% !TEX program = xelatex
%
% Brezis《泛函分析、Sobolev 空间与偏微分方程》中译本
% 原书: Haim Brezis, Functional Analysis, Sobolev Spaces and
%       Partial Differential Equations, Universitext, Springer, 2011.
% 模板: 李文威《代数学方法》AJbook.cls (lww 模板)
%
% 编译: bash _build.sh   (xelatex x2 -> texindy -> xelatex)
% 本译本不使用 biblatex/biber, 理由见《数学译本制作规范》§6.

\documentclass[
	11pt,
	geometry=a4,
	coverpage=cover,
	fontsetup=font-setup-open,
	titlesetup=titles-setup,
]{AJbook}

% ==================== 类未加载的宏包 ====================
\usepackage{mycommand}      % \R \CC \Z \N \Q \norm \Ker \Image \supp \lrangle 等
\usepackage{myarrows}       % 可伸缩箭头 (依赖 tikz-cd)
\usepackage{tikz-cd}        % 交换图
\usepackage{mathrsfs}       % \mathscr
\usepackage{enumitem}       % 列表间距与自定义标签
\usepackage{xurl}           % 长 URL 折行
\usepackage{hologo}         % \hologo{XeLaTeX}
\usepackage{chngcntr}       % \counterwithout (图全书连续编号)
\usepackage{caption}        % 图题注格式 (原书 "Fig. 3")
\usepackage{longtable}      % 跨页表格 (符号表)
\usetikzlibrary{arrows.meta,calc,positioning,decorations.markings,patterns,shapes.geometric}

\setcounter{MaxMatrixCols}{20}
\setlength{\emergencystretch}{2.5em}

% ==================== 函数分析专用记号 ====================
\newcommand{\abs}[1]{\left|#1\right|}
\newcommand{\inner}[2]{\left(#1,#2\right)}          % Hilbert 内积
\newcommand{\dual}[2]{\left\langle #1,#2\right\rangle} % 对偶配对
\newcommand{\weakto}{\rightharpoonup}                % 弱收敛
\newcommand{\weakstar}{\overset{*}{\rightharpoonup}} % 弱*收敛
% "弱*" 中的星号必须紧贴"弱"字: xeCJK 默认在汉字与数学之间插胶,
% 直接用 弱$^\star$ 会印成 "弱 ⋆ 拓扑". 此处局部关掉 CJK 胶.
\newcommand{\wstar}{\texorpdfstring{{\xeCJKsetup{CJKecglue={}}弱$^\star$}}{弱*}}
\newcommand{\embeds}{\hookrightarrow}                % 连续嵌入
\newcommand{\embedsd}{\hookrightarrow\joinrel\relbar\joinrel\rightarrow} % 紧嵌入
\DeclareMathOperator{\dist}{dist}
\DeclareMathOperator{\esssup}{ess\,sup}
\DeclareMathOperator{\essinf}{ess\,inf}
\DeclareMathOperator{\codim}{codim}
\DeclareMathOperator{\Int}{Int}
\DeclareMathOperator{\osc}{osc}
\DeclareMathOperator{\ind}{ind}      % Fredholm 指标 (第 6 章评注)
\newcommand{\defeq}{\coloneqq}

% 原书专用的算子记号
\DeclareMathOperator{\Lip}{Lip}
\newcommand{\LL}{\mathcal{L}}          % 有界线性算子空间
\newcommand{\KK}{\mathcal{K}}          % 紧算子空间
\newcommand{\Cc}{C_c}
\newcommand{\Ccinf}{C_c^\infty}
\newcommand{\Wmp}{W^{m,p}}
\newcommand{\Wonep}{W_0^{1,p}}
\newcommand{\Wone}{W^{1,p}}

% ==================== 编号体系 (严格保真原书体例) ====================
% 原书体例:
%   Theorem / Proposition / Corollary 共用一个按章连续计数器: 1.1, 1.2, ...
%   Lemma 独立按章编号: Lemma 1.1, 1.2, ...
%   Remark 独立按章编号, 章内不带小数点: Remark 1, 2, ...
%   Example 独立按章编号, 章内不带小数点: Example 1, 2, ...
%   Definition 无编号, 行内粗体 "Definition." 起头
%   公式按章从 (1) 重新开始, 章内连续
%   习题 1.1--1.27 按章编号, 编号外加方框
%   图全书连续 Fig. 1--9
%   解答按习题号, 解答内公式 (S1), (S2), ... 每题重新开始
%   Problem 全书连续 1--51
%
% 原书另有 "•" 起头的重点结论 (Theorem/Proposition/Corollary/Lemma/Remark 各有),
% 圆点排在标题行左侧, 编号照常递增. 用一个可选参数控制是否加圆点.

\newcounter{cstmt}[chapter]   % 定理/命题/推论共用
\renewcommand{\thecstmt}{\thechapter.\arabic{cstmt}}
\newcounter{clemma}[chapter]
\renewcommand{\theclemma}{\thechapter.\arabic{clemma}}
\newcounter{cremark}[chapter]
\renewcommand{\thecremark}{\arabic{cremark}}
% 例: 按节重新编号 (原书 §3.2 有例 1、2, §3.7 又从例 1 开始)
\newcounter{cexample}[section]
\renewcommand{\thecexample}{\arabic{cexample}}

% 星号标记的结论: 原书把 "⋆" 紧贴在 "Theorem" 前, 不插入空隙.
% xeCJK 默认在汉字与数学间加胶, 故局部关掉.

\newcommand{\cstmtStar}[1]{{\xeCJKsetup{CJKecglue={}}$\star$#1}}

\newcommand{\cstmtHead}[2]{%
	\par\addvspace{2mm}%
	\noindent{\sffamily\bfseries\thmheiti #1}%
	\ifstrempty{#2}{}{\enspace(#2)}%
	{\sffamily\bfseries\thmheiti .}%
	\enspace\ignorespaces
}
\newcommand{\cstmtTail}{\par\addvspace{2mm}}

% 类在它自己的 \AtBeginDocument 里用 ntheorem 建了 theorem/definition/remark/example 等,
% 那个钩子比本文件注册得早, 所以本文件的钩子后执行, 此时环境已存在, 用 \renewenvironment 接管.
\AtBeginDocument{%
	% 无编号定义: "定义" + 空格 + 宋体正文
	\renewenvironment{definition}{%
		\par\addvspace{2mm}\noindent
		{\sffamily\bfseries\thmheiti 定义.}\enspace\songti\ignorespaces
	}{\cstmtTail}%
	% 定理/命题/推论: 共用计数器, 楷体正文, 可选圆点/星号
	\renewenvironment{theorem}[1][]{\refstepcounter{cstmt}\cstmtHead{定理~\thecstmt}{#1}\kaishu}{\cstmtTail}%
	\renewenvironment{proposition}[1][]{\refstepcounter{cstmt}\cstmtHead{命题~\thecstmt}{#1}\kaishu}{\cstmtTail}%
	\renewenvironment{corollary}[1][]{\refstepcounter{cstmt}\cstmtHead{推论~\thecstmt}{#1}\kaishu}{\cstmtTail}%
	\newenvironment{keytheorem}[1][]{\refstepcounter{cstmt}\cstmtHead{\textbullet\enspace 定理~\thecstmt}{#1}\kaishu}{\cstmtTail}%
	\newenvironment{keyproposition}[1][]{\refstepcounter{cstmt}\cstmtHead{\textbullet\enspace 命题~\thecstmt}{#1}\kaishu}{\cstmtTail}%
	\newenvironment{keycorollary}[1][]{\refstepcounter{cstmt}\cstmtHead{\textbullet\enspace 推论~\thecstmt}{#1}\kaishu}{\cstmtTail}%
	\newenvironment{startheorem}[1][]{\refstepcounter{cstmt}\cstmtHead{\cstmtStar{定理~\thecstmt}}{#1}\kaishu}{\cstmtTail}%
	\newenvironment{starproposition}[1][]{\refstepcounter{cstmt}\cstmtHead{\cstmtStar{命题~\thecstmt}}{#1}\kaishu}{\cstmtTail}%
	\newenvironment{starcorollary}[1][]{\refstepcounter{cstmt}\cstmtHead{\cstmtStar{推论~\thecstmt}}{#1}\kaishu}{\cstmtTail}%
	% 引理: 独立计数器
	\renewenvironment{lemma}[1][]{\refstepcounter{clemma}\cstmtHead{引理~\theclemma}{#1}\kaishu}{\cstmtTail}%
	\newenvironment{keylemma}[1][]{\refstepcounter{clemma}\cstmtHead{\textbullet\enspace 引理~\theclemma}{#1}\kaishu}{\cstmtTail}%
	\newenvironment{starlemma}[1][]{\refstepcounter{clemma}\cstmtHead{\cstmtStar{引理~\theclemma}}{#1}\kaishu}{\cstmtTail}%
	% 注记 / 例: 独立计数器, 章内无小数点, 宋体正文
	\renewenvironment{remark}[1][]{\refstepcounter{cremark}\cstmtHead{注记~\thecremark}{#1}\songti}{\cstmtTail}%
	\newenvironment{keyremark}[1][]{\refstepcounter{cremark}\cstmtHead{\textbullet\enspace 注记~\thecremark}{#1}\songti}{\cstmtTail}%
	\newenvironment{starremark}[1][]{\refstepcounter{cremark}\cstmtHead{\cstmtStar{注记~\thecremark}}{#1}\songti}{\cstmtTail}%
	\renewenvironment{example}[1][]{\refstepcounter{cexample}\cstmtHead{例~\thecexample}{#1}\songti}{\cstmtTail}%
	\newenvironment{starexample}[1][]{\refstepcounter{cexample}\cstmtHead{\cstmtStar{例~\thecexample}}{#1}\songti}{\cstmtTail}%
	% 原书用 "Notation." 起头的行内小节, 与 Definition 同级
	\newenvironment{notation}{%
		\par\addvspace{2mm}\noindent
		{\sffamily\bfseries\thmheiti 记号.}\enspace\songti\ignorespaces
	}{\cstmtTail}%
	% 每章末尾 "Comments on Chapter N" 内的带号粗体小标题 ("1. Generalizations...")
	% 原书有两种排法: 带小标题的 (第 1 章) 编号单独成行, 其后另起小标题;
	% 不带小标题的 (第 2、3 章) 编号与正文同行. 分别用 \cmtitem 与 \cmtnum.
	\newcommand{\cmtitem}[1]{%
		\par\addvspace{0.7\baselineskip}\noindent
		{\sffamily\bfseries #1}\par\nobreak\addvspace{0.25\baselineskip}\noindent\ignorespaces
	}%
	\newcommand{\cmtnum}[1]{%
		\par\addvspace{0.7\baselineskip}\noindent
		{\sffamily\bfseries #1}\enspace\ignorespaces
	}%
	% 其下的字母小标题 "(a) The Krein--Milman theorem."
	\newcommand{\cmtsub}[1]{%
		\par\addvspace{0.45\baselineskip}\noindent
		{\sffamily\bfseries #1}\par\nobreak\addvspace{0.15\baselineskip}\noindent\ignorespaces
	}%
	% 评注里 "字母 + 结论" 同行起排的情形 ("(a) Theorem 5.12 (...)." 第 5 章评注 1):
	% 字母排在定理编号之前, 不另起一行.
	\newenvironment{cmtlettertheorem}[2]{%
		\refstepcounter{cstmt}%
		\par\addvspace{2mm}\noindent
		{\sffamily\bfseries\thmheiti (#1)\enspace 定理~\thecstmt}%
		\ifstrempty{#2}{}{\enspace(#2)}%
		{\sffamily\bfseries\thmheiti .}\enspace\kaishu\ignorespaces
	}{\cstmtTail}%
		% 证明环境 (统一黑体标题 + 空心方格证毕符号 \openbox)
		\renewenvironment{proof}[1][]{%
			\par\addvspace{2mm}\noindent
			{\sffamily\bfseries\thmheiti 证明\ifstrempty{#1}{}{~(#1)}.}%
			\enspace\ignorespaces
		}{%
			\hfill\openbox\par\addvspace{2mm}%
		}%
		% 超链接锚点唯一化 (同一编号在不同章重复时必须区分)
	\renewcommand{\theHcstmt}{\thechapter.\arabic{cstmt}}%
	\renewcommand{\theHclemma}{\thechapter.\arabic{clemma}}%
	\renewcommand{\theHcremark}{\thechapter.\arabic{cremark}}%
	\renewcommand{\theHcexample}{\thechapter.\thesection.\arabic{cexample}}%
	\renewcommand{\theHequation}{\thechapter.\arabic{equation}}%
}

% 公式: 按章从 (1) 重新开始, 章内连续
% \counterwithin* 只重置计数器, 不改写 \theequation, 故显示仍为 (1)(2)...
\counterwithin*{equation}{chapter}
\renewcommand{\theequation}{\arabic{equation}}

% 脚注: 原书每章从 1 重新编号 (第 1 章 1--7, 第 3 章 1--3 ...)
\counterwithin*{footnote}{chapter}
% 原书脚注: 编号上标后留一个空隙再接正文, 续行齐左边界 (无悬挂缩进).
\makeatletter
\renewcommand\@makefntext[1]{\noindent\@makefnmark\enspace #1}
\makeatother

% 带星号的节: 原书把 "⋆" 印在节号之前 (⋆2.4 ...), 但目录与天眉不带星号.
% 做法: 给节标题的编号前缀加一个可开关的星号; 星号不进入 \thesection,
% 因此 \addcontentsline 与 \sectionmark 都取不到它.
\newcommand{\AJsecstar}{}
\titleformat{name=\section}
	{\filleft\normalfont\sffamily\bfseries\CJKfamily{sectionfont}}
	{}
	{0mm}
	{ \settowidth{\BoxTtlwidth}{\Huge \AJsecstar\thesection \hspace{0.7em} \Large #1}
		\ifdim \BoxTtlwidth < \textwidth
			\MakeSectBox{\Huge \AJsecstar\thesection \hspace{0.7em} \Large #1}\vskip-18pt%
		\else
			\Huge \underline{\AJsecstar\thesection} \hspace{0.7em} \Large #1%
		\fi }
	[\gdef\AJsecstar{}]
\newcommand{\starsection}[1]{\gdef\AJsecstar{$\star$}\section{#1}}

% 原书的行内粗体小标题 ("Definition of the adjoint A*." / "Examples" 等)
\newcommand{\runin}[1]{%
	\par\addvspace{2mm}\noindent{\sffamily\bfseries #1}\enspace\ignorespaces
}

% 图: 全书连续编号
% 原书题注形如 "Fig. 3"(编号前有 "Fig." 前缀), 故题注用 "图 3" 而非 "图3".
\counterwithout{figure}{chapter}
\renewcommand{\thefigure}{\arabic{figure}}
\AtBeginDocument{\renewcommand{\theHfigure}{\arabic{figure}}}
\DeclareCaptionLabelFormat{brezisfig}{图~#2}
\captionsetup[figure]{labelformat=brezisfig, labelsep=quad}

% ==================== 习题 / 解答 / 问题 ====================
% 习题: 章.序, 编号外加方框 (原书体例)
% 类在 \AtBeginDocument 里用 ntheorem 定义了 exercise (共用 theorem 计数器);
% 本文件注册的钩子晚于类, 因此在这里解除并重建.
\AtBeginDocument{%
	\let\exercise\relax \let\endexercise\relax
}
\newcounter{brexer}[chapter]
\renewcommand{\thebrexer}{\thechapter.\arabic{brexer}}
\providecommand{\theHexer}{}
\renewcommand{\theHexer}{\thechapter.\arabic{brexer}}
\makeatletter
\AtBeginDocument{%
	\newenvironment{exercise}[1][]{%
		\par\medskip\noindent
		\refstepcounter{brexer}%
		\hypertarget{exer:\thebrexer}{}%
		\@ifundefined{r@hassol:\thebrexer}{%
			\fbox{\small\bfseries\thebrexer}%
		}{%
			\hyperlink{sol:\thebrexer}{\fbox{\small\bfseries\thebrexer}}%
		}%
		\ifstrempty{#1}{\enspace\ignorespaces}%
			{\enspace{\itshape #1}\par\nobreak\ignorespaces}%
	}{\par\medskip}%
}

% 解答: 与习题同号, 但不重复占计数器. 双向超链接回原题
\newcommand{\solutionhead}[1]{%
	\par\addvspace{0.8\baselineskip}%
	\noindent\hypertarget{sol:#1}{}%
	\label{hassol:#1}%
	\hyperlink{exer:#1}{{\sffamily\bfseries #1}}%
	\par\nobreak\addvspace{0.2\baselineskip}%
}

% 解答内的公式: (S1), (S2), ... 每题重新开始
\newcounter{sequation}
\renewcommand{\thesequation}{S\arabic{sequation}}
\newcommand{\resetsequation}{\setcounter{sequation}{0}}

% Problem: 全书连续, 双向超链接到部分问题解答
\newcounter{bproblem}
\renewcommand{\thebproblem}{\arabic{bproblem}}
\newenvironment{problem}[1][]{%
	\par\bigskip\noindent
	\refstepcounter{bproblem}%
	\hypertarget{prob:\thebproblem}{}%
	\@ifundefined{r@hassolprob:\thebproblem}{%
		{\sffamily\bfseries 问题~\thebproblem}%
	}{%
		\hyperlink{solprob:\thebproblem}{{\sffamily\bfseries 问题~\thebproblem}}%
	}%
	\ifstrempty{#1}{}{\quad(#1)}%
	\par\nobreak\smallskip\noindent\ignorespaces
}{\par\bigskip}

% 问题解答标题: 设置锚点 solprob:N 与标记 hassolprob:N,
% 点击标题文字可跳回正文原题 prob:N
\newcommand{\problemsolutionhead}[1]{%
	\par\addvspace{0.8\baselineskip}%
	\noindent\hypertarget{solprob:#1}{}%
	\label{hassolprob:#1}%
	\hyperlink{prob:#1}{{\sffamily\bfseries 问题~#1}}%
	\par\nobreak\addvspace{0.2\baselineskip}%
}
\makeatother

% ==================== 参考文献 (原书体例, 不用 biblatex) ====================
% 原书: 按作者分组, 组内 [1] [2] 编号, 正文写 "J. Diestel [1, Chapter 1]".
% \bref{Key}{n}             -> 可点击 [n]
% \brefcite{Key}{n}{后缀}    -> 可点击 [n 后缀]
\newcommand{\bref}[2]{\hyperlink{ref:#1-#2}{[#2]}}
\newcommand{\brefcite}[3]{\hyperlink{ref:#1-#2}{[#2#3]}}
\newcommand{\brefentry}[3]{%
	\par\addvspace{0.35\baselineskip}%
	\noindent\hypertarget{ref:#1-#2}{}\hangindent=1.5em\hangafter=1
	#3~\bref{#1}{#2}\par
}

% ==================== 索引 ====================
\usepackage[xindy, noautomatic]{imakeidx}
\makeindex[program=truexindy, title={名词索引暨英汉对照}, columns=2, intoc=true,%
	options=-M texindy -I xelatex -C utf8]

% ==================== 超链接 ====================
\usepackage[bookmarksnumbered=true,hyperindex=true]{hyperref}
\hypersetup{
	pdfauthor={Haim Brezis},
	pdftitle={泛函分析、Sobolev 空间与偏微分方程},
	pdfsubject={Functional Analysis, Sobolev Spaces and Partial Differential Equations},
	colorlinks=true,
	linkcolor=blue,
	citecolor=red,
	urlcolor=teal
}

\begin{document}

\frontmatter
% chapters/frontmatter.tex
% 前置部分: 献词 / 前言 / 使用说明 / 致谢
% 源: 原书 PDF 第 5--8 页 (印刷页 v--ix)

\thispagestyle{empty}
\vspace*{0.22\textheight}
\begin{center}
	\large\itshape
	献给 Felix Browder\\
	——我的良师与挚友,\\
	是他教会我以泛函分析学者的眼光\\
	欣赏偏微分方程
\end{center}
\cleardoublepage

\chapter*{译者序}
\markboth{译者序}{译者序}
\addcontentsline{toc}{chapter}{译者序}
\enlargethispage{2\baselineskip}

海姆·布雷齐斯 (Haïm Brezis, 1944--2024) 教授是当代偏微分方程与非线性泛函分析大师.
本书由其在巴黎大学与罗格斯大学讲授的讲义凝练而成,
将抽象泛函分析与现代 PDE (尤其是 Sobolev 空间与变分法) 浑然融合, 是享誉全球的经典教材.
为方便国内学子研读, 我们依据 Springer 2011 年英文版进行了全本中文重排与翻译.
针对电子化研读体验, 本次重排特别做了如下现代化改进与体例说明:

\begin{enumerate}[label=\arabic*., leftmargin=2.2em, itemsep=0.15em, topsep=0.15em]
	\item \textbf{习题与解答的双向交互跳转}:
	      本书收录了极其丰富的习题与综合性问题.为了让读者在自学和演算时获得最佳阅读体验,
	      我们在 PDF 电子版中全面实现了\emph{双向可逆超链接} (Bidirectional Hyperlinks):
	      \begin{itemize}[label=$\circ$, leftmargin=1.5em, itemsep=0.05em, topsep=0.05em]
		      \item 正文各章末尾的习题中, 凡在书末「部分习题解答」中附有解答的题目,
		            其方框题号 (如 \fbox{\small\bfseries 1.1}) 均内置交互跳转链接, 点击即可直接跳至书末对应解答;
		            在书末解答处, 点击加粗题号同样可瞬间跳回正文原题所在页面.
		      \item 在后置的 51 道大型「问题」与「部分问题解答」之间,
		            亦同样支持点击问题标题实现原题与参考解答之间的双向无缝跳转.
	      \end{itemize}
	\item \textbf{参考文献引用体例说明}:
	      原书参考文献列表严格遵循欧洲古典分析专著的编排传统:
	      文献按作者姓氏首字母升序排列, 但\emph{每位作者名下的著作各自独立按 $[1],[2],\dots$ 重新编号} (而非全书统一的单一大流水号).
	      因此, 正文中出现的引用形式为 ``作者名 [编号]'' (例如 ``H.~L.~Royden [1]'', ``E.~Lieb--M.~Loss [1]'', ``J.~Dieudonné [1]'').
	      文中的每一个 $[1]$ 或 $[2]$ 均已建立精确的超链接锚点, 点击即可直达书末该作者名下的具体出版条目.
	\item \textbf{图表矢量重绘与高精度交叉引用}:
	      原书中全部 9 幅几何原理示意图均采用 TikZ 矢量代码重新精细绘制, 避免了扫描位图的失真与模糊;
	      定理、引理、推论、小节编号 ($\S$) 以及公式引用均全面绑定超链接 (\verb|\ref| 与 \verb|\eqref|),
	      确保读者可快速穿透引用链条.
	\item \textbf{翻译原则与严谨性}:
	      本译本坚持严谨忠实原则, 绝不擅自删节、合并或跳跃原文论证细节.
	      对于原版英文印刷中已发现的个别符号笔误, 译本均在 LaTeX 源文件中以注释形式详实记录并予以严格核对.
\end{enumerate}

\vspace{0.6\baselineskip}
\begin{flushright}
	译者\\
	2026 年 10 月
\end{flushright}

\cleardoublepage

\chapter*{前言}
\markboth{前言}{前言}
\addcontentsline{toc}{chapter}{前言}

本书源于我在巴黎大学多年讲授的一门课程.它面向在实分析方面已有良好基础的学生
(例如 G.\,B.~Folland [2], A.\,W.~Knapp [1] 和 H.\,L.~Royden [1] 的教科书中所讲述的内容).
我构想了一套混合了两个不同``世界''之内容的方案: 泛函分析 (FA) 与偏微分方程 (PDE).
第一部分处理 FA 与算子理论中的抽象结果.第二部分研究具有特定可微性质的函数空间
(一个或多个实变量的函数), 即著名的 Sobolev 空间, 它位于现代 PDE 理论的核心.
我说明如何把 FA 中的抽象结果用于求解 PDE.Sobolev 空间出现在纯数学与应用数学的广泛问题中.
它们出现在例如微分几何, 调和分析, 工程, 力学和物理中产生的线性与非线性 PDE 里.
它们属于任何分析方向研究生的工具箱.

遗憾的是, FA 与 PDE 虽然联系紧密, 却常常分属两门课程.FA 中处理的许多问题都源自 PDE
(从历史的角度看, 可参见例如 J.~Dieudonné [1] 与 H.~Brezis--F.~Browder [1]).
专论 FA 的书 (甚至多卷本巨著) 数量众多.处理 PDE 的教科书也不少.
然而, 面向研究生的综合性讲述却很少见, 我试图填补这一空缺.
常常被数学中最抽象的构造所吸引的学生, 通常也为 FA 的优雅所吸引.
另一方面, 他们又对 PDE 中没完没了的公式以及数不清的下标感到畏惧.
我尝试先分析一维 PDE (即 ODE——常微分方程) 这一简单情形, 从而呈现一条从 FA 到 PDE 的
``平滑''过渡; 对初学者而言, 一维情形看起来要容易把握得多.按这种做法, 我所阐述的技巧
对 ODE 而言可能过于精细, 但它们后来成为 PDE 理论的基石.
这种安排使学生此后处理精细的高维 PDE 时容易得多.

本书的一个较早版本最初于 1983 年以法文出版, 此后有多种译本, 在世界范围内颇受欢迎,
并被许多欧洲大学采用为教科书.法文本的一个缺陷是缺少习题.
本书含有丰富的习题.我打算在未来的版本中再增加一些.
我也概述了一些较新的进展, 特别是在非线性 PDE 方向上的进展.

\section*{使用说明}
\addcontentsline{toc}{section}{使用说明}

\begin{enumerate}[label=\arabic*., leftmargin=2.2em, itemsep=0.35em]
	\item 以圆点符号 \textbullet\ 起头的陈述或段落极为重要, 务必透彻掌握, 才能理解后文.
	\item 标有星号 $\star$ 的结果初学者可以跳过; 它们只对高年级读者有意义.
	\item 在每一章中, 我把命题, 定理与推论连续编号 (例如命题~
\ref{prop:3.6} 之后是定理~
\ref{thm:3.7}, 推论~
\ref{cor:3.8},
	      等等).只有注记与引理另行单独编号.
	\item 为简化表述, 我假定所有向量空间都是 $\mathbb{R}$ 上的.多数结果对 $\mathbb{C}$ 上的向量空间
	      仍然成立.我在第~
\ref{chap:11}~章加了一节简短的讨论, 说明其异同.
	\item 许多章之后附有大量习题.部分解答放在书末.更精细的问题另设一节``问题'',
	      其后是``问题的部分解答''.这些问题通常需要用到若干章的知识.
	      我在每个问题的开头标明了所涉及的章.有些习题与问题阐述的是本章正文中未加细节
	      或未给证明的结果.
\end{enumerate}

\section*{致谢}
\addcontentsline{toc}{section}{致谢}

在本书的写作过程中, 我得到了两位挚友兼前同事的诸多鼓励: Ph.~Ciarlet 与 H.~Berestycki.
我非常感谢 G.~Tronel, M.~Comte, Th.~Gallouet, S.~Guerre-Delabrière, O.~Kavian, S.~Kichenassamy 以及已故的 Th.~Lachand-Robert, 他们与我分享了与学生打交道的``实战经验''.
S.~Antman, D.~Kinderlehrer 和 Y.~Li 向我解释了美国学生的背景与``品味''.
C.~Jones 惠允我使用他为个人之需所准备的, 法文原书若干章的英译稿.
我要感谢 A.~Ponce, H.-M.~Nguyen, H.~Castro 和 H.~Wang, 他们仔细校阅了本书的部分内容.
我有幸在罗格斯大学得到两位出色的助手, 本书的大部分内容由她们录入:
现已退休的 Barbara Miller, 以及现任的 Barbara Mastrian.对她们始终如一的奉献与专业帮助,
我的赞美与感激之词不足以表达.她们总能为 PDE 公式中棘手的繁难之处找到令人满意的处理办法.
没有她们的热情与耐心, 本书绝无可能完成.一如既往, 与 Springer 的 Ann Kostant 在这一项目上的合作
令我十分愉快.过去我多次有机会体会到她对数学界长期不渝的投入.

作者部分得到 NSF 基金 DMS-0802958 的资助.

\vspace{1.5\baselineskip}
\begin{flushright}
	Haim Brezis\\
	罗格斯大学\\
	2010 年 3 月
\end{flushright}

\cleardoublepage


\mainmatter
% chapters/ch01.tex
% 第~
% \ref{chap:1}~章 Hahn--Banach 定理. 共轭凸函数理论引论
% 源: 原书 PDF 第 12--40 页 (印刷页 1--29)

\chapter{Hahn--Banach 定理. 共轭凸函数理论引论}
\label{chap:1}

\index{bansuisuanzi@伴随算子 (adjoint)}
\index{Hahn--Banach theorem@Hahn--Banach 定理 (Hahn--Banach theorem)}
\section{Hahn--Banach 定理的解析形式: 线性泛函的延拓}
\label{sec:1.1}

设 $E$ 为 $\R$ 上的向量空间. 回忆\emph{泛函}是指定义在 $E$ 上 (或 $E$ 的某个子空间上),
取值于 $\R$ 的函数. 本节的主要结果, 是关于把定义在 $E$ 的线性子空间上的线性泛函
延拓为定义在整个 $E$ 上的线性泛函.

\index{HahnBanachdingli@Hahn--Banach 定理!jiexixingshi@解析形式 (Hahn--Banach theorem, analytic form)}
\index{Helly@Helly (Helly)}
\begin{theorem}[Helly, Hahn--Banach 解析形式]
\label{thm:1.1}
设 $p:E\to\R$ 为满足\footnote{满足 \eqref{eq:1.1} 与 \eqref{eq:1.2} 的函数 $p$ 有时称为
\emph{\index{Minkowskifanhan@Minkowski 泛函 (Minkowski functional)}
Minkowski 泛函}(Minkowski functional).}
\begin{align}
p(\lambda x) &= \lambda p(x) && \forall x\in E \quad\text{及}\quad \forall\lambda>0,
  \label{eq:1.1}\\
p(x+y) &\le p(x)+p(y) && \forall x,y\in E
  \label{eq:1.2}
\end{align}
的函数.设 $G\subset E$ 为线性子空间, $g:G\to\R$ 为满足
\begin{equation}
g(x)\le p(x)\quad \forall x\in G
  \label{eq:1.3}
\end{equation}
的线性泛函.在这些假设下, 存在定义在整个 $E$ 上的线性泛函 $f$, 它延拓 $g$, 即
$g(x)=f(x)\ \forall x\in G$, 并且满足
\begin{equation}
f(x)\le p(x)\quad \forall x\in E.
  \label{eq:1.4}
\end{equation}
\end{theorem}

定理  \ref{thm:1.1} 的证明依赖于 \index{Zornyinli@Zorn 引理 (Zorn's lemma)}
Zorn 引理, 这是关于有序集的一条著名而非常有用的性质.
在叙述 Zorn 引理之前, 我们须澄清一些概念.设 $P$ 为带 (偏) 序关系 $\le$ 的集合.
若对 $Q$ 中任意一对 $(a,b)$, 或者 $a\le b$, 或者 $b\le a$ (或二者兼有!), 则称子集 $Q\subset P$
为\emph{全序}的.设 $Q\subset P$ 为 $P$ 的子集; 若对每个 $a\in Q$ 都有 $a\le c$, 则称 $c\in P$
为 $Q$ 的\emph{上界}.若不存在 $x\in P$ 使得 $m\le x$ (除非 $x=m$), 则称 $m\in P$ 为 $P$ 的
\emph{极大元}.注意 $P$ 的极大元不必是 $P$ 的上界.

若 $P$ 中每个全序子集 $Q$ 都有上界, 则称 $P$ 是\emph{归纳的}.

\begin{keylemma}[Zorn]
\label{lem:1.1}
每个非空且归纳的有序集都有极大元.
\end{keylemma}

Zorn 引理可由选择公理推出, 但我们在此不讨论它的推导; 可参见例如 J.~Dugundji
\bref{Dugundji}{1}, N.~Dunford--J.~T.~Schwartz \bref{DunfordSchwartz}{1} (第 1 卷,
定理 1.2.7), E.~Hewitt--K.~Stromberg \bref{HewittStromberg}{1}, S.~Lang \bref{Lang}{1}
以及 A.~Knapp \bref{Knapp}{1}.

\begin{remark}
\label{rem:1.1}
Zorn 引理在分析中有许多重要的应用.它是在证明某些看似天真的存在性命题时的基本工具,
例如``每个向量空间都有基''(见习题  \ref{exer:1.5}) 和``任何向量空间上都存在非平凡的
线性泛函''.多数分析学家并不知道如何证明 Zorn 引理; 但对分析学家而言,
理解 Zorn 引理的陈述并能够正确地使用它, 却是相当要紧的!
\end{remark}

\begin{proof}[引理 \ref{lem:1.2}]
% 原书此处误作 "Proof of Lemma 1.2", 实际证明的是定理 1.1 (真正的引理~
% \ref{lem:1.2} 的证明见
% 印刷页 6). 两份 OCR 与 PDF 文本层一致, 属原书排印之误. 依技能规范: 保留原样并记录,
% 不擅自更正. 见《项目交接文档》"原书疑似错误"条目.
考虑集合
\begin{equation*}
P=\left\{h:D(h)\subset E\to\R\ \middle|\
\begin{array}{l}
D(h)\ \text{是}\ E\ \text{的线性子空间},\\
h\ \text{线性},\quad G\subset D(h),\\
h\ \text{延拓}\ g,\ \text{且}\ h(x)\le p(x)\quad\forall x\in D(h)
\end{array}
\right\}.
\end{equation*}
在 $P$ 上定义序关系
\begin{equation*}
(h_1\le h_2)\iff\bigl(D(h_1)\subset D(h_2)\ \text{且}\ h_2\ \text{延拓}\ h_1\bigr).
\end{equation*}
显然 $P$ 非空, 因为 $g\in P$.我们断言 $P$ 是归纳的.事实上, 设 $Q\subset P$ 为全序子集;
把 $Q$ 写成 $Q=(h_i)_{i\in I}$, 并置
\begin{equation*}
D(h)=\bigcup_{i\in I}D(h_i),\qquad
h(x)=h_i(x)\quad \text{若}\ x\in D(h_i)\ \text{对某个}\ i.
\end{equation*}
容易看出 $h$ 的定义是有意义的, $h\in P$, 且 $h$ 是 $Q$ 的上界.于是我们可以应用 Zorn 引理,
从而得到 $P$ 中的一个极大元 $f$.我们断言 $D(f)=E$, 这就完成了定理  \ref{thm:1.1} 的证明.

反设 $D(f)\ne E$.取 $x_0\notin D(f)$; 置 $D(h)=D(f)+\R x_0$, 并对每个 $x\in D(f)$ 置
$h(x+tx_0)=f(x)+t\alpha\ (t\in\R)$, 其中常数 $\alpha\in\R$ 将被选得使 $h\in P$.
我们必须保证
\begin{equation*}
f(x)+t\alpha\le p(x+tx_0)\quad \forall x\in D(f)\quad\text{及}\quad \forall t\in\R.
\end{equation*}
由 \eqref{eq:1.1} 只需验证
\begin{equation*}
\begin{cases}
f(x)+\alpha\le p(x+x_0) & \forall x\in D(f),\\
f(x)-\alpha\le p(x-x_0) & \forall x\in D(f).
\end{cases}
\end{equation*}
换言之, 我们必须找到某个 $\alpha$ 满足
\begin{equation*}
\sup_{y\in D(f)}\{f(y)-p(y-x_0)\}\le\alpha\le
\inf_{x\in D(f)}\{p(x+x_0)-f(x)\}.
\end{equation*}
这样的 $\alpha$ 存在, 因为
\begin{equation*}
f(y)-p(y-x_0)\le p(x+x_0)-f(x)\quad \forall x\in D(f),\ \forall y\in D(f);
\end{equation*}
事实上, 由 \eqref{eq:1.2} 可得
\begin{equation*}
f(x)+f(y)\le p(x+y)\le p(x+x_0)+p(y-x_0).
\end{equation*}
于是 $f\le h$; 但这是不可能的, 因为 $f$ 是极大元且 $h\ne f$.
\end{proof}

现在描述定理  \ref{thm:1.1} 在 $E$ 为赋范向量空间 (n.v.s.), 范数为 $\norm{\cdot}$ 的情形下的
一些简单应用.

\begin{notation}
记 $E^\star$ 为 $E$ 的\emph{\index{duioukongjian@对偶空间 (dual space)}
对偶空间}, 即 $E$ 上一切连续线性泛函构成的空间;
$E^\star$ 上的 (对偶) 范数定义为
\begin{equation}
\norm{f}_{E^\star}=\sup_{\substack{x\in E\\ \norm{x}\le 1}}\abs{f(x)}
=\sup_{\substack{x\in E\\ \norm{x}\le 1}}f(x).
  \label{eq:1.5}
\end{equation}
在不致混淆时, 我们也用 $\norm{f}$ 代替 $\norm{f}_{E^\star}$.

给定 $f\in E^\star$ 与 $x\in E$, 我们常写 $\dual{f}{x}$ 代替 $f(x)$;
并称 $\dual{\cdot}{\cdot}$ 为对偶 $E^\star,E$ 的\emph{标量积}.

众所周知 $E^\star$ 是 Banach 空间, 即 $E^\star$ 是完备的 (即使 $E$ 不完备);
这可由 $\R$ 完备这一事实推出.
\end{notation}

\begin{keycorollary}
\label{cor:1.2}
设 $G\subset E$ 为线性子空间.若 $g:G\to\R$ 为连续线性泛函, 则存在 $f\in E^\star$
延拓 $g$, 且
\begin{equation*}
\norm{f}_{E^\star}=\sup_{\substack{x\in G\\ \norm{x}\le 1}}\abs{g(x)}=\norm{g}_{G^\star}.
\end{equation*}
\end{keycorollary}

\begin{proof}
在定理  \ref{thm:1.1} 中取 $p(x)=\norm{g}_{G^\star}\norm{x}$.
\end{proof}

\begin{keycorollary}
\label{cor:1.3}
对每个 $x_0\in E$ 都存在 $f_0\in E^\star$ 使得
\begin{equation*}
\norm{f_0}=\norm{x_0}\quad\text{且}\quad \dual{f_0}{x_0}=\norm{x_0}^2.
\end{equation*}
\end{keycorollary}

\begin{proof}
在推论  \ref{cor:1.2} 中取 $G=\R x_0$ 及 $g(tx_0)=t\norm{x_0}^2$, 于是
$\norm{g}_{G^\star}=\norm{x_0}$.
\end{proof}

\begin{remark}
\label{rem:1.2}
推论  \ref{cor:1.3} 给出的元素 $f_0$ 一般并不唯一 (试着构造一个例子, 或见习题
 \ref{exer:1.2}).然而, 若 $E^\star$ 是\index{yangetufanshu@严格凸范数 (strictly convex norm)}
严格凸的\footnote{若
$\norm{tx+(1-t)y}<1,\ \forall t\in(0,1),\ \forall x,y$ 满足 $\norm{x}=\norm{y}=1$ 且
$x\ne y$, 则称赋范空间是\emph{严格凸}的; 见习题  \ref{exer:1.26}.}
——例如 $E$ 为 Hilbert 空间 (见第~
\ref{chap:5}~章), 或 $E=L^p(\Omega)$ 且 $1<p<\infty$
(见第~
\ref{chap:4}~章)——则 $f_0$ 是唯一的.一般地, 我们对每个 $x_0\in E$ 置
\begin{equation*}
F(x_0)=\bigl\{f_0\in E^\star;\ \norm{f_0}=\norm{x_0}\ \text{且}\ \dual{f_0}{x_0}=\norm{x_0}^2\bigr\}.
\end{equation*}
这个 (多值的) 映射 $x_0\mapsto F(x_0)$ 称为从 $E$ 到 $E^\star$ 的\emph{\index{duiouyingshe@对偶映射 (duality map)}
对偶映射};
它的某些性质在习题  \ref{exer:1.1},  \ref{exer:1.2},  \ref{exer:3.28} 与问题 13 中有描述.
\end{remark}

\begin{keycorollary}
\label{cor:1.4}
对每个 $x\in E$ 都有
\begin{equation}
\norm{x}=\sup_{\substack{f\in E^\star\\ \norm{f}\le 1}}\abs{\dual{f}{x}}
=\max_{\substack{f\in E^\star\\ \norm{f}\le 1}}\abs{\dual{f}{x}}.
  \label{eq:1.6}
\end{equation}
\end{keycorollary}

\begin{proof}
我们总可假设 $x\ne 0$.显然
\begin{equation*}
\sup_{\substack{f\in E^\star\\ \norm{f}\le 1}}\abs{\dual{f}{x}}\le\norm{x}.
\end{equation*}
另一方面, 由推论  \ref{cor:1.3} 知存在 $f_0\in E^\star$ 使得 $\norm{f_0}=\norm{x}$ 且
$\dual{f_0}{x}=\norm{x}^2$.置 $f_1=f_0/\norm{x}$, 于是 $\norm{f_1}=1$ 且
$\dual{f_1}{x}=\norm{x}$.
\end{proof}

\begin{remark}
\label{rem:1.3}
公式 \eqref{eq:1.5}——它是一个\emph{定义}——不应与公式 \eqref{eq:1.6} 混淆,
后者是一个\emph{命题}.一般地, \eqref{eq:1.5} 中的 ``sup'' 未必取到; 例如见习题
 \ref{exer:1.3}.然而, 若 $E$ 是自反 Banach 空间 (见第~
\ref{chap:3}~章), 则 \eqref{eq:1.5} 中的
``sup'' 可以取到; R.~C.~James 的一个深刻结果断言其逆亦真: 若 $E$ 是 Banach 空间,
且对每个 $f\in E^\star$, \eqref{eq:1.5} 中的 sup 都能取到, 则 $E$ 自反; 参见例如
J.~Diestel \brefcite{Diestel}{1}{, 第 1 章} 或 R.~Holmes \bref{Holmes}{1}.
\end{remark}

\section{Hahn--Banach 定理的几何形式: 凸集的分离}
\label{sec:1.2}

我们先给出关于超平面的一些预备事实.以下 $E$ 均指赋范向量空间.

\begin{definition}
\emph{仿射超平面}是指 $E$ 的如下形式的子集
\begin{equation*}
H=\{x\in E;\ f(x)=\alpha\},
\end{equation*}
其中 $f$ 是不恒为零的线性泛函\footnote{我们不假定 $f$ 连续 (在每个无穷维赋范空间中都存在
不连续的线性泛函; 见习题  \ref{exer:1.5}).}, $\alpha\in\R$ 为给定常数.
我们写 $H=[f=\alpha]$, 并称 $f=\alpha$ 是 $H$ 的方程.
\end{definition}

\begin{proposition}
\label{prop:1.5}
超平面 $H=[f=\alpha]$ 是闭的, 当且仅当 $f$ 连续.
\end{proposition}

\begin{proof}
显然, 若 $f$ 连续则 $H$ 闭.反之, 设 $H$ 闭.$H$ 的补集 $H^c$ 是开的且非空
(因为 $f$ 不恒为零).取 $x_0\in H^c$, 于是 $f(x_0)\ne\alpha$, 例如 $f(x_0)<\alpha$.
取 $r>0$ 使得 $B(x_0,r)\subset H^c$, 其中
\begin{equation*}
B(x_0,r)=\{x\in E;\ \norm{x-x_0}<r\}.
\end{equation*}
我们断言
\begin{equation}
f(x)<\alpha\quad \forall x\in B(x_0,r).
  \label{eq:1.7}
\end{equation}
事实上, 反设对某个 $x_1\in B(x_0,r)$ 有 $f(x_1)>\alpha$.线段
\begin{equation*}
\{x_t=(1-t)x_0+tx_1;\ t\in[0,1]\}
\end{equation*}
包含在 $B(x_0,r)$ 中, 从而 $f(x_t)\ne\alpha,\ \forall t\in[0,1]$; 另一方面, 对某个
$t\in[0,1]$ 有 $f(x_t)=\alpha$, 即 $t=\dfrac{f(x_1)-\alpha}{f(x_1)-f(x_0)}$,
矛盾, 故 \eqref{eq:1.7} 得证.由 \eqref{eq:1.7} 推出
\begin{equation*}
f(x_0+rz)<\alpha\quad \forall z\in B(0,1).
\end{equation*}
因此 $f$ 连续, 且 $\norm{f}\le\dfrac{1}{r}\bigl(\alpha-f(x_0)\bigr)$.
\end{proof}

\begin{definition}
设 $A,B$ 为 $E$ 的两个子集.若
\begin{equation*}
f(x)\le\alpha\quad \forall x\in A\quad\text{且}\quad f(x)\ge\alpha\quad \forall x\in B,
\end{equation*}
则称超平面 $H=[f=\alpha]$ \emph{分离} $A$ 与 $B$.

若存在 $\varepsilon>0$ 使得
\begin{equation*}
f(x)\le\alpha-\varepsilon\ \ \forall x\in A\quad\text{且}\quad
f(x)\ge\alpha+\varepsilon\ \ \forall x\in B,
\end{equation*}
则称 $H$ \emph{严格分离} $A$ 与 $B$.

从几何上说, 分离意味着 $A$ 位于 $H$ 所决定的半空间之一中, 而 $B$ 位于另一个半空间中;
见图  \ref{fig:1}.

最后回忆, 子集 $A\subset E$ 称为\emph{凸}的, 若
\begin{equation*}
tx+(1-t)y\in A\quad \forall x,y\in A,\quad \forall t\in[0,1].
\end{equation*}
\end{definition}

\index{tujifenli@凸集分离 (separation of convex sets)}
\begin{keytheorem}[Hahn--Banach, 第一几何形式]
\label{thm:1.6}
设 $A\subset E$ 与 $B\subset E$ 为两个非空凸子集, 且 $A\cap B=\emptyset$.
假设其中一个是开的.则存在分离 $A$ 与 $B$ 的闭超平面.
\end{keytheorem}

定理  \ref{thm:1.6} 的证明依赖于下面两条引理.

\begin{figure}[htbp]
	\centering
	\begin{tikzpicture}[scale=0.95]
		% 超平面 H (斜线)
		\draw (0.2,-0.6) -- (4.6,3.0) node[above right, inner sep=2pt] {$H$};
		% 凸集 A (H 左上方)
		\begin{scope}[shift={(1.05,2.35)}]
			\draw plot[smooth cycle, tension=0.75] coordinates {
				(-0.85,0.05) (-0.62,0.42) (-0.18,0.55) (0.16,0.36)
				(0.44,0.44) (0.72,0.16) (0.58,-0.16) (0.72,-0.42)
				(0.30,-0.58) (-0.10,-0.50) (-0.40,-0.62) (-0.70,-0.36)};
		\end{scope}
		\node at (1.05,2.35) {$A$};
		% 凸集 B (H 右下方)
		\begin{scope}[shift={(3.35,0.72)}]
			\draw plot[smooth cycle, tension=0.75] coordinates {
				(-0.62,0.30) (-0.30,0.48) (0.02,0.34) (0.26,0.50)
				(0.58,0.30) (0.50,-0.02) (0.66,-0.34) (0.30,-0.50)
				(-0.06,-0.40) (-0.36,-0.52) (-0.60,-0.18)};
		\end{scope}
		\node at (3.35,0.72) {$B$};
	\end{tikzpicture}
	\caption{}
	\label{fig:1}
\end{figure}

\begin{lemma}
\label{lem:1.2}
设 $C\subset E$ 为含 $0\in C$ 的开凸集.对每个 $x\in E$ 置
\begin{equation}
p(x)=\inf\{\alpha>0;\ \alpha^{-1}x\in C\}
  \label{eq:1.8}
\end{equation}
($p$ 称为 $C$ 的\emph{规范函数}(gauge) 或 $C$ 的 Minkowski 泛函).
则 $p$ 满足 \eqref{eq:1.1}, \eqref{eq:1.2} 以及下述性质:
\begin{align}
&\text{存在常数 } M \text{ 使得}\quad 0\le p(x)\le M\norm{x}\quad \forall x\in E,
  \label{eq:1.9}\\
&C=\{x\in E;\ p(x)<1\}.
  \label{eq:1.10}
\end{align}
\end{lemma}

\begin{proof}
显然 \eqref{eq:1.1} 成立.

\emph{\eqref{eq:1.9} 的证明.} 取 $r>0$ 使得 $B(0,r)\subset C$; 显然有
\begin{equation*}
p(x)\le\frac{1}{r}\norm{x}\quad \forall x\in E.
\end{equation*}

\emph{\eqref{eq:1.10} 的证明.} 先设 $x\in C$; 因 $C$ 开, 故当 $\varepsilon>0$ 充分小时
$(1+\varepsilon)x\in C$, 于是 $p(x)\le\dfrac{1}{1+\varepsilon}<1$.反之, 若 $p(x)<1$,
则存在 $\alpha\in(0,1)$ 使 $\alpha^{-1}x\in C$, 从而
$x=\alpha(\alpha^{-1}x)+(1-\alpha)0\in C$.

\emph{\eqref{eq:1.2} 的证明.} 设 $x,y\in E$, $\varepsilon>0$.利用 \eqref{eq:1.1} 与 \eqref{eq:1.10} 得
$\dfrac{x}{p(x)+\varepsilon}\in C$ 且 $\dfrac{y}{p(y)+\varepsilon}\in C$.于是
\begin{equation*}
t\frac{x}{p(x)+\varepsilon}+(1-t)\frac{y}{p(y)+\varepsilon}\in C
\quad \text{对一切}\ t\in[0,1].
\end{equation*}
取 $t=\dfrac{p(x)+\varepsilon}{p(x)+p(y)+2\varepsilon}$, 即得
$\dfrac{x+y}{p(x)+p(y)+2\varepsilon}\in C$.再次利用 \eqref{eq:1.1} 与 \eqref{eq:1.10}, 便推出
$p(x+y)<p(x)+p(y)+2\varepsilon,\ \forall\varepsilon>0$.
\end{proof}

\begin{lemma}
\label{lem:1.3}
设 $C\subset E$ 为非空开凸集, $x_0\in E$ 且 $x_0\notin C$.则存在 $f\in E^\star$
使得 $f(x)<f(x_0)\ \forall x\in C$.特别地, 超平面 $[f=f(x_0)]$ 分离 $\{x_0\}$ 与 $C$.
\end{lemma}

\begin{proof}
经平移, 我们总可假设 $0\in C$.于是可引入 $C$ 的规范函数 $p$ (见引理  \ref{lem:1.2}).
考虑线性子空间 $G=\R x_0$ 以及由
\begin{equation*}
g(tx_0)=t,\qquad t\in\R
\end{equation*}
定义的线性泛函 $g:G\to\R$.显然
\begin{equation*}
g(x)\le p(x)\quad \forall x\in G
\end{equation*}
(分 $t>0$ 与 $t\le 0$ 两种情形考虑).由定理  \ref{thm:1.1}, 存在 $E$ 上的线性泛函 $f$
延拓 $g$ 且满足
\begin{equation*}
f(x)\le p(x)\quad \forall x\in E.
\end{equation*}
特别地, $f(x_0)=1$, 且由 \eqref{eq:1.9} 知 $f$ 连续.由 \eqref{eq:1.10} 推出对每个 $x\in C$ 有 $f(x)<1$.
\end{proof}

\begin{proof}[定理 \ref{thm:1.6}]
置 $C=A-B$, 于是 $C$ 凸 (请验证!), $C$ 开 (因为 $C=\bigcup_{y\in B}(A-y)$),
且 $0\notin C$ (因为 $A\cap B=\emptyset$).由引理  \ref{lem:1.3}, 存在 $f\in E^\star$
使得
\begin{equation*}
f(z)<0\quad \forall z\in C,
\end{equation*}
即
\begin{equation*}
f(x)<f(y)\quad \forall x\in A,\ \forall y\in B.
\end{equation*}
取常数 $\alpha$ 满足
\begin{equation*}
\sup_{x\in A}f(x)\le\alpha\le\inf_{y\in B}f(y).
\end{equation*}
显然, 超平面 $[f=\alpha]$ 分离 $A$ 与 $B$.
\end{proof}

\begin{keytheorem}[Hahn--Banach, 第二几何形式]
\label{thm:1.7}
设 $A\subset E$ 与 $B\subset E$ 为两个非空凸子集, 且 $A\cap B=\emptyset$.
假设 $A$ 闭而 $B$ 紧.则存在严格分离 $A$ 与 $B$ 的闭超平面.
\end{keytheorem}

\begin{proof}
置 $C=A-B$, 于是 $C$ 凸, 闭 (请验证!), 且 $0\notin C$.因此存在 $r>0$ 使得
$B(0,r)\cap C=\emptyset$.由定理  \ref{thm:1.6}, 存在分离 $B(0,r)$ 与 $C$ 的闭超平面.
因此存在 $f\in E^\star$, $f\not\equiv 0$, 使得
\begin{equation*}
f(x-y)\le f(rz)\quad \forall x\in A,\ \forall y\in B,\ \forall z\in B(0,1).
\end{equation*}
于是 $f(x-y)\le-r\norm{f}\ \ \forall x\in A,\ \forall y\in B$.置
$\varepsilon=\tfrac{1}{2}r\norm{f}>0$, 便得到
\begin{equation*}
f(x)+\varepsilon\le f(y)-\varepsilon\quad \forall x\in A,\ \forall y\in B.
\end{equation*}
取 $\alpha$ 使得
\begin{equation*}
\sup_{x\in A}f(x)+\varepsilon\le\alpha\le\inf_{y\in B}f(y)-\varepsilon,
\end{equation*}
即知超平面 $[f=\alpha]$ 严格分离 $A$ 与 $B$.
\end{proof}

\begin{remark}
\label{rem:1.4}
设 $A\subset E$ 与 $B\subset E$ 为两个非空凸集, 且 $A\cap B=\emptyset$.
若不作进一步的假设, 一般不可能用闭超平面分离 $A$ 与 $B$.甚至可以构造 $A$ 与 $B$
都闭的这样的例子 (见习题  \ref{exer:1.14}).然而, 若 $E$ 是有限维的, 则总能分离任意两个
非空凸集 $A$ 与 $B$, 只要 $A\cap B=\emptyset$ (不需要任何进一步的假设!); 见习题
 \ref{exer:1.9}.
\end{remark}

我们用一条非常有用的结论来结束本节:

\begin{keycorollary}
\label{cor:1.8}
设 $F\subset E$ 为线性子空间且 $\overline{F}\ne E$.则存在 $f\in E^\star$, $f\not\equiv 0$,
使得
\begin{equation*}
\dual{f}{x}=0\quad \forall x\in F.
\end{equation*}
\end{keycorollary}

\begin{proof}
取 $x_0\in E$ 且 $x_0\notin\overline{F}$.在定理  \ref{thm:1.7} 中取
$A=\overline{F}$, $B=\{x_0\}$, 得到严格分离 $\overline{F}$ 与 $\{x_0\}$ 的闭超平面
$[f=\alpha]$.于是
\begin{equation*}
\dual{f}{x}<\alpha<\dual{f}{x_0}\quad \forall x\in\overline{F}.
\end{equation*}
由此推出 $\dual{f}{x}=0\ \forall x\in F$, 因为对每个 $\lambda\in\R$ 都有
$\lambda\dual{f}{x}<\alpha$.
\end{proof}

\begin{keyremark}
\label{rem:1.5}
推论  \ref{cor:1.8} 在证明线性子空间 $F\subset E$ 稠密时经常被用到.
只需证明 $E$ 上每个在 $F$ 上为零的连续线性泛函必在整个 $E$ 上为零.
\end{keyremark}

\section{双对偶 $E^{\star\star}$. 正交关系}
\label{sec:1.3}

设 $E$ 为赋范向量空间, $E^\star$ 为其对偶空间, 范数为
\begin{equation*}
\norm{f}_{E^\star}=\sup_{\substack{x\in E\ \norm{x}\le 1}}\abs{\dual{f}{x}}.
\end{equation*}
\emph{双对偶} $E^{\star\star}$ 是 $E^\star$ 的对偶, 其范数为
\begin{equation*}
\norm{\xi}_{E^{\star\star}}=\sup_{\substack{f\in E^\star\ \norm{f}\le 1}}\abs{\dual{\xi}{f}}
\qquad (\xi\in E^{\star\star}).
\end{equation*}
存在\emph{\index{dianfandanshe@典范单射 (canonical injection)}
典范单射} $J:E\to E^{\star\star}$, 定义如下: 给定 $x\in E$, 映射
$f\mapsto\dual{f}{x}$ 是 $E^\star$ 上的连续线性泛函; 因此它是 $E^{\star\star}$ 中的一个元素,
我们把它记为 $Jx$.\footnote{$J$ 不应与注记 \ref{rem:1.2} 中定义的对偶映射
$F:E\to E^\star$ 相混淆.} 我们有
\begin{equation*}
\dual{Jx}{f}_{E^{\star\star},E^\star}=\dual{f}{x}_{E^\star,E}
\qquad \forall x\in E,\ \forall f\in E^\star.
\end{equation*}
显然 $J$ 是线性的, 且 $J$ 是等距的, 即 $\norm{Jx}_{E^{\star\star}}=\norm{x}_{E}$;
事实上, 我们有
\begin{equation*}
\norm{Jx}_{E^{\star\star}}=\sup_{\substack{f\in E^\star\ \norm{f}\le 1}}\abs{\dual{Jx}{f}}
=\sup_{\substack{f\in E^\star\ \norm{f}\le 1}}\abs{\dual{f}{x}}=\norm{x}
\end{equation*}
(由推论  \ref{cor:1.4}).

可能发生 $J$ 不是从 $E$ 到 $E^{\star\star}$ 的满射的情形 (见第 3, 4 章).
然而, 借助 $J$ 把 $E$ 与 $E^{\star\star}$ 的一个子空间等同起来, 是方便的.
若 $J$ 恰为满射, 则称 $E$ 是\emph{自反}的, 并把 $E^{\star\star}$ 与 $E$ 等同 (见第~
\ref{chap:3}~章).

\begin{notation}
若 $M\subset E$ 为线性子空间, 我们置
\begin{equation*}
M^\perp=\{f\in E^\star;\ \dual{f}{x}=0\ \ \forall x\in M\}.
\end{equation*}
若 $N\subset E^\star$ 为线性子空间, 我们置
\begin{equation*}
N^\perp=\{x\in E;\ \dual{f}{x}=0\ \ \forall f\in N\}.
\end{equation*}
注意——按定义——$N^\perp$ 是 $E$ 的子集而非 $E^{\star\star}$ 的子集.显然 $M^\perp$
(相应地 $N^\perp$) 是 $E^\star$ (相应地 $E$) 的闭线性子空间.我们说 $M^\perp$
(相应地 $N^\perp$) 是与 $M$ (相应地 $N$) \emph{正交}的空间.
\end{notation}

\begin{proposition}
\label{prop:1.9}
设 $M\subset E$ 为线性子空间.则
\begin{equation*}
(M^\perp)^\perp=\overline{M}.
\end{equation*}
设 $N\subset E^\star$ 为线性子空间.则
\begin{equation*}
(N^\perp)^\perp\supset\overline{N}.
\end{equation*}
\end{proposition}

\begin{proof}
显然 $M\subset(M^\perp)^\perp$, 且因 $(M^\perp)^\perp$ 闭, 故
$\overline{M}\subset(M^\perp)^\perp$.反之, 我们证明 $(M^\perp)^\perp\subset\overline{M}$.
反设存在 $x_0\in(M^\perp)^\perp$ 使得 $x_0\notin\overline{M}$.由定理  \ref{thm:1.7},
存在严格分离 $\{x_0\}$ 与 $\overline{M}$ 的闭超平面.于是存在 $f\in E^\star$ 与
$\alpha\in\R$ 使得
\begin{equation*}
\dual{f}{x}<\alpha<\dual{f}{x_0}\quad \forall x\in M.
\end{equation*}
因 $M$ 是线性空间, 故推出 $\dual{f}{x}=0\ \forall x\in M$, 且 $\dual{f}{x_0}>0$.
因此 $f\in M^\perp$, 从而 $\dual{f}{x_0}=0$, 矛盾.

同样显然 $N\subset(N^\perp)^\perp$, 从而 $\overline{N}\subset(N^\perp)^\perp$.
\end{proof}

\begin{remark}
\label{rem:1.6}
可能发生 $(N^\perp)^\perp$ 严格大于 $\overline{N}$ 的情形 (见习题  \ref{exer:1.16}).
不过, 试着去证明 $(N^\perp)^\perp=\overline{N}$ 并看清论证在哪里断裂, 是有启发的.
设 $f_0\in E^\star$ 使得 $f_0\in(N^\perp)^\perp$ 且 $f_0\notin\overline{N}$.
在 $E^\star$ 中应用 Hahn--Banach 定理, 我们可以严格分离 $\{f_0\}$ 与 $\overline{N}$.
于是存在 $\xi\in E^{\star\star}$ 使得 $\dual{\xi}{f_0}>0$.但我们无法导出矛盾, 因为
$\xi\notin N^\perp$——除非我们碰巧 (纯属巧合!) 知道 $\xi\in E$, 更确切地说
$\xi=Jx_0$ 对某个 $x_0\in E$ 成立.特别地, 若 $E$ 自反, 则 $(N^\perp)^\perp=\overline{N}$
确实成立.在一般情形下可以证明, $(N^\perp)^\perp$ 与 $N$ 在\wstar{}拓扑
$\sigma(E^\star,E)$ 下的闭包重合 (见第~
\ref{chap:3}~章).
\end{remark}

\section{共轭凸函数理论速览}
\label{sec:1.4}

我们从关于下半连续函数与凸函数的一些基本事实开始.本节考虑定义在集合 $E$ 上, 取值于 $(-\infty,+\infty]$ 的函数 $\varphi$, 于是 $\varphi$ 可以取到值 $+\infty$
(但 $-\infty$ 被排除).记 $D(\varphi)$ 为 $\varphi$ 的\emph{定义域}, 即
\begin{equation*}
D(\varphi)=\{x\in E;\ \varphi(x)<+\infty\}.
\end{equation*}

\begin{notation}
$\varphi$ 的\emph{\index{shangtu@上图 (epigraph)}
上图}是集合\footnote{我们强调 $\R=(-\infty,\infty)$, 因此 $\lambda$
不取到值 $\infty$.}
\begin{equation*}
\operatorname{epi}\varphi=\{[x,\lambda]\in E\times\R;\ \varphi(x)\le\lambda\}.
\end{equation*}
\end{notation}

现在设 $E$ 为拓扑空间.我们回忆下述内容.

\begin{definition}
\index{xiabanlianxu@下半连续 (lower semicontinuous)}
若对每个 $\lambda\in\R$, 集合
\begin{equation*}
[\varphi\le\lambda]=\{x\in E;\ \varphi(x)\le\lambda\}
\end{equation*}
是闭的, 则称函数 $\varphi:E\to(-\infty,+\infty]$ 是\emph{下半连续}的 (l.s.c.).
\end{definition}

关于下半连续函数, 有如下一些众所周知的基本事实 (可参见例如 G.~Choquet
\bref{Choquet}{1}, J.~Dixmier \bref{Dixmier}{1}, J.~R.~Munkres \bref{Munkres}{1}, H.~L.~Royden \bref{Royden}{1}):

\begin{enumerate}[label=\arabic*., leftmargin=2.2em, itemsep=0.25em]
	\item 若 $\varphi$ 下半连续, 则 $\operatorname{epi}\varphi$ 在 $E\times\R$ 中闭;
	      反之亦然.
	\item 若 $\varphi$ 下半连续, 则对每个 $x\in E$ 与每个 $\varepsilon>0$, 存在 $x$ 的邻域
	      $V$ 使得
	      \begin{equation*}
	      \varphi(y)\ge\varphi(x)-\varepsilon\quad \forall y\in V;
	      \end{equation*}
	      反之亦然.

	      特别地, 若 $\varphi$ 下半连续, 则对 $E$ 中每个满足 $x_n\to x$ 的序列 $(x_n)$, 有
	      \begin{equation*}
	      \liminf_{n\to\infty}\varphi(x_n)\ge\varphi(x);
	      \end{equation*}
	      反之, 若 $E$ 是度量空间, 上述命题的逆也成立.
	\item 若 $\varphi_1$ 与 $\varphi_2$ 下半连续, 则 $\varphi_1+\varphi_2$ 下半连续.
	\item 若 $(\varphi_i)_{i\in I}$ 是一族下半连续函数, 则它们的上包络也下半连续, 即由
	      \begin{equation*}
	      \varphi(x)=\sup_{i\in I}\varphi_i(x)
	      \end{equation*}
	      定义的函数 $\varphi$ 下半连续.
	\item 若 $E$ 紧且 $\varphi$ 下半连续, 则 $\inf_E\varphi$ 可以取到.

	      (若 $E$ 是紧度量空间, 可以用极小化序列来论证.对一般的紧拓扑空间,
	      对适当取值的 $\lambda$ 考虑集合 $[\varphi\le\lambda]$.)
\end{enumerate}

现在设 $E$ 为向量空间.回忆下述定义.

\begin{definition}
若
\begin{equation*}
\varphi(tx+(1-t)y)\le t\varphi(x)+(1-t)\varphi(y)
\qquad \forall x,y\in E,\ \forall t\in(0,1),
\end{equation*}
则称函数 $\varphi:E\to(-\infty,+\infty]$ 是\emph{凸}的.
\end{definition}

我们将用到凸函数的一些基本性质:

\begin{enumerate}[label=\arabic*., leftmargin=2.2em, itemsep=0.25em]
	\item 若 $\varphi$ 是凸函数, 则 $\operatorname{epi}\varphi$ 是 $E\times\R$ 中的凸集;
	      反之亦然.
	\item 若 $\varphi$ 是凸函数, 则对每个 $\lambda\in\R$, 集合 $[\varphi\le\lambda]$ 是凸的;
	      但逆命题不成立.
	\item 若 $\varphi_1$ 与 $\varphi_2$ 凸, 则 $\varphi_1+\varphi_2$ 凸.
	\item 若 $(\varphi_i)_{i\in I}$ 是一族凸函数, 则上包络 $\sup_i\varphi_i$ 是凸的.
\end{enumerate}

以下我们假设 $E$ 是赋范向量空间.

\begin{definition}
设 $\varphi:E\to(-\infty,+\infty]$ 满足 $\varphi\not\equiv+\infty$
(即 $D(\varphi)\ne\emptyset$).定义\emph{\index{gongehanshu@共轭函数 (conjugate function)}
共轭函数}
$\varphi^\star:E^\star\to(-\infty,+\infty]$ 为\footnote{$\varphi^\star$ 有时称为 $\varphi$
的\emph{Legendre 变换}.}
\begin{equation*}
\varphi^\star(f)=\sup_{x\in E}\{\dual{f}{x}-\varphi(x)\}\qquad (f\in E^\star).
\end{equation*}
注意 $\varphi^\star$ 在 $E^\star$ 上凸且下半连续.事实上, 对每个固定的 $x\in E$,
函数 $f\mapsto\dual{f}{x}-\varphi(x)$ 在 $E^\star$ 上凸且连续 (从而下半连续).
于是这些函数 (当 $x$ 遍历 $E$ 时) 的上包络凸且下半连续.
\end{definition}

\begin{remark}
\label{rem:1.7}
显然有不等式
\begin{equation}
\dual{f}{x}\le\varphi(x)+\varphi^\star(f)\qquad \forall x\in E,\ \forall f\in E^\star,
  \label{eq:1.11}
\end{equation}
它有时称为\emph{Young 不等式}.当然, 按我们对 $\varphi^\star$ 的定义, 这一事实是显然的!
Young 不等式的经典形式 (见第~
\ref{chap:4}~章定理~
\ref{thm:4.6} 的证明) 断言
\begin{equation}
ab\le\frac{1}{p}a^p+\frac{1}{p'}b^{p'}\qquad \forall a,b\ge 0
  \label{eq:1.12}
\end{equation}
其中 $1<p<\infty$ 且 $\dfrac{1}{p}+\dfrac{1}{p'}=1$.不等式 \eqref{eq:1.12} 是
\eqref{eq:1.11} 的一个特例, 取 $E=E^\star=\R$ 且
$\varphi(t)=\dfrac{1}{p}\abs{t}^p$, $\varphi^\star(s)=\dfrac{1}{p'}\abs{s}^{p'}$
(见习题  \ref{exer:1.18}, 第 (h) 问).
\end{remark}

\begin{figure}[htbp]
	\centering
	\begin{tikzpicture}[scale=1.0]
		% 坐标轴
		\draw[-{Stealth[length=2mm]}] (-0.35,0) -- (6.3,0) node[right] {$E$};
		\draw[-{Stealth[length=2mm]}] (0,-0.4) -- (0,4.5) node[above] {$\R$};
		% 上图 A = epi phi (抛物体)
		\fill[gray!25] plot[domain=1.05:4.45, samples=60, smooth]
			(\x, {0.42*(\x-2.75)^2+1.35}) -- (4.45,4.35) -- (1.05,4.35) -- cycle;
		\draw[thick] plot[domain=1.05:4.45, samples=60, smooth]
			(\x, {0.42*(\x-2.75)^2+1.35});
		\node at (2.75,3.55) {$A=\operatorname{epi}\varphi$};
		% 水平阴影线
		\foreach \yy in {1.62,1.78,1.94,2.10,2.26,2.42}{
			\draw[thin] (1.35,\yy) -- (4.15,\yy);
		}
		% 分离超平面 H
		\draw[thick] (0.55,0.75) -- (6.0,3.95) node[right] {$H$};
		% 点 B = [x0, lambda0]
		\fill (2.55,0.72) circle (1.6pt);
		\node[right] at (2.62,0.72) {$[x_0,\lambda_0]=B$};
		\draw (2.55,0.02) -- (2.55,-0.10) node[below] {$x_0$};
		\draw (-0.10,0.72) -- (0.10,0.72) node[left] {$\lambda_0$};
	\end{tikzpicture}
	\caption{}
	\label{fig:2}
\end{figure}

\begin{proposition}
\label{prop:1.10}
设 $\varphi:E\to(-\infty,+\infty]$ 凸且下半连续, 且 $\varphi\not\equiv+\infty$.
则 $\varphi^\star\not\equiv+\infty$; 特别地, $\varphi$ 被某个仿射连续函数从下方界住.
\end{proposition}

\begin{proof}
取 $x_0\in D(\varphi)$ 与 $\lambda_0<\varphi(x_0)$.在空间 $E\times\R$ 中应用定理
 \ref{thm:1.7} (Hahn--Banach, 第二几何形式), 取 $A=\operatorname{epi}\varphi$ 与
$B=\{[x_0,\lambda_0]\}$.于是存在 $E\times\R$ 中的闭超平面 $H=[\Phi=\alpha]$,
它严格分离 $A$ 与 $B$; 见图  \ref{fig:2}.注意函数 $x\in E\mapsto\Phi([x,0])$ 是 $E$ 上的
连续线性泛函, 故对某个 $f\in E^\star$ 有 $\Phi([x,0])=\dual{f}{x}$.令
$k=\Phi([0,1])$, 则
\begin{equation*}
\Phi([x,\lambda])=\dual{f}{x}+k\lambda\qquad \forall[x,\lambda]\in E\times\R.
\end{equation*}
写出 $\Phi>\alpha$ 于 $A$ 上以及 $\Phi<\alpha$ 于 $B$ 上, 即得
\begin{equation*}
\dual{f}{x}+k\lambda>\alpha,\qquad \forall[x,\lambda]\in\operatorname{epi}\varphi,
\end{equation*}
以及
\begin{equation*}
\dual{f}{x_0}+k\lambda_0<\alpha.
\end{equation*}
特别地, 我们有
\begin{equation}
\dual{f}{x}+k\varphi(x)>\alpha\qquad \forall x\in D(\varphi)
  \label{eq:1.13}
\end{equation}
从而
\begin{equation*}
\dual{f}{x_0}+k\varphi(x_0)>\alpha>\dual{f}{x_0}+k\lambda_0.
\end{equation*}
由此推出 $k>0$.由 \eqref{eq:1.13} 我们有
\begin{equation*}
\left\langle-\frac{1}{k}f,x\right\rangle-\varphi(x)<-\frac{\alpha}{k}
\qquad \forall x\in D(\varphi),
\end{equation*}
因此 $\varphi^\star\!\left(-\dfrac{1}{k}f\right)<+\infty$.
\end{proof}

若迭代 $\star$ 运算, 我们得到定义在 $E^{\star\star}$ 上的函数 $\varphi^{\star\star}$.
但我们选择把 $\varphi^{\star\star}$ 限制到 $E$ 上, 即定义
\begin{equation*}
\varphi^{\star\star}(x)=\sup_{f\in E^\star}\{\dual{f}{x}-\varphi^\star(f)\}
\qquad (x\in E).
\end{equation*}

\index{FenchelMoreaudingli@Fenchel--Moreau 定理 (Fenchel--Moreau theorem)}
\begin{keytheorem}[Fenchel--Moreau]
\label{thm:1.11}
设 $\varphi:E\to(-\infty,+\infty]$ 凸, 下半连续, 且 $\varphi\not\equiv+\infty$.
则 $\varphi^{\star\star}=\varphi$.
\end{keytheorem}

\begin{proof}
分两步进行:

\emph{第一步:} 我们附加假设 $\varphi\ge 0$, 并断言 $\varphi^{\star\star}=\varphi$.

首先, 显然 $\varphi^{\star\star}\le\varphi$, 因为
$\dual{f}{x}-\varphi^\star(f)\le\varphi(x)\ \forall x\in E$ 且 $\forall f\in E^\star$.
为证明 $\varphi^{\star\star}=\varphi$, 我们用反证法, 设对某个 $x_0\in E$ 有
$\varphi^{\star\star}(x_0)<\varphi(x_0)$.可能 $\varphi(x_0)=+\infty$, 但
$\varphi^{\star\star}(x_0)$ 总是有限的.在空间 $E\times\R$ 中应用定理  \ref{thm:1.7}
(Hahn--Banach, 第二几何形式), 取 $A=\operatorname{epi}\varphi$ 与
$B=[x_0,\varphi^{\star\star}(x_0)]$.于是, 如同命题  \ref{prop:1.10} 的证明中那样,
存在 $f\in E^\star$, $k\in\R$ 与 $\alpha\in\R$ 使得
\begin{align}
\dual{f}{x}+k\lambda &> \alpha && \forall[x,\lambda]\in\operatorname{epi}\varphi,
  \label{eq:1.14}\\
\dual{f}{x_0}+k\varphi^{\star\star}(x_0) &< \alpha.
  \label{eq:1.15}
\end{align}
由此推出 $k\ge 0$ (固定 $x\in D(\varphi)$ 并令 $\lambda\to+\infty$ 于 \eqref{eq:1.14}).
[这里不能像命题  \ref{prop:1.10} 的证明中那样断言 $k>0$; 可能 $k=0$, 这对应于
$E\times\R$ 中一个``竖直的''超平面 $H$.]

设 $\varepsilon>0$; 因 $\varphi\ge 0$, 由 \eqref{eq:1.14} 有
\begin{equation*}
\dual{f}{x}+(k+\varepsilon)\varphi(x)\ge\alpha\qquad \forall x\in D(\varphi).
\end{equation*}
因此
\begin{equation*}
\varphi^\star\!\left(-\frac{f}{k+\varepsilon}\right)\le-\frac{\alpha}{k+\varepsilon}.
\end{equation*}
由 $\varphi^{\star\star}(x_0)$ 的定义,
\begin{equation*}
\varphi^{\star\star}(x_0)\ge\left\langle-\frac{f}{k+\varepsilon},x_0\right\rangle
-\varphi^\star\!\left(-\frac{f}{k+\varepsilon}\right)
\ge\left\langle-\frac{f}{k+\varepsilon},x_0\right\rangle+\frac{\alpha}{k+\varepsilon}.
\end{equation*}
于是
\begin{equation*}
\dual{f}{x_0}+(k+\varepsilon)\varphi^{\star\star}(x_0)\ge\alpha\qquad \forall\varepsilon>0,
\end{equation*}
这与 \eqref{eq:1.15} 矛盾.

\emph{第二步:} 一般情形.

固定某个 $f_0\in D(\varphi^\star)$ (由命题  \ref{prop:1.10} 知
$D(\varphi^\star)\ne\emptyset$), 并定义
\begin{equation*}
\tilde\varphi(x)=\varphi(x)-\dual{f_0}{x}+\varphi^\star(f_0),
\end{equation*}
于是 $\tilde\varphi$ 凸, 下半连续, $\tilde\varphi\not\equiv+\infty$, 且
$\tilde\varphi\ge 0$.由第一步知 $(\tilde\varphi)^{\star\star}=\tilde\varphi$.
现在计算 $(\tilde\varphi)^\star$ 与 $(\tilde\varphi)^{\star\star}$.我们有
\begin{equation*}
(\tilde\varphi)^\star(f)=\varphi^\star(f+f_0)-\varphi^\star(f_0)
\end{equation*}
以及
\begin{equation*}
(\tilde\varphi)^{\star\star}(x)=\varphi^{\star\star}(x)-\dual{f_0}{x}+\varphi^\star(f_0).
\end{equation*}
写出 $(\tilde\varphi)^{\star\star}=\tilde\varphi$, 即得 $\varphi^{\star\star}=\varphi$.
\end{proof}

让我们考察一些例子.

\begin{example}
考虑 $\varphi(x)=\norm{x}$.容易验证
\begin{equation*}
\varphi^\star(f)=
\begin{cases}
0 & \text{若}\ \norm{f}\le 1,\\
+\infty & \text{若}\ \norm{f}>1.
\end{cases}
\end{equation*}
于是
\begin{equation*}
\varphi^{\star\star}(x)=\sup_{\substack{f\in E^\star\\ \norm{f}\le 1}}\dual{f}{x}.
\end{equation*}
写出等式 $\varphi^{\star\star}=\varphi$, 我们再次得到推论  \ref{cor:1.4} 的一部分.
\end{example}

\begin{example}
给定非空集合 $K\subset E$, 置
\begin{equation*}
I_K(x)=
\begin{cases}
0 & \text{若}\ x\in K,\\
+\infty & \text{若}\ x\notin K.
\end{cases}
\end{equation*}
函数 $I_K$ 称为 $K$ 的\emph{\index{shixinghanshu@示性函数 (indicator function)}
示性函数}(indicator function)
(不应与 $K$ 的\emph{特征函数} $\chi_K$ 相混淆, 后者在 $K$ 上为 $1$, 在 $K$ 外为 $0$).
注意 $I_K$ 是凸函数当且仅当 $K$ 是凸集, $I_K$ 下半连续当且仅当 $K$ 闭.
共轭函数 $(I_K)^\star$ 称为 $K$ 的\emph{\index{zhichenghanshu@支撑函数 (supporting function)}
支撑函数}.

容易看出, 若 $K=M$ 是线性子空间, 则 $(I_M)^\star=I_{M^\perp}$ 且
$(I_M)^{\star\star}=I_{(M^\perp)^\perp}$.假设 $M$ 是闭线性空间, 写出
$(I_M)^{\star\star}=I_M$, 即得 $(M^\perp)^\perp=M$.在某种意义上,
定理  \ref{thm:1.11} 可以看作命题  \ref{prop:1.9} 的对偶.
\end{example}

我们用共轭函数的另一条有用性质来结束本章.

\begin{startheorem}[Fenchel--Rockafellar]
\label{thm:1.12}
设 $\varphi,\psi:E\to(-\infty,+\infty]$ 为两个凸函数.假设存在
$x_0\in D(\varphi)\cap D(\psi)$ 使得 $\varphi$ 在 $x_0$ 处连续.则
\begin{align*}
\inf_{x\in E}\{\varphi(x)+\psi(x)\}
&= \sup_{f\in E^\star}\{-\varphi^\star(-f)-\psi^\star(f)\}\\
&= \max_{f\in E^\star}\{-\varphi^\star(-f)-\psi^\star(f)\}
 = -\min_{f\in E^\star}\{\varphi^\star(-f)+\psi^\star(f)\}.
\end{align*}
\end{startheorem}

定理  \ref{thm:1.12} 的证明依赖于下述引理.

\begin{lemma}
\label{lem:1.4}
设 $C\subset E$ 为凸集, 则 $\Int C$ 是凸的.\footnote{如通常那样, $\Int C$ 表示 $C$
的内部.} 此外, 若 $\Int C\ne\emptyset$, 则
\begin{equation*}
\overline{C}=\overline{\Int C}.
\end{equation*}
\end{lemma}

引理  \ref{lem:1.4} 的证明可参见例如习题  \ref{exer:1.7}.

\begin{proof}[定理 \ref{thm:1.12}]
置
\begin{align*}
a &= \inf_{x\in E}\{\varphi(x)+\psi(x)\},\\
b &= \sup_{f\in E^\star}\{-\varphi^\star(-f)-\psi^\star(f)\}.
\end{align*}
显然 $b\le a$.若 $a=-\infty$, 则定理  \ref{thm:1.12} 的结论显然成立.
故以下可假设 $a\in\R$.置 $C=\operatorname{epi}\varphi$, 于是 $\Int C\ne\emptyset$
(因为 $\varphi$ 在 $x_0$ 处连续).我们在定理  \ref{thm:1.6} (Hahn--Banach,
第一几何形式) 中取 $A=\Int C$ 与
\begin{equation*}
B=\{[x,\lambda]\in E\times\R;\ \lambda\le a-\psi(x)\}.
\end{equation*}
于是 $A$ 与 $B$ 是非空凸集.此外 $A\cap B=\emptyset$; 事实上, 若
$[x,\lambda]\in A$, 则 $\lambda>\varphi(x)$; 另一方面 $\varphi(x)\ge a-\psi(x)$
(由 $a$ 的定义), 故 $[x,\lambda]\notin B$.

于是存在分离 $A$ 与 $B$ 的闭超平面 $H$.由此推出 $H$ 也分离 $\overline{A}$ 与 $B$.
但由引理  \ref{lem:1.4} 知 $\overline{A}=\overline{C}$.因此存在 $f\in E^\star$, $k\in\R$ 与 $\alpha\in\R$, 使得 $E\times\R$ 中的超平面 $H=[\Phi=\alpha]$ 分离
$\overline{C}$ 与 $B$, 其中
\begin{equation*}
\Phi([x,\lambda])=\dual{f}{x}+k\lambda\qquad \forall[x,\lambda]\in E\times\R.
\end{equation*}
于是有
\begin{align}
\dual{f}{x}+k\lambda &\ge \alpha && \forall[x,\lambda]\in\overline{C},
  \label{eq:1.16}\\
\dual{f}{x}+k\lambda &\le \alpha && \forall[x,\lambda]\in B.
  \label{eq:1.17}
\end{align}

取 $x=x_0$ 并令 $\lambda\to+\infty$ 于 \eqref{eq:1.16}, 可见 $k\ge 0$.我们断言
\begin{equation}
k>0.
  \label{eq:1.18}
\end{equation}
反设 $k=0$; 由此推出 $\norm{f}\ne 0$ (因为 $\Phi\not\equiv 0$).由 \eqref{eq:1.16}
与 \eqref{eq:1.17} 我们有
\begin{equation*}
\dual{f}{x}\ge\alpha\ \ \forall x\in D(\varphi),\qquad
\dual{f}{x}\le\alpha\ \ \forall x\in D(\psi).
\end{equation*}
但对某个充分小的 $\varepsilon_0>0$ 有 $B(x_0,\varepsilon_0)\subset D(\varphi)$, 故
\begin{equation*}
\dual{f}{x_0+\varepsilon_0z}\ge\alpha\qquad \forall z\in B(0,1),
\end{equation*}
这蕴含 $\dual{f}{x_0}\ge\alpha+\varepsilon_0\norm{f}$.另一方面, 因
$x_0\in D(\psi)$, 有 $\dual{f}{x_0}\le\alpha$; 于是得到 $\norm{f}=0$, 矛盾,
\eqref{eq:1.18} 获证.

由 \eqref{eq:1.16} 与 \eqref{eq:1.17} 得
\begin{equation*}
\varphi^\star\!\left(-\frac{f}{k}\right)\le-\frac{\alpha}{k}
\qquad\text{以及}\qquad
\psi^\star\!\left(\frac{f}{k}\right)\le\frac{\alpha}{k}-a,
\end{equation*}
从而
\begin{equation*}
-\varphi^\star\!\left(-\frac{f}{k}\right)-\psi^\star\!\left(\frac{f}{k}\right)\ge a.
\end{equation*}
另一方面, 由 $b$ 的定义有
\begin{equation*}
-\varphi^\star\!\left(-\frac{f}{k}\right)-\psi^\star\!\left(\frac{f}{k}\right)\le b.
\end{equation*}
于是得到
\begin{equation*}
a=b=-\varphi^\star\!\left(-\frac{f}{k}\right)-\psi^\star\!\left(\frac{f}{k}\right).
\end{equation*}
\end{proof}

\begin{example}
设 $K$ 为非空凸集.我们断言, 对每个 $x_0\in E$ 有
\begin{equation}
\dist(x_0,K)=\inf_{x\in K}\norm{x-x_0}
=\max_{\substack{f\in E^\star\\ \norm{f}\le 1}}\{\dual{f}{x_0}-I_K^\star(f)\}.
  \label{eq:1.19}
\end{equation}
事实上, 我们有
\begin{equation*}
\inf_{x\in K}\norm{x-x_0}=\inf_{x\in E}\{\varphi(x)+\psi(x)\},
\end{equation*}
其中 $\varphi(x)=\norm{x-x_0}$, $\psi(x)=I_K(x)$.应用定理  \ref{thm:1.12}, 即得
\eqref{eq:1.19}.在 $K=M$ 为线性子空间这一特殊情形, 我们得到关系式
\begin{equation*}
\dist(x_0,M)=\inf_{x\in M}\norm{x-x_0}
=\max_{\substack{f\in M^\perp\\ \norm{f}\le 1}}\dual{f}{x_0}.
\end{equation*}
\end{example}

\begin{remark}
\label{rem:1.8}
在 $\inf_{x\in K}\norm{x-x_0}$ 取不到的情形, 关系式 \eqref{eq:1.19} 可能为我们提供
一些有用的信息 (例如见习题  \ref{exer:1.17}).极小曲面理论提供了一个有趣的场景:
其中原问题 (即 $\inf_{x\in E}\{\varphi(x)+\psi(x)\}$) 不必有解, 而对偶问题
(即 $\max_{f\in E^\star}\{-\varphi^\star(-f)-\psi^\star(f)\}$) 有解; 见
I.~Ekeland--R.~Temam \bref{EkelandTemam}{1}.
\end{remark}

\begin{example}
设 $\varphi:E\to\R$ 凸且连续, $M\subset E$ 为线性子空间.则
\begin{equation*}
\inf_{x\in M}\varphi(x)=-\min_{f\in M^\perp}\varphi^\star(f).
\end{equation*}
只需在定理  \ref{thm:1.12} 中取 $\psi=I_M$.
\end{example}

\section*{第 1 章评注}
\addcontentsline{toc}{section}{第 1 章评注}
\markright{第 1 章评注}

\cmtitem{1. Hahn--Banach 定理的推广与变形.}
Hahn--Banach 定理的第一几何形式 (定理  \ref{thm:1.6}) 在一般拓扑向量空间中仍然成立.
第二几何形式 (定理  \ref{thm:1.7}) 在局部凸空间中成立——这类空间起着重要的作用,
例如在\emph{广义函数理论}中 (参见例如 L.~Schwartz \bref{Schwartz}{1} 与
F.~Treves \bref{Treves}{1}).有兴趣的读者可参阅例如 N.~Bourbaki \bref{Bourbaki}{1}, J.~Kelley--I.~Namioka \bref{KelleyNamioka}{1}, G.~Choquet \brefcite{Choquet}{2}{}
(第 2 卷), A.~Taylor--D.~Lay \bref{TaylorLay}{1} 以及 A.~Knapp \bref{Knapp}{2}.

\cmtitem{2. Hahn--Banach 定理的应用.}
Hahn--Banach 定理的应用范围广泛而多样.这里给出两个例子:

\cmtsub{(a) Krein--Milman 定理.}
Hahn--Banach 定理的第二几何形式是证明 Krein--Milman 定理的基本要素.
在叙述这一结果之前, 我们需要一些定义.设 $E$ 为赋范向量空间, $A$ 为 $E$ 的子集.
$A$ 的\emph{凸包}, 记作 $\operatorname{conv}A$, 是包含 $A$ 的最小凸集.显然
$\operatorname{conv}A$ 由 $A$ 中元素的全体有限凸组合构成, 即
\begin{equation*}
\operatorname{conv}A=\left\{\sum_{i\in I}t_ia_i;\ I\ \text{有限},\ a_i\in A\ \forall i,\
t_i\ge 0\ \forall i,\ \text{且}\ \sum_{i\in I}t_i=1\right\}.
\end{equation*}
$A$ 的\emph{闭凸包}, 记作 $\overline{\operatorname{conv}}A$, 是 $\operatorname{conv}A$ 的闭包.
给定凸集 $K\subset E$, 若点 $x\in K$ 不能写成 $K$ 中两点 $x_0,x_1$ 的凸组合, 即
$x\ne(1-t)x_0+tx_1$ (其中 $t\in(0,1)$, $x_0\ne x_1$), 则称 $x$ 为\emph{极端点}.

\begin{keytheorem}[Krein--Milman]
\label{thm:1.13}
设 $K\subset E$ 为紧凸集.则 $K$ 与其极端点集的闭凸包重合.
\end{keytheorem}

Krein--Milman 定理本身有无数的应用与推广 (例如 Choquet 积分表示定理, Bochner 定理, Bernstein 定理等).关于这一广阔的课题, 可参见例如 N.~Bourbaki \bref{Bourbaki}{1}, G.~Choquet \brefcite{Choquet}{2}{} (第 2 卷), R.~Phelps \bref{Phelps}{1}, C.~Dellacherie--P.~A.~Meyer \brefcite{DellacherieMeyer}{1}{} (第 10 章), N.~Dunford--J.~T.~Schwartz \brefcite{DunfordSchwartz}{1}{} (第 1 卷), W.~Rudin \bref{Rudin}{1}, R.~Larsen \bref{Larsen}{1}, J.~Kelley--I.~Namioka \bref{KelleyNamioka}{1}, R.~Edwards \bref{Edwards}{1}.
Y.~Pinchover 给出的一个对 PDE 的有趣应用载于 S.~Agmon \bref{Agmon}{2}.
Krein--Milman 定理的证明见问题 1.

\cmtsub{(b) 在偏微分方程理论中.}
我们仅举一例: 具常系数的\emph{一般}微分算子 $P(D)$ 存在基本解 (Malgrange--Ehrenpreis
定理), 这依赖于 Hahn--Banach 定理的解析形式; 参见例如 L.~Hörmander
\bref{Hormander}{1}, \bref{Hormander}{2}, K.~Yosida \bref{Yosida}{1}, W.~Rudin \bref{Rudin}{1}, F.~Treves \bref{Treves}{2}, M.~Reed--B.~Simon \brefcite{ReedSimon}{1}{} (第 2 卷).

本着同样的精神, 还可提及 P.~Lax 的方法对 Laplacian 的 Green 函数存在性的证明;
见 P.~Lax \brefcite{Lax}{1}{} (第~\ref{sec:9.5}~节) 与 P.~Garabedian \bref{Garabedian}{1}.
给定任意 $f\in L^N(\Omega)$, 方程 $\operatorname{div}u=f$ 在 $\Omega\subset\R^N$ 中
存在解 $u\in L^\infty(\Omega)$ 的证明依赖于 Hahn--Banach 定理 (见
J.~Bourgain--H.~Brezis \bref{BourgainBrezis}{1}, \bref{BourgainBrezis}{2}).
出人意料的是, 通过 Hahn--Banach 定理得到的 $u$ \emph{非线性}地依赖于 $f$.
事实上, 不存在从 $L^N$ 到 $L^\infty$ 的有界线性算子由 $f$ 给出 $u$.
这说明在 Hahn--Banach 定理的证明中使用 Zorn 引理 (以及背后的选择公理)
可能是微妙的事, 并且可能破坏问题的线性特征.有时无法绕开这一障碍.

\cmtitem{3. 凸函数.}
凸分析与对偶原理是近年来大为扩展并日益流行的课题; 参见例如 J.~J.~Moreau
\bref{Moreau}{1}, R.~T.~Rockafellar \bref{Rockafellar}{1}, \bref{Rockafellar}{2}, I.~Ekeland--R.~Temam \bref{EkelandTemam}{1}, I.~Ekeland--T.~Turnbull \bref{EkelandTurnbull}{1}, F.~Clarke \bref{Clarke}{1}, J.~P.~Aubin--I.~Ekeland \bref{AubinEkeland}{1}, J.~B.~Hiriart-Urruty--C.~Lemaréchal \bref{HiriartUrrutyLemarechal}{1}.
在应用方面, 可提及以下几点:

\cmtsub{(a)}
博弈论, 经济学, 最优化, 凸规划; 见 J.~P.~Aubin \bref{Aubin}{1}, \bref{Aubin}{2}, \bref{Aubin}{3}, J.~P.~Aubin--I.~Ekeland \bref{AubinEkeland}{1}, S.~Karlin
\bref{Karlin}{1}, A.~Balakrishnan \bref{Balakrishnan}{1}, V.~Barbu--I.~Precupanu \bref{BarbuPrecupanu}{1}, J.~Franklin \bref{Franklin}{1}, J.~Stoer--C.~Witzgall \bref{StoerWitzgall}{1}.

\cmtsub{(b)}
力学; 见 J.~J.~Moreau \bref{Moreau}{2}, P.~Germain \bref{Germain}{1}, \bref{Germain}{2}, G.~Duvaut--J.~L.~Lions \bref{DuvautLions}{1}, R.~Temam--G.~Strang \bref{TemamStrang}{1} 以及 P.~Germain 在此文之后的评注, H.~D.~Bui \bref{Bui}{1} 及其中的大量参考文献.还请注意 J.~F.~Toland
\bref{Toland}{1}, \bref{Toland}{2}, \bref{Toland}{3} (关于旋转链的研究), A.~Damlamian \bref{Damlamian}{1} (关于等离子体物理中出现的一个问题)
以及 G.~Auchmuty \bref{Auchmuty}{1} 对 (非凸) 对偶的使用.

\cmtsub{(c)}
单调算子与非线性半群理论; 见 H.~Brezis \bref{Brezis}{1}, F.~Browder
\bref{Browder}{1}, V.~Barbu \bref{Barbu}{1} 以及 R.~Phelps \bref{Phelps}{2}.

\cmtsub{(d)}
涉及 Hamilton 系统周期解与非线性振动弦的变分问题; 见 F.~Clarke, I.~Ekeland, J.~M.~Lasry, H.~Brezis, J.~M.~Coron, L.~Nirenberg 的近期工作 (我们引用例如
F.~Clarke--I.~Ekeland \bref{ClarkeEkeland}{1}, H.~Brezis--J.~M.~Coron--L.~Nirenberg \bref{BrezisCoronNirenberg}{1}, H.~Brezis \bref{Brezis}{2}, J.~P.~Aubin--I.~Ekeland \bref{AubinEkeland}{1}, I.~Ekeland \bref{Ekeland}{1}, 以及它们所列的文献).

\cmtsub{(e)}
概率论中的大偏差理论; 见例如 R.~Azencott 等 \bref{Azencott}{1}, D.~W.~Stroock \bref{Stroock}{1}.

\cmtsub{(f)}
偏微分方程与复分析理论; 见 L.~Hörmander \bref{Hormander}{3}.

\cmtitem{4. 有界线性算子的延拓.}
设 $E$ 与 $F$ 为两个 Banach 空间, $G\subset E$ 为闭子空间.设
$S:G\to F$ 为有界线性算子.人们可以问: 能否把 $S$ 延拓为有界线性算子
$T:E\to F$? 注意推论  \ref{cor:1.2} 只在 $F=\R$ 时解决了这一问题.一般地,
答案是否定的 (即使 $E$ 与 $F$ 都是自反空间; 见习题  \ref{exer:1.27}),
只有某些特殊情形例外; 例如:

\cmtsub{(a)}
若 $\dim F<\infty$.可在 $F$ 中选取一组基, 并对 $S$ 的每个分量应用推论  \ref{cor:1.2}.

\cmtsub{(b)}
若 $G$ 有拓扑补空间 (见第~\ref{sec:2.4}~节).特别地, 若 $\dim G<\infty$ 或
$\codim G<\infty$, 或 $E$ 是 Hilbert 空间, 则此条件成立.

人们也可以问: 是否存在\emph{保持范数}的延拓 $T$, 即
$\norm{T}_{\LL(E,F)}=\norm{S}_{\LL(G,F)}$? 答案只在某些例外情形下是肯定的;
见 L.~Nachbin \bref{Nachbin}{1}, J.~Kelley \bref{Kelley}{1} 以及习题  \ref{exer:5.15}.

\section*{第 1 章习题}
\addcontentsline{toc}{section}{第 1 章习题}
\markright{第 1 章习题}

\begin{exercise}[对偶映射的性质.]
\label{exer:1.1}
设 $E$ 为赋范向量空间.对每个 $x\in E$, \emph{对偶映射} $F$ 定义为
\begin{equation*}
F(x)=\bigl\{f\in E^\star;\ \norm{f}=\norm{x}\ \text{且}\ \dual{f}{x}=\norm{x}^2\bigr\}.
\end{equation*}

\begin{enumerate}[label=\arabic*., leftmargin=2.2em, itemsep=0.25em]
	\item 证明
	      \begin{equation*}
	      F(x)=\bigl\{f\in E^\star;\ \norm{f}\le\norm{x}\ \text{且}\ \dual{f}{x}=\norm{x}^2\bigr\},
	      \end{equation*}
	      并由此推出 $F(x)$ 非空, 闭且凸.
	\item 证明: 若 $E^\star$ 严格凸, 则 $F(x)$ 只含一个点.
	\item 证明
	      \begin{equation*}
	      F(x)=\Bigl\{f\in E^\star;\ \tfrac12\norm{y}^2-\tfrac12\norm{x}^2\ge\dual{f}{y-x}
	      \quad \forall y\in E\Bigr\}.
	      \end{equation*}
	\item 由此推出
	      \begin{equation*}
	      \dual{F(x)-F(y)}{x-y}\ge 0\quad \forall x,y\in E,
	      \end{equation*}
	      更精确地说,
	      \begin{equation*}
	      \dual{f-g}{x-y}\ge 0\quad \forall x,y\in E,\quad \forall f\in F(x),\quad \forall g\in F(y).
	      \end{equation*}
	      证明事实上还有
	      \begin{equation*}
	      \dual{f-g}{x-y}\ge\bigl(\norm{x}-\norm{y}\bigr)^2
	      \quad \forall x,y\in E,\quad \forall f\in F(x),\quad \forall g\in F(y).
	      \end{equation*}
	\item 再假设 $E^\star$ 严格凸, 并设 $x,y\in E$ 满足
	      \begin{equation*}
	      \dual{F(x)-F(y)}{x-y}=0.
	      \end{equation*}
	      证明 $Fx=Fy$.
\end{enumerate}
\end{exercise}

\begin{exercise}
\label{exer:1.2}
设 $E$ 为 $n$ 维向量空间, $(e_i)_{1\le i\le n}$ 为 $E$ 的一组基.给定 $x\in E$, 写
$x=\sum_{i=1}^n x_ie_i$, 其中 $x_i\in\R$; 给定 $f\in E^\star$, 置 $f_i=\dual{f}{e_i}$.

\begin{enumerate}[label=\arabic*., leftmargin=2.2em, itemsep=0.25em]
	\item 在 $E$ 上考虑范数
	      \begin{equation*}
	      \norm{x}_1=\sum_{i=1}^n\abs{x_i}.
	      \end{equation*}
	      \begin{enumerate}[label=(\alph*), leftmargin=2.0em, itemsep=0.2em]
	      	\item 用 $f_i$ 显式地算出 $f\in E^\star$ 的\index{duioufanshu@对偶范数 (dual norm)}
对偶范数 $\norm{f}_{E^\star}$.
	      	\item 对每个 $x\in E$ 显式地确定集合 $F(x)$ (对偶映射).
	      \end{enumerate}
	\item 同样的问题, 但 $E$ 上赋以范数
	      \begin{equation*}
	      \norm{x}_\infty=\max_{1\le i\le n}\abs{x_i}.
	      \end{equation*}
	\item 同样的问题, 但 $E$ 上赋以范数
	      \begin{equation*}
	      \norm{x}_2=\Bigl(\sum_{i=1}^n\abs{x_i}^2\Bigr)^{1/2},
	      \end{equation*}
	      更一般地, 赋以范数
	      \begin{equation*}
	      \norm{x}_p=\Bigl(\sum_{i=1}^n\abs{x_i}^p\Bigr)^{1/p},
	      \end{equation*}
	      其中 $p\in(1,\infty)$.
\end{enumerate}
\end{exercise}

\begin{exercise}
\label{exer:1.3}
设 $E=\{u\in C([0,1];\R);\ u(0)=0\}$, 赋以通常的范数
\begin{equation*}
\norm{u}=\max_{t\in[0,1]}\abs{u(t)}.
\end{equation*}
考虑线性泛函
\begin{equation*}
f:\ u\in E\mapsto f(u)=\int_0^1 u(t)\,dt.
\end{equation*}

\begin{enumerate}[label=\arabic*., leftmargin=2.2em, itemsep=0.25em]
	\item 证明 $f\in E^\star$ 并算出 $\norm{f}_{E^\star}$.
	\item 能否找到某个 $u\in E$, 使得 $\norm{u}=1$ 且 $f(u)=\norm{f}_{E^\star}$?
\end{enumerate}
\end{exercise}

\begin{exercise}
\label{exer:1.4}
考虑空间 $E=c_0$ (趋于零的数列), 赋以通常的范数 (见第~\ref{sec:11.3}~节).对 $E$ 中的每个元素
$u=(u_1,u_2,u_3,\dots)$ 定义
\begin{equation*}
f(u)=\sum_{n=1}^\infty\frac{1}{2^n}u_n.
\end{equation*}

\begin{enumerate}[label=\arabic*., leftmargin=2.2em, itemsep=0.25em]
	\item 验证 $f$ 是 $E$ 上的连续线性泛函, 并算出 $\norm{f}_{E^\star}$.
	\item 能否找到某个 $u\in E$, 使得 $\norm{u}=1$ 且 $f(u)=\norm{f}_{E^\star}$?
\end{enumerate}
\end{exercise}

\begin{exercise}
\label{exer:1.5}
设 $E$ 为无限维赋范向量空间.

\begin{enumerate}[label=\arabic*., leftmargin=2.2em, itemsep=0.25em]
	\item 证明 (利用 Zorn 引理) 存在 $E$ 中的\emph{代数基} $(e_i)_{i\in I}$, 使得
	      $\norm{e_i}=1\ \forall i\in I$.
	      回忆: 代数基 (或称 Hamel 基) 是 $E$ 中的子集 $(e_i)_{i\in I}$, 使得每个
	      $x\in E$ 都可唯一地写成
	      \begin{equation*}
	      x=\sum_{i\in J}x_ie_i\quad \text{其中 } J\subset I,\ J \text{ 有限}.
	      \end{equation*}
	\item 构造一个不连续的线性泛函 $f:E\to\R$.
	\item 再假设 $E$ 是 Banach 空间, 证明 $I$ 不可数.
	      [提示: 利用 \index{Bairegangdingli@Baire 纲定理 (Baire category theorem)}
Baire 纲定理 (定理  \ref{thm:2.1}).]
\end{enumerate}
\end{exercise}

\begin{exercise}
\label{exer:1.6}
设 $E$ 为赋范向量空间, $H\subset E$ 为一张超平面.设 $V\subset E$ 为包含 $H$ 的仿射子空间.

\begin{enumerate}[label=\arabic*., leftmargin=2.2em, itemsep=0.25em]
	\item 证明或者 $V=H$, 或者 $V=E$.
	\item 由此推出 $H$ 或者在 $E$ 中闭, 或者在 $E$ 中稠密.
\end{enumerate}
\end{exercise}

\begin{exercise}
\label{exer:1.7}
设 $E$ 为赋范向量空间, $C\subset E$ 为凸集.

\begin{enumerate}[label=\arabic*., leftmargin=2.2em, itemsep=0.25em]
	\item 证明 $\overline{C}$ 与 $\Int C$ 都是凸的.
	\item 给定 $x\in C$ 与 $y\in\Int C$, 证明 $tx+(1-t)y\in\Int C\ \forall t\in(0,1)$.
	\item 由此推出: 只要 $\Int C\ne\emptyset$, 就有 $\overline{C}=\overline{\Int C}$.
\end{enumerate}
\end{exercise}

\begin{exercise}
\label{exer:1.8}
设 $E$ 为赋范向量空间, 范数为 $\norm{\cdot}$.设 $C\subset E$ 为开凸集且 $0\in C$.
以 $p$ 记 $C$ 的规范函数 (见引理  \ref{lem:1.2}).

\begin{enumerate}[label=\arabic*., leftmargin=2.2em, itemsep=0.25em]
	\item 假设 $C$ 对称 (即 $-C=C$) 且有界, 证明 $p$ 是一个与 $\norm{\cdot}$ 等价的范数.
	\item 设 $E=C([0,1];\R)$, 赋以通常的范数
	      \begin{equation*}
	      \norm{u}=\max_{t\in[0,1]}\abs{u(t)}.
	      \end{equation*}
	      令
	      \begin{equation*}
	      C=\Bigl\{u\in E;\ \int_0^1\abs{u(t)}^2dt<1\Bigr\}.
	      \end{equation*}
	      验证 $C$ 是凸的, 对称的, 且 $0\in C$.$C$ 在 $E$ 中有界吗? 算出 $C$ 的规范函数
	      $p$, 并证明 $p$ 是 $E$ 上的范数.$p$ 与 $\norm{\cdot}$ 等价吗?
\end{enumerate}
\end{exercise}

\begin{exercise}[有限维空间中的 Hahn--Banach 定理.]
\label{exer:1.9}
设 $E$ 为有限维赋范空间.设 $C\subset E$ 为非空凸集且 $0\notin C$.我们断言: 总存在
一张超平面把 $C$ 与 $\{0\}$ 分离.

[注意每张超平面都是闭的 (为什么?).本题的要点在于: 对 $C$ 不需要任何附加假设.]

\begin{enumerate}[label=\arabic*., leftmargin=2.2em, itemsep=0.25em]
	\item 设 $(x_n)_{n\ge1}$ 为 $C$ 的可数子集且在 $C$ 中稠密 (它为什么存在?).
	      对每个 $n$ 令
	      \begin{equation*}
	      C_n=\operatorname{conv}\{x_1,x_2,\dots,x_n\}
	      =\Bigl\{x=\sum_{i=1}^n t_ix_i;\ t_i\ge 0\ \forall i\ \text{且}\ \sum_{i=1}^n t_i=1\Bigr\}.
	      \end{equation*}
	      验证 $C_n$ 是紧的, 且 $\bigcup_{n=1}^\infty C_n$ 在 $C$ 中稠密.
	\item 证明存在某个 $f_n\in E^\star$ 使得
	      \begin{equation*}
	      \norm{f_n}=1\quad \text{且}\quad \dual{f_n}{x}\ge 0\quad \forall x\in C_n.
	      \end{equation*}
	\item 由此推出存在某个 $f\in E^\star$ 使得
	      \begin{equation*}
	      \norm{f}=1\quad \text{且}\quad \dual{f}{x}\ge 0\quad \forall x\in C.
	      \end{equation*}
	      得出结论.
	\item 设 $A,B\subset E$ 为非空不交凸集.证明存在一张超平面 $H$ 把 $A$ 与 $B$ 分离.
\end{enumerate}
\end{exercise}

\begin{exercise}
\label{exer:1.10}
设 $E$ 为赋范向量空间, $I$ 为任意指标集.固定 $E$ 中的子集 $(x_i)_{i\in I}$ 与 $\R$ 中的
子集 $(\alpha_i)_{i\in I}$.证明下述两条性质等价:
\begin{align*}
	\text{(A)}\qquad & \text{存在某个 } f\in E^\star \text{ 使得 } \dual{f}{x_i}=\alpha_i
	  \quad \forall i\in I.\\[0.5em]
	\text{(B)}\qquad & \left\{\begin{aligned}
	  &\text{存在常数 } M\ge 0, \text{ 使得对每个有限子集 } J\subset I
	  \text{ 以及任意一组实数 } (\beta_i)_{i\in J},\\
	  &\text{都有 }\ \Bigl|\sum_{i\in J}\beta_i\alpha_i\Bigr|
	  \le M\Bigl\|\sum_{i\in J}\beta_ix_i\Bigr\|.
	\end{aligned}\right.
\end{align*}
注意: 在 (B)$\Rightarrow$(A) 的证明中, 可以找到某个 $f\in E^\star$ 满足
$\norm{f}_{E^\star}\le M$.

[提示: 先在 $(x_i)_{i\in I}$ 张成的线性空间上试着定义 $f$.]
\end{exercise}

\begin{exercise}
\label{exer:1.11}
设 $E$ 为赋范向量空间, $M>0$.固定 $E^\star$ 中的 $n$ 个元素 $(f_i)_{1\le i\le n}$ 与 $n$ 个
实数 $(\alpha_i)_{1\le i\le n}$.证明下述两条性质等价:
\begin{align*}
	\text{(A)}\qquad & \left\{\begin{aligned}
	  &\forall\varepsilon>0\ \exists x_\varepsilon\in E \text{ 使得 }\\
	  &\norm{x_\varepsilon}\le M+\varepsilon \quad\text{且}\quad
	  \dual{f_i}{x_\varepsilon}=\alpha_i \quad \forall i=1,2,\dots,n.
	\end{aligned}\right.\\[0.5em]
	\text{(B)}\qquad & \Bigl|\sum_{i=1}^n\beta_i\alpha_i\Bigr|
	  \le M\Bigl\|\sum_{i=1}^n\beta_if_i\Bigr\|
	  \quad \forall\beta_1,\beta_2,\dots,\beta_n\in\R.
\end{align*}
[提示: 证明 (B)$\Rightarrow$(A) 时, 先考虑 $f_i$ 线性无关的情形, 并仿照引理  \ref{lem:3.3}
的证明.]

比较习题  \ref{exer:1.10},  \ref{exer:1.11} 与引理  \ref{lem:3.3}.
\end{exercise}

\begin{exercise}
\label{exer:1.12}
设 $E$ 为向量空间.固定 $E$ 上的 $n$ 个线性泛函 $(f_i)_{1\le i\le n}$ 与 $n$ 个实数
$(\alpha_i)_{1\le i\le n}$.证明下述两条性质等价:
\begin{align*}
	\text{(A)}\qquad & \text{存在某个 } x\in E \text{ 使得 } f_i(x)=\alpha_i
	  \quad \forall i=1,2,\dots,n.\\[0.2em]
	\text{(B)}\qquad & \left\{\begin{aligned}
	  &\text{对任意一组满足 } \sum_{i=1}^n\beta_if_i=0 \text{ 的实数 }
	  \beta_1,\beta_2,\dots,\beta_n,\\
	  &\text{也必有 } \sum_{i=1}^n\beta_i\alpha_i=0.
	\end{aligned}\right.
\end{align*}
\end{exercise}

\begin{exercise}
\label{exer:1.13}
设 $E=\R^n$, 并令
\begin{equation*}
P=\{x\in\R^n;\ x_i\ge 0\quad \forall i=1,2,\dots,n\}.
\end{equation*}
设 $M$ 为 $E$ 的线性子空间且 $M\cap P=\{0\}$.证明存在 $E$ 中的一张超平面 $H$, 使得
\begin{equation*}
M\subset H\quad \text{且}\quad H\cap P=\{0\}.
\end{equation*}
[提示: 先证明 $M^\perp\cap\Int P\ne\emptyset$.]
\end{exercise}

\begin{exercise}
\label{exer:1.14}
设 $E=\ell^1$ (见第~\ref{sec:11.3}~节), 并考虑两个集合
\begin{equation*}
X=\{x=(x_n)_{n\ge1}\in E;\ x_{2n}=0\ \forall n\ge1\}
\end{equation*}
与
\begin{equation*}
Y=\Bigl\{y=(y_n)_{n\ge1}\in E;\ y_{2n}=\frac{1}{2^n}y_{2n-1}\ \forall n\ge1\Bigr\}.
\end{equation*}

\begin{enumerate}[label=\arabic*., leftmargin=2.2em, itemsep=0.25em]
	\item 验证 $X$ 与 $Y$ 都是闭线性空间, 且 $\overline{X+Y}=E$.
	\item 令 $c\in E$ 由下式定义:
	      \begin{equation*}
	      \begin{cases}
	      	c_{2n-1}=0 & \forall n\ge1,\\[0.2em]
	      	c_{2n}=\dfrac{1}{2^n} & \forall n\ge1.
	      \end{cases}
	      \end{equation*}
	      验证 $c\notin X+Y$.
	\item 令 $Z=X-c$, 验证 $Y\cap Z=\emptyset$.$E$ 中是否存在把 $Y$ 与 $Z$ 分离的闭超平面?
	      与定理  \ref{thm:1.7} 及习题  \ref{exer:1.9} 作比较.
	\item 在 $E=\ell^p$ ($1<p<\infty$) 以及 $E=c_0$ 中问同样的问题.
\end{enumerate}
\end{exercise}

\begin{exercise}
\label{exer:1.15}
设 $E$ 为赋范向量空间, $C\subset E$ 为凸集且 $0\in C$.置
\begin{align*}
\text{(A)}\qquad C^\star&=\{f\in E^\star;\ \dual{f}{x}\le1\quad \forall x\in C\},\\
\text{(B)}\qquad C^{\star\star}&=\{x\in E;\ \dual{f}{x}\le1\quad \forall f\in C^\star\}.
\end{align*}

\begin{enumerate}[label=\arabic*., leftmargin=2.2em, itemsep=0.25em]
	\item 证明 $C^{\star\star}=\overline{C}$.
	\item 若 $C$ 是线性空间, $C^\star$ 是什么?
\end{enumerate}
\end{exercise}

\begin{exercise}
\label{exer:1.16}
设 $E=\ell^1$, 从而 $E^\star=\ell^\infty$ (见第~\ref{sec:11.3}~节).把 $N=c_0$ 看作 $E^\star$ 的闭子空间.
确定
\begin{equation*}
N^\perp=\{x\in E;\ \dual{f}{x}=0\quad \forall f\in N\}
\end{equation*}
与
\begin{equation*}
N^{\perp\perp}=\{f\in E^\star;\ \dual{f}{x}=0\quad \forall x\in N^\perp\}.
\end{equation*}
验证 $N^{\perp\perp}\ne N$.
\end{exercise}

\begin{exercise}
\label{exer:1.17}
设 $E$ 为赋范向量空间, $f\in E^\star$ 且 $f\ne0$.设 $M$ 为超平面 $[f=0]$.

\begin{enumerate}[label=\arabic*., leftmargin=2.2em, itemsep=0.25em]
	\item 确定 $M^\perp$.
	\item 证明对每个 $x\in E$,
	      \begin{equation*}
	      \dist(x,M)=\inf_{y\in M}\norm{x-y}=\frac{\abs{\dual{f}{x}}}{\norm{f}}.
	      \end{equation*}
	      [找一个直接的方法, 或利用第~\ref{sec:1.4}~节例 3.]
	\item 现在假设 $E=\{u\in C([0,1];\R);\ u(0)=0\}$ 且
	      \begin{equation*}
	      \dual{f}{u}=\int_0^1 u(t)\,dt,\quad u\in E.
	      \end{equation*}
	      证明 $\dist(u,M)=\bigl|\int_0^1u(t)dt\bigr|\ \forall u\in E$.
	      证明对任何 $u\in E\setminus M$, 下确界 $\inf_{v\in M}\norm{u-v}$ 都取不到.
\end{enumerate}
\end{exercise}

\begin{exercise}
\label{exer:1.18}
验证下面定义的函数 $\varphi:\R\to(-\infty,+\infty]$ 是凸的, 下半连续的, 并确定共轭函数
$\varphi^\star$.画出它们的图像并标出各自的上图.

\begin{enumerate}[label=(\alph*), leftmargin=2.0em, itemsep=0.3em]
	\item $\varphi(x)=ax+b$, 其中 $a,b\in\R$.
	\item $\varphi(x)=e^x$.
	\item
	      $\varphi(x)=\begin{cases}0 & \text{若 } \abs{x}\le1,\\ +\infty & \text{若 } \abs{x}>1.\end{cases}$
	\item
	      $\varphi(x)=\begin{cases}0 & \text{若 } x=0,\\ +\infty & \text{若 } x\ne0.\end{cases}$
	\item
	      $\varphi(x)=\begin{cases}-\log x & \text{若 } x>0,\\ +\infty & \text{若 } x\le0.\end{cases}$
	\item
	      $\varphi(x)=\begin{cases}-(1-x^2)^{1/2} & \text{若 } \abs{x}\le1,\\
	      +\infty & \text{若 } \abs{x}>1.\end{cases}$
	\item
	      $\varphi(x)=\begin{cases}\tfrac12\abs{x}^2 & \text{若 } \abs{x}\le1,\\
	      \abs{x}-\tfrac12 & \text{若 } \abs{x}>1.\end{cases}$
	\item $\varphi(x)=\dfrac{1}{p}\abs{x}^p$, 其中 $1<p<\infty$.
	\item $\varphi(x)=x^+=\max\{x,0\}$.
	\item
	      $\varphi(x)=\begin{cases}\dfrac{1}{p}x^p & \text{若 } x\ge0,\ \text{其中 }
	      1<p<+\infty,\\ +\infty & \text{若 } x<0.\end{cases}$
	\item
	      $\varphi(x)=\begin{cases}-\dfrac{1}{p}x^p & \text{若 } x\ge0,\ \text{其中 }
	      0<p<1,\\ +\infty & \text{若 } x<0.\end{cases}$
	\item $\varphi(x)=\dfrac{1}{p}\bigl[(\abs{x}-1)^+\bigr]^p$, 其中 $1<p<\infty$.
\end{enumerate}
\end{exercise}

\begin{exercise}
\label{exer:1.19}
设 $E$ 为赋范向量空间.

\begin{enumerate}[label=\arabic*., leftmargin=2.2em, itemsep=0.25em]
	\item 设 $\varphi,\psi:E\to(-\infty,+\infty]$ 为两个函数且 $\varphi\le\psi$.证明
	      $\psi^\star\le\varphi^\star$.
	\item 设 $F:\R\to(-\infty,+\infty]$ 为凸的, 下半连续的函数, 满足 $F(0)=0$ 且
	      $F(t)\ge0\ \forall t\in\R$.令 $\varphi(x)=F(\norm{x})$.
	      证明 $\varphi$ 凸且下半连续, 并且 $\varphi^\star(f)=F^\star(\norm{f})\ \forall f\in E^\star$.
\end{enumerate}
\end{exercise}

\begin{exercise}
\label{exer:1.20}
设 $E=\ell^p$, $1\le p<\infty$ (见第~\ref{sec:11.3}~节).验证下面定义的函数
$\varphi:E\to(-\infty,+\infty]$ 是凸的, 下半连续的, 并确定 $\varphi^\star$.
对 $x=(x_1,x_2,\dots,x_n,\dots)$ 置
\begin{enumerate}[label=(\alph*), leftmargin=2.0em, itemsep=0.3em]
	\item
	      $\varphi(x)=\begin{cases}\sum_{k=1}^{+\infty}k\abs{x_k}^2
	      & \text{若 } \sum_{k=1}^\infty k\abs{x_k}^2<+\infty,\\ +\infty & \text{其他情形}.\end{cases}$
	\item $\varphi(x)=\sum_{k=2}^{+\infty}\abs{x_k}^k$.
	      (验证对每个 $x\in E$ 都有 $\varphi(x)<\infty$.)
	\item
	      $\varphi(x)=\begin{cases}\sum_{k=1}^{+\infty}\abs{x_k}
	      & \text{若 } \sum_{k=1}^\infty\abs{x_k}<+\infty,\\ +\infty & \text{其他情形}.\end{cases}$
\end{enumerate}
\end{exercise}

\begin{exercise}
\label{exer:1.21}
设 $E=E^\star=\R^2$, 并令
\begin{equation*}
C=\{[x_1,x_2];\ x_1\ge0,\ x_2\ge0\}.
\end{equation*}
在 $E$ 上定义函数
\begin{equation*}
\varphi(x)=\begin{cases}-\sqrt{x_1x_2} & \text{若 } x\in C,\\ +\infty & \text{若 } x\notin C.\end{cases}
\end{equation*}

\begin{enumerate}[label=\arabic*., leftmargin=2.2em, itemsep=0.25em]
	\item 证明 $\varphi$ 在 $E$ 上凸且下半连续.
	\item 确定 $\varphi^\star$.
	\item 考虑集合 $D=\{[x_1,x_2];\ x_1=0\}$ 与函数 $\psi=I_D$.算出下列表达式的值:
	      \begin{equation*}
	      \inf_{x\in E}\{\varphi(x)+\psi(x)\}
	      \quad\text{与}\quad
	      \sup_{f\in E^\star}\{-\varphi^\star(-f)-\psi^\star(f)\}.
	      \end{equation*}
	\item 与定理  \ref{thm:1.12} 的结论作比较, 并解释其中的差别.
\end{enumerate}
\end{exercise}

\begin{exercise}
\label{exer:1.22}
设 $E$ 为赋范向量空间, $A\subset E$ 为非空闭集.令
\begin{equation*}
\varphi(x)=\dist(x,A)=\inf_{a\in A}\norm{x-a}.
\end{equation*}

\begin{enumerate}[label=\arabic*., leftmargin=2.2em, itemsep=0.25em]
	\item 验证 $\abs{\varphi(x)-\varphi(y)}\le\norm{x-y}\ \forall x,y\in E$.
	\item 假设 $A$ 是凸的, 证明 $\varphi$ 是凸的.
	\item 反之, 假设 $\varphi$ 是凸的, 证明 $A$ 是凸的.
	\item 对不一定凸的任意 $A$, 证明 $\varphi^\star=(I_A)^\star+I_{B_{E^\star}}$.
\end{enumerate}
\end{exercise}

\index{xiajuanji@下卷积 (inf-convolution)}
\begin{exercise}[下卷积.]
\label{exer:1.23}
设 $E$ 为赋范向量空间.给定两个函数 $\varphi,\psi:E\to(-\infty,+\infty]$, 按下述方式定义
$\varphi$ 与 $\psi$ 的\emph{下卷积}: 对每个 $x\in E$, 令
\begin{equation*}
(\varphi\nabla\psi)(x)=\inf_{y\in E}\{\varphi(x-y)+\psi(y)\}.
\end{equation*}
注意以下几点:
\begin{enumerate}[label=(\roman*), leftmargin=2.4em, itemsep=0.15em]
	\item $(\varphi\nabla\psi)(x)$ 可能取值 $\pm\infty$;
	\item $(\varphi\nabla\psi)(x)<+\infty$ 当且仅当 $x\in D(\varphi)+D(\psi)$.
\end{enumerate}

\begin{enumerate}[label=\arabic*., leftmargin=2.2em, itemsep=0.25em]
	\item 假设 $D(\varphi^\star)\cap D(\psi^\star)\ne\emptyset$, 证明 $(\varphi\nabla\psi)$ 不取值
	      $-\infty$, 并且
	      \begin{equation*}
	      (\varphi\nabla\psi)^\star=\varphi^\star+\psi^\star.
	      \end{equation*}
	\item 假设 $D(\varphi)\cap D(\psi)\ne\emptyset$, 证明在 $E^\star$ 上
	      \begin{equation*}
	      (\varphi+\psi)^\star\le(\varphi^\star\nabla\psi^\star).
	      \end{equation*}
	\item 假设 $\varphi$ 与 $\psi$ 都是凸的, 且存在 $x_0\in D(\varphi)\cap D(\psi)$ 使 $\varphi$
	      在 $x_0$ 处连续.证明在 $E^\star$ 上
	      \begin{equation*}
	      (\varphi+\psi)^\star=(\varphi^\star\nabla\psi^\star).
	      \end{equation*}
	\item 假设 $\varphi$ 与 $\psi$ 都是凸的, 下半连续的, 且 $D(\varphi)\cap D(\psi)\ne\emptyset$.
	      证明在 $E$ 上
	      \begin{equation*}
	      (\varphi^\star\nabla\psi^\star)^\star=(\varphi+\psi).
	      \end{equation*}
\end{enumerate}
给定函数 $\varphi:E\to(-\infty,+\infty]$, 置
\begin{equation*}
\operatorname{epi}_{s}\varphi=\{[x,\lambda]\in E\times\R;\ \varphi(x)<\lambda\}.
\end{equation*}

\begin{enumerate}[label=\arabic*., leftmargin=2.2em, itemsep=0.25em, start=5]
	\item 验证 $\varphi$ 是凸的当且仅当 $\operatorname{epi}_{s}\varphi$ 是 $E\times\R$ 中的凸集.
	\item 设 $\varphi,\psi:E\to(-\infty,+\infty]$ 为函数且 $D(\varphi^\star)\cap D(\psi^\star)\ne\emptyset$.
	      证明
	      \begin{equation*}
	      \operatorname{epi}_{s}(\varphi\nabla\psi)
	      =(\operatorname{epi}_{s}\varphi)+(\operatorname{epi}_{s}\psi).
	      \end{equation*}
	\item 由此推出: 若 $\varphi,\psi:E\to(-\infty,+\infty]$ 为凸函数且
	      $D(\varphi^\star)\cap D(\psi^\star)\ne\emptyset$, 则 $(\varphi\nabla\psi)$ 是凸函数.
\end{enumerate}
\end{exercise}

\begin{exercise}[用下卷积作正则化.]
\label{exer:1.24}
设 $E$ 为赋范向量空间, $\varphi:E\to(-\infty,+\infty]$ 为凸的, 下半连续的函数且
$\varphi\not\equiv+\infty$.我们的目标是构造一列函数 $(\varphi_n)$, 使得:
\begin{enumerate}[label=(\roman*), leftmargin=2.4em, itemsep=0.15em]
	\item 对每个 $n$, $\varphi_n:E\to(-\infty,+\infty)$ 是凸的, 连续的;
	\item 对每个 $x$, 数列 $(\varphi_n(x))_n$ 单调不减且收敛于 $\varphi(x)$.
\end{enumerate}
为此, 令
\begin{equation*}
\varphi_n(x)=\inf_{y\in E}\{n\norm{x-y}+\varphi(y)\}.
\end{equation*}

\begin{enumerate}[label=\arabic*., leftmargin=2.2em, itemsep=0.25em]
	\item 证明存在足够大的 $N$, 使得当 $n\ge N$ 时, 对一切 $x\in E$, $\varphi_n(x)$ 都有限.
	      以下取 $n\ge N$.
	\item 证明 $\varphi_n$ 是凸的 (见习题  \ref{exer:1.23}), 并且
	      \begin{equation*}
	      \abs{\varphi_n(x_1)-\varphi_n(x_2)}\le n\norm{x_1-x_2}
	      \quad \forall x_1,x_2\in E.
	      \end{equation*}
	\item 确定 $(\varphi_n)^\star$.
	\item 验证 $\varphi_n(x)\le\varphi(x)\ \forall x\in E$, $\forall n$.证明对每个 $x\in E$,
	      数列 $(\varphi_n(x))_n$ 单调不减.
	\item 给定 $x\in D(\varphi)$, 取 $y_n\in E$ 使得
	      \begin{equation*}
	      \varphi_n(x)\le n\norm{x-y_n}+\varphi(y_n)\le\varphi_n(x)+\frac{1}{n}.
	      \end{equation*}
	      证明 $\lim_{n\to\infty}y_n=x$, 并由此推出 $\lim_{n\to\infty}\varphi_n(x)=\varphi(x)$.
	\item 对 $x\notin D(\varphi)$, 证明 $\lim_{n\to\infty}\varphi_n(x)=+\infty$.
	      [提示: 用反证法.]
\end{enumerate}
\end{exercise}

\begin{exercise}[半标量积.]
\label{exer:1.25}
设 $E$ 为赋范向量空间.

\begin{enumerate}[label=\arabic*., leftmargin=2.2em, itemsep=0.25em]
	\item 设 $\varphi:E\to(-\infty,+\infty)$ 为凸函数.给定 $x,y\in E$, 考虑函数
	      \begin{equation*}
	      h(t)=\frac{\varphi(x+ty)-\varphi(x)}{t},\quad t>0.
	      \end{equation*}
	      验证 $h$ 在 $(0,+\infty)$ 上单调不减, 并由此推出
	      \begin{equation*}
	      \lim_{t\downarrow0}h(t)=\inf_{t>0}h(t)
	      \end{equation*}
	      在 $[-\infty,+\infty)$ 中存在.用
	      \begin{equation*}
	      [x,y]=\inf_{t>0}\frac{1}{2t}\bigl[\norm{x+ty}^2-\norm{x}^2\bigr]
	      \end{equation*}
	      定义\emph{半标量积} $[x,y]$.
	\item 证明 $\abs{[x,y]}\le\norm{x}\norm{y}\ \forall x,y\in E$.
	\item 证明
	      \begin{equation*}
	      [x,\lambda x+\mu y]=\lambda\norm{x}^2+\mu[x,y]
	      \quad \forall x,y\in E,\quad \forall\lambda\in\R,\quad \forall\mu\ge0
	      \end{equation*}
	      以及
	      \begin{equation*}
	      [\lambda x,\mu y]=\lambda\mu[x,y]
	      \quad \forall x,y\in E,\quad \forall\lambda\ge0,\quad \forall\mu\ge0.
	      \end{equation*}
	\item 证明对每个 $x\in E$, 函数 $y\mapsto[x,y]$ 是凸的.证明函数
	      $G(x,y)=-[x,y]$ 在 $E\times E$ 上下半连续.
	\item 证明
	      \begin{equation*}
	      [x,y]=\max_{f\in F(x)}\dual{f}{y}\quad \forall x,y\in E,
	      \end{equation*}
	      其中 $F$ 表示对偶映射 (见推论  \ref{cor:1.3} 之后的注记 \ref{rem:1.2} 以及
	      习题  \ref{exer:1.1}).
	      [提示: 置 $\alpha=[x,y]$, 并对如下定义的函数 $\varphi$ 与 $\psi$ 应用
	      定理  \ref{thm:1.12}:
	      \begin{equation*}
	      \varphi(z)=\tfrac12\norm{x+z}^2-\tfrac12\norm{x}^2,\quad z\in E,
	      \end{equation*}
	      以及
	      \begin{equation*}
	      \psi(z)=\begin{cases}-t\alpha & \text{当 } z=ty\ \text{且}\ t\ge0,\\
	      +\infty & \text{其他情形}.\end{cases}]
	      \end{equation*}
	\item 对 $E=\R^n$ 赋以范数 $\norm{x}_p$ ($1\le p\le\infty$), 显式地确定 $[x,y]$
	      (见第~\ref{sec:11.3}~节).
	      [提示: 利用习题  \ref{exer:1.2} 的结果.]
\end{enumerate}
\end{exercise}

\begin{exercise}[严格凸的范数与函数.]
\label{exer:1.26}
设 $E$ 为赋范向量空间.若
\begin{equation*}
\norm{tx+(1-t)y}<1,\quad
\forall x,y\in E\ \text{且}\ x\ne y,\ \norm{x}=\norm{y}=1,\quad \forall t\in(0,1),
\end{equation*}
则称范数 $\norm{\cdot}$ 是\emph{严格凸}的 (或称空间 $E$ 是严格凸的).
若
\begin{equation*}
\varphi(tx+(1-t)y)<t\varphi(x)+(1-t)\varphi(y)
\quad \forall x,y\in E\ \text{且}\ x\ne y,\quad \forall t\in(0,1),
\end{equation*}
则称函数 $\varphi:E\to(-\infty,+\infty]$ 是\emph{严格凸}的.

\begin{enumerate}[label=\arabic*., leftmargin=2.2em, itemsep=0.25em]
	\item 证明范数 $\norm{\cdot}$ 严格凸当且仅当函数 $\varphi(x)=\norm{x}^2$ 严格凸.
	\item 对 $\varphi(x)=\norm{x}^p$ ($1<p<\infty$) 问同样的问题.
\end{enumerate}
\end{exercise}

\begin{exercise}
\label{exer:1.27}
设 $E$ 与 $F$ 为两个 Banach 空间, $G\subset E$ 为闭子空间.设 $T:G\to F$ 为连续线性映射.
本题的目的是说明: 有时 $T$ 不能延拓为连续线性映射 $\widetilde{T}:E\to F$.为此, 设 $E$
为 Banach 空间, $G\subset E$ 为不承认任何补空间的闭子空间 (见第~
\ref{chap:2}~章注记 8).令 $F=G$,
$T=I$ (恒等映射).证明 $T$ 不能延拓.

[提示: 用反证法.]

与推论  \ref{cor:1.2} 的结论作比较.

\end{exercise}

% chapters/ch02.tex
% 第~
\ref{chap:2}~章 一致有界原理与\index{bituxiangdingli@闭图像定理 (closed graph theorem)}
闭图像定理
% 源: 原书 PDF 第 41--64 页 (印刷页 31--54)

\chapter{一致有界原理与闭图像定理}
\label{chap:2}

\section{Baire 纲定理}
\label{sec:2.1}

下面的经典结果在第~
\ref{chap:2}~章各定理的证明中起着本质作用.

\begin{keytheorem}[Baire]
\label{thm:2.1}
设 $X$ 为完备度量空间, $(X_n)_{n\ge1}$ 为 $X$ 中的一列闭子集.假设对每个 $n\ge1$
都有
\begin{equation*}
\Int X_n=\emptyset.
\end{equation*}
则
\begin{equation*}
\Int\Bigl(\bigcup_{n=1}^\infty X_n\Bigr)=\emptyset.
\end{equation*}
\end{keytheorem}

\begin{remark}
\label{rem:2.1}
Baire 纲定理常以下述形式使用.设 $X$ 为非空完备度量空间, $(X_n)_{n\ge1}$ 为一列
闭子集, 使得
\begin{equation*}
\bigcup_{n=1}^\infty X_n=X.
\end{equation*}
则存在某个 $n_0$, 使 $\Int X_{n_0}\ne\emptyset$.
\end{remark}

\begin{proof}
置 $O_n=X_n^c$, 则对每个 $n\ge1$, $O_n$ 是 $X$ 中的开稠集.我们要证明
$G=\bigcap_{n=1}^\infty O_n$ 在 $X$ 中稠密.设 $\omega$ 为 $X$ 中的非空开集;
我们将证明 $\omega\cap G\ne\emptyset$.

如通常那样, 置
\begin{equation*}
B(x,r)=\{y\in X;\ d(y,x)<r\}.
\end{equation*}
任取 $x_0\in\omega$ 与 $r_0>0$, 使得
\begin{equation*}
B(x_0,r_0)\subset\omega.
\end{equation*}
然后取 $x_1\in B(x_0,r_0)\cap O_1$ 与 $r_1>0$, 使得
\begin{equation*}
\left\{\begin{aligned}
	&B(x_1,r_1)\subset B(x_0,r_0)\cap O_1,\\
	&0<r_1<\frac{r_0}{2},
\end{aligned}\right.
\end{equation*}
这总是可以做到的, 因为 $O_1$ 是开稠集.用归纳法构造两列 $(x_n)$ 与 $(r_n)$, 使得
\begin{equation*}
\left\{\begin{aligned}
	&B(x_{n+1},r_{n+1})\subset B(x_n,r_n)\cap O_{n+1},\quad \forall n\ge0,\\
	&0<r_{n+1}<\frac{r_n}{2}.
\end{aligned}\right.
\end{equation*}
由此可知 $(x_n)$ 是 Cauchy 列; 设 $x_n\to\ell$.

由于对每个 $n\ge0$ 与每个 $p\ge0$ 都有 $x_{n+p}\in B(x_n,r_n)$, 令 $p\to\infty$
取极限即得
\begin{equation*}
\ell\in\overline{B(x_n,r_n)},\quad \forall n\ge0.
\end{equation*}
特别地, $\ell\in\omega\cap G$.
\end{proof}

\section{一致有界原理}
\label{sec:2.2}

\begin{notation}
设 $E$ 与 $F$ 为两个赋范向量空间.以 $\LL(E,F)$ 记从 $E$ 到 $F$ 的连续 (即有界)
线性算子构成的空间, 赋以范数
\begin{equation*}
\norm{T}_{\LL(E,F)}=\sup_{\substack{x\in E\\ \norm{x}\le1}}\norm{Tx}.
\end{equation*}
如通常那样, 用 $\LL(E)$ 代替 $\LL(E,E)$.
\end{notation}

\index{BanachSteinhausdingli@Banach--Steinhaus 定理 (Banach--Steinhaus theorem)}
\begin{keytheorem}[Banach--Steinhaus, 一致有界原理]
\label{thm:2.2}
设 $E$ 与 $F$ 为两个 Banach 空间, $(T_i)_{i\in I}$ 为一族 (不必可数) 从 $E$ 到 $F$ 的
连续线性算子.假设
\begin{equation}
\sup_{i\in I}\norm{T_ix}<\infty\quad \forall x\in E.
\label{eq:2.1}
\end{equation}
则
\begin{equation}
\sup_{i\in I}\norm{T_i}_{\LL(E,F)}<\infty.
\label{eq:2.2}
\end{equation}
换言之, 存在常数 $c$ 使得
\begin{equation*}
\norm{T_ix}\le c\norm{x}\quad \forall x\in E,\ \forall i\in I.
\end{equation*}
\end{keytheorem}

\begin{remark}
\label{rem:2.2}
定理  \ref{thm:2.2} 的结论颇为出人意料.由逐点估计竟能推出整体 (一致) 估计.
\end{remark}

\begin{proof}
对每个 $n\ge1$, 令
\begin{equation*}
X_n=\{x\in E;\ \forall i\in I,\ \norm{T_ix}\le n\},
\end{equation*}
则 $X_n$ 是闭的, 且由 \eqref{eq:2.1} 有
\begin{equation*}
\bigcup_{n=1}^\infty X_n=E.
\end{equation*}
由 Baire 纲定理, 存在某个 $n_0\ge1$ 使 $\Int(X_{n_0})\ne\emptyset$.取 $x_0\in E$ 与
$r>0$ 使 $B(x_0,r)\subset X_{n_0}$.我们有
\begin{equation*}
\norm{T_i(x_0+rz)}\le n_0\quad \forall i\in I,\ \forall z\in B(0,1).
\end{equation*}
由此得到
\begin{equation*}
r\norm{T_i}_{\LL(E,F)}\le n_0+\norm{T_ix_0},
\end{equation*}
从而 \eqref{eq:2.2} 成立.
\end{proof}

\begin{remark}
\label{rem:2.3}
回忆: 一般说来, 连续映射的逐点极限不必连续.线性性假设在定理  \ref{thm:2.2} 中起着
本质作用.然而要注意: 在定理  \ref{thm:2.2} 的框架下, 并不能推出
$\norm{T_n-T}_{\LL(E,F)}\to0$.
\end{remark}

下面是一致有界原理的几个直接推论.

\begin{corollary}
\label{cor:2.3}
设 $E$ 与 $F$ 为两个 Banach 空间, $(T_n)$ 为从 $E$ 到 $F$ 的一列连续线性算子, 使得
对每个 $x\in E$, $T_nx$ 当 $n\to\infty$ 时收敛, 其极限记为 $Tx$.则
\begin{enumerate}[label=(\alph*), leftmargin=2.0em, itemsep=0.15em]
	\item $\sup_n\norm{T_n}_{\LL(E,F)}<\infty$,
	\item $T\in\LL(E,F)$,
	\item $\norm{T}_{\LL(E,F)}\le\liminf_{n\to\infty}\norm{T_n}_{\LL(E,F)}$.
\end{enumerate}
\end{corollary}

\begin{proof}
(a) 直接由定理  \ref{thm:2.2} 得到, 因而存在常数 $c$ 使得
\begin{equation*}
\norm{T_nx}\le c\norm{x}\quad \forall n,\ \forall x\in E.
\end{equation*}
取极限得
\begin{equation*}
\norm{Tx}\le c\norm{x}\quad \forall x\in E.
\end{equation*}
$T$ 显然是线性的, 故得 (b).

最后, 我们有
\begin{equation*}
\norm{T_nx}\le\norm{T_n}_{\LL(E,F)}\norm{x}\quad \forall x\in E,
\end{equation*}
于是 (c) 直接可得.
\end{proof}

\begin{keycorollary}
\label{cor:2.4}
设 $G$ 为 Banach 空间, $B$ 为 $G$ 的子集.假设
\begin{equation}
\text{对每个 } f\in G^\star,\ \text{集合 } f(B)=\{\dual{f}{x};\ x\in B\}
\text{ 有界 (在 } \R \text{ 中)}.
\label{eq:2.3}
\end{equation}
则
\begin{equation}
B \text{ 有界}.
\label{eq:2.4}
\end{equation}
\end{keycorollary}

\begin{proof}
我们对 $E=G^\star$, $F=\R$, $I=B$ 应用定理  \ref{thm:2.2}.对每个 $b\in B$, 置
\begin{equation*}
T_b(f)=\dual{f}{b},\quad f\in E=G^\star,
\end{equation*}
于是由 \eqref{eq:2.3},
\begin{equation*}
\sup_{b\in B}\abs{T_b(f)}<\infty\quad \forall f\in E.
\end{equation*}
由定理  \ref{thm:2.2}, 存在常数 $c$ 使得
\begin{equation*}
\abs{\dual{f}{b}}\le c\norm{f}\quad \forall f\in G^\star,\ \forall b\in B.
\end{equation*}
因此 (利用推论  \ref{cor:1.4}) 得到
\begin{equation*}
\norm{b}\le c\quad \forall b\in B.
\end{equation*}
\end{proof}

\begin{remark}
\label{rem:2.4}
推论  \ref{cor:2.4} 说的是: 要证明集合 $B$ 有界, 只须通过有界线性泛函去``观察''
$B$.这在有限维空间中是一种熟悉的做法, 那里线性泛函就是相对于某组基的分量.在
某种意义下, 推论  \ref{cor:2.4} 在无限维空间中取代了分量的用法.有时人们把推论
 \ref{cor:2.4} 的结论表述为: ``弱有界'' $\Longleftrightarrow$ ``强有界'' (见第~
\ref{chap:3}~章).
\end{remark}

下面是与推论  \ref{cor:2.4} 对偶的命题:

\begin{corollary}
\label{cor:2.5}
设 $G$ 为 Banach 空间, $B^\star$ 为 $G^\star$ 的子集.假设
\begin{equation}
\text{对每个 } x\in G,\ \text{集合 } \dual{B^\star}{x}=\{\dual{f}{x};\ f\in B^\star\}
\text{ 有界 (在 } \R \text{ 中)}.
\label{eq:2.5}
\end{equation}
则
\begin{equation}
B^\star \text{ 有界}.
\label{eq:2.6}
\end{equation}
\end{corollary}

\begin{proof}
对 $E=G$, $F=\R$, $I=B^\star$ 应用定理  \ref{thm:2.2}.对每个 $b\in B^\star$, 置
$T_b(x)=\dual{b}{x}$ ($x\in G=E$).我们得到: 存在常数 $c$ 使得
\begin{equation*}
\abs{\dual{b}{x}}\le c\norm{x}\quad \forall b\in B^\star,\ \forall x\in G.
\end{equation*}
由对偶范数的定义即得
\begin{equation*}
\norm{b}\le c\quad \forall b\in B^\star.
\end{equation*}
\end{proof}

\section{\index{kaiyingshedingli@开映射定理 (open mapping theorem)}
开映射定理与闭图像定理}
\label{sec:2.3}

下面是两个属于 Banach 的基本结果.

\begin{keytheorem}[开映射定理]
\label{thm:2.6}
设 $E$ 与 $F$ 为两个 Banach 空间, $T$ 为从 $E$ 到 $F$ 的连续线性算子且为满射
(= 映上).则存在常数 $c>0$ 使得
\begin{equation}
T(B_E(0,1))\supset B_F(0,c).
\label{eq:2.7}
\end{equation}
\end{keytheorem}

\begin{remark}
\label{rem:2.5}
性质 \eqref{eq:2.7} 蕴含: $T$ 把 $E$ 中任何开集的像都是 $F$ 中的开集 (这正是本
定理名称的由来!).事实上, 设 $U$ 在 $E$ 中开, 证明 $T(U)$ 开.任取
$y_0\in T(U)$, 则 $y_0=Tx_0$ 对某个 $x_0\in U$ 成立.取 $r>0$ 使
$B(x_0,r)\subset U$, 即 $x_0+B(0,r)\subset U$.于是
\begin{equation*}
y_0+T(B(0,r))\subset T(U).
\end{equation*}
利用 \eqref{eq:2.7} 得
\begin{equation*}
T(B(0,r))\supset B(0,rc),
\end{equation*}
从而
\begin{equation*}
B(y_0,rc)\subset T(U).
\end{equation*}
\end{remark}

定理  \ref{thm:2.6} 的几个重要推论如下.

\begin{keycorollary}
\index{youjienidingli@有界逆定理 (bounded inverse theorem)}
\label{cor:2.7}
设 $E$ 与 $F$ 为两个 Banach 空间, $T$ 为从 $E$ 到 $F$ 的连续线性算子且为双射,
即既单射 (= 一对一) 又满射.则 $T^{-1}$ 也连续 (从 $F$ 到 $E$).
\end{keycorollary}

\begin{proof}[推论 \ref{cor:2.7}]
性质 \eqref{eq:2.7} 与 $T$ 的单射性蕴含: 若 $x\in E$ 满足 $\norm{Tx}<c$, 则
$\norm{x}<1$.由齐次性得
\begin{equation*}
\norm{x}\le\frac{1}{c}\norm{Tx}\quad \forall x\in E,
\end{equation*}
因此 $T^{-1}$ 连续.
\end{proof}

\begin{corollary}
\label{cor:2.8}
设 $E$ 为赋以两个范数 $\norm{\cdot}_1$ 与 $\norm{\cdot}_2$ 的向量空间.假设 $E$ 对
这两个范数都是 Banach 空间, 且存在常数 $C\ge0$ 使得
\begin{equation*}
\norm{x}_2\le C\norm{x}_1\quad \forall x\in E.
\end{equation*}
则这两个范数等价, 即存在常数 $c>0$ 使得
\begin{equation*}
\norm{x}_1\le c\norm{x}_2\quad \forall x\in E.
\end{equation*}
\end{corollary}

\begin{proof}[推论 \ref{cor:2.8}]
对 $E=(E,\norm{\cdot}_1)$, $F=(E,\norm{\cdot}_2)$, $T=I$ 应用推论  \ref{cor:2.7}.
\end{proof}

\begin{proof}[定理 \ref{thm:2.6}]
我们把论证分为两步:

\emph{第一步.} 设 $T$ 为从 $E$ 映上 $F$ 的线性满射.则存在常数 $c>0$ 使得
\begin{equation}
T(B(0,1))\supset B(0,2c).
\label{eq:2.8}
\end{equation}

\emph{证明.} 置 $X_n=nT(B(0,1))$.因 $T$ 是满射, 有 $\bigcup_{n=1}^\infty X_n=F$;
由 Baire 纲定理, 存在某个 $n_0$ 使 $\Int(X_{n_0})\ne\emptyset$.由此
\begin{equation*}
\Int\bigl[T(B(0,1))\bigr]\ne\emptyset.
\end{equation*}
取 $c>0$ 与 $y_0\in F$ 使得
\begin{equation}
B(y_0,4c)\subset T(B(0,1)).
\label{eq:2.9}
\end{equation}
特别地, $y_0\in T(B(0,1))$; 又由对称性,
\begin{equation}
-y_0\in T(B(0,1)).
\label{eq:2.10}
\end{equation}
把 \eqref{eq:2.9} 与 \eqref{eq:2.10} 相加得
\begin{equation*}
B(0,4c)\subset T(B(0,1))+T(B(0,1)).
\end{equation*}
另一方面, 因 $T(B(0,1))$ 是凸的, 有
\begin{equation*}
T(B(0,1))+T(B(0,1))=2T(B(0,1)),
\end{equation*}
于是 \eqref{eq:2.8} 成立.

\emph{第二步.} 设 $T$ 为从 $E$ 到 $F$ 的连续线性算子且满足 \eqref{eq:2.8}.则
\begin{equation}
T(B(0,1))\supset B(0,c).
\label{eq:2.11}
\end{equation}

\emph{证明.} 任取 $y\in F$ 且 $\norm{y}<c$.目标是找 $x\in E$ 使得
$\norm{x}<1$ 且 $Tx=y$.由 \eqref{eq:2.8} 知
\begin{equation}
\forall\varepsilon>0\quad\exists z\in E\ \text{使}\ \norm{z}<\frac12\
\text{且}\ \norm{y-Tz}<\varepsilon.
\label{eq:2.12}
\end{equation}
取 $\varepsilon=c/2$, 找到 $z_1\in E$ 使得
\begin{equation*}
\norm{z_1}<\frac12\quad\text{且}\quad \norm{y-Tz_1}<\frac{c}{2}.
\end{equation*}
对 $y-Tz_1$ (代替 $y$) 施以同样的构造, 取 $\varepsilon=c/4$, 找到 $z_2\in E$ 使得
\begin{equation*}
\norm{z_2}<\frac14\quad\text{且}\quad \norm{(y-Tz_1)-Tz_2}<\frac{c}{4}.
\end{equation*}
如此继续, 用归纳法得到一列 $(z_n)$ 使得
\begin{equation*}
\norm{z_n}<\frac{1}{2^n}\quad\text{且}\quad
\norm{y-T(z_1+z_2+\cdots+z_n)}<\frac{c}{2^n}\quad \forall n.
\end{equation*}
由此可知数列 $x_n=z_1+z_2+\cdots+z_n$ 是 Cauchy 列.设 $x_n\to x$, 显然
$\norm{x}<1$ 且 $y=Tx$ (因 $T$ 连续).
\end{proof}

\begin{keytheorem}[闭图像定理]
\label{thm:2.9}
设 $E$ 与 $F$ 为两个 Banach 空间.设 $T$ 为从 $E$ 到 $F$ 的线性算子.假设 $T$ 的
图像 $G(T)$ 在 $E\times F$ 中闭.则 $T$ 连续.
\end{keytheorem}

\begin{remark}
\label{rem:2.6}
逆命题显然成立, 因为任何连续映射 (线性或非线性) 的图像都是闭的.
\end{remark}

\begin{proof}[定理 \ref{thm:2.9}]
在 $E$ 上考虑两个范数
\begin{equation*}
\norm{x}_1=\norm{x}_E+\norm{Tx}_F\quad\text{与}\quad \norm{x}_2=\norm{x}_E
\end{equation*}
(范数 $\norm{\cdot}_1$ 称为\emph{图像范数}).

利用 $G(T)$ 闭这一假设, 容易验证 $E$ 对范数 $\norm{\cdot}_1$ 是 Banach 空间.
另一方面, $E$ 对范数 $\norm{\cdot}_2$ 也是 Banach 空间, 且
$\norm{\cdot}_2\le\norm{\cdot}_1$.由推论  \ref{cor:2.8}, 这两个范数等价, 因而存在
常数 $c>0$ 使 $\norm{x}_1\le c\norm{x}_2$.于是 $\norm{Tx}_F\le c\norm{x}_E$.
\end{proof}

\starsection{补子空间. 线性算子的右可逆性与左可逆性}
\label{sec:2.4}

我们先给出 Banach 空间中闭子空间的若干几何性质, 它们可由开映射定理推出.

\begin{startheorem}
\label{thm:2.10}
设 $E$ 为 Banach 空间.假设 $G$ 与 $L$ 为两个闭线性子空间, 且 $G+L$ 闭.则存在
常数 $C\ge0$ 使得
\begin{equation}
\left\{\begin{aligned}
	&\text{每个 } z\in G+L \text{ 都可分解为}\\
	&z=x+y,\quad \text{其中 } x\in G,\ y\in L,\
	\norm{x}\le C\norm{z},\ \norm{y}\le C\norm{z}.
\end{aligned}\right.
\label{eq:2.13}
\end{equation}
\end{startheorem}

\begin{proof}
考虑乘积空间 $G\times L$, 赋以范数
\begin{equation*}
\norm{[x,y]}=\norm{x}+\norm{y},
\end{equation*}
以及赋以 $E$ 的范数的空间 $G+L$.

映射 $T:G\times L\to G+L$, $T[x,y]=x+y$, 是连续, 线性且满射的.由开映射定理,
存在常数 $c>0$, 使得每个满足 $\norm{z}<c$ 的 $z\in G+L$ 都可写成 $z=x+y$, 其中
$x\in G$, $y\in L$, 且 $\norm{x}+\norm{y}<1$.由齐次性, 每个 $z\in G+L$ 都可写成
\begin{equation*}
z=x+y\quad \text{其中 } x\in G,\ y\in L,\ \text{且}\
\norm{x}+\norm{y}\le(1/c)\norm{z}.
\end{equation*}
\end{proof}

\begin{starcorollary}
\label{cor:2.11}
在与定理  \ref{thm:2.10} 相同的假设下, 存在常数 $C$ 使得
\begin{equation}
\dist(x,G\cap L)\le C\{\dist(x,G)+\dist(x,L)\}\quad \forall x\in E.
\label{eq:2.14}
\end{equation}
\end{starcorollary}

\begin{proof}
给定 $x\in E$ 与 $\varepsilon>0$, 存在 $a\in G$ 与 $b\in L$ 使得
\begin{equation*}
\norm{x-a}\le\dist(x,G)+\varepsilon,\qquad
\norm{x-b}\le\dist(x,L)+\varepsilon.
\end{equation*}
把性质 \eqref{eq:2.13} 用于 $z=a-b$, 知存在 $a'\in G$ 与 $b'\in L$ 使得
\begin{equation*}
a-b=a'+b',\qquad \norm{a'}\le C\norm{a-b},\qquad \norm{b'}\le C\norm{a-b}.
\end{equation*}
于是 $a-a'\in G\cap L$, 且
\begin{align*}
\dist(x,G\cap L)&\le\norm{x-(a-a')}\le\norm{x-a}+\norm{a'}\\
&\le\norm{x-a}+C\norm{a-b}\le\norm{x-a}+C\bigl(\norm{x-a}+\norm{x-b}\bigr)\\
&\le(1+C)\dist(x,G)+\dist(x,L)+(1+2C)\varepsilon.
\end{align*}
最后令 $\varepsilon\to0$ 即得 \eqref{eq:2.14}.
\end{proof}

\begin{remark}
\label{rem:2.7}
推论  \ref{cor:2.11} 的逆命题也成立: 若 $G$ 与 $L$ 是两个闭线性子空间且
\eqref{eq:2.14} 成立, 则 $G+L$ 闭 (见习题  \ref{exer:2.16}).
\end{remark}

\begin{definition}
设 $G\subset E$ 为 Banach 空间 $E$ 的闭子空间.称子空间 $L\subset E$ 为 $G$ 的
\emph{拓扑\index{bukongjian@补空间 (complementary subspaces)}
补空间}, 或简称为 $G$ 的\emph{补空间}, 如果
\begin{enumerate}[label=(\roman*), leftmargin=2.4em, itemsep=0.15em]
	\item $L$ 闭,
	\item $G\cap L=\{0\}$ 且 $G+L=E$.
\end{enumerate}
此时也说 $G$ 与 $L$ 是 $E$ 的\emph{互补子空间}.若如此, 则每个 $z\in E$ 都可唯一地
写成 $z=x+y$, 其中 $x\in G$, $y\in L$.

由定理  \ref{thm:2.10} 可知, 投影算子 $z\mapsto x$ 与 $z\mapsto y$ 都是连续线性算子.
(这一性质也可以作为互补子空间的定义.)
\end{definition}

\subsection*{例}

\begin{enumerate}[label=\arabic*., leftmargin=2.2em, itemsep=0.5em]
	\item \emph{每个有限维子空间} $G$ 都有补空间.事实上, 设 $e_1,e_2,\dots,e_n$ 为
	      $G$ 的一组基.每个 $x\in G$ 可写成 $x=\sum_{i=1}^n x_ie_i$.置
	      $\varphi_i(x)=x_i$.利用 Hahn--Banach 定理 (解析形式) ——更确切地说, 推论
	       \ref{cor:1.2}——每个 $\varphi_i$ 都可延拓为定义在 $E$ 上的连续线性泛函
	      $\widetilde{\varphi}_i$.容易验证
	      $L=\bigcap_{i=1}^n(\widetilde{\varphi}_i)^{-1}(0)$ 是 $G$ 的补空间.
	\item \emph{每个有限余维的闭子空间} $G$ 都有补空间.只须取任一有限维空间 $L$,
	      使得 $G\cap L=\{0\}$ 且 $G+L=E$ ($L$ 因有限维而闭).
\end{enumerate}

下面是这类情形的一个典型例子.设 $N\subset E^\star$ 为 $p$ 维子空间.则
\begin{equation*}
G=\{x\in E;\ \dual{f}{x}=0\ \forall f\in N\}=N^\perp
\end{equation*}
是闭的, 且余维为 $p$.事实上, 设 $f_1,f_2,\dots,f_p$ 为 $N$ 的一组基.则存在
$e_1,e_2,\dots,e_p\in E$ 使得
\begin{equation*}
\dual{f_i}{e_j}=\delta_{ij}\quad \forall i,j=1,2,\dots,p.
\end{equation*}
[考虑由
\begin{equation}
\Phi(x)=\bigl(\dual{f_1}{x},\dual{f_2}{x},\dots,\dual{f_p}{x}\bigr)
\label{eq:2.15}
\end{equation}
定义的映射 $\Phi:E\to\R^p$, 注意 $\Phi$ 是满射; 否则由 Hahn--Banach 定理 (第二
几何形式), 会存在某个 $\alpha=(\alpha_1,\alpha_2,\dots,\alpha_p)\ne0$, 使得
\begin{equation*}
\alpha\cdot\Phi(x)=\Bigl\langle\sum_{i=1}^p\alpha_if_i,\ x\Bigr\rangle=0
\quad \forall x\in E,
\end{equation*}
这是荒谬的.]

容易验证向量组 $(e_i)_{1\le i\le p}$ 线性无关, 且由这些 $e_i$ 张成的空间是 $G$ 的
补空间.$N^\perp$ 的余维等于 $N$ 的维数这一事实的另一证明见第~
\ref{chap:11}~章 (命题~
\ref{prop:11.11}).

\begin{enumerate}[label=\arabic*., leftmargin=2.2em, itemsep=0.5em, start=3]
	\item 在 Hilbert 空间中, 每个闭子空间都有补空间 (见第~\ref{sec:5.2}~节).
\end{enumerate}

\begin{remark}
\label{rem:2.8}
重要的是要知道: 某些闭子空间 (即使在自反 Banach 空间中) \emph{没有}补空间.
事实上, J.~Lindenstrauss 与 L.~Tzafriri \bref{LindenstraussTzafriri}{1} 的一个著名
结果表明: 在每个不同构于 Hilbert 空间的 Banach 空间中, 都存在\emph{没有}补空间的
闭子空间.
\end{remark}

\begin{definition}
设 $T\in\LL(E,F)$.称算子 $S\in\LL(F,E)$ 为 $T$ 的\emph{右逆}, 如果
$T\circ S=I_F$.称算子 $S\in\LL(F,E)$ 为 $T$ 的\emph{左逆}, 如果
$S\circ T=I_E$.
\end{definition}

下面的结果给出这类逆存在的充分必要条件.

\begin{startheorem}
\label{thm:2.12}
假设 $T\in\LL(E,F)$ 是满射.则下述性质等价:
\begin{enumerate}[label=(\roman*), leftmargin=2.4em, itemsep=0.15em]
	\item $T$ 有右逆.
	\item $N(T)=T^{-1}(0)$ 在 $E$ 中有补空间.
\end{enumerate}
\end{startheorem}

\begin{proof}
(i) $\Rightarrow$ (ii). 设 $S$ 为 $T$ 的右逆.容易看出 (请验证) $R(S)=S(F)$ 是
$N(T)$ 在 $E$ 中的补空间.

(ii) $\Rightarrow$ (i). 设 $L$ 为 $N(T)$ 的补空间.设 $P$ 为从 $E$ 到 $L$ 的
(连续) 投影算子.给定 $f\in F$, 以 $x$ 记方程 $Tx=f$ 的任一解.置 $Sf=Px$.
注意 $S$ 与 $x$ 的选取无关.容易验证 $S\in\LL(F,E)$ 且 $T\circ S=I_F$.
\end{proof}

\begin{remark}
\label{rem:2.9}
由注记 \ref{rem:2.8} 与定理  \ref{thm:2.12}, 很容易构造没有右逆的满射算子 $T$.
事实上, 设 $G\subset E$ 为没有补空间的闭子空间, 令 $F=E/G$, 并令 $T$ 为从 $E$ 到
$F$ 的典范投影 (商空间的定义与性质见第~\ref{sec:11.2}~节).
\end{remark}

\begin{startheorem}
\label{thm:2.13}
假设 $T\in\LL(E,F)$ 是单射.则下述性质等价:
\begin{enumerate}[label=(\roman*), leftmargin=2.4em, itemsep=0.15em]
	\item $T$ 有左逆.
	\item $R(T)=T(E)$ 闭, 且在 $F$ 中有补空间.
\end{enumerate}
\end{startheorem}

\begin{proof}
(i) $\Rightarrow$ (ii). 容易验证 $R(T)$ 闭, 且 $N(S)$ 是 $R(T)$ 的补空间
[写 $f=TSf+(f-TSf)$].

(ii) $\Rightarrow$ (i). 设 $P$ 为从 $F$ 到 $R(T)$ 的连续投影算子.设 $f\in F$;
因 $Pf\in R(T)$, 存在唯一的 $x\in E$ 使 $Tx=Pf$.置 $Sf=x$.显然
$S\circ T=I_E$; 此外, 由推论  \ref{cor:2.7}, $S$ 连续.
\end{proof}

\starsection{再论正交性}
\label{sec:2.5}

有一些简单的公式, 给出和空间或交空间的正交表达.

\begin{proposition}
\label{prop:2.14}
设 $G$ 与 $L$ 为 $E$ 中的两个闭子空间.则
\begin{equation}
G\cap L=(G^\perp+L^\perp)^\perp,
\label{eq:2.16}
\end{equation}
\begin{equation}
G^\perp\cap L^\perp=(G+L)^\perp.
\label{eq:2.17}
\end{equation}
\end{proposition}

\begin{proof}[\eqref{eq:2.16}]
显然 $G\cap L\subset(G^\perp+L^\perp)^\perp$; 事实上, 若 $x\in G\cap L$ 且
$f\in G^\perp+L^\perp$, 则 $\dual{f}{x}=0$.反之, 我们有
$G^\perp\subset G^\perp+L^\perp$, 从而 $(G^\perp+L^\perp)^\perp\subset
G^{\perp\perp}=G$ (注意: 若 $N_1\subset N_2$, 则 $N_2^\perp\subset N_1^\perp$);
同理 $(G^\perp+L^\perp)^\perp\subset L$.因此 $(G^\perp+L^\perp)^\perp\subset
G\cap L$.
\end{proof}

\begin{proof}[\eqref{eq:2.17}]
用与 \eqref{eq:2.16} 的证明相同的论证.
\end{proof}

\begin{corollary}
\label{cor:2.15}
设 $G$ 与 $L$ 为 $E$ 中的两个闭子空间.则
\begin{equation}
(G\cap L)^\perp\supset G^\perp+L^\perp,
\label{eq:2.18}
\end{equation}
\begin{equation}
(G^\perp\cap L^\perp)^\perp=\overline{G+L}.
\label{eq:2.19}
\end{equation}
\end{corollary}

\begin{proof}
利用命题  \ref{prop:1.9} 与  \ref{prop:2.14}.
\end{proof}

下面是一个更深刻的结果.

\begin{startheorem}
\label{thm:2.16}
设 $G$ 与 $L$ 为 Banach 空间 $E$ 中的两个闭子空间.则下述性质等价:
\begin{enumerate}[label=(\alph*), leftmargin=2.0em, itemsep=0.15em]
	\item $G+L$ 在 $E$ 中闭,
	\item $G^\perp+L^\perp$ 在 $E^\star$ 中闭,
	\item $G+L=(G^\perp\cap L^\perp)^\perp$,
	\item $G^\perp+L^\perp=(G\cap L)^\perp$.
\end{enumerate}
\end{startheorem}

\begin{proof}
(a) $\Longleftrightarrow$ (c) 由 \eqref{eq:2.19} 得出.(d) $\Rightarrow$ (b) 是
显然的.

剩下的是蕴含关系 (a) $\Rightarrow$ (d) 与 (b) $\Rightarrow$ (a).

(a) $\Rightarrow$ (d). 由 \eqref{eq:2.18}, 只须证明
$(G\cap L)^\perp\subset G^\perp+L^\perp$.给定 $f\in(G\cap L)^\perp$, 考虑如下
定义的泛函 $\varphi:G+L\to\R$.对每个 $x\in G+L$, 写 $x=a+b$, 其中 $a\in G$, $b\in L$.置
\begin{equation*}
\varphi(x)=\dual{f}{a}.
\end{equation*}
显然 $\varphi$ 与 $x$ 的分解无关, 且 $\varphi$ 线性.另一方面, 由定理
 \ref{thm:2.10}, 我们可以选取 $x$ 的分解使得 $\norm{a}\le C\norm{x}$, 于是
\begin{equation*}
\abs{\varphi(x)}\le C\norm{x}\quad \forall x\in G+L.
\end{equation*}
把 $\varphi$ 延拓为整个 $E$ 上的连续线性泛函 $\widetilde{\varphi}$ (见推论
 \ref{cor:1.2}).于是
\begin{equation*}
f=(f-\widetilde{\varphi})+\widetilde{\varphi}
\quad\text{其中}\quad f-\widetilde{\varphi}\in G^\perp,\quad
\widetilde{\varphi}\in L^\perp.
\end{equation*}

(b) $\Rightarrow$ (a). 由推论  \ref{cor:2.11} 知存在常数 $C$ 使得
\begin{equation}
\dist(f,G^\perp\cap L^\perp)\le C\{\dist(f,G^\perp)+\dist(f,L^\perp)\}
\quad \forall f\in E^\star.
\label{eq:2.20}
\end{equation}
另一方面, 我们有
\begin{equation}
\dist(f,G^\perp)=\sup_{\substack{x\in G\\ \norm{x}\le1}}\dual{f}{x}
\quad \forall f\in E^\star.
\label{eq:2.21}
\end{equation}
[对 $\varphi(x)=I_{B_E}(x)-\dual{f}{x}$ 与 $\psi(x)=I_G(x)$ 应用定理
 \ref{thm:1.12}, 其中 $B_E=\{x\in E;\ \norm{x}\le1\}$.]

同理,
\begin{equation}
\dist(f,L^\perp)=\sup_{\substack{x\in L\\ \norm{x}\le1}}\dual{f}{x}
\quad \forall f\in E^\star,
\label{eq:2.22}
\end{equation}
并且 (由 \eqref{eq:2.17})
\begin{equation}
\dist(f,G^\perp\cap L^\perp)=\dist(f,(G+L)^\perp)
=\sup_{\substack{x\in G+L\\ \norm{x}\le1}}\dual{f}{x}\quad \forall f\in E^\star.
\label{eq:2.23}
\end{equation}
合并 \eqref{eq:2.20}, \eqref{eq:2.21}, \eqref{eq:2.22} 与 \eqref{eq:2.23} 得
\begin{equation}
\sup_{\substack{x\in G+L\\ \norm{x}\le1}}\dual{f}{x}
\le C\Bigl\{\sup_{\substack{x\in G\\ \norm{x}\le1}}\dual{f}{x}
+\sup_{\substack{x\in L\\ \norm{x}\le1}}\dual{f}{x}\Bigr\}
\quad \forall f\in E^\star.
\label{eq:2.24}
\end{equation}
由 \eqref{eq:2.24} 推出
\begin{equation}
\overline{B_G+B_L}\supset\frac{1}{C}B_{\overline{G+L}}.
\label{eq:2.25}
\end{equation}
事实上, 用反证法, 若存在 $x_0\in\overline{G+L}$ 满足 $\norm{x_0}\le1/C$ 且
$x_0\notin\overline{B_G+B_L}$, 则在 $E$ 中存在一张闭超平面把 $\{x_0\}$ 与
$\overline{B_G+B_L}$ 严格分离.于是存在 $f_0\in E^\star$ 与 $\alpha\in\R$ 使得
\begin{equation*}
\dual{f_0}{x}<\alpha<\dual{f_0}{x_0}\quad \forall x\in\overline{B_G+B_L}.
\end{equation*}
因此
\begin{equation*}
\sup_{\substack{x\in G\\ \norm{x}\le1}}\dual{f_0}{x}
+\sup_{\substack{x\in L\\ \norm{x}\le1}}\dual{f_0}{x}
\le\alpha<\dual{f_0}{x_0},
\end{equation*}
这与 \eqref{eq:2.24} 矛盾, 故 \eqref{eq:2.25} 得证.

最后, 考虑赋以范数
\begin{equation*}
\norm{[x,y]}=\max\{\norm{x},\norm{y}\}
\end{equation*}
的空间 $X=G\times L$, 以及赋以 $E$ 的范数的空间 $Y=\overline{G+L}$.由
$T([x,y])=x+y$ 定义的映射 $T:X\to Y$ 线性且连续.由 \eqref{eq:2.25} 知
\begin{equation*}
\overline{T(B_X)}\supset\frac{1}{C}B_Y.
\end{equation*}
利用定理  \ref{thm:2.6} (开映射定理) 证明中的第二步, 得到
\begin{equation*}
T(B_X)\supset\frac{1}{2C}B_Y.
\end{equation*}
由此可知 $T$ 从 $X$ 映上 $Y$, 即 $G+L=\overline{G+L}$.
\end{proof}

\section{\index{wujiexianxingsuanzi@无界线性算子 (unbounded linear operator)}
无界线性算子引论. \index{bansuisuanzi@伴随算子 (adjoint operator)}
伴随算子的定义}
\label{sec:2.6}

\begin{definition}
设 $E$ 与 $F$ 为两个 Banach 空间.从 $E$ 到 $F$ 的\emph{无界线性算子}是指定义在
线性子空间 $D(A)\subset E$ 上, 取值于 $F$ 的线性映射 $A:D(A)\subset E\to F$.集合
$D(A)$ 称为 $A$ 的\emph{定义域}.

若 $D(A)=E$ 且存在常数 $c\ge0$ 使
\begin{equation*}
\norm{Au}\le c\norm{u}\quad \forall u\in E,
\end{equation*}
则称 $A$ 是\emph{有界的} (或\emph{连续的}).

有界算子的范数定义为
\begin{equation*}
\norm{A}_{\LL(E,F)}=\operatorname*{Sup}_{u\ne0}\frac{\norm{Au}}{\norm{u}}.
\end{equation*}
\end{definition}

\begin{remark}
\label{rem:2.10}
当然可能出现这样的情况: 一个``无界线性算子''实际上是有界的.这个术语稍欠一致,
但通行且不会引起混淆.
\end{remark}

以下是一些重要的定义与进一步的记号:
\begin{equation*}
\boxed{\text{$A$ 的图像}=G(A)=\{[u,Au];\ u\in D(A)\}\subset E\times F,}
\end{equation*}
\begin{equation*}
\boxed{\text{$A$ 的值域}=R(A)=\{Au;\ u\in D(A)\}\subset F,}
\end{equation*}
\begin{equation*}
\boxed{\text{$A$ 的核}=N(A)=\{u\in D(A);\ Au=0\}\subset E.}
\end{equation*}
若 $G(A)$ 在 $E\times F$ 中闭, 则称映射 $A$ 是\emph{闭的}.

\begin{keyremark}
\label{rem:2.11}
要证明算子 $A$ 是闭的, 一般按如下步骤进行.取 $D(A)$ 中的序列 $(u_n)$, 使
$u_n\to u$ (在 $E$ 中) 且 $Au_n\to f$ (在 $F$ 中).然后验证两件事:
\begin{enumerate}[label=(\alph*), leftmargin=2.0em, itemsep=0.15em]
	\item $u\in D(A)$,
	\item $f=Au$.
\end{enumerate}
注意: 只考虑满足 $u_n\to0$ (在 $E$ 中) 且 $Au_n\to f$ (在 $F$ 中) 的序列
(并证明 $f=0$) 是\emph{不够}的.
\end{keyremark}

\begin{remark}
\label{rem:2.12}
若 $A$ 闭, 则 $N(A)$ 闭; 然而 $R(A)$ 不必闭.
\end{remark}

\begin{remark}
\label{rem:2.13}
实践中, 大多数无界算子是\emph{闭}的且是\emph{稠定}的, 即 $D(A)$ 在 $E$ 中稠密.
\end{remark}

\runin{伴随算子 $A^\star$ 的定义.}
设 $A:D(A)\subset E\to F$ 为稠定的无界线性算子.我们按下述方式引入无界算子
$A^\star:D(A^\star)\subset F^\star\to E^\star$.首先定义它的定义域:
\begin{equation*}
D(A^\star)=\{v\in F^\star;\ \exists c\ge0\ \text{使得}\
\abs{\dual{v}{Au}}\le c\norm{u}\ \forall u\in D(A)\}.
\end{equation*}
显然 $D(A^\star)$ 是 $F^\star$ 的线性子空间.现在定义 $A^\star v$.给定
$v\in D(A^\star)$, 考虑由
\begin{equation*}
g(u)=\dual{v}{Au}\quad \forall u\in D(A)
\end{equation*}
定义的映射 $g:D(A)\to\R$.我们有
\begin{equation*}
\abs{g(u)}\le c\norm{u}\quad \forall u\in D(A).
\end{equation*}
由 Hahn--Banach 定理 (解析形式; 见定理  \ref{thm:1.1}), 存在线性映射
$f:E\to\R$, 它延拓 $g$, 且满足
\begin{equation*}
\abs{f(u)}\le c\norm{u}\quad \forall u\in E.
\end{equation*}
于是 $f\in E^\star$.注意 $g$ 的延拓是唯一的, 因为 $D(A)$ 在 $E$ 中稠密.

置
\begin{equation*}
A^\star v=f.
\end{equation*}
这个无界线性算子 $A^\star:D(A^\star)\subset F^\star\to E^\star$ 称为 $A$ 的
\emph{伴随算子}.简言之, $A$ 与 $A^\star$ 之间的基本关系为
\begin{equation*}
\dual{v}{Au}_{F^\star,F}=\dual{A^\star v}{u}_{E^\star,E}
\quad \forall u\in D(A),\ \forall v\in D(A^\star).
\end{equation*}

\begin{remark}
\label{rem:2.14}
延拓 $g$ 不必借助 Hahn--Banach 定理.只须用经典的连续性延拓即可: 因为 $D(A)$
稠密, $g$ 在 $D(A)$ 上一致连续, 而 $\R$ 完备 (见例如 H.~L.~Royden
\bref{Royden}{1} (第~
\ref{chap:7}~章命题 11) 或 J.~Dugundji \bref{Dugundji}{1} (第 XIV 章
定理~
\ref{thm:5.2})).
\end{remark}

\begin{starremark}
\label{rem:2.15}
可能出现 $D(A^\star)$ 在 $F^\star$ 中不稠密的情形 (即使 $A$ 闭); 但这是一种相当
病态的情形 (见习题  \ref{exer:2.22}).下述断言总是成立: 若 $A$ 闭, 则
$D(A^\star)$ 对于第~
\ref{chap:3}~章定义的\wstar{}拓扑 $\sigma(F^\star,F)$ 在 $F^\star$ 中
稠密 (见问题 9).特别地, 若 $F$ 自反, 则 $D(A^\star)$ 对通常的 (范数) 拓扑在
$F^\star$ 中稠密 (见定理  \ref{thm:3.24}).
\end{starremark}

\begin{remark}
\label{rem:2.16}
若 $A$ 是有界算子, 则 $A^\star$ 也是有界算子 (从 $F^\star$ 到 $E^\star$), 并且
\begin{equation*}
\norm{A^\star}_{\LL(F^\star,E^\star)}=\norm{A}_{\LL(E,F)}.
\end{equation*}
事实上, 显然 $D(A^\star)=F^\star$.由基本关系,
\begin{equation*}
\abs{\dual{A^\star v}{u}}\le\norm{A}\norm{u}\norm{v}
\quad \forall u\in E,\ \forall v\in F^\star,
\end{equation*}
由此推出 $\norm{A^\star v}\le\norm{A}\norm{v}$, 从而 $\norm{A^\star}\le\norm{A}$.

我们还有
\begin{equation*}
\abs{\dual{v}{Au}}\le\norm{A^\star}\norm{u}\norm{v}
\quad \forall u\in E,\ \forall v\in F^\star,
\end{equation*}
由此 (由推论  \ref{cor:1.4}) 推出 $\norm{Au}\le\norm{A^\star}\norm{u}$, 从而
$\norm{A}\le\norm{A^\star}$.
\end{remark}

\begin{proposition}
\label{prop:2.17}
设 $A:D(A)\subset E\to F$ 为稠定的无界线性算子.则 $A^\star$ 闭, 即
$G(A^\star)$ 在 $F^\star\times E^\star$ 中闭.
\end{proposition}

\begin{proof}
设 $v_n\in D(A^\star)$ 使 $v_n\to v$ (在 $F^\star$ 中) 且 $A^\star v_n\to f$
(在 $E^\star$ 中).须验证 (a) $v\in D(A^\star)$ 与 (b) $A^\star v=f$.

我们有
\begin{equation*}
\dual{v_n}{Au}=\dual{A^\star v_n}{u}\quad \forall u\in D(A).
\end{equation*}
取极限得
\begin{equation*}
\dual{v}{Au}=\dual{f}{u}\quad \forall u\in D(A).
\end{equation*}
因此 $v\in D(A^\star)$ (因为
$\abs{\dual{v}{Au}}\le\norm{f}\norm{u}\ \forall u\in D(A)$), 且 $A^\star v=f$.
\end{proof}

\emph{$A$ 与 $A^\star$ 的图像由一个非常简单的正交关系相联系:} 考虑由
\begin{equation*}
I([v,f])=[-f,v]
\end{equation*}
定义的同构 $I:F^\star\times E^\star\to E^\star\times F^\star$.设
$A:D(A)\subset E\to F$ 为稠定的无界线性算子.则
\begin{equation*}
\boxed{I[G(A^\star)]=G(A)^\perp.}
\end{equation*}
事实上, 设 $[v,f]\in F^\star\times E^\star$, 则
\begin{align*}
[v,f]\in G(A^\star)
&\Longleftrightarrow \dual{f}{u}=\dual{v}{Au}\quad \forall u\in D(A)\\
&\Longleftrightarrow -\dual{f}{u}+\dual{v}{Au}=0\quad \forall u\in D(A)\\
&\Longleftrightarrow [-f,v]\in G(A)^\perp.
\end{align*}

下面是值域与核之间的若干标准正交关系:

\begin{corollary}
\label{cor:2.18}
设 $A:D(A)\subset E\to F$ 为稠定且闭的无界线性算子.则
\begin{align*}
\text{(i)}\qquad &N(A)=R(A^\star)^\perp,\\
\text{(ii)}\qquad &N(A^\star)=R(A)^\perp,\\
\text{(iii)}\qquad &N(A)^\perp\supset\overline{R(A^\star)},\\
\text{(iv)}\qquad &N(A^\star)^\perp=\overline{R(A)}.
\end{align*}
\end{corollary}

\begin{proof}
注意 (iii) 与 (iv) 直接由 (i), (ii) 结合命题  \ref{prop:1.9} 得到.(i) 与 (ii)
有简单直接的证明 (见习题  \ref{exer:2.18}).不过, 用下述办法把这些事实与命题
 \ref{prop:2.14} 联系起来是有启发性的.考虑空间 $X=E\times F$, 从而
$X^\star=E^\star\times F^\star$, 以及 $X$ 的子空间
\begin{equation*}
G=G(A)\quad\text{与}\quad L=E\times\{0\}.
\end{equation*}
很容易验证
\begin{equation}
N(A)\times\{0\}=G\cap L,
\label{eq:2.26}
\end{equation}
\begin{equation}
E\times R(A)=G+L,
\label{eq:2.27}
\end{equation}
\begin{equation}
\{0\}\times N(A^\star)=G^\perp\cap L^\perp,
\label{eq:2.28}
\end{equation}
\begin{equation}
R(A^\star)\times F^\star=G^\perp+L^\perp.
\label{eq:2.29}
\end{equation}

\emph{(i) 的证明.} 由 \eqref{eq:2.29} 有
\begin{equation*}
R(A^\star)^\perp\times\{0\}=(G^\perp+L^\perp)^\perp=G\cap L
\quad\text{(由 \eqref{eq:2.16})}
=N(A)\times\{0\}\quad\text{(由 \eqref{eq:2.26})}.
\end{equation*}

\emph{(ii) 的证明.} 由 \eqref{eq:2.27} 有
\begin{equation*}
\{0\}\times R(A)^\perp=(G+L)^\perp=G^\perp\cap L^\perp
\quad\text{(由 \eqref{eq:2.17})}
=\{0\}\times N(A^\star)\quad\text{(由 \eqref{eq:2.28})}.
\end{equation*}
\end{proof}

\begin{remark}
\label{rem:2.17}
即使 $A$ 是有界线性算子, 也可能出现 $N(A)^\perp\ne R(A^\star)$ (见习题
 \ref{exer:2.23}).不过, 下述断言总是成立: $N(A)^\perp$ 是 $R(A^\star)$ 关于
\wstar{}拓扑 $\sigma(E^\star,E)$ 的闭包 (见问题 9).特别地, 若 $E$ 自反, 则
$N(A)^\perp=R(A^\star)$.
\end{remark}

\starsection{闭值域算子的刻画. 满射算子的刻画}
\label{sec:2.7}

关于闭值域算子的主要结果如下.

\begin{startheorem}
\label{thm:2.19}
设 $A:D(A)\subset E\to F$ 为稠定且闭的无界线性算子.则下述性质等价:
\begin{enumerate}[label=(\roman*), leftmargin=2.4em, itemsep=0.15em]
	\item $R(A)$ 闭,
	\item $R(A^\star)$ 闭,
	\item $R(A)=N(A^\star)^\perp$,
	\item $R(A^\star)=N(A)^\perp$.
\end{enumerate}
\end{startheorem}

\begin{proof}
沿用推论  \ref{cor:2.18} 证明中的记号, 有
\begin{align*}
\text{(i)}&\Longleftrightarrow G+L\ \text{在 } X \text{ 中闭 (见 \eqref{eq:2.27})},\\
\text{(ii)}&\Longleftrightarrow G^\perp+L^\perp\ \text{在 } X^\star \text{ 中闭
(见 \eqref{eq:2.29})},\\
\text{(iii)}&\Longleftrightarrow G+L=(G^\perp\cap L^\perp)^\perp
\ \text{(见 \eqref{eq:2.27} 与 \eqref{eq:2.28})},\\
\text{(iv)}&\Longleftrightarrow (G\cap L)^\perp=G^\perp+L^\perp
\ \text{(见 \eqref{eq:2.26} 与 \eqref{eq:2.29})}.
\end{align*}
结论于是由定理  \ref{thm:2.16} 得到.
\end{proof}

\begin{remark}
\label{rem:2.18}
设 $A:D(A)\subset E\to F$ 为闭的无界线性算子.则 $R(A)$ 闭当且仅当存在常数 $C$
使得
\begin{equation*}
\dist(u,N(A))\le C\norm{Au}\quad \forall u\in D(A);
\end{equation*}
见习题  \ref{exer:2.14}.
\end{remark}

下一个结果给出满射算子的一个有用刻画.

\begin{startheorem}
\label{thm:2.20}
设 $A:D(A)\subset E\to F$ 为稠定且闭的无界线性算子.则下述性质等价:
\begin{enumerate}[label=(\alph*), leftmargin=2.0em, itemsep=0.15em]
	\item $A$ 是满射, 即 $R(A)=F$,
	\item 存在常数 $C$ 使得
	      \begin{equation*}
	      \norm{v}\le C\norm{A^\star v}\quad \forall v\in D(A^\star),
	      \end{equation*}
	\item $N(A^\star)=\{0\}$ 且 $R(A^\star)$ 闭.
\end{enumerate}
\end{startheorem}

\begin{remark}
\label{rem:2.19}
蕴含关系 (b) $\Rightarrow$ (a) 在实践中有时很有用, 可用来断定算子 $A$ 是满射.
做法如下.假设 $v$ 满足 $A^\star v=f$, 试证 $\norm{v}\le C\norm{f}$ (其中 $C$
与 $f$ 无关).这称为\emph{先验估计}方法.人们并不关心方程 $A^\star v=f$ 是否
有解; 而是先验地假定 $v$ 已给定, 再设法估计其范数.
\end{remark}

\begin{proof}
(a) $\Rightarrow$ (b). 置
\begin{equation*}
B^\star=\{v\in D(A^\star);\ \norm{A^\star v}\le1\}.
\end{equation*}
由齐次性, 只须证明 $B^\star$ 有界.为此——由推论  \ref{cor:2.5} (\index{yizhiyoujieyuanli@一致有界原理 (uniform boundedness principle)}
一致有界原理)
——只须证明: 对任意给定的 $f_0\in F$, 集合 $\dual{B^\star}{f_0}$ 有界 (在 $\R$
中).因 $A$ 是满射, 存在 $u_0\in D(A)$ 使 $Au_0=f_0$.对每个 $v\in B^\star$ 有
\begin{equation*}
\dual{v}{f_0}=\dual{v}{Au_0}=\dual{A^\star v}{u_0},
\end{equation*}
从而 $\abs{\dual{v}{f_0}}\le\norm{u_0}$.

(b) $\Rightarrow$ (c). 设 $f_n=A^\star v_n\to f$.对 $v_n-v_m$ 使用 (b) 可知
$(v_n)$ 是 Cauchy 列, 于是 $v_n\to v$.因 $A^\star$ 闭 (由命题  \ref{prop:2.17}),
得 $A^\star v=f$.

(c) $\Rightarrow$ (a). 因 $R(A^\star)$ 闭, 由定理  \ref{thm:2.19} 得
$R(A)=N(A^\star)^\perp=F$.
\end{proof}

还有一个``对偶''命题.

\begin{startheorem}
\label{thm:2.21}
设 $A:D(A)\subset F$ 为稠定且闭的无界线性算子.则下述性质等价:
\begin{enumerate}[label=(\alph*), leftmargin=2.0em, itemsep=0.15em]
	\item $A^\star$ 是满射, 即 $R(A^\star)=E^\star$,
	\item 存在常数 $C$ 使得
	      \begin{equation*}
	      \norm{u}\le C\norm{Au}\quad \forall u\in D(A),
	      \end{equation*}
	\item $N(A)=\{0\}$ 且 $R(A)$ 闭.
\end{enumerate}
\end{startheorem}

\begin{proof}
证明与定理  \ref{thm:2.20} 类似, 留作练习.
\end{proof}

\begin{remark}
\label{rem:2.20}
若假设 $\dim E<\infty$ 或 $\dim F<\infty$, 则下述各条等价:
\begin{align*}
&A \text{ 满射} \Longleftrightarrow A^\star \text{ 单射},\\
&A^\star \text{ 满射} \Longleftrightarrow A \text{ 单射},
\end{align*}
这确实是有限维空间中线性算子的经典结果.这些等价关系成立的原因是 $R(A)$ 与
$R(A^\star)$ 都是有限维的 (从而闭).

在一般情形下, 只有蕴含关系
\begin{align*}
&A \text{ 满射} \Rightarrow A^\star \text{ 单射},\\
&A^\star \text{ 满射} \Rightarrow A \text{ 单射}.
\end{align*}
逆命题不成立, 从下面这个简单的例子可以看出.设 $E=F=\ell^2$; 对每个
$x\in\ell^2$ 写 $x=(x_n)_{n\ge1}$, 并置 $Ax=\bigl(\frac1n x_n\bigr)_{n\ge1}$.
容易看出 $A$ 是有界算子且 $A^\star=A$; $A^\star$ (相应地 $A$) 是单射, 但 $A$
(相应地 $A^\star$) 不是满射; $R(A)$ (相应地 $R(A^\star)$) 稠密但不闭.
\end{remark}

\section*{第 2 章评注}
\addcontentsline{toc}{section}{第 2 章评注}
\markright{第 2 章评注}

\cmtnum{1.}
可以显式地写出一些没有补空间的简单闭子空间.例如 $c_0$ 是 $\ell^\infty$ 中
没有补空间的闭子空间; 见例如 C.~DeVito \bref{DeVito}{1} (记号 $c_0$ 与
$\ell^\infty$ 的说明见第~\ref{sec:11.3}~节).其他例子见 W.~Rudin \bref{Rudin}{1} ($L^1$
的一个子空间), G.~Köthe \bref{Kothe}{1} 以及 B.~Beauzamy \bref{Beauzamy}{1}
($\ell^p$ 的一个子空间, $p\ne2$).

\cmtnum{2.}
第~
\ref{chap:2}~章的大部分结果都可推广到 \emph{Fréchet 空间} (可度量化且完备的局部凸
空间).可能的推广方式很多; 见例如 H.~Schaefer \bref{Schaefer}{1}, J.~Horváth \bref{Horvath}{1}, R.~Edwards \bref{Edwards}{1}, F.~Treves
\bref{Treves}{1}, \bref{Treves}{3}, G.~Köthe \bref{Kothe}{1}.这些推广的动机
来自\emph{广义函数理论} (见 L.~Schwartz \bref{Schwartz}{1}), 其中许多重要的空间
\emph{不是} Banach 空间.关于在偏微分方程理论中的应用, 读者可参阅
L.~Hörmander \bref{Hormander}{1} 或 F.~Treves \bref{Treves}{1}, \bref{Treves}{2},
\bref{Treves}{3}.

\cmtnum{3.}
第~\ref{sec:2.5}~节的结果在 T.~Kato \bref{Kato}{1} 中有各种推广.

\section*{第 2 章习题}
\addcontentsline{toc}{section}{第 2 章习题}
\markright{第 2 章习题}

\begin{exercise}[凸函数的连续性.]
\label{exer:2.1}
设 $E$ 为 Banach 空间, $\varphi:E\to(-\infty,+\infty]$ 为凸的 l.s.c. 函数.
假设 $x_0\in\Int D(\varphi)$.

\begin{enumerate}[label=\arabic*., leftmargin=2.2em, itemsep=0.25em]
	\item 证明存在两个常数 $R>0$ 与 $M$, 使得
	      \begin{equation*}
	      \varphi(x)\le M\quad \forall x\in E\ \text{满足}\ \norm{x-x_0}\le R.
	      \end{equation*}
	      [提示: 取适当的 $\rho>0$, 考虑集合
	      \begin{equation*}
	      F_n=\{x\in E;\ \norm{x-x_0}\le\rho\ \text{且}\ \varphi(x)\le n\}.]
	      \end{equation*}
	\item 证明: $\forall r<R$, $\exists L\ge0$ 使得
	      \begin{equation*}
	      \abs{\varphi(x_1)-\varphi(x_2)}\le L\norm{x_1-x_2}
	      \quad \forall x_1,x_2\in E\ \text{满足}\
	      \norm{x_i-x_0}\le r,\ i=1,2.
	      \end{equation*}
	      更精确地说, 可以取 $L=\dfrac{2[M-\varphi(x_0)]}{R-r}$.
\end{enumerate}
\end{exercise}

\begin{exercise}
\label{exer:2.2}
设 $E$ 为向量空间, $p:E\to\R$ 为具有下述三条性质的函数:
\begin{enumerate}[label=(\roman*), leftmargin=2.4em, itemsep=0.15em]
	\item $p(x+y)\le p(x)+p(y)\quad \forall x,y\in E$,
	\item 对每个固定的 $x\in E$, 函数 $\lambda\mapsto p(\lambda x)$ 是从 $\R$ 到
	      $\R$ 的连续函数,
	\item 只要 $E$ 中的序列 $(y_n)$ 满足 $p(y_n)\to0$, 就有
	      $p(\lambda y_n)\to0$ 对每个 $\lambda\in\R$ 成立.
\end{enumerate}
假设 $(x_n)$ 为 $E$ 中的序列, 满足 $p(x_n)\to0$, 而 $(\alpha_n)$ 为 $\R$ 中的
有界序列.证明 $p(0)=0$ 且 $p(\alpha_nx_n)\to0$.

[提示: 给定 $\varepsilon>0$, 考虑集合
\begin{equation*}
F_n=\{\lambda\in\R;\ \abs{p(\lambda x_k)}\le\varepsilon,\ \forall k\ge n\}.]
\end{equation*}

并由此推出: 若 $(x_n)$ 为 $E$ 中的序列, 满足 $p(x_n-x)\to0$ 对某个 $x\in E$
成立, 而 $(\alpha_n)$ 为 $\R$ 中的序列, 满足 $\alpha_n\to\alpha$, 则
$p(\alpha_nx_n)\to p(\alpha x)$.
\end{exercise}

\begin{exercise}
\label{exer:2.3}
设 $E$ 与 $F$ 为两个 Banach 空间, $(T_n)$ 为 $\LL(E,F)$ 中的序列.假设对每个
$x\in E$, $T_nx$ 当 $n\to\infty$ 时收敛, 其极限记为 $Tx$.证明: 若 $x_n\to x$
(在 $E$ 中), 则 $T_nx_n\to Tx$ (在 $F$ 中).
\end{exercise}

\begin{exercise}
\label{exer:2.4}
设 $E$ 与 $F$ 为两个 Banach 空间, $a:E\times F\to\R$ 为满足下述条件的双线性
形式:
\begin{enumerate}[label=(\roman*), leftmargin=2.4em, itemsep=0.15em]
	\item 对每个固定的 $x\in E$, 映射 $y\mapsto a(x,y)$ 连续;
	\item 对每个固定的 $y\in F$, 映射 $x\mapsto a(x,y)$ 连续.
\end{enumerate}
证明存在常数 $C\ge0$ 使得
\begin{equation*}
\abs{a(x,y)}\le C\norm{x}\norm{y}\quad \forall x\in E,\ \forall y\in F.
\end{equation*}
[提示: 引入线性算子 $T:E\to F^\star$, 并借助推论  \ref{cor:2.5} 证明 $T$ 有界.]
\end{exercise}

\begin{exercise}
\label{exer:2.5}
设 $E$ 为 Banach 空间, $\varepsilon_n$ 为一列正数, 满足 $\lim\varepsilon_n=0$.
又设 $(f_n)$ 为 $E^\star$ 中的序列, 具有下述性质:
\begin{equation*}
\left\{\begin{aligned}
	&\exists r>0,\quad \forall x\in E\ \text{且}\ \norm{x}<r,\ \exists C(x)\in\R\
	\text{使得}\\
	&\dual{f_n}{x}\le\varepsilon_n\norm{f_n}+C(x)\quad \forall n.
\end{aligned}\right.
\end{equation*}
证明 $(f_n)$ 有界.

[提示: 引入 $g_n=f_n/(1+\varepsilon_n\norm{f_n})$.]
\end{exercise}

\begin{exercise}[局部有界的非线性单调算子.]
\label{exer:2.6}
设 $E$ 为 Banach 空间, $D(A)$ 为 $E$ 的任意子集.称 (非线性) 映射
$A:D(A)\subset E\to E^\star$ 为\emph{单调}的, 如果它满足
\begin{equation*}
\dual{Ax-Ay}{x-y}\ge0\quad \forall x,y\in D(A).
\end{equation*}

\begin{enumerate}[label=\arabic*., leftmargin=2.2em, itemsep=0.25em]
	\item 设 $x_0\in\Int D(A)$.证明存在两个常数 $R>0$ 与 $C$, 使得
	      \begin{equation*}
	      \norm{Ax}\le C\quad \forall x\in D(A)\ \text{满足}\ \norm{x-x_0}<R.
	      \end{equation*}
	      [提示: 用反证法, 构造 $D(A)$ 中的序列 $(x_n)$, 使得 $x_n\to x_0$ 且
	      $\norm{Ax_n}\to\infty$.取 $r>0$ 使 $B(x_0,r)\subset D(A)$.利用 $A$
	      在 $x_n$ 处与在 $(x_0+x)$ 处 (其中 $\norm{x}<r$) 的单调性.应用习题
	       \ref{exer:2.5}.]
	\item 对点 $x_0\in\Int[\operatorname{conv}D(A)]$ 证明同样的结论.
	\item 把第 1 问的结论推广到 $A$ 为多值的情形, 即对每个 $x\in D(A)$, $Ax$
	      是 $E^\star$ 的非空子集; 单调性定义如下:
	      \begin{equation*}
	      \dual{f-g}{x-y}\ge0\quad \forall x,y\in D(A),\ \forall f\in Ax,\
	      \forall g\in Ay.
	      \end{equation*}
\end{enumerate}
\end{exercise}

\begin{exercise}
\label{exer:2.7}
设 $\alpha=(\alpha_n)$ 为给定的实数列, $1\le p\le\infty$.假设对 $\ell^p$ 中
的每个元素 $x=(x_n)$ 都有 $\sum\abs{\alpha_n}\abs{x_n}<\infty$ ($\ell^p$ 空间的
定义见第~\ref{sec:11.3}~节).

证明 $\alpha\in\ell^{p'}$.
\end{exercise}

\begin{exercise}
\label{exer:2.8}
设 $E$ 为 Banach 空间, $T:E\to E^\star$ 为满足
\begin{equation*}
\dual{Tx}{x}\ge0\quad \forall x\in E
\end{equation*}
的线性算子.证明 $T$ 是有界算子.

[有两种方法: (i) 利用习题  \ref{exer:2.6}; 或 (ii) 应用闭图像定理.]
\end{exercise}

\begin{exercise}
\label{exer:2.9}
设 $E$ 为 Banach 空间, $T:E\to E^\star$ 为满足
\begin{equation*}
\dual{Tx}{y}=\dual{Ty}{x}\quad \forall x,y\in E
\end{equation*}
的线性算子.证明 $T$ 是有界算子.
\end{exercise}

\begin{exercise}
\label{exer:2.10}
设 $E$ 与 $F$ 为两个 Banach 空间, $T\in\LL(E,F)$ 为满射.

\begin{enumerate}[label=\arabic*., leftmargin=2.2em, itemsep=0.25em]
	\item 设 $M$ 为 $E$ 的任意子集.证明: $T(M)$ 在 $F$ 中闭当且仅当
	      $M+N(T)$ 在 $E$ 中闭.
	\item 由此推出: 若 $M$ 为 $E$ 中的闭向量空间且 $\dim N(T)<\infty$, 则
	      $T(M)$ 闭.
\end{enumerate}
\end{exercise}

\begin{exercise}
\label{exer:2.11}
设 $E$ 为 Banach 空间, $F=\ell^1$, $T\in\LL(E,F)$ 为满射.证明存在
$S\in\LL(F,E)$ 使得 $T\circ S=I_F$, 即 $S$ 是 $T$ 的右逆.

[提示: 不要用定理  \ref{thm:2.12}; 试着利用 $\ell^1$ 的典范基显式地定义 $S$.]
\end{exercise}

\begin{exercise}
\label{exer:2.12}
设 $E$ 与 $F$ 为两个 Banach 空间, 其范数分别为 $\norm{\cdot}_E$ 与
$\norm{\cdot}_F$.设 $T\in\LL(E,F)$ 使得 $R(T)$ 闭且 $\dim N(T)<\infty$.以
$\abs{\cdot}$ 记 $E$ 上的另一个范数, 它比 $\norm{\cdot}_E$ 弱, 即
$\abs{x}\le M\norm{x}_E\ \forall x\in E$.

证明存在常数 $C$ 使得
\begin{equation*}
\norm{x}_E\le C\bigl(\norm{Tx}_F+\abs{x}\bigr)\quad \forall x\in E.
\end{equation*}
[提示: 用反证法.]
\end{exercise}

\begin{exercise}
\label{exer:2.13}
设 $E$ 与 $F$ 为两个 Banach 空间.证明集合
\begin{equation*}
\Omega=\{T\in\LL(E,F);\ T\ \text{有左逆}\}
\end{equation*}
在 $\LL(E,F)$ 中开.

[提示: 先证明集合
\begin{equation*}
\mathcal{O}=\{T\in\LL(E,F);\ T\ \text{是双射}\}
\end{equation*}
在 $\LL(E,F)$ 中开.]
\end{exercise}

\begin{exercise}
\label{exer:2.14}
设 $E$ 与 $F$ 为两个 Banach 空间.

\begin{enumerate}[label=\arabic*., leftmargin=2.2em, itemsep=0.25em]
	\item 设 $T\in\LL(E,F)$.证明 $R(T)$ 闭当且仅当存在常数 $C$ 使得
	      \begin{equation*}
	      \dist(x,N(T))\le C\norm{Tx}\quad \forall x\in E.
	      \end{equation*}
	      [提示: 利用商空间 $E/N(T)$; 见第~\ref{sec:11.2}~节.]
	\item 设 $A:D(A)\subset E\to F$ 为闭的无界算子.证明 $R(A)$ 闭当且仅当
	      存在常数 $C$ 使得
	      \begin{equation*}
	      \dist(u,N(A))\le C\norm{Au}\quad \forall u\in D(A).
	      \end{equation*}
	      [提示: 考虑算子 $T:E_0\to F$, 其中 $E_0=D(A)$ 赋以图像范数, 且
	      $T=A$.]
\end{enumerate}
\end{exercise}

\begin{exercise}
\label{exer:2.15}
设 $E_1$, $E_2$ 与 $F$ 为三个 Banach 空间.设 $T_1\in\LL(E_1,F)$,
$T_2\in\LL(E_2,F)$ 使得
\begin{equation*}
R(T_1)\cap R(T_2)=\{0\}\quad\text{且}\quad R(T_1)+R(T_2)=F.
\end{equation*}
证明 $R(T_1)$ 与 $R(T_2)$ 都闭.

[提示: 对由 $T(x_1,x_2)=T_1x_1+T_2x_2$ 定义的映射
$T:E_1\times E_2\to F$ 应用习题  \ref{exer:2.10}.]
\end{exercise}

\begin{exercise}
\label{exer:2.16}
设 $E$ 为 Banach 空间, $G$ 与 $L$ 为 $E$ 的两个闭子空间.假设存在常数 $C$
使得
\begin{equation*}
\dist(x,G\cap L)\le C\dist(x,L),\quad \forall x\in G.
\end{equation*}
证明 $G+L$ 闭.
\end{exercise}

\begin{exercise}
\label{exer:2.17}
设 $E=C([0,1])$, 赋以通常的范数.考虑由
\begin{equation*}
D(A)=C^1([0,1])\quad\text{与}\quad Au=u'=\frac{\dd u}{\dd t}
\end{equation*}
定义的算子 $A:D(A)\subset E\to E$.

\begin{enumerate}[label=\arabic*., leftmargin=2.2em, itemsep=0.25em]
	\item 验证 $\overline{D(A)}=E$.
	\item $A$ 是否闭?
	\item 考虑由
	      \begin{equation*}
	      D(B)=C^2([0,1])\quad\text{与}\quad Bu=u'=\frac{\dd u}{\dd t}
	      \end{equation*}
	      定义的算子 $B:D(B)\subset E\to E$.$B$ 是否闭?
\end{enumerate}
\end{exercise}

\begin{exercise}
\label{exer:2.18}
设 $E$ 与 $F$ 为两个 Banach 空间, $A:D(A)\subset E\to F$ 为稠定的无界算子.

\begin{enumerate}[label=\arabic*., leftmargin=2.2em, itemsep=0.25em]
	\item 证明 $N(A^\star)=R(A)^\perp$ 且 $N(A)\subset R(A^\star)^\perp$.
	\item 假设 $A$ 还是闭的, 证明 $N(A)=R(A^\star)^\perp$.
\end{enumerate}
[试着找出直接的论证, 不要依赖推论  \ref{cor:2.18} 的证明.第 2 问用反证法:
假设存在 $u\in R(A^\star)^\perp$ 使 $[u,0]\notin G(A)$, 并应用 Hahn--Banach 定理.]
\end{exercise}

\begin{exercise}
\label{exer:2.19}
设 $E$ 为 Banach 空间, $A:D(A)\subset E\to E^\star$ 为稠定的无界算子.

\begin{enumerate}[label=\arabic*., leftmargin=2.2em, itemsep=0.25em]
	\item 假设存在常数 $C$ 使得
	      \begin{equation}
	      \dual{Au}{u}\ge-C\norm{Au}^2\quad \forall u\in D(A).
	      \tag{1}
	      \end{equation}
	      证明 $N(A)\subset N(A^\star)$.
	\item 反之, 假设 $N(A)\subset N(A^\star)$.又假设 $A$ 闭且 $R(A)$ 闭.
	      证明存在常数 $C$ 使 (1) 成立.
\end{enumerate}
\end{exercise}

\begin{exercise}
\label{exer:2.20}
设 $E$ 与 $F$ 为两个 Banach 空间.设 $T\in\LL(E,F)$, $A:D(A)\subset E\to F$
为稠定且闭的无界算子.考虑由
\begin{equation*}
D(B)=D(A),\qquad B=A+T
\end{equation*}
定义的算子 $B:D(B)\subset E\to F$.

\begin{enumerate}[label=\arabic*., leftmargin=2.2em, itemsep=0.25em]
	\item 证明 $B$ 闭.
	\item 证明 $D(B^\star)=D(A^\star)$ 且 $B^\star=A^\star+T^\star$.
\end{enumerate}
\end{exercise}

\begin{exercise}
\label{exer:2.21}
设 $E$ 为无限维 Banach 空间.固定元素 $a\in E$, $a\ne0$, 以及不连续的线性
泛函 $f:E\to\R$ (这样的泛函存在; 见习题  \ref{exer:1.5}).考虑由
\begin{equation*}
D(A)=E,\qquad Ax=x-f(x)a
\end{equation*}
定义的算子 $A:E\to E$.

\begin{enumerate}[label=\arabic*., leftmargin=2.2em, itemsep=0.25em]
	\item 确定 $N(A)$ 与 $R(A)$.
	\item $A$ 是否闭?
	\item 确定 $A^\star$ (仔细地定义 $D(A^\star)$).
	\item 确定 $N(A^\star)$ 与 $R(A^\star)$.
	\item 比较 $N(A)$ 与 $R(A^\star)^\perp$, 以及 $N(A^\star)$ 与 $R(A)^\perp$.
	\item 与习题  \ref{exer:2.18} 的结果作比较.
\end{enumerate}
\end{exercise}

\begin{exercise}
\label{exer:2.22}
本习题的目的是构造一个无界算子 $A:D(A)\subset E\to E$, 它是稠定, 闭的, 且
满足 $\overline{D(A^\star)}\ne E^\star$.

设 $E=\ell^1$, 从而 $E^\star=\ell^\infty$.考虑由
\begin{equation*}
D(A)=\bigl\{u=(u_n)\in\ell^1;\ (nu_n)\in\ell^1\bigr\}
\quad\text{与}\quad Au=(nu_n)
\end{equation*}
定义的算子 $A:D(A)\subset E\to E$.

\begin{enumerate}[label=\arabic*., leftmargin=2.2em, itemsep=0.25em]
	\item 验证 $A$ 是稠定且闭的.
	\item 确定 $D(A^\star)$, $A^\star$ 与 $\overline{D(A^\star)}$.
\end{enumerate}
\end{exercise}

\begin{exercise}
\label{exer:2.23}
设 $E=\ell^1$, 从而 $E^\star=\ell^\infty$.考虑由
\begin{equation*}
Tu=\Bigl(\frac1n u_n\Bigr)_{n\ge1}
\quad\text{对每个 } u=(u_n)_{n\ge1}\in\ell^1
\end{equation*}
定义的算子 $T\in\LL(E,E)$.

确定 $N(T)$, $N(T)^\perp$, $T^\star$, $R(T^\star)$ 与
$\overline{R(T^\star)}$.与推论  \ref{cor:2.18} 作比较.
\end{exercise}

\begin{exercise}
\label{exer:2.24}
设 $E$, $F$ 与 $G$ 为三个 Banach 空间.设 $A:D(A)\subset E\to F$ 为稠定的
无界算子.设 $T\in\LL(F,G)$, 考虑由 $D(B)=D(A)$ 与 $B=T\circ A$ 定义的算子
$B:D(B)\subset E\to G$.

\begin{enumerate}[label=\arabic*., leftmargin=2.2em, itemsep=0.25em]
	\item 确定 $B^\star$.
	\item (用一个例子) 证明: 即使 $A$ 闭, $B$ 也\emph{不必}闭.
\end{enumerate}
\end{exercise}

\begin{exercise}
\label{exer:2.25}
设 $E$, $F$ 与 $G$ 为三个 Banach 空间.

\begin{enumerate}[label=\arabic*., leftmargin=2.2em, itemsep=0.25em]
	\item 设 $T\in\LL(E,F)$, $S\in\LL(F,G)$.证明
	      \begin{equation*}
	      (S\circ T)^\star=T^\star\circ S^\star.
	      \end{equation*}
	\item 假设 $T\in\LL(E,F)$ 是双射.证明 $T^\star$ 是双射, 且
	      $(T^\star)^{-1}=(T^{-1})^\star$.
\end{enumerate}
\end{exercise}

\begin{exercise}
\label{exer:2.26}
设 $E$ 与 $F$ 为两个 Banach 空间, $T\in\LL(E,F)$.设
$\psi:F\to(-\infty,+\infty]$ 为凸函数.假设在 $R(T)$ 中存在某个元素, 使 $\psi$
在该处有限且连续.置
\begin{equation*}
\varphi(x)=\psi(Tx),\qquad x\in E.
\end{equation*}
证明: 对每个 $f\in F^\star$,
\begin{equation*}
\varphi^\star(T^\star f)
=\inf_{g\in N(T^\star)}\psi^\star(f-g)
=\min_{g\in N(T^\star)}\psi^\star(f-g).
\end{equation*}
\end{exercise}

\begin{exercise}
\label{exer:2.27}
设 $E$, $F$ 为两个 Banach 空间, $T\in\LL(E,F)$.假设 $R(T)$ 具有有限余维,
即存在 $F$ 的有限维子空间 $X$, 使得 $X+R(T)=F$ 且 $X\cap R(T)=\{0\}$.

证明 $R(T)$ 闭.
\end{exercise}

% chapters/ch03.tex
% 第~
\ref{chap:3}~章 \index{ruotuopu@弱拓扑 (weak topology)}
弱拓扑. \index{zifankongjian@自反空间 (reflexive space)}
自反空间. \index{kefenkongjian@可分空间 (separable space)}
可分空间. \index{yizhitukongjian@一致凸空间 (uniformly convex space)}
一致凸性
% 源: 原书 PDF 第 65--97 页 (印刷页 55--87)

\chapter{弱拓扑. 自反空间. 可分空间. 一致凸性}
\label{chap:3}

\section{使一族映射连续的, 最粗的拓扑}
\label{sec:3.1}

本章开头先回顾拓扑学中一个熟知的概念.设 $X$ 为一个集合 (不带任何结构),
$(Y_i)_{i\in I}$ 为一族\emph{拓扑空间}.给定一族映射 $(\varphi_i)_{i\in I}$, 其中对
每个 $i\in I$, $\varphi_i$ 都把 $X$ 映入 $Y_i$.我们考虑下述问题:

\runin{问题 1.}
在 $X$ 上构造一个拓扑, 使所有映射 $(\varphi_i)_{i\in I}$ 都连续.如果可能,
找出一个\emph{最经济}的拓扑 $\mathcal{T}$, 即开集最少的拓扑.

注意: 若给 $X$ 配上离散拓扑 (即 $X$ 的每个子集都是开集), 则每个映射
$\varphi_i$ 都连续; 当然, 这个拓扑远非``最便宜''的, 事实上它是最贵的! 我们将会
看到, 总存在一个 (唯一的) ``最便宜''的拓扑 $\mathcal{T}$, 使每个映射 $\varphi_i$ 都
连续.它称为与族 $(\varphi_i)_{i\in I}$ 相伴的\emph{最粗}或\emph{最弱}拓扑 (有时
也称为\emph{初始拓扑}).

若 $\omega_i\subset Y_i$ 是任一开集, 则 $\varphi_i^{-1}(\omega_i)$ 必须是 $\mathcal{T}$
中的开集.当 $\omega_i$ 取遍 $Y_i$ 的开集族, $i$ 取遍 $I$ 时, 我们得到一族 $X$ 的
子集, 其中每一个在拓扑 $\mathcal{T}$ 下都必须是开集.把这族子集记为
$(U_\lambda)_{\lambda\in\Lambda}$.当然, 这族子集不必构成一个拓扑.因此我们引出:

\runin{问题 2.}
给定集合 $X$ 以及 $X$ 中的一族子集 $(U_\lambda)_{\lambda\in\Lambda}$, 在 $X$ 上构造
最便宜的拓扑 $\mathcal{T}$, 使 $U_\lambda$ 对一切 $\lambda\in\Lambda$ 都是开集.

换言之, 我们要找出 $X$ 的子集构成的最便宜的一族 $\mathcal{F}$, 它对有限交与任意并
稳定\footnote{意思是: $\mathcal{F}$ 中有限多个集合的交以及 $\mathcal{F}$ 中任意多个
集合的并都仍属于 $\mathcal{F}$.}, 并且具有性质: 对每个 $\lambda\in\Lambda$ 都有
$U_\lambda\in\mathcal{F}$.构造如下.首先考虑 $(U_\lambda)_{\lambda\in\Lambda}$ 中
集合的有限交, 即 $\bigcap_{\lambda\in\Gamma}U_\lambda$, 其中 $\Gamma\subset\Lambda$
为有限集.这样得到 $X$ 的一族新子集, 记作 $\Phi$, 它包含
$(U_\lambda)_{\lambda\in\Lambda}$, 并且对有限交稳定.然而它不必对任意并稳定.因此
我们接着考虑由 $\Phi$ 中元素作任意并所得的一族 $\mathcal{F}$.显然 $\mathcal{F}$
对任意并稳定.$\mathcal{F}$ 是否对有限交稳定则并不显然; 但事实上我们有下述结果:

\begin{lemma}
\label{lem:3.1}
族 $\mathcal{F}$ 对有限交稳定.
\end{lemma}

引理  \ref{lem:3.1} 的证明——一道令人愉快的集合论练习——留给读者; 参见例如
G.~Folland \bref{Folland}{2}.现在显然, 上述构造给出了具有所需性质的最便宜的
拓扑.

\begin{remark}
\label{rem:3.1}
在 $\mathcal{F}$ 的构造中不能交换运算的次序.同样自然的是先取任意并, 再取有限交.
这样得到的一族对有限交稳定; 但它对任意并并不稳定.必须再考虑一次任意并, 过程
才稳定下来.
\end{remark}

总结上述讨论可知: 拓扑 $\mathcal{T}$ 的开集, 是先取形如
$\varphi_i^{-1}(\omega_i)$ 的集合的有限交, 再取任意并而得到的.由此可知, 对每个
$x\in X$, 取形如 $\bigcap_{\text{有限}}\varphi_i^{-1}(V_i)$ 的集合 (其中 $V_i$ 是
$\varphi_i(x)$ 在 $Y_i$ 中的邻域) 便得到 $x$ 在拓扑 $\mathcal{T}$ 下的一个邻域基.
回忆: 在拓扑空间中, 点 $x$ 的\emph{邻域基}是指 $x$ 的一族邻域, 使得 $x$ 的每个
邻域都包含该族中的某个邻域.

以下总给 $X$ 配上与族 $(\varphi_i)_{i\in I}$ 相伴的最弱拓扑 $\mathcal{T}$.拓扑
$\mathcal{T}$ 有两条简单的性质:

\begin{keyproposition}
\label{prop:3.1}
设 $(x_n)$ 为 $X$ 中的序列.则 $x_n\to x$ (在 $\mathcal{T}$ 中) 当且仅当对每个
$i\in I$ 都有 $\varphi_i(x_n)\to\varphi_i(x)$.
\end{keyproposition}

\begin{proof}
若 $x_n\to x$, 则对每个 $i$ 都有 $\varphi_i(x_n)\to\varphi_i(x)$, 因为每个
$\varphi_i$ 关于 $\mathcal{T}$ 连续.反之, 设 $U$ 为 $x$ 的邻域.由前面的讨论, 总可
设 $U$ 具有形式 $U=\bigcap_{i\in J}\varphi_i^{-1}(V_i)$, 其中 $J\subset I$ 为有限集.
对每个 $i\in J$ 存在整数 $N_i$, 使得当 $n\ge N_i$ 时 $\varphi_i(x_n)\in V_i$.
于是当 $n\ge N=\max_{i\in J}N_i$ 时 $x_n\in U$.
\end{proof}

\begin{keyproposition}
\label{prop:3.2}
设 $Z$ 为拓扑空间, $\psi$ 为从 $Z$ 到 $X$ 的映射.则 $\psi$ 连续当且仅当对每个
$i\in I$, $\varphi_i\circ\psi$ 都是从 $Z$ 到 $Y_i$ 的连续映射.
\end{keyproposition}

\begin{proof}
若 $\psi$ 连续, 则对每个 $i\in I$, $\varphi_i\circ\psi$ 也连续.反之, 须证明对每个
开集 $U\subset X$, $\psi^{-1}(U)$ 在 $Z$ 中开.但已知 $U$ 具有形式
$U=\bigcup_{\text{任意}}\bigcap_{\text{有限}}\varphi_i^{-1}(\omega_i)$, 其中
$\omega_i$ 在 $Y_i$ 中开.于是
\begin{equation*}
\psi^{-1}(U)
=\bigcup_{\text{任意}}\bigcap_{\text{有限}}\psi^{-1}\bigl[\varphi_i^{-1}(\omega_i)\bigr]
=\bigcup_{\text{任意}}\bigcap_{\text{有限}}(\varphi_i\circ\psi)^{-1}(\omega_i),
\end{equation*}
它在 $Z$ 中开, 因为每个映射 $\varphi_i\circ\psi$ 都连续.
\end{proof}

\section{弱拓扑 $\sigma(E,E^\star)$ 的定义与初等性质}
\label{sec:3.2}

设 $E$ 为 Banach 空间, $f\in E^\star$.以 $\varphi_f:E\to\R$ 记线性泛函
$\varphi_f(x)=\dual{f}{x}$.当 $f$ 取遍 $E^\star$ 时, 我们得到一族从 $E$ 到 $\R$
的映射 $(\varphi_f)_{f\in E^\star}$.现在暂且不管 $E$ 上通常的拓扑 (与 $\norm{\cdot}$
相联系的拓扑), 而在集合 $E$ 上定义一个新拓扑如下:

\begin{definition}
$E$ 上的\emph{弱拓扑} $\sigma(E,E^\star)$ 是与族 $(\varphi_f)_{f\in E^\star}$
相伴的最粗拓扑 (按第  \ref{sec:3.1} 节的意义, 其中 $X=E$, 每个 $Y_i=\R$, 且
$I=E^\star$).
\end{definition}

注意每个映射 $\varphi_f$ 关于通常拓扑连续, 因此弱拓扑比通常拓扑更弱.

\begin{proposition}
\label{prop:3.3}
弱拓扑 $\sigma(E,E^\star)$ 是 Hausdorff 拓扑.
\end{proposition}

\begin{proof}
给定 $x_1,x_2\in E$, $x_1\ne x_2$, 须找出弱拓扑 $\sigma(E,E^\star)$ 中的两个开集
$O_1$ 与 $O_2$, 使 $x_1\in O_1$, $x_2\in O_2$, 且 $O_1\cap O_2=\emptyset$.
由 Hahn--Banach 定理 (第二几何形式), 存在一个闭超平面把 $\{x_1\}$ 与 $\{x_2\}$
严格分离.于是存在 $f\in E^\star$ 与 $\alpha\in\R$, 使
\begin{equation*}
\dual{f}{x_1}<\alpha<\dual{f}{x_2}.
\end{equation*}
置
\begin{equation*}
O_1=\{x\in E;\ \dual{f}{x}<\alpha\}=\varphi_f^{-1}\bigl((-\infty,\alpha)\bigr),
\end{equation*}
\begin{equation*}
O_2=\{x\in E;\ \dual{f}{x}>\alpha\}=\varphi_f^{-1}\bigl((\alpha,+\infty)\bigr).
\end{equation*}
显然 $O_1$ 与 $O_2$ 关于 $\sigma(E,E^\star)$ 开, 且满足所需性质.
\end{proof}

\begin{keyproposition}
\label{prop:3.4}
设 $x_0\in E$; 给定 $\varepsilon>0$ 与 $E^\star$ 中的\emph{有限}集
$\{f_1,f_2,\dots,f_k\}$, 考虑
\begin{equation*}
V=V(f_1,f_2,\dots,f_k;\varepsilon)
=\{x\in E;\ \abs{\dual{f_i}{x-x_0}}<\varepsilon\ \ \forall i=1,2,\dots,k\}.
\end{equation*}
则 $V$ 是 $x_0$ 在拓扑 $\sigma(E,E^\star)$ 下的邻域.此外, 令 $\varepsilon$, $k$ 与
$E^\star$ 中的 $f_i$ 变化, 便得到 $x_0$ 关于 $\sigma(E,E^\star)$ 的一个邻域基.
\end{keyproposition}

\begin{proof}
显然 $V=\bigcap_{i=1}^k\varphi_{f_i}^{-1}\bigl((a_i-\varepsilon,a_i+\varepsilon)\bigr)$,
其中 $a_i=\dual{f_i}{x_0}$, 它关于拓扑 $\sigma(E,E^\star)$ 开且包含 $x_0$.反之,
设 $U$ 为 $x_0$ 关于 $\sigma(E,E^\star)$ 的邻域.由第  \ref{sec:3.1} 节的讨论知,
存在包含 $x_0$ 的开集 $W$, $W\subset U$, 具有形式
$W=\bigcap_{\text{有限}}\varphi_{f_i}^{-1}(\omega_i)$, 其中 $\omega_i$ 是
$a_i=\dual{f_i}{x_0}$ 在 $\R$ 中的邻域.于是存在 $\varepsilon>0$, 使对每个 $i$ 都有
$(a_i-\varepsilon,a_i+\varepsilon)\subset\omega_i$.由此 $x_0\in V\subset W\subset U$.
\end{proof}

\begin{notation}
若 $E$ 中的序列 $(x_n)$ 在弱拓扑 $\sigma(E,E^\star)$ 下收敛于 $x$, 我们写
\begin{equation*}
\boxed{x_n\weakto x.}
\end{equation*}
为避免混淆, 有时也说``$x_n\weakto x$ 在 $\sigma(E,E^\star)$ 中弱收敛''.为了完全
清楚, 有时也用``$x_n\to x$ 强收敛''来强调强收敛, 意即 $\norm{x_n-x}\to0$.
\end{notation}

\begin{keyproposition}
\label{prop:3.5}
设 $(x_n)$ 为 $E$ 中的序列.则
\begin{enumerate}[label=(\roman*), leftmargin=2.4em, itemsep=0.15em]
	\item $[x_n\weakto x\ \text{在}\ \sigma(E,E^\star)\ \text{中弱收敛}]
	      \iff[\dual{f}{x_n}\to\dual{f}{x}\ \forall f\in E^\star]$.
	\item 若 $x_n\to x$ 强收敛, 则 $x_n\weakto x$ 在 $\sigma(E,E^\star)$ 中弱收敛.
	\item 若 $x_n\weakto x$ 在 $\sigma(E,E^\star)$ 中弱收敛, 则 $(\norm{x_n})$ 有界,
	      且 $\norm{x}\le\liminf\norm{x_n}$.
	\item 若 $x_n\weakto x$ 在 $\sigma(E,E^\star)$ 中弱收敛, 且 $f_n\to f$ 在
	      $E^\star$ 中强收敛 (即 $\norm{f_n-f}_{E^\star}\to0$), 则
	      $\dual{f_n}{x_n}\to\dual{f}{x}$.
\end{enumerate}
\end{keyproposition}

\begin{proof}
(i) 由命题  \ref{prop:3.1} 与弱拓扑 $\sigma(E,E^\star)$ 的定义即得.

(ii) 由 (i) 即得, 因为
$\abs{\dual{f}{x_n}-\dual{f}{x}}\le\norm{f}\norm{x_n-x}$; 由弱拓扑比强拓扑更弱这一
事实也可看出.

(iii) 由一致有界原理 (见推论  \ref{cor:2.4}) 即得, 因为对每个 $f\in E^\star$, 集合
$(\dual{f}{x_n})_n$ 有界.在不等式
\begin{equation*}
\abs{\dual{f}{x_n}}\le\norm{f}\norm{x_n}
\end{equation*}
中取极限, 得
\begin{equation*}
\abs{\dual{f}{x}}\le\norm{f}\liminf\norm{x_n},
\end{equation*}
从而 (由推论  \ref{cor:1.4})
\begin{equation*}
\norm{x}=\sup_{\norm{f}\le1}\abs{\dual{f}{x}}\le\liminf\norm{x_n}.
\end{equation*}

(iv) 由不等式
\begin{equation*}
\abs{\dual{f_n}{x_n}-\dual{f}{x}}
\le\abs{\dual{f_n-f}{x_n}}+\abs{\dual{f}{x_n-x}}
\le\norm{f_n-f}\norm{x_n}+\abs{\dual{f}{x_n-x}},
\end{equation*}
结合 (i) 与 (iii) 即得.
\end{proof}

\begin{keyproposition}
\label{prop:3.6}
当 $E$ 为有限维时, 弱拓扑 $\sigma(E,E^\star)$ 与通常拓扑相同.特别地, 序列
$(x_n)$ 弱收敛当且仅当它强收敛.
\end{keyproposition}

\begin{proof}
由于弱拓扑的开集总比强拓扑少, 只须验证每个强开集都是弱开集.设 $x_0\in E$,
$U$ 为 $x_0$ 在强拓扑中的邻域.须找出 $x_0$ 在弱拓扑 $\sigma(E,E^\star)$ 中的
邻域 $V$, 使 $V\subset U$.换言之, 须找出 $E^\star$ 中的
$f_1,f_2,\dots,f_k$ 与 $\varepsilon>0$, 使
\begin{equation*}
V=\{x\in E;\ \abs{\dual{f_i}{x-x_0}}<\varepsilon\ \ \forall i=1,2,\dots,k\}\subset U.
\end{equation*}
取 $r>0$ 使 $B(x_0,r)\subset U$.在 $E$ 中取基 $e_1,e_2,\dots,e_k$, 使
$\norm{e_i}=1$, $\forall i$.每个 $x\in E$ 都有分解
$x=\sum_{i=1}^k x_ie_i$, 而映射 $x\mapsto x_i$ 是 $E$ 上的连续线性泛函, 记为
$f_i$.对每个 $x\in V$ 有
\begin{equation*}
\norm{x-x_0}\le\sum_{i=1}^k\abs{\dual{f_i}{x-x_0}}<k\varepsilon.
\end{equation*}
取 $\varepsilon=r/k$, 即得 $V\subset U$.
\end{proof}

\begin{remark}
\label{rem:3.2}
弱拓扑 $\sigma(E,E^\star)$ 中的开集 (相应地, 闭集) 总是强拓扑中的开集 (相应地,
闭集).在任何无穷维空间中, 弱拓扑都严格粗于强拓扑; 即存在强拓扑中的开集
(相应地, 闭集), 它们在弱拓扑中不是开集 (相应地, 闭集).下面给出两个例子:
\end{remark}

\begin{example}
\label{exa:3.1}
单位球面 $S=\{x\in E;\ \norm{x}=1\}$ (其中 $E$ 为无穷维) 在弱拓扑
$\sigma(E,E^\star)$ 中绝不闭.更确切地说, 有
\begin{equation}
\overline{S}^{\sigma(E,E^\star)}=B_E,
\label{eq:3.1}
\end{equation}
其中 $\overline{S}^{\sigma(E,E^\star)}$ 表示 $S$ 在拓扑 $\sigma(E,E^\star)$ 中的闭包,
而 $B_E$ (第~
\ref{chap:2}~章中已定义) 表示 $E$ 中的闭单位球,
\begin{equation*}
B_E=\{x\in E;\ \norm{x}\le1\}.
\end{equation*}

先验证每个满足 $\norm{x_0}<1$ 的 $x_0\in E$ 都属于
$\overline{S}^{\sigma(E,E^\star)}$.事实上, 设 $V$ 为 $x_0$ 在
$\sigma(E,E^\star)$ 中的邻域.须证明 $V\cap S\ne\emptyset$.由命题  \ref{prop:3.4},
总可设 $V$ 具有形式
\begin{equation*}
V=\{x\in E;\ \abs{\dual{f_i}{x-x_0}}<\varepsilon\ \ \forall i=1,2,\dots,k\},
\end{equation*}
其中 $\varepsilon>0$, $f_1,f_2,\dots,f_k\in E^\star$.取 $y_0\in E$, $y_0\ne0$, 使
\begin{equation*}
\dual{f_i}{y_0}=0\quad\forall i=1,2,\dots,k.
\end{equation*}
[这样的 $y_0$ 存在; 否则映射 $\varphi:E\to\R^k$,
$\varphi(x)=(\dual{f_i}{x})_{1\le i\le k}$ 将是单射, 而 $\varphi$ 将是 $E$ 到
$\varphi(E)$ 的同构, 从而 $\dim E\le k$, 与 $E$ 为无穷维的假设矛盾.]\footnote{这一
构造的几何解释如下: 当 $E$ 为无穷维时, $x_0$ 在拓扑 $\sigma(E,E^\star)$ 中的每个
邻域 $V$ 都包含一条过 $x_0$ 的直线, 甚至包含一个过 $x_0$ 的``巨大''仿射空间.}
函数 $g(t)=\norm{x_0+ty_0}$ 在 $[0,\infty)$ 上连续, 且 $g(0)<1$,
$\lim_{t\to+\infty}g(t)=+\infty$.因此存在 $t_0>0$, 使 $\norm{x_0+t_0y_0}=1$.
于是 $x_0+t_0y_0\in V\cap S$, 从而得到
\begin{equation*}
S\subset B_E\subset\overline{S}^{\sigma(E,E^\star)}.
\end{equation*}

为完成 \eqref{eq:3.1} 的证明, 只须知道 $B_E$ 在拓扑 $\sigma(E,E^\star)$ 中闭.
但
\begin{equation*}
B_E=\bigcap_{\substack{f\in E^\star\\ \norm{f}\le1}}\{x\in E;\ \abs{\dual{f}{x}}\le1\},
\end{equation*}
它是弱闭集的交.
\end{example}

\begin{example}
\label{exa:3.2}
单位球 $U=\{x\in E;\ \norm{x}<1\}$ (其中 $E$ 为无穷维) 在弱拓扑
$\sigma(E,E^\star)$ 中绝不\emph{开}.反设 $U$ 弱开.则其补集
$U^c=\{x\in E;\ \norm{x}\ge1\}$ 弱闭.于是 $S=B_E\cap U^c$ 也弱闭; 这与例
\ref{exa:3.1} 矛盾.
\end{example}

\begin{starremark}
\label{rem:3.3}
在无穷维空间中, 弱拓扑绝不可度量化, 即 $E$ 上不存在度量 (从而也不存在范数) 在
$E$ 上诱导出弱拓扑 $\sigma(E,E^\star)$; 见习题  \ref{exer:3.8}.然而, 我们后面
将会看到 (定理  \ref{thm:3.29}), 若 $E^\star$ 可分, 则可在 $E$ 上定义一个范数,
它在 $E$ 的\emph{有界}集上诱导出弱拓扑 $\sigma(E,E^\star)$.
\end{starremark}

\begin{starremark}
\label{rem:3.4}
通常, 在无穷维空间中存在弱收敛而不强收敛的序列.例如, 若 $E^\star$ 可分或 $E$
自反, 则可构造 $E$ 中的序列 $(x_n)$, 使 $\norm{x_n}=1$ 且 $x_n\weakto0$ 弱收敛
(见习题  \ref{exer:3.22}).然而, 也存在无穷维空间, 使每个弱收敛序列都强收敛.
例如 $\ell^1$ 就有这个不寻常的性质 (见问题 8).这样的空间相当``稀少'', 且有些
``病态''.这一奇怪的事实与注记 \ref{rem:3.2} 并不矛盾; 注记 \ref{rem:3.2} 断言:
在无穷维空间中, 弱拓扑与强拓扑总是不同的, 弱拓扑严格粗于强拓扑.请记住: 两个
\emph{度量} (或可度量化) 空间若有相同的收敛序列, 则它们的拓扑相同; 然而, 两个
\emph{拓扑}空间若有相同的收敛序列, 它们的拓扑却不必相同.
\end{starremark}

\section{弱拓扑, 凸集与线性算子}
\label{sec:3.3}

每个弱闭集都是强闭集, 而其逆在无穷维空间中不成立 (见注记 \ref{rem:3.2}).然而,
知道下述事实非常有用: 对凸集而言, 弱闭 $=$ 强闭:

\begin{keytheorem}
\label{thm:3.7}
设 $C$ 为 $E$ 的凸子集.则 $C$ 在弱拓扑 $\sigma(E,E^\star)$ 中闭当且仅当它在强
拓扑中闭.
\end{keytheorem}

\begin{proof}
设 $C$ 强闭, 要证 $C$ 弱闭.我们将验证 $C$ 的补集 $C^c$ 在弱拓扑中开.为此,
设 $x_0\notin C$.由 Hahn--Banach 定理, 存在一个闭超平面把 $\{x_0\}$ 与 $C$
严格分离.于是存在 $f\in E^\star$ 与 $\alpha\in\R$, 使
\begin{equation*}
\dual{f}{x_0}<\alpha<\dual{f}{y}\quad\forall y\in C.
\end{equation*}
置
\begin{equation*}
V=\{x\in E;\ \dual{f}{x}<\alpha\};
\end{equation*}
于是 $x_0\in V$, $V\cap C=\emptyset$ (即 $V\subset C^c$), 且 $V$ 在弱拓扑中开.
\end{proof}

\begin{corollary}[Mazur]
\label{cor:3.8}
假设 $(x_n)$ 弱收敛于 $x$.则存在由 $x_n$ 的凸组合构成的序列 $(y_n)$, 它强收敛
于 $x$.
\end{corollary}

\begin{proof}
置 $C=\operatorname{conv}\bigl(\bigcup_{p=1}^\infty\{x_p\}\bigr)$, 即 $x_n$ 的凸包.
由于 $x$ 属于 $\bigcup_{p=1}^\infty\{x_p\}$ 的弱闭包, 它自然属于 $C$ 的弱闭包.
由定理  \ref{thm:3.7}, $x\in\overline{C}$ (即 $C$ 的强闭包), 结论由此得出.
\end{proof}

\begin{remark}
\label{rem:3.5}
推论  \ref{cor:3.8} 有一些变形 (见习题  \ref{exer:3.4} 与  \ref{exer:5.24}).此外,
注意定理  \ref{thm:3.7} 的证明表明: 每个闭凸集 $C$ 都等于所有包含 $C$ 的闭半空间
之交.
\end{remark}

\begin{keycorollary}
\label{cor:3.9}
假设 $\varphi:E\to(-\infty,+\infty]$ 凸, 且在强拓扑中下半连续.则 $\varphi$ 在弱
拓扑 $\sigma(E,E^\star)$ 中下半连续.
\end{keycorollary}

\begin{proof}
对每个 $\lambda\in\R$, 集合
\begin{equation*}
A=\{x\in E;\ \varphi(x)\le\lambda\}
\end{equation*}
凸且强闭.由定理  \ref{thm:3.7} 它弱闭, 因此 $\varphi$ 弱下半连续.
\end{proof}

\begin{keyremark}
\label{rem:3.6}
在实践中要证明一个函数在弱拓扑中下半连续可能相当困难.推论  \ref{cor:3.9} 常按
下述方式使用:
\begin{equation*}
\boxed{\varphi\ \text{凸且强连续}\ \Rightarrow\ \varphi\ \text{弱下半连续}.}
\end{equation*}
例如, 函数 $\varphi(x)=\norm{x}$ 凸且强连续; 因此它弱下半连续.特别地, 若
$x_n\weakto x$ 弱收敛, 则推出 $\norm{x}\le\liminf\norm{x_n}$ (另见命题
 \ref{prop:3.5}).
\end{keyremark}

\begin{theorem}
\label{thm:3.10}
设 $E$ 与 $F$ 为两个 Banach 空间, $T$ 为从 $E$ 到 $F$ 的线性算子.假设 $T$ 在强
拓扑中连续.则 $T$ 是从 $E$ 弱 $\sigma(E,E^\star)$ 到 $F$ 弱 $\sigma(F,F^\star)$
的连续算子, 反之亦然.
\end{theorem}

\begin{proof}
由命题  \ref{prop:3.2}, 只须验证对每个 $f\in F^\star$, 映射
$x\mapsto\dual{f}{Tx}$ 是从 $E$ 弱 $\sigma(E,E^\star)$ 到 $\R$ 的连续映射.但
$x\mapsto\dual{f}{Tx}$ 是 $E$ 上的连续线性泛函, 因此它在弱拓扑
$\sigma(E,E^\star)$ 中也连续.

反之, 设 $T$ 是从 $E$ 弱到 $F$ 弱的连续算子.则 $G(T)$ 在赋以乘积拓扑
$\sigma(E,E^\star)\times\sigma(F,F^\star)$ 的 $E\times F$ 中闭, 而该乘积拓扑显然
就是 $\sigma(E\times F,(E\times F)^\star)$.于是 $G(T)$ 强闭 (任何弱闭集都强闭).
借助闭图像定理 (定理  \ref{thm:2.9}), 我们断定 $T$ 是从 $E$ 强到 $F$ 强的连续
算子.
\end{proof}

\begin{remark}
\label{rem:3.7}
上述论证还表明: 若线性算子 $T$ 是从 $E$ 强到 $F$ 弱的连续算子, 则 $T$ 是从 $E$
强到 $F$ 强的连续算子.因此, 对线性算子而言, 下述连续性都是同一回事:
$S\to S$, $W\to W$, $S\to W$ ($S=$ 强, $W=$ 弱).另一方面, 很少有线性算子是
$W\to S$ 连续的; 这种情况成立当且仅当 $T$ 是 $S\to S$ 连续的, 并且
$\dim R(T)<\infty$ (见习题  \ref{exer:6.7}).

此外, 注意一般说来, 从 $E$ 强到 $F$ 强连续的非线性映射, 不必是从 $E$ 弱到 $F$ 弱
连续的 (见例如习题  \ref{exer:4.20}).这是非线性问题中困难的一个主要来源.
\end{remark}

\section{\texorpdfstring{\wstar{}拓扑 $\sigma(E^\star,E)$}{弱*拓扑 $\sigma(E^\star,E)$}}
\label{sec:3.4}

到目前为止, $E^\star$ 上有两个拓扑:
\begin{enumerate}[label=(\alph*), leftmargin=2.0em, itemsep=0.15em]
	\item 与 $E^\star$ 的范数相联系的通常 (强) 拓扑;
	\item 在 $E^\star$ 上施行第  \ref{sec:3.3} 节的构造所得的弱拓扑
	      $\sigma(E^\star,E^{\star\star})$.
\end{enumerate}
现在要在 $E^\star$ 上定义第三个拓扑, 称为\emph{\wstar{}拓扑}, 记作
$\sigma(E^\star,E)$ (这里的 $^\star$ 提醒我们: 这个拓扑只在共轭空间上定义).
对每个 $x\in E$ 考虑线性泛函 $\varphi_x:E^\star\to\R$,
$f\mapsto\varphi_x(f)=\dual{f}{x}$.当 $x$ 取遍 $E$ 时, 我们得到一族从 $E^\star$
到 $\R$ 的映射 $(\varphi_x)_{x\in E}$.

\begin{definition}
\emph{\wstar{}拓扑} $\sigma(E^\star,E)$ 是 $E^\star$ 上与族 $(\varphi_x)_{x\in E}$
相伴的最粗拓扑 (按第  \ref{sec:3.1} 节的意义, 其中 $X=E^\star$, 所有 $Y_i=\R$,
且 $I=E$).
\end{definition}

由于 $E\subset E^{\star\star}$, 显然拓扑 $\sigma(E^\star,E)$ 比拓扑
$\sigma(E^\star,E^{\star\star})$ 更粗; 即拓扑 $\sigma(E^\star,E)$ 的开集
(相应地, 闭集) 比拓扑 $\sigma(E^\star,E^{\star\star})$ 少, 而后者又比强拓扑的开集
(相应地, 闭集) 少.

\begin{remark}
\label{rem:3.8}
读者大概会问, 为什么人们对弱拓扑如此狂热! 原因如下: 更粗的拓扑有更多的紧集.
例如 $E^\star$ 中的闭单位球 $B_{E^\star}$ 在强拓扑中绝不紧 (除非
$\dim E<\infty$; 见定理  \ref{thm:6.5}), 但它在\wstar{}拓扑中总是紧的 (见定理
 \ref{thm:3.16}).了解了紧集的基本作用——例如在极小化这样的存在性机制中的作用——
就很容易理解\wstar{}拓扑的重要性.
\end{remark}

\begin{proposition}
\label{prop:3.11}
\wstar{}拓扑是 Hausdorff 拓扑.
\end{proposition}

\begin{proof}
给定 $f_1,f_2\in E^\star$, $f_1\ne f_2$, 存在 $x\in E$ 使
$\dual{f_1}{x}\ne\dual{f_2}{x}$ (这里不用 Hahn--Banach 定理, 只用到 $f_1\ne f_2$
这一事实).例如设 $\dual{f_1}{x}<\dual{f_2}{x}$, 取 $\alpha$ 使
\begin{equation*}
\dual{f_1}{x}<\alpha<\dual{f_2}{x}.
\end{equation*}
置
\begin{equation*}
O_1=\{f\in E^\star;\ \dual{f}{x}<\alpha\}=\varphi_x^{-1}\bigl((-\infty,\alpha)\bigr),
\end{equation*}
\begin{equation*}
O_2=\{f\in E^\star;\ \dual{f}{x}>\alpha\}=\varphi_x^{-1}\bigl((\alpha,+\infty)\bigr).
\end{equation*}
则 $O_1$ 与 $O_2$ 是 $\sigma(E^\star,E)$ 中的开集, 且 $f_1\in O_1$, $f_2\in O_2$,
$O_1\cap O_2=\emptyset$.
\end{proof}

\begin{proposition}
\label{prop:3.12}
设 $f_0\in E^\star$; 给定 $E$ 中的\emph{有限}集 $\{x_1,x_2,\dots,x_k\}$ 与
$\varepsilon>0$, 考虑
\begin{equation*}
V=V(x_1,x_2,\dots,x_k;\varepsilon)
=\bigl\{f\in E^\star;\ \abs{\dual{f-f_0}{x_i}}<\varepsilon\ \ \forall i=1,2,\dots,k\bigr\}.
\end{equation*}
则 $V$ 是 $f_0$ 在拓扑 $\sigma(E^\star,E)$ 下的邻域.此外, 令 $\varepsilon$, $k$ 与
$E$ 中的 $x_i$ 变化, 便得到 $f_0$ 关于 $\sigma(E^\star,E)$ 的一个邻域基.
\end{proposition}

\begin{proof}
与命题  \ref{prop:3.4} 的证明相同.
\end{proof}

\begin{notation}
若 $E^\star$ 中的序列 $(f_n)$ 在\wstar{}拓扑下收敛于 $f$, 我们写
\begin{equation*}
\boxed{f_n\weakstar f.}
\end{equation*}
为避免混淆, 有时也说``$f_n\weakstar f$ 在 $\sigma(E^\star,E)$ 中'', ``$f_n\weakto f$ 在 $\sigma(E^\star,E^{\star\star})$ 中''以及``$f_n\to f$ 强收敛''.
\end{notation}

\begin{keyproposition}
\label{prop:3.13}
设 $(f_n)$ 为 $E^\star$ 中的序列.则
\begin{enumerate}[label=(\roman*), leftmargin=2.4em, itemsep=0.15em]
	\item $[f_n\weakstar f\ \text{在}\ \sigma(E^\star,E)\ \text{中}]
	      \iff[\dual{f_n}{x}\to\dual{f}{x},\ \forall x\in E]$.
	\item 若 $f_n\to f$ 强收敛, 则 $f_n\weakto f$ 在
	      $\sigma(E^\star,E^{\star\star})$ 中.

	      若 $f_n\weakto f$ 在 $\sigma(E^\star,E^{\star\star})$ 中, 则
	      $f_n\weakstar f$ 在 $\sigma(E^\star,E)$ 中.
	\item 若 $f_n\weakstar f$ 在 $\sigma(E^\star,E)$ 中, 则 $(\norm{f_n})$ 有界,
	      且 $\norm{f}\le\liminf\norm{f_n}$.
	\item 若 $f_n\weakstar f$ 在 $\sigma(E^\star,E)$ 中, 且 $x_n\to x$ 在 $E$ 中强
	      收敛, 则 $\dual{f_n}{x_n}\to\dual{f}{x}$.
\end{enumerate}
\end{keyproposition}

\begin{proof}
照抄命题  \ref{prop:3.5} 的证明.
\end{proof}

\begin{remark}
\label{rem:3.9}
假设 $f_n\weakstar f$ 在 $\sigma(E^\star,E)$ 中 (甚至在
$\sigma(E^\star,E^{\star\star})$ 中弱收敛), 且 $x_n\weakto x$ 在
$\sigma(E,E^\star)$ 中.一般不能断定 $\dual{f_n}{x_n}\to\dual{f}{x}$ (在 Hilbert
空间中很容易构造出例子).
\end{remark}

\begin{remark}
\label{rem:3.10}
当 $E$ 为有限维空间时, $E^\star$ 上的三个拓扑 (强, 弱, \wstar{}) 重合.事实上,
典范单射 $J:E\to E^{\star\star}$ (见第  \ref{sec:1.3} 节) 是满射 (因为
$\dim E=\dim E^{\star\star}$), 因此 $\sigma(E^\star,E)=\sigma(E^\star,E^{\star\star})$.
\end{remark}

\begin{starproposition}
\label{prop:3.14}
设 $\varphi:E^\star\to\R$ 为关于\wstar{}拓扑连续的线性泛函.则存在 $x_0\in E$,
使
\begin{equation*}
\varphi(f)=\dual{f}{x_0}\quad\forall f\in E^\star.
\end{equation*}
\end{starproposition}

其证明依赖下述有用的代数引理:

\begin{lemma}
\label{lem:3.2}
设 $X$ 为向量空间, $\varphi,\varphi_1,\varphi_2,\dots,\varphi_k$ 为 $X$ 上的
$(k+1)$ 个线性泛函, 使得
\begin{equation}
[\varphi_i(v)=0\ \ \forall i=1,2,\dots,k]\ \Rightarrow\ [\varphi(v)=0].
\label{eq:3.2}
\end{equation}
则存在常数 $\lambda_1,\lambda_2,\dots,\lambda_k\in\R$, 使
$\varphi=\sum_{i=1}^k\lambda_i\varphi_i$.
\end{lemma}


\begin{proof}[引理 \ref{lem:3.2}]
考虑由
\begin{equation*}
F(u)=[\varphi(u),\varphi_1(u),\varphi_2(u),\dots,\varphi_k(u)]
\end{equation*}
定义的映射 $F:X\to\R^{k+1}$.由假设 \eqref{eq:3.2} 可知,
$a=[1,0,0,\dots,0]$ 不属于 $R(F)$.于是可用 $\R^{k+1}$ 中的某个超平面把
$\{a\}$ 与 $R(F)$ 严格分离; 即存在常数
$\lambda,\lambda_1,\lambda_2,\dots,\lambda_k$ 与 $\alpha$, 使
\begin{equation*}
\lambda<\alpha<\lambda\varphi(u)+\sum_{i=1}^k\lambda_i\varphi_i(u)
\quad\forall u\in X.
\end{equation*}
由此推出
\begin{equation*}
\lambda\varphi(u)+\sum_{i=1}^k\lambda_i\varphi_i(u)=0\quad\forall u\in X,
\end{equation*}
并且 $\lambda<0$ (从而 $\lambda\ne0$).
\end{proof}

\begin{proof}[命题 \ref{prop:3.14}]
由于 $\varphi$ 关于\wstar{}拓扑连续, 存在 $0$ 关于 $\sigma(E^\star,E)$ 的邻域
$V$, 使
\begin{equation*}
\abs{\varphi(f)}<1\quad\forall f\in V.
\end{equation*}
总可设
\begin{equation*}
V=\{f\in E^\star;\ \abs{\dual{f}{x_i}}<\varepsilon\ \ \forall i=1,2,\dots,k\},
\end{equation*}
其中 $x_i\in E$, $\varepsilon>0$.特别地,
\begin{equation*}
[\dual{f}{x_i}=0\ \ \forall i=1,2,\dots,k]\ \Rightarrow\ [\varphi(f)=0].
\end{equation*}
由引理  \ref{lem:3.2} 得
\begin{equation*}
\varphi(f)=\sum_{i=1}^k\lambda_i\dual{f}{x_i}
=\Bigl\langle f,\sum_{i=1}^k\lambda_ix_i\Bigr\rangle\quad\forall f\in E^\star.
\end{equation*}
\end{proof}

\begin{starcorollary}
\label{cor:3.15}
假设 $H$ 是 $E^\star$ 中的超平面, 且在 $\sigma(E^\star,E)$ 中闭.则 $H$ 具有形式
\begin{equation*}
H=\{f\in E^\star;\ \dual{f}{x_0}=\alpha\},
\end{equation*}
其中 $x_0\in E$, $x_0\ne0$, $\alpha\in\R$.
\end{starcorollary}

\begin{proof}
$H$ 可写成
\begin{equation*}
H=\{f\in E^\star;\ \varphi(f)=\alpha\},
\end{equation*}
其中 $\varphi$ 是 $E^\star$ 上的线性泛函, $\varphi\not\equiv0$.取
$f_0\notin H$, 并取 $f_0$ 关于拓扑 $\sigma(E^\star,E)$ 的邻域 $V$, 使
$V\subset H^c$.可设
\begin{equation*}
V=\{f\in E^\star;\ \abs{\dual{f-f_0}{x_i}}<\varepsilon\ \ \forall i=1,2,\dots,k\}.
\end{equation*}
由于 $V$ 凸, 我们有
\begin{equation}
\varphi(f)<\alpha\quad\forall f\in V
\label{eq:3.3}
\end{equation}
或
\begin{equation*}
\varphi(f)>\alpha\quad\forall f\in V.
\tag{3$'$}
\label{eq:3.3p}
\end{equation*}
例如设 \eqref{eq:3.3} 成立, 则得
\begin{equation*}
\varphi(g)<\alpha-\varphi(f_0)\quad\forall g\in W=V-f_0,
\end{equation*}
而由于 $-W=W$, 我们推出
\begin{equation}
\abs{\varphi(g)}\le\abs{\alpha-\varphi(f_0)}\quad\forall g\in W.
\label{eq:3.4}
\end{equation}
由 \eqref{eq:3.4} 可知 $\varphi$ 在 $0$ 处关于拓扑 $\sigma(E^\star,E)$ 连续 (因为
$W$ 是 $0$ 的邻域).应用命题  \ref{prop:3.14}, 我们断定存在 $x_0\in E$, 使
\begin{equation*}
\varphi(f)=\dual{f}{x_0}\quad\forall f\in E^\star.
\end{equation*}
\end{proof}

\begin{remark}
\label{rem:3.11}
假设典范单射 $J:E\to E^{\star\star}$ 不是满射.则拓扑 $\sigma(E^\star,E)$ 严格
粗于拓扑 $\sigma(E^\star,E^{\star\star})$.例如, 取 $\xi\in E^{\star\star}$,
$\xi\notin J(E)$.则集合
\begin{equation*}
H=\{f\in E^\star;\ \dual{\xi}{f}=0\}
\end{equation*}
在 $\sigma(E^\star,E^{\star\star})$ 中闭, 但由推论  \ref{cor:3.15} 可知它在
$\sigma(E^\star,E)$ 中不闭.从这个例子还可知: 在强拓扑中闭的凸集, 在\wstar{}
拓扑中不必闭.$E^\star$ 中的闭凸集有两类:
\begin{enumerate}[label=(\alph*), leftmargin=2.0em, itemsep=0.15em]
	\item 强闭的凸集 ($=$ 在拓扑 $\sigma(E^\star,E^{\star\star})$ 中闭, 由定理
	       \ref{thm:3.7});
	\item 在 $\sigma(E^\star,E)$ 中闭的凸集.
\end{enumerate}
\end{remark}

\index{BanachAlaogluBourbakidingli@Banach--Alaoglu--Bourbaki 定理 (Banach--Alaoglu--Bourbaki theorem)}
\begin{keytheorem}[Banach--Alaoglu--Bourbaki]
\label{thm:3.16}
闭单位球
\begin{equation*}
B_{E^\star}=\{f\in E^\star;\ \norm{f}\le1\}
\end{equation*}
在\wstar{}拓扑 $\sigma(E^\star,E)$ 中紧.
\end{keytheorem}

\begin{remark}
\label{rem:3.12}
$B_{E^\star}$ 的紧性是\wstar{}拓扑最本质的性质; 另见注记 \ref{rem:3.8}.
\end{remark}

\begin{proof}
考虑 Cartesian 积 $Y=\R^E$, 它由所有从 $E$ 到 $\R$ 的映射组成; 把 $Y$ 的元素记作
$\omega=(\omega_x)_{x\in E}$, 其中 $\omega_x\in\R$.空间 $Y$ 赋以标准\emph{乘积
拓扑} (见例如 H.~L.~Royden \bref{Royden}{1}, J.~R.~Munkres \bref{Munkres}{1}, A.~Knapp \bref{Knapp}{1} 或 J.~Dixmier \bref{Dixmier}{1}), 即 $Y$ 上与映射族
$\omega\mapsto\omega_x$ (当 $x$ 取遍 $E$ 时) 相伴的最粗拓扑; 这当然与逐点收敛
拓扑相同 (见例如 J.~R.~Munkres \bref{Munkres}{1}).以下总给 $E^\star$ 配上
\wstar{}拓扑 $\sigma(E^\star,E)$.

由于 $E^\star$ 由从 $E$ 到 $\R$ 的特殊映射 (即连续线性映射) 组成, 我们可以把
$E^\star$ 看作 $Y$ 的子集.更确切地说, 设 $\Phi:E^\star\to Y$ 为从 $E^\star$ 到
$Y$ 的典范单射, 使得 $\Phi(f)=(\omega_x)_{x\in E}$, 其中 $\omega_x=\dual{f}{x}$.
显然 $\Phi$ 是从 $E^\star$ 到 $Y$ 的连续映射 (用命题  \ref{prop:3.2}, 并注意对每个
固定的 $x\in E$, 映射 $f\in E^\star\mapsto(\Phi(f))_x=\dual{f}{x}$ 连续).逆映射
$\Phi^{-1}$ 也是从 $\Phi(E^\star)$ (赋以 $Y$ 的拓扑) 到 $E^\star$ 的连续映射:
事实上, 再次用命题  \ref{prop:3.2}, 只须验证对每个固定的 $x\in E$, 映射
$\omega\mapsto\dual{\Phi^{-1}(\omega)}{x}$ 在 $\Phi(E^\star)$ 上连续; 这是显然的,
因为 $\dual{\Phi^{-1}(\omega)}{x}=\omega_x$ (注意 $\omega=\Phi(f)$ 对某个
$f\in E^\star$ 成立, 且 $\dual{\Phi^{-1}(\omega)}{x}=\dual{f}{x}=\omega_x$).
换言之, $\Phi$ 是从 $E^\star$ 到 $\Phi(E^\star)$ 的同胚.另一方面, 显然
$\Phi(B_{E^\star})=K$, 其中 $K$ 由
\begin{equation*}
K=\Bigl\{\omega\in Y\ \Big|\ \abs{\omega_x}\le\norm{x},\
\omega_{x+y}=\omega_x+\omega_y\
\text{且}\ \omega_{\lambda x}=\lambda\omega_x\ \forall\lambda\in\R,\ \forall x,y\in E\Bigr\}
\end{equation*}
定义.为完成定理  \ref{thm:3.16} 的证明, 只须验证 $K$ 是 $Y$ 的紧子集.把 $K$
写成 $K=K_1\cap K_2$, 其中
\begin{equation*}
K_1=\{\omega\in Y;\ \abs{\omega_x}\le\norm{x}\ \ \forall x\in E\}
\end{equation*}
且
\begin{equation*}
K_2=\bigl\{\omega\in Y;\ \omega_{x+y}=\omega_x+\omega_y\
\text{且}\ \omega_{\lambda x}=\lambda\omega_x\ \forall\lambda\in\R,\ \forall x,y\in E\bigr\}.
\end{equation*}
集合 $K_1$ 也可写成紧区间的乘积
\begin{equation*}
K_1=\prod_{x\in E}[-\norm{x},+\norm{x}].
\end{equation*}
回忆: (任意) 紧空间的乘积紧——这是 Tychonoff 的一个深刻定理; 见例如
H.~L.~Royden \bref{Royden}{1}, G.~B.~Folland \bref{Folland}{2}, J.~R.~Munkres \bref{Munkres}{1}, A.~Knapp \bref{Knapp}{1} 或 J.~Dixmier
\bref{Dixmier}{1}.因此 $K_1$ 紧.另一方面, $K_2$ 在 $Y$ 中闭; 事实上, 对每个固定
的 $\lambda\in\R$, $x,y\in E$, 集合
\begin{equation*}
A_{x,y}=\{\omega\in Y;\ \omega_{x+y}-\omega_x-\omega_y=0\},
\end{equation*}
\begin{equation*}
B_{\lambda,x}=\{\omega\in Y;\ \omega_{\lambda x}-\lambda\omega_x=0\}
\end{equation*}
在 $Y$ 中闭 (因为映射 $\omega\mapsto\omega_{x+y}-\omega_x-\omega_y$ 与
$\omega\mapsto\omega_{\lambda x}-\lambda\omega_x$ 在 $Y$ 上连续), 而 $K_2$ 可写成
\begin{equation*}
K_2=\Bigl(\bigcap_{x,y\in E}A_{x,y}\Bigr)\cap
\Bigl(\bigcap_{\substack{x\in E\ \lambda\in\R}}B_{\lambda,x}\Bigr).
\end{equation*}
最后, $K$ 紧, 因为它是紧集 ($K_1$) 与闭集 ($K_2$) 的交.
\end{proof}
\section{自反空间}
\label{sec:3.5}

\begin{definition}
设 $E$ 为 Banach 空间, $J:E\to E^{\star\star}$ 为从 $E$ 到 $E^{\star\star}$ 的
典范单射 (见第  \ref{sec:1.3} 节).若 $J$ 是满射, 即 $J(E)=E^{\star\star}$, 则称
空间 $E$ 是\emph{自反的}.
\end{definition}

当 $E$ 自反时, 通常把 $E^{\star\star}$ 与 $E$ 等同起来.

\begin{remark}
\label{rem:3.13}
分析中许多重要的空间都是自反的.显然, 有限维空间自反 (因为
$\dim E=\dim E^\star=\dim E^{\star\star}$).我们将在第~
\ref{chap:4}~章 (另见第~
\ref{chap:11}~章) 看到,
当 $1<p<\infty$ 时 $L^p$ (以及 $\ell^p$) 空间自反.在第~
\ref{chap:5}~章我们将看到 Hilbert
空间自反.然而, 分析中同样重要的一些空间不是自反的; 例如:
\begin{itemize}[leftmargin=1.6em, itemsep=0.15em]
	\item $L^1$ 与 $L^\infty$ (以及 $\ell^1$, $\ell^\infty$) 不自反 (见第~
\ref{chap:4}~章与
	      第~
\ref{chap:11}~章);
	\item $C(K)$, 即无穷紧度量空间 $K$ 上的连续函数空间, 不自反 (见习题
	       \ref{exer:3.25}).
\end{itemize}
\end{remark}

\begin{starremark}
\label{rem:3.14}
在上述定义中使用 $J$ 是本质的.R.~C.~James \bref{James}{1} 构造了一个引人注目的
非自反空间的例子, 它具有下述性质: 存在从 $E$ 到 $E^{\star\star}$ 的满射等距.
\end{starremark}

我们的下一个结果描述自反空间的一个基本性质:

\begin{keytheorem}[Kakutani]
\label{thm:3.17}
设 $E$ 为 Banach 空间.则 $E$ 自反当且仅当
\begin{equation*}
B_E=\{x\in E;\ \norm{x}\le1\}
\end{equation*}
在弱拓扑 $\sigma(E,E^\star)$ 中紧.
\end{keytheorem}

\begin{proof}
先设 $E$ 自反, 于是 $J(B_E)=B_{E^{\star\star}}$.由定理  \ref{thm:3.16} 已知
$B_{E^{\star\star}}$ 在拓扑 $\sigma(E^{\star\star},E^\star)$ 中紧.因此只须验证
$J^{-1}$ 是从赋以 $\sigma(E^{\star\star},E^\star)$ 的 $E^{\star\star}$ 到赋以
$\sigma(E,E^\star)$ 的 $E$ 的连续映射.由命题  \ref{prop:3.2}, 我们只须证明对每个
固定的 $f\in E^\star$, 映射 $\xi\mapsto\dual{f}{J^{-1}\xi}$ 在赋以
$\sigma(E^{\star\star},E^\star)$ 的 $E^{\star\star}$ 上连续.但
$\dual{f}{J^{-1}\xi}=\dual{\xi}{f}$, 而映射 $\xi\mapsto\dual{\xi}{f}$ 确实在
$E^{\star\star}$ 上关于拓扑 $\sigma(E^{\star\star},E^\star)$ 连续.因此我们证明了
$B_E$ 在 $\sigma(E,E^\star)$ 中紧.

逆命题更为精细, 它依赖下述两个引理:
\end{proof}

\begin{lemma}[Helly]
\label{lem:3.3}
设 $E$ 为 Banach 空间.给定 $E^\star$ 中的 $f_1,f_2,\dots,f_k$ 以及 $\R$ 中的
$\gamma_1,\gamma_2,\dots,\gamma_k$.则下述性质等价:
\begin{enumerate}[label=(\roman*), leftmargin=2.4em, itemsep=0.15em]
	\item $\forall\varepsilon>0$, $\exists x_\varepsilon\in E$, 使
	      $\norm{x_\varepsilon}\le1$ 且
	      \begin{equation*}
	      \abs{\dual{f_i}{x_\varepsilon}-\gamma_i}<\varepsilon
	      \quad\forall i=1,2,\dots,k,
	      \end{equation*}
	\item $\abs{\sum_{i=1}^k\beta_i\gamma_i}
	      \le\norm{\sum_{i=1}^k\beta_i f_i}
	      \quad\forall\beta_1,\beta_2,\dots,\beta_k\in\R$.
\end{enumerate}
\end{lemma}

\begin{proof}
(i) $\Rightarrow$ (ii). 固定 $\beta_1,\beta_2,\dots,\beta_k\in\R$, 并置
$S=\sum_{i=1}^k\abs{\beta_i}$.由 (i) 得
\begin{equation*}
\abs{\sum_{i=1}^k\beta_i\dual{f_i}{x_\varepsilon}
-\sum_{i=1}^k\beta_i\gamma_i}\le\varepsilon S,
\end{equation*}
从而
\begin{equation*}
\abs{\sum_{i=1}^k\beta_i\gamma_i}
\le\norm{\sum_{i=1}^k\beta_i f_i}\norm{x_\varepsilon}+\varepsilon S
\le\norm{\sum_{i=1}^k\beta_i f_i}+\varepsilon S.
\end{equation*}
由于这对每个 $\varepsilon>0$ 都成立, 便得到 (ii).

(ii) $\Rightarrow$ (i). 置 $\gamma=(\gamma_1,\gamma_2,\dots,\gamma_k)\in\R^k$,
考虑由
\begin{equation*}
\varphi(x)=(\dual{f_1}{x},\dots,\dual{f_k}{x})
\end{equation*}
定义的映射 $\varphi:E\to\R^k$.性质 (i) 说的正是
$\gamma\in\overline{\varphi(B_E)}$.反设 (i) 不成立, 即
$\gamma\notin\overline{\varphi(B_E)}$.于是 $\{\gamma\}$ 与
$\overline{\varphi(B_E)}$ 可用 $\R^k$ 中某个超平面严格分离; 即存在
$\beta=(\beta_1,\beta_2,\dots,\beta_k)\in\R^k$ 与 $\alpha\in\R$, 使
\begin{equation*}
\beta\cdot\varphi(x)<\alpha<\beta\cdot\gamma\quad\forall x\in B_E.
\end{equation*}
由此推出
\begin{equation*}
\Bigl\langle\sum_{i=1}^k\beta_i f_i,x\Bigr\rangle<\alpha<\sum_{i=1}^k\beta_i\gamma_i
\quad\forall x\in B_E,
\end{equation*}
于是
\begin{equation*}
\norm{\sum_{i=1}^k\beta_i f_i}\le\alpha<\sum_{i=1}^k\beta_i\gamma_i,
\end{equation*}
这与 (ii) 矛盾.
\end{proof}

\begin{lemma}[Goldstine]
\label{lem:3.4}
设 $E$ 为任一 Banach 空间.则 $J(B_E)$ 在 $B_{E^{\star\star}}$ 中关于拓扑
$\sigma(E^{\star\star},E^\star)$ 稠密; 从而 $J(E)$ 在 $E^{\star\star}$ 中关于拓扑
$\sigma(E^{\star\star},E^\star)$ 稠密.
\end{lemma}

\begin{proof}
设 $\xi\in B_{E^{\star\star}}$, $V$ 为 $\xi$ 关于拓扑
$\sigma(E^{\star\star},E^\star)$ 的邻域.须证明
$V\cap J(B_E)\ne\emptyset$.如通常那样, 可设 $V$ 具有形式
\begin{equation*}
V=\bigl\{\eta\in E^{\star\star};\ \abs{\dual{\eta-\xi}{f_i}}<\varepsilon
\ \ \forall i=1,2,\dots,k\bigr\},
\end{equation*}
其中 $f_1,f_2,\dots,f_k$ 为 $E^\star$ 中给定的元素, $\varepsilon>0$.须找出
$x\in B_E$, 使 $J(x)\in V$, 即
\begin{equation*}
\abs{\dual{f_i}{x}-\dual{\xi}{f_i}}<\varepsilon\quad\forall i=1,2,\dots,k.
\end{equation*}
置 $\gamma_i=\dual{\xi}{f_i}$.由引理  \ref{lem:3.3}, 只须验证
\begin{equation*}
\abs{\sum_{i=1}^k\beta_i\gamma_i}\le\norm{\sum_{i=1}^k\beta_i f_i},
\end{equation*}
这是显然的, 因为
$\sum_{i=1}^k\beta_i\gamma_i=\bigl\langle\xi,\sum_{i=1}^k\beta_i f_i\bigr\rangle$
且 $\norm{\xi}\le1$.
\end{proof}

\begin{remark}
\label{rem:3.15}
注意 $J(B_E)$ 在 $B_{E^{\star\star}}$ 中关于强拓扑闭.事实上, 若
$\xi_n=J(x_n)\to\xi$, 则看出 $(x_n)$ 是 $B_E$ 中的 Cauchy 列 (因为 $J$ 是等距),
从而 $x_n\to x$, 于是 $\xi=Jx$.由此可知, 除非 $J(B_E)=B_{E^{\star\star}}$
(即 $E$ 自反), $J(B_E)$ 在 $B_{E^{\star\star}}$ 中关于强拓扑不稠密.
\end{remark}

\begin{remark}
\label{rem:3.16}
引理  \ref{lem:3.4} 的另一证明见问题 9 (它基于 $E^{\star\star}$ 中 Hahn--Banach
定理的一个变形).
\end{remark}

\begin{proof}[定理 \ref{thm:3.17} 的续证]
典范单射 $J:E\to E^{\star\star}$ 总是从 $\sigma(E,E^\star)$ 到
$\sigma(E^{\star\star},E^\star)$ 连续的, 因为对每个固定的 $f\in E^\star$, 映射
$x\mapsto\dual{Jx}{f}=\dual{f}{x}$ 关于 $\sigma(E,E^\star)$ 连续.假设 $B_E$ 在
拓扑 $\sigma(E,E^\star)$ 中紧, 我们推出 $J(B_E)$ 在 $E^{\star\star}$ 中关于拓扑
$\sigma(E^{\star\star},E^\star)$ 紧——从而闭.另一方面, 由引理  \ref{lem:3.4},
$J(B_E)$ 在同一拓扑下在 $B_{E^{\star\star}}$ 中稠密.于是
$J(B_E)=B_{E^{\star\star}}$, 从而 $J(E)=E^{\star\star}$.
\end{proof}

关于自反空间的紧性, 我们还有下述两个结果:

\begin{keytheorem}
\label{thm:3.18}
假设 $E$ 为自反 Banach 空间, $(x_n)$ 为 $E$ 中的有界序列.则存在子列
$(x_{n_k})$, 它在弱拓扑 $\sigma(E,E^\star)$ 中收敛.
\end{keytheorem}

逆命题也成立, 即:

\begin{startheorem}[Eberlein--\v{S}mulian]
\label{thm:3.19}
假设 $E$ 为 Banach 空间, 使 $E$ 中每个有界序列都有弱收敛 (在
$\sigma(E,E^\star)$ 中) 的子列.则 $E$ 自反.
\end{startheorem}

定理  \ref{thm:3.18} 的证明需要绕道可分空间, 将在第  \ref{sec:3.6} 节给出.定理
 \ref{thm:3.19} 的证明相当精细, 从略; 见例如 R.~Holmes \bref{Holmes}{1}, K.~Yosida \bref{Yosida}{1}, N.~Dunford--J.~T.~Schwartz \bref{DunfordSchwartz}{1}, J.~Diestel \bref{Diestel}{2} 或问题 10.

\begin{remark}
\label{rem:3.17}
为弄清定理  \ref{thm:3.17},  \ref{thm:3.18} 与  \ref{thm:3.19} 之间的联系, 回忆
下述事实是有用的:
\begin{enumerate}[label=(\roman*), leftmargin=2.4em, itemsep=0.2em]
	\item 若 $X$ 为度量空间, 则
	      \begin{equation*}
	      [X\ \text{紧}]\ \iff\ [X\ \text{中每个序列都有收敛子列}].
	      \end{equation*}
	\item 存在紧拓扑空间 $X$ 以及 $X$ 中没有任何收敛子列的序列.典型的例子是
	      $X=B_{E^\star}$, 它在拓扑 $\sigma(E^\star,E)$ 中紧; 当 $E=\ell^\infty$ 时
	      很容易构造 $X$ 中没有收敛子列的序列 (见习题  \ref{exer:3.18}).
	\item 若拓扑空间 $X$ 具有``每个序列都有收敛子列''的性质, 则 $X$ 不必紧.
\end{enumerate}
\end{remark}

下面是自反空间的另一些性质.

\begin{keyproposition}
\label{prop:3.20}
假设 $E$ 为自反 Banach 空间, $M\subset E$ 为 $E$ 的闭线性子空间.则 $M$ 自反.
\end{keyproposition}

\begin{proof}
空间 $M$ (赋以 $E$ 的范数) 先验地有两个不同的弱拓扑:
\begin{enumerate}[label=(\alph*), leftmargin=2.0em, itemsep=0.15em]
	\item 由 $\sigma(E,E^\star)$ 诱导的拓扑;
	\item 它自己的弱拓扑 $\sigma(M,M^\star)$.
\end{enumerate}
事实上, 这两个拓扑相同 (因为由 Hahn--Banach 定理, $M$ 上每个连续线性泛函都是
$E$ 上某个连续线性泛函在 $M$ 上的限制).由定理  \ref{thm:3.17}, 须验证 $B_M$ 在
拓扑 $\sigma(M,M^\star)$ 中紧, 或在拓扑 $\sigma(E,E^\star)$ 中紧.然而, $B_E$ 在
拓扑 $\sigma(E,E^\star)$ 中紧, 而 $M$ 在拓扑 $\sigma(E,E^\star)$ 中闭 (由定理
 \ref{thm:3.7}).因此 $B_M$ 在拓扑 $\sigma(E,E^\star)$ 中紧.
\end{proof}

\begin{corollary}
\label{cor:3.21}
Banach 空间 $E$ 自反当且仅当其共轭空间 $E^\star$ 自反.
\end{corollary}

\begin{proof}
$E$ 自反 $\Rightarrow$ $E^\star$ 自反.证明的想法很简单, 因为粗略地说, 我们有
$E^{\star\star}=E\Rightarrow E^{\star\star\star}=E^\star$.更确切地说, 设 $J$ 为
从 $E$ 到 $E^{\star\star}$ 的典范同构.给定 $\varphi\in E^{\star\star\star}$.
映射 $x\mapsto\dual{\varphi}{Jx}$ 是 $E$ 上的连续线性泛函.记它为 $f\in E^\star$,
于是
\begin{equation*}
\dual{\varphi}{Jx}=\dual{f}{x}\quad\forall x\in E.
\end{equation*}
但我们也有
\begin{equation*}
\dual{\varphi}{Jx}=\dual{Jx}{f}\quad\forall x\in E.
\end{equation*}
由于 $J$ 满射, 我们推出
\begin{equation*}
\dual{\varphi}{\xi}=\dual{\xi}{f}\quad\forall\xi\in E^{\star\star},
\end{equation*}
这正好说明从 $E^\star$ 到 $E^{\star\star\star}$ 的典范单射是满射.

$E^\star$ 自反 $\Rightarrow$ $E$ 自反.由上一段我们已知 $E^{\star\star}$ 自反.
由于 $J(E)$ 是 $E^{\star\star}$ 中强闭的子空间, 由命题  \ref{prop:3.20} 推得
$J(E)$ 自反.因此 $E$ 自反.\footnote{显然, 若 $E$ 与 $F$ 为 Banach 空间, 且 $T$
是从 $E$ 到 $F$ 的线性满射等距, 则 $E$ 自反当且仅当 $F$ 自反.当然, 这与注记
\ref{rem:3.14} 并不矛盾!}
\end{proof}

\begin{keycorollary}
\label{cor:3.22}
设 $E$ 为自反 Banach 空间, $K\subset E$ 为 $E$ 的有界闭凸子集.则 $K$ 在拓扑
$\sigma(E,E^\star)$ 中紧.
\end{keycorollary}

\begin{proof}
$K$ 关于拓扑 $\sigma(E,E^\star)$ 闭 (由定理  \ref{thm:3.7}).另一方面, 存在常数
$m$ 使 $K\subset mB_E$, 而 $mB_E$ 在 $\sigma(E,E^\star)$ 中紧 (由定理
 \ref{thm:3.17}).
\end{proof}

\begin{keycorollary}
\label{cor:3.23}
设 $E$ 为自反 Banach 空间, $A\subset E$ 为 $E$ 的非空闭凸子集.设
$\varphi:A\to(-\infty,+\infty]$ 为凸下半连续函数, 且 $\varphi\not\equiv+\infty$, 并
设
\begin{equation}
\lim_{\substack{x\in A\\ \norm{x}\to\infty}}\varphi(x)=+\infty
\label{eq:3.5}
\end{equation}
(若 $A$ 有界则无此假设).则 $\varphi$ 在 $A$ 上达到最小值, 即存在 $x_0\in A$, 使
\begin{equation*}
\varphi(x_0)=\min_A\varphi.
\end{equation*}
\end{keycorollary}

\begin{proof}
任取 $a\in A$ 使 $\varphi(a)<+\infty$, 并考虑集合
\begin{equation*}
\tilde{A}=\{x\in A;\ \varphi(x)\le\varphi(a)\}.
\end{equation*}
则 $\tilde{A}$ 闭, 凸且有界 (由 \eqref{eq:3.5}), 从而它在拓扑 $\sigma(E,E^\star)$
中紧 (由推论  \ref{cor:3.22}).另一方面, $\varphi$ 在拓扑 $\sigma(E,E^\star)$ 中也
下半连续 (由推论  \ref{cor:3.9}).于是 $\varphi$ 在 $\tilde{A}$ 上达到最小值 (见
第~
\ref{chap:1}~章下半连续函数定义之后的性质 5), 即存在 $x_0\in\tilde{A}$, 使
\begin{equation*}
\varphi(x_0)\le\varphi(x)\quad\forall x\in\tilde{A}.
\end{equation*}
若 $x\in A\setminus\tilde{A}$, 则 $\varphi(x_0)\le\varphi(a)<\varphi(x)$; 因此
\begin{equation*}
\varphi(x_0)\le\varphi(x)\quad\forall x\in A.
\end{equation*}
\end{proof}

\begin{remark}
\label{rem:3.18}
推论  \ref{cor:3.23} 是自反空间与凸函数在变分法与最优化中的许多问题里如此重要的
主要原因.
\end{remark}

\begin{theorem}
\label{thm:3.24}
设 $E$ 与 $F$ 为两个自反 Banach 空间.设 $A:D(A)\subset E\to F$ 为稠定且闭的
无界线性算子.则 $D(A^\star)$ 在 $F^\star$ 中稠密.于是 $A^{\star\star}$ 有定义
($A^{\star\star}:D(A^{\star\star})\subset E^{\star\star}\to F^{\star\star}$),
它也可看作从 $E$ 到 $F$ 的无界算子.并且有
\begin{equation*}
\boxed{A^{\star\star}=A.}
\end{equation*}
\end{theorem}

\begin{proof}
\emph{1. $D(A^\star)$ 在 $F^\star$ 中稠密.} 设 $\varphi$ 为 $F^\star$ 上的连续
线性泛函, 它在 $D(A^\star)$ 上为零.由推论  \ref{cor:1.8}, 只须证明 $\varphi\equiv0$
在 $F^\star$ 上.由于 $F$ 自反, $\varphi\in F$, 且有
\begin{equation}
\dual{w}{\varphi}=0\quad\forall w\in D(A^\star).
\label{eq:3.6}
\end{equation}
若 $\varphi\ne0$, 则 $[0,\varphi]\notin G(A)$ (在 $E\times F$ 中).于是可用
$E\times F$ 中的闭超平面把 $[0,\varphi]$ 与 $G(A)$ 严格分离; 即存在
$[f,v]\in E^\star\times F^\star$ 与 $\alpha\in\R$, 使
\begin{equation*}
\dual{f}{u}+\dual{v}{Au}<\alpha<\dual{v}{\varphi}\quad\forall u\in D(A).
\end{equation*}
由此推出
\begin{equation*}
\dual{f}{u}+\dual{v}{Au}=0\quad\forall u\in D(A)
\end{equation*}
且
\begin{equation*}
\dual{v}{\varphi}\ne0.
\end{equation*}
于是 $v\in D(A^\star)$, 而在 \eqref{eq:3.6} 中取 $w=v$ 便导致矛盾.

\emph{2. $A^{\star\star}=A$.} 回忆 (见第  \ref{sec:2.6} 节)
\begin{equation*}
I[G(A^\star)]=G(A)^\perp
\end{equation*}
且
\begin{equation*}
I[G(A^{\star\star})]=G(A^\star)^\perp.
\end{equation*}
由此
\begin{equation*}
G(A^{\star\star})=G(A)^{\perp\perp}=G(A),
\end{equation*}
因为 $A$ 闭.
\end{proof}

\section{可分空间}
\label{sec:3.6}

\begin{definition}
若度量空间 $E$ 存在可数稠密子集 $D\subset E$, 则称 $E$ 是\emph{可分的}.
\end{definition}

分析中许多重要的空间都是可分的.显然, 有限维空间可分.我们将在第~
\ref{chap:4}~章 (另见
第~
\ref{chap:11}~章) 看到, 当 $1\le p<\infty$ 时 $L^p$ (以及 $\ell^p$) 空间可分.此外,
$C(K)$, 即紧度量空间 $K$ 上的连续函数空间, 也可分 (见问题 24).然而,
$L^\infty$ 与 $\ell^\infty$ 不可分 (见第~
\ref{chap:4}~章与第~
\ref{chap:11}~章).

\begin{proposition}
\label{prop:3.25}
设 $E$ 为可分度量空间, $F\subset E$ 为任一子集.则 $F$ 也可分.
\end{proposition}

\begin{proof}
设 $(u_n)$ 为 $E$ 的可数稠密子集.设 $(r_m)$ 为任一趋于 $0$ 的正数列.只要
$B(u_n,r_m)\cap F$ 非空, 就任取一点 $a_{m,n}\in B(u_n,r_m)\cap F$.集合
$(a_{m,n})$ 可数且在 $F$ 中稠密.
\end{proof}

\begin{theorem}
\label{thm:3.26}
设 $E$ 为 Banach 空间, 且 $E^\star$ 可分.则 $E$ 可分.
\end{theorem}

\begin{remark}
\label{rem:3.19}
逆命题不成立.我们将在第~
\ref{chap:4}~章看到, $E=L^1$ 可分, 但其共轭空间
$E^\star=L^\infty$ 不可分.
\end{remark}

\begin{proof}
设 $(f_n)_{n\ge1}$ 在 $E^\star$ 中可数稠密.由于
\begin{equation*}
\norm{f_n}=\sup_{\substack{x\in E\\ \norm{x}\le1}}\dual{f_n}{x},
\end{equation*}
我们可以找到 $x_n\in E$, 使
\begin{equation*}
\norm{x_n}=1\quad\text{且}\quad\dual{f_n}{x_n}\ge\frac12\norm{f_n}.
\end{equation*}
以 $L_0$ 记 $(x_n)_{n\ge1}$ 在 $\Q$ 上张成的向量空间; 即 $L_0$ 由 $(x_n)_{n\ge1}$
的, 系数在 $\Q$ 中的所有有限线性组合组成.我们断言 $L_0$ 可数.事实上, 对每个
整数 $n$, 设 $\Phi_n$ 为 $(x_k)_{1\le k\le n}$ 在 $\Q$ 上张成的向量空间.显然
$\Phi_n$ 可数, 且 $L_0=\bigcup_{n\ge1}\Phi_n$.

以 $L$ 记 $(x_n)_{n\ge1}$ 在 $\R$ 上张成的向量空间.当然, $L_0$ 是 $L$ 的稠密
子集.我们断言 $L$ 是 $E$ 的稠密子空间——这就完成了证明 ($L_0$ 将是 $E$ 的可数
稠密子集).设 $f\in E^\star$ 为在 $L$ 上为零的连续线性泛函; 由推论
 \ref{cor:1.8}, 须证明 $f=0$.给定任意 $\varepsilon>0$, 存在整数 $N$, 使
$\norm{f-f_N}<\varepsilon$.我们有
\begin{equation*}
\frac12\norm{f_N}\le\dual{f_N}{x_N}=\dual{f_N-f}{x_N}<\varepsilon
\end{equation*}
(因为 $\dual{f}{x_N}=0$).于是 $\norm{f}\le\norm{f-f_N}+\norm{f_N}<3\varepsilon$.
因此 $f=0$.
\end{proof}

\begin{corollary}
\label{cor:3.27}
设 $E$ 为 Banach 空间.则
\begin{equation*}
[E\ \text{自反且可分}]\ \iff\ [E^\star\ \text{自反且可分}].
\end{equation*}
\end{corollary}

\begin{proof}
由推论  \ref{cor:3.21} 与定理  \ref{thm:3.26} 已知
\begin{equation*}
[E^\star\ \text{自反且可分}]\ \Rightarrow\ [E\ \text{自反且可分}].
\end{equation*}
反之, 若 $E$ 自反且可分, 则 $E^{\star\star}=J(E)$ 也自反且可分; 因此 $E^\star$
自反且可分.
\end{proof}

可分性与弱拓扑的可度量化密切相关.回忆: 若拓扑空间 $X$ 上存在一个度量诱导 $X$
的拓扑, 则称 $X$ 是\emph{可度量化的}.

\begin{theorem}
\label{thm:3.28}
设 $E$ 为可分 Banach 空间.则 $B_{E^\star}$ 在\wstar{}拓扑
$\sigma(E^\star,E)$ 中可度量化.

反之, 若 $B_{E^\star}$ 在 $\sigma(E^\star,E)$ 中可度量化, 则 $E$ 可分.
\end{theorem}

还有一个``对偶''的命题.

\begin{theorem}
\label{thm:3.29}
设 $E$ 为 Banach 空间, 且 $E^\star$ 可分.则 $B_E$ 在弱拓扑
$\sigma(E,E^\star)$ 中可度量化.

反之, 若 $B_E$ 在 $\sigma(E,E^\star)$ 中可度量化, 则 $E^\star$ 可分.
\end{theorem}

\begin{proof}[定理 \ref{thm:3.28}]
设 $(x_n)_{n\ge1}$ 为 $B_E$ 的可数稠密子集.对每个 $f\in E^\star$ 置
\begin{equation*}
[f]=\sum_{n=1}^\infty\frac{1}{2^n}\abs{\dual{f}{x_n}}.
\end{equation*}
显然, $[\cdot]$ 是 $E^\star$ 上的范数, 且 $[f]\le\norm{f}$.令
$d(f,g)=[f-g]$ 为相应的度量.我们将证明: $d$ 在 $B_{E^\star}$ 上诱导的拓扑与
$\sigma(E^\star,E)$ 限制在 $B_{E^\star}$ 上的拓扑相同.

(a) 设 $f_0\in B_{E^\star}$, $V$ 为 $f_0$ 关于 $\sigma(E^\star,E)$ 的邻域.须找出
$r>0$, 使
\begin{equation*}
U=\{f\in B_{E^\star};\ d(f,f_0)<r\}\subset V.
\end{equation*}
如通常那样, 可设 $V$ 具有形式
\begin{equation*}
V=\{f\in B_{E^\star};\ \abs{\dual{f-f_0}{y_i}}<\varepsilon
\ \ \forall i=1,2,\dots,k\},
\end{equation*}
其中 $\varepsilon>0$, $y_1,y_2,\dots,y_k\in E$.不妨设
$\norm{y_i}\le1$ 对每个 $i=1,2,\dots,k$ 成立.对每个 $i$ 存在整数 $n_i$, 使
\begin{equation*}
\norm{y_i-x_{n_i}}<\varepsilon/4
\end{equation*}
(因为集合 $(x_n)_{n\ge1}$ 在 $B_E$ 中稠密).取 $r>0$ 充分小, 使
\begin{equation*}
2^{n_i}r<\varepsilon/2\quad\forall i=1,2,\dots,k.
\end{equation*}
我们断言: 对这样的 $r$, $U\subset V$.事实上, 若 $d(f,f_0)<r$, 则
\begin{equation*}
\frac{1}{2^{n_i}}\abs{\dual{f-f_0}{x_{n_i}}}<r\quad\forall i=1,2,\dots,k,
\end{equation*}
从而对每个 $i=1,2,\dots,k$,
\begin{equation*}
\abs{\dual{f-f_0}{y_i}}
=\abs{\dual{f-f_0}{y_i-x_{n_i}}+\dual{f-f_0}{x_{n_i}}}
<\frac{\varepsilon}{2}+\frac{\varepsilon}{2}.
\end{equation*}
由此 $f\in V$.

(b) 设 $f_0\in B_{E^\star}$.给定 $r>0$, 须找出 $f_0$ 关于 $\sigma(E^\star,E)$ 的
邻域 $V$, 使
\begin{equation*}
V\subset U=\{f\in B_{E^\star};\ d(f,f_0)<r\}.
\end{equation*}
我们取 $V$ 为
\begin{equation*}
V=\{f\in B_{E^\star};\ \abs{\dual{f-f_0}{x_i}}<\varepsilon
\ \ \forall i=1,2,\dots,k\},
\end{equation*}
其中 $\varepsilon$ 与 $k$ 待定, 使得 $V\subset U$.对 $f\in V$ 有
\begin{equation*}
\begin{aligned}
d(f,f_0)&=\sum_{n=1}^k\frac{1}{2^n}\abs{\dual{f-f_0}{x_n}}
+\sum_{n=k+1}^\infty\frac{1}{2^n}\abs{\dual{f-f_0}{x_n}}\\
&<\varepsilon+2\sum_{n=k+1}^\infty\frac{1}{2^n}
=\varepsilon+\frac{1}{2^{k-1}}.
\end{aligned}
\end{equation*}
因此只须取 $\varepsilon=\frac r2$, 并取 $k$ 充分大, 使
$\frac{1}{2^{k-1}}<\frac r2$.
\end{proof}

\begin{proof}[定理 \ref{thm:3.28} 的逆命题]
$\star$ 反过来, 设 $B_{E^\star}$ 在 $\sigma(E^\star,E)$ 中可度量化, 我们来证明
$E$ 可分.置
\begin{equation*}
U_n=\{f\in B_{E^\star};\ d(f,0)<1/n\},
\end{equation*}
并设 $V_n$ 为 $0$ 在 $\sigma(E^\star,E)$ 中的邻域, 使 $V_n\subset U_n$.可设 $V_n$
具有形式
\begin{equation*}
V_n=\{f\in B_{E^\star};\ \abs{\dual{f}{x}}<\varepsilon_n\ \ \forall x\in\Phi_n\},
\end{equation*}
其中 $\varepsilon_n>0$, $\Phi_n$ 为 $E$ 的有限子集.置
\begin{equation*}
D=\bigcup_{n=1}^\infty\Phi_n,
\end{equation*}
则 $D$ 可数.

我们断言: 由 $D$ 张成的向量空间在 $E$ 中稠密 (由此推出 $E$ 可分).事实上, 设
$f\in E^\star$ 满足 $\dual{f}{x}=0$, $\forall x\in D$.则 $f\in V_n$,
$\forall n$, 从而 $f\in U_n$, $\forall n$, 于是 $f=0$.
\end{proof}

\begin{proof}[定理 \ref{thm:3.29}]
蕴含关系
\begin{equation*}
[E^\star\ \text{可分}]\ \Rightarrow\ [B_E\ \text{在}\ \sigma(E,E^\star)\ \text{中可度量化}]
\end{equation*}
的证明与上面完全相同——只需把 $E$ 与 $E^\star$ 的角色互换.逆命题的证明更为
精细 (请找出上面的证明在哪里失效); 我们引用 N.~Dunford--J.~T.~Schwartz
\bref{DunfordSchwartz}{1} 或习题  \ref{exer:3.24}.
\end{proof}

\begin{remark}
\label{rem:3.20}
应当再次强调 (见注记 \ref{rem:3.3}): 在无穷维空间中, 整个 $E$ (相应地,
$E^\star$) 上的弱拓扑 $\sigma(E,E^\star)$ (相应地, \wstar{}拓扑
$\sigma(E^\star,E)$) 不可度量化; 见习题  \ref{exer:3.8}.特别地, 范数
$[\cdot]$ 在整个 $E^\star$ 上诱导的拓扑与\wstar{}拓扑并不相同.
\end{remark}

\begin{corollary}
\label{cor:3.30}
设 $E$ 为可分 Banach 空间, $(f_n)$ 为 $E^\star$ 中的有界序列.则存在子列
$(f_{n_k})$, 它在\wstar{}拓扑 $\sigma(E^\star,E)$ 中收敛.
\end{corollary}

\begin{proof}
不妨设 $\norm{f_n}\le1$ 对一切 $n$ 成立.集合 $B_{E^\star}$ 关于拓扑
$\sigma(E^\star,E)$ 紧且可度量化 (由定理  \ref{thm:3.16} 与  \ref{thm:3.28}).
结论由此得出.
\end{proof}

现在可以回到定理  \ref{thm:3.18} 的证明了:

\begin{proof}[定理 \ref{thm:3.18}]
设 $M_0$ 为由 $x_n$ 张成的向量空间, $M=\overline{M_0}$.显然 $M$ 可分 (见定理
 \ref{thm:3.26} 的证明).此外, $M$ 自反 (由命题  \ref{prop:3.20}).于是 $B_M$ 在
弱拓扑 $\sigma(M,M^\star)$ 中紧且可度量化, 因为 $M^\star$ 可分 (这里用到推论
 \ref{cor:3.27} 与定理  \ref{thm:3.29}).因此可以找到子列 $(x_{n_k})$, 它在
$\sigma(M,M^\star)$ 中弱收敛, 从而 $(x_{n_k})$ 也在 $\sigma(E,E^\star)$ 中弱收敛
(如命题  \ref{prop:3.20} 的证明中那样).
\end{proof}

\section{一致凸空间}
\label{sec:3.7}

\begin{definition}
若 Banach 空间满足
\begin{equation*}
\forall\varepsilon>0\ \exists\delta>0\ \text{使得}\
\Bigl[x,y\in E,\ \norm{x}\le1,\ \norm{y}\le1\ \text{且}\ \norm{x-y}>\varepsilon
\Bigr]\ \Rightarrow\ \Bigl\|\frac{x+y}{2}\Bigr\|<1-\delta,
\end{equation*}
则称它是\emph{一致凸的}.
\end{definition}

一致凸性是单位球的一个\emph{几何}性质: 若把一根长为 $\varepsilon>0$ 的直尺在单位
球中滑动, 则它的中点必须落在某个半径为 $(1-\delta)$ 的球内, 其中 $\delta>0$.特别
地, 单位\emph{球面}必须是``圆的'', 不能包含任何线段.

\begin{example}
\label{exa:3.3}
设 $E=\R^2$.范数 $\norm{x}_2=\bigl[\abs{x_1}^2+\abs{x_2}^2\bigr]^{1/2}$ 一致凸,
而范数 $\norm{x}_1=\abs{x_1}+\abs{x_2}$ 与
$\norm{x}_\infty=\max(\abs{x_1},\abs{x_2})$ 不是一致凸的.只要盯着单位球看,
如图  \ref{fig:3} 所示, 这一点就一目了然.
\end{example}

\begin{figure}[htbp]
	\centering
	\begin{tikzpicture}[scale=1.05]
		% 左图: l2 范数的单位圆
		\begin{scope}
			\draw[-{Stealth[length=2mm]}] (-1.55,0) -- (1.75,0);
			\draw[-{Stealth[length=2mm]}] (0,-1.55) -- (0,1.75);
			\draw[thick] (0,0) circle (1.0);
			\fill (0,0) circle (1.5pt);
			\node[below=6mm, align=center] at (0,-1.0) {\small 关于范数 $\norm{\cdot}_2$ 的 $E$ 的单位球};
		\end{scope}
		% 右图: l1 范数的单位菱形
		\begin{scope}[xshift=6.0cm]
			\draw[-{Stealth[length=2mm]}] (-1.55,0) -- (1.75,0);
			\draw[-{Stealth[length=2mm]}] (0,-1.55) -- (0,1.75);
			\draw[thick] (1,0) -- (0,1) -- (-1,0) -- (0,-1) -- cycle;
			\fill (0,0) circle (1.5pt);
			\fill (1,0) circle (1.5pt);
			\fill (-1,0) circle (1.5pt);
			\fill (0,1) circle (1.5pt);
			\fill (0,-1) circle (1.5pt);
			\node[below=6mm, align=center] at (0,-1.0) {\small 关于范数 $\norm{\cdot}_1$ 的 $E$ 的单位球};
		\end{scope}
	\end{tikzpicture}
	\caption{}
	\label{fig:3}
\end{figure}

\begin{example}
\label{exa:3.4}
我们将在第~
\ref{chap:4}~章与第~
\ref{chap:5}~章看到, 当 $1<p<\infty$ 时 $L^p$ 空间一致凸, 而 Hilbert
空间也一致凸.
\end{example}

\index{MilmanPettisdingli@Milman--Pettis 定理 (Milman--Pettis theorem)}
\begin{keytheorem}[Milman--Pettis]
\label{thm:3.31}
每个一致凸 Banach 空间都自反.
\end{keytheorem}

\begin{remark}
\label{rem:3.21}
一致凸性是范数的一个\emph{几何}性质; 等价的范数不必一致凸.另一方面, 自反性是
一个\emph{拓扑}性质: 自反空间关于等价范数仍然自反.定理  \ref{thm:3.31} 的一个
引人注目的特点是: 一个几何性质蕴含了一个拓扑性质.一致凸性常被用作证明自反性的
工具; 但它不是终极工具——存在一些古怪的自反空间, 它们不允许任何一致凸的等价
范数!
\end{remark}

\begin{proof}
设 $\xi\in E^{\star\star}$ 且 $\norm{\xi}=1$.须证明 $\xi\in J(B_E)$.由于
$J(B_E)$ 在 $E^{\star\star}$ 中关于强拓扑闭, 只须证明
\begin{equation}
\forall\varepsilon>0\ \exists x\in B_E\ \text{使得}\ \norm{\xi-J(x)}\le\varepsilon.
\label{eq:3.7}
\end{equation}
固定 $\varepsilon>0$, 设 $\delta>0$ 为一致凸性的模.选取某个 $f\in E^\star$, 使
$\norm{f}=1$ 且
\begin{equation}
\dual{\xi}{f}>1-(\delta/2)
\label{eq:3.8}
\end{equation}
(这是可能的, 因为 $\norm{\xi}=1$).置
\begin{equation*}
V=\{\eta\in E^{\star\star};\ \abs{\dual{\eta-\xi}{f}}<\delta/2\},
\end{equation*}
于是 $V$ 是 $\xi$ 在拓扑 $\sigma(E^{\star\star},E^\star)$ 中的一个邻域.由于
$J(B_E)$ 关于 $\sigma(E^{\star\star},E^\star)$ 在 $B_{E^{\star\star}}$ 中稠密
(引理  \ref{lem:3.4}), 我们知道 $V\cap J(B_E)\ne\emptyset$, 从而存在 $x\in B_E$
使 $J(x)\in V$.我们断言这个 $x$ 满足 \eqref{eq:3.7}.反设
$\norm{\xi-Jx}>\varepsilon$, 即 $\xi\in(Jx+\varepsilon B_{E^{\star\star}})^c=W$.
集合 $W$ 也是 $\xi$ 在拓扑 $\sigma(E^{\star\star},E^\star)$ 中的一个邻域 (因为
$B_{E^{\star\star}}$ 在 $\sigma(E^{\star\star},E^\star)$ 中闭).再次利用引理
 \ref{lem:3.4}, 我们知道 $V\cap W\cap J(B_E)\ne\phi$, 即存在 $y\in B_E$ 使
$J(y)\in V\cap W$.由 $J(x),J(y)\in V$, 得
\begin{equation*}
\abs{\dual{f}{x}-\dual{\xi}{f}}<\delta/2
\end{equation*}
以及
\begin{equation*}
\abs{\dual{f}{y}-\dual{\xi}{f}}<\delta/2.
\end{equation*}
把这两个不等式相加, 得到
\begin{equation*}
2\dual{\xi}{f}<\dual{f}{x+y}+\delta\le\norm{x+y}+\delta.
\end{equation*}
结合 \eqref{eq:3.8}, 得到
\begin{equation*}
\Bigl\|\frac{x+y}{2}\Bigr\|>1-\delta.
\end{equation*}
由一致凸性推出 $\norm{x-y}\le\varepsilon$; 这是荒谬的, 因为 $J(y)\in W$
(即 $\norm{x-y}>\varepsilon$).
\end{proof}

最后给出一致凸空间的一个有用性质.

\begin{proposition}
\index{RadonRieszxingzhi@Radon--Riesz 性质 (Radon--Riesz property)}
\label{prop:3.32}
假设 $E$ 是一致凸 Banach 空间.设 $(x_n)$ 为 $E$ 中的序列, 使
$x_n\weakto x$ 弱收敛于 $\sigma(E,E^\star)$, 且
\begin{equation*}
\limsup\norm{x_n}\le\norm{x}.
\end{equation*}
则 $x_n\to x$ 强收敛.
\end{proposition}

\begin{proof}
总可设 $x\ne0$ (否则结论显然).置
\begin{equation*}
\lambda_n=\max(\norm{x_n},\norm{x}),\quad
y_n=\lambda_n^{-1}x_n,\quad y=\norm{x}^{-1}x,
\end{equation*}
于是 $\lambda_n\to\norm{x}$ 且 $y_n\weakto y$ 弱收敛于 $\sigma(E,E^\star)$.
由此
\begin{equation*}
\norm{y}\le\liminf\norm{(y_n+y)/2}
\end{equation*}
(见命题  \ref{prop:3.5}(iii)).另一方面, $\norm{y}=1$ 且 $\norm{y_n}\le1$, 因此
事实上 $\norm{(y_n+y)/2}\to1$.由一致凸性推出 $\norm{y_n-y}\to0$, 从而
$x_n\to x$ 强收敛.
\end{proof}

\section*{第 3 章评注}
\addcontentsline{toc}{section}{第 3 章评注}
\markright{第 3 章评注}

\cmtnum{1.}
拓扑 $\sigma(E,E^\star)$, $\sigma(E^\star,E)$ 等都是局部凸拓扑.作为这样的拓扑,
它们享有局部凸空间的一切性质; 例如 Hahn--Banach 定理 (几何形式), Krein--Milman 定理等仍然成立; 参见例如 N.~Bourbaki \bref{Bourbaki}{1}, A.~Knapp \bref{Knapp}{1}, 以及问题 9.

\cmtnum{2.}
下面是\wstar{}拓扑的另一个值得提及的引人注目的性质.

\begin{startheorem}[Banach--Dieudonn\'e--Krein--\v{S}mulian]
\label{thm:3.33}
设 $E$ 为 Banach 空间, $C\subset E^\star$ 凸.假设对每个整数 $n$, 集合
$C\cap(nB_{E^\star})$ 关于拓扑 $\sigma(E^\star,E)$ 闭.则 $C$ 关于拓扑
$\sigma(E^\star,E)$ 闭.
\end{startheorem}

其证明可在例如 N.~Bourbaki \bref{Bourbaki}{1}, R.~Larsen \bref{Larsen}{1}, R.~Holmes \bref{Holmes}{1}, N.~Dunford--J.~T.~Schwartz \bref{DunfordSchwartz}{1}, H.~Schaefer \bref{Schaefer}{1} 以及问题 11 中找到.上述文献也包含了大量与
Eberlein--\v{S}mulian 定理 (定理  \ref{thm:3.19}) 相关的材料.

\cmtnum{3.}
\emph{对偶向量空间}理论——它推广了对偶 $\dual{E}{E^\star}$——在四十年代末与
五十年代初非常流行, 尤其是在与广义函数理论的联系中.若在 $X\times Y$ 上存在一个
分离点的双线性形式 $\langle\,,\rangle$ (即 $\forall x\ne0\ \exists y$ 使
$\dual{x}{y}\ne0$, 且 $\forall y\ne0\ \exists x$ 使 $\dual{x}{y}\ne0$), 则称
两个向量空间 $X$ 与 $Y$ 处于对偶之中.在 $X$ (或 $Y$) 上可以定义许多拓扑, 例如
弱拓扑 $\sigma(X,Y)$, Mackey 拓扑 $\tau(X,Y)$ 以及强拓扑 $\beta(X,Y)$.这些
拓扑在非 Banach 空间——例如广义函数理论中使用的空间——里很有意义.关于这一
主题, 读者可参阅例如 N.~Bourbaki \bref{Bourbaki}{1}, H.~Schaefer
\bref{Schaefer}{1}, G.~K\"othe \bref{Kothe}{1}, F.~Treves \bref{Treves}{1}, J.~Kelley--I.~Namioka \bref{KelleyNamioka}{1}, R.~Edwards \bref{Edwards}{1}, J.~Horv\'ath \bref{Horvath}{1} 等.

\cmtnum{4.}
可分性, 自反性与一致凸性这些性质也与函数 $x\mapsto\norm{x}$ 的可微性性质有关
(参见例如 J.~Diestel \bref{Diestel}{1}, B.~Beauzamy \bref{Beauzamy}{1} 以及
问题 13).关于具有良好几何性质的等价范数的存在性, 已有大量研究.例如, 怎样判断
一个 Banach 空间是否允许一个等价的\emph{一致凸}范数? 这一信息有多大用处? (这样的
空间称为\emph{超自反的}; 参见例如 J.~Diestel \bref{Diestel}{1} 或 B.~Beauzamy
\bref{Beauzamy}{1}.)\ Banach 空间的几何自六十年代初以来蓬勃发展, 已成为一个活跃
的领域, 与 A.~Dvoretzky, A.~Grothendieck, R.~C.~James, J.~Lindenstrauss, V.~Milman, L.~Tzafriri (以及他们在以色列的学派), A.~Pelczynski, P.~Enflo, L.~Schwartz (以及他包括 G.~Pisier, B.~Maurey, B.~Beauzamy 在内的学派), W.~B.~Johnson, H.~P.~Rosenthal, J.~Bourgain, D.~Preiss, M.~Talagrand, T.~Gowers 等许多人的名字联系在一起.关于这一主题, 读者可参阅 B.~Beauzamy
\bref{Beauzamy}{1}, J.~Diestel \bref{Diestel}{1}, \bref{Diestel}{2}, J.~Lindenstrauss--L.~Tzafriri \bref{LindenstraussTzafriri}{2}, L.~Schwartz
\bref{Schwartz}{2}, R.~Deville--G.~Godefroy--V.~Zizler \bref{Deville}{1}, Y.~Benyamini--J.~Lindenstrauss \bref{Benyamini}{1}, F.~Albiac--N.~Kalton
\bref{Albiac}{1}, A.~Pietsch \bref{Pietsch}{1} 等著作.

\section*{第 3 章习题}
\addcontentsline{toc}{section}{第 3 章习题}
\markright{第 3 章习题}

\begin{exercise}
\label{exer:3.1}
设 $E$ 为 Banach 空间, $A\subset E$ 为在弱拓扑 $\sigma(E,E^\star)$ 中紧的子集.
证明 $A$ 有界.
\end{exercise}

\begin{exercise}
\label{exer:3.2}
设 $E$ 为 Banach 空间, $(x_n)$ 为 $E$ 中的序列, 使 $x_n\weakto x$ 在弱拓扑
$\sigma(E,E^\star)$ 中弱收敛.置
\begin{equation*}
\sigma_n=\frac1n(x_1+x_2+\cdots+x_n).
\end{equation*}
证明 $\sigma_n\weakto x$ 在弱拓扑 $\sigma(E,E^\star)$ 中弱收敛.
\end{exercise}

\begin{exercise}
\label{exer:3.3}
设 $E$ 为 Banach 空间, $A\subset E$ 为凸子集.证明 $A$ 在强拓扑中的闭包与它在
弱拓扑 $\sigma(E,E^\star)$ 中的闭包相同.
\end{exercise}

\begin{exercise}
\label{exer:3.4}
设 $E$ 为 Banach 空间, $(x_n)$ 为 $E$ 中的序列, 使 $x_n\weakto x$ 在弱拓扑
$\sigma(E,E^\star)$ 中弱收敛.

\begin{enumerate}[label=\arabic*., leftmargin=2.2em, itemsep=0.25em]
	\item 证明存在 $E$ 中的序列 $(y_n)$, 使得
	      \begin{equation*}
	      \text{(a)}\quad y_n\in\operatorname{conv}\Bigl(\bigcup_{i=n}^\infty\{x_i\}\Bigr)
	      \quad\forall n
	      \end{equation*}
	      并且
	      \begin{equation*}
	      \text{(b)}\quad y_n\to x\quad\text{强收敛}.
	      \end{equation*}
	\item 证明存在 $E$ 中的序列 $(z_n)$, 使得
	      \begin{equation*}
	      \text{(a$'$)}\quad z_n\in\operatorname{conv}\Bigl(\bigcup_{i=1}^n\{x_i\}\Bigr)
	      \quad\forall n
	      \end{equation*}
	      并且
	      \begin{equation*}
	      \text{(b$'$)}\quad z_n\to x\quad\text{强收敛}.
	      \end{equation*}
\end{enumerate}
\end{exercise}

\begin{exercise}
\label{exer:3.5}
设 $E$ 为 Banach 空间, $K\subset E$ 为在强拓扑中紧的子集.设 $(x_n)$ 为 $K$ 中的
序列, 使 $x_n\weakto x$ 弱收敛于 $\sigma(E,E^\star)$.证明 $x_n\to x$ 强收敛.

[提示: 用反证法.]
\end{exercise}

\begin{exercise}
\label{exer:3.6}
设 $X$ 为拓扑空间, $E$ 为 Banach 空间.设 $u,v:X\to E$ 为两个连续映射, 其中
$E$ 配上弱拓扑 $\sigma(E,E^\star)$.

\begin{enumerate}[label=\arabic*., leftmargin=2.2em, itemsep=0.25em]
	\item 证明映射 $x\mapsto u(x)+v(x)$ 是从 $X$ 到配上 $\sigma(E,E^\star)$ 的
	      $E$ 的连续映射.
	\item 设 $a:X\to\R$ 为连续函数.证明映射 $x\mapsto a(x)u(x)$ 是从 $X$ 到
	      配上 $\sigma(E,E^\star)$ 的 $E$ 的连续映射.
\end{enumerate}
\end{exercise}

\begin{exercise}
\label{exer:3.7}
设 $E$ 为 Banach 空间, $A\subset E$ 为在弱拓扑 $\sigma(E,E^\star)$ 中闭的子集.
设 $B\subset E$ 为在弱拓扑 $\sigma(E,E^\star)$ 中紧的子集.

\begin{enumerate}[label=\arabic*., leftmargin=2.2em, itemsep=0.25em]
	\item 证明 $A+B$ 在 $\sigma(E,E^\star)$ 中闭.
	\item 再假设 $A$ 与 $B$ 凸, 非空且不相交.证明存在一个闭超平面严格分离 $A$
	      与 $B$.
\end{enumerate}
\end{exercise}

\begin{exercise}
\label{exer:3.8}
设 $E$ 为无穷维 Banach 空间.我们的目的是证明配上弱拓扑的 $E$ 不可度量化.
反设 $E$ 上存在一个度量 $d(x,y)$, 它在 $E$ 上诱导出与 $\sigma(E,E^\star)$ 相同的
拓扑.

\begin{enumerate}[label=\arabic*., leftmargin=2.2em, itemsep=0.25em]
	\item 对每个整数 $k\ge1$, 用 $V_k$ 记 $0$ 在拓扑 $\sigma(E,E^\star)$ 中的
	      一个邻域, 使得
	      \begin{equation*}
	      V_k\subset\Bigl\{x\in E;\ d(x,0)<\frac1k\Bigr\}.
	      \end{equation*}
	      证明存在 $E^\star$ 中的序列 $(f_n)$, 使得每个 $g\in E^\star$ 都是
	      这些 $f_n$ 的 (有限) 线性组合.

	      [提示: 利用引理  \ref{lem:3.2}.]
	\item 推出 $E^\star$ 是有限维的.

	      [提示: 像习题 1.5 那样利用 Baire 纲定理.]
	\item 得出结论.
	\item 用类似的方法证明配上\wstar{}拓扑 $\sigma(E^\star,E)$ 的 $E^\star$
	      不可度量化.
\end{enumerate}
\end{exercise}

\begin{exercise}
\label{exer:3.9}
设 $E$ 为 Banach 空间, $M\subset E$ 为线性子空间, $f_0\in E^\star$.证明存在
某个 $g_0\in M^\perp$, 使得
\begin{equation*}
\inf_{g\in M^\perp}\norm{f_0-g}=\norm{f_0-g_0}.
\end{equation*}
建议两种方法:

\begin{enumerate}[label=\arabic*., leftmargin=2.2em, itemsep=0.25em]
	\item 利用定理  \ref{thm:1.12}.
	\item 利用\wstar{}拓扑 $\sigma(E^\star,E)$.
\end{enumerate}
\end{exercise}

\begin{exercise}
\label{exer:3.10}
设 $E$ 与 $F$ 为两个 Banach 空间.设 $T\in\LL(E,F)$, 从而
$T^\star\in\LL(F^\star,E^\star)$.证明 $T^\star$ 是从配上 $\sigma(F^\star,F)$
的 $F^\star$ 到配上 $\sigma(E^\star,E)$ 的 $E^\star$ 的连续映射.
\end{exercise}

\begin{exercise}
\label{exer:3.11}
设 $E$ 为 Banach 空间, $A:E\to E^\star$ 为定义在 $D(A)=E$ 上的单调映射;
见习题 2.6.假设对每个 $x,y\in E$, 映射
\begin{equation*}
t\in\R\mapsto\dual{A(x+ty)}{y}
\end{equation*}
在 $t=0$ 处连续.证明 $A$ 是从强拓扑下的 $E$ 到配上 $\sigma(E^\star,E)$ 的
$E^\star$ 的连续映射.
\end{exercise}

\begin{exercise}
\label{exer:3.12}
设 $E$ 为 Banach 空间, $x_0\in E$.设 $\varphi:E\to(-\infty,+\infty]$ 为凸的
l.s.c. 函数, $\varphi\not\equiv+\infty$.

\begin{enumerate}[label=\arabic*., leftmargin=2.2em, itemsep=0.25em]
	\item 证明下列性质等价:
	      \begin{equation*}
	      \text{(A)}\quad
	      \begin{aligned}
	      &\exists R,\ \exists M<+\infty\ \text{使得}\ \varphi(x)\le M,\\
	      &\forall x\in E\ \text{满足}\ \norm{x-x_0}\le R,
	      \end{aligned}
	      \end{equation*}
	      \begin{equation*}
	      \text{(B)}\quad
	      \lim_{\substack{f\in E^\star\\ \norm{f}\to\infty}}
	      \{\varphi^\star(f)-\dual{f}{x_0}\}=+\infty.
	      \end{equation*}
	\item 假设 (A) 或 (B) 成立, 证明
	      \begin{equation*}
	      \inf_{f\in E^\star}\{\varphi^\star(f)-\dual{f}{x_0}\}
	      \end{equation*}
	      可以达到.

	      [提示: 利用\wstar{}拓扑 $\sigma(E^\star,E)$ 或定理  \ref{thm:1.12}.]

	      这个 $\inf$ 的值是多少?
\end{enumerate}
\end{exercise}

\begin{exercise}
\label{exer:3.13}
设 $E$ 为 Banach 空间.设 $(x_n)$ 为 $E$ 中的序列, $x\in E$.置
\begin{equation*}
K_n=\overline{\operatorname{conv}\Bigl(\bigcup_{i=n}^\infty\{x_i\}\Bigr)}.
\end{equation*}

\begin{enumerate}[label=\arabic*., leftmargin=2.2em, itemsep=0.25em]
	\item 证明: 若 $x_n\weakto x$ 弱收敛于 $\sigma(E,E^\star)$, 则
	      \begin{equation*}
	      \bigcap_{n=1}^\infty K_n=\{x\}.
	      \end{equation*}
	\item 假设 $E$ 自反.证明: 若 $(x_n)$ 有界且
	      $\bigcap_{n=1}^\infty K_n=\{x\}$, 则 $x_n\weakto x$ 弱收敛于
	      $\sigma(E,E^\star)$.
	\item 假设 $E$ 有限维且 $\bigcap_{n=1}^\infty K_n=\{x\}$.证明
	      $x_n\to x$.

	      [注意这里并没有假设 $(x_n)$ 有界.]
	\item 在 $\ell^p$ ($1<p<\infty$, 见第~
\ref{chap:11}~章) 中构造一个序列 $(x_n)$, 使得
	      $\bigcap_{n=1}^\infty K_n=\{x\}$, 而 $(x_n)$ \emph{不}有界.

	      [第 3 问与第 4 问的结果得自 Guy Amram 与 Daniel Baffet.]
\end{enumerate}
\end{exercise}

\begin{exercise}
\label{exer:3.14}
设 $E$ 为自反 Banach 空间, $I$ 为指标集.考虑 $E^\star$ 中的一族 $(f_i)_{i\in I}$
与 $\R$ 中的一族 $(\alpha_i)_{i\in I}$.设 $M>0$.证明下列性质等价:
\begin{equation*}
\text{(A)}\quad
\begin{cases}
\text{存在某个}\ x\in E,\ \norm{x}\le M,\ \text{使得}\ \dual{f_i}{x}=\alpha_i\\
\text{对每个}\ i\in I.
\end{cases}
\end{equation*}
\begin{equation*}
\text{(B)}\quad
\begin{cases}
\text{对}\ \R\ \text{中任一满足}\ J\subset I,\ J\ \text{有限的族}\
(\beta_i)_{i\in J},\ \text{都有}\\
\abs{\sum_{i\in J}\beta_i\alpha_i}\le M\bigl\|\sum_{i\in J}\beta_if_i\bigr\|.
\end{cases}
\end{equation*}
与习题 1.10, 1.11 及引理  \ref{lem:3.3} 比较.
\end{exercise}

\begin{exercise}[凸集上测度的重心.]
\label{exer:3.15}
设 $E$ 为自反 Banach 空间, $K\subset E$ 有界, 闭且凸.下文给 $K$ 配上
$\sigma(E,E^\star)$, 于是 $K$ 紧.设 $F=C(K)$ 配上通常的范数.固定某个
$\mu\in F^\star$ 且 $\norm{\mu}=1$, 并假设 $\mu\ge0$, 即
\begin{equation*}
\dual{\mu}{u}\ge0\quad\forall u\in C(K),\ u\ge0\ \text{在}\ K\ \text{上}.
\end{equation*}
证明存在唯一的元素 $x_0\in K$, 使得
\begin{equation}
\dual{\mu}{f|_K}=\dual{f}{x_0}\quad\forall f\in E^\star.
\tag{1}
\end{equation}
[提示: 先找出某个满足 (1) 的 $x_0\in E$, 然后借助 Hahn--Banach 定理证明
$x_0\in K$.]
\end{exercise}

\begin{exercise}
\label{exer:3.16}
设 $E$ 为 Banach 空间.

\begin{enumerate}[label=\arabic*., leftmargin=2.2em, itemsep=0.25em]
	\item 设 $(f_n)$ 为 $E^\star$ 中的序列, 使得对每个 $x\in E$,
	      $\dual{f_n}{x}$ 收敛到一个极限.证明存在某个 $f\in E^\star$, 使得
	      $f_n\weakstar f$ 在 $\sigma(E^\star,E)$ 中\wstar{}收敛.
	\item 这里假设 $E$ 自反.设 $(x_n)$ 为 $E$ 中的序列, 使得对每个
	      $f\in E^\star$, $\dual{f}{x_n}$ 收敛到一个极限.证明存在某个 $x\in E$,
	      使得 $x_n\weakto x$ 在 $\sigma(E,E^\star)$ 中弱收敛.
	\item 在非自反空间 $E$ 中构造一个例子, 使第 2 问的结论失效.

	      [提示: 取 $E=c_0$ (见第~\ref{sec:11.3}~节) 以及
	      $x_n=(1,1,\dots,1_{(n)},0,0,\dots)$.]
\end{enumerate}
\end{exercise}

\begin{exercise}
\label{exer:3.17}
\begin{enumerate}[label=\arabic*., leftmargin=2.2em, itemsep=0.25em]
	\item 设 $(x_n)$ 为 $\ell^p$ 中的序列, $1\le p\le\infty$.假设 $x_n\weakto x$
	      在 $\sigma(\ell^p,\ell^{p'})$ 中弱收敛, 证明:
	      \begin{enumerate}[label=(\alph*), leftmargin=2.0em, itemsep=0.15em]
	      	\item $(x_n)$ 在 $\ell^p$ 中有界,
	      	\item 对每个 $i$ 有 $x^n_i\to x_i$ ($n\to\infty$), 其中
	      	      $x_n=(x^n_1,x^n_2,\dots,x^n_i,\dots)$,
	      	      $x=(x_1,x_2,\dots,x_i,\dots)$.
	      \end{enumerate}
	\item 反过来, 设 $(x_n)$ 为 $\ell^p$ 中的序列, $1<p\le\infty$.假设 (a) 与
	      (b) 成立 (对某个记作 $x_i$ 的极限).证明 $x\in\ell^p$ 且 $x_n\weakto x$
	      在 $\sigma(\ell^p,\ell^{p'})$ 中弱收敛.
\end{enumerate}
\end{exercise}

\begin{exercise}
\label{exer:3.18}
对每个整数 $n\ge1$, 置
\begin{equation*}
e_n=(0,0,\dots,1_{(n)},0,\dots).
\end{equation*}
\begin{enumerate}[label=\arabic*., leftmargin=2.2em, itemsep=0.25em]
	\item 证明 $e_n\weakto0$ 在弱拓扑 $\sigma(\ell^p,\ell^{p'})$ 中弱收敛,
	      $1<p\le\infty$.
	\item 证明不存在子列 $(e_{n_k})$ 在 $\ell^1$ 中关于 $\sigma(\ell^1,\ell^\infty)$
	      收敛.
	\item 构造一个 Banach 空间 $E$ 与 $E^\star$ 中的序列 $(f_n)$ 的例子, 使得
	      $\norm{f_n}=1\ \forall n$, 且 $(f_n)$ 没有在 $\sigma(E^\star,E)$ 中收敛的
	      子列.这与 $B_{E^\star}$ 在拓扑 $\sigma(E^\star,E)$ 中的紧性矛盾吗?

	      [提示: 取 $E=\ell^\infty$.]
\end{enumerate}
\end{exercise}

\begin{exercise}
\label{exer:3.19}
设 $E=\ell^p$ 且 $F=\ell^q$, 其中 $1<p<\infty$, $1<q<\infty$.设
$a:\R\to\R$ 为连续函数, 使得
\begin{equation*}
\abs{a(t)}\le C\abs{t}^{p/q}\quad\forall t\in\R.
\end{equation*}
给定
\begin{equation*}
x=(x_1,x_2,\dots,x_i,\dots)\in\ell^p,
\end{equation*}
置
\begin{equation*}
Ax=\bigl(a(x_1),a(x_2),\dots,a(x_i),\dots\bigr).
\end{equation*}
\begin{enumerate}[label=\arabic*., leftmargin=2.2em, itemsep=0.25em]
	\item 证明 $Ax\in\ell^q$, 且映射 $x\mapsto Ax$ 是从 $\ell^p$ (强) 到
	      $\ell^q$ (强) 的连续映射.
	\item 证明: 若 $(x^n)$ 为 $\ell^p$ 中的序列, 使 $x^n\weakto x$ 在
	      $\sigma(\ell^p,\ell^{p'})$ 中弱收敛, 则 $Ax^n\weakto Ax$ 在
	      $\sigma(\ell^q,\ell^{q'})$ 中弱收敛.
	\item 推出 $A$ 是从配上 $\sigma(E,E^\star)$ 的 $B_E$ 到配上
	      $\sigma(F,F^\star)$ 的 $F$ 的连续映射.
\end{enumerate}
\end{exercise}

\begin{exercise}
\label{exer:3.20}
设 $E$ 为 Banach 空间.

\begin{enumerate}[label=\arabic*., leftmargin=2.2em, itemsep=0.25em]
	\item 证明存在紧拓扑空间 $K$ 以及从 $E$ 到配上通常范数的 $C(K)$ 的等距
	      同构.

	      [提示: 取 $K=B_{E^\star}$ 配上 $\sigma(E^\star,E)$.]
	\item 假设 $E$ 可分, 证明存在从 $E$ 到 $\ell^\infty$ 的等距同构.
\end{enumerate}
\end{exercise}

\begin{exercise}
\label{exer:3.21}
设 $E$ 为可分 Banach 空间, $(f_n)$ 为 $E^\star$ 中的有界序列.直接证明——不利用
$E^\star$ 的可度量化性——存在一个子列 $(f_{n_k})$ 在 $\sigma(E^\star,E)$ 中
收敛.

[提示: 利用对角线过程.]
\end{exercise}

\begin{exercise}
\label{exer:3.22}
设 $E$ 为无穷维 Banach 空间, 满足下列假设之\emph{一}:
\begin{enumerate}[label=(\alph*), leftmargin=2.0em, itemsep=0.15em]
	\item $E^\star$ 可分,
	\item $E$ 自反.
\end{enumerate}
证明存在 $E$ 中的序列 $(x_n)$, 使得
\begin{equation*}
\norm{x_n}=1\quad\forall n
\qquad\text{且}\qquad
x_n\weakto0\ \text{弱收敛于}\ \sigma(E,E^\star).
\end{equation*}
\end{exercise}

\begin{exercise}
\label{exer:3.23}
若 $E$ 自反, 定理  \ref{thm:2.16} 的证明会变得容易得多.特别地, 找出
(b) $\Rightarrow$ (a) 的一个简单证明.
\end{exercise}

\begin{exercise}
\label{exer:3.24}
本习题的目的是勾勒定理  \ref{thm:3.29} 证明的一部分, 即: 若 Banach 空间 $E$
使得 $B_E$ 关于 $\sigma(E,E^\star)$ 可度量化, 则 $E^\star$ 可分.设 $d(x,y)$
为 $B_E$ 上的度量, 它在 $B_E$ 上诱导出与 $\sigma(E,E^\star)$ 相同的拓扑.置
\begin{equation*}
U_n=\Bigl\{x\in B_E;\ d(x,0)<\frac1n\Bigr\}.
\end{equation*}
设 $V_n$ 为 $0$ 关于 $\sigma(E,E^\star)$ 的邻域, 使 $V_n\subset U_n$.不妨设
$V_n$ 具有形式
\begin{equation*}
V_n=\{x\in E;\ \abs{\dual{f}{x}}<\varepsilon_n\quad\forall f\in\Phi_n\},
\end{equation*}
其中 $\varepsilon_n>0$, $\Phi_n\subset E^\star$ 是某个有限子集.置
$D=\bigcup_{n=1}^\infty\Phi_n$, 并设 $F$ 为 $D$ 张成的向量空间.我们断言
$F$ 在 $E^\star$ 中关于强拓扑稠密.反设 $\overline{F}\ne E^\star$.

\begin{enumerate}[label=\arabic*., leftmargin=2.2em, itemsep=0.25em]
	\item 证明存在某个 $\xi\in E^{\star\star}$ 与某个 $f_0\in E^\star$, 使得
	      \begin{equation*}
	      \dual{\xi}{f_0}>1,\qquad
	      \dual{\xi}{f}=0\quad\forall f\in F,\qquad
	      \text{且}\qquad \norm{\xi}=1.
	      \end{equation*}
	\item 置
	      \begin{equation*}
	      W=\Bigl\{x\in B_E;\ \abs{\dual{f_0}{x}}<\frac12\Bigr\}.
	      \end{equation*}
	      证明存在某个整数 $n_0\ge1$, 使 $V_{n_0}\subset W$.
	\item 证明存在 $x_1\in B_E$, 使得
	      \begin{equation*}
	      \begin{cases}
	      \abs{\dual{f}{x_1}-\dual{\xi}{f}}<\varepsilon_{n_0}
	      &\forall f\in\Phi_{n_0},\\[2pt]
	      \abs{\dual{f_0}{x_1}-\dual{\xi}{f_0}}<\dfrac12.
	      \end{cases}
	      \end{equation*}
	\item 推出 $x_1\in V_{n_0}$ 且 $\dual{f_0}{x_1}>\frac12$.
	\item 得出结论.
\end{enumerate}
\end{exercise}

\begin{exercise}
\label{exer:3.25}
设 $K$ 为紧度量空间, 且不是有限集.证明 $C(K)$ 不自反.

[提示: 设 $(a_n)$ 为 $K$ 中的序列, 使 $a_n\to a$ 且 $a_n\ne a\ \forall n$.
考虑线性泛函 $f(u)=\sum_{n=1}^\infty\frac1{2^n}u(a_n)$, $u\in C(K)$, 并按
习题 1.3 与 1.4 的方式论证.]
\end{exercise}

\begin{exercise}
\label{exer:3.26}
设 $F$ 为可分 Banach 空间, $(a_n)$ 为 $B_F$ 的稠密子集.考虑线性算子
$T:\ell^1\to F$,
\begin{equation*}
Tx=\sum_{i=1}^\infty x_ia_i
\qquad\text{其中}\quad x=(x_1,x_2,\dots,x_n,\dots)\in\ell^1.
\end{equation*}

\begin{enumerate}[label=\arabic*., leftmargin=2.2em, itemsep=0.25em]
	\item 证明 $T$ 有界且满射.
\end{enumerate}
下文再假设 $F$ 是无穷维的, 且 $F^\star$ 可分.

\begin{enumerate}[label=\arabic*., leftmargin=2.2em, itemsep=0.25em, start=2]
	\item 证明 $T$ 没有右逆.

	      [提示: 利用习题  \ref{exer:3.22} 与问题 8 的结果.]
	\item 推出 $N(T)$ 在 $\ell^1$ 中没有补空间.
	\item 确定 $E^\star$.
\end{enumerate}
\end{exercise}

\begin{exercise}
\label{exer:3.27}
设 $E$ 为可分 Banach 空间, 范数为 $\norm{\cdot}$.$E^\star$ 上的对偶范数也记作
$\norm{\cdot}$.本习题的目的是在 $E$ 上构造一个等价范数, 它严格凸, 且其对偶范数
也严格凸.

设 $(a_n)\subset B_E$ 为 $B_E$ 关于强拓扑的稠密子集.设
$(b_n)\subset B_{E^\star}$ 为 $B_{E^\star}$ 的可数子集, 它在 $B_{E^\star}$ 中
关于\wstar{}拓扑 $\sigma(E^\star,E)$ 稠密.这样的集合为何存在?

给定 $f\in E^\star$, 置
\begin{equation*}
\|f\|_1=\Bigl\{\|f\|^2+\sum_{n=1}^\infty\frac{1}{2^n}\abs{\dual{f}{a_n}}^2
\Bigr\}^{1/2}.
\end{equation*}
\begin{enumerate}[label=\arabic*., leftmargin=2.2em, itemsep=0.25em]
	\item 证明 $\|\cdot\|_1$ 是与 $\|\cdot\|$ 等价的范数.
	\item 证明 $\|\cdot\|_1$ 严格凸.

	      [提示: 利用习题 1.26.]
\end{enumerate}
给定 $x\in E$, 置
\begin{equation*}
\|x\|_2=\Bigl\{\|x\|_1^2+\sum_{n=1}^\infty\frac{1}{2^n}
\abs{\dual{b_n}{x}}^2\Bigr\}^{1/2},
\end{equation*}
其中 $\|x\|_1=\sup_{\|f\|_1\le1}\dual{f}{x}$.

\begin{enumerate}[label=\arabic*., leftmargin=2.2em, itemsep=0.25em, start=3]
	\item 证明 $\|\cdot\|_2$ 是与 $\|\cdot\|$ 等价的严格凸范数.
	\item 证明 $\|\cdot\|_2$ 的对偶范数也严格凸.

	      [提示: 利用习题 1.23 第 3 问的结果.]
	\item 基于问题 4 的结果找出另一种方法.
\end{enumerate}
\end{exercise}

\begin{exercise}
\label{exer:3.28}
设 $E$ 为一致凸 Banach 空间.用 $F$ 记从 $E$ 到 $E^\star$ 的 (多值) 对偶映射,
见推论  \ref{cor:1.3} 之后的注记 2 以及习题 1.1.

证明: 对每个 $f\in E^\star$, 存在唯一的 $x\in E$, 使得 $f\in Fx$.
\end{exercise}

\begin{exercise}
\label{exer:3.29}
设 $E$ 为一致凸 Banach 空间.

\begin{enumerate}[label=\arabic*., leftmargin=2.2em, itemsep=0.25em]
	\item 证明 $\forall M>0$, $\forall\varepsilon>0$, $\exists\delta>0$, 使得
	      \begin{equation*}
	      \Bigl\|\frac{x+y}{2}\Bigr\|^2\le\frac12\norm{x}^2+\frac12\norm{y}^2-\delta
	      \end{equation*}
	      \begin{equation*}
	      \forall x,y\in E\ \text{满足}\ \norm{x}\le M,\ \norm{y}\le M
	      \ \text{且}\ \norm{x-y}>\varepsilon.
	      \end{equation*}
	      [提示: 用反证法.]
	\item 当 $\|\cdot\|^2$ 换成 $\|\cdot\|^p$ ($1<p<\infty$) 时, 问同样的问题.
\end{enumerate}
\end{exercise}

\begin{exercise}
\label{exer:3.30}
设 $E$ 为范数 $\norm{\cdot}$ 的 Banach 空间.假设 $E$ 上存在一个等价的范数,
记作 $\abs{\cdot}$, 它一致凸.

证明: 对任意 $k>1$, 存在 $E$ 上的一致凸范数 $|\!|\!|\cdot|\!|\!|$, 使得
\begin{equation*}
\norm{x}\le|\!|\!|x|\!|\!|\le k\norm{x}\quad\forall x\in E.
\end{equation*}
[提示: 置 $|\!|\!|x|\!|\!|^2=\norm{x}^2+\alpha\abs{x}^2$, 其中 $\alpha>0$ 足够
小, 并利用习题  \ref{exer:3.29}.]

例: $E=\R^n$.
\end{exercise}

\begin{exercise}
\label{exer:3.31}
设 $E$ 为一致凸 Banach 空间.

\begin{enumerate}[label=\arabic*., leftmargin=2.2em, itemsep=0.25em]
	\item 证明
	      \begin{equation*}
	      \forall\varepsilon>0,\quad
	      \forall\alpha\in\Bigl(0,\frac12\Bigr),\quad
	      \exists\delta>0
	      \end{equation*}
	      使得
	      \begin{equation*}
	      \norm{tx+(1-t)y}\le1-\delta
	      \end{equation*}
	      \begin{equation*}
	      \forall t\in[\alpha,1-\alpha],\quad\forall x,y\in E
	      \ \text{满足}\ \norm{x}\le1,\ \norm{y}\le1\ \text{且}\
	      \norm{x-y}\ge\varepsilon.
	      \end{equation*}
	      [提示: 若 $\alpha\le t\le\frac12$, 则写
	      $tx+(1-t)y=\frac12(y+z)$.]
	\item 推出 $E$ 严格凸.
\end{enumerate}
\end{exercise}

\begin{exercise}[一致凸 Banach 空间中闭凸集上的投影.]
\label{exer:3.32}
设 $E$ 为一致凸 Banach 空间, $C\subset E$ 为非空闭凸集.

\begin{enumerate}[label=\arabic*., leftmargin=2.2em, itemsep=0.25em]
	\item 证明对每个 $x\in E$,
	      \begin{equation*}
	      \inf_{y\in C}\norm{x-y}
	      \end{equation*}
	      由 $C$ 中某个唯一点达到, 记作 $P_Cx$.
	\item 证明 $C$ 中每个极小化序列 $(y_n)$ 都强收敛到 $P_Cx$.
	\item 证明映射 $x\mapsto P_Cx$ 是从强拓扑下的 $E$ 到强拓扑下的 $E$ 的连续
	      映射.
	\item 更精确地, 证明 $P_C$ 在 $E$ 的有界子集上一致连续.

	      [提示: 利用习题  \ref{exer:3.29}.]
\end{enumerate}
设 $\varphi:E\to(-\infty,+\infty]$ 为凸的 l.s.c. 函数,
$\varphi\not\equiv+\infty$.

\begin{enumerate}[label=\arabic*., leftmargin=2.2em, itemsep=0.25em, start=5]
	\item 证明对每个 $x\in E$ 与每个整数 $n\ge1$,
	      \begin{equation*}
	      \inf_{y\in E}\bigl\{n\norm{x-y}^2+\varphi(y)\bigr\}
	      \end{equation*}
	      由某个唯一点达到, 记作 $y_n$.
	\item 证明 $y_n\to P_Cx$ ($n\to\infty$), 其中
	      $C=\overline{D(\varphi)}$.
\end{enumerate}
\end{exercise}

% chapters/ch04.tex
% 第~
% \ref{chap:4}~章 $L^p$ 空间
% 源: 原书 PDF 第 98--139 页 (印刷页 89--130)

\chapter{$L^p$ 空间}
\label{chap:4}

设 $(\Omega,\mathcal{M},\mu)$ 表示一个\emph{测度空间}, 即 $\Omega$ 是一个集合, 并且

\begin{enumerate}[label=(\roman*), leftmargin=2.6em, itemsep=0.25em]
	\item $\mathcal{M}$ 是 $\Omega$ 中的一个 \emph{$\sigma$-代数}, 即 $\mathcal{M}$ 是由
	      $\Omega$ 的子集组成的一个族, 使得:
	      \begin{enumerate}[label=(\alph*), leftmargin=2.4em, itemsep=0.15em]
	      	\item $\emptyset\in\mathcal{M}$,
	      	\item $A\in\mathcal{M}\ \Rightarrow\ A^c\in\mathcal{M}$,
	      	\item $\bigcup_{n=1}^\infty A_n\in\mathcal{M}$, 只要 $A_n\in\mathcal{M}\ \forall n$;
	      \end{enumerate}
	\item $\mu$ 是一个\emph{测度}, 即 $\mu:\mathcal{M}\to[0,\infty]$ 满足
	      \begin{enumerate}[label=(\alph*), leftmargin=2.4em, itemsep=0.15em]
	      	\item $\mu(\emptyset)=0$,
	      	\item $\mu\Bigl(\bigcup_{n=1}^\infty A_n\Bigr)=\sum_{n=1}^\infty\mu(A_n)$,
	      	      只要 $(A_n)$ 是 $\mathcal{M}$ 的成员构成的互不相交的可数族.
	      \end{enumerate}
\end{enumerate}

$\mathcal{M}$ 的成员称为\emph{可测集}.有时我们用 $|A|$ 代替 $\mu(A)$.我们还将
假设——尽管这并非本质的——以下条件成立:

\begin{enumerate}[label=(\roman*), leftmargin=2.6em, start=3, itemsep=0.25em]
	\item $\Omega$ 是 \emph{$\sigma$-有限的}, 即存在 $\mathcal{M}$ 中的可数族
	      $(\Omega_n)$, 使得 $\Omega=\bigcup_{n=1}^\infty\Omega_n$ 且
	      $\mu(\Omega_n)<\infty\ \forall n$.
\end{enumerate}

满足 $\mu(E)=0$ 的集合 $E\in\mathcal{M}$ 称为\emph{零测集}.若一个性质在 $\Omega$
上除一个零测集外处处成立, 则说该性质\emph{几乎处处}成立 (或对\emph{几乎所有的}
$x\in\Omega$ 成立), 简记为 a.e..

我们假定读者熟悉可测函数与可积函数 $f:\Omega\to\R$ 的概念; 参见例如
H.~L.~Royden \bref{Royden}{1}, G.~B.~Folland \bref{Folland}{2}, A.~Knapp \bref{Knapp}{1}, D.~L.~Cohn \bref{Cohn}{1}, A.~Friedman \bref{Friedman}{3}, W.~Rudin \bref{Rudin}{2}, P.~Halmos \bref{Halmos}{1}, E.~Hewitt--K.~Stromberg \bref{HewittStromberg}{1}, R.~Wheeden--A.~Zygmund \bref{WheedenZygmund}{1}, J.~Neveu \bref{Neveu}{1}, P.~Malliavin \bref{Malliavin}{1}, A.~J.~Weir \bref{Weir}{1}, A.~Kolmogorov--S.~Fomin \bref{KolmogorovFomin}{1}, I.~Fonseca--G.~Leoni \bref{FonsecaLeoni}{1}.我们用 $L^1(\Omega,\mu)$, 或简记为
$L^1(\Omega)$ (或干脆记作 $L^1$), 表示由 $\Omega$ 到 $\R$ 的可积函数构成的
空间.

我们常把 $\int_\Omega f\,d\mu$ 写成 $\int f$, 并采用记号
\begin{equation*}
\norm{f}_{L^1}=\norm{f}_1=\int_\Omega\abs{f}\,d\mu=\int\abs{f}.
\end{equation*}
按惯例, 我们把几乎处处相等的两个函数视为同一函数.下面回顾一些基本事实.

\section{人人都必须知道的一些积分结果}
\label{sec:4.1}

\begin{keytheorem}[单调收敛定理, Beppo Levi]
\label{thm:4.1}
设 $(f_n)$ 为 $L^1$ 中的函数序列, 满足
\begin{enumerate}[label=(\alph*), leftmargin=2.2em, itemsep=0.1em]
	\item $f_1\le f_2\le\cdots\le f_n\le f_{n+1}\le\cdots$ 在 $\Omega$ 上几乎处处成立,
	\item $\sup_n\int f_n<\infty$.
\end{enumerate}
则 $f_n(x)$ 在 $\Omega$ 上几乎处处收敛于一个有限极限, 记作 $f(x)$; 函数
$f$ 属于 $L^1$, 且 $\norm{f_n-f}_1\to0$.
\end{keytheorem}

\begin{keytheorem}[控制收敛定理, Lebesgue]
\label{thm:4.2}
设 $(f_n)$ 为 $L^1$ 中的函数序列, 满足
\begin{enumerate}[label=(\alph*), leftmargin=2.2em, itemsep=0.1em]
	\item $f_n(x)\to f(x)$ 在 $\Omega$ 上几乎处处成立,
	\item 存在函数 $g\in L^1$, 使得对一切 $n$ 都有 $\abs{f_n(x)}\le g(x)$
	      在 $\Omega$ 上几乎处处成立.
\end{enumerate}
则 $f\in L^1$, 且 $\norm{f_n-f}_1\to0$.
\end{keytheorem}

\begin{lemma}[Fatou 引理]
\label{lem:4.1}
设 $(f_n)$ 为 $L^1$ 中的函数序列, 满足
\begin{enumerate}[label=(\alph*), leftmargin=2.2em, itemsep=0.1em]
	\item 对一切 $n$ 都有 $f_n\ge0$ 几乎处处成立,
	\item $\sup_n\int f_n<\infty$.
\end{enumerate}
对几乎所有的 $x\in\Omega$, 令 $f(x)=\liminf_{n\to\infty}f_n(x)\le+\infty$.则
$f\in L^1$, 且
\begin{equation*}
\int f\le\liminf_{n\to\infty}\int f_n.
\end{equation*}
\end{lemma}

一个基本的例子是 $\Omega=\R^N$, $\mathcal{M}$ 由 Lebesgue 可测集组成,
$\mu$ 为 $\R^N$ 上的 Lebesgue 测度.

\begin{notation}
我们用 $C_c(\R^N)$ 表示 $\R^N$ 上具有紧支集的连续函数构成的空间, 即
\begin{equation*}
C_c(\R^N)=\{f\in C(\R^N);\ f(x)=0\ \ \forall x\in\R^N\setminus K,\
\text{其中}\ K\ \text{紧}\}.
\end{equation*}
\end{notation}

\begin{theorem}[稠密性]
\label{thm:4.3}
空间 $C_c(\R^N)$ 在 $L^1(\R^N)$ 中稠密; 即
\begin{equation*}
\forall f\in L^1(\R^N)\ \ \forall\varepsilon>0\ \ \exists f_1\in C_c(\R^N)\
\text{使得}\ \norm{f-f_1}_1\le\varepsilon.
\end{equation*}
\end{theorem}

设 $(\Omega_1,\mathcal{M}_1,\mu_1)$ 与 $(\Omega_2,\mathcal{M}_2,\mu_2)$ 为两个
$\sigma$-有限的测度空间.按标准方式可以在 Cartesian 积
$\Omega=\Omega_1\times\Omega_2$ 上定义测度空间的结构
$(\Omega,\mathcal{M},\mu)$.

\begin{theorem}[Tonelli]
\label{thm:4.4}
设 $F(x,y):\Omega_1\times\Omega_2\to\R$ 为可测函数, 满足
\begin{equation*}
\text{(a)}\ \ \int_{\Omega_2}\abs{F(x,y)}\,d\mu_2<\infty\ \ \text{对几乎所有的}
\ x\in\Omega_1,
\end{equation*}
以及
\begin{equation*}
\text{(b)}\ \ \int_{\Omega_1}d\mu_1\int_{\Omega_2}\abs{F(x,y)}\,d\mu_2<\infty.
\end{equation*}
则 $F\in L^1(\Omega_1\times\Omega_2)$.
\end{theorem}

\begin{theorem}[Fubini]
\label{thm:4.5}
假设 $F\in L^1(\Omega_1\times\Omega_2)$.则对几乎所有的 $x\in\Omega_1$,
$F(x,y)\in L^1_y(\Omega_2)$, 且 $\int_{\Omega_2}F(x,y)\,d\mu_2\in
L^1_x(\Omega_1)$.同样, 对几乎所有的 $y\in\Omega_2$,
$F(x,y)\in L^1_x(\Omega_1)$, 且 $\int_{\Omega_1}F(x,y)\,d\mu_1\in
L^1_y(\Omega_2)$.此外还有
\begin{equation*}
\int_{\Omega_1}d\mu_1\int_{\Omega_2}F(x,y)\,d\mu_2
=\int_{\Omega_2}d\mu_2\int_{\Omega_1}F(x,y)\,d\mu_1
=\iint_{\Omega_1\times\Omega_2}F(x,y)\,d\mu_1d\mu_2.
\end{equation*}
\end{theorem}

\section{$L^p$ 空间的定义与初等性质}
\label{sec:4.2}

\begin{definition}
设 $p\in\R$, $1<p<\infty$; 令
\begin{equation*}
L^p(\Omega)=\bigl\{f:\Omega\to\R;\ f\ \text{可测且}\ \abs{f}^p\in L^1(\Omega)\bigr\}
\end{equation*}
并置
\begin{equation*}
\norm{f}_{L^p}=\norm{f}_p=\Bigl(\int_\Omega\abs{f(x)}^p\,d\mu\Bigr)^{1/p}.
\end{equation*}
\end{definition}

稍后我们将验证 $\norm{\cdot}_p$ 是一个范数.

\begin{definition}
令
\begin{equation*}
L^\infty(\Omega)=\Bigl\{f:\Omega\to\R\ \Big|\ f\ \text{可测, 且存在常数}\ C\
\text{使得}\ \abs{f(x)}\le C\ \text{在}\ \Omega\ \text{上几乎处处成立}\Bigr\},
\end{equation*}
并置
\begin{equation*}
\norm{f}_{L^\infty}=\norm{f}_\infty=\inf\{C;\ \abs{f(x)}\le C\
\text{在}\ \Omega\ \text{上几乎处处成立}\}.
\end{equation*}
\end{definition}

下面的注记说明 $\norm{\cdot}_\infty$ 是一个范数:

\begin{remark}
\label{rem:4.1}
若 $f\in L^\infty$, 则
\begin{equation*}
\abs{f(x)}\le\norm{f}_\infty\quad\text{在}\ \Omega\ \text{上几乎处处成立}.
\end{equation*}
事实上, 存在序列 $C_n$, 使 $C_n\to\norm{f}_\infty$, 且对每个 $n$ 都有
$\abs{f(x)}\le C_n$ 在 $\Omega$ 上几乎处处成立.于是存在零测集 $E_n$, 使得
$\abs{f(x)}\le C_n$ 对一切 $x\in\Omega\setminus E_n$ 成立.令
$E=\bigcup_{n=1}^\infty E_n$, 则 $|E|=0$, 且
\begin{equation*}
\abs{f(x)}\le C_n\quad\forall n,\ \ \forall x\in\Omega\setminus E;
\end{equation*}
由此推出 $\abs{f(x)}\le\norm{f}_\infty\ \forall x\in\Omega\setminus E$.
\end{remark}

\begin{notation}
设 $1\le p\le\infty$; 我们用 $p'$ 表示\emph{共轭指数},
\begin{equation*}
\boxed{\frac1p+\frac1{p'}=1.}
\end{equation*}
\end{notation}

\index{Hlderbudengshi@Hölder 不等式 (Holder inequality)}
\begin{keytheorem}[Hölder 不等式]
\label{thm:4.6}
假设 $f\in L^p$, $g\in L^{p'}$, 其中 $1\le p\le\infty$.则 $fg\in L^1$, 且
\begin{equation}
\boxed{\int\abs{fg}\le\norm{f}_p\norm{g}_{p'}.}
\label{eq:4.1}
\end{equation}
\end{keytheorem}

\begin{proof}
若 $p=1$ 或 $p=\infty$, 结论显然; 因此设 $1<p<\infty$.我们回顾 Young 不等式:\footnote{有时方便使用
的形式是 $ab\le\varepsilon a^p+C_\varepsilon b^{p'}$, 其中
$C_\varepsilon=\varepsilon^{-1/(p-1)}$.}
\begin{equation}
\boxed{ab\le\frac1pa^p+\frac1{p'}b^{p'}\quad\forall a\ge0,\ \ \forall b\ge0.}
\label{eq:4.2}
\end{equation}
不等式 \eqref{eq:4.2} 是函数 $\log$ 在 $(0,\infty)$ 上凹这一事实的直接推论:
\begin{equation*}
\log\Bigl(\frac1pa^p+\frac1{p'}b^{p'}\Bigr)\ge\frac1p\log a^p+\frac1{p'}\log b^{p'}
=\log ab.
\end{equation*}
我们有
\begin{equation*}
\abs{f(x)g(x)}\le\frac1p\abs{f(x)}^p+\frac1{p'}\abs{g(x)}^{p'}\quad
\text{a.e.}\ x\in\Omega.
\end{equation*}
由此推出 $fg\in L^1$, 且
\begin{equation}
\int\abs{fg}\le\frac1p\norm{f}_p^p+\frac1{p'}\norm{g}_{p'}^{p'}.
\label{eq:4.3}
\end{equation}
在 \eqref{eq:4.3} 中把 $f$ 换成 $\lambda f$ ($\lambda>0$), 得到
\begin{equation}
\int\abs{fg}\le\frac{\lambda^{p-1}}{p}\norm{f}_p^p
+\frac1{\lambda^{p'}p'}\norm{g}_{p'}^{p'}.
\label{eq:4.4}
\end{equation}
选取 $\lambda=\norm{f}_p^{-1}\norm{g}_{p'}^{p'/p}$ (使 \eqref{eq:4.4} 的右端达到
最小), 即得 \eqref{eq:4.1}.
\end{proof}

\begin{remark}
\label{rem:4.2}
记住 Hölder 不等式的下述推广是有益的: 假设 $f_1,f_2,\dots,f_k$ 是满足
\begin{equation*}
f_i\in L^{p_i},\ \ 1\le i\le k,\quad\text{其中}\
\frac1p=\frac1{p_1}+\frac1{p_2}+\cdots+\frac1{p_k}\le1
\end{equation*}
的函数.则乘积 $f=f_1f_2\cdots f_k$ 属于 $L^p$, 且
\begin{equation*}
\norm{f}_p\le\norm{f_1}_{p_1}\norm{f_2}_{p_2}\cdots\norm{f_k}_{p_k}.
\end{equation*}
特别地, 若 $f\in L^p\cap L^q$, 其中 $1\le p\le q\le\infty$, 则对一切满足
$p\le r\le q$ 的 $r$ 都有 $f\in L^r$, 且下述``插值不等式''成立:
\begin{equation*}
\norm{f}_r\le\norm{f}_p^\alpha\norm{f}_q^{1-\alpha},\quad
\text{其中}\ \frac1r=\frac\alpha p+\frac{1-\alpha}q,\ 0\le\alpha\le1;
\end{equation*}
参见习题 4.4.
\end{remark}

\begin{theorem}
\label{thm:4.7}
对任意的 $p$, $1\le p\le\infty$, $L^p$ 是向量空间, 且 $\norm{\cdot}_p$ 是一个范数.
\end{theorem}

\begin{proof}
$p=1$ 与 $p=\infty$ 的情形是清楚的.因此设 $1<p<\infty$, 并设
$f,g\in L^p$.我们有
\begin{equation*}
\abs{f(x)+g(x)}^p\le\bigl(\abs{f(x)}+\abs{g(x)}\bigr)^p
\le2^p\bigl(\abs{f(x)}^p+\abs{g(x)}^p\bigr).
\end{equation*}
从而 $f+g\in L^p$.另一方面,
\begin{equation*}
\norm{f+g}_p^p=\int\abs{f+g}^{p-1}\abs{f+g}
\le\int\abs{f+g}^{p-1}\abs{f}+\int\abs{f+g}^{p-1}\abs{g}.
\end{equation*}
但 $\abs{f+g}^{p-1}\in L^{p'}$, 由 Hölder 不等式得到
\begin{equation*}
\norm{f+g}_p^p\le\norm{f+g}_p^{p-1}\bigl(\norm{f}_p+\norm{g}_p\bigr),
\end{equation*}
即 $\norm{f+g}_p\le\norm{f}_p+\norm{g}_p$.
\end{proof}

\index{FischerRieszdingli@Fischer--Riesz 定理 (Fischer--Riesz theorem)}
\begin{keytheorem}[Fischer--Riesz]
\label{thm:4.8}
对任意的 $p$, $1\le p\le\infty$, $L^p$ 是 Banach 空间.
\end{keytheorem}

\begin{proof}
我们区分 $p=\infty$ 与 $1\le p<\infty$ 两种情形.

\emph{情形 1: $p=\infty$.}设 $(f_n)$ 是 $L^\infty$ 中的 Cauchy 序列.给定整数
$k\ge1$, 存在整数 $N_k$, 使得 $\norm{f_m-f_n}_\infty\le\frac1k$ 对一切
$m,n\ge N_k$ 成立.于是存在零测集 $E_k$, 使得
\begin{equation}
\abs{f_m(x)-f_n(x)}\le\frac1k\quad\forall x\in\Omega\setminus E_k,\ \
\forall m,n\ge N_k.
\label{eq:4.5}
\end{equation}
再令 $E=\bigcup_kE_k$——于是 $E$ 是零测集——我们看出对一切
$x\in\Omega\setminus E$, 序列 $f_n(x)$ 是 Cauchy 序列 (在 $\R$ 中).因此
$f_n(x)\to f(x)$ 对一切 $x\in\Omega\setminus E$ 成立.在 \eqref{eq:4.5} 中令
$m\to\infty$ 取极限, 得到
\begin{equation*}
\abs{f(x)-f_n(x)}\le\frac1k\quad\text{对一切}\ x\in\Omega\setminus E,\ \
\forall n\ge N_k.
\end{equation*}
我们由此断定 $f\in L^\infty$, 且 $\norm{f-f_n}_\infty\le\frac1k$
$\forall n\ge N_k$; 因此 $f_n\to f$ 在 $L^\infty$ 中.

\emph{情形 2: $1\le p<\infty$.}设 $(f_n)$ 是 $L^p$ 中的 Cauchy 序列.为得出结论,
只须证明某个子列在 $L^p$ 中收敛.

我们抽取子列 $(f_{n_k})$, 使得
\begin{equation*}
\norm{f_{n_{k+1}}-f_{n_k}}_p\le\frac1{2^k}\quad\forall k\ge1.
\end{equation*}
[做法如下: 取 $n_1$ 使 $\norm{f_m-f_n}_p\le\frac12$ 对一切 $m,n\ge n_1$ 成立;
再取 $n_2\ge n_1$ 使 $\norm{f_m-f_n}_p\le\frac1{2^2}$ 对一切 $m,n\ge n_2$
成立, 等等.] 我们断言 $f_{n_k}$ 在 $L^p$ 中收敛.为简化记号, 把 $f_{n_k}$
写成 $f_k$, 于是有
\begin{equation}
\norm{f_{k+1}-f_k}_p\le\frac1{2^k}\quad\forall k\ge1.
\label{eq:4.6}
\end{equation}
令
\begin{equation*}
g_n(x)=\sum_{k=1}^n\abs{f_{k+1}(x)-f_k(x)},
\end{equation*}
于是
\begin{equation*}
\norm{g_n}_p\le1.
\end{equation*}
由单调收敛定理, $g_n(x)$ 在 $\Omega$ 上几乎处处趋于一个有限极限, 记作 $g(x)$,
且 $g\in L^p$.另一方面, 对 $m\ge n\ge2$ 有
\begin{equation*}
\abs{f_m(x)-f_n(x)}\le\abs{f_m(x)-f_{m-1}(x)}+\cdots+\abs{f_{n+1}(x)-f_n(x)}
\le g(x)-g_{n-1}(x).
\end{equation*}
由此推出 $f_n(x)$ 在 $\Omega$ 上几乎处处是 Cauchy 序列, 且收敛于一个有限极限,
记作 $f(x)$.在 $\Omega$ 上几乎处处有
\begin{equation}
\abs{f(x)-f_n(x)}\le g(x)\quad\text{对}\ n\ge2,
\label{eq:4.7}
\end{equation}
特别地 $f\in L^p$.最后, 由控制收敛定理推出 $\norm{f_n-f}_p\to0$, 因为
$\abs{f_n(x)-f(x)}^p\to0$ 几乎处处成立, 且
$\abs{f_n-f}^p\le g^p\in L^1$.
\end{proof}

\begin{theorem}
\label{thm:4.9}
设 $(f_n)$ 为 $L^p$ 中的序列, $f\in L^p$ 满足 $\norm{f_n-f}_p\to0$.
则存在子列 $(f_{n_k})$ 与函数 $h\in L^p$, 使得
\begin{enumerate}[label=(\alph*), leftmargin=2.2em, itemsep=0.1em]
	\item $f_{n_k}(x)\to f(x)$ 在 $\Omega$ 上几乎处处成立,
	\item $\abs{f_{n_k}(x)}\le h(x)\ \forall k$, 在 $\Omega$ 上几乎处处成立.
\end{enumerate}
\end{theorem}

\begin{proof}
当 $p=\infty$ 时结论显然.因此设 $1\le p<\infty$.由于 $(f_n)$ 是 Cauchy 序列,
我们可以回到定理  \ref{thm:4.8} 的证明, 考虑一个子列 $(f_{n_k})$——记作
$(f_k)$——它满足 \eqref{eq:4.6}, 并且 $f_k(x)$ 几乎处处收敛于一个极限\footnote{先验地,
应当把 $f$ 与 $f^\star$ 区分开: 按假设 $f_n\to f$ 在 $L^p$ 中, 而另一方面
$f_{n_k}(x)\to f^\star(x)$ 几乎处处成立.} $f^\star(x)$, 其中 $f^\star\in L^p$.
此外, 由 \eqref{eq:4.7} 有 $\abs{f^\star(x)-f_k(x)}\le g(x)\ \forall k$, 在
$\Omega$ 上几乎处处成立, 其中 $g\in L^p$.由控制收敛定理知 $f_k\to f^\star$ 在
$L^p$ 中, 从而 $f=f^\star$ 几乎处处成立.此外还有
$\abs{f_k(x)}\le\abs{f^\star(x)}+g(x)$, 结论随之成立.
\end{proof}

\section{自反性. 可分性. $L^p$ 的对偶}
\label{sec:4.3}

我们分别考虑以下三种情形:
\begin{enumerate}[label=(\Alph*), leftmargin=2.6em, itemsep=0.1em]
	\item $1<p<\infty$,
	\item $p=1$,
	\item $p=\infty$.
\end{enumerate}

\subsection*{A. $1<p<\infty$ 时 $L^p(\Omega)$ 的研究}

这种情形最为``有利'': $L^p$ 是自反的, 可分的, 且 $L^p$ 的对偶是 $L^{p'}$.

\begin{keytheorem}
\label{thm:4.10}
对任意的 $p$, $1<p<\infty$, $L^p$ 是自反的.
\end{keytheorem}

证明分三步.

\runin{第 1 步 (Clarkson 第一不等式).}
设 $2\le p<\infty$.我们断言
\begin{equation}
\Bigl\|\frac{f+g}2\Bigr\|_p^p+\Bigl\|\frac{f-g}2\Bigr\|_p^p
\le\frac12\bigl(\norm{f}_p^p+\norm{g}_p^p\bigr)\quad\forall f,g\in L^p.
\label{eq:4.8}
\end{equation}

\emph{\eqref{eq:4.8} 的证明.}显然只须证明
\begin{equation*}
\Bigl|\frac{a+b}2\Bigr|^p+\Bigl|\frac{a-b}2\Bigr|^p
\le\frac12\bigl(\abs{a}^p+\abs{b}^p\bigr)\quad\forall a,b\in\R.
\end{equation*}
首先注意
\begin{equation*}
\alpha^p+\beta^p\le(\alpha^2+\beta^2)^{p/2}\quad\forall\alpha,\beta\ge0
\end{equation*}
(由齐性可设 $\beta=1$, 并注意函数 $(x^2+1)^{p/2}-x^p-1$ 在 $[0,\infty)$ 上
递增).取 $\alpha=\bigl|\frac{a+b}2\bigr|$, $\beta=\bigl|\frac{a-b}2\bigr|$,
得到
\begin{equation*}
\begin{aligned}
\Bigl|\frac{a+b}2\Bigr|^p+\Bigl|\frac{a-b}2\Bigr|^p
&\le\Bigl[\Bigl|\frac{a+b}2\Bigr|^2+\Bigl|\frac{a-b}2\Bigr|^2\Bigr]^{p/2}
=\Bigl(\frac{a^2}2+\frac{b^2}2\Bigr)^{p/2}\\
&\le\frac12\bigl(\abs{a}^p+\abs{b}^p\bigr)
\end{aligned}
\end{equation*}
(最后一个不等式来自函数 $x\mapsto\abs{x}^{p/2}$ 的凸性, 因为 $p\ge2$).

\runin{第 2 步: $2\le p<\infty$ 时 $L^p$ 一致凸, 从而自反.}
事实上, 设 $\varepsilon>0$, $f,g\in L^p$, 满足 $\norm{f}_p\le1$, $\norm{g}_p\le1$, 且 $\norm{f-g}_p>\varepsilon$.由 \eqref{eq:4.8} 推出
\begin{equation*}
\Bigl\|\frac{f+g}2\Bigr\|_p^p<1-\Bigl(\frac\varepsilon2\Bigr)^p,
\end{equation*}
从而 $\bigl\|\frac{f+g}2\bigr\|_p<1-\delta$, 其中
$\delta=1-\bigl[1-\bigl(\frac\varepsilon2\bigr)^p\bigr]^{1/p}>0$.因此 $L^p$ 一致凸,
由定理  \ref{thm:3.31} 知它自反.

\runin{第 3 步: $1<p\le2$ 时 $L^p$ 自反.}

\begin{proof}
设 $1<p<\infty$.考虑如下定义的算子 $T:L^p\to(L^{p'})^\star$: 固定
$u\in L^p$; 映射 $f\in L^{p'}\mapsto\int uf$ 是 $L^{p'}$ 上的连续线性泛函,
因此它定义了 $(L^{p'})^\star$ 中的一个元素, 记作 $Tu$, 使得
\begin{equation*}
\dual{Tu}{f}=\int uf\quad\forall f\in L^{p'}.
\end{equation*}
我们断言
\begin{equation}
\norm{Tu}_{(L^{p'})^\star}=\norm{u}_p\quad\forall u\in L^p.
\label{eq:4.9}
\end{equation}
事实上, 由 Hölder 不等式有
\begin{equation*}
\abs{\dual{Tu}{f}}\le\norm{u}_p\norm{f}_{p'}\quad\forall f\in L^{p'},
\end{equation*}
因此 $\norm{Tu}_{(L^{p'})^\star}\le\norm{u}_p$.

另一方面, 令
\begin{equation*}
f_0(x)=\abs{u(x)}^{p-2}u(x)
\end{equation*}
(若 $u(x)=0$ 则 $f_0(x)=0$).显然
\begin{equation*}
f_0\in L^{p'},\quad\norm{f_0}_{p'}=\norm{u}_p^{p-1},
\quad\text{且}\quad\dual{Tu}{f_0}=\norm{u}_p^p;
\end{equation*}
于是
\begin{equation}
\norm{Tu}_{(L^{p'})^\star}\ge\frac{\dual{Tu}{f_0}}{\norm{f_0}_{p'}}
=\norm{u}_p.
\label{eq:4.10}
\end{equation}
因此我们证明了 $T$ 是从 $L^p$ 到 $(L^{p'})^\star$ 的等距映射, 这意味着
$T(L^p)$ 是 $(L^{p'})^\star$ 的闭子空间 (因为 $L^p$ 是 Banach 空间).

现在设 $1<p\le2$.由于 $L^{p'}$ 自反 (由第 2 步), 故 $(L^{p'})^\star$ 也自反
(推论  \ref{cor:3.21}).由命题  \ref{prop:3.20} 断定 $T(L^p)$ 自反, 从而 $L^p$
也自反.
\end{proof}

\begin{remark}
\label{rem:4.3}
事实上, $1<p\le2$ 时 $L^p$ 也是一致凸的.这是 Clarkson 第二不等式的推论, 该不等式
对 $1<p\le2$ 成立:
\begin{equation*}
\Bigl\|\frac{f+g}2\Bigr\|_{p'}^p+\Bigl\|\frac{f-g}2\Bigr\|_{p'}^p
\le\Bigl(\frac12\norm{f}_p^p+\frac12\norm{g}_p^p\Bigr)^{1/(p-1)}
\quad\forall f,g\in L^p.
\end{equation*}
这个不等式比 Clarkson 第一不等式更难证明 (参见例如问题 20 或
E.~Hewitt--K.~Stromberg \bref{HewittStromberg}{1}).显然, 它蕴涵 $1<p\le2$ 时
$L^p$ 一致凸; 另一种方法参见 C.~Morawetz \bref{Morawetz}{1} (习题 4.12) 或
J.~Diestel \bref{Diestel}{1}.
\end{remark}

\begin{keytheorem}[Riesz 表示定理]
\label{thm:4.11}
设 $1<p<\infty$, $\varphi\in(L^p)^\star$.则存在唯一的函数 $u\in L^{p'}$,
使得
\begin{equation*}
\dual{\varphi}{f}=\int uf\quad\forall f\in L^p.
\end{equation*}
此外,
\begin{equation*}
\norm{u}_{p'}=\norm{\varphi}_{(L^p)^\star}.
\end{equation*}
\end{keytheorem}

\begin{remark}
\label{rem:4.4}
定理  \ref{thm:4.11} 非常重要.它说明 $1<p<\infty$ 时 $L^p$ 上的每个连续线性泛函
都可以``具体地''表示为一个积分.映射 $\varphi\mapsto u$ 是一个线性满等距, 它使
我们能够把``抽象''空间 $(L^p)^\star$ 与 $L^{p'}$ 等同起来.

在下文中, 我们将系统地作等同
\begin{equation*}
\boxed{(L^p)^\star=L^{p'}.}
\end{equation*}
\end{remark}

\begin{proof}
考虑由 $\dual{Tu}{f}=\int uf\ \forall u\in L^{p'},\ \forall f\in L^p$ 定义的
算子 $T:L^{p'}\to(L^p)^\star$.定理  \ref{thm:4.10} 证明 (第 3 步) 中所用的论证
表明
\begin{equation*}
\norm{Tu}_{(L^p)^\star}=\norm{u}_{p'}\quad\forall u\in L^{p'}.
\end{equation*}
我们断言 $T$ 是满射.事实上, 令 $E=T(L^{p'})$.由于 $E$ 是闭子空间, 只须证明
$E$ 在 $(L^p)^\star$ 中稠密.设 $h\in(L^p)^{\star\star}$ 满足
$\dual{h}{Tu}=0\ \forall u\in L^{p'}$.由于 $L^p$ 自反, $h\in L^p$, 且满足
$\int uh=0\ \forall u\in L^{p'}$.取 $u=\abs{h}^{p-2}h$, 即得 $h=0$.
\end{proof}

\begin{theorem}
\label{thm:4.12}
对任意的 $p$, $1\le p<\infty$, 空间 $C_c(\R^N)$ 在 $L^p(\R^N)$ 中稠密.
\end{theorem}

在证明定理  \ref{thm:4.12} 之前, 先引入一些记号.

\begin{notation}
\emph{截断}运算 $T_n:\R\to\R$ 定义为
\begin{equation*}
T_nr=\begin{cases}
r & \text{若}\ \abs{r}\le n,\\[2pt]
\dfrac{nr}{\abs{r}} & \text{若}\ \abs{r}>n.
\end{cases}
\end{equation*}
给定集合 $E\subset\Omega$, 定义\emph{特征函数}\footnote{不要与第~
\ref{chap:1}~章引入的示性函数
$I_E$ 混淆.} $\chi_E$ 为
\begin{equation*}
\chi_E(x)=\begin{cases}
1 & \text{若}\ x\in E,\\
0 & \text{若}\ x\in\Omega\setminus E.
\end{cases}
\end{equation*}
\end{notation}

\begin{proof}
首先我们断言: 给定 $f\in L^p(\R^N)$ 与 $\varepsilon>0$, 存在函数
$g\in L^\infty(\R^N)$ 与 $\R^N$ 中的紧集 $K$, 使得 $g$ 在 $K$ 之外为零, 且
\begin{equation}
\norm{f-g}_p<\varepsilon.
\label{eq:4.11}
\end{equation}
事实上, 设 $\chi_n$ 为 $B(0,n)$ 的特征函数, 令 $f_n=\chi_nT_nf$.由控制收敛定理
看出 $\norm{f_n-f}_p\to0$, 因此可以取 $n$ 足够大的 $g=f_n$.其次, 给定
$\delta>0$, 由定理  \ref{thm:4.3} 存在 $g_1\in C_c(\R^N)$, 使得
\begin{equation*}
\norm{g-g_1}_1<\delta.
\end{equation*}
我们总可以假设 $\norm{g_1}_\infty\le\norm{g}_\infty$; 否则用 $T_ng_1$ 代替
$g_1$, 其中 $n=\norm{g}_\infty$.最后有
\begin{equation*}
\norm{g-g_1}_p\le\norm{g-g_1}_1^{1/p}\norm{g-g_1}_\infty^{1-(1/p)}
\le\delta^{1/p}\bigl(2\norm{g}_\infty\bigr)^{1-(1/p)}.
\end{equation*}
选取 $\delta>0$ 足够小, 使
\begin{equation*}
\delta^{1/p}\bigl(2\norm{g}_\infty\bigr)^{1-(1/p)}<\varepsilon,
\end{equation*}
即得结论.
\end{proof}

\begin{definition}
若存在 $\mathcal{M}$ 中成员构成的可数族 $(E_n)$, 使得由 $(E_n)$ 生成的
$\sigma$-代数与 $\mathcal{M}$ 重合 (即 $\mathcal{M}$ 是包含一切 $E_n$ 的最小
$\sigma$-代数), 则称测度空间 $\Omega$ 是\emph{可分的}.
\end{definition}

\begin{example}
测度空间 $\Omega=\R^N$ 是可分的.事实上, 对 $(E_n)$ 可以取任意一族可数个开集,
使得 $\R^N$ 中每个开集都能写成这些 $E_n$ 的并.更一般地, 若 $\Omega$ 是可分度量
空间, 而 $\mathcal{M}$ 由 Borel 集组成 (即 $\mathcal{M}$ 是由 $\Omega$ 中的开集
生成的 $\sigma$-代数), 则 $\Omega$ 是可分测度空间.
\end{example}

\begin{theorem}
\label{thm:4.13}
假设 $\Omega$ 是可分测度空间.则对任意的 $p$, $1\le p<\infty$,
$L^p(\Omega)$ 是可分的.
\end{theorem}

我们只考虑 $\Omega=\R^N$ 的情形, 因为一般情形稍嫌麻烦.注意, 作为推论,
对任意可测集 $\Omega\subset\R^N$, $L^p(\Omega)$ 也是可分的.事实上, 存在从
$L^p(\Omega)$ 到 $L^p(\R^N)$ 的典范等距 (在 $\Omega$ 外作零延拓); 因此可以把
$L^p(\Omega)$ 与 $L^p(\R^N)$ 的一个子空间等同起来, 从而 $L^p(\Omega)$ 可分
(由命题  \ref{prop:3.25}).

\begin{proof}[定理 \ref{thm:4.13}, $\Omega=\R^N$ 的情形]
用 $\mathcal{R}$ 表示 $\R^N$ 中形如
$R=\prod_{k=1}^N(a_k,b_k)$ ($a_k,b_k\in\Q$) 的集合构成的可数族.用 $E$ 表示由
函数族 $(\chi_R)_{R\in\mathcal{R}}$ 在 $\Q$ 上生成的向量空间, 即 $E$ 由函数
$\chi_R$ 的, 系数为有理数的有限线性组合组成, 因此 $E$ 是可数的.

我们断言 $E$ 在 $L^p(\R^N)$ 中稠密.事实上, 给定 $f\in L^p(\R^N)$ 与
$\varepsilon>0$, 存在 $f_1\in C_c(\R^N)$, 使得 $\norm{f-f_1}_p<\varepsilon$.
取 $R\in\mathcal{R}$ 为任一包含 $\operatorname{supp}f_1$ ($f_1$ 的支集) 的方体.
给定 $\delta>0$, 容易构造 $f_2\in E$, 使得 $\norm{f_1-f_2}_\infty<\delta$, 且
$f_2$ 在 $R$ 之外为零: 只须把 $R$ 分成 $\mathcal{R}$ 中的小方体, 使 $f_1$ 在
每个小方体上的振幅 (即 $\sup-\inf$) 小于 $\delta$.因此
$\norm{f_1-f_2}_p\le\norm{f_1-f_2}_\infty\abs{R}^{1/p}<\delta\abs{R}^{1/p}$.
只要选取 $\delta>0$ 使 $\delta\abs{R}^{1/p}<\varepsilon$, 便得到
$\norm{f-f_2}_p<2\varepsilon$.
\end{proof}

\subsection*{B. $L^1(\Omega)$ 的研究}

我们先描述 $L^1(\Omega)$ 的对偶空间.

\begin{keytheorem}[Riesz 表示定理]
\label{thm:4.14}
设 $\varphi\in(L^1)^\star$.则存在唯一的函数 $u\in L^\infty$, 使得
\begin{equation*}
\dual{\varphi}{f}=\int uf\quad\forall f\in L^1.
\end{equation*}
此外,
\begin{equation*}
\norm{u}_\infty=\norm{\varphi}_{(L^1)^\star}.
\end{equation*}
\end{keytheorem}

\begin{keyremark}
\label{rem:4.5}
定理  \ref{thm:4.14} 断言 $L^1$ 上的每个连续线性泛函都可以``具体地''表示为一个
积分.映射 $\varphi\mapsto u$ 是一个线性满等距, 它使我们能够把``抽象''空间
$(L^1)^\star$ 与 $L^\infty$ 等同起来.在下文中, 我们将系统地作等同
\begin{equation*}
\boxed{(L^1)^\star=L^\infty.}
\end{equation*}
\end{keyremark}

\begin{proof}
设 $(\Omega_n)$ 为 $\Omega$ 中的可测集序列, 使得
$\Omega=\bigcup_{n=1}^\infty\Omega_n$ 且 $\abs{\Omega_n}<\infty\ \forall n$.令
$\chi_n=\chi_{\Omega_n}$.

$u$ 的唯一性是显然的.事实上, 设 $u\in L^\infty$ 满足
\begin{equation*}
\int uf=0\quad\forall f\in L^1.
\end{equation*}
取 $f=\chi_n\operatorname{sign}u$ (本书通篇采用约定 $\operatorname{sign}0=0$),
即得 $u=0$ 在 $\Omega_n$ 上几乎处处成立, 从而 $u=0$ 在 $\Omega$ 上几乎处处成立.

现在证明 $u$ 的存在性.首先, 我们构造一个函数 $\theta\in L^2(\Omega)$, 使得
\begin{equation*}
\theta(x)\ge\varepsilon_n>0\quad\forall x\in\Omega_n.
\end{equation*}
这样的函数 $\theta$ 的存在性是清楚的.事实上, 令 $\theta$ 在 $\Omega_1$ 上等于
$\alpha_1$, 在 $\Omega_2\setminus\Omega_1$ 上等于 $\alpha_2$, …, 在
$\Omega_n\setminus\Omega_{n-1}$ 上等于 $\alpha_n$, 等等, 并调整常数
$\alpha_n>0$ 使 $\theta\in L^2$.

映射 $f\in L^2(\Omega)\mapsto\dual{\varphi}{\theta f}$ 是 $L^2(\Omega)$ 上的连续
线性泛函.由定理  \ref{thm:4.11} (取 $p=2$) 存在函数 $v\in L^2(\Omega)$, 使得
\begin{equation}
\dual{\varphi}{\theta f}=\int vf\quad\forall f\in L^2(\Omega).
\label{eq:4.12}
\end{equation}
令 $u(x)=v(x)/\theta(x)$.显然 $u$ 有定义, 因为 $\theta>0$ 在 $\Omega$ 上; 此外
$u$ 可测, 且 $u\chi_n\in L^2(\Omega)$.我们断言 $u$ 具有一切所需的性质.有
\begin{equation}
\dual{\varphi}{\chi_ng}=\int u\chi_ng\quad\forall g\in L^\infty(\Omega)\ \ \forall n.
\label{eq:4.13}
\end{equation}
事实上, 只须在 \eqref{eq:4.12} 中取 $f=\chi_ng/\theta$ (注意 $f\in L^2(\Omega)$,
因为 $f$ 在 $\Omega_n$ 上有界且在 $\Omega_n$ 外为零).

其次, 我们断言 $u\in L^\infty(\Omega)$, 且
\begin{equation}
\norm{u}_\infty\le\norm{\varphi}_{(L^1)^\star}.
\label{eq:4.14}
\end{equation}
固定任一常数 $C>\norm{\varphi}_{(L^1)^\star}$, 并令
\begin{equation*}
A=\{x\in\Omega;\ \abs{u(x)}>C\}.
\end{equation*}
我们来验证 $A$ 是零测集.事实上, 在 \eqref{eq:4.13} 中取
$g=\chi_A\operatorname{sign}u$, 得到
\begin{equation*}
\int_{A\cap\Omega_n}\abs{u}\le\norm{\varphi}_{(L^1)^\star}\abs{A\cap\Omega_n},
\end{equation*}
从而
\begin{equation*}
C\abs{A\cap\Omega_n}\le\norm{\varphi}_{(L^1)^\star}\abs{A\cap\Omega_n}.
\end{equation*}
由此推出 $\abs{A\cap\Omega_n}=0\ \forall n$, 从而 $A$ 是零测集.\eqref{eq:4.14}
的证明至此完成.

最后, 我们断言
\begin{equation}
\dual{\varphi}{h}=\int uh\quad\forall h\in L^1(\Omega).
\label{eq:4.15}
\end{equation}
事实上, 只须在 \eqref{eq:4.13} 中取 $g=T_nh$ ($h$ 的截断), 并注意
$\chi_nT_nh\to h$ 在 $L^1(\Omega)$ 中.

为完成定理  \ref{thm:4.14} 的证明, 只剩验证
$\norm{u}_\infty=\norm{\varphi}_{(L^1)^\star}$.由 \eqref{eq:4.15} 有
\begin{equation*}
\abs{\dual{\varphi}{h}}\le\norm{u}_\infty\norm{h}_1\quad\forall h\in L^1(\Omega),
\end{equation*}
因此 $\norm{\varphi}_{(L^1)^\star}\le\norm{u}_\infty$.再结合 \eqref{eq:4.14}
即得结论.
\end{proof}

\begin{keyremark}
\label{rem:4.6}
除 $\Omega$ 由有限个原子组成这一平凡情形 (此时 $L^1(\Omega)$ 是有限维的) 之外,
空间 $L^1(\Omega)$ 绝不是自反的.事实上, 反设 $L^1(\Omega)$ 自反, 并考虑两种
情形:
\begin{enumerate}[label=(\roman*), leftmargin=2.4em, itemsep=0.1em]
	\item $\forall\varepsilon>0\ \exists\omega\subset\Omega$ 可测, 使
	      $0<\mu(\omega)<\varepsilon$.
	\item $\exists\varepsilon>0$, 使得对一切满足 $\mu(\omega)>0$ 的可测集
	      $\omega\subset\Omega$ 都有 $\mu(\omega)\ge\varepsilon$.
\end{enumerate}
在情形 (i) 中, 存在可测集的递减序列 $(\omega_n)$, 使得
$\mu(\omega_n)>0\ \forall n$ 且 $\mu(\omega_n)\to0$ [先任取序列
$(\omega'_k)$, 使 $0<\mu(\omega'_k)<1/2^k$, 再令
$\omega_n=\bigcup_{k=n}^\infty\omega'_k$].令 $\chi_n=\chi_{\omega_n}$, 并定义
$u_n=\chi_n/\norm{\chi_n}_1$.由于 $\norm{u_n}_1=1$, 存在子列——仍记作 $u_n$
——以及某个 $u\in L^1$, 使得 $u_n\weakto u$ 在弱拓扑 $\sigma(L^1,L^\infty)$ 中
(由定理  \ref{thm:3.18}), 即
\begin{equation}
\int u_n\varphi\to\int u\varphi\quad\forall\varphi\in L^\infty.
\label{eq:4.16}
\end{equation}
另一方面, 对固定的 $j$ 与 $n>j$ 有 $\int u_n\chi_j=1$.令 $n\to\infty$ 取极限,
得到 $\int u\chi_j=1\ \forall j$.最后注意 (由控制收敛定理)
$\int u\chi_j\to0$ 当 $j\to\infty$——矛盾.

在情形 (ii) 中, 空间 $\Omega$ 是纯原子的, 且由可数个互不相同的原子 $(a_n)$ 的并
组成 (除非只有有限个原子!).此时 $L^1(\Omega)$ 同构于 $\ell^1$, 故只须证明
$\ell^1$ 不自反.考虑典范基:
\begin{equation*}
e_n=(0,0,\dots,\underset{(n)}1,0,0\dots).
\end{equation*}
假设 $\ell^1$ 自反, 则存在子列 $(e_{n_k})$ 与某个 $x\in\ell^1$, 使得
$e_{n_k}\weakto x$ 在弱拓扑 $\sigma(\ell^1,\ell^\infty)$ 中, 即
\begin{equation*}
\dual{\varphi}{e_{n_k}}\xrightarrow[k\to\infty]{}\dual{\varphi}{x}
\quad\forall\varphi\in\ell^\infty.
\end{equation*}
取
\begin{equation*}
\varphi=\varphi_j=(0,0,\dots,\underset{(j)}1,1,1,\dots),
\end{equation*}
便得到 $\dual{\varphi_j}{x}=1\ \forall j$.另一方面, 当 $j\to\infty$ 时
$\dual{\varphi_j}{x}\to0$ (因为 $x\in\ell^1$)——矛盾.
\end{keyremark}

\subsection*{C. $L^\infty$ 的研究}

我们已经知道 (定理  \ref{thm:4.14}) $L^\infty=(L^1)^\star$.作为对偶空间,
$L^\infty$ 享有若干良好的性质.特别地, 有:

\begin{enumerate}[label=(\roman*), leftmargin=2.6em, itemsep=0.2em]
	\item 闭单位球 $B_{L^\infty}$ 在\wstar{}拓扑 $\sigma(L^\infty,L^1)$ 中紧
	      (由定理  \ref{thm:3.16}).
	\item 若 $\Omega$ 是 $\R^N$ 中的可测子集, $(f_n)$ 是 $L^\infty(\Omega)$ 中的
	      有界序列, 则存在子列 $(f_{n_k})$ 与某个 $f\in L^\infty(\Omega)$, 使得
	      $f_{n_k}\weakstar f$ 在\wstar{}拓扑 $\sigma(L^\infty,L^1)$ 中 (这是推论
	       \ref{cor:3.30} 与定理  \ref{thm:4.13} 的推论).
\end{enumerate}

然而, 除 $\Omega$ 由有限个原子组成这一平凡情形外, $L^\infty(\Omega)$ 不自反;
否则 $L^1(\Omega)$ 将自反 (由推论  \ref{cor:3.21}), 而我们知道 $L^1$ 不自反
(注记 \ref{rem:4.6}).由此推出 $L^\infty$ 的对偶空间 $(L^\infty)^\star$ 包含
$L^1$ (因为 $L^\infty=(L^1)^\star$), 且 $(L^\infty)^\star$ 严格大于 $L^1$.
换言之, 存在 $L^\infty$ 上的连续线性泛函 $\varphi$, 它们不能表示为
\begin{equation*}
\dual{\varphi}{f}=\int uf\quad\forall f\in L^\infty
\end{equation*}
的形式, 其中 $u\in L^1$.

事实上, 我们来描述这样一个泛函的``具体''例子.令
$\varphi_0:C_c(\R^N)\to\R$ 定义为
\begin{equation*}
\varphi_0(f)=f(0)\quad\text{对}\ f\in C_c(\R^N).
\end{equation*}
显然 $\varphi_0$ 是 $C_c(\R^N)$ 上关于 $\norm{\cdot}_\infty$ 范数的连续线性
泛函.由 Hahn--Banach 定理, 可以把 $\varphi_0$ 延拓为 $L^\infty(\R^N)$ 上的
连续线性泛函 $\varphi$, 并且有
\begin{equation}
\dual{\varphi}{f}=f(0)\quad\forall f\in C_c(\R^N).
\label{eq:4.17}
\end{equation}
我们来验证不存在函数 $u\in L^1(\R^N)$ 使得
\begin{equation}
\dual{\varphi}{f}=\int uf\quad\forall f\in L^\infty(\R^N).
\label{eq:4.18}
\end{equation}
反设这样的函数 $u$ 存在.由 \eqref{eq:4.17} 与 \eqref{eq:4.18} 推出
\begin{equation*}
\int uf=0\quad\forall f\in C_c(\R^N)\ \text{且}\ f(0)=0.
\end{equation*}
对 $\Omega=\R^N\setminus\{0\}$ 应用推论  \ref{cor:4.24}, 看出 $u=0$ 在
$\R^N\setminus\{0\}$ 上几乎处处成立, 从而 $u=0$ 在 $\R^N$ 上几乎处处成立.
由 \eqref{eq:4.18} 断定
\begin{equation*}
\dual{\varphi}{f}=0\quad\forall f\in L^\infty(\R^N),
\end{equation*}
这与 \eqref{eq:4.17} 矛盾.

\begin{starremark}
\label{rem:4.7}
$L^\infty$ 的对偶空间与 $L^1$ 并不重合, 但我们仍可以问: $(L^\infty)^\star$ 究竟
是什么样子? 为此, 把 $L^\infty(\Omega;\CC)$ 看作交换 $C^\star$-代数是有益的
(参见例如 W.~Rudin \bref{Rudin}{1}).由 Gelfand 定理,
$L^\infty(\Omega;\CC)$ 同构且等距于某个紧拓扑空间 $K$ 上连续复值函数构成的空间
$C(K;\CC)$ ($K$ 是代数 $L^\infty$ 的谱; 除 $\Omega$ 由有限个原子组成的情形外,
$K$ 不可度量化).因此 $(L^\infty(\Omega;\CC))^\star$ 可以等同于 $K$ 上的复值
Radon 测度构成的空间, 而 $(L^\infty(\Omega;\R))^\star$ 可以等同于 $K$ 上的实值
Radon 测度构成的空间; 更多细节见本章末评注 3, W.~Rudin \bref{Rudin}{1} 以及
K.~Yosida \bref{Yosida}{1} (第 118 页).
\end{starremark}

\begin{remark}
\label{rem:4.8}
除 $\Omega$ 由有限个原子组成的情形外, 空间 $L^\infty(\Omega)$ 不可分.为证明
这一事实, 使用下述结论是方便的.
\end{remark}

\begin{lemma}
\label{lem:4.2}
设 $E$ 为 Banach 空间.假设存在族 $(O_i)_{i\in I}$, 使得
\begin{enumerate}[label=(\roman*), leftmargin=2.4em, itemsep=0.1em]
	\item 对每个 $i\in I$, $O_i$ 是 $E$ 的非空开子集,
	\item 若 $i\ne j$ 则 $O_i\cap O_j=\emptyset$,
	\item $I$ 不可数.
\end{enumerate}
则 $E$ 不可分.
\end{lemma}

\begin{proof}[引理 \ref{lem:4.2}]
反设 $E$ 可分.用 $(u_n)_{n\in\N}$ 表示 $E$ 中一个可数稠密集.对每个
$i\in I$, 集合 $O_i\cap(u_n)_{n\in\N}\ne\emptyset$, 故可以选取 $n(i)$ 使
$u_{n(i)}\in O_i$.映射 $i\mapsto n(i)$ 是单射; 事实上, 若 $n(i)=n(j)$, 则
$u_{n(i)}=u_{n(j)}\in O_i\cap O_j$, 从而 $i=j$.因此 $I$ 可数——矛盾.
\end{proof}

现在证明 $L^\infty(\Omega)$ 不可分.我们断言存在 $\Omega$ 中可测集构成的不可数族
$(\omega_i)_{i\in I}$, 其中各集合互不相同, 即当 $i\ne j$ 时对称差
$\omega_i\triangle\omega_j$ 的测度为正.然后对由
\begin{equation*}
O_i=\{f\in L^\infty(\Omega);\ \norm{f-\chi_{\omega_i}}_\infty<1/2\}
\end{equation*}
定义的族 $(O_i)_{i\in I}$ 应用引理  \ref{lem:4.2} 即得结论 (注意: 若 $\omega$ 与
$\omega'$ 互不相同, 则 $\norm{\chi_\omega-\chi_{\omega'}}_\infty=1$).当 $\Omega$
是 $\R^N$ 中的开集时, 不可数族 $(\omega_i)$ 的存在性是显然的, 因为可以考虑所有
球 $B(x_0,r)$, 其中 $x_0\in\Omega$ 而 $r>0$ 足够小.

当 $\Omega$ 是一般的测度空间时, 把 $\Omega$ 分成它的原子部分 $\Omega_a$ 与非原子
(= 弥散) 部分 $\Omega_d$; 然后区分两种情形:
\begin{enumerate}[label=(\roman*), leftmargin=2.4em, itemsep=0.1em]
	\item $\Omega_d$ 不是零测集.
	\item $\Omega_d$ 是零测集.
\end{enumerate}
在情形 (i) 中, 对每个实数 $t$, $0<t<\mu(\Omega_d)$, 存在可测集 $\omega$ 使
$\mu(\omega)=t$; 参见例如 P.~Halmos \bref{Halmos}{1}, A.~J.~Weir \bref{Weir}{1}
或 J.~Neveu \bref{Neveu}{1}.这样就得到一个不可数的, 互不相同的可测集族.

在情形 (ii) 中, $\Omega$ 由可数个互不相同的原子 $(a_n)$ 的并组成 (除非 $\Omega$
由有限个原子组成).对任意整数集合 $A\subset\N$, 定义
$\omega_A=\bigcup_{n\in A}a_n$.显然 $(\omega_A)$ 是一个不可数的, 互不相同的
可测集族.

下表总结了 $\Omega$ 为 $\R^N$ 的可测子集时空间 $L^p(\Omega)$ 的主要性质:

\begin{center}
\renewcommand{\arraystretch}{1.35}
\begin{tabular}{|c|c|c|c|}
\hline
 & 自反 & 可分 & 对偶空间 \\
\hline
$L^p$, $1<p<\infty$ & 是 & 是 & $L^{p'}$ \\
\hline
$L^1$ & 否 & 是 & $L^\infty$ \\
\hline
$L^\infty$ & 否 & 否 & 严格大于 $L^1$ \\
\hline
\end{tabular}
\end{center}

\section{卷积与正则化}
\label{sec:4.4}

我们先定义函数 $f\in L^1(\R^N)$ 与函数 $g\in L^p(\R^N)$ 的卷积乘积.

\begin{keytheorem}[Young]
\label{thm:4.15}
设 $f\in L^1(\R^N)$, $g\in L^p(\R^N)$, 其中 $1\le p\le\infty$.则对几乎所有的
$x\in\R^N$, 函数 $y\mapsto f(x-y)g(y)$ 在 $\R^N$ 上可积, 且我们定义
\begin{equation*}
\boxed{(f\star g)(x)=\int_{\R^N}f(x-y)g(y)\,dy.}
\end{equation*}
此外 $f\star g\in L^p(\R^N)$, 且
\begin{equation*}
\boxed{\norm{f\star g}_p\le\norm{f}_1\norm{g}_p.}
\end{equation*}
\end{keytheorem}

\begin{proof}
当 $p=\infty$ 时结论显然.考虑两种情形: (i) $p=1$, (ii) $1<p<\infty$.

\emph{情形 (i): $p=1$.}令 $F(x,y)=f(x-y)g(y)$.对几乎所有的 $y\in\R^N$ 有
\begin{equation*}
\int_{\R^N}\abs{F(x,y)}\,dx=\abs{g(y)}\int_{\R^N}\abs{f(x-y)}\,dx
=\abs{g(y)}\norm{f}_1<\infty,
\end{equation*}
并且
\begin{equation*}
\int_{\R^N}dy\int_{\R^N}\abs{F(x,y)}\,dx=\norm{g}_1\norm{f}_1<\infty.
\end{equation*}
由 Tonelli 定理 (定理  \ref{thm:4.4}) 推出 $F\in L^1(\R^N\times\R^N)$.应用
Fubini 定理 (定理  \ref{thm:4.5}), 看出
\begin{equation*}
\int_{\R^N}\abs{F(x,y)}\,dy<\infty\quad\text{对几乎所有的}\ x\in\R^N,
\end{equation*}
并且
\begin{equation*}
\int_{\R^N}dx\int_{\R^N}\abs{F(x,y)}\,dy
=\int_{\R^N}dy\int_{\R^N}\abs{F(x,y)}\,dx=\norm{f}_1\norm{g}_1.
\end{equation*}
这正是 $p=1$ 时定理  \ref{thm:4.15} 的结论.

\emph{情形 (ii): $1<p<\infty$.}由情形 (i) 知, 对几乎所有的固定 $x\in\R^N$,
函数 $y\mapsto\abs{f(x-y)}\abs{g(y)}^p$ 在 $\R^N$ 上可积, 即
\begin{equation*}
\abs{f(x-y)}^{1/p}\abs{g(y)}\in L^p_y(\R^N).
\end{equation*}
由于 $\abs{f(x-y)}^{1/p'}\in L^{p'}_y(\R^N)$, 由 Hölder 不等式推出
\begin{equation*}
\abs{f(x-y)}\abs{g(y)}=\abs{f(x-y)}^{1/p'}\abs{f(x-y)}^{1/p}\abs{g(y)}
\in L^1_y(\R^N),
\end{equation*}
且
\begin{equation*}
\int_{\R^N}\abs{f(x-y)}\abs{g(y)}\,dy
\le\norm{f}_1^{1/p'}\Bigl(\int_{\R^N}\abs{f(x-y)}\abs{g(y)}^p\,dy\Bigr)^{1/p},
\end{equation*}
即
\begin{equation*}
\abs{(f\star g)(x)}^p\le\norm{f}_1^{p/p'}\bigl(\abs{f}\star\abs{g}^p\bigr)(x).
\end{equation*}
由情形 (i) 断定 $f\star g\in L^p(\R^N)$, 且
\begin{equation*}
\norm{f\star g}_p^p\le\norm{f}_1^{p/p'}\norm{f}_1\norm{g}_p^p,
\end{equation*}
即 $\norm{f\star g}_p\le\norm{f}_1\norm{g}_p$.
\end{proof}

\begin{notation}
给定 $\R^N$ 上的函数 $f$, 令 $\check f(x)=f(-x)$.
\end{notation}

\begin{proposition}
\label{prop:4.16}
设 $f\in L^1(\R^N)$, $g\in L^p(\R^N)$, $h\in L^{p'}(\R^N)$.则
\begin{equation*}
\int_{\R^N}(f\star g)h=\int_{\R^N}g(\check f\star h).
\end{equation*}
\end{proposition}

\begin{proof}
函数 $F(x,y)=f(x-y)g(y)h(x)$ 属于 $L^1(\R^N\times\R^N)$, 因为由定理
 \ref{thm:4.15} 与 Hölder 不等式有
\begin{equation*}
\int\abs{h(x)}\,dx\int\abs{f(x-y)}\abs{g(y)}\,dy<\infty.
\end{equation*}
因此
\begin{equation*}
\begin{aligned}
\int(f\star g)(x)h(x)\,dx
&=\int dx\int F(x,y)\,dy=\int dy\int F(x,y)\,dx\\
&=\int g(y)(\check f\star h)(y)\,dy.
\end{aligned}
\end{equation*}
\end{proof}

\runin{支集与卷积.}
函数 $f$ 的支集的概念是标准的: $\operatorname{supp}f$ 是使 $f$ 为零的最大开集的
补集; 换言之, $\operatorname{supp}f$ 是集合 $\{x;\ f(x)\ne0\}$ 的闭包.在处理
等价类 (例如空间 $L^p$) 时, 这个概念并不合适.我们需要一个内蕴的定义, 即当
$f_1=f_2$ 几乎处处成立时, $\operatorname{supp}f_1$ 与 $\operatorname{supp}f_2$
应当相同 (或只相差一个零测集).读者容易承认, 通常的概念对 $\R$ 上的
$f=\chi_\Q$ 没有意义.在下面的命题中我们引入合适的定义.

\begin{proposition}[以及支集的定义]
\label{prop:4.17}
设 $f:\R^N\to\R$ 为任一函数.考虑 $\R^N$ 上满足下述条件的开集构成的族
$(\omega_i)_{i\in I}$: 对每个 $i\in I$, $f=0$ 在 $\omega_i$ 上几乎处处成立.令
$\omega=\bigcup_{i\in I}\omega_i$.则 $f=0$ 在 $\omega$ 上几乎处处成立.

按定义, $\operatorname{supp}f$ 是 $\omega$ 在 $\R^N$ 中的补集.
\end{proposition}

\begin{remark}
\label{rem:4.9}
(a) 假设 $f_1=f_2$ 在 $\R^N$ 上几乎处处成立; 显然
$\operatorname{supp}f_1=\operatorname{supp}f_2$.因此我们可以谈论 $f\in L^p$ 的
$\operatorname{supp}f$——无须指明在等价类中选取了哪个代表元.

(b) 若 $f$ 是 $\R^N$ 上的连续函数, 则容易验证 $\operatorname{supp}f$ 的新定义与
通常的定义一致.
\end{remark}

\begin{proof}[命题 \ref{prop:4.17}]
由于集合 $I$ 不必可数, 所以 $f=0$ 在 $\omega$ 上几乎处处成立这一点并不显然.
但可以如下归到可数情形.存在 $\R^N$ 中开集构成的可数族 $(O_n)$, 使得 $\R^N$ 中
每个开集都是某些 $O_n$ 的并.写 $\omega_i=\bigcup_{n\in A_i}O_n$, 且
$\omega=\bigcup_{n\in B}O_n$, 其中 $B=\bigcup_{i\in I}A_i$.由于对每个
$n\in B$, $f=0$ 在 $O_n$ 上几乎处处成立, 断定 $f=0$ 在 $\omega$ 上几乎处处
成立.
\end{proof}

\begin{keyproposition}
\label{prop:4.18}
设 $f\in L^1(\R^N)$, $g\in L^p(\R^N)$, 其中 $1\le p\le\infty$.则
\begin{equation*}
\boxed{\operatorname{supp}(f\star g)\subset\overline{\operatorname{supp}f
+\operatorname{supp}g}.}
\end{equation*}
\end{keyproposition}

\begin{proof}
固定 $x\in\R^N$, 使函数 $y\mapsto f(x-y)g(y)$ 可积 (见定理  \ref{thm:4.15}).有
\begin{equation*}
(f\star g)(x)=\int f(x-y)g(y)\,dy
=\int_{(x-\operatorname{supp}f)\cap\operatorname{supp}g}f(x-y)g(y)\,dy.
\end{equation*}
若 $x\notin\operatorname{supp}f+\operatorname{supp}g$, 则
$(x-\operatorname{supp}f)\cap\operatorname{supp}g=\emptyset$, 从而
$(f\star g)(x)=0$.于是 $(f\star g)(x)=0$ 在
$(\operatorname{supp}f+\operatorname{supp}g)^c$ 上几乎处处成立.特别地,
\begin{equation*}
(f\star g)(x)=0\quad\text{a.e. 在}\
\operatorname{Int}\bigl[(\operatorname{supp}f+\operatorname{supp}g)^c\bigr]
\ \text{上},
\end{equation*}
因此
\begin{equation*}
\operatorname{supp}(f\star g)\subset\overline{\operatorname{supp}f
+\operatorname{supp}g}.
\end{equation*}
\end{proof}

\begin{keyremark}
\label{rem:4.10}
若 $f$ 与 $g$ 都具有紧支集, 则 $f\star g$ 也具有紧支集.然而, 若两者中只有一个
具有紧支集, 则 $f\star g$ 未必具有紧支集.
\end{keyremark}

\begin{definition}
设 $\Omega\subset\R^N$ 为开集, $1\le p\le\infty$.若对每个紧集 $K\subset\Omega$
都有 $f\chi_K\in L^p(\Omega)$, 则称函数 $f:\Omega\to\R$ 属于
$L^p_{\mathrm{loc}}(\Omega)$.

注意: 若 $f\in L^p_{\mathrm{loc}}(\Omega)$, 则 $f\in L^1_{\mathrm{loc}}(\Omega)$.
\end{definition}

\begin{proposition}
\label{prop:4.19}
设 $f\in C_c(\R^N)$, $g\in L^1_{\mathrm{loc}}(\R^N)$.则 $(f\star g)(x)$ 对每个
$x\in\R^N$ 都有定义, 并且 $(f\star g)\in C(\R^N)$.
\end{proposition}

\begin{proof}
注意对每个 $x\in\R^N$, 函数 $y\mapsto f(x-y)g(y)$ 在 $\R^N$ 上可积, 因此
$(f\star g)(x)$ 对每个 $x\in\R^N$ 都有定义.

设 $x_n\to x$, 并取 $\R^N$ 中固定的紧集 $K$, 使
$(x_n-\operatorname{supp}f)\subset K\ \forall n$.于是
$f(x_n-y)=0\ \forall n,\ \forall y\notin K$.由 $f$ 的一致连续性推出
\begin{equation*}
\abs{f(x_n-y)-f(x-y)}\le\varepsilon_n\chi_K(y)\quad\forall n,\ \forall y\in\R^N,
\end{equation*}
其中 $\varepsilon_n\to0$.于是
\begin{equation*}
\abs{(f\star g)(x_n)-(f\star g)(x)}\le\varepsilon_n\int_K\abs{g(y)}\,dy
\xrightarrow{\ }0.
\end{equation*}
\end{proof}

\begin{notation}
设 $\Omega\subset\R^N$ 为开集.

$C(\Omega)$ 是 $\Omega$ 上连续函数构成的空间.

$C^k(\Omega)$ 是 $\Omega$ 上 $k$ 次连续可微函数构成的空间 ($k\ge1$ 为整数).

$C^\infty(\Omega)=\bigcap_kC^k(\Omega)$.

$C_c(\Omega)$ 是 $\Omega$ 上在 $\Omega$ 内具有紧支集的连续函数构成的空间, 即在某个
紧集 $K\subset\Omega$ 之外为零的函数构成的空间.

$C_c^k(\Omega)=C^k(\Omega)\cap C_c(\Omega)$.

$C_c^\infty(\Omega)=C^\infty(\Omega)\cap C_c(\Omega)$

(有些作者用 $\mathcal{D}(\Omega)$ 或 $C_0^\infty(\Omega)$ 代替
$C_c^\infty(\Omega)$).

若 $f\in C^1(\Omega)$, 其梯度定义为
\begin{equation*}
\nabla f=\Bigl(\frac{\partial f}{\partial x_1},
\frac{\partial f}{\partial x_2},\dots,\frac{\partial f}{\partial x_N}\Bigr).
\end{equation*}
若 $f\in C^k(\Omega)$, 而 $\alpha=(\alpha_1,\alpha_2,\dots,\alpha_N)$ 是长度
$\abs{\alpha}=\alpha_1+\alpha_2+\cdots+\alpha_N$ 小于 $k$ 的多重指标, 则写
\begin{equation*}
D^\alpha f=\frac{\partial^{\alpha_1}}{\partial x_1^{\alpha_1}}
\frac{\partial^{\alpha_2}}{\partial x_2^{\alpha_2}}\cdots
\frac{\partial^{\alpha_N}}{\partial x_N^{\alpha_N}}f.
\end{equation*}
\end{notation}

\begin{keyproposition}
\label{prop:4.20}
设 $f\in C_c^k(\R^N)$ ($k\ge1$), $g\in L^1_{\mathrm{loc}}(\R^N)$.则
$f\star g\in C^k(\R^N)$, 且
\begin{equation*}
\boxed{D^\alpha(f\star g)=(D^\alpha f)\star g\quad\forall\alpha,\
\abs{\alpha}\le k.}
\end{equation*}
特别地, 若 $f\in C_c^\infty(\R^N)$, $g\in L^1_{\mathrm{loc}}(\R^N)$, 则
$f\star g\in C^\infty(\R^N)$.
\end{keyproposition}

\begin{proof}
由归纳法只须考虑 $k=1$ 的情形.给定 $x\in\R^N$, 我们断言 $f\star g$ 在 $x$ 处
可微, 且
\begin{equation*}
\nabla(f\star g)(x)=(\nabla f)\star g(x).
\end{equation*}
设 $h\in\R^N$, $\abs{h}<1$.则对一切 $y\in\R^N$,
\begin{equation*}
\begin{aligned}
&\abs{f(x+h-y)-f(x-y)-h\cdot\nabla f(x-y)}\\
&\quad=\Bigl|\int_0^1\bigl[h\cdot\nabla f(x+sh-y)-h\cdot\nabla f(x-y)\bigr]ds
\Bigr|\le\abs{h}\varepsilon(\abs{h}),
\end{aligned}
\end{equation*}
其中 $\varepsilon(\abs{h})\to0$ 当 $\abs{h}\to0$ (因为 $\nabla f$ 在 $\R^N$ 上
一致连续).

取 $\R^N$ 中固定的紧集 $K$, 使其足够大, 以致
$x+B(0,1)-\operatorname{supp}f\subset K$.有
\begin{equation*}
f(x+h-y)-f(x-y)-h\cdot\nabla f(x-y)=0\quad\forall y\notin K,\
\forall h\in B(0,1),
\end{equation*}
因此
\begin{equation*}
\abs{f(x+h-y)-f(x-y)-h\cdot\nabla f(x-y)}\le\abs{h}\varepsilon(\abs{h})\chi_K(y)
\quad\forall y\in\R^N,\ \forall h\in B(0,1).
\end{equation*}
于是对 $h\in B(0,1)$,
\begin{equation*}
\abs{(f\star g)(x+h)-(f\star g)(x)-h\cdot(\nabla f\star g)(x)}
\le\abs{h}\varepsilon(\abs{h})\int_K\abs{g(y)}\,dy.
\end{equation*}
由此推出 $f\star g$ 在 $x$ 处可微, 且
$\nabla(f\star g)(x)=(\nabla f)\star g(x)$.
\end{proof}

\subsection*{磨光核}

\begin{definition}
\emph{\index{moguanghe@磨光核 (mollifier)}
磨光核序列} $(\rho_n)_{n\ge1}$ 是指 $\R^N$ 上满足下述条件的任一函数序列:
\begin{equation*}
\rho_n\in C_c^\infty(\R^N),\quad
\operatorname{supp}\rho_n\subset\overline{B(0,1/n)},\quad
\int\rho_n=1,\quad\rho_n\ge0\ \text{在}\ \R^N\ \text{上}.
\end{equation*}
在下文中, 我们将系统地使用记号 $(\rho_n)$ 表示一个磨光核序列.
\end{definition}

从一个单一的函数 $\rho\in C_c^\infty(\R^N)$ 出发 (要求
$\operatorname{supp}\rho\subset\overline{B(0,1)}$, $\rho\ge0$ 在 $\R^N$ 上, 且
$\rho$ 不恒为零) 容易生成一个磨光核序列——例如函数
\begin{equation*}
\rho(x)=\begin{cases}
e^{1/(\abs{x}^2-1)} & \text{若}\ \abs{x}<1,\\
0 & \text{若}\ \abs{x}>1.
\end{cases}
\end{equation*}
令 $\rho_n(x)=Cn^N\rho(nx)$, 其中 $C=1/\int\rho$, 便得到磨光核序列.

\begin{proposition}
\label{prop:4.21}
假设 $f\in C(\R^N)$.则 $(\rho_n\star f)\xrightarrow[n\to\infty]{}f$ 在 $\R^N$
的紧集上一致收敛.
\end{proposition}

\begin{proof}\footnote{用卷积作正则化的技巧最初由 Leray 与 Friedrichs 引入.}
设 $K\subset\R^N$ 为固定的紧集.给定 $\varepsilon>0$, 存在 $\delta>0$ (依赖于
$K$ 与 $\varepsilon$), 使得
\begin{equation*}
\abs{f(x-y)-f(x)}<\varepsilon\quad\forall x\in K,\ \forall y\in B(0,\delta).
\end{equation*}
对 $x\in\R^N$ 有
\begin{equation*}
\begin{aligned}
(\rho_n\star f)(x)-f(x)
&=\int\bigl[f(x-y)-f(x)\bigr]\rho_n(y)\,dy\\
&=\int_{B(0,1/n)}\bigl[f(x-y)-f(x)\bigr]\rho_n(y)\,dy.
\end{aligned}
\end{equation*}
当 $n>1/\delta$ 且 $x\in K$ 时得到
\begin{equation*}
\abs{(\rho_n\star f)(x)-f(x)}\le\varepsilon\int\rho_n=\varepsilon.
\end{equation*}
\end{proof}

\begin{keytheorem}
\label{thm:4.22}
假设 $f\in L^p(\R^N)$, $1\le p<\infty$.则
$(\rho_n\star f)\xrightarrow[n\to\infty]{}f$ 在 $L^p(\R^N)$ 中.
\end{keytheorem}

\begin{proof}
给定 $\varepsilon>0$, 固定函数 $f_1\in C_c(\R^N)$, 使
$\norm{f-f_1}_p<\varepsilon$ (见定理  \ref{thm:4.12}).由命题  \ref{prop:4.21} 知
$(\rho_n\star f_1)\to f_1$ 在 $\R^N$ 的每个紧集上一致收敛.另一方面, 由命题
 \ref{prop:4.18} 有
\begin{equation*}
\operatorname{supp}(\rho_n\star f_1)\subset B(0,1/n)+\operatorname{supp}f_1
\subset B(0,1)+\operatorname{supp}f_1,
\end{equation*}
这是一个固定的紧集.由此推出
\begin{equation*}
\norm{(\rho_n\star f_1)-f_1}_p\xrightarrow[n\to\infty]{}0.
\end{equation*}
最后, 写
\begin{equation*}
(\rho_n\star f)-f=\bigl[\rho_n\star(f-f_1)\bigr]
+\bigl[(\rho_n\star f_1)-f_1\bigr]+[f_1-f],
\end{equation*}
于是
\begin{equation*}
\norm{(\rho_n\star f)-f}_p\le2\norm{f-f_1}_p+\norm{(\rho_n\star f_1)-f_1}_p
\end{equation*}
(由定理  \ref{thm:4.15}).我们断定
\begin{equation*}
\limsup_{n\to\infty}\norm{(\rho_n\star f)-f}_p\le2\varepsilon
\quad\forall\varepsilon>0,
\end{equation*}
因此 $\lim_{n\to\infty}\norm{(\rho_n\star f)-f}_p=0$.
\end{proof}

\begin{keycorollary}
\label{cor:4.23}
设 $\Omega\subset\R^N$ 为开集.则对任意的 $1\le p<\infty$,
$C_c^\infty(\Omega)$ 在 $L^p(\Omega)$ 中稠密.
\end{keycorollary}

\begin{proof}
给定 $f\in L^p(\Omega)$, 令
\begin{equation*}
\bar f(x)=\begin{cases}
f(x) & \text{若}\ x\in\Omega,\\
0 & \text{若}\ x\in\R^N\setminus\Omega,
\end{cases}
\end{equation*}
于是 $\bar f\in L^p(\R^N)$.

取 $\R^N$ 中紧集构成的序列 $(K_n)$, 使得
\begin{equation*}
\bigcup_{n=1}^\infty K_n=\Omega
\quad\text{且}\quad\dist(K_n,\Omega^c)\ge2/n\quad\forall n.
\end{equation*}
[例如可以取
$K_n=\{x\in\R^N;\ \abs{x}\le n\ \text{且}\ \dist(x,\Omega^c)\ge2/n\}$.]
令 $g_n=\chi_{K_n}\bar f$, $f_n=\rho_n\star g_n$, 于是
\begin{equation*}
\operatorname{supp}f_n\subset B(0,1/n)+K_n\subset\Omega.
\end{equation*}
由此推出 $f_n\in C_c^\infty(\Omega)$.另一方面有
\begin{equation*}
\begin{aligned}
\norm{f_n-f}_{L^p(\Omega)}=\norm{f_n-\bar f}_{L^p(\R^N)}
&\le\norm{(\rho_n\star g_n)-(\rho_n\star\bar f)}_{L^p(\R^N)}
+\norm{(\rho_n\star\bar f)-\bar f}_{L^p(\R^N)}\\
&\le\norm{g_n-\bar f}_{L^p(\R^N)}
+\norm{(\rho_n\star\bar f)-\bar f}_{L^p(\R^N)}.
\end{aligned}
\end{equation*}
最后注意 $\norm{g_n-\bar f}_{L^p(\R^N)}\to0$ (由控制收敛定理),
而 $\norm{(\rho_n\star\bar f)-\bar f}_{L^p(\R^N)}\to0$ (由定理
 \ref{thm:4.22}).我们断定 $\norm{f_n-f}_{L^p(\Omega)}\to0$.
\end{proof}

\begin{corollary}
\label{cor:4.24}
设 $\Omega\subset\R^N$ 为开集, $u\in L^1_{\mathrm{loc}}(\Omega)$ 满足
\begin{equation*}
\int uf=0\quad\forall f\in C_c^\infty(\Omega).
\end{equation*}
则 $u=0$ 在 $\Omega$ 上几乎处处成立.
\end{corollary}

\begin{proof}
设 $g\in L^\infty(\R^N)$, 其支集 $\operatorname{supp}g$ 是包含在 $\Omega$ 中的
紧集.令 $g_n=\rho_n\star g$, 则当 $n$ 足够大时 $g_n\in C_c^\infty(\Omega)$.
因此有
\begin{equation}
\int ug_n=0\quad\forall n.
\label{eq:4.19}
\end{equation}
由于 $g_n\to g$ 在 $L^1(\R^N)$ 中 (由定理  \ref{thm:4.22}), 存在子列——仍记作
$g_n$——使得 $g_n\to g$ 在 $\R^N$ 上几乎处处成立 (见定理  \ref{thm:4.9}).此外,
$\norm{g_n}_{L^\infty(\R^N)}\le\norm{g}_{L^\infty(\R^N)}$.在 \eqref{eq:4.19}
中取极限 (由控制收敛定理), 得到
\begin{equation}
\int ug=0.
\label{eq:4.20}
\end{equation}

设 $K$ 为包含在 $\Omega$ 中的紧集.取函数 $g$ 为
\begin{equation*}
g=\begin{cases}
\operatorname{sign}u & \text{在}\ K\ \text{上},\\
0 & \text{在}\ \R^N\setminus K\ \text{上}.
\end{cases}
\end{equation*}
由 \eqref{eq:4.20} 推出 $\int_K\abs{u}=0$, 从而 $u=0$ 在 $K$ 上几乎处处成立.
由于这对任一紧集 $K\subset\Omega$ 都成立, 断定 $u=0$ 在 $\Omega$ 上几乎处处成立.
\end{proof}

\section{$L^p$ 中强紧性的判别准则}
\label{sec:4.5}

判断 $L^p(\Omega)$ 中的一族函数是否在 $L^p(\Omega)$ 中 (就强拓扑而言) 具有紧闭包,
是重要的.我们回顾: Ascoli--Arzel\`a 定理回答了 $C(K)$——紧度量空间 $K$ 上取值于
$\R$ 的连续函数空间——中的同一问题.

\begin{keytheorem}[Ascoli--Arzel\`a]
\label{thm:4.25}
设 $K$ 为紧度量空间, $H$ 为 $C(K)$ 的有界子集.假设 $H$ 一致等度连续, 即
\begin{equation}
\forall\varepsilon>0\ \exists\delta>0\ \text{使得}\ d(x_1,x_2)<\delta
\Rightarrow\abs{f(x_1)-f(x_2)}<\varepsilon\quad\forall f\in H.
\label{eq:4.21}
\end{equation}
则 $H$ 在 $C(K)$ 中的闭包紧.
\end{keytheorem}

Ascoli--Arzel\`a 定理的证明可参见例如 W.~Rudin \bref{Rudin}{1}, \bref{Rudin}{2}, A.~Knapp \bref{Knapp}{1}, J.~Dixmier \bref{Dixmier}{1}, A.~Friedman \bref{Friedman}{3}, G.~Choquet \bref{Choquet}{1}, K.~Yosida \bref{Yosida}{1}, H.~L.~Royden \bref{Royden}{1}, J.~R.~Munkres \bref{Munkres}{1}, G.~B.~Folland \bref{Folland}{2} 等.

\begin{notation}[函数的平移]
令 $(\tau_hf)(x)=f(x+h)$, $x\in\R^N$, $h\in\R^N$.
\end{notation}

下面这个定理及其推论是 Ascoli--Arzel\`a 定理的``$L^p$ 版本''.

\begin{keytheorem}[Kolmogorov--M. Riesz--Fr\'echet]
\label{thm:4.26}
设 $\mathcal{F}$ 为 $L^p(\R^N)$ 的有界集, $1\le p<\infty$.假设\footnote{假设
\eqref{eq:4.22} 应与 \eqref{eq:4.21} 比较.它是一个``积分型''等度连续性假设.}
\begin{equation}
\lim_{\abs{h}\to0}\norm{\tau_hf-f}_p=0\quad\text{关于}\ f\in\mathcal{F}\
\text{一致},
\label{eq:4.22}
\end{equation}
即 $\forall\varepsilon>0\ \exists\delta>0$ 使得
$\norm{\tau_hf-f}_p<\varepsilon\ \forall f\in\mathcal{F},\ \forall h\in\R^N$,
$\abs{h}<\delta$.

则对任一测度有限的 $\R^N$ 的可测子集 $\Omega$, $\mathcal{F}|_\Omega$ 在
$L^p(\Omega)$ 中的闭包紧.

[这里 $\mathcal{F}|_\Omega$ 表示 $\mathcal{F}$ 中函数在 $\Omega$ 上的限制.]
\end{keytheorem}

证明分四步.

\runin{第 1 步.}我们断言
\begin{equation}
\norm{(\rho_n\star f)-f}_{L^p(\R^N)}\le\varepsilon\quad\forall f\in\mathcal{F},\
\forall n>1/\delta.
\label{eq:4.23}
\end{equation}
事实上, 由 Hölder 不等式有
\begin{equation*}
\begin{aligned}
\abs{(\rho_n\star f)(x)-f(x)}
&\le\int\abs{f(x-y)-f(x)}\rho_n(y)\,dy\\
&\le\Bigl(\int\abs{f(x-y)-f(x)}^p\rho_n(y)\,dy\Bigr)^{1/p}.
\end{aligned}
\end{equation*}
于是
\begin{equation*}
\begin{aligned}
\int\abs{(\rho_n\star f)(x)-f(x)}^p\,dx
&\le\iint\abs{f(x-y)-f(x)}^p\rho_n(y)\,dx\,dy\\
&=\int_{B(0,1/n)}\rho_n(y)\,dy\int\abs{f(x-y)-f(x)}^p\,dx\le\varepsilon^p,
\end{aligned}
\end{equation*}
只要 $1/n<\delta$.

\runin{第 2 步.}我们断言
\begin{equation}
\norm{\rho_n\star f}_{L^\infty(\R^N)}\le C_n\norm{f}_{L^p(\R^N)}
\quad\forall f\in\mathcal{F}
\label{eq:4.24}
\end{equation}
且
\begin{equation}
\abs{(\rho_n\star f)(x_1)-(\rho_n\star f)(x_2)}\le C_n\norm{f}_p\abs{x_1-x_2}
\quad\forall f\in\mathcal{F},\ \forall x_1,x_2\in\R^N,
\label{eq:4.25}
\end{equation}
其中 $C_n$ 只依赖于 $n$.

不等式 \eqref{eq:4.24} 由 Hölder 不等式得到, 其中 $C_n=\norm{\rho_n}_{p'}$.
另一方面, $\nabla(\rho_n\star f)=(\nabla\rho_n)\star f$, 因此
\begin{equation*}
\norm{\nabla(\rho_n\star f)}_{L^\infty(\R^N)}
\le\norm{\nabla\rho_n}_{L^{p'}(\R^N)}\norm{f}_{L^p(\R^N)}.
\end{equation*}
于是得到 \eqref{eq:4.25}, 其中 $C_n=\norm{\nabla\rho_n}_{L^{p'}(\R^N)}$.

\runin{第 3 步.}给定 $\varepsilon>0$ 与测度有限的 $\Omega\subset\R^N$, 存在
$\Omega$ 的有界可测子集 $\omega$, 使得
\begin{equation}
\norm{f}_{L^p(\Omega\setminus\omega)}<\varepsilon\quad\forall f\in\mathcal{F}.
\label{eq:4.26}
\end{equation}
事实上, 写
\begin{equation*}
\norm{f}_{L^p(\Omega\setminus\omega)}
\le\norm{f-(\rho_n\star f)}_{L^p(\R^N)}
+\norm{\rho_n\star f}_{L^p(\Omega\setminus\omega)}.
\end{equation*}
由 \eqref{eq:4.24}, 只须选取 $\omega$ 使 $\abs{\Omega\setminus\omega}$ 足够小.

\runin{第 4 步: 结论.}由于 $L^p(\Omega)$ 完备, 只须 (参见例如 A.~Knapp
\bref{Knapp}{1} 或 J.~R.~Munkres \bref{Munkres}{1} 第~\ref{sec:7.3}~节) 证明
$\mathcal{F}|_\Omega$ 全有界, 即给定任一 $\varepsilon>0$, 存在
$\mathcal{F}|_\Omega$ 的有限个半径为 $\varepsilon$ 的球覆盖.给定
$\varepsilon>0$, 取有界可测集 $\omega$ 使 \eqref{eq:4.26} 成立, 并取
$n>1/\delta$.族 $H=(\rho_n\star\mathcal{F})|_{\bar\omega}$ 满足
Ascoli--Arzel\`a 定理的全部假设 (由第 2 步).因此 $H$ 在 $C(\bar\omega)$ 中的
闭包紧; 从而 $H$ 在 $L^p(\omega)$ 中的闭包也紧.于是可以用有限个
$L^p(\omega)$ 中半径为 $\varepsilon$ 的球覆盖 $H$, 即
\begin{equation*}
H\subset\bigcup_iB(g_i,\varepsilon)\quad\text{其中}\ g_i\in L^p(\omega).
\end{equation*}
考虑由
\begin{equation*}
\bar g_i=\begin{cases}
g_i & \text{在}\ \omega\ \text{上},\\
0 & \text{在}\ \Omega\setminus\omega\ \text{上}
\end{cases}
\end{equation*}
定义的函数 $\bar g_i:\Omega\to\R$, 以及 $L^p(\Omega)$ 中的球
$B(\bar g_i,3\varepsilon)$.

我们断言这些球覆盖 $\mathcal{F}|_\Omega$.事实上, 给定 $f\in\mathcal{F}$, 存在
某个 $i$ 使
\begin{equation*}
\norm{(\rho_n\star f)-g_i}_{L^p(\omega)}<\varepsilon.
\end{equation*}
由于
\begin{equation*}
\norm{f-\bar g_i}_{L^p(\Omega)}^p=\int_{\Omega\setminus\omega}\abs{f}^p
+\int_\omega\abs{f-g_i}^p,
\end{equation*}
由 \eqref{eq:4.26} 有
\begin{equation*}
\begin{aligned}
\norm{f-\bar g_i}_{L^p(\Omega)}
&\le\varepsilon+\norm{f-g_i}_{L^p(\omega)}\\
&\le\varepsilon+\norm{f-(\rho_n\star f)}_{L^p(\R^N)}
+\norm{(\rho_n\star f)-g_i}_{L^p(\omega)}<3\varepsilon.
\end{aligned}
\end{equation*}
我们断定 $\mathcal{F}|_\Omega$ 在 $L^p(\Omega)$ 中的闭包紧.

\begin{remark}
\label{rem:4.11}
当试图证明 $L^p(\Omega)$ ($\Omega$ 有界) 中的一族函数 $\mathcal{F}$ 在
$L^p(\Omega)$ 中具有紧闭包时, 通常方便的做法是先把函数延拓到整个 $\R^N$ 上, 然后
应用定理  \ref{thm:4.26}, 再考虑它们在 $\Omega$ 上的限制.
\end{remark}

\begin{remark}
\label{rem:4.12}
在定理  \ref{thm:4.26} 的假设下, 一般不能断定 $\mathcal{F}$ 本身在
$L^p(\R^N)$ 中具有紧闭包 (可构造例子, 或参见习题 4.33).还需要一个附加假设;
我们下面来叙述它:
\end{remark}

\begin{corollary}
\label{cor:4.27}
设 $\mathcal{F}$ 为 $L^p(\R^N)$ 的有界集, $1\le p<\infty$.假设
\eqref{eq:4.22} 成立, 并且
\begin{equation}
\begin{cases}
\forall\varepsilon>0\ \exists\Omega\subset\R^N\ \text{有界, 可测, 使得}\\
\norm{f}_{L^p(\R^N\setminus\Omega)}<\varepsilon\quad\forall f\in\mathcal{F}.
\end{cases}
\label{eq:4.27}
\end{equation}
则 $\mathcal{F}$ 在 $L^p(\R^N)$ 中的闭包紧.
\end{corollary}

\begin{proof}
给定 $\varepsilon>0$, 取有界可测集 $\Omega\subset\R^N$ 使 \eqref{eq:4.27} 成立.
由定理  \ref{thm:4.26} 知 $\mathcal{F}|_\Omega$ 在 $L^p(\Omega)$ 中的闭包紧.
于是可以用有限个 $L^p(\Omega)$ 中半径为 $\varepsilon$ 的球覆盖
$\mathcal{F}|_\Omega$, 即
\begin{equation*}
\mathcal{F}|_\Omega\subset\bigcup_iB(g_i,\varepsilon)\quad\text{其中}\
g_i\in L^p(\Omega).
\end{equation*}
令
\begin{equation*}
\bar g_i(x)=\begin{cases}
g_i(x) & \text{在}\ \Omega\ \text{中},\\
0 & \text{在}\ \R^N\setminus\Omega\ \text{上}.
\end{cases}
\end{equation*}
显然 $\mathcal{F}$ 被 $L^p(\R^N)$ 中的球 $B(\bar g_i,2\varepsilon)$ 覆盖.
\end{proof}

\begin{remark}
\label{rem:4.13}
推论  \ref{cor:4.27} 的逆命题也成立 (见习题 4.34).因此我们得到了
$L^p(\R^N)$ 中紧集的完全刻画.
\end{remark}

最后给出定理  \ref{thm:4.26} 的一个有用应用:

\begin{corollary}
\label{cor:4.28}
设 $G$ 为 $L^1(\R^N)$ 中固定的函数, 并令
\begin{equation*}
\mathcal{F}=G\star B,
\end{equation*}
其中 $B$ 是 $L^p(\R^N)$ 中的有界集, $1\le p<\infty$.则对任一测度有限的可测集
$\Omega$, $\mathcal{F}|_\Omega$ 在 $L^p(\Omega)$ 中的闭包紧.
\end{corollary}

\begin{proof}
显然 $\mathcal{F}$ 在 $L^p(\R^N)$ 中有界.另一方面, 若写 $f=G\star u$, 其中
$u\in B$, 则
\begin{equation*}
\norm{\tau_hf-f}_p=\norm{(\tau_hG-G)\star u}_p\le C\norm{\tau_hG-G}_1,
\end{equation*}
再借助下述引理即得结论:
\end{proof}

\begin{lemma}
\label{lem:4.3}
设 $G\in L^q(\R^N)$, $1\le q<\infty$.则
\begin{equation*}
\lim_{h\to0}\norm{\tau_hG-G}_q=0.
\end{equation*}
\end{lemma}

\begin{proof}
给定 $\varepsilon>0$, 由定理  \ref{thm:4.12} 存在函数 $G_1\in C_c(\R^N)$, 使得
$\norm{G-G_1}_q<\varepsilon$.写
\begin{equation*}
\begin{aligned}
\norm{\tau_hG-G}_q
&\le\norm{\tau_hG-\tau_hG_1}_q+\norm{\tau_hG_1-G_1}_q+\norm{G_1-G}_q\\
&\le2\varepsilon+\norm{\tau_hG_1-G_1}_q.
\end{aligned}
\end{equation*}
由于 $\lim_{h\to0}\norm{\tau_hG_1-G_1}_q=0$, 看出
\begin{equation*}
\limsup_{h\to0}\norm{\tau_hG-G}_q\le2\varepsilon\quad\forall\varepsilon>0.
\end{equation*}
\end{proof}

\section*{第 4 章评注}
\addcontentsline{toc}{section}{第 4 章评注}
\markright{第 4 章评注}

\cmtitem{1. Egorov 定理.}
积分论的一些基本结果已在 \S\ref{sec:4.1} 中回顾.有一个未被提及的有用结果如下.

\begin{startheorem}[Egorov]
\label{thm:4.29}
假设 $\Omega$ 是测度有限的测度空间.设 $(f_n)$ 为 $\Omega$ 上的可测函数序列,
使得
\begin{equation*}
f_n(x)\to f(x)\quad\text{a.e. 在}\ \Omega\ \text{上}\quad
(\text{其中}\ \abs{f(x)}<\infty\ \text{a.e.}).
\end{equation*}
则 $\forall\varepsilon>0\ \exists A\subset\Omega$ 可测, 使 $\abs{\Omega\setminus A}
<\varepsilon$ 且 $f_n\to f$ 在 $A$ 上一致收敛.
\end{startheorem}

其证明见\index{Vitalishouliandingli@Vitali 收敛定理 (Vitali convergence theorem)}
习题 4.14, P.~Halmos \bref{Halmos}{1}, G.~B.~Folland \bref{Folland}{2}, E.~Hewitt--K.~Stromberg \bref{HewittStromberg}{1}, R.~Wheeden--A.~Zygmund \bref{WheedenZygmund}{1}, K.~Yosida \bref{Yosida}{1}, A.~Friedman \bref{Friedman}{3} 等.

\cmtitem{2. $L^1$ 中的弱紧集.}
由于 $L^1$ 不自反, $L^1$ 的有界集关于弱拓扑 $\sigma(L^1,L^\infty)$ 并不起重要
作用.下述结果给出了 $L^1$ 中弱紧集的一个有用刻画.

\begin{startheorem}[Dunford--Pettis]
\label{thm:4.30}
设 $\mathcal{F}$ 为 $L^1(\Omega)$ 的有界集.则 $\mathcal{F}$ 在弱拓扑
$\sigma(L^1,L^\infty)$ 中具有紧闭包, 当且仅当 $\mathcal{F}$ 等度可积, 即
\begin{equation*}
\text{(a)}\ \begin{cases}
\forall\varepsilon>0\ \exists\delta>0\ \text{使得}\\
\displaystyle\int_A\abs{f}<\varepsilon\quad
\forall A\subset\Omega\ \text{可测, 且}\ \abs{A}<\delta,\ \forall f\in\mathcal{F},
\end{cases}
\end{equation*}
并且
\begin{equation*}
\text{(b)}\ \begin{cases}
\forall\varepsilon>0\ \exists\omega\subset\Omega\ \text{可测, 且}\
\abs{\omega}<\infty\ \text{使得}\\
\displaystyle\int_{\Omega\setminus\omega}\abs{f}<\varepsilon\quad\forall f\in\mathcal{F}.
\end{cases}
\end{equation*}
\end{startheorem}

定理  \ref{thm:4.30} 的证明与讨论见问题 23, 或 N.~Dunford--J.~T.~Schwartz
\bref{DunfordSchwartz}{1}, B.~Beauzamy \bref{Beauzamy}{1}, J.~Diestel \bref{Diestel}{2}, I.~Fonseca--G.~Leoni \bref{FonsecaLeoni}{1};
关于概率论方面, 还可见 J.~Neveu \bref{Neveu}{1}, C.~Dellacherie--P.~A.~Meyer \bref{DellacherieMeyer}{1}; 另见习题 4.36.

\cmtitem{3. Radon 测度.}
正如刚才所指出的, $L^1$ 的有界集不具有紧性性质.为克服这种紧性的缺失, 有时把
$L^1$ 嵌入一个更大的空间——Radon 测度空间——是非常有用的.

例如, 假设 $\Omega$ 是 $\R^N$ 中带 Lebesgue 测度的有界开集.考虑空间
$E=C(\overline\Omega)$, 配以范数 $\norm{u}=\sup_{x\in\overline\Omega}\abs{u(x)}$.
它的对偶空间记作 $\mathcal{M}(\overline\Omega)$, 称为 $\overline\Omega$ 上的
\emph{Radon 测度}空间.$\mathcal{M}(\overline\Omega)$ 上的\wstar{}拓扑有时称为
``淡''拓扑.

我们把 $L^1(\Omega)$ 与 $\mathcal{M}(\overline\Omega)$ 的一个子空间等同起来.为此
引入如下定义的映射 $L^1(\Omega)\to\mathcal{M}(\overline\Omega)$: 给定
$f\in L^1(\Omega)$, 映射 $u\in C(\overline\Omega)\mapsto\int_\Omega fu\,dx$ 是
$C(\overline\Omega)$ 上的连续线性泛函, 记作 $Tf$, 于是
\begin{equation*}
\dual{Tf}{u}_{E^\star,E}=\int_\Omega fu\,dx\quad\forall u\in E.
\end{equation*}
显然 $T$ 是线性的, 而且 $T$ 是等距, 因为
\begin{equation*}
\norm{Tf}_{\mathcal{M}(\overline\Omega)}
=\sup_{\substack{u\in E\\ \norm{u}\le1}}\int_\Omega fu=\norm{f}_1
\end{equation*}
(见习题 4.26).

利用 $T$ 可以把 $L^1(\Omega)$ 与 $\mathcal{M}(\overline\Omega)$ 的一个子空间
等同起来.由于 $\mathcal{M}(\overline\Omega)$ 是可分空间 $C(\overline\Omega)$ 的
对偶空间, 它在\wstar{}拓扑中具有某些紧性性质.特别地, 若 $(f_n)$ 是
$L^1(\Omega)$ 中的有界序列, 则存在子列 $(f_{n_k})$ 与 Radon 测度 $\mu$, 使得
$f_{n_k}\overset{\star}{\rightharpoonup}\mu$ 在\wstar{}拓扑
$\sigma(E^\star,E)$ 中, 即
\begin{equation*}
\int_\Omega f_{n_k}u\to\dual{\mu}{u}\quad\forall u\in C(\overline\Omega).
\end{equation*}
例如, $L^1$ 中的序列可以关于\wstar{}拓扑收敛于一个 Dirac 测度.Radon 测度的
一些进一步的性质见问题 24.

``测度''这一术语由下述结果提供依据, 它把上述定义与集合论意义下标准的测度概念联系
起来:

\begin{theorem}[Riesz 表示定理]
\label{thm:4.31}
设 $\mu$ 为 $\Omega$ 上的 Radon 测度.则存在 $\Omega$ 上唯一的带号 Borel 测度
$\nu$ (即定义在 $\Omega$ 的 Borel 集上的测度), 使得
\begin{equation*}
\dual{\mu}{u}=\int_\Omega u\,d\nu\quad\forall u\in C(\overline\Omega).
\end{equation*}
\end{theorem}

通常方便的做法是把空间 $E=C(\overline\Omega)$ 换成子空间
\begin{equation*}
E_0=\{f\in C(\overline\Omega);\ f=0\ \text{在}\ \overline\Omega\ \text{的边界上}\}.
\end{equation*}
$E_0$ 的对偶记作 $\mathcal{M}(\Omega)$ (以区别于 $\mathcal{M}(\overline\Omega)$).
Riesz 表示定理仍然成立, 只需附加条件 $\abs{\nu}(\overline\Omega\
\text{的边界})=0$.

关于这一广阔而经典的题材, 参见例如 H.~L.~Royden \bref{Royden}{1}, W.~Rudin \bref{Rudin}{2}, G.~B.~Folland \bref{Folland}{2}, A.~Knapp \bref{Knapp}{1}, P.~Malliavin \bref{Malliavin}{1}, P.~Halmos \bref{Halmos}{1}, I.~Fonseca--G.~Leoni \bref{FonsecaLeoni}{1}.

\cmtitem{4. 向量值函数的 Bochner 积分.}
设 $\Omega$ 为测度空间, $E$ 为 Banach 空间.空间 $L^p(\Omega;E)$ 由定义在
$\Omega$ 上, 取值于 $E$ 的, 在某种适当意义下可测且满足
$\int_\Omega\norm{f(x)}^p\,d\mu<\infty$ 的全部函数 $f$ 组成 (当 $p=\infty$ 时作
通常的修改).\S\ref{sec:4.2} 与 \S\ref{sec:4.3} 中描述的大部分性质在对 $E$ 附加一些假设后仍然成立.
例如, 若 $E$ 自反且 $1<p<\infty$, 则 $L^p(\Omega;E)$ 自反, 其对偶空间为
$L^{p'}(\Omega;E^\star)$.更多细节见 K.~Yosida \bref{Yosida}{1}, D.~L.~Cohn \bref{Cohn}{1}, E.~Hille \bref{Hille}{1}, B.~Beauzamy \bref{Beauzamy}{1}, L.~Schwartz \bref{Schwartz}{3}.
当 $\Omega$ 是 $\R$ 中的区间时, 空间 $L^p(\Omega;E)$ 在演化方程的研究中非常
有用 (见第~
\ref{chap:10}~章).

\cmtitem{5. 插值理论.}
开创插值理论的最引人注目的结果是下述定理.

\begin{theorem}[Schur, M. Riesz, Thorin]
\label{thm:4.32}
假设 $\Omega$ 是测度空间, $\abs{\Omega}<\infty$, 且 $T:L^1(\Omega)\to
L^1(\Omega)$ 是有界线性算子, 其范数为 $M_1=\norm{T}_{\mathcal{L}(L^1,L^1)}$.
此外假设 $T:L^\infty(\Omega)\to L^\infty(\Omega)$ 是有界线性算子, 其范数为
$M_\infty=\norm{T}_{\mathcal{L}(L^\infty,L^\infty)}$.

则对一切 $1<p<\infty$, $T$ 是从 $L^p(\Omega)$ 到 $L^p(\Omega)$ 的有界算子,
且其范数 $M_p$ 满足
\begin{equation*}
M_p\le M_1^{1/p}M_\infty^{1/p'}.
\end{equation*}
\end{theorem}

插值理论最初由 I.~Schur, M.~Riesz, G.~O.~Thorin, J.~Marcinkiewicz 与
A.~Zygmund 发现.许多作者作出了决定性的贡献, 其中包括 J.-L.~Lions, J.~Peetre, A.~P.~Calder\'on, E.~Stein 与 E.~Gagliardo.它已成为调和分析
(参见例如 E.~Stein--G.~Weiss \bref{SteinWeiss}{1}, E.~Stein \bref{Stein}{1}, C.~Sadosky \bref{Sadosky}{1}) 与偏微分方程 (参见例如
J.-L.~Lions--E.~Magenes \bref{LionsMagenes}{1}) 中的有用工具.关于这些问题,
还可见 G.~B.~Folland \bref{Folland}{2}, N.~Dunford--J.~T.~Schwartz
\bref{DunfordSchwartz}{1} (第 1 卷第 520 页), J.~Bergh--J.~L\"ofstr\"om
\bref{BerghLofstrom}{1}, M.~Reed--B.~Simon \bref{ReedSimon}{1} (第 2 卷第 27
页) 以及问题 22.

\cmtitem{6. Young 不等式.}
下面是定理  \ref{thm:4.15} 的一个推广.

\begin{theorem}[Young]
\label{thm:4.33}
假设 $f\in L^p(\R^N)$, $g\in L^q(\R^N)$, 其中 $1\le p\le\infty$,
$1\le q\le\infty$, 且
\begin{equation*}
\frac1r=\frac1p+\frac1q-1\ge0.
\end{equation*}
则 $f\star g\in L^r(\R^N)$, 且 $\norm{f\star g}_r\le\norm{f}_p\norm{g}_q$.
\end{theorem}

其证明参见例如习题 4.30.

\cmtnum{7.}
卷积的概念——它可推广到广义函数 (见 L.~Schwartz \bref{Schwartz}{1} 或
A.~Knapp \bref{Knapp}{2})——在偏微分方程理论中起着基本作用.例如, 方程
$P(D)u=f$ 在 $\R^N$ 中 (其中 $P(D)$ 是任一常系数微分算子) 有形式为
$u=E\star f$ 的解, 其中 $E$ 是 $P(D)$ 的基本解 (Malgrange--Ehrenpreis 定理;
另见第~
\ref{chap:1}~章评注 2b).特别地, 方程 $\Delta u=f$ 在 $\R^3$ 中有形式为
$u=E\star f$ 的解, 其中 $E(x)=-(4\pi\abs{x})^{-1}$.

\section*{第 4 章习题}
\addcontentsline{toc}{section}{第 4 章习题}
\markright{第 4 章习题}

除另有说明外, $\Omega$ 表示一个 $\sigma$-有限测度空间.

\begin{exercise}
\label{exer:4.1}
设 $\alpha>0$, $\beta>0$.令
\begin{equation*}
f(x)=\bigl(1+\abs{x}^\alpha\bigr)^{-1}
\bigl(1+\bigl|\log\abs{x}\bigr|^\beta\bigr)^{-1},\quad x\in\R^N.
\end{equation*}
$f$ 在什么条件下属于 $L^p(\R^N)$?
\end{exercise}

\begin{exercise}
\label{exer:4.2}
假设 $\abs{\Omega}<\infty$, $1\le p\le q\le\infty$.证明 $L^q(\Omega)\subset
L^p(\Omega)$ 且嵌入连续.更确切地说, 证明
\begin{equation*}
\norm{f}_p\le\abs{\Omega}^{\frac1p-\frac1q}\norm{f}_q\quad\forall f\in L^q(\Omega).
\end{equation*}
[提示: 使用 Hölder 不等式.]
\end{exercise}

\begin{exercise}
\label{exer:4.3}
\begin{enumerate}[label=\arabic*., leftmargin=2.2em, itemsep=0.35em]
	\item 设 $f,g\in L^p(\Omega)$, $1\le p\le\infty$.证明
	      \begin{equation*}
	      h(x)=\max\{f(x),g(x)\}\in L^p(\Omega).
	      \end{equation*}
	\item 设 $(f_n)$ 与 $(g_n)$ 为 $L^p(\Omega)$ 中的两个序列, $1\le p\le\infty$,
	      使得 $f_n\to f$ 在 $L^p(\Omega)$ 中, $g_n\to g$ 在 $L^p(\Omega)$ 中.
	      令 $h_n=\max\{f_n,g_n\}$, 证明 $h_n\to h$ 在 $L^p(\Omega)$ 中.
	\item 设 $(f_n)$ 为 $L^p(\Omega)$ 中的序列, $1\le p<\infty$; $(g_n)$ 为
	      $L^\infty(\Omega)$ 中的有界序列.假设 $f_n\to f$ 在 $L^p(\Omega)$ 中,
	      且 $g_n\to g$ 几乎处处.证明 $f_ng_n\to fg$ 在 $L^p(\Omega)$ 中.
\end{enumerate}
\end{exercise}

\begin{exercise}
\label{exer:4.4}
\begin{enumerate}[label=\arabic*., leftmargin=2.2em, itemsep=0.35em]
	\item 设 $f_1,f_2,\dots,f_k$ 是 $k$ 个函数, 满足 $f_i\in L^{p_i}(\Omega)\ \forall i$,
	      其中 $1\le p_i\le\infty$, 且 $\sum_{i=1}^k\frac1{p_i}\le1$.令
	      \begin{equation*}
	      f(x)=\prod_{i=1}^kf_i(x).
	      \end{equation*}
	      证明 $f\in L^p(\Omega)$, 其中 $\frac1p=\sum_{i=1}^k\frac1{p_i}$, 且
	      \begin{equation*}
	      \norm{f}_p\le\prod_{i=1}^k\norm{f_i}_{p_i}.
	      \end{equation*}
	      [提示: 先对 $k=2$ 证明, 再用归纳法.]
	\item 由此推出: 若 $f\in L^p(\Omega)\cap L^q(\Omega)$, 其中
	      $1\le p\le\infty$, $1\le q\le\infty$, 则对 $p$ 与 $q$ 之间的每个 $r$
	      都有 $f\in L^r(\Omega)$.更确切地说, 写
	      \begin{equation*}
	      \frac1r=\frac\alpha p+\frac{1-\alpha}q,\quad\text{其中}\ \alpha\in[0,1],
	      \end{equation*}
	      并证明
	      \begin{equation*}
	      \norm{f}_r\le\norm{f}_p^\alpha\norm{f}_q^{1-\alpha}.
	      \end{equation*}
\end{enumerate}
\end{exercise}

\begin{exercise}
\label{exer:4.5}
设 $1\le p<\infty$, $1\le q\le\infty$.
\begin{enumerate}[label=\arabic*., leftmargin=2.2em, itemsep=0.35em]
	\item 证明 $L^1(\Omega)\cap L^\infty(\Omega)$ 是 $L^p(\Omega)$ 的稠密子集.
	\item 证明集合
	      \begin{equation*}
	      \bigl\{f\in L^p(\Omega)\cap L^q(\Omega);\ \norm{f}_q\le1\bigr\}
	      \end{equation*}
	      在 $L^p(\Omega)$ 中闭.
	\item 设 $(f_n)$ 为 $L^p(\Omega)\cap L^q(\Omega)$ 中的序列, $f\in L^p(\Omega)$.
	      假设 $f_n\to f$ 在 $L^p(\Omega)$ 中, 且 $\norm{f_n}_q\le C$.证明
	      $f\in L^r(\Omega)$, 且对 $p$ 与 $q$ 之间的每个 $r$ ($r\ne q$) 都有
	      $f_n\to f$ 在 $L^r(\Omega)$ 中.
\end{enumerate}
\end{exercise}

\begin{exercise}
\label{exer:4.6}
假设 $\abs{\Omega}<\infty$.
\begin{enumerate}[label=\arabic*., leftmargin=2.2em, itemsep=0.35em]
	\item 设 $f\in L^\infty(\Omega)$.证明 $\lim_{p\to\infty}\norm{f}_p=
	      \norm{f}_\infty$.
	\item 设 $f\in\bigcap_{1\le p<\infty}L^p(\Omega)$, 并假设存在常数 $C$, 使得
	      \begin{equation*}
	      \norm{f}_p\le C\quad\forall1\le p<\infty.
	      \end{equation*}
	      证明 $f\in L^\infty(\Omega)$.
	\item 对 $\Omega=(0,1)$, 构造一个函数 $f\in\bigcap_{1\le p<\infty}L^p(\Omega)$
	      的例子, 使得 $f\notin L^\infty(\Omega)$.
\end{enumerate}
\end{exercise}

\begin{exercise}
\label{exer:4.7}
设 $1\le q\le p\le\infty$.设 $a(x)$ 为 $\Omega$ 上的可测函数.假设对每个函数
$u\in L^p(\Omega)$ 都有 $au\in L^q(\Omega)$.

证明 $a\in L^r(\Omega)$, 其中
\begin{equation*}
r=\begin{cases}
\dfrac{pq}{p-q} & \text{若}\ p<\infty,\\[8pt]
q & \text{若}\ p=\infty.
\end{cases}
\end{equation*}
[提示: 使用闭图像定理.]
\end{exercise}

\begin{exercise}
\label{exer:4.8}
设 $X\subset L^1(\Omega)$ 是 $L^1(\Omega)$ 中的闭向量空间.假设
\begin{equation*}
X\subset\bigcup_{1<q\le\infty}L^q(\Omega).
\end{equation*}
\begin{enumerate}[label=\arabic*., leftmargin=2.2em, itemsep=0.35em]
	\item 证明存在某个 $p>1$, 使得 $X\subset L^p(\Omega)$.
	      [提示: 对每个整数 $n\ge1$ 考虑集合
	      \begin{equation*}
	      X_n=\bigl\{f\in X\cap L^{1+(1/n)}(\Omega);\
	      \norm{f}_{1+(1/n)}\le n\bigr\}.]
	      \end{equation*}
	\item 证明存在常数 $C$, 使得
	      \begin{equation*}
	      \norm{f}_p\le C\norm{f}_1\quad\forall f\in X.
	      \end{equation*}
\end{enumerate}
\end{exercise}

\begin{exercise}[Jensen 不等式]
\label{exer:4.9}
假设 $\abs{\Omega}<\infty$.设 $j:\R\to(-\infty,+\infty]$ 为凸的, 下半连续的
函数, $j\not\equiv+\infty$.设 $f\in L^1(\Omega)$ 满足 $f(x)\in D(j)$ 几乎处处,
且 $j(f)\in L^1(\Omega)$.证明
\begin{equation*}
j\Bigl(\frac1{\abs{\Omega}}\int_\Omega f\Bigr)
\le\frac1{\abs{\Omega}}\int_\Omega j(f).
\end{equation*}
\end{exercise}

\begin{exercise}[凸被积函数]
\label{exer:4.10}
假设 $\abs{\Omega}<\infty$.设 $1\le p<\infty$, $j:\R\to\R$ 为凸的连续函数.
考虑由
\begin{equation*}
J(u)=\begin{cases}
\displaystyle\int_\Omega j(u(x))\,dx & \text{若}\ j(u)\in L^1(\Omega),\\[8pt]
+\infty & \text{若}\ j(u)\notin L^1(\Omega)
\end{cases}
\end{equation*}
定义的函数 $J:L^p(\Omega)\to(-\infty,+\infty]$.
\begin{enumerate}[label=\arabic*., leftmargin=2.2em, itemsep=0.35em]
	\item 证明 $J$ 是凸的.
	\item 证明 $J$ 下半连续.
	      [提示: 先考虑 $j\ge0$ 的情形, 并使用 Fatou 引理.]
	\item 证明共轭函数 $J^\star:L^{p'}(\Omega)\to(-\infty,+\infty]$ 由
	      \begin{equation*}
	      J^\star(f)=\begin{cases}
	      \displaystyle\int_\Omega j^\star(f(x))\,dx & \text{若}\ j^\star(f)\in L^1(\Omega),\\[8pt]
	      +\infty & \text{若}\ j^\star(f)\notin L^1(\Omega)
	      \end{cases}
	      \end{equation*}
	      给出.
	      [提示: 当 $1<p<\infty$ 时考虑
	      $J_n(u)=J(u)+\frac1n\int\abs{u}^p$, 并求出 $J_n^\star$.]
	\item 用 $\partial j$ (相应地, $\partial J$) 表示 $j$ (相应地, $J$) 的次微分
	      (见问题 2).设 $u\in L^p(\Omega)$, $f\in L^{p'}(\Omega)$; 证明
	      \begin{equation*}
	      f\in\partial J(u)\iff f(x)\in\partial j(u(x))\quad\text{a.e. 在}\
	      \Omega\ \text{上}.
	      \end{equation*}
\end{enumerate}
\end{exercise}

\begin{exercise}[空间 $L^\alpha(\Omega)$, $0<\alpha<1$]
\label{exer:4.11}
设 $0<\alpha<1$.令
\begin{equation*}
L^\alpha(\Omega)=\bigl\{u:\Omega\to\R;\ u\ \text{可测且}\
\abs{u}^\alpha\in L^1(\Omega)\bigr\}
\end{equation*}
并置
\begin{equation*}
[u]_\alpha=\Bigl(\int\abs{u}^\alpha\Bigr)^{1/\alpha}.
\end{equation*}
\begin{enumerate}[label=\arabic*., leftmargin=2.2em, itemsep=0.35em]
	\item 验证 $L^\alpha$ 是向量空间, 但 $[\cdot]_\alpha$ 不是范数.更确切地说,
	      证明: 若 $u,v\in L^\alpha(\Omega)$, $u\ge0$ 几乎处处, $v\ge0$ 几乎处处,
	      则
	      \begin{equation*}
	      [u+v]_\alpha\ge[u]_\alpha+[v]_\alpha.
	      \end{equation*}
	\item 证明
	      \begin{equation*}
	      [u+v]_\alpha^\alpha\le[u]_\alpha^\alpha+[v]_\alpha^\alpha
	      \quad\forall u,v\in L^\alpha(\Omega).
	      \end{equation*}
\end{enumerate}
\end{exercise}

\begin{exercise}[$1<p\le2$ 时 $L^p$ 一致凸 (C. Morawetz 的方法)]
\label{exer:4.12}
\begin{enumerate}[label=\arabic*., leftmargin=2.2em, itemsep=0.35em]
	\item 设 $1<p<\infty$.证明存在只依赖于 $p$ 的常数 $C$, 使得
	      \begin{equation*}
	      \abs{a-b}^p\le C\bigl(\abs{a}^p+\abs{b}^p\bigr)^{1-s}
	      \Bigl(\abs{a}^p+\abs{b}^p-2\Bigl|\frac{a+b}2\Bigr|^p\Bigr)^s
	      \quad\forall a,b\in\R,
	      \end{equation*}
	      其中 $s=p/2$.
	\item 由此推出 $1<p\le2$ 时 $L^p(\Omega)$ 一致凸.
	      [提示: 使用第 1 问与 Hölder 不等式.]
\end{enumerate}
\end{exercise}

\begin{exercise}
\label{exer:4.13}
\begin{enumerate}[label=\arabic*., leftmargin=2.2em, itemsep=0.35em]
	\item 验证
	      \begin{equation*}
	      \bigl|\abs{a+b}-\abs{a}-\abs{b}\bigr|\le2\abs{b}\quad\forall a,b\in\R.
	      \end{equation*}
	\item 设 $(f_n)$ 为 $L^1(\Omega)$ 中的序列, 满足
	      \begin{enumerate}[label=(\roman*), leftmargin=2.4em, itemsep=0.1em]
	      	\item $f_n(x)\to f(x)$ 几乎处处,
	      	\item $(f_n)$ 在 $L^1(\Omega)$ 中有界, 即 $\norm{f_n}_1\le M\ \forall n$.
	      \end{enumerate}
	      证明 $f\in L^1(\Omega)$, 且
	      \begin{equation*}
	      \lim_{n\to\infty}\int\bigl\{\abs{f_n}-\abs{f_n-f}\bigr\}=\int\abs{f}.
	      \end{equation*}
	      [提示: 在第 1 问中取 $a=f_n-f$, $b=f$, 并考虑序列
	      $\varphi_n=\bigl|\abs{f_n}-\abs{f_n-f}-\abs{f}\bigr|$.]
	\item 设 $(f_n)$ 为 $L^1(\Omega)$ 中的序列, $f$ 为 $L^1(\Omega)$ 中的函数, 满足
	      \begin{enumerate}[label=(\roman*), leftmargin=2.4em, itemsep=0.1em]
	      	\item $f_n(x)\to f(x)$ 几乎处处,
	      	\item $\norm{f_n}_1\to\norm{f}$.
	      \end{enumerate}
	      证明 $\norm{f_n-f}_1=0$.
\end{enumerate}
\end{exercise}

\begin{exercise}[Egorov 与 Vitali 定理]
\label{exer:4.14}
假设 $\abs{\Omega}<\infty$.设 $(f_n)$ 为可测函数序列, 使得 $f_n\to f$ 几乎处处
(其中 $\abs{f}<\infty$ 几乎处处).
\begin{enumerate}[label=\arabic*., leftmargin=2.2em, itemsep=0.35em]
	\item 设 $\alpha>0$ 固定.证明
	      \begin{equation*}
	      \operatorname{meas}\bigl[\abs{f_n-f}>\alpha\bigr]
	      \xrightarrow[n\to\infty]{}0.
	      \end{equation*}
	\item 更确切地说, 令
	      \begin{equation*}
	      S_n(\alpha)=\bigcup_{k\ge n}\bigl[\abs{f_k-f}>\alpha\bigr].
	      \end{equation*}
	      证明 $\abs{S_n(\alpha)}\xrightarrow[n\to\infty]{}0$.
	\item (Egorov). 证明
	      \begin{equation*}
	      \begin{cases}
	      \forall\delta>0\ \ \exists A\subset\Omega\ \text{可测, 使得}\\
	      \abs{A}<\delta\ \text{且}\ f_n\to f\ \text{在}\ \Omega\setminus A\
	      \text{上一致收敛}.
	      \end{cases}
	      \end{equation*}
	      [提示: 给定整数 $m\ge1$, 借助第 2 问证明存在可测集
	      $\Omega_m\subset\Omega$, 使 $\abs{\Omega_m}<\delta/2^m$, 且存在整数
	      $N_m$ 使得
	      \begin{equation*}
	      \abs{f_k(x)-f(x)}<\frac1m\quad\forall k\ge N_m,\
	      \forall x\in\Omega\setminus\Omega_m.]
	      \end{equation*}
	\item (Vitali). 设 $(f_n)$ 为 $L^p(\Omega)$ 中的序列, $1\le p<\infty$.假设
	      \begin{enumerate}[label=(\roman*), leftmargin=2.4em, itemsep=0.1em]
	      	\item $\forall\varepsilon>0\ \exists\delta>0$ 使得
	      	      $\int_A\abs{f_n}^p<\varepsilon\ \forall n$ 且对一切满足
	      	      $\abs{A}<\delta$ 的可测集 $A\subset\Omega$;
	      	\item $f_n\to f$ 几乎处处.
	      \end{enumerate}
	      证明 $f\in L^p(\Omega)$, 且 $f_n\to f$ 在 $L^p(\Omega)$ 中.
\end{enumerate}
\end{exercise}

\begin{exercise}
\label{exer:4.15}
设 $\Omega=(0,1)$.
\begin{enumerate}[label=\arabic*., leftmargin=2.2em, itemsep=0.35em]
	\item 考虑由 $f_n(x)=ne^{-nx}$ 定义的函数序列 $(f_n)$.证明
	      \begin{enumerate}[label=(\roman*), leftmargin=2.4em, itemsep=0.1em]
	      	\item $f_n\to0$ 几乎处处;
	      	\item $f_n$ 在 $L^1(\Omega)$ 中有界;
	      	\item $f_n\not\to0$ 在 $L^1(\Omega)$ 中强收敛;
	      	\item $f_n\not\weakto0$ 弱收敛于 $\sigma(L^1,L^\infty)$.
	      \end{enumerate}
	      更确切地说, 不存在弱收敛于 $\sigma(L^1,L^\infty)$ 的子列.
	\item 设 $1<p<\infty$, 考虑由 $g_n(x)=n^{1/p}e^{-nx}$ 定义的函数序列
	      $(g_n)$.证明
	      \begin{enumerate}[label=(\roman*), leftmargin=2.4em, itemsep=0.1em]
	      	\item $g_n\to0$ 几乎处处;
	      	\item $(g_n)$ 在 $L^p(\Omega)$ 中有界;
	      	\item $g_n\not\to0$ 在 $L^p(\Omega)$ 中强收敛;
	      	\item $g_n\weakto0$ 弱收敛于 $\sigma(L^p,L^{p'})$.
	      \end{enumerate}
\end{enumerate}
\end{exercise}

\begin{exercise}
\label{exer:4.16}
设 $1<p<\infty$.设 $(f_n)$ 为 $L^p(\Omega)$ 中的序列, 满足
\begin{enumerate}[label=(\roman*), leftmargin=2.4em, itemsep=0.1em]
	\item $f_n$ 在 $L^p(\Omega)$ 中有界,
	\item $f_n\to f$ 在 $\Omega$ 上几乎处处.
\end{enumerate}
\begin{enumerate}[label=\arabic*., leftmargin=2.2em, itemsep=0.35em]
	\item 证明 $f_n\weakto f$ 弱收敛于 $\sigma(L^p,L^{p'})$.
	      [提示: 先证明: 若 $f_n\weakto\tilde f$ 弱收敛于
	      $\sigma(L^p,L^{p'})$ 且 $f_n\to f$ 几乎处处, 则 $f=\tilde f$ 几乎处处
	      (使用习题  \ref{exer:3.4}).]
	\item 若把假设 (ii) 换成
	      \begin{enumerate}[label=(\roman*)$'$, leftmargin=3.0em, itemsep=0.1em]
	      	\item $\norm{f_n-f}_1\to0$,
	      \end{enumerate}
	      结论相同.
	\item 现在假设 (i), (ii) 成立, 且 $\abs{\Omega}<\infty$.证明对每个满足
	      $1\le q<p$ 的 $q$ 都有 $\norm{f_n-f}_q\to0$.
	      [提示: 引入截断函数 $T_kf_n$, 或使用 Egorov 定理.]
\end{enumerate}
\end{exercise}

\begin{exercise}[Brezis--Lieb 引理]
\label{exer:4.17}
设 $1<p<\infty$.
\begin{enumerate}[label=\arabic*., leftmargin=2.2em, itemsep=0.35em]
	\item 证明存在常数 $C$ (依赖于 $p$), 使得
	      \begin{equation*}
	      \bigl|\abs{a+b}^p-\abs{a}^p-\abs{b}^p\bigr|
	      \le C\bigl(\abs{a}^{p-1}\abs{b}+\abs{a}\abs{b}^{p-1}\bigr)
	      \quad\forall a,b\in\R.
	      \end{equation*}
	\item 设 $(f_n)$ 为 $L^p(\Omega)$ 中的有界序列, 使得 $f_n\to f$ 在 $\Omega$ 上
	      几乎处处.证明 $f\in L^p(\Omega)$, 且
	      \begin{equation*}
	      \lim_{n\to\infty}\int_\Omega\bigl\{\abs{f_n}^p-\abs{f_n-f}^p\bigr\}
	      =\int_\Omega\abs{f}^p.
	      \end{equation*}
	      [提示: 在第 1 问中取 $a=f_n-f$, $b=f$.注意由习题  \ref{exer:4.16},
	      $\abs{f_n-f}\weakto0$ 在 $L^p$ 中弱收敛, 而 $\abs{f_n-f}^{p-1}
	      \weakto0$ 在 $L^{p'}$ 中弱收敛.]
	\item 由此推出: 若 $(f_n)$ 为 $L^p(\Omega)$ 中的序列, 满足
	      \begin{enumerate}[label=(\roman*), leftmargin=2.4em, itemsep=0.1em]
	      	\item $f_n(x)\to f(x)$ 几乎处处,
	      	\item $\norm{f_n}_p\to\norm{f}_p$,
	      \end{enumerate}
	      则 $\norm{f_n-f}_p\to0$.
	\item 为第 3 问找出另一种方法.
\end{enumerate}
\end{exercise}

\begin{exercise}[Rademacher 函数]
\label{exer:4.18}
设 $1\le p\le\infty$, $f\in L^p_{\mathrm{loc}}(\R)$.假设 $f$ 是 $T$-周期的,
即 $f(x+T)=f(x)$ 对几乎所有的 $x\in\R$ 成立.令
\begin{equation*}
\bar f=\frac1T\int_0^Tf(t)\,dt.
\end{equation*}
考虑由
\begin{equation*}
u_n(x)=f(nx),\quad x\in(0,1)
\end{equation*}
定义的 $L^p(0,1)$ 中的序列 $(u_n)$.
\begin{enumerate}[label=\arabic*., leftmargin=2.2em, itemsep=0.35em]
	\item 证明 $u_n\weakto\bar f$ 在 $L^p(0,1)$ 中关于拓扑 $\sigma(L^p,L^{p'})$
	      弱收敛.
	\item 求 $\lim_{n\to\infty}\norm{u_n-\bar f}_p$.
	\item 考察下述例子:
	      \begin{enumerate}[label=(\roman*), leftmargin=2.4em, itemsep=0.1em]
	      	\item $u_n(x)=\sin nx$,
	      	\item $u_n(x)=f(nx)$, 其中 $f$ 是 1-周期的, 且
	      	      \begin{equation*}
	      	      f(x)=\begin{cases}
	      	      \alpha & \text{若}\ x\in(0,1/2),\\
	      	      \beta & \text{若}\ x\in(1/2,1).
	      	      \end{cases}
	      	      \end{equation*}
	      \end{enumerate}
	      例 (ii) 中的函数称为 \emph{Rademacher 函数}.
\end{enumerate}
\end{exercise}

\begin{exercise}
\label{exer:4.19}
\begin{enumerate}[label=\arabic*., leftmargin=2.2em, itemsep=0.35em]
	\item 设 $(f_n)$ 为 $L^p(\Omega)$ 中的序列, $1<p<\infty$, $f\in L^p(\Omega)$.
	      假设
	      \begin{enumerate}[label=(\roman*), leftmargin=2.4em, itemsep=0.1em]
	      	\item $f_n\weakto f$ 弱收敛于 $\sigma(L^p,L^{p'})$,
	      	\item $\norm{f_n}_p\to\norm{f}_p$.
	      \end{enumerate}
	      证明 $f_n\to f$ 在 $L^p(\Omega)$ 中强收敛.
	\item 构造 $L^1(0,1)$ 中的序列 $(f_n)$, $f_n\ge0$, 使得:
	      \begin{enumerate}[label=(\roman*), leftmargin=2.4em, itemsep=0.1em]
	      	\item $f_n\weakto f$ 弱收敛于 $\sigma(L^1,L^\infty)$,
	      	\item $\norm{f_n}_1\to\norm{f}_1$,
	      	\item $\norm{f_n-f}_1\not\to0$.
	      \end{enumerate}
	      与习题  \ref{exer:4.13} 的结果以及命题  \ref{prop:3.32} 作比较.
\end{enumerate}
\end{exercise}

\begin{exercise}
\label{exer:4.20}
假设 $\abs{\Omega}<\infty$.设 $1\le p<\infty$, $1\le q<\infty$.

设 $a:\R\to\R$ 为连续函数, 使得
\begin{equation*}
\abs{a(t)}\le C\bigl\{\abs{t}^{p/q}+1\bigr\}\quad\forall t\in\R.
\end{equation*}
考虑由
\begin{equation*}
(Au)(x)=a(u(x)),\quad x\in\Omega
\end{equation*}
定义的 (非线性) 映射 $A:L^p(\Omega)\to L^q(\Omega)$.
\begin{enumerate}[label=\arabic*., leftmargin=2.2em, itemsep=0.35em]
	\item 证明 $A$ 从 $L^p(\Omega)$ 强拓扑到 $L^q(\Omega)$ 强拓扑连续.
	\item 取 $\Omega=(0,1)$, 并假设对每个满足 $u_n\weakto u$ 弱收敛于
	      $\sigma(L^p,L^{p'})$ 的序列 $(u_n)$ 都有 $Au_n\weakto Au$ 弱收敛于
	      $\sigma(L^q,L^{q'})$.证明 $a$ 是仿射函数.
	      [提示: 使用 Rademacher 函数; 见习题  \ref{exer:4.18}.]
\end{enumerate}
\end{exercise}

\begin{exercise}
\label{exer:4.21}
给定函数 $u_0:\R\to\R$, 令 $u_n(x)=u_0(x+n)$.
\begin{enumerate}[label=\arabic*., leftmargin=2.2em, itemsep=0.35em]
	\item 假设 $u_0\in L^p(\R)$, $1<p<\infty$.证明 $u_n\weakto0$ 在 $L^p(\R)$
	      中关于弱拓扑 $\sigma(L^p,L^{p'})$ 弱收敛.
	\item 假设 $u_0\in L^\infty(\R)$, 且 $u_0(x)\to0$ 当 $\abs{x}\to\infty$ 在下述
	      弱意义下成立: 对每个 $\delta>0$, 集合 $[\abs{u_0}>\delta]$ 的测度有限.
	      证明 $u_n\overset{\star}{\rightharpoonup}0$ 在 $L^\infty(\R)$ 中关于
	      \wstar{}拓扑 $\sigma(L^\infty,L^1)$ 弱收敛.
	\item 取 $u_0=\chi_{(0,1)}$.证明不存在子列 $(u_{n_k})$ 在 $L^1(\R)$ 中关于
	      $\sigma(L^1,L^\infty)$ 收敛.
\end{enumerate}
\end{exercise}

\begin{exercise}
\label{exer:4.22}
\begin{enumerate}[label=\arabic*., leftmargin=2.2em, itemsep=0.35em]
	\item 设 $(f_n)$ 为 $L^p(\Omega)$ 中的序列, $1<p\le\infty$, $f\in L^p(\Omega)$.
	      证明下述性质等价:
	      \begin{equation*}
	      \text{(A)}\ \ f_n\weakto f\ \text{在}\ \sigma(L^p,L^{p'})\ \text{中},
	      \end{equation*}
	      \begin{equation*}
	      \text{(B)}\ \begin{cases}
	      \norm{f_n}_p\le C\quad\text{且}\\
	      \displaystyle\int_Ef_n\to\int_Ef\quad\forall E\subset\Omega\ \text{可测, 且}\
	      \abs{E}<\infty.
	      \end{cases}
	      \end{equation*}
	\item 若 $p=1$ 且 $\abs{\Omega}<\infty$, 证明 (A) $\iff$ (B).
	\item 假设 $p=1$ 且 $\abs{\Omega}=\infty$.证明 (A) $\Rightarrow$ (B).
	      构造一个例子说明一般地 (B) $\not\Rightarrow$ (A).
	      [提示: 使用习题  \ref{exer:4.21} 第 3 问.]
	\item 设 $(f_n)$ 为 $L^1(\Omega)$ 中的序列, $f\in L^1(\Omega)$, 且
	      $\abs{\Omega}=\infty$.假设
	      \begin{enumerate}[label=(\alph*), leftmargin=2.4em, itemsep=0.1em]
	      	\item $f_n\ge0\ \forall n$, 且 $f\ge0$ 在 $\Omega$ 上几乎处处,
	      	\item $\int_\Omega f_n\to\int_\Omega f$,
	      	\item $\int_Ef_n\to\int_Ef\ \forall E\subset\Omega$ 可测, 且
	      	      $\abs{E}<\infty$.
	      \end{enumerate}
	      证明 $f_n\weakto f$ 在 $L^1(\Omega)$ 中弱收敛于 $\sigma(L^1,L^\infty)$.
	      [提示: 证明 $\int_Ff_n\to\int_Ff\ \forall F\subset\Omega$ 可测, 且
	      $\abs{F}\le\infty$.]
\end{enumerate}
\end{exercise}

\begin{exercise}
\label{exer:4.23}
设 $f:\Omega\to\R$ 为可测函数, $1\le p\le\infty$.本习题的目的是证明集合
\begin{equation*}
C=\bigl\{u\in L^p(\Omega);\ u\ge f\ \text{a.e.}\bigr\}
\end{equation*}
在 $L^p(\Omega)$ 中关于拓扑 $\sigma(L^p,L^{p'})$ 闭.
\begin{enumerate}[label=\arabic*., leftmargin=2.2em, itemsep=0.35em]
	\item 先假设 $1\le p<\infty$.证明 $C$ 在强 $L^p$ 拓扑中凸且闭.由此推出
	      $C$ 在 $\sigma(L^p,L^{p'})$ 中闭.
	\item 取 $p=\infty$, 证明
	      \begin{equation*}
	      C=\Bigl\{u\in L^\infty(\Omega)\ \Big|\ \int u\phi\ge\int f\phi\
	      \forall\phi\in L^1(\Omega)\ \text{满足}\ f\phi\in L^1(\Omega)\
	      \text{且}\ \phi\ge0\ \text{a.e.}\Bigr\}.
	      \end{equation*}
	      [提示: 先假设 $f\in L^\infty(\Omega)$; 一般情形下引入集合
	      $\omega_n=[\abs{f}<n]$.]
	\item 由此推出: 当 $p=\infty$ 时, $C$ 在 $\sigma(L^\infty,L^1)$ 中闭.
	\item 设 $f_1,f_2\in L^\infty(\Omega)$, $f_1\le f_2$ 几乎处处.证明集合
	      \begin{equation*}
	      C=\bigl\{u\in L^\infty(\Omega);\ f_1\le u\le f_2\ \text{a.e.}\bigr\}
	      \end{equation*}
	      在 $L^\infty(\Omega)$ 中关于拓扑 $\sigma(L^\infty,L^1)$ 紧.
\end{enumerate}
\end{exercise}

\begin{exercise}
\label{exer:4.24}
设 $u\in L^\infty(\R^N)$.设 $(\rho_n)$ 为磨光核序列.设 $(\zeta_n)$ 为
$L^\infty(\R^N)$ 中的序列, 使得
\begin{equation*}
\norm{\zeta_n}_\infty\le1\quad\forall n
\quad\text{且}\quad\zeta_n\to\zeta\ \text{a.e. 在}\ \R^N\ \text{上}.
\end{equation*}
令
\begin{equation*}
v_n=\rho_n\star(\zeta_nu)\quad\text{且}\quad v=\zeta u.
\end{equation*}
\begin{enumerate}[label=\arabic*., leftmargin=2.2em, itemsep=0.35em]
	\item 证明 $v_n\overset{\star}{\rightharpoonup}v$ 在 $L^\infty(\R^N)$ 中关于
	      \wstar{}拓扑 $\sigma(L^\infty,L^1)$ 弱收敛.
	\item 证明对每个球 $B$ 都有 $\int_B\abs{v_n-v}\to0$.
\end{enumerate}
\end{exercise}

\begin{exercise}[$L^\infty(\Omega)$ 中函数的正则化]
\label{exer:4.25}
设 $\Omega\subset\R^N$ 为开集.
\begin{enumerate}[label=\arabic*., leftmargin=2.2em, itemsep=0.35em]
	\item 设 $u\in L^\infty(\Omega)$.证明存在 $C_c^\infty(\Omega)$ 中的序列
	      $(u_n)$, 使得
	      \begin{enumerate}[label=(\alph*), leftmargin=2.4em, itemsep=0.1em]
	      	\item $\norm{u_n}_\infty\le\norm{u}_\infty\ \forall n$,
	      	\item $u_n\to u$ 在 $\Omega$ 上几乎处处,
	      	\item $u_n\overset{\star}{\rightharpoonup}u$ 在 $L^\infty(\Omega)$ 中
	      	      关于\wstar{}拓扑 $\sigma(L^\infty,L^1)$ 弱收敛.
	      \end{enumerate}
	\item 若 $u\ge0$ 在 $\Omega$ 上几乎处处, 证明还可以要求
	      \begin{enumerate}[label=(\alph*), leftmargin=2.4em, start=4, itemsep=0.1em]
	      	\item $u_n\ge0$ 在 $\Omega$ 上 $\forall n$.
	      \end{enumerate}
	\item 由此推出 $C_c^\infty(\Omega)$ 在 $L^\infty(\Omega)$ 中关于拓扑
	      $\sigma(L^\infty,L^1)$ 稠密.
\end{enumerate}
\end{exercise}

\begin{exercise}
\label{exer:4.26}
设 $\Omega\subset\R^N$ 为开集, $f\in L^1_{\mathrm{loc}}(\Omega)$.
\begin{enumerate}[label=\arabic*., leftmargin=2.2em, itemsep=0.35em]
	\item 证明 $f\in L^1(\Omega)$ 当且仅当
	      \begin{equation*}
	      A=\sup\Bigl\{\int f\phi;\ \phi\in C_c(\Omega),\
	      \norm{\phi}_\infty\le1\Bigr\}<\infty.
	      \end{equation*}
	      若 $f\in L^1(\Omega)$, 证明 $A=\norm{f}_1$.
	\item 证明 $f^+\in L^1(\Omega)$ 当且仅当
	      \begin{equation*}
	      B=\sup\Bigl\{\int f\phi;\ \phi\in C_c(\Omega),\
	      \norm{\phi}_\infty\le1\ \text{且}\ \phi\ge0\Bigr\}<\infty.
	      \end{equation*}
	      若 $f^+\in L^1(\Omega)$, 证明 $B=\norm{f^+}_1$.
	\item 把 $C_c(\Omega)$ 换成 $C_c^\infty(\Omega)$ 时, 同样的问题.
	\item 由此推出
	      \begin{equation*}
	      \Bigl\{\int f\phi=0\ \forall\phi\in C_c^\infty(\Omega)\Bigr\}
	      \Longrightarrow\bigl[f=0\ \text{a.e.}\bigr]
	      \end{equation*}
	      以及
	      \begin{equation*}
	      \Bigl\{\int f\phi\ge0\ \forall\phi\in C_c^\infty(\Omega),\
	      \phi\ge0\Bigr\}\Longrightarrow\bigl[f\ge0\ \text{a.e.}\bigr].
	      \end{equation*}
\end{enumerate}
\end{exercise}

\begin{exercise}
\label{exer:4.27}
设 $\Omega\subset\R^N$ 为开集.设 $u,v\in L^1_{\mathrm{loc}}(\Omega)$, 且 $u\ne0$
在一个测度为正的集合上几乎处处成立.假设
\begin{equation*}
\Bigl\{\phi\in C_c^\infty(\Omega)\ \text{且}\ \int u\phi>0\Bigr\}
\Longrightarrow\Bigl\{\int v\phi\ge0\Bigr\}.
\end{equation*}
证明存在常数 $\lambda\ge0$, 使得 $v=\lambda u$.
\end{exercise}

\begin{exercise}
\label{exer:4.28}
设 $\rho\in L^1(\R^N)$, $\int\rho=1$.令 $\rho_n(x)=n^N\rho(nx)$.设
$f\in L^p(\R^N)$, $1\le p<\infty$.证明 $\rho_n\star f\to f$ 在 $L^p(\R^N)$ 中.
\end{exercise}

\begin{exercise}
\label{exer:4.29}
设 $K\subset\R^N$ 为紧子集.证明存在 $C_c^\infty(\R^N)$ 中的函数序列 $(u_n)$,
使得
\begin{enumerate}[label=(\alph*), leftmargin=2.4em, itemsep=0.15em]
	\item $0\le u_n\le1$ 在 $\R^N$ 上,
	\item $u_n=1$ 在 $K$ 上,
	\item $\operatorname{supp}u_n\subset K+B(0,1/n)$,
	\item $\abs{D^\alpha u_n(x)}\le C_\alpha n^{\abs{\alpha}}\ \forall x\in\R^N$,
	      对每个多重指标 $\alpha$ (其中 $C_\alpha$ 只依赖于 $\alpha$ 而不依赖于
	      $n$).
\end{enumerate}
[提示: 设 $\chi_n$ 为 $K+B(0,1/2n)$ 的特征函数; 取
$u_n=\rho_{2n}\star\chi_n$.]
\end{exercise}

\begin{exercise}[Young 不等式]
\label{exer:4.30}
设 $1\le p\le\infty$, $1\le q\le\infty$, 满足 $\frac1p+\frac1q\ge1$.令
$\frac1r=\frac1p+\frac1q-1$, 于是 $1\le r\le\infty$.设
$f\in L^p(\R^N)$, $g\in L^q(\R^N)$.
\begin{enumerate}[label=\arabic*., leftmargin=2.2em, itemsep=0.35em]
	\item 证明对几乎所有的 $x\in\R^N$, 函数 $y\mapsto f(x-y)g(y)$ 在 $\R^N$ 上
	      可积.
	      [提示: 令 $\alpha=p/q'$, $\beta=q/p'$, 并写
	      \begin{equation*}
	      \abs{f(x-y)g(y)}=\abs{f(x-y)}^\alpha\abs{g(y)}^\beta
	      \bigl[\abs{f(x-y)}^{1-\alpha}\abs{g(y)}^{1-\beta}\bigr].]
	      \end{equation*}
	\item 令
	      \begin{equation*}
	      (f\star g)(x)=\int_{\R^N}f(x-y)g(y)\,dy.
	      \end{equation*}
	      证明 $f\star g\in L^r(\R^N)$, 且 $\norm{f\star g}_r\le\norm{f}_p
	      \norm{g}_q$.
	\item 这里假设 $\frac1p+\frac1q=1$.证明
	      \begin{equation*}
	      f\star g\in C(\R^N)\cap L^\infty(\R^N),
	      \end{equation*}
	      并且, 若 $1<p<\infty$, 则当 $\abs{x}\to\infty$ 时
	      $(f\star g)(x)\to0$.
\end{enumerate}
\end{exercise}

\begin{exercise}
\label{exer:4.31}
设 $f\in L^p(\R^N)$, $1\le p<\infty$.对每个 $r>0$ 令
\begin{equation*}
f_r(x)=\frac1{\abs{B(x,r)}}\int_{B(x,r)}f(y)\,dy,\quad x\in\R^N.
\end{equation*}
\begin{enumerate}[label=\arabic*., leftmargin=2.2em, itemsep=0.35em]
	\item 证明 $f_r\in L^p(\R^N)\cap C(\R^N)$, 且当 $\abs{x}\to\infty$ 时
	      $f_r(x)\to0$ ($r$ 固定).
	\item 证明当 $r\to0$ 时 $f_r\to f$ 在 $L^p(\R^N)$ 中.
	      [提示: 把 $f_r$ 写成 $\phi_r\star f$ 的形式, 其中 $\phi_r$ 取法适当.]
\end{enumerate}
\end{exercise}

\begin{exercise}
\label{exer:4.32}
\begin{enumerate}[label=\arabic*., leftmargin=2.2em, itemsep=0.35em]
	\item 设 $f,g\in L^1(\R^N)$, $h\in L^p(\R^N)$, $1\le p\le\infty$.证明
	      $f\star g=g\star f$ 且 $(f\star g)\star h=f\star(g\star h)$.
	\item 设 $f\in L^1(\R^N)$.假设 $f\star\phi=0\ \forall\phi\in
	      C_c^\infty(\R^N)$.证明 $f=0$ 在 $\R^N$ 上几乎处处成立.对
	      $f\in L^1_{\mathrm{loc}}(\R^N)$ 也有同样的问题.
	\item 设 $a\in L^1(\R^N)$ 为固定函数.考虑由
	      \begin{equation*}
	      T_a(u)=a\star u
	      \end{equation*}
	      定义的算子 $T_a:L^2(\R^N)\to L^2(\R^N)$.验证 $T_a$ 有界, 且
	      $\norm{T_a}_{\mathcal{L}(L^2)}\le\norm{a}_{L^1(\R^N)}$.计算
	      $T_a\circ T_b$, 并证明 $T_a\circ T_b=T_b\circ T_a\
	      \forall a,b\in L^1(\R^N)$.确定 $(T_a)^\star$, $T_a\circ(T_a)^\star$
	      与 $(T_a)^\star\circ T_a$.在关于 $a$ 的什么条件下有
	      $(T_a)^\star=T_a$?
\end{enumerate}
\end{exercise}

\begin{exercise}
\label{exer:4.33}
固定函数 $\phi\in C_c(\R)$, $\phi\not\equiv0$, 并考虑函数族
\begin{equation*}
\mathcal{F}=\bigcup_{n=1}^\infty\{\phi_n\},
\end{equation*}
其中 $\phi_n(x)=\phi(x+n)$, $x\in\R$.
\begin{enumerate}[label=\arabic*., leftmargin=2.2em, itemsep=0.35em]
	\item 假设 $1\le p<\infty$.证明 $\forall\varepsilon>0\ \exists\delta>0$,
	      使得
	      \begin{equation*}
	      \norm{\tau_hf-f}_p<\varepsilon\quad\forall f\in\mathcal{F}\
	      \text{且}\ \forall h\in\R,\ \abs{h}<\delta.
	      \end{equation*}
	\item 证明 $\mathcal{F}$ 在 $L^p(\R)$ 中不具有紧闭包.
\end{enumerate}
\end{exercise}

\begin{exercise}
\label{exer:4.34}
设 $1\le p<\infty$, $\mathcal{F}\subset L^p(\R^N)$ 为 $L^p(\R^N)$ 的紧子集.
\begin{enumerate}[label=\arabic*., leftmargin=2.2em, itemsep=0.35em]
	\item 证明 $\mathcal{F}$ 在 $L^p(\R^N)$ 中有界.
	\item 证明 $\forall\varepsilon>0\ \exists\delta>0$, 使得
	      \begin{equation*}
	      \norm{\tau_hf-f}_p<\varepsilon\quad\forall f\in\mathcal{F}\
	      \text{且}\ \forall h\in\R^N,\ \abs{h}<\delta.
	      \end{equation*}
	\item 证明 $\forall\varepsilon>0\ \exists\Omega\subset\R^N$ 有界, 开, 使得
	      \begin{equation*}
	      \norm{f}_{L^p(\R^N\setminus\Omega)}<\varepsilon\quad\forall f\in\mathcal{F}.
	      \end{equation*}
	      与推论  \ref{cor:4.27} 作比较.
\end{enumerate}
\end{exercise}

\begin{exercise}
\label{exer:4.35}
固定函数 $G\in L^p(\R^N)$, $1\le p<\infty$, 并令 $\mathcal{F}=G\star B$, 其中
$B$ 是 $L^1(\R^N)$ 中的有界集.

证明对任一测度有限的可测集 $\Omega\subset\R^N$, $\mathcal{F}|_\Omega$ 在
$L^p(\Omega)$ 中的闭包紧.与推论  \ref{cor:4.28} 作比较.
\end{exercise}

\begin{exercise}[等度可积族]
\label{exer:4.36}
若子集 $\mathcal{F}\subset L^1(\Omega)$ 满足下述性质, 则称它是\emph{等度可积}
的:\footnote{可以证明: 若测度空间 $\Omega$ 是弥散的 (即 $\Omega$ 没有原子), 则
(a) 由 (b) 与 (c) 推出.例如考虑带 Lebesgue 测度的 $\Omega=\R^N$.}
\begin{enumerate}[label=(\alph*), leftmargin=2.4em, itemsep=0.15em]
	\item $\mathcal{F}$ 在 $L^1(\Omega)$ 中有界,
	\item $\begin{cases}
	      \forall\varepsilon>0\ \exists\delta>0\ \text{使得}\\
	      \displaystyle\int_E\abs{f}<\varepsilon\quad\forall f\in\mathcal{F},\
	      \forall E\subset\Omega\ \text{可测, 且}\ \abs{E}<\delta,
	      \end{cases}$
	\item $\begin{cases}
	      \forall\varepsilon>0\ \exists\omega\subset\Omega\ \text{可测, 且}\
	      \abs{\omega}<\infty\ \text{使得}\\
	      \displaystyle\int_{\Omega\setminus\omega}\abs{f}<\varepsilon\quad
	      \forall f\in\mathcal{F}.
	      \end{cases}$
\end{enumerate}
设 $(\Omega_n)$ 为 $\Omega$ 中的非减可测集序列, $\abs{\Omega_n}<\infty\ \forall n$,
且 $\Omega=\bigcup_n\Omega_n$.
\begin{enumerate}[label=\arabic*., leftmargin=2.2em, itemsep=0.35em]
	\item 证明 $\mathcal{F}$ 等度可积当且仅当
	      \begin{equation*}
	      \lim_{t\to\infty}\sup_{f\in\mathcal{F}}\int_{[\abs{f}>t]}\abs{f}=0
	      \end{equation*}
	      且
	      \begin{equation*}
	      \lim_{n\to\infty}\sup_{f\in\mathcal{F}}\int_{\Omega\setminus\Omega_n}
	      \abs{f}=0.
	      \end{equation*}
	\item 证明: 若 $\mathcal{F}\subset L^1(\Omega)$ 紧, 则 $\mathcal{F}$ 等度
	      可积.逆命题是否成立?
\end{enumerate}
\end{exercise}

\begin{exercise}
\label{exer:4.37}
固定函数 $f\in L^1(\R)$, 使得
\begin{equation*}
\int_{-\infty}^{+\infty}f(t)\,dt=0\quad\text{且}\quad
\int_0^{+\infty}f(t)\,dt>0,
\end{equation*}
并对 $x\in I=(-1,+1)$ 令 $u_n(x)=nf(nx)$.
\begin{enumerate}[label=\arabic*., leftmargin=2.2em, itemsep=0.35em]
	\item 证明
	      \begin{equation*}
	      \lim_{n\to\infty}\int_Iu_n(x)\phi(x)\,dx=0\quad
	      \forall\phi\in C([-1,+1]).
	      \end{equation*}
	\item 验证序列 $(u_n)$ 在 $L^1(I)$ 中有界.证明 $(u_n)$ 的任一子列都不等度
	      可积.
	\item 证明不存在函数 $u\in L^1(I)$, 使得沿某个子列 $(u_{n_k})$ 有
	      \begin{equation*}
	      \lim_{k\to\infty}\int_Iu_{n_k}(x)\phi(x)\,dx
	      =\int_Iu(x)\phi(x)\,dx\quad\forall\phi\in L^\infty(I).
	      \end{equation*}
	\item 与 Dunford--Pettis 定理作比较 (见问题 23 的 A3 问).
	\item 证明存在子列 $(u_{n_k})$, 使得当 $k\to\infty$ 时 $u_{n_k}(x)\to0$
	      在 $I$ 上几乎处处成立.
	      [提示: 计算 $\int_{[n^{-1/2}<\abs{x}<1]}\abs{u_n(x)}\,dx$, 并应用定理
	       \ref{thm:4.9}.]
\end{enumerate}
\end{exercise}

\begin{exercise}
\label{exer:4.38}
令 $I=(0,1)$, 并考虑由
\begin{equation*}
u_n(x)=\begin{cases}
n & \text{若}\ x\in\displaystyle\bigcup_{j=0}^{n-1}
\Bigl(\frac jn,\frac jn+\frac1{n^2}\Bigr),\\[10pt]
0 & \text{其他情形}
\end{cases}
\end{equation*}
定义的 $L^1(I)$ 中的函数序列 $(u_n)$.
\begin{enumerate}[label=\arabic*., leftmargin=2.2em, itemsep=0.35em]
	\item 验证 $\abs{\operatorname{supp}u_n}=\frac1n$, $\norm{u_n}_1=1$.
	\item 证明
	      \begin{equation*}
	      \lim_{n\to+\infty}\int_Iu_n(x)\phi(x)\,dx=\int_I\phi(x)\,dx
	      \quad\forall\phi\in C([0,1]).
	      \end{equation*}
	      [提示: 先考虑 $\phi\in C^1([0,1])$ 的情形.]
	\item 证明 $(u_n)$ 的任一子列都不等度可积.
	\item 证明不存在函数 $u\in L^1(I)$, 使得沿某个子列 $(u_{n_k})$ 有
	      \begin{equation*}
	      \lim_{k\to\infty}\int_Iu_{n_k}(x)\phi(x)\,dx
	      =\int_Iu(x)\phi(x)\,dx\quad\forall\phi\in L^\infty(I).
	      \end{equation*}
	      [提示: 取进一步的子列 $(u_{n'_k})$, 使
	      $\sum_k\abs{\operatorname{supp}u_{n'_k}}<1$.]
	\item 证明存在子列 $(u_{n_k})$, 使得当 $k\to\infty$ 时 $u_{n_k}(x)\to0$
	      在 $I$ 上几乎处处成立.
\end{enumerate}
\end{exercise}

% chapters/ch05.tex
% 第~
\ref{chap:5}~章 \index{Hilbertkongjian@Hilbert 空间 (Hilbert space)}
Hilbert 空间
% 源: 原书 PDF 第 140--165 页 (印刷页 131--156)

\chapter{Hilbert 空间}
\label{chap:5}

\section{定义与初等性质. \index{daotujishangdetouying@到凸集上的投影 (projection on convex set)}
到闭凸集上的投影}
\label{sec:5.1}

\begin{definition}
设 $H$ 为向量空间.\index{biaoliangji@标量积 (scalar product)}
\emph{标量积} $(u,v)$ 是 $H\times H$ 上取值于 $\R$ 的双线性
形式 (即由 $H\times H$ 到 $\R$ 的, 对两个变量都线性的映射), 使得
\begin{align*}
(u,v)&=(v,u) &&\forall u,v\in H\quad(\text{对称性}),\\
(u,u)&\ge0 &&\forall u\in H\quad(\text{正性}),\\
(u,u)&\ne0 &&\forall u\ne0\quad(\text{确定性}).
\end{align*}
\end{definition}

我们回忆, 标量积满足 Cauchy--Schwarz 不等式
\begin{equation*}
\abs{(u,v)}\le(u,u)^{1/2}(v,v)^{1/2}\quad\forall u,v\in H.
\end{equation*}
[记住下面一点有时是有用的: Cauchy--Schwarz 不等式的证明并不需要假设
$(u,u)\ne0\ \forall u\ne0$.] 由 Cauchy--Schwarz 不等式推出, 量
\begin{equation*}
\boxed{\abs{u}=(u,u)^{1/2}}
\end{equation*}
是一个范数——我们将常用 $\abs{\cdot}$ (而不是 $\norm{\cdot}$) 表示由标量积
导出的范数.事实上, 我们有
\begin{equation*}
\abs{u+v}^2=(u+v,u+v)=\abs{u}^2+(u,v)+(v,u)+\abs{v}^2
\le\abs{u}^2+2\abs{u}\abs{v}+\abs{v}^2,
\end{equation*}
从而 $\abs{u+v}\le\abs{u}+\abs{v}$.

我们回忆\index{pingxingsibianxingfaze@平行四边形法则 (parallelogram law)}
经典的\emph{平行四边形法则}:
\begin{equation}
\abs{\frac{a+b}{2}}^2+\abs{\frac{a-b}{2}}^2
=\frac12\bigl(\abs{a}^2+\abs{b}^2\bigr)\quad\forall a,b\in H.
\label{eq:5.1}
\end{equation}

\begin{definition}
\emph{Hilbert 空间}是赋以一个标量积的向量空间 $H$, 使得 $H$ 关于范数
$\abs{\cdot}$ 完备.
\end{definition}

下文 $H$ 总表示一个 Hilbert 空间.

\runin{基本例子.}
赋以标量积
\begin{equation*}
(u,v)=\int_\Omega u(x)v(x)\,d\mu
\end{equation*}
的 $L^2(\Omega)$ 是 Hilbert 空间.特别地, $\ell^2$ 是 Hilbert 空间.第~
\ref{chap:8}~章与
第~
\ref{chap:9}~章中研究的 Sobolev 空间 $H^1$ 是 Hilbert 空间的另一个例子; 它以
$L^2(\Omega)$ 为``模型''.

\begin{keyproposition}
\label{prop:5.1}
$H$ 一致凸, 从而自反.
\end{keyproposition}

\begin{proof}
设 $\varepsilon>0$, $u,v\in H$ 满足 $\abs{u}\le1$, $\abs{v}\le1$ 且
$\abs{u-v}>\varepsilon$.由平行四边形法则有
\begin{equation*}
\abs{\frac{u+v}{2}}^2<1-\frac{\varepsilon^2}{4}
\quad\text{从而}\quad
\abs{\frac{u+v}{2}}<1-\delta,
\quad\text{其中}\ \delta=1-\Bigl(1-\frac{\varepsilon^2}{4}\Bigr)^{1/2}>0.
\end{equation*}
\end{proof}

\begin{keytheorem}[到闭凸集上的投影]
\label{thm:5.2}
设 $K\subset H$ 为非空闭凸集.则对每个 $f\in H$ 存在唯一的元素 $u\in K$,
使得
\begin{equation}
\abs{f-u}=\min_{v\in K}\abs{f-v}=\dist(f,K).
\label{eq:5.2}
\end{equation}
此外, $u$ 由下述性质\emph{刻画}:
\begin{equation}
u\in K\quad\text{且}\quad(f-u,v-u)\le0\quad\forall v\in K.
\label{eq:5.3}
\end{equation}
\end{keytheorem}

\begin{notation}
上述元素 $u$ 称为 $f$ 在 $K$ 上的\emph{投影}, 记作
\begin{equation*}
\boxed{u=P_Kf.}
\end{equation*}
\end{notation}

不等式 \eqref{eq:5.3} 说的是: 向量 $\overrightarrow{uf}$ 与任一向量
$\overrightarrow{uv}$ ($v\in K$) 的标量积 $\le0$, 即这两个向量所成的角
$\theta\ge\pi/2$; 见图  \ref{fig:4}.

\begin{figure}[tbp]
	\centering
	\begin{tikzpicture}[scale=1.0]
		\draw[thick] (0,0) ellipse (2.4 and 1.2);
		\coordinate (u) at (1.62,0.88);
		\coordinate (v) at (-0.35,0.88);
		\coordinate (f) at (2.95,1.78);
		\draw[-{Stealth[length=2mm]}] (u) -- (v);
		\draw[-{Stealth[length=2mm]}] (u) -- (f);
		\fill (u) circle (1.3pt);
		\draw (u) ++(34:0.62) arc[start angle=34, end angle=180, radius=0.62];
		\node at ($(u)+(107:0.95)$) {$\theta$};
		\node at (-0.55,-0.35) {$K$};
		\node[right] at (u) {$u=P_Kf$};
		\node[above left] at (v) {$v$};
		\node[right] at (f) {$f$};
	\end{tikzpicture}
	\caption{}
	\label{fig:4}
\end{figure}

\begin{proof}
\emph{(a) 存在性.} 我们将给出两个不同的证明:

\begin{enumerate}[label=\arabic*., leftmargin=2.2em, itemsep=0.3em]
	\item 函数 $\varphi(v)=\abs{f-v}$ 凸, 连续, 且
	      $\lim_{\abs{v}\to\infty}\varphi(v)=+\infty$.由于 $H$ 自反, 由推论
	       \ref{cor:3.23} 推出 $\varphi$ 在 $K$ 上达到最小值.
	\item 第二个证明\emph{不}依赖自反空间与一致凸空间的理论.它是直接的论证.
	      设 $(v_n)$ 为 \eqref{eq:5.2} 的\emph{极小化序列}, 即 $v_n\in K$ 且
	      \begin{equation*}
	      d_n=\abs{f-v_n}\to d=\inf_{v\in K}\abs{f-v}.
	      \end{equation*}
	      我们断言 $(v_n)$ 是 Cauchy 序列.事实上, 对
	      $a=f-v_n$, $b=f-v_m$ 应用平行四边形法则, 得到
	      \begin{equation*}
	      \abs{f-\frac{v_n+v_m}{2}}^2+\abs{\frac{v_n-v_m}{2}}^2
	      =\frac12\bigl(d_n^2+d_m^2\bigr).
	      \end{equation*}
	      但 $\dfrac{v_n+v_m}{2}\in K$, 从而
	      $\abs{f-\frac{v_n+v_m}{2}}\ge d$.由此推出
	      \begin{equation*}
	      \abs{\frac{v_n-v_m}{2}}^2\le\frac12\bigl(d_n^2+d_m^2\bigr)-d^2
	      \quad\text{且}\quad\lim_{m,n\to\infty}\abs{v_n-v_m}=0.
	      \end{equation*}
	      因此序列 $(v_n)$ 收敛于某个极限 $u\in K$, 且 $d=\abs{f-u}$.
\end{enumerate}

\emph{(b) \eqref{eq:5.2} 与 \eqref{eq:5.3} 的等价性.} 假设 $u\in K$ 满足 \eqref{eq:5.2}, 并设 $w\in K$.我们有
\begin{equation*}
v=(1-t)u+tw\in K\quad\forall t\in[0,1],
\end{equation*}
从而
\begin{equation*}
\abs{f-u}\le\abs{f-[(1-t)u+tw]}=\abs{(f-u)-t(w-u)}.
\end{equation*}
因此
\begin{equation*}
\abs{f-u}^2\le\abs{f-u}^2-2t(f-u,w-u)+t^2\abs{w-u}^2,
\end{equation*}
这蕴含 $2(f-u,w-u)\le t\abs{w-u}^2\ \forall t\in(0,1]$.令 $t\to0$ 即得 \eqref{eq:5.3}.

反之, 假设 $u$ 满足 \eqref{eq:5.3}.则我们有
\begin{equation*}
\abs{u-f}^2-\abs{v-f}^2=2(f-u,v-u)-\abs{u-v}^2\le0\quad\forall v\in K;
\end{equation*}
这蕴含 \eqref{eq:5.2}.

\emph{(c) 唯一性.} 假设 $u_1$ 与 $u_2$ 满足 \eqref{eq:5.3}.我们有
\begin{align}
(f-u_1,v-u_1)&\le0\quad\forall v\in K,\label{eq:5.4}\\
(f-u_2,v-u_2)&\le0\quad\forall v\in K.\label{eq:5.5}
\end{align}
在 \eqref{eq:5.4} 中取 $v=u_2$, 在 \eqref{eq:5.5} 中取 $v=u_1$, 并把相应的两个不等式相加, 得到
$\abs{u_1-u_2}^2\le0$.
\end{proof}

\begin{remark}
\label{rem:5.1}
极小化问题与不等式组相联系, 这并不令人意外.让我们回忆一个熟知的例子.
设 $F:\R\to\R$ 为可微函数, 并设 $u\in[0,1]$ 是 $F$ 在 $[0,1]$ 上达到最小值的
点.则或者 $u\in(0,1)$ 且 $F'(u)=0$, 或者 $u=0$ 且 $F'(u)\le0$, 或者 $u=1$
且 $F'(u)=1$.这三种情形可概括为: $u\in[0,1]$ 且
$F'(u)(v-u)\le0\ \forall v\in[0,1]$; 另见习题 5.10.
% 原书此处印作 "u = 1 and F'(u) = 1", 与极小性不符 (应为 F'(u) >= 0 或
% F'(u) <= 0 的端情形). 两份 OCR 与 PDF 文本层一致, 属原书排印之误.
% 依技能规范: 保留原样并记录, 不擅自更正. 见《项目交接文档》"原书疑似错误".
\end{remark}

\begin{remark}
\label{rem:5.2}
设 $K\subset E$ 为一致凸 Banach 空间 $E$ 中的非空闭凸集.则对每个 $f\in E$
存在唯一的元素 $u\in E$, 使得
\begin{equation*}
\norm{f-u}=\min_{v\in K}\norm{f-v}=\dist(f,K);
\end{equation*}
见习题 3.32.
\end{remark}

\begin{proposition}
\label{prop:5.3}
设 $K\subset H$ 为非空闭凸集.则 $P_K$ 不增大距离, 即
\begin{equation*}
\abs{P_Kf_1-P_Kf_2}\le\abs{f_1-f_2}\quad\forall f_1,f_2\in H.
\end{equation*}
\end{proposition}

\begin{proof}
令 $u_1=P_Kf_1$, $u_2=P_Kf_2$.我们有
\begin{align}
(f_1-u_1,v-u_1)&\le0\quad\forall v\in K,\label{eq:5.6}\\
(f_2-u_2,v-u_2)&\le0\quad\forall v\in K.\label{eq:5.7}
\end{align}
在 \eqref{eq:5.6} 中取 $v=u_2$, 在 \eqref{eq:5.5} 中取 $v=u_1$, 并把相应的两个不等式相加, 得到
\begin{equation*}
\abs{u_1-u_2}^2\le(f_1-f_2,u_1-u_2).
\end{equation*}
由此推出 $\abs{u_1-u_2}\le\abs{f_1-f_2}$.
% 原书此处印作 "v = u1 in (5)", 应为 (7). 依技能规范保留原样并记录.
\end{proof}

\begin{corollary}
\label{cor:5.4}
假设 $M\subset H$ 为闭线性子空间.设 $f\in H$.则 $u=P_Mf$ 由下式刻画:
\begin{equation}
\boxed{u\in M\quad\text{且}\quad(f-u,v)=0\quad\forall v\in M.}
\label{eq:5.8}
\end{equation}
此外, $P_M$ 是线性算子, 称为\emph{正交投影}.
\end{corollary}

\begin{proof}
由 \eqref{eq:5.3} 我们有
\begin{equation*}
(f-u,v-u)\le0\quad\forall v\in M,
\end{equation*}
从而
\begin{equation*}
(f-u,tv-u)\le0\quad\forall v\in M,\ \forall t\in\R.
\end{equation*}
由此推出 \eqref{eq:5.8} 成立.

反之, 若 $u$ 满足 \eqref{eq:5.8}, 则我们有
\begin{equation*}
(f-u,v-u)=0\quad\forall v\in M.
\end{equation*}
$P_M$ 的线性性是显然的.
\end{proof}
\section{Hilbert 空间的对偶空间}
\label{sec:5.2}

在 Hilbert 空间中, 写出连续线性泛函是非常容易的.任取 $f\in H$; 则映射
$u\mapsto(f,u)$ 是 $H$ 上的连续线性泛函.一个值得注意的事实是: $H$ 上的
\emph{一切}连续线性泛函都以这种方式得到:

\index{RieszFrchetbiaoshidingli@Riesz--Fréchet 表示定理 (Riesz--Frechet representation theorem)}
\begin{keytheorem}[Riesz--Fréchet 表示定理]
\label{thm:5.5}
给定任一 $\varphi\in H^\star$, 存在唯一的 $f\in H$, 使得
\begin{equation*}
\dual{\varphi}{u}=(f,u)\quad\forall u\in H.
\end{equation*}
此外,
\begin{equation*}
\abs{f}=\norm{\varphi}_{H^\star}.
\end{equation*}
\end{keytheorem}

\begin{proof}
我们再一次给出两个证明:

\begin{enumerate}[label=\arabic*., leftmargin=2.2em, itemsep=0.3em]
	\item 第一个证明与定理  \ref{thm:4.11} 的证明几乎完全相同.考虑如下定义的
	      映射 $T:H\to H^\star$: 给定任一 $f\in H$, 映射 $u\mapsto(f,u)$ 是
	      $H$ 上的连续线性泛函.它定义了 $H^\star$ 中的一个元素, 记作 $Tf$,
	      使得
	      \begin{equation*}
	      \dual{Tf}{u}=(f,u)\quad\forall u\in H.
	      \end{equation*}
	      显然 $\norm{Tf}_{H^\star}=\abs{f}$.于是 $T$ 是由 $H$ \emph{到}
	      $T(H)$ (它是 $H^\star$ 的闭子空间) 上的线性等距.为完成证明, 只须
	      证明 $T(H)$ 在 $H^\star$ 中稠密.假设 $h$ 是 $H^\star$ 上的连续线性
	      泛函, 且它在 $T(H)$ 上为零.由于 $H$ 自反, $h$ 属于 $H$ 且满足
	      $\dual{Tf}{h}=0\ \forall f\in H$.由此推出 $(f,h)=0\ \forall f\in H$,
	      从而 $h=0$.
	\item 第二个证明是更直接的论证, 它避免了任何对自反性的使用.令
	      $M=\varphi^{-1}(\{0\})$, 于是 $M$ 是 $H$ 的闭子空间.我们总可以假设
	      $M\ne H$ (否则 $\varphi\equiv0$, 定理  \ref{thm:5.5} 的结论是显然的
	      ——只须取 $f=0$).我们断言: 存在某个元素 $g\in H$, 使得
	      \begin{equation*}
	      \abs{g}=1\quad\text{且}\quad(g,v)=0\quad\forall v\in M
	      \quad(\text{从而}\ g\notin M).
	      \end{equation*}
	      事实上, 取 $g_0\in H$ 且 $g_0\notin M$.令 $g_1=P_Mg_0$.则
	      \begin{equation*}
	      g=(g_0-g_1)/\abs{g_0-g_1}
	      \end{equation*}
	      具有所需性质.

	      给定任一 $u\in H$, 令
	      \begin{equation*}
	      v=u-\lambda g\quad\text{其中}\
	      \lambda=\frac{\dual{\varphi}{u}}{\dual{\varphi}{g}}.
	      \end{equation*}
	      注意 $v$ 是有定义的, 因为 $\dual{\varphi}{g}\ne0$; 此外 $v\in M$,
	      因为 $\dual{\varphi}{v}=0$.由此推出 $(g,v)=0$, 即
	      \begin{equation*}
	      \dual{\varphi}{u}=\dual{\varphi}{g}(g,u)\quad\forall u\in H,
	      \end{equation*}
	      取 $f=\dual{\varphi}{g}g$ 即完成证明.
\end{enumerate}
\end{proof}

\begin{keyremark}
\label{rem:5.3}
\emph{$H$ 与 $H^\star$: 等同还是不等同? 三元组 $V\subset H\subset V^\star$.}
定理  \ref{thm:5.5} 断言存在由 $H$ 到 $H^\star$ 上的一个典范等距.因此把
$H$ 与 $H^\star$ 等同起来是``合理的''.我们将\emph{常常}这样做, 但\emph{并非
总是}这样做.下面是一个典型的, 在许多应用中出现的, 应当谨慎对待等同的情形.
设 $H$ 是带有标量积 $(\ ,\ )$ 与相应范数 $\abs{\cdot}$ 的 Hilbert 空间.设
$V\subset H$ 是 $H$ 中的稠密线性子空间.设 $V$ 有自己的范数 $\norm{\cdot}$,
且 $V$ 关于 $\norm{\cdot}$ 是 Banach 空间.设嵌入 $V\subset H$ 连续, 即
\begin{equation*}
\abs{v}\le C\norm{v}\quad\forall v\in V.
\end{equation*}
[例如 $H=L^2(0,1)$, $V=L^p(0,1)$ 且 $p>2$, 或 $V=C([0,1])$.]

存在一个典范映射 $T:H^\star\to V^\star$, 它不过是把 $H$ 上的连续线性泛函
$\varphi$ 限制到 $V$ 上, 即
\begin{equation*}
\dual{T\varphi}{v}_{V^\star,V}=\dual{\varphi}{v}_{H^\star,H}.
\end{equation*}
容易看出 $T$ 具有下述性质:
\begin{enumerate}[label=(\roman*), leftmargin=2.6em, itemsep=0.1em]
	\item $\norm{T\varphi}_{V^\star}\le C\abs{\varphi}_{H^\star}
	      \quad\forall\varphi\in H^\star$,
	\item $T$ 是单射,
	\item 若 $V$ 自反, 则 $R(T)$ 在 $V^\star$ 中稠密.\footnote{然而一般说来
	      $T$ \emph{不是}满射.}
\end{enumerate}

把 $H^\star$ 与 $H$ 等同起来, 并把 $T$ 看作由 $H^\star$ 到 $V^\star$ 的典范
嵌入, 通常写作
\begin{equation}
\boxed{V\subset H\simeq H^\star\subset V^\star,}
\label{eq:5.9}
\end{equation}
其中一切嵌入都连续且稠密 (只要 $V$ 自反).人们说 $H$ 是\emph{枢轴空间}.
注意标量积 $\dual{\cdot}{\cdot}_{V^\star,V}$ 与 $(\ ,\ )$ 只要都有意义就
一致, 即
\begin{equation*}
\dual{f}{v}_{V^\star,V}=(f,v)\quad\forall f\in H,\ \forall v\in V.
\end{equation*}
\end{keyremark}

如果 $V$ 本身是带有自己的标量积 $((\ ,\ ))$ 的 Hilbert 空间, 而该标量积与
范数 $\norm{\cdot}$ 相对应, 则情形变得更微妙.当然, 我们可以借助
$((\ ,\ ))$ 把 $V$ 与 $V^\star$ 等同起来.然而 \eqref{eq:5.9} 就变得荒谬了.这表明
不能同时把 $V$ 与 $H$ 同它们的对偶空间等同起来: 必须作出选择.通常的习惯是
把 $H^\star$ 与 $H$ 等同, 写出 \eqref{eq:5.9}, 而不把 $V^\star$ 与 $V$ 等同 [自然, 仍
存在由 $V$ 到 $V^\star$ 的等距, 但它不被看作恒等映射].下面是一个极有启发的
例子.

令
\begin{equation*}
H=\ell^2=\Bigl\{u=(u_n)_{n\ge1};\ \sum_{n=1}^\infty u_n^2<\infty\Bigr\},
\end{equation*}
赋以标量积 $(u,v)=\sum_{n=1}^\infty u_nv_n$.

令
\begin{equation*}
V=\Bigl\{u=(u_n)_{n\ge1};\ \sum_{n=1}^\infty n^2u_n^2<\infty\Bigr\},
\end{equation*}
赋以标量积 $((u,v))=\sum_{n=1}^\infty n^2u_nv_n$.

显然 $V\subset H$ 且嵌入连续, 并且 $V$ 在 $H$ 中稠密.这里我们把 $H^\star$ 与
$H$ 等同, 而把 $V^\star$ 与空间
\begin{equation*}
V^\star=\Bigl\{f=(f_n)_{n\ge1};\ \sum_{n=1}^\infty\frac{1}{n^2}f_n^2
<\infty\Bigr\}
\end{equation*}
等同, 它比 $H$ 大.标量积 $\dual{\cdot}{\cdot}_{V^\star,V}$ 由
\begin{equation*}
\dual{f}{v}_{V^\star,V}=\sum_{n=1}^\infty f_nv_n
\end{equation*}
给出, 而 Riesz--Fréchet 同构 $T:V\to V^\star$ 由
\begin{equation*}
u=(u_n)_{n\ge1}\mapsto Tu=(n^2u_n)_{n\ge1}
\end{equation*}
给出.

\begin{remark}
\label{rem:5.4}
不借助一致凸空间的理论, 也容易证明 Hilbert 空间自反.只须两次使用
Riesz--Fréchet 同构 (先由 $H$ 到 $H^\star$, 再由 $H^\star$ 到 $H^{\star\star}$).
\end{remark}

\begin{remark}
\label{rem:5.5}
设 $H$ 是与它的对偶空间 $H^\star$ 等同的 Hilbert 空间.设 $M$ 为 $H$ 的
子空间.我们已经在 \S\ref{sec:1.3} 中把 $M^\perp$ 定义为 $H^\star$ 的子空间.现在可以
把它看作 $H$ 的子空间, 即
\begin{equation*}
M^\perp=\{u\in H;\ (u,v)=0\ \forall v\in M\}.
\end{equation*}
显然有 $M\cap M^\perp=\{0\}$.此外, 若 $M$ 闭, 则还有
$M+M^\perp=H$.事实上, 每个 $f\in H$ 都可以写成
\begin{equation*}
f=(P_Mf)+(f-P_Mf)
\end{equation*}
而 $f-P_Mf\in M^\perp$; 更确切地说, $f-P_Mf=P_{M^\perp}f$.

由此推出, 在 Hilbert 空间中每个闭子空间都有补空间 (\S\ref{sec:2.4} 意义下).
\end{remark}
\section{\index{Stampacchiadingli@Stampacchia 定理 (Stampacchia theorem)}
Stampacchia 定理与 \index{LaxMilgramdingli@Lax--Milgram 定理 (Lax--Milgram theorem)}
Lax--Milgram 定理}
\label{sec:5.3}

\begin{definition}
称双线性形式 $a:H\times H\to\R$ 为
\begin{enumerate}[label=(\roman*), leftmargin=2.6em, itemsep=0.25em]
	\item \emph{连续}的, 如果存在常数 $C$, 使得
	      \begin{equation*}
	      \abs{a(u,v)}\le C\abs{u}\abs{v}\quad\forall u,v\in H;
	      \end{equation*}
	\item \emph{强制}的, 如果存在常数 $\alpha>0$, 使得
	      \begin{equation*}
	      a(v,v)\ge\alpha\abs{v}^2\quad\forall v\in H.
	      \end{equation*}
\end{enumerate}
\end{definition}

\begin{theorem}[Stampacchia]
\label{thm:5.6}
假设 $a(u,v)$ 是 $H$ 上的连续强制双线性形式.设 $K\subset H$ 为非空闭凸子集.
则对任意给定的 $\varphi\in H^\star$, 存在唯一的元素 $u\in K$, 使得
\begin{equation}
a(u,v-u)\ge\dual{\varphi}{v-u}\quad\forall v\in K.
\label{eq:5.10}
\end{equation}
此外, 若 $a$ 对称, 则 $u$ 由下述性质刻画:
\begin{equation}
\boxed{u\in K\quad\text{且}\quad
\frac12a(u,u)-\dual{\varphi}{u}
=\min_{v\in K}\Bigl\{\frac12a(v,v)-\dual{\varphi}{v}\Bigr\}.}
\label{eq:5.11}
\end{equation}
\end{theorem}

定理  \ref{thm:5.6} 的证明依赖于下面这个非常经典的结果.

\begin{keytheorem}[Banach 不动点定理——压缩映射原理]
\label{thm:5.7}
设 $X$ 为非空完备度量空间, $S:X\to X$ 为严格压缩, 即
\begin{equation*}
d(Sv_1,Sv_2)\le k\,d(v_1,v_2)\quad\forall v_1,v_2\in X,\ \text{其中}\ k<1.
\end{equation*}
则 $S$ 有唯一不动点, $u=Su$.
\end{keytheorem}

证明见例如 T.~M.~Apostol \bref{Apostol}{1}, G.~Choquet \bref{Choquet}{1}, A.~Friedman \bref{Friedman}{3}.

\begin{proof}[定理 \ref{thm:5.6}]
由 Riesz--Fréchet 表示定理 (定理  \ref{thm:5.5}) 我们知道, 存在唯一的 $f\in H$,
使得
\begin{equation*}
\dual{\varphi}{v}=(f,v)\quad\forall v\in H.
\end{equation*}
另一方面, 若固定 $u\in H$, 则映射 $v\mapsto a(u,v)$ 是 $H$ 上的连续线性泛函.
再一次使用 Riesz--Fréchet 表示定理, 我们找到 $H$ 中某个唯一的元素, 记作
$Au$, 使得 $a(u,v)=(Au,v)\ \forall v\in H$.显然 $A$ 是由 $H$ 到 $H$ 的线性
算子, 满足
\begin{align}
\abs{Au}&\le C\abs{u}\quad\forall u\in H,\label{eq:5.12}\\
(Au,u)&\ge\alpha\abs{u}^2\quad\forall u\in H.\label{eq:5.13}
\end{align}
问题 \eqref{eq:5.10} 归结为寻找某个 $u\in K$, 使得
\begin{equation}
(Au,v-u)\ge(f,v-u)\quad\forall v\in K.
\label{eq:5.14}
\end{equation}
设 $\rho>0$ 为常数 (稍后确定).注意 \eqref{eq:5.14} 等价于
\begin{equation}
(\rho f-\rho Au+u-u,v-u)\le0\quad\forall v\in K,
\label{eq:5.15}
\end{equation}
即
\begin{equation*}
u=P_K(\rho f-\rho Au+u).
\end{equation*}
对每个 $v\in K$, 令 $Sv=P_K(\rho f-\rho Av+v)$.我们断言: 若 $\rho>0$ 选得
适当, 则 $S$ 是严格压缩.事实上, 由于 $P_K$ 不增大距离 (见命题  \ref{prop:5.3}),
我们有
\begin{equation*}
\abs{Sv_1-Sv_2}\le\abs{(v_1-v_2)-\rho(Av_1-Av_2)},
\end{equation*}
从而
\begin{align*}
\abs{Sv_1-Sv_2}^2
&=\abs{v_1-v_2}^2-2\rho(Av_1-Av_2,v_1-v_2)+\rho^2\abs{Av_1-Av_2}^2\\
&\le\abs{v_1-v_2}^2\bigl(1-2\rho\alpha+\rho^2C^2\bigr).
\end{align*}
选取 $\rho>0$ 使得 $k^2=1-2\rho\alpha+\rho^2C^2<1$ (即
$0<\rho<2\alpha/C^2$), 即知 $S$ 有唯一不动点.\footnote{若须数值地计算该不动点,
则取 $\rho=\alpha/C^2$ 是有益的, 这样可使 $k$ 最小, 从而加速 $S$ 的迭代的收敛.}

现在假设形式 $a(u,v)$ 也是对称的.则 $a(u,v)$ 在 $H$ 上定义了一个新的标量积;
相应的范数 $a(u,u)^{1/2}$ 与原范数 $\abs{u}$ 等价.由此推出 $H$ 关于这个新标量积
也是 Hilbert 空间.利用 Riesz--Fréchet 定理, 我们现在可以借助新的标量积来表示
泛函 $\varphi$, 即存在某个唯一的元素 $g\in H$, 使得
\begin{equation*}
\dual{\varphi}{v}=a(g,v)\quad\forall v\in H.
\end{equation*}
问题 \eqref{eq:5.10} 归结为寻找某个 $u\in K$, 使得
\begin{equation}
a(g-u,v-u)\le0\quad\forall v\in K.
\label{eq:5.16}
\end{equation}
\\eqref{eq:5.16} 的解是一位老朋友: $u$ 不过是 $g$ 关于新标量积 $a$ 在 $K$ 上的投影.我们还
知道 (由定理  \ref{thm:5.2}) $u$ 是 $K$ 中使
\begin{equation*}
\min_{v\in K}a(g-v,g-v)^{1/2}
\end{equation*}
达到的唯一元素.这相当于在 $K$ 上极小化函数
\begin{equation*}
v\mapsto a(g-v,g-v)=a(v,v)-2a(g,v)+a(g,g)
=a(v,v)-2\dual{\varphi}{v}+a(g,g),
\end{equation*}
或等价地极小化函数
\begin{equation*}
v\mapsto\frac12a(v,v)-\dual{\varphi}{v}.
\end{equation*}
\end{proof}

\begin{remark}
\label{rem:5.6}
容易验证: 若 $a(u,v)$ 是满足
\begin{equation*}
a(v,v)\ge0\quad\forall v\in H
\end{equation*}
的双线性形式, 则函数 $v\mapsto a(v,v)$ 凸.
\end{remark}

\begin{keycorollary}[Lax--Milgram]
\label{cor:5.8}
假设 $a(u,v)$ 是 $H$ 上的连续强制双线性形式.则对任意给定的
$\varphi\in H^\star$, 存在唯一的元素 $u\in H$, 使得
\begin{equation}
a(u,v)=\dual{\varphi}{v}\quad\forall v\in H.
\label{eq:5.17}
\end{equation}
此外, 若 $a$ 对称, 则 $u$ 由下述性质刻画:
\begin{equation}
\boxed{u\in H\quad\text{且}\quad
\frac12a(u,u)-\dual{\varphi}{u}
=\min_{v\in H}\Bigl\{\frac12a(v,v)-\dual{\varphi}{v}\Bigr\}.}
\label{eq:5.18}
\end{equation}
\end{keycorollary}

\begin{proof}
在定理  \ref{thm:5.6} 中取 $K=H$, 并仿照推论  \ref{cor:5.4} 的证明进行论证.
\end{proof}

\begin{remark}
\label{rem:5.7}
Lax--Milgram 定理是求解线性椭圆型偏微分方程的一个非常简单而有效的工具 (见
第~
\ref{chap:8}~章与第~
\ref{chap:9}~章).注意方程 \eqref{eq:5.17} 与极小化问题 \eqref{eq:5.18} 之间的联系是很有意思的.
当这类问题在力学或物理学中出现时, 它们往往有自然的解释: 最小作用量原理, 能量的极小化, 等等.用变分法的语言来说, \eqref{eq:5.17} 是与极小化问题 \eqref{eq:5.18} 相联系的
Euler 方程.粗略地说, \eqref{eq:5.17} 说的是``$F'(u)=0$'', 其中 $F$ 是函数
$F(v)=\frac12a(v,v)-\dual{\varphi}{v}$.
\end{remark}

\begin{remark}
\label{rem:5.8}
有一个直接而初等的论证可以证明 \eqref{eq:5.17} 有唯一解.事实上, 这相当于证明
\begin{equation*}
\forall f\in H\ \ \exists u\in H\ \text{唯一, 使得}\ Au=f,
\end{equation*}
即 $A$ 是由 $H$ 到 $H$ 上的双射.这是下述事实的平凡推论:
\begin{enumerate}[label=(\alph*), leftmargin=2.2em, itemsep=0.2em]
	\item $A$ 是单射 (因为 $A$ 强制),
	\item $R(A)$ 闭, 因为 $\alpha\abs{v}\le\abs{Av}\ \forall v\in H$
	      (强制性的推论),
	\item $R(A)$ 稠密; 事实上, 假设 $v\in H$ 满足
	      \begin{equation*}
	      (Au,v)=0\quad\forall u\in H,
	      \end{equation*}
	      则 $v=0$.
\end{enumerate}
\end{remark}

\section{Hilbert 和. 正交基}
\label{sec:5.4}

\begin{definition}
设 $(E_n)_{n\ge1}$ 为 $H$ 的一列闭子空间.若
\begin{enumerate}[label=(\alph*), leftmargin=2.2em, itemsep=0.25em]
	\item 空间 $E_n$ 两两正交, 即
	      \begin{equation*}
	      (u,v)=0\quad\forall u\in E_n,\ \forall v\in E_m,\ m\ne n,
	      \end{equation*}
	\item 由 $\bigcup_{n=1}^\infty E_n$ 张成的线性空间在 $H$ 中稠密,\footnote{由
	      $E_n$ 张成的线性空间按代数意义理解, 即属于诸空间 $(E_n)$ 的元素的有
	      限线性组合.}
\end{enumerate}
则说 $H$ 是诸 $E_n$ 的\emph{Hilbert 和}, 并记作 $H=\bigoplus_nE_n$.
\end{definition}

\begin{keytheorem}
\label{thm:5.9}
假设 $H$ 是诸 $E_n$ 的 Hilbert 和.给定 $u\in H$, 令
\begin{equation*}
u_n=P_{E_n}u\quad\text{及}\quad S_n=\sum_{k=1}^n u_k.
\end{equation*}
则我们有
\begin{equation}
\lim_{n\to\infty}S_n=u
\label{eq:5.19}
\end{equation}
以及
\begin{equation}
\sum_{k=1}^\infty\abs{u_k}^2=\abs{u}^2
\quad(\text{Bessel--Parseval 恒等式}).
\label{eq:5.20}
\end{equation}
\end{keytheorem}

使用下述引理是方便的.

\begin{lemma}
\label{lem:5.1}
假设 $(v_n)$ 为 $H$ 中任一序列, 使得
\begin{align}
(v_m,v_n)&=0\quad\forall m\ne n,\label{eq:5.21}\\
\sum_{k=1}^\infty\abs{v_k}^2&<\infty.\label{eq:5.22}
\end{align}
令
\begin{equation*}
S_n=\sum_{k=1}^n v_k.
\end{equation*}
则
\begin{equation*}
S=\lim_{n\to\infty}S_n
\end{equation*}
存在, 并且
\begin{equation}
\abs{S}^2=\sum_{k=1}^\infty\abs{v_k}^2.
\label{eq:5.23}
\end{equation}
\end{lemma}

\begin{proof}[引理 \ref{lem:5.1}]
注意对 $m>n$ 有
\begin{equation*}
\abs{S_m-S_n}^2=\sum_{k=n+1}^m\abs{v_k}^2.
\end{equation*}
由此推出 $S_n$ 是 Cauchy 序列, 从而 $S=\lim_{n\to\infty}S_n$ 存在.另一方面,
我们有
\begin{equation*}
\abs{S_n}^2=\sum_{k=1}^n\abs{v_k}^2.
\end{equation*}
令 $n\to\infty$ 即得 \eqref{eq:5.23}.
\end{proof}

\begin{proof}[定理 \ref{thm:5.9}]
由于 $u_n=P_{E_n}u$, 我们有 (由 \eqref{eq:5.8})
\begin{equation}
(u-u_n,v)=0\quad\forall v\in E_n,
\label{eq:5.24}
\end{equation}
特别地,
\begin{equation*}
(u,u_n)=\abs{u_n}^2.
\end{equation*}
把这些等式相加, 得到
\begin{equation*}
(u,S_n)=\sum_{k=1}^n\abs{u_k}^2.
\end{equation*}
但我们还有
\begin{equation}
\sum_{k=1}^n\abs{u_k}^2=\abs{S_n}^2,
\label{eq:5.25}
\end{equation}
从而得到
\begin{equation*}
(u,S_n)=\abs{S_n}^2.
\end{equation*}
由此推出 $\abs{S_n}\le\abs{u}$, 因此 $\sum_{k=1}^n\abs{u_k}^2\le\abs{u}^2$.

于是我们可以应用引理  \ref{lem:5.1}, 断言 $S=\lim_{n\to\infty}S_n$ 存在.让我们
\emph{不}假设 (b) 而确定 $S$.设 $F$ 为由诸 $E_n$ 张成的线性空间.我们断言
\begin{equation}
S=P_{\overline{F}}u.
\label{eq:5.26}
\end{equation}
事实上, 我们有
\begin{equation*}
(u-S_n,v)=0\quad\forall v\in E_m,\ m\le n
\end{equation*}
(只须写出 $u-S_n=(u-u_m)-\sum_{k\ne m}u_k$).令 $n\to\infty$ 得到
\begin{equation*}
(u-S,v)=0\quad\forall v\in E_m,\ \forall m,
\end{equation*}
从而
\begin{equation*}
(u-S,v)=0\quad\forall v\in F,
\end{equation*}
这蕴含
\begin{equation*}
(u-S,v)=0\quad\forall v\in\overline{F}.
\end{equation*}
另一方面, $S_n\in F\ \forall n$, 而取极限得 $S\in\overline{F}$.这就证明了 \eqref{eq:5.26}.
当然, 若 (b) 成立, 则 $\overline{F}=H$, 从而 $S=u$.在 \eqref{eq:5.25} 中令
$n\to\infty$ 取极限, 即得 \eqref{eq:5.20}.
\end{proof}

\begin{definition}
称 $H$ 中的序列 $(e_n)_{n\ge1}$ 为 $H$ 的\emph{正交基} (或\emph{Hilbert
基},\footnote{不要与代数 (= Hamel) 基混淆; 后者是 $H$ 中的族 $(e_i)_{i\in I}$,
使得每个 $u\in H$ 都能唯一地写成诸 $e_i$ 的有限线性组合 (见习题 1.5).}或
在没有混淆时简称为\emph{基}\footnote{有些作者说 $(e_n)$ 是完全正交系.}), 如果
它满足下述性质:
\begin{enumerate}[label=(\roman*), leftmargin=2.6em, itemsep=0.2em]
	\item $\abs{e_n}=1\ \forall n$ 且 $(e_m,e_n)=0\ \forall m\ne n$,
	\item 由诸 $e_n$ 张成的线性空间在 $H$ 中稠密.
\end{enumerate}
\end{definition}

\begin{keycorollary}
\label{cor:5.10}
设 $(e_n)$ 为正交基.则对每个 $u\in H$, 我们有
\begin{equation*}
u=\sum_{k=1}^\infty(u,e_k)e_k,
\quad\text{即}\quad u=\lim_{n\to\infty}\sum_{k=1}^n(u,e_k)e_k
\end{equation*}
以及
\begin{equation*}
\abs{u}^2=\sum_{k=1}^\infty\abs{(u,e_k)}^2.
\end{equation*}
反之, 给定任一序列 $(\alpha_n)\in\ell^2$, 级数 $\sum_{k=1}^\infty\alpha_ke_k$
收敛到某个元素 $u\in H$, 使得 $(u,e_k)=\alpha_k\ \forall k$ 且
$\abs{u}^2=\sum_{k=1}^\infty\alpha_k^2$.
\end{keycorollary}

\begin{proof}
注意 $H$ 是诸空间 $E_n=\R e_n$ 的 Hilbert 和, 且 $P_{E_n}u=(u,e_n)e_n$.
使用定理  \ref{thm:5.9} 与引理  \ref{lem:5.1}.
\end{proof}

\begin{remark}
\label{rem:5.9}
一般说来, 定理  \ref{thm:5.9} 中的级数 $\sum u_k$ 与推论  \ref{cor:5.10} 中的
级数 $\sum(u,e_k)e_k$ \emph{不是}绝对收敛的, 即可能出现
$\sum_{k=1}^\infty\abs{u_k}=\infty$ 或
$\sum_{k=1}^\infty\abs{(u,e_k)}=\infty$.
\end{remark}

\begin{keytheorem}
\label{thm:5.11}
每个可分 Hilbert 空间都有正交基.
\end{keytheorem}

\begin{proof}
设 $(v_n)$ 为 $H$ 的可数稠密子集.用 $F_k$ 表示由
$\{v_1,v_2,\dots,v_k\}$ 张成的线性空间.序列 $(F_k)$ 是有限维空间的不减序列,
使得 $\bigcup_{k=1}^\infty F_k$ 在 $H$ 中稠密.在 $F_1$ 中任取一个单位向量
$e_1$.若 $F_2\ne F_1$, 则在 $F_2$ 中存在某个向量 $e_2$, 使得 $\{e_1,e_2\}$ 是
$F_2$ 的正交基.重复同样的构造, 即得 $H$ 的一个正交基.
\end{proof}

\begin{remark}
\label{rem:5.10}
定理  \ref{thm:5.11} 与推论  \ref{cor:5.10} 合起来表明: 一切可分 Hilbert 空间都
与空间 $\ell^2$ 同构且等距.尽管这个结果看似惊人, 考虑诸如 $L^2(\Omega)$
(或 Sobolev 空间 $H^1(\Omega)$ 等) 这样的其他 Hilbert 空间仍然是十分重要的.
原因是: 许多很好的线性 (或非线性) 算子, 一旦用基写出来, 可能显得很可怕.
\end{remark}

\begin{remark}
\label{rem:5.11}
若 $H$ 是不可分的 Hilbert 空间——这是一种相当不寻常的情形——则仍可 (借助
Zorn 引理) 证明不可数正交基 $(e_i)_{i\in I}$ 的存在性; 见例如 W.~Rudin
\bref{Rudin}{2}, A.~E.~Taylor--D.~C.~Lay \bref{TaylorLay}{1}, G.~B.~Folland \bref{Folland}{2}, G.~Choquet \bref{Choquet}{1}.
\end{remark}
\section*{第 5 章评注}
\addcontentsline{toc}{section}{第 5 章评注}
\markright{第 5 章评注}

\cmtitem{1. Hilbert 空间的刻画.}
有时知道一个给定的范数 $\norm{\cdot}$ 在向量空间 $E$ 上是否为 Hilbert 范数
是有用的, 即是否存在 $E$ 上的标量积 $(\ ,\ )$, 使得
$\norm{u}=(u,u)^{1/2}\ \forall u\in E$.已知若干种判别准则:

\begin{cmtlettertheorem}{a}{Fréchet--von Neumann--Jordan}
\label{thm:5.12}
假设范数 $\norm{\cdot}$ 满足平行四边形法则 \eqref{eq:5.1}.则 $\norm{\cdot}$ 是 Hilbert
范数.
\end{cmtlettertheorem}

证明见 K.~Yosida \bref{Yosida}{1} 或习题 5.1.

\begin{cmtlettertheorem}{b}{Kakutani}
\label{thm:5.13}
假设 $E$ 是 $\dim E\ge3$ 的赋范空间.假设每个二维子空间 $F$ 都有一个范数为
1 的投影算子 (即存在有界线性投影算子 $P:E\to F$, 使得 $Pu=u\ \forall u\in F$
且 $\norm{P}\le1$).\footnote{我们指出, 每个一维子空间总有范数为 1 的投影算子.
(使用 Hahn--Banach 定理.)}则 $\norm{\cdot}$ 是 Hilbert 范数.
\end{cmtlettertheorem}

\begin{cmtlettertheorem}{c}{de Figueiredo--Karlovitz}
\label{thm:5.14}
设 $E$ 是 $\dim E\ge3$ 的赋范空间.考虑单位球上的径向投影, 即
\begin{equation*}
Tu=\begin{cases}
u & \text{若}\ \norm{u}\le1,\\
u/\norm{u} & \text{若}\ \norm{u}>1.
\end{cases}
\end{equation*}
假设\footnote{可以证明, 在任意赋范空间中, $T$ 满足
$\norm{Tu-Tv}\le2\norm{u-v}\ \forall u,v\in E$, 且常数 2 不能改进;
见习题 5.6.}
\begin{equation*}
\norm{Tu-Tv}\le\norm{u-v}\quad\forall u,v\in E.
\end{equation*}
则 $\norm{\cdot}$ 是 Hilbert 范数.
\end{cmtlettertheorem}

最后, 让我们回忆一个已经提到过的结果 (注记 2.8).

\begin{cmtlettertheorem}{d}{Lindenstrauss--Tzafriri}
\label{thm:5.15}
假设 $E$ 是 Banach 空间, 使得每个闭子空间都有补空间.\footnote{这等价于说
每个闭子空间都有一个有界投影算子 $P$.注意这里——与定理  \ref{thm:5.13}
相反——我们不假设 $\norm{P}\le1$.}则 $E$ 可 Hilbert 化, 即存在一个等价的
Hilbert 范数.
\end{cmtlettertheorem}

\cmtitem{2. 变分不等式.}
Stampacchia 定理是变分不等式理论的出发点 (见例如
D.~Kinderlehrer--G.~Stampacchia \bref{KinderlehrerStampacchia}{1}), 该理论在
力学与物理学 (见例如 G.~Duvaut--J.~L.~Lions \bref{DuvautLions}{1}), 自由边界
值问题 (见例如 C.~Baiocchi--A.~Capelo \bref{BaiocchiCapelo}{1} 与
A.~Friedman \bref{Friedman}{4}), 最优控制 (见例如 J.-L.~Lions
\bref{Lions}{2} 与 V.~Barbu \bref{Barbu}{2}), 随机控制 (见
A.~Bensoussan--J.-L.~Lions \bref{BensoussanLions}{1}) 中有大量应用.

\cmtitem{3. 与单调算子相联系的非线性方程.}
Stampacchia 定理与 Lax--Milgram 定理可以推广到某些类非线性算子.例如,
让我们提到下述结果.

\begin{theorem}[Minty--Browder]
\label{thm:5.16}
设 $E$ 为自反 Banach 空间.设 $A:E\to E^\star$ 为连续非线性映射, 使得
\begin{equation*}
\dual{Av_1-Av_2}{v_1-v_2}>0\quad\forall v_1,v_2\in E,\ v_1\ne v_2,
\end{equation*}
且
\begin{equation*}
\lim_{\norm{v}\to\infty}\frac{\dual{Av}{v}}{\norm{v}}=\infty.
\end{equation*}
则对每个 $f\in E^\star$ 存在方程 $Au=f$ 的唯一解 $u\in E$.
\end{theorem}

有兴趣的读者可在 F.~Browder \bref{Browder}{1} 与 J.-L.~Lions
\bref{Lions}{3} 中找到定理  \ref{thm:5.16} 的证明以及许多推广与应用; 另见
问题 31.

\cmtitem{4. 特殊的正交基. Fourier 级数. 小波.}
在第~
\ref{chap:6}~章中我们将介绍一种构造正交基的非常有力的方法, 即取紧自伴算子的
特征向量.在实践中, 人们常常使用 $L^2(\Omega)$ 的由微分算子的特征函数组成的
特殊基 (见 \S\ref{sec:8.6} 与 \S\ref{sec:9.8}).由
\begin{equation*}
e_n(x)=\sqrt{2/\pi}\sin nx,\quad n\ge1,
\end{equation*}
或
\begin{equation*}
e_n(x)=\sqrt{2/\pi}\cos nx,\quad n\ge0,
\end{equation*}
定义的 $L^2(0,\pi)$ 上的正交基很受喜爱, 因为它导致 Fourier 级数与调和分析
——一个本身就很重要的领域; 见例如 J.~M.~Ash \bref{Ash}{1}, H.~Dym--H.~P.~McKean \bref{DymMcKean}{1}, Y.~Katznelson
\bref{Katznelson}{1}, C.~S.~Rees--S.~M.~Shah--C.~V.~Stanojevic
\bref{ReesShahStanojevic}{1}.

这里有一个使分析学家困惑了数十年的问题.给定 $u\in L^2(0,\pi)$, 考虑它的
Fourier 级数 $S_n=\sum_{k=1}^n(u,e_k)e_k$.已知 (见推论  \ref{cor:5.10})
$S_n\to u$ 在 $L^2(0,\pi)$ 中.由此推出子列 $S_{n_k}\to u$ 在 $(0,\pi)$ 上
几乎处处 (见定理  \ref{thm:4.9}).但是否可以说整个序列 $S_n\to u$ 在
$(0,\pi)$ 上几乎处处? 答案由下述非常深刻的结果给出:

\begin{theorem}[Carleson]
\label{thm:5.17}
若 $u\in L^2(0,\pi)$, 则 $S_n\to u$ 几乎处处.
\end{theorem}

$L^2(0,1)$ 或 $L^2(\R)$ 的其他经典基与 Bessel, Legendre, Hermite, Laguerre, Chebyshev, Jacobi 等人的名字相联系.我们把有兴趣的读者引向
R.~Courant--D.~Hilbert \bref{CourantHilbert}{1} 第 1 卷以及
R.~Dautray--J.-L.~Lions \bref{DautrayLions}{1} 第 VIII 章; 另见第~
\ref{chap:8}~章末尾的
评注 (Sturm--Liouville 算子的谱性质).近来, 由阶梯函数组成的 $L^2(0,1)$ 的
Haar 基与 Walsh 基也引起了很大兴趣; 见例如习题 5.31, 5.32, G.~Alexits
\bref{Alexits}{1}, H.~F.~Harmuth \bref{Harmuth}{1}.

小波理论提供了一种非常重要而优美的新型基.它是分解函数, 信号, 语音, 图像等的
有力工具.有兴趣的读者可以查阅 Y.~Meyer \bref{Meyer}{1}, \bref{Meyer}{2}, \bref{Meyer}{3}, R.~Coifman 与 Y.~Meyer \bref{CoifmanMeyer}{1}, I.~Daubechies \bref{Daubechies}{1}, G.~David \bref{David}{1}, C.~K.~Chui \bref{Chui}{1}, M.~B.~Ruskai 等 \bref{RuskaiEtAl}{1}, J.~J.~Benedetto--M.~W.~Frazier \bref{BenedettoFrazier}{1}, G.~Kaiser
\bref{Kaiser}{1}, J.~P.~Kahane--P.~G.~Lemarié-Rieusset
\bref{KahaneLemarie}{1}, S.~Mallat \bref{Mallat}{1}, G.~Bachman--L.~Narici--E.~Beckenstein \bref{BachmanNariciBeckenstein}{1}, T.~F.~Chan--J.~Shen \bref{ChanShen}{1}, P.~Wojtaszczyk \bref{Wojtaszczyk}{1}, E.~Hernandez--G.~Weiss \bref{HernandezWeiss}{1} 等人的近著及其中的参考文献.

\cmtitem{5. Banach 空间中的 Schauder 基.}
设 $E$ 为 Banach 空间.称序列 $(e_n)_{n\ge1}$ 为\emph{Schauder 基}, 如果对每个
$u\in E$ 存在唯一的序列 $(\alpha_n)_{n\ge1}\subset\R$, 使得
$u=\sum_{k=1}^\infty\alpha_ke_k$ (即
$u=\lim_{n\to\infty}\sum_{k=1}^n\alpha_ke_k$).这样的基在 Banach 空间的几何
理论中起着重要作用 (见例如 B.~Beauzamy \bref{Beauzamy}{1}, J.~Lindenstrauss--L.~Tzafriri \bref{LindenstraussTzafriri}{2}, J.~Diestel \bref{Diestel}{2}, R.~C.~James \bref{James}{2}).分析中使用的
一切经典 (可分) Banach 空间都有 Schauder 基 (见例如 I.~Singer
\bref{Singer}{1}).这一事实使 Banach 猜测每个可分 Banach 空间都有基.经过
数十年徒劳的努力之后, P.~Enflo \bref{Enflo}{1} 发现了反例.甚至可以在
$\ell^p$ ($1<p<\infty$, $p\ne2$) 中构造没有 Schauder 基的闭子空间 (见
J.~Lindenstrauss--L.~Tzafriri \bref{LindenstraussTzafriri}{2}).A.~Szankowski
\bref{Szankowski}{1} 发现了另一个惊人的例子: 当 $H$ 为无限维可分 Hilbert
空间时, $\mathcal{L}(H)$ (配以通常的范数) 没有 Schauder 基.在第~
\ref{chap:6}~章中我们
将看到, 关于紧算子的一个相关问题也有否定的答案.
\section*{第 5 章习题}
\addcontentsline{toc}{section}{第 5 章习题}
\markright{第 5 章习题}

下文 $H$ 总表示赋以标量积 $(\ ,\ )$ 及相应范数 $\abs{\cdot}$ 的 Hilbert 空间.

\begin{exercise}[平行四边形法则.]
\label{exer:5.1}
假设 $E$ 是赋以范数 $\norm{\cdot}$ 的向量空间, 该范数满足平行四边形法则, 即
\begin{equation*}
\norm{a+b}^2+\norm{a-b}^2=2\bigl(\norm{a}^2+\norm{b}^2\bigr)
\quad\forall a,b\in E.
\end{equation*}
我们的目的是证明由
\begin{equation*}
(u,v)=\frac12\bigl(\norm{u+v}^2-\norm{u}^2-\norm{v}^2\bigr),
\quad u,v\in E,
\end{equation*}
定义的量是标量积, 且 $(u,u)=\norm{u}^2$.
\begin{enumerate}[label=\arabic*., leftmargin=2.2em, itemsep=0.3em]
	\item 验证
	      \begin{equation*}
	      (u,v)=(v,u),\quad(-u,v)=-(u,v)\quad\text{且}\quad(u,2v)=2(u,v)
	      \quad\forall u,v\in E.
	      \end{equation*}
	\item 证明
	      \begin{equation*}
	      (u+v,w)=(u,w)+(v,w)\quad\forall u,v,w\in E.
	      \end{equation*}
	      [提示: 依次对 (i) $a=u$, $b=v$; (ii) $a=u+w$, $b=v+w$;
	      (iii) $a=u+v+w$, $b=w$ 使用平行四边形法则.]
	\item 证明 $(\lambda u,v)=\lambda(u,v)\ \forall\lambda\in\R,\ \forall u,v\in E$.

	      [提示: 先考虑 $\lambda\in\N$ 的情形, 再考虑 $\lambda\in\Q$, 最后
	      考虑 $\lambda\in\R$.]
	\item 作出结论.
\end{enumerate}
\end{exercise}

\begin{exercise}[当 $p\ne2$ 时 $L^p$ 不是 Hilbert 空间.]
\label{exer:5.2}
设 $\Omega$ 为测度空间, 并假设存在可测集 $A\subset\Omega$, 使得
$0<\abs{A}<\abs{\Omega}$.

证明: 对任意 $1\le p\le\infty$, $p\ne2$, 范数 $\norm{\cdot}_p$ 都不满足
平行四边形法则.

[提示: 使用支集互不相交的函数.]
\end{exercise}

\begin{exercise}
\label{exer:5.3}
设 $(u_n)$ 为 $H$ 中的序列, $(t_n)$ 为 $(0,\infty)$ 中的序列, 使得
\begin{equation*}
(t_nu_n-t_mu_m,u_n-u_m)\le0\quad\forall m,n.
\end{equation*}
\begin{enumerate}[label=\arabic*., leftmargin=2.2em, itemsep=0.3em]
	\item 假设序列 $(t_n)$ 不减 (可能无界).证明序列 $(u_n)$ 收敛.

	      [提示: 证明序列 $(\abs{u_n})$ 不增.]
	\item 假设序列 $(t_n)$ 不增.证明下述二者必居其一:
	      \begin{enumerate}[label=(\roman*), leftmargin=2.4em, itemsep=0.1em]
	      	\item 或者 $\abs{u_n}\to\infty$,
	      	\item 或者 $(u_n)$ 收敛.
	      \end{enumerate}
	      若 $t_n\to t>0$, 证明 $(u_n)$ 收敛; 而若 $t_n\to0$, 证明情形 (i)
	      与 (ii) 都可能出现.
\end{enumerate}
\end{exercise}

\begin{exercise}
\label{exer:5.4}
设 $K\subset H$ 为非空闭凸集.设 $f\in H$, $u=P_Kf$.证明
\begin{equation*}
\abs{v-u}^2\le\abs{v-f}^2-\abs{u-f}^2\quad\forall v\in K.
\end{equation*}
由此推出
\begin{equation*}
\abs{v-u}\le\abs{v-f}\quad\forall v\in K.
\end{equation*}
给出几何解释.
\end{exercise}

\begin{exercise}
\label{exer:5.5}
\begin{enumerate}[label=\arabic*., leftmargin=2.2em, itemsep=0.3em]
	\item 设 $(K_n)$ 为 $H$ 中不增的闭凸集序列, 使得
	      $\bigcap_nK_n\ne\emptyset$.证明对每个 $f\in H$, 序列
	      $u_n=P_{K_n}f$ (强) 收敛到一个极限, 并确定该极限.
	\item 设 $(K_n)$ 为 $H$ 中不增的非空闭凸集序列.证明对每个 $f\in H$,
	      序列 $u_n=P_{K_n}f$ (强) 收敛到一个极限, 并确定该极限.

	      设 $\varphi:H\to\R$ 为下方有界的连续函数.证明序列
	      $\alpha_n=\inf_{K_n}\varphi$ 收敛, 并确定该极限.
\end{enumerate}
\end{exercise}

\begin{exercise}[到单位球上的径向投影.]
\label{exer:5.6}
设 $E$ 为赋以范数 $\norm{\cdot}$ 的向量空间.令
\begin{equation*}
Tu=\begin{cases}
u & \text{若}\ \norm{u}\le1,\\
u/\norm{u} & \text{若}\ \norm{u}>1.
\end{cases}
\end{equation*}
\begin{enumerate}[label=\arabic*., leftmargin=2.2em, itemsep=0.3em]
	\item 证明 $\norm{Tu-Tv}\le2\norm{u-v}\quad\forall u,v\in E$.
	\item 证明一般说来常数 2 不能改进.

	      [提示: 取 $E=\R^2$, 范数 $\norm{u}=\abs{u_1}+\abs{u_2}$.]
	\item 若 $\norm{\cdot}$ 是 Hilbert 范数, 情形如何?
\end{enumerate}
\end{exercise}

\begin{exercise}[到凸锥上的投影.]
\label{exer:5.7}
设 $K\subset H$ 为顶点在原点的凸锥, 即
\begin{equation*}
0\in K\quad\text{且}\quad\lambda u+\mu v\in K
\quad\forall\lambda,\mu>0,\ \forall u,v\in K;
\end{equation*}
此外假设 $K$ 闭.

给定 $f\in H$, 证明 $u=P_Kf$ 由下述性质刻画:
\begin{equation*}
u\in K,\quad(f-u,v)\le0\quad\forall v\in K,\quad\text{且}\quad(f-u,u)=0.
\end{equation*}
\end{exercise}

\begin{exercise}
\label{exer:5.8}
设 $\Omega$ 为测度空间, $h:\Omega\to[0,+\infty)$ 为可测函数.令
\begin{equation*}
K=\{u\in L^2(\Omega);\ \abs{u(x)}\le h(x)\ \text{a.e. 在}\ \Omega\ \text{上}\}.
\end{equation*}
验证 $K$ 是 $H=L^2(\Omega)$ 中的非空闭凸集.确定 $P_K$.
\end{exercise}

\begin{exercise}
\label{exer:5.9}
设 $A\subset H$ 与 $B\subset H$ 为两个非空闭凸集, 使得 $A\cap B=\emptyset$ 且
$B$ 有界.令
\begin{equation*}
C=A-B.
\end{equation*}
\begin{enumerate}[label=\arabic*., leftmargin=2.2em, itemsep=0.3em]
	\item 证明 $C$ 闭且凸.
	\item 令 $u=P_C0$, 并把 $u$ 写成 $u=a_0-b_0$, 其中 $a_0\in A$, $b_0\in B$ (这是可能的, 因为 $u\in C$).证明
	      $\abs{a_0-b_0}=\dist(A,B)=\inf_{a\in A,b\in B}\abs{a-b}$.
	      确定 $P_Ab_0$ 与 $P_Ba_0$.
	\item 设 $a_1\in A$, $b_1\in B$ 是另一对使得
	      $\abs{a_1-b_1}=\dist(A,B)$ 的元素.证明 $u=a_1-b_1$.
	      画出一些图形, 其中对 $[a_0,b_0]$ 唯一 (相应地, 不唯一).
	\item 对 Hilbert 空间的情形, 给出 Hahn--Banach 定理第二几何形式的一个
	      简单证明.
\end{enumerate}
\end{exercise}

\begin{exercise}
\label{exer:5.10}
设 $F:H\to\R$ 为 $C^1$ 类凸函数.设 $K\subset H$ 凸, $u\in H$.证明下述性质
等价:
\begin{enumerate}[label=(\roman*), leftmargin=2.4em, itemsep=0.15em]
	\item $F(u)\le F(v)\quad\forall v\in K$,
	\item $(F'(u),v-u)\ge0\quad\forall v\in K$.
\end{enumerate}
例: $F(v)=\abs{v-f}^2$, 其中 $f\in H$ 给定.
\end{exercise}

\begin{exercise}
\label{exer:5.11}
设 $M\subset H$ 为不化为 $\{0\}$ 的闭线性子空间.设 $f\in H$, $f\notin M^\perp$.
\begin{enumerate}[label=\arabic*., leftmargin=2.2em, itemsep=0.3em]
	\item 证明
	      \begin{equation*}
	      m=\inf_{\substack{u\in M\\\abs{u}=1}}(f,u)
	      \end{equation*}
	      被唯一地达到.
	\item 设 $\varphi_1,\varphi_2,\varphi_3\in H$ 给定, $E$ 表示由
	      $\{\varphi_1,\varphi_2,\varphi_3\}$ 张成的线性空间.在下列情形确定
	      $m$:
	      \begin{enumerate}[label=(\roman*), leftmargin=2.4em, itemsep=0.1em]
	      	\item $M=E$,
	      	\item $M=E^\perp$.
	      \end{enumerate}
	\item 考察 $H=L^2(0,1)$, $\varphi_1(t)=t$, $\varphi_2(t)=t^2$, $\varphi_3(t)=t^3$ 的情形.
\end{enumerate}
\end{exercise}

\begin{exercise}[预 Hilbert 空间的完备化.]
\label{exer:5.12}
设 $E$ 为赋以标量积 $(\ ,\ )$ 的向量空间.不假设 $E$ 关于范数
$\abs{u}=(u,u)^{1/2}$ 完备 ($E$ 称为\emph{预 Hilbert 空间}).

回忆对偶空间 $E^\star$ 赋以对偶范数 $\norm{f}_{E^\star}$ 后是完备的.设
$T:E\to E^\star$ 为如下定义的映射:
\begin{equation*}
\dual{Tu}{v}_{E^\star,E}=(u,v)\quad\forall u,v\in E.
\end{equation*}
验证 $T$ 是线性等距.$T$ 是满射吗?

我们的目的是证明 $R(T)$ 在 $E^\star$ 中稠密, 且 $\norm{\cdot}_{E^\star}$ 是
Hilbert 范数.

\begin{enumerate}[label=\arabic*., leftmargin=2.2em, itemsep=0.3em]
	\item 把 $E$ 的标量积转移到 $R(T)$ 上, 并把它延拓到 $\overline{R(T)}$.
	      所得的标量积记作 $((f,g))$, 其中 $f,g\in\overline{R(T)}$.

	      验证相应的范数 $((f,f))^{1/2}$ 在 $R(T)$ 上与
	      $\norm{f}_{E^\star}$ 一致.证明
	      \begin{equation*}
	      \dual{f}{v}=((f,Tv))\quad\forall v\in E,\ \forall f\in R(T).
	      \end{equation*}
	\item 证明 $\overline{R(T)}=E^\star$.

	      [提示: 给定 $f\in E^\star$, 把 $f$ 转移为 $\overline{R(T)}$ 上的
	      线性泛函, 并使用 $\overline{R(T)}$ 中的 Riesz--Fréchet 表示定理.]

	      由此推出 $E^\star$ 关于范数 $\norm{\cdot}_{E^\star}$ 是 Hilbert
	      空间.
	\item 作出结论: $E$ 的完备化可以等同于 $E^\star$.(完备化的定义见例如
	      A.~Friedman \bref{Friedman}{3}.)
\end{enumerate}
\end{exercise}

\begin{exercise}
\label{exer:5.13}
设 $E$ 为赋以范数 $\norm{\cdot}_E$ 的向量空间.对偶范数记作
$\norm{\cdot}_{E^\star}$.回忆 (\emph{多值}) 对偶映射定义为
\begin{equation*}
F(u)=\{f\in E^\star;\ \norm{f}_{E^\star}=\norm{u}_E
\ \text{且}\ \dual{f}{u}=\norm{u}_E^2\}.
\end{equation*}
\begin{enumerate}[label=\arabic*., leftmargin=2.2em, itemsep=0.3em]
	\item 假设 $F$ 满足下述性质:
	      \begin{equation*}
	      F(u)+F(v)\subset F(u+v)\quad\forall u,v\in E.
	      \end{equation*}
	      证明范数 $\norm{\cdot}_E$ 由某个标量积导出.

	      [提示: 使用习题 5.1.]
	\item 反之, 若范数 $\norm{\cdot}_E$ 由某个标量积导出, 关于 $F$ 能说些
	      什么?

	      [提示: 使用习题 5.12 与 1.1.]
\end{enumerate}
\end{exercise}

\begin{exercise}
\label{exer:5.14}
设 $a:H\times H\to\R$ 为双线性连续形式, 使得
\begin{equation*}
a(v,v)\ge0\quad\forall v\in H.
\end{equation*}
证明函数 $v\mapsto F(v)=a(v,v)$ 凸, 属于 $C^1$ 类, 并确定它的微分.
\end{exercise}

\begin{exercise}
\label{exer:5.15}
设 $G\subset H$ 为 Hilbert 空间 $H$ 的线性子空间; $G$ 赋以 $H$ 的范数.设
$F$ 为 Banach 空间.设 $S:G\to F$ 为有界线性算子.证明存在有界线性算子
$T:H\to F$, 它延拓 $S$, 且使得
\begin{equation*}
\norm{T}_{\mathcal{L}(H,F)}=\norm{S}_{\mathcal{L}(G,F)}.
\end{equation*}
\end{exercise}

\begin{exercise}[三元组 $V\subset H\subset V^\star$.]
\label{exer:5.16}
设 $H$ 为赋以标量积 $(\ ,\ )$ 及相应范数 $\abs{\cdot}$ 的 Hilbert 空间.设
$V\subset H$ 为在 $V$ 中稠密的线性子空间.
% 原书此处印作 "dense in V", 按上下文 (并与 5.2 节注记 3 的同一表述对照)
% 应为 "dense in H". 依技能规范保留原样并记录, 不擅自更正.
假设 $V$ 有自己的范数
$\norm{\cdot}$, 且 $V$ 关于 $\norm{\cdot}$ 是 Banach 空间.还假设嵌入
$V\subset H$ 连续, 即 $\abs{v}\le C\norm{v}\ \forall v\in V$.考虑由
\begin{equation*}
\dual{Tu}{v}_{V^\star,V}=(u,v)\quad\forall u\in H,\ \forall v\in V
\end{equation*}
定义的算子 $T:H\to V^\star$.
\begin{enumerate}[label=\arabic*., leftmargin=2.2em, itemsep=0.3em]
	\item 证明 $\norm{Tu}_{V^\star}\le C\abs{u}\ \forall u\in H$.
	\item 证明 $T$ 是单射.
	\item 证明若 $V$ 自反, 则 $R(T)$ 在 $V^\star$ 中稠密.
	\item 给定 $f\in V^\star$, 证明 $f\in R(T)$ 当且仅当存在常数 $a\ge0$,
	      使得 $\abs{\dual{f}{v}_{V^\star,V}}\le a\abs{v}\ \forall v\in V$.
\end{enumerate}
\end{exercise}

\begin{exercise}
\label{exer:5.17}
设 $M,N\subset H$ 为两个闭线性子空间.假设
$(u,v)=0\ \forall u\in M,\ \forall v\in N$.证明 $M+N$ 闭.
\end{exercise}

\begin{exercise}
\label{exer:5.18}
设 $E$ 为 Banach 空间, $H$ 为 Hilbert 空间.设
$T\in\mathcal{L}(E,H)$.证明下述性质等价:
\begin{enumerate}[label=(\roman*), leftmargin=2.4em, itemsep=0.15em]
	\item $T$ 有左逆,
	\item 存在常数 $C$, 使得 $\norm{u}\le C\abs{Tu}\ \forall u\in E$.
\end{enumerate}
\end{exercise}

\begin{exercise}
\label{exer:5.19}
设 $(u_n)$ 为 $H$ 中的序列, 使得 $u_n\weakto u$ 弱收敛.假设
$\limsup\abs{u_n}\le\abs{u}$.不依赖命题  \ref{prop:3.32} 证明
$u_n\to u$ 强收敛.
\end{exercise}

\begin{exercise}
\label{exer:5.20}
假设 $S\in\mathcal{L}(H)$ 满足 $(Su,u)\ge0\ \forall u\in H$.
\begin{enumerate}[label=\arabic*., leftmargin=2.2em, itemsep=0.3em]
	\item 证明 $N(S)=R(S)^\perp$.
	\item 证明对每个 $t>0$, $I+tS$ 是双射.
	\item 证明
	      \begin{equation*}
	      \lim_{t\to+\infty}(I+tS)^{-1}f=P_{N(S)}f\quad\forall f\in H.
	      \end{equation*}
	      [提示: 有两种方法:
	      \begin{enumerate}[label=(\alph*), leftmargin=2.2em, itemsep=0.1em]
	      	\item 分别考虑 $f\in N(S)$ 与 $f\in R(S)$ 的情形.
	      	\item 使用弱收敛.
	      \end{enumerate}]
\end{enumerate}
\end{exercise}

\begin{exercise}[线性压缩的迭代. Kakutani--Yosida 遍历定理.]
\label{exer:5.21}
设 $T\in\mathcal{L}(H)$ 满足 $\norm{T}\le1$.给定 $f\in H$ 与整数 $n\ge1$,
令
\begin{equation*}
\sigma_n(f)=\frac1n\bigl(f+Tf+T^2f+\cdots+T^{n-1}f\bigr)
\end{equation*}
以及
\begin{equation*}
\mu_n(f)=\Bigl(\frac{I+T}{2}\Bigr)^nf.
\end{equation*}
我们的目的是证明
\begin{equation*}
\lim_{n\to\infty}\sigma_n(f)=\lim_{n\to\infty}\mu_n(f)=P_{N(I-T)}f.
\end{equation*}
\begin{enumerate}[label=\arabic*., leftmargin=2.2em, itemsep=0.3em]
	\item 验证 $N(I-T)=R(I-T)^\perp$.
	\item 假设 $f\in R(I-T)$.证明存在常数 $C$, 使得
	      $\abs{\sigma_n(f)}\le C/n\ \forall n\ge1$.
	\item 由此推出: 对每个 $f\in H$, 有
	      \begin{equation*}
	      \lim_{n\to\infty}\sigma_n(f)=P_{N(I-T)}f.
	      \end{equation*}
	\item 令 $S=\frac12(I+T)$.证明
	      \begin{equation}
	      \abs{u-Su}^2+\abs{Su}^2\le\abs{u}^2\quad\forall u\in H.
	      \tag{1}
	      \end{equation}
	      由此推出
	      \begin{equation*}
	      \sum_{i=0}^\infty\abs{S^iu-S^{i+1}u}^2\le\abs{u}^2\quad\forall u\in H
	      \end{equation*}
	      以及
	      \begin{equation*}
	      \abs{S^n(u-Su)}\le\frac{\abs{u}}{\sqrt{n+1}}
	      \quad\forall u\in H,\ \forall n\ge1.
	      \end{equation*}
	\item 假设 $f\in R(I-T)$.证明存在常数 $C$, 使得
	      $\abs{\mu_n(f)}\le C/\sqrt{n}\ \forall n\ge1$.
	\item 由此推出: 对每个 $f\in H$, 有
	      \begin{equation*}
	      \lim_{n\to\infty}\mu_n(f)=P_{N(I-T)}f.
	      \end{equation*}
\end{enumerate}
\end{exercise}

\begin{exercise}
\label{exer:5.22}
设 $C\subset H$ 为非空闭凸集, $T:C\to C$ 为非线性压缩, 即
\begin{equation*}
\abs{Tu-Tv}\le\abs{u-v}\quad\forall u,v\in C.
\end{equation*}
\begin{enumerate}[label=\arabic*., leftmargin=2.2em, itemsep=0.3em]
	\item 设 $(u_n)$ 为 $C$ 中的序列, 使得
	      \begin{equation*}
	      u_n\weakto u\ \text{弱收敛}\quad\text{且}\quad(u_n-Tu_n)\to f
	      \ \text{强收敛}.
	      \end{equation*}
	      证明 $u-Tu=f$.

	      [提示: 先考虑 $C=H$ 的情形, 并使用不等式
	      $((u-Tu)-(v-Tv),u-v)\ge0\ \forall u,v$.]
	\item 由此推出: 若 $C$ 有界且 $T(C)\subset C$, 则 $T$ 有不动点.

	      [提示: 考虑 $T_\varepsilon u=(1-\varepsilon)Tu+\varepsilon a$,
	      其中 $a\in C$ 固定, $\varepsilon>0$, $\varepsilon\to0$.]
\end{enumerate}
\end{exercise}

\begin{exercise}[Zarantonello 不等式.]
\label{exer:5.23}
设 $T:H\to H$ 为 (非线性) 压缩.假设 $\alpha_1,\alpha_2,\dots,\alpha_n\in\R$
满足 $\alpha_i\ge0\ \forall i$ 且 $\sum_{i=1}^n\alpha_i=1$.假设
$u_1,u_2,\dots,u_n\in H$, 并令
\begin{equation*}
\sigma=\sum_{i=1}^n\alpha_iu_i.
\end{equation*}
证明
\begin{equation*}
\abs{T\sigma-\sum_{i=1}^n\alpha_iTu_i}^2
\le\frac12\sum_{i,j=1}^n\alpha_i\alpha_j
\bigl(\abs{u_i-u_j}^2-\abs{Tu_i-Tu_j}^2\bigr).
\end{equation*}
[提示: 写出
\begin{equation*}
\abs{T\sigma-\sum_{i=1}^n\alpha_iTu_i}^2
=\sum_{i,j=1}^n\alpha_i\alpha_j(T\sigma-Tu_i,T\sigma-Tu_j),
\end{equation*}
并使用恒等式 $(a,b)=\frac12\bigl(\abs{a}^2+\abs{b}^2-\abs{a-b}^2\bigr)$.]

当 $T$ 是等距 (即 $\abs{Tu-Tv}=\abs{u-v}\ \forall u,v\in H$) 时能推出什么?
\end{exercise}

\begin{exercise}[Banach--Saks 性质.]
\label{exer:5.24}
\begin{enumerate}[label=\arabic*., leftmargin=2.2em, itemsep=0.3em]
	\item 假设 $(u_n)$ 为 $H$ 中的序列, 使得 $u_n\weakto0$ 弱收敛.用归纳法
	      构造子列 $(u_{n_j})$, 使得 $u_{n_1}=u_1$ 且
	      \begin{equation*}
	      \abs{(u_{n_j},u_{n_k})}\le\frac1k
	      \quad\forall k\ge2,\ \forall j=1,2,\dots,k-1.
	      \end{equation*}
	      由此推出: 由 $\sigma_p=\frac1p\sum_{j=1}^pu_{n_j}$ 定义的序列
	      $(\sigma_p)$ 当 $p\to\infty$ 时强收敛于 0.

	      [提示: 估计 $\abs{\sigma_p}^2$.]
	\item 假设 $(u_n)$ 为 $H$ 中的有界序列.证明存在子列 $(u_{n_j})$, 使得
	      序列 $\sigma_p=\frac1p\sum_{j=1}^pu_{n_j}$ 当 $p\to\infty$ 时强
	      收敛到一个极限.

	      与推论  \ref{cor:3.8} 及习题 3.4 比较.
\end{enumerate}
\end{exercise}

\begin{exercise}[Opial 引理的若干变体.]
\label{exer:5.25}
设 $K\subset H$ 为非空闭凸集.设 $(u_n)$ 为 $H$ 中的序列, 使得对每个
$v\in K$, 序列 $(\abs{u_n-v})$ 不增.
\begin{enumerate}[label=\arabic*., leftmargin=2.2em, itemsep=0.3em]
	\item 验证序列 $(\dist(u_n,K))$ 不增.
	\item 证明序列 $(P_Ku_n)$ 强收敛到一个极限, 记作 $\ell$.

	      [提示: 使用习题 5.4.]
	\item 这里假设序列 $(u_n)$ 具有性质
	      \begin{equation*}
	      \text{(P)}\qquad
	      \begin{cases}
	      \text{每当某个子列 }(u_{n_k})\text{ 弱收敛到某个极限 }u\in H\text{ 时,}\\
	      \text{都有 }u\in K.
	      \end{cases}
	      \end{equation*}
	      证明 $u_n\weakto\ell$ 弱收敛.
	\item 这里假设 $\bigcup_{\lambda>0}\lambda(K-K)=H$.证明存在某个 $u\in H$,
	      使得 $u_n\weakto u$ 弱收敛且 $P_Ku=\ell$.
	\item 这里假设 $\Int K\ne\emptyset$.证明存在某个 $u\in H$, 使得
	      $u_n\to u$ 强收敛.

	      [提示: 先考虑 $K$ 是单位球的情形, 再考虑一般情形.]
	\item 令 $\sigma_n=\frac1n(u_1+u_2+\cdots+u_n)$, 并假设序列 $(\sigma_n)$
	      满足性质 (P).证明 $\sigma_n\weakto\ell$ 弱收敛.
\end{enumerate}
\end{exercise}

\begin{exercise}
\label{exer:5.26}
假设 $(e_n)$ 是 $H$ 的正交基.
\begin{enumerate}[label=\arabic*., leftmargin=2.2em, itemsep=0.3em]
	\item 验证 $e_n\weakto0$ 弱收敛.

	      设 $(a_n)$ 为 $\R$ 中的有界序列, 令
	      $u_n=\frac1n\sum_{i=1}^na_ie_i$.
	\item 证明 $\abs{u_n}\to0$.
	\item 证明 $\sqrt{n}\,u_n\weakto0$ 弱收敛.
\end{enumerate}
\end{exercise}

\begin{exercise}
\label{exer:5.27}
设 $D\subset H$ 为子集, 使得由 $D$ 张成的线性空间在 $H$ 中稠密.设
$(E_n)_{n\ge1}$ 为 $H$ 中两两正交的闭子空间序列.假设
\begin{equation*}
\sum_{n=1}^\infty\abs{P_{E_n}u}^2=\abs{u}^2\quad\forall u\in D.
\end{equation*}
证明 $H$ 是诸 $E_n$ 的 Hilbert 和.
\end{exercise}

\begin{exercise}
\label{exer:5.28}
假设 $H$ 可分.
\begin{enumerate}[label=\arabic*., leftmargin=2.2em, itemsep=0.3em]
	\item 设 $V\subset H$ 为在 $H$ 中稠密的线性子空间.证明 $V$ 包含 $H$ 的
	      一个正交基.
	\item 设 $(e_n)_{n\ge1}$ 为 $H$ 中的正交序列, 即 $(e_i,e_j)=\delta_{ij}$.
	      证明存在 $H$ 的一个正交基, 它包含 $\bigcup_{n=1}^\infty\{e_n\}$.
\end{enumerate}
\end{exercise}

\begin{exercise}[Grothendieck 的一个引理.]
\label{exer:5.29}
设 $\Omega$ 为 $\abs{\Omega}<\infty$ 的测度空间.设 $E$ 为 $L^p(\Omega)$
($1\le p<\infty$) 的闭子空间.假设 $E\subset L^\infty(\Omega)$.我们的目的
是证明 $\dim E<\infty$.
\begin{enumerate}[label=\arabic*., leftmargin=2.2em, itemsep=0.3em]
	\item 证明存在常数 $C$, 使得
	      \begin{equation*}
	      \norm{u}_\infty\le C\norm{u}_p\quad\forall u\in E.
	      \end{equation*}
	      [提示: 使用推论~
\ref{cor:2.8}.]
	\item 证明存在常数 $M$, 使得
	      \begin{equation*}
	      \norm{u}_\infty\le M\norm{u}_2\quad\forall u\in E.
	      \end{equation*}
	      [提示: 区分 $1\le p\le2$ 与 $2<p<\infty$ 两种情形.]
	\item 由此推出 $E$ 是 $L^2(\Omega)$ 的闭子空间.

	      下文假设 $\dim E=\infty$.设 $(e_n)_{n\ge1}$ 为 $E$ (赋以 $L^2$
	      标量积) 中的正交序列.
	\item 固定任一整数 $k\ge1$.证明存在零测集 $\omega\subset\Omega$, 使得
	      \begin{equation*}
	      \Bigl|\sum_{i=1}^k\alpha_ie_i(x)\Bigr|
	      \le M\Bigl(\sum_{i=1}^k\alpha_i^2\Bigr)^{1/2}
	      \quad\forall x\in\Omega\setminus\omega,\
	      \forall\alpha=(\alpha_1,\alpha_2,\dots,\alpha_k)\in\R^k.
	      \end{equation*}
	      [提示: 先考虑 $\alpha\in\Q^k$ 的情形.]
	\item 由此推出 $\sum_{i=1}^k\abs{e_i(x)}^2\le M^2
	      \ \forall x\in\Omega\setminus\omega$.
	\item 作出结论.
\end{enumerate}
\end{exercise}

\begin{exercise}
\label{exer:5.30}
设 $(e_n)_{n\ge1}$ 为 $H=L^2(0,1)$ 中的正交序列.设 $p(t)$ 为 $H$ 中给定
的函数.
\begin{enumerate}[label=\arabic*., leftmargin=2.2em, itemsep=0.3em]
	\item 证明对每个 $t\in[0,1]$, 有
	      \begin{equation}
	      \sum_{n=1}^\infty\Bigl|\int_0^tp(s)e_n(s)\,ds\Bigr|^2
	      \le\int_0^t\abs{p(s)}^2\,ds.
	      \tag{1}
	      \end{equation}
	\item 由此推出
	      \begin{equation}
	      \sum_{n=1}^\infty\int_0^1\Bigl|\int_0^tp(s)e_n(s)\,ds\Bigr|^2dt
	      \le\int_0^1\abs{p(t)}^2(1-t)\,dt.
	      \tag{2}
	      \end{equation}
	\item 现在假设 $(e_n)_{n\ge1}$ 是 $H$ 的正交基.证明 (1) 与 (2) 变成
	      等式.
	\item 反之, 假设 (2) 中等号成立, 且 $p(t)\ne0$ 几乎处处.证明
	      $(e_n)_{n\ge1}$ 是正交基.
\end{enumerate}
例: $p\equiv1$.
\end{exercise}

\begin{exercise}[Haar 基.]
\label{exer:5.31}
给定整数 $n\ge1$, 把它写成 $n=k+2^p$, 其中 $p\ge0$, $k\ge0$ 是由条件
$k\le2^p-1$ 唯一确定的整数.考虑由下式定义在 $(0,1)$ 上的函数
\begin{equation*}
\varphi_n(t)=\begin{cases}
2^{p/2} & \text{若}\ k2^{-p}<t<\bigl(k+\frac12\bigr)2^{-p},\\
-2^{p/2} & \text{若}\ \bigl(k+\frac12\bigr)2^{-p}<t<(k+1)2^{-p},\\
0 & \text{其他情形}.
\end{cases}
\end{equation*}
令 $\varphi_0\equiv1$, 证明 $(\varphi_n)_{n\ge0}$ 是 $L^2(0,1)$ 的正交基.
\end{exercise}

\begin{exercise}[Rademacher 系与 Walsh 基.]
\label{exer:5.32}
对每个整数 $i\ge0$, 考虑由 $r_i(t)=(-1)^{[2^it]}$ 定义在 $(0,1)$ 上的
函数 $r_i(t)$ (如通常那样, $[x]$ 表示不超过 $x$ 的最大整数).
\begin{enumerate}[label=\arabic*., leftmargin=2.2em, itemsep=0.3em]
	\item 验证 $(r_i)_{i\ge0}$ 是 $L^2(0,1)$ 中的正交序列 (称为
	      \emph{Rademacher 系}).
	\item $(r_i)_{i\ge0}$ 是正交基吗?

	      [提示: 考虑函数 $u=r_1r_2$.]
	\item 给定整数 $n\ge0$, 考虑它的二进制表示
	      $n=\sum_{i=0}^\ell\alpha_i2^i$, 其中 $\alpha_i\in\{0,1\}$.令
	      \begin{equation*}
	      w_n(t)=\prod_{i=0}^\ell r_{i+1}(t)^{\alpha_i}.
	      \end{equation*}
	      证明 $(w_n)_{n\ge0}$ 是 $L^2(0,1)$ 的正交基 (称为
	      \emph{Walsh 基}).注意 $(r_i)_{i\ge0}$ 是 $(w_n)_{n\ge0}$ 的
	      子集.
\end{enumerate}
\end{exercise}

% chapters/ch06.tex
% chapters/ch06.tex
% 第~
\ref{chap:6}~章 \index{jinsuanzi@紧算子 (compact operator)}
紧算子. 自伴紧算子的\index{pufenjie@谱分解 (spectral decomposition)}
谱分解
% 源: 原书 PDF 第 166--188 页 (印刷页 157--179)

\chapter{紧算子. 自伴紧算子的谱分解}
\label{chap:6}

\section{定义. 初等性质. 伴随算子}
\label{sec:6.1}

本章自始至终 (除非另有说明) 用 $E$ 与 $F$ 表示两个 Banach 空间.

\begin{definition}
设 $T\in\LL(E,F)$ 为有界算子.若 $T(B_E)$ 在 $F$ 中 (在强拓扑下) 有紧的
闭包, 则称 $T$ 是\emph{紧算子}.
\end{definition}

由 $E$ 到 $F$ 的一切紧算子构成之集记作 $\KK(E,F)$.为简单起见, 记
$\KK(E)=\KK(E,E)$.

\begin{theorem}
\label{thm:6.1}
集 $\KK(E,F)$ 是 $\LL(E,F)$ 的闭线性子空间 (在由范数
$\norm{\cdot}_{\LL(E,F)}$ 所诱导的拓扑下).
\end{theorem}

\begin{proof}
显然, 两个紧算子之和是紧算子.设 $(T_n)$ 为紧算子序列, $T$ 为有界算子,
使得 $\norm{T_n-T}_{\LL(E,F)}\to0$.我们断言 $T$ 是紧算子.由于 $F$ 完备,
只须验证: 对每个 $\varepsilon>0$, 存在 $T(B_E)$ 的半径 $\varepsilon$ 的有限覆盖
(见例如 J.~R.~Munkres \bref{Munkres}{1} 第~\ref{sec:7.3}~节).取整数 $n$, 使
$\norm{T_n-T}_{\LL(E,F)}<\varepsilon/2$.由于 $T_n(B_E)$ 有紧的闭包, 存在
$T_n(B_E)$ 的半径 $\varepsilon/2$ 的有限覆盖, 即
$T_n(B_E)\subset\bigcup_{i\in I}B(f_i,\varepsilon/2)$.由此推出
\begin{equation*}
T(B_E)\subset\bigcup_{i\in I}B(f_i,\varepsilon).
\end{equation*}
\end{proof}

\begin{definition}
设 $T\in\LL(E,F)$.若 $T$ 的值域 $R(T)$ 是有限维的, 则称 $T$ 是
\emph{有限秩算子}.
\end{definition}

显然, 任何有限秩算子都是紧算子, 因此我们有下述结论.

\begin{corollary}
\label{cor:6.2}
设 $(T_n)$ 为有限秩算子序列, $T\in\LL(E,F)$ 满足
$\norm{T_n-T}_{\LL(E,F)}\to0$.则 $T\in\KK(E,F)$.
\end{corollary}

\begin{starremark}
\label{rem:6.1}
著名的「逼近问题」(Banach, Grothendieck) 讨论的是推论  \ref{cor:6.2} 的逆:
给定紧算子 $T$, 是否总存在有限秩算子序列 $(T_n)$, 使得
$\norm{T_n-T}_{\LL(E,F)}\to0$? 这个问题长期悬而未决, 直到 1972 年
P.~Enflo \bref{Enflo}{1} 找到了反例.原来的构造相当复杂, 其后又找到了更简单的
例子, 例如取 $F$ 为 $\ell^p$ 的某个闭子空间 (对任意 $1<p<\infty$, $p\ne2$).
有兴趣的读者可在 J.~Lindenstrauss--L.~Tzafriri \bref{LindenstraussTzafriri}{2}
中找到关于逼近问题的详细讨论.注意, 在某些特殊情形下逼近问题的答案是肯定的
——例如当 $F$ 是 Hilbert 空间时.事实上, 令 $K=T(B_E)$.给定 $\varepsilon>0$,
存在 $K$ 的半径 $\varepsilon$ 的有限覆盖, 即
$K\subset\bigcup_{i\in I}B(f_i,\varepsilon)$.设 $G$ 为诸 $f_i$ 张成的向量空间,
并令 $T_\varepsilon=P_GT$, 于是 $T_\varepsilon$ 是有限秩的.我们断言
$\norm{T_\varepsilon-T}_{\LL(E,F)}<2\varepsilon$.对每个 $x\in B_E$, 存在
$i_0\in I$, 使得
\begin{equation}
\norm{Tx-f_{i_0}}<\varepsilon.
\label{eq:6.1}
\end{equation}
从而
\begin{equation*}
\norm{P_GTx-P_Gf_{i_0}}<\varepsilon,
\end{equation*}
即
\begin{equation}
\norm{P_GTx-f_{i_0}}<\varepsilon.
\label{eq:6.2}
\end{equation}
合并 \eqref{eq:6.1} 与 \eqref{eq:6.2}, 得到
\begin{equation*}
\norm{P_GTx-Tx}<2\varepsilon\quad\forall x\in B_E,
\end{equation*}
即
\begin{equation*}
\norm{T_\varepsilon-T}_{\LL(E,F)}<2\varepsilon.
\end{equation*}
[更一般地可以看出: 若 $F$ 有 Schauder 基, 则对每个空间 $E$ 以及由 $E$ 到
$F$ 的每个紧算子, 逼近问题的答案都是肯定的.]

联系到逼近问题, 我们提一种在非线性分析中非常有用的技巧, 即用有限秩的
非线性映射来逼近一个连续映射 (线性或非线性).设 $X$ 为拓扑空间, $F$ 为
Banach 空间, $T:X\to F$ 为连续映射, 使得 $T(X)$ 在 $F$ 中有紧的闭包.
我们断言: 对每个 $\varepsilon>0$, 存在有限秩的连续映射 $T_\varepsilon:X\to F$,
使得
\begin{equation}
\norm{T_\varepsilon(x)-T(x)}<\varepsilon\quad\forall x\in X.
\label{eq:6.3}
\end{equation}
事实上, 由于 $K=T(X)$ 紧, 存在 $K$ 的有限覆盖, 即
$K\subset\bigcup_{i\in I}B(f_i,\varepsilon/2)$.令
\begin{equation*}
T_\varepsilon(x)=\frac{\sum_{i\in I}q_i(x)f_i}{\sum_{i\in I}q_i(x)},
\quad\text{其中}\ q_i(x)=\max\{\varepsilon-\norm{Tx-f_i},0\};
\end{equation*}
显然 $T_\varepsilon$ 满足 \eqref{eq:6.3}.

这种逼近方式非常有用, 例如可用它由 Brouwer 不动点定理推出 Schauder 不动点
定理 (见例如 K.~Deimling \bref{Deimling}{1}, A.~Granas--J.~Dugundji \bref{GranasDugundji}{1}, J.~Franklin \bref{Franklin}{1}, 以及习题  \ref{exer:6.26}).类似的构造与
Schauder 不动点定理相结合, 还被 Lomonosov 以出人意料的方式用来证明一大类
线性算子存在非平凡不变子空间 (见例如 C.~Pearcy \bref{Pearcy}{1}, N.~Akhiezer--I.~Glazman \bref{AkhiezerGlazman}{1}, A.~Granas--J.~Dugundji \bref{GranasDugundji}{1}, 以及问题
42).另一个基于 Schauder 不动点定理而有简单证明的线性结果
是 Krein--Rutman 定理 (见定理  \ref{thm:6.13} 与问题 41).
\end{starremark}

\begin{proposition}
\label{prop:6.3}
设 $E$, $F$, $G$ 为三个 Banach 空间.设 $T\in\LL(E,F)$,
$S\in\KK(F,G)$ [相应地, $T\in\KK(E,F)$, $S\in\LL(F,G)$].则
$S\circ T\in\KK(E,G)$.
\end{proposition}

证明是显然的.

\begin{theorem}[Schauder]
\label{thm:6.4}
若 $T\in\KK(E,F)$, 则 $T^\star\in\KK(F^\star,E^\star)$.反之亦然.
\end{theorem}

\begin{proof}
我们要证明 $T^\star(B_{F^\star})$ 在 $E^\star$ 中有紧的闭包.设 $(v_n)$ 为
$B_{F^\star}$ 中的序列.我们断言 $(T^\star(v_n))$ 有收敛子列.令
$K=T(B_E)$; 这是一个紧度量空间.考虑由
\begin{equation*}
H=\{\phi_n:x\in K\mapsto\dual{v_n}{x};\ n=1,2,\dots\}
\end{equation*}
定义的集 $H\subset C(K)$.Ascoli--Arzelà 定理 (定理  \ref{thm:4.25}) 的假设
成立.于是存在子列 (仍记作 $\phi_{n_k}$), 它在 $K$ 上一致收敛于某个连续函数
$\phi\in C(K)$.特别地, 我们有
\begin{equation*}
\sup_{u\in B_E}\abs{\dual{v_{n_k}}{Tu}-\phi(Tu)}\xrightarrow[k\to\infty]{}0.
\end{equation*}
从而
\begin{equation*}
\sup_{u\in B_E}\abs{\dual{v_{n_k}}{Tu}-\dual{v_{n_\ell}}{Tu}}
\xrightarrow[k,\ell\to\infty]{}0,
\end{equation*}
即 $\norm{T^\star v_{n_k}-T^\star v_{n_\ell}}_{E^\star}
\xrightarrow[k,\ell\to\infty]{}0$.于是 $T^\star v_{n_k}$ 在 $E^\star$ 中
收敛.

反之, 设 $T^\star\in\KK(F^\star,E^\star)$.由第一部分, 我们已经知道
$T^{\star\star}\in\KK(E^{\star\star},F^{\star\star})$.特别地,
$T^{\star\star}(B_E)$ 在 $F^{\star\star}$ 中有紧的闭包.但
$T(B_E)=T^{\star\star}(B_E)$, 且 $F$ 在 $F^{\star\star}$ 中闭.因此
$T(B_E)$ 在 $F$ 中有紧的闭包.
\end{proof}

\begin{remark}
\label{rem:6.2}
设 $E$ 与 $F$ 为两个 Banach 空间, $T\in\KK(E,F)$.若 $(u_n)$ 在 $E$ 中
弱收敛于 $u$, 则 $(Tu_n)$ 强收敛于 $Tu$.若 $E$ 自反, 逆命题也成立
(见习题  \ref{exer:6.7}).
\end{remark}

\section{Riesz--Fredholm 理论}
\label{sec:6.2}

我们首先给出一些有用的预备结果.

\begin{lemma}[Riesz 引理]
\label{lem:6.1}
设 $E$ 为赋范向量空间 (n.v.s.), $M\subset E$ 为闭线性子空间, 且 $M\ne E$.
则
\begin{equation*}
\forall\varepsilon>0\ \ \exists u\in E\ \text{使得}\ \norm{u}=1\
\text{且}\ \dist(u,M)\ge1-\varepsilon.
\end{equation*}
\end{lemma}

\begin{proof}
取 $v\in E$, $v\notin M$.由于 $M$ 闭, 故
\begin{equation*}
d=\dist(v,M)>0.
\end{equation*}
任取 $m_0\in M$, 使得
\begin{equation*}
d\le\norm{v-m_0}\le d/(1-\varepsilon).
\end{equation*}
则
\begin{equation*}
u=\frac{v-m_0}{\norm{v-m_0}}
\end{equation*}
满足所需性质.事实上, 对每个 $m\in M$, 我们有
\begin{equation*}
\norm{u-m}=\left\|\frac{v-m_0}{\norm{v-m_0}}-m\right\|
\ge\frac{d}{\norm{v-m_0}}\ge1-\varepsilon,
\end{equation*}
因为 $m_0+\norm{v-m_0}m\in M$.
\end{proof}

\begin{remark}
\label{rem:6.3}
若 $M$ 是有限维的 (或更一般地, 若 $M$ 自反), 则在引理  \ref{lem:6.1} 中可
取 $\varepsilon=0$.但一般说来这并不成立 (见习题  \ref{exer:1.17}).
\end{remark}

\begin{keytheorem}[Riesz]
\label{thm:6.5}
设 $E$ 为赋范向量空间, 且 $B_E$ 紧.则 $E$ 是有限维的.
\end{keytheorem}

\begin{proof}
反设 $E$ 是无穷维的.于是存在 $E$ 的有限维子空间序列 $(E_n)$, 使得
$E_{n-1}\subset E_n$ 且 $E_{n-1}\ne E_n$.由引理  \ref{lem:6.1}, 存在序列
$(u_n)$, $u_n\in E_n$, 使得 $\norm{u_n}=1$ 且
$\dist(u_n,E_{n-1})\ge1/2$.特别地, 当 $m<n$ 时
$\norm{u_n-u_m}\ge1/2$.于是 $(u_n)$ 没有收敛子列, 这与 $B_E$ 紧的假设
矛盾.
\end{proof}

\begin{keytheorem}[Fredholm 二择一定理]
\label{thm:6.6}
设 $T\in\KK(E)$.则
\begin{enumerate}[label=(\alph*), leftmargin=2.2em, itemsep=0.1em]
	\item $N(I-T)$ 是有限维的,
	\item $R(I-T)$ 闭, 且更确切地, $R(I-T)=N(I-T^\star)^\perp$,
	\item $N(I-T)=\{0\}\iff R(I-T)=E$,
	\item $\dim N(I-T)=\dim N(I-T^\star)$.
\end{enumerate}
\end{keytheorem}

\begin{remark}
\label{rem:6.4}
Fredholm 二择一定理讨论方程 $u-Tu=f$ 的可解性.它说:
\begin{itemize}[leftmargin=1.6em, itemsep=0.15em]
	\item 或者对每个 $f\in E$, 方程 $u-Tu=f$ 有唯一解,
	\item 或者齐次方程 $u-Tu=0$ 有 $n$ 个线性无关解, 而此时非齐次方程
	      $u-Tu=f$ 可解当且仅当 $f$ 满足 $n$ 个正交条件, 即
	      $f\in N(I-T^\star)^\perp$.
\end{itemize}
\end{remark}

\begin{remark}
\label{rem:6.5}
性质 (c) 在有限维空间中是我们熟悉的.若 $\dim E<\infty$, 则由 $E$ 到自身的
线性算子单射 (= 一一) 当且仅当它满射 (= 到上).然而在无穷维空间中, 一个有界
算子可以单射而不满射, 反之亦然, 例如 $\ell^2$ 中的右移位 (相应地, 左移位)
(见注记 \ref{rem:6.6}).因此, 断言 (c) 是形如 $I-T$ (其中 $T\in\KK(E)$)
的算子的一个值得注意的性质.
\end{remark}

\begin{proof}
\begin{enumerate}[label=(\alph*), leftmargin=2.2em, itemsep=0.3em]
	\item 令 $E_1=N(I-T)$.则 $B_{E_1}\subset T(B_E)$, 从而 $B_{E_1}$ 紧.
	      由定理  \ref{thm:6.5}, $E_1$ 必为有限维的.
	\item 设 $f_n=u_n-Tu_n\to f$.我们要证明 $f\in R(I-T)$.令
	      $d_n=\dist(u_n,N(I-T))$.由于 $N(I-T)$ 有限维, 存在
	      $v_n\in N(I-T)$, 使得 $d_n=\norm{u_n-v_n}$.我们有
	      \begin{equation}
	      f_n=(u_n-v_n)-T(u_n-v_n).
	      \label{eq:6.4}
	      \end{equation}
	      我们断言 $\norm{u_n-v_n}$ 保持有界.若不然, 则存在子列使得
	      $\norm{u_{n_k}}-v_{n_k}\to\infty$.令
	      $w_n=(u_n-v_n)/\norm{u_n-v_n}$.由 \eqref{eq:6.4} 看出
	      $w_{n_k}-Tw_{n_k}\to0$.再取子列 (仍简记作 $w_{n_k}$), 可设
	      $Tw_{n_k}\to z$.于是 $w_{n_k}\to z$ 且 $z\in N(I-T)$, 从而
	      $\dist(w_{n_k},N(I-T))\to0$.另一方面,
	      \begin{equation*}
	      \dist(w_n,N(I-T))
	      =\frac{\dist(u_n,N(I-T))}{\norm{u_n-v_n}}=1
	      \end{equation*}
	      (因为 $v_n\in N(I-T)$); 矛盾.

	      于是 $\norm{u_n-v_n}$ 保持有界, 而 $T$ 是紧算子, 故可抽出子列使得
	      $T(u_{n_k}-v_{n_k})$ 收敛于某个极限 $\ell$.由 \eqref{eq:6.4}
	      推出 $u_{n_k}-v_{n_k}\to f+\ell$.令 $g=f+\ell$, 则 $g-Tg=f$,
	      即 $f\in R(I-T)$.这就完成了算子 $(I-T)$ 有闭值域这一事实的证明.
	      因此可以应用定理  \ref{thm:2.19}, 推出
	      \begin{equation*}
	      R(I-T)=N(I-T^\star)^\perp,\qquad
	      R(I-T^\star)=N(I-T)^\perp.
	      \end{equation*}
	\item 我们先证蕴涵 $\Rightarrow$.反设
	      \begin{equation*}
	      E_1=R(I-T)\ne E.
	      \end{equation*}
	      则 $E_1$ 是 Banach 空间且 $T(E_1)\subset E_1$.于是
	      $T|_{E_1}\in\KK(E_1)$, 且 $E_2=(I-T)(E_1)$ 是 $E_1$ 的闭子空间.
	      又 $E_2\ne E_1$ (因为 $(I-T)$ 单射).令 $E_n=(I-T)^n(E)$, 得到
	      一列 (严格) 递减的闭子空间.用 Riesz 引理可构造序列 $(u_n)$, 使得
	      $u_n\in E_n$, $\norm{u_n}=1$ 且 $\dist(u_n,E_{n+1})\ge1/2$.我们有
	      \begin{equation*}
	      Tu_n-Tu_m=-(u_n-Tu_n)+(u_m-Tu_m)+(u_n-u_m).
	      \end{equation*}
	      注意若 $n>m$, 则 $E_{n+1}\subset E_n\subset E_{m+1}\subset E_m$,
	      因此
	      \begin{equation*}
	      -(u_n-Tu_n)+(u_m-Tu_m)+u_n\in E_{m+1}.
	      \end{equation*}
	      由此推出 $\norm{Tu_n-Tu_m}\ge\dist(u_m,E_{m+1})\ge1/2$.这不可能,
	      因为 $T$ 是紧算子.于是我们证明了 $R(I-T)=E$.

	      反之, 设 $R(I-T)=E$.由推论  \ref{cor:2.18}, 我们知道
	      $N(I-T^\star)=R(I-T)^\perp=\{0\}$.由于 $T^\star\in\KK(E^\star)$,
	      可对之应用上一步, 推出 $R(I-T^\star)=E^\star$.再次使用推论
	       \ref{cor:2.18}, 得到 $N(I-T)=R(I-T^\star)^\perp=\{0\}$.
	\item 令 $d=\dim N(I-T)$, $d^\star=\dim N(I-T^\star)$.我们先证
	      $d^\star\le d$.若不然, 设 $d<d^\star$.由于 $N(I-T)$ 有限维,
	      它在 $E$ 中有补空间 (见第~\ref{sec:2.4}~节例 1).于是存在由 $E$ 到 $N(I-T)$
	      的连续投影算子 $P$.另一方面, $R(I-T)=N(I-T^\star)^\perp$ 有有限
	      余维 $d^\star$ (见第~\ref{sec:2.4}~节例 2), 从而它在 $E$ 中有维数为 $d^\star$
	      的补空间, 记作 $F$.由于 $d<d^\star$, 存在线性映射
	      $\Lambda:N(I-T)\to F$, 它单射但不满射.令 $S=T+\Lambda\circ P$.则
	      $S\in\KK(E)$, 因为 $\Lambda\circ P$ 是有限秩的.

	      我们断言 $N(I-S)=\{0\}$.事实上, 若
	      \begin{equation*}
	      0=u-Su=(u-Tu)-(\Lambda\circ Pu),
	      \end{equation*}
	      则
	      \begin{equation*}
	      u-Tu=0\quad\text{且}\quad \Lambda\circ Pu=0,
	      \end{equation*}
	      即 $u\in N(I-T)$ 且 $\Lambda u=0$.因此 $u=0$.

	      对算子 $S$ 应用 (c), 得到 $R(I-S)=E$.这是荒谬的, 因为存在
	      $f\in F$ 使 $f\notin R(\Lambda)$, 于是方程 $u-Su=f$ 无解.

	      于是我们证明了 $d^\star\le d$.把这一事实应用于 $T^\star$, 得到
	      \begin{equation*}
	      \dim N(I-T^{\star\star})\le\dim N(I-T^\star)\le\dim N(I-T).
	      \end{equation*}
	      但 $N(I-T^{\star\star})\supset N(I-T)$, 因此 $d=d^\star$.
\end{enumerate}
\end{proof}

\section{紧算子的谱}
\label{sec:6.3}

下面是一些重要的定义.

\begin{definition}
设 $T\in\LL(E)$.\emph{\index{yujieji@预解集 (resolvent set)}
预解集} (记作 $\rho(T)$) 定义为
\begin{equation*}
\rho(T)=\{\lambda\in\R;\ (T-\lambda I)\ \text{是由}\ E\ \text{到}\ E\ \text{的双射}\}.
\end{equation*}
\emph{谱} (记作 $\sigma(T)$) 是预解集的补集, 即 $\sigma(T)=\R\setminus\rho(T)$.
若 $N(T-\lambda I)\ne\{0\}$, 则称实数 $\lambda$ 是 $T$ 的\emph{\index{tezhengzhi@特征值 (eigenvalue)}
特征值};
$N(T-\lambda I)$ 是对应的\emph{特征空间}.全部特征值之集记作 $EV(T)$.
\footnote{有些作者用 $\sigma_p(T)$ (点谱) 代替 $EV(T)$.}
\end{definition}

记住下面一点是有用的: 若 $\lambda\in\rho(T)$, 则
$(T-\lambda I)^{-1}\in\LL(E)$ (见推论  \ref{cor:2.7}).

\begin{remark}
\label{rem:6.6}
显然 $EV(T)\subset\sigma(T)$.一般说来, 这个包含可以是严格的:
\footnote{当然, 若 $E$ 有限维, 则 $EV(T)=\sigma(T)$.}
可能存在某个 $\lambda$, 使得
\begin{equation*}
N(T-\lambda I)=\{0\}\quad\text{且}\quad R(T-\lambda I)\ne E
\end{equation*}
(这样的 $\lambda$ 属于谱但不是特征值).例如考虑 $E=\ell^2$ 中的右移位,
即 $Tu=(0,u_1,u_2,\dots)$, 其中 $u=(u_1,u_2,u_3,\dots)$.则
$0\in\sigma(T)$, 而 $0\notin EV(T)$.事实上, 此时 $EV(T)=\emptyset$, 而
$\sigma(T)=[-1,+1]$ (见习题  \ref{exer:6.18}).当然, 在有限维或无穷维空间中都
可能出现 $EV(T)=\sigma(T)=\emptyset$; 例如考虑 $\R^2$ 中旋转 $\pi/2$ 的
变换, 或在 $\ell^2$ 中考虑算子 $Tu=(-u_2,u_1,-u_4,u_3,\dots)$.若在 $\CC$
上的向量空间中讨论 (见第~\ref{sec:11.4}~节), 情况就完全不同了; 在 $\CC$ 上的空间中,
特征值与谱的研究要有趣得多.众所周知, 在 $\CC$ 上的有限维空间中,
$EV(T)=\sigma(T)\ne\emptyset$ (它们是特征多项式的根).在 $\CC$ 上的无穷维
空间中, 一个非平凡的结果断言 $\sigma(T)$ 总是非空的 (见第~\ref{sec:11.4}~节).
然而, $EV(T)=\emptyset$ 是可能发生的 (例如取 $E=\ell^2$ 中的右移位).
\end{remark}

\begin{proposition}
\label{prop:6.7}
有界算子 $T$ 的谱 $\sigma(T)$ 紧, 且
$\sigma(T)\subset[-\norm{T},+\norm{T}]$.
\end{proposition}

\begin{proof}
设 $\lambda\in\R$ 满足 $\abs{\lambda}>\norm{T}$.我们将证明 $T-\lambda I$ 是
双射, 从而 $\sigma(T)\subset[-\norm{T},+\norm{T}]$.给定 $f\in E$, 方程
$Tu-\lambda u=f$ 有唯一解, 因为它可写成 $u=\lambda^{-1}(Tu-f)$, 而压缩映射
原理 (定理  \ref{thm:5.7}) 适用.

下面证明 $\rho(T)$ 是开集.设 $\lambda_0\in\rho(T)$.给定 $\lambda\in\R$
(接近 $\lambda_0$) 与 $f\in E$, 我们试图求解
\begin{equation}
Tu-\lambda u=f.
\label{eq:6.5}
\end{equation}
方程 \eqref{eq:6.5} 可写成
\begin{equation*}
Tu-\lambda_0u=f+(\lambda-\lambda_0)u,
\end{equation*}
即
\begin{equation}
u=(T-\lambda_0I)^{-1}\bigl[f+(\lambda-\lambda_0)u\bigr].
\label{eq:6.6}
\end{equation}
再次应用压缩映射原理, 可见当
$\abs{\lambda-\lambda_0}\,\norm{(T-\lambda_0I)^{-1}}<1$ 时 \eqref{eq:6.6}
有解.
\end{proof}

\begin{keytheorem}
\label{thm:6.8}
设 $T\in\KK(E)$ 且 $\dim E=\infty$, 则我们有:
\begin{enumerate}[label=(\alph*), leftmargin=2.2em, itemsep=0.1em]
	\item $0\in\sigma(T)$,
	\item $\sigma(T)\setminus\{0\}=EV(T)\setminus\{0\}$,
	\item 下述情形之一成立:
	      \begin{itemize}[leftmargin=1.6em, itemsep=0.15em]
	      	\item $\sigma(T)=\{0\}$,
	      	\item $\sigma(T)\setminus\{0\}$ 是有限集,
	      	\item $\sigma(T)\setminus\{0\}$ 是收敛于 $0$ 的序列.
	      \end{itemize}
\end{enumerate}
\end{keytheorem}

\begin{proof}
\begin{enumerate}[label=(\alph*), leftmargin=2.2em, itemsep=0.3em]
	\item 若不然, 设 $0\notin\sigma(T)$.则 $T$ 是双射, 而
	      $I=T\circ T^{-1}$ 是紧的.于是 $B_E$ 紧, 由定理  \ref{thm:6.5} 得
	      $\dim E<\infty$; 矛盾.
	\item 设 $\lambda\in\sigma(T)$, $\lambda\ne0$.我们要证明 $\lambda$ 是
	      特征值.若不然, 设 $N(T-\lambda I)=\{0\}$.则由定理  \ref{thm:6.6}(c),
	      我们知道 $R(T-\lambda I)=E$, 因此 $\lambda\in\rho(T)$; 矛盾.
\end{enumerate}

为证明断言 (c), 我们将使用下述引理.

\begin{lemma}
\label{lem:6.2}
设 $T\in\KK(E)$, $(\lambda_n)_{n\ge1}$ 为互不相同的实数序列, 使得
\begin{equation*}
\lambda_n\to\lambda\quad\text{且}\quad
\lambda_n\in\sigma(T)\setminus\{0\}\ \ \forall n.
\end{equation*}
则 $\lambda=0$.换言之, $\sigma(T)\setminus\{0\}$ 的所有点都是孤立点.
\end{lemma}

\begin{proof}
我们知道 $\lambda_n\in EV(T)$; 取 $e_n\ne0$ 使 $(T-\lambda_nI)e_n=0$.令
$E_n$ 为由 $\{e_1,e_2,\dots,e_n\}$ 张成的空间.我们断言对一切 $n$ 都有
$E_n\subset E_{n+1}$, $E_n\ne E_{n+1}$.只须验证对一切 $n$, 向量
$e_1,e_2,\dots,e_n$ 线性无关.对 $n$ 归纳证明.设此结论对 $n$ 成立, 并设
$e_{n+1}=\sum_{i=1}^n\alpha_ie_i$.则
\begin{equation*}
Te_{n+1}=\sum_{i=1}^n\alpha_i\lambda_ie_i
=\sum_{i=1}^n\alpha_i\lambda_{n+1}e_i.
\end{equation*}
由此推出 $\alpha_i(\lambda_i-\lambda_{n+1})=0$, $i=1,2,\dots,n$, 从而
$\alpha_i=0$, $i=1,2,\dots,n$; 矛盾.于是我们证明了
$E_n\subset E_{n+1}$, $E_n\ne E_{n+1}$ 对一切 $n$ 成立.

应用 Riesz 引理 (引理  \ref{lem:6.1}), 可构造序列 $(u_n)_{n\ge1}$, 使得
$u_n\in E_n$, $\norm{u_n}=1$ 且 $\dist(u_n,E_{n-1})\ge1/2$ 对一切 $n\ge2$
成立.当 $2\le m<n$ 时, 有
\begin{equation*}
E_{m-1}\subset E_m\subset E_{n-1}\subset E_n.
\end{equation*}
另一方面, 显然 $(T-\lambda_nI)E_n\subset E_{n-1}$.于是
\begin{equation*}
\left\|\frac{Tu_n}{\lambda_n}-\frac{Tu_m}{\lambda_m}\right\|
=\left\|\frac{(Tu_n-\lambda_nu_n)}{\lambda_n}
-\frac{(Tu_m-\lambda_mu_m)}{\lambda_m}+u_n-u_m\right\|
\ge\dist(u_n,E_{n-1})\ge1/2.
\end{equation*}
若 $\lambda_n\to\lambda$ 且 $\lambda\ne0$, 则矛盾, 因为 $(Tu_n)$ 有收敛子列.
\end{proof}

定理  \ref{thm:6.8} 的证明 (续).对每个整数 $n\ge1$, 集
\begin{equation*}
\sigma(T)\cap\{\lambda\in\R;\ \abs{\lambda}\ge1/n\}
\end{equation*}
或空或有限 (若它含无穷多个互不相同的点, 则存在子列收敛于某个满足
$\abs{\lambda}\ge1/n$ 的 $\lambda$——因为 $\sigma(T)$ 紧——这与引理
 \ref{lem:6.2} 矛盾).因此, 若 $\sigma(T)\setminus\{0\}$ 含无穷多个互不相同的
点, 我们可把它们排成一个趋于 $0$ 的序列.
\end{proof}

\begin{remark}
\label{rem:6.7}
给定任何收敛于 $0$ 的序列 $(\alpha_n)$, 存在紧算子 $T$ 使得
$\sigma(T)=(\alpha_n)\cup\{0\}$.在 $\ell^2$ 中, 只须考虑由
$Tu=(\alpha_1u_1,\alpha_2u_2,\dots,\alpha_nu_n,\dots)$ (其中
$u=(u_1,u_2,\dots,u_n,\dots)$) 定义的乘法算子 $T$.注意 $T$ 是紧的,
因为 $T$ 是有限秩算子的极限.更确切地说, 令
$T_nu=(\alpha_1u_1,\alpha_2u_2,\dots,\alpha_nu_n,0,0,\dots)$; 则
$\norm{T_n-T}\to0$.在这个例子中我们还看到, $0$ 可以属于也可以不属于
$EV(T)$.另一方面, 若 $0\in EV(T)$, 则对应的特征空间, 即 $N(T)$, 可以是
有限维或无穷维的.
\end{remark}

\section{自伴紧算子的谱分解}
\label{sec:6.4}

下文假设 $E=H$ 是 Hilbert 空间, $T\in\LL(H)$.把 $H^\star$ 与 $H$ 等同,
我们可把 $T^\star$ 看作由 $H$ 到自身的算子.

\begin{definition}
若 $T^\star=T$, 即
\begin{equation*}
(Tu,v)=(u,Tv)\quad\forall u,v\in H,
\end{equation*}
则称有界算子 $T\in\LL(H)$ 是\emph{自伴}的.
\end{definition}

\begin{proposition}
\label{prop:6.9}
设 $T\in\LL(H)$ 为自伴算子.令
\begin{equation*}
m=\inf_{\substack{u\in H\\ \abs{u}=1}}(Tu,u)
\quad\text{与}\quad
M=\sup_{\substack{u\in H\\ \abs{u}=1}}(Tu,u).
\end{equation*}
则 $\sigma(T)\subset[m,M]$, $m\in\sigma(T)$, $M\in\sigma(T)$.此外,
$\norm{T}=\max\{\abs{m},\abs{M}\}$.
\end{proposition}

\begin{proof}
设 $\lambda>M$; 我们将证明 $\lambda\in\rho(T)$.我们有
\begin{equation*}
(Tu,u)\le M\abs{u}^2\quad\forall u\in H,
\end{equation*}
从而
\begin{equation*}
(\lambda u-Tu,u)\ge(\lambda-M)\abs{u}^2=\alpha\abs{u}^2
\quad\forall u\in H,\ \text{其中}\ \alpha>0.
\end{equation*}
应用 Lax--Milgram 定理 (推论  \ref{cor:5.8}), 推出 $\lambda I-T$ 是双射,
因此 $\lambda\in\rho(T)$.类似地, 任何 $\lambda<m$ 都属于 $\rho(T)$, 因此
$\sigma(T)\subset[m,M]$.

下面证明 $M\in\sigma(T)$ (证明 $m\in\sigma(T)$ 类似).双线性形式
$a(u,v)=(Mu-Tu,v)$ 对称且满足
\begin{equation*}
a(v,v)\ge0\quad\forall v\in H.
\end{equation*}
因此它满足 Cauchy--Schwarz 不等式
\begin{equation*}
\abs{a(u,v)}\le a(u,u)^{1/2}a(v,v)^{1/2}\quad\forall u,v\in H,
\end{equation*}
即
\begin{equation*}
\abs{(Mu-Tu,v)}\le(Mu-Tu,u)^{1/2}(Mv-Tv,v)^{1/2}
\quad\forall u,v\in H.
\end{equation*}
由此推出
\begin{equation}
\abs{Mu-Tu}\le C(Mu-Tu,u)^{1/2}\quad\forall u\in H.
\label{eq:6.7}
\end{equation}
由 $M$ 的定义, 存在序列 $(u_n)$ 使得 $\abs{u_n}=1$ 且
$(Tu_n,u_n)\to M$.由 \eqref{eq:6.7} 推出 $\abs{Mu_n-Tu_n}\to0$, 从而
$M\in\sigma(T)$ (因为若 $M\in\rho(T)$, 则
$u_n=(MI-T)^{-1}(Mu_n-Tu_n)\to0$, 这不可能).

最后, 我们证明 $\norm{T}=\mu$, 其中 $\mu=\max\{\abs{m},\abs{M}\}$.对
$u,v\in H$ 写出
\begin{align*}
(T(u+v),u+v)&=(Tu,u)+(Tv,v)+2(Tu,v),\\
(T(u-v),u-v)&=(Tu,u)+(Tv,v)-2(Tu,v).
\end{align*}
于是
\begin{equation*}
4(Tu,v)=(T(u+v),u+v)-(T(u-v),u-v)
\le M\abs{u+v}^2-m\abs{u-v}^2,
\end{equation*}
从而
\begin{equation*}
4\abs{(Tu,v)}\le\mu\bigl(\abs{u+v}^2+\abs{u-v}^2\bigr)
=2\mu\bigl(\abs{u}^2+\abs{v}^2\bigr).
\end{equation*}
用 $\alpha>0$ 代替 $v$ 为 $\alpha v$, 得到
\begin{equation*}
4\abs{(Tu,v)}\le2\mu\left(\frac{\abs{u}^2}{\alpha}+\alpha\abs{v}^2\right).
\end{equation*}
下一步对 $\alpha$ 极小化右端, 即取 $\alpha=\abs{u}/\abs{v}$, 然后得到
\begin{equation*}
\abs{(Tu,v)}\le\mu\abs{u}\,\abs{v}\quad\forall u,v,
\end{equation*}
从而 $\norm{T}\le\mu$.另一方面, 显然 $\abs{(Tu,u)}\le\norm{T}\abs{u}^2$,
故 $\abs{m}\le\norm{T}$ 且 $\abs{M}\le\norm{T}$, 从而 $\mu\le\norm{T}$.
\end{proof}

\begin{corollary}
\label{cor:6.10}
设 $T\in\LL(H)$ 为自伴算子, 且 $\sigma(T)=\{0\}$.则 $T=0$.
\end{corollary}

我们最后的结论是一个基本结果.它断言每个紧自伴算子都可以在某个适当的基下
对角化.

\begin{keytheorem}
\label{thm:6.11}
设 $H$ 为可分 Hilbert 空间, $T$ 为紧自伴算子.则存在由 $T$ 的特征向量构成的
Hilbert 基.
\end{keytheorem}

\begin{proof}
设 $(\lambda_n)_{n\ge1}$ 为 $T$ 的全部 (互不相同的) 非零特征值.令
\begin{equation*}
\lambda_0=0,\qquad E_0=N(T),\qquad E_n=N(T-\lambda_nI).
\end{equation*}
回忆
\begin{equation*}
0\le\dim E_0\le\infty\quad\text{且}\quad0<\dim E_n<\infty.
\end{equation*}
我们断言 $H$ 是诸 $E_n$ ($n=0,1,2,\dots$) 的 Hilbert 和 (按第~\ref{sec:5.4}~节的
意义):

(i) 空间 $(E_n)_{n\ge0}$ 两两正交.

事实上, 若 $u\in E_m$, $v\in E_n$ 且 $m\ne n$, 则
\begin{equation*}
Tu=\lambda_mu\quad\text{且}\quad Tv=\lambda_nv,
\end{equation*}
从而
\begin{equation*}
(Tu,v)=\lambda_m(u,v)=(u,Tv)=\lambda_n(u,v).
\end{equation*}
因此 $(u,v)=0$.

(ii) 设 $F$ 为诸空间 $(E_n)_{n\ge0}$ 张成的向量空间.我们将证明 $F$ 在 $H$
中稠密.

显然 $T(F)\subset F$.由此推出 $T(F^\perp)\subset F^\perp$; 事实上, 给定
$u\in F^\perp$, 我们有
\begin{equation*}
(Tu,v)=(u,Tv)=0\quad\forall v\in F,
\end{equation*}
故 $Tu\in F^\perp$.$T$ 在 $F^\perp$ 上的限制记作 $T_0$.它是 $F^\perp$ 上的
紧自伴算子.我们断言 $\sigma(T_0)=\{0\}$.若不然, 设某个 $\lambda\ne0$ 属于
$\sigma(T_0)$.由于 $\lambda\in EV(T_0)$, 存在 $u\in F^\perp$, $u\ne0$, 使得
$T_0u=\lambda u$.于是 $\lambda$ 是 $T$ 的特征值之一, 设为
$\lambda=\lambda_n$ ($n\ge1$).从而 $u\in E_n\subset F$.由于
$u\in F^\perp\cap F$, 推出 $u=0$; 矛盾.应用推论  \ref{cor:6.10}, 推出
$T_0=0$, 即 $T$ 在 $F^\perp$ 上为零.由此推出 $F^\perp\subset N(T)$.另一方面,
$N(T)\subset F$, 因而 $F^\perp\subset F$.这意味着 $F^\perp=\{0\}$, 故 $F$
在 $H$ 中稠密.

最后, 在每个子空间 $(E_n)_{n\ge0}$ 中选取一个 Hilbert 基 ($E_0$ 的这种基的
存在性由定理  \ref{thm:5.11} 得出; 对其余的 $E_n$, $n\ge1$, 这是显然的,
因为它们有限维).这些基的并显然是 $H$ 的由 $T$ 的特征向量构成的 Hilbert 基.
\end{proof}

\begin{remark}
\label{rem:6.8}
设 $T$ 为紧自伴算子.由前面的分析, 我们可以把任何 $u\in H$ 写成
\begin{equation*}
u=\sum_{n=0}^\infty u_n\quad\text{其中}\ u_n\in E_n.
\end{equation*}
则 $Tu=\sum_{n=1}^\infty\lambda_nu_n$.给定整数 $k\ge1$, 令
\begin{equation*}
T_ku=\sum_{n=1}^k\lambda_nu_n.
\end{equation*}
显然 $T_k$ 是有限秩算子, 且
\begin{equation*}
\norm{T_k-T}\le\sup_{n\ge k+1}\abs{\lambda_n}\to0
\quad\text{当}\ k\to\infty.
\end{equation*}
回忆: 事实上在 Hilbert 空间中, 每个紧算子——不必自伴——都是有限秩算子序列的
极限 (见注记 \ref{rem:6.1}).
\end{remark}

\section*{第 6 章评注}
\addcontentsline{toc}{section}{第 6 章评注}
\markright{第 6 章评注}

\cmtitem{$\star$1. Fredholm 算子.}
定理  \ref{thm:6.6} 是通向 Fredholm 算子理论的第一步.给定两个 Banach 空间
$E$ 与 $F$, 若 $A\in\LL(E,F)$ 满足
\begin{enumerate}[label=(\roman*), leftmargin=2.6em, itemsep=0.1em]
	\item $N(A)$ 是有限维的,
	\item $R(A)$ 闭且余维有限,\footnote{设 $A\in\LL(E,F)$ 满足 $N(A)$ 有限维
	      且 $R(A)$ 余维有限 (即存在有限维空间 $G\subset F$, 使得
	      $R(A)+G=F$).则可推出 $R(A)$ 闭 (见习题  \ref{exer:2.27}).}
\end{enumerate}
则称 $A$ 是\emph{Fredholm 算子} (或 \emph{Noether 算子})——记作
$A\in\Phi(E,F)$.$A$ 的\emph{指标}定义为
\begin{equation*}
\ind A=\dim N(A)-\codim R(A).
\end{equation*}
例如, $A=I-T$ (其中 $T\in\KK(E)$) 是指标为零的 Fredholm 算子; 这由定理
 \ref{thm:6.6} 得出.

Fredholm 算子的主要性质如下:

\begin{enumerate}[label=(\alph*), leftmargin=2.2em, itemsep=0.3em]
	\item Fredholm 算子类 $\Phi(E,F)$ 是 $\LL(E,F)$ 的开子集, 且映射
	      $A\mapsto\ind A$ 连续; 因此它在 $\Phi(E,F)$ 的每个连通分支上为常数.
	\item 每个算子 $A\in\Phi(E,F)$ 模有限秩算子可逆, 即存在算子
	      $B\in\LL(F,E)$, 使得
	      \begin{equation*}
	      (A\circ B-I_F)\quad\text{与}\quad(B\circ A-I_E)
	      \end{equation*}
	      都是有限秩算子.

	      反之, 设 $A\in\LL(E,F)$, 并设存在 $B\in\LL(F,E)$, 使得
	      \begin{equation*}
	      A\circ B-I_F\in\KK(F)\quad\text{且}\quad
	      B\circ A-I_E\in\KK(E).
	      \end{equation*}
	      则 $A\in\Phi(E,F)$.
	\item 若 $A\in\Phi(E,F)$ 且 $T\in\KK(E,F)$, 则 $A+T\in\Phi(E,F)$ 且
	      $\ind(A+T)=\ind A$.
	\item 若 $A\in\Phi(E,F)$ 且 $B\in\Phi(F,G)$, 则 $B\circ A\in\Phi(E,G)$ 且
	      $\ind(B\circ A)=\ind(A)+\ind(B)$.
\end{enumerate}

关于这个问题, 见例如 T.~Kato \bref{Kato}{1}, M.~Schechter \bref{Schechter}{1}, S.~Lang \bref{Lang}{1}, A.~E.~Taylor--D.~C.~Lay \bref{TaylorLay}{1}, P.~Lax \bref{Lax}{1}, L.~Hörmander \bref{Hormander}{2} (第 3 卷), 以及问题
38.

\cmtitem{$\star$2. Hilbert--Schmidt 算子.}
设 $H$ 为可分 Hilbert 空间.若有 $H$ 中的 Hilbert 基 $(e_n)$, 使得
$\norm{T}_{HS}^2=\sum\abs{Te_n}^2<\infty$, 则称有界算子 $T\in\LL(H)$ 为
\emph{Hilbert--Schmidt 算子}.可以证明这个定义与基的选取无关, 且
$\norm{\cdot}_{HS}$ 是一个范数.每个 Hilbert--Schmidt 算子都是紧的.
Hilbert--Schmidt 算子起着重要作用, 尤其是由于下述结果:

\begin{theorem}
\label{thm:6.12}
设 $H=L^2(\Omega)$, $K(x,y)\in L^2(\Omega\times\Omega)$.则算子
\begin{equation*}
u\mapsto(Ku)(x)=\int_\Omega K(x,y)u(y)\,dy
\end{equation*}
是 Hilbert--Schmidt 算子.

反之, $L^2(\Omega)$ 上的每个 Hilbert--Schmidt 算子都具有上述形式, 其中
$K(x,y)\in L^2(\Omega\times\Omega)$ 由算子唯一确定.
\end{theorem}

关于这个问题, 见例如 A.~Balakrishnan \bref{Balakrishnan}{1}, N.~Dunford--J.~T.~Schwartz \bref{DunfordSchwartz}{1} 第 2 卷, 以及问题 40.

\cmtitem{3. 特征值的重数.}
设 $T\in\KK(E)$, $\lambda\in\sigma(T)\setminus\{0\}$.可以证明序列
$N((T-\lambda I)^k)$, $k=1,2,\dots$, 严格递增到某个有限的 $p$ 之后保持
不变 (见例如 A.~E.~Taylor--D.~C.~Lay \bref{TaylorLay}{1}, E.~Kreyszig \bref{Kreyszig}{1}, 以及问题 36).这个整数 $p$ 称为
$(T-\lambda I)$ 的\emph{升数}.$N(T-\lambda I)$ 的维数被某些作者称为
$\lambda$ 的\emph{几何重数}, $N((T-\lambda I)^p)$ 的维数称为 $\lambda$ 的
\emph{代数重数}; 若 $E$ 是 Hilbert 空间且 $T$ 自伴, 二者相等 (见问题 36).

\cmtitem{4. 谱分析.}
设 $H$ 为 Hilbert 空间, $T\in\LL(H)$ 为自伴算子 (可以不是紧的).
有一种称为 $T$ 的\emph{谱族}的构造, 它推广了第~\ref{sec:6.4}~节的谱分解.它特别允许
定义\emph{泛函演算}, 即对任何连续函数 $f$ 赋予量 $f(T)$ 意义.它还可以推广到
无界算子与非自伴算子, 只要假设 $T$ 是\emph{正规}的, 即 $TT^\star=T^\star T$.
谱分析是一个广阔的领域, 尤其是在 $\CC$ 上的 Banach 空间中 (见第~\ref{sec:11.4}~节),
有许多应用与分支.初等的介绍见例如 W.~Rudin \bref{Rudin}{1}, E.~Kreyszig \bref{Kreyszig}{1}, A.~Friedman \bref{Friedman}{3}, K.~Yosida \bref{Yosida}{1}.更完整的叙述见例如 M.~Reed--B.~Simon
\bref{ReedSimon}{1}, T.~Kato \bref{Kato}{1}, R.~Dautray--J.-L.~Lions \bref{DautrayLions}{1} 第 VIII 与 IX 章, N.~Dunford--J.~T.~Schwartz \bref{DunfordSchwartz}{1} 第 2 卷, N.~Akhiezer--I.~Glazman \bref{AkhiezerGlazman}{1}, A.~E.~Taylor--D.~C.~Lay \bref{TaylorLay}{1}, J.~Weidmann \bref{Weidmann}{1}, J.~B.~Conway \bref{Conway}{1}, P.~Lax \bref{Lax}{1}, 以及
M.~Schechter \bref{Schechter}{2}.

\cmtnum{5. 极小极大原理.} Courant--Fischer 的极小极大公式提供了计算特征值的
一种非常有用的方法; 见例如 R.~Courant--D.~Hilbert \bref{CourantHilbert}{1}, P.~Lax \bref{Lax}{1}, 以及问题 37.H.~Weinberger \bref{Weinberger}{2} 的
专著包含这个专题的大量发展.

\cmtitem{6. Krein--Rutman 定理.}
下述结果在研究二阶椭圆算子的谱性质时有重要应用 (见第~
\ref{chap:9}~章).

\begin{startheorem}[Krein--Rutman]
\label{thm:6.13}
设 $E$ 为 Banach 空间, $P$ 为顶点在 $0$ 的凸锥, 即
$\lambda x+\mu y\in P\ \forall\lambda\ge0$, $\forall\mu\ge0$,
$\forall x\in P$, $\forall y\in P$.假设 $P$ 闭, $\Int P\ne\emptyset$,
且 $P\ne E$.设 $T\in\KK(E)$ 满足 $T(P\setminus\{0\})\subset\Int P$.则存在
$x_0\in\Int P$ 与 $\lambda_0>0$, 使得 $Tx_0=\lambda_0x_0$; 此外, $\lambda_0$
是唯一对应于 $P$ 中特征向量的特征值, 即若 $Tx=\lambda x$, $x\in P$,
$x\ne0$, 则 $\lambda=\lambda_0$ 且 $x=mx_0$ (其中 $m>0$).最后,
\begin{equation*}
\lambda_0=\max\{\abs{\lambda};\ \lambda\in\sigma(T)\},
\end{equation*}
且 $\lambda$ 的重数 (几何重数与代数重数) 等于一.
\end{startheorem}

问题 41 中给出的证明属于 P.~Rabinowitz \bref{Rabinowitz}{2}.上述
Krein--Rutman 定理的变体可在例如 H.~Schaefer \bref{Schaefer}{1}, R.~Nussbaum \bref{Nussbaum}{1}, F.~F.~Bonsall \bref{Bonsall}{1} 与
J.~F.~Toland \bref{Toland}{4} 中找到.

\section*{第 6 章习题}
\addcontentsline{toc}{section}{第 6 章习题}
\markright{第 6 章习题}

\begin{exercise}
\label{exer:6.1}
设 $E=\ell^p$, $1\le p\le\infty$ (见第~\ref{sec:11.3}~节).设 $(\lambda_n)$ 为 $\R$
中的有界序列, 考虑由
\begin{equation*}
Tx=(\lambda_1x_1,\lambda_2x_2,\dots,\lambda_nx_n,\dots)
\end{equation*}
定义的算子 $T\in\LL(E)$, 其中 $x=(x_1,x_2,\dots,x_n,\dots)$.证明 $T$ 是由
$E$ 到 $E$ 的紧算子当且仅当 $\lambda_n\to0$.
\end{exercise}

\begin{exercise}
\label{exer:6.2}
设 $E$ 与 $F$ 为两个 Banach 空间, $T\in\LL(E,F)$.
\begin{enumerate}[label=\arabic*., leftmargin=2.2em, itemsep=0.3em]
	\item 假设 $E$ 自反.证明 $T(B_E)$ 是 (强) 闭的.
	\item 假设 $E$ 自反且 $T\in\KK(E,F)$.证明 $T(B_E)$ 紧.
	\item 设 $E=F=C([0,1])$, $Tu(t)=\int_0^tu(s)\,ds$.验证 $T\in\KK(E)$.
	      证明 $T(B_E)$ 不闭.
\end{enumerate}
\end{exercise}

\begin{exercise}
\label{exer:6.3}
设 $E$ 与 $F$ 为两个 Banach 空间, $T\in\KK(E,F)$.假设 $\dim E=\infty$.
证明存在 $E$ 中的序列 $(u_n)$, 使得 $\norm{u_n}_E=1$ 且
$\norm{Tu_n}_F\to0$.
[提示: 用反证法.]
\end{exercise}

\begin{exercise}
\label{exer:6.4}
设 $1\le p<\infty$.验证 $\ell^p\subset c_0$ 且嵌入连续 (关于 $\ell^p$ 与
$c_0$ 的定义, 见第~\ref{sec:11.3}~节).这个嵌入是紧的吗?
[提示: 使用 $\ell^p$ 的典范基 $(e_n)$.]
\end{exercise}

\begin{exercise}
\label{exer:6.5}
设 $(\lambda_n)$ 为正数序列, 满足 $\lim_{n\to\infty}\lambda_n=+\infty$.
设 $V$ 为满足
\begin{equation*}
\sum_{n=1}^\infty\lambda_n\abs{u_n}^2<\infty
\end{equation*}
的序列 $(u_n)_{n\ge1}$ 构成的空间.$V$ 赋以标量积
\begin{equation*}
((u,v))=\sum_{n=1}^\infty\lambda_nu_nv_n.
\end{equation*}
证明 $V$ 是 Hilbert 空间, 且 $V\subset\ell^2$ 的嵌入是紧的.
\end{exercise}

\begin{exercise}
\label{exer:6.6}
设 $1\le q\le p\le\infty$.证明由 $L^p(0,1)$ 到 $L^q(0,1)$ 的典范嵌入
连续但不紧.
[提示: 使用 Rademacher 函数; 见习题  \ref{exer:4.18}.]
\end{exercise}

\begin{exercise}
\label{exer:6.7}
设 $E$ 与 $F$ 为两个 Banach 空间, $T\in\LL(E,F)$.考虑下述性质:
\begin{equation*}
\text{(P)}\qquad
\begin{cases}
\text{对 }E\text{ 中每个弱收敛序列 }(u_n),\\
u_n\weakto u\ \text{蕴含}\ Tu_n\to Tu\ \text{在 }F\text{ 中强收敛}.
\end{cases}
\end{equation*}
\begin{equation*}
\text{(Q)}\qquad
\begin{cases}
T\text{ 是由 }E\text{ (赋以弱拓扑 }\sigma(E,E^\star)\text{) 到 }F\\
\text{(赋以强拓扑) 的连续算子}.
\end{cases}
\end{equation*}
\begin{enumerate}[label=\arabic*., leftmargin=2.2em, itemsep=0.3em]
	\item 证明
	      \begin{equation*}
	      \text{(Q)}\iff T\ \text{是有限秩算子}.
	      \end{equation*}
	\item 证明 $T\in\KK(E,F)\Rightarrow\text{(P)}$.
	\item 假设 $E=\ell^1$ 或 $F=\ell^1$.证明每个算子 $T\in\LL(E,F)$ 都满足
	      (P).[提示: 使用问题 8 的一个结果.]
\end{enumerate}

下文假设 $E$ 自反.

\begin{enumerate}[label=\arabic*., leftmargin=2.2em, start=4, itemsep=0.3em]
	\item 证明 $T\in\KK(E,F)\iff\text{(P)}$.
	\item 由此推出每个算子 $T\in\LL(E,\ell^1)$ 都是紧的.
	\item 证明每个算子 $T\in\LL(c_0,E)$ 都是紧的.
	      [提示: 考虑伴随算子 $T^\star$.]
\end{enumerate}
\end{exercise}

\begin{exercise}
\label{exer:6.8}
设 $E$ 与 $F$ 为两个 Banach 空间, $T\in\KK(E,F)$.假设 $R(T)$ 闭.
\begin{enumerate}[label=\arabic*., leftmargin=2.2em, itemsep=0.3em]
	\item 证明 $T$ 是有限秩算子.
	      [提示: 使用开映射定理, 即定理  \ref{thm:2.6}.]
	\item 此外假设 $\dim N(T)<\infty$.证明 $\dim E<\infty$.
\end{enumerate}
\end{exercise}

\begin{exercise}
\label{exer:6.9}
设 $E$ 与 $F$ 为两个 Banach 空间, $T\in\LL(E,F)$.
\begin{enumerate}[label=\arabic*., leftmargin=2.2em, itemsep=0.3em]
	\item 证明下述三个性质等价:\footnote{投影算子是指满足 $P^2=P$ 的
	      算子 $P$.}
	      \begin{equation*}
	      \text{(A)}\qquad
	      \begin{cases}
	      \dim N(T)<\infty\ \text{且}\ R(T)\ \text{闭}.
	      \end{cases}
	      \end{equation*}
	      \begin{equation*}
	      \text{(B)}\qquad
	      \begin{cases}
	      \text{存在有限秩投影算子 }P\in\LL(E)\text{ 与常数 }C,\\
	      \text{使得}\ \norm{u}_E\le C(\norm{Tu}_F+\norm{Pu}_E)
	      \quad\forall u\in E.
	      \end{cases}
	      \end{equation*}
	      \begin{equation*}
	      \text{(C)}\qquad
	      \begin{cases}
	      \text{存在 Banach 空间 }G\text{, 算子 }Q\in\KK(E,G)\text{ 与常数 }C,\\
	      \text{使得}\ \norm{u}_E\le C(\norm{Tu}_F+\norm{Qu}_G)
	      \quad\forall u\in E.
	      \end{cases}
	      \end{equation*}
	      [提示: 当 $\dim N(T)<\infty$ 时, 考虑 $N(T)$ 的补空间;
	      见第~\ref{sec:2.4}~节.] 与习题  \ref{exer:2.12} 作比较.
	\item 假设 $T$ 满足 (A).证明对每个 $S\in\KK(E,F)$, $(T+S)$ 也满足 (A).
	\item 证明满足 (A) 的全部算子 $T\in\LL(E,F)$ 之集在 $\LL(E,F)$ 中是开集.
	\item 设 $F_0$ 为 $F$ 的闭线性子空间, $S\in\KK(F_0,F)$.证明
	      $(I+S)(F_0)$ 是 $F$ 的闭子空间.
\end{enumerate}
\end{exercise}

\begin{exercise}
\label{exer:6.10}
设 $Q(t)=\sum_{k=1}^pa_kt^k$ 为多项式, 且 $Q(1)\ne0$.设 $E$ 为 Banach
空间, $T\in\LL(E)$.假设 $Q(T)\in\KK(E)$.
\begin{enumerate}[label=\arabic*., leftmargin=2.2em, itemsep=0.3em]
	\item 证明 $\dim N(I-T)<\infty$, 且 $R(I-T)$ 闭.更一般地, 证明对每个闭
	      子空间 $E_0\subset E$, $(I-T)(E_0)$ 闭.
	      [提示: 写出 $Q(1)-Q(t)=\tilde Q(t)(1-t)$, 其中 $\tilde Q$ 为某个
	      多项式, 并应用习题  \ref{exer:6.9}.]
	\item 证明 $N(I-T)=\{0\}\iff R(I-T)=E$.
	\item 证明 $\dim N(I-T)=\dim N(I-T^\star)$.
	      [第 2, 3 问的提示: 使用与定理  \ref{thm:6.6} 的证明相同的方法.]
\end{enumerate}
\end{exercise}

\begin{exercise}
\label{exer:6.11}
设 $K$ 为紧度量空间, $E=C(K;\R)$ 赋以通常的范数
$\norm{u}=\max_{x\in K}\abs{u(x)}$.设 $F\subset E$ 为闭子空间.假设每个函数
$u\in F$ 都是 Hölder 连续的, 即
\begin{equation*}
\forall u\in F\ \ \exists\alpha\in(0,1]\ \text{与}\ \exists L
\ \text{使得}\ \abs{u(x)-u(y)}\le Ld(x,y)^\alpha
\quad\forall x,y\in K.
\end{equation*}
本题的目的是证明 $F$ 是有限维的.
\begin{enumerate}[label=\arabic*., leftmargin=2.2em, itemsep=0.3em]
	\item 证明存在常数 $\gamma\in(0,1]$ 与 $C\ge0$ (二者都与 $u$ 无关),
	      使得
	      \begin{equation*}
	      \abs{u(x)-u(y)}\le C\norm{u}d(x,y)^\gamma
	      \quad\forall u\in F,\ \forall x,y\in K.
	      \end{equation*}
	      [提示: 应用 Baire 纲定理 (定理  \ref{thm:2.1}), 取
	      $F_n=\{u\in F;\ \abs{u(x)-u(y)}\le nd(x,y)^{1/n}
	      \ \forall x,y\in K\}$.]
	\item 证明 $B_F$ 紧, 并作出结论.
\end{enumerate}
\end{exercise}

\begin{exercise}
\label{exer:6.12}
\emph{J.-L.~Lions 的一个引理.}
设 $X$, $Y$, $Z$ 为三个 Banach 空间, 范数分别为 $\norm{\cdot}_X$, $\norm{\cdot}_Y$, $\norm{\cdot}_Z$.假设 $X\subset Y$ 的嵌入是紧的,
$Y\subset Z$ 的嵌入是连续的.证明
\begin{equation*}
\forall\varepsilon>0\ \ \exists C_\varepsilon\ge0\ \text{使得}\
\norm{u}_Y\le\varepsilon\norm{u}_X+C_\varepsilon\norm{u}_Z
\quad\forall u\in X.
\end{equation*}
[提示: 用反证法.]

\emph{应用.} 证明
\begin{equation*}
\forall\varepsilon>0\ \ \exists C_\varepsilon\ge0\ \text{使得}\
\max_{[0,1]}\abs{u}\le\varepsilon\max_{[0,1]}\abs{u'}+C_\varepsilon\norm{u}_{L^1}
\quad\forall u\in C^1([0,1]).
\end{equation*}
\end{exercise}

\begin{exercise}
\label{exer:6.13}
设 $E$ 与 $F$ 为两个 Banach 空间, 范数分别为 $\norm{\cdot}_E$ 与
$\norm{\cdot}_F$.假设 $E$ 自反.设 $T\in\KK(E,F)$.考虑 $E$ 上的另一个
范数 $\abs{\cdot}$, 它比 $\norm{\cdot}_E$ 弱, 即
$\abs{u}\le C\norm{u}_E\ \forall u\in E$.证明
\begin{equation*}
\forall\varepsilon>0\ \ \exists C_\varepsilon\ge0\ \text{使得}\
\norm{Tu}_F\le\varepsilon\norm{u}_E+C_\varepsilon\abs{u}
\quad\forall u\in E.
\end{equation*}
证明当 $E$ 不自反时结论可能不成立.
[提示: 取 $E=C([0,1])$, $F=\R$, $\norm{u}=\norm{u}_{L^\infty}$,
$\abs{u}=\norm{u}_{L^1}$.]
\end{exercise}

\begin{exercise}
\label{exer:6.14}
设 $E$ 为 Banach 空间, $T\in\LL(E)$ 且 $\norm{T}<1$.
\begin{enumerate}[label=\arabic*., leftmargin=2.2em, itemsep=0.3em]
	\item 证明 $(I-T)$ 是双射, 且
	      \begin{equation*}
	      \norm{(I-T)^{-1}}\le\frac{1}{1-\norm{T}}.
	      \end{equation*}
	\item 令 $S_n=I+T+\cdots+T^{n-1}$.证明
	      \begin{equation*}
	      \norm{S_n-(I-T)^{-1}}\le\frac{\norm{T}^n}{1-\norm{T}}.
	      \end{equation*}
\end{enumerate}
\end{exercise}

\begin{exercise}
\label{exer:6.15}
设 $E$ 为 Banach 空间, $T\in\LL(E)$.
\begin{enumerate}[label=\arabic*., leftmargin=2.2em, itemsep=0.3em]
	\item 设 $\lambda\in\R$ 满足 $\abs{\lambda}>\norm{T}$.证明
	      \begin{equation*}
	      \norm{I+\lambda(T-\lambda I)^{-1}}\le\frac{\norm{T}}{\abs{\lambda}-\norm{T}}.
	      \end{equation*}
	\item 设 $\lambda\in\rho(T)$.验证
	      \begin{equation*}
	      (T-\lambda I)^{-1}T=T(T-\lambda I)^{-1},
	      \end{equation*}
	      并证明
	      \begin{equation*}
	      \dist(\lambda,\sigma(T))\ge\frac{1}{\norm{(T-\lambda I)^{-1}}}.
	      \end{equation*}
	\item 假设 $0\in\rho(T)$.证明
	      \begin{equation*}
	      \sigma(T^{-1})=\frac{1}{\sigma(T)}.
	      \end{equation*}
\end{enumerate}

下文假设 $1\in\rho(T)$; 令
\begin{equation*}
U=(T+I)(T-I)^{-1}=(T-I)^{-1}(T+I).
\end{equation*}

\begin{enumerate}[label=\arabic*., leftmargin=2.2em, start=4, itemsep=0.3em]
	\item 验证 $1\in\rho(U)$, 并用 $T$ 给出 $(U-I)^{-1}$ 的简单表达式.
	\item 证明 $T=(U+I)(U-I)^{-1}$.
	\item 考虑函数 $f(t)=\frac{t+1}{t-1}$, $t\in\R$.证明
	      \begin{equation*}
	      \sigma(U)=f(\sigma(T)).
	      \end{equation*}
\end{enumerate}
\end{exercise}

\begin{exercise}
\label{exer:6.16}
设 $E$ 为 Banach 空间, $T\in\LL(E)$.
\begin{enumerate}[label=\arabic*., leftmargin=2.2em, itemsep=0.3em]
	\item 假设 $T^2=I$.证明 $\sigma(T)\subset\{-1,+1\}$, 并对
	      $\lambda\ne\pm1$ 确定 $(T-\lambda I)^{-1}$.
	\item 更一般地, 假设存在整数 $n\ge2$ 使得 $T^n=I$.证明
	      $\sigma(T)\subset\{-1,+1\}$, 并对 $\lambda\ne\pm1$ 确定
	      $(T-\lambda I)^{-1}$.
	\item 假设存在整数 $n\ge2$ 使得 $T^n=0$.证明 $\sigma(T)=\{0\}$, 并对
	      $\lambda\ne0$ 确定 $(T-\lambda I)^{-1}$.
	\item 假设存在整数 $n\ge2$ 使得 $\norm{T^n}<1$.证明 $I-T$ 是双射, 并用
	      $(I-T^n)^{-1}$ 与 $T$ 的迭代给出 $(I-T)^{-1}$ 的表达式.
\end{enumerate}
\end{exercise}

\begin{exercise}
\label{exer:6.17}
设 $E=\ell^p$, $1\le p\le\infty$, $(\lambda_n)$ 为 $\R$ 中的有界序列.
考虑由
\begin{equation*}
Mx=(\lambda_1x_1,\lambda_2x_2,\dots,\lambda_nx_n,\dots)
\end{equation*}
定义的乘法算子 $M\in\LL(E)$, 其中 $x=(x_1,x_2,\dots,x_n,\dots)$.确定
$EV(M)$ 与 $\sigma(M)$.
\end{exercise}

\begin{exercise}
\label{exer:6.18}
\emph{移位算子的谱性质.}
$E=\ell^2$ 中的元素记作 $x=(x_1,x_2,\dots,x_n,\dots)$.考虑算子
\begin{equation*}
S_rx=(0,x_1,x_2,\dots,x_{n-1},\dots),\qquad
S_\ell x=(x_2,x_3,x_4,\dots,x_{n+1},\dots),
\end{equation*}
它们分别称为\emph{右移位}与\emph{左移位}.
\begin{enumerate}[label=\arabic*., leftmargin=2.2em, itemsep=0.3em]
	\item 确定 $\norm{S_r}$ 与 $\norm{S_\ell}$.$S_r$ 或 $S_\ell$ 属于
	      $\KK(E)$ 吗?
	\item 证明 $EV(S_r)=\emptyset$.
	\item 证明 $\sigma(S_r)=[-1,+1]$.
	\item 证明 $EV(S_\ell)=(-1,+1)$.确定对应的特征空间.
	\item 证明 $\sigma(S_\ell)=[-1,+1]$.
	\item 确定 $S_r^\star$ 与 $S_\ell^\star$.
	\item 证明对每个 $\lambda\in(-1,+1)$, 空间 $R(S_r-\lambda I)$ 与
	      $R(S_\ell-\lambda I)$ 都闭.给出这些空间的显式表示.
	      [提示: 应用定理  \ref{thm:2.19} 与  \ref{thm:2.20}.]
	\item 证明空间 $R(S_r\pm I)$ 与 $R(S_\ell\pm I)$ 稠密但不闭.
\end{enumerate}

考虑由
\begin{equation*}
Mx=(\alpha_1x_1,\alpha_2x_2,\dots,\alpha_nx_n,\dots)
\end{equation*}
定义的乘法算子 $M$, 其中 $(\alpha_n)$ 是 $\R$ 中的有界序列.

\begin{enumerate}[label=\arabic*., leftmargin=2.2em, start=9, itemsep=0.3em]
	\item 确定 $EV(S_r\circ M)$.
	\item 假设当 $n\to\infty$ 时 $\alpha_n\to\alpha$.证明
	      \begin{equation*}
	      \sigma(S_r\circ M)=[-\abs{\alpha},+\abs{\alpha}].
	      \end{equation*}
	      [提示: 应用定理  \ref{thm:6.6}.]
	\item 假设对每个整数 $n$, $\alpha_{2n}=a$, $\alpha_{2n+1}=b$, 其中
	      $a\ne b$.确定 $\sigma(S_r\circ M)$.
	      [提示: 计算 $(S_r\circ M)^2$, 并应用习题  \ref{exer:6.16} 第 4 问.]
\end{enumerate}
\end{exercise}

\begin{exercise}
\label{exer:6.19}
设 $E$ 为 Banach 空间, $T\in\LL(E)$.
\begin{enumerate}[label=\arabic*., leftmargin=2.2em, itemsep=0.3em]
	\item 证明 $\sigma(T^\star)=\sigma(T)$.
	\item 举例说明 $EV(T)$ 与 $EV(T^\star)$ 之间一般没有包含关系.
	      [提示: 考虑右移位与左移位.]
\end{enumerate}
\end{exercise}

\begin{exercise}
\label{exer:6.20}
设 $E=L^p(0,1)$, $1\le p<\infty$.给定 $u\in E$, 令
\begin{equation*}
Tu(x)=\int_0^xu(t)\,dt.
\end{equation*}
\begin{enumerate}[label=\arabic*., leftmargin=2.2em, itemsep=0.3em]
	\item 证明 $T\in\KK(E)$.
	\item 确定 $EV(T)$ 与 $\sigma(T)$.
	\item 当 $\lambda\in\rho(T)$ 时, 给出 $(T-\lambda I)^{-1}$ 的显式公式.
	\item 确定 $T^\star$.
\end{enumerate}
\end{exercise}

\begin{exercise}
\label{exer:6.21}
设 $V$ 与 $H$ 为两个 Banach 空间, 范数分别为 $\norm{\cdot}$ 与
$\abs{\cdot}$, 满足
\begin{equation*}
V\subset H\ \text{的嵌入是紧的}.
\end{equation*}
设 $p(u)$ 为 $V$ 上的半范数, 使得 $p(u)+\abs{u}$ 是 $V$ 上与
$\norm{\cdot}$ 等价的范数.令
\begin{equation*}
N=\{u\in V;\ p(u)=0\},
\end{equation*}
并对 $u\in V$ 令 $\dist(u,N)=\inf_{v\in N}\norm{u-v}$.
\begin{enumerate}[label=\arabic*., leftmargin=2.2em, itemsep=0.3em]
	\item 证明 $N$ 是有限维空间.
	      [提示: 考虑 $N$ 中赋以范数 $\abs{\cdot}$ 的单位球.]
	\item 证明存在常数 $K_1>0$, 使得
	      \begin{equation*}
	      p(u)\le K_1\dist(u,N)\quad\forall u\in V.
	      \end{equation*}
	\item 证明存在常数 $K_2>0$, 使得
	      \begin{equation*}
	      K_2\dist(u,N)\le p(u)\quad\forall u\in V.
	      \end{equation*}
	      [提示: 用反证法.假设存在 $V$ 中的序列 $(u_n)$, 使得
	      $\dist(u_n,N)=1\ \forall n$ 且 $p(u_n)\to0$.]
\end{enumerate}
\end{exercise}

\begin{exercise}
\label{exer:6.22}
设 $E$ 为 Banach 空间, $T\in\LL(E)$.给定多项式
$Q(t)=\sum_{k=0}^pa_kt^k$ ($a_k\in\R$), 令
$Q(T)=\sum_{k=0}^pa_kT^k$.
\begin{enumerate}[label=\arabic*., leftmargin=2.2em, itemsep=0.3em]
	\item 证明 $Q(EV(T))\subset EV(Q(T))$.
	\item 证明 $Q(\sigma(T))\subset\sigma(Q(T))$.
	\item 在 $E=\R^2$ 中构造一个例子, 使上述包含是严格的.
\end{enumerate}

下文假设 $E$ 是 Hilbert 空间 (与其对偶空间 $H^\star$ 等同), 且
$T^\star=T$.

\begin{enumerate}[label=\arabic*., leftmargin=2.2em, start=4, itemsep=0.3em]
	\item 这里假设多项式 $Q$ 没有实根, 即
	      $Q(t)\ne0\ \forall t\in\R$.证明 $Q(T)$ 是双射.
	      [提示: 先考虑 $Q$ 是二次多项式的情形, 更确切地说,
	      $Q(t)=t^2+1$.]
	\item 由此推出对每个多项式 $Q$, 都有
	      \begin{enumerate}[label=(\roman*), leftmargin=2.4em, itemsep=0.1em]
	      	\item $Q(EV(T))=EV(Q(T))$,
	      	\item $Q(\sigma(T))=\sigma(Q(T))$.
	      \end{enumerate}
	      [提示: 写出
	      $Q(t)-\lambda=(t-t_1)(t-t_2)\cdots(t-t_q)\tilde Q(t)$, 其中
	      $t_1,t_2,\dots,t_q$ 是 $Q(t)-\lambda$ 的实根, $\tilde Q$ 无实根.]
\end{enumerate}
\end{exercise}

\begin{exercise}
\label{exer:6.23}
\emph{谱半径.}
设 $E$ 为 Banach 空间, $T\in\LL(E)$.令
\begin{equation*}
a_n=\log\norm{T^n},\quad n\ge1.
\end{equation*}
\begin{enumerate}[label=\arabic*., leftmargin=2.2em, itemsep=0.3em]
	\item 验证
	      \begin{equation*}
	      a_{i+j}\le a_i+a_j\quad\forall i,j\ge1.
	      \end{equation*}
	\item 由此推出
	      \begin{equation*}
	      \lim_{n\to+\infty}(a_n/n)\ \text{存在, 且等于}\ \inf_{m\ge1}(a_m/m).
	      \end{equation*}
	      [提示: 固定整数 $m\ge1$.对任意整数 $n\ge1$, 写出
	      $n=mq+r$, 其中 $q=[\frac nm]$ 是不超过 $n/m$ 的最大整数,
	      $0\le r<m$.注意 $a_n\le\frac nm a_m+a_r$.]
	\item 作出结论: $r(T)=\lim_{n\to\infty}\norm{T^n}^{1/n}$ 存在, 且
	      $r(T)\le\norm{T}$.在 $E=\R^2$ 中构造一个例子, 使得 $r(T)=0$
	      而 $\norm{T}=1$.
\end{enumerate}

数 $r(T)$ 称为 $T$ 的\emph{谱半径}.

\begin{enumerate}[label=\arabic*., leftmargin=2.2em, start=4, itemsep=0.3em]
	\item 证明 $\sigma(T)\subset[-r(T),+r(T)]$.由此推出: 若
	      $\sigma(T)\ne\emptyset$, 则
	      $\max\{\abs{\lambda};\ \lambda\in\sigma(T)\}\le r(T)$.
	      [提示: 注意若 $\lambda\in\sigma(T)$, 则
	      $\lambda^n\in\sigma(T^n)$; 见习题  \ref{exer:6.22}.]
	\item 在 $E=\R^3$ 中构造一个例子, 使得 $\sigma(T)=\{0\}$, 而
	      $r(T)=1$.
\end{enumerate}

下文取 $E=L^p(0,1)$, $1\le p\le\infty$.考虑由
\begin{equation*}
Tu(t)=\int_0^tu(s)\,ds
\end{equation*}
定义的算子 $T\in\LL(E)$.

\begin{enumerate}[label=\arabic*., leftmargin=2.2em, start=6, itemsep=0.3em]
	\item 用归纳法证明: 当 $n\ge2$ 时,
	      \begin{equation*}
	      (T^nu)(t)=\frac{1}{(n-1)!}\int_0^t(t-\tau)^{n-1}u(\tau)\,d\tau.
	      \end{equation*}
	\item 由此推出 $\norm{T^n}\le\frac{1}{n!}$.
	      [提示: 使用卷积积的一个不等式.]
	\item 证明 $T$ 的谱半径为 $0$.
	      [提示: 使用 Stirling 公式.]
	\item 证明 $\sigma(T)=\{0\}$.与习题  \ref{exer:6.20} 作比较.
\end{enumerate}
\end{exercise}

\begin{exercise}
\label{exer:6.24}
假设 $T\in\LL(H)$ 自伴.
\begin{enumerate}[label=\arabic*., leftmargin=2.2em, itemsep=0.3em]
	\item 证明下述性质等价:
	      \begin{enumerate}[label=(\roman*), leftmargin=2.4em, itemsep=0.1em]
	      	\item $(Tu,u)\ge0\ \forall u\in H$,
	      	\item $\sigma(T)\subset[0,\infty)$.
	      \end{enumerate}
	      [提示: 应用命题  \ref{prop:6.9}.]
	\item 证明下述性质等价:
	      \begin{enumerate}[label=(\roman*), leftmargin=2.4em, start=3, itemsep=0.1em]
	      	\item $\norm{T}\le1$ 且 $(Tu,u)\ge0\ \forall u\in H$,
	      	\item $0\le(Tu,u)\le\abs{u}^2\ \forall u\in H$,
	      	\item $\sigma(T)\subset[0,1]$,
	      	\item $(Tu,u)\ge\abs{Tu}^2\ \forall u\in H$.
	      \end{enumerate}
	      [提示: 为证 (v)$\Rightarrow$(vi), 对 $(T+\varepsilon I)^{-1}$
	      ($\varepsilon>0$) 应用命题  \ref{prop:6.9}.]
	\item 证明下述性质等价:
	      \begin{enumerate}[label=(\roman*), leftmargin=2.4em, start=7, itemsep=0.1em]
	      	\item $(Tu,u)\le\abs{Tu}^2\ \forall u\in H$,
	      	\item $(0,1)\subset\rho(T)$.
	      \end{enumerate}
	      [提示: 引入 $U=2T-I$.]
\end{enumerate}
\end{exercise}

\begin{exercise}
\label{exer:6.25}
设 $E$ 为 Banach 空间, $K\in\KK(E)$.证明存在
$M\in\LL(E)$, $\tilde M\in\LL(E)$ 以及有限秩投影算子 $P$, $\tilde P$,
使得
\begin{enumerate}[label=(\roman*), leftmargin=2.6em, itemsep=0.15em]
	\item $M\circ(I+K)=I-P$,
	\item $(I+K)\circ\tilde M=I-\tilde P$.
\end{enumerate}
[提示: 设 $X$ 为 $N(I+K)$ 在 $E$ 中的补空间.则 $(I+K)|_X$ 是由 $X$ 到
$R(I+K)$ 的双射.记其逆为 $M$.设 $Q$ 为 $E$ 到 $R(I+K)$ 的投影算子,
并令 $\tilde M=M\circ Q$.证明 (i) 与 (ii) 成立.]
\end{exercise}

\begin{exercise}
\label{exer:6.26}
\emph{从 Brouwer 不动点定理到 Schauder 不动点定理.}
本题假设下述结果已知 (证明见例如 K.~Deimling \bref{Deimling}{1}, A.~Granas--J.~Dugundji \bref{GranasDugundji}{1}, 或
L.~Nirenberg \bref{Nirenberg}{2}).

\smallskip
\noindent\emph{定理 (Brouwer).} 设 $F$ 为有限维空间, $Q\subset F$ 为非空紧凸集.
设 $f:Q\to Q$ 为连续映射.则 $f$ 有不动点, 即存在 $p\in Q$ 使得
$f(p)=p$.

\smallskip
我们的目标是证明下述结果.

\smallskip
\noindent\emph{定理 (Schauder).} 设 $E$ 为 Banach 空间, $C$ 为 $E$ 中的非空
闭凸集.设 $F:C\to C$ 为连续映射, 且 $F(C)\subset K$, 其中 $K$ 是 $C$ 的紧
子集.则 $F$ 在 $K$ 中有不动点.

\begin{enumerate}[label=\arabic*., leftmargin=2.2em, itemsep=0.3em]
	\item 给定 $\varepsilon>0$, 考虑 $K$ 的有限覆盖, 即
	      $K\subset\bigcup_{i\in I}B(y_i,\varepsilon/2)$, 其中 $I$ 有限,
	      $y_i\in K\ \forall i\in I$.定义函数
	      $q(x)=\sum_{i\in I}q_i(x)$, 其中
	      \begin{equation*}
	      q_i(x)=\max\{\varepsilon-\norm{Fx-y_i},0\}.
	      \end{equation*}
	      验证 $q$ 在 $C$ 上连续, 且 $q(x)\ge\varepsilon/2\ \forall x\in C$.
	\item 令
	      \begin{equation*}
	      F_\varepsilon(x)=\frac{\sum_{i\in I}q_i(x)y_i}{q(x)},\quad x\in C.
	      \end{equation*}
	      证明 $F_\varepsilon:C\to C$ 连续, 且
	      \begin{equation*}
	      \norm{F_\varepsilon(x)-F(x)}\le\varepsilon,\quad\forall x\in C.
	      \end{equation*}
	\item 证明 $F_\varepsilon$ 有不动点 $x_\varepsilon\in C$.
	      [提示: 令 $Q=\operatorname{conv}\bigl(\bigcup_{i\in I}\{y_i\}\bigr)$.
	      验证 $F_\varepsilon|_Q$ 在 $Q$ 中有不动点 $x_\varepsilon\in Q$.]
	\item 证明对某个序列 $\varepsilon_n\to0$, $(x_{\varepsilon_n})$ 收敛于
	      某个极限 $x\in C$.证明 $F(x)=x$.
\end{enumerate}
\end{exercise}


% chapters/ch07.tex
% 第~
\ref{chap:7}~章 \index{HilleYosidadingli@Hille--Yosida 定理 (Hille--Yosida theorem)}
Hille--Yosida 定理
% 源: 原书 PDF 第 189--207 页 (印刷页 181--199)

\chapter{Hille--Yosida 定理}
\label{chap:7}

\section{\index{jidadandiaosuanzi@极大单调算子 (maximal monotone operator)}
极大单调算子的定义与初等性质}
\label{sec:7.1}

本章自始至终用 $H$ 表示 Hilbert 空间.

\begin{definition}
称无界线性算子 $A:D(A)\subset H\to H$ 是\emph{单调的},\footnote{有些作者称
$A$ 是\emph{增生算子} (accretive), 或称 $-A$ 是\emph{耗散算子}
(dissipative).} 若它满足
\begin{equation*}
(Av,v)\ge0\quad\forall v\in D(A).
\end{equation*}
若此外还有 $R(I+A)=H$, 即
\begin{equation*}
\forall f\in H\ \ \exists u\in D(A)\ \text{使得}\ u+Au=f,
\end{equation*}
则称 $A$ 是\emph{极大单调的}.
\end{definition}

\begin{proposition}
\label{prop:7.1}
设 $A$ 为极大单调算子.则
\begin{enumerate}[label=(\alph*), leftmargin=2.2em, itemsep=0.1em]
	\item $D(A)$ 在 $H$ 中稠密,
	\item $A$ 是闭算子,
	\item 对每个 $\lambda>0$, $(I+\lambda A)$ 是由 $D(A)$ 到 $H$ 的双射,
	      $(I+\lambda A)^{-1}$ 是有界算子, 且
	      $\norm{(I+\lambda A)^{-1}}_{\LL(H)}\le1$.
\end{enumerate}
\end{proposition}

\begin{proof}
\begin{enumerate}[label=(\alph*), leftmargin=2.2em, itemsep=0.3em]
	\item 设 $f\in H$ 满足 $(f,v)=0\ \forall v\in D(A)$.我们断言 $f=0$.
	      事实上, 存在某个 $v_0\in D(A)$ 使得 $v_0+Av_0=f$.我们有
	      \begin{equation*}
	      0=(f,v_0)=\abs{v_0}^2+(Av_0,v_0)\ge\abs{v_0}^2.
	      \end{equation*}
	      于是 $v_0=0$, 从而 $f=0$.
	\item 首先注意: 给定任何 $f\in H$, 存在唯一的 $u\in D(A)$ 使得
	      $u+Au=f$, 因为若 $u'$ 是另一个解, 则
	      \begin{equation*}
	      u-u'+A(u-u')=0.
	      \end{equation*}
	      与 $(u-u')$ 作标量积并利用单调性, 可见 $u-u'=0$.其次注意
	      $\abs{u}\le\abs{f}$, 因为
	      $\abs{u}^2+(Au,u)=(f,u)\ge\abs{u}^2$.因此映射 $f\mapsto u$
	      (记作 $(I+A)^{-1}$) 是由 $H$ 到自身的有界线性算子, 且
	      $\norm{(I+A)^{-1}}_{\LL(H)}\le1$.下面证明 $A$ 是闭算子.设
	      $(u_n)$ 为 $D(A)$ 中的序列, 使得 $u_n\to u$ 且 $Au_n\to f$.
	      我们要验证 $u\in D(A)$ 且 $Au=f$.但 $u_n+Au_n\to u+f$, 因此
	      \begin{equation*}
	      u_n=(I+A)^{-1}(u_n+Au_n)\to(I+A)^{-1}(u+f).
	      \end{equation*}
	      于是 $u=(I+A)^{-1}(u+f)$, 即 $u\in D(A)$ 且 $u+Au=u+f$.
	\item 我们将证明: 若对某个 $\lambda_0>0$ 有 $R(I+\lambda_0A)=H$, 则对每个
	      $\lambda>\lambda_0/2$ 都有 $R(I+\lambda A)=H$.首先注意——如 (b) 中
	      那样——对每个 $f\in H$ 存在唯一的 $u\in D(A)$ 使得
	      $u+\lambda_0Au=f$.此外, 映射 $f\mapsto u$ (记作
	      $(I+\lambda_0A)^{-1}$) 是有界线性算子, 且
	      $\norm{(I+\lambda_0A)^{-1}}_{\LL(H)}\le1$.我们试图求解方程
	      \begin{equation}
	      u+\lambda Au=f\quad\text{其中}\ \lambda>0.
	      \label{eq:7.1}
	      \end{equation}
	      方程 \eqref{eq:7.1} 可写成
	      \begin{equation*}
	      u+\lambda_0Au=\frac{\lambda_0}{\lambda}f
	      +\left(1-\frac{\lambda_0}{\lambda}\right)u,
	      \end{equation*}
	      或者等价地
	      \begin{equation}
	      u=(I+\lambda_0A)^{-1}
	      \left[\frac{\lambda_0}{\lambda}f
	      +\left(1-\frac{\lambda_0}{\lambda}\right)u\right].
	      \label{eq:7.2}
	      \end{equation}
	      若 $\left|1-\frac{\lambda_0}{\lambda}\right|<1$, 即
	      $\lambda>\lambda_0/2$, 则可应用压缩映射原理 (定理  \ref{thm:5.7}),
	      推出 \eqref{eq:7.2} 有解.

	      结论 (c) 容易由归纳法得出: 由于 $I+A$ 满射, 故对每个
	      $\lambda>1/2$, $I+\lambda A$ 满射, 从而对每个 $\lambda>1/4$
	      也满射, 依此类推.
\end{enumerate}
\end{proof}

\begin{remark}
\label{rem:7.1}
若 $A$ 极大单调, 则对每个 $\lambda>0$, $\lambda A$ 也极大单调.然而, 若
$A$ 与 $B$ 都是极大单调算子, 则定义在 $D(A)\cap D(B)$ 上的 $A+B$ 不必是
极大单调的.
\end{remark}

\begin{definition}
设 $A$ 为极大单调算子.对每个 $\lambda>0$, 令
\begin{equation*}
J_\lambda=(I+\lambda A)^{-1}
\quad\text{与}\quad
A_\lambda=\frac1\lambda(I-J_\lambda);
\end{equation*}
$J_\lambda$ 称为 $A$ 的\emph{预解式}, $A_\lambda$ 称为 $A$ 的
\emph{\index{Yosidabijin@Yosida 逼近 (Yosida approximation)}
Yosida 逼近} (或\emph{正则化}).记住
$\norm{J_\lambda}_{\LL(H)}\le1$.
\end{definition}

\begin{proposition}
\label{prop:7.2}
设 $A$ 为极大单调算子.则
\begin{align*}
A_\lambda v&=A(J_\lambda v)
&&\forall v\in H\ \text{及}\ \forall\lambda>0,
\tag*{(a$_1$)}\\
A_\lambda v&=J_\lambda(Av)
&&\forall v\in D(A)\ \text{及}\ \forall\lambda>0,
\tag*{(a$_2$)}\\
\abs{A_\lambda v}&\le\abs{Av}
&&\forall v\in D(A)\ \text{及}\ \forall\lambda>0,
\tag*{(b)}\\
\lim_{\lambda\to0}J_\lambda v&=v
&&\forall v\in H,
\tag*{(c)}\\
\lim_{\lambda\to0}A_\lambda v&=Av
&&\forall v\in D(A),
\tag*{(d)}\\
(A_\lambda v,v)&\ge0
&&\forall v\in H\ \text{及}\ \forall\lambda>0,
\tag*{(e)}\\
\abs{A_\lambda v}&\le(1/\lambda)\abs{v}
&&\forall v\in H\ \text{及}\ \forall\lambda>0.
\tag*{(f)}
\end{align*}
\end{proposition}

\begin{proof}
(a$_1$) 可写成 $v=(J_\lambda v)+\lambda A(J_\lambda v)$, 这正是 $J_\lambda v$
的定义.

(a$_2$) 由 (a$_1$) 我们有
\begin{equation*}
A_\lambda v+A(v-J_\lambda v)=Av,
\end{equation*}
即
\begin{equation*}
A_\lambda v+\lambda A(A_\lambda v)=Av,
\end{equation*}
这意味着 $A_\lambda v=(I+\lambda A)^{-1}Av$.

(b) 容易由 (a$_2$) 得出.

(c) 先设 $v\in D(A)$.则
\begin{equation*}
\abs{v-J_\lambda v}=\lambda\abs{A_\lambda v}\le\lambda\abs{Av}
\quad\text{(由 (b))},
\end{equation*}
从而 $\lim_{\lambda\to0}J_\lambda v=v$.

现在设 $v$ 是 $H$ 中的一般元素.给定任何 $\varepsilon>0$, 存在某个
$v_1\in D(A)$ 使得 $\abs{v-v_1}\le\varepsilon$ (因为由命题  \ref{prop:7.1},
$D(A)$ 在 $H$ 中稠密).我们有
\begin{equation*}
\abs{J_\lambda v-v}\le\abs{J_\lambda v-J_\lambda v_1}
+\abs{J_\lambda v_1-v_1}+\abs{v_1-v}
\le2\abs{v-v_1}+\abs{J_\lambda v_1-v_1}
\le2\varepsilon+\abs{J_\lambda v_1-v_1}.
\end{equation*}
于是
\begin{equation*}
\limsup_{\lambda\to0}\abs{J_\lambda v-v}\le2\varepsilon
\quad\forall\varepsilon>0,
\end{equation*}
从而
\begin{equation*}
\lim_{\lambda\to0}\abs{J_\lambda v-v}=0.
\end{equation*}

(d) 这是 (a$_2$) 与 (c) 的推论.

(e) 我们有
\begin{equation*}
(A_\lambda v,v)=(A_\lambda v,v-J_\lambda v)+(A_\lambda v,J_\lambda v)
=\lambda\abs{A_\lambda v}^2+(A(J_\lambda v),J_\lambda v),
\end{equation*}
从而
\begin{equation}
(A_\lambda v,v)\ge\lambda\abs{A_\lambda v}^2.
\label{eq:7.3}
\end{equation}

(f) 这是 \eqref{eq:7.3} 与 Cauchy--Schwarz 不等式的推论.
\end{proof}

\begin{remark}
\label{rem:7.2}
命题  \ref{prop:7.2} 说明 $(A_\lambda)_{\lambda>0}$ 是一族有界算子, 当
$\lambda\to0$ 时它们「逼近」无界算子 $A$.这种逼近方式将被非常频繁地使用.
当然, 一般说来 $\norm{A_\lambda}_{\LL(H)}$ 当 $\lambda\to0$ 时会「爆炸」.
\end{remark}

\section{发展方程 \texorpdfstring{$du/dt+Au=0$}{du/dt+Au=0} 在
\texorpdfstring{$[0,+\infty)$}{[0,+∞)} 上的解.
\texorpdfstring{$u(0)=u_0$}{u(0)=u0}. 存在性与唯一性}
\label{sec:7.2}

我们先从一个非常经典的结果开始:

\begin{keytheorem}[Cauchy, Lipschitz, Picard]
\label{thm:7.3}
设 $E$ 为 Banach 空间, $F:E\to E$ 为 Lipschitz 映射, 即存在常数 $L$,
使得
\begin{equation*}
\norm{Fu-Fv}\le L\norm{u-v}\quad\forall u,v\in E.
\end{equation*}
则给定任何 $u_0\in E$, 存在唯一的解 $u\in C^1([0,+\infty);E)$ 满足
\begin{equation}
\begin{cases}
\dfrac{du}{dt}(t)=Fu(t) & \text{在}\ [0,+\infty)\ \text{上},\\
u(0)=u_0.
\end{cases}
\label{eq:7.4}
\end{equation}
$u_0$ 称为\emph{初值}.
\end{keytheorem}

\begin{proof}
\emph{存在性.} 求解 \eqref{eq:7.4} 等价于求某个
$u\in C([0,+\infty);E)$ 满足积分方程
\begin{equation}
u(t)=u_0+\int_0^tF(u(s))\,ds.
\label{eq:7.5}
\end{equation}
给定 $k>0$ (稍后确定), 令
\begin{equation*}
X=\Bigl\{u\in C([0,+\infty);E);\ \sup_{t\ge0}e^{-kt}\norm{u(t)}<\infty\Bigr\}.
\end{equation*}
容易验证 $X$ 关于范数
\begin{equation*}
\norm{u}_X=\sup_{t\ge0}e^{-kt}\norm{u(t)}
\end{equation*}
是 Banach 空间.对每个 $u\in X$, 由
\begin{equation*}
(\mathcal{T}u)(t)=u_0+\int_0^tF(u(s))\,ds
\end{equation*}
定义的函数 $\mathcal{T}u$ 也属于 $X$.此外, 我们有
\begin{equation*}
\norm{\mathcal{T}u-\mathcal{T}v}_X\le\frac Lk\norm{u-v}_X
\quad\forall u,v\in X.
\end{equation*}
取 $k>L$, 可见 $\mathcal{T}$ 在 $X$ 中有 (唯一) 不动点 $u$, 它是
\eqref{eq:7.5} 的解.

\emph{唯一性.} 设 $u$ 与 $u'$ 是 \eqref{eq:7.4} 的两个解, 令
\begin{equation*}
\phi(t)=\norm{u(t)-u'(t)}.
\end{equation*}
由 \eqref{eq:7.5} 推出
\begin{equation*}
\phi(t)\le L\int_0^t\phi(s)\,ds\quad\forall t\ge0,
\end{equation*}
从而 $\phi\equiv0$.
\end{proof}

前一个定理在研究常微分方程时极其有用.然而在研究偏微分方程时它用处不大.
我们下一个结果在求解发展型偏微分方程时是非常强有力的工具; 见第~
\ref{chap:10}~章.

\begin{keytheorem}[Hille--Yosida]
\label{thm:7.4}
设 $A$ 为极大单调算子.则给定任何 $u_0\in D(A)$, 存在唯一的函数\footnote{空间
$D(A)$ 赋以图范数 $\abs{v}+\abs{Av}$, 或与之等价的 Hilbert 范数
$(\abs{v}^2+\abs{Av}^2)^{1/2}$.}
\begin{equation*}
u\in C^1([0,+\infty);H)\cap C([0,+\infty);D(A))
\end{equation*}
满足
\begin{equation}
\begin{cases}
\dfrac{du}{dt}+Au=0 & \text{在}\ [0,+\infty)\ \text{上},\\
u(0)=u_0.
\end{cases}
\label{eq:7.6}
\end{equation}
此外,
\begin{equation*}
\abs{u(t)}\le\abs{u_0}
\quad\text{且}\quad
\left|\frac{du}{dt}(t)\right|=\abs{Au(t)}\le\abs{Au_0}
\quad\forall t\ge0.
\end{equation*}
\end{keytheorem}

\begin{remark}
\label{rem:7.3}
定理  \ref{thm:7.4} 的主要意义在于, 它把「发展问题」的研究归结为「定态方程」
$u+Au=f$ 的研究 (假定我们已经知道 $A$ 单调, 而这在实践中容易验证).
\end{remark}

\begin{proof}
证明分为六步.

\emph{第 1 步: 唯一性.} 设 $u$ 与 $u'$ 是 \eqref{eq:7.6} 的两个解.我们有
\begin{equation*}
\left(\frac{d}{dt}(u-u'),(u-u')\right)
=-(A(u-u'),u-u')\le0.
\end{equation*}
但\footnote{记住: 若 $\phi\in C^1([0,+\infty);H)$, 则
$\abs{\phi}^2\in C^1([0,+\infty);\R)$, 且
$\dfrac{d}{dt}\abs{\phi}^2=2\left(\dfrac{d\phi}{dt},\phi\right)$.}
\begin{equation*}
\frac12\frac{d}{dt}\abs{u(t)-u'(t)}^2
=\left(\frac{d}{dt}(u(t)-u'(t)),u(t)-u'(t)\right).
\end{equation*}
因此函数 $t\mapsto\abs{u(t)-u'(t)}$ 在 $[0,+\infty)$ 上不增.由于
$\abs{u(0)-u'(0)}=0$, 推出
\begin{equation*}
\abs{u(t)-u'(t)}=0\quad\forall t\ge0.
\end{equation*}

证明存在性的主要想法是: 在 \eqref{eq:7.6} 中用 $A_\lambda$ 代替 $A$, 对
近似问题应用定理  \ref{thm:7.3}, 然后利用各种与 $\lambda$ 无关的估计令
$\lambda\to0$ 取极限.于是设 $u_\lambda$ 为问题
\begin{equation}
\begin{cases}
\dfrac{du_\lambda}{dt}+A_\lambda u_\lambda=0
& \text{在}\ [0,+\infty)\ \text{上},\\
u_\lambda(0)=u_0\in D(A)
\end{cases}
\label{eq:7.7}
\end{equation}
的解.

\emph{第 2 步:} 我们有估计
\begin{equation}
\abs{u_\lambda(t)}\le\abs{u_0}
\qquad\forall t\ge0,\ \forall\lambda>0,
\label{eq:7.8}
\end{equation}
\begin{equation}
\left|\frac{du_\lambda}{dt}(t)\right|=\abs{A_\lambda u_\lambda(t)}
\le\abs{Au_0}
\qquad\forall t\ge0,\ \forall\lambda>0.
\label{eq:7.9}
\end{equation}
它们直接由下述引理以及 $\abs{A_\lambda u_0}\le\abs{Au_0}$ 得出.

\begin{lemma}
\label{lem:7.1}
设 $w\in C^1([0,+\infty);H)$ 满足
\begin{equation}
\frac{dw}{dt}+A_\lambda w=0\quad\text{在}\ [0,+\infty)\ \text{上}.
\label{eq:7.10}
\end{equation}
则函数 $t\mapsto\abs{w(t)}$ 与 $t\mapsto
\left|\dfrac{dw}{dt}(t)\right|=\abs{A_\lambda w(t)}$ 在 $[0,+\infty)$ 上不增.
\end{lemma}

\begin{proof}
我们有
\begin{equation*}
\left(\frac{dw}{dt},w\right)+(A_\lambda w,w)=0.
\end{equation*}
由命题  \ref{prop:7.2}(e) 我们知道 $(A_\lambda w,w)\ge0$, 从而
$\frac12\frac{d}{dt}\abs{w}^2\le0$, 故 $\abs{w(t)}$ 不增.另一方面, 由于
$A_\lambda$ 是线性有界算子, 由 \eqref{eq:7.10} 可 (用归纳法) 推出
$w\in C^\infty([0,+\infty);H)$, 并且
\begin{equation*}
\frac{d}{dt}\left(\frac{dw}{dt}\right)
+A_\lambda\left(\frac{dw}{dt}\right)=0.
\end{equation*}
把上述事实应用于 $\dfrac{dw}{dt}$, 可见
$\left|\dfrac{dw}{dt}(t)\right|$ 不增.事实上, 对任意阶 $k$, 函数
$\left|\dfrac{d^kw}{dt^k}(t)\right|$ 都不增.
\end{proof}

\emph{第 3 步:} 这里我们将证明: 对每个 $t\ge0$, 当 $\lambda\to0$ 时
$u_\lambda(t)$ 收敛于某个极限, 记作 $u(t)$.此外, 这个收敛在每个有界区间
$[0,T]$ 上是一致的.

对每个 $\lambda,\mu>0$, 我们有
\begin{equation*}
\frac{du_\lambda}{dt}-\frac{du_\mu}{dt}
+A_\lambda u_\lambda-A_\mu u_\mu=0,
\end{equation*}
从而
\begin{equation}
\frac12\frac{d}{dt}\abs{u_\lambda(t)-u_\mu(t)}^2
+(A_\lambda u_\lambda(t)-A_\mu u_\mu(t),u_\lambda(t)-u_\mu(t))=0.
\label{eq:7.11}
\end{equation}
为简单起见略去 $t$, 我们写出
\begin{align}
&(A_\lambda u_\lambda-A_\mu u_\mu,u_\lambda-u_\mu)\notag\\
&\quad=(A_\lambda u_\lambda-A_\mu u_\mu,
u_\lambda-J_\lambda u_\lambda+J_\lambda u_\lambda
-J_\mu u_\mu+J_\mu u_\mu-u_\mu)\notag\\
&\quad=(A_\lambda u_\lambda-A_\mu u_\mu,\lambda A_\lambda u_\lambda
-\mu A_\mu u_\mu)+(A(J_\lambda u_\lambda-J_\mu u_\mu),
J_\lambda u_\lambda-J_\mu u_\mu)\notag\\
&\quad\ge(A_\lambda u_\lambda-A_\mu u_\mu,\lambda A_\lambda u_\lambda
-\mu A_\mu u_\mu).
\label{eq:7.12}
\end{align}
由 \eqref{eq:7.9}, \eqref{eq:7.11} 与 \eqref{eq:7.12} 推出
\begin{equation*}
\frac12\frac{d}{dt}\abs{u_\lambda-u_\mu}^2
\le2(\lambda+\mu)\abs{Au_0}^2.
\end{equation*}
对此不等式积分, 得到
\begin{equation*}
\abs{u_\lambda(t)-u_\mu(t)}^2\le4(\lambda+\mu)t\abs{Au_0}^2,
\end{equation*}
即
\begin{equation}
\abs{u_\lambda(t)-u_\mu(t)}\le2\sqrt{(\lambda+\mu)t}\,\abs{Au_0}.
\label{eq:7.13}
\end{equation}
由此推出, 对每个固定的 $t\ge0$, 当 $\lambda\to0$ 时 $u_\lambda(t)$ 是
Cauchy 列, 从而收敛于某个极限, 记作 $u(t)$.在 \eqref{eq:7.13} 中令
$\mu\to0$ 取极限, 我们有
\begin{equation*}
\abs{u_\lambda(t)-u(t)}\le2\sqrt{\lambda t}\,\abs{Au_0}.
\end{equation*}
因此, 收敛在 $t$ 上于每个有界区间 $[0,T]$ 一致, 故
$u\in C([0,+\infty);H)$.

\emph{第 4 步:} 这里再假设 $u_0\in D(A^2)$, 即 $u_0\in D(A)$ 且
$Au_0\in D(A)$, 我们证明当 $\lambda\to0$ 时
$\dfrac{du_\lambda}{dt}(t)$ 收敛于某个极限, 且收敛在每个有界区间 $[0,T]$ 上
一致.

令 $v_\lambda=\dfrac{du_\lambda}{dt}$, 于是
$\dfrac{dv_\lambda}{dt}+A_\lambda v_\lambda=0$.沿用第 3 步的论证, 可见
\begin{equation}
\frac12\frac{d}{dt}\abs{v_\lambda-v_\mu}^2
\le(\abs{A_\lambda v_\lambda}+\abs{A_\mu v_\mu})
(\lambda\abs{A_\lambda v_\lambda}+\mu\abs{A_\mu v_\mu}).
\label{eq:7.14}
\end{equation}
由引理  \ref{lem:7.1}, 我们有
\begin{equation}
\abs{A_\lambda v_\lambda(t)}\le\abs{A_\lambda v_\lambda(0)}
=\abs{A_\lambda A_\lambda u_0},
\label{eq:7.15}
\end{equation}
类似地
\begin{equation}
\abs{A_\mu v_\mu(t)}\le\abs{A_\mu v_\mu(0)}
=\abs{A_\mu A_\mu u_0}.
\label{eq:7.16}
\end{equation}
最后, 由于 $Au_0\in D(A)$, 我们得到
\begin{equation*}
A_\lambda A_\lambda u_0=J_\lambda AJ_\lambda Au_0
=J_\lambda J_\lambda AAu_0=J_\lambda^2A^2u_0,
\end{equation*}
从而
\begin{equation}
\abs{A_\lambda A_\lambda u_0}\le\abs{A^2u_0},
\qquad
\abs{A_\mu A_\mu u_0}\le\abs{A^2u_0}.
\label{eq:7.17}
\end{equation}
合并 \eqref{eq:7.14}, \eqref{eq:7.15}, \eqref{eq:7.16} 与
\eqref{eq:7.17}, 我们得到
\begin{equation*}
\frac12\frac{d}{dt}\abs{v_\lambda-v_\mu}^2
\le2(\lambda+\mu)\abs{A^2u_0}^2.
\end{equation*}
如第 3 步那样, 我们得出结论: 当 $\lambda\to0$ 时
$v_\lambda(t)=\dfrac{du_\lambda}{dt}(t)$ 收敛于某个极限, 且收敛在每个有界
区间 $[0,T]$ 上一致.

\emph{第 5 步:} 假设 $u_0\in D(A^2)$, 这里证明 $u$ 是 \eqref{eq:7.6} 的解.

由第 3 步与第 4 步, 我们知道对一切 $T<\infty$,
\begin{equation*}
\begin{cases}
u_\lambda(t)\to u(t), & \text{当}\ \lambda\to0,\ \text{在}\ [0,T]\ \text{上一致},\\
\dfrac{du_\lambda}{dt}(t)\ \text{当}\ \lambda\to0\ \text{时收敛},
& \text{在}\ [0,T]\ \text{上一致}.
\end{cases}
\end{equation*}
由此容易推出 $u\in C^1([0,+\infty);H)$, 且当 $\lambda\to0$ 时
$\dfrac{du_\lambda}{dt}(t)\to\dfrac{du}{dt}(t)$ 在 $[0,T]$ 上一致.把
\eqref{eq:7.7} 改写为
\begin{equation}
\frac{du_\lambda}{dt}(t)+A(J_\lambda u_\lambda(t))=0.
\label{eq:7.18}
\end{equation}
注意当 $\lambda\to0$ 时 $J_\lambda u_\lambda(t)\to u(t)$, 因为
\begin{equation*}
\abs{J_\lambda u_\lambda(t)-u(t)}
\le\abs{J_\lambda u_\lambda(t)-J_\lambda u(t)}
+\abs{J_\lambda u(t)-u(t)}
\le\abs{u_\lambda(t)-u(t)}+\abs{J_\lambda u(t)-u(t)}\to0.
\end{equation*}
利用 $A$ 有闭图像这一事实, 由 \eqref{eq:7.18} 推出
$u(t)\in D(A)\ \forall t\ge0$, 且
\begin{equation*}
\frac{du}{dt}(t)+Au(t)=0.
\end{equation*}
最后, 由于 $u\in C^1([0,+\infty);H)$, 函数 $t\mapsto Au(t)$ 由
$[0,+\infty)$ 到 $H$ 连续, 故 $u\in C([0,+\infty);D(A))$.于是我们得到了
\eqref{eq:7.6} 的一个解, 且此外还满足
\begin{equation*}
\abs{u(t)}\le\abs{u_0}\quad\forall t\ge0
\quad\text{与}\quad
\left|\frac{du}{dt}(t)\right|=\abs{Au(t)}\le\abs{Au_0}
\quad\forall t\ge0.
\end{equation*}

\emph{第 6 步:} 这里我们完成定理的证明.

我们将使用下述引理.

\begin{lemma}
\label{lem:7.2}
设 $u_0\in D(A)$.则 $\forall\varepsilon>0$, $\exists u'_0\in D(A^2)$,
使得 $\abs{u_0-u'_0}<\varepsilon$ 且 $\abs{Au_0-Au'_0}<\varepsilon$.换言之,
$D(A^2)$ 在 $D(A)$ 中稠密 (关于图范数).
\end{lemma}

\begin{proof}[引理 \ref{lem:7.2}]
令 $u'_0=J_\lambda u_0$, 其中 $\lambda>0$ 取得适当小 (稍后确定).我们有
\begin{equation*}
u'_0\in D(A)\quad\text{且}\quad u'_0+\lambda Au'_0=u_0.
\end{equation*}
于是 $Au'_0\in D(A)$, 即 $u'_0\in D(A^2)$.另一方面, 由命题
 \ref{prop:7.2}, 我们知道
\begin{equation*}
\lim_{\lambda\to0}\abs{J_\lambda u_0-u_0}=0,\qquad
\lim_{\lambda\to0}\abs{J_\lambda Au_0-Au_0}=0,
\end{equation*}
且
\begin{equation*}
J_\lambda Au_0=AJ_\lambda u_0.
\end{equation*}
取 $\lambda>0$ 充分小即得所需结论.
\end{proof}

现在我们回到定理  \ref{thm:7.4} 的证明.给定 $u_0\in D(A)$, 我们 (用引理
 \ref{lem:7.2}) 构造 $D(A^2)$ 中的序列 $(u_{0n})$, 使得 $u_{0n}\to u_0$ 且
$Au_{0n}\to Au_0$.由第 5 步, 我们知道问题
\begin{equation}
\begin{cases}
\dfrac{du_n}{dt}+Au_n=0 & \text{在}\ [0,+\infty)\ \text{上},\\
u_n(0)=u_{0n}
\end{cases}
\label{eq:7.19}
\end{equation}
有解 $u_n$.对一切 $t\ge0$, 我们有
\begin{equation*}
\abs{u_n(t)-u_m(t)}\le\abs{u_{0n}-u_{0m}}
\xrightarrow[m,n\to\infty]{}0,
\end{equation*}
\begin{equation*}
\left|\frac{du_n}{dt}(t)-\frac{du_m}{dt}(t)\right|
\le\abs{Au_{0n}-Au_{0m}}
\xrightarrow[m,n\to\infty]{}0.
\end{equation*}
因此
\begin{equation*}
u_n(t)\to u(t)\ \text{在}\ [0,+\infty)\ \text{上一致},
\qquad
\frac{du_n}{dt}(t)\to\frac{du}{dt}(t)
\ \text{在}\ [0,+\infty)\ \text{上一致},
\end{equation*}
其中 $u\in C^1([0,+\infty);H)$.在 \eqref{eq:7.19} 中取极限——利用 $A$ 是
闭算子这一事实——可见 $u(t)\in D(A)$ 且 $u$ 满足 \eqref{eq:7.6}.由
\eqref{eq:7.6} 推出 $u\in C([0,+\infty);D(A))$.
\end{proof}

\begin{remark}
\label{rem:7.4}
设 $u_\lambda$ 为 \eqref{eq:7.7} 的解:
\begin{enumerate}[label=(\alph*), leftmargin=2.2em, itemsep=0.2em]
	\item 假设 $u_0\in D(A)$.我们知道 (由第 3 步), 当 $\lambda\to0$ 时,
	      对每个 $t\ge0$, $u_\lambda(t)$ 收敛于某个极限 $u(t)$.可以直接
	      证明 $u\in C^1([0,+\infty);H)\cap C([0,+\infty);D(A))$, 且它满足
	      \eqref{eq:7.6}.
	\item 只假设 $u_0\in H$.仍可证明当 $\lambda\to0$ 时, 对每个 $t\ge0$,
	      $u_\lambda(t)$ 收敛于某个极限, 记作 $u(t)$.但可能出现这样的情形:
	      这个极限 $u(t)$ 对一切 $t>0$ 都不属于 $D(A)$, 且 $u(t)$ 在
	      $[0,+\infty)$ 上处处不可微.因此 $u(t)$ 不是 \eqref{eq:7.6} 的
	      「古典」解.事实上, 对这样的 $u_0$, 问题 \eqref{eq:7.6} 没有古典
	      解.尽管如此, 我们仍可把 $u(t)$ 看作 \eqref{eq:7.6} 的「广义」解.
	      我们将在第~\ref{sec:7.4}~节看到, 当 $A$ 自伴时这种情况不会发生: 此时对每个
	      $u_0\in H$ (即使 $u_0\notin D(A)$), $u(t)$ 都是 \eqref{eq:7.6}
	      的「古典」解.
\end{enumerate}
\end{remark}

\begin{starremark}[压缩半群]
\label{rem:7.5}
对每个 $t\ge0$, 考虑线性映射 $u_0\in D(A)\mapsto u(t)\in D(A)$, 其中
$u(t)$ 是定理  \ref{thm:7.4} 给出的 \eqref{eq:7.6} 的解.由于
$\abs{u(t)}\le\abs{u_0}$ 且 $D(A)$ 在 $H$ 中稠密, 我们可以按连续性把这个
映射延拓为由 $H$ 到自身的有界算子, 记作 $S_A(t)$.\footnote{或者也可以
用注记 \ref{rem:7.4}(b) 直接在 $H$ 上把 $S_A(t)$ 定义为映射
$u_0\in H\mapsto u(t)\in H$.} 容易验证 $S_A(t)$ 满足下述性质:
\begin{enumerate}[label=(\alph*), leftmargin=2.2em, itemsep=0.2em]
	\item 对每个 $t\ge0$, $S_A(t)\in\LL(H)$ 且
	      $\norm{S_A(t)}_{\LL(H)}\le1$,
	\item
	      \begin{equation*}
	      \begin{cases}
	      S_A(t_1+t_2)=S_A(t_1)\circ S_A(t_2) & \forall t_1,t_2\ge0,\\
	      S_A(0)=I,
	      \end{cases}
	      \end{equation*}
	\item
	      \begin{equation*}
	      \lim_{\substack{t\to0\\ t>0}}\abs{S_A(t)u_0-u_0}=0
	      \quad\forall u_0\in H.
	      \end{equation*}
\end{enumerate}
这样一族依赖于参数 $t\ge0$, 由 $H$ 到自身的算子 $\{S(t)\}_{t\ge0}$ 若满足
(a), (b), (c), 则称为\emph{连续压缩半群}.

Hille 与 Yosida 的一个值得注意的结果断言: 反之, 给定 $H$ 上的连续压缩半群
$S(t)$, 存在唯一的极大单调算子 $A$, 使得
$S(t)=S_A(t)\ \forall t\ge0$.这就在极大单调算子与连续压缩半群之间建立了
一一对应.(证明见第~
\ref{chap:7}~章评注中引用的文献.)
\end{starremark}

\begin{keyremark}
\label{rem:7.6}
设 $A$ 为极大单调算子, $\lambda\in\R$.问题
\begin{equation*}
\begin{cases}
\dfrac{du}{dt}+Au+\lambda u=0 & \text{在}\ [0,+\infty)\ \text{上},\\
u(0)=u_0
\end{cases}
\end{equation*}
可用下述简单办法归结为问题 \eqref{eq:7.6}.令
\begin{equation*}
v(t)=e^{\lambda t}u(t).
\end{equation*}
则 $v$ 满足
\begin{equation*}
\begin{cases}
\dfrac{dv}{dt}+Av=0 & \text{在}\ [0,+\infty)\ \text{上},\\
v(0)=u_0.
\end{cases}
\end{equation*}
\end{keyremark}

\section{正则性}
\label{sec:7.3}

这里我们将证明: 只要对初值 $u_0$ 作附加假设, 定理  \ref{thm:7.4} 中得到的
\eqref{eq:7.6} 的解 $u$ 就比 $C^1([0,+\infty);H)\cap
C([0,+\infty);D(A))$ 更正则.为此, 我们用归纳法定义空间
\begin{equation*}
D(A^k)=\{v\in D(A^{k-1});\ Av\in D(A^{k-1})\},
\end{equation*}
其中 $k$ 为任何整数, $k\ge2$.容易看出 $D(A^k)$ 关于标量积
\begin{equation*}
(u,v)_{D(A^k)}=\sum_{j=0}^k(A^ju,A^jv)
\end{equation*}
是 Hilbert 空间; 相应的范数为
\begin{equation*}
\abs{u}_{D(A^k)}=\left(\sum_{j=0}^k\abs{A^ju}^2\right)^{1/2}.
\end{equation*}

\begin{theorem}
\label{thm:7.5}
假设对某个整数 $k\ge2$, $u_0\in D(A^k)$.则定理  \ref{thm:7.4} 中得到的
问题 \eqref{eq:7.6} 的解 $u$ 满足
\begin{equation*}
u\in C^{k-j}([0,+\infty);D(A^j))\quad\forall j=0,1,\dots,k.
\end{equation*}
\end{theorem}

\begin{proof}
先设 $k=2$.考虑赋以标量积 $(u,v)_{D(A)}$ 的 Hilbert 空间 $H_1=D(A)$.
容易验证由
\begin{equation*}
\begin{cases}
D(A_1)=D(A^2),\\
A_1u=Au & \text{对}\ u\in D(A_1)
\end{cases}
\end{equation*}
定义的算子 $A_1:D(A_1)\subset H_1\to H_1$ 在 $H_1$ 中极大单调.把定理
 \ref{thm:7.4} 应用于空间 $H_1$ 中的算子 $A_1$, 可见存在函数
\begin{equation*}
u\in C^1([0,+\infty);H_1)\cap C([0,+\infty);D(A_1))
\end{equation*}
使得
\begin{equation*}
\begin{cases}
\dfrac{du}{dt}+A_1u=0 & \text{在}\ [0,+\infty)\ \text{上},\\
u(0)=u_0.
\end{cases}
\end{equation*}
特别地, $u$ 满足 \eqref{eq:7.6}; 由唯一性, 这个 $u$ 就是 \eqref{eq:7.6}
的解.只须再验证 $u\in C^2([0,+\infty);H)$.由于
\begin{equation*}
A\in\LL(H_1,H)\quad\text{且}\quad u\in C([0,+\infty);H_1),
\end{equation*}
推出 $Au\in C^1([0,+\infty);H)$ 且
\begin{equation}
\frac{d}{dt}(Au)=A\left(\frac{du}{dt}\right).
\label{eq:7.20}
\end{equation}
应用 \eqref{eq:7.6}, 可见 $\dfrac{du}{dt}\in C^1([0,+\infty);H)$, 即
$u\in C^2([0,+\infty);H)$, 且
\begin{equation}
\frac{d}{dt}\left(\frac{du}{dt}\right)
+A\left(\frac{du}{dt}\right)=0\quad\text{在}\ [0,+\infty)\ \text{上}.
\label{eq:7.21}
\end{equation}

现在我们转向一般情形 $k\ge3$.我们对 $k$ 归纳: 假设结论对直到 $k-1$ 的
阶数都成立, 并设 $u_0\in D(A^k)$.由前面的分析, 我们知道 \eqref{eq:7.6}
的解 $u$ 属于 $C^2([0,+\infty);H)\cap C^1([0,+\infty);D(A))$, 且 $u$ 满足
\eqref{eq:7.21}.令
\begin{equation*}
v=\frac{du}{dt},
\end{equation*}
我们有
\begin{equation*}
v\in C^1([0,+\infty);H)\cap C([0,+\infty);D(A)),
\end{equation*}
\begin{equation*}
\begin{cases}
\dfrac{dv}{dt}+Av=0 & \text{在}\ [0,+\infty)\ \text{上},\\
v(0)=-Au_0.
\end{cases}
\end{equation*}
换言之, $v$ 是对应于初值 $v_0=-Au_0$ 的 \eqref{eq:7.6} 的解.由于
$v_0\in D(A^{k-1})$, 由归纳假设我们知道
\begin{equation}
v\in C^{k-1-j}([0,+\infty);D(A^j))\quad\forall j=0,1,\dots,k-1,
\label{eq:7.22}
\end{equation}
即
\begin{equation*}
u\in C^{k-j}([0,+\infty);D(A^j))\quad\forall j=0,1,\dots,k-1.
\end{equation*}
只须再验证
\begin{equation}
u\in C([0,+\infty);D(A^k)).
\label{eq:7.23}
\end{equation}
在 \eqref{eq:7.22} 中取 $j=k-1$, 可见
\begin{equation}
\frac{du}{dt}\in C([0,+\infty);D(A^{k-1})).
\label{eq:7.24}
\end{equation}
由 \eqref{eq:7.24} 与方程 \eqref{eq:7.6} 推出
\begin{equation*}
Au\in C([0,+\infty);D(A^{k-1})),
\end{equation*}
即 \eqref{eq:7.23}.
\end{proof}

\section{自伴情形}
\label{sec:7.4}

设 $A:D(A)\subset H\to H$ 为无界线性算子, $\overline{D(A)}=H$.把
$H^\star$ 与 $H$ 等同, 我们可以把 $A^\star$ 看作 $H$ 中的无界线性算子.

\begin{definition}
称
\begin{itemize}[leftmargin=1.6em, itemsep=0.15em]
	\item $A$ 是\emph{对称的}, 若 $(Au,v)=(u,Av)\ \forall u,v\in D(A)$,
	\item $A$ 是\emph{自伴的}, 若 $D(A^\star)=D(A)$ 且 $A^\star=A$.
\end{itemize}
\end{definition}

\begin{remark}
\label{rem:7.7}
对有界算子, 对称算子与自伴算子的概念一致.然而, 若 $A$ 无界, 则对称算子与
自伴算子之间有微妙的差别.显然, 任何自伴算子都是对称的.逆命题不成立: 算子
$A$ 对称当且仅当 $A\subset A^\star$, 即 $D(A)\subset D(A^\star)$ 且在
$D(A)$ 上 $A^\star=A$.可能出现 $A$ 对称而 $D(A)\ne D(A^\star)$ 的情形.
我们下一个结果表明: 若 $A$ 极大单调, 则
\begin{equation*}
(A\ \text{对称})\iff(A\ \text{自伴}).
\end{equation*}
\end{remark}

\begin{proposition}
\label{prop:7.6}
设 $A$ 为极大单调的对称算子.则 $A$ 自伴.
\end{proposition}

\begin{proof}
令 $J_1=(I+A)^{-1}$.我们先证明 $J_1$ 自伴.由于 $J_1\in\LL(H)$, 只须
验证
\begin{equation}
(J_1u,v)=(u,J_1v)\quad\forall u,v\in H.
\label{eq:7.25}
\end{equation}
令 $u_1=J_1u$, $v_1=J_1v$, 于是
\begin{equation*}
u_1+Au_1=u,\qquad v_1+Av_1=v.
\end{equation*}
由于按假设 $(u_1,Av_1)=(Au_1,v_1)$, 推出 $(u_1,v)=(u,v_1)$, 即
\eqref{eq:7.25}.

设 $u\in D(A^\star)$, 令 $f=u+A^\star u$.我们有
\begin{equation*}
(f,v)=(u,v+Av)\quad\forall v\in D(A),
\end{equation*}
即
\begin{equation*}
(f,J_1w)=(u,w)\quad\forall w\in H.
\end{equation*}
因此 $u=J_1f$, 从而 $u\in D(A)$.这就证明了 $D(A^\star)=D(A)$, 因此 $A$
自伴.
\end{proof}

\begin{remark}
\label{rem:7.8}
必须注意: 若 $A$ 是单调算子 (甚至是对称单调算子), 则 $A^\star$ 不必单调.
然而可以证明下述性质等价:
\begin{equation*}
A\ \text{极大单调}\iff A^\star\ \text{极大单调}
\iff A\ \text{闭},\ D(A)\ \text{稠密},\ A\ \text{与}\ A^\star\
\text{单调}.
\end{equation*}
这个结果的一个更一般的版本见问题 16.
\end{remark}

\begin{keytheorem}
\label{thm:7.7}
设 $A$ 为自伴的极大单调算子.则对每个 $u_0\in H$, 存在唯一的函数\footnote{
让我们强调定理  \ref{thm:7.4} 与定理  \ref{thm:7.7} 的区别.这里
$u_0\in H$ (而不是 $u_0\in D(A)$); 结论是 \eqref{eq:7.6} 有解, 且该解在
远离 $t=0$ 处光滑.然而, $\left|\dfrac{du}{dt}(t)\right|$ 当 $t\to0$ 时可能
「爆炸」.}
\begin{equation*}
u\in C([0,+\infty);H)\cap C^1((0,+\infty);H)
\cap C((0,+\infty);D(A))
\end{equation*}
使得
\begin{equation*}
\begin{cases}
\dfrac{du}{dt}+Au=0 & \text{在}\ (0,+\infty)\ \text{上},\\
u(0)=u_0.
\end{cases}
\end{equation*}
此外, 我们有
\begin{equation*}
\abs{u(t)}\le\abs{u_0}
\quad\text{且}\quad
\left|\frac{du}{dt}(t)\right|=\abs{Au(t)}\le\frac1t\abs{u_0}
\quad\forall t>0,
\end{equation*}
\begin{equation}
u\in C^k((0,+\infty);D(A^\ell))\quad\forall k,\ell\ \text{为整数}.
\label{eq:7.26}
\end{equation}
\end{keytheorem}

\begin{proof}
\emph{唯一性.} 设 $u$ 与 $u'$ 是两个解.由 $A$ 的单调性, 可见
$\phi(t)=\abs{u(t)-u'(t)}^2$ 在 $(0,+\infty)$ 上不增.另一方面, $\phi$ 在
$[0,+\infty)$ 上连续且 $\phi(0)=0$.因此 $\phi\equiv0$.

\emph{存在性.} 证明分为两步:

\emph{第 1 步.} 先设 $u_0\in D(A^2)$, 并设 $u$ 是定理  \ref{thm:7.4} 给出的
\eqref{eq:7.6} 的解.我们断言
\begin{equation}
\left|\frac{du}{dt}(t)\right|\le\frac1t\abs{u_0}\quad\forall t>0.
\label{eq:7.27}
\end{equation}
如命题  \ref{prop:7.6} 的证明中那样, 我们有
\begin{equation*}
J_\lambda^\star=J_\lambda\quad\text{与}\quad A_\lambda^\star=A_\lambda
\quad\forall\lambda>0.
\end{equation*}
我们回到定理  \ref{thm:7.4} 的证明中引入的近似问题:
\begin{equation}
\frac{du_\lambda}{dt}+A_\lambda u_\lambda=0\ \text{在}\ [0,+\infty)\ \text{上},
\qquad u_\lambda(0)=u_0.
\label{eq:7.28}
\end{equation}
把 \eqref{eq:7.28} 与 $u_\lambda$ 作标量积并在 $[0,T]$ 上积分, 得到
\begin{equation}
\frac12\abs{u_\lambda(T)}^2+\int_0^T(A_\lambda u_\lambda,u_\lambda)\,dt
=\frac12\abs{u_0}^2.
\label{eq:7.29}
\end{equation}
把 \eqref{eq:7.28} 与 $t\dfrac{du_\lambda}{dt}$ 作标量积并在 $[0,T]$ 上
积分, 得到
\begin{equation}
\int_0^T\left|\frac{du_\lambda}{dt}(t)\right|^2t\,dt
+\int_0^T\left(A_\lambda u_\lambda(t),
\frac{du_\lambda}{dt}(t)\right)t\,dt=0.
\label{eq:7.30}
\end{equation}
但
\begin{equation*}
\frac{d}{dt}(A_\lambda u_\lambda,u_\lambda)
=\left(A_\lambda\frac{du_\lambda}{dt},u_\lambda\right)
+\left(A_\lambda u_\lambda,\frac{du_\lambda}{dt}\right)
=2\left(A_\lambda u_\lambda,\frac{du_\lambda}{dt}\right),
\end{equation*}
因为 $A_\lambda^\star=A_\lambda$.对 \eqref{eq:7.30} 中的第二个积分作分部
积分, 得到
\begin{align}
\int_0^T\left(A_\lambda u_\lambda(t),
\frac{du_\lambda}{dt}(t)\right)t\,dt
&=\frac12\int_0^T\frac{d}{dt}\bigl[(A_\lambda u_\lambda,u_\lambda)\bigr]t\,dt
\notag\\
&=\frac12(A_\lambda u_\lambda(T),u_\lambda(T))T
-\frac12\int_0^T(A_\lambda u_\lambda,u_\lambda)\,dt.
\label{eq:7.31}
\end{align}
另一方面, 由于函数 $t\mapsto\left|\dfrac{du_\lambda}{dt}(t)\right|$ 不增
(由引理  \ref{lem:7.1}), 我们有
\begin{equation}
\int_0^T\left|\frac{du_\lambda}{dt}(t)\right|^2t\,dt
\ge\left|\frac{du_\lambda}{dt}(T)\right|^2\frac{T^2}{2}.
\label{eq:7.32}
\end{equation}
合并 \eqref{eq:7.29}, \eqref{eq:7.30}, \eqref{eq:7.31} 与
\eqref{eq:7.32}, 我们得到
\begin{equation*}
\frac12\abs{u_\lambda(T)}^2+T(A_\lambda u_\lambda(T),u_\lambda(T))
+T^2\left|\frac{du_\lambda}{dt}(T)\right|^2\le\frac12\abs{u_0}^2;
\end{equation*}
特别地推出
\begin{equation}
\left|\frac{du_\lambda}{dt}(T)\right|\le\frac1T\abs{u_0}
\quad\forall T>0.
\label{eq:7.33}
\end{equation}
最后, 在 \eqref{eq:7.33} 中令 $\lambda\to0$ 取极限.这就完成了
\eqref{eq:7.27} 的证明, 因为
$\dfrac{du_\lambda}{dt}\to\dfrac{du}{dt}$ (见定理  \ref{thm:7.4} 证明中的
第 5 步).

\emph{第 2 步.} 现在假设 $u_0\in H$.设 $(u_{0n})$ 为 $D(A^2)$ 中的序列,
使得 $u_{0n}\to u_0$ (回忆 $D(A^2)$ 在 $D(A)$ 中稠密, 而 $D(A)$ 在 $H$ 中
稠密; 因此 $D(A^2)$ 在 $H$ 中稠密).设 $u_n$ 为
\begin{equation*}
\begin{cases}
\dfrac{du_n}{dt}+Au_n=0 & \text{在}\ [0,+\infty)\ \text{上},\\
u_n(0)=u_{0n}
\end{cases}
\end{equation*}
的解.由定理  \ref{thm:7.4}, 我们知道
\begin{equation*}
\abs{u_n(t)-u_m(t)}\le\abs{u_{0n}-u_{0m}}
\quad\forall m,n,\ \forall t\ge0,
\end{equation*}
而由第 1 步
\begin{equation*}
\left|\frac{du_n}{dt}(t)-\frac{du_m}{dt}(t)\right|
\le\frac1t\abs{u_{0n}-u_{0m}}\quad\forall m,n,\ \forall t>0.
\end{equation*}
由此推出 $u_n$ 在 $[0,+\infty)$ 上一致收敛于某个极限 $u(t)$, 且
$\dfrac{du_n}{dt}(t)$ 在每个区间 $[\delta,+\infty)$ ($\delta>0$) 上一致收敛
于 $\dfrac{du}{dt}(t)$.极限函数 $u$ 满足
\begin{equation*}
u\in C([0,+\infty);H)\cap C^1((0,+\infty);H),
\qquad u(t)\in D(A)\ \forall t>0
\end{equation*}
且
\begin{equation*}
\frac{du}{dt}(t)+Au(t)=0\quad\forall t>0
\end{equation*}
(这里用到 $A$ 闭这一事实).

现在我们转向 \eqref{eq:7.26} 的证明.我们将对 $k\ge2$ 归纳证明
\begin{equation}
u\in C^{k-j}((0,+\infty);D(A^j))\quad\forall j=0,1,\dots,k.
\label{eq:7.34}
\end{equation}
假设 \eqref{eq:7.34} 对直到 $k-1$ 的阶数都成立.特别地, 我们有
\begin{equation}
u\in C((0,+\infty);D(A^{k-1})).
\label{eq:7.35}
\end{equation}
为证明 \eqref{eq:7.34}, 只须 (由定理  \ref{thm:7.5}) 验证
\begin{equation}
u\in C((0,+\infty);D(A^k)).
\label{eq:7.36}
\end{equation}
考虑 Hilbert 空间 $\tilde H=D(A^{k-1})$ 以及由
\begin{equation*}
\begin{cases}
D(\tilde A)=D(A^k),\\
\tilde A=A
\end{cases}
\end{equation*}
定义的算子 $\tilde A:D(\tilde A)\subset\tilde H\to\tilde H$.容易看出
$\tilde A$ 在 $\tilde H$ 中极大单调且对称; 因此它自伴.把定理  \ref{thm:7.7}
的第一个断言应用于空间 $\tilde H$ 中的算子 $\tilde A$, 可见对任何
$v_0\in\tilde H$, 问题
\begin{equation}
\begin{cases}
\dfrac{dv}{dt}+Av=0 & \text{在}\ (0,+\infty)\ \text{上},\\
v(0)=v_0
\end{cases}
\label{eq:7.37}
\end{equation}
有唯一解 $v$.此外,
\begin{equation*}
v\in C([0,+\infty);\tilde H)\cap C^1((0,+\infty);\tilde H)
\cap C((0,+\infty);D(\tilde A)).
\end{equation*}
取 $v_0=u(\varepsilon)$ ($\varepsilon>0$)——由 \eqref{eq:7.35} 我们已经知道
$v_0\in\tilde H$——我们推出 $u\in C((\varepsilon,+\infty);D(A^k))$, 这就
完成了 \eqref{eq:7.36} 的证明.
\end{proof}

\section*{第 7 章评注}
\addcontentsline{toc}{section}{第 7 章评注}
\markright{第 7 章评注}

\cmtitem{1. Banach 空间中的 Hille--Yosida 定理.}
Hille--Yosida 定理可以推广到 Banach 空间.确切的陈述如下.设 $E$ 为 Banach
空间, $A:D(A)\subset E\to E$ 为无界线性算子.若 $\overline{D(A)}=E$ 且对
每个 $\lambda>0$, $I+\lambda A$ 是由 $D(A)$ 到 $E$ 的双射, 并且
$\norm{(I+\lambda A)^{-1}}_{\LL(E)}\le1$, 则称 $A$ 是 \emph{$m$-增生}的.

\begin{theorem}[Hille--Yosida]
\label{thm:7.8}
设 $A$ 为 $m$-增生算子.则给定任何 $u_0\in D(A)$, 存在唯一的函数
\begin{equation*}
u\in C^1([0,+\infty);E)\cap C([0,+\infty);D(A))
\end{equation*}
使得
\begin{equation}
\begin{cases}
\dfrac{du}{dt}+Au=0 & \text{在}\ [0,+\infty)\ \text{上},\\
u(0)=u_0.
\end{cases}
\label{eq:7.38}
\end{equation}
此外,
\begin{equation*}
\norm{u(t)}\le\norm{u_0}
\quad\text{且}\quad
\left\|\frac{du}{dt}(t)\right\|=\norm{Au(t)}\le\norm{Au_0}
\quad\forall t\ge0.
\end{equation*}
把映射 $u_0\mapsto u(t)$ 按连续性延拓到整个 $E$, 记作 $S_A(t)$.它是 $E$
上的连续压缩半群.反之, 给定任何连续压缩半群 $S(t)$, 存在唯一的
$m$-增生算子 $A$, 使得 $S(t)=S_A(t)\ \forall t\ge0$.
\end{theorem}

关于证明, 见例如 P.~Lax \bref{Lax}{1}, A.~Pazy \bref{Pazy}{1}, J.~Goldstein \bref{Goldstein}{1}, E.~Davies \bref{Davies}{1}, \bref{Davies}{2}, K.~Yosida \bref{Yosida}{1}, M.~Reed--B.~Simon
\bref{ReedSimon}{1} 第 2 卷, H.~Tanabe \bref{Tanabe}{1}, N.~Dunford--J.~T.~Schwartz \bref{DunfordSchwartz}{1} 第 1 卷, M.~Schechter \bref{Schechter}{1}, A.~Friedman \bref{Friedman}{2}, R.~Dautray--J.-L.~Lions \bref{DautrayLions}{1} 第 XVII 章, A.~Balakrishnan \bref{Balakrishnan}{1}, T.~Kato \bref{Kato}{1}, W.~Rudin \bref{Rudin}{1}.这些文献对半群理论作了大量展开.

\cmtitem{2. 指数公式.}
求解 \eqref{eq:7.38} 有许多迭代技巧.这里提一种基本方法.

\begin{theorem}
\label{thm:7.9}
假设 $A$ 为 $m$-增生算子.则对每个 $u_0\in D(A)$, \eqref{eq:7.38} 的解
$u$ 由「指数公式」给出:
\begin{equation}
u(t)=\lim_{n\to+\infty}
\left[\left(I+\frac tnA\right)^{-1}\right]^nu_0.
\label{eq:7.39}
\end{equation}
\end{theorem}

证明见例如 K.~Yosida \bref{Yosida}{1} 与 A.~Pazy \bref{Pazy}{1}.公式
\eqref{eq:7.39} 用数值分析的语言来说, 对应于 \eqref{eq:7.38} 的隐式时间
离散格式的收敛性 (见例如 K.~W.~Morton--D.~F.~Mayers
\bref{MortonMayers}{1}).更确切地说, 把区间 $[0,t]$ 分成 $n$ 个等长的小区间
$\Delta t=t/n$, 并从 $u^0$ 出发递推求解方程
\begin{equation*}
\frac{u^{j+1}-u^j}{\Delta t}+Au^{j+1}=0,
\quad j=0,1,\dots,n-1.
\end{equation*}
换言之, $u^n$ 由
\begin{equation*}
u^n=(I+\Delta tA)^{-n}u_0
=\left(I+\frac tnA\right)^{-n}u_0
\end{equation*}
给出.当 $n\to\infty$ (即 $\Delta t\to0$) 时, $u^n$ 收敛于 $u(t)$ 是
「符合直觉的」.

\cmtnum{3.} 定理  \ref{thm:7.7} 是解析半群理论的第一步.关于这个专题, 见
例如 K.~Yosida \bref{Yosida}{1}, T.~Kato \bref{Kato}{1}, M.~Reed--B.~Simon \bref{ReedSimon}{1} 第 2 卷, A.~Friedman
\bref{Friedman}{2}, A.~Pazy \bref{Pazy}{1}, 以及 H.~Tanabe
\bref{Tanabe}{1}.

\cmtitem{4. 非齐次方程. 非线性方程.}
在 Banach 空间 $E$ 中考虑问题
\begin{equation}
\begin{cases}
\dfrac{du}{dt}(t)+Au(t)=f(t) & \text{在}\ [0,T]\ \text{上},\\
u(0)=u_0.
\end{cases}
\label{eq:7.40}
\end{equation}
下述结论成立.

\begin{theorem}
\label{thm:7.10}
假设 $A$ 为 $m$-增生算子.则对每个 $u_0\in D(A)$ 与每个
$f\in C^1([0,T];E)$, 存在 \eqref{eq:7.40} 的唯一解 $u$, 满足
\begin{equation*}
u\in C^1([0,T];E)\cap C([0,T];D(A)).
\end{equation*}
此外, $u$ 由公式
\begin{equation}
u(t)=S_A(t)u_0+\int_0^tS_A(t-s)f(s)\,ds
\label{eq:7.41}
\end{equation}
给出, 其中 $S_A(t)$ 是评注 1 中引入的半群.
\end{theorem}

注意, 若只假设 $f\in L^1((0,T);E)$, 公式 \eqref{eq:7.41} 仍有意义, 并
给出 \eqref{eq:7.40} 的广义解.关于这些问题, 见例如 T.~Kato
\bref{Kato}{1}, A.~Pazy \bref{Pazy}{1}, R.~H.~Martin \bref{Martin}{1}, H.~Tanabe \bref{Tanabe}{1}.

在物理应用中会遇到许多形如
\begin{equation*}
\frac{du}{dt}+Au=F(u)
\end{equation*}
的「半线性」方程, 其中 $F$ 是由 $E$ 到 $E$ 的非线性映射.关于这些问题,
见例如 R.~H.~Martin \bref{Martin}{1}, Th.~Cazenave--A.~Haraux \bref{CazenaveHaraux}{1}, 以及第~
\ref{chap:10}~章的评注.

我们还要提到, 第~
\ref{chap:7}~章的一些结果有非线性推广.考虑非线性的 $m$-增生算子
$A:D(A)\subset E\to E$ 是有用的.关于这个专题, 见例如 H.~Brezis
\bref{Brezis}{1} 与 V.~Barbu \bref{Barbu}{1}.

% chapters/ch08.tex
% chapters/ch08.tex
% 第~
% \ref{chap:8}~章 Sobolev 空间与一维边值问题的变分表述
% 源: 原书 PDF 第 208--267 页 (印刷页 201--260)

\chapter{Sobolev 空间与一维边值问题的变分表述}
\label{chap:8}

\section{动机}
\label{sec:8.1}

考虑如下问题.给定 $f\in C([a,b])$, 求一个函数 $u$ 满足
\begin{equation}
\begin{cases}
-u''+u=f & \text{在}\ [a,b]\ \text{上},\\
u(a)=u(b)=0.
\end{cases}
\label{eq:8.1}
\end{equation}
\eqref{eq:8.1} 的\emph{古典解}——或\emph{强解}——是指满足 \eqref{eq:8.1} 的通常
意义的 $C^2$ 函数.众所周知, 通过一个非常简单的计算就能显式地解出
\eqref{eq:8.1}, 但我们回避这一点, 以便在这个初等例子上说明方法.

用 $\varphi\in C^1([a,b])$ 乘 \eqref{eq:8.1} 并分部积分; 我们得到
\begin{equation}
\int_a^b u'\varphi'+\int_a^b u\varphi=\int_a^b f\varphi
\quad\forall\varphi\in C^1([a,b]),\ \varphi(a)=\varphi(b)=0.
\label{eq:8.2}
\end{equation}
注意只要 $u\in C^1([a,b])$, \eqref{eq:8.2} 就有意义 (而 \eqref{eq:8.1} 要求
$u$ 有二阶导数); 事实上, 只须知道 $u,u'\in L^1(a,b)$, 其中 $u'$ 的意义尚待明确.
我们 (暂时) 说, 满足 \eqref{eq:8.2} 的 $C^1$ 函数 $u$ 是 \eqref{eq:8.1} 的
\emph{弱解}.

下面的程序概括了偏微分方程理论中\emph{变分方法}的主要步骤:

\runin{步骤 A.} 明确\emph{弱解}的概念.这涉及 \emph{Sobolev 空间}, 它们是我们的
\emph{基本工具}.

\runin{步骤 B.} 通过 Lax--Milgram 定理用变分方法建立\emph{弱解的存在性与唯一性}.

\runin{步骤 C.} 证明弱解属于 $C^2$ 类 (例如): 这是一个\emph{正则性}结果.

\runin{步骤 D.} 通过证明任何 $C^2$ 的弱解都是古典解, 从而回到古典解.

完成步骤 D 非常简单.事实上, 设 $u\in C^2([a,b])$, $u(a)=u(b)=0$, 且 $u$ 满足
\eqref{eq:8.2}.对 \eqref{eq:8.2} 分部积分, 得到
\begin{equation*}
\int_a^b(-u''+u-f)\varphi=0
\quad\forall\varphi\in C^1([a,b]),\ \varphi(a)=\varphi(b)=0,
\end{equation*}
从而
\begin{equation*}
\int_a^b(-u''+u-f)\varphi=0
\quad\forall\varphi\in C_c^1((a,b)).
\end{equation*}
由此推出 (见推论  \ref{cor:4.24}) $-u''+u=f$ 在 $(a,b)$ 上几乎处处成立, 于是
在 $[a,b]$ 上处处成立, 因为 $u\in C^2([a,b])$.
% 原书此处印作 "see Corollary 4.15"; 第~
% \ref{chap:4}~章并无编号 4.15 的推论, 引理~
% \ref{lem:8.1} 的证明
% (印刷页 205) 处又引作 "Corollary 4.24", 所指的显然是同一条 (u 的积分为零 ⇒ u=0
% a.e.). 依技能规范保留原样并记录, 不擅自更正.

\section{\index{yiweiSobolevkongjian@一维 Sobolev 空间 (Sobolev space in dimension 1)}
Sobolev 空间 $W^{1,p}(I)$}
\label{sec:8.2}

设 $I=(a,b)$ 为开区间, 可以无界, 并设 $p\in\R$ 满足 $1\le p\le\infty$.

\begin{definition}
\emph{Sobolev 空间} $W^{1,p}(I)$\footnote{若无混淆, 我们将写 $W^{1,p}$ 而不写
$W^{1,p}(I)$, 写 $H^1$ 而不写 $H^1(I)$.} 定义为
\begin{equation*}
W^{1,p}(I)=\left\{u\in L^p(I);\ \exists g\in L^p(I)\ \text{使得}\
\int_I u\varphi'=-\int_I g\varphi\ \ \forall\varphi\in C_c^1(I)\right\}.
\end{equation*}
我们记
\begin{equation*}
H^1(I)=W^{1,2}(I).
\end{equation*}
对 $u\in W^{1,p}(I)$, 我们记\footnote{注意这有意义: 由推论  \ref{cor:4.24},
$g$ 几乎处处唯一确定.} $u'=g$.
\end{definition}

\begin{remark}
\label{rem:8.1}
在 $W^{1,p}$ 的定义中, 我们称 $\varphi$ 为\emph{检验函数}.我们本来同样可以
用 $C_c^\infty(I)$ 作为检验函数的类, 因为若 $\varphi\in C_c^1(I)$, 则当 $n$
足够大时 $\rho_n\star\varphi\in C_c^\infty(I)$, 且 $\rho_n\star\varphi\to\varphi$
在 $C^1$ 中 (见第~\ref{sec:4.4}~节; 当然, $\varphi$ 在 $I$ 外延拓为 $0$).
\end{remark}

\begin{remark}
\label{rem:8.2}
显然, 若 $u\in C^1(I)\cap L^p(I)$ 且 $u'\in L^p(I)$ (这里 $u'$ 是 $u$ 的通常
导数), 则 $u\in W^{1,p}(I)$.此外, $u$ 的通常导数与它在 $W^{1,p}$ 意义下的
导数一致——所以记号是相容的! 特别地, 若 $I$ 有界, 则
$C^1(\bar I)\subset W^{1,p}(I)$ 对一切 $1\le p\le\infty$ 成立.
\end{remark}

\runin{例.} 设 $I=(-1,+1)$.作为练习证明下述事实:

\noindent (i) 函数 $u(x)=\abs{x}$ 对每个 $1\le p\le\infty$ 都属于 $W^{1,p}(I)$,
且 $u'=g$, 其中
\begin{equation*}
g(x)=
\begin{cases}
+1 & \text{若}\ 0<x<1,\\
-1 & \text{若}\ -1<x<0.
\end{cases}
\end{equation*}
更一般地, $\bar I$ 上连续且在 $\bar I$ 上分段 $C^1$ 的函数对一切
$1\le p\le\infty$ 都属于 $W^{1,p}(I)$.

\noindent (ii) 上述函数 $g$ 对任何 $1\le p\le\infty$ 都不属于 $W^{1,p}(I)$.

\begin{starremark}
\label{rem:8.3}
也可以用分布的语言来定义 $W^{1,p}$ (见 L.~Schwartz \bref{Schwartz}{1} 或
A.~Knapp \bref{Knapp}{2}).一切函数 $u\in L^p(I)$ 都有分布意义下的导数; 这个
导数是巨大的分布空间 $\mathcal{D}'(I)$ 中的元素.我们说 $u\in W^{1,p}$, 如果
这个分布导数恰好落在 $L^p$ 中, 而 $L^p$ 是 $\mathcal{D}'(I)$ 的子空间.当
$I=\R$ 且 $p=2$ 时, Sobolev 空间也可以用 Fourier 变换来定义; 参见例如
J.~L.~Lions--E.~Magenes \bref{LionsMagenes}{1}, P.~Malliavin
\bref{Malliavin}{1}, H.~Triebel \bref{Triebel}{1}, L.~Grafakos
\bref{Grafakos}{1}.这里不采取这个观点.
\end{starremark}

\begin{notation}
空间 $W^{1,p}$ 配以范数
\begin{equation*}
\norm{u}_{W^{1,p}}=\norm{u}_{L^p}+\norm{u'}_{L^p},
\end{equation*}
有时 (若 $1<p<\infty$) 也配以等价范数
$(\norm{u}_{L^p}^p+\norm{u'}_{L^p}^p)^{1/p}$.空间 $H^1$ 配以标量积
\begin{equation*}
(u,v)_{H^1}=(u,v)_{L^2}+(u',v')_{L^2}=\int_a^b(uv+u'v')
\end{equation*}
及相应的范数
\begin{equation*}
\norm{u}_{H^1}=(\norm{u}_{L^2}^2+\norm{u'}_{L^2}^2)^{1/2}.
\end{equation*}
\end{notation}

\begin{proposition}
\label{prop:8.1}
空间 $W^{1,p}$ 对 $1\le p\le\infty$ 是 Banach 空间.当 $1<p<\infty$ 时它自反,\footnote{这个性质是 $W^{1,p}$ 的一大优势.在变分法问题中, $W^{1,p}$ 优于
$C^1$, 因为后者不自反.极小元的存在性在自反空间中容易建立 (参见例如推论
 \ref{cor:3.23}).} 当 $1\le p<\infty$ 时可分.空间 $H^1$ 是可分 Hilbert 空间.
\end{proposition}

\begin{proof}
\begin{enumerate}[label=(\alph*), leftmargin=2.2em, itemsep=0.3em]
	\item 设 $(u_n)$ 为 $W^{1,p}$ 中的 Cauchy 序列; 则 $(u_n)$ 与 $(u_n')$ 是
	      $L^p$ 中的 Cauchy 序列.由此推出 $u_n$ 在 $L^p$ 中收敛到某个极限 $u$,
	      且 $u_n'$ 在 $L^p$ 中收敛到某个极限 $g$.我们有
	      \begin{equation*}
	      \int_I u_n\varphi'=-\int_I u_n'\varphi
	      \quad\forall\varphi\in C_c^1(I),
	      \end{equation*}
	      取极限得
	      \begin{equation*}
	      \int_I u\varphi'=-\int_I g\varphi
	      \quad\forall\varphi\in C_c^1(I).
	      \end{equation*}
	      于是 $u\in W^{1,p}$, $u'=g$, 且 $\norm{u_n-u}_{W^{1,p}}\to0$.
	\item 当 $1<p<\infty$ 时 $W^{1,p}$ 自反.显然, 乘积空间
	      $E=L^p(I)\times L^p(I)$ 自反.由 $Tu=[u,u']$ 定义的算子
	      $T:W^{1,p}\to E$ 是从 $W^{1,p}$ 到 $E$ 的等距映射.由于 $W^{1,p}$
	      是 Banach 空间, $T(W^{1,p})$ 是 $E$ 的闭子空间.由此推出
	      $T(W^{1,p})$ 自反 (见命题  \ref{prop:3.20}).因此 $W^{1,p}$ 也自反.
	\item 当 $1\le p<\infty$ 时 $W^{1,p}$ 可分.显然, 乘积空间
	      $E=L^p(I)\times L^p(I)$ 可分.于是 $T(W^{1,p})$ 也可分 (由命题
	       \ref{prop:3.25}).因此 $W^{1,p}$ 可分.
\end{enumerate}
\end{proof}

\begin{remark}
\label{rem:8.4}
记住下面这个事实是方便的, 我们在命题  \ref{prop:8.1} 的证明中已用过它: 设
$(u_n)$ 为 $W^{1,p}$ 中的序列, 使得 $u_n\to u$ 在 $L^p$ 中且 $(u_n')$ 在
$L^p$ 中收敛到某个极限; 则 $u\in W^{1,p}$ 且 $\norm{u_n-u}_{W^{1,p}}\to0$.
事实上, 当 $1<p\le\infty$ 时, 只须知道 $u_n\to u$ 在 $L^p$ 中且
$\norm{u_n'}_{L^p}$ 保持有界, 便可断言 $u\in W^{1,p}$ (见习题 8.2).
\end{remark}

粗略地说, $W^{1,p}$ 中的函数是 $L^p$ 函数的原函数.更精确地说, 我们有下述结果:

\begin{theorem}
\label{thm:8.2}
设 $u\in W^{1,p}(I)$, $1\le p\le\infty$, $I$ 有界或无界; 则存在函数
$\tilde u\in C(\bar I)$, 使得
\begin{equation*}
u=\tilde u\quad\text{在}\ I\ \text{上几乎处处成立},
\end{equation*}
且
\begin{equation*}
\tilde u(x)-\tilde u(y)=\int_y^x u'(t)\,dt\quad\forall x,y\in\bar I.
\end{equation*}
\end{theorem}

\begin{remark}
\label{rem:8.5}
让我们强调定理  \ref{thm:8.2} 的内容.首先注意, 若一个函数 $u$ 属于 $W^{1,p}$,
则一切满足 $v=u$ 在 $I$ 上几乎处处成立的函数 $v$ 也属于 $W^{1,p}$ (这直接由
$W^{1,p}$ 的定义得出).定理  \ref{thm:8.2} 断言, 每个 $u\in W^{1,p}$ 在 $\bar I$
上有一个 (且只有一个) 连续代表元, 即存在 $\bar I$ 上的连续函数属于 $u$ 的等价类
(若 $v=u$ 几乎处处成立, 则 $v\sim u$).在有用的时候,\footnote{例如, 为了给
每个 $x\in\bar I$ 处的 $u(x)$ 赋予意义.} 我们用它的连续代表元代替 $u$.为简化
记号, 我们也把连续代表元写成 $u$.最后我们指出, 性质 ``$u$ 有连续代表元'' 与
``$u$ 几乎处处连续'' 并不相同.
\end{remark}

\begin{remark}
\label{rem:8.6}
由定理  \ref{thm:8.2} 推出: 若 $u\in W^{1,p}$ 且 $u'\in C(\bar I)$ (即 $u'$ 在
$\bar I$ 上有连续代表元), 则 $u\in C^1(\bar I)$; 更精确地说,
$\tilde u\in C^1(\bar I)$, 但如上所述, 我们不再区分 $u$ 与 $\tilde u$.
\end{remark}

在定理  \ref{thm:8.2} 的证明中我们将用到下述引理:

\begin{lemma}
\label{lem:8.1}
设 $f\in L^1_{\mathrm{loc}}(I)$ 满足
\begin{equation}
\int_I f\varphi'=0\quad\forall\varphi\in C_c^1(I).
\label{eq:8.3}
\end{equation}
则存在常数 $C$, 使得 $f=C$ 在 $I$ 上几乎处处成立.
\end{lemma}

\begin{proof}
固定函数 $\psi\in C_c(I)$, 使得 $\int_I\psi=1$.对任何函数 $w\in C_c(I)$,
存在 $\varphi\in C_c^1(I)$ 使得
\begin{equation*}
\varphi'=w-\left(\int_I w\right)\psi.
\end{equation*}
事实上, 函数 $h=w-(\int_I w)\psi$ 连续, 在 $I$ 中具有紧支集, 且
$\int_I h=0$.因此 $h$ 在 $I$ 中有一个 (唯一的) 具紧支集的原函数.由
\eqref{eq:8.3} 推出
\begin{equation*}
\int_I f\left(w-\left(\int_I w\right)\psi\right)=0
\quad\forall w\in C_c(I),
\end{equation*}
即
\begin{equation*}
\int_I\left(f-\left(\int_I f\psi\right)\right)w=0
\quad\forall w\in C_c(I),
\end{equation*}
于是 (由推论  \ref{cor:4.24}) $f-(\int_I f\psi)=0$ 在 $I$ 上几乎处处成立,
即 $f=C$ 在 $I$ 上几乎处处成立, 其中 $C=\int_I f\psi$.
\end{proof}

\begin{lemma}
\label{lem:8.2}
设 $g\in L^1_{\mathrm{loc}}(I)$; 对固定的 $y_0\in I$, 令
\begin{equation*}
v(x)=\int_{y_0}^x g(t)\,dt,\quad x\in I.
\end{equation*}
则 $v\in C(I)$ 且
\begin{equation*}
\int_I v\varphi'=-\int_I g\varphi\quad\forall\varphi\in C_c^1(I).
\end{equation*}
\end{lemma}

\begin{proof}
我们有
\begin{equation*}
\begin{aligned}
\int_I v\varphi'
&=\int_I\left(\int_{y_0}^x g(t)\,dt\right)\varphi'(x)\,dx\\
&=-\int_a^{y_0}dx\int_x^{y_0}g(t)\varphi'(x)\,dt
+\int_{y_0}^b dx\int_{y_0}^x g(t)\varphi'(x)\,dt.
\end{aligned}
\end{equation*}
由 Fubini 定理,
\begin{equation*}
\begin{aligned}
\int_I v\varphi'
&=-\int_a^{y_0}g(t)\,dt\int_a^t\varphi'(x)\,dx
+\int_{y_0}^b g(t)\,dt\int_t^b\varphi'(x)\,dx\\
&=-\int_I g(t)\varphi(t)\,dt.
\end{aligned}
\end{equation*}
\end{proof}

\begin{proof}[定理 \ref{thm:8.2}]
固定 $y_0\in I$, 令 $\bar u(x)=\int_{y_0}^x u'(t)\,dt$.由引理  \ref{lem:8.2}
有
\begin{equation*}
\int_I\bar u\varphi'=-\int_I u'\varphi\quad\forall\varphi\in C_c^1(I).
\end{equation*}
于是 $\int_I(u-\bar u)\varphi'=0\ \forall\varphi\in C_c^1(I)$.由引理
 \ref{lem:8.1} 推出 $u-\bar u=C$ 在 $I$ 上几乎处处成立.函数
$\tilde u(x)=\bar u(x)+C$ 具有所需的性质.
\end{proof}

\begin{remark}
\label{rem:8.7}
引理  \ref{lem:8.2} 表明: 函数 $g\in L^p$ 的原函数 $v$ 属于 $W^{1,p}$, 只要
我们还知道 $v\in L^p$; 当 $I$ 有界时这总是成立的.
\end{remark}

\begin{proposition}
\label{prop:8.3}
设 $u\in L^p$, $1<p\le\infty$.下述性质等价:
\begin{enumerate}[label=(\roman*), leftmargin=2.6em, itemsep=0.2em]
	\item $u\in W^{1,p}$,
	\item 存在常数 $C$, 使得
	      \begin{equation*}
	      \abs{\int_I u\varphi'}\le C\norm{\varphi}_{L^{p'}(I)}
	      \quad\forall\varphi\in C_c^1(I).
	      \end{equation*}
\end{enumerate}
此外, 在 (ii) 中可以取 $C=\norm{u'}_{L^p(I)}$.
\end{proposition}

\begin{proof}
(i) $\Rightarrow$ (ii).这是显然的.

(ii) $\Rightarrow$ (i).线性泛函
\begin{equation*}
\varphi\in C_c^1(I)\mapsto\int_I u\varphi'
\end{equation*}
定义在 $L^{p'}$ 的一个稠密子空间上 (因为 $p'<\infty$), 且关于 $L^{p'}$ 范数连续.
因此它延拓为定义在整个 $L^{p'}$ 上的有界线性泛函 $F$ (应用 Hahn--Banach 定理,
或直接用连续性延拓).由 Riesz 表示定理 (定理  \ref{thm:4.11} 与
 \ref{thm:4.14}), 存在 $g\in L^p$ 使得
\begin{equation*}
\dual{F}{\varphi}=\int_I g\varphi\quad\forall\varphi\in L^{p'}.
\end{equation*}
特别地,
\begin{equation*}
\int_I u\varphi'=\int_I g\varphi\quad\forall\varphi\in C_c^1(I),
\end{equation*}
从而 $u\in W^{1,p}$.
\end{proof}

\begin{starremark}
\label{rem:8.8}[绝对连续函数与有界变差函数]
当 $p=1$ 时, 蕴涵 (i) $\Rightarrow$ (ii) 仍然成立, 但逆命题不成立.为说明这
一点, 设 $I$ 有界.满足 $p=1$ 的 (i) 的函数, 即 $W^{1,1}(I)$ 中的函数, 称为
\emph{绝对连续函数}.它们也由下述性质刻画:
\begin{equation*}
(AC)\qquad
\begin{cases}
\forall\varepsilon>0,\ \exists\delta>0\ \text{使得对每个由互不相交区间}
\ (a_k,b_k)\subset I\\
\text{组成的有限序列, 只要}\ \sum\abs{b_k-a_k}<\delta,\ \text{就有}\
\sum\abs{u(b_k)-u(a_k)}<\varepsilon.
\end{cases}
\end{equation*}
另一方面, 满足 $p=1$ 的 (ii) 的函数称为\emph{有界变差函数}; 这些函数可以用许多
不同的方式刻画:
\begin{enumerate}[label=(\alph*), leftmargin=2.2em, itemsep=0.15em]
	\item 它们是 $I$ 上两个有界单调不减函数 (可以间断) 之差,
	\item 它们是满足下述性质的函数 $u$:
	      \begin{equation*}
	      (BV)\qquad
	      \begin{cases}
	      \text{存在常数}\ C\ \text{使得}\\
	      \sum_{i=0}^{k-1}\abs{u(t_{i+1})-u(t_i)}\le C
	      \ \text{对}\ I\ \text{中一切}\ t_0<t_1<\cdots<t_k,
	      \end{cases}
	      \end{equation*}
	\item 它们是 $L^1(I)$ 中这样的函数 $u$: 其分布导数是有限测度.
\end{enumerate}
注意有界变差函数不必有连续代表元.关于这一课题, 参见例如 E.~Hewitt--
K.~Stromberg \bref{HewittStromberg}{1}, A.~Kolmogorov--S.~Fomin
\bref{KolmogorovFomin}{1}, S.~Chae \bref{Chae}{1}, H.~Royden
\bref{Royden}{1}, G.~Folland \bref{Folland}{2}, G.~Buttazzo--M.~Giaquinta--
S.~Hildebrandt \bref{ButtazzoGiaquintaHildebrandt}{1}, W.~Rudin
\bref{Rudin}{2}, R.~Wheeden--A.~Zygmund \bref{WheedenZygmund}{1} 以及
A.~Knapp \bref{Knapp}{1}.
\end{starremark}
\begin{proposition}
\label{prop:8.4}
$L^\infty(I)$ 中的函数 $u$ 属于 $W^{1,\infty}(I)$ 当且仅当存在常数 $C$, 使得
\begin{equation*}
\abs{u(x)-u(y)}\le C\abs{x-y}\quad\text{对}\ I\ \text{中几乎所有的}\ x,y.
\end{equation*}
\end{proposition}

\begin{proof}
若 $u\in W^{1,\infty}(I)$, 我们可以应用定理  \ref{thm:8.2} 推出
\begin{equation*}
\abs{u(x)-u(y)}\le\norm{u'}_{L^\infty}\abs{x-y}
\quad\text{对}\ I\ \text{中几乎所有的}\ x,y.
\end{equation*}
反之, 设 $\varphi\in C_c^1(I)$.对 $h\in\R$, 当 $\abs{h}$ 足够小时, 我们有
\begin{equation*}
\int_I[u(x+h)-u(x)]\varphi(x)\,dx=\int_I u(x)[\varphi(x-h)-\varphi(x)]\,dx
\end{equation*}
(当 $h$ 小时这些积分有意义, 因为 $\varphi$ 的支集在 $I$ 的某个紧子集中).利用
关于 $u$ 的假设, 得到
\begin{equation*}
\abs{\int_I u(x)[\varphi(x-h)-\varphi(x)]\,dx}\le C\abs{h}\norm{\varphi}_{L^1}.
\end{equation*}
除以 $\abs{h}$ 并令 $h\to0$, 我们得到
\begin{equation*}
\abs{\int_I u\varphi'}\le C\norm{\varphi}_{L^1}
\quad\forall\varphi\in C_c^1(I).
\end{equation*}
现在可以应用命题  \ref{prop:8.3}, 断言 $u\in W^{1,\infty}$.
\end{proof}

命题  \ref{prop:8.4} 的 $L^p$ 版本如下:

\begin{proposition}
\label{prop:8.5}
设 $u\in L^p(\R)$, $1<p<\infty$.下述性质等价:
\begin{enumerate}[label=(\roman*), leftmargin=2.6em, itemsep=0.2em]
	\item $u\in W^{1,p}(\R)$,
	\item 存在常数 $C$, 使得对一切 $h\in\R$,
	      \begin{equation*}
	      \norm{\tau_hu-u}_{L^p(\R)}\le C\abs{h}.
	      \end{equation*}
\end{enumerate}
此外, 在 (ii) 中可以取 $C=\norm{u'}_{L^p(\R)}$.

回忆 $(\tau_hu)(x)=u(x+h)$.
\end{proposition}

\begin{proof}
(i) $\Rightarrow$ (ii).(当 $p=1$ 时这个蕴涵也成立.) 由定理  \ref{thm:8.2},
对 $\R$ 中一切 $x$ 与 $h$,
\begin{equation*}
u(x+h)-u(x)=\int_x^{x+h}u'(t)\,dt=h\int_0^1u'(x+sh)\,ds.
\end{equation*}
于是
\begin{equation*}
\abs{u(x+h)-u(x)}\le\abs{h}\int_0^1\abs{u'(x+sh)}\,ds.
\end{equation*}
应用 Hölder 不等式, 得到
\begin{equation*}
\abs{u(x+h)-u(x)}^p\le\abs{h}^p\int_0^1\abs{u'(x+sh)}^p\,ds.
\end{equation*}
由此推出
\begin{equation*}
\begin{aligned}
\int_\R\abs{u(x+h)-u(x)}^p\,dx
&\le\abs{h}^p\int_\R dx\int_0^1\abs{u'(x+sh)}^p\,ds\\
&\le\abs{h}^p\int_0^1ds\int_\R\abs{u'(x+sh)}^p\,dx.
\end{aligned}
\end{equation*}
但当 $0<s<1$ 时,
\begin{equation*}
\int_\R\abs{u'(x+sh)}^p\,dx=\int_\R\abs{u'(y)}^p\,dy,
\end{equation*}
由此即可推出 (ii).

(ii) $\Rightarrow$ (i).设 $\varphi\in C_c^1(\R)$.对一切 $h\in\R$ 有
\begin{equation*}
\int_\R[u(x+h)-u(x)]\varphi(x)\,dx=\int_\R u(x)[\varphi(x-h)-\varphi(x)]\,dx.
\end{equation*}
利用 Hölder 不等式与 (ii), 得到
\begin{equation*}
\abs{\int_\R[u(x+h)-u(x)]\varphi(x)\,dx}\le C\abs{h}\norm{\varphi}_{L^{p'}(\R)},
\end{equation*}
从而
\begin{equation*}
\abs{\int_\R u(x)[\varphi(x-h)-\varphi(x)]\,dx}\le C\abs{h}\norm{\varphi}_{L^{p'}(\R)}.
\end{equation*}
除以 $\abs{h}$ 并令 $h\to0$, 我们得到
\begin{equation*}
\abs{\int_\R u\varphi'}\le C\norm{\varphi}_{L^{p'}(\R)}.
\end{equation*}
我们再次应用命题  \ref{prop:8.3}, 断言 $u\in W^{1,p}(\R)$.
\end{proof}

某些基本的分析运算只对定义在整个 $\R$ 上的函数才有意义 (例如卷积与 Fourier
变换).因此, 能够把一个函数 $u\in W^{1,p}(I)$ 延拓成函数
$\bar u\in W^{1,p}(\R)$ 是有用的.\footnote{若把 $u$ 在 $I$ 外延拓为 $0$, 则所得
函数一般不属于 $W^{1,p}(\R)$ (见注记 \ref{rem:8.5} 与第~\ref{sec:8.3}~节).} 下面的结果
处理这一点.

\begin{theorem}[延拓算子]
\label{thm:8.6}
设 $1\le p\le\infty$.存在有界线性算子 $P:W^{1,p}(I)\to W^{1,p}(\R)$, 称为
\emph{延拓算子}, 它满足下述性质:
\begin{enumerate}[label=(\roman*), leftmargin=2.6em, itemsep=0.2em]
	\item $Pu|_I=u\quad\forall u\in W^{1,p}(I)$,
	\item $\norm{Pu}_{L^p(\R)}\le C\norm{u}_{L^p(I)}\quad\forall u\in W^{1,p}(I)$,
	\item $\norm{Pu}_{W^{1,p}(\R)}\le C\norm{u}_{W^{1,p}(I)}\quad\forall u\in W^{1,p}(I)$,
\end{enumerate}
其中 $C$ 只依赖于 $\abs{I}\le\infty$.\footnote{在 (ii) 中可取 $C=4$, 在 (iii) 中
可取 $C=4(1+1/\abs{I})$.}
\end{theorem}

\begin{proof}
从 $I=(0,\infty)$ 的情形开始, 我们证明\emph{反射延拓}
\begin{equation*}
(Pu)(x)=u^\star(x)=
\begin{cases}
u(x) & \text{若}\ x\ge0,\\
u(-x) & \text{若}\ x<0
\end{cases}
\end{equation*}
可用.显然有
\begin{equation*}
\norm{u^\star}_{L^p(\R)}\le2\norm{u}_{L^p(I)}.
\end{equation*}
令
\begin{equation*}
v(x)=
\begin{cases}
u'(x) & \text{若}\ x>0,\\
-u'(-x) & \text{若}\ x<0,
\end{cases}
\end{equation*}
容易验证 $v\in L^p(\R)$ 且
\begin{equation*}
u^\star(x)-u^\star(0)=\int_0^x v(t)\,dt\quad\forall x\in\R.
\end{equation*}
由此推出 $u^\star\in W^{1,p}(\R)$ (见注记 \ref{rem:8.7}) 且
$\norm{u^\star}_{W^{1,p}(\R)}\le2\norm{u}_{W^{1,p}(I)}$.

\begin{figure}[tbp]
	\centering
	\begin{tikzpicture}[scale=1.0]
		\draw[-{Stealth[length=2mm]}] (-0.35,0) -- (7.6,0) node[right] {$x$};
		\draw[-{Stealth[length=2mm]}] (0,-0.3) -- (0,2.9) node[above] {$\eta$};
		\draw[thick] (0,2.35) -- (2.7,2.35)
			.. controls (3.45,2.35) and (3.4,0.35) .. (4.15,0.35)
			-- (7.2,0.35);
		\foreach \x/\lab in {0/0, 2.2/{$\frac14$}, 3.55/{$\frac12$}, 4.9/{$\frac34$}, 6.2/{$1$}}
			\draw (\x,0) -- (\x,-0.09) node[below] {\small \lab};
		\node[left] at (-0.08,2.35) {\small $1$};
	\end{tikzpicture}
	\caption{}
	\label{fig:5}
\end{figure}
现在考虑 $I$ 有界的情形; 不失一般性, 可设 $I=(0,1)$.固定函数
$\eta\in C^1(\R)$, $0\le\eta\le1$, 使得
\begin{equation*}
\eta(x)=
\begin{cases}
1 & \text{若}\ x<1/4,\\
0 & \text{若}\ x>3/4.
\end{cases}
\end{equation*}
见图  \ref{fig:5}.

给定 $(0,1)$ 上的函数 $f$, 令
\begin{equation*}
\tilde f(x)=
\begin{cases}
f(x) & \text{若}\ 0<x<1,\\
0 & \text{若}\ x>1.
\end{cases}
\end{equation*}
我们将需要下述引理.
\end{proof}


\begin{lemma}
\label{lem:8.3}
设 $u\in W^{1,p}(I)$.则
\begin{equation*}
\eta\tilde u\in W^{1,p}(0,\infty)
\quad\text{且}\quad
(\eta\tilde u)'=\eta'\tilde u+\eta u'.
\end{equation*}
\end{lemma}

\begin{proof}
设 $\varphi\in C_c^1((0,\infty))$; 则
\begin{equation*}
\begin{aligned}
\int_0^\infty\eta\tilde u\varphi'
&=\int_0^1\eta u\varphi'=\int_0^1u[(\eta\varphi)'-\eta'\varphi]\\
&=-\int_0^1u'\eta\varphi-\int_0^1u\eta'\varphi
\qquad\text{因为}\ \eta\varphi\in C_c^1((0,1))\\
&=-\int_0^\infty(u'\eta+\tilde u\eta')\varphi.
\end{aligned}
\end{equation*}
\end{proof}

\begin{proof}[定理 \ref{thm:8.6}, 续]
给定 $u\in W^{1,p}(I)$, 写
\begin{equation*}
u=\eta u+(1-\eta)u.
\end{equation*}
函数 $\eta u$ 先由 $\eta\tilde u$ 延拓到 $(0,\infty)$ (由引理  \ref{lem:8.3}),
再用反射延拓到 $\R$.这样我们得到函数 $v_1\in W^{1,p}(\R)$, 它延拓 $\eta u$,
且
\begin{equation*}
\norm{v_1}_{L^p(\R)}\le2\norm{u}_{L^p(I)},\qquad
\norm{v_1}_{W^{1,p}(\R)}\le C\norm{u}_{W^{1,p}(I)}
\end{equation*}
(其中 $C$ 依赖于 $\norm{\eta'}_{L^\infty}$).

对 $(1-\eta)u$ 同样处理, 即先把 $(1-\eta)u$ 在 $(-\infty,0)$ 上延拓为 $0$
从而延拓到 $(-\infty,1)$, 再用反射延拓到 $\R$ (这次关于点 $1$ 而非点 $0$
反射).这样我们得到函数 $v_2\in W^{1,p}(\R)$, 它延拓 $(1-\eta)u$ 且满足
\begin{equation*}
\norm{v_2}_{L^p(\R)}\le2\norm{u}_{L^p(I)},\qquad
\norm{v_2}_{W^{1,p}(\R)}\le C\norm{u}_{W^{1,p}(I)}.
\end{equation*}
于是 $Pu=v_1+v_2$ 满足定理的条件.
\end{proof}

\index{Poincarbudengshi@Poincaré 不等式 (Poincare inequality)}
$C^1$ 函数的某些性质对 $W^{1,p}$ 函数仍然成立 (参见例如推论  \ref{cor:8.10} 与
 \ref{cor:8.11}).借助基于下述结果的稠密性论证来建立这些性质是方便的.

\begin{keytheorem}[稠密性]
\label{thm:8.7}
设 $u\in W^{1,p}(I)$, $1\le p<\infty$.则存在 $C_c^\infty(\R)$ 中的序列
$(u_n)$, 使得 $u_n|_I\to u$ 在 $W^{1,p}(I)$ 中.
\end{keytheorem}

\begin{remark}
\label{rem:8.9}
一般说来, 不存在 $C_c^\infty(I)$ 中的序列 $(u_n)$ 使得 $u_n\to u$ 在
$W^{1,p}(I)$ 中 (见第~\ref{sec:8.3}~节).这与 $L^p$ 空间形成对照: 回忆对每个函数
$u\in L^p(I)$, 存在 $C_c^\infty(I)$ 中的序列 $(u_n)$ 使得 $u_n\to u$ 在
$L^p(I)$ 中 (见推论  \ref{cor:4.23}).
\end{remark}

\begin{proof}
我们总可以假设 $I=\R$; 否则用定理  \ref{thm:8.6} 把 $u$ 延拓为
$W^{1,p}(\R)$ 中的函数.我们使用卷积 (使函数成为 $C^\infty$) 与截断 (使其支集
紧) 的基本技巧.

\runin{(a) 卷积.} 我们将需要下述引理.

\begin{lemma}
\label{lem:8.4}
设 $\rho\in L^1(\R)$, $v\in W^{1,p}(\R)$, $1\le p\le\infty$.则
$\rho\star v\in W^{1,p}(\R)$ 且 $(\rho\star v)'=\rho\star v'$.
\end{lemma}

\begin{proof}
首先设 $\rho$ 有紧支集.我们已经知道 (定理  \ref{thm:4.15})
$\rho\star v\in L^p(\R)$.设 $\varphi\in C_c^1(\R)$; 由命题  \ref{prop:4.16}
与  \ref{prop:4.20} 有
\begin{equation*}
\int(\rho\star v)\varphi'=\int v(\check\rho\star\varphi')
=\int v(\check\rho\star\varphi)'=-\int v'(\check\rho\star\varphi)
=-\int(\rho\star v')\varphi,
\end{equation*}
由此推出
\begin{equation*}
\rho\star v\in W^{1,p}\quad\text{且}\quad(\rho\star v)'=\rho\star v'.
\end{equation*}
若 $\rho$ 没有紧支集, 引入 $C_c(\R)$ 中的序列 $(\rho_n)$ 使得
$\rho_n\to\rho$ 在 $L^1(\R)$ 中 (见推论  \ref{cor:4.23}).由上述结果, 得到
\begin{equation*}
\rho_n\star v\in W^{1,p}(\R)\quad\text{且}\quad
(\rho_n\star v)'=\rho_n\star v'.
\end{equation*}
但 $\rho_n\star v\to\rho\star v$ 在 $L^p(\R)$ 中且
$\rho_n\star v'\to\rho\star v'$ 在 $L^p(\R)$ 中 (由定理  \ref{thm:4.15}).
借助注记 \ref{rem:8.4} 我们断言
\begin{equation*}
\rho\star v\in W^{1,p}(\R)\quad\text{且}\quad(\rho\star v)'=\rho\star v'.
\end{equation*}
\end{proof}

\runin{(b) 截断.} 固定函数 $\zeta\in C_c^\infty(\R)$, 使得 $0\le\zeta\le1$ 且
\begin{equation*}
\zeta(x)=
\begin{cases}
1 & \text{若}\ \abs{x}<1,\\
0 & \text{若}\ \abs{x}\ge2.
\end{cases}
\end{equation*}
定义序列
\begin{equation}
\zeta_n(x)=\zeta(x/n)\quad n=1,2,\dots.
\label{eq:8.4}
\end{equation}
由控制收敛定理容易推出: 若函数 $f$ 属于 $L^p(\R)$, $1\le p<\infty$, 则
$\zeta_nf\to f$ 在 $L^p(\R)$ 中.

\runin{(c) 结论.} 选取一列磨光核 $(\rho_n)$.我们断言序列
$u_n=\zeta_n(\rho_n\star u)$ 在 $W^{1,p}(\R)$ 中收敛到 $u$.首先,
$\norm{u_n-u}_p\to0$.事实上, 写
\begin{equation*}
u_n-u=\zeta_n((\rho_n\star u)-u)+(\zeta_nu-u),
\end{equation*}
从而
\begin{equation*}
\norm{u_n-u}_p\le\norm{\rho_n\star u-u}_p+\norm{\zeta_nu-u}_p\to0.
\end{equation*}
其次, 由引理  \ref{lem:8.4} 有
\begin{equation*}
u_n'=\zeta_n'(\rho_n\star u)+\zeta_n(\rho_n\star u').
\end{equation*}
因此
\begin{equation*}
\begin{aligned}
\norm{u_n'-u'}_p
&\le\norm{\zeta_n'(\rho_n\star u)}_p+\norm{\zeta_n(\rho_n\star u')-u'}_p\\
&\le\frac{C}{n}\norm{u}_p+\norm{\rho_n\star u'-u'}_p+\norm{\zeta_nu'-u'}_p\to0,
\end{aligned}
\end{equation*}
其中 $C=\norm{\zeta'}_\infty$.
\end{proof}

下一个结果是 Sobolev 不等式 (也称 \emph{Sobolev 嵌入}) 的一个重要原型.

\begin{keytheorem}
\index{Sobolevqianru@Sobolev 嵌入 (Sobolev embedding)}
\label{thm:8.8}
存在常数 $C$ (只依赖于 $\abs{I}\le\infty$), 使得
\begin{equation}
\norm{u}_{L^\infty(I)}\le C\norm{u}_{W^{1,p}(I)}
\quad\forall u\in W^{1,p}(I),\ \forall 1\le p\le\infty.
\label{eq:8.5}
\end{equation}
换言之, $W^{1,p}(I)\subset L^\infty(I)$ 且嵌入连续, 对一切
$1\le p\le\infty$ 成立.此外, 若 $I$ 有界, 则
\begin{equation}
\text{嵌入}\ W^{1,p}(I)\subset C(\bar I)\ \text{对一切}\ 1<p\le\infty\ \text{是紧的},
\label{eq:8.6}
\end{equation}
\begin{equation}
\text{嵌入}\ W^{1,1}(I)\subset L^q(I)\ \text{对一切}\ 1\le q<\infty\ \text{是紧的}.
\label{eq:8.7}
\end{equation}
\end{keytheorem}

\begin{proof}
我们首先对 $I=\R$ 证明 \eqref{eq:8.5}; 一般情形则由延拓定理 (定理
 \ref{thm:8.6}) 推出.设 $v\in C_c^1(\R)$; 若 $1\le p<\infty$, 令
$G(s)=\abs{s}^{p-1}s$.函数 $w=G(v)$ 属于 $C_c^1(\R)$, 且
\begin{equation*}
w'=G'(v)v'=p\abs{v}^{p-1}v'.
\end{equation*}
于是对 $x\in\R$ 有
\begin{equation*}
G(v(x))=\int_{-\infty}^x p\abs{v(t)}^{p-1}v'(t)\,dt,
\end{equation*}
由 Hölder 不等式
\begin{equation*}
\abs{v(x)}^p\le p\norm{v}_p^{p-1}\norm{v'}_p,
\end{equation*}
由此断言
\begin{equation}
\norm{v}_\infty\le C\norm{v}_{W^{1,p}}\quad\forall v\in C_c^1(\R),
\label{eq:8.8}
\end{equation}
其中 $C$ 是普适常数 (与 $p$ 无关).\footnote{注意
$p^{1/p}\le e^{1/e}\ \forall p\ge1$.}

现在用稠密性论证.设 $u\in W^{1,p}(\R)$; 存在序列
$(u_n)\subset C_c^1(\R)$ 使得 $u_n\to u$ 在 $W^{1,p}(\R)$ 中 (由定理
 \ref{thm:8.7}).应用 \eqref{eq:8.8}, 可见 $(u_n)$ 是 $L^\infty(\R)$ 中的
Cauchy 序列.于是 $u_n\to u$ 在 $L^\infty(\R)$ 中, 我们得到 \eqref{eq:8.5}.

\emph{\eqref{eq:8.6} 的证明.} 设 $H$ 为 $W^{1,p}(I)$ 中的单位球,
$1<p\le\infty$.对 $u\in H$ 有
\begin{equation*}
\abs{u(x)-u(y)}=\abs{\int_y^x u'(t)\,dt}
\le\norm{u'}_p\abs{x-y}^{1/p'}\le\abs{x-y}^{1/p'}
\quad\forall x,y\in I.
\end{equation*}
于是由 Ascoli--Arzelà 定理 (定理  \ref{thm:4.25}) 知 $H$ 在 $C(\bar I)$ 中的
闭包紧.

\emph{\eqref{eq:8.7} 的证明.} 设 $H$ 为 $W^{1,1}(I)$ 中的单位球.设 $P$ 为
定理  \ref{thm:8.6} 中的延拓算子, 令 $F=P(H)$, 于是 $H=F|_I$.我们通过应用
定理  \ref{thm:4.26} 来证明 $H$ 在 $L^q(I)$ 中的闭包紧 (对一切
$1\le q<\infty$).显然, $F$ 在 $W^{1,1}(\R)$ 中有界; 因此 $F$ 在 $L^q(\R)$
中也有界, 因为它在 $L^1(\R)$ 与 $L^\infty(\R)$ 中都有界.现在验证第~
\ref{chap:4}~章的条件
\eqref{eq:4.22}, 即
\begin{equation*}
\lim_{h\to0}\norm{\tau_hf-f}_q=0\quad\text{对}\ f\in F\ \text{一致}.
\end{equation*}
由命题  \ref{prop:8.5}, 对每个 $f\in F$,
\begin{equation*}
\norm{\tau_hf-f}_{L^1(\R)}\le\abs{h}\norm{f'}_{L^1(\R)}\le C\abs{h},
\end{equation*}
因为 $F$ 是 $W^{1,1}(\R)$ 的有界子集.因此
\begin{equation*}
\norm{\tau_hf-f}_{L^q(\R)}^q\le(2\norm{f}_{L^\infty(\R)})^{q-1}
\norm{\tau_hf-f}_{L^1(\R)}\le C\abs{h},
\end{equation*}
从而
\begin{equation*}
\norm{\tau_hf-f}_{L^q(\R)}\le C\abs{h}^{1/q},
\end{equation*}
其中 $C$ 与 $f$ 无关.由于 $q\ne\infty$, 所需结论随之成立.
\end{proof}

\begin{remark}
\label{rem:8.10}
嵌入 $W^{1,1}(I)\subset C(\bar I)$ 连续, 但从不紧, 即使 $I$ 是有界区间;
读者应自己给出论证或参见习题 8.2.然而, 若 $(u_n)$ 是 $W^{1,1}(I)$ 中的有界
序列 ($I$ 有界或无界), 则存在子列 $(u_{n_k})$ 使得 $u_{n_k}(x)$ 对一切 $x\in I$
收敛 (这是 Helly 选择定理; 参见例如 A.~Kolmogorov--S.~Fomin
\bref{KolmogorovFomin}{1} 与习题 8.3).当 $I$ 无界且 $1<p\le\infty$ 时, 我们
知道嵌入 $W^{1,p}(I)\subset L^\infty(I)$ 连续; 这个嵌入从不紧——同样, 请给出
论证或参见习题 8.4.然而, 若 $(u_n)$ 在 $W^{1,p}(I)$ 中有界,
$1<p\le\infty$, 则存在子列 $(u_{n_k})$ 与某个 $u\in W^{1,p}(I)$, 使得
$u_{n_k}\to u$ 在 $L^\infty(J)$ 中, 对 $I$ 的每个有界子集 $J$.
\end{remark}

\begin{remark}
\label{rem:8.11}
设 $I$ 为有界区间, $1\le p\le\infty$, $1\le q\le\infty$.由定理
 \ref{thm:8.2} 与 \eqref{eq:8.5} 容易证明, 范数
\begin{equation*}
|||u|||=\norm{u'}_p+\norm{u}_q
\end{equation*}
与 $W^{1,p}(I)$ 的范数等价.
\end{remark}

\begin{remark}
\label{rem:8.12}
设 $I$ 为无界区间.若 $u\in W^{1,p}(I)$, 则 $u\in L^q(I)$ 对一切
$q\in[p,\infty]$, 因为
\begin{equation*}
\int_I\abs{u}^q\le\norm{u}_\infty^{q-p}\norm{u}_p^p.
\end{equation*}
但一般说来 $u\notin L^q(I)$ 对 $q\in[1,p)$ (见习题 8.1).
\end{remark}

\begin{corollary}
\label{cor:8.9}
设 $I$ 为无界区间, $u\in W^{1,p}(I)$, $1\le p<\infty$.则
\begin{equation}
\lim_{\substack{x\in I\\ \abs{x}\to\infty}}u(x)=0.
\label{eq:8.9}
\end{equation}
\end{corollary}

\begin{proof}
由定理  \ref{thm:8.7}, 存在 $C_c^1(\R)$ 中的序列 $(u_n)$ 使得
$u_n|_I\to u$ 在 $W^{1,p}(I)$ 中.由 \eqref{eq:8.5} 推出
$\norm{u_n-u}_{L^\infty(I)}\to0$.由此得到 \eqref{eq:8.9}.事实上, 给定
$\varepsilon>0$, 选取 $n$ 足够大使得 $\norm{u_n-u}_{L^\infty(I)}<\varepsilon$.
当 $\abs{x}$ 足够大时 $u_n(x)=0$ (因为 $u_n\in C_c^1(\R)$), 从而
$\abs{u(x)}<\varepsilon$.
\end{proof}

\begin{corollary}[乘积的求导]
\label{cor:8.10}
设 $u,v\in W^{1,p}(I)$, $1\le p\le\infty$.\footnote{注意这个结果与
$L^p$ 函数的性质形成对照: 一般说来, 若 $u,v\in L^p$, 则乘积 $uv$ 不属于
$L^p$.我们说 $W^{1,p}(I)$ 是一个 Banach 代数.} 则
\begin{equation}
uv\in W^{1,p}(I)\quad\text{且}\quad
(uv)'=u'v+uv'.
\label{eq:8.10}
\end{equation}
此外, 分部积分公式成立:
\begin{equation}
\int_y^x u'v=u(x)v(x)-u(y)v(y)-\int_y^x uv'\quad\forall x,y\in\bar I.
\label{eq:8.11}
\end{equation}
\end{corollary}

\begin{proof}
首先回忆 $u\in L^\infty$ (由定理  \ref{thm:8.8}), 从而 $uv\in L^p$.为证明
$(uv)'\in L^p$, 我们从 $1\le p<\infty$ 的情形开始.设 $(u_n)$ 与 $(v_n)$ 为
$C_c^1(\R)$ 中的序列, 使得 $u_n|_I\to u$ 且 $v_n|_I\to v$ 在 $W^{1,p}(I)$
中.于是 $u_n|_I\to u$ 且 $v_n|_I\to v$ 在 $L^\infty(I)$ 中 (再次由定理
 \ref{thm:8.8}).由此推出 $u_nv_n|_I\to uv$ 在 $L^\infty(I)$ 中, 也在 $L^p(I)$
中.我们有
\begin{equation*}
(u_nv_n)'=u_n'v_n+u_nv_n'\to u'v+uv'\quad\text{在}\ L^p(I)\ \text{中}.
\end{equation*}
再次对序列 $(u_nv_n)$ 应用注记 \ref{rem:8.4}, 我们断言 $uv\in W^{1,p}(I)$
且 \eqref{eq:8.10} 成立.对 \eqref{eq:8.10} 积分即得 \eqref{eq:8.11}.

现在转向 $p=\infty$ 的情形; 设 $u,v\in W^{1,\infty}(I)$.于是
$uv\in L^\infty(I)$ 且 $u'v+uv'\in L^\infty(I)$.只须验证
\begin{equation*}
\int_I uv\varphi'=-\int_I(u'v+uv')\varphi\quad\forall\varphi\in C_c^1(I).
\end{equation*}
为此, 固定有界开区间 $J\subset I$ 使得 $\operatorname{supp}\varphi\subset J$.
于是 $u,v\in W^{1,p}(J)$ 对一切 $p<\infty$, 且由上述结果知
\begin{equation*}
\int_J uv\varphi'=-\int_J(u'v+uv')\varphi,
\end{equation*}
即
\begin{equation*}
\int_I uv\varphi'=-\int_I(u'v+uv')\varphi.
\end{equation*}
\end{proof}

\begin{corollary}[复合的求导]
\label{cor:8.11}
设 $G\in C^1(\R)$ 满足\footnote{当 $I$ 有界时 (或当 $I$ 无界但 $p=\infty$ 时)
这一限制是不必要的.当 $I$ 无界且 $1\le p<\infty$ 时它是本质的.}
$G(0)=0$, 并设 $u\in W^{1,p}(I)$, $1\le p\le\infty$.则
\begin{equation*}
G\circ u\in W^{1,p}(I)\quad\text{且}\quad(G\circ u)'=(G'\circ u)u'.
\end{equation*}
\end{corollary}

\begin{proof}
令 $M=\norm{u}_\infty$.由于 $G(0)=0$, 存在常数 $C$ 使得
$\abs{G(s)}\le C\abs{s}$ 对一切 $s\in[-M,+M]$.于是 $\abs{G\circ u}\le C\abs{u}$;
由此推出 $G\circ u\in L^p(I)$.类似地, $(G'\circ u)u'\in L^p(I)$.只须验证
\begin{equation}
\int_I(G\circ u)\varphi'=-\int_I(G'\circ u)u'\varphi
\quad\forall\varphi\in C_c^1(I).
\label{eq:8.12}
\end{equation}
首先设 $1\le p<\infty$.则存在 $C_c^1(\R)$ 中的序列 $(u_n)$ 使得
$u_n|_I\to u$ 在 $W^{1,p}(I)$ 中, 也在 $L^\infty(I)$ 中.于是
$(G\circ u_n)|_I\to G\circ u$ 在 $L^\infty(I)$ 中, 且
$(G'\circ u_n)u_n'|_I\to(G'\circ u)u'$ 在 $L^p(I)$ 中.显然 (由 $C^1$ 函数的
标准法则) 有
\begin{equation*}
\int_I(G\circ u_n)\varphi'=-\int_I(G'\circ u_n)u_n'\varphi
\quad\forall\varphi\in C_c^1(I),
\end{equation*}
由此推出 \eqref{eq:8.12}.对 $p=\infty$ 的情形, 按推论  \ref{cor:8.10} 的证明
中同样的方式进行.
\end{proof}

\subsection*{Sobolev 空间 $W^{m,p}$}

\begin{definition}
给定整数 $m\ge2$ 与实数 $1\le p\le\infty$, 我们用归纳法定义空间
\begin{equation*}
W^{m,p}(I)=\{u\in W^{m-1,p}(I);\ u'\in W^{m-1,p}(I)\}.
\end{equation*}
我们也记
\begin{equation*}
H^m(I)=W^{m,2}(I).
\end{equation*}
\end{definition}

容易证明: $u\in W^{m,p}(I)$ 当且仅当存在 $m$ 个函数
$g_1,g_2,\dots,g_m\in L^p(I)$, 使得
\begin{equation*}
\int_I uD^j\varphi=(-1)^j\int_I g_j\varphi
\quad\forall\varphi\in C_c^\infty(I),\ \forall j=1,2,\dots,m,
\end{equation*}
其中 $D^j\varphi$ 表示 $\varphi$ 的第 $j$ 阶导数.当 $u\in W^{m,p}(I)$ 时,
我们于是可以考虑 $u$ 的逐次导数: $u'=g_1$, $(u')'=g_2$, $\dots$, 直到 $m$
阶.它们记作 $Du,D^2u,\dots,D^mu$.空间 $W^{m,p}(I)$ 配以范数
\begin{equation*}
\norm{u}_{W^{m,p}}=\norm{u}_p+\sum_{\alpha=1}^m\norm{D^\alpha u}_p,
\end{equation*}
空间 $H^m(I)$ 配以标量积
\begin{equation*}
(u,v)_{H^m}=(u,v)_{L^2}+\sum_{\alpha=1}^m(D^\alpha u,D^\alpha v)_{L^2}
=\int_I uv+\sum_{\alpha=1}^m\int_I D^\alpha u\,D^\alpha v.
\end{equation*}

可以证明范数 $\norm{\cdot}_{W^{m,p}}$ 与范数
\begin{equation*}
|||u|||=\norm{u}_p+\norm{D^mu}_p
\end{equation*}
等价.更精确地说, 可以证明: 对每个整数 $j$, $1\le j\le m-1$, 以及对每个
$\varepsilon>0$, 存在常数 $C$ (依赖于 $\varepsilon$ 与 $\abs{I}\le\infty$),
使得
\begin{equation*}
\norm{D^ju}_p\le\varepsilon\norm{D^mu}_p+C\norm{u}_p
\quad\forall u\in W^{m,p}(I)
\end{equation*}
(参见例如 R.~Adams \bref{Adams}{1}, 或习题 8.6 中 $\abs{I}<\infty$ 的情形).

读者可以把对 $W^{1,p}$ 证明的一切性质推广到空间 $W^{m,p}$; 例如, 若 $I$ 有界,
则 $W^{m,p}(I)\subset C^{m-1}(\bar I)$, 嵌入连续 (相应地, 当
$1<p\le\infty$ 时嵌入紧).
\section{空间 $W_0^{1,p}$}
\label{sec:8.3}

\begin{definition}
给定 $1\le p<\infty$, 用 $W_0^{1,p}(I)$ 表示 $C_c^1(I)$ 在 $W^{1,p}(I)$ 中的
闭包.\footnote{我们不对 $p=\infty$ 定义 $W_0^{1,p}$.} 令
\begin{equation*}
H_0^1(I)=W_0^{1,2}(I).
\end{equation*}
空间 $W_0^{1,p}(I)$ 配以 $W^{1,p}(I)$ 的范数, 空间 $H_0^1$ 配以 $H^1$ 的
标量积.\footnote{若无混淆, 我们常写 $W_0^{1,p}$ 与 $H_0^1$ 而不写
$W_0^{1,p}(I)$ 与 $H_0^1(I)$.}
\end{definition}

空间 $W_0^{1,p}$ 是可分 Banach 空间.此外, 当 $p>1$ 时它自反.空间 $H_0^1$
是可分 Hilbert 空间.

\begin{remark}
\label{rem:8.13}
当 $I=\R$ 时, 我们知道 $C_c^1(\R)$ 在 $W^{1,p}(\R)$ 中稠密 (见定理
 \ref{thm:8.7}), 因此 $W_0^{1,p}(\R)=W^{1,p}(\R)$.
\end{remark}

\begin{remark}
\label{rem:8.14}
利用磨光核序列 $(\rho_n)$, 容易验证下述事实:
\begin{enumerate}[label=(\roman*), leftmargin=2.6em, itemsep=0.2em]
	\item $C_c^\infty(I)$ 在 $W_0^{1,p}(I)$ 中稠密.
	\item 若 $u\in W^{1,p}(I)\cap C_c(I)$, 则 $u\in W_0^{1,p}(I)$.
\end{enumerate}
\end{remark}

下一个结果给出 $W_0^{1,p}(I)$ 中函数的基本刻画.

\begin{keytheorem}
\label{thm:8.12}
设 $u\in W^{1,p}(I)$.则 $u\in W_0^{1,p}(I)$ 当且仅当 $u=0$ 在 $\partial I$ 上.
\end{keytheorem}

\begin{remark}
\label{rem:8.15}
定理  \ref{thm:8.12} 说明了空间 $W_0^{1,p}(I)$ 所起的中心作用.微分方程 (或
偏微分方程) 常与边界条件耦合, 即在 $\partial I$ 上规定 $u$ 的值.
\end{remark}

\begin{proof}
若 $u\in W_0^{1,p}$, 则存在 $C_c^1(I)$ 中的序列 $(u_n)$ 使得 $u_n\to u$ 在
$W^{1,p}(I)$ 中.因此 $u_n\to u$ 在 $\bar I$ 上一致, 从而 $u=0$ 在 $\partial I$
上.

反之, 设 $u\in W^{1,p}(I)$ 满足 $u=0$ 在 $\partial I$ 上.固定任一函数
$G\in C^1(\R)$, 使得
\begin{equation*}
G(t)=
\begin{cases}
0 & \text{若}\ \abs{t}\le1,\\
t & \text{若}\ \abs{t}\ge2,
\end{cases}
\qquad\text{且}\qquad
\abs{G(t)}\le\abs{t}\quad\forall t\in\R.
\end{equation*}
令 $u_n=(1/n)G(nu)$, 于是 $u_n\in W^{1,p}(I)$ (由推论  \ref{cor:8.11}).另一方面,
\begin{equation*}
\operatorname{supp}u_n\subset\{x\in I;\ \abs{u(x)}\ge1/n\},
\end{equation*}
因此 $\operatorname{supp}u_n$ 包含在 $I$ 的某个紧子集中 (利用 $u=0$ 在
$\partial I$ 上以及当 $\abs{x}\to\infty$, $x\in I$ 时 $u(x)\to0$ 这一事实).
于是 $u_n\in W_0^{1,p}(I)$ (见注记 \ref{rem:8.14}).最后, 由控制收敛定理容易
验证 $u_n\to u$ 在 $W^{1,p}(I)$ 中.于是 $u\in W_0^{1,p}(I)$.
\end{proof}

\begin{remark}
\label{rem:8.16}
我们再提一下 $W_0^{1,p}$ 函数的另外两个刻画:
\begin{enumerate}[label=(\roman*), leftmargin=2.6em, itemsep=0.25em]
	\item 设 $1\le p<\infty$, $u\in L^p(I)$.定义 $\bar u$ 为
	      \begin{equation*}
	      \bar u(x)=
	      \begin{cases}
	      u(x) & \text{若}\ x\in I,\\
	      0 & \text{若}\ x\in\R\setminus I.
	      \end{cases}
	      \end{equation*}
	      则 $u\in W_0^{1,p}(I)$ 当且仅当 $\bar u\in W^{1,p}(\R)$.
	\item 设 $1<p<\infty$, $u\in L^p(I)$.则 $u$ 属于 $W_0^{1,p}(I)$ 当且仅当
	      存在常数 $C$, 使得
	      \begin{equation*}
	      \abs{\int_I u\varphi'}\le C\norm{\varphi}_{L^{p'}(I)}
	      \quad\forall\varphi\in C_c^1(\R).
	      \end{equation*}
\end{enumerate}
\end{remark}

\begin{keyproposition}[Poincaré 不等式]
\label{prop:8.13}
设 $I$ 为有界区间.则存在常数 $C$ (依赖于 $\abs{I}<\infty$), 使得
\begin{equation}
\norm{u}_{W^{1,p}(I)}\le C\norm{u'}_{L^p(I)}\quad\forall u\in W_0^{1,p}(I).
\label{eq:8.13}
\end{equation}
换言之, 在 $W_0^{1,p}$ 上, 量 $\norm{u'}_{L^p(I)}$ 是与 $W^{1,p}$ 范数等价的
一个范数.
\end{keyproposition}

\begin{proof}
设 $u\in W_0^{1,p}(I)$ (其中 $I=(a,b)$).由于 $u(a)=0$, 我们有
\begin{equation*}
\abs{u(x)}=\abs{u(x)-u(a)}=\abs{\int_a^x u'(t)\,dt}\le\norm{u'}_{L^1}.
\end{equation*}
于是 $\norm{u}_{L^\infty(I)}\le\norm{u'}_{L^1(I)}$, 而 \eqref{eq:8.13} 由
Hölder 不等式推出.
\end{proof}

\begin{remark}
\label{rem:8.17}
若 $I$ 有界, 则表达式 $(u',v')_{L^2}=\int u'v'$ 定义了 $H_0^1$ 上的一个标量积,
且相应的范数, 即 $\norm{u'}_{L^2}$, 与 $H^1$ 范数等价.
\end{remark}

\begin{remark}
\label{rem:8.18}
给定整数 $m\ge2$ 与实数 $1\le p<\infty$, 空间 $W_0^{m,p}(I)$ 定义为
$C_c^m(I)$ 在 $W^{m,p}(I)$ 中的闭包.可以证明 (见习题 8.9)
\begin{equation*}
W_0^{m,p}(I)=\{u\in W^{m,p}(I);\ u=Du=\cdots=D^{m-1}u=0
\ \text{在}\ \partial I\ \text{上}\}.
\end{equation*}
必须注意
\begin{equation*}
W_0^{2,p}(I)=\{u\in W^{2,p}(I);\ u=Du=0\ \text{在}\ \partial I\ \text{上}\}
\end{equation*}
与
\begin{equation*}
W^{2,p}(I)\cap W_0^{1,p}(I)=\{u\in W^{2,p}(I);\ u=0\ \text{在}\ \partial I\ \text{上}\}
\end{equation*}
之间的区别.
\end{remark}

\subsection*{$\star$ $W_0^{1,p}(I)$ 的对偶空间}

\begin{notation}
空间 $W_0^{1,p}(I)$ ($1\le p<\infty$) 的对偶空间记作 $W^{-1,p'}(I)$, 空间
$H_0^1(I)$ 的对偶空间记作 $H^{-1}(I)$.
\end{notation}

沿用第~
\ref{chap:5}~章注记 \ref{rem:5.3} 的做法, 我们把 $L^2$ 与其对偶等同起来, 但不把
$H_0^1$ 与其对偶等同起来.我们有包含关系
\begin{equation*}
H_0^1\subset L^2\subset H^{-1},
\end{equation*}
其中这些嵌入连续且稠密 (即值域稠密).

若 $I$ 是有界区间, 我们有
\begin{equation*}
W_0^{1,p}\subset L^2\subset W^{-1,p'}\quad\text{对一切}\ 1\le p<\infty,
\end{equation*}
嵌入连续 (当 $1<p<\infty$ 时嵌入稠密).

若 $I$ 无界, 我们只有
\begin{equation*}
W_0^{1,p}\subset L^2\subset W^{-1,p'}\quad\text{对一切}\ 1\le p\le2,
\end{equation*}
嵌入连续 (见注记 \ref{rem:8.12}).

$W^{-1,p'}$ 中的元素可以借助 $L^{p'}$ 中的函数来表示; 确切地说, 我们有下述结果.

\begin{proposition}
\label{prop:8.14}
设 $F\in W^{-1,p'}(I)$.则存在两个函数 $f_0,f_1\in L^{p'}(I)$, 使得
\begin{equation*}
\dual{F}{u}=\int_I f_0u+\int_I f_1u'\quad\forall u\in W_0^{1,p}(I)
\end{equation*}
且
\begin{equation*}
\norm{F}_{W^{-1,p'}}=\max\{\norm{f_0}_{p'},\norm{f_1}_{p'}\}.
\end{equation*}
当 $I$ 有界时, 可以取 $f_0=0$.
\end{proposition}

\begin{proof}
考虑乘积空间 $E=L^p(I)\times L^p(I)$, 配以范数
$\norm{h}=\norm{h_0}_p+\norm{h_1}_p$, 其中 $h=[h_0,h_1]$.

映射 $T:u\in W_0^{1,p}(I)\mapsto[u,u']\in E$ 是从 $W_0^{1,p}(I)$ 到 $E$ 的
等距映射.令 $G=T(W_0^{1,p}(I))$, 配以 $E$ 的范数, 并令
$S=T^{-1}:G\to W_0^{1,p}(I)$.映射 $h\in G\mapsto\dual{F}{Sh}$ 是 $G$ 上的
连续线性泛函.由 Hahn--Banach 定理, 它可以延拓为 $E$ 上的连续线性泛函
$\Phi$, 且 $\norm{\Phi}_{E^\star}=\norm{F}$.由 Riesz 表示定理, 我们知道存在
两个函数 $f_0,f_1\in L^{p'}(I)$, 使得
\begin{equation*}
\dual{\Phi}{h}=\int_I f_0h_0+\int_I f_1h_1
\quad\forall h=[h_0,h_1]\in E.
\end{equation*}
容易验证 $\norm{\Phi}_{E^\star}=\max\{\norm{f_0}_{p'},\norm{f_1}_{p'}\}$.
此外, 我们有
\begin{equation*}
\dual{\Phi}{Tu}=\dual{F}{u}=\int_I f_0u+\int_I f_1u'
\quad\forall u\in W_0^{1,p}.
\end{equation*}
当 $I$ 有界时, 空间 $W_0^{1,p}(I)$ 可以配以范数 $\norm{u'}_p$ (见命题
 \ref{prop:8.13}).我们用 $E=L^p(I)$ 与
$T:u\in W_0^{1,p}(I)\mapsto u'\in L^p(I)$ 重复同样的论证.
\end{proof}

\begin{remark}
\label{rem:8.19}
函数 $f_0$ 与 $f_1$ 不由 $F$ 唯一确定.
\end{remark}

\begin{remark}
\label{rem:8.20}
元素 $F\in W^{-1,p'}(I)$ 通常与分布 $f_0-f_1'$ 等同起来 (按定义, 分布
$f_0-f_1'$ 是线性泛函
$u\mapsto\int_I f_0u+\int_I f_1u'$, 定义在 $C_c^\infty(I)$ 上).
\end{remark}

\begin{remark}
\label{rem:8.21}
命题  \ref{prop:8.14} 的第一个断言对 $W^{1,p}$ ($1\le p<\infty$) 上的连续线性
泛函也成立, 即 $W^{1,p}$ 上的每个连续线性泛函 $F$ 都可以表示为
\begin{equation*}
\dual{F}{u}=\int_I f_0u+\int_I f_1u'\quad\forall u\in W^{1,p},
\end{equation*}
其中 $f_0,f_1\in L^{p'}$ 为某两个函数.
\end{remark}

\section{一些边值问题的例子}
\label{sec:8.4}

考虑问题
\begin{equation}
\begin{cases}
-u''+u=f & \text{在}\ I=(0,1)\ \text{上},\\
u(0)=u(1)=0,
\end{cases}
\label{eq:8.14}
\end{equation}
其中 $f$ 是给定函数 (例如属于 $C(\bar I)$, 或者更一般地属于 $L^2(I)$).边界
条件 $u(0)=u(1)=0$ 称为 (\emph{齐次}) \emph{Dirichlet 边界条件}.

\begin{definition}
\eqref{eq:8.14} 的\emph{古典解}是指满足 \eqref{eq:8.14} 的通常意义的
$C^2(\bar I)$ 函数.\eqref{eq:8.14} 的\emph{弱解}是指满足
\begin{equation}
\int_I u'v'+\int_I uv=\int_I fv\quad\forall v\in H_0^1(I)
\label{eq:8.15}
\end{equation}
的函数 $u\in H_0^1(I)$.
\end{definition}

让我们把 \S \ref{sec:8.1} 中概述的程序``付诸实施'':

\runin{步骤 A.} \emph{每个古典解都是弱解.}由分部积分 (其合理性已在推论
 \ref{cor:8.10} 中说明) 这是显然的.

\runin{步骤 B.} \emph{弱解的存在性与唯一性.}这就是下述结果的内容.

\begin{keyproposition}
\label{prop:8.15}
给定任意 $f\in L^2(I)$, 存在唯一的解 $u\in H_0^1$ 满足 \eqref{eq:8.15}.
此外, $u$ 由
\begin{equation*}
\boxed{\min_{v\in H_0^1}\left\{\frac12\int_I(v'^2+v^2)-\int_I fv\right\}}
\end{equation*}
得到; 这就是\emph{Dirichlet 原理}.
\end{keyproposition}

\begin{proof}
我们在 Hilbert 空间 $H=H_0^1(I)$ 中对双线性形式
\begin{equation*}
a(u,v)=\int_I u'v'+\int_I uv=(u,v)_{H^1}
\end{equation*}
以及线性泛函 $\varphi:v\mapsto\int_I fv$ 应用 Lax--Milgram 定理 (推论
 \ref{cor:5.8}).
\end{proof}

\begin{remark}
\label{rem:8.22}
给定 $F\in H^{-1}(I)$, 由 Riesz--Fréchet 表示定理 (定理  \ref{thm:5.5}) 我们
知道存在唯一的 $u\in H_0^1(I)$ 使得
\begin{equation*}
(u,v)_{H^1}=\dual{F}{v}_{H^{-1},H_0^1}\quad\forall v\in H_0^1.
\end{equation*}
映射 $F\mapsto u$ 是从 $H^{-1}$ 到 $H_0^1$ 上的 Riesz--Fréchet 同构.函数 $u$
与 \eqref{eq:8.14} 在 \eqref{eq:8.15} 意义下的弱解重合.
\end{remark}

\subsection*{步骤 C 与步骤 D. 弱解的正则性.古典解的复原}

首先注意, 若 $f\in L^2$ 且 $u\in H_0^1$ 是 \eqref{eq:8.14} 的弱解, 则
$u\in H^2$.事实上, 我们有
\begin{equation*}
\int_I u'v'=\int_I(f-u)v\quad\forall v\in C_c^1(I),
\end{equation*}
因此 $u'\in H^1$ (由 $H^1$ 的定义以及 $f-u\in L^2$), 即 $u\in H^2$.

此外, 若假设 $f\in C(\bar I)$, 则弱解 $u$ 属于 $C^2(\bar I)$.事实上,
$(u')'\in C(\bar I)$, 从而 $u'\in C^1(\bar I)$ (见注记 \ref{rem:8.6}).从弱解
$u\in C^2(\bar I)$ 到古典解的过渡已在 \S \ref{sec:8.1} 中完成.

\begin{remark}
\label{rem:8.23}
若 $f\in H^k(I)$, 其中 $k$ 为整数且 $k\ge1$, 则容易验证 (用归纳法)
\eqref{eq:8.15} 的解 $u$ 属于 $H^{k+2}(I)$.
\end{remark}

上面描述的方法是极其灵活的, 可以适应大量的问题.我们举出经常遇到的几个
例子.在每个问题中, 精确地确定函数空间并找出适当的弱表述都是关键.

\begin{example}[非齐次 Dirichlet 条件]
\label{ex:8.1}
考虑问题
\begin{equation}
\begin{cases}
-u''+u=f & \text{在}\ I=(0,1)\ \text{上},\\
u(0)=\alpha,\ u(1)=\beta,
\end{cases}
\label{eq:8.16}
\end{equation}
其中 $\alpha,\beta\in\R$ 给定, $f$ 是给定函数.
\end{example}

\begin{keyproposition}
\label{prop:8.16}
给定 $\alpha,\beta\in\R$ 与 $f\in L^2(I)$, 存在唯一的函数 $u\in H^2(I)$
满足 \eqref{eq:8.16}.此外, $u$ 由
\begin{equation*}
\min_{\substack{v\in H^1(I)\\ v(0)=\alpha,\ v(1)=\beta}}
\left\{\frac12\int_I(v'^2+v^2)-\int_I fv\right\}
\end{equation*}
得到.若另外 $f\in C(\bar I)$, 则 $u\in C^2(\bar I)$.
\end{keyproposition}

\begin{proof}
我们给出两种可能的途径:

\emph{方法 1.} 固定任一光滑函数\footnote{例如取 $u_0$ 为仿射函数.} $u_0$,
使得 $u_0(0)=\alpha$, $u_0(1)=\beta$.引入新未知函数 $\tilde u=u-u_0$.于是
$\tilde u$ 满足
\begin{equation*}
\begin{cases}
-\tilde u''+\tilde u=f+u_0''-u_0 & \text{在}\ I\ \text{上},\\
\tilde u(0)=\tilde u(1)=0.
\end{cases}
\end{equation*}
我们化归为关于 $\tilde u$ 的前一个问题.

\emph{方法 2.} 在空间 $H^1(I)$ 中考虑闭凸集
\begin{equation*}
K=\{v\in H^1(I);\ v(0)=\alpha\ \text{且}\ v(1)=\beta\}.
\end{equation*}
若 $u$ 是 \eqref{eq:8.16} 的古典解, 则
\begin{equation*}
\int_I u'(v-u)'+\int_I u(v-u)=\int_I f(v-u)\quad\forall v\in K.
\end{equation*}
特别地,
\begin{equation}
\int_I u'(v-u)'+\int_I u(v-u)\ge\int_I f(v-u)\quad\forall v\in K.
\label{eq:8.17}
\end{equation}
我们现在可以引用 Stampacchia 定理 (定理  \ref{thm:5.6}): 存在唯一的函数
$u\in K$ 满足 \eqref{eq:8.17}, 并且 $u$ 由
\begin{equation*}
\min_{v\in K}\left\{\frac12\int_I(v'^2+v^2)-\int_I fv\right\}
\end{equation*}
得到.为从 \eqref{eq:8.16} 恢复出古典解, 在 \eqref{eq:8.17} 中取
$v=u\pm w$, 其中 $w\in H_0^1$, 得到
\begin{equation*}
\int_I u'w'+\int_I uw=\int_I fw\quad\forall w\in H_0^1.
\end{equation*}
由此推出 (与上面一样) $u\in H^2(I)$.若 $f\in C(\bar I)$, 则与齐次情形相同的
论证说明 $u\in C^2(\bar I)$.
\end{proof}

\begin{starexample}[Sturm--Liouville 问题]
\label{ex:8.2}
考虑问题
\begin{equation}
\begin{cases}
-(pu')'+qu=f & \text{在}\ I=(0,1)\ \text{上},\\
u(0)=u(1)=0,
\end{cases}
\label{eq:8.18}
\end{equation}
其中 $p\in C^1(\bar I)$, $q\in C(\bar I)$, $f\in L^2(I)$ 给定, 且
\begin{equation*}
p(x)\ge\alpha>0\quad\forall x\in I.
\end{equation*}
若 $u$ 是 \eqref{eq:8.18} 的古典解, 则
\begin{equation*}
\int_I pu'v'+\int_I quv=\int_I fv\quad\forall v\in H_0^1(I).
\end{equation*}
我们取 $H_0^1(I)$ 作为函数空间, 并取
\begin{equation*}
a(u,v)=\int_I pu'v'+\int_I quv
\end{equation*}
作为 $H_0^1$ 上的对称连续双线性形式.若 $q\ge0$ 在 $I$ 上, 则由 Poincaré
不等式 (命题  \ref{prop:8.13}) 这个形式是强制的.于是由 Lax--Milgram 定理,
存在唯一的 $u\in H_0^1$ 使得
\begin{equation}
a(u,v)=\int_I fv\quad\forall v\in H_0^1(I).
\label{eq:8.19}
\end{equation}
此外, $u$ 由
\begin{equation*}
\min_{v\in H_0^1(I)}\left\{\frac12\int_I(pv'^2+qv^2)-\int_I fv\right\}
\end{equation*}
得到.由 \eqref{eq:8.19} 显然有 $pu'\in H^1$; 于是 (由推论  \ref{cor:8.10})
$u'=(1/p)(pu')\in H^1$, 从而 $u\in H^2$.最后, 若 $f\in C(\bar I)$, 则
$pu'\in C^1(\bar I)$, 于是 $u'\in C^1(\bar I)$, 即 $u\in C^2(\bar I)$.步骤 D
可以照搬, 我们断言 $u$ 是 \eqref{eq:8.18} 的古典解.

现在考虑更一般的问题
\begin{equation}
\begin{cases}
-(pu')'+ru'+qu=f & \text{在}\ I=(0,1)\ \text{上},\\
u(0)=u(1)=0.
\end{cases}
\label{eq:8.20}
\end{equation}
关于 $p$, $q$, $f$ 的假设同上, 而 $r\in C(\bar I)$.若 $u$ 是 \eqref{eq:8.20}
的古典解, 则
\begin{equation*}
\int_I pu'v'+\int_I ru'v+\int_I quv=\int_I fv\quad\forall v\in H_0^1.
\end{equation*}
我们取 $H_0^1(I)$ 作为函数空间, 并取
\begin{equation*}
a(u,v)=\int_I pu'v'+\int_I ru'v+\int_I quv
\end{equation*}
作为连续双线性形式.这个形式不是对称的.在某些情形下它是强制的; 例如
\begin{enumerate}[label=(\roman*), leftmargin=2.4em, itemsep=0.2em]
	\item 若 $q\ge1$ 且 $r^2<4\alpha$;
	\item 或者若 $q\ge1$ 且 $r\in C^1(\bar I)$, $r'\le2$; 这里我们用到
	      \begin{equation*}
	      \int rv'v=-\frac12\int r'v^2\quad\forall v\in H_0^1.
	      \end{equation*}
\end{enumerate}
这时可以应用 Lax--Milgram 定理, 但没有与之直接对应的极小化问题.下面是一个
使我们能恢复对称双线性形式的技巧.引入 $r/p$ 的一个原函数 $R$, 并令
$\zeta=e^{-R}$.方程 \eqref{eq:8.20} 乘以 $\zeta$ 后可以写成
\begin{equation*}
-\zeta pu''-\zeta p'u'+\zeta ru'+\zeta qu=\zeta f,
\end{equation*}
或者 (因为 $\zeta'p+\zeta r=0$)
\begin{equation*}
-(\zeta pu')'+\zeta qu=\zeta f.
\end{equation*}
在 $H_0^1$ 上定义对称连续双线性形式
\begin{equation*}
a(u,v)=\int_I\zeta pu'v'+\int_I\zeta quv.
\end{equation*}
当 $q\ge0$ 时这个形式是强制的, 于是存在唯一的 $u\in H_0^1$ 使得
\begin{equation*}
a(u,v)=\int_I\zeta fv\quad\forall v\in H_0^1.
\end{equation*}
此外, $u$ 由
\begin{equation*}
\min_{v\in H_0^1(I)}\left\{\frac12\int_I(\zeta pv'^2+\zeta qv^2)
-\int_I\zeta fv\right\}
\end{equation*}
得到.容易验证 $u\in H^2$; 并且若 $f\in C(\bar I)$, 则 $u\in C^2(\bar I)$ 是
\eqref{eq:8.20} 的古典解.
\end{starexample}

\begin{example}[齐次 Neumann 条件]
\label{ex:8.3}
考虑问题
\begin{equation}
\begin{cases}
-u''+u=f & \text{在}\ I=(0,1)\ \text{上},\\
u'(0)=u'(1)=0.
\end{cases}
\label{eq:8.21}
\end{equation}
\end{example}

\begin{keyproposition}
\label{prop:8.17}
给定 $f\in L^2(I)$, 存在唯一的函数 $u\in H^2(I)$ 满足 \eqref{eq:8.21}.
\footnote{注意 $u\in H^2(I)\Rightarrow u\in C^1(\bar I)$, 因此条件
$u'(0)=u'(1)=0$ 有意义.若只知道 $u\in H^1$, 它就没有意义.}此外, $u$ 由
\begin{equation*}
\min_{v\in H^1(I)}\left\{\frac12\int_I(v'^2+v^2)-\int_I fv\right\}
\end{equation*}
得到.若另外 $f\in C(\bar I)$, 则 $u\in C^2(\bar I)$.
\end{keyproposition}

\begin{proof}
若 $u$ 是 \eqref{eq:8.21} 的古典解, 则
\begin{equation}
\int_I u'v'+\int_I uv=\int_I fv\quad\forall v\in H^1(I).
\label{eq:8.22}
\end{equation}
我们取 $H^1(I)$ 作为函数空间: 由于 $u(0)$ 与 $u(1)$ 先验未知, 像上面那样在
$H_0^1$ 中工作是毫无意义的.我们对双线性形式
$a(u,v)=\int_I u'v'+\int_I uv$ 与线性泛函 $\varphi:v\mapsto\int_I fv$ 应用
Lax--Milgram 定理.这样我们得到满足 \eqref{eq:8.22} 的唯一函数
$u\in H^1(I)$.由 \eqref{eq:8.22} 如上推出 $u\in H^2(I)$.再次利用
\eqref{eq:8.22}, 我们得到
\begin{equation}
\int_I(-u''+u-f)v+u'(1)v(1)-u'(0)v(0)=0\quad\forall v\in H^1(I).
\label{eq:8.23}
\end{equation}
在 \eqref{eq:8.23} 中先取 $v\in H_0^1$, 得到 $-u''+u=f$ 几乎处处成立.回到
\eqref{eq:8.23}, 剩下
\begin{equation*}
u'(1)v(1)-u'(0)v(0)=0\quad\forall v\in H^1(I).
\end{equation*}
由于 $v(0)$ 与 $v(1)$ 任意, 我们推出 $u'(0)=u'(1)=0$.
\end{proof}

\begin{example}[非齐次 Neumann 条件]
\label{ex:8.4}
考虑问题
\begin{equation}
\begin{cases}
-u''+u=f & \text{在}\ I=(0,1)\ \text{上},\\
u'(0)=\alpha,\ u'(1)=\beta,
\end{cases}
\label{eq:8.24}
\end{equation}
其中 $\alpha,\beta\in\R$ 给定, $f$ 是给定函数.
\end{example}

\begin{proposition}
\label{prop:8.18}
给定任意 $f\in L^2(I)$ 与 $\alpha,\beta\in\R$, 存在唯一的函数 $u\in H^2(I)$
满足 \eqref{eq:8.24}.此外, $u$ 由
\begin{equation*}
\min_{v\in H^1(I)}\left\{\frac12\int_I(v'^2+v^2)-\int_I fv+\alpha v(0)-\beta v(1)\right\}
\end{equation*}
得到.若另外 $f\in C(\bar I)$, 则 $u\in C^2(\bar I)$.
\end{proposition}

\begin{proof}
若 $u$ 是 \eqref{eq:8.24} 的古典解, 则
\begin{equation*}
\int_I u'v'+\int_I uv=\int_I fv-\alpha v(0)+\beta v(1)\quad\forall v\in H^1(I).
\end{equation*}
我们取 $H^1(I)$ 作为函数空间, 并对双线性形式
$a(u,v)=\int_I u'v'+\int_I uv$ 与线性泛函
\begin{equation*}
\varphi:v\mapsto\int_I fv-\alpha v(0)+\beta v(1)
\end{equation*}
应用 Lax--Milgram 定理.这个线性泛函连续 (由定理  \ref{thm:8.8}).然后照例 3
的做法证明 $u\in H^2(I)$ 且 $u'(0)=\alpha$, $u'(1)=\beta$.
\end{proof}

\begin{example}[混合边界条件]
\label{ex:8.5}
考虑问题
\begin{equation}
\begin{cases}
-u''+u=f & \text{在}\ I=(0,1)\ \text{上},\\
u(0)=0,\ u'(1)=0.
\end{cases}
\label{eq:8.25}
\end{equation}
若 $u$ 是 \eqref{eq:8.25} 的古典解, 则
\begin{equation}
\int_I u'v'+\int_I uv=\int_I fv\quad\forall v\in H^1(I)\ \text{且}\ v(0)=0.
\label{eq:8.26}
\end{equation}
适当的空间是
\begin{equation*}
H=\{v\in H^1(I);\ v(0)=0\},
\end{equation*}
配以 $H^1$ 标量积.其余部分留作习题.
\end{example}

\begin{example}[Robin 条件, 或``第三类''边界条件]
\label{ex:8.6}
考虑问题
\begin{equation}
\begin{cases}
-u''+u=f & \text{在}\ I=(0,1)\ \text{上},\\
u'(0)=ku(0),\ u(1)=0,
\end{cases}
\label{eq:8.27}
\end{equation}
其中 $k\in\R$ 给定.\footnote{更一般地, 可以处理边界条件
$\alpha_0u'(0)+\beta_0u(0)=0$, $\alpha_1u'(1)+\beta_1u(1)=0$, 只要对常数
$\alpha_0,\beta_0,\alpha_1,\beta_1$ 加适当的条件.}

若 $u$ 是 \eqref{eq:8.27} 的古典解, 则
\begin{equation*}
\int_I u'v'+\int_I uv+ku(0)v(0)=\int_I fv
\quad\forall v\in H^1(I)\ \text{且}\ v(1)=0.
\end{equation*}
应用 Lax--Milgram 定理的适当空间是 Hilbert 空间
\begin{equation*}
H=\{v\in H^1(I);\ v(1)=0\},
\end{equation*}
配以 $H^1$ 标量积.双线性形式
\begin{equation*}
a(u,v)=\int_I u'v'+\int_I uv+ku(0)v(0)
\end{equation*}
对称且连续.它在 $k\ge0$ 时强制.\footnote{若 $k<0$ 且 $\abs{k}$ 足够小,
形式 $a(u,v)$ 仍然是强制的.另一方面, 显式计算表明存在某个负的 $k$ 值以及
(光滑) 函数 $f$, 使得 \eqref{eq:8.27} 无解 (见习题 8.21).}
\end{example}

\begin{example}[周期边界条件]
\label{ex:8.7}
考虑问题
\begin{equation}
\begin{cases}
-u''+u=f & \text{在}\ I=(0,1)\ \text{上},\\
u(0)=u(1),\ u'(0)=u'(1).
\end{cases}
\label{eq:8.28}
\end{equation}
若 $u$ 是 \eqref{eq:8.28} 的古典解, 则
\begin{equation}
\int_I u'v'+\int_I uv=\int_I fv\quad\forall v\in H^1(I)\ \text{且}\ v(0)=v(1).
\label{eq:8.29}
\end{equation}
应用 Lax--Milgram 定理的适当框架是 Hilbert 空间
\begin{equation*}
H=\{v\in H^1(I);\ v(0)=v(1)\},
\end{equation*}
配以双线性形式 $a(u,v)=\int_I u'v'+\int_I uv$.当 $f\in L^2(I)$ 时我们得到
\eqref{eq:8.28} 的一个解 $u\in H^2(I)$.若另外 $f\in C(\bar I)$, 则这个解
是古典解.
\end{example}

\begin{example}[$\R$ 上的一个边值问题]
\label{ex:8.8}
考虑问题
\begin{equation}
\begin{cases}
-u''+u=f & \text{在}\ \R\ \text{上},\\
u(x)\to0 & \text{当}\ \abs{x}\to\infty,
\end{cases}
\label{eq:8.30}
\end{equation}
其中 $f$ 给定于 $L^2(\R)$.\eqref{eq:8.30} 的\emph{古典解}是指满足
\eqref{eq:8.30} 的通常意义的 $C^2(\R)$ 函数.\eqref{eq:8.30} 的\emph{弱解}
是指满足
\begin{equation*}
\int_\R u'v'+\int_\R uv=\int_\R fv\quad\forall v\in H^1(\R)
\end{equation*}
的函数 $u\in H^1(\R)$.

我们首先必须证明任何古典解都是弱解; 让我们先验证 $u\in H^1(\R)$.像定理
 \ref{thm:8.7} 的证明中那样取截断函数序列 $(\zeta_n)$.用 $\zeta_nu$ 乘
\eqref{eq:8.30} 并分部积分, 我们得到
\begin{equation*}
\int_\R u'(\zeta_nu'+\zeta_n'u)+\int_\R\zeta_nu^2=\int_\R\zeta_nfu,
\end{equation*}
由此推出
\begin{equation}
\int_\R\zeta_n(u'^2+u^2)=\int_\R\zeta_nfu+\frac12\int_\R\zeta_n''u^2.
\label{eq:8.31}
\end{equation}
但是
\begin{equation*}
\frac12\int_\R\zeta_n''u^2\le\frac{C}{n^2}\int_{n<\abs{x}<2n}u^2
\quad\text{其中}\ C=\norm{\zeta''}_{L^\infty(\R)},
\end{equation*}
并且当 $n\to\infty$ 时 $\frac{1}{n^2}\int_{n<\abs{x}<2n}u^2\to0$, 因为
$u(x)\to0$ ($\abs{x}\to\infty$).把不等式
\begin{equation*}
\int_\R\zeta_nfu\le\frac12\int_\R\zeta_nu^2+\frac12\int_\R\zeta_nf^2
\end{equation*}
代入 \eqref{eq:8.31}, 我们看到当 $n\to\infty$ 时
$\int\zeta_n(u'^2+u^2)$ 保持有界, 因此 $u\in H^1(\R)$.

假设 $u$ 是 \eqref{eq:8.30} 的古典解, 我们有
\begin{equation*}
\int_\R u'v'+\int_\R uv=\int_\R fv\quad\forall v\in C_c^1(\R).
\end{equation*}
由稠密性 (以及 $u\in H^1(\R)$), 这对每个 $v\in H^1(\R)$ 都成立.因此 $u$
是 \eqref{eq:8.30} 的弱解.

为得到弱解的存在性与唯一性, 只须在 Hilbert 空间 $H^1(\R)$ 中应用
Lax--Milgram 定理.容易验证弱解 $u$ 属于 $H^2(\R)$, 并且若进一步
$f\in C(\R)$, 则 $u\in C^2(\R)$.我们断言 (利用推论  \ref{cor:8.9}):
给定 $f\in L^2(\R)\cap C(\R)$, 问题 \eqref{eq:8.30} 有唯一的古典解 (它
并且属于 $H^2(\R)$).
\end{example}

\begin{remark}
\label{rem:8.24}
问题
\begin{equation*}
\begin{cases}
-u''=f & \text{在}\ \R\ \text{上},\\
u(x)\to0 & \text{当}\ \abs{x}\to\infty,
\end{cases}
\end{equation*}
不能用上述技巧处理, 因为双线性形式 $a(u,v)=\int u'v'$ 在 $H^1(\R)$ 中不是
强制的.事实上, 即使 $f$ 是具紧支集的光滑函数, 这个问题也未必有解 (为什么?).
\end{remark}

\begin{remark}
\label{rem:8.25}
另一方面, 同样的方法适用于问题
\begin{equation*}
\begin{cases}
-u''+u=f & \text{在}\ I=(0,+\infty)\ \text{上},\\
u(0)=0 & \text{且}\ u(x)\to0\ \text{当}\ x\to+\infty,
\end{cases}
\end{equation*}
其中 $f$ 给定于 $L^2(0,+\infty)$.
\end{remark}

\section{最大值原理}
\label{sec:8.5}

下面是一个非常有用, 称为\emph{最大值原理}的性质.

\begin{keytheorem}
\index{yiweijidazhiyuanli@一维极大值原理 (maximum principle in dimension 1)}
\label{thm:8.19}
设 $f\in L^2(I)$, $I=(0,1)$, 并设 $u\in H^2(I)$ 是 Dirichlet 问题
\begin{equation}
\begin{cases}
-u''+u=f & \text{在}\ I\ \text{上},\\
u(0)=\alpha,\ u(1)=\beta
\end{cases}
\label{eq:8.32}
\end{equation}
的解.则对每个 $x\in I$,\footnote{$\sup_I f$ 与 $\inf_I f$ 分别指 $f$ 的
本质上确界 (可能为 $+\infty$) 与本质下确界 (可能为 $-\infty$).回忆
$\operatorname{ess\,sup}f=\inf\{C;\ f(x)\le C\ \text{a.e.}\}$,
$\operatorname{ess\,inf}f=-\operatorname{ess\,sup}(-f)$.}
\begin{equation}
\min\{\alpha,\beta,\inf_I f\}\le u(x)\le\max\{\alpha,\beta,\sup_I f\}.
\label{eq:8.33}
\end{equation}
\end{keytheorem}

\begin{proof}[(用 Stampacchia 的截断方法)]
我们有
\begin{equation}
\int_I u'v'+\int_I uv=\int_I fv\quad\forall v\in H_0^1(I).
\label{eq:8.34}
\end{equation}
固定任一函数 $G\in C^1(\R)$, 使得
\begin{enumerate}[label=(\roman*), leftmargin=2.6em, itemsep=0.2em]
	\item $G$ 在 $(0,+\infty)$ 上严格递增,
	\item 当 $t\in(-\infty,0]$ 时 $G(t)=0$.
\end{enumerate}
令 $K=\max\{\alpha,\beta,\sup_I f\}$, 并假设 $K<\infty$.我们将证明
$u\le K$ 在 $I$ 上成立.函数 $v=G(u-K)$ 属于 $H^1(I)$, 甚至属于
$H_0^1(I)$, 因为
\begin{equation*}
u(0)-K=\alpha-K\le0\quad\text{且}\quad u(1)-K=\beta-K\le0.
\end{equation*}
把 $v$ 代入 \eqref{eq:8.34}, 我们得到
\begin{equation*}
\int_I u'^2G'(u-K)+\int_I uG(u-K)=\int_I fG(u-K),
\end{equation*}
即
\begin{equation*}
\int_I u'^2G'(u-K)+\int_I(u-K)G(u-K)=\int_I(f-K)G(u-K).
\end{equation*}
但 $(f-K)\le0$ 且 $G(u-K)\ge0$, 由此推出 $(f-K)G(u-K)\le0$, 于是
\begin{equation*}
\int_I(u-K)G(u-K)\le0.
\end{equation*}
由于 $tG(t)\ge0$ 对一切 $t\in\R$ 成立, 上述不等式推出
$(u-K)G(u-K)=0$ 几乎处处成立.于是 $u\le K$ 几乎处处成立, 从而在 $I$ 上
处处成立, 因为 $u$ 连续.$u$ 的下界可通过把上述上界应用于 $-u$ 得到.
\end{proof}

\begin{remark}
\label{rem:8.26}
当 $f\in C(\bar I)$ 时, $u\in C^2(\bar I)$, 可以用另一种方法建立
\eqref{eq:8.33}: 最大值原理的\emph{古典方法}.设 $x_0\in\bar I$ 是 $u$ 在
$\bar I$ 上取到最大值的点.若 $x_0=0$ 或 $x_0=1$, 结论是显然的.否则
$0<x_0<1$, 于是 $u'(x_0)=0$, $u''(x_0)\le0$.由方程
\eqref{eq:8.32} 推出
% 原书此处印作 "From equation (33)"; (33) 是结论不等式, 而 u(x_0)=f(x_0)+u''(x_0)
% 实际来自方程 (32) (-u''+u=f).依技能规范保留原样并记录, 不擅自更正.
\begin{equation*}
u(x_0)=f(x_0)+u''(x_0)\le f(x_0)\le K,
\end{equation*}
从而 $u\le K$ 在 $I$ 上成立.
\end{remark}

下面是定理  \ref{thm:8.19} 的一些直接推论.

\begin{keycorollary}
\label{cor:8.20}
设 $u$ 是 \eqref{eq:8.34} 的解.
\begin{enumerate}[label=(\roman*), leftmargin=2.6em, itemsep=0.25em]
	\item 若 $u\ge0$ 在 $\partial I$ 上且 $f\ge0$ 在 $I$ 上, 则 $u\ge0$ 在
	      $I$ 上.
	\item 若 $u=0$ 在 $\partial I$ 上且 $f\in L^\infty(I)$, 则
	      $\norm{u}_{L^\infty(I)}\le\norm{f}_{L^\infty(I)}$.
	\item 若 $f=0$ 在 $I$ 上, 则
	      $\norm{u}_{L^\infty(I)}\le\norm{u}_{L^\infty(\partial I)}$.
\end{enumerate}
\end{keycorollary}

对于 Neumann 条件的情形, 我们有类似的结果.

\begin{proposition}
\label{prop:8.21}
设 $f\in L^2(I)$, $I=(0,1)$, 并设 $u\in H^2(I)$ 是问题
\begin{equation*}
\begin{cases}
-u''+u=f & \text{在}\ I\ \text{上},\\
u'(0)=u'(1)=0
\end{cases}
\end{equation*}
的解.则对每个 $x\in\bar I$,
\begin{equation}
\inf_I f\le u(x)\le\sup_I f.
\label{eq:8.35}
\end{equation}
\end{proposition}

\begin{proof}
我们有
\begin{equation}
\int_I u'v'+\int_I uv=\int_I fv\quad\forall v\in H^1(I).
\label{eq:8.36}
\end{equation}
取 $K=\sup_I f$ 以及上面同一个函数 $G$, 把 $v=G(u-K)$ 代入
\eqref{eq:8.36}.然后照定理  \ref{thm:8.19} 的证明进行.
\end{proof}

\begin{remark}
\label{rem:8.27}
若 $f\in C(\bar I)$, 则 $u\in C^2(\bar I)$, 并且可以沿注记 \ref{rem:8.26}
的思路建立 \eqref{eq:8.35}.注意若 $u$ 在 $\partial I$ 上取到最大值, 比如在
$0$ 处, 则 $u''(0)\le0$ (把 $u$ 关于 $0$ 向左作反射延拓, 并利用
$u'(0)=0$ 这一事实).
\end{remark}

\begin{remark}
\label{rem:8.28}
设 $f\in L^2(\R)$, 并设 $u\in H^2(\R)$ 是例 \ref{ex:8.8} 中讨论的问题
\begin{equation*}
\begin{cases}
-u''+u=f & \text{在}\ \R\ \text{上},\\
u(x)\to0 & \text{当}\ \abs{x}\to\infty
\end{cases}
\end{equation*}
的解.则对一切 $x\in\R$,
\begin{equation*}
\inf_\R f\le u(x)\le\sup_\R f.
\end{equation*}
\end{remark}

\section{特征函数与谱分解}
\label{sec:8.6}

下面是基本结果.

\begin{keytheorem}
\label{thm:8.22}
设 $p\in C^1(\bar I)$, $I=(0,1)$, 且 $p\ge\alpha>0$ 在 $I$ 上; 设
$q\in C(\bar I)$.则存在实数序列 $(\lambda_n)$ 与 $L^2(I)$ 的 Hilbert 基
$(e_n)$, 使得 $e_n\in C^2(\bar I)$ 对一切 $n$ 成立, 且
\begin{equation}
\begin{cases}
-(pe_n')'+qe_n=\lambda_ne_n & \text{在}\ I\ \text{上},\\
e_n(0)=e_n(1)=0.
\end{cases}
\label{eq:8.37}
\end{equation}
此外, 当 $n\to+\infty$ 时 $\lambda_n\to+\infty$.
\end{keytheorem}

我们说 $(\lambda_n)$ 是带 Dirichlet 边界条件的微分算子
$Au=-(pu')'+qu$ 的\emph{特征值}, $(e_n)$ 是相应的\emph{特征函数}.

\begin{proof}
我们总可以假设 $q\ge0$, 因为若不然, 取任一常数 $C$ 使得 $q+C\ge0$, 这
相当于在 \eqref{eq:8.37} 中把 $\lambda_n$ 换成 $\lambda_n+C$.对每个
$f\in L^2(I)$, 存在唯一的 $u\in H^2(I)\cap H_0^1(I)$ 满足
\begin{equation}
\begin{cases}
-(pu')'+qu=f & \text{在}\ I\ \text{上},\\
u(0)=u(1)=0.
\end{cases}
\label{eq:8.38}
\end{equation}
记 $T$ 为映射 $f\mapsto u$, 把它看作从 $L^2(I)$ 到 $L^2(I)$ 的算子.
\footnote{我们也可以把 $T$ 看作从 $H_0^1$ 到 $H_0^1$ 的算子 (见
\S\ref{sec:9.8} 注记 28).}我们断言 $T$ 自伴且紧.

首先是紧性.由 \eqref{eq:8.38} 我们有
\begin{equation*}
\int_I pu'^2+\int_I qu^2=\int_I fu,
\end{equation*}
于是 $\alpha\norm{u'}_{L^2}^2\le\norm{f}_{L^2}\norm{u}_{L^2}$.由此推出
$\norm{u}_{H^1}\le C\norm{f}_{L^2}$, 其中 $C$ 是只依赖于 $\alpha$ 的常数.
这可以写成
\begin{equation*}
\norm{Tf}_{H^1}\le C\norm{f}_{L^2}\quad\forall f\in L^2(I).
\end{equation*}
由于 $H^1(I)$ 到 $L^2(I)$ 的嵌入是紧的 (因为 $I$ 有界), 我们推出 $T$ 是
从 $L^2(I)$ 到 $L^2(I)$ 的紧算子.

其次, 我们证明 $T$ 自伴, 即
\begin{equation*}
\int_I(Tf)g=\int_I f(Tg)\quad\forall f,g\in L^2(I).
\end{equation*}
事实上, 令 $u=Tf$ 且 $v=Tg$, 我们有
\begin{equation}
-(pu')'+qu=f
\label{eq:8.39}
\end{equation}
以及
\begin{equation}
-(pv')'+qv=g.
\label{eq:8.40}
\end{equation}
把 \eqref{eq:8.39} 乘以 $v$, 把 \eqref{eq:8.40} 乘以 $u$, 然后积分, 我们得到
\begin{equation*}
\int_I pu'v'+\int_I quv=\int_I fv=\int_I gu,
\end{equation*}
这就是所要的结论.

最后, 我们注意
\begin{equation}
\int_I(Tf)f=\int_I uf=\int_I pu'^2+qu^2\ge0\quad\forall f\in L^2(I),
\label{eq:8.41}
\end{equation}
并且 $N(T)=\{0\}$, 因为 $Tf=0$ 蕴含 $u=0$, 从而 $f=0$.

应用定理  \ref{thm:6.11}, 我们知道 $L^2(I)$ 有一个由 $T$ 的特征向量组成的
Hilbert 基 $(e_n)_{n\ge1}$, 相应的特征值为 $(\mu_n)_{n\ge1}$.我们有
$\mu_n>0$ 对一切 $n$ 成立 (由 \eqref{eq:8.41} 有 $\mu_n\ge0$, 而
$\mu_n\ne0$, 因为 $N(T)=\{0\}$).我们还知道 $\mu_n\to0$.由
$Te_n=\mu_ne_n$ 得到
\begin{equation*}
-(pe_n')'+qe_n=\lambda_ne_n\quad\text{其中}\ \lambda_n=1/\mu_n,
\qquad e_n(0)=e_n(1)=0.
\end{equation*}
此外, $e_n\in C^2(\bar I)$, 因为 $f=\lambda_ne_n\in C(\bar I)$ (事实上,
若 $p,q\in C^\infty(\bar I)$, 则 $e_n\in C^\infty(\bar I)$).
\end{proof}

\begin{example}
若 $p\equiv1$ 且 $q\equiv0$, 则
\begin{equation*}
e_n(x)=\sqrt2\sin(n\pi x)\quad\text{且}\quad\lambda_n=n^2\pi^2,
\quad n=1,2,\dots.
\end{equation*}
\end{example}

\begin{remark}
\label{rem:8.29}
对于同一个微分算子, 特征值与特征函数随边界条件而变.作为习题, 试确定
算子 $Au=-u''$ 在例 \ref{ex:8.3}, \ref{ex:8.5}, \ref{ex:8.6} 与
\ref{ex:8.7} 的边界条件下的特征值.
\end{remark}

\begin{remark}
\label{rem:8.30}
$I$ 有界这一假设在证明算子 $T$ 的紧性时起本质作用.当 $I$ 无界时, 定理
 \ref{thm:8.22} 的结论一般不再成立;\footnote{在某些情形下, 对 $p$ 与 $q$
作适当假设, 定理  \ref{thm:8.22} 的结论在无界区间上仍然成立 (见问题 51).}
人们反而会遇到非常有趣的连续谱现象——关于这个课题, 可参看例如
M.~Reed--B.~Simon \bref{ReedSimon}{1}.在习题 8.38 中我们确定算子
$T:f\mapsto u$ 的特征值与谱, 其中 $u\in H^2(\R)$ 是问题 \eqref{eq:8.30}
的解: $T$ 是从 $L^2(\R)$ 到自身的自伴有界算子, 但它不紧.
\end{remark}

\section*{第 8 章评注}
\addcontentsline{toc}{section}{第 8 章评注}

\cmtitem{1. 一些进一步的不等式.}

让我们举出一些涉及 Sobolev 范数的非常有用的不等式:

(i) \emph{Poincaré--Wirtinger 不等式.}

设 $I$ 为有界区间.给定 $u\in L^2(I)$, 令
$\bar u=\frac{1}{\abs{I}}\int_I u$ (这是 $u$ 在 $I$ 上的平均值).我们有
\begin{equation*}
\norm{u-\bar u}_\infty\le\norm{u'}_1\quad\forall u\in W^{1,1}(I)
\end{equation*}
(见问题 47).

(ii) \emph{Hardy 不等式.}

设 $I=(0,1)$, $u\in W_0^{1,p}(I)$, $1<p<\infty$.则函数
\begin{equation*}
v(x)=\frac{u(x)}{x(1-x)}
\end{equation*}
属于 $L^p(I)$, 并且进一步
\begin{equation*}
\norm{v}_p\le C_p\norm{u'}_p\quad\forall u\in W_0^{1,p}(I)
\end{equation*}
(见习题 8.8).

(iii) \emph{Gagliardo--Nirenberg 插值不等式.}

设 $I$ 为有界区间, $1\le r\le\infty$, $1\le q\le p\le\infty$.则存在
常数 $C$ 使得
\begin{equation}
\norm{u}_p\le C\norm{u}_q^{1-a}\norm{u}_{W^{1,r}}^a
\quad\forall u\in W^{1,r}(I),
\label{eq:8.42}
\end{equation}
其中 $0\le a\le1$ 由
$a(\frac1q-\frac1r+1)=\frac1q-\frac1p$ 定义 (见习题 8.15).特别地, 由
不等式 \eqref{eq:8.42} 推出, 若 $p<\infty$ (或者甚至 $p=\infty$ 但
$r>1$), 则
\begin{equation}
\forall\varepsilon>0\ \exists C_\varepsilon>0\ \text{使得}\
\norm{u}_p\le\varepsilon\norm{u}_{W^{1,r}}+C_\varepsilon\norm{u}_q
\quad\forall u\in W^{1,r}(I).
\label{eq:8.43}
\end{equation}
也可以用直接的``紧性方法''建立 \eqref{eq:8.43}; 见习题 8.5.其它更一般的
不等式可在 L.~Nirenberg \bref{Nirenberg}{1} 中找到 (另见
A.~Friedman \bref{Friedman}{2}).特别地, 我们提请注意不等式
\begin{equation*}
\norm{u'}_p\le C\norm{u}_{W^{2,r}}^{1/2}\norm{u}_q^{1/2}
\quad\forall u\in W^{2,r}(I),
\end{equation*}
其中 $p$ 是 $q$ 与 $r$ 的调和平均, 即 $\frac1p=\frac12(\frac1q+\frac1r)$.

\cmtitem{2. Hilbert--Schmidt 算子.}

可以证明, 把 $L^2(I)$ 中每个 $f$ 对应到问题
\begin{equation*}
\begin{cases}
-(pu')'+qu=f & \text{在}\ I=(0,1)\ \text{上},\\
u(0)=u(1)=0
\end{cases}
\end{equation*}
的唯一解 $u$ 的算子 $T:f\mapsto u$ (假设 $p\ge\alpha>0$ 且 $q\ge0$) 是从
$L^2(I)$ 到 $L^2(I)$ 的 Hilbert--Schmidt 算子; 见习题 8.37.

\cmtitem{3. Sturm--Liouville 算子的谱性质.}

关于有界区间 $I$ 上带 Dirichlet 条件的 Sturm--Liouville 算子
$Au=-(pu')'+qu$ 的谱性质, 人们知道许多.其中我们提一下:
\begin{enumerate}[label=(\roman*), leftmargin=2.6em, itemsep=0.25em]
	\item 每个特征值的重数为 1: 这时说每个特征值是\emph{单}的.
	\item 若特征值 $(\lambda_n)$ 按递增顺序排列, 则相应于 $\lambda_n$ 的
	      特征函数 $e_n(x)$ 在 $I$ 上恰有 $(n-1)$ 个零点; 特别地, 第一个
	      特征函数 $e_1(x)$ 在 $I$ 上保持定号, 通常取 $e_1>0$ 在 $I$ 上.
	\item 当 $n\to\infty$ 时, 商 $\lambda_n/n^2$ 收敛到一个正的极限.
\end{enumerate}
其中一些性质在习题 8.33, 8.42 与问题 49 中讨论.感兴趣的读者还可以查阅
Weinberger \bref{Weinberger}{1}, M.~Protter--H.~Weinberger
\bref{ProtterWeinberger}{1}, E.~Coddington--N.~Levinson
\bref{CoddingtonLevinson}{1}, Ph.~Hartman \bref{Hartman}{1},
S.~Agmon \bref{Agmon}{1}, R.~Courant--D.~Hilbert \bref{CourantHilbert}{1},
第 1 卷, E.~Ince \bref{Ince}{1}, Y.~Pinchover--J.~Rubinstein
\bref{PinchoverRubinstein}{1}, A.~Zettl \bref{Zettl}{1} 以及
G.~Buttazzo--M.~Giaquinta--S.~Hildebrandt \bref{ButtazzoGiaquintaHildebrandt}{1}.

著名的 Gelfand--Levitan 理论研究一个重要的``反''问题: 仅由 Sturm--Liouville
算子 $Au=-u''+q(x)u$ 的谱的知识, 能够获取关于函数 $q(x)$ 的哪些信息?
这个问题因其众多应用而备受关注; 可参看例如 B.~Levitan \bref{Levitan}{1},
以及第~
\ref{chap:9}~章评注 13.

\section*{第 8 章习题}
\addcontentsline{toc}{section}{第 8 章习题}
\markright{第 8 章习题}

\begin{exercise}
\label{exer:8.1}
考虑函数
\begin{equation*}
u(x)=(1+x^2)^{-\alpha/2}\bigl(\log(2+x^2)\bigr)^{-1},\quad x\in\R,
\end{equation*}
其中 $0<\alpha<1$.验证 $u\in W^{1,p}(\R)$ 对一切 $p\in[1/\alpha,\infty]$
成立, 而 $u\notin L^q(\R)$ 对一切 $q\in[1,1/\alpha)$ 成立.
\end{exercise}

\begin{exercise}
\label{exer:8.2}
设 $I=(0,1)$.
\begin{enumerate}[label=\arabic*., leftmargin=2.2em, itemsep=0.3em]
	\item 假设 $(u_n)$ 是 $W^{1,p}(I)$ 中的有界序列, $1<p\le\infty$.证明
	      存在子序列 $(u_{n_k})$ 与某个 $u\in W^{1,p}(I)$, 使得
	      $\norm{u_{n_k}-u}_{L^\infty}\to0$.此外, 若 $1<p<\infty$, 则
	      $u_{n_k}'\weakto u'$ 在 $L^p(I)$ 中弱收敛; 若 $p=\infty$, 则
	      $u_{n_k}'\weakstar u'$ 在 $\sigma(L^\infty,L^1)$ 中弱$^\star$收敛.
	\item 构造 $W^{1,1}(I)$ 中的有界序列 $(u_n)$, 它没有在 $L^\infty(I)$ 中
	      收敛的子序列.
\end{enumerate}
[提示: 考虑由
\begin{equation*}
u_n(x)=
\begin{cases}
0 & \text{若}\ x\in[0,\tfrac12],\\
n(x-\tfrac12) & \text{若}\ x\in(\tfrac12,\tfrac12+\tfrac1n),\\
1 & \text{若}\ x\in[\tfrac12+\tfrac1n,1],
\end{cases}
\end{equation*}
定义的序列 $(u_n)$, 其中 $n\ge2$.]
\end{exercise}

\begin{exercise}
\label{exer:8.3}
\emph{Helly 选择定理.}设 $(u_n)$ 是 $W^{1,1}(0,1)$ 中的有界序列.目标是
证明存在子序列 $(u_{n_k})$, 使得 $u_{n_k}(x)$ 对每个 $x\in[0,1]$ 都收敛到
一个极限.
\begin{enumerate}[label=\arabic*., leftmargin=2.2em, itemsep=0.3em]
	\item 证明总可以进一步假设
	      \begin{equation}
	      \forall n,\ u_n\ \text{在}\ [0,1]\ \text{上单调不减}.
	      \tag{1}
	      \end{equation}
	      [提示: 考虑序列 $v_n(x)=\int_0^x\abs{u_n'(t)}\,dt$ 与
	      $w_n=v_n-u_n$.]
\end{enumerate}
以下假设 (1) 成立.
\begin{enumerate}[label=\arabic*., leftmargin=2.2em, itemsep=0.3em, start=2]
	\item 证明存在子序列 $(u_{n_k})$ 与可测集 $E\subset[0,1]$,
	      $\abs{E}=0$, 使得对每个 $x\in[0,1]\setminus E$, $u_{n_k}(x)$
	      收敛到一个极限, 记作 $u(x)$.
	      [提示: 利用 $W^{1,1}\subset L^1$ 的嵌入是紧的这一事实.]
	\item 证明 $u$ 在 $[0,1]\setminus E$ 上单调不减, 并推出存在可数集
	      $D\subset(0,1)$ 与单调不减函数
	      $\bar u:(0,1)\setminus D\to\R$, 使得
	      $\bar u(x+0)=\bar u(x-0)$ 对一切 $x\in(0,1)\setminus D$ 成立, 且
	      $\bar u(x)=u(x)$ 对一切 $x\in(0,1)\setminus(D\cup E)$ 成立.
	\item 证明 $u_{n_k}(x)\to\bar u(x)$ 对一切 $x\in(0,1)\setminus D$ 成立.
	\item 从序列 $(u_{n_k})$ 中构造一个对每个 $x\in[0,1]$ 都收敛的子序列.
	      [提示: 用对角线过程.]
\end{enumerate}
\end{exercise}

\begin{exercise}
\label{exer:8.4}
固定函数 $\phi\in C_c^\infty(\R)$, $\phi\not\equiv0$, 并令
$u_n(x)=\phi(x+n)$.设 $1\le p\le\infty$.
\begin{enumerate}[label=\arabic*., leftmargin=2.2em, itemsep=0.3em]
	\item 验证 $(u_n)$ 在 $W^{1,p}(\R)$ 中有界.
	\item 证明对任何 $1\le q\le\infty$, 都不存在在 $L^q(\R)$ 中强收敛的
	      子序列 $(u_{n_k})$.
	\item 证明对一切 $p\in(1,\infty)$, $u_n\weakto0$ 在 $W^{1,p}(\R)$ 中
	      弱收敛.
\end{enumerate}
\end{exercise}

\begin{exercise}
\label{exer:8.5}
设 $p>1$.
\begin{enumerate}[label=\arabic*., leftmargin=2.2em, itemsep=0.3em]
	\item 证明对一切 $\varepsilon>0$ 存在 $C=C(\varepsilon,p)$, 使得
	      \begin{equation}
	      \norm{u}_{L^\infty(0,1)}\le\varepsilon\norm{u'}_{L^p(0,1)}
	      +C\norm{u}_{L^1(0,1)}\quad\forall u\in W^{1,p}(0,1).
	      \tag{1}
	      \end{equation}
	      [提示: 取 $X=W^{1,p}(0,1)$, $Y=L^\infty(0,1)$,
	      $Z=L^1(0,1)$, 应用习题 6.12.]
	\item 证明当 $p=1$ 时 (1) 不成立.
	      [提示: 取 $u(x)=x^n$ 并令 $n\to\infty$.]
	\item 设 $1\le q<\infty$.证明对一切 $\varepsilon>0$ 存在
	      $C=C(\varepsilon,q)$, 使得
	      \begin{equation}
	      \norm{u}_{L^q(0,1)}\le\varepsilon\norm{u'}_{L^1(0,1)}
	      +C\norm{u}_{L^1(0,1)}\quad\forall u\in W^{1,1}(0,1).
	      \tag{2}
	      \end{equation}
\end{enumerate}
\end{exercise}

\begin{exercise}
\label{exer:8.6}
设 $I=(0,1)$, $p>1$.
\begin{enumerate}[label=\arabic*., leftmargin=2.2em, itemsep=0.3em]
	\item 验证 $W^{2,p}(I)\subset C^1(\bar I)$, 且嵌入是紧的.
	\item 由此推出: 对一切 $\varepsilon>0$, 存在 $C=C(\varepsilon,p)$, 使得
	      \begin{equation*}
	      \norm{u'}_{L^\infty(I)}+\norm{u}_{L^\infty(I)}
	      \le\varepsilon\norm{u''}_{L^p(I)}+C\norm{u}_{L^1(I)}
	      \quad\forall u\in W^{2,p}(I).
	      \end{equation*}
	\item 设 $1\le q<\infty$.证明对一切 $\varepsilon>0$ 存在
	      $C=C(\varepsilon,q)$, 使得
	      \begin{equation*}
	      \norm{u'}_{L^q(I)}+\norm{u}_{L^\infty(I)}
	      \le\varepsilon\norm{u''}_{L^1(I)}+C\norm{u}_{L^1(I)}
	      \quad\forall u\in W^{2,1}(I).
	      \end{equation*}
\end{enumerate}
更一般地, 设 $m\ge2$ 为整数.
\begin{enumerate}[label=\arabic*., leftmargin=2.2em, itemsep=0.3em, start=4]
	\item 证明对一切 $\varepsilon>0$ 存在 $C=C(\varepsilon,m,p)$, 使得
	      \begin{equation*}
	      \sum_{j=0}^{m-1}\norm{D^ju}_{L^\infty(I)}
	      \le\varepsilon\norm{D^mu}_{L^p(I)}+C\norm{u}_{L^1(I)}
	      \quad\forall u\in W^{m,p}(I).
	      \end{equation*}
	\item 设 $1\le q<\infty$.证明对一切 $\varepsilon>0$ 存在
	      $C=C(\varepsilon,q)$, 使得
	      \begin{equation*}
	      \norm{D^{(m-1)}u}_{L^q(I)}
	      +\sum_{j=0}^{m-2}\norm{D^ju}_{L^\infty(I)}
	      \le\varepsilon\norm{D^mu}_{L^1(I)}+C\norm{u}_{L^1(I)}
	      \quad\forall u\in W^{m,1}(I).
	      \end{equation*}
\end{enumerate}
\end{exercise}

\begin{exercise}
\label{exer:8.7}
设 $I=(0,1)$.给定定义在 $I$ 上的函数 $u$, 令
\begin{equation*}
\bar u(x)=
\begin{cases}
u(x) & \text{若}\ x\in I,\\
0 & \text{若}\ x\in\R,\ x\notin I.
\end{cases}
\end{equation*}
\begin{enumerate}[label=\arabic*., leftmargin=2.2em, itemsep=0.3em]
	\item 假设 $u\in W_0^{1,p}(I)$, $1\le p<\infty$.证明
	      $\bar u\in W^{1,p}(\R)$.
	\item 反之, 设 $u\in L^p(I)$, $1\le p<\infty$, 且 $\bar u\in W^{1,p}(\R)$.
	      证明 $u\in W_0^{1,p}(I)$.
	\item 设 $u\in L^p(I)$, $1<p<\infty$.证明 $u\in W_0^{1,p}(I)$ 当且仅当
	      存在常数 $C$, 使得
	      \begin{equation*}
	      \abs{\int_\R u\varphi'}\le C\norm{\varphi}_{L^{p'}(\R)}
	      \quad\forall\varphi\in C_c^1(\R).
	      \end{equation*}
\end{enumerate}
\end{exercise}

\begin{exercise}
\label{exer:8.8}
\begin{enumerate}[label=\arabic*., leftmargin=2.2em, itemsep=0.3em]
	\item 设 $u\in W^{1,p}(0,1)$, $1<p<\infty$.证明若 $u(0)=0$, 则
	      $\frac{u(x)}{x}\in L^p(0,1)$ 且
	      \begin{equation*}
	      \left\|\frac{u(x)}{x}\right\|_{L^p(0,1)}
	      \le\frac{p}{p-1}\norm{u'}_{L^p(0,1)}.
	      \end{equation*}
	      [提示: 用问题 34 的 C 部分.]
	\item 反之, 设 $u\in W^{1,p}(0,1)$, $1\le p<\infty$, 且
	      $\frac{u(x)}{x}\in L^p(0,1)$.证明 $u(0)=0$.
	      [提示: 用反证法.]
	\item 设 $u(x)=(1+\abs{\log x})^{-1}$.验证 $u\in W^{1,1}(0,1)$,
	      $u(0)=0$, 但 $\frac{u(x)}{x}\notin L^1(0,1)$.
	\item 假设 $u\in W^{1,p}(0,1)$, $1\le p<\infty$, 且 $u(0)=0$.固定任一
	      函数 $\zeta\in C^\infty(\R)$, 使得 $\zeta(x)=0$ 对一切
	      $x\in(-\infty,1]$ 成立, 且 $\zeta(x)=1$ 对一切
	      $x\in[2,+\infty)$ 成立.令 $\zeta_n(x)=\zeta(nx)$,
	      $u_n(x)=\zeta_n(x)u(x)$, $n=1,2,\dots$.验证
	      $u_n\in W^{1,p}(0,1)$, 并证明当 $n\to\infty$ 时
	      $u_n\to u$ 在 $W^{1,p}(0,1)$ 中.
	      [提示: 分别考虑 $p=1$ 与 $p>1$ 两种情形.]
\end{enumerate}
\end{exercise}

\begin{exercise}
\label{exer:8.9}
令 $I=(0,1)$.
\begin{enumerate}[label=\arabic*., leftmargin=2.2em, itemsep=0.3em]
	\item 设 $u\in W^{2,p}(I)$, $1<p<\infty$.假设 $u(0)=u'(0)=0$.证明
	      $\frac{u(x)}{x^2}\in L^p(I)$ 且 $\frac{u'(x)}{x}\in L^p(I)$,
	      并且
	      \begin{equation}
	      \left\|\frac{u(x)}{x^2}\right\|_{L^p(I)}
	      +\left\|\frac{u'(x)}{x}\right\|_{L^p(I)}
	      \le C_p\norm{u''}_{L^p(I)}.
	      \tag{1}
	      \end{equation}
	      [提示: 参看习题 8.8.]
	\item 由此推出 $v(x)=\frac{u(x)}{x}\in W^{1,p}(I)$ 且 $v(0)=0$.
	\item 设 $u$ 如第 1 问.令 $u_n=\zeta_nu$, 其中 $\zeta_n$ 在习题 8.8
	      第 4 问中定义.验证 $u_n\in W^{2,p}(I)$, 且当 $n\to\infty$ 时
	      $u_n\to u$ 在 $W^{2,p}(I)$ 中.
	\item 更一般地, 设 $m\ge1$ 为整数, $1<p<\infty$.假设 $u\in X_m$, 其中
	      \begin{equation*}
	      X_m=\{u\in W^{m,p}(I);\ u(0)=Du(0)=\cdots=D^{m-1}u(0)=0\}.
	      \end{equation*}
	      证明 $\frac{u(x)}{x^m}\in L^p(I)$ 且 $\frac{u(x)}{x^{m-1}}\in X_1$.
	      [提示: 对 $m$ 用归纳法.]
	\item 假设 $u\in X_m$, 证明
	      \begin{equation*}
	      v=\frac{D^ju(x)}{x^{m-j-k}}\in X_k
	      \quad\forall j,k\ \text{为整数},\ j\ge0,\ k\ge1,\ j+k\le m-1.
	      \end{equation*}
	\item 设 $u$ 如第 4 问, $\zeta_n$ 如第 3 问.证明
	      $\zeta_nu\in W^{m,p}(I)$, 且当 $n\to\infty$ 时
	      $\zeta_nu\to u$ 在 $W^{m,p}(I)$ 中.
	\item 在 $p>1$ 的情形下给出第~
\ref{chap:8}~章注记 \ref{rem:8.18} 的一个证明.
	\item 现在设 $u\in W^{2,1}(I)$, $u(0)=u'(0)=0$.令
	      \begin{equation*}
	      v(x)=
	      \begin{cases}
	      \frac{u(x)}{x} & \text{若}\ x\in(0,1],\\
	      0 & \text{若}\ x=0.
	      \end{cases}
	      \end{equation*}
	      验证 $v\in C([0,1])$.证明 $v\in W^{1,1}(I)$.
	      [提示: 注意
	      $v'(x)=\frac{1}{x^2}\int_0^xu''(t)t\,dt$.]
	\item 构造一个函数 $u\in W^{2,1}(I)$ 的例子, 它满足
	      $u(0)=u'(0)=0$, 但 $\frac{u(x)}{x^2}\notin L^1(I)$ 且
	      $\frac{u'(x)}{x}\notin L^1(I)$.
	      [提示: 用习题 8.8 的第 3 问.]
	\item 设 $u$ 如第 8 问, $\zeta_n$ 如第 3 问.验证
	      $u_n=\zeta_nu\in W^{2,1}(I)$, 且当 $n\to\infty$ 时
	      $u_n\to u$ 在 $W^{2,1}(I)$ 中.
	\item 在 $m=2$, $p=1$ 的情形下给出第~
\ref{chap:8}~章注记 \ref{rem:8.18} 的一个
	      证明.
	\item 把第 8--11 问推广到 $W^{m,1}(I)$, $m\ge2$.
\end{enumerate}
[第 8 问的结果及其对 $m\ge2$ 的推广属于 Hernan Castro 与 Hui Wang.]
\end{exercise}

\begin{exercise}
\label{exer:8.10}
设 $I=(0,1)$.设 $u\in W^{1,p}(I)$, $1\le p<\infty$.我们的目标是证明
$u'=0$ 在集合 $E=\{x\in I;\ u(x)=0\}$ 上几乎处处成立.固定函数
$G\in C^1(\R,\R)$, 使得 $\abs{G(t)}\le1$ 对一切 $t\in\R$ 成立,
$\abs{G'(t)}\le C$ 对一切 $t\in\R$ 成立 (其中 $C$ 为某常数), 且
\begin{equation*}
G(t)=
\begin{cases}
1 & \text{若}\ t\ge1,\\
t & \text{若}\ \abs{t}\le1/2,\\
-1 & \text{若}\ t\le-1.
\end{cases}
\end{equation*}
令
\begin{equation*}
v_n(x)=\frac1nG(nu(x)).
\end{equation*}
\begin{enumerate}[label=\arabic*., leftmargin=2.2em, itemsep=0.3em]
	\item 验证当 $n\to\infty$ 时 $\norm{v_n}_{L^\infty(I)}\to0$.
	\item 证明 $v_n\in W^{1,p}(I)$ 并算出 $v_n'$.
	\item 由此推出 $\abs{v_n'}$ 被 $L^p(I)$ 中一个固定函数所控制.
	\item 证明当 $n\to\infty$ 时 $v_n'(x)\to f(x)$ 在 $I$ 上几乎处处成立,
	      并确定 $f$.
	      [提示: 分别考虑 $x\in E$ 与 $x\notin E$ 两种情形.]
	\item 由此推出 $v_n'\to f$ 在 $L^p(I)$ 中.
	\item 证明 $f=0$ 在 $I$ 上几乎处处成立, 并断言 $u'=0$ 在 $E$ 上几乎
	      处处成立.
\end{enumerate}
\end{exercise}

\begin{exercise}
\label{exer:8.11}
设 $F\in C(\R,\R)$, 并假设 $F\in C^1(\R\setminus\{0\})$, 且
$\abs{F'(t)}\le C$ 对一切 $t\in\R\setminus\{0\}$ 成立 (其中 $C$ 为某常数).
设 $1\le p<\infty$.目标是证明: 对每个 $u\in W^{1,p}(0,1)$,
$v=F(u)$ 属于 $W^{1,p}(0,1)$, 且
\begin{equation*}
v'(x)=
\begin{cases}
F'(u(x))u'(x) & \text{在}\ [u(x)\ne0]\ \text{上几乎处处},\\
0 & \text{在}\ [u(x)=0]\ \text{上几乎处处}.
\end{cases}
\end{equation*}
\begin{enumerate}[label=\arabic*., leftmargin=2.2em, itemsep=0.3em]
	\item 构造 $C^\infty(\R)$ 中的序列 $(F_n)$, 使得
	      $\norm{F_n'}_{L^\infty(\R)}\le C$ 对一切 $n$ 成立, 且
	      \begin{equation*}
	      \begin{cases}
	      F_n\to F & \text{在}\ \R\ \text{的紧子集上一致},\\
	      F_n'\to F' & \text{在}\ \R\setminus\{0\}\ \text{的紧子集上一致}.
	      \end{cases}
	      \end{equation*}
	\item 验证 $v_n=F_n(u)\in W^{1,p}(0,1)$, 且 $v_n'=F_n'(u)u'$.
	\item 证明 $v_n\to v=F(u)$ 在 $C([0,1])$ 中, 且 $v_n'$ 在 $L^p(0,1)$
	      中收敛到 $f$, 其中
	      \begin{equation*}
	      f(x)=
	      \begin{cases}
	      F'(u(x))u'(x) & \text{在}\ [u(x)\ne0]\ \text{上几乎处处},\\
	      0 & \text{在}\ [u(x)=0]\ \text{上几乎处处}.
	      \end{cases}
	      \end{equation*}
	      [提示: 应用控制收敛定理与习题 8.10.]
	\item 由此推出 $v\in W^{1,p}(0,1)$ 且 $v'=f$.证明
	      $v_n\to v$ 在 $W^{1,p}(0,1)$ 中.
	\item 设 $(u_k)$ 是 $W^{1,p}(0,1)$ 中的序列, 使得 $u_k\to u$ 在
	      $W^{1,p}(0,1)$ 中.证明 $F(u_k)\to F(u)$ 在 $W^{1,p}$ 中.
	      [提示: 应用定理~
\ref{thm:4.9} 并过渡到子序列 (仍记作 $u_k$), 可以假设
	      $u_k'\to u'$ 在 $(0,1)$ 上几乎处处成立, 且
	      $\abs{u_k'}\le g$ 对一切 $k$ 成立, 其中 $g\in L^p(0,1)$ 为某函数.
	      令 $w_k=F(u_k)$, 验证 $w_k'\to f$ 在 $(0,1)$ 上几乎处处成立,
	      其中 $f$ 如第 3 问所定义.由此推出 $w_k\to F(u)$ 在
	      $W^{1,p}(0,1)$ 中.断言原序列 $F(u_k)$ 收敛到 $F(u)$ 在
	      $W^{1,p}(0,1)$ 中.]
	\item 应用: 取 $F(t)=t^+=\max\{t,0\}$.验证
	      $u^+\in W^{1,p}(0,1)$ 对一切 $u\in W^{1,p}(0,1)$ 成立.算出
	      $(u^+)'$.
\end{enumerate}
\end{exercise}

\begin{exercise}
\label{exer:8.12}
设 $I=(0,1)$, $1\le p\le\infty$.令
\begin{equation*}
B_p=\{u\in W^{1,p}(I);\ \norm{u}_{L^p(I)}+\norm{u'}_{L^p(I)}\le1\}.
\end{equation*}
\begin{enumerate}[label=\arabic*., leftmargin=2.2em, itemsep=0.3em]
	\item 证明当 $1<p\le\infty$ 时 $B_p$ 是 $L^p(I)$ 的闭子集; 更确切
	      地说, $B_p$ 在 $L^p(I)$ 中紧.
	\item 证明 $B_1$ 不是 $L^1(I)$ 的闭子集.
\end{enumerate}
\end{exercise}

\begin{exercise}
\label{exer:8.13}
设 $1\le p<\infty$, $u\in W^{1,p}(\R)$.令
\begin{equation*}
D_hu(x)=\frac1h\bigl(u(x+h)-u(x)\bigr),\quad x\in\R,\ h>0.
\end{equation*}
证明当 $h\to0$ 时 $D_hu\to u'$ 在 $L^p(\R)$ 中.
[提示: 利用 $C_c^1(\R)$ 在 $W^{1,p}(\R)$ 中稠密这一事实.]
\end{exercise}

\begin{exercise}
\label{exer:8.14}
设 $u\in C^1((0,1))$.证明下列条件等价:
\begin{enumerate}[label=(\alph*), leftmargin=2.4em, itemsep=0.2em]
	\item $u\in W^{1,1}(0,1)$,
	\item $u'\in L^1(0,1)$ (其中 $u'$ 表示 $u$ 的通常意义的导数),
	\item $u\in BV(0,1)$ (关于 $BV$ 的定义见注记 \ref{rem:8.8}).
\end{enumerate}
验证函数 $u(x)=x\sin(1/x)$ ($0<x\le1$), $u(0)=0$, 在 $[0,1]$ 上连续,
但 $u\notin W^{1,1}(0,1)$.
\end{exercise}

\begin{exercise}
\label{exer:8.15}
\emph{Gagliardo--Nirenberg 不等式 (第一种形式).}设 $I=(0,1)$.
\begin{enumerate}[label=\arabic*., leftmargin=2.2em, itemsep=0.3em]
	\item 设 $1\le q<\infty$, $1<r\le\infty$.证明
	      \begin{equation}
	      \norm{u}_{L^\infty(I)}\le C\norm{u}_{W^{1,r}(I)}^a
	      \norm{u}_{L^q(I)}^{1-a}\quad\forall u\in W^{1,r}(I)
	      \tag{1}
	      \end{equation}
	      其中 $C=C(q,r)$ 为某常数, 而 $0<a<1$ 由
	      \begin{equation}
	      a\left(\frac1q+1-\frac1r\right)=\frac1q
	      \tag{2}
	      \end{equation}
	      定义.
	      [提示: 先考虑 $u(0)=0$ 的情形, 并写出
	      $G(u(x))=\int_0^xG'(u(t))u'(t)\,dt$, 其中
	      $G(t)=\abs{t}^{\alpha-1}t$, $\alpha=\frac1a$.当
	      $u(0)\ne0$ 时, 把上述不等式应用于 $\zeta u$, 其中
	      $\zeta\in C^1([0,1])$, $\zeta(0)=0$, 且当 $t\in[\frac12,1]$
	      时 $\zeta(t)=1$.]
	\item 设 $1\le q<p<\infty$, $1\le r\le\infty$.证明
	      \begin{equation}
	      \norm{u}_{L^p(I)}\le C\norm{u}_{W^{1,r}(I)}^b
	      \norm{u}_{L^q(I)}^{1-b}\quad\forall u\in W^{1,r}(I)
	      \tag{3}
	      \end{equation}
	      其中 $C=C(p,q,r)$ 为某常数, 而 $0<b<1$ 由
	      \begin{equation}
	      b\left(\frac1q+1-\frac1r\right)=\frac1q-\frac1p
	      \tag{4}
	      \end{equation}
	      定义.
	      [提示: 写出
	      $\norm{u}_{L^p(I)}^p=\int_I\abs{u}^q\abs{u}^{p-q}
	      \le\norm{u}_{L^q(I)}^q\norm{u}_{L^\infty}^{p-q}$, 并在
	      $r>1$ 时使用 (1).]
	\item 在与第 2 问相同的假设下, 证明
	      \begin{equation}
	      \norm{u}_{L^p(I)}\le C\norm{u'}_{L^r(I)}^b
	      \norm{u}_{L^q(I)}^{1-b}
	      \quad\forall u\in W^{1,r}(I)\ \text{且}\ \int_Iu=0.
	      \tag{5}
	      \end{equation}
\end{enumerate}
\end{exercise}

\begin{exercise}
\label{exer:8.16}
设 $E=L^p(0,1)$, $1\le p<\infty$.考虑由
\begin{equation*}
D(A)=\{u\in W^{1,p}(0,1),\ u(0)=0\},\qquad Au=u'
\end{equation*}
定义的无界算子 $A:D(A)\subset E\to E$.
\begin{enumerate}[label=\arabic*., leftmargin=2.2em, itemsep=0.3em]
	\item 验证 $D(A)$ 在 $E$ 中稠密, 且 $A$ 是闭算子 (即 $G(A)$ 在
	      $E\times E$ 中闭).
	\item 确定 $R(A)$ 与 $N(A)$.
	\item 算出 $A^\star$.验证当 $1<p<\infty$ 时 $D(A^\star)$ 在
	      $E^\star=L^{p'}(0,1)$ 中稠密, 但当 $p=1$ 时 $D(A^\star)$ 在
	      $E^\star=L^\infty(0,1)$ 中不稠密.
	\item 对由
	      \begin{equation*}
	      D(\tilde A)=W_0^{1,p}(0,1),\qquad \tilde Au=u'
	      \end{equation*}
	      定义的算子 $\tilde A$ 回答同样的问题.
\end{enumerate}
\end{exercise}

\begin{exercise}
\label{exer:8.17}
设 $H=L^2(0,1)$, 并设 $A:D(A)\subset H\to H$ 是由 $Au=u''$ 定义的无界
算子, 其定义域将在下面明确.在下列各情形下确定 $A^\star$, $D(A^\star)$,
$N(A)$ 与 $N(A^\star)$:
\begin{enumerate}[label=\arabic*., leftmargin=2.2em, itemsep=0.3em]
	\item $D(A)=\{u\in H^2(0,1);\ u(0)=u(1)=0\}$.
	\item $D(A)=H^2(0,1)$.
	\item $D(A)=\{u\in H^2(0,1);\ u(0)=u(1)=u'(0)=u'(1)=0\}$.
	\item $D(A)=\{u\in H^2(0,1);\ u(0)=u(1)\}$.
\end{enumerate}
对算子 $Au=u''-xu'$ 回答同样的问题.
\end{exercise}

\begin{exercise}
\label{exer:8.18}
验证映射 $u\mapsto u(0)$ 是从 $H^1(0,1)$ 到 $\R$ 的连续线性泛函.由此
推出存在唯一的 $v_0\in H^1(0,1)$, 使得
\begin{equation*}
u(0)=\int_0^1(u'v_0'+uv_0)\quad\forall u\in H^1(0,1).
\end{equation*}
证明 $v_0$ 是某个带适当边界条件的微分方程的解.显式算出 $v_0$.
[提示: 考虑 \S \ref{sec:8.4} 中的例 \ref{ex:8.4}.]
\end{exercise}

\begin{exercise}
\label{exer:8.19}
设 $H=L^2(0,1)$, 考虑由
\begin{equation*}
\varphi(u)=
\begin{cases}
\frac12\int_0^1u'^2 & \text{若}\ u\in H^1(0,1),\\
+\infty & \text{若}\ u\in L^2(0,1)\ \text{且}\ u\notin H^1(0,1)
\end{cases}
\end{equation*}
定义的函数 $\varphi:H\to(-\infty,+\infty]$.
\begin{enumerate}[label=\arabic*., leftmargin=2.2em, itemsep=0.3em]
	\item 验证 $\varphi$ 凸且下半连续.
	\item 对每个 $f\in H$ 计算 $\varphi^\star(f)$.
	      [提示: 先证明若 $\int_0^1f\ne0$, 则
	      $\varphi^\star(f)=+\infty$.然后假设
	      $\int_0^1f=0$, 并令 $F(x)=\int_0^xf(t)\,dt$.注意
	      $\int_0^1fv=-\int_0^1Fv'$ 对一切 $v\in H^1(0,1)$ 成立.]
\end{enumerate}
\end{exercise}

\begin{exercise}
\label{exer:8.20}
令
\begin{equation*}
V=\{v\in H^1(0,1);\ v(0)=0\}.
\end{equation*}
\begin{enumerate}[label=\arabic*., leftmargin=2.2em, itemsep=0.3em]
	\item 给定 $f\in L^2(0,1)$, 使得 $\frac1xf(x)\in L^2(0,1)$, 证明存在
	      唯一的 $u\in V$ 满足
	      \begin{equation}
	      \int_0^1u'(x)v'(x)\,dx+\int_0^1\frac{u(x)v(x)}{x^2}\,dx
	      =\int_0^1\frac{f(x)v(x)}{x^2}\,dx\quad\forall v\in V.
	      \tag{1}
	      \end{equation}
	      [提示: 用习题 8.8 的第 1 问.]
	\item 与 (1) 对应的极小化问题是什么?
\end{enumerate}
以下假设 $\frac{1}{x^2}f(x)\in L^2(0,1)$.
\begin{enumerate}[label=\arabic*., leftmargin=2.2em, itemsep=0.3em, start=3]
	\item 在 (1) 中取 $v(x)=\frac{u(x)}{(x+\varepsilon)^2}$, $\varepsilon>0$,
	      并推出
	      \begin{equation}
	      \int_0^1\left|\frac{d}{dx}\left(\frac{u(x)}{x+\varepsilon}\right)
	      \right|^2dx\le\int_0^1\frac{f(x)}{x^2}\frac{u(x)}{(x+\varepsilon)^2}\,dx.
	      \tag{2}
	      \end{equation}
	\item 证明 $\frac{u(x)}{x^2}\in L^2(0,1)$,
	      $\frac{u(x)}{x}\in H^1(0,1)$, 且 $\frac{u'(x)}{x}\in L^2(0,1)$.
	      [提示: 再次使用习题 8.8 的第 1 问, 并令 $\varepsilon\to0$.]
	\item 由此推出 $u\in H^2(0,1)$, 并且
	      \begin{equation}
	      \begin{cases}
	      -u''(x)+\frac{u(x)}{x^2}=\frac{f(x)}{x^2} & \text{在}\ (0,1)\ \text{上几乎处处},\\
	      u(0)=u'(0)=0\ \text{且}\ u'(1)=0.
	      \end{cases}
	      \tag{3}
	      \end{equation}
	\item 反之, 假设函数 $u\in H^2(0,1)$ 满足
	      $\frac{u(x)}{x^2}\in L^2(0,1)$ 与 (3).证明 (1) 成立.
\end{enumerate}
\end{exercise}

\begin{exercise}
\label{exer:8.21}
假设 $p\in C^1([0,1])$, $p(x)\ge\alpha>0$ 对一切 $x\in[0,1]$ 成立, 且
$q\in C([0,1])$, $q(x)\ge0$ 对一切 $x\in[0,1]$ 成立.设 $v_0\in C^2([0,1])$
是
\begin{equation}
\begin{cases}
-(pv_0')'+qv_0=0 & \text{在}\ [0,1]\ \text{上},\\
v_0(0)=1,\ v_0(1)=0
\end{cases}
\tag{1}
\end{equation}
的唯一解.令 $k_0=v_0'(0)$.
\begin{enumerate}[label=\arabic*., leftmargin=2.2em, itemsep=0.3em]
	\item 验证 $k_0\le-\alpha/p(0)$.
	      [提示: 用 $v_0$ 乘方程 (1) 并分部积分.利用
	      $1\le\norm{v_0'}_1\le\norm{v_0'}_2$.]
\end{enumerate}
现在研究问题
\begin{equation}
\begin{cases}
-(pu')'+qu=f & \text{在}\ (0,1)\ \text{上},\\
u'(0)=ku(0),\ u(1)=0,
\end{cases}
\tag{2}
\end{equation}
其中 $k\in\R$ 固定, $f\in L^2(0,1)$ 给定.
\begin{enumerate}[label=\arabic*., leftmargin=2.2em, itemsep=0.3em, start=2]
	\item 假设 $k=k_0$.证明
	      \begin{equation*}
	      [(2)\ \text{有解}\ u\in H^2(0,1)]\iff\left[\int_0^1fv_0=0\right].
	      \end{equation*}
	      $u$ 唯一吗?
	\item 现在假设 $k\ne k_0$.证明对每个 $f\in L^2(0,1)$, 问题 (2) 都有
	      唯一的解 $u\in H^2(0,1)$.
	      [提示: 利用习题 8.5, 找常数 $K$, 使得双线性形式
	      $a(u,v)=\int_0^1(pu'v'+quv+Kuv)+p(0)ku(0)v(0)$ 在 $H^1$ 上
	      强制.把 (2) 写成 $u=T(f+Ku)$ 的形式, 其中 $T$ 为某个适当的
	      紧算子.然后应用 Fredholm 二择一定理的断言 (c).]
\end{enumerate}
\end{exercise}

\begin{exercise}
\label{exer:8.22}
令
\begin{equation*}
K=\{\rho\in H^1(0,1);\ \rho\ge0\ \text{在}\ (0,1)\ \text{上},\
\sqrt\rho\in H^1(0,1)\}.
\end{equation*}
\begin{enumerate}[label=\arabic*., leftmargin=2.2em, itemsep=0.3em]
	\item 构造一个函数 $\rho\in H^1(0,1)$ 的例子, 它在 $(0,1)$ 上
	      $\rho\ge0$, 但 $\rho\notin K$.
	\item 给定 $\rho\in H^1(0,1)$, $\rho\ge0$ 在 $(0,1)$ 上, 令
	      \begin{equation*}
	      \mu=
	      \begin{cases}
	      \frac12\frac{\rho'}{\sqrt\rho} & \text{在}\ [\rho>0]\ \text{上},\\
	      0 & \text{在}\ [\rho=0]\ \text{上}.
	      \end{cases}
	      \end{equation*}
	      证明 $\rho\in K\iff\mu\in L^2$, 并且这时 $(\sqrt\rho)'=\mu$.
	      [提示: 考虑 $\rho_\varepsilon=\rho+\varepsilon$.]
	\item 证明 $K$ 是以 $0$ 为顶点的凸锥.
	\item 证明函数 $\rho\in K\mapsto\norm{(\sqrt\rho)'}_{L^2}^2$ 凸.
\end{enumerate}
\end{exercise}

\begin{exercise}
\label{exer:8.23}
设 $I=(0,1)$, 并固定常数 $k>0$.
\begin{enumerate}[label=\arabic*., leftmargin=2.2em, itemsep=0.3em]
	\item 给定 $f\in L^1(I)$, 证明存在唯一的 $u\in H_0^1(I)$ 满足
	      \begin{equation}
	      \int_Iu'v'+k\int_Iuv=\int_Ifv\quad\forall v\in H_0^1(I).
	      \tag{1}
	      \end{equation}
	\item 证明 $u\in W^{2,1}(I)$.
	\item 证明
	      \begin{equation*}
	      \norm{u}_{L^1(I)}\le\frac1k\norm{f}_{L^1(I)}.
	      \end{equation*}
	      [提示: 固定函数 $\gamma\in C^1(\R,\R)$, 使得 $\gamma'(t)\ge0$
	      对一切 $t\in\R$ 成立, $\gamma(0)=0$, $\gamma(t)=+1$ 对一切
	      $t\ge1$ 成立, $\gamma(t)=-1$ 对一切 $t\le-1$ 成立.在 (1) 中
	      取 $v=\gamma(nu)$ 并令 $n\to\infty$.]
	\item 现在假设 $f\in L^p(I)$, $1<p<\infty$.证明存在与 $k$ 及 $p$
	      无关的常数 $\delta>0$, 使得
	      \begin{equation*}
	      \norm{u}_{L^p(I)}\le\frac{1}{k+\delta/pp'}\norm{f}_{L^p(I)}.
	      \end{equation*}
	      [提示: 当 $2\le p<\infty$ 时, 在 (1) 中取
	      $v=\gamma(u)$, 其中 $\gamma(t)=\abs{t}^{p-1}\operatorname{sign}t$.
	      当 $1<p<2$ 时, 用对偶性.]
	\item 证明若 $f\in L^\infty(I)$, 则
	      \begin{equation*}
	      \norm{u}_{L^\infty(I)}\le C_k\norm{f}_{L^\infty(I)},
	      \end{equation*}
	      并求出最佳常数 $C_k$.
	      [提示: 显式算出对应于 $f\equiv1$ 的 (1) 的解 $u$.]
\end{enumerate}
\end{exercise}

\begin{exercise}
\label{exer:8.24}
设 $I=(0,1)$.
\begin{enumerate}[label=\arabic*., leftmargin=2.2em, itemsep=0.3em]
	\item 证明对每个 $\varepsilon>0$ 存在常数 $C_\varepsilon$, 使得
	      \begin{equation*}
	      \abs{u(1)}^2\le\varepsilon\norm{u'}_{L^2(I)}^2
	      +C_\varepsilon\norm{u}_{L^2(I)}^2\quad\forall u\in H^1(I).
	      \end{equation*}
	      [提示: 用习题 8.5, 或者直接写出
	      $u^2(1)=u^2(x)+2\int_x^1u(t)u'(t)\,dt$.]
	\item 证明若常数 $k>0$ 足够大, 则对每个 $f\in L^2(I)$ 存在唯一的
	      $u\in H^2(I)$ 满足
	      \begin{equation}
	      \begin{cases}
	      -u''+ku=f & \text{在}\ (0,1)\ \text{上},\\
	      u'(0)=0\ \text{且}\ u'(1)=u(1).
	      \end{cases}
	      \tag{1}
	      \end{equation}
	      问题 (1) 的弱表述是什么? 相应的极小化问题是什么?
	\item 假设 $k$ 足够大.设 $T$ 为算子 $T:f\mapsto u$, 其中 $u$ 是 (1)
	      的解.证明 $T$ 是 $L^2(I)$ 中的自伴紧算子.
	\item 由此推出存在 $\R$ 中的序列 $(\lambda_n)$, 满足
	      $\abs{\lambda_n}\to\infty$, 以及 $C^2(I)$ 中的函数序列 $(u_n)$,
	      使得 $\norm{u_n}_{L^2(I)}=1$ 且
	      \begin{equation*}
	      \begin{cases}
	      -u_n''=\lambda_nu_n & \text{在}\ (0,1)\ \text{上},\\
	      u_n'(0)=0\ \text{且}\ u_n'(1)=u_n(1).
	      \end{cases}
	      \end{equation*}
	      证明 $\lambda_n\to+\infty$.
	\item 设 $\Lambda$ 为一切使存在 $u\not\equiv0$ 满足
	      \begin{equation*}
	      \begin{cases}
	      -u''=\lambda u & \text{在}\ (0,1)\ \text{上},\\
	      u'(0)=0\ \text{且}\ u'(1)=u(1)
	      \end{cases}
	      \end{equation*}
	      的 $\lambda\in\R$ 的集合.确定 $\Lambda$ 中的正元素.证明 $\Lambda$
	      中恰有一个负值 (记作 $\lambda_0$).
	      [提示: 不要试图显式算出 $\Lambda$; 而是利用两条曲线的交点.]
	\item 当 $k=\abs{\lambda_0}$ 时, 第 2 问中的情形如何?
\end{enumerate}
\end{exercise}

\begin{exercise}
\label{exer:8.25}
设 $I=(0,2)$, $V=H^1(I)$.考虑双线性形式
\begin{equation*}
a(u,v)=\int_0^2u'(t)v'(t)\,dt
+\left(\int_0^1u(t)\,dt\right)\left(\int_0^1v(t)\,dt\right).
\end{equation*}
\begin{enumerate}[label=\arabic*., leftmargin=2.2em, itemsep=0.3em]
	\item 验证 $a(u,v)$ 是连续对称双线性形式, 且 $a(u,u)=0$ 蕴含 $u=0$.
	\item 证明 $a$ 强制.
	      [提示: 用反证法, 假设存在 $H^1(I)$ 中的序列 $(u_n)$, 使得
	      $a(u_n,u_n)\to0$ 且 $\norm{u_n}_{H^1}=1$.取子序列
	      $(u_{n_k})$, 使它在 $H^1(I)$ 中弱收敛且在 $L^2(I)$ 中强收敛到
	      某个极限 $u$.证明 $u=0$.]
	\item 由此推出对每个 $f\in L^2(I)$ 存在唯一的 $u\in H^1(I)$ 满足
	      \begin{equation}
	      a(u,v)=\int_0^2fv\quad\forall v\in H^1(I).
	      \tag{1}
	      \end{equation}
	      相应的极小化问题是什么?
	\item 证明 (1) 的解属于 $H^2(I)$ (特别地 $u\in C^1(I)$).确定 $u$
	      所满足的方程与边界条件.
	      [提示: 方便的做法是令 $g=(\int_0^1u)\chi$, 其中 $\chi$ 是
	      $(0,1)$ 的特征函数.]
	\item 假设 $f\in C(I)$, 设 $u$ 是 (1) 的解.证明对每个 $p<\infty$,
	      $u$ 属于 $W^{2,p}(I)$.证明 $u\in C^2(I)$ 当且仅当
	      $\int_If=0$.
	\item 当 $f$ 为常数时, 显式确定 (1) 的解 $u$.
	\item 令 $u=Tf$, 其中 $u$ 是 (1) 的解, $f\in L^2(I)$.验证 $T$ 是
	      从 $L^2(I)$ 到自身的自伴紧算子.
	\item 研究 $T$ 的特征值.
\end{enumerate}
\end{exercise}

\begin{exercise}
\label{exer:8.26}
从 Hilbert 空间 $H$ 到自身的具有界线性算子 $S$ 称为\emph{非负}的, 记作
$S\ge0$, 若
\begin{equation*}
(Sf,f)\ge0\quad\forall f\in H.
\end{equation*}
令 $I=(0,1)$.假设 $p\in C^1(I)$ 与 $q\in C(I)$ 满足
\begin{equation*}
p(x)\ge\alpha>0\quad\text{且}\quad q(x)\ge\alpha>0\quad\forall x\in I.
\end{equation*}
回忆给定 $f\in H=L^2(I)$, 方程
\begin{equation}
-(pu')'+qu=f\quad\text{在}\ I\ \text{上}
\tag{1}
\end{equation}
带 Dirichlet 边界条件
\begin{equation}
u(0)=u(1)=0
\tag{2}
\end{equation}
有唯一的解 $u\in H^2(I)$.这个解记作 $u_D=S_Df$, 其中 $S_D$ 看作从 $H$ 到
自身的具有界线性算子.类似地, (1) 带 Neumann 边界条件
\begin{equation}
u'(0)=u'(1)=0
\tag{3}
\end{equation}
有唯一的解 $u\in H^2(I)$.这个解记作 $u_N=S_Nf$, 其中 $S_N$ 也看作从 $H$
到自身的具有界线性算子.
\begin{enumerate}[label=\arabic*., leftmargin=2.2em, itemsep=0.3em]
	\item 证明 $S_D\ge0$ 且 $S_N\ge0$.
	\item 回忆与 Dirichlet 条件及 Neumann 条件相应的极小化原理.由此推出
	      \begin{equation}
	      \frac12\int_Ip(u_N')^2+qu_N^2-\int_Ifu_N
	      \le\frac12\int_Ip(u_D')^2+qu_D^2-\int_Ifu_D.
	      \tag{4}
	      \end{equation}
	      [提示: 利用 $H_0^1\subset H^1$.]
	\item 证明 $S_N-S_D\ge0$.
	      [提示: 利用 (4), 并把 (1) 分别乘以 $u_D$ 与 $u_N$.]
\end{enumerate}
给定实数 $k\ge0$, 考虑与边界条件
\begin{equation}
p(0)u'(0)=ku(0)\quad\text{且}\quad u(1)=0
\tag{5}
\end{equation}
相联系的方程 (1).
\begin{enumerate}[label=\arabic*., leftmargin=2.2em, itemsep=0.3em, start=4]
	\item 验证问题 (1) 与 (5) 有唯一的解, 记作 $u_k=S_kf$.相应的
	      极小化原理是什么?
	\item 证明 $S_k\ge0$.
	\item 设 $k_1\ge k_2\ge0$.证明 $S_{k_2}-S_{k_1}\ge0$.
\end{enumerate}
\end{exercise}

\begin{exercise}
\label{exer:8.27}
设 $I=(-1,+1)$.考虑 $H_0^1(I)$ 上由
\begin{equation*}
a(u,v)=\int_I(u'v'+uv-\lambda u(0)v)
\end{equation*}
定义的双线性形式, 其中 $\lambda\in\R$ 固定.
\begin{enumerate}[label=\arabic*., leftmargin=2.2em, itemsep=0.3em]
	\item 验证 $a(u,v)$ 是 $H_0^1(I)$ 上的连续双线性形式.
	\item 证明若 $\abs{\lambda}<\sqrt2$, 则双线性形式 $a$ 强制.
	      [提示: 验证 $\abs{u(0)}\le\norm{u'}_{L^2}$ 对一切
	      $u\in H_0^1(I)$ 成立.]
	\item 由此推出若 $\abs{\lambda}<\sqrt2$, 则对每个 $f\in L^2(I)$ 存在
	      唯一的解 $u\in H^2(I)\cap H_0^1(I)$, 满足
	      \begin{equation}
	      \begin{cases}
	      -u''+u-\lambda u(0)=f & \text{在}\ I\ \text{上},\\
	      u(-1)=u(1)=0.
	      \end{cases}
	      \tag{1}
	      \end{equation}
	\item 证明存在唯一的值 $\lambda=\lambda_0\in\R$ (要求显式确定), 使得
	      问题
	      \begin{equation}
	      \begin{cases}
	      -u''+u=\lambda u(0) & \text{在}\ I\ \text{上},\\
	      u(-1)=u(1)=0
	      \end{cases}
	      \tag{2}
	      \end{equation}
	      有解 $u\not\equiv0$.
	      [提示: 方便的做法是引入问题
	      \begin{equation}
	      \begin{cases}
	      -\varphi''+\varphi=1 & \text{在}\ I\ \text{上},\\
	      \varphi(-1)=\varphi(1)=0
	      \end{cases}
	      \tag{3}
	      \end{equation}
	      的唯一解 $\varphi$.显式算出 $\varphi$.]
	\item 证明若 $\lambda\ne\lambda_0$, 则对每个 $f\in L^2(I)$, (1) 都有
	      唯一的解 $u\in H^2(I)\cap H_0^1(I)$.
	      [提示: 考虑线性算子 $S:g\mapsto v$, 其中 $g\in L^2(I)$, 而
	      $v\in H^2(I)\cap H_0^1(I)$ 是问题
	      \begin{equation}
	      \begin{cases}
	      -v''+v=g & \text{在}\ I\ \text{上},\\
	      v(-1)=v(1)=0
	      \end{cases}
	      \tag{4}
	      \end{equation}
	      的唯一解.把 (1) 写成 $u-\lambda u(0)\varphi=Sf$.]
	\item 当 $\lambda=\lambda_0$ 时, 完整分析问题 (1).
	      [提示: 找出 $Sf$ 所应满足的简单充要条件, 使问题 (1) 有解.]
\end{enumerate}
\end{exercise}

\begin{exercise}
\label{exer:8.28}
设 $H=L^2(0,1)$ 配以其通常的标量积.考虑由
\begin{equation*}
(Tf)(x)=x\int_x^1f(t)\,dt+\int_0^xtf(t)\,dt,
\quad 0\le x\le1,
\end{equation*}
定义的算子 $T:H\to H$.
\begin{enumerate}[label=\arabic*., leftmargin=2.2em, itemsep=0.3em]
	\item 验证 $T$ 是有界算子.
	\item 证明 $T$ 是紧算子.
	\item 证明 $T$ 自伴.
	\item 证明 $(Tf,f)\ge0$ 对一切 $f\in H$ 成立, 且 $(Tf,f)=0$ 蕴含
	      $f=0$.
	\item 令 $u=Tf$.证明 $u\in H^2(0,1)$ 并算出 $u''$.验证
	      $u(0)=u'(1)=0$.
	\item 确定 $T$ 的谱与特征值.仔细考察 $\lambda=0$ 的情形.
\end{enumerate}
以下令
\begin{equation*}
e_k(x)=\sqrt2\sin\left(\left(k+\frac12\right)\pi x\right),
\quad k=0,1,2,\dots.
\end{equation*}
\begin{enumerate}[label=\arabic*., leftmargin=2.2em, itemsep=0.3em, start=7]
	\item 验证 $(e_k)$ 是 $H$ 的标准正交基.
	      [提示: 用第 6 问.]
	\item 由此推出由
	      \begin{equation*}
	      \tilde e_k(x)=\sqrt2\cos\left(\left(k+\frac12\right)\pi x\right),
	      \quad k=0,1,2,\dots
	      \end{equation*}
	      定义的序列 $(\tilde e_k)$ 也是 $H$ 的标准正交基.
	      [提示: 考虑 $e_k(1-x)$.]
\end{enumerate}
给定 $f\in H$, 用 $(\alpha_k(f))$ 表示 $f$ 在基 $(e_k)$ 中的分量.
\begin{enumerate}[label=\arabic*., leftmargin=2.2em, itemsep=0.3em, start=9]
	\item 计算下列函数的 $\alpha_k(f)$:
	      \begin{enumerate}[label=(\alph*), leftmargin=2.2em, itemsep=0.2em]
	      	\item $f_1(x)=\chi_{[a,b]}(x)=
	      	      \begin{cases}1 & \text{若}\ x\in[a,b],\\
	      	      0 & \text{若}\ x\notin[a,b],\end{cases}$
	      	      其中 $0\le a<b\le1$.
	      	\item $f_2(x)=x$.
	      	\item $f_3(x)=x^2$.
	      \end{enumerate}
\end{enumerate}
最后, 我们打算利用分量 $\alpha_k(f)$ 来刻画属于 $H^1(0,1)$ 的函数
$f\in L^2(0,1)$.
\begin{enumerate}[label=\arabic*., leftmargin=2.2em, itemsep=0.3em, start=10]
	\item 假设 $f\in H^1(0,1)$.证明存在常数 $a\in\R$ (依赖于 $f$),
	      使得 $(k\alpha_k(f)+a)\in\ell^2$, 即
	      \begin{equation}
	      \sum_{k=0}^\infty\abs{k\alpha_k(f)+a}^2<\infty.
	      \tag{1}
	      \end{equation}
	      [提示: 在计算 $\alpha_k(f)$ 时作一次分部积分.]
	\item 反之, 假设 $f\in L^2(0,1)$, 且对某个 $a\in\R$, (1) 成立.
	      证明 $f\in H^1(0,1)$.
	      [提示: 令 $\tilde f=f+\frac{a\pi}{\sqrt2}$, 并令
	      $\tilde f_n(x)=\sum_{k=0}^{n-1}\alpha_k(\tilde f)e_k(x)$.验证
	      当 $n\to\infty$ 时 $\norm{\tilde f_n'}_{L^2}$ 保持有界.]
\end{enumerate}
\end{exercise}

\begin{exercise}
\label{exer:8.29}
令
\begin{equation*}
a(u,v)=\int_0^1(u'v'+uv)+(u(1)-u(0))(v(1)-v(0))
\quad\forall u,v\in H^1(0,1).
\end{equation*}
\begin{enumerate}[label=\arabic*., leftmargin=2.2em, itemsep=0.3em]
	\item 验证 $a$ 是 $H^1(0,1)$ 上的连续强制双线性形式.
	\item 由此推出对每个 $f\in L^2(0,1)$, 存在唯一的 $u\in H^2(0,1)$ 满足
	      \begin{equation}
	      a(u,v)=\int_0^1fv\quad\forall v\in H^1(0,1).
	      \tag{1}
	      \end{equation}
	\item 验证 $u$ 满足
	      \begin{equation}
	      \begin{cases}
	      -u''+u=f & \text{在}\ (0,1)\ \text{上},\\
	      u'(0)=u(0)-u(1),\\
	      u'(1)=u(0)-u(1).
	      \end{cases}
	      \tag{2}
	      \end{equation}
	      证明 (2) 的任何解 $u\in H^2(0,1)$ 都满足 (1).
\end{enumerate}
设 $T:L^2(0,1)\to L^2(0,1)$ 为由 $Tf=u$ 定义的算子.
\begin{enumerate}[label=\arabic*., leftmargin=2.2em, itemsep=0.3em, start=4]
	\item 验证 $T$ 自伴且紧.
	\item 证明若 $f\ge0$ 在 $(0,1)$ 上几乎处处成立, 则 $u=Tf\ge0$ 在
	      $(0,1)$ 上.
	\item 验证 $(Tf,f)_{L^2}\ge0$ 对一切 $f\in L^2(0,1)$ 成立.
	\item 确定 $\mathrm{EV}(T)$.
\end{enumerate}
\end{exercise}

\begin{exercise}
\label{exer:8.30}
设 $k\in\R$, $k\ne1$, 考虑空间
\begin{equation*}
V=\{v\in H^1(0,1);\ v(0)=kv(1)\}
\end{equation*}
与双线性形式
\begin{equation*}
a(u,v)=\int_0^1(u'v'+uv)
-\left(\int_0^1u\right)\left(\int_0^1v\right)\quad\forall u,v\in V.
\end{equation*}
\begin{enumerate}[label=\arabic*., leftmargin=2.2em, itemsep=0.3em]
	\item 验证 $V$ 是 $H^1(0,1)$ 的闭子空间.以下 $V$ 配以由 $H^1$ 标量积
	      诱导的 Hilbert 结构.
	\item 证明 $a$ 是 $V$ 上的连续强制双线性形式.
	      [提示: 证明存在常数 $C$, 使得
	      $\norm{v}_{L^\infty(0,1)}\le C\norm{v'}_{L^2(0,1)}$ 对一切
	      $v\in V$ 成立.]
	\item 由此推出对每个 $f\in L^2(0,1)$ 存在问题
	      \begin{equation}
	      u\in V\quad\text{且}\quad
	      a(u,v)=\int_0^1fv\quad\forall v\in V
	      \tag{1}
	      \end{equation}
	      的唯一解.
	\item 证明 (1) 的解 $u$ 属于 $H^2(0,1)$, 并且满足
	      \begin{equation}
	      \begin{cases}
	      -u''+u-\int_0^1u=f & \text{在}\ (0,1)\ \text{上},\\
	      u(0)=ku(1)\ \text{且}\ u'(1)=ku'(0).
	      \end{cases}
	      \tag{2}
	      \end{equation}
	\item 反之, 证明任何满足 (2) 的函数 $u\in H^2(0,1)$ 都是 (1) 的解.
	\item 设 $k_n\in\R$, $k_n\ne1$ 对一切 $n$ 成立, 是收敛到 $k\ne1$ 的
	      序列.令
	      \begin{equation*}
	      V_n=\{v\in H^1(0,1);\ v(0)=k_nv(1)\}.
	      \end{equation*}
	      给定 $f\in L^2(0,1)$, 设 $u_n$ 是
	      \begin{equation}
	      u_n\in V_n\quad\text{且}\quad
	      a(u_n,v)=\int_0^1fv\quad\forall v\in V_n
	      \tag{$1_n$}
	      \end{equation}
	      的解.证明当 $n\to\infty$ 时 $u_n\to u$ 在 $H^1(0,1)$ 中, 其中
	      $u$ 是 (1) 的解.由此推出 $u_n\to u$ 在 $H^2(0,1)$ 中.
	      [提示: 验证由
	      \begin{equation*}
	      u^{(n)}(x)=u(x)+\frac{k-k_n}{k_n-1}u(1)
	      \end{equation*}
	      定义的函数 $u^{(n)}$ 属于 $V_n$, 并在 $H^1(0,1)$ 中收敛到 $u$.
	      证明
	      $a(u_n-u^{(n)},u_n-u^{(n)})=(k_n-k)u'(0)(u_n(1)-u^{(n)}(1))$.]
	\item 若 $k_n$ 收敛到 $1$, 序列 $(u_n)$ 的情形如何?
	\item 考虑由 $Tf=u$ 定义的算子 $T:L^2(0,1)\to L^2(0,1)$, 其中 $u$
	      是 (1) 的解.证明 $T$ 自伴且紧.研究 $\mathrm{EV}(T)$.
\end{enumerate}
\end{exercise}

\begin{exercise}
\label{exer:8.31}
考虑 $(0,1)$ 上带 Neumann 边界条件 $u'(0)=u'(1)=0$ 的 Sturm--Liouville
算子 $Au=-u''+u$.
\begin{enumerate}[label=\arabic*., leftmargin=2.2em, itemsep=0.3em]
	\item 算出 $A$ 的特征值及相应的特征函数.
	\item 给定 $f\in L^2$ 满足 $\int_0^1f=0$, 设 $u$ 是
	      \begin{equation*}
	      \begin{cases}
	      -u''+u=f & \text{在}\ (0,1)\ \text{上},\\
	      u'(0)=u'(1)=0
	      \end{cases}
	      \end{equation*}
	      的解.证明
	      \begin{equation*}
	      \norm{u}_{L^2(0,1)}\le\frac{1}{(1+\pi^2)}\norm{f}_{L^2(0,1)}.
	      \end{equation*}
	      [提示: 应用问题 49 的第 6 问.]
	\item 设 $(u_n)$ 是由
	      \begin{equation*}
	      \begin{cases}
	      -u_n''+u_n=u_{n-1} & \text{在}\ (0,1)\ \text{上},\\
	      u_n'(0)=u_n'(1)=0
	      \end{cases}
	      \end{equation*}
	      归纳定义的序列, 从某个 $u_0\in L^2(0,1)$ 出发.证明
	      \begin{equation*}
	      \norm{u_n-\bar u_0}_{L^2(0,1)}
	      \le\frac{1}{(1+\pi^2)^n}\norm{u_0-\bar u_0}_{L^2(0,1)}
	      \quad\forall n,
	      \end{equation*}
	      其中 $\bar u_0=\int_0^1u_0$.
\end{enumerate}
\end{exercise}

\begin{exercise}
\label{exer:8.32}
令
\begin{equation*}
V=\{v\in H^1(0,1);\ v(1)=0\}.
\end{equation*}
又令
\begin{equation*}
H=\{f\ \text{在}\ (0,1)\ \text{上可测, 且}\ xf(x)\in L^2(0,1)\}.
\end{equation*}
\begin{enumerate}[label=\arabic*., leftmargin=2.2em, itemsep=0.3em]
	\item 证明 $H$ 配以标量积
	      \begin{equation*}
	      (f,g)_H=\int_0^1f(x)g(x)x^2\,dx
	      \end{equation*}
	      是 Hilbert 空间.
	\item 给定 $f\in H$ 与 $\varepsilon>0$, 验证存在唯一的 $u\in V$ 满足
	      \begin{equation*}
	      \int_0^1u'(x)v'(x)(x^2+\varepsilon)\,dx
	      +\int_0^1u(x)v(x)x^2\,dx
	      =\int_0^1f(x)v(x)x^2\,dx\quad\forall v\in V.
	      \end{equation*}
	      这个 $u$ 记作 $u_\varepsilon$.
	\item 证明 $u_\varepsilon\in H^2(0,1)$, 并且满足
	      \begin{equation}
	      \begin{cases}
	      -((x^2+\varepsilon)u_\varepsilon')'+x^2u_\varepsilon=x^2f
	      & \text{在}\ (0,1)\ \text{上},\\
	      u_\varepsilon'(0)=0\ \text{且}\ u_\varepsilon(1)=0.
	      \end{cases}
	      \tag{1}
	      \end{equation}
	\item 由此推出
	      \begin{equation}
	      \abs{u_\varepsilon'(x)}\le\frac1x\int_0^xt\abs{f(t)-u_\varepsilon(t)}\,dt
	      \quad\forall x\in[0,1].
	      \tag{2}
	      \end{equation}
	\item 证明当 $\varepsilon\to0$ 时, $xu_\varepsilon(x)$ 与
	      $u_\varepsilon'(x)$ 在 $L^2(0,1)$ 中保持有界.
	      [提示: 用习题 8.8 的第 1 问.]
	\item 令 $\varepsilon\to0$ 取极限, 并断言存在唯一的 $u\in V$ 满足
	      \begin{equation}
	      \int_0^1(u'(x)v'(x)+u(x)v(x))x^2\,dx
	      =\int_0^1f(x)v(x)x^2\,dx\quad\forall v\in V.
	      \tag{3}
	      \end{equation}
\end{enumerate}
考虑由 $Tf=u$ 定义的算子 $T:H\to H$, 其中 $u$ 是 (3) 的解.
\begin{enumerate}[label=\arabic*., leftmargin=2.2em, itemsep=0.3em, start=7]
	\item 验证 $T$ 是从 $H$ 到自身的自伴紧算子.
	\item 确定 $T$ 的全部特征值.
	      [提示: 寻找形如 $\frac1x\sin kx$ 的特征函数, 其中 $k$ 取适当的值.]
\end{enumerate}
\end{exercise}

\begin{exercise}
\label{exer:8.33}
\emph{特征值的单性.}考虑 $I=(0,1)$ 上的 Sturm--Liouville 算子
\begin{equation*}
Au=-(pu')'+qu,
\end{equation*}
其中 $p\in C^1([0,1])$, $p\ge\alpha>0$ 在 $I$ 上, $q\in C([0,1])$.(不作
进一步的假设; 特别地, 相应的双线性形式
$a(u,v)=\int_0^1(pu'v'+quv)$ 不必强制.)令
\begin{equation*}
N=\{u\in H^2(0,1);\ a(u,v)=0\ \forall v\in H_0^1(0,1)\}.
\end{equation*}
\begin{enumerate}[label=\arabic*., leftmargin=2.2em, itemsep=0.3em]
	\item 证明存在唯一的 $U\in N$ 满足 $U(0)=1$ 且 $U'(0)=0$.
	      [提示: 把方程 $Au=0$ 写成一阶微分方程组, 并应用定理  \ref{thm:7.3}
	      (Cauchy--Lipschitz--Picard).]
	\item 证明 $\dim N=2$.
	      [提示: 考虑满足 $V(0)=0$ 与 $V'(0)=1$ 的唯一 $V\in N$.然后把
	      任何 $u\in N$ 写成 $u=u(0)U+u'(0)V$.]
	\item 令 $N_0=\{u\in N;\ u(0)=0\}$.验证 $\dim N_0=1$.
	\item 令 $N_{00}=\{u\in N;\ u(0)=u(1)=0\}$.证明
	      \begin{equation*}
	      \dim N_{00}=1\iff 0\ \text{是}\ A\ \text{带零 Dirichlet 条件的特征值}.
	      \end{equation*}
	      否则 $\dim N_{00}=0$.
	\item 由此推出 $A$ 带零 Dirichlet 条件的一切特征值都是单的.(作为对照,
	      与周期边界条件相联系的特征值可以有重数 2; 见习题 8.34.)
	\item 把上述结果推广到把条件 $p\in C^1([0,1])$ 换成 $p\in C([0,1])$,
	      并把 $N$ 换成
	      \begin{equation*}
	      \tilde N=\{u\in H^1(0,1);\ a(u,v)=0\ \forall v\in H_0^1(0,1)\}
	      \end{equation*}
	      的情形.
\end{enumerate}
\end{exercise}

\begin{exercise}
\label{exer:8.34}
考虑问题
\begin{equation}
\begin{cases}
-u''+u=f & \text{在}\ (0,1)\ \text{上},\\
u(0)=u(1)\ \text{且}\ u'(1)-u'(0)=k,
\end{cases}
\tag{1}
\end{equation}
其中 $k\in\R$ 与 $f(x)$ 给定.
\begin{enumerate}[label=\arabic*., leftmargin=2.2em, itemsep=0.3em]
	\item 求出问题 (1) 的弱表述.
	\item 证明对每个 $f\in L^2(0,1)$ 与每个 $k\in\R$, (1) 都有唯一的弱解
	      $u\in H^1(0,1)$.相应的极小化问题是什么?
	\item 证明弱解 $u$ 属于 $H^2(0,1)$ 且满足 (1).验证若
	      $f\in C([0,1])$, 则 $u\in C^2([0,1])$.
	\item 证明若 $f\le0$ 在 $(0,1)$ 上且 $k\le0$, 则 $u\le0$ 在 $(0,1)$ 上.
	\item 取 $k=0$, 考虑由 $Tf=u$ 定义的算子
	      $T:L^2(0,1)\to L^2(0,1)$.验证 $T$ 是自伴紧算子.算出 $T$ 的
	      特征值, 并证明除第一个之外, 每个特征值 $\lambda$ 的重数都是 2
	      (即 $\dim N(T-\lambda I)=2$).
\end{enumerate}
\emph{注.} 注意作为对照, 带 Dirichlet 边界条件 ($u(0)=u(1)=0$) 的
Sturm--Liouville 算子的每个特征值都是单的 (即相应的特征空间维数为 1).
\end{exercise}

\begin{exercise}
\label{exer:8.35}
固定两个函数 $a,b\in C([0,1])$, 考虑问题
\begin{equation}
\begin{cases}
-u''+au'+bu=f & \text{在}\ (0,1)\ \text{上},\\
u(0)=u(1)=0,
\end{cases}
\tag{1}
\end{equation}
其中 $f\in L^2(0,1)$.

给定 $g\in L^2(0,1)$, 设 $v\in H^2(0,1)$ 是
\begin{equation}
\begin{cases}
-v''=g & \text{在}\ (0,1)\ \text{上},\\
v(0)=v(1)=0
\end{cases}
\tag{2}
\end{equation}
的唯一解.令 $Sg=v$, 于是 $S:L^2(0,1)\to H^2(0,1)$.
\begin{enumerate}[label=\arabic*., leftmargin=2.2em, itemsep=0.3em]
	\item 验证问题 (1) 等价于
	      \begin{equation}
	      \begin{cases}
	      u\in H^1(0,1),\\
	      u=S(f-au'-bu).
	      \end{cases}
	      \tag{3}
	      \end{equation}
\end{enumerate}
考虑由
\begin{equation*}
Tu=-S(au'+bu),\quad u\in H^1(0,1)
\end{equation*}
定义的算子 $T:H^1(0,1)\to H^1(0,1)$.
\begin{enumerate}[label=\arabic*., leftmargin=2.2em, itemsep=0.3em, start=2]
	\item 证明 $T$ 是紧算子.
	\item 证明问题 (1) 对每个 $f\in L^2(0,1)$ 都有解, 当且仅当 (1) 在
	      $f=0$ 时唯一的解是 $u=0$.
	\item 假设 $b\ge0$ 在 $(0,1)$ 上.证明 (1) 在 $f=0$ 时唯一的解是
	      $u=0$.
	      [提示: 固定常数 $k>0$, 使得 $k^2-ka-b>0$ 在 $[0,1]$ 上.令
	      $u_\varepsilon(x)=u(x)+\varepsilon e^{kx}$, $\varepsilon>0$.对
	      $u_\varepsilon$ 施行最大值原理的``古典方法''; 见第~
\ref{chap:8}~章注记
	      \ref{rem:8.26}.由此推出 $u(x)\le\varepsilon e^k$ 对一切
	      $x\in[0,1]$, 一切 $\varepsilon>0$ 成立.]
	      断言对每个 $f\in L^2(0,1)$, 问题 (1) 都有唯一的解
	      $u\in H^2(0,1)$.
	\item 验证一般情形下 (对 $a$ 与 $b$ 不作任何假设), 问题 (1) 在 $f=0$
	      时的解空间维数为 0 或 1.若这个维数为 1, 证明问题 (1) 有解当且
	      仅当 $\int_0^1f\varphi_0=0$ 对某个待定的函数
	      $\varphi_0\not\equiv0$ 成立.
	      [提示: 用习题 8.33 与 Fredholm 二择一定理.]
\end{enumerate}
\end{exercise}

\begin{exercise}
\label{exer:8.36}
设 $I=(0,1)$.给定 $I$ 上的两个函数 $f_1,f_2$, 考虑方程组
\begin{equation}
\begin{cases}
-u_1''+u_2=f_1 & \text{在}\ I\ \text{上},\\
-u_2''-u_1=f_2 & \text{在}\ I\ \text{上},
\end{cases}
\tag{1}
\end{equation}
其中 $u_1,u_2$ 是未知函数.

\noindent\textbf{-A-}\quad 本部分规定 Dirichlet 条件
\begin{equation}
u_1(0)=u_2(0)=u_1(1)=u_2(1)=0.
\tag{2}
\end{equation}
\begin{enumerate}[label=\arabic*., leftmargin=2.2em, itemsep=0.3em]
	\item 为问题 (1)--(2) 定义适当的弱解概念.证明对每对
	      $f=[f_1,f_2]\in L^2(I)\times L^2(I)$, 存在唯一的弱解
	      $u=[u_1,u_2]\in H_0^1(I)\times H_0^1(I)$.
	\item 验证 $u_1,u_2\in H^2(I)$.
	\item 证明若 $f=[f_1,f_2]\in C(\bar I)\times C(\bar I)$, 则
	      $u=[u_1,u_2]\in C^2(\bar I)\times C^2(\bar I)$, 且 $u$ 是
	      (1)--(2) 的古典解.
	\item 考虑由 $Tf=u$ 定义的从 $L^2(I)\times L^2(I)$ 到自身的算子
	      $T$.验证 $T$ 紧.
	\item 证明 $\mathrm{EV}(T)=\varnothing$.
	\item $T$ 满射吗? 断言 $\sigma(T)=\{0\}$.
	\item $T$ 自伴吗? 算出 $T^\star$.
\end{enumerate}

\noindent\textbf{-B-}\quad 本部分规定 Neumann 条件
\begin{equation}
u_1'(0)=u_1'(1)=u_2'(0)=u_2'(1)=0.
\tag{3}
\end{equation}
\begin{enumerate}[label=\arabic*., leftmargin=2.2em, itemsep=0.3em]
	\item 为问题 (1)--(3) 定义适当的弱解概念.验证 Lax--Milgram 定理
	      不适用.
\end{enumerate}
给定 $\varepsilon>0$, 考虑方程组
\begin{equation}
\begin{cases}
-u_1''+u_2+\varepsilon u_1=f_1 & \text{在}\ I\ \text{上},\\
-u_2''-u_1+\varepsilon u_2=f_2 & \text{在}\ I\ \text{上}.
\end{cases}
\tag{$1_\varepsilon$}
\end{equation}
\begin{enumerate}[label=\arabic*., leftmargin=2.2em, itemsep=0.3em, start=2]
	\item 证明对每个 $f=[f_1,f_2]\in L^2(I)\times L^2(I)$, 问题
	      $(1_\varepsilon)$--(3) 都有唯一的弱解
	      $u_\varepsilon=[u_1^\varepsilon,u_2^\varepsilon]
	      \in H^1(I)\times H^1(I)$.
	\item 证明
	      \begin{equation*}
	      \norm{u_1^\varepsilon}_{L^2(I)}^2+\norm{u_2^\varepsilon}_{L^2(I)}^2
	      \le\norm{f_1}_{L^2(I)}^2+\norm{f_2}_{L^2(I)}^2.
	      \end{equation*}
	\item 由此推出当 $\varepsilon\to0$ 时,
	      $u_\varepsilon=[u_1^\varepsilon,u_2^\varepsilon]$ 在
	      $H^2(I)\times H^2(I)$ 中保持有界.
	\item 证明对每个 $f=[f_1,f_2]\in L^2(I)\times L^2(I)$, (1)--(3) 都有
	      唯一的解 $u=[u_1,u_2]\in H^2(I)\times H^2(I)$.
\end{enumerate}
\end{exercise}

\begin{exercise}
\label{exer:8.37}
\begin{enumerate}[label=\arabic*., leftmargin=2.2em, itemsep=0.3em]
	\item 证明从 $H^1(0,1)$ 到 $L^2(0,1)$ 的恒等算子是 Hilbert--Schmidt
	      算子 (见问题 40).
	      [提示: 写出 $u(x)=u(0)+\int_0^xu'(t)\,dt$, 并应用问题 40 中的
	      第 A3, A6 与 B4 问.]
	\item 考虑 Sturm--Liouville 问题
	      \begin{equation*}
	      \begin{cases}
	      -(pu')'+qu=\lambda u & \text{在}\ (0,1)\ \text{上},\\
	      u(0)=u(1)=0
	      \end{cases}
	      \end{equation*}
	      的特征值 $(\lambda_n)$.回忆在定理  \ref{thm:8.22} 的假设下
	      $\lambda_n\to+\infty$.于是存在某个整数 $N$, 使得
	      $\lambda_n>0$ 对一切 $n\ge N$ 成立.证明
	      \begin{equation*}
	      \sum_{n=N}^{+\infty}\frac{1}{\lambda_n^2}<\infty.
	      \end{equation*}
\end{enumerate}
\emph{注.} 习题 8.42 描述了一个精确得多的估计.
\end{exercise}

\begin{exercise}
\label{exer:8.38}
\emph{具有连续谱的算子的例子.}给定 $f\in L^2(\R)$, 设 $u\in H^1(\R)$ 是
问题
\begin{equation*}
-u''+u=f\quad\text{在}\ \R\ \text{上}
\end{equation*}
的唯一 (弱) 解, 其意义为
\begin{equation*}
\int_\R u'v'+\int_\R uv=\int_\R fv\quad\forall v\in H^1(\R).
\end{equation*}
令 $u=Tf$, 并把 $T$ 看作从 $H=L^2(\R)$ 到自身的具有界线性算子.
\begin{enumerate}[label=\arabic*., leftmargin=2.2em, itemsep=0.3em]
	\item 验证 $T^\star=T$ ($H$ 与其对偶空间等同), 且 $\norm{T}\le1$.
	\item 证明 $\mathrm{EV}(T)=\varnothing$.$T$ 是紧算子吗?
	\item 设 $\lambda\in(-\infty,0)$; 验证 $\lambda\in\rho(T)$.
	\item 设 $\lambda\in(1,+\infty)$; 验证 $\lambda\in\rho(T)$.
	\item 由此推出 $\sigma(T)\subset[0,1]$.
	\item $T$ 满射吗? 断言 $0\in\sigma(T)$.
	\item $(T-I)$ 满射吗? 断言 $1\in\sigma(T)$ 且 $\norm{T}=1$.
	\item 设 $\lambda\in(0,1)$.$(T-\lambda I)$ 满射吗?
	\item 断言 $\sigma(T)=[0,1]$.
\end{enumerate}
\end{exercise}

\begin{exercise}
\label{exer:8.39}
给定 $f\in L^2(0,1)$, 考虑由
\begin{equation*}
\varphi(v)=\frac12\int_0^1v'^2+\frac14\int_0^1v^4-\int_0^1fv,
\quad v\in H^1(0,1),
\end{equation*}
定义的函数 $\varphi:H^1(0,1)\to\R$.
\begin{enumerate}[label=\arabic*., leftmargin=2.2em, itemsep=0.3em]
	\item 验证 $\varphi$ 在 $H^1(0,1)$ 上凸且连续.
	\item 证明当 $\norm{v}_{H^1}\to\infty$ 时 $\varphi(v)\to+\infty$.
	\item 由此推出存在唯一的 $u\in H^1(0,1)$, 使得
	      \begin{equation}
	      \varphi(u)=\min_{v\in H^1(0,1)}\varphi(v).
	      \tag{1}
	      \end{equation}
	\item 证明
	      \begin{equation}
	      \int_0^1(u'v'+u^3v)=\int_0^1fv\quad\forall v\in H^1(0,1).
	      \tag{2}
	      \end{equation}
	      [提示: 写出 $\varphi(u)\le\varphi(u+\varepsilon v)$ 对一切
	      $v\in H^1(0,1)$ 成立, 并显式算出
	      $\varphi(u+\varepsilon v)$.]
	\item 证明 $u\in H^2(0,1)$, 且 $u$ 满足
	      \begin{equation}
	      \begin{cases}
	      -u''+u^3=f & \text{在}\ (0,1)\ \text{上几乎处处},\\
	      u'(0)=u'(1)=0.
	      \end{cases}
	      \tag{3}
	      \end{equation}
	\item 反之, 证明 (3) 的任何解都满足 (1).由此推出 (3) 有唯一解.
	\item 证明若 $f\ge0$ 在 $(0,1)$ 上几乎处处成立, 则 $u\ge0$ 在
	      $(0,1)$ 上.
	      [提示: 用与 \S \ref{sec:8.5} 相同的技巧.]
	\item 证明若 $f\in L^\infty(0,1)$, 则
	      \begin{equation*}
	      \norm{u}_{L^\infty(0,1)}^3\le\norm{f}_{L^\infty(0,1)}.
	      \end{equation*}
	      [提示: 像定理  \ref{thm:8.19} 的证明那样论证, 用
	      $G(u-K^{1/3})$ 作试验函数, 其中 $K=\norm{f}_{L^\infty}$.]
	\item 当 $\varphi(v)$ 被
	      \begin{equation*}
	      \psi(v)=\frac12\int_0^1v'^2+\frac14v^4(0)-\int_0^1fv,
	      \quad v\in H^1(0,1)
	      \end{equation*}
	      替换时, 情形如何?
\end{enumerate}
\end{exercise}

\begin{exercise}
\label{exer:8.40}
设 $j\in C^1(\R,\R)$ 是凸函数, 满足
\begin{equation}
j(t)\ge\abs{t}-C\quad\forall t\in\R,\ \text{对某个}\ C\in\R,
\tag{1}
\end{equation}
以及
\begin{equation}
-1<j'(t)<+1\quad\forall t\in\R.
\tag{2}
\end{equation}
[一个值得记住的好例子是 $j(t)=(1+t^2)^{1/2}$.]

给定 $f\in L^2(0,1)$, 考虑由
\begin{equation*}
\varphi(v)=\frac12\int_0^1v'^2+\int_0^1j(v)-\int_0^1fv,
\quad v\in H^1(0,1),
\end{equation*}
定义的函数 $\varphi:H^1(0,1)\to\R$.
\begin{enumerate}[label=\arabic*., leftmargin=2.2em, itemsep=0.3em]
	\item 证明若 $\abs{\int_0^1f}>1$, 则
	      \begin{equation*}
	      \inf_{v\in H^1(0,1)}\varphi(v)=-\infty.
	      \end{equation*}
	      [提示: 取 $v=$ 常数.]
	\item 证明若 $\abs{\int_0^1f}\le1$, 则
	      \begin{equation*}
	      \inf_{v\in H^1(0,1)}\varphi(v)>-\infty.
	      \end{equation*}
	      [提示: 写出 $\int_0^1fv=\int_0^1\bar f(v-\bar v)+\bar f\bar v$, 其中
	      $\bar v=\int_0^1v$, $\bar f=\int_0^1f$, 并使用
	      Poincaré--Wirtinger 不等式; 见第~
\ref{chap:8}~章评注 1 与问题 47.]
	\item 证明若 $\abs{\int_0^1f}<1$, 则
	      \begin{equation*}
	      \lim_{\norm{v}_{H^1(0,1)}\to\infty}\varphi(v)=+\infty.
	      \end{equation*}
	      由此推出 $\inf_{v\in H^1(0,1)}\varphi(v)$ 可以达到, 并且每个极小
	      点 $u\in H^1(0,1)$ 都满足
	      \begin{equation*}
	      \int_0^1u'w'+\int_0^1j'(u)w=\int_0^1fw\quad\forall w\in H^1(0,1).
	      \end{equation*}
	      证明 $u\in H^2(0,1)$, 并且是问题
	      \begin{equation}
	      \begin{cases}
	      -u''+j'(u)=f & \text{在}\ (0,1)\ \text{上},\\
	      u'(0)=u'(1)=0
	      \end{cases}
	      \tag{3}
	      \end{equation}
	      的解.
	\item 假设 $\abs{\int_0^1f}=1$.证明
	      $\inf_{v\in H^1(0,1)}\varphi(v)$ 不能达到.
	      [提示: 用反证法.若下确界在某个 $u$ 处达到, 则 $u$ 满足 (3).
	      在 $(0,1)$ 上积分 (3).]
	\item 当 $\abs{\int_0^1f}=1$ 时, $\varphi$ 的极小化序列 $(u_n)$ 的
	      情形如何?
	      [提示: 证明 $u_n=\bar u_n+(u_n-\bar u_n)$, 其中
	      $\abs{\bar u_n}\to\infty$, 而当 $n\to\infty$ 时
	      $\norm{u_n-\bar u_n}_{H^1}\le C$.]
\end{enumerate}
\end{exercise}

\begin{exercise}
\label{exer:8.41}
设 $q\in C([0,1])$, 并假设双线性形式
\begin{equation*}
a(u,v)=\int_0^1(u'v'+quv),\quad u,v\in H_0^1(0,1),
\end{equation*}
在 $H_0^1(0,1)$ 上强制.空间 $H_0^1(0,1)$ 配以标量积 $a(u,v)$ (现记作
$(u,v)_H$), 范数 $\abs{u}_H=(u,u)_H^{1/2}$.
\begin{enumerate}[label=\arabic*., leftmargin=2.2em, itemsep=0.3em]
	\item 证明
	      \begin{equation}
	      \alpha=\sup\left\{\int_0^1\abs{u}^4;\ u\in H_0^1(0,1)\
	      \text{且}\ \abs{u}_H=1\right\}>0
	      \tag{1}
	      \end{equation}
	      可由某个 $u_0$ 达到.
	      [提示: 考虑在 $H_0^1(0,1)$ 中弱收敛到某个 $u_0$ 的极大化序列
	      $(u_n)$.验证 $\alpha=\int_0^1\abs{u_0}^4$ 且
	      $\abs{u_0}_H\le1$.通过引入 $u_0/\abs{u_0}_H$ 证明
	      $\abs{u_0}_H=1$.]
	\item 证明可以假设 $u_0\ge0$ 在 $[0,1]$ 上.
	      [提示: 把 $u_0$ 换成 $\abs{u_0}$, 并应用习题 8.11.]
	\item 证明 $u_0$ 属于 $H^2(0,1)$, 并且满足
	      \begin{equation}
	      \begin{cases}
	      -u_0''+qu_0=\frac1\alpha u_0^3 & \text{在}\ (0,1)\ \text{上},\\
	      u_0(0)=u_0(1)=0.
	      \end{cases}
	      \tag{2}
	      \end{equation}
	      [提示: 写出 $\norm{w_\varepsilon}_{L^4(0,1)}
	      \le\norm{u_0}_{L^4(0,1)}$, 其中
	      $w_\varepsilon=\frac{u_0+\varepsilon v}{\abs{u_0+\varepsilon v}_H}$,
	      $v\in H_0^1(0,1)$, 而 $\varepsilon>0$ 足够小.然后对
	      $\norm{w_\varepsilon}_{L^4(0,1)}^4$ 与
	      $\abs{u_0+\varepsilon v}_H^2$ 作 Taylor 展开, 令
	      $\varepsilon\to0$.]
	\item 由此推出 $u_0(x)>0$ 对一切 $x\in(0,1)$ 成立.
	      [提示: 用强最大值原理; 见问题 45.]
	\item 设 $u_1$ 是 (1) 的任一极大化函数.证明或者
	      $u_1(x)>0$ 对一切 $x\in(0,1)$ 成立, 或者 $u_1(x)<0$ 对一切
	      $x\in(0,1)$ 成立.
	      [提示: 验证 $\abs{u_1(x)}>0$ 对一切 $x\in(0,1)$ 成立.]
	\item 由此推出问题
	      \begin{equation}
	      \begin{cases}
	      -u''+qu=u^3 & \text{在}\ (0,1)\ \text{上},\\
	      u>0 & \text{在}\ (0,1)\ \text{上},\\
	      u(0)=u(1)=0
	      \end{cases}
	      \tag{3}
	      \end{equation}
	      存在解 $u\in C^2([0,1])$.
	      [提示: 取 $u=ku_0$, 其中 $k>0$ 为某个适当的常数.]
	\item 现在假设 $a$ 不强制, 更确切地说, 存在某个
	      $v_1\in H_0^1(0,1)$, 使得 $v_1\not\equiv0$ 且
	      $a(v_1,v_1)\le0$.证明问题 (3) 无解.
	      [提示: 验证 $\lambda_1\le0$, 其中 $\lambda_1$ 是
	      $Au=-u''+qu$ 的第一个特征值; 然后把 (3) 乘以相应的特征函数
	      $\varphi_1$.]
\end{enumerate}
\end{exercise}

\begin{exercise}
\label{exer:8.42}
\emph{Sturm--Liouville 特征值的渐近行为.}考虑 $I=(0,L)$ 上的算子
$Av=-v''+a(x)v$, 带零 Dirichlet 条件, 其中 $a\in C([0,L])$.
\begin{enumerate}[label=\arabic*., leftmargin=2.2em, itemsep=0.3em]
	\item 用 $(\lambda_n)$ 表示 $A$ 的特征值序列.证明
	      \begin{equation*}
	      \left|\lambda_n-\frac{\pi^2n^2}{L^2}\right|
	      \le\norm{a}_{L^\infty(0,L)}\quad\forall n.
	      \end{equation*}
	      [提示: 考虑对应于 $a\equiv0$ 的算子 $A_0$ 的特征值, 并使用
	      Courant--Fischer 极小极大原理 (见问题 49).]
\end{enumerate}
现在考虑一般的 Sturm--Liouville 算子
\begin{equation*}
Bu=-(pu')'+qu\quad\text{在}\ (0,1)\ \text{上},
\end{equation*}
带零 Dirichlet 条件.假设 $p\in C^2([0,1])$, $p\ge\alpha>0$ 在 $(0,1)$ 上,
$q\in C([0,1])$.令 $L=\int_0^1p(t)^{-1/2}\,dt$, 并引入新变量
$x=\int_0^tp(t)^{-1/2}\,dt$, 使得当 $0<t<1$ 时 $0<x<L$.给定函数
$u\in C^2([0,1])$, 令
\begin{equation*}
v(x)=p^{1/4}(t)u(t),\quad 0<t<1.
\end{equation*}
\begin{enumerate}[label=\arabic*., leftmargin=2.2em, itemsep=0.3em, start=2]
	\item 证明 $u$ 在 $(0,1)$ 上满足 $-(pu')'+qu=\mu u$ 当且仅当 $v$ 在
	      $(0,L)$ 上满足 $-v''+av=\mu v$, 其中 $a\in C([0,1])$ 只依赖于
	      $p$ 与 $q$.
	      [提示: 经过一些繁复的计算后, 证明
	      $a(x)=q(t)+\frac14p''(t)-\frac{(p'(t))^2}{16p(t)}$.]
	\item 由此推出算子 $B$ 的特征值 $(\mu_n)$ 满足
	      \begin{equation*}
	      \left|\mu_n-\frac{\pi^2n^2}{L^2}\right|\le C.
	      \end{equation*}
\end{enumerate}
\end{exercise}

% chapters/ch09.tex
% chapters/ch09.tex
% chapters/ch09.tex
% 第~
% \ref{chap:9}~章 Sobolev 空间与 N 维椭圆边值问题的变分表述
% 源: 原书 PDF 第 269--329 页 (印刷页 263--323)

\chapter{Sobolev 空间与 N 维椭圆边值问题的变分表述}
\label{chap:9}

\section{\index{NweiSobolevkongjian@N 维 Sobolev 空间 (Sobolev space in dimension N)}
Sobolev 空间 $W^{1,p}(\Omega)$ 的定义与初等性质}
\label{sec:9.1}

设 $\Omega\subset\R^N$ 为开集, $p\in\R$ 满足 $1\le p\le\infty$.

\begin{definition}
\emph{Sobolev 空间} $W^{1,p}(\Omega)$ 定义为\footnote{若无混淆, 我们常写
$W^{1,p}$ 而不写 $W^{1,p}(\Omega)$.}
\begin{equation*}
W^{1,p}(\Omega)=\left\{u\in L^p(\Omega)\ \middle|\
\begin{array}{l}
\exists g_1,g_2,\dots,g_N\in L^p(\Omega)\ \text{使得}\\
\displaystyle\int_\Omega u\frac{\partial\varphi}{\partial x_i}
=-\int_\Omega g_i\varphi\quad\forall\varphi\in C_c^\infty(\Omega),\
\forall i=1,2,\dots,N
\end{array}
\right\}.
\end{equation*}
我们令
\begin{equation*}
H^1(\Omega)=W^{1,2}(\Omega).
\end{equation*}
对 $u\in W^{1,p}(\Omega)$, 我们定义\footnote{这个定义是有意义的: 由推论
 \ref{cor:4.24}, $g_i$ (几乎处处) 唯一.}
$\dfrac{\partial u}{\partial x_i}=g_i$, 并记
\begin{equation*}
\nabla u=\operatorname{grad}u=
\left(\frac{\partial u}{\partial x_1},\frac{\partial u}{\partial x_2},
\dots,\frac{\partial u}{\partial x_N}\right).
\end{equation*}
空间 $W^{1,p}(\Omega)$ 配以范数
\begin{equation*}
\norm{u}_{W^{1,p}}=\norm{u}_p+\sum_{i=1}^N
\norm{\frac{\partial u}{\partial x_i}}_p,
\end{equation*}
有时 (若 $1\le p<\infty$) 也配以等价范数
$(\norm{u}_p^p+\sum_{i=1}^N\norm{\frac{\partial u}{\partial x_i}}_p^p)^{1/p}$.
空间 $H^1(\Omega)$ 配以标量积
\begin{equation*}
(u,v)_{H^1}=(u,v)_{L^2}+\sum_{i=1}^N
\left(\frac{\partial u}{\partial x_i},\frac{\partial v}{\partial x_i}\right)_{L^2}
=\int_\Omega uv+\sum_{i=1}^N\frac{\partial u}{\partial x_i}
\frac{\partial v}{\partial x_i}.
\end{equation*}
相应的范数
\begin{equation*}
\norm{u}_{H^1}=\left(\norm{u}_2^2+\sum_{i=1}^N
\norm{\frac{\partial u}{\partial x_i}}_2^2\right)^{1/2}
\end{equation*}
与 $W^{1,2}$ 范数等价.
\end{definition}

\begin{keyproposition}
\label{prop:9.1}
$W^{1,p}(\Omega)$ 对每个 $1\le p\le\infty$ 都是 Banach 空间.当
$1<p<\infty$ 时它自反, 当 $1\le p<\infty$ 时它可分.$H^1(\Omega)$ 是可分
Hilbert 空间.
\end{keyproposition}

\begin{proof}
仿照命题  \ref{prop:8.1} 的证明, 取算子 $Tu=[u,\nabla u]$.
\end{proof}

\begin{remark}
\label{rem:9.1}
在 $W^{1,p}$ 的定义中, 我们本来同样可以用 $C_c^\infty(\Omega)$ 作为检验函数
$\varphi$ 的类 (而不用 $C_c^1(\Omega)$); 为证明这一点, 利用一列磨光核
$(\rho_n)$.
\end{remark}

\begin{remark}
\label{rem:9.2}
显然, 若 $u\in C^1(\Omega)\cap L^p(\Omega)$ 且
$\frac{\partial u}{\partial x_i}\in L^p(\Omega)$ 对一切 $i=1,2,\dots,N$
成立 (这里 $\frac{\partial u}{\partial x_i}$ 指 $u$ 的通常偏导数), 则
$u\in W^{1,p}(\Omega)$.此外, 通常的偏导数与 $W^{1,p}$ 意义下的偏导数一致,
所以记号是相容的.特别地, 若 $\Omega$ 有界, 则
$C^1(\overline\Omega)\subset W^{1,p}(\Omega)$ 对一切 $1\le p\le\infty$
成立.反之可以证明: 若 $u\in W^{1,p}(\Omega)$ (其中 $1\le p\le\infty$) 且
$\frac{\partial u}{\partial x_i}\in C(\Omega)$ 对一切 $i=1,2,\dots,N$ 成立
(这里 $\frac{\partial u}{\partial x_i}$ 指 $W^{1,p}$ 意义下的偏导数), 则
$u\in C^1(\Omega)$; 更精确地说, 存在函数 $\tilde u\in C^1(\Omega)$ 使得
$u=\tilde u$ 几乎处处成立.
\end{remark}

\begin{starremark}
\label{rem:9.3}
对每个 $u\in L^1_{\mathrm{loc}}(\Omega)$, 分布论赋予
$\frac{\partial u}{\partial x_i}$ 一个意义 ($\frac{\partial u}{\partial x_i}$
是巨大的分布空间 $\mathcal{D}'(\Omega)$ 中的元素, 这个空间特别地包含
$L^1_{\mathrm{loc}}(\Omega)$).用分布的语言可以说, $W^{1,p}(\Omega)$ 就是
那些使得所有偏导数 $\frac{\partial u}{\partial x_i}$ ($1\le i\le N$, 分布
意义下的) 都属于 $L^p(\Omega)$ 的函数 $u\in L^p(\Omega)$ 的集合.

当 $\Omega=\R^N$ 且 $p=2$ 时, 也可以用 Fourier 变换来定义 Sobolev 空间;
参见例如 J.-L.~Lions--E.~Magenes \bref{LionsMagenes}{1}, P.~Malliavin \bref{Malliavin}{1}, H.~Triebel \bref{Triebel}{1}, L.~Grafakos \bref{Grafakos}{1}.这里不采取这个观点.
\end{starremark}

\begin{remark}
\label{rem:9.4}
记住下述事实是有用的:

\noindent (i) 设 $(u_n)$ 为 $W^{1,p}$ 中的序列, 使得 $u_n\to u$ 在 $L^p$ 中,
且 $(\nabla u_n)$ 在 $(L^p)^N$ 中收敛到某个极限.则 $u\in W^{1,p}$ 且
$\norm{u_n-u}_{W^{1,p}}\to0$.当 $1<p\le\infty$ 时, 只须知道 $u_n\to u$
在 $L^p$ 中且 $(\nabla u_n)$ 在 $(L^p)^N$ 中有界, 就能推出 $u\in W^{1,p}$
(为什么?).

\noindent (ii) 给定定义在 $\Omega$ 上的函数 $f$, 我们记 $\bar f$ 为它在
$\Omega$ 外的延拓, 即
\begin{equation*}
\bar f(x)=
\begin{cases}
f(x) & \text{若}\ x\in\Omega,\\
0 & \text{若}\ x\in\R^N\setminus\Omega.
\end{cases}
\end{equation*}
设 $u\in W^{1,p}(\Omega)$, $\alpha\in C_c^1(\Omega)$.则\footnote{注意: 一般
说来 $\bar u\notin W^{1,p}(\R^N)$ (为什么?).}
\begin{equation*}
\alpha u\in W^{1,p}(\R^N)
\quad\text{且}\quad
\frac{\partial}{\partial x_i}(\alpha u)=\alpha\frac{\partial u}{\partial x_i}
+\frac{\partial\alpha}{\partial x_i}u.
\end{equation*}
事实上, 设 $\varphi\in C_c^1(\R^N)$; 我们有
\begin{align*}
\int_{\R^N}\alpha u\frac{\partial\varphi}{\partial x_i}
&=\int_\Omega\alpha u\frac{\partial\varphi}{\partial x_i}
=\int_\Omega u\left(\frac{\partial}{\partial x_i}(\alpha\varphi)
-\frac{\partial\alpha}{\partial x_i}\varphi\right)\\
&=-\int_\Omega\left(\frac{\partial u}{\partial x_i}\alpha\varphi
+u\frac{\partial\alpha}{\partial x_i}\varphi\right)
=-\int_{\R^N}\left(\alpha\frac{\partial u}{\partial x_i}
+\frac{\partial\alpha}{\partial x_i}u\right)\varphi.
\end{align*}
如果把 $\alpha\in C_c^1(\Omega)$ 换成
$\alpha\in C^1(\R^N)\cap L^\infty(\R^N)$, 其中
$\nabla\alpha\in L^\infty(\R^N)^N$ 且
$\supp\alpha\subset\R^N\setminus(\partial\Omega)$, 结论仍然成立.

下面是对一般开集 $\Omega$ 成立的第一个稠密性结果; 我们稍后在附加假设下
(推论  \ref{cor:9.8}) 建立一个更精确的结果.我们需要下述概念.
\end{remark}

\begin{definition}
设 $\Omega\subset\R^N$ 为开集.我们说 $\R^N$ 中的开集 $\omega$ \emph{强含于}
$\Omega$, 并记 $\omega\subset\subset\Omega$, 如果 $\overline\omega\subset\Omega$
且 $\overline\omega$ 紧.\footnote{$\overline\omega$ 表示 $\omega$ 在 $\R^N$
中的闭包.}
\end{definition}

\begin{keytheorem}[Friedrichs]
\label{thm:9.2}
设 $u\in W^{1,p}(\Omega)$, $1\le p<\infty$.则存在 $C_c^\infty(\R^N)$ 中的
序列 $(u_n)$ 使得
\begin{equation}
u_n|_\Omega\to u\quad\text{在}\ L^p(\Omega)\ \text{中}
\label{eq:9.1}
\end{equation}
且
\begin{equation}
\nabla u_n|_\omega\to\nabla u|_\omega
\quad\text{在}\ L^p(\omega)^N\ \text{中对一切}\ \omega\subset\subset\Omega.
\label{eq:9.2}
\end{equation}
若 $\Omega=\R^N$ 且 $u\in W^{1,p}(\R^N)$, $1\le p<\infty$, 则存在
$C_c^\infty(\R^N)$ 中的序列 $(u_n)$ 使得
\begin{equation*}
u_n\to u\quad\text{在}\ L^p(\R^N)\ \text{中}
\quad\text{且}\quad
\nabla u_n\to\nabla u\quad\text{在}\ L^p(\R^N)^N\ \text{中}.
\end{equation*}
\end{keytheorem}

在证明中我们要用到下述引理.

\begin{lemma}
\label{lem:9.1}
设 $\rho\in L^1(\R^N)$, $v\in W^{1,p}(\R^N)$, $1\le p\le\infty$.则
\begin{equation*}
\rho\star v\in W^{1,p}(\R^N)
\quad\text{且}\quad
\frac{\partial}{\partial x_i}(\rho\star v)=\rho\star\frac{\partial v}{\partial x_i}
\quad\forall i=1,2,\dots,N.
\end{equation*}
\end{lemma}

\begin{proof}[引理 \ref{lem:9.1}]
仿照引理  \ref{lem:8.4} 的证明.
\end{proof}

\begin{proof}[定理 \ref{thm:9.2}]
令
\begin{equation*}
\bar u(x)=
\begin{cases}
u(x) & \text{若}\ x\in\Omega,\\
0 & \text{若}\ x\in\R^N\setminus\Omega,
\end{cases}
\end{equation*}
并令 $v_n=\rho_n\star\bar u$ (其中 $(\rho_n)$ 是一列磨光核).我们知道
(见第~\ref{sec:4.4}~节) $v_n\in C^\infty(\R^N)$ 且 $v_n\to\bar u$ 在 $L^p(\R^N)$ 中.
我们断言: 对一切 $\omega\subset\subset\Omega$,
$\nabla v_n|_\omega\to\nabla u|_\omega$ 在 $L^p(\omega)^N$ 中.事实上,
给定 $\omega\subset\subset\Omega$, 固定函数 $\alpha\in C_c^1(\Omega)$,
$0\le\alpha\le1$, 使得 $\alpha=1$ 在 $\omega$ 的某个邻域上.当 $n$ 足够大时
我们有
\begin{equation}
\rho_n\star(\alpha u)=\rho_n\star\bar u\quad\text{在}\ \omega\ \text{上},
\label{eq:9.3}
\end{equation}
因为
\begin{align*}
\supp(\rho_n\star\alpha u-\rho_n\star\bar u)
&=\supp\bigl(\rho_n\star(1-\bar\alpha)\bar u\bigr)\\
&\subset\supp\rho_n+\supp\bigl((1-\bar\alpha)\bar u\bigr)
\subset B\left(0,\frac1n\right)+\supp(1-\bar\alpha)
\subset(\omega)^c
\end{align*}
对足够大的 $n$ 成立.由引理  \ref{lem:9.1} 与注记 \ref{rem:9.4}(ii) 我们有
\begin{equation*}
\frac{\partial}{\partial x_i}(\rho_n\star\alpha u)
=\rho_n\star\left(\alpha\frac{\partial u}{\partial x_i}
+\frac{\partial\alpha}{\partial x_i}u\right).
\end{equation*}
由此推出
\begin{equation*}
\frac{\partial}{\partial x_i}(\rho_n\star\alpha u)
\to\alpha\frac{\partial u}{\partial x_i}
+\frac{\partial\alpha}{\partial x_i}u\quad\text{在}\ L^p(\R^N)\ \text{中},
\end{equation*}
特别地
\begin{equation*}
\frac{\partial}{\partial x_i}(\rho_n\star\alpha u)
\to\frac{\partial u}{\partial x_i}\quad\text{在}\ L^p(\omega)\ \text{中}.
\end{equation*}
由 \eqref{eq:9.3} 我们有
\begin{equation*}
\frac{\partial}{\partial x_i}(\rho_n\star\bar u)
\to\frac{\partial u}{\partial x_i}\quad\text{在}\ L^p(\omega)\ \text{中}.
\end{equation*}
最后, 我们用一列截断函数 $(\zeta_n)$ 乘序列 $(v_n)$, 如定理  \ref{thm:8.7}
的证明中那样.\footnote{本章中我们总是用 $(\zeta_n)$ 表示一列截断函数, 即
固定函数 $\zeta\in C_c^\infty(\R^N)$, $0\le\zeta\le1$, 使得
\begin{equation*}
\zeta(x)=
\begin{cases}
1 & \text{若}\ \abs{x}\le1,\\
0 & \text{若}\ \abs{x}\ge2,
\end{cases}
\end{equation*}
并令 $\zeta_n(x)=\zeta(x/n)$, $n=1,2,\dots$.} 容易验证序列
$u_n=\zeta_nv_n$ 具有所要的性质, 即 $u_n\in C_c^\infty(\R^N)$,
$u_n\to u$ 在 $L^p(\Omega)$ 中, 且对每个 $\omega\subset\subset\Omega$,
$\nabla u_n\to\nabla u$ 在 $(L^p(\omega))^N$ 中.

当 $\Omega=\R^N$ 时, 序列 $u_n=\zeta_n(\rho_n\star u)$ 具有所要的性质.
\end{proof}

\begin{starremark}
\label{rem:9.5}
可以证明 (Meyers--Serrin 定理): 若 $u\in W^{1,p}(\Omega)$, $1\le p<\infty$,
则存在 $C^\infty(\Omega)\cap W^{1,p}(\Omega)$ 中的序列 $(u_n)$ 使得
$u_n\to u$ 在 $W^{1,p}(\Omega)$ 中; 这个结果的证明相当精细 (参见例如
R.~Adams \bref{Adams}{1} 或 A.~Friedman \bref{Friedman}{2}).一般说来,
若 $\Omega$ 是任意开集且 $u\in W^{1,p}(\Omega)$, 则未必存在
$C_c^1(\R^N)$ 中的序列 $(u_n)$ 使得 $u_n|_\Omega\to u$ 在 $W^{1,p}(\Omega)$
中.请把 Meyers--Serrin 定理 (它对任意开集成立) 与推论  \ref{cor:9.8}
(它假设 $\Omega$ 是正则的) 作比较.
\end{starremark}

下面是 $W^{1,p}$ 函数的一个简单刻画:

\begin{proposition}
\label{prop:9.3}
设 $u\in L^p(\Omega)$, $1<p\le\infty$.下述性质等价:

\noindent (i) $u\in W^{1,p}(\Omega)$;

\noindent (ii) 存在常数 $C$ 使得
\begin{equation*}
\left|\int_\Omega u\frac{\partial\varphi}{\partial x_i}\right|
\le C\norm{\varphi}_{L^{p'}(\Omega)}
\quad\forall\varphi\in C_c^\infty(\Omega),\
\forall i=1,2,\dots,N;
\end{equation*}

\noindent (iii) 存在常数 $C$ 使得对一切 $\omega\subset\subset\Omega$ 与一切
$h\in\R^N$, $\abs{h}<\dist(\omega,\partial\Omega)$, 有
\begin{equation*}
\norm{\tau_hu-u}_{L^p(\omega)}\le C\abs{h}.
\end{equation*}
(注意: 对 $x\in\omega$ 与 $\abs{h}<\dist(\omega,\partial\Omega)$,
$\tau_hu(x)=u(x+h)$ 有意义.)

此外, 在 (ii) 与 (iii) 中都可以取 $C=\norm{\nabla u}_{L^p(\Omega)}$.
若 $\Omega=\R^N$, 我们有
\begin{equation*}
\norm{\tau_hu-u}_{L^p(\R^N)}\le\abs{h}\norm{\nabla u}_{L^p(\R^N)}.
\end{equation*}
\end{proposition}

\begin{proof}
\noindent (i)$\Rightarrow$(ii). 显然.

\noindent (ii)$\Rightarrow$(i). 仿照命题  \ref{prop:8.3} 的证明.

\noindent (i)$\Rightarrow$(iii). 先设 $u\in C_c^\infty(\R^N)$.设 $h\in\R^N$,
令
\begin{equation*}
v(t)=u(x+th),\quad t\in\R.
\end{equation*}
则 $v'(t)=h\cdot\nabla u(x+th)$, 于是
\begin{equation*}
u(x+h)-u(x)=v(1)-v(0)=\int_0^1v'(t)\,dt
=\int_0^1h\cdot\nabla u(x+th)\,dt.
\end{equation*}
由此推出, 当 $1\le p<\infty$ 时,
\begin{equation*}
\abs{\tau_hu(x)-u(x)}^p\le\abs{h}^p\int_0^1
\abs{\nabla u(x+th)}^p\,dt,
\end{equation*}
且
\begin{align*}
\int_\omega\abs{\tau_hu(x)-u(x)}^p\,dx
&\le\abs{h}^p\int_\omega dx\int_0^1\abs{\nabla u(x+th)}^p\,dt\\
&=\abs{h}^p\int_0^1dt\int_\omega\abs{\nabla u(x+th)}^p\,dx
=\abs{h}^p\int_0^1dt\int_{\omega+th}\abs{\nabla u(y)}^p\,dy.
\end{align*}
若 $\abs{h}<\dist(\omega,\partial\Omega)$, 则存在开集
$\omega'\subset\subset\Omega$ 使得 $\omega+th\subset\omega'$ 对一切
$t\in[0,1]$ 成立, 于是
\begin{equation}
\norm{\tau_hu-u}_{L^p(\omega)}^p\le\abs{h}^p\int_{\omega'}
\abs{\nabla u}^p.
\label{eq:9.4}
\end{equation}
这就完成了 $u\in C_c^\infty(\R^N)$, $1\le p<\infty$ 情形的 (iii) 的证明.
现在设 $u\in W^{1,p}(\Omega)$, $1\le p<\infty$.由定理  \ref{thm:9.2}, 存在
$C_c^\infty(\R^N)$ 中的序列 $(u_n)$ 使得 $u_n\to u$ 在 $L^p(\Omega)$ 中,
且对一切 $\omega\subset\subset\Omega$, $\nabla u_n\to\nabla u$ 在
$L^p(\omega)^N$ 中.把 \eqref{eq:9.4} 用于 $(u_n)$ 并取极限, 即得对每个
$u\in W^{1,p}(\Omega)$ ($1\le p<\infty$) 成立的 (iii).当 $p=\infty$ 时,
对上文 (关于 $p<\infty$) 的结论令 $p\to\infty$.

\noindent (iii)$\Rightarrow$(ii). 设 $\varphi\in C_c^\infty(\Omega)$, 考虑开集
$\omega$ 使得 $\supp\varphi\subset\omega\subset\subset\Omega$.设 $h\in\R^N$,
$\abs{h}<\dist(\omega,\partial\Omega)$.由 (iii) 我们有
\begin{equation*}
\left|\int_\Omega(\tau_hu-u)\varphi\right|
\le C\abs{h}\norm{\varphi}_{L^{p'}(\Omega)}.
\end{equation*}
另一方面, 由于
\begin{equation*}
\int_\Omega(u(x+h)-u(x))\varphi(x)\,dx
=\int_\Omega u(y)\bigl(\varphi(y-h)-\varphi(y)\bigr)\,dy,
\end{equation*}
由此推出
\begin{equation*}
\left|\int_\Omega u(y)\frac{\varphi(y-h)-\varphi(y)}{\abs{h}}\,dy\right|
\le C\norm{\varphi}_{L^{p'}(\Omega)}.
\end{equation*}
取 $h=te_i$, $t\in\R$, 并令 $t\to0$, 即得 (ii).
\end{proof}

\begin{starremark}
\label{rem:9.6}
当 $p=1$ 时, 下述蕴涵仍然成立:
\begin{equation*}
\text{(i)}\Rightarrow\text{(ii)}\iff\text{(iii)}.
\end{equation*}
满足 $p=1$ 的 (ii) (或 (iii)) 的函数称为\emph{有界变差函数} (用分布的语言,
有界变差函数就是一个 $L^1$ 函数, 其所有一阶导数 (分布意义下的) 都是有界
测度).这个空间在许多应用中起着重要作用.人们在极小曲面理论中 (参见例如
E.~Giusti \bref{Giusti}{1} 以及 E.~DeGiorgi, M.~Miranda 等人的著作及其中
所引文献), 在弹性与塑性问题中 (有限变形函数, 参见例如
R.~Temam--G.~Strang \bref{TemamStrang}{2} 以及所引的 P.~Suquet 的工作), 在允许间断解的一阶拟线性方程即所谓守恒律中 (参见例如
A.~I.~Volpert \bref{Volpert}{1} 与 A.~Bressan \bref{Bressan}{1}),
都会遇到有界变差函数 (或具有类似性质的函数).关于这个广阔的课题, 还可参看
L.~Ambrosio--N.~Fusco--D.~Pallara \bref{Ambrosio}{1} 的书以及本章末的
评注 16.
\end{starremark}

\begin{remark}
\label{rem:9.7}
命题  \ref{prop:9.3} ((i)$\Rightarrow$(iii)) 表明, 任何
$u\in W^{1,\infty}(\Omega)$ 在 $\Omega$ 上都有连续的代表元.更精确地说,
若 $\Omega$ 连通, 则
\begin{equation}
\abs{u(x)-u(y)}\le\norm{\nabla u}_{L^\infty(\Omega)}
\dist_\Omega(x,y)\quad\forall x,y\in\Omega
\label{eq:9.5}
\end{equation}
(对上述连续代表元 $u$ 成立), 其中 $\dist_\Omega(x,y)$ 表示 $\Omega$ 中从
$x$ 到 $y$ 的测地距离; 特别地, 若 $\Omega$ 凸, 则
$\dist_\Omega(x,y)=\abs{x-y}$.由此还可推出: 若
$u\in W^{1,p}(\Omega)$ (其中 $1\le p\le\infty$, $\Omega$ 为某开集) 且
$\nabla u=0$ 在 $\Omega$ 上几乎处处成立, 则 $u$ 在 $\Omega$ 的每个连通分支上
为常数.
\end{remark}

\begin{proposition}[乘积的微分]
\label{prop:9.4}
设 $u,v\in W^{1,p}(\Omega)\cap L^\infty(\Omega)$, $1\le p\le\infty$.则
$uv\in W^{1,p}(\Omega)\cap L^\infty(\Omega)$ 且
\begin{equation*}
\frac{\partial}{\partial x_i}(uv)=\frac{\partial u}{\partial x_i}v
+u\frac{\partial v}{\partial x_i},\quad i=1,2,\dots,N.
\end{equation*}
\end{proposition}

\begin{proof}
如推论  \ref{cor:8.10} 的证明中那样, 只须考虑 $1\le p<\infty$ 的情形.由
定理  \ref{thm:9.2}, 存在 $C_c^\infty(\R^N)$ 中的序列 $(u_n),(v_n)$ 使得
\begin{equation*}
u_n\to u,\quad v_n\to v\quad\text{在}\ L^p(\Omega)\ \text{中且几乎处处在}\
\Omega\ \text{上},
\end{equation*}
\begin{equation*}
\nabla u_n\to\nabla u,\quad\nabla v_n\to\nabla v\quad\text{在}
\ L^p(\omega)^N\ \text{中对一切}\ \omega\subset\subset\Omega.
\end{equation*}
检查定理  \ref{thm:9.2} 的证明, 我们容易看出还有
\begin{equation*}
\norm{u_n}_{L^\infty(\R^N)}\le\norm{u}_{L^\infty(\Omega)}
\quad\text{且}\quad
\norm{v_n}_{L^\infty(\R^N)}\le\norm{v}_{L^\infty(\Omega)}.
\end{equation*}
另一方面,
\begin{equation*}
\int_\Omega u_nv_n\frac{\partial\varphi}{\partial x_i}
=-\int_\Omega\left(\frac{\partial u_n}{\partial x_i}v_n
+u_n\frac{\partial v_n}{\partial x_i}\right)\varphi
\quad\forall\varphi\in C_c^1(\Omega).
\end{equation*}
取极限 (利用控制收敛定理), 这成为
\begin{equation*}
\int_\Omega uv\frac{\partial\varphi}{\partial x_i}
=-\int_\Omega\left(\frac{\partial u}{\partial x_i}v
+u\frac{\partial v}{\partial x_i}\right)\varphi
\quad\forall\varphi\in C_c^1(\Omega).
\end{equation*}
\end{proof}

\begin{proposition}[复合的微分]
\label{prop:9.5}
设 $G\in C^1(\R)$ 满足 $G(0)=0$ 且 $\abs{G'(s)}\le M$ 对一切 $s\in\R$
成立, 其中 $M$ 为某常数.设 $u\in W^{1,p}(\Omega)$, $1\le p\le\infty$.则
\begin{equation*}
G\circ u\in W^{1,p}(\Omega)
\quad\text{且}\quad
\frac{\partial}{\partial x_i}(G\circ u)=(G'\circ u)
\frac{\partial u}{\partial x_i},\quad i=1,2,\dots,N.
\end{equation*}
\end{proposition}

\begin{proof}
我们有 $\abs{G(s)}\le M\abs{s}$ 对一切 $s\in\R$ 成立, 于是
$\abs{G\circ u}\le M\abs{u}$; 因此 $G\circ u\in L^p(\Omega)$, 并且
$(G'\circ u)\frac{\partial u}{\partial x_i}\in L^p(\Omega)$.剩下要验证
\begin{equation}
\int_\Omega(G\circ u)\frac{\partial\varphi}{\partial x_i}
=-\int_\Omega(G'\circ u)\frac{\partial u}{\partial x_i}\varphi
\quad\forall\varphi\in C_c^1(\Omega).
\label{eq:9.6}
\end{equation}
当 $1\le p<\infty$ 时, 取 $C_c^\infty(\R^N)$ 中的序列 $(u_n)$ 使得
$u_n\to u$ 在 $L^p(\Omega)$ 中且几乎处处在 $\Omega$ 上,
$\nabla u_n\to\nabla u$ 在 $L^p(\omega)^N$ 中对一切
$\omega\subset\subset\Omega$ 成立 (定理  \ref{thm:9.2}).我们有
\begin{equation*}
\int_\Omega(G\circ u_n)\frac{\partial\varphi}{\partial x_i}
=-\int_\Omega(G'\circ u_n)\frac{\partial u_n}{\partial x_i}\varphi
\quad\forall\varphi\in C_c^1(\Omega).
\end{equation*}
但 $G\circ u_n\to G\circ u$ 在 $L^p(\Omega)$ 中, 且
$(G'\circ u_n)\frac{\partial u_n}{\partial x_i}\to
(G'\circ u)\frac{\partial u}{\partial x_i}$ 在 $L^p(\omega)$ 中对一切
$\omega\subset\subset\Omega$ 成立 (由控制收敛), 于是 \eqref{eq:9.6}
得证.当 $p=\infty$ 时, 固定开集 $\Omega'$ 使得
$\supp\varphi\subset\Omega'\subset\subset\Omega$.则
$u\in W^{1,p}(\Omega')$ 对一切 $p<\infty$ 成立, 且 \eqref{eq:9.6} 由上述
结果推出.
\end{proof}

\begin{proposition}[变量替换公式]
\label{prop:9.6}
设 $\Omega$ 与 $\Omega'$ 为 $\R^N$ 中的两个开集, $H:\Omega'\to\Omega$ 为
双射, $x=H(y)$, 使得 $H\in C^1(\Omega')$, $H^{-1}\in C^1(\Omega)$,
$\operatorname{Jac}H\in L^\infty(\Omega')$,
$\operatorname{Jac}H^{-1}\in L^\infty(\Omega)$.\footnote{$\operatorname{Jac}H$
表示 Jacobi 矩阵 $\frac{\partial H_i}{\partial y_j}$; 于是它是
$L^\infty(\Omega')^{N\times N}$ 中的函数.} 设 $u\in W^{1,p}(\Omega)$,
$1\le p\le\infty$.则 $u\circ H\in W^{1,p}(\Omega')$ 且
\begin{equation*}
\frac{\partial}{\partial y_j}u(H(y))=\sum_i
\frac{\partial u}{\partial x_i}(H(y))\frac{\partial H_i}{\partial y_j}(y)
\quad\forall j=1,2,\dots,N.
\end{equation*}
\end{proposition}

\begin{proof}
当 $1\le p<\infty$ 时, 取 $C_c^\infty(\R^N)$ 中的序列 $(u_n)$ 使得
$u_n\to u$ 在 $L^p(\Omega)$ 中, 且 $\nabla u_n\to\nabla u$ 在
$L^p(\omega)^N$ 中对一切 $\omega\subset\subset\Omega$ 成立.于是
$u_n\circ H\to u\circ H$ 在 $L^p(\Omega')$ 中, 且
\begin{equation*}
\left(\frac{\partial u_n}{\partial x_i}\circ H\right)
\frac{\partial H_i}{\partial y_j}
\to\left(\frac{\partial u}{\partial x_i}\circ H\right)
\frac{\partial H_i}{\partial y_j}
\quad\text{在}\ L^p(\omega')\ \text{中对一切}\ \omega'\subset\subset\Omega'.
\end{equation*}
给定 $\psi\in C_c^1(\Omega')$, 我们有
\begin{equation*}
\int_{\Omega'}(u_n\circ H)\frac{\partial\psi}{\partial y_j}\,dy
=-\int_{\Omega'}\sum_i\left(\frac{\partial u_n}{\partial x_i}\circ H\right)
\frac{\partial H_i}{\partial y_j}\psi\,dy.
\end{equation*}
取极限即得所要结果.当 $p=\infty$ 时, 按命题  \ref{prop:9.5} 证明的末尾
同样的方式处理.
\end{proof}

\subsection*{空间 $W^{m,p}(\Omega)$}

设 $m\ge2$ 为整数, $p$ 为实数, $1\le p\le\infty$.我们归纳地定义
\begin{equation*}
W^{m,p}(\Omega)=\left\{u\in W^{m-1,p}(\Omega);\
\frac{\partial u}{\partial x_i}\in W^{m-1,p}(\Omega)\
\forall i=1,2,\dots,N\right\}.
\end{equation*}
等价地, 这些集合也可以这样引入:
\begin{equation*}
W^{m,p}(\Omega)=\left\{u\in L^p(\Omega)\ \middle|\
\begin{array}{l}
\forall\alpha,\ \abs{\alpha}\le m,\ \exists g_\alpha\in L^p(\Omega)\
\text{使得}\\
\displaystyle\int_\Omega uD^\alpha\varphi=(-1)^{\abs{\alpha}}
\int_\Omega g_\alpha\varphi\quad\forall\varphi\in C_c^\infty(\Omega)
\end{array}
\right\},
\end{equation*}
其中我们采用标准的多重指标记号
$\alpha=(\alpha_1,\alpha_2,\dots,\alpha_N)$, $\alpha_i\ge0$ 为整数,
\begin{equation*}
\abs{\alpha}=\sum_{i=1}^N\alpha_i
\quad\text{且}\quad
D^\alpha\varphi=\frac{\partial^{\abs{\alpha}}\varphi}
{\partial x_1^{\alpha_1}\partial x_2^{\alpha_2}\cdots\partial x_N^{\alpha_N}}.
\end{equation*}
我们令 $D^\alpha u=g_\alpha$.空间 $W^{m,p}(\Omega)$ 配以范数
\begin{equation*}
\norm{u}_{W^{m,p}}=\sum_{0\le\abs{\alpha}\le m}\norm{D^\alpha u}_p
\end{equation*}
是 Banach 空间.

空间 $H^m(\Omega)=W^{m,2}(\Omega)$ 配以标量积
\begin{equation*}
(u,v)_{H^m}=\sum_{0\le\abs{\alpha}\le m}(D^\alpha u,D^\alpha v)_{L^2}
\end{equation*}
是 Hilbert 空间.

\begin{remark}
\label{rem:9.8}
可以证明: 若 $\Omega$ ``足够光滑''且 $\Gamma=\partial\Omega$ 有界, 则
$W^{m,p}(\Omega)$ 上的范数等价于范数
\begin{equation*}
\norm{u}_p+\sum_{\abs{\alpha}=m}\norm{D^\alpha u}_p.
\end{equation*}
更精确地说, 可以证明: 对每个多重指标 $\alpha$, $0<\abs{\alpha}<m$, 以及
每个 $\varepsilon>0$, 存在常数 $C$ (依赖于 $\Omega,\varepsilon,\alpha$)
使得
\begin{equation*}
\norm{D^\alpha u}_p\le\varepsilon\sum_{\abs{\beta}=m}\norm{D^\beta u}_p
+C\norm{u}_p\quad\forall u\in W^{m,p}(\Omega)
\end{equation*}
(参见例如 R.~Adams \bref{Adams}{1}).
\end{remark}

\section{延拓算子}
\label{sec:9.2}

通过先考虑 $\Omega=\R^N$ 的情形来建立 $W^{1,p}(\Omega)$ 中函数的性质, 往往
是方便的 (例如见第~\ref{sec:9.3}~节的结果).因此, 能够把一个函数
$u\in W^{1,p}(\Omega)$ 延拓为函数 $\tilde u\in W^{1,p}(\R^N)$ 是有用的.
这并非总是可能的 (在一般区域 $\Omega$ 中).然而, 若 $\Omega$ ``光滑'',
这样的延拓可以构造出来.让我们先确切说明光滑开集的概念.

\begin{notation}
给定 $x\in\R^N$, 写
\begin{equation*}
x=(x',x_N),\quad x'\in\R^{N-1},\quad
x'=(x_1,x_2,\dots,x_{N-1}),
\end{equation*}
并令
\begin{equation*}
\abs{x'}=\left(\sum_{i=1}^{N-1}x_i^2\right)^{1/2}.
\end{equation*}
我们定义
\begin{align*}
&\R^N_+=\{x=(x',x_N);\ x_N>0\},\\
&Q=\{x=(x',x_N);\ \abs{x'}<1\ \text{且}\ \abs{x_N}<1\},\\
&Q_+=Q\cap\R^N_+,\\
&Q_0=\{x=(x',0);\ \abs{x'}<1\}.
\end{align*}
\end{notation}

\begin{definition}
我们说开集 $\Omega$ 属于 $C^1$ 类, 如果对每个 $x\in\partial\Omega=\Gamma$, 存在
$x$ 在 $\R^N$ 中的邻域 $U$ 以及双射 $H:Q\to U$, 使得
\begin{equation*}
H\in C^1(Q),\quad H^{-1}\in C^1(U),\quad
H(Q_+)=U\cap\Omega,\quad H(Q_0)=U\cap\Gamma.
\end{equation*}
映射 $H$ 称为\emph{局部坐标卡}.
\end{definition}

\begin{theorem}
\index{yantuosuanzi@延拓算子 (extension operator)}
\label{thm:9.7}
设 $\Omega$ 属于 $C^1$ 类且 $\Gamma$ 有界 (或者 $\Omega=\R^N_+$).则存在线性
延拓算子
\begin{equation*}
P:W^{1,p}(\Omega)\to W^{1,p}(\R^N)\quad(1\le p\le\infty),
\end{equation*}
使得对一切 $u\in W^{1,p}(\Omega)$,
\begin{equation*}
\begin{array}{ll}
Pu|_\Omega=u, & \text{(i)}\\[2pt]
\norm{Pu}_{L^p(\R^N)}\le C\norm{u}_{L^p(\Omega)}, & \text{(ii)}\\[2pt]
\norm{Pu}_{W^{1,p}(\R^N)}\le C\norm{u}_{W^{1,p}(\Omega)}, & \text{(iii)}
\end{array}
\end{equation*}
其中 $C$ 只依赖于 $\Omega$.
\end{theorem}

我们先证明一个关于反射延拓的简单而基本的引理.

\begin{lemma}
\label{lem:9.2}
给定 $u\in W^{1,p}(Q_+)$, $1\le p\le\infty$, 定义 $Q$ 上的函数
$u^\star$ 为反射延拓, 即
\begin{equation*}
u^\star(x',x_N)=
\begin{cases}
u(x',x_N) & \text{若}\ x_N>0,\\
u(x',-x_N) & \text{若}\ x_N<0.
\end{cases}
\end{equation*}
则 $u^\star\in W^{1,p}(Q)$ 且
\begin{equation*}
\norm{u^\star}_{L^p(Q)}\le2\norm{u}_{L^p(Q_+)},\quad
\norm{u^\star}_{W^{1,p}(Q)}\le2\norm{u}_{W^{1,p}(Q_+)}.
\end{equation*}
\end{lemma}

\begin{proof}
事实上, 我们将证明
\begin{equation}
\frac{\partial u^\star}{\partial x_i}
=\left(\frac{\partial u}{\partial x_i}\right)^\star
\quad\text{对}\ 1\le i\le N-1
\label{eq:9.7}
\end{equation}
且
\begin{equation}
\frac{\partial u^\star}{\partial x_N}
=\left(\frac{\partial u}{\partial x_N}\right)^{\square},
\label{eq:9.8}
\end{equation}
其中 $(\frac{\partial u}{\partial x_i})^\star$ 表示 $\frac{\partial u}
{\partial x_i}$ 的反射延拓, 且只要 $f$ 定义在 $Q_+$ 上, 我们就令
\begin{equation*}
f^{\square}(x',x_N)=
\begin{cases}
f(x',x_N) & \text{若}\ x_N>0,\\
-f(x',-x_N) & \text{若}\ x_N<0.
\end{cases}
\end{equation*}
我们将使用 $C^\infty(\R)$ 中的一列函数 $(\eta_k)$, 它定义为
\begin{equation*}
\eta_k(t)=\eta(kt),\quad t\in\R,\ k=1,2,\dots,
\end{equation*}
其中 $\eta$ 是任一固定函数, $\eta\in C^\infty(\R)$, 满足
\begin{equation*}
\eta(t)=
\begin{cases}
0 & \text{若}\ t<1/2,\\
1 & \text{若}\ t>1.
\end{cases}
\end{equation*}

\noindent\emph{\eqref{eq:9.7} 的证明.} 设 $\varphi\in C_c^1(Q)$.对
$1\le i\le N-1$, 我们有
\begin{equation}
\int_Q u^\star\frac{\partial\varphi}{\partial x_i}
=\int_{Q_+}u\frac{\partial\psi}{\partial x_i},
\label{eq:9.9}
\end{equation}
其中 $\psi(x',x_N)=\varphi(x',x_N)+\varphi(x',-x_N)$.函数 $\psi$ 一般
不属于 $C_c^1(Q_+)$, 因此不能作为检验函数 (在 $W^{1,p}$ 的定义中).另一方
面, $\eta_k(x_N)\psi(x',x_N)\in C_c^1(Q_+)$, 于是
\begin{equation*}
\int_{Q_+}u\frac{\partial}{\partial x_i}(\eta_k\psi)
=-\int_{Q_+}\frac{\partial u}{\partial x_i}\eta_k\psi.
\end{equation*}
由于 $\frac{\partial}{\partial x_i}(\eta_k\psi)=\eta_k
\frac{\partial\psi}{\partial x_i}$, 我们有
\begin{equation}
\int_{Q_+}u\eta_k\frac{\partial\psi}{\partial x_i}
=-\int_{Q_+}\frac{\partial u}{\partial x_i}\eta_k\psi.
\label{eq:9.10}
\end{equation}
在 \eqref{eq:9.10} 中令 $k\to\infty$ 取极限 (利用控制收敛), 得到
\begin{equation}
\int_{Q_+}u\frac{\partial\psi}{\partial x_i}
=-\int_{Q_+}\frac{\partial u}{\partial x_i}\psi.
\label{eq:9.11}
\end{equation}
结合 \eqref{eq:9.9} 与 \eqref{eq:9.11}, 我们得到
\begin{equation*}
\int_Q u^\star\frac{\partial\varphi}{\partial x_i}
=-\int_{Q_+}\frac{\partial u}{\partial x_i}\psi
=-\int_Q\left(\frac{\partial u}{\partial x_i}\right)^\star\varphi,
\end{equation*}
由此推出 \eqref{eq:9.7}.

\noindent\emph{\eqref{eq:9.8} 的证明.} 对每个 $\varphi\in C_c^1(Q)$, 我们有
\begin{equation}
\int_Q u^\star\frac{\partial\varphi}{\partial x_N}
=\int_{Q_+}u\frac{\partial\chi}{\partial x_N},
\label{eq:9.12}
\end{equation}
其中 $\chi(x',x_N)=\varphi(x',x_N)-\varphi(x',-x_N)$.注意
$\chi(x',0)=0$, 于是存在常数 $M$ 使得
$\abs{\chi(x',x_N)}\le M\abs{x_N}$ 在 $Q$ 上成立.由于
$\eta_k\chi\in C_c^1(Q_+)$, 我们有
\begin{equation}
\int_{Q_+}u\frac{\partial}{\partial x_N}(\eta_k\chi)
=-\int_{Q_+}\frac{\partial u}{\partial x_N}\eta_k\chi.
\label{eq:9.13}
\end{equation}
但
\begin{equation}
\frac{\partial}{\partial x_N}(\eta_k\chi)=\eta_k\frac{\partial\chi}
{\partial x_N}+k\eta'(kx_N)\chi.
\label{eq:9.14}
\end{equation}
我们断言
\begin{equation}
\int_{Q_+}uk\eta'(kx_N)\chi\to0\quad\text{当}\ k\to\infty.
\label{eq:9.15}
\end{equation}
事实上, 我们有
\begin{equation*}
\left|\int_{Q_+}uk\eta'(kx_N)\chi\right|
\le kMC\int_{0<x_N<1/k}\abs{u}x_N\,dx
\le MC\int_{0<x_N<1/k}\abs{u}\,dx,
\end{equation*}
其中 $C=\sup_{t\in[0,1]}\abs{\eta'(t)}$, 由此推出 \eqref{eq:9.15}.

由 \eqref{eq:9.13}, \eqref{eq:9.14} 与 \eqref{eq:9.15} 推出
\begin{equation*}
\int_{Q_+}u\frac{\partial\chi}{\partial x_N}
=-\int_{Q_+}\frac{\partial u}{\partial x_N}\chi.
\end{equation*}
最后, 我们有
\begin{equation}
\int_{Q_+}\frac{\partial u}{\partial x_N}\chi
=\int_Q\left(\frac{\partial u}{\partial x_N}\right)^{\square}\varphi.
\label{eq:9.16}
\end{equation}
结合 \eqref{eq:9.12} 与 \eqref{eq:9.16}, 得到 \eqref{eq:9.8}.这就完成了
引理  \ref{lem:9.2} 的证明.
\end{proof}

引理  \ref{lem:9.2} 的结论在把 $Q_+$ 换成 $\R^N_+$ 后仍然成立 (证明不变).
这就建立了 $\Omega=\R^N_+$ 情形的定理  \ref{thm:9.7}.

\begin{starremark}
\label{rem:9.9}
引理  \ref{lem:9.2} 对某些不属于 $C^1$ 类的开集 $\Omega$ 给出了延拓算子的
非常简单的构造.例如考虑正方形
\begin{equation*}
\Omega=\{x\in\R^2;\ 0<x_1<1,\ 0<x_2<1\}.
\end{equation*}
设 $u\in W^{1,p}(\Omega)$.通过四次相继的反射 (见图  \ref{fig:6}), 我们得到 $u$ 在
\begin{equation*}
\widetilde\Omega=\{x\in\R^2;\ -1<x_1<3,\ -1<x_2<3\}
\end{equation*}
中的延拓 $\tilde u\in W^{1,p}(\widetilde\Omega)$.

\begin{figure}[tbp]
	\centering
	\begin{tikzpicture}[scale=0.95]
		\draw[dashed] (0,0) rectangle (4,4);
		\draw[dashed] (2,0) -- (2,4);
		\draw[dashed] (1,0) -- (1,2);
		\draw[dashed] (0,2) -- (2,2);
		\draw (1,1) rectangle (2,2);
		\node at (1.28,1.48) {$\Omega$};
		\node at (0.45,1.15) {$2$};
		\node at (1.35,0.55) {$1$};
		\node at (1.12,3.15) {$3$};
		\node at (2.85,2.15) {$4$};
	\end{tikzpicture}
	\caption{}
	\label{fig:6}
\end{figure}

然后固定任一函数 $\psi\in C_c^1(\widetilde\Omega)$ 使得 $\psi=1$ 在
$\Omega$ 上.用 $Pu$ 表示把 $\psi\tilde u$ 在 $\widetilde\Omega$ 外延拓为
$0$ 所得的 $\R^2$ 上的函数.容易证明算子
$P:W^{1,p}(\Omega)\to W^{1,p}(\R^2)$ 满足 (i), (ii) 与 (iii).
\end{starremark}

下述引理非常有用.

\begin{lemma}[单位分解]
\label{lem:9.3}
设 $K$ 为 $\R^N$ 的紧子集, $U_1,U_2,\dots,U_k$ 为 $K$ 的开覆盖, 即
$K\subset\bigcup_{i=1}^kU_i$.则存在函数
$\theta_0,\theta_1,\theta_2,\dots,\theta_k\in C^\infty(\R^N)$ 使得
\begin{equation*}
0\le\theta_i\le1\quad\forall i=0,1,2,\dots,k
\quad\text{且}\quad\sum_{i=0}^k\theta_i=1\ \text{在}\ \R^N\ \text{上},
\end{equation*}
\begin{equation*}
\begin{array}{ll}
\supp\theta_i\ \text{紧且}\ \supp\theta_i\subset U_i
\quad\forall i=1,2,\dots,k, & \text{(i)}\\[2pt]
\supp\theta_0\subset\R^N\setminus K. & \text{(ii)}
\end{array}
\end{equation*}
若 $\Omega$ 为有界开集且 $K=\partial\Omega$, 则
$\theta_0|_\Omega\in C_c^\infty(\Omega)$.
\end{lemma}

\begin{proof}
这个引理是经典的; 类似的陈述可在例如 S.~Agmon \bref{Agmon}{1}, R.~Adams \bref{Adams}{1}, G.~Folland \bref{Folland}{1}, P.~Malliavin \bref{Malliavin}{1} 中找到.
\end{proof}

\begin{proof}[定理 \ref{thm:9.7}]
我们用局部坐标卡把 $\Gamma=\partial\Omega$ ``拉直'', 并使用单位分解.\footnote{下文中我们将经常使用这个技巧, 把在 $\R^N_+$ (或 $Q_+$) 上证明
的结果转移到光滑开集 $\Omega$ 上的同样结论.}更精确地说, 由于 $\Gamma$ 紧
且属于 $C^1$ 类, 存在 $\R^N$ 中的开集 $(U_i)_{1\le i\le k}$ 使得
$\Gamma\subset\bigcup_{i=1}^kU_i$, 以及双射 $H_i:Q\to U_i$ 使得
\begin{equation*}
H_i\in C^1(Q),\quad H_i^{-1}\in C^1(U_i),\quad
H_i(Q_+)=U_i\cap\Omega,\quad H_i(Q_0)=U_i\cap\Gamma.
\end{equation*}
考虑引理  \ref{lem:9.3} 中引入的函数
$\theta_0,\theta_1,\theta_2,\dots,\theta_k$.给定
$u\in W^{1,p}(\Omega)$, 写
\begin{equation*}
u=\sum_{i=0}^k\theta_iu=\sum_{i=0}^ku_i,
\end{equation*}
其中 $u_i=\theta_iu$.

现在我们把每个函数 $u_i$ 延拓到 $\R^N$, 对 $u_0$ 与 $(u_i)_{1\le i\le k}$
分别处理.

\noindent (a) \emph{$u_0$ 的延拓.} 我们用
\begin{equation*}
\bar u_0(x)=
\begin{cases}
u_0(x) & \text{若}\ x\in\Omega,\\
0 & \text{若}\ x\in\R^N\setminus\Omega
\end{cases}
\end{equation*}
定义 $u_0$ 到 $\R^N$ 的延拓.回忆
$\theta_0\in C^1(\R^N)\cap L^\infty(\R^N)$,
$\nabla\theta_0\in L^\infty(\R^N)$, 因为
$\nabla\theta_0=-\sum_{i=1}^k\nabla\theta_i$ 有紧支集, 且
$\supp\theta_0\subset\R^N\setminus\Gamma$.由此推出 (由注记
\ref{rem:9.4}(ii)) $\bar u_0\in W^{1,p}(\R^N)$ 且
\begin{equation*}
\frac{\partial}{\partial x_i}\bar u_0=\theta_0
\frac{\partial u}{\partial x_i}+\frac{\partial\theta_0}{\partial x_i}\bar u.
\end{equation*}
于是
\begin{equation*}
\norm{\bar u_0}_{W^{1,p}(\R^N)}\le C\norm{u}_{W^{1,p}(\Omega)}.
\end{equation*}

\noindent (b) \emph{$u_i$ 的延拓, $1\le i\le k$.} 考虑 $u$ 在
$U_i\cap\Omega$ 上的限制, 并借助 $H_i$ 把这个函数``转移''到 $Q_+$ 上.
更精确地说, 令 $v_i(y)=u(H_i(y))$, $y\in Q_+$.我们知道 (命题
 \ref{prop:9.6}) $v_i\in W^{1,p}(Q_+)$.然后用反射 (引理  \ref{lem:9.2})
定义 $v_i$ 在 $Q$ 上的延拓; 记之为 $v_i^\star$.我们知道
$v_i^\star\in W^{1,p}(Q)$.用 $H_i^{-1}$ 把 $v_i^\star$ ``转回''到 $U_i$
上, 记之为 $w_i$:
\begin{equation*}
w_i(x)=v_i^\star[H_i^{-1}(x)]\quad\text{对}\ x\in U_i.
\end{equation*}
则 $w_i\in W^{1,p}(U_i)$, $w_i=u$ 在 $U_i\cap\Omega$ 上, 且
\begin{equation*}
\norm{w_i}_{W^{1,p}(U_i)}\le C\norm{u}_{W^{1,p}(U_i\cap\Omega)}.
\end{equation*}
最后, 对 $x\in\R^N$ 令
\begin{equation*}
\hat u_i(x)=
\begin{cases}
\theta_i(x)w_i(x) & \text{若}\ x\in U_i,\\
0 & \text{若}\ x\in\R^N\setminus U_i,
\end{cases}
\end{equation*}
于是 $\hat u_i\in W^{1,p}(\R^N)$ (见注记 \ref{rem:9.4}(ii)),
$\hat u_i=u_i$ 在 $\Omega$ 上, 且
\begin{equation*}
\norm{\hat u_i}_{W^{1,p}(\R^N)}\le C\norm{u}_{W^{1,p}(U_i\cap\Omega)}.
\end{equation*}

\noindent (c) \emph{结论.} 算子
$Pu=\bar u_0+\sum_{i=1}^N\hat u_i$ 具有所有所要的性质.
\end{proof}

\begin{keycorollary}[稠密性]
\label{cor:9.8}
设 $\Omega$ 属于 $C^1$ 类, $u\in W^{1,p}(\Omega)$, $1\le p<\infty$.则存在
$C_c^\infty(\R^N)$ 中的序列 $(u_n)$ 使得 $u_n|_\Omega\to u$ 在
$W^{1,p}(\Omega)$ 中.换言之, $C_c^\infty(\R^N)$ 中函数在 $\Omega$ 上的限制
构成 $W^{1,p}(\Omega)$ 的稠密子空间.
\end{keycorollary}

\begin{proof}
先设 $\Gamma$ 有界.则由定理  \ref{thm:9.7} 存在延拓算子 $P$.序列\footnote{按
惯例, $(\rho_n)$ 是一列磨光核, $(\zeta_n)$ 是一列截断函数, 如定理
 \ref{thm:9.2} 的证明中那样.}
$\zeta_n(\rho_n\star Pu)$ 在 $W^{1,p}(\R^N)$ 中收敛到 $Pu$, 于是它解决了
问题.当 $\Gamma$ 无界时, 我们先考虑序列 $\zeta_nu$.给定 $\varepsilon>0$,
固定 $n_0$ 使得 $\norm{\zeta_{n_0}u-u}_{W^{1,p}}<\varepsilon$.然后可以构造
$\zeta_{n_0}u$ 的延拓 $v\in W^{1,p}(\R^N)$ (因为这只涉及 $\Gamma$ 与一个
大球的交集).最后任取 $w\in C_c^\infty(\R^N)$ 使得
$\norm{w-v}_{W^{1,p}(\R^N)}<\varepsilon$.
\end{proof}

\section{Sobolev 不等式}
\label{sec:9.3}

在第~
\ref{chap:8}~章中我们看到, 若 $\Omega$ 是一维的, 则对一切 $1\le p\le\infty$,
$W^{1,p}(\Omega)\subset L^\infty(\Omega)$ 且嵌入连续.在维数 $N\ge2$ 时,
这个包含只对 $p>N$ 成立; 当 $p\le N$ 时可以构造出不属于 $L^\infty$ 的
$W^{1,p}$ 函数 (见注记 \ref{rem:9.16}).然而, 一个本质上属于 Sobolev 的
重要结果表明: 若 $1\le p<N$, 则 $W^{1,p}(\Omega)\subset L^{p^\star}(\Omega)$
且嵌入连续, 其中某个 $p^\star\in(p,+\infty)$.这个结果通常称为
\emph{Sobolev 嵌入定理}.我们先考虑下述情形:

\runin{A. $\Omega=\R^N$ 的情形.}

\begin{keytheorem}[Sobolev, Gagliardo, Nirenberg]
\index{Sobolevbudengshi@Sobolev 不等式 (Sobolev inequality)}
\label{thm:9.9}
设 $1\le p<N$.则
\begin{equation*}
W^{1,p}(\R^N)\subset L^{p^\star}(\R^N),
\end{equation*}
其中 $p^\star$ 由
\begin{equation*}
\frac1{p^\star}=\frac1p-\frac1N
\end{equation*}
给出, 且存在常数\footnote{可以取 $C(p,N)=(N-1)p/(N-p)$, 但这个常数不是最优
的.最优常数是已知的 (但它并不简单!), 参见 Th.~Aubin \bref{Aubin}{1}, G.~Talenti \bref{Talenti}{1} 与 E.~Lieb \bref{Lieb}{1}.}
$C=C(p,N)$ 使得
\begin{equation}
\norm{u}_{p^\star}\le C\norm{\nabla u}_p
\quad\forall u\in W^{1,p}(\R^N).
\label{eq:9.17}
\end{equation}
\end{keytheorem}

\begin{remark}
\label{rem:9.10}
$p^\star$ 的值可以通过一个非常简单的伸缩论证得到 (物理学家所钟爱的伸缩论证
有时以最少的努力给出有用的信息).事实上, 假设存在常数 $C$ 与 $q$
($1\le q\le\infty$) 使得
\begin{equation}
\norm{u}_q\le C\norm{\nabla u}_p\quad\forall u\in C_c^\infty(\R^N).
\label{eq:9.18}
\end{equation}
则必有 $q=p^\star$.为看出这一点, 固定任一函数
$u\in C_c^\infty(\R^N)$, 把 $u_\lambda(x)=u(\lambda x)$ 代入
\eqref{eq:9.18}.我们得到
\begin{equation*}
\norm{u}_q\le C\lambda^{\left(1+\frac Nq-\frac Np\right)}
\norm{\nabla u}_p\quad\forall\lambda>0,
\end{equation*}
由此推出 $1+\frac Nq-\frac Np=0$, 即 $q=p^\star$ (只要 $u$ 不恒为零).
\end{remark}

定理  \ref{thm:9.9} 的证明依赖于下述引理:

\begin{lemma}
\label{lem:9.4}
设 $N\ge2$, $f_1,f_2,\dots,f_N\in L^{N-1}(\R^{N-1})$.对 $x\in\R^N$ 与
$1\le i\le N$, 令
\begin{equation*}
\tilde x_i=(x_1,x_2,\dots,x_{i-1},x_{i+1},\dots,x_N)\in\R^{N-1},
\end{equation*}
即从列表中删去 $x_i$.则函数
\begin{equation*}
f(x)=f_1(\tilde x_1)f_2(\tilde x_2)\cdots f_N(\tilde x_N),\quad x\in\R^N,
\end{equation*}
属于 $L^1(\R^N)$, 且
\begin{equation*}
\norm{f}_{L^1(\R^N)}\le\prod_{i=1}^N\norm{f_i}_{L^{N-1}(\R^{N-1})}.
\end{equation*}
\end{lemma}

\begin{proof}
$N=2$ 的情形是平凡的 (为什么?).让我们考虑 $N=3$ 的情形.我们有
\begin{align*}
\int_\R\abs{f(x)}\,dx_3
&=\abs{f_3(x_1,x_2)}\int_\R\abs{f_1(x_2,x_3)}\abs{f_2(x_1,x_3)}\,dx_3\\
&\le\abs{f_3(x_1,x_2)}\left(\int_\R\abs{f_1(x_2,x_3)}^2dx_3\right)^{1/2}
\left(\int_\R\abs{f_2(x_1,x_3)}^2dx_3\right)^{1/2}
\end{align*}
(由 Cauchy--Schwarz).再用一次 Cauchy--Schwarz 得到
\begin{equation*}
\int_{\R^3}\abs{f(x)}\,dx\le\norm{f_3}_{L^2(\R^2)}
\norm{f_1}_{L^2(\R^2)}\norm{f_2}_{L^2(\R^2)}.
\end{equation*}
一般情形用归纳法得到——先假设结论对 $N$ 成立, 再推出它对 $N+1$ 成立.
固定 $x_{N+1}\in\R$; 由 Hölder 不等式,
\begin{align*}
&\int_{\R^N}\abs{f(x)}\,dx_1dx_2\cdots dx_N\\
&\quad\le\norm{f_{N+1}}_{L^N(\R^N)}
\left(\int\abs{f_1f_2\cdots f_N}^{N'}dx_1dx_2\cdots dx_N\right)^{1/N'}
\end{align*}
(其中 $N'=N/(N-1)$).对函数 $\abs{f_1}^{N'},\abs{f_2}^{N'},\dots,
\abs{f_N}^{N'}$ 应用归纳假设, 我们得到
\begin{equation*}
\int_{\R^N}\abs{f_1}^{N'}\cdots\abs{f_N}^{N'}\,dx_1\cdots dx_N
\le\prod_{i=1}^N\norm{f_i}_{L^N(\R^{N-1})}^{N'},
\end{equation*}
由此推出
\begin{equation*}
\int_{\R^N}\abs{f(x)}\,dx_1\cdots dx_N
\le\norm{f_{N+1}}_{L^N(\R^N)}\prod_{i=1}^N
\norm{f_i}_{L^N(\R^{N-1})}.
\end{equation*}
现在让 $x_{N+1}$ 变化.每个函数
$x_{N+1}\mapsto\norm{f_i}_{L^N(\R^{N-1})}$ 都属于 $L^N(\R)$,
$1\le i\le N$.因此它们的乘积
$\prod_{i=1}^N\norm{f_i}_{L^N(\R^{N-1})}$ 属于 $L^1(\R)$ (见第~
\ref{chap:4}~章中
Hölder 不等式后的注记 2), 并且
\begin{equation*}
\int_{\R^{N+1}}\abs{f(x)}\,dx_1dx_2\cdots dx_Ndx_{N+1}
\le\prod_{i=1}^{N+1}\norm{f_i}_{L^N(\R^N)}.
\end{equation*}
\end{proof}

\begin{proof}[定理 \ref{thm:9.9}]
我们从 $p=1$ 且 $u\in C_c^1(\R^N)$ 的情形开始.我们有
\begin{align*}
\abs{u(x_1,x_2,\dots,x_N)}
&=\left|\int_{-\infty}^{x_1}\frac{\partial u}{\partial x_1}
(t,x_2,\dots,x_N)\,dt\right|\\
&\le\int_{-\infty}^{+\infty}\left|\frac{\partial u}{\partial x_1}
(t,x_2,\dots,x_N)\right|dt,
\end{align*}
类似地, 对每个 $1\le i\le N$,
\begin{equation*}
\abs{u(x_1,x_2,\dots,x_N)}\le\int_{-\infty}^{+\infty}
\left|\frac{\partial u}{\partial x_i}(x_1,x_2,\dots,x_{i-1},t,
x_{i+1},\dots,x_N)\right|dt
\overset{\text{def}}{=}f_i(\tilde x_i).
\end{equation*}
于是
\begin{equation*}
\abs{u(x)}^N\le\prod_{i=1}^Nf_i(\tilde x_i).
\end{equation*}
由引理  \ref{lem:9.4} 推出
\begin{equation*}
\int_{\R^N}\abs{u(x)}^{N/(N-1)}dx\le\prod_{i=1}^N
\norm{f_i}_{L^1(\R^{N-1})}^{1/(N-1)}
=\prod_{i=1}^N\left\|\frac{\partial u}{\partial x_i}\right\|_{L^1(\R^N)}^{1/(N-1)}.
\end{equation*}
因此我们有
\begin{equation}
\norm{u}_{L^{N/(N-1)}(\R^N)}\le\prod_{i=1}^N
\left\|\frac{\partial u}{\partial x_i}\right\|_{L^1(\R^N)}^{1/N}.
\label{eq:9.19}
\end{equation}
这就完成了 $p=1$ 且 $u\in C_c^1(\R^N)$ 情形下 \eqref{eq:9.17} 的证明.
现在转向 $1<p<N$ 的情形, 仍设 $u\in C_c^1(\R^N)$.设 $m\ge1$; 把
\eqref{eq:9.19} 用于 $\abs{u}^{m-1}u$ 而不是 $u$.我们得到
\begin{align}
\norm{u}_{mN/(N-1)}^m
&\le m\prod_{i=1}^N\left\|\abs{u}^{m-1}
\frac{\partial u}{\partial x_i}\right\|_1^{1/N}\notag\\
&\le m\norm{u}_{p'(m-1)}^{m-1}\prod_{i=1}^N
\left\|\frac{\partial u}{\partial x_i}\right\|_p^{1/N}.
\label{eq:9.20}
\end{align}
然后取 $m$ 使得 $mN/(N-1)=p'(m-1)$, 这给出
$m=(N-1)p^\star/N$ (由于 $1<p<N$, 有 $m\ge1$).我们得到
\begin{equation*}
\norm{u}_{p^\star}\le m\prod_{i=1}^N
\left\|\frac{\partial u}{\partial x_i}\right\|_p^{1/N},
\end{equation*}
从而
\begin{equation*}
\norm{u}_{p^\star}\le C\norm{\nabla u}_p\quad\forall u\in C^1(\R^N).
\end{equation*}
为完成证明, 设 $u\in W^{1,p}(\R^N)$, 并取 $C_c^1(\R^N)$ 中的序列
$(u_n)$ 使得 $u_n\to u$ 在 $W^{1,p}(\R^N)$ 中.必要时抽取子序列, 还可以假设
$u_n\to u$ 几乎处处成立.我们有
\begin{equation*}
\norm{u_n}_{p^\star}\le C\norm{\nabla u_n}_p.
\end{equation*}
由 Fatou 引理\footnote{也可以注意到 $(u_n)$ 是 $L^{p^\star}$ 中的 Cauchy
序列来得出结论.}推出
\begin{equation*}
u\in L^{p^\star}\quad\text{且}\quad\norm{u}_{p^\star}\le C\norm{\nabla u}_p.\textquoteleft
\end{equation*}
% 原书此处不等式末尾多印了一个孤立的左单引号 (见印刷页 281); 依技能规范
% 保留原样并记录, 不擅自更正.
\end{proof}

\begin{keycorollary}
\label{cor:9.10}
设 $1\le p<N$.则
\begin{equation*}
W^{1,p}(\R^N)\subset L^q(\R^N)\quad\forall q\in[p,p^\star]
\end{equation*}
且嵌入连续.
\end{keycorollary}

\begin{proof}
给定 $q\in[p,p^\star]$, 我们写
\begin{equation*}
\frac1q=\frac\alpha p+\frac{1-\alpha}{p^\star}
\quad\text{其中}\ \alpha\in[0,1].
\end{equation*}
我们知道 (见第~
\ref{chap:4}~章注记 2)
\begin{equation*}
\norm{u}_q\le\norm{u}_p^\alpha\norm{u}_{p^\star}^{1-\alpha}
\le\norm{u}_p+\norm{u}_{p^\star}
\end{equation*}
(由 Young 不等式).利用定理  \ref{thm:9.9}, 我们推出
\begin{equation*}
\norm{u}_q\le C\norm{u}_{W^{1,p}}\quad\forall u\in W^{1,p}(\R^N).
\end{equation*}
\end{proof}

\begin{keycorollary}[极限情形 $p=N$]
\label{cor:9.11}
我们有
\begin{equation*}
W^{1,p}(\R^N)\subset L^q(\R^N)\quad\forall q\in[N,+\infty).
\end{equation*}
\end{keycorollary}

\begin{proof}
设 $u\in C_c^1(\R^N)$; 把 \eqref{eq:9.20} 用于 $p=N$, 我们得到
\begin{equation*}
\norm{u}_{mN/(N-1)}^m\le m\norm{u}_{(m-1)N/(N-1)}^{m-1}
\norm{\nabla u}_N\quad\forall m\ge1,
\end{equation*}
借助 Young 不等式, 我们有
\begin{equation}
\norm{u}_{mN/(N-1)}\le C\bigl(\norm{u}_{(m-1)N/(N-1)}
+\norm{\nabla u}_N\bigr)\quad\forall m\ge1.
\label{eq:9.21}
\end{equation}
在 \eqref{eq:9.21} 中先取 $m=N$; 它成为
\begin{equation*}
\norm{u}_{N^2/(N-1)}\le C\norm{u}_{W^{1,N}},
\end{equation*}
由插值不等式 (见第~
\ref{chap:4}~章注记 2) 我们有
\begin{equation}
\norm{u}_q\le C\norm{u}_{W^{1,N}}
\label{eq:9.22}
\end{equation}
对一切满足 $N\le q\le N^2/(N-1)$ 的 $q$ 成立.用 $m=N+1$, $m=N+2$ 等
反复迭代这个论证, 我们得到
\begin{equation}
\norm{u}_q\le C\norm{u}_{W^{1,N}}\quad\forall u\in C^1(\R^N)
\label{eq:9.23}
\end{equation}
对一切 $q\in[N,+\infty)$ 成立, 其中常数 $C$ 依赖于 $q$ 与 $N$.\footnote{当
$q\to+\infty$ 时这个常数``爆破''.}不等式 \eqref{eq:9.23} 由稠密性延拓到
$W^{1,N}$.
\end{proof}

\begin{keytheorem}[Morrey]
\index{Morreybudengshi@Morrey 不等式 (Morrey inequality)}
\label{thm:9.12}
设 $p>N$.则
\begin{equation}
W^{1,p}(\R^N)\subset L^\infty(\R^N)
\label{eq:9.24}
\end{equation}
且嵌入连续.此外, 对一切 $u\in W^{1,p}(\R^N)$, 我们有
\begin{equation}
\abs{u(x)-u(y)}\le C\abs{x-y}^\alpha\norm{\nabla u}_p
\quad\text{对几乎一切}\ x,y\in\R^N,
\label{eq:9.25}
\end{equation}
其中 $\alpha=1-(N/p)$, $C$ 为常数 (只依赖于 $p$ 与 $N$).
\end{keytheorem}

\begin{remark}
\label{rem:9.11}
不等式 \eqref{eq:9.25} 蕴含存在函数 $\tilde u\in C(\R^N)$ 使得
$u=\tilde u$ 在 $\R^N$ 上几乎处处成立.(事实上, 设 $A\subset\R^N$ 为零测集
使得 \eqref{eq:9.25} 对 $x,y\in\R^N\setminus A$ 成立; 由于
$\R^N\setminus A$ 在 $\R^N$ 中稠密, 函数 $u|_{\R^N\setminus A}$ 有 (唯一的)
连续延拓到 $\R^N$.)换言之, 每个 $u\in W^{1,p}(\R^N)$ ($p>N$) 都有连续
代表元.在有用的时候, 我们用它的连续代表元代替 $u$, 并且也用 $u$ 表示这个
连续代表元.
\end{remark}

\begin{proof}[定理 \ref{thm:9.12}]
我们先对 $u\in C_c^1(\R^N)$ 建立 \eqref{eq:9.25}.设 $Q$ 为包含 $0$ 的开
立方体, 其边——长为 $r$——平行于坐标轴.对 $x\in Q$ 我们有
\begin{equation*}
u(x)-u(0)=\int_0^1\frac{d}{dt}u(tx)\,dt,
\end{equation*}
于是
\begin{equation}
\abs{u(x)-u(0)}\le\int_0^1\sum_{i=1}^N\abs{x_i}
\left|\frac{\partial u}{\partial x_i}(tx)\right|dt
\le r\sum_{i=1}^N\int_0^1\left|\frac{\partial u}{\partial x_i}(tx)\right|dt.
\label{eq:9.26}
\end{equation}
令
\begin{equation*}
\bar u=\frac{1}{\abs{Q}}\int_Qu(x)\,dx=(\text{$u$ 在 $Q$ 上的平均值}).
\end{equation*}
在 $Q$ 上积分 \eqref{eq:9.26}, 我们得到
\begin{align*}
\abs{\bar u-u(0)}
&\le\frac{r}{\abs{Q}}\int_Qdx\sum_{i=1}^N\int_0^1
\left|\frac{\partial u}{\partial x_i}(tx)\right|dt\\
&=\frac{1}{r^{N-1}}\int_0^1dt\int_Q\sum_{i=1}^N
\left|\frac{\partial u}{\partial x_i}(tx)\right|dx\\
&=\frac{1}{r^{N-1}}\int_0^1dt\int_{tQ}\sum_{i=1}^N
\left|\frac{\partial u}{\partial x_i}(y)\right|\frac{dy}{t^N}.
\end{align*}
于是由 Hölder 不等式, 我们有
\begin{equation*}
\int_{tQ}\left|\frac{\partial u}{\partial x_i}(y)\right|dy
\le\left(\int_Q\left|\frac{\partial u}{\partial x_i}\right|^p\right)^{1/p}
\abs{tQ}^{1/p'}
\end{equation*}
(因为当 $t\in(0,1)$ 时 $tQ\subset Q$).由此推出
\begin{equation*}
\abs{\bar u-u(0)}\le\frac{1}{r^{N-1}}\norm{\nabla u}_{L^p(Q)}
r^{N/p'}\int_0^1\frac{t^{N/p'}}{t^N}\,dt
=\frac{r^{1-(N/p)}}{1-(N/p)}\norm{\nabla u}_{L^p(Q)}.
\end{equation*}
由平移, 这个不等式对一切边——长为 $r$——平行于坐标轴的立方体 $Q$ 仍然成立.
于是我们有
\begin{equation}
\abs{\bar u-u(x)}\le\frac{r^{1-(N/p)}}{1-(N/p)}\norm{\nabla u}_{L^p(Q)}
\quad\forall x\in Q.
\label{eq:9.27}
\end{equation}
把这两个不等式相加 (并利用三角不等式), 我们得到
\begin{equation}
\abs{u(x)-u(y)}\le\frac{2r^{1-(N/p)}}{1-(N/p)}
\norm{\nabla u}_{L^p(Q)}\quad\forall x,y\in Q.
\label{eq:9.28}
\end{equation}
给定 $\R^N$ 中任意两点 $x,y$, 存在这样的立方体 $Q$, 其边长
$r=2\abs{x-y}$ 且包含 $x$ 与 $y$.这就对 $u\in C_c^1(\R^N)$ 推出了
\eqref{eq:9.25}.对一般函数 $u\in W^{1,p}(\R^N)$, 我们使用
$C_c^1(\R^N)$ 中的序列 $(u_n)$ 使得 $u_n\to u$ 在 $W^{1,p}(\R^N)$ 中且
$u_n\to u$ 几乎处处成立.

现在我们证明 \eqref{eq:9.24}.设 $u\in C_c^1(\R^N)$, $x\in\R^N$, $Q$ 为
边长为 $r=1$ 且包含 $x$ 的立方体.由 \eqref{eq:9.27} 与 Hölder 不等式,
我们有
\begin{equation*}
\abs{u(x)}\le\abs{\bar u}+C\norm{\nabla u}_{L^p(Q)}
\le C\norm{u}_{W^{1,p}(Q)}\le C\norm{u}_{W^{1,p}(\R^N)},
\end{equation*}
其中 $C$ 只依赖于 $p$ 与 $N$.于是
\begin{equation*}
\norm{u}_{L^\infty(\R^N)}\le C\norm{u}_{W^{1,p}(\R^N)}
\quad\forall u\in C_c^1(\R^N).
\end{equation*}
对一般函数 $u\in W^{1,p}(\R^N)$, 我们使用标准的稠密性论证.
\end{proof}

\begin{remark}
\label{rem:9.12}
由 \eqref{eq:9.24} 推出: 若 $u\in W^{1,p}(\R^N)$, $N<p<\infty$, 则
\begin{equation*}
\lim_{\abs{x}\to\infty}u(x)=0.
\end{equation*}
事实上, 存在 $C_c^1(\R^N)$ 中的序列 $(u_n)$ 使得 $u_n\to u$ 在
$W^{1,p}(\R^N)$ 中.由 \eqref{eq:9.24}, $u$ 也是 $u_n$ 在 $\R^N$ 上的
一致极限.
\end{remark}

\begin{keycorollary}
\label{cor:9.13}
设 $m\ge1$ 为整数, $p\in[1,+\infty)$.我们有
\begin{equation*}
W^{m,p}(\R^N)\subset L^q(\R^N),
\end{equation*}
其中
\begin{equation*}
\frac1q=\frac1p-\frac mN\quad\text{若}\ \frac1p-\frac mN>0,
\end{equation*}
\begin{equation*}
W^{m,p}(\R^N)\subset L^q(\R^N)\quad\forall q\in[p,+\infty)
\quad\text{若}\ \frac1p-\frac mN=0,
\end{equation*}
\begin{equation*}
W^{m,p}(\R^N)\subset L^\infty(\R^N)\quad\text{若}\
\frac1p-\frac mN<0,
\end{equation*}
且所有这些嵌入都连续.此外, 若 $m-(N/p)>0$ 不是整数, 令\footnote{$[\ ]$
表示整数部分.}
\begin{equation*}
k=[m-(N/p)]\quad\text{且}\quad
\theta=m-(N/p)-k\ (0<\theta<1).
\end{equation*}
则对一切 $u\in W^{m,p}(\R^N)$,
\begin{equation*}
\norm{D^\alpha u}_{L^\infty(\R^N)}\le C\norm{u}_{W^{m,p}(\R^N)}
\quad\forall\alpha,\ \abs{\alpha}\le k,
\end{equation*}
并且\footnote{这蕴含对一切 $\abs{\alpha}<k$, $D^\alpha u$ 是 Lipschitz
连续的, 即
\begin{equation*}
\abs{D^\alpha u(x)-D^\alpha u(y)}\le C\norm{u}_{W^{m,p}}\abs{x-y}
\quad\text{对几乎一切}\ x,y\in\R^N.
\end{equation*}}
\begin{equation*}
\abs{D^\alpha u(x)-D^\alpha u(y)}\le C\norm{u}_{W^{m,p}(\R^N)}
\abs{x-y}^\theta\quad\text{对几乎一切}\ x,y\in\R^N,\
\forall\alpha,\ \abs{\alpha}=k.
\end{equation*}
特别地, $W^{m,p}(\R^N)\subset C^k(\R^N)$.\footnote{这应当理解为在连续
代表元的选择意义下成立.}
\end{keycorollary}

\begin{proof}
所有这些结果都由定理  \ref{thm:9.9}, 推论  \ref{cor:9.11} 与定理
 \ref{thm:9.12} 反复应用得到.
\end{proof}

\begin{remark}
\label{rem:9.13}
$p=1$ 且 $m=N$ 的情形是特殊的.我们有 $W^{N,1}\subset L^\infty$.
(但一般说来, 当 $p>1$ 且 $m=N/p$ 时 $W^{m,p}\subset L^\infty$ 并不成立.)
事实上, 对 $u\in C_c^\infty(\R^N)$, 我们有
\begin{equation*}
u(x_1,x_2,\dots,x_N)=\int_{-\infty}^{x_1}\int_{-\infty}^{x_2}
\cdots\int_{-\infty}^{x_N}\frac{\partial^Nu}
{\partial x_1\partial x_2\cdots\partial x_N}(t_1,t_2,\dots,t_N)
\,dt_1dt_2\cdots dt_N,
\end{equation*}
于是
\begin{equation}
\norm{u}_{L^\infty}\le C\norm{u}_{W^{N,1}}
\quad\forall u\in C_c^\infty(\R^N).
\label{eq:9.29}
\end{equation}
一般函数 $u\in W^{N,1}$ 的情形由稠密性推出.
\end{remark}

现在让我们转向下述情形.

\runin{B. $\Omega\subset\R^N$ 的情形.}

这里我们假设 $\Omega$ 是 $C^1$ 类开集且 $\Gamma$ 有界, 或者 $\Omega=\R^N_+$.

\begin{keycorollary}
\label{cor:9.14}
设 $1\le p\le\infty$.我们有
\begin{equation*}
W^{1,p}(\Omega)\subset L^{p^\star}(\Omega),
\end{equation*}
其中
\begin{equation*}
\frac1{p^\star}=\frac1p-\frac1N,\quad\text{若}\ p<N,
\end{equation*}
\begin{equation*}
W^{1,p}(\Omega)\subset L^q(\Omega)\quad\forall q\in[p,+\infty),
\quad\text{若}\ p=N,
\end{equation*}
\begin{equation*}
W^{1,p}(\Omega)\subset L^\infty(\Omega),\quad\text{若}\ p>N,
\end{equation*}
且所有这些嵌入都连续.此外, 若 $p>N$, 则对一切
$u\in W^{1,p}(\Omega)$,
\begin{equation*}
\abs{u(x)-u(y)}\le C\norm{u}_{W^{1,p}}\abs{x-y}^\alpha
\quad\text{对几乎一切}\ x,y\in\Omega,
\end{equation*}
其中 $\alpha=1-(N/p)$, $C$ 只依赖于 $\Omega$, $p$ 与 $N$.特别地,
$W^{1,p}(\Omega)\subset C(\overline\Omega)$.\footnote{再强调一次, 这应当
理解为在连续代表元的选择意义下成立.}
\end{keycorollary}

\begin{proof}
考虑延拓算子 $P:W^{1,p}(\Omega)\to W^{1,p}(\R^N)$ (见定理
 \ref{thm:9.7}), 然后应用定理  \ref{thm:9.9}, 推论  \ref{cor:9.11} 与定理
 \ref{thm:9.12}.
\end{proof}

\begin{keycorollary}
\label{cor:9.15}
把 $\R^N$ 换成 $\Omega$ 后, 推论  \ref{cor:9.13} 的结论仍然成立.\footnote{确切
地说, 若 $m-(N/p)>0$ 不是整数, 则 $W^{m,p}(\Omega)\subset C^k(\overline\Omega)$,
其中 $k=[m-(N/p)]$, 且
\begin{equation*}
C^k(\overline\Omega)=\{u\in C^k(\Omega);\ D^\alpha u\ \text{对一切}
\ \abs{\alpha}\le k\ \text{在}\ \overline\Omega\ \text{上有连续延拓}\}.
\end{equation*}}
\end{keycorollary}

\begin{proof}
反复应用推论  \ref{cor:9.14}.\footnote{另一种做法是把推论  \ref{cor:9.13}
与延拓算子 $P:W^{m,p}(\Omega)\to W^{m,p}(\R^N)$ 结合起来使用, 但这需要
一个额外的假设: 为构造这个延拓算子, $\Omega$ 必须属于 $C^m$ 类.}
\end{proof}

\begin{keytheorem}[Rellich--Kondrachov]
\label{thm:9.16}
设 $\Omega$ 有界且属于 $C^1$ 类.则下述嵌入是紧的:
\begin{equation*}
W^{1,p}(\Omega)\subset L^q(\Omega)\quad\forall q\in[1,p^\star),
\end{equation*}
其中
\begin{equation*}
\frac1{p^\star}=\frac1p-\frac1N,\quad\text{若}\ p<N,
\end{equation*}
\begin{equation*}
W^{1,p}(\Omega)\subset L^q(\Omega)\quad\forall q\in[p,+\infty),
\quad\text{若}\ p=N,
\end{equation*}
\begin{equation*}
W^{1,p}(\Omega)\subset C(\overline\Omega),\quad\text{若}\ p>N.
\end{equation*}
特别地, 对一切 $p$ (以及一切 $N$), $W^{1,p}(\Omega)\subset L^p(\Omega)$
是紧嵌入.
\end{keytheorem}

\begin{proof}
$p>N$ 的情形由推论  \ref{cor:9.14} 与 Ascoli--Arzelà 定理推出.$p=N$ 的
情形归结为 $p<N$ 的情形.因此只剩下 $p<N$ 的情形.

设 $H$ 为 $W^{1,p}(\Omega)$ 中的单位球.设 $P$ 为定理  \ref{thm:9.7} 的延拓
算子.令 $F=P(H)$, 于是 $H=F|_\Omega$.为证明对 $q\in[1,p^\star)$, $H$
在 $L^q(\Omega)$ 中有紧闭包, 我们引用定理  \ref{thm:4.26}.由于 $\Omega$
有界, 我们总可以假设 $q\ge p$.显然 $F$ 在 $W^{1,p}(\R^N)$ 中有界, 因此
由推论  \ref{cor:9.10} 它在 $L^q(\R^N)$ 中也有界.我们要验证
\begin{equation*}
\lim_{\abs{h}\to0}\norm{\tau_hf-f}_{L^q(\R^N)}=0
\quad\text{关于}\ f\in F\ \text{一致}.
\end{equation*}
由命题  \ref{prop:9.3} 我们有
\begin{equation*}
\norm{\tau_hf-f}_{L^p(\R^N)}\le\abs{h}\norm{\nabla f}_{L^p(\R^N)}
\quad\forall f\in F.
\end{equation*}
由于 $p\le q<p^\star$, 我们可以写
\begin{equation*}
\frac1q=\frac\alpha p+\frac{1-\alpha}{p^\star}
\quad\text{其中}\ \alpha\in(0,1].
\end{equation*}
借助插值不等式 (见第~
\ref{chap:4}~章注记 2), 我们有
\begin{align*}
\norm{\tau_hf-f}_{L^q(\R^N)}
&\le\norm{\tau_hf-f}_{L^p(\R^N)}^\alpha
\norm{\tau_hf-f}_{L^{p^\star}(\R^N)}^{1-\alpha}\\
&\le\abs{h}^\alpha\norm{\nabla f}_{L^p(\R^N)}^\alpha
\bigl(2\norm{f}_{L^{p^\star}(\R^N)}\bigr)^{1-\alpha}\le C\abs{h}^\alpha,
\end{align*}
其中 $C$ 与 $F$ 无关 (因为 $F$ 在 $W^{1,p}(\R^N)$ 中有界).所要结论由此
推出.
\end{proof}

\begin{remark}
\label{rem:9.14}
定理  \ref{thm:9.16} 在下述意义下是``几乎最优的'':

\noindent (i) 若 $\Omega$ 无界, 嵌入 $W^{1,p}(\Omega)\subset L^p(\Omega)$
一般不是紧的.\footnote{关于无界区域上这个嵌入的紧性, 参见 R.~Adams
\bref{Adams}{1} 第 167 页的详细讨论.}

\noindent (ii) 嵌入 $W^{1,p}(\Omega)\subset L^{p^\star}(\Omega)$ 从来不是
紧的, 即使 $\Omega$ 有界且光滑.
\end{remark}

\begin{starremark}
\label{rem:9.15}
设 $\Omega$ 为 $C^1$ 类有界开集.则只要
\begin{equation*}
1\le q\le p^\star\ \text{若}\ 1\le p<N,\quad
1\le q<\infty\ \text{若}\ p=N,\quad
1\le q\le\infty\ \text{若}\ p>N,
\end{equation*}
范数
\begin{equation*}
|||u|||=\norm{\nabla u}_p+\norm{u}_q
\end{equation*}
就与 $W^{1,p}$ 范数等价.
\end{starremark}

\begin{starremark}[极限情形 $p=N$]
\label{rem:9.16}
设 $\Omega$ 为 $C^1$ 类有界开集, $u\in W^{1,N}(\Omega)$.则一般说来
$u\notin L^\infty(\Omega)$.例如, 若
\begin{equation*}
\Omega=\{x\in\R^N;\ \abs{x}<1/2\},
\end{equation*}
则函数
\begin{equation*}
u(x)=(\log 1/\abs{x})^\alpha
\quad\text{其中}\ 0<\alpha<1-(1/N)
\end{equation*}
属于 $W^{1,N}(\Omega)$, 但由于 $x=0$ 处的奇异性它不是有界的.尽管如此, 我们
有 Trudinger 不等式
\begin{equation*}
\int_\Omega e^{\abs{u}^{N/(N-1)}}<\infty\quad\forall u\in W^{1,N}(\Omega)
\end{equation*}
(参见例如 R.~Adams \bref{Adams}{1} 或 D.~Gilbarg--N.~Trudinger
\bref{GilbargTrudinger}{1}).
\end{starremark}

\section{空间 $W_0^{1,p}(\Omega)$}
\label{sec:9.4}

\begin{definition}
设 $1\le p<\infty$; $W_0^{1,p}(\Omega)$ 表示 $C_c^1(\Omega)$ 在
$W^{1,p}(\Omega)$ 中的闭包.令\footnote{有歧义时, 我们将写 $W_0^{1,p}$, $H_0^1$ 而不写 $W_0^{1,p}(\Omega)$, $H_0^1(\Omega)$.}
\begin{equation*}
H_0^1(\Omega)=W_0^{1,2}(\Omega).
\end{equation*}
空间 $W_0^{1,p}$ 配以 $W^{1,p}$ 范数是可分 Banach 空间; 当 $1<p<\infty$ 时
它自反.$H_0^1$ 配以 $H^1$ 标量积是 Hilbert 空间.
\end{definition}

\begin{starremark}
\label{rem:9.17}
由于 $C_c^1(\R^N)$ 在 $W^{1,p}(\R^N)$ 中稠密, 我们有
\begin{equation*}
W_0^{1,p}(\R^N)=W^{1,p}(\R^N).
\end{equation*}
相反, 若 $\Omega\subset\R^N$ 且 $\Omega\ne\R^N$, 则一般说来
$W_0^{1,p}(\Omega)\ne W^{1,p}(\Omega)$.然而, 若 $\R^N\setminus\Omega$
``足够薄''且 $p<N$, 则 $W_0^{1,p}(\Omega)=W^{1,p}(\Omega)$.例如, 若
$\Omega=\R^N\setminus\{0\}$ 且 $N\ge2$, 则可以证明
$H_0^1(\Omega)=H^1(\Omega)$.
\end{starremark}

\begin{remark}
\label{rem:9.18}
利用一列磨光核容易验证 $C_c^\infty(\Omega)$ 在 $W_0^{1,p}(\Omega)$ 中
稠密.换言之, 在 $W_0^{1,p}(\Omega)$ 的定义中本来同样可以用
$C_c^\infty(\Omega)$ 代替 $C_c^1(\Omega)$.

$W_0^{1,p}(\Omega)$ 中的函数``大致''就是那些``在 $\Gamma=\partial\Omega$
上为零''的 $W^{1,p}(\Omega)$ 函数.要确切说明这一点是微妙的, 因为函数
$u\in W^{1,p}(\Omega)$ 只在几乎处处意义下定义 (而 $\Gamma$ 的测度为零!),
并且 $u$ 未必有连续代表元.\footnote{尽管如此, 若
$u\in W^{1,p}(\Omega)$ (当 $\Omega$ 正则时), 可以赋予 $u|_\Gamma$ 一个
意义, 并且可以证明例如 $u|_\Gamma\in L^p(\Gamma)$.这依赖于迹理论 (见本章
末的评注).}下述刻画表明我们``确实''得到了``在 $\Gamma$ 上为零''的函数.
我们先从一个简单事实开始:
\end{remark}

\begin{lemma}
\label{lem:9.5}
设 $u\in W^{1,p}(\Omega)$, $1\le p<\infty$, 并假设 $\supp u$ 是 $\Omega$
的紧子集.则 $u\in W_0^{1,p}(\Omega)$.
\end{lemma}

\begin{proof}
固定开集 $\omega$ 使得 $\supp u\subset\omega\subset\subset\Omega$, 并选取
$\alpha\in C_c^1(\omega)$ 使得 $\alpha=1$ 在 $\supp u$ 上; 于是
$\alpha u=u$.另一方面 (定理  \ref{thm:9.2}), 存在 $C_c^\infty(\R^N)$ 中的
序列 $(u_n)$ 使得 $u_n\to u$ 在 $L^p(\Omega)$ 中且
$\nabla u_n\to\nabla u$ 在 $L^p(\omega)^N$ 中.由此推出
$\alpha u_n\to\alpha u$ 在 $W^{1,p}(\Omega)$ 中.于是
$\alpha u$ 属于 $W_0^{1,p}(\Omega)$, $u$ 也如此.
\end{proof}

\begin{theorem}
\label{thm:9.17}
设 $\Omega$ 属于 $C^1$ 类.设\footnote{回忆若 $p>N$, 则
$u\in W^{1,p}(\Omega)\Rightarrow u\in C(\overline\Omega)$ (见推论
 \ref{cor:9.14}).}
\begin{equation*}
u\in W^{1,p}(\Omega)\cap C(\overline\Omega),
\end{equation*}
其中 $1\le p<\infty$.则下述性质等价:

\noindent (i) $u=0$ 在 $\Gamma$ 上;

\noindent (ii) $u\in W_0^{1,p}(\Omega)$.
\end{theorem}

\begin{proof}
\noindent (i)$\Rightarrow$(ii). 先设 $\supp u$ 有界.固定函数
$G\in C^1(\R)$ 使得
\begin{equation*}
\abs{G(t)}\le\abs{t}\quad\forall t\in\R
\quad\text{且}\quad
G(t)=
\begin{cases}
0 & \text{若}\ \abs{t}\le1,\\
t & \text{若}\ \abs{t}\ge2.
\end{cases}
\end{equation*}
则 $u_n=(1/n)G(nu)$ 属于 $W^{1,p}$ (由命题  \ref{prop:9.5}).容易验证
(利用控制收敛) $u_n\to u$ 在 $W^{1,p}$ 中.另一方面,
\begin{equation*}
\supp u_n\subset\{x\in\Omega;\ \abs{u(x)}\ge1/n\},
\end{equation*}
于是 $\supp u_n$ 是含于 $\Omega$ 的紧集.由引理  \ref{lem:9.5},
$u_n\in W_0^{1,p}$, 由此推出 $u\in W_0^{1,p}$.在 $\supp u$ 无界的一般
情形, 考虑序列 $(\zeta_nu)$ (其中 $(\zeta_n)$ 是定理  \ref{thm:9.2} 的证明
中的一列截断函数).由上述结果 $\zeta_nu\in W_0^{1,p}$, 且由于
$\zeta_nu\to u$ 在 $W^{1,p}$ 中, 我们推出 $u\in W_0^{1,p}$.

\noindent (ii)$\Rightarrow$(i). 利用局部坐标卡, 这归结为下述问题.设
$u\in W_0^{1,p}(Q_+)\cap C(\overline{Q_+})$; 证明 $u=0$ 在 $Q_0$ 上.

设 $(u_n)$ 为 $C_c^1(Q_+)$ 中的序列使得 $u_n\to u$ 在 $W^{1,p}(Q_+)$ 中.
对 $(x',x_N)\in Q_+$, 我们有
\begin{equation*}
\abs{u_n(x',x_N)}\le\int_0^{x_N}\left|\frac{\partial u_n}{\partial x_N}
(x',t)\right|dt,
\end{equation*}
于是对 $0<\varepsilon<1$,
\begin{equation*}
\frac1\varepsilon\int_{\abs{x'}<1}\int_0^\varepsilon
\abs{u_n(x',x_N)}\,dx'\,dx_N
\le\int_{\abs{x'}<1}\int_0^\varepsilon
\left|\frac{\partial u_n}{\partial x_N}(x',t)\right|dx'\,dt.
\end{equation*}
取极限, 当 $n\to\infty$ 时 ($\varepsilon>0$ 固定) 我们得到
\begin{equation*}
\frac1\varepsilon\int_{\abs{x'}<1}\int_0^\varepsilon
\abs{u(x',x_N)}\,dx'\,dx_N
\le\int_{\abs{x'}<1}\int_0^\varepsilon
\left|\frac{\partial u}{\partial x_N}(x',t)\right|dx'\,dt.
\end{equation*}
最后, 令 $\varepsilon\to0$, 我们得到
\begin{equation*}
\int_{\abs{x'}<1}\abs{u(x',0)}\,dx'=0
\end{equation*}
(因为 $u\in C(\overline{Q_+})$ 且 $\frac{\partial u}{\partial x_N}
\in L^1(Q_+)$).于是 $u=0$ 在 $Q_0$ 上.
\end{proof}

\begin{remark}
\label{rem:9.19}
在 (i)$\Rightarrow$(ii) 的证明中我们没有用到 $\Omega$ 的光滑性.然而,
逆命题 (ii)$\Rightarrow$(i) 需要对 $\Omega$ 加光滑性假设 (例如考虑
$\Omega=\R^N\setminus\{0\}$, $N\ge2$, $p\le N$).
\end{remark}

下面是 $W_0^{1,p}$ 的另一个刻画.

\begin{proposition}
\label{prop:9.18}
设 $\Omega$ 属于 $C^1$ 类.设 $u\in L^p(\Omega)$, $1<p<\infty$.下述性质
等价:

\noindent (i) $u\in W_0^{1,p}(\Omega)$;

\noindent (ii) 存在常数 $C$ 使得
\begin{equation*}
\left|\int_\Omega u\frac{\partial\varphi}{\partial x_i}\right|
\le C\norm{\varphi}_{L^{p'}(\Omega)}
\quad\forall\varphi\in C_c^1(\R^N),\ \forall i=1,2,\dots,N;
\end{equation*}

\noindent (iii) 函数
\begin{equation*}
\bar u(x)=
\begin{cases}
u(x) & \text{若}\ x\in\Omega,\\
0 & \text{若}\ x\in\R^N\setminus\Omega
\end{cases}
\end{equation*}
属于 $W^{1,p}(\R^N)$, 并且此时
$\frac{\partial\bar u}{\partial x_i}=\frac{\partial u}{\partial x_i}$.
\end{proposition}

\begin{proof}
\noindent (i)$\Rightarrow$(ii). 设 $(u_n)$ 为 $C_c^1(\Omega)$ 中的序列
使得 $u_n\to u$ 在 $W^{1,p}$ 中.对 $\varphi\in C_c^1(\R^N)$ 我们有
\begin{equation*}
\left|\int_\Omega u_n\frac{\partial\varphi}{\partial x_i}\right|
=\left|\int_\Omega\frac{\partial u_n}{\partial x_i}\varphi\right|
\le\left\|\frac{\partial u_n}{\partial x_i}\right\|_p
\norm{\varphi}_{p'}.
\end{equation*}
取极限即得 (ii).

\noindent (ii)$\Rightarrow$(iii). 设 $\varphi\in C_c^1(\R^N)$; 我们有
\begin{equation*}
\left|\int_{\R^N}\bar u\frac{\partial\varphi}{\partial x_i}\right|
=\left|\int_\Omega u\frac{\partial\varphi}{\partial x_i}\right|
\le C\norm{\varphi}_{L^{p'}(\Omega)}
\le C\norm{\varphi}_{L^{p'}(\R^N)}.
\end{equation*}
于是 $\bar u\in W^{1,p}(\R^N)$ (由命题  \ref{prop:9.3}).

\noindent (iii)$\Rightarrow$(i). 总可以假设 $\Omega$ 有界 (否则考虑序列
$(\zeta_nu)$).利用局部坐标卡与单位分解, 这归结为下述问题.设
$u\in L^p(Q_+)$ 使得函数
\begin{equation*}
\bar u(x)=
\begin{cases}
u(x) & \text{若}\ x\in Q,\ x_N>0,\\
0 & \text{若}\ x\in Q,\ x_N<0
\end{cases}
\end{equation*}
属于 $W^{1,p}(Q)$; 证明
\begin{equation*}
\alpha u\in W_0^{1,p}(Q_+)\quad\forall\alpha\in C_c^1(Q).
\end{equation*}
设 $(\rho_n)$ 为一列磨光核使得
\begin{equation*}
\supp\rho_n\subset\left\{x\in\R^N;\ \frac1{2n}<x_N<\frac1n\right\};
\end{equation*}
例如可以取 $\rho_n(x)=n^N\rho(nx)$, 其中
$\supp\rho\subset\{x\in\R^N;\ (1/2)<x_N<1\}$.于是
$\rho_n\star(\alpha\bar u)\to\alpha\bar u$ 在 $W^{1,p}(\R^N)$ 中 (注意
$\alpha\bar u$ 在 $Q$ 外延拓为 $0$ 后属于 $W^{1,p}(\R^N)$).另一方面,
\begin{equation*}
\supp(\rho_n\star\alpha\bar u)\subset\supp\rho_n+
\supp(\alpha\bar u)\subset Q_+
\end{equation*}
对足够大的 $n$ 成立.由此推出
\begin{equation*}
\rho_n\star(\alpha\bar u)\in C_c^1(Q_+),
\end{equation*}
于是 $\alpha u\in W_0^{1,p}(Q_+)$.
\end{proof}

\begin{remark}
\label{rem:9.20}
推论  \ref{cor:9.14} 的证明使用了延拓算子, 因此必须假设 $\Omega$ 光滑.
若把 $W^{1,p}(\Omega)$ 换成 $W_0^{1,p}(\Omega)$, 则可以使用在 $\Omega$ 外
零延拓的典范延拓, 它对任意区域 $\Omega$ 都有效 (在命题  \ref{prop:9.18}
的证明中, 蕴涵 (i)$\Rightarrow$(iii) 没有用到 $\Omega$ 的任何光滑性假设).
特别地, 推论  \ref{cor:9.14} 的结论对任意开集 $\Omega$ 的
$W_0^{1,p}(\Omega)$ 成立.类似地, 定理  \ref{thm:9.16} 的结论对任意有界
开集 $\Omega$ 的 $W_0^{1,p}(\Omega)$ 成立.由定理  \ref{thm:9.9} 还可以推出:
若 $\Omega$ 是任意开集且 $1\le p<N$, 则
\begin{equation}
\norm{u}_{L^{p^\star}(\Omega)}\le C(p,N)\norm{\nabla u}_{L^p(\Omega)}
\quad\forall u\in W_0^{1,p}(\Omega).
\label{eq:9.30}
\end{equation}
\end{remark}

\begin{keycorollary}[Poincaré 不等式]
\label{cor:9.19}
设 $1\le p<\infty$, $\Omega$ 为有界开集.则存在常数 $C$ (依赖于 $\Omega$
与 $p$) 使得
\begin{equation*}
\norm{u}_{L^p(\Omega)}\le C\norm{\nabla u}_{L^p(\Omega)}
\quad\forall u\in W_0^{1,p}(\Omega).
\end{equation*}
特别地, 表达式 $\norm{\nabla u}_{L^p(\Omega)}$ 是 $W_0^{1,p}(\Omega)$ 上的
范数, 且与范数 $\norm{u}_{W^{1,p}}$ 等价; 在 $H_0^1(\Omega)$ 上, 表达式
$\sum_{i=1}^N\int_\Omega\frac{\partial u}{\partial x_i}
\frac{\partial v}{\partial x_i}$ 是标量积, 它诱导的范数为
$\norm{\nabla u}_{L^2}$, 且与范数 $\norm{u}_{H^1}$ 等价.
\end{keycorollary}

\begin{remark}
\label{rem:9.21}
若 $\Omega$ 有有限测度, 或者 $\Omega$ 在某个轴上的投影有界, Poincaré
不等式仍然成立.
\end{remark}

\begin{remark}
\label{rem:9.22}
对每个整数 $m\ge1$ 与 $1\le p<\infty$, 定义 $W_0^{m,p}(\Omega)$ 为
$C_c^m(\Omega)$ 在 $W^{m,p}(\Omega)$ 中的闭包.``大致''说来, 函数
$u$ 属于 $W_0^{m,p}(\Omega)$, 如果 $u\in W^{m,p}(\Omega)$ 且对一切满足
$\abs{\alpha}\le m-1$ 的多重指标 $\alpha$, $D^\alpha u=0$ 在 $\Gamma$ 上.
重要的是要注意到, 当 $m\ge2$ 时 $W_0^{m,p}(\Omega)$ 与
$W^{m,p}(\Omega)\cap W_0^{1,p}(\Omega)$ 是有区别的.
\end{remark}

\subsection*{$W_0^{1,p}(\Omega)$ 的对偶空间}

\begin{notation}
我们用 $W^{-1,p'}(\Omega)$ 表示 $W_0^{1,p}(\Omega)$ 的对偶空间
($1\le p<\infty$), 用 $H^{-1}(\Omega)$ 表示 $H_0^1(\Omega)$ 的对偶.
$L^2(\Omega)$ 的对偶与 $L^2(\Omega)$ 等同, 但我们不把 $H_0^1(\Omega)$ 与
它的对偶等同 (见第~
\ref{chap:5}~章注记 3).我们有包含关系
\begin{equation*}
H_0^1(\Omega)\subset L^2(\Omega)\subset H^{-1}(\Omega),
\end{equation*}
其中这些嵌入连续且稠密.

若 $\Omega$ 有界, 则
\begin{equation*}
W_0^{1,p}(\Omega)\subset L^2(\Omega)\subset W^{-1,p'}(\Omega)
\quad\text{若}\ 2N/(N+2)\le p<\infty,
\end{equation*}
且嵌入连续且稠密.若 $\Omega$ 无界, 同样的结论成立, 但只对范围
$2N/(N+2)\le p\le2$.
\end{notation}

$W^{-1,p'}$ 中的元素由下述结果完全刻画:

\begin{proposition}
\label{prop:9.20}
设 $F\in W^{-1,p'}(\Omega)$.则存在函数
$f_0,f_1,f_2,\dots,f_N\in L^{p'}(\Omega)$ 使得
\begin{equation*}
\dual{F}{v}=\int_\Omega f_0v+\sum_{i=1}^N\int_\Omega
f_i\frac{\partial v}{\partial x_i}\quad\forall v\in W_0^{1,p}(\Omega),
\end{equation*}
且
\begin{equation*}
\norm{F}=\max_{0\le i\le N}\norm{f_i}_{p'}.
\end{equation*}
若 $\Omega$ 有界, 我们可以取 $f_0=0$.
\end{proposition}

\begin{proof}
仿照命题  \ref{prop:8.14} 的证明.
\end{proof}

\section{\index{tuoyuanbianzhiwenti@椭圆边值问题 (elliptic boundary value problem)}
一些边值问题的变分表述}
\label{sec:9.5}

现在我们将利用前面的框架来研究一些二阶椭圆型偏微分方程 (= PDE).

\begin{example}[Laplace 算子的齐次 Dirichlet 问题]
\label{ex:9.1}
设 $\Omega\subset\R^N$ 为有界开集.我们寻找函数 $u:\Omega\to\R$ 满足
\begin{equation}
\begin{cases}
-\Delta u+u=f & \text{在}\ \Omega\ \text{中},\\
u=0 & \text{在}\ \Gamma=\partial\Omega\ \text{上},
\end{cases}
\label{eq:9.31}
\end{equation}
其中
\begin{equation*}
\Delta u=\sum_{i=1}^N\frac{\partial^2u}{\partial x_i^2}
=\text{$u$ 的 Laplace 算子},
\end{equation*}
而 $f$ 是 $\Omega$ 上给定的函数.边界条件 $u=0$ 在 $\Gamma$ 上称为
(\emph{齐次}) \emph{Dirichlet 条件}.
\end{example}

\begin{definition}
\eqref{eq:9.31} 的\emph{古典解}是满足 \eqref{eq:9.31} (在通常意义下) 的
函数 $u\in C^2(\Omega)$.\eqref{eq:9.31} 的\emph{弱解}是满足
\begin{equation}
\int_\Omega\nabla u\cdot\nabla v+\int_\Omega uv=\int_\Omega fv
\quad\forall v\in H_0^1(\Omega)
\label{eq:9.32}
\end{equation}
的函数 $u\in H_0^1(\Omega)$, 其中
$\nabla u\cdot\nabla v=\sum_{i=1}^N\frac{\partial u}{\partial x_i}
\frac{\partial v}{\partial x_i}$.
\end{definition}

我们实施第~
\ref{chap:8}~章中描述的程序.

\runin{步骤 A: 每个古典解都是弱解.}

事实上, $u\in H^1(\Omega)\cap C(\overline\Omega)$ 且 $u=0$ 在 $\Gamma$ 上,
于是由定理  \ref{thm:9.17} (另见注记 \ref{rem:9.19})
$u\in H_0^1(\Omega)$.另一方面, 若 $v\in C_c^1(\Omega)$ 我们有
\begin{equation*}
\int_\Omega\nabla u\cdot\nabla v+\int_\Omega uv=\int_\Omega fv,
\end{equation*}
由稠密性, 它对一切 $v\in H_0^1(\Omega)$ 仍然成立.

\runin{步骤 B: 弱解的存在性与唯一性.}

这就是下述基本结果的内容.

\begin{keytheorem}[Dirichlet, Riemann, Poincaré, Hilbert]
\label{thm:9.21}
给定任意 $f\in L^2(\Omega)$, \eqref{eq:9.31} 存在唯一的弱解
$u\in H_0^1(\Omega)$.此外, $u$ 由
\begin{equation*}
\min_{v\in H_0^1(\Omega)}\left\{\frac12\int_\Omega
(\abs{\nabla v}^2+\abs{v}^2)-\int_\Omega fv\right\}
\end{equation*}
得到.这就是\emph{Dirichlet 原理}.
\end{keytheorem}

\begin{proof}
在 Hilbert 空间 $H=H_0^1(\Omega)$ 中对双线性形式
\begin{equation*}
a(u,v)=\int_\Omega(\nabla u\cdot\nabla v+uv)
\end{equation*}
与线性泛函 $\varphi:v\mapsto\int_\Omega fv$ 应用 Lax--Milgram 定理.
\end{proof}

\runin{步骤 C: \index{ruojiedezhengzexing@弱解的正则性 (regularity of weak solutions)}
弱解的正则性.}

这个问题是微妙的.我们将在第~\ref{sec:9.6}~节中讨论它.

\runin{步骤 D: 复原古典解.}

假设 \eqref{eq:9.31} 的弱解 $u\in H_0^1(\Omega)$ 属于 $C^2(\Omega)$,
并假设 $\Omega$ 属于 $C^1$ 类.则 $u=0$ 在 $\Gamma$ 上 (由定理
 \ref{thm:9.17}).另一方面, 我们有
\begin{equation*}
\int_\Omega(-\Delta u+u)v=\int_\Omega fv\quad\forall v\in C_c^1(\Omega),
\end{equation*}
于是 $-\Delta u+u=f$ 在 $\Omega$ 上几乎处处成立 (由推论  \ref{cor:4.24}).
事实上, 由于 $u\in C^2(\Omega)$, $-\Delta u+u=f$ 在 $\Omega$ 上处处成立;
于是 $u$ 是古典解.

我们现在描述另外一些例子.在每个情形中, 精确地指明函数空间与适当的弱
表述都是至关重要的.

\begin{example}[非齐次 Dirichlet 条件]
\label{ex:9.2}
设 $\Omega\subset\R^N$ 为有界开集.我们寻找函数 $u:\Omega\to\R$ 满足
\begin{equation}
\begin{cases}
-\Delta u+u=f & \text{在}\ \Omega\ \text{中},\\
u=g & \text{在}\ \Gamma\ \text{上},
\end{cases}
\label{eq:9.33}
\end{equation}
其中 $f$ 在 $\Omega$ 上给定, $g$ 在 $\Gamma$ 上给定.假设存在函数
$\tilde g\in H^1(\Omega)\cap C(\overline\Omega)$ 使得\footnote{例如, 若
$\Omega$ 属于 $C^1$ 类且 $g\in C^1(\Gamma)$, 这个假设就满足.若 $\Omega$
足够正则, 则不必假设 $\tilde g\in C(\overline\Omega)$.应用迹理论 (见本章
末的评注), 只须知道 $\tilde g\in H^1(\Omega)$, 即
$g\in H^{1/2}(\Gamma)$.} $\tilde g=g$ 在 $\Gamma$ 上, 并考虑集合
\begin{equation*}
K=\{v\in H^1(\Omega);\ v-\tilde g\in H_0^1(\Omega)\}.
\end{equation*}
由定理  \ref{thm:9.17} 推出 $K$ 与 $\tilde g$ 的选取无关, 只依赖于 $g$.
$K$ 是 $H^1(\Omega)$ 中的非空闭凸集.
\end{example}

\begin{definition}
\eqref{eq:9.33} 的\emph{古典解}是满足 \eqref{eq:9.33} 的函数
$u\in C^2(\Omega)$.\eqref{eq:9.33} 的\emph{弱解}是满足
\begin{equation}
\int_\Omega(\nabla u\cdot\nabla v+uv)=\int_\Omega fv
\quad\forall v\in H_0^1(\Omega)
\label{eq:9.34}
\end{equation}
的函数 $u\in K$.
\end{definition}

如上所述, 任何古典解都是弱解.

\begin{keyproposition}
\label{prop:9.22}
给定任意 $f\in L^2(\Omega)$, \eqref{eq:9.33} 存在唯一的弱解
$u\in K$.此外, $u$ 由
\begin{equation*}
\min_{v\in K}\left\{\frac12\int_\Omega
(\abs{\nabla v}^2+v^2)-\int_\Omega fv\right\}
\end{equation*}
得到.
\end{keyproposition}

\begin{proof}
我们断言: $u\in K$ 是 \eqref{eq:9.33} 的弱解当且仅当
\begin{equation}
\int_\Omega\nabla u\cdot(\nabla v-\nabla u)+\int_\Omega u(v-u)
\ge\int_\Omega f(v-u)\quad\forall v\in K.
\label{eq:9.35}
\end{equation}
事实上, 若 $u$ 是 \eqref{eq:9.33} 的弱解, 显然 \eqref{eq:9.35} 成立, 甚至
取等号.反之, 若 $u\in K$ 满足 \eqref{eq:9.35}, 在 \eqref{eq:9.35} 中取
$v=u\pm w$, 其中 $w\in H_0^1(\Omega)$, 就得到 \eqref{eq:9.34}.然后我们可以
应用 Stampacchia 定理 (定理  \ref{thm:5.6}) 来完成证明.
\end{proof}

正则性与古典解的复原的研究遵循与例 \ref{ex:9.1} 相同的模式.

\begin{example}[一般二阶椭圆方程]
\label{ex:9.3}
设 $\Omega\subset\R^N$ 为有界开集.给定函数
$a_{ij}(x)\in C^1(\Omega)$, $1\le i,j\le N$, 满足椭圆性条件
\begin{equation}
\sum_{i,j=1}^Na_{ij}(x)\xi_i\xi_j\ge\alpha\abs{\xi}^2,
\quad\forall x\in\Omega,\ \forall\xi\in\R^N,
\label{eq:9.36}
\end{equation}
其中 $\alpha>0$.还给定函数 $a_0\in C(\Omega)$.我们寻找函数
$u:\Omega\to\R$ 满足
\begin{equation}
\begin{cases}
\displaystyle-\sum_{i,j=1}^N\frac{\partial}{\partial x_j}
\left(a_{ij}\frac{\partial u}{\partial x_i}\right)+a_0u=f
& \text{在}\ \Omega\ \text{中},\\[2mm]
u=0 & \text{在}\ \Gamma\ \text{上}.
\end{cases}
\label{eq:9.37}
\end{equation}
\eqref{eq:9.37} 的\emph{古典解}是满足 \eqref{eq:9.37} (在通常意义下) 的
函数 $u\in C^2(\Omega)$.\eqref{eq:9.37} 的\emph{弱解}是满足
\begin{equation}
\int_\Omega\sum_{i,j=1}^Na_{ij}\frac{\partial u}{\partial x_i}
\frac{\partial v}{\partial x_j}+\int_\Omega a_0uv=\int_\Omega fv
\quad\forall v\in H_0^1(\Omega)
\label{eq:9.38}
\end{equation}
的函数 $u\in H_0^1(\Omega)$.

如上所述, 任何古典解都是弱解.若 $a_0(x)\ge0$ 在 $\Omega$ 上, 则对一切
$f\in L^2(\Omega)$ 存在唯一的弱解 $u\in H_0^1$: 只须在空间
$H=H_0^1$ 中对连续双线性形式
\begin{equation*}
a(u,v)=\int_\Omega\sum_{i,j=1}^Na_{ij}\frac{\partial u}{\partial x_i}
\frac{\partial v}{\partial x_j}+\int_\Omega a_0uv
\end{equation*}
应用 Lax--Milgram 定理.$a(\ ,\ )$ 的强制性来自椭圆性假设, 假设
$a_0\ge0$ 以及 Poincaré 不等式.若矩阵 $(a_{ij})$ 还是对称的, 则形式
$a(\ ,\ )$ 对称, 且 $u$ 由
\begin{equation*}
\min_{v\in H_0^1}\left\{\frac12\int_\Omega\left(
\sum_{i,j=1}^Na_{ij}\frac{\partial v}{\partial x_i}
\frac{\partial v}{\partial x_j}+a_0v^2\right)-\int_\Omega fv\right\}
\end{equation*}
得到.

现在我们考虑一个更一般的问题: 寻找函数 $u:\Omega\to\R$ 满足
\begin{equation}
\begin{cases}
\displaystyle-\sum_{i,j}\frac{\partial}{\partial x_j}
\left(a_{ij}\frac{\partial u}{\partial x_i}\right)
+\sum_ia_i\frac{\partial u}{\partial x_i}+a_0u=f
& \text{在}\ \Omega\ \text{中},\\[2mm]
u=0 & \text{在}\ \Gamma\ \text{上},
\end{cases}
\label{eq:9.39}
\end{equation}
其中函数 $(a_{ij})\in L^\infty(\Omega)$ 满足椭圆性条件, 而函数
$(a_i)_{0\le i\le N}$ 给定在 $L^\infty(\Omega)$ 中.\eqref{eq:9.39} 的
\emph{弱解}是使得
\begin{equation}
\int_\Omega\sum_{i,j}a_{ij}\frac{\partial u}{\partial x_i}
\frac{\partial v}{\partial x_j}+\int_\Omega\sum_ia_i
\frac{\partial u}{\partial x_i}v+\int_\Omega a_0uv=\int_\Omega fv
\quad\forall v\in H_0^1
\label{eq:9.40}
\end{equation}
的函数 $u\in H_0^1$.相应的连续双线性形式为
\begin{equation}
a(u,v)=\int_\Omega\sum_{i,j}a_{ij}\frac{\partial u}{\partial x_i}
\frac{\partial v}{\partial x_j}+\int_\Omega\sum_ia_i
\frac{\partial u}{\partial x_i}v+\int_\Omega a_0uv.
\label{eq:9.41}
\end{equation}
一般说来这个形式不是对称的;\footnote{在维数 $N$ 中没有像一维那样已知的
办法把它化为对称情形.}在某些情形下它是强制的: 这时可以用 Lax--Milgram
定理得到弱解的存在性与唯一性.在一般情形——即使没有强制性——我们仍有下述
结果.
\end{example}

\begin{theorem}
\label{thm:9.23}
若 $f=0$, 则 \eqref{eq:9.40} 的解 $u\in H_0^1$ 的集合是有限维向量空间,
设其维数为 $d$.此外, 存在 $d$ 维子空间 $F\subset L^2(\Omega)$ 使得\footnote{
换言之, \eqref{eq:9.40} 有解当且仅当 $f$ 满足 $d$ 个正交性条件.}
\begin{equation*}
[\eqref{eq:9.40}\ \text{有解}]
\iff\left[\int_\Omega fv=0\quad\forall v\in F\right].
\end{equation*}
\end{theorem}

\begin{remark}
\label{rem:9.23}
假设与 \eqref{eq:9.40} 相联系的齐次方程 (即 $f=0$ 的情形) 以 $u=0$ 为
唯一解.则对每个 $f\in L^2$, \eqref{eq:9.40} 存在唯一的解
$u\in H_0^1$.\footnote{注意椭圆问题中解的存在性与唯一性之间的密切关系.
这个值得注意的关系是 Fredholm 二择一定理 (定理  \ref{thm:6.6}) 的推论.}
特别地, 若 $a_0\ge0$ 在 $\Omega$ 上, 则可以用极大值原理型的方法证明
$f=0\Rightarrow u=0$.于是我们推出, 只在假设 $a_0\ge0$ 在 $\Omega$ 上
(而对 $a_i$, $1\le i\le N$ 不作任何假设) 的条件下, 对每个 $f\in L^2$,
\eqref{eq:9.40} 都存在唯一的解 $u\in H_0^1$; 参见例如
D.~Gilbarg--N.~Trudinger \bref{GilbargTrudinger}{1}.
\end{remark}

\begin{proof}[定理 \ref{thm:9.23}]
固定 $\lambda>0$, 使得双线性形式
\begin{equation*}
a(u,v)+\lambda\int_\Omega uv
\end{equation*}
在 $H_0^1$ 上强制.对每个 $f\in L^2$ 存在唯一的 $u\in H_0^1$ 满足
\begin{equation*}
a(u,\varphi)+\lambda\int_\Omega u\varphi=\int_\Omega f\varphi
\quad\forall\varphi\in H_0^1.
\end{equation*}
令 $u=Tf$, 于是 $T:L^2\to L^2$ 是紧线性算子 (由于 $\Omega$ 有界, 嵌入
$H_0^1\subset L^2$ 是紧的; 见定理  \ref{thm:9.16} 与注记 \ref{rem:9.20}).
方程 \eqref{eq:9.40} 等价于
\begin{equation}
u=T(f+\lambda u).
\label{eq:9.42}
\end{equation}
令 $v=f+\lambda u$ 为新未知量, \eqref{eq:9.42} 成为
\begin{equation}
v-\lambda Tv=f.
\label{eq:9.43}
\end{equation}
结论由 Fredholm 二择一定理推出.
\end{proof}

\begin{example}[齐次 Neumann 问题]
\label{ex:9.4}
设 $\Omega\subset\R^N$ 为 $C^1$ 类有界区域.我们寻找函数
$u:\Omega\to\R$ 满足
\begin{equation}
\begin{cases}
-\Delta u+u=f & \text{在}\ \Omega\ \text{中},\\[1mm]
\dfrac{\partial u}{\partial n}=0 & \text{在}\ \Gamma\ \text{上},
\end{cases}
\label{eq:9.44}
\end{equation}
其中 $f$ 在 $\Omega$ 上给定; $\frac{\partial u}{\partial n}$ 表示 $u$ 的
外法向导数, 即 $\frac{\partial u}{\partial n}=\nabla u\cdot n$, 其中
$n$ 是 $\Gamma$ 的单位外法向量.边界条件 $\frac{\partial u}{\partial n}=0$
在 $\Gamma$ 上称为 (\emph{齐次}) \emph{Neumann 条件}.
\end{example}

\begin{definition}
\eqref{eq:9.44} 的\emph{古典解}是满足 \eqref{eq:9.44} 的函数
$u\in C^2(\overline\Omega)$.\eqref{eq:9.44} 的\emph{弱解}是满足
\begin{equation}
\int_\Omega\nabla u\cdot\nabla v+\int_\Omega uv=\int_\Omega fv
\quad\forall v\in H^1(\Omega)
\label{eq:9.45}
\end{equation}
的函数 $u\in H^1(\Omega)$.
\end{definition}

\runin{步骤 A: 每个古典解都是弱解.}

回忆由 Green 公式我们有
\begin{equation}
\int_\Omega(\Delta u)v=\int_\Gamma\frac{\partial u}{\partial n}v\,d\sigma
-\int_\Omega\nabla u\cdot\nabla v
\quad\forall u\in C^2(\overline\Omega),\ \forall v\in C^1(\overline\Omega),
\label{eq:9.46}
\end{equation}
其中 $d\sigma$ 是 $\Gamma$ 上的面测度.若 $u$ 是 \eqref{eq:9.44} 的古典
解, 则 $u\in H^1(\Omega)$, 且我们有
\begin{equation*}
\int_\Omega\nabla u\cdot\nabla v+\int_\Omega uv=\int_\Omega fv
\quad\forall v\in C^1(\overline\Omega).
\end{equation*}
由稠密性 (推论  \ref{cor:9.8}) 我们推出
\begin{equation*}
\int_\Omega\nabla u\cdot\nabla v+\int_\Omega uv=\int_\Omega fv
\quad\forall v\in H^1(\Omega).
\end{equation*}

\runin{步骤 B: 弱解的存在性与唯一性.}

\begin{proposition}
\label{prop:9.24}
对每个 $f\in L^2(\Omega)$, \eqref{eq:9.44} 存在唯一的弱解
$u\in H^1(\Omega)$.此外, $u$ 由
\begin{equation*}
\min_{v\in H^1(\Omega)}\left\{\frac12\int_\Omega
(\abs{\nabla v}^2+v^2)-\int_\Omega fv\right\}
\end{equation*}
得到.
\end{proposition}

\begin{proof}
在 $H=H^1(\Omega)$ 中应用 Lax--Milgram 定理.
\end{proof}

\runin{步骤 C: 弱解的正则性.}

这将在第~\ref{sec:9.6}~节中讨论.

\runin{步骤 D: 复原古典解.}

若 $u\in C^2(\overline\Omega)$ 是 \eqref{eq:9.44} 的弱解, 则由
\eqref{eq:9.46} 我们有
\begin{equation}
\int_\Omega(-\Delta u+u)v+\int_\Gamma\frac{\partial u}{\partial n}
v\,d\sigma=\int_\Omega fv\quad\forall v\in C^1(\overline\Omega).
\label{eq:9.47}
\end{equation}
在 \eqref{eq:9.47} 中先取 $v\in C_c^1(\Omega)$, 推出
\begin{equation*}
-\Delta u+u=f\quad\text{在}\ \Omega\ \text{中}.
\end{equation*}
然后回到 \eqref{eq:9.47}, 取 $v\in C^1(\overline\Omega)$; 得到
\begin{equation*}
\int_\Gamma\frac{\partial u}{\partial n}v\,d\sigma=0
\quad\forall v\in C^1(\overline\Omega),
\end{equation*}
因此 $\frac{\partial u}{\partial n}=0$ 在 $\Gamma$ 上.

\begin{example}[无界区域]
\label{ex:9.5}
当 $\Omega$ 是 $\R^N$ 中的无界开集时, 除了 $\Gamma=\partial\Omega$ 上的
通常边界条件之外, 还要在无穷远处加一个边界条件, 例如
$u(x)\to0$ 当 $\abs{x}\to\infty$.在弱解的水平上, 这``翻译''为条件
$u\in H^1$.当然, 必须首先证明: 若 $u$ 是古典解且
$u(x)\to0$ 当 $\abs{x}\to\infty$, 则 $u$ 必属于 $H^1$ (见第~
\ref{chap:8}~章例
\ref{ex:8.8} 中的讨论).下面是几个典型的例子:

\noindent (a) $\Omega=\R^N$; 给定 $f\in L^2(\R^N)$, 方程
\begin{equation*}
-\Delta u+u=f\quad\text{在}\ \R^N\ \text{中}
\end{equation*}
在下述意义下有唯一的弱解:
\begin{equation*}
u\in H^1(\R^N)\quad\text{且}\quad
\int_{\R^N}\nabla u\cdot\nabla v+\int_{\R^N}uv=\int_{\R^N}fv
\quad\forall v\in H^1(\R^N).
\end{equation*}

\noindent (b) $\Omega=\R^N_+$; 给定 $f\in L^2(\R^N_+)$, 问题
\begin{equation*}
\begin{cases}
-\Delta u+u=f & \text{在}\ \R^N_+\ \text{中},\\
u(x',0)=0 & \text{对}\ x'\in\R^{N-1}
\end{cases}
\end{equation*}
在下述意义下有唯一的弱解:
\begin{equation*}
u\in H_0^1(\Omega)\quad\text{且}\quad
\int_\Omega\nabla u\cdot\nabla v+\int_\Omega uv=\int_\Omega fv
\quad\forall v\in H_0^1(\Omega).
\end{equation*}

\noindent (c) $\Omega=\R^N_+$; 给定 $f\in L^2(\R^N_+)$, 问题
\begin{equation*}
\begin{cases}
-\Delta u+u=f & \text{在}\ \R^N_+\ \text{中},\\
\dfrac{\partial u}{\partial x_N}(x',0)=0 & \text{对}\ x'\in\R^{N-1}
\end{cases}
\end{equation*}
在下述意义下有唯一的弱解:
\begin{equation*}
u\in H^1(\Omega)\quad\text{且}\quad
\int_\Omega\nabla u\cdot\nabla v+\int_\Omega uv=\int_\Omega fv
\quad\forall v\in H^1(\Omega).
\end{equation*}
\end{example}


\section{弱解的正则性}
\label{sec:9.6}

\begin{definition}
设 $m\ge1$ 为整数.我们说开集 $\Omega$ 属于 $C^m$ 类, 如果对每个
$x\in\Gamma$ 存在 $x$ 在 $\R^N$ 中的邻域 $U$ 与双射 $H:Q\to U$ 使得
\begin{equation*}
H\in C^m(Q),\quad H^{-1}\in C^m(U),\quad
H(Q_+)=U\cap\Omega,\quad H(Q_0)=U\cap\Gamma.
\end{equation*}
若 $\Omega$ 对一切 $m$ 都属于 $C^m$ 类, 则称 $\Omega$ 属于 $C^\infty$ 类.
\end{definition}

主要的正则性结果如下.

\begin{keytheorem}[Dirichlet 问题的正则性]
\label{thm:9.25}
设 $\Omega$ 是 $C^2$ 类开集, $\Gamma$ 有界 (或者 $\Omega=\R^N_+$).设
$f\in L^2(\Omega)$, $u\in H_0^1(\Omega)$ 满足
\begin{equation}
\int_\Omega\nabla u\nabla\varphi+\int_\Omega u\varphi
=\int_\Omega f\varphi\quad\forall\varphi\in H_0^1(\Omega).
\label{eq:9.48}
\end{equation}
则 $u\in H^2(\Omega)$ 且 $\norm{u}_{H^2}\le C\norm{f}_{L^2}$, 其中 $C$ 是
只依赖于 $\Omega$ 的常数.此外, 若 $\Omega$ 属于 $C^{m+2}$ 类且
$f\in H^m(\Omega)$, 则
\begin{equation*}
u\in H^{m+2}(\Omega)\quad\text{且}\quad
\norm{u}_{H^{m+2}}\le C\norm{f}_{H^m}.
\end{equation*}
特别地, 若 $f\in H^m(\Omega)$ 且 $m>N/2$, 则 $u\in C^2(\Omega)$.最后, 若
$\Omega$ 属于 $C^\infty$ 类且 $f\in C^\infty(\Omega)$, 则
$u\in C^\infty(\Omega)$.
\end{keytheorem}

\begin{theorem}[Neumann 问题的正则性]
\label{thm:9.26}
在与定理  \ref{thm:9.25} 相同的假设下, 对 Neumann 问题的解, 即对满足
\begin{equation}
\int_\Omega\nabla u\cdot\nabla\varphi+\int_\Omega u\varphi
=\int_\Omega f\varphi\quad\forall\varphi\in H^1(\Omega)
\label{eq:9.49}
\end{equation}
的 $u\in H^1(\Omega)$, 可以得到相同的结论.
\end{theorem}

\begin{remark}
\label{rem:9.24}
对与一般二阶椭圆算子相联系的 Dirichlet (或 Neumann) 问题的解, 也可以
得到相同的结论, 即若 $u\in H_0^1(\Omega)$ 满足
\begin{equation*}
\int_\Omega\sum_{i,j}a_{ij}\frac{\partial u}{\partial x_i}
\frac{\partial\varphi}{\partial x_j}+\int_\Omega\sum_ia_i
\frac{\partial u}{\partial x_i}\varphi+\int_\Omega a_0u\varphi
=\int_\Omega f\varphi\quad\forall\varphi\in H_0^1(\Omega),
\end{equation*}
则\footnote{若 $\Omega$ 无界, 我们还必须假设
$D^\alpha a_{ij}\in L^\infty(\Omega)$ $\forall\alpha$, $\abs{\alpha}\le m+1$
以及 $D^\alpha a_i\in L^\infty(\Omega)$ $\forall\alpha$, $\abs{\alpha}\le m$.}
\begin{equation*}
[f\in L^2(\Omega),\ a_{ij}\in C^1(\Omega)\ \text{且}\ a_i\in C(\Omega)]
\Rightarrow u\in H^2(\Omega),
\end{equation*}
并且对 $m\ge1$,
\begin{equation*}
[f\in H^m(\Omega),\ a_{ij}\in C^{m+1}(\Omega)\ \text{且}\ a_i\in C^m(\Omega)]
\Rightarrow u\in H^{m+2}(\Omega).
\end{equation*}
\end{remark}

我们只证明定理  \ref{thm:9.25}; 定理  \ref{thm:9.26} 的证明完全类似.
证明的主要思想如下.我们先考虑 $\Omega=\R^N$ 的情形, 然后考虑
$\Omega=\R^N_+$ 的情形.对一般的区域 $\Omega$ 我们分两步进行:

\noindent 1. \emph{内部正则性}, 即 $u$ 在每个区域 $\omega\subset\subset\Omega$ 上
正则.这里, 证明遵循与 $\Omega=\R^N$ 相同的模式.

\noindent 2. \emph{边界正则性}, 即 $u$ 在边界的某个邻域上正则.这里, 在局部
坐标卡中, 证明类似于 $\Omega=\R^N_+$ 的情形.

我们建议读者在着手处理一般情形之前先好好研究 $\Omega=\R^N$ 与
$\Omega=\R^N_+$ 这两种情形.本节的安排如下:

\noindent A. $\Omega=\R^N$ 的情形.

\noindent B. $\Omega=\R^N_+$ 的情形.

\noindent C. 一般情形:

\noindent C1. 内部估计.

\noindent C2. 边界附近的估计.

证明的本质要素是 L.~Nirenberg 的\emph{平移方法}\footnote{也称为\emph{差商}
技巧.}.

\subsection*{A. $\Omega=\R^N$ 的情形}

\noindent\emph{记号.} 给定 $h\in\R^N$, $h\neq0$, 令
\begin{equation*}
D_hu=\frac1{\abs{h}}(\tau_hu-u),
\quad\text{即}\quad
D_hu(x)=\frac{u(x+h)-u(x)}{\abs{h}}.
\end{equation*}
在 \eqref{eq:9.48} 中取 $\varphi=D_{-h}(D_hu)$.这是允许的, 因为
$\varphi\in H^1(\R^N)$ (由于 $u\in H^1(\R^N)$); 我们得到
\begin{equation*}
\int\abs{\nabla D_hu}^2+\int\abs{D_hu}^2=\int fD_{-h}(D_hu),
\end{equation*}
从而
\begin{equation}
\norm{D_hu}_{H^1}^2\le\norm{f}_2\norm{D_{-h}(D_hu)}_2.
\label{eq:9.50}
\end{equation}
另一方面, 回忆 (命题  \ref{prop:9.3})
\begin{equation}
\norm{D_{-h}v}_2\le\norm{\nabla v}_2\quad\forall v\in H^1.
\label{eq:9.51}
\end{equation}
把它用于 $v=D_hu$, 得到
\begin{equation*}
\norm{D_hu}_{H^1}^2\le\norm{f}_2\norm{\nabla(D_hu)}_2,
\end{equation*}
从而
\begin{equation*}
\norm{D_hu}_{H^1}\le\norm{f}_2.
\end{equation*}
特别地,
\begin{equation}
\left\|D_h\frac{\partial u}{\partial x_i}\right\|_2\le\norm{f}_2
\quad\forall i=1,2,\dots,N.
\label{eq:9.52}
\end{equation}
再次应用命题  \ref{prop:9.3}, 我们看到
$\frac{\partial u}{\partial x_i}\in H^1$, 从而 $u\in H^2$.

我们现在证明 $f\in H^1\Rightarrow u\in H^3$.记 $Du$ 为导数
$\frac{\partial u}{\partial x_i}$ ($1\le i\le N$) 中的任何一个.我们已经知道
$Du\in H^1$.还须证明 $Du\in H^2$.为此只须验证
\begin{equation}
\int\nabla(Du)\cdot\nabla\varphi+\int(Du)\varphi=\int(Df)\varphi
\quad\forall\varphi\in H^1
\label{eq:9.53}
\end{equation}
(然后我们就可以对 $Du$ 应用前面的分析, 得到 $Du\in H^2$).若
$\varphi\in C_c^\infty(\R^N)$, 我们可以在 \eqref{eq:9.48} 中把 $\varphi$ 换成
$D\varphi$; 它成为
\begin{equation*}
\int\nabla u\cdot\nabla(D\varphi)+\int uD\varphi=\int fD\varphi,
\end{equation*}
从而
\begin{equation*}
\int\nabla(Du)\cdot\nabla\varphi+\int(Du)\varphi=\int(Df)\varphi
\quad\forall\varphi\in C_c^\infty(\R^N).
\end{equation*}
这蕴含 \eqref{eq:9.53}, 因为 $C_c^\infty(\R^N)$ 在 $H^1(\R^N)$ 中稠密
(命题 9.2).
% 原书此处印作 "(Proposition 9.2)" (印刷页 301); 本章的命题编号为 9.1, 9.3,
% 9.4, ... 并无编号 9.2 的命题 (9.2 是定理 9.2, Friedrichs 的逼近定理, 其
% Ω=R^N 的情形恰好给出这一稠密性). 依技能规范保留原样并记录, 不擅自更正,
% 故此处不用 \ref 以免产生未定义引用.

为证明 $f\in H^m\Rightarrow u\in H^{m+2}$, 只须对 $m$ 作归纳并应用
\eqref{eq:9.53}.

\subsection*{B. $\Omega=\R^N_+$ 的情形}

我们再次使用平移, 但只在切向进行, 即沿方向
$h\in\R^{N-1}\times\{0\}$: 我们说 $h$ \emph{平行于边界}, 并记作
$h\parallel\Gamma$.关键的一点是
\begin{equation*}
u\in H_0^1(\Omega)\ \Rightarrow\ \tau_hu\in H_0^1(\Omega)
\quad\text{若}\ h\parallel\Gamma.
\end{equation*}
换言之, $H_0^1(\Omega)$ 在切向平移下不变.

我们取 $h\parallel\Gamma$, 并在 \eqref{eq:9.48} 中代入
$\varphi=D_{-h}(D_hu)$; 得到
\begin{equation*}
\int\abs{\nabla(D_hu)}^2+\int\abs{D_hu}^2=\int fD_{-h}(D_hu),
\end{equation*}
即
\begin{equation}
\norm{D_hu}_{H^1}^2\le\norm{f}_2\norm{D_{-h}(D_hu)}_2.
\label{eq:9.54}
\end{equation}
现在我们使用下述引理.

\begin{lemma}
\label{lem:9.6}
我们有
\begin{equation*}
\norm{D_hv}_{L^2(\Omega)}\le\norm{\nabla v}_{L^2(\Omega)}
\quad\forall v\in H^1(\Omega),\ \forall h\parallel\Gamma.
\end{equation*}
\end{lemma}

\begin{proof}
先取 $v\in C_c^1(\R^N)$, 仿照命题  \ref{prop:9.3} 的证明 (注意
$\Omega+th=\Omega$ 对一切 $t$ 与一切 $h\parallel\Gamma$ 成立).对一般的
$v\in H^1(\Omega)$ 用稠密性论证.
\end{proof}

结合 \eqref{eq:9.54} 与引理  \ref{lem:9.6}, 我们得到
\begin{equation}
\norm{D_hu}_{H^1}\le\norm{f}_2\quad\forall h\parallel\Gamma.
\label{eq:9.55}
\end{equation}
设 $1\le j\le N$, $1\le k\le N-1$, $h=\abs{h}e_k$, 并设
$\varphi\in C_c^\infty(\Omega)$.我们有
\begin{equation*}
\int D_h\left(\frac{\partial u}{\partial x_j}\right)\varphi
=-\int uD_{-h}\left(\frac{\partial\varphi}{\partial x_j}\right),
\end{equation*}
而由 \eqref{eq:9.55},
\begin{equation*}
\left|\int uD_{-h}\left(\frac{\partial\varphi}{\partial x_j}\right)\right|
\le\norm{f}_2\norm{\varphi}_2.
\end{equation*}
令 $h\to0$ 取极限, 它成为
\begin{equation}
\left|\int u\frac{\partial^2\varphi}{\partial x_j\partial x_k}\right|
\le\norm{f}_2\norm{\varphi}_2
\quad\forall1\le j\le N,\ \forall1\le k\le N-1.
\label{eq:9.56}
\end{equation}
最后, 我们断言
\begin{equation}
\left|\int u\frac{\partial^2\varphi}{\partial x_N^2}\right|
\le\norm{f}_2\norm{\varphi}_2\quad\forall\varphi\in C_c^\infty(\Omega).
\label{eq:9.57}
\end{equation}
为证明 \eqref{eq:9.57}, 我们回到方程 \eqref{eq:9.48}, 推出
\begin{equation*}
\left|\int u\frac{\partial^2\varphi}{\partial x_N^2}\right|
\le\sum_{i=1}^{N-1}\left|\int u\frac{\partial^2\varphi}{\partial x_i^2}\right|
+\left|\int(f-u)\varphi\right|\le C\norm{f}_2\norm{\varphi}_2
\end{equation*}
(这里用了 \eqref{eq:9.56}).结合 \eqref{eq:9.56} 与 \eqref{eq:9.57}, 我们得到
\begin{equation*}
\left|\int u\frac{\partial^2\varphi}{\partial x_j\partial x_k}\right|
\le C\norm{f}_2\norm{\varphi}_2
\quad\forall\varphi\in C_c^\infty(\Omega),\ \forall1\le j,k\le N.
\end{equation*}
于是 $u\in H^2(\Omega)$, 因为存在函数 $f_{jk}\in L^2(\Omega)$ 使得
\begin{equation*}
\int u\frac{\partial^2\varphi}{\partial x_j\partial x_k}
=\int f_{jk}\varphi\quad\forall\varphi\in C_c^\infty(\Omega)
\end{equation*}
(如同命题  \ref{prop:8.3} 的证明).

最后我们证明 $f\in H^m(\Omega)\Rightarrow u\in H^{m+2}(\Omega)$.用 $Du$
表示任一切向导数 $Du=\frac{\partial u}{\partial x_j}$, $1\le j\le N-1$.
我们先建立下述结果.

\begin{lemma}
\label{lem:9.7}
设 $u\in H^2(\Omega)\cap H_0^1(\Omega)$ 满足 \eqref{eq:9.48}.则
$Du\in H_0^1(\Omega)$, 并且
\begin{equation}
\int\nabla(Du)\cdot\nabla\varphi+\int(Du)\varphi=\int(Df)\varphi
\quad\forall\varphi\in H_0^1(\Omega).
\label{eq:9.58}
\end{equation}
\end{lemma}

\begin{proof}
唯一微妙之处在于证明 $Du\in H_0^1(\Omega)$; 至于 \eqref{eq:9.58}, 它可由
\eqref{eq:9.48} 取 $D\varphi$ 代替 $\varphi$ (其中
$\varphi\in C_c^\infty(\Omega)$) 然后利用稠密性得到.设 $h=\abs{h}e_j$,
$1\le j\le N-1$, 于是 $D_hu\in H_0^1$ (因为 $H_0^1$ 在切向平移下不变).
由引理  \ref{lem:9.6} 我们有
\begin{equation*}
\norm{D_hu}_{H^1}\le\norm{u}_{H^2}.
\end{equation*}
于是存在序列 $h_n\to0$ 使得 $D_{h_n}u$ 弱收敛于 $H_0^1$ 中某个 $g$
(因为 $H_0^1$ 是 Hilbert 空间).特别地, $D_{h_n}u\rightharpoonup g$ 在
$L^2$ 中弱收敛.对 $\varphi\in C_c^\infty(\Omega)$ 我们有
\begin{equation*}
\int(D_hu)\varphi=\int u(D_{-h}\varphi),
\end{equation*}
令 $h_n\to0$ 取极限, 得到
\begin{equation*}
\int g\varphi=-\int u\frac{\partial\varphi}{\partial x_j}
\quad\forall\varphi\in C_c^\infty(\Omega).
\end{equation*}
因此 $\frac{\partial u}{\partial x_j}=g\in H_0^1(\Omega)$.
\end{proof}

\runin{证明 $f\in H^m\Rightarrow u\in H^{m+2}$.} 对 $m$ 作归纳.假设结论
到 $m$ 阶成立, 并设 $f\in H^{m+1}$.我们已经知道 $u\in H^{m+2}$; 又 $Du$
(任一切向导数) 属于 $H_0^1(\Omega)$ 且满足 \eqref{eq:9.58}.把归纳假设用于
$Du$ 与 $Df$, 我们看到 $Du\in H^{m+2}$.为完成证明, 例如只须验证
$\frac{\partial^2u}{\partial x_N^2}\in H^{m+1}$.为此我们再一次回到方程
\eqref{eq:9.48}, 把它写成
\begin{equation*}
\frac{\partial^2u}{\partial x_N^2}
=-\sum_{i=1}^{N-1}\frac{\partial^2u}{\partial x_i^2}+u-f\in H^{m+1}.
\end{equation*}

\subsection*{C. 一般情形}

我们只证明 $f\in L^2(\Omega)\Rightarrow u\in H^2(\Omega)$; 蕴含关系
$f\in H^m\Rightarrow u\in H^{m+2}$ 如同情形 A 与 B 那样对 $m$ 作归纳.
为简化叙述, 我们假设 $\Omega$ 有界.我们使用单位分解, 并像定理  \ref{thm:9.7}
的证明那样写 $u=\sum_{i=0}^k\theta_iu$.

\runin{C1. 内部估计.} 我们断言 $\theta_0u\in H^2(\Omega)$.由于
$\theta_0|_\Omega\in C_c^\infty(\Omega)$, 把 $\theta_0u$ 在 $\Omega$ 外延拓为
$0$ 所得的函数属于 $H^1(\R^N)$ (见注记 \ref{rem:9.4}(ii)).容易验证
$\theta_0u$ 是 $\R^N$ 中下述方程的弱解:
\begin{equation*}
-\Delta(\theta_0u)+\theta_0u=\theta_0f-2\nabla\theta_0\cdot\nabla u
-(\Delta\theta_0)u\defeq g,
\end{equation*}
其中 $g\in L^2(\R^N)$.由情形 A 推出 $\theta_0u\in H^2(\R^N)$ 且
\begin{equation*}
\norm{\theta_0u}_{H^2}\le C(\norm{f}_2+\norm{u}_{H^1})
\le C'\norm{f}_2
\end{equation*}
(因为由 \eqref{eq:9.48} 有 $\norm{u}_{H^1}\le\norm{f}_2$).

\runin{C2. 边界附近的估计.} 我们断言 $\theta_iu\in H^2(\Omega)$,
$1\le i\le k$.回忆我们有双射 $H:Q\to U_i$ 使得
\begin{equation*}
H\in C^2(Q),\quad J=H^{-1}\in C^2(U_i),\quad
H(Q_+)=\Omega\cap U_i,\quad H(Q_0)=\Gamma\cap U_i.
\end{equation*}
我们写 $x=H(y)$, $y=H^{-1}(x)=J(x)$.容易验证
$v=\theta_iu\in H_0^1(\Omega\cap U_i)$, 且 $v$ 是 $\Omega\cap U_i$ 中下述
方程的弱解:
\begin{equation*}
-\Delta v=\theta_if-\theta_iu-2\nabla\theta_i\cdot\nabla u
-(\Delta\theta_i)u\defeq g,
\end{equation*}
其中 $g\in L^2(\Omega\cap U_i)$ 且 $\norm{g}_2\le C\norm{f}_2$.更精确地说,
\begin{equation}
\int_{\Omega\cap U_i}\nabla v\cdot\nabla\varphi\,dx
=\int_{\Omega\cap U_i}g\varphi\,dx
\quad\forall\varphi\in H_0^1(\Omega\cap U_i).
\label{eq:9.59}
\end{equation}
现在我们把 $v|_{\Omega\cap U_i}$ 转移到 $Q_+$ 上.令
\begin{equation*}
w(y)=v(H(y))\quad\text{对}\ y\in Q_+,
\quad\text{即}\quad
w(Jx)=v(x)\quad\text{对}\ x\in\Omega\cap U_i.
\end{equation*}
下述基本引理表明方程 \eqref{eq:9.59} 变成 $w$ 在 $Q_+$ 上的一个二阶椭圆
方程.\footnote{更一般地, 如果我们从 $v$ 的椭圆方程出发, 就会得到 $w$ 的椭圆
方程: 椭圆性条件在变量替换下保持不变.}

\begin{lemma}
\label{lem:9.8}
在上述记号下, $w$ 属于 $H_0^1(Q_+)$ 且满足
\begin{equation}
\sum_{k,\ell=1}^N\int_{Q_+}a_{k\ell}\frac{\partial w}{\partial y_k}
\frac{\partial\psi}{\partial y_\ell}\,dy
=\int_{Q_+}\tilde g\psi\,dy\quad\forall\psi\in H_0^1(Q_+),
\label{eq:9.60}
\end{equation}
其中\footnote{$\operatorname{detJac}H$ 表示 Jacobi 行列式, 即 Jacobi 矩阵
$\operatorname{Jac}H=\left(\frac{\partial H_i}{\partial y_j}\right)$ 的行列式.}
$\tilde g=(g\circ H)\abs{\operatorname{detJac}H}\in L^2(Q_+)$, 且函数
$a_{k\ell}\in C^1(Q_+)$ 满足椭圆性条件 \eqref{eq:9.36}.
\end{lemma}

\begin{proof}
设 $\psi\in H_0^1(Q_+)$, 并对 $x\in\Omega\cap U_i$ 令 $\varphi(x)=\psi(Jx)$.
则 $\varphi\in H_0^1(\Omega\cap U_i)$ 且
\begin{equation*}
\frac{\partial v}{\partial x_j}=\sum_k\frac{\partial w}{\partial y_k}
\frac{\partial J_k}{\partial x_j},\qquad
\frac{\partial\varphi}{\partial x_j}=\sum_\ell\frac{\partial\psi}{\partial y_\ell}
\frac{\partial J_\ell}{\partial x_j}.
\end{equation*}
于是由积分中通常的变量替换公式,
\begin{equation*}
\begin{aligned}
\int_{\Omega\cap U_i}\nabla v\cdot\nabla\varphi\,dx
&=\int_{\Omega\cap U_i}\sum_{j,k,\ell}\frac{\partial J_k}{\partial x_j}
\frac{\partial J_\ell}{\partial x_j}\frac{\partial w}{\partial y_k}
\frac{\partial\psi}{\partial y_\ell}\,dx\\
&=\int_{Q_+}\sum_{j,k,\ell}\frac{\partial J_k}{\partial x_j}
\frac{\partial J_\ell}{\partial x_j}\frac{\partial w}{\partial y_k}
\frac{\partial\psi}{\partial y_\ell}\abs{\operatorname{detJac}H}\,dy.
\end{aligned}
\end{equation*}
于是
\begin{equation}
\int_{\Omega\cap U_i}\nabla v\cdot\nabla\varphi\,dx
=\int_{Q_+}\sum_{k,\ell}a_{k\ell}\frac{\partial w}{\partial y_k}
\frac{\partial\psi}{\partial y_\ell}\,dy,
\label{eq:9.61}
\end{equation}
其中
\begin{equation*}
a_{k\ell}=\sum_j\frac{\partial J_k}{\partial x_j}
\frac{\partial J_\ell}{\partial x_j}\abs{\operatorname{detJac}H}.
\end{equation*}
我们注意到 $a_{k\ell}\in C^1(Q_+)$, 且椭圆性条件成立, 因为对一切
$\xi\in\R^N$ 有
\begin{equation*}
\sum_{k,\ell}a_{k\ell}\xi_k\xi_\ell
=\abs{\operatorname{detJac}H}\sum_j
\left|\sum_k\frac{\partial J_k}{\partial x_j}\xi_k\right|^2
\ge\alpha\abs{\xi}^2,
\end{equation*}
其中 $\alpha>0$, 因为 Jacobi 矩阵 $\operatorname{Jac}H$ 与
$\operatorname{Jac}J$ 都非奇异.

另一方面, 我们有
\begin{equation}
\int_{\Omega\cap U_i}g\varphi\,dx
=\int_{Q_+}(g\circ H)\psi\abs{\operatorname{detJac}H}\,dy.
\label{eq:9.62}
\end{equation}
结合 \eqref{eq:9.59}, \eqref{eq:9.61} 与 \eqref{eq:9.62} 即得
\eqref{eq:9.60}.引理  \ref{lem:9.8} 证毕.
\end{proof}

我们现在回到边界估计的证明, 并证明
$w\in H^2(Q_+)$\footnote{下文中我们用 $C$ 表示只依赖于 $a_{k\ell}$ 的各种
常数.} 且 $\norm{w}_{H^2}\le C\norm{\tilde g}_2$.回到
$\Omega\cap U_i$, 这将推出 $\theta_iu$ 属于 $H^2(\Omega\cap U_i)$, 从而
实际上属于 $H^2(\Omega)$ 且 $\norm{\theta_iu}_{H^2}\le C\norm{f}_2$.

如同情形 B ($\Omega=\R^N_+$), 我们使用切向平移.在 \eqref{eq:9.60} 中取
$\psi=D_{-h}(D_hw)$, 其中 $h\parallel Q_0$, 且 $\abs{h}$ 足够小使得
$\psi\in H_0^1(Q_+)$.\footnote{回忆
$\operatorname{supp}w\subset\{(x',x_N);\ \abs{x'}<1-\delta\ \text{且}\
0\le x_N<1-\delta\}$ 对某个 $\delta>0$ 成立.} 我们得到
\begin{equation}
\sum_{k,\ell}\int_{Q_+}D_h\left(a_{k\ell}\frac{\partial w}{\partial y_k}\right)
\frac{\partial}{\partial y_\ell}(D_hw)
=\int_{Q_+}\tilde gD_{-h}(D_hw).
\label{eq:9.63}
\end{equation}
但
\begin{equation}
\int_{Q_+}\tilde gD_{-h}(D_hw)\le\norm{\tilde g}_2
\norm{D_{-h}(D_hw)}_2\le\norm{\tilde g}_2\norm{\nabla D_hw}_2
\label{eq:9.64}
\end{equation}
(由引理  \ref{lem:9.6}).另一方面, 写
\begin{equation*}
D_h\left(a_{k\ell}\frac{\partial w}{\partial y_k}\right)(y)
=a_{k\ell}(y+h)\frac{\partial}{\partial y_k}D_hw(y)
+(D_ha_{k\ell}(y))\frac{\partial w}{\partial y_k}(y),
\end{equation*}
于是
\begin{equation}
\sum_{k,\ell}\int_{Q_+}D_h\left(a_{k\ell}\frac{\partial w}{\partial y_k}\right)
\frac{\partial}{\partial y_\ell}(D_hw)
\ge\alpha\norm{\nabla(D_hw)}_2^2-C\norm{w}_{H^1}\norm{\nabla D_hw}_2.
\label{eq:9.65}
\end{equation}
结合 \eqref{eq:9.64} 与 \eqref{eq:9.65}, 我们得到
\begin{equation}
\norm{\nabla D_hw}_2\le C(\norm{w}_{H^1}+\norm{\tilde g}_2)
\le C\norm{\tilde g}_2
\label{eq:9.66}
\end{equation}
(注意由 \eqref{eq:9.60} 与 Poincaré 不等式有
$\norm{w}_{H^1}\le C\norm{\tilde g}_2$).由 \eqref{eq:9.66} 推出——如同
情形 B 那样——
\begin{equation}
\left|\int_{Q_+}\frac{\partial w}{\partial y_k}
\frac{\partial\psi}{\partial y_\ell}\right|
\le C\norm{\tilde g}_2\norm{\psi}_2
\quad\forall\psi\in C_c^1(Q_+),\ \forall(k,\ell)\neq(N,N).
\label{eq:9.67}
\end{equation}
为推出 $w\in H^2(Q_+)$ (且 $\norm{w}_{H^2}\le C\norm{\tilde g}_2$),
还须证明
\begin{equation}
\left|\int_{Q_+}\frac{\partial w}{\partial y_N}
\frac{\partial\psi}{\partial y_N}\right|
\le C\norm{\tilde g}_2\norm{\psi}_2\quad\forall\psi\in C_c^1(Q_+).
\label{eq:9.68}
\end{equation}
为此我们回到方程, 在其中把 $\psi$ 换成 $(1/a_{NN})\psi$
($\psi\in C_c^1(Q_+)$); 这是允许的, 因为 $a_{NN}\in C^1(Q_+)$ 且
$a_{NN}\ge\alpha>0$.方程成为
\begin{equation*}
\int a_{NN}\frac{\partial w}{\partial y_N}
\frac{\partial}{\partial y_N}\left(\frac{\psi}{a_{NN}}\right)
=\int\frac{\tilde g}{a_{NN}}\psi
-\sum_{(k,\ell)\neq(N,N)}\int a_{k\ell}\frac{\partial w}{\partial y_k}
\frac{\partial}{\partial y_\ell}\left(\frac{\psi}{a_{NN}}\right),
\end{equation*}
即
\begin{equation}
\begin{aligned}
\int\frac{\partial w}{\partial y_N}\frac{\partial\psi}{\partial y_N}
={}&\int\frac1{a_{NN}}\frac{\partial a_{NN}}{\partial y_N}
\frac{\partial w}{\partial y_N}\psi+\int\frac{\tilde g}{a_{NN}}\psi\\
&+\sum_{(k,\ell)\neq(N,N)}\int\frac{\partial w}{\partial y_k}
\frac{\partial a_{k\ell}}{\partial y_\ell}\frac{\psi}{a_{NN}}
-\sum_{(k,\ell)\neq(N,N)}\int\frac{\partial w}{\partial y_k}
\frac{\partial}{\partial y_\ell}\left(\frac{a_{k\ell}\psi}{a_{NN}}\right).
\end{aligned}
\label{eq:9.69}
\end{equation}
结合 \eqref{eq:9.67}\footnote{我们对 $(a_{k\ell}a_{NN})\psi$ 代替 $\psi$
使用 \eqref{eq:9.67}.} 与 \eqref{eq:9.69}, 我们得到
\begin{equation*}
\left|\int_{Q_+}\frac{\partial w}{\partial y_N}
\frac{\partial\psi}{\partial y_N}\right|
\le C(\norm{w}_{H^1}+\norm{\tilde g}_2)\norm{\psi}_2
\quad\forall\psi\in C_c^1(Q_+).
\end{equation*}
这就建立了 \eqref{eq:9.68}, 并完成了边界附近的估计.

\begin{remark}
\label{rem:9.25}
设 $\Omega$ 为任意开集, $u\in H^1(\Omega)$ 使得
\begin{equation*}
\int_\Omega\nabla u\cdot\nabla\varphi=\int_\Omega f\varphi
\quad\forall\varphi\in C_c^\infty(\Omega).
\end{equation*}
我们假设 $f\in H^m(\Omega)$.则对每个 $\theta\in C_c^\infty(\Omega)$ 有
$\theta u\in H^{m+2}(\Omega)$: 我们说 $u\in H^{m+2}_{\mathrm{loc}}(\Omega)$.
为证明这一点, 只须像情形 C1 那样进行并对 $m$ 作归纳.特别地,
$f\in C^\infty(\Omega)\Rightarrow u\in C^\infty(\Omega)$.\footnote{但一般
说来, 我们不能断言例如 $u\in C(\Omega)$ (即使 $\Omega$ 与 $f$ 都非常光滑),
因为没有给定任何边界条件.}

对\emph{很弱}意义下的解也有同样的结论, 即 $u\in L^2(\Omega)$ 使得
\begin{equation*}
-\int_\Omega u\Delta\varphi=\int_\Omega f\varphi
\quad\forall\varphi\in C_c^\infty(\Omega).
\end{equation*}
(证明要稍微微妙一些; 参见例如 S.~Agmon \bref{Agmon}{1}.) 我们要强调椭圆
问题中正则性结果的\emph{局部}性质.更精确地说, 设 $f\in L^2(\Omega)$,
$u\in H_0^1(\Omega)$ 是
\begin{equation*}
\int_\Omega\nabla u\cdot\nabla\varphi+\int_\Omega u\varphi
=\int_\Omega f\varphi\quad\forall\varphi\in H_0^1(\Omega)
\end{equation*}
的唯一弱解.固定 $\omega\subset\subset\Omega$; 则 $u|_\omega$ 依赖于 $f$
在\emph{整个} $\Omega$ 上的值——而不仅是 $f$ 在 $\omega$ 上的值.\footnote{例如,
若 $f\ge0$ 在 $\Omega$ 上, $f=0$ 在 $\omega$ 上, 而 $f>0$ 在 $\Omega$ 的某个
开子集上, 则 $u>0$ 在 $\Omega$ 上 (从而在 $\omega$ 上); 见本章末评注中的强
极大值原理.} 相反, $u|_\omega$ 的\emph{正则性}只依赖于 $f|_\omega$ 的正则性.
例如, 即使 $f$ 在 $\omega$ 外非常不规则, 也有
$f\in C^\infty(\omega)\Rightarrow u\in C^\infty(\omega)$.这个性质称为
\emph{亚椭圆性}.
\end{remark}

\begin{remark}
\label{rem:9.26}
从某种观点看, 正则性结果有点令人惊讶.事实上, 对 $\Delta u$ 即对导数之和
$\sum_k\frac{\partial^2u}{\partial x_k^2}$ 所作的假设, 迫使我们对\emph{每个}
导数 $\frac{\partial^2u}{\partial x_i\partial x_j}$ 分别得到同一性质的结论.
\end{remark}

\section{极大值原理}
\label{sec:9.7}

极大值原理是一个非常有用的工具, 它有许多种表述形式.这里我们介绍几种
简单的形式.

设 $\Omega$ 为 $\R^N$ 的一般开子集.

\begin{keytheorem}[Dirichlet 问题的极大值原理]
\index{Nweijidazhiyuanli@N 维极大值原理 (maximum principle in dimension N)}
\label{thm:9.27}
假设\footnote{若 $\Omega$ 属于 $C^1$ 类, 则可以借助迹理论 (它赋予
$u|_\Gamma$ 以意义; 见本章末的评注) 去掉假设 $u\in C(\overline\Omega)$;
此外若 $u\in H_0^1(\Omega)$, 也可以去掉假设 $u\in C(\overline\Omega)$.}
$f\in L^2(\Omega)$ 且 $u\in H^1(\Omega)\cap C(\overline\Omega)$ 满足
\begin{equation}
\int_\Omega\nabla u\cdot\nabla\varphi+\int_\Omega u\varphi
=\int_\Omega f\varphi\quad\forall\varphi\in H_0^1(\Omega).
\label{eq:9.70}
\end{equation}
则对一切 $x\in\Omega$,
\begin{equation*}
\min\left\{\inf_\Gamma u,\ \inf_\Omega f\right\}\le u(x)
\le\max\left\{\sup_\Gamma u,\ \sup_\Omega f\right\}.
\end{equation*}
(这里以及下文中, $\sup$ 与 $\inf$ 分别表示本性上确界与本性下确界.)
\end{keytheorem}

\begin{proof}
我们使用 Stampacchia 的截断方法.固定函数 $G\in C^1(\R)$ 使得
\begin{equation*}
\begin{array}{ll}
\text{(i)}\ \abs{G'(s)}\le M & \forall s\in\R,\\[1mm]
\text{(ii)}\ G\ \text{在}\ (0,+\infty)\ \text{上严格递增}, &\\[1mm]
\text{(iii)}\ G(s)=0 & \forall s\le0.
\end{array}
\end{equation*}
令
\begin{equation*}
K=\max\left\{\sup_\Gamma u,\ \sup_\Omega f\right\},
\end{equation*}
并假设 $K<\infty$ (否则无须证明).令 $v=G(u-K)$.我们区分两种情形:

\noindent (a) $\abs{\Omega}<\infty$ 的情形.

则 $v\in H^1(\Omega)$ (由命题  \ref{prop:9.5} 用于函数
$t\mapsto G(t-K)-G(-K)$).另一方面, $v\in H_0^1(\Omega)$, 因为
$v\in C(\overline\Omega)$ 且 $v=0$ 在 $\Gamma$ 上 (见定理  \ref{thm:9.17}).
把这个 $v$ 代入 \eqref{eq:9.70}, 并像定理  \ref{thm:8.19} 的证明那样进行.
% 原书此处印作 "as in the proof of Theorem 8.18" (印刷页 308); 第~
% \ref{chap:8}~章并无编号
% 8.18 的定理 (8.18 是命题, 讲二阶边值问题的存在唯一性), 用 Stampacchia 截断
% 方法证明极大值原理的是定理~
% \ref{thm:8.19}. 依技能规范保留原样并记录, 不擅自更正.

\noindent (b) $\abs{\Omega}=\infty$ 的情形.

此时 $K\ge0$ (因为 $f(x)\le K$ 在 $\Omega$ 上几乎处处成立且 $f\in L^2$
蕴含 $K\ge0$).固定 $K'>K$.由命题  \ref{prop:9.5} 用于函数
$t\mapsto G(t-K')$, 我们看到 $v=G(u-K')\in H^1(\Omega)$.此外
$v\in C(\overline\Omega)$ 且 $v=0$ 在 $\Gamma$ 上; 因此 $v\in H_0^1(\Omega)$.
把这个 $v$ 代入 \eqref{eq:9.70}, 我们有
\begin{equation}
\int_\Omega\abs{\nabla u}^2G'(u-K')+\int_\Omega uG(u-K')
=\int_\Omega fG(u-K').
\label{eq:9.71}
\end{equation}
另一方面, $G(u-K')\in L^1(\Omega)$, 因为\footnote{因为
$G(u-K')-G(-K')\le M\abs{u}$ 且 $G(-K')=0$ (由于 $-K'<0$).}
\begin{equation*}
0\le G(u-K')\le M\abs{u},
\end{equation*}
而在集合 $[u\ge K']=\{x\in\Omega;\ u(x)\ge K'\}$ 上有
\begin{equation*}
K'\int_{[u\ge K']}\abs{u}\le\int_\Omega u^2<\infty.
\end{equation*}
由 \eqref{eq:9.71} 我们推出
\begin{equation*}
\int_\Omega(u-K')G(u-K')\le\int_\Omega(f-K')G(u-K')\le0.
\end{equation*}
于是 $u\le K'$ 在 $\Omega$ 上几乎处处成立, 从而 $u\le K$ 在 $\Omega$ 上几乎
处处成立 (因为 $K'>K$ 是任意的).
\end{proof}

\begin{keycorollary}
\label{cor:9.28}
设 $f\in L^2(\Omega)$, $u\in H^1(\Omega)\cap C(\overline\Omega)$\footnote{同
上, 在某些情形下可以去掉假设 $u\in C(\overline\Omega)$.} 满足
\eqref{eq:9.70}.则
\begin{equation}
[u\ge0\ \text{在}\ \Gamma\ \text{上且}\ f\ge0\ \text{在}\ \Omega\ \text{中}]
\Rightarrow[u\ge0\ \text{在}\ \Omega\ \text{中}],
\label{eq:9.72}
\end{equation}
\begin{equation}
\norm{u}_{L^\infty(\Omega)}\le\max\left\{
\norm{u}_{L^\infty(\Gamma)},\ \norm{f}_{L^\infty(\Omega)}\right\}.
\label{eq:9.73}
\end{equation}
特别地,
\begin{equation}
\text{若}\ f=0\ \text{在}\ \Omega\ \text{中, 则}\
\norm{u}_{L^\infty(\Omega)}\le\norm{u}_{L^\infty(\Gamma)},
\label{eq:9.74}
\end{equation}
\begin{equation}
\text{若}\ u=0\ \text{在}\ \Gamma\ \text{上, 则}\
\norm{u}_{L^\infty(\Omega)}\le\norm{f}_{L^\infty(\Omega)}.
\label{eq:9.75}
\end{equation}
\end{keycorollary}

\begin{remark}
\label{rem:9.27}
若 $\Omega$ 有界且 $u$ 是方程
\begin{equation}
-\Delta u+u=f\quad\text{在}\ \Omega\ \text{中}
\label{eq:9.76}
\end{equation}
的古典解, 则可以对定理  \ref{thm:9.27} 给出另一种证明.事实上, 设
$x_0\in\overline\Omega$ 使得 $u(x_0)=\max_{\overline\Omega}u$.

\noindent (i) 若 $x_0\in\Gamma$, 则 $u(x_0)\le\sup_\Gamma u\le K$.

\noindent (ii) 若 $x_0\in\Omega$, 则 $\nabla u(x_0)=0$ 且
$\frac{\partial^2u}{\partial x_i^2}(x_0)\le0$ 对一切 $1\le i\le N$ 成立,
从而 $\Delta u(x_0)\le0$.由此, 利用方程 \eqref{eq:9.76} 我们得到
\begin{equation*}
u(x_0)=f(x_0)+\Delta u(x_0)\le f(x_0)\le K.
\end{equation*}

这个方法的好处是它适用于一般的二阶椭圆方程.例如, 定理  \ref{thm:9.27} 的
结论对
\begin{equation}
-\sum_{i,j=1}^N\frac{\partial}{\partial x_j}
\left(a_{ij}\frac{\partial u}{\partial x_i}\right)
+\sum_{i=1}^Na_i\frac{\partial u}{\partial x_i}+u=f
\quad\text{在}\ \Omega\ \text{中}
\label{eq:9.77}
\end{equation}
也成立.注意若 $x_0\in\Omega$, 则
\begin{equation*}
\sum_{i,j=1}^Na_{ij}(x_0)\frac{\partial^2u}{\partial x_i\partial x_j}(x_0)\le0;
\end{equation*}
事实上, 通过 (依赖于 $x_0$ 的) 坐标变换可以把问题化为矩阵 $a_{ij}(x_0)$
为对角形的情形.定理  \ref{thm:9.27} 的结论对 \eqref{eq:9.77} 的弱解仍然
成立, 但证明要微妙一些; 参见 D.~Gilbarg--N.~Trudinger
\bref{GilbargTrudinger}{1}.
\end{remark}

\begin{proposition}
\label{prop:9.29}
设函数 $a_{ij}\in L^\infty(\Omega)$ 满足椭圆性条件 \eqref{eq:9.36}, 且
$a_i,a_0\in L^\infty(\Omega)$, $a_0\ge0$ 在 $\Omega$ 上.设
$f\in L^2(\Omega)$, $u\in H^1(\Omega)\cap C(\overline\Omega)$\footnote{同
上, 在某些情形下可以去掉假设 $u\in C(\overline\Omega)$.} 使得
\begin{equation}
\int_\Omega\sum_{i,j}a_{ij}\frac{\partial u}{\partial x_i}
\frac{\partial\varphi}{\partial x_j}
+\int_\Omega\sum_ia_i\frac{\partial u}{\partial x_i}\varphi
+\int_\Omega a_0u\varphi=\int_\Omega f\varphi
\quad\forall\varphi\in H_0^1(\Omega).
\label{eq:9.78}
\end{equation}
则
\begin{equation}
[u\ge0\ \text{在}\ \Gamma\ \text{上且}\ f\ge0\ \text{在}\ \Omega\ \text{中}]
\Rightarrow[u\ge0\ \text{在}\ \Omega\ \text{中}].
\label{eq:9.79}
\end{equation}
设另外 $a_0\equiv0$ 且 $\Omega$ 有界.则
\begin{equation}
[f\ge0\ \text{在}\ \Omega\ \text{中}]
\Rightarrow\left[u\ge\inf_\Gamma u\ \text{在}\ \Omega\ \text{中}\right],
\label{eq:9.80}
\end{equation}
且
\begin{equation}
[f=0\ \text{在}\ \Omega\ \text{中}]
\Rightarrow\left[\inf_\Gamma u\le\sup_\Gamma u\ \text{在}\ \Omega\
\text{中}\right].
\label{eq:9.81}
\end{equation}
\end{proposition}

\begin{proof}
我们在 $a_i\equiv0$ ($1\le i\le N$) 的情形下证明这个结果; 一般情形更为
微妙 (参见 D.~Gilbarg--N.~Trudinger \bref{GilbargTrudinger}{1}, 定理 8.1).
证明 \eqref{eq:9.79} 相当于证明
\begin{equation*}
[u\le0\ \text{在}\ \Gamma\ \text{上且}\ f\le0\ \text{在}\ \Omega\ \text{中}]
\Rightarrow[u\le0\ \text{在}\ \Omega\ \text{中}].
\tag*{(79$'$)}
\end{equation*}
在 \eqref{eq:9.78} 中取 $\varphi=G(u)$, 其中 $G$ 如定理  \ref{thm:9.27} 的
证明中那样; 我们于是得到
\begin{equation*}
\int_\Omega\sum_{i,j}a_{ij}\frac{\partial u}{\partial x_i}
\frac{\partial u}{\partial x_j}G'(u)\le0,
\end{equation*}
从而
\begin{equation*}
\int_\Omega\abs{\nabla u}^2G'(u)\le0.
\end{equation*}
令 $H(t)=\int_0^t[G'(s)]^{1/2}ds$, 于是 $H(u)\in H_0^1(\Omega)$ 且
\begin{equation*}
\abs{\nabla H(u)}^2=\abs{\nabla u}^2G'(u)=0.
\end{equation*}
由此推出\footnote{注意若 $f\in W_0^{1,p}(\Omega)$, $1\le p<\infty$ 且
$\nabla f=0$ 在 $\Omega$ 中, 则 $f=0$ 在 $\Omega$ 中.事实上, 设 $\bar f$
为 $f$ 在 $\Omega$ 外按 $0$ 延拓所得; 则 $\bar f\in W^{1,p}(\R^N)$ 且
$\nabla\bar f=\nabla f=0$ (见命题  \ref{prop:9.18}).于是 $\bar f$ 为常数
(见注记 \ref{rem:9.7}), 而由 $\bar f\in L^p(\R^N)$ 得
$\bar f\equiv0$.} $H(u)=0$ 在 $\Omega$ 中, 从而 $u\le0$ 在 $\Omega$ 中.

现在我们以下述形式证明 \eqref{eq:9.80}:
\begin{equation*}
[f\le0\ \text{在}\ \Omega\ \text{中}]
\Rightarrow\left[u\le\sup_\Gamma u\ \text{在}\ \Omega\ \text{中}\right].
\tag*{(80$'$)}
\end{equation*}
令 $K=\sup_\Gamma u$; 则 $(u-K)$ 满足 \eqref{eq:9.78}, 因为 $a_0\equiv0$
且 $(u-K)\in H^1(\Omega)$ (由于 $\Omega$ 有界).应用
(79$'$), 我们得到 $u-K\le0$ 在 $\Omega$ 中, 即
(80$'$).最后, \eqref{eq:9.81} 由 \eqref{eq:9.80} 与
(80$'$) 推出.
\end{proof}

\begin{proposition}[Neumann 问题的极大值原理]
\label{prop:9.30}
设 $f\in L^2(\Omega)$, $u\in H^1(\Omega)$ 使得
\begin{equation*}
\int_\Omega\nabla u\cdot\nabla\varphi+\int_\Omega u\varphi
=\int_\Omega f\varphi\quad\forall\varphi\in H^1(\Omega).
\end{equation*}
则对几乎一切 $x\in\Omega$,
\begin{equation*}
\inf_\Omega f\le u(x)\le\sup_\Omega f.
\end{equation*}
\end{proposition}

\begin{proof}
仿照定理  \ref{thm:9.27} 的证明.
\end{proof}

\section{特征函数与谱分解}
\label{sec:9.8}

本节中我们假设 $\Omega$ 是有界开集.

\begin{keytheorem}
\label{thm:9.31}
存在 $L^2(\Omega)$ 的 Hilbert 基 $(e_n)_{n\ge1}$ 以及实数列
$(\lambda_n)_{n\ge1}$, 使得 $\lambda_n>0$ $\forall n$ 且
$\lambda_n\to+\infty$, 并且
\begin{equation}
e_n\in H_0^1(\Omega)\cap C^\infty(\Omega),
\label{eq:9.82}
\end{equation}
\begin{equation}
-\Delta e_n=\lambda_ne_n\quad\text{在}\ \Omega\ \text{中}.
\label{eq:9.83}
\end{equation}
我们说 $\lambda_n$ 是 $-\Delta$ (带 Dirichlet 边界条件) 的\emph{特征值},
$e_n$ 是相应的\emph{特征函数}.
\end{keytheorem}

\begin{proof}
给定 $f\in L^2(\Omega)$, 令 $u=Tf$ 为问题
\begin{equation}
\int_\Omega\nabla u\cdot\nabla\varphi=\int_\Omega f\varphi
\quad\forall\varphi\in H_0^1(\Omega)
\label{eq:9.84}
\end{equation}
的唯一解 $u\in H_0^1(\Omega)$.我们把 $T$ 看作从 $L^2(\Omega)$ 到
$L^2(\Omega)$ 的算子.则 $T$ 是自伴紧算子 (重复定理  \ref{thm:8.22} 的证明,
并利用 $H_0^1(\Omega)\subset L^2(\Omega)$ 为紧嵌入这一事实).另一方面,
% 原书此处印作 "repeat the proof of Theorem 8.21" (印刷页 311); 第~
% \ref{chap:8}~章中
% 8.21 是命题 (Neumann 情形的极大值原理), 讲特征函数与 Hilbert 基的是定理
% 8.22. 依技能规范保留原样并记录, 不擅自更正.
$N(T)=\{0\}$ 且 $(Tf,f)_{L^2}\ge0$ $\forall f\in L^2$.我们推出 (应用定理
 \ref{thm:6.11}) $L^2$ 存在由 $T$ 的特征函数组成的 Hilbert 基 $(e_n)$,
相应的特征值 $(\mu_n)$ 满足 $\mu_n>0$ $\forall n$ 且 $\mu_n\to0$.于是
$e_n\in H_0^1(\Omega)$ 且
\begin{equation*}
\int_\Omega\nabla e_n\cdot\nabla\varphi
=\frac1{\mu_n}\int_\Omega e_n\varphi\quad\forall\varphi\in H_0^1(\Omega).
\end{equation*}
换言之, $e_n$ 是 \eqref{eq:9.83} 的弱解, 其中 $\lambda_n=1/\mu_n$.由第~\ref{sec:9.6}~节的正则性结果 (见注记 \ref{rem:9.25}) 我们知道
$e_n\in H^2(\omega)$ 对每个 $\omega\subset\subset\Omega$ 成立.由此推出
$e_n\in H^4(\omega)$ 对每个 $\omega\subset\subset\Omega$ 成立, 然后
$e_n\in H^6(\omega)$ 对每个 $\omega\subset\subset\Omega$ 成立, 等等.于是
$e_n\in\cap_{m\ge1}H^m(\omega)$ 对一切 $\omega\subset\subset\Omega$ 成立.
因此 $e_n\in C^\infty(\omega)$ 对一切 $\omega\subset\subset\Omega$ 成立, 即
$e_n\in C^\infty(\Omega)$.
\end{proof}

\begin{remark}
\label{rem:9.28}
在定理  \ref{thm:9.31} 的假设下, 序列 $(e_n/\sqrt{\lambda_n})$ 是赋以标量积
$\int_\Omega\nabla u\cdot\nabla v$ 的 $H_0^1(\Omega)$ 的 Hilbert 基, 而
$(e_n/\sqrt{\lambda_n+1})$ 是赋以标量积
$\int_\Omega(\nabla u\cdot\nabla v+uv)$ 的 $H_0^1(\Omega)$ 的 Hilbert 基.
事实上, 显然 $(e_n/\sqrt{\lambda_n})$ 在 $H_0^1(\Omega)$ 中正交规范 (利用
\eqref{eq:9.83}).还须验证 $e_n$ 张成的向量空间在 $H_0^1(\Omega)$ 中稠密.
为此, 设 $f\in H_0^1(\Omega)$ 使得 $(f,e_n)_{H_0^1}=0$ $\forall n$.我们要
证明 $f=0$.由 \eqref{eq:9.83} 有 $\lambda_n\int_\Omega e_nf=0$ $\forall n$,
从而 $f=0$ (因为 $(e_n)$ 是 $L^2(\Omega)$ 的 Hilbert 基).
\end{remark}

\begin{remark}
\label{rem:9.29}
在定理  \ref{thm:9.31} 的假设下 (对一般的有限区域 $\Omega$), 可以证明
$e_n\in L^\infty(\Omega)$.另一方面, 若 $\Omega$ 属于 $C^\infty$ 类, 则
$e_n\in C^\infty(\overline\Omega)$; 这容易由定理  \ref{thm:9.25} 推出.
\end{remark}

\begin{remark}
\label{rem:9.30}
设 $a_{ij}\in L^\infty(\Omega)$ 满足椭圆性条件 \eqref{eq:9.36}, 且
$a_0\in L^\infty(\Omega)$.则存在 $L^2(\Omega)$ 的 Hilbert 基 $(e_n)$ 以及
实数列 $(\lambda_n)$, $\lambda_n\to+\infty$, 使得 $e_n\in H_0^1(\Omega)$ 且
\begin{equation*}
\int_\Omega\sum_{i,j}a_{ij}\frac{\partial e_n}{\partial x_i}
\frac{\partial\varphi}{\partial x_j}+\int_\Omega a_0e_n\varphi
=\lambda_n\int_\Omega e_n\varphi\quad\forall\varphi\in H_0^1(\Omega).
\end{equation*}
\end{remark}

\section*{第 9 章评注}
\addcontentsline{toc}{section}{第 9 章评注}

本章是 Sobolev 空间理论与椭圆方程理论的引论.希望深入钻研这个广阔课题的
读者可以查阅一份内容丰富的文献目录; 我们谨援引其中一部分:
S.~Agmon \bref{Agmon}{1}, L.~Bers--F.~John--M.~Schechter
\bref{BersJohnSchechter}{1}, J.-L.~Lions \bref{Lions}{1},
J.-L.~Lions--E.~Magenes \bref{LionsMagenes}{1}, A.~Friedman
\bref{Friedman}{2}, M.~Miranda \bref{Miranda}{1}, G.~Folland
\bref{Folland}{1}, F.~Treves \bref{Treves}{4}, R.~Adams \bref{Adams}{1},
D.~Gilbarg--N.~Trudinger \bref{GilbargTrudinger}{1},
G.~Stampacchia \bref{Stampacchia}{1}, R.~Courant--D.~Hilbert
\bref{CourantHilbert}{1} 第 2 卷, H.~Weinberger \bref{Weinberger}{1},
L.~Nirenberg \bref{Nirenberg}{1}, E.~Giusti \bref{Giusti}{2},
L.~C.~Evans \bref{Evans}{1}, M.~Giaquinta \bref{Giaquinta}{1},
E.~Lieb--M.~Loss \bref{Lieb}{1}, M.~Taylor \bref{Taylor}{1},
W.~Ziemer \bref{Ziemer}{1}, O.~Ladyzhenskaya--N.~Uraltseva
\bref{Ladyzhenskaya}{1}, N.~Krylov \bref{Krylov}{1}, \bref{Krylov}{2},
V.~Maz'ja \bref{Mazja}{1}, C.~Morrey \bref{Morrey}{1},
Y.~Z.~Chen--L.~C.~Wu \bref{ChenWu}{1}, E.~DiBenedetto
\bref{DiBenedetto}{1}, Q.~Han--F.~H.~Lin \bref{HanLin}{1}, J.~Jost
\bref{Jost}{1}, W.~Strauss \bref{Strauss}{1}, 以及这些著作中援引的文献.

\cmtnum{1.} 在第~
\ref{chap:9}~章中我们常常假设 $\Omega$ 属于 $C^1$ 类; 这排除了例如带
``角''的区域.在各种情形下都可以减弱这个假设, 用某种``异乎寻常''的条件来
代替它: $\Omega$ 分片属于 $C^1$ 类, $\Omega$ 是 Lipschitz 区域, $\Omega$
具有锥性质, $\Omega$ 具有线段性质, 等等; 参见例如 R.~Adams \bref{Adams}{1}
与 S.~Agmon \bref{Agmon}{1}.

\cmtnum{2.} 定理  \ref{thm:9.7} (延拓算子的存在性) 可以借助反射延拓技巧的
适当推广, 移植到空间 $W^{m,p}(\Omega)$ ($\Omega$ 属于 $C^m$ 类); 参见例如
R.~Adams \bref{Adams}{1} 与 S.~Agmon \bref{Agmon}{1}.

\cmtnum{3.} 一些涉及 Sobolev 范数的非常有用的不等式.

\cmtnum{$\bullet$\ A. Poincaré--Wirtinger 不等式.}

设 $\Omega$ 为 $C^1$ 类连通开集, $1\le p\le\infty$.则存在常数 $C$ 使得
\begin{equation*}
\norm{u-\bar u}_p\le C\norm{\nabla u}_p\quad\forall u\in W^{1,p}(\Omega),
\end{equation*}
其中 $\bar u=\frac1{\abs{\Omega}}\int_\Omega u$.由此, 利用 Sobolev 不等式
推出: 若 $p<N$, 则
\begin{equation*}
\norm{u-\bar u}_{p^\star}\le C\norm{\nabla u}_p
\quad\forall u\in W^{1,p}(\Omega).
\end{equation*}

\cmtnum{$\bullet$\ B. Hardy 不等式.}

设 $\Omega$ 为 $C^1$ 类有界开集, $1<p<\infty$.令
$d(x)=\dist(x,\Gamma)$.则存在常数 $C$ 使得
\begin{equation*}
\left\|\frac ud\right\|_p\le C\norm{\nabla u}_p
\quad\forall u\in W_0^{1,p}(\Omega).
\end{equation*}
反之,
\begin{equation*}
[u\in W^{1,p}(\Omega)\ \text{且}\ (u/d)\in L^p(\Omega)]
\Rightarrow[u\in W_0^{1,p}(\Omega)];
\end{equation*}
见 J.-L.~Lions--E.~Magenes \bref{LionsMagenes}{1}.

\cmtnum{$\bullet$\ C. Gagliardo--Nirenberg 插值不等式.}

这里我们只举出几个在应用中经常遇到的例子.一般情形见
L.~Nirenberg \bref{Nirenberg}{1} 或 A.~Friedman \bref{Friedman}{2}.

为确定起见, 设 $\Omega\subset\R^N$ 为正则有界开集.

\noindent\emph{例 1.} 设 $u\in L^p(\Omega)\cap W^{2,r}(\Omega)$, 其中
$1\le p\le\infty$, $1\le r\le\infty$.则 $u\in W^{1,q}(\Omega)$, 其中 $q$ 是
$p$ 与 $r$ 的调和平均, 即 $\frac1q=\frac12(\frac1p+\frac1r)$, 并且
\begin{equation*}
\norm{Du}_{L^q}\le C\norm{u}_{W^{2,r}}^{1/2}\norm{u}_{L^p}^{1/2}.
\end{equation*}
特例:

\noindent (a) $p=\infty$, 从而 $q=2r$.我们有
\begin{equation*}
\norm{Du}_{L^q}\le C\norm{u}_{W^{2,r}}^{1/2}\norm{u}_{L^\infty}^{1/2}.
\end{equation*}
除别的用处之外, 这个不等式可以用来证明 $W^{2,r}\cap L^\infty$ 是一个代数,
也就是说,
\begin{equation*}
u,v\in W^{2,r}\cap L^\infty\Rightarrow uv\in W^{2,r}\cap L^\infty
\end{equation*}
(这个性质对 $W^{m,r}\cap L^\infty$ ($m$ 为整数, $m\ge2$) 仍然成立).

\noindent (b) $p=q=r$.我们有
\begin{equation*}
\norm{Du}_{L^p}\le C\norm{u}_{W^{2,p}}^{1/2}\norm{u}_{L^p}^{1/2},
\end{equation*}
由此特别推出
\begin{equation*}
\norm{Du}_{L^p}\le\varepsilon\norm{D^2u}_{L^p}
+C_\varepsilon\norm{u}_{L^p}\quad\forall\varepsilon>0.
\end{equation*}

\noindent\emph{例 2.} 设 $1\le q\le p<\infty$.则
\begin{equation}
\norm{u}_{L^p}\le C\norm{u}_{L^q}^{1-a}\norm{u}_{W^{1,N}}^a
\quad\forall u\in W^{1,N}(\Omega),\ \text{其中}\ a=1-(q/p).
\label{eq:9.85}
\end{equation}
我们指出经常用到的一个特例
\begin{equation*}
N=2,\quad p=4,\quad q=2,\quad\text{且}\ a=1/2,
\end{equation*}
也就是说,
\begin{equation*}
\norm{u}_{L^4}\le C\norm{u}_{L^2}^{1/2}\norm{u}_{H^1}^{1/2}
\quad\forall u\in H^1(\Omega).
\end{equation*}
在这方面我们注意到, 还有通常的插值不等式 (第~
\ref{chap:4}~章注记 2)
\begin{equation*}
\norm{u}_{L^p}\le\norm{u}_{L^q}^{1-a}\norm{u}_{L^\infty}^a
\quad\text{其中}\ a=1-(q/p),
\end{equation*}
但它不能推出 \eqref{eq:9.85}, 因为 $W^{1,N}$ 不包含在 $L^\infty$ 中.

\noindent\emph{例 3.} 设 $1\le q\le p\le\infty$ 且 $r>N$.则
\begin{equation}
\norm{u}_{L^p}\le C\norm{u}_{L^q}^{1-a}\norm{u}_{W^{1,r}}^a
\quad\forall u\in W^{1,r}(\Omega),
\label{eq:9.86}
\end{equation}
其中 $a=\left(\frac1q-\frac1p\right)\Big/\left(\frac1q+\frac1N-\frac1r\right)$.

\cmtnum{$\bullet$\ 4.} 下述性质有时是有用的.设 $u\in W^{1,p}(\Omega)$,
$1\le p\le\infty$, $\Omega$ 为任意开集.则 $\nabla u=0$ 在集合
$\{x\in\Omega;\ u(x)=k\}$ 上几乎处处成立, 其中 $k$ 为任意常数.

\cmtnum{$\star$5.} 当 $p>N$ 时, $W^{1,p}(\Omega)$ 中的函数在 $\Omega$ 中
几乎处处按通常意义可微.更精确地说, 设 $u\in W^{1,p}(\Omega)$, $p>N$.则
存在零测集 $A\subset\Omega$ 使得
\begin{equation*}
\lim_{h\to0}\frac{u(x+h)-u(x)-h\cdot\nabla u(x)}{\abs{h}}=0
\quad\forall x\in\Omega\setminus A.
\end{equation*}
当 $u\in W^{1,p}(\Omega)$ 且 $p\le N$ ($N>1$) 时这个性质不成立.关于这个
问题可查阅 E.~Stein \bref{Stein}{1} (第~
\ref{chap:8}~章).

\cmtnum{6.} \emph{分数阶 Sobolev 空间.}

可以定义一族介于 $L^p(\Omega)$ 与 $W^{1,p}(\Omega)$ 之间的空间.更精确地说,
若 $0<s<1$ ($s\in\R$) 且 $1\le p<\infty$, 令
\begin{equation*}
W^{s,p}(\Omega)=\left\{u\in L^p(\Omega);\
\frac{\abs{u(x)-u(y)}}{\abs{x-y}^{s+(N/p)}}\in L^p(\Omega\times\Omega)
\right\},
\end{equation*}
赋以自然的范数.令 $H^s(\Omega)=W^{s,2}(\Omega)$.关于这些空间的研究, 参见
例如 R.~Adams \bref{Adams}{1}, J.-L.~Lions--E.~Magenes
\bref{LionsMagenes}{1}, P.~Malliavin \bref{Malliavin}{1},
H.~Triebel \bref{Triebel}{1} 以及 L.~Grafakos \bref{Grafakos}{1}.
空间 $W^{s,p}(\Omega)$ 也可以定义为 $W^{1,p}$ 与 $L^p$ 之间的插值空间;
当 $p=2$ 且 $\Omega=\R^N$ 时还可以用 Fourier 变换来定义.最后, 对实数
$s$ (非整数) 且 $s>1$, 我们如下定义 $W^{s,p}(\Omega)$: 写
$s=m+\sigma$, 其中 $m$ 是 $s$ 的整数部分, 并令
\begin{equation*}
W^{s,p}(\Omega)=\{u\in W^{m,p}(\Omega);\
D^\alpha u\in W^{\sigma,p}(\Omega)\ \forall\alpha,\ \abs{\alpha}=m\}.
\end{equation*}
借助局部坐标卡还可以定义 $W^{s,p}(\Gamma)$, 其中 $\Gamma$ 是光滑流形
(例如正则开集的边界).这些空间在迹理论中起重要作用 (见评注 7).

\cmtnum{$\bullet$\ 7.} \emph{迹理论.}

设 $1\le p<\infty$.我们先从一个基本引理开始.

\begin{lemma}
\label{lem:9.9}
设 $\Omega=\R^N_+$.则存在常数 $C$ 使得
\begin{equation*}
\left(\int_{\R^{N-1}}\abs{u(x',0)}^p\,dx'\right)^{1/p}
\le C\norm{u}_{W^{1,p}(\Omega)}\quad\forall u\in C_c^1(\R^N).
\end{equation*}
\end{lemma}

\begin{proof}
令 $G(t)=\abs{t}^{p-1}t$, 并设 $u\in C_c^1(\R^N)$.我们有
\begin{equation*}
\begin{aligned}
G(u(x',0))&=-\int_0^{+\infty}\frac{\partial}{\partial x_N}
G(u(x',x_N))\,dx_N\\
&=-\int_0^{+\infty}G'(u(x',x_N))
\frac{\partial u}{\partial x_N}(x',x_N)\,dx_N.
\end{aligned}
\end{equation*}
于是
\begin{equation*}
\begin{aligned}
\abs{u(x',0)}^p
&\le p\int_0^\infty\abs{u(x',x_N)}^{p-1}
\left|\frac{\partial u}{\partial x_N}(x',x_N)\right|dx_N\\
&\le C\left(\int_0^\infty\abs{u(x',x_N)}^p\,dx_N
+\int_0^\infty\left|\frac{\partial u}{\partial x_N}(x',x_N)\right|^p
dx_N\right),
\end{aligned}
\end{equation*}
对 $x'\in\R^{N-1}$ 积分即得结论.
\end{proof}

由引理  \ref{lem:9.9} 可以推出: 由
$u\mapsto u|_\Gamma$ ($\Gamma=\partial\Omega=\R^{N-1}\times\{0\}$) 定义在
$C_c^1(\R^N)$ 上, 取值于 $L^p(\Gamma)$ 的映射, 按稠密性可以延拓为从
$W^{1,p}(\Omega)$ 到 $L^p(\Gamma)$ 的有界线性算子.这个算子按定义就是 $u$ 在
$\Gamma$ 上的\emph{迹}; 它也记作 $u|_\Gamma$.

我们指出, $L^p(\R^N_+)$ 与 $W^{1,p}(\R^N_+)$ 之间有一个本质区别:
$L^p(\R^N_+)$ 中的函数在 $\Gamma$ 上没有迹.利用局部坐标卡, 容易想象当
$\Omega$ 是 $\R^N$ 中的正则开集 (例如 $\Omega$ 属于 $C^1$ 类且 $\Gamma$ 有界)
时, 如何对函数 $u\in W^{1,p}(\Omega)$ 定义它在 $\Gamma=\partial\Omega$ 上的
迹.在这种情形下 $u|_\Gamma\in L^p(\Gamma)$ (关于面测度 $d\sigma$).迹的最
重要的性质如下:

\noindent (i) 若 $u\in W^{1,p}(\Omega)$, 则实际上
$u|_\Gamma\in W^{1-(1/p),p}(\Gamma)$ 且
\begin{equation*}
\norm{u|_\Gamma}_{W^{1-(1/p),p}(\Gamma)}\le C\norm{u}_{W^{1,p}(\Omega)}
\quad\forall u\in W^{1,p}(\Omega).
\end{equation*}
此外, \index{jisuanzi@迹算子 (trace operator)}
迹算子 $u\mapsto u|_\Gamma$ 是从 $W^{1,p}(\Omega)$ 到
$W^{1-(1/p),p}(\Gamma)$ 上的满射.

\noindent (ii) 迹算子的核是 $W_0^{1,p}(\Omega)$, 即
\begin{equation*}
W_0^{1,p}(\Omega)=\{u\in W^{1,p}(\Omega);\ u|_\Gamma=0\}.
\end{equation*}

\noindent (iii) 我们有 Green 公式
\begin{equation*}
\int_\Omega\frac{\partial u}{\partial x_i}v
=-\int_\Omega u\frac{\partial v}{\partial x_i}
+\int_\Gamma uv(\vec n\cdot\vec e_i)\,d\sigma\quad\forall u,v\in H^1(\Omega),
\end{equation*}
其中 $\vec n$ 是 $\Gamma$ 的单位外法向量.注意面积分是有意义的, 因为
$u,v\in L^2(\Gamma)$.

同样地, 对函数 $u\in W^{2,p}(\Omega)$ 我们可以谈论 $\frac{\partial u}{\partial n}$:
令 $\frac{\partial u}{\partial n}=(\nabla u)|_\Gamma\cdot\vec n$, 这是有意义的,
因为 $(\nabla u)|_\Gamma\in L^p(\Gamma)^N$, 且
$\frac{\partial u}{\partial n}\in L^p(\Gamma)$ (实际上
$\frac{\partial u}{\partial n}\in W^{1-(1/p),p}(\Gamma)$).Green 公式也成立:
\begin{equation*}
-\int_\Omega(\Delta u)v=\int_\Omega\nabla u\cdot\nabla v
-\int_\Gamma\frac{\partial u}{\partial n}v\,d\sigma\quad\forall u,v\in H^2(\Omega).
\end{equation*}

\noindent (iv) 算子
$u\mapsto\{u|_\Gamma,\frac{\partial u}{\partial n}\}$ 是从 $W^{2,p}(\Omega)$
到 $W^{2-(1/p),p}(\Gamma)\times W^{1-(1/p),p}(\Gamma)$ 上的有界线性满射.
关于这些问题, $p=2$ 的情形见 J.-L.~Lions--E.~Magenes
\bref{LionsMagenes}{1} (关于 $p\neq2$ 的情形见其中援引的文献).

\cmtnum{8.} \emph{2$m$ 阶算子与椭圆组.}

第~
\ref{chap:9}~章中证明的存在性与正则性结果可以推广到 $2m$ 阶椭圆算子与椭圆组.
\footnote{但极大值原理不能推广, 除非在非常特殊的情形.} 本质要素之一是
G\aa rding 不等式.关于这些问题, 参见 S.~Agmon \bref{Agmon}{1},
J.-L.~Lions--E.~Magenes \bref{LionsMagenes}{1},
S.~Agmon--A.~Douglis--L.~Nirenberg \bref{AgmonDouglisNirenberg}{1}.
$2m$ 阶算子与某些组在力学与物理学中起重要作用.我们特别指出双调和算子
$\Delta^2$ (板的弯曲理论), 弹性力学组与 Stokes 组 (流体力学); 参见例如
Ph.~Ciarlet \bref{Ciarlet}{1}, G.~Duvaut--J.-L.~Lions
\bref{DuvautLions}{1}, R.~Temam \bref{Temam}{1},
J.~Ne\v cas--L.~Hlava\v cek \bref{NecasHlavacek}{1}, M.~Gurtin
\bref{Gurtin}{1}.

\cmtnum{9.} \emph{$L^p$ 与 $C^{0,\alpha}$ 空间中的正则性.}

第~
\ref{chap:9}~章中关于 $p=2$ 证明的正则性定理可以推广到 $p\neq2$ 的情形.

\begin{keytheorem}[Agmon--Douglis--Nirenberg]
\label{thm:9.32}
设 $\Omega$ 属于 $C^2$ 类且 $\Gamma$ 有界, $1<p<\infty$.则对一切
$f\in L^p(\Omega)$, 方程
\begin{equation}
-\Delta u+u=f\quad\text{在}\ \Omega\ \text{中}
\label{eq:9.87}
\end{equation}
存在唯一的解 $u\in W^{2,p}(\Omega)\cap W_0^{1,p}(\Omega)$.此外, 若
$\Omega$ 属于 $C^{m+2}$ 类且 $f\in W^{m,p}(\Omega)$ ($m\ge1$ 为整数), 则
\begin{equation*}
u\in W^{m+2,p}(\Omega)\quad\text{且}\quad
\norm{u}_{W^{m+2,p}}\le C\norm{f}_{W^{m,p}}.
\end{equation*}
把 \eqref{eq:9.87} 换成带光滑系数的二阶椭圆方程, 有类似的结果.定理
 \ref{thm:9.32} 的证明比 $p=2$ 的情形 (定理  \ref{thm:9.25}) 复杂得多.
``古典''的方法本质上依赖于两个要素:
\end{keytheorem}

\noindent (a) 利用基本解给出 $u$ 的显式表示公式.例如, 若 $\Omega=\R^3$,
则 \eqref{eq:9.87} 的解由 $u=G\star f$ 给出, 其中
$G(x)=\frac{c}{\abs{x}}e^{-\abs{x}}$.于是形式上有
$\frac{\partial^2u}{\partial x_i\partial x_j}
=\frac{\partial^2G}{\partial x_i\partial x_j}\star f$;``不幸的是''
$\frac{\partial^2G}{\partial x_i\partial x_j}$ 不属于 $L^1(\R^3)$,\footnote{
但几乎属于!} 因为在 $x=0$ 处有奇性, 因此不能对卷积积应用初等估计 (例如定理
4.15).

\noindent (b) 为克服这个困难, 人们使用 Calderón--Zygmund 的 $L^p$ 奇异积分
理论 (参见例如 E.~Stein \bref{Stein}{1} 与 L.~Bers--F.~John--M.~Schechter
\bref{BersJohnSchechter}{1}).

\noindent\emph{注意:} 定理  \ref{thm:9.32} 的结论对 $p=1$ 与 $p=\infty$
不成立.

在 H\"older 空间框架中的另一个基本正则性结果如下.\footnote{回忆当
$0<\alpha<1$ 且 $m$ 为整数时,
\begin{equation*}
C^{0,\alpha}(\Omega)=\left\{u\in C(\Omega);\
\sup_{\substack{x,y\in\Omega\\x\neq y}}
\frac{\abs{u(x)-u(y)}}{\abs{x-y}^\alpha}<\infty\right\},
\end{equation*}
且 $C^{m,\alpha}(\Omega)=\{u\in C(\Omega);\ D^\beta u\in C^{0,\alpha}(\Omega)\
\forall\beta,\ \abs{\beta}=m\}$.}

\begin{keytheorem}[Schauder]
\label{thm:9.33}
设 $\Omega$ 有界且属于 $C^{2,\alpha}$ 类, $0<\alpha<1$.则对每个
$f\in C^{0,\alpha}(\Omega)$, 问题
\begin{equation}
\begin{cases}
-\Delta u+u=f & \text{在}\ \Omega\ \text{中},\\
u=0 & \text{在}\ \Gamma\ \text{上}
\end{cases}
\label{eq:9.88}
\end{equation}
存在唯一的解 $u\in C^{2,\alpha}(\Omega)$.此外, 若 $\Omega$ 属于
$C^{m+2,\alpha}$ 类 ($m\ge1$ 为整数) 且 $f\in C^{m,\alpha}(\Omega)$, 则
$u\in C^{m+2,\alpha}(\Omega)$ 且
$\norm{u}_{C^{m+2,\alpha}}\le C\norm{f}_{C^{m,\alpha}}$.
\end{keytheorem}

把 \eqref{eq:9.88} 换成带光滑系数的二阶椭圆算子, 有类似的结果.定理
 \ref{thm:9.33} 的证明——如同定理  \ref{thm:9.32} 的证明一样——依赖于 $u$ 的
显式表示以及 H\"older, Korn, Lichtenstein, Giraud 的 $C^{0,\alpha}$ 空间
奇异积分理论.关于这个课题, 参见
S.~Agmon--A.~Douglis--L.~Nirenberg \bref{AgmonDouglisNirenberg}{1},
L.~Bers--F.~John--M.~Schechter \bref{BersJohnSchechter}{1},
C.~Morrey \bref{Morrey}{1}, D.~Gilbarg--N.~Trudinger
\bref{GilbargTrudinger}{1}.另一种避免奇异积分理论的方法由 Campanato 与
Stampacchia 提出 (参见例如 Y.~Z.~Chen--L.~C.~Wu \bref{ChenWu}{1} 与
E.~Giusti \bref{Giusti}{2}).其它初等方法由 A.~Brandt \bref{Brandt}{1}
(仅基于极大值原理) 与 L.~Simon \bref{Simon}{2} 发展起来.

设 $\Omega$ 为有界正则开集, $f\in C(\Omega)$.由定理  \ref{thm:9.32},
存在 $u\in W^{2,p}(\Omega)\cap W_0^{1,p}(\Omega)$ (对一切 $1<p<\infty$),
它是 \eqref{eq:9.87} 的唯一解.特别地, $u\in C^{1,\alpha}(\Omega)$ 对一切
$0<\alpha<1$ 成立 (由 Morrey 定理 (定理  \ref{thm:9.12})).一般说来, $u$
不属于 $C^2$, 甚至不属于 $W^{2,\infty}$.这就解释了为什么人们常常避免在
空间 $L^1(\Omega)$, $L^\infty(\Omega)$ 与 $C(\Omega)$ 中工作——对于这些
空间我们没有最优的正则性结果.

定理  \ref{thm:9.32} 与  \ref{thm:9.33} 可以推广到 $2m$ 阶椭圆算子与椭圆组;
见 S.~Agmon--A.~Douglis--L.~Nirenberg \bref{AgmonDouglisNirenberg}{1}.
最后我们指出, 从另一个方向看, 带不连续系数的二阶椭圆方程是大量工作的
主题.我们举出下述著名结果为例.

\begin{keytheorem}[De Giorgi, Nash, Stampacchia]
\label{thm:9.34}
设 $\Omega\subset\R^N$, $N\ge2$, 为有界正则开集.假设函数
$a_{ij}\in L^\infty(\Omega)$ 满足椭圆性条件 \eqref{eq:9.36}.设
$f\in L^p(\Omega)$, $p>N/2$, 且 $u\in H_0^1(\Omega)$ 使得
\begin{equation*}
\int_\Omega\sum_{i,j}a_{ij}\frac{\partial u}{\partial x_i}
\frac{\partial\varphi}{\partial x_j}=\int_\Omega f\varphi
\quad\forall\varphi\in H_0^1(\Omega).
\end{equation*}
则 $u\in C^{0,\alpha}(\Omega)$, 其中 $0<\alpha<1$ 依赖于 $\Omega$, $a_{ij}$
与 $p$.
\end{keytheorem}

关于这些问题, 参见 G.~Stampacchia \bref{Stampacchia}{1},
D.~Gilbarg--N.~Trudinger \bref{GilbargTrudinger}{1},
O.~Ladyzhenskaya--N.~Uraltseva \bref{Ladyzhenskaya}{1} 以及
E.~Giusti \bref{Giusti}{2}.

\cmtnum{10.} \emph{变分方法的一些缺点以及如何绕开它们!}

变分方法非常容易地给出弱解的存在性.它并不总是适用, 但可以补充完善.
我们指出两个例子.设 $\Omega\subset\R^N$ 为有界正则开集.

\noindent (a) \emph{对偶方法.} 设 $f\in L^1(\Omega)$——甚至设 $f$ 是
$\Omega$ 上的一个 (Radon) 测度——并寻找问题
\begin{equation}
\begin{cases}
-\Delta u+u=f & \text{在}\ \Omega\ \text{中},\\
u=0 & \text{在}\ \Gamma\ \text{上}
\end{cases}
\label{eq:9.89}
\end{equation}
的解.只要 $N>1$, 线性泛函 $\varphi\mapsto\int_\Omega f\varphi$ 对每个
$\varphi\in H_0^1(\Omega)$ 就没有定义, 因此变分方法无效.另一方面, 可以使用
下述技巧.用 $T:L^2(\Omega)\to L^2(\Omega)$ 记算子 $f\mapsto u$ (其中 $u$ 是
\eqref{eq:9.89} 的解, 它对 $f\in L^2(\Omega)$ 存在).我们知道 $T$ 是自伴的.
另一方面 (定理  \ref{thm:9.32}), 对 $2\le p<\infty$ 有
$T:L^p(\Omega)\to W^{2,p}(\Omega)$, 而由 Sobolev 与 Morrey 定理, 当
$p>N/2$ 时 $T:L^p(\Omega)\to C^0(\Omega)$.由对偶性推出, 当 $p>N/2$ 时
\begin{equation*}
T^\star:\mathcal M(\Omega)=C^0(\Omega)^\star\to L^{p'}(\Omega).
\end{equation*}
由于 $T$ 在 $L^2$ 中自伴, $T^\star$ 是 $T$ 的延拓; 于是可以把
$u=T^\star f$ 看作 \eqref{eq:9.89} 的广义解.事实上, 若
$f\in L^1(\Omega)$, 则 $u=T^\star f\in L^q(\Omega)$ 对一切
$q<N/(N-2)$ 成立; $u$ 是 \eqref{eq:9.89} 在
\begin{equation*}
-\int_\Omega u\Delta\varphi+\int_\Omega u\varphi=\int_\Omega f\varphi
\quad\forall\varphi\in C^2(\Omega),\ \varphi=0\ \text{在}\ \Gamma\ \text{上}
\end{equation*}
意义下的唯一 (\emph{很}) 弱解.按照同样的想法, 可以研究 $f$ 给定于
$H^{-m}(\Omega)$ 时的 \eqref{eq:9.89}; 见 J.-L.~Lions--E.~Magenes
\bref{LionsMagenes}{1}.

\noindent (b) \emph{稠密性方法.} 设 $g\in C(\Gamma)$ 并寻找问题
\begin{equation}
\begin{cases}
-\Delta u+u=0 & \text{在}\ \Omega\ \text{中},\\
u=g & \text{在}\ \Gamma\ \text{上}
\end{cases}
\label{eq:9.90}
\end{equation}
的解.一般说来, 若 $g\in C(\Gamma)$, 并不存在函数 $\tilde g\in H^1(\Omega)$
使得 $\tilde g|_\Gamma=g$ (见评注 7, 并注意 $C(\Gamma)$ 不包含在
$H^{1/2}(\Gamma)$ 中).因此不可能在 $H^1(\Omega)$ 中寻找 \eqref{eq:9.90} 的
解: 变分方法无效.然而, 我们有下述结果.

\begin{keytheorem}
\label{thm:9.35}
\eqref{eq:9.90} 存在唯一的解 $u\in C(\overline\Omega)\cap C^\infty(\Omega)$.
\end{keytheorem}

\begin{proof}
固定 $\tilde g\in C_c(\R^N)$ 使得 $\tilde g|_\Gamma=g$; 由
Tietze--Urysohn 定理 (参见例如 J.~Dieudonn\'e \bref{Dieudonne}{1},
J.~Dugundji \bref{Dugundji}{1}, J.~Munkres \bref{Munkres}{1})
这样的 $\tilde g$ 存在.设 $(\tilde g_n)$ 是 $C_c^\infty(\R^N)$ 中的序列,
使得 $\tilde g_n\to g$ 在 $\R^N$ 上一致收敛.令
$g_n=\tilde g_n|_\Gamma$.应用变分方法与正则性结果, 我们看到问题
\begin{equation*}
\begin{cases}
-\Delta u_n+u_n=0 & \text{在}\ \Omega\ \text{中},\\
u_n=g_n & \text{在}\ \Gamma\ \text{上}
\end{cases}
\end{equation*}
存在古典解 $u_n\in C^2(\Omega)$.由极大值原理 (推论  \ref{cor:9.28}) 我们有
\begin{equation*}
\norm{u_m-u_n}_{L^\infty(\Omega)}\le\norm{g_m-g_n}_{L^\infty(\Gamma)}.
\end{equation*}
于是 $(u_n)$ 是 $C(\overline\Omega)$ 中的 Cauchy 序列, 且
$u_n\to u$ 在 $C(\overline\Omega)$ 中.显然有
\begin{equation*}
\int_\Omega u(-\Delta\varphi+\varphi)=0\quad\forall\varphi\in C_c^\infty(\Omega),
\end{equation*}
从而 $u\in C^\infty(\Omega)$ (见注记 \ref{rem:9.25}).于是
$u\in C(\overline\Omega)\cap C^\infty(\Omega)$ 满足 \eqref{eq:9.90}.
\eqref{eq:9.90} 的解的唯一性由极大值原理推出 (见注记 \ref{rem:9.27}).
\end{proof}

\begin{starremark}
\label{rem:9.31}
在定理  \ref{thm:9.35} 中, 假设 $\Omega$ 足够光滑是本质的.当 $\Omega$ 有
``病态''边界时, 我们会遇到位势论中的问题 (正则点, Wiener 判别法等).

求解 \eqref{eq:9.90} 的另一种方法是在位势论中很经典的 Perron 方法.定义
\begin{equation*}
u(x)=\sup\{v(x);\ v\in C(\overline\Omega)\cap C^2(\Omega),\
-\Delta v+v\le0\ \text{在}\ \Omega\ \text{中},\ v\le g\ \text{在}\
\Gamma\ \text{上}\},
\end{equation*}
并 (直接) 证明 $u$ 满足 \eqref{eq:9.90}.满足
$-\Delta v+v\le0$ 在 $\Omega$ 中且 $v\le g$ 在 $\Gamma$ 上的函数 $v$
称为 \eqref{eq:9.90} 的\emph{下解}.
\end{starremark}

\cmtnum{11.} \emph{强极大值原理.}

当 $u$ 是古典解时, 我们可以加强命题  \ref{prop:9.29} 的结论.更精确地说,
设 $\Omega$ 为连通的有界正则开集, $a_{ij}\in C^1(\Omega)$ 满足椭圆性条件
\eqref{eq:9.36}, $a_i,a_0\in C(\Omega)$, $a_0\ge0$ 在 $\Omega$ 上.

\begin{keytheorem}[Hopf]
\index{qiangjidazhiyuanli@强极大值原理 (strong maximum principle)}
\label{thm:9.36}
设 $u\in C(\overline\Omega)\cap C^2(\Omega)$ 满足
\begin{equation}
-\sum_{i,j}\frac{\partial}{\partial x_j}
\left(a_{ij}\frac{\partial u}{\partial x_i}\right)
+\sum_ia_i\frac{\partial u}{\partial x_i}+a_0u=f
\quad\text{在}\ \Omega\ \text{中}.
\label{eq:9.91}
\end{equation}
假设 $f\ge0$ 在 $\Omega$ 上.若存在 $x_0\in\Omega$ 使得
$u(x_0)=\min_{\overline\Omega}u$, 且 $u(x_0)\le0$,\footnote{若 $a_0=0$,
则假设 $u(x_0)\le0$ 是多余的.} 则 $u$ 在 $\Omega$ 中为常数 (并且进一步
$f=0$ 在 $\Omega$ 中).
\end{keytheorem}

关于证明, 参见例如 L.~Bers--F.~John--M.~Schechter
\bref{BersJohnSchechter}{1}, D.~Gilbarg--N.~Trudinger
\bref{GilbargTrudinger}{1}, M.~Protter--H.~Weinberger
\bref{ProtterWeinberger}{1} 以及 P.~Pucci--J.~Serrin
\bref{PucciSerrin}{1}.

\begin{keycorollary}
\label{cor:9.37}
设 $u\in C(\overline\Omega)\cap C^2(\Omega)$ 满足 \eqref{eq:9.91},
其中 $f\ge0$ 在 $\Omega$ 上.假设 $u\ge0$ 在 $\Gamma$ 上.则
\begin{itemize}
	\item 或者 $u>0$ 在 $\Omega$ 中,
	\item 或者 $u\equiv0$ 在 $\Omega$ 中.
\end{itemize}
\end{keycorollary}

关于与极大值原理相关的其它结果 (Harnack 不等式等), 参见例如
G.~Stampacchia \bref{Stampacchia}{1},
D.~Gilbarg--N.~Trudinger \bref{GilbargTrudinger}{1},
M.~Protter--H.~Weinberger \bref{ProtterWeinberger}{1},
R.~Sperb \bref{Sperb}{1} 以及 P.~Pucci--J.~Serrin
\bref{PucciSerrin}{1}.

\cmtnum{12.} \emph{Laplace--Beltrami 算子.}

定义在 Riemann 流形 (带边界或不带边界) 上的椭圆算子, 特别是
Laplace--Beltrami 算子, 在微分几何与物理学中起重要作用; 参见例如
Y.~Choquet--C.~Dewitt--M.~Dillard \bref{ChoquetDewittDillard}{1}.

\cmtnum{13.} \emph{谱性质. 反问题.}

二阶椭圆算子的特征值与特征函数有许多值得注意的性质.这里我们举出其中
一些.设 $\Omega\subset\R^N$ 为连通的有界正则开集.设
$a_{ij}\in C^1(\Omega)$ 满足椭圆性条件 \eqref{eq:9.36}, $a_0\in C(\Omega)$.
设 $A$ 为算子
\begin{equation*}
Au=-\sum_{i,j}\frac{\partial}{\partial x_j}
\left(a_{ij}\frac{\partial u}{\partial x_i}\right)+a_0u
\end{equation*}
并带齐次 Dirichlet 条件 ($u=0$ 在 $\Gamma$ 上).用 $(\lambda_n)$ 记 $A$
的特征值按递增顺序排列的序列, 当 $n\to\infty$ 时
$\lambda_n\to+\infty$.则第一个特征值 $\lambda_1$ 的重数为 1 (这时说
$\lambda_1$ 是\emph{单}特征值),\footnote{当维数 $N\ge2$ 时, 其余特征值
的重数可以大于 1.} 且我们可以选取相应的特征函数 $e_1$ 使 $e_1>0$ 在
$\Omega$ 上; 这由 Krein--Rutman 定理推出 (见第~
\ref{chap:6}~章评注与问题 41).此外,
可以证明当 $n\to\infty$ 时 $\lambda_n\sim cn^{2/N}$, 其中 $c>0$; 见
S.~Agmon \bref{Agmon}{1}.

$\Omega$ 的几何性质\footnote{特别是当 $\Omega$ 是不带边界的 Riemann 流形
且 $A$ 是 Laplace--Beltrami 算子时.}与 $A$ 的谱之间的关系是大量研究的
主题; 参见例如 M.~Kac \bref{Kac}{1}, Marcel Berger \bref{Berger}{1},
R.~Osserman \bref{Osserman}{1}, I.~M.~Singer \bref{Singer}{1},
P.~B\'erard \bref{Berard}{1}, I.~Chavel \bref{Chavel}{1}.谱几何的目标
是仅从谱 $(\lambda_n)$ 的知识出发, ``还原''关于 $\Omega$ 的最大量的信息.

一个极其简单的问题是下述问题.设 $\Omega_1$ 与 $\Omega_2$ 是 $\R^2$ 中
两个有界区域; 假设算子 $-\Delta$ (带 Dirichlet 边界条件) 在 $\Omega_1$ 与
$\Omega_2$ 上的特征值相同.$\Omega_1$ 与 $\Omega_2$ 是否等距? 这个问题被
M.~Kac 戏称为: ``人能听出鼓的形状吗?''\footnote{因为附着在边界 $\Gamma$ 上
的膜振动的谐波是函数 $e_n(x)\sin\sqrt{\lambda_n}t$, 其中 $(\lambda_n,e_n)$
是 $-\Delta$ 带 Dirichlet 边界条件的特征值与特征函数.} 已知若 $\Omega_1$
是圆盘, 答案是肯定的.1991 年, C.~Gordon--L.~Webb--S.~Wolpert
\bref{GordonWebbWolpert}{1} 对带角的区域给出了否定的答案.Kac 问题对光滑
区域仍未解决.

另一类重要的``反问题''是仅从边界上的测量 (例如 Dirichlet-to-Neumann 映射)
或``无穷远处''的测量 (逆散射) 来确定偏微分方程中的系数与参数, 或者确定
内部物体的形状与特征.这类问题出现在许多领域 (医学成像, 地震学等); 参见
例如 G.~Uhlmann \bref{Uhlmann}{1}, C.~B.~Croke 等 \bref{CrokeEtAl}{1}.

\cmtnum{14.} \emph{退化椭圆问题.}

考虑形如
\begin{equation*}
\begin{cases}
\displaystyle-\sum_{i,j}\frac{\partial}{\partial x_j}
\left(a_{ij}\frac{\partial u}{\partial x_i}\right)
+\sum_ia_i\frac{\partial u}{\partial x_i}+a_0u=f
& \text{在}\ \Omega\ \text{中},\\[2mm]
\text{边界条件给定在}\ \Gamma\ \text{上, 或给定在}\ \Gamma\ \text{的一部分上}
\end{cases}
\end{equation*}
的问题, 其中函数 $a_{ij}$ 不满足椭圆性条件 \eqref{eq:9.36}, 而只满足
\begin{equation}
\sum_{i,j}a_{ij}(x)\xi_i\xi_j\ge0\quad\forall x\in\Omega,\ \forall\xi\in\R^N.
\tag*{(36$'$)}
\end{equation}
参见例如 J.~Kohn--L.~Nirenberg \bref{KohnNirenberg}{1},
M.~S.~Baouendi--C.~Goulaouic \bref{BaouendiGoulaouic}{1},
O.~Oleinik--E.~Radkevitch \bref{OleinikRadkevitch}{1} 的著作.

\cmtnum{15.} \emph{非线性椭圆问题.}

这是一个由几何, 力学, 物理学, 最优控制, 概率论等中数不清的问题所激发的
庞大研究领域.自 20 世纪 30 年代初 Leray 与 Schauder 的开创性工作以来,
它有了惊人的发展.我们区分几个类别:

\cmtsub{(a) 半线性问题.} 例如, 这包括形如
\begin{equation}
\begin{cases}
-\Delta u=f(x,u) & \text{在}\ \Omega\ \text{中},\\
u=0 & \text{在}\ \Gamma\ \text{上}
\end{cases}
\label{eq:9.92}
\end{equation}
的问题, 其中 $f(x,u)$ 是给定的函数.

这一类中包括所谓的\emph{分支问题}, 其中人们研究问题
\begin{equation*}
\begin{cases}
-\Delta u=f_\lambda(x,u) & \text{在}\ \Omega\ \text{中},\\
u=0 & \text{在}\ \Gamma\ \text{上}
\end{cases}
\tag*{(91$'$)}
\end{equation*}
的解集 $(\lambda,u)$ 的结构, 其中 $\lambda$ 是变化的参数.

\cmtsub{(b) 拟线性问题.} 考虑形如
\begin{equation}
\begin{cases}
\displaystyle-\sum_{i,j}\frac{\partial}{\partial x_j}
\left(a_{ij}(x,u,\nabla u)\frac{\partial u}{\partial x_i}\right)
=f(x,u,\nabla u) & \text{在}\ \Omega\ \text{中},\\[2mm]
u=0 & \text{在}\ \Gamma\ \text{上}
\end{cases}
\label{eq:9.93}
\end{equation}
的问题, 其中函数 $a_{ij}(x,u,p)$ 是椭圆的, 但可能是退化的; 例如我们有
\begin{equation*}
\sum_{i,j}a_{ij}(x,u,p)\xi_i\xi_j\ge\alpha(u,p)\abs{\xi}^2
\quad\forall x\in\Omega,\ \forall\xi\in\R^N,\ \forall u\in\R,\
\forall p\in\R^N,
\end{equation*}
其中 $\alpha(u,p)>0$ $\forall u\in\R$, $\forall p\in\R^N$, 但
$\alpha(u,p)$ 不能被一个常数 $\alpha>0$ 一致地控制在下.特别地, 著名的
\emph{极小曲面方程}属于这一类, 其中
$a_{ij}(x,u,p)=\delta_{ij}(1+\abs{p}^2)^{-1/2}$.更一般地, 人们考虑形如
\begin{equation}
F(x,u,Du,D^2u)=0
\label{eq:9.94}
\end{equation}
的\emph{完全非线性}椭圆问题, 其中矩阵
$\frac{\partial F}{\partial q_{ij}}(x,u,p,q)$ 是椭圆的 (可能是退化的).
例如, \emph{Monge--Amp\`ere 方程}属于这一类.

\cmtsub{(c) 自由边界问题.} 问题是要在一个\emph{事先未给定}的开集 $\Omega$
中求解一个线性椭圆方程.$\Omega$ 未知这一事实常常通过在 $\Gamma$ 上给定
\emph{两个}边界条件 (例如 Dirichlet 条件与 Neumann 条件) 来``补偿''.
问题在于同时找出开集 $\Omega$ 与函数 $u$, 使得 $\ldots$.

\runin{技巧:}

\noindent (a) 对问题 \eqref{eq:9.92} 或 (92$'$) 有几种技巧可用:

\begin{itemize}
	\item \emph{单调性方法}, 见 F.~Browder \bref{Browder}{1} 与
	      J.-L.~Lions \bref{Lions}{3}.
	\item \emph{拓扑方法} (Schauder 不动点定理, Leray--Schauder 度理论
	      等); 见 J.~T.~Schwartz \bref{Schwartz}{1},
	      M.~Krasnoselskii \bref{Krasnoselskii}{1}, 以及
	      L.~Nirenberg \bref{Nirenberg}{2}, \bref{Nirenberg}{3}.
	\item \emph{变分方法} (临界点理论, 极小极大技巧, Morse 理论等); 见
	      P.~Rabinowitz \bref{Rabinowitz}{1}, \bref{Rabinowitz}{2},
	      Melvyn Berger \bref{Berger}{1},
	      M.~Krasnoselskii \bref{Krasnoselskii}{1},
	      L.~Nirenberg \bref{Nirenberg}{3},
	      J.~Mawhin--M.~Willem \bref{MawhinWillem}{1},
	      M.~Willem \bref{Willem}{1}, M.~Struwe \bref{Struwe}{1}.
\end{itemize}

一般性的综述参见例如 A.~Ambrosetti--G.~Prodi \bref{AmbrosettiProdi}{1}
与 E.~Zeidler \bref{Zeidler}{1} 的书.

\noindent (b) 求解 \eqref{eq:9.93} 型问题可能要涉及精细的估计技巧;\footnote{
例如极小曲面方程就是这种情形.} 见 E.~De Giorgi,
O.~Ladyzhenskaya--N.~Uraltseva \bref{Ladyzhenskaya}{1},
J.~Serrin \bref{Serrin}{1}, E.~Bombieri \bref{Bombieri}{1} 以及
D.~Gilbarg--N.~Trudinger \bref{GilbargTrudinger}{1} 的著作.关于完全非线性
方程特别是 Monge--Amp\`ere 方程, 近来也取得了重要进展; 参见例如
S.~T.~Yau \bref{Yau}{1}, L.~Caffarelli--L.~Nirenberg--J.~Spruck
\bref{CaffarelliNirenbergSpruck}{1} 以及 X.~Cabr\'e--L.~Caffarelli
\bref{CabreCaffarelli}{1}.

\noindent (c) 关于自由边界问题, 近年来出现了许多新结果, 它们常常与变分
不等式理论相联系; 参见例如 D.~Kinderlehrer--G.~Stampacchia
\bref{KinderlehrerStampacchia}{1}, C.~Baiocchi--A.~Capelo
\bref{BaiocchiCapelo}{1}, A.~Friedman \bref{Friedman}{4},
J.~Crank \bref{Crank}{1} 以及 L.~Caffarelli--S.~Salsa
\bref{CaffarelliSalsa}{1}.

\cmtnum{16.} \emph{几何测度论.}

这个领域处于几何与偏微分方程的交叉处, 自 20 世纪 60 年代以来有了广泛的
发展, 其起点是 H.~Federer, E.~De Giorgi, A.~I.~Volpert 与 F.~Almgren 的
基本贡献, 它们与变分法, 等周不等式等中产生的问题相联系.它在物理问题中有
大量应用, 例如相变, 力学中的断裂, 图像处理中的边缘检测, 液晶中的线涡, 超导体与超流体.空间 $BV$ (有界变差函数) 在这些问题中起特殊的作用.我们
参见例如 L.~Ambrosio--N.~Fusco--D.~Pallara \bref{Ambrosio}{1},
L.~Simon \bref{Simon}{1}, L.~C.~Evans--R.~Gariepy \bref{EvansGariepy}{1}
以及 F.~H.~Lin--X.~P.~Yang \bref{LinYang}{1}.

% chapters/ch10.tex
% 第~
% \ref{chap:10}~章 发展方程: 热方程与波动方程
% 源: 原书 PDF 第 329--351 页 (印刷页 325--347)

\chapter{发展方程: 热方程与波动方程}
\label{chap:10}

\section{热方程: 存在性, 唯一性与正则性}
\label{sec:10.1}

\begin{notation}
设 $\Omega\subset\R^N$ 为边界为 $\Gamma$ 的开集.令
\begin{equation*}
Q=\Omega\times(0,+\infty),\qquad
\Sigma=\Gamma\times(0,+\infty);
\end{equation*}
$\Sigma$ 称为柱体 $Q$ 的\emph{侧边界}.见图  \ref{fig:7}.
\end{notation}

考虑下述问题: 求函数 $u(x,t):\Omega\times[0,+\infty)\to\R$ 使得
\begin{equation}
\frac{\partial u}{\partial t}-\Delta u=0\quad\text{在}\ Q\ \text{中},
\label{eq:10.1}
\end{equation}
\begin{equation}
u=0\quad\text{在}\ \Sigma\ \text{上},
\label{eq:10.2}
\end{equation}
\begin{equation}
u(x,0)=u_0(x)\quad\text{在}\ \Omega\ \text{上},
\label{eq:10.3}
\end{equation}

\begin{figure}[tbp]
	\centering
	\begin{tikzpicture}[scale=0.95]
		\draw[-{Stealth}] (-0.5,0) -- (5.2,0) node[below left] {$x$};
		\draw[-{Stealth}] (0,-0.5) -- (0,4.4) node[above right] {$t$};
		\draw (1.4,0) -- (1.4,3.6);
		\draw (3.6,0) -- (3.6,3.6);
		\draw (1.4,3.6) .. controls (2.1,4.3) and (2.9,4.3) .. (3.6,3.6);
		\draw[dashed] (1.4,3.6) .. controls (2.1,2.9) and (2.9,2.9) .. (3.6,3.6);
		\draw (1.4,0.35) .. controls (2.1,1.05) and (2.9,1.05) .. (3.6,0.35);
		\draw[dashed] (1.4,0.35) .. controls (2.1,-0.35) and (2.9,-0.35) .. (3.6,0.35);
		\draw (1.4,0) -- (1.4,-0.12);
		\draw (3.6,0) -- (3.6,-0.12);
		\node at (2.5,2.4) {$Q$};
		\node at (2.5,-0.35) {$\Omega$};
		\node[right] at (3.62,2.4) {$\Sigma$};
	\end{tikzpicture}
	\caption{}
	\label{fig:7}
\end{figure}

其中 $\Delta=\sum_{i=1}^N\dfrac{\partial^2}{\partial x_i^2}$ 表示关于空间变量
$x$ 的 Laplace 算子, $t$ 是时间变量, $u_0(x)$ 是给定的函数, 称为
\emph{初值} (或 \emph{Cauchy 数据}).

\index{refangcheng@热方程 (heat equation)}
方程 \eqref{eq:10.1} 称为\emph{热方程}, 因为它刻画区域 $\Omega$ 中时刻 $t$
的温度分布 $u$.热方程及其各种变形出现在许多扩散现象中\footnote{热的扩散
只是众多例子之一.} (见本章末评注).热方程是\emph{抛物型}方程的最简单
例子.\footnote{关于偏微分方程分为``椭圆型'', ``抛物型'', ``双曲型''三类
的传统分类, 参见例如 R.~Courant--D.~Hilbert \bref{CourantHilbert}{1}.}

方程 \eqref{eq:10.2} 是 (齐次) \emph{Dirichlet 边界条件}; 可以把它换成
Neumann 条件
\begin{equation*}
\frac{\partial u}{\partial n}=0\quad\text{在}\ \Sigma\ \text{上}
\tag*{(2$'$)}
\end{equation*}
($n$ 是 $\Gamma$ 的外单位法向量), 或第~
\ref{chap:8}~章与第~
\ref{chap:9}~章中遇到的任何边界条件.
条件 \eqref{eq:10.2} 对应于假设边界 $\Gamma$ 保持零温度; 条件 $(2')$
对应于假设通过 $\Gamma$ 的热流为零.我们通过把 $u(x,t)$ 看作定义在
$[0,+\infty)$ 上, 取值于空间 $H$ 的函数来求解问题 \eqref{eq:10.1}--
\eqref{eq:10.3}, 其中 $H$ 是只依赖于 $x$ 的函数空间: 例如 $H=L^2(\Omega)$,
或 $H=H_0^1(\Omega)$.当我们只写 $u(t)$ 时, 意思是 $u(t)$ 是 $H$ 的元素,
即函数 $x\mapsto u(x,t)$.这种观点使我们能够把 Hille--Yosida 定理与第 8, 9
章的结果结合起来, 非常容易地求解问题 \eqref{eq:10.1}--\eqref{eq:10.3}.

为简化起见, 我们在整个第~
\ref{chap:10}~章中假设 $\Omega$ 属于 $C^\infty$ 类且
$\Gamma$ 有界 (但若只关心弱解, 这个假设可以大大减弱).

\begin{keytheorem}
\label{thm:10.1}
假设 $u_0\in L^2(\Omega)$.则存在唯一的函数 $u(x,t)$ 满足
\eqref{eq:10.1}, \eqref{eq:10.2}, \eqref{eq:10.3} 以及
\begin{equation}
u\in C([0,\infty);L^2(\Omega))\cap C((0,\infty);H^2(\Omega)\cap
H_0^1(\Omega)),
\label{eq:10.4}
\end{equation}
\begin{equation}
u\in C^1((0,\infty);L^2(\Omega)).
\label{eq:10.5}
\end{equation}
此外,
\begin{equation*}
u\in C^\infty(\overline\Omega\times[\varepsilon,\infty))
\quad\forall\varepsilon>0.
\end{equation*}
最后, $u\in L^2(0,\infty;H_0^1(\Omega))$ 且\footnote{依照上面的讨论, 我们
采用下述记号:
$\abs{u(T)}_{L^2(\Omega)}=\int_\Omega\abs{u(x,T)}^2dx$ 以及
$\abs{\nabla u(t)}^2_{L^2(\Omega)}
=\sum_{i=1}^N\int_\Omega\abs{\dfrac{\partial u}{\partial x_i}(x,t)}^2dx$.}
\begin{equation}
\frac12\abs{u(T)}_{L^2(\Omega)}^2
+\int_0^T\abs{\nabla u(t)}^2_{L^2(\Omega)}dt
=\frac12\abs{u_0}_{L^2(\Omega)}^2\quad\forall T>0.
\label{eq:10.6}
\end{equation}
\end{keytheorem}

\begin{proof}
我们在 $H=L^2(\Omega)$ 中应用 Hille--Yosida 理论 (但也可以作其他选择; 见
定理  \ref{thm:10.2} 的证明).考虑由
\begin{equation*}
D(A)=H^2(\Omega)\cap H_0^1(\Omega),\qquad Au=-\Delta u
\end{equation*}
定义的无界算子 $A:D(A)\subset H\to H$.

注意, 边界条件 \eqref{eq:10.2} 已被纳入 $A$ 的定义域中, 这一点很重要.
我们断言 $A$ 是自伴极大单调算子.于是可以应用定理  \ref{thm:7.7}, 推出
\eqref{eq:10.1}--\eqref{eq:10.3} 存在唯一的满足 \eqref{eq:10.4} 与
\eqref{eq:10.5} 的解.

(i) \emph{$A$ 单调.} 对每个 $u\in D(A)$, 有
\begin{equation*}
(Au,u)_{L^2}=\int_\Omega(-\Delta u)u=\int_\Omega\abs{\nabla u}^2\ge0.
\end{equation*}

(ii) \emph{$A$ 极大单调.} 须验证 $R(I+A)=H=L^2$.但我们已知 (见定理
 \ref{thm:9.25}): 对每个 $f\in L^2$, 方程 $u-\Delta u=f$ 存在唯一的解
$u\in H^2\cap H_0^1$.

(iii) \emph{$A$ 自伴.} 由命题  \ref{prop:7.6}, 只须验证 $A$ 对称.对每个
$u,v\in D(A)$, 有
\begin{equation*}
(Au,v)_{L^2}=\int_\Omega(-\Delta u)v=\int_\Omega\nabla u\cdot\nabla v
\end{equation*}
以及
\begin{equation*}
(u,Av)_{L^2}=\int_\Omega u(-\Delta v)=\int_\Omega\nabla u\cdot\nabla v,
\end{equation*}
从而 $(Au,v)=(u,Av)$.

其次, 由定理  \ref{thm:9.25} 推出: 对每个整数 $\ell$, 有
$D(A^\ell)\subset H^{2\ell}(\Omega)$, 且嵌入连续.更精确地说,
\begin{equation*}
D(A^\ell)=\{u\in H^{2\ell}(\Omega);\ u=\Delta u=\dots=\Delta^{\ell-1}u=0
\ \text{在}\ \Gamma\ \text{上}\}.
\end{equation*}
由定理  \ref{thm:7.7} 我们知道, \eqref{eq:10.1}--\eqref{eq:10.3} 的解 $u$
满足
\begin{equation*}
u\in C^k((0,\infty);D(A^\ell))\quad\forall k,\ \forall\ell,
\end{equation*}
从而
\begin{equation*}
u\in C^k((0,\infty);H^{2\ell}(\Omega))\quad\forall k,\ \forall\ell.
\end{equation*}
由此推出 (借助推论  \ref{cor:9.15})
\begin{equation*}
u\in C^k((0,\infty);C^k(\overline\Omega))\quad\forall k.
\end{equation*}
现在我们转向 \eqref{eq:10.6} 的证明.形式上, 我们把 \eqref{eq:10.1} 乘以
$u$ 并在 $\Omega\times(0,T)$ 上积分.但必须小心, 因为 $u(t)$ 在 $(0,\infty)$
上可微, 而在 $[0,\infty)$ 上不可微.考虑函数
$\varphi(t)=\dfrac12\abs{u(t)}^2_{L^2(\Omega)}$.它在 $(0,\infty)$ 上属于
$C^1$ 类 (由 \eqref{eq:10.5}), 且对 $t>0$,
\begin{equation*}
\varphi'(t)=\left(u(t),\frac{du}{dt}(t)\right)_{L^2}
=(u(t),\Delta u(t))_{L^2}=-\int_\Omega\abs{\nabla u(t)}^2.
\end{equation*}
因此, 对 $0<\varepsilon<T<\infty$, 我们得到
\begin{equation*}
\varphi(T)-\varphi(\varepsilon)
=\int_\varepsilon^T\varphi'(t)dt
=-\int_\varepsilon^T\abs{\nabla u(t)}^2_{L^2}dt.
\end{equation*}
最后令 $\varepsilon\to0$.由于 $\varphi(\varepsilon)\to\frac12\abs{u_0}^2$
(因为 $u\in C([0,\infty];L^2(\Omega))$), 我们得到
$u\in L^2(0,\infty;H_0^1(\Omega))$, 且 \eqref{eq:10.6} 成立.

若对 $u_0$ 作附加假设, 则解 $u$ 直到 $t=0$ 都会变得更正则 (回忆: 在
$t=0$ 之外, 定理  \ref{thm:10.1} 总保证 $u$ 光滑, 即
$u\in C^\infty(\overline\Omega\times[\varepsilon,\infty))\ \forall\varepsilon>0$).
\end{proof}

\begin{theorem}
\label{thm:10.2}
(a) 若 $u_0\in H_0^1(\Omega)$, 则 \eqref{eq:10.1}--\eqref{eq:10.3} 的解
$u$ 满足
\begin{equation*}
u\in C([0,\infty);H_0^1(\Omega))\cap L^2(0,\infty;H^2(\Omega))
\end{equation*}
且
\begin{equation*}
\frac{\partial u}{\partial t}\in L^2(0,\infty;L^2(\Omega)).
\end{equation*}
此外, 我们有
\begin{equation}
\int_0^T\left\|\frac{\partial u}{\partial t}(t)\right\|^2_{L^2(\Omega)}dt
+\frac12\abs{\nabla u(T)}^2_{L^2(\Omega)}
=\frac12\abs{\nabla u_0}^2_{L^2(\Omega)}.
\label{eq:10.7}
\end{equation}

(b) 若 $u_0\in H^2(\Omega)\cap H_0^1(\Omega)$, 则
\begin{equation*}
u\in C([0,\infty);H^2(\Omega))\cap L^2(0,\infty;H^3(\Omega))
\end{equation*}
且
\begin{equation*}
\frac{\partial u}{\partial t}\in L^2(0,\infty;H_0^1(\Omega)).
\end{equation*}

(c) 若 $u_0\in H^k(\Omega)\ \forall k$ 且满足所谓的\emph{相容性条件}
\begin{equation}
u_0=\Delta u_0=\dots=\Delta^ju_0=\dots=0\quad\text{在}\ \Gamma\ \text{上}
\label{eq:10.8}
\end{equation}
(对每个整数 $j$), 则 $u\in C^\infty(\overline\Omega\times[0,\infty))$.
\end{theorem}

\begin{proof}[(a)]
这里我们在空间 $H_1=H_0^1(\Omega)$ 中工作, 它配以标量积
\begin{equation*}
(u,v)_{H_1}=\int_\Omega\nabla u\cdot\nabla v+\int_\Omega uv.
\end{equation*}
在 $H_1$ 中考虑由
\begin{equation*}
D(A_1)=\{u\in H^3(\Omega)\cap H_0^1(\Omega);\ \Delta u\in H_0^1(\Omega)\},
\qquad A_1u=-\Delta u
\end{equation*}
定义的无界算子 $A_1:D(A_1)\subset H_1\to H_1$.

我们断言 $A_1$ 极大单调且自伴.

(i) \emph{$A_1$ 单调.} 对每个 $u\in D(A_1)$, 有
\begin{equation*}
(A_1u,u)_{H_1}=\int_\Omega\nabla(-\Delta u)\cdot\nabla u
+\int_\Omega(-\Delta u)u
=\int_\Omega\abs{\Delta u}^2+\int_\Omega\abs{\nabla u}^2\ge0.
\end{equation*}

(ii) \emph{$A_1$ 极大单调.} 由定理  \ref{thm:9.25} 我们知道: 对每个
$f\in H^1(\Omega)$, 问题
\begin{equation*}
\begin{cases}
u-\Delta u=f & \text{在}\ \Omega\ \text{中},\\
u=0 & \text{在}\ \Gamma\ \text{上}
\end{cases}
\end{equation*}
的解 $u\in H_0^1(\Omega)$ 属于 $H^3(\Omega)$.若此外 $f\in H_0^1(\Omega)$,
则 $\Delta u\in H_0^1(\Omega)$, 从而 $u\in D(A_1)$.

(iii) \emph{$A_1$ 对称.} 对每个 $u,v\in D(A_1)$, 有
\begin{align*}
(A_1u,v)_{H_1}
&=\int_\Omega\nabla(-\Delta u)\cdot\nabla v+\int_\Omega(-\Delta u)v\\
&=\int_\Omega\Delta u\,\Delta v+\int_\Omega\nabla u\cdot\nabla v
=(u,A_1v)_{H_1}.
\end{align*}

应用定理  \ref{thm:7.7} 可见: 若 $u_0\in H_0^1(\Omega)$, 则存在
\eqref{eq:10.1}--\eqref{eq:10.3} 的解 $u$ (由唯一性, 它与定理
 \ref{thm:10.1} 中得到的解相同), 使得
\begin{equation*}
u\in C([0,\infty);H_0^1(\Omega)).
\end{equation*}
最后, 令 $\varphi(t)=\dfrac12\abs{\nabla u(t)}^2_{L^2(\Omega)}$.这个函数在
$(0,\infty)$ 上属于 $C^\infty$ 类, 且
\begin{equation*}
\varphi'(t)=\left(\nabla u(t),\nabla\frac{du}{dt}(t)\right)_{L^2}
=\left(-\Delta u(t),\frac{du}{dt}(t)\right)_{L^2}
=-\left\|\frac{du}{dt}(t)\right\|^2_{L^2}.
\end{equation*}
由此推出, 对 $0<\varepsilon<T<\infty$, 有
\begin{equation*}
\varphi(T)-\varphi(\varepsilon)
+\int_\varepsilon^T\left\|\frac{du}{dt}(t)\right\|^2_{L^2}dt=0.
\end{equation*}
当 $\varepsilon\to0$ 时,
$\varphi(\varepsilon)\to\frac12\abs{\nabla u_0}^2_{L^2}$, 于是容易得到结论.
\end{proof}

\begin{proof}[(b)]
这里我们在空间 $H_2=H^2(\Omega)\cap H_0^1(\Omega)$ 中工作, 它配以标量积
\begin{equation*}
(u,v)_{H_2}=(\Delta u,\Delta v)_{L^2}+(u,v)_{L^2}
\end{equation*}
(相应的范数与通常的 $H^2$ 范数等价; 为什么?).在 $H_2$ 中考虑由
\begin{equation*}
D(A_2)=\{u\in H^4(\Omega);\ u\in H_0^1(\Omega)\ \text{且}\
\Delta u\in H_0^1(\Omega)\},\qquad A_2u=-\Delta u
\end{equation*}
定义的无界算子 $A_2:D(A_2)\subset H_2\to H_2$.

容易证明 $A_2$ 是 $H_2$ 中的自伴极大单调算子.\footnote{更一般地, 若
$A:D(A)\subset H\to H$ 是自伴极大单调算子, 则可以考虑 Hilbert 空间
$\tilde H=D(A)$, 它配以标量积
$(u,v)_{\tilde H}=(Au,Av)+(u,v)$.此时由 $D(\tilde A)=D(A^2)$ 与
$\tilde Au=Au$ 定义的算子 $\tilde A:D(\tilde A)\subset\tilde H\to\tilde H$
是 $\tilde H$ 中的自伴极大单调算子.} 因此我们可以对 $H_2$ 中的 $A_2$
应用定理  \ref{thm:7.7}.最后, 令
$\varphi(t)=\dfrac12\abs{\Delta u(t)}^2_{L^2}$.这个函数在 $(0,\infty)$ 上
属于 $C^\infty$ 类, 且
\begin{equation*}
\varphi'(t)=\left(\Delta u(t),\Delta\frac{du}{dt}(t)\right)_{L^2}
=(\Delta u(t),\Delta^2u(t))_{L^2}=-\abs{\nabla\Delta u(t)}^2_{L^2}.
\end{equation*}
于是, 对 $0<\varepsilon<T<\infty$, 有
\begin{equation*}
\frac12\abs{\Delta u(T)}^2_{L^2}-\frac12\abs{\Delta u(\varepsilon)}^2_{L^2}
+\int_\varepsilon^T\abs{\nabla\Delta u(t)}^2_{L^2}dt=0.
\end{equation*}
取极限 $\varepsilon\to0$, 可见 $u\in L^2(0,\infty;H^3(\Omega))$ (为什么?),
且 (由方程 \eqref{eq:10.1}) $\dfrac{du}{dt}\in L^2(0,\infty;H^1(\Omega))$.
\end{proof}

\begin{proof}[(c)]
在空间 $H=L^2(\Omega)$ 中, 考虑由
\begin{equation*}
D(A)=H^2(\Omega)\cap H_0^1(\Omega),\qquad Au=-\Delta u
\end{equation*}
定义的算子 $A:D(A)\subset H\to H$.

应用定理  \ref{thm:7.5}, 我们知道: 若 $u_0\in D(A^k)$, $k\ge1$, 则
\begin{equation*}
u\in C^{k-j}([0,\infty);D(A^j))\quad\forall j=0,1,\dots,k.
\end{equation*}
假设 \eqref{eq:10.8} 恰好说明 $u_0\in D(A^k)$ 对每个整数 $k\ge1$ 成立.
因此我们有
\begin{equation*}
u\in C^{k-j}([0,\infty);D(A^j))\quad\forall k\ge1,\
\forall j=0,1,\dots,k.
\end{equation*}
由此推出 (如同定理  \ref{thm:10.1} 的证明)
$u\in C^\infty(\overline\Omega\times[0,\infty))$.
\end{proof}

\begin{keyremark}
\label{rem:10.1}
定理  \ref{thm:10.1} 表明热方程对初值 $u_0$ 有很强的\emph{光滑化效应}.
注意, 即使初值不连续, 解 $u(x,t)$ 对每个 $t>0$ 关于 $x$ 也是 $C^\infty$ 的.
特别地, 这个效应意味着热方程是\emph{时间不可逆}的.一般说来, 人们不能
求解带有``终值''数据
\begin{equation}
u(x,T)=u_T(x)\quad\text{在}\ \Omega\ \text{上}
\label{eq:10.11}
\end{equation}
的反向问题
\begin{equation}
\frac{\partial u}{\partial t}-\Delta u=0\quad\text{在}\ \Omega\times(0,T)\ \text{中},
\label{eq:10.9}
\end{equation}
\begin{equation}
u=0\quad\text{在}\ \Gamma\times(0,T)\ \text{上},
\label{eq:10.10}
\end{equation}
我们必然要假设
\begin{equation*}
u_T\in C^\infty(\overline\Omega)
\quad\text{且}\quad
\Delta^ju_T=0\ \text{在}\ \Gamma\ \text{上}\ \forall j\ge0.
\end{equation*}
但即使作了这个假设, 反向问题 \eqref{eq:10.9}, \eqref{eq:10.10}, \eqref{eq:10.11} 也未必有解.不应把这个与问题 $(9')$, \eqref{eq:10.10}, \eqref{eq:10.11} 混淆, 其中
\begin{equation*}
-\frac{\partial u}{\partial t}-\Delta u=0
\quad\text{在}\ \Omega\times(0,T)\ \text{中},
\tag*{(9$'$)}
\end{equation*}
它对任何数据 $u_T\in L^2(\Omega)$ 总有唯一解 (把 $t$ 换成 $T-t$ 并应用
定理  \ref{thm:10.1}).
\end{keyremark}

\begin{remark}
\label{rem:10.2}
若把 Dirichlet 条件换成 Neumann 条件, 前面的结果 (作一些小的修改后) 仍然
成立.
\end{remark}

\begin{remark}
\label{rem:10.3}
当 $\Omega$ 有界时, 问题 \eqref{eq:10.1}--\eqref{eq:10.3} 也可以通过
在 $L^2(\Omega)$ 的一个 Hilbert 基中分解来求解.为此, 非常方便的是选取
$L^2(\Omega)$ 的一个由 $-\Delta$ 的特征函数 (带零 Dirichlet 条件) 组成的基
$(e_i(x))_{i\ge1}$, 即
\begin{equation*}
\begin{cases}
-\Delta e_i=\lambda_ie_i & \text{在}\ \Omega\ \text{中},\\
e_i=0 & \text{在}\ \Gamma\ \text{上}
\end{cases}
\end{equation*}
(见 \S \ref{sec:9.8}).我们以级数\footnote{由于显而易见的原因, 这个方法也
称为``分离变量法''或 Fourier 方法.事实上, Fourier 正是在研究一个空间变量
的热方程时发现了 Fourier 级数.}
\begin{equation}
u(x,t)=\sum_{i=1}^\infty a_i(t)e_i(x)
\label{eq:10.12}
\end{equation}
的形式寻求 \eqref{eq:10.1}--\eqref{eq:10.3} 的解.我们立即看出, 函数
$a_i(t)$ 必须满足
\begin{equation*}
a_i'(t)+\lambda_ia_i(t)=0,
\end{equation*}
从而 $a_i(t)=a_i(0)e^{-\lambda_it}$.常数 $a_i(0)$ 由关系式
\begin{equation}
u_0(x)=\sum_{i=1}^\infty a_i(0)e_i(x)
\label{eq:10.13}
\end{equation}
确定.换言之, \eqref{eq:10.1}--\eqref{eq:10.3} 的解 $u$ 由
\begin{equation}
u(x,t)=\sum_{i=1}^\infty a_i(0)e^{-\lambda_it}e_i(x)
\label{eq:10.14}
\end{equation}
给出, 其中常数 $a_i(0)$ 是 $u_0(x)$ 在基 $(e_i)$ 中的分量, 即
$a_i(0)=\int_\Omega u_0e_i$.

关于这个级数的收敛性 (以及由此得到的 $u$ 的正则性), 我们参见
H.~Weinberger \bref{Weinberger}{1}.注意这个方法与求解线性微分方程组
\begin{equation*}
\frac{du}{dt}+Mu=0
\end{equation*}
时所用的标准技巧之间的类似之处, 其中 $u(t)$ 取值于一个有限维向量空间,
$M$ 是对称矩阵.当然, 主要差别来自问题 \eqref{eq:10.1}--\eqref{eq:10.3}
对应的是一个无穷维方程组这一事实.
\end{remark}

\begin{remark}
\label{rem:10.4}
\index{xiangrongxingtiaojian@相容性条件 (compatibility conditions)}
相容性条件 \eqref{eq:10.8} 看上去也许有些神秘, 但事实上它们是自然的.
它们是使 \eqref{eq:10.1}--\eqref{eq:10.3} 有直到 $t=0$ 都光滑的解, 即
$u\in C^\infty(\overline\Omega\times[0,\infty))$ 的\emph{必要}条件 (假设
$u_0\in C^\infty(\overline\Omega)$ 且 $u_0=0$ 在 $\partial\Omega$ 上, 并不能
保证直到 $t=0$ 的光滑性).事实上, 设 $u\in C^\infty(\overline\Omega\times
[0,\infty))$ 满足 \eqref{eq:10.1}--\eqref{eq:10.3}.则显然
\begin{equation}
\frac{\partial^ju}{\partial t^j}=0
\quad\text{在}\ \Gamma\times(0,\infty)\ \text{上}\quad\forall j,
\label{eq:10.15}
\end{equation}
由连续性, 我们还有
\begin{equation*}
\frac{\partial^ju}{\partial t^j}=0
\quad\text{在}\ \Gamma\times[0,\infty)\ \text{上}\quad\forall j.
\end{equation*}
另一方面,
\begin{equation*}
\frac{\partial^2u}{\partial t^2}
=\Delta\left(\frac{\partial u}{\partial t}\right)=\Delta^2u
\quad\text{在}\ Q\ \text{中},
\end{equation*}
从而由归纳法,
\begin{equation*}
\frac{\partial^ju}{\partial t^j}=\Delta^ju\quad\text{在}\ Q\ \text{中}
\quad\forall j.
\end{equation*}
再次由连续性, 我们有
\begin{equation}
\frac{\partial^ju}{\partial t^j}=\Delta^ju
\quad\text{在}\ \Omega\times[0,\infty)\ \text{中}.
\label{eq:10.16}
\end{equation}
在 $\Gamma\times\{0\}$ 上比较 \eqref{eq:10.15} 与 \eqref{eq:10.16}, 我们
得到 \eqref{eq:10.8}.
\end{remark}

\begin{remark}
\label{rem:10.5}
当然, 若我们在定理  \ref{thm:10.2} 的 (b) 与 (c) 两种情形之间作中间假设,
则 $u$ 在 $t=0$ 附近的正则性结果有许多变体.
\end{remark}
\section{极大值原理}
\label{sec:10.2}

主要结果如下.

\begin{keytheorem}
\label{thm:10.3}
假设 $u_0\in L^2(\Omega)$, 设 $u$ 为 \eqref{eq:10.1}--\eqref{eq:10.3} 的解.
则对所有 $(x,t)\in Q$, 有
\begin{equation*}
\min\left\{0,\inf_\Omega u_0\right\}\le u(x,t)\le\max\left\{0,
\sup_\Omega u_0\right\}.
\end{equation*}
\end{keytheorem}

\begin{proof}
与椭圆情形一样, 我们使用 Stampacchia 的截断方法.令
\begin{equation*}
K=\max\left\{0,\sup_\Omega u_0\right\}
\end{equation*}
并假设 $K<+\infty$.固定一个如定理  \ref{thm:9.27} 的证明中的函数 $G$,
并令
\begin{equation*}
H(s)=\int_0^sG(\sigma)d\sigma,\quad s\in\R.
\end{equation*}
容易验证, 由
\begin{equation*}
\varphi(t)=\int_\Omega H(u(x,t)-K)dx
\end{equation*}
定义的函数 $\varphi$ 具有下述性质:
\begin{equation}
\varphi\in C([0,\infty);\R),\quad\varphi(0)=0,\quad
\varphi\ge0\ \text{在}\ [0,\infty)\ \text{上},
\label{eq:10.17}
\end{equation}
\begin{equation}
\varphi\in C^1((0,\infty);\R),
\label{eq:10.18}
\end{equation}
且
\begin{align*}
\varphi'(t)&=\int_\Omega G(u(x,t)-K)\frac{\partial u}{\partial t}(x,t)dx
=\int_\Omega G(u(x,t)-K)\Delta u(x,t)dx\\
&=-\int_\Omega G'(u-K)\abs{\nabla u}^2dx\le0,
\end{align*}
因为 $G(u(x,t)-K)\in H_0^1(\Omega)$ 对每个 $t>0$ 成立.由此推出
$\varphi\equiv0$, 从而对每个 $t>0$, $u(x,t)\le K$ 在 $\Omega$ 上几乎
处处成立.
\end{proof}

\begin{corollary}
\label{cor:10.4}
设 $u_0\in L^2(\Omega)$.\eqref{eq:10.1}--\eqref{eq:10.3} 的解 $u$ 具有
下述性质:

(i) 若 $u_0\ge0$ 在 $\Omega$ 上几乎处处成立, 则 $u\ge0$ 在 $Q$ 中.

(ii) 若 $u_0\in L^\infty(\Omega)$, 则 $u\in L^\infty(Q)$ 且
\begin{equation}
\norm{u}_{L^\infty(Q)}\le\norm{u_0}_{L^\infty(\Omega)}.
\label{eq:10.19}
\end{equation}
\end{corollary}

\begin{corollary}
\label{cor:10.5}
设 $u_0\in C(\overline\Omega)\cap L^2(\Omega)$ 且 $u_0=0$ 在 $\Gamma$ 上.
\footnote{若 $\Omega$ 无界, 我们还假设当 $\abs{x}\to\infty$ 时
$u_0(x)\to0$.} 则 \eqref{eq:10.1}--\eqref{eq:10.3} 的解 $u$ 属于
$C(\overline Q)$.
\end{corollary}

\begin{proof}[推论 \ref{cor:10.5}]
设 $(u_{0n})$ 为 $C_c^\infty(\Omega)$ 中的函数序列, 使得
$u_{0n}\to u_0$ 在 $L^\infty(\Omega)$ 中且在 $L^2(\Omega)$ 中 (这样的序列
容易构造).由定理  \ref{thm:10.2}, 对应于初值 $u_{0n}$ 的
\eqref{eq:10.1}--\eqref{eq:10.3} 的解 $u_n$ 属于 $C^\infty(\overline Q)$.
另一方面 (定理  \ref{thm:7.7}), 我们知道
\begin{equation*}
\abs{u_n(t)-u(t)}_{L^2(\Omega)}\le\abs{u_{0n}-u_0}_{L^2(\Omega)}
\quad\forall t\ge0.
\end{equation*}
由 \eqref{eq:10.19} 我们有
\begin{equation*}
\norm{u_n-u_m}_{L^\infty(Q)}\le\norm{u_{0n}-u_{0m}}_{L^\infty(\Omega)}.
\end{equation*}
因此, 序列 $(u_n)$ 在 $Q$ 上一致收敛于 $u$, 从而 $u\in C(\overline Q)$.
\end{proof}

与椭圆情形一样, 极大值原理还有另一种途径.为简单起见, 这里假设
$\Omega$ 有界.设 $u(x,t)$ 是满足\footnote{注意, 我们没有给定任何边界条件或
任何初值.}
\begin{equation}
u\in C(\overline\Omega\times[0,T]),
\label{eq:10.20}
\end{equation}
\begin{equation}
u\ \text{在}\ \Omega\times(0,T)\ \text{中关于}\ t\ \text{属于}\ C^1\
\text{类, 关于}\ x\ \text{属于}\ C^2\ \text{类},
\label{eq:10.21}
\end{equation}
\begin{equation}
\frac{\partial u}{\partial t}-\Delta u\le0\quad\text{在}\
\Omega\times(0,T)\ \text{中}
\label{eq:10.22}
\end{equation}
的函数.

\begin{theorem}
\label{thm:10.6}
假设 \eqref{eq:10.20}, \eqref{eq:10.21} 与 \eqref{eq:10.22} 成立.则
\begin{equation}
\max_{\overline\Omega\times[0,T]}u=\max_Pu,
\label{eq:10.23}
\end{equation}
其中 $P=(\overline\Omega\times\{0\})\cup(\Gamma\times[0,T])$ 称为柱体
$\Omega\times(0,T)$ 的\emph{\index{paowubianjie@抛物边界 (parabolic boundary)}
抛物边界}.
\end{theorem}

\begin{proof}
令 $v(x,t)=u(x,t)+\varepsilon\abs{x}^2$, 其中 $\varepsilon>0$, 于是
\begin{equation}
\frac{\partial v}{\partial t}-\Delta v\le-2\varepsilon N<0
\quad\text{在}\ \Omega\times(0,T)\ \text{中}.
\label{eq:10.24}
\end{equation}
我们断言
\begin{equation*}
\max_{\overline\Omega\times[0,T]}v=\max_Pv.
\end{equation*}
若不然, 则存在点 $(x_0,t_0)\in\overline\Omega\times[0,T]$ 使得
$(x_0,t_0)\notin P$ 且
\begin{equation*}
\max_{\overline\Omega\times[0,T]}v=v(x_0,t_0).
\end{equation*}
由于 $x_0\in\Omega$ 且 $0<t_0\le T$, 我们有
\begin{equation}
\Delta v(x_0,t_0)\le0
\label{eq:10.25}
\end{equation}
以及
\begin{equation}
\frac{\partial v}{\partial t}(x_0,t_0)\ge0
\label{eq:10.26}
\end{equation}
(若 $t_0<T$, 则 $\dfrac{\partial v}{\partial t}(x_0,t_0)=0$; 若
$t_0=T$, 则 $\dfrac{\partial v}{\partial t}(x_0,t_0)\ge0$).\footnote{为
稳妥起见, 应当在 $\Omega\times(0,T')$ 中工作, 其中 $T'<T$, 然后令
$T'\to T$, 因为 $v$ 只在 $\Omega\times(0,T)$ 中关于 $t$ 属于 $C^1$ 类.}
结合 \eqref{eq:10.25} 与 \eqref{eq:10.26}, 我们得到
$\left(\dfrac{\partial v}{\partial t}-\Delta v\right)(x_0,t_0)\ge0$, 这与
\eqref{eq:10.24} 矛盾.因此我们有
\begin{equation*}
\max_{\overline\Omega\times[0,T]}v=\max_Pv\le\max_Pu+\varepsilon C,
\end{equation*}
其中 $C=\sup_{x\in\Omega}\abs{x}^2$.由于 $u\le v$, 我们推出
\begin{equation*}
\max_{\overline\Omega\times[0,T]}u\le\max_Pu+\varepsilon C
\quad\forall\varepsilon>0.
\end{equation*}
这就完成了 \eqref{eq:10.23} 的证明.
\end{proof}

\section{波动方程}
\label{sec:10.3}

设 $\Omega\subset\R^N$ 为开集.如上, 我们令
\begin{equation*}
Q=\Omega\times(0,\infty),\qquad\Sigma=\Gamma\times(0,\infty).
\end{equation*}
考虑下述问题: 求函数 $u(x,t):\Omega\times[0,\infty)\to\R$ 满足
\begin{equation}
\frac{\partial^2u}{\partial t^2}-\Delta u=0\quad\text{在}\ Q\ \text{中},
\label{eq:10.27}
\end{equation}
\begin{equation}
u=0\quad\text{在}\ \Sigma\ \text{上},
\label{eq:10.28}
\end{equation}
\begin{equation}
u(x,0)=u_0(x)\quad\text{在}\ \Omega\ \text{上},
\label{eq:10.29}
\end{equation}
\begin{equation}
\frac{\partial u}{\partial t}(x,0)=v_0(x)\quad\text{在}\ \Omega\ \text{上},
\label{eq:10.30}
\end{equation}
其中 $\Delta=\sum_{i=1}^N\dfrac{\partial^2}{\partial x_i^2}$ 表示关于空间
变量 $x$ 的 Laplace 算子, $t$ 是时间变量, $u_0,v_0$ 是给定的函数.

\index{bodongfangcheng@波动方程 (wave equation)}
方程 \eqref{eq:10.27} 称为\emph{波动方程}.算子
$\left(\dfrac{\partial^2}{\partial t^2}-\Delta\right)$ 常记作 $\square$,
\index{dAlembertsuanzi@d'Alembert 算子 (d'Alembertian)}
称为 \emph{d'Alembert 算子}.波动方程是\emph{双曲型}方程的典型例子.

当 $N=1$ 且 $\Omega=(0,1)$ 时, 方程 \eqref{eq:10.27} 刻画在没有外力作用时
弦的微小\footnote{完整的方程是一个非常困难的非线性方程; 方程
\eqref{eq:10.27} 是它在平衡位置附近的线性化版本.}振动.对每个 $t$, 函数
$x\in\Omega\mapsto u(x,t)$ 的图形表示弦在时刻 $t$ 的构形.当 $N=2$ 时,
方程 \eqref{eq:10.27} 刻画弹性膜的微小振动.对每个 $t$, 函数
$x\in\Omega\mapsto u(x,t)$ 的图形表示膜在时刻 $t$ 的构形.更一般地, 方程
\eqref{eq:10.27} 刻画波 (声波, 电磁波等) 在某个均匀弹性介质
$\Omega\subset\R^N$ 中的传播.

方程 \eqref{eq:10.28} 是 (齐次) Dirichlet 边界条件; 可以把它换成 Neumann
条件或第~
\ref{chap:8}~章或第~
\ref{chap:9}~章中遇到的任何边界条件.条件 $u=0$ 在 $\Sigma$ 上意味着
弦 (或膜) 在 $\Gamma$ 上固定, 而 Neumann 条件则说弦在其端点处自由.方程
\eqref{eq:10.29} 与 \eqref{eq:10.30} 表示系统的初始状态: 初始构形 (也说
初始位移) 由 $u_0$ 描述, 初速度由 $v_0$ 描述.数据 $(u_0,v_0)$ 通常称为
\emph{Cauchy 数据}.

为简化起见, 我们在本节中始终假设 $\Omega$ 属于 $C^\infty$ 类, 且 $\Gamma$
有界.

\begin{keytheorem}[存在性与唯一性]
\label{thm:10.7}
假设 $u_0\in H^2(\Omega)\cap H_0^1(\Omega)$ 且 $v_0\in H_0^1(\Omega)$.则
\eqref{eq:10.27}--\eqref{eq:10.30} 存在唯一的解 $u$, 满足
\begin{equation}
u\in C([0,\infty);H^2(\Omega)\cap H_0^1(\Omega))\cap
C^1([0,\infty);H_0^1(\Omega))\cap C^2([0,\infty);L^2(\Omega)).
\label{eq:10.31}
\end{equation}
此外,\footnote{我们采用与前面各节相同的记号, 即
$\left\|\dfrac{\partial u}{\partial t}(t)\right\|^2_{L^2(\Omega)}
=\int_\Omega\abs{\dfrac{\partial u}{\partial t}(x,t)}^2dx$,
$\norm{\nabla u(t)}^2_{L^2(\Omega)}
=\sum_{i=1}^N\int_\Omega\abs{\dfrac{\partial u}{\partial x_i}(x,t)}^2dx$.}
\begin{equation}
\left\|\frac{\partial u}{\partial t}(t)\right\|^2_{L^2(\Omega)}
+\abs{\nabla u(t)}^2_{L^2(\Omega)}
=\abs{v_0}^2_{L^2(\Omega)}+\abs{\nabla u_0}^2_{L^2(\Omega)}
\quad\forall t\ge0.
\label{eq:10.32}
\end{equation}
\end{keytheorem}

\begin{remark}
\label{rem:10.6}
方程 \eqref{eq:10.32} 是一条\emph{守恒律}, 它断言系统的能量不随时间改变.
\end{remark}

在证明定理  \ref{thm:10.7} 之前, 我们先提一个正则性结果.

\begin{theorem}[正则性]
\label{thm:10.8}
假设初值满足
\begin{equation*}
u_0\in H^k(\Omega),\quad v_0\in H^k(\Omega)\quad\forall k,
\end{equation*}
以及相容性条件
\begin{align*}
\Delta^ju_0&=0\quad\text{在}\ \Gamma\ \text{上}\quad
\forall j\ge0,\ j\ \text{为整数},\\
\Delta^jv_0&=0\quad\text{在}\ \Gamma\ \text{上}\quad
\forall j\ge0,\ j\ \text{为整数}.
\end{align*}
则 \eqref{eq:10.27}--\eqref{eq:10.30} 的解 $u$ 属于
$C^\infty(\overline\Omega\times[0,\infty))$.
\end{theorem}

\begin{proof}[定理 \ref{thm:10.7}]
与 \S \ref{sec:10.1} 一样, 我们把 $u(x,t)$ 看作定义在 $[0,\infty)$ 上的
向量值函数; 更精确地说, 对每个 $t\ge0$, $u(t)$ 表示映射
$x\mapsto u(x,t)$.我们把 \eqref{eq:10.27} 写成一阶方程组的形式:\footnote{
这是标准的手法, 即把一个 $k$ 阶微分方程写成一个由 $k$ 个一阶方程组成的
方程组.}
\begin{equation}
\begin{cases}
\dfrac{\partial u}{\partial t}-v=0 & \text{在}\ Q\ \text{中},\\[2mm]
\dfrac{\partial v}{\partial t}-\Delta u=0 & \text{在}\ Q\ \text{中},
\end{cases}
\label{eq:10.33}
\end{equation}
并令 $U=\binom{u}{v}$, 于是 \eqref{eq:10.33} 变成
\begin{equation}
\frac{dU}{dt}+AU=0,
\label{eq:10.34}
\end{equation}
其中
\begin{equation}
AU=\begin{pmatrix}0&-I\\-\Delta&0\end{pmatrix}U
=\begin{pmatrix}0&-I\\-\Delta&0\end{pmatrix}
\begin{pmatrix}u\\v\end{pmatrix}
=\begin{pmatrix}-v\\-\Delta u\end{pmatrix}.
\label{eq:10.35}
\end{equation}
现在我们在空间 $H=H_0^1(\Omega)\times L^2(\Omega)$ 中应用 Hille--Yosida
理论, 它配以标量积
\begin{equation*}
(U_1,U_2)=\int_\Omega\nabla u_1\cdot\nabla u_2dx+\int_\Omega u_1u_2dx
+\int_\Omega v_1v_2dx,
\end{equation*}
其中 $U_1=\binom{u_1}{v_1}$, $U_2=\binom{u_2}{v_2}$.

考虑由 \eqref{eq:10.35} 与
\begin{equation*}
D(A)=(H^2(\Omega)\cap H_0^1(\Omega))\times H_0^1(\Omega)
\end{equation*}
定义的无界算子 $A:D(A)\subset H\to H$.

注意, 边界条件 \eqref{eq:10.28} 已被纳入空间 $H$ 中.条件
$v=\dfrac{\partial u}{\partial t}=0$ 在 $\Sigma$ 上是 \eqref{eq:10.28} 的
直接推论.

我们断言 $A+I$ 在 $H$ 中极大单调:

(i) \emph{$A+I$ 单调}; 事实上, 若 $U=\binom{u}{v}\in D(A)$, 则
\begin{align*}
(AU,U)_H+\abs{U}^2_H
&=-\int_\Omega\nabla v\cdot\nabla u-\int_\Omega uv
+\int_\Omega(-\Delta u)v+\int_\Omega u^2+\int_\Omega\abs{\nabla u}^2
+\int_\Omega v^2\\
&=-\int_\Omega uv+\int_\Omega u^2+\int_\Omega v^2+\int_\Omega\abs{\nabla u}^2
\ge0.
\end{align*}

(ii) \emph{$A+I$ 极大单调.} 这归结为证明 $A+2I$ 是满射.给定
$F=\binom{f}{g}\in H$, 我们必须求解方程 $AU+2U=F$, 即方程组
\begin{equation}
\begin{cases}
-v+2u=f & \text{在}\ \Omega\ \text{中},\\
-\Delta u+2v=g & \text{在}\ \Omega\ \text{中},
\end{cases}
\label{eq:10.36}
\end{equation}
其中
\begin{equation*}
u\in H^2(\Omega)\cap H_0^1(\Omega)\quad\text{且}\quad
v\in H_0^1(\Omega).
\end{equation*}
由 \eqref{eq:10.36} 推出
\begin{equation}
-\Delta u+4u=2f+g.
\label{eq:10.37}
\end{equation}
方程 \eqref{eq:10.37} 有唯一的解 $u\in H^2(\Omega)\cap H_0^1(\Omega)$
(由定理  \ref{thm:9.25}).然后取 $v=2u-f$ 即得到 $v\in H_0^1(\Omega)$.
这样就解出了 \eqref{eq:10.36}.

应用 Hille--Yosida 定理 (定理  \ref{thm:7.4}) 与注记 \ref{rem:7.7}, 我们
看到问题
\begin{equation}
\begin{cases}
\dfrac{dU}{dt}+AU=0 & \text{在}\ [0,\infty)\ \text{上},\\
U(0)=U_0
\end{cases}
\label{eq:10.38}
\end{equation}
存在唯一的解, 满足
\begin{equation}
U\in C^1([0,\infty);H)\cap C([0,\infty);D(A)),
\label{eq:10.39}
\end{equation}
因为 $U_0=\binom{u_0}{v_0}\in D(A)$.由 \eqref{eq:10.39} 我们推出
\eqref{eq:10.31}.

为证明 \eqref{eq:10.32}, 只须把 \eqref{eq:10.27} 乘以
$\dfrac{\partial u}{\partial t}$ 并在 $\Omega$ 上积分.注意
\begin{equation*}
\int_\Omega\frac{\partial^2u}{\partial t^2}
\frac{\partial u}{\partial t}dx
=\frac12\frac{\partial}{\partial t}\int_\Omega
\abs{\frac{\partial u}{\partial t}(x,t)}^2dx
\end{equation*}
以及
\begin{equation*}
\int_\Omega(-\Delta u)\frac{\partial u}{\partial t}dx
=\int_\Omega\nabla u\cdot\frac{\partial}{\partial t}(\nabla u)dx
=\frac12\frac{\partial}{\partial t}\int_\Omega\abs{\nabla u}^2dx.
\end{equation*}
\end{proof}

\begin{remark}
\label{rem:10.7}
当 $\Omega$ 有界时, 我们可以在 $H_0^1(\Omega)$ 上使用标量积
$\int\nabla u_1\cdot\nabla u_2$ (见推论  \ref{cor:9.19}), 并在
$H=H_0^1(\Omega)\times L^2(\Omega)$ 上使用标量积
\begin{equation*}
(U_1,U_2)=\int_\Omega\nabla u_1\cdot\nabla u_2+\int_\Omega v_1v_2,
\end{equation*}
其中 $U_1=\begin{pmatrix}u_1\\v_1\end{pmatrix}$,
$U_2=\begin{pmatrix}u_2\\v_2\end{pmatrix}$.

用这个标量积, 我们有
\begin{equation*}
(AU,U)=-\int_\Omega\nabla v\cdot\nabla u+\int_\Omega(-\Delta u)v=0
\quad\forall U=\begin{pmatrix}u\\v\end{pmatrix}\in D(A).
\end{equation*}
容易验证:

(i) $A$ 与 $-A$ 极大单调,

(ii) $A^\star=-A$.
\end{remark}

因此我们也可以求解问题
\begin{equation*}
\frac{dU}{dt}-AU=0\quad\text{在}\ [0,+\infty)\ \text{上},\quad U(0)=U_0,
\end{equation*}
或等价地
\begin{equation*}
\frac{dU}{dt}+AU=0\quad\text{在}\ (-\infty,0]\ \text{上},\quad U(0)=U_0
\end{equation*}
(只须把 $t$ 换成 $-t$).\footnote{换言之, 时间是可逆的; 从这个观点看,
波动方程与热方程 (对热方程时间不可逆) 之间有本质的差别.} 关系式
\eqref{eq:10.32} 可以写成
\begin{equation*}
\abs{U(t)}_H=\abs{U_0}_H\quad\forall t\in\R.
\end{equation*}
我们说单参数族 $\{U(t)\}_{t\in\R}$ 是 $H$ 上的\emph{等距群}.

\begin{keyremark}
\label{rem:10.8}
与热方程相反, 波动方程对初值\emph{没有}光滑化效应.为说明这一点, 只须考虑
$\Omega=\R$ 的情形.此时 \eqref{eq:10.27}--\eqref{eq:10.30} 有一个非常
简单的显式解, 即
\begin{equation}
u(x,t)=\frac12(u_0(x+t)+u_0(x-t))
+\frac12\int_{x-t}^{x+t}v_0(s)ds.
\label{eq:10.40}
\end{equation}
特别地, 若 $v_0=0$, 我们有
\begin{equation*}
u(x,t)=\frac12(u_0(x+t)+u_0(x-t)).
\end{equation*}
显然 $u$ 并不比 $u_0$ 更正则.我们甚至可以更精确些.假设
$u_0\in C^\infty(\R\setminus\{x_0\})$.则 $u(x,t)$ 在 $\R\times\R$ 上
属于 $C^\infty$ 类, 但在直线 $x+t=x_0$ 与 $x-t=x_0$ 上除外.这些直线
称为过点 $(x_0,0)$ 的\emph{特征线}.我们说奇性沿特征线传播.
\end{keyremark}

\begin{remark}
\label{rem:10.9}
当 $\Omega$ 有界时, 问题 \eqref{eq:10.27}--\eqref{eq:10.30} 可以像热方程
那样通过在一个 Hilbert 基中分解来求解.非常方便的是在 $L^2(\Omega)$ 的
由 $-\Delta$ 的特征函数 (带 Dirichlet 条件) 组成的基 $(e_i)$ 中工作, 即
$-\Delta e_i=\lambda_ie_i$ 在 $\Omega$ 中, $e_i=0$ 在 $\Gamma$ 上; 回忆
$\lambda_i>0$.我们以级数
\begin{equation}
u(x,t)=\sum_ia_i(t)e_i(x)
\label{eq:10.41}
\end{equation}
的形式寻求 \eqref{eq:10.27}--\eqref{eq:10.30} 的解.我们立即看出, 函数
$a_i(t)$ 必须满足
\begin{equation*}
a_i''(t)+\lambda_ia_i(t)=0,
\end{equation*}
从而
\begin{equation*}
a_i(t)=a_i(0)\cos(\sqrt{\lambda_i}t)
+\frac{a_i'(0)}{\sqrt{\lambda_i}}\sin(\sqrt{\lambda_i}t).
\end{equation*}
常数 $a_i(0)$ 与 $a_i'(0)$ 由关系式
\begin{equation*}
u_0(x)=\sum_ia_i(0)e_i(x)
\quad\text{与}\quad
v_0(x)=\sum_ia_i'(0)e_i(x)
\end{equation*}
确定.换言之, $a_i(0)$ 与 $a_i'(0)$ 是 $u_0$ 与 $v_0$ 在基 $(e_i)$ 中的
分量.关于这个级数的收敛性, 参见例如 H.~Weinberger \bref{Weinberger}{1}.
\end{remark}

\begin{proof}[定理 \ref{thm:10.8}]
我们使用与定理  \ref{thm:10.7} 的证明相同的记号.容易通过对 $k$ 作归纳
证明
\begin{equation*}
D(A^k)=\left\{\begin{pmatrix}u\\v\end{pmatrix}\ \middle|\
\begin{array}{l}
u\in H^{k+1}(\Omega)\ \text{且}\ \Delta^ju=0\ \text{在}\ \Gamma\ \text{上}
\ \forall j,\ 0\le j\le[k/2]\\[1mm]
v\in H^k(\Omega)\ \text{且}\ \Delta^jv=0\ \text{在}\ \Gamma\ \text{上}
\ \forall j,\ 0\le j\le[(k+1)/2]-1
\end{array}
\right\}.
\end{equation*}
特别地, $D(A^k)\subset H^{k+1}(\Omega)\times H^k(\Omega)$, 且嵌入连续.
应用定理  \ref{thm:7.5}, 我们看到: 若
$U_0=\binom{u_0}{v_0}\in D(A^k)$, 则 \eqref{eq:10.38} 的解 $U$ 满足
\begin{equation*}
U\in C^{k-j}([0,\infty);D(A^j))\quad\forall j=0,1,\dots,k.
\end{equation*}
于是 $u\in C^{k-j}([0,\infty);H^{j+1}(\Omega))\ \forall j=0,1,\dots,k$.
借助推论  \ref{cor:9.15}, 我们在定理  \ref{thm:10.8} 的假设下 (即
$U_0\in D(A^k)\ \forall k$) 得出
\begin{equation*}
u\in C^k(\overline\Omega\times[0,\infty))\quad\forall k.
\end{equation*}
\end{proof}

\begin{remark}
\label{rem:10.10}
定理  \ref{thm:10.8} 中引入的相容性条件, 是问题
\eqref{eq:10.27}--\eqref{eq:10.30} 有解
$u\in C^\infty(\overline\Omega\times[0,\infty))$ 的\emph{充要}条件.证明
与注记 \ref{rem:10.4} 相同.
\end{remark}

\begin{remark}
\label{rem:10.11}
\S \ref{sec:10.3} 中介绍的方法也可以用来求解 \emph{Klein--Gordon} 方程
\begin{equation*}
\frac{\partial^2u}{\partial t^2}-\Delta u+m^2u=0\quad\text{在}\ Q\ \text{中},
\quad m>0.
\tag*{(27$'$)}
\end{equation*}
注意, $(27')$ 不能通过形如
$v(x,t)=e^{\lambda t}u(x,t)$ 的未知函数变换化为 \eqref{eq:10.27}.
\end{remark}
\section*{第 10 章评注}
\addcontentsline{toc}{section}{第 10 章评注}

\cmtitem{\emph{关于热方程的评注}}

\cmtnum{1.} \emph{J.-L.~Lions 的途径.}

下述结果使我们能够在一个非常一般的框架中证明抛物型问题弱解的存在性与唯一性.
这个定理可以看作 Lax--Milgram 定理的抛物型对应物.设 $H$ 是带有标量积
$(\ ,\ )$ 与范数 $\abs{\cdot}$ 的 Hilbert 空间.对偶空间 $H^\star$ 与 $H$
等同.

设 $V$ 是另一个 Hilbert 空间, 范数为 $\norm{\cdot}$.我们假设
$V\subset H$, 且嵌入稠密, 连续, 于是
\begin{equation*}
V\subset H\subset V^\star
\end{equation*}
(见注记 \ref{rem:5.3}).

固定 $T>0$; 对几乎所有的 $t\in[0,T]$, 给定双线性形式
$a(t;u,v):V\times V\to\R$, 它满足下述性质:

(i) 对每个 $u,v\in V$, 函数 $t\mapsto a(t;u,v)$ 可测,

(ii) 对几乎所有的 $t\in[0,T]$, $\abs{a(t;u,v)}\le M\norm{u}\norm{v}\
\forall u,v\in V$,

(iii) 对几乎所有的 $t\in[0,T]$, $a(t;v,v)\ge\alpha\norm{v}^2-C\abs{v}^2\
\forall v\in V$,

其中 $\alpha>0$, $M$ 与 $C$ 为常数.

\begin{theorem}[J.-L.~Lions]
\label{thm:10.9}
给定 $f\in L^2(0,T;V^\star)$ 与 $u_0\in H$, 存在唯一的函数 $u$ 满足
\begin{equation*}
u\in L^2(0,T;V)\cap C([0,T];H),\qquad
\frac{du}{dt}\in L^2(0,T;V^\star),
\end{equation*}
\begin{equation*}
\dual{\frac{du}{dt}(t)}{v}+a(t;u(t),v)=\dual{f(t)}{v}
\quad\text{对几乎所有的}\ t\in(0,T),\ \forall v\in V,
\end{equation*}
以及
\begin{equation*}
u(0)=u_0.
\end{equation*}
\end{theorem}

证明参见例如 J.-L.~Lions--E.~Magenes \bref{LionsMagenes}{1}.

\emph{应用.} $H=L^2(\Omega)$, $V=H_0^1(\Omega)$,
\begin{equation*}
a(t;u,v)=\sum_{i,j}\int_\Omega a_{ij}(x,t)\frac{\partial u}{\partial x_i}
\frac{\partial v}{\partial x_j}dx
+\sum_i\int_\Omega a_i(x,t)\frac{\partial u}{\partial x_i}v
+\int_\Omega a_0(x,t)uv\,dx,
\end{equation*}
其中 $a_{ij},a_i,a_0\in L^\infty(\Omega\times(0,T))$ 且
\begin{equation}
\sum_{i,j}a_{ij}(x,t)\xi_i\xi_j\ge\alpha\abs{\xi}^2
\quad\text{对几乎所有的}\ (x,t)\in\Omega\times(0,T),\ \forall\xi\in\R^N,
\label{eq:10.42}
\end{equation}
$\alpha>0$.这样我们就得到问题
\begin{equation}
\begin{cases}
\displaystyle\frac{\partial u}{\partial t}
-\sum_{i,j}\frac{\partial}{\partial x_j}
\left(a_{ij}\frac{\partial u}{\partial x_i}\right)
+\sum_ia_i\frac{\partial u}{\partial x_i}+a_0u=f
& \text{在}\ \Omega\times(0,T)\ \text{中},\\[2mm]
u=0 & \text{在}\ \Gamma\times(0,T)\ \text{上},\\
u(x,0)=u_0(x) & \text{在}\ \Omega\ \text{上}
\end{cases}
\label{eq:10.43}
\end{equation}
的弱解.在对数据作附加假设时, \eqref{eq:10.43} 的解有更好的正则性; 见
下面的评注.

\cmtnum{2.} \emph{$C^\infty$ 正则性.}

这里我们假设 $\Omega$ 有界且属于 $C^\infty$ 类.设
$a_{ij},a_i,a_0\in C^\infty(\overline\Omega\times[0,T])$ 满足
\eqref{eq:10.42}.

\begin{theorem}
\label{thm:10.10}
假设 $u_0\in L^2(\Omega)$ 且 $f\in C^\infty(\overline\Omega\times[0,T])$.
则 \eqref{eq:10.43} 的解 $u$ 对每个 $\varepsilon>0$ 属于
$C^\infty(\overline\Omega\times[\varepsilon,T])$.若此外
$u_0\in C^\infty(\overline\Omega)$ 且 $\{f,u_0\}$ 在
$\Gamma\times\{0\}$ 上满足适当的相容性条件,\footnote{我们不显式写出这些
关系式; 它们是 \eqref{eq:10.8} 的自然推广 (另见注记 \ref{rem:10.4}).}
则 $u\in C^\infty(\overline\Omega\times[0,T])$.
\end{theorem}

证明参见例如 J.-L.~Lions--E.~Magenes \bref{LionsMagenes}{1},
A.~Friedman \bref{Friedman}{1}, \bref{Friedman}{2}, 以及
O.~Ladyzhenskaya--V.~Solonnikov--N.~Uraltseva
\bref{LadyzhenskayaSolonnikovUraltseva}{1}; 它基于与第~
\ref{chap:7}~章和
\S \ref{sec:10.1} 中给出的估计非常相似的估计.

我们要提到, 还有一种抽象理论把 Hille--Yosida 理论推广到形如
$\dfrac{du}{dt}(t)+A(t)u(t)=f(t)$ 的问题, 其中对每个 $t$, $A(t)$ 是极大
单调算子.这个理论由 T.~Kato, H.~Tanabe, P.~E.~Sobolevskii 等人发展起来.
处理它比 Hille--Yosida 理论在技术上更复杂; 参见
A.~Friedman \bref{Friedman}{2}, H.~Tanabe \bref{Tanabe}{1}, 以及
K.~Yosida \bref{Yosida}{1}.

\cmtnum{3.} \emph{$L^p$ 与 $C^{0,\alpha}$ 正则性.}

考虑问题\footnote{当然, 我们也可以给定非齐次 Dirichlet 条件
$u(x,t)=g(x,t)$ 在 $\Gamma\times(0,T)$ 上, 但为简单起见我们只处理
$g=0$ 的情形.}
\begin{equation}
\begin{cases}
\dfrac{\partial u}{\partial t}-\Delta u=f & \text{在}\
\Omega\times(0,T)\ \text{中},\\[2mm]
u=0 & \text{在}\ \Gamma\times(0,T)\ \text{上},\\
u(x,0)=u_0(x) & \text{在}\ \Omega\ \text{上}.
\end{cases}
\label{eq:10.44}
\end{equation}
为方便起见, 假设 $\Omega$ 有界且属于 $C^\infty$ 类.我们先从一个简单的
结果开始.

\begin{theorem}[$L^2$ 正则性]
\label{thm:10.11}
给定 $f\in L^2(\Omega\times(0,T))$ 与 $u_0\in H_0^1(\Omega)$, \eqref{eq:10.44}
存在唯一的解, 满足
\begin{equation*}
u\in C([0,T];H_0^1(\Omega))\cap L^2(0,T;H^2(\Omega)\cap H_0^1(\Omega))
\end{equation*}
且
\begin{equation*}
\frac{\partial u}{\partial t}\in L^2(0,T;L^2(\Omega)).
\end{equation*}
\end{theorem}

证明是容易的; 参见例如 J.-L.~Lions--E.~Magenes \bref{LionsMagenes}{1}.
更一般地, 在 $L^p$ 空间中我们有下述结果.

\begin{theorem}[$L^p$ 正则性]
\label{thm:10.12}
给定 $f\in L^p(\Omega\times(0,T))$, 其中 $1<p<\infty$, 且
$u_0=0$,\footnote{为简化起见.} 则 \eqref{eq:10.44} 存在唯一的解, 满足
\begin{equation*}
u,\ \frac{\partial u}{\partial t},\
\frac{\partial u}{\partial x_i},\
\frac{\partial^2u}{\partial x_i\partial x_j}
\in L^p(\Omega\times(0,T))\quad\forall i,j.
\end{equation*}
\end{theorem}

\begin{theorem}[Hölder 正则性]
\label{thm:10.13}
设 $0<\alpha<1$.假设\footnote{即
$\abs{f(x_1,t_1)-f(x_2,t_2)}\le C(\abs{x_1-x_2}^2+\abs{t_1-t_2})^{\alpha/2}
\ \forall x_1,x_2,t_1,t_2$.}
$f\in C^{\alpha,\alpha/2}(\overline\Omega\times[0,T])$ 且
$u_0\in C^{2+\alpha}(\overline\Omega)$ 满足自然的相容性条件
\begin{equation*}
u_0=0\ \text{在}\ \Gamma\ \text{上}
\quad\text{且}\quad
-\Delta u_0=f(x,0)\ \text{在}\ \Gamma\ \text{上}.
\end{equation*}
则 \eqref{eq:10.44} 存在唯一的解 $u$, 使得
\begin{equation*}
u,\ \frac{\partial u}{\partial t},\
\frac{\partial u}{\partial x_i},\
\frac{\partial^2u}{\partial x_i\partial x_j}
\in C^{\alpha,\alpha/2}(\overline\Omega\times[0,T])\quad\forall i,j.
\end{equation*}
\end{theorem}

定理  \ref{thm:10.12} 与  \ref{thm:10.13} 的证明是精细的, 但定理
 \ref{thm:10.12} 中 $p=2$ 的情形除外.与椭圆情形一样 (见第~
\ref{chap:9}~章末的评注),
它们依赖于下述两点:

(i) 一个显式的表示公式, 它涉及 $\dfrac{\partial}{\partial t}-\Delta$ 的
基本解.例如, 若 $\Omega=\R^N$ 且 $f=0$, 则
\begin{equation}
u(x,t)=\int_{\R^N}E(x-y,t)u_0(y)dy=E\star u_0,
\label{eq:10.45}
\end{equation}
其中 $\star$ 只表示关于空间变量 $x$ 的卷积, $E$ 是\emph{热核},
$E(x,t)=(4\pi t)^{-N/2}e^{-\abs{x}^2/4t}$; 参见例如
G.~Folland \bref{Folland}{1}.

(ii) 奇异积分的技术.

关于这个课题, 参见例如 O.~Ladyzhenskaya--V.~Solonnikov--N.~Uraltseva
\bref{LadyzhenskayaSolonnikovUraltseva}{1}, A.~Friedman
\bref{Friedman}{1}, N.~Krylov \bref{Krylov}{1}, \bref{Krylov}{2},
P.~Grisvard \bref{Grisvard}{1} (第 9 节),
D.~Stroock--S.~Varadhan \bref{StroockVaradhan}{1},
A.~Brandt \bref{Brandt}{2}, B.~Knerr \bref{Knerr}{1}, 以及
L.~Simon \bref{Simon}{2} 为 Hölder 正则性设计了更初等的论证.

应当记住的一般``哲学''如下: 若 $u$ 是 \eqref{eq:10.44} 在 $u_0=0$ 时的解,
则 $\dfrac{\partial u}{\partial t}$ 与 $\Delta u$ 都与 $f$ 有相同的正则性.

最后, 我们要提到: 若把 $\Delta$ 换成
\begin{equation*}
\sum_{i,j}\frac{\partial}{\partial x_j}
\left(a_{ij}(x,t)\frac{\partial u}{\partial x_i}\right)
+\sum_ia_i(x,t)\frac{\partial u}{\partial x_i}+a_0(x,t)u,
\end{equation*}
其中系数光滑且满足
\begin{equation}
\sum_{i,j}a_{ij}(x,t)\xi_i\xi_j\ge\nu\abs{\xi}^2
\quad\forall x,t,\ \forall\xi\in\R^N,\ \nu>0,
\label{eq:10.46}
\end{equation}
则定理  \ref{thm:10.11},  \ref{thm:10.12} 与  \ref{thm:10.13} 的结论仍然
成立.在系数不规则 (即 $a_{ij}\in L^\infty(\Omega\times(0,T))$ 且满足
\eqref{eq:10.46}) 的情形, Nash--Moser 的一个困难结果表明: 存在某个
$\alpha>0$ 使得 $u\in C^{\alpha,\alpha/2}(\overline\Omega\times[0,T])$;
参见例如 O.~Ladyzhenskaya--V.~Solonnikov--N.~Uraltseva
\bref{LadyzhenskayaSolonnikovUraltseva}{1}.

\cmtnum{4.} \emph{抛物型方程的一些例子.}

线性与非线性抛物型方程 (以及方程组) 出现在许多领域: 力学, 物理学, 化学, 生物学, 最优控制, 概率论, 金融, 图像处理等.我们举出一些例子:

(i) \emph{Navier--Stokes 方程组}:
\begin{equation}
\frac{\partial u_i}{\partial t}-\Delta u_i+\sum_ju_j\frac{\partial u_i}
{\partial x_j}=f_i+\frac{\partial p}{\partial x_i}
\quad\text{在}\ \Omega\times(0,T)\ \text{中},\ 1\le i\le N,
\label{eq:10.47}
\end{equation}
\begin{equation}
\operatorname{div}u=\sum_{i=1}^N\frac{\partial u_i}{\partial x_i}=0
\quad\text{在}\ \Omega\times(0,T)\ \text{中},
\label{eq:10.48}
\end{equation}
\begin{equation}
u=0\quad\text{在}\ \Gamma\times(0,T)\ \text{上},
\label{eq:10.49}
\end{equation}
\begin{equation}
u(x,0)=u_0(x)\quad\text{在}\ \Omega\ \text{上},
\label{eq:10.50}
\end{equation}
它在流体力学中起中心作用; 参见例如 R.~Temam \bref{Temam}{1} 及其中的
参考文献.

(ii) \emph{反应--扩散方程组.} 这些是形如
\begin{equation*}
\begin{cases}
\dfrac{\partial u}{\partial t}-M\Delta u=f(u)
& \text{在}\ \Omega\times(0,T)\ \text{中}\\
\text{加上边界条件与初值}
\end{cases}
\end{equation*}
的非线性抛物型方程或方程组, 其中 $u(x,t)$ 取值于 $\R^m$, $M$ 是一个
$m\times m$ (对角) 矩阵, $f$ 是由 $\R^m$ 到 $\R^m$ 的非线性映射.这些
方程组被用来模拟各种领域中出现的现象: 化学, 生物学, 神经生理学, 流行病学, 燃烧, 群体遗传学, 生态学, 地质学等; 参见例如 P.~Fife \bref{Fife}{1}
及其大量的参考文献.反应--扩散方程的解表现出范围广泛的行为, 包括行波的
形成与自组织斑图的形成.

(iii) \emph{自由边界问题.} 例如, Stefan 问题描述冰水混合物的演化; 参见
例如 E.~Magenes \bref{Magenes}{1} 的综述文章与 A.~Friedman
\bref{Friedman}{4} 的书.

(iv) 扩散方程在概率论 (Brown 运动, Markov 过程, 扩散过程, 随机微分方程等)
中起中心作用; 参见例如 D.~Stroock--S.~Varadhan \bref{StroockVaradhan}{1}.

(v) 许多其他半线性抛物型问题的例子见 D.~D.~Henry \bref{Henry}{1} 与
Th.~Cazenave--A.~Haraux \bref{CazenaveHaraux}{1}.

(vi) 热方程的一个有趣的应用与 Atiyah--Singer 指标有关; 参见例如
P.~Gilkey \bref{Gilkey}{1}.

(vii) 更精致的非线性扩散方程用于图像处理 (Perona--Malik 模型的变体).
G.~Perelman 最近对著名的 Poincaré 猜想的解决, 依赖于 R.~Hamilton 对
Ricci 流的细致研究, 而 Ricci 流是一种非线性热方程.

\cmtnum{5.} 关于抛物型方程极大值原理的进一步结果, 参见例如
A.~Friedman \bref{Friedman}{1}, M.~Protter--H.~Weinberger
\bref{ProtterWeinberger}{1}, R.~Sperb \bref{Sperb}{1}.例如, 若 $u$ 是
\eqref{eq:10.1}--\eqref{eq:10.3} 在 $u_0\ge0$ 且 $u_0\not\equiv0$ 时的解,
则 $u(x,t)>0\ \forall x\in\Omega,\ \forall t>0$.当 $\Omega=\R^N$ 时, 这
容易由显式表示公式 \eqref{eq:10.45} 推出.

\cmtitem{\emph{关于波动方程的评注}}

\cmtnum{6.} \emph{波动方程的弱解.}

波动方程弱解的存在性与唯一性有一个一般的抽象框架.设 $V$ 与 $H$ 是两个
Hilbert 空间, 使得 $V\subset H\subset V^\star$ (如评注 1 中那样).对每个
$t\in[0,T]$, 给定对称连续双线性形式 $a(t;u,v):V\times V\to\R$, 使得

(i) 函数 $t\mapsto a(t;u,v)$ 属于 $C^1$ 类 $\forall u,v\in V$,

(ii) $a(t;v,v)\ge\alpha\norm{v}^2-C\abs{v}^2\ \forall t\in[0,T],\
\forall v\in V$, $\alpha>0$.

\begin{theorem}[J.-L.~Lions]
\label{thm:10.14}
给定 $f\in L^2(0,T;H)$, $u_0\in V$ 与 $v_0\in H$, 存在唯一的函数 $u$ 满足
\begin{equation*}
u\in C([0,T];V),\qquad
\frac{du}{dt}\in C([0,T];H),\qquad
\frac{d^2u}{dt^2}\in L^2(0,T;V^\star),
\end{equation*}
\begin{equation*}
\dual{\frac{d^2u}{dt^2}(t)}{v}+a(t;u(t),v)=\dual{f(t)}{v}
\quad\text{对几乎所有的}\ t\in(0,T),\ \forall v\in V,
\end{equation*}
\begin{equation*}
u(0)=u_0\quad\text{且}\quad\frac{du}{dt}(0)=v_0.
\end{equation*}
\end{theorem}

证明参见例如 J.-L.~Lions--E.~Magenes \bref{LionsMagenes}{1}.

\emph{应用.} 令 $H=L^2(\Omega)$, $V=H_0^1(\Omega)$,
\begin{equation*}
a(t;u,v)=\int_\Omega\sum_{i,j}a_{ij}(x,t)\frac{\partial u}{\partial x_i}
\frac{\partial v}{\partial x_j}dx+\int_\Omega a_0(x,t)uv\,dx,
\end{equation*}
其中 \eqref{eq:10.42} 成立且
\begin{equation*}
a_{ij},\ \frac{\partial a_{ij}}{\partial t},\ a_0,\
\frac{\partial a_0}{\partial t}\in L^\infty(\Omega\times(0,T)),\qquad
a_{ij}=a_{ji}\ \forall i,j.
\end{equation*}
则问题
\begin{equation*}
\begin{cases}
\displaystyle\frac{\partial^2u}{\partial t^2}
-\sum_{i,j}\frac{\partial}{\partial x_j}
\left(a_{ij}\frac{\partial u}{\partial x_i}\right)+a_0u=f
& \text{在}\ \Omega\times(0,T)\ \text{中},\\
\eqref{eq:10.28},~\eqref{eq:10.29},~\eqref{eq:10.30}
\end{cases}
\end{equation*}
存在唯一的弱解.注意, 对初值的假设 ($u_0\in H_0^1(\Omega)$ 且
$v_0\in L^2(\Omega)$) 比定理  \ref{thm:10.7} 中的假设更弱.在对 $f$, $u_0$, $v_0$ (正则性与相容性条件) 以及 $a_{ij}$, $a_0$ 作附加假设时, 可以得到
$u$ 的更好正则性.

\cmtnum{7.} 波动方程的 $L^p$ 理论是精细的, 在过去 30 年里得到了广泛的
研究.Strichartz 估计是一个重要的工具; 参见例如
S.~Klainerman \bref{Klainerman}{1}.

\cmtnum{8.} \emph{极大值原理.}

极大值原理的一些非常特殊的形式对波动方程成立; 参见例如
M.~Protter--H.~Weinberger \bref{ProtterWeinberger}{1}.例如, 设 $u$ 是
\eqref{eq:10.27}--\eqref{eq:10.30} 的解.

(i) 若 $\Omega=\R$, $u_0\ge0$ 且 $v_0\ge0$, 则 $u\ge0$.

(ii) 若 $\Omega=\R^2$, $u_0=0$ 且 $v_0\ge0$, 则 $u\ge0$.

断言 (i) 由表示公式 \eqref{eq:10.40} 推出.在 $\R^N$ 中有类似但更复杂的
公式; 参见例如 S.~Mizohata \bref{Mizohata}{1}, G.~Folland
\bref{Folland}{1}, H.~Weinberger \bref{Weinberger}{1},
R.~Courant--D.~Hilbert \bref{CourantHilbert}{1}, 以及
S.~Mikhlin \bref{Mikhlin}{1}.它蕴含 (ii).

然而, 我们要提醒读者注意下述几点:

(iii) 若 $\Omega=(0,1)$, $u_0\ge0$ 且 $v_0=0$, 则一般不能推出 $u\ge0$.

(iv) 若 $\Omega=\R^2$, $u_0\ge0$ 且 $v_0=0$, 则一般不能推出 $u\ge0$.

J.~Mawhin--R.~Ortega--A.~M.~Robles--P\'erez \bref{MawhinOrtegaRoblesPerez}{1}
最近为电报方程 (它与波动方程相似) 设计了一种不寻常的极大值原理形式.

\cmtnum{9.} \emph{\index{yilaiquyu@依赖区域 (domain of dependence)}
依赖区域. 波的传播. \index{Huygensyuanli@Huygens 原理 (Huygens' principle)}
Huygens 原理.}

热方程与波动方程之间有本质的差别:

(i) 对热方程, 初值的一个小扰动\footnote{即在小的区域内局部化的扰动.}
会立即在所有地方被感觉到, 即 $\forall x\in\Omega,\ \forall t>0$.例如,
我们已经看到, 若 $u_0\ge0$ 且 $u_0\not\equiv0$, 则
$u(x,t)>0\ \forall x\in\Omega,\ \forall t>0$.我们说热以\emph{无穷大速度}
传播.\footnote{从物理上说, 这并不现实! 然而, 表示公式 \eqref{eq:10.45}
表明: 当初值的扰动局部化在 $x=x_0$ 附近时, 若 $t$ 小且 $\abs{x-x_0}$ 大,
则它在点 $(x,t)$ 处的效应可以忽略.}

(ii) 对波动方程, 情形完全不同.例如假设 $\Omega=\R$.显式公式
\eqref{eq:10.40} 表明 $u(\bar x,\bar t)$ 只依赖于 $u_0$ 与 $v_0$ 在区间
$[\bar x-\bar t,\bar x+\bar t]$ 上的值; 见图  \ref{fig:8}.

\begin{figure}[tbp]
	\centering
	\begin{tikzpicture}[scale=1.05]
		\draw[-{Stealth}] (-0.4,0) -- (6.2,0) node[below] {$x$};
		\draw[-{Stealth}] (0,-0.4) -- (0,3.8) node[right] {$t$};
		\draw (0,1.6) -- (-0.08,1.6) node[left] {$\bar t$};
		\draw (1.1,-0.3) -- (3.0,1.6);
		\draw (4.9,-0.3) -- (3.0,1.6);
		\fill (3.0,1.6) circle (1.8pt);
		\node[right] at (3.15,1.5) {$(\bar x,\bar t)$};
		\foreach \xpos in {1.4,1.55,...,4.6}{
			\draw (\xpos-0.08,-0.12) -- (\xpos+0.08,0.18);
		}
		\draw (3.0,0.15) -- (3.0,-0.15);
		\node[below] at (0.9,-0.25) {$\bar x-\bar t$};
		\node[below] at (3.0,-0.15) {$\bar x$};
		\node[below] at (5.1,-0.25) {$\bar x+\bar t$};
	\end{tikzpicture}
	\caption{}
	\label{fig:8}
\end{figure}

我们说 $x$ 轴上的区间 $[\bar x-\bar t,\bar x+\bar t]$ 是点
$(\bar x,\bar t)$ 的\emph{依赖区域}.对 $\Omega=\R^N$ ($N\ge2$) 也成立
同样的事实: $u(\bar x,\bar t)$ 只依赖于 $u_0$ 与 $v_0$ 在球
$\{x\in\R^N;\ \abs{x-\bar x}\le\bar t\}$ 上的值.这个位于超平面
$\R^N\times\{0\}$ 中的球称为点 $(\bar x,\bar t)$ 的依赖区域.从几何上看,
它是锥
\begin{equation*}
\{(x,t)\in\R^N\times\R;\ \abs{x-\bar x}\le\bar t-t\ \text{且}\
t\le\bar t\}
\end{equation*}
与超平面 $\R^N\times\{0\}$ 的交.物理解释是: 波以小于 1 的速度传播.
\footnote{出现速度 1 是因为我们对波动方程作了标准化.有些读者可能更愿意
使用方程 $\dfrac{\partial^2u}{\partial t^2}-c^2\Delta u=0$.} 在时刻
$t=0$ 局部化于区域\footnote{即 $u_0$ 与 $v_0$ 的支集在 $D$ 中.} $D$ 中的
信号, 只有在时刻 $t\ge\dist(x,D)$ 之后才能在点 $x\in\R^N$ 处被感觉到
(当 $t<\dist(x,D)$ 时 $u(x,t)=0$).

当 $N>1$ 为奇数时, 例如 $N=3$, 还有一个更引人注目的效应:
$u(\bar x,\bar t)$ 只依赖于 $u_0$ 与 $v_0$\footnote{以及它们的某些导数.}
在球面 $\{x\in\R^N;\ \abs{x-\bar x}=t\}$ 上的值.这就是 \emph{Huygens
原理}.从物理上说, 它表明在时刻 $t=0$ 局部化于区域 $D$ 中的信号, 只在
时间区间 $[t_1,t_2]$ 内于点 $x\in\R^N$ 处被观察到, 其中
$t_1=\inf_{y\in D}\dist(x,y)$ 且 $t_2=\sup_{y\in D}\dist(x,y)$.在时刻
$t_2$ 之后, 信号在点 $x$ 处就不再被感觉到.

另一方面, 若维数 $N$ 为偶数 (例如 $N=2$), 则信号在 $x$ 处对所有时刻
$t>t_1$ 都会持续存在.\footnote{这个效应随时间衰减, 但不会完全消失.}

\emph{对音乐的一个应用.} 位于 $\R^3$ 中, 与乐器\footnote{尺寸很小的
乐器.}相距 $d$ 的听者, 在时刻 $t$ 听到的是在时刻 $t-d$ 演奏的音符,
此外什么也听不到!\footnote{而在 $\R^2$ 中, 他会听到在时间区间 $[0,t-d]$
内演奏的所有音符的一个加权平均.}

关于 Huygens 原理的更多细节, 读者可以查阅 R.~Courant--D.~Hilbert
\bref{CourantHilbert}{1}, G.~Folland \bref{Folland}{1},
P.~Garabedian \bref{Garabedian}{1}, 以及 S.~Mikhlin
\bref{Mikhlin}{1}.

% chapters/ch11.tex
% 第~
% \ref{chap:11}~章 杂记与补充
% 源: 原书 PDF 第 352--373 页 (印刷页 349--370)

\chapter{杂记与补充}
\label{chap:11}

为了使全书的叙述更加紧凑, 本章收集了未纳入正文的若干补充内容.它们与第 1--7
章相联系.其中一些证明非常简略.有几个证明被略去, 有兴趣的读者可参阅参考文献.

\section{有限维与有限余维空间}
\label{sec:11.1}

众所周知, 任何 $p$ 维的有限维空间 $X$ 都同构于 $\R^p$.特别地, $X$ 完备,
$X$ 上的所有范数都等价, 且闭单位球 $B_X$ 是紧的.

\begin{proposition}
\label{prop:11.1}
设 $E$ 为 Banach 空间, $X\subset E$ 为有限维空间.则 $X$ 是闭的.
\end{proposition}

\begin{proof}
设 $(x_n)$ 为 $X$ 中的序列, 使得 $x_n\to x$ 在 $E$ 中.则 $(x_n)$ 是 $X$
中的 Cauchy 序列, 从而 $(x_n)$ 收敛于 $X$ 中的极限.因此 $x\in X$.
\end{proof}

\begin{proposition}
\label{prop:11.2}
假设 $X$ 有限维, $F$ 为 Banach 空间.则任何线性算子 $T:X\to F$ 必有界.
\end{proposition}

\begin{proof}
设 $(e_i)$ 为 $X$ 的基, 并把 $x$ 写成 $x=\sum_{i=1}^px_ie_i$.则
$Tx=\sum_{i=1}^px_iTe_i$, 从而
$\norm{Tx}\le\sum_{i=1}^p\abs{x_i}\norm{Te_i}\le(\max_i\norm{Te_i})\sum_{i=1}^p\abs{x_i}
\le C\norm{x}$.
\end{proof}

特别地, $X$ 上的所有线性泛函都连续.有限维空间 $X$ 的对偶空间 $X^\star$
也是有限维的, 且 $\dim X^\star=\dim X$.更精确地说, 若 $(e_i)$ 是 $X$ 的
基, 则把 $x\in X$ 写为 $x=\sum_{i=1}^px_ie_i$, 并令 $f_i(x)=x_i$,
$i=1,2,\dots,p$.显然泛函 $(f_i)$ 在 $X^\star$ 中线性无关且它们生成
$X^\star$.因此它们构成 $X^\star$ 的基.较不显然的是下述结论:

\begin{proposition}
\label{prop:11.3}
假设 $X$ 为 Banach 空间 (其中 $\dim X\le\infty$), 使得 $X^\star$ 有限维.
则 $X$ 有限维且 $\dim X=\dim X^\star$.
\end{proposition}

\begin{proof}
我们需要 Hahn--Banach 定理, 更精确地说是推论  \ref{cor:1.4}.设
$J:X\to X^{\star\star}$ 为 \S \ref{sec:1.3} 中定义的典范单射.由于
$\dim X^\star<\infty$, 由上面的讨论推出 $\dim X^{\star\star}<\infty$.但
$X$ 同构于 $J(X)\subset X^{\star\star}$, 从而
$\dim X\le\dim X^{\star\star}=\dim X^\star$.因此 $\dim X<\infty$, 再次
由上面的讨论推出 $\dim X^\star=\dim X$.
\end{proof}

\begin{proposition}
\label{prop:11.4}
设 $E$ 为 Banach 空间, $M\subset E$ 为闭子空间.假设 $X\subset E$ 为有限维
子空间.则 $(M+X)$ 是闭的.此外, $(M+X)$ 在 $E$ 中有补空间当且仅当 $M$ 有
补空间.
\end{proposition}

[警告: 回忆一般说来, 两个闭子空间的和不必是闭的; 参见例如习题 1.14.]

\begin{proof}
先额外假设 $M\cap X=\{0\}$.把 $u_n$ 写成 $u_n=x_n+y_n$, 其中 $x_n\in X$,
$y_n\in M$, 且 $u_n\to u$ 在 $E$ 中.我们断言 $(x_n)$ 有界.若不然, 则对
某个子序列 $n_k\to\infty$, 有 $\norm{x_{n_k}}\to\infty$.抽取进一步的
子序列, 可以假设在 $X$ 中 $\dfrac{x_{n_k}}{\norm{x_{n_k}}}\to\xi$, 满足
$\norm{\xi}=1$ (这里我们用到 $\dim X<\infty$ 的事实).从而
$\dfrac{y_{n_k}}{\norm{x_{n_k}}}=\dfrac{u_{n_k}}{\norm{x_{n_k}}}
-\dfrac{x_{n_k}}{\norm{x_{n_k}}}\to-\xi$; 此外, $\xi\in M$ (因为 $M$ 是
闭的).因此 $\xi\in M\cap X$, 我们必有 $\xi=0$.矛盾.于是我们证明了
$(x_n)$ 有界.抽取子序列, 可以假设 $x_{n_k}\to x$ 在 $X$ 中.则
$y_n\to u-x\in M$ (因为 $M$ 是闭的).因此 $u\in(M+X)$, 这完成了
$M\cap X=\{0\}$ 时 $(M+X)$ 闭的证明.

在一般情形下, 设 $\widetilde X$ 为 $(M\cap X)$ 在 $X$ 中的补空间 (这是
有限维的事情).显然 $\widetilde X$ 有限维, $M\cap\widetilde X=\{0\}$, 且
$M+\widetilde X=M+X$.我们已经证明了 $(M+\widetilde X)$ 是闭的, 从而
$(M+X)$ 也是闭的.

现在设 $M$ 在 $E$ 中有补空间, 记作 $N$.设 $P_M$ 与 $P_N$ 分别为到 $M$
与 $N$ 上的投影算子.由于 $P_N(X)$ 维数有限, 它在 $N$ 中有补空间, 记作
$\widetilde N$ (见 \S \ref{sec:2.4}).我们断言 $\widetilde N$ 是 $(M+X)$
在 $E$ 中的补空间.

首先我们有
\begin{equation*}
(M+X)\cap\widetilde N=\{0\}.
\end{equation*}
事实上, 若 $\tilde n=m+x$, 其中 $\tilde n\in\widetilde N$, $m\in M$,
$x\in X$, 则
\begin{equation*}
\tilde n=P_N\tilde n=P_N(m+x)=P_Nx\in P_N(X),
\end{equation*}
从而 $\tilde n\in\widetilde N\cap P_N(X)=\{0\}$.

其次, 我们有
\begin{equation*}
(M+X)+\widetilde N=E.
\end{equation*}
事实上, 任何 $\xi\in E$ 可以写成
\begin{equation*}
\xi=P_M\xi+P_N\xi,
\end{equation*}
而 $P_N\xi$ 可以进一步分解为
\begin{equation*}
P_N\xi=P_Nx+\tilde n,
\end{equation*}
其中 $x\in X$, $\tilde n\in\widetilde N$.但 $x=P_Mx+P_Nx$, 从而
\begin{equation*}
P_N\xi=(x-P_Mx)+\tilde n,
\end{equation*}
因此
\begin{equation*}
\xi=P_M\xi+x-P_Mx+\tilde n\in(M+X)+\widetilde N.
\end{equation*}

反之, 假设 $(M+X)$ 在 $E$ 中有补空间, 记作 $W$.如上, 设 $\widetilde X$
为 $(M\cap X)$ 在 $X$ 中的补空间.我们断言 $(W+\widetilde X)$ 是 $M$ 的
补空间.

首先我们有
\begin{equation*}
(W+\widetilde X)\cap M=\{0\}.
\end{equation*}
事实上, 若 $m\in M$ 可以写成 $m=w+\tilde x$, 其中 $w\in W$,
$\tilde x\in\widetilde X$, 则 $w=m-\tilde x$, 从而
$w\in(M+X)\cap W=\{0\}$.因此
$m=\tilde x\in(M\cap X)\cap\widetilde X=\{0\}$.

最后, 我们验证
\begin{equation*}
(W+\widetilde X)+M=E.
\end{equation*}
事实上, 只须验证
\begin{equation*}
(M+\widetilde X)=(M+X)
\end{equation*}
(因为 $W+(M+X)=E$).显然 $M+\widetilde X\subset M+X$ (因为
$\widetilde X\subset X$).反之, 任何 $x\in X$ 可以写成 $x=x_1+\tilde x$,
其中 $x_1\in M\cap X$, $\tilde x\in\widetilde X$.因此
$M+X\subset M+\widetilde X$.
\end{proof}

设 $M$ 为 Banach 空间 $E$ 的子空间.回忆: 若存在有限维空间 $X\subset E$
使得 $M+X=E$, 则称 $M$ 具有\emph{有限余维}.我们总可以假设 $M\cap X=\{0\}$
(否则在 $X$ 中选取 $M\cap X$ 的补空间).$M$ 的余维数按定义是这样的 $X$ 的
维数, 记作 $\operatorname{codim}M$ (它与 $X$ 的具体选取无关); 它与
$\dim(E/M)$ 一致.

[警告: 有限余维的子空间不必是闭的.例如, 若 $\dim E=\infty$, 取 $E$ 上的
任何不连续的线性泛函 $f$ (见习题 1.5).则 $M=f^{-1}(\{0\})$ 余维数为 1,
但 $M$ 不是闭的 (由命题  \ref{prop:1.5}); 事实上, $M$ 在 $E$ 中稠密.]

\begin{proposition}
\label{prop:11.5}
设 $E$ 为 Banach 空间, $M$ 为 $E$ 的有限余维的闭子空间.则包含 $M$ 的
任何子空间 $\widetilde M\subset E$ 必为闭的.
\end{proposition}

\begin{proof}
空间 $M$ 在 $\widetilde M$ 中有代数补空间, 记作 $X$.显然 $\dim X<\infty$,
且 $\widetilde M=X+M$.应用命题  \ref{prop:11.4}, 可知 $\widetilde M$ 是闭的.
\end{proof}

\begin{proposition}
\label{prop:11.6}
设 $E$ 为 Banach 空间, $M$ 为 $E$ 的有限余维的闭子空间.设 $D$ 为 $E$ 的
稠密子空间.则存在 $M$ 的补空间 $X$ 满足 $X\subset D$.
\end{proposition}

\begin{proof}
设 $d$ 为 $M$ 在 $E$ 中的余维数.若 $d=0$, 则 $M=E$, 可以取 $X=\{0\}$.
因此可以假设 $d\ge1$.固定任何满足 $x_1\notin M$ 的 $x_1\in D$; 这是可能
的, 否则 $D\subset M$ 蕴含 $E=\overline D\subset M\ne E$, 矛盾.令
$M_1=M+\R x_1$.则 $M_1$ 是闭的 (由命题  \ref{prop:11.4}), 且
$\operatorname{codim}M_1=d-1$.重复这一构造 $(d-1)$ 次, 即得到维数为 $d$
的子空间 $X\subset D$, 使得 $M+X=E$ 且 $M\cap X=\{0\}$.
\end{proof}

\begin{proposition}
\label{prop:11.7}
设 $E$ 为 Banach 空间, $G,L\subset E$ 为闭子空间.假设存在有限维空间
$X_1,X_2\subset E$ 使得
\begin{equation}
G+L+X_1=E,
\label{eq:11.1}
\end{equation}
\begin{equation}
G\cap L\subset X_2.
\label{eq:11.2}
\end{equation}
则 $G$ (相应地 $L$) 有补空间.
\end{proposition}

\begin{proof}
我们把证明分成两步.

\emph{第 1 步: 当 $X_2=\{0\}$ 时命题  \ref{prop:11.7} 的结论成立.}

设 $\widetilde X_1$ 为 $(G+L)\cap X_1$ 在 $X_1$ 中的补空间.由命题
 \ref{prop:11.4} 我们已知 $(L+\widetilde X_1)$ 是闭的.我们断言
$(L+\widetilde X_1)$ 是 $G$ 的补空间.

首先, 我们有
\begin{equation*}
G+(L+\widetilde X_1)=E.
\end{equation*}
事实上, 任何 $\xi\in E$ 可以写成 $\xi=g+\ell+h$, 其中 $g\in G$, $\ell\in L$,
$h\in X_1$, 而 $h$ 可以进一步分解为 $h=h_1+h_2$, 其中
$h_1\in(G+L)\cap X_1$, $h_2\in\widetilde X_1$.因此
$\xi\in G+L+\widetilde X_1$.

其次, 我们有
\begin{equation*}
G\cap(L+\widetilde X_1)=\{0\}.
\end{equation*}
事实上, 假设 $g=\ell+\tilde x_1$, 其中 $g\in G$, $\ell\in L$,
$\tilde x_1\in\widetilde X_1$.则 $\tilde x_1=g-\ell$, 从而
$\tilde x_1\in(G+L)\cap\widetilde X_1=\{0\}$.因此
$g=\ell\in G\cap L=\{0\}$ (这是第 1 步中的假设).

\emph{第 2 步: 一般情形.}

设 $\widetilde G$ 为 $(G\cap L)$ 在 $G$ 中的补空间, $\widetilde L$ 为
$(G\cap L)$ 在 $L$ 中的补空间 (注意, 由于 $G\cap L$ 有限维,
$\widetilde G$ 与 $\widetilde L$ 存在; 见 \S \ref{sec:2.4}).我们断言
\begin{equation}
(\widetilde G+\widetilde L)+(X_1+X_2)=E
\label{eq:11.3}
\end{equation}
以及
\begin{equation}
(\widetilde G\cap\widetilde L)=\{0\}.
\label{eq:11.4}
\end{equation}
这将完成本命题的证明.事实上, 由第 1 步推出 $\widetilde G$ 有补空间.
因此由命题  \ref{prop:11.4}, $G=\widetilde G+(G\cap L)$ 也有补空间.

\emph{\eqref{eq:11.3} 的验证.} 任何 $\xi\in E$ 可以写成
\begin{equation*}
\xi=g+\ell+x_1,\quad g\in G,\ \ell\in L,\ x_1\in X_1.
\end{equation*}
但 $g=\tilde g+h_1$, 其中 $\tilde g\in\widetilde G$, $h_1\in G\cap L$;
同理 $\ell=\tilde\ell+h_2$, 其中 $\tilde\ell\in\widetilde L$,
$h_2\in G\cap L$.因此
\begin{equation*}
\xi=(\tilde g+\tilde\ell)+x_1+(h_1+h_2)\in(\widetilde G+\widetilde L)+(X_1+X_2).
\end{equation*}

\emph{\eqref{eq:11.4} 的验证.} 设 $g\in\widetilde G\cap\widetilde L$.则
$g\in(G\cap L)\cap\widetilde L$ (因为 $\widetilde G\subset G$ 且
$\widetilde L\subset L$).但 $(G\cap L)\cap\widetilde L=\{0\}$.
\end{proof}
\section{商空间}
\label{sec:11.2}

设 $E$ 为 Banach 空间, $M$ 为闭子空间.我们在 $E$ 上考虑由 $x\sim y$
当且仅当 $x-y\in M$ 定义的等价关系.全体等价类的集合是一个向量空间, 记作
$E/M$, \index{shangkongjian@商空间 (quotient space)}
称为 $E$ (模 $M$) 的\emph{商空间}.把每个 $x\in E$ 映到其等价类
$[x]$ 的典范映射记作 $\pi:E\to E/M$.显然 $\pi$ 是满射线性算子.商空间
$E/M$ \index{shangfanshu@商范数 (quotient norm)}
配以\emph{商范数}
\begin{equation*}
\norm{[x]}_{E/M}=\norm{\pi(x)}_{E/M}
=\inf_{\substack{y\in E\\ y\in[x]}}\norm{y}
=\inf_{m\in M}\norm{x-m}.
\end{equation*}
容易看出 $\norm{[x]}_{E/M}$ 是 $E/M$ 上的范数 (为验证
$\norm{[x]}_{E/M}=0\Rightarrow[x]=0$, 须用到 $M$ 闭的事实).此外,
$\pi:E\to E/M$ 是有界算子, 且 $\norm{\pi}\le1$.若无混淆, 我们简写
$\norm{\cdot}$ 而不写 $\norm{\cdot}_{E/M}$.

\begin{proposition}
\label{prop:11.8}
配以范数 $\norm{\cdot}_{E/M}$ 的商空间 $E/M$ 是 Banach 空间.
\end{proposition}

\begin{proof}
设 $(\pi(x_k))$ 为 $E/M$ 中的 Cauchy 序列.我们须证明 $(\pi(x_k))$ 收敛;
由于 $(\pi(x_k))$ 是 Cauchy 序列, 只须证明某个子序列收敛.抽取子序列 (仍
记作 $(x_k)$), 可以假设对所有 $k$,
$\norm{\pi(x_{k+1})-\pi(x_k)}<\dfrac{1}{2^k}$ (见定理  \ref{thm:4.8} 的
证明).因此存在 $M$ 中的序列 $(m_k)$ 使得
$\norm{x_{k+1}-x_k-m_k}<\dfrac{1}{2^k}$.令 $m_k=\mu_{k+1}-\mu_k$, 其中
$\mu_1=0$, 且 $\mu_k\in M\ \forall k$.由于 $(x_k-\mu_k)$ 是 $E$ 中的
Cauchy 序列, 它在 $E$ 中收敛于极限 $\ell$.因此
$\pi(x_k)=\pi(x_k-\mu_k)$ 也在 $E/M$ 中收敛 (收敛于 $\pi(\ell)$).
\end{proof}

\begin{proposition}
\label{prop:11.9}
设 $M$ 为 $E$ 的闭子空间, 设 $\pi^\star:(E/M)^\star\to E^\star$ 为
$\pi:E\to E/M$ 的对偶算子.则 $R(\pi^\star)=M^\perp$, 更精确地说,
$\pi^\star$ 是由 $(E/M)^\star$ 到 $M^\perp$ 上的双射, 满足
\begin{equation*}
\norm{\pi^\star(\xi)}_{E^\star}=\norm{\xi}_{(E/M)^\star}
\quad\forall\xi\in(E/M)^\star.
\end{equation*}
特别地, $(E/M)^\star$ 与 $M^\perp$ 同构且等距.
\end{proposition}

\begin{proof}
设 $\xi\in(E/M)^\star$ 且 $x\in E$, 记
\begin{equation*}
\dual{\pi^\star(\xi)}{x}=\dual{\xi}{\pi(x)}.
\end{equation*}
若 $x\in M$, 则 $\pi(x)=0$, 从而 $\dual{\pi^\star(\xi)}{x}=0\ \forall x\in M$,
即 $\pi^\star(\xi)\in M^\perp$.

反之, 设 $f\in M^\perp$; 我们须证明对某个 $\xi\in(E/M)^\star$, 有
$f=\pi^\star(\xi)$.给定 $y\in E/M$, 把 $y$ 写成 $y=\pi(x)$, 其中 $x\in E$,
然后定义 $\xi(y)=\dual{f}{x}$.注意这个定义不依赖于 $x$ 的具体选取, 因为
$f\in M^\perp$.显然 $\xi$ 关于 $y$ 线性, 且对一切 $m\in M$, 有
$\abs{\xi(y)}\le\norm{f}_{E^\star}\norm{x-m}$.对所有 $m\in M$ 取下确界,
给出
$\abs{\xi(y)}\le\norm{f}_{E^\star}\norm{\pi(x)}_{E/M}
=\norm{f}_{E^\star}\norm{y}_{E/M}$.因此 $\xi\in(E/M)^\star$, 且显然
$\dual{\pi^\star(\xi)}{x}=\dual{\xi}{\pi(x)}=\dual{f}{x}\ \forall x\in E$,
即 $\pi^\star(\xi)=f$.此外,
$\norm{\xi}_{(E/M)^\star}\le\norm{f}_{E^\star}=\norm{\pi^\star(\xi)}_{E^\star}$.

另一方面, 我们知道
\begin{equation*}
\norm{\pi^\star(\xi)}_{E^\star}\le\norm{\xi}_{(E/M)^\star},
\end{equation*}
因为 $\norm{\pi^\star}=\norm{\pi}\le1$.从而
\begin{equation*}
\norm{\pi^\star(\xi)}_{E^\star}=\norm{\xi}_{(E/M)^\star}
\quad\forall\xi\in(E/M)^\star.
\end{equation*}
\end{proof}

设 $F,G$ 为 Banach 空间, $T\in\LL(F,G)$.考虑 $F$ 的闭子空间 $N(T)$,
商空间 $F/N(T)$, 以及典范映射 $\pi:F\to F/N(T)$.算子 $T$ 可以分解为
$T=\widetilde T\circ\pi$, 其中 $\widetilde T:F/N(T)\to G$; 事实上, 给定
$y\in F/N(T)$, 把 $y$ 写成 $y=\pi(x)$, 其中 $x\in F$, 并令
$\widetilde Ty=Tx$.显然 $\widetilde T$ 良定 (不依赖于 $x$ 的选取), 并且
是由 $F/N(T)$ 到 $R(T)$ 上的双射; 此外, $\norm{\widetilde T}=\norm{T}$.

现在考虑这个框架的一个特殊情形.设 $M$ 为 Banach 空间 $E$ 的闭子空间.
设 $T:E^\star\to M^\star$ 定义为
\begin{equation*}
T(f)=f|_M\quad\forall f\in E^\star.
\end{equation*}
则 $N(T)=M^\perp$, 且 (由 Hahn--Banach 定理) $R(T)=M^\star$.把上述结论
应用于 $F=E^\star$ 与 $G=M^\star$, 我们得到双射算子
$\widetilde T:E^\star/M^\perp\to M^\star$, 使得 $\widetilde T\circ\pi=T$.

\begin{proposition}
\label{prop:11.10}
对任何 Banach 空间 $E$ 与 $E$ 的任何闭子空间 $M$, 算子 $\widetilde T$ 是由
$E^\star/M^\perp$ 到 $M^\star$ 上的双射等距.
\end{proposition}

\begin{proof}
只须证明 $\widetilde T$ 是等距.给定任何 $f\in E^\star$, 考虑 $M$ 上的泛函
$f|_M$.由推论  \ref{cor:1.2} 我们知道, 存在泛函 $\tilde f\in E^\star$
使得 $\tilde f|_M=f|_M$ 且
$\norm{\tilde f}_{E^\star}=\norm{f|_M}_{M^\star}=\norm{T(f)}_{M^\star}$.
由于 $f-\tilde f\in M^\perp$, 我们有
\begin{equation*}
\norm{\pi(f)}_{E^\star/M^\perp}=\dist(f,M^\perp)
\le\norm{f-(f-\tilde f)}_{E^\star}=\norm{\tilde f}_{E^\star}
=\norm{Tf}_{M^\star}.
\end{equation*}
因此我们证明了
\begin{equation*}
\norm{\pi(f)}_{E^\star/M^\perp}\le\norm{Tf}_{M^\star}\quad\forall f\in E^\star.
\end{equation*}
但 $T=\widetilde T\circ\pi$, 从而
\begin{equation*}
\norm{\pi(f)}_{E^\star/M^\perp}\le\norm{\widetilde T(\pi(f))}_{M^\star}
\quad\forall f\in E^\star,
\end{equation*}
即
\begin{equation*}
\norm{y}_{E^\star/M^\perp}\le\norm{\widetilde T(y)}_{M^\star}
\quad\forall y\in E^\star/M^\perp.
\end{equation*}
另一方面, 显然有
\begin{equation*}
\norm{(\widetilde T\circ\pi)(f)}_{M^\star}=\norm{T(f)}_{M^\star}\le\norm{f}_{E^\star}
\quad\forall f\in E^\star.
\end{equation*}
把 $f$ 换成 $(f-g)$ (其中 $g\in M^\perp$) 并在 $g\in M^\perp$ 上取下确界,
得到
\begin{equation*}
\norm{\widetilde T(\pi(f))}_{M^\star}\le\norm{\pi(f)}_{E^\star/M^\perp}
\quad\forall f\in E^\star,
\end{equation*}
即
\begin{equation*}
\norm{\widetilde T(y)}_{M^\star}\le\norm{y}_{E^\star/M^\perp}
\quad\forall y\in E^\star/M^\perp.
\end{equation*}
由此得出 $\widetilde T$ 是等距.
\end{proof}

商空间 $E/M$ 继承了空间 $E$ 的许多性质, 例如自反性与一致凸性.

\begin{proposition}
\label{prop:11.11}
假设 $E$ 为自反 Banach 空间, $M$ 为闭子空间.则 $E/M$ 自反.
\end{proposition}

\begin{proof}
我们知道 $E^\star$ 自反 (见推论  \ref{cor:3.21}), 从而 $M^\perp$ 也自反
(作为 $E^\star$ 的闭子空间; 见命题  \ref{prop:3.20}).另一方面, $M^\perp$
同构于 $(E/M)^\star$ (由命题  \ref{prop:11.9}).因此 $(E/M)^\star$ 自反,
再次由推论  \ref{cor:3.21}, $E/M$ 也自反.
\end{proof}

\begin{proposition}
\label{prop:11.12}
假设 $E$ 为一致凸 Banach 空间, $M$ 为闭子空间.则 $E/M$ 一致凸.
\end{proposition}

\begin{proof}
设 $\pi(x),\pi(y)\in E/M$ 满足 $\norm{\pi(x)}\le1$, $\norm{\pi(y)}\le1$,
且 $\norm{\pi(x)-\pi(y)}>\varepsilon$.由于 $E$ 自反, 我们知道 (见推论
 \ref{cor:3.23}) 存在 $m_1\in M$ 与 $m_2\in M$ 使得
$\norm{x-m_1}\le1$ 且 $\norm{y-m_2}\le1$.此外, 对所有 $m\in M$,
$\norm{(x-y)-m}>\varepsilon$.由 $E$ 的一致凸性推出
\begin{equation*}
\norm{\frac{(x-m_1)+(y-m_2)}{2}}<1-\delta,
\end{equation*}
从而
\begin{equation*}
\norm{\frac{\pi(x)+\pi(y)}{2}}<1-\delta.
\end{equation*}
\end{proof}

\begin{proposition}
\label{prop:11.13}
设 $E$ 为 Banach 空间, $M\subset E$ 为闭子空间.则

(a) $\dim M<\infty$ 当且仅当 $\operatorname{codim}M^\perp<\infty$, 且在这时
$\dim M=\operatorname{codim}M^\perp$,

(b) $\operatorname{codim}M<\infty$ 当且仅当 $\dim M^\perp<\infty$, 且在这时
$\operatorname{codim}M=\dim M^\perp$.
\end{proposition}

\begin{proof}
(a) 由命题  \ref{prop:11.10} 我们知道, $E^\star/M^\perp$ 总是与 $M^\star$
同构.因此
$\dim M^\star<\infty\iff\dim(E^\star/M^\perp)<\infty$.由命题
 \ref{prop:11.3} 我们知道,
$\dim M<\infty\iff\dim M^\star<\infty$ 且此时 $\dim M=\dim M^\star$.
另一方面,
$\dim(E^\star/M^\perp)<\infty\iff\operatorname{codim}M^\perp<\infty$.
因此 $\dim M<\infty\iff\operatorname{codim}M^\perp<\infty$ 且
$\dim M=\dim M^\star=\dim(E^\star/M^\perp)=\operatorname{codim}M^\perp$.

(b) 命题  \ref{prop:11.9} 给出
$\dim M^\perp<\infty\iff\dim(E/M)^\star<\infty$.再次利用命题
 \ref{prop:11.3}, 这等价于 $\dim(E/M)<\infty$, 即
$\operatorname{codim}M<\infty$.此时
$\dim M^\perp=\dim(E/M)^\star=\dim(E/M)=\operatorname{codim}M$.
\end{proof}

一个``对偶''断言部分成立.

\begin{proposition}
\label{prop:11.14}
设 $N\subset E^\star$ 为闭子空间.则 $\dim N<\infty$ 当且仅当
$\operatorname{codim}N^\perp<\infty$, 且在这时
$\dim N=\operatorname{codim}N^\perp$.还成立
$\dim N^\perp\le\operatorname{codim}N$, 但可能出现
$\dim N^\perp<\operatorname{codim}N<\infty$ 的情形.
\end{proposition}

\begin{proof}
回忆
\begin{equation*}
N^\perp=\{x\in E;\ \dual{f}{x}=0\ \forall f\in N\}.
\end{equation*}
显然 $N\subset N^{\perp\perp}$; 但可能出现 $N\ne N^{\perp\perp}$ (见第 1
章注记 6).例如, 取 $\xi\in E^{\star\star}\setminus E$, 并令
$N=\xi^{-1}(\{0\})=\{f\in E^\star;\ \dual{\xi}{f}=0\}$.则 $N$ 是 $E^\star$
的余维数为 1 的闭子空间 (即 $N$ 是超平面).然而, $N^\perp=\{0\}$ (因为
由引理  \ref{lem:3.2}, $N$ 在 $E^{\star\star}$ 中的正交补为 $\R\xi$, 从而
$N$ 在 $E$ 中的正交补 $N^\perp$ 简化为 $\{0\}$).在这种情形下,
$N^\perp=\{0\}$, 从而 $N^{\perp\perp}=E^\star$, 且 $\dim N^\perp=0$,
而 $\operatorname{codim}N=1$.

现在我们回到一般情形.由于 $N\subset N^{\perp\perp}$, 我们有
$\operatorname{codim}N^{\perp\perp}\le\operatorname{codim}N\le\infty$.令
$M=N^\perp\subset E$.由命题  \ref{prop:11.10} 我们有
\begin{equation*}
\operatorname{codim}M^\perp=\dim(E^\star/M^\perp)=\dim M^\star,
\end{equation*}
从而 $\operatorname{codim}N^{\perp\perp}=\dim M\le\infty$.因此
\begin{equation*}
\dim N^\perp\le\operatorname{codim}N\le\infty.
\end{equation*}

现在证明 $\dim N<\infty\Rightarrow\operatorname{codim}N^\perp<\infty$ 且
$\operatorname{codim}N^\perp=\dim N$.

我们先断言 $N^{\perp\perp}=N$.我们已知 $N\subset N^{\perp\perp}$.设
$f_1,f_2,\dots,f_p$ 为 $N$ 的基, 设 $f\in N^{\perp\perp}$.由于在
$N^\perp=\{x\in E;\ \dual{f_i}{x}=0\ \forall i\}$ 上 $f=0$, 我们可以应用
引理  \ref{lem:3.2} 并得出 $f=\sum\lambda_if_i$.因此
$N^{\perp\perp}\subset N$.如上, 令 $M=N^\perp$.由于 $\dim M^\perp<\infty$,
由命题  \ref{prop:11.13} 我们推出 $\operatorname{codim}M<\infty$ 且
$\operatorname{codim}M=\dim M^\perp$, 即
$\operatorname{codim}N^\perp=\dim N$.

反之, 假设 $\operatorname{codim}N^\perp<\infty$, 再次令 $M=N^\perp$, 于是
$\operatorname{codim}M<\infty$.再次应用命题  \ref{prop:11.13} 给出
$\dim M^\perp<\infty$, 即 $\dim N^{\perp\perp}<\infty$.由于
$N\subset N^{\perp\perp}$, 我们推出 $\dim N<\infty$, 这回到了前面的情形.
因此 $\dim N=\operatorname{codim}N^\perp$.
\end{proof}
\section{若干经典的序列空间}
\label{sec:11.3}

给定序列 $x=(x_1,x_2,\dots,x_k,\dots)$, 令
\begin{align*}
\norm{x}_p&=\left(\sum_{k=1}^\infty\abs{x_k}^p\right)^{1/p},\quad 1\le p<\infty,\\
\norm{x}_\infty&=\sup_k\abs{x_k},
\end{align*}
并考虑相应的空间
\begin{align*}
\ell^p&=\{x;\ \norm{x}_p<\infty\},\quad 1\le p<\infty,\\
\ell^\infty&=\{x;\ \norm{x}_\infty<\infty\},
\end{align*}
它们关于 $\ell^p$ (相应地 $\ell^\infty$) 范数是 Banach 空间.这可以直接
证明 (且相当容易); 也可以应用定理  \ref{thm:4.8} 于配以计数测度
$\mu(E)=E\subset\N$ 中的点数的 $\Omega=\N$.下面提到的许多性质都是第~
\ref{chap:4}~章
中一般结果的推论.为了方便读者, 我们也给出一些直接的证明.

$\ell^\infty$ 有两个有趣的子空间:
\begin{align*}
c&=\left\{x;\ \lim_{k\to\infty}x_k\ \text{存在}\right\},\\
c_0&=\left\{x;\ \lim_{k\to\infty}x_k=0\right\}.
\end{align*}
它们都配以 $\ell^\infty$ 范数.显然 $c_0\subset c\subset\ell^\infty$, 其中
$c_0$ 在 $c$ 中闭, $c$ 在 $\ell^\infty$ 中闭.

Hölder 不等式采取如下形式:
\begin{equation}
\left|\sum_{k=1}^\infty x_ky_k\right|\le\norm{x}_p\norm{y}_{p'}
\quad\forall x\in\ell^p,\ \forall y\in\ell^{p'}\ \text{满足}\
\frac1p+\frac{1}{p'}=1.
\label{eq:11.5}
\end{equation}
空间 $\ell^2$ 是配以标量积
\begin{equation*}
(x,y)=\sum_{k=1}^\infty x_ky_k
\end{equation*}
的 Hilbert 空间.显然 $\ell^p\subset c_0$, 满足
\begin{equation*}
\norm{x}_\infty\le\norm{x}_p\quad\forall p,\ 1\le p<\infty,\ \forall x\in\ell^p,
\end{equation*}
这给出当 $1\le p\le q\le\infty$ 时 $\ell^p\subset\ell^q$, 满足
\begin{equation*}
\norm{x}_q\le\norm{x}_p\quad\forall x\in\ell^p.
\end{equation*}

\begin{proposition}
\label{prop:11.15}
对 $1<p<\infty$, 空间 $\ell^p$ 是自反的, 甚至是一致凸的.
\end{proposition}

\begin{proof}
应用定理  \ref{thm:4.10} 与习题 4.12 于 $\Omega=\N$.
\end{proof}

\begin{proposition}
\label{prop:11.16}
空间 $c$, $c_0$ 以及 $\ell^p$ ($1\le p<\infty$) 是可分的.
\end{proposition}

\begin{proof}
令
\begin{equation*}
D=\{x=(x_k);\ x_k\in\Q\ \forall k,\ \text{且对足够大的}\ k\ \text{有}\ x_k=0\}.
\end{equation*}
显然 $D$ 是可数的; 此外, 当 $1\le p<\infty$ 时 $D$ 在 $\ell^p$ 中稠密, 且
在 $c_0$ 中稠密.集合 $D+\lambda(1,1,1,\dots)$ ($\lambda\in\Q$) 是可数的,
且在 $c$ 中稠密.
\end{proof}

\begin{proposition}
\label{prop:11.17}
空间 $\ell^\infty$ 不可分.
\end{proposition}

\begin{proof}
假设 $A\subset\ell^\infty$ 是可数的.我们将验证 $A$ 不能在 $\ell^\infty$
中稠密.记 $A=(a_k)$, 其中每个 $a_k\in\ell^\infty$, 从而
$a_k=(a_k^1,a_k^2,\dots)$.对每个整数 $k$, 令
\begin{equation*}
b_k=\begin{cases}
a_k^k+1 & \text{若}\ \abs{a_k^k}\le1,\\
0 & \text{若}\ \abs{a_k^k}>1.
\end{cases}
\end{equation*}
注意 $b=(b_k)\in\ell^\infty$ 且 $\abs{b_k-a_k^k}\ge1\ \forall k$.因此,
\begin{equation*}
\norm{b-a_k}_\infty\ge\abs{b_k-a_k^k}\ge1\quad\forall k,
\end{equation*}
从而 $b\notin\overline A$.
\end{proof}

\begin{proposition}
\label{prop:11.18}
设 $1\le p<\infty$.给定任何 $\varphi\in(\ell^p)^\star$, 存在唯一的
$u\in\ell^{p'}$ 使得
\begin{equation*}
\dual{\varphi}{x}=\sum_{k=1}^\infty u_kx_k\quad\forall x\in\ell^p.
\end{equation*}
此外,
\begin{equation*}
\norm{u}_{p'}=\norm{\varphi}_{(\ell^p)^\star}.
\end{equation*}
\end{proposition}

\begin{proof}
令 $e_k=(0,0,\dots,\underset{(k)}{1},0,0,\dots)$.令 $u_k=\varphi(e_k)$.
我们断言 $u=(u_k)\in\ell^{p'}$ 且
\begin{equation}
\norm{u}_{p'}\le\norm{\varphi}_{(\ell^p)^\star}.
\label{eq:11.6}
\end{equation}
当 $p=1$ 时不等式 \eqref{eq:11.6} 是显然的, 因为
\begin{equation*}
\abs{u_k}\le\norm{\varphi}_{(\ell^1)^\star}\norm{e_k}_1
\le\norm{\varphi}_{(\ell^1)^\star}\quad\forall k.
\end{equation*}
现在转向 $1<p<\infty$ 的情形.固定整数 $N$.则对每个
$x=(x_1,x_2,\dots,x_N,0,0,\dots)$, 我们有
\begin{equation}
\sum_{k=1}^N u_kx_k=\varphi\left(\sum_{k=1}^N x_ke_k\right)
\le\norm{\varphi}_{(\ell^p)^\star}\norm{x}_p.
\label{eq:11.7}
\end{equation}
选取 $x_k=\abs{u_k}^{p'-2}u_k$ 给出
\begin{equation*}
\left(\sum_{k=1}^N\abs{u_k}^{p'}\right)^{1/p'}
\le\norm{\varphi}_{(\ell^p)^\star}.
\end{equation*}
令 $N\to\infty$, 可见 $u\in\ell^{p'}$ 且 \eqref{eq:11.6} 成立.此外,
\begin{equation*}
\varphi(x)=\sum_{k=1}^\infty u_kx_k\quad\forall x\in D,
\end{equation*}
其中 $D$ 是命题  \ref{prop:11.16} 的证明中定义的.由于 $D$ 在 $\ell^p$
中稠密, 我们得到
\begin{equation*}
\varphi(x)=\sum_{k=1}^\infty u_kx_k\quad\forall x\in\ell^p.
\end{equation*}
由 Hölder 不等式得出
\begin{equation*}
\abs{\varphi(x)}\le\norm{u}_{p'}\norm{x}_p\quad\forall x\in\ell^p,
\end{equation*}
从而 $\norm{\varphi}_{(\ell^p)^\star}\le\norm{u}_{p'}$.结合 \eqref{eq:11.6},
我们得到 $\norm{\varphi}_{(\ell^p)^\star}=\norm{u}_{p'}$.

$u$ 的唯一性是显然的.
\end{proof}

\begin{proposition}
\label{prop:11.19}
给定任何 $\varphi\in(c_0)^\star$, 存在唯一的 $u\in\ell^1$ 使得
\begin{equation*}
\dual{\varphi}{x}=\sum_{k=1}^\infty u_kx_k\quad\forall x\in c_0.
\end{equation*}
此外,
\begin{equation*}
\norm{u}_1=\norm{\varphi}_{(c_0)^\star}.
\end{equation*}
\end{proposition}

\begin{proof}
这是命题  \ref{prop:11.18} 的证明的一个容易的改编 (取 $p=\infty$ 且
$p'=1$); 证明的最后部分成立是因为 $D$ 在 $c_0$ 中稠密 (但在 $\ell^\infty$
中不稠密).
\end{proof}

\begin{proposition}
\label{prop:11.20}
给定 $\varphi\in(c)^\star$, 存在唯一的对 $(u,\lambda)\in\ell^1\times\R$
使得
\begin{equation*}
\dual{\varphi}{x}=\sum_{k=1}^\infty u_kx_k+\lambda\lim_{k\to\infty}x_k
\quad\forall x\in c.
\end{equation*}
此外,
\begin{equation*}
\norm{u}_1+\abs{\lambda}=\norm{\varphi}_{(c)^\star}.
\end{equation*}
\end{proposition}

\begin{proof}
把命题  \ref{prop:11.19} 应用于 $\varphi|_{c_0}$, 我们找到某个
$u\in\ell^1$ 使得
\begin{equation*}
\varphi(y)=\sum_{k=1}^\infty u_ky_k\quad\forall y\in c_0.
\end{equation*}
若 $x\in c$, 把 $x$ 写成 $x=y+ae$, 其中 $e=(1,1,1,\dots)$,
$a=\lim_{k\to\infty}x_k$, 且 $y\in c_0$.则
\begin{equation*}
\varphi(x)=\sum_{k=1}^\infty u_ky_k+a\varphi(e)
=\sum_{k=1}^\infty u_k(x_k-a)+a\varphi(e)
=\sum_{k=1}^\infty u_kx_k+\lambda a,
\end{equation*}
其中 $\lambda=\varphi(e)-\sum_{k=1}^\infty u_k$.

反之, 给定任何 $u\in\ell^1$ 与 $\lambda\in\R$, 泛函
\begin{equation}
\varphi(x)=\sum_{k=1}^\infty u_kx_k+\lambda\lim_{k\to\infty}x_k,
\quad x\in c,
\label{eq:11.8}
\end{equation}
定义了 $(c)^\star$ 的一个元素.我们断言
\begin{equation}
\norm{\varphi}_{(c)^\star}=\norm{u}_1+\abs{\lambda}.
\label{eq:11.9}
\end{equation}
显然
\begin{equation}
\norm{\varphi}_{(c)^\star}\le\norm{u}_1+\abs{\lambda}.
\label{eq:11.10}
\end{equation}
在 \eqref{eq:11.8} 中选取 $x=(x_k)$, 其中 $N$ 为固定的整数, 且
\begin{equation*}
x_k=\begin{cases}
\operatorname{sign}(u_k), & 1\le k\le N,\\
\operatorname{sign}(\lambda), & k>N,
\end{cases}
\end{equation*}
给出
\begin{equation*}
\varphi(x)=\sum_{k=1}^N\abs{u_k}+\operatorname{sign}(\lambda)
\sum_{k=N+1}^\infty u_k+\abs{\lambda}\le\norm{\varphi}_{(c)^\star}.
\end{equation*}
令 $N\to\infty$, 我们得到
\begin{equation*}
\norm{u}_1+\abs{\lambda}\le\norm{\varphi}_{(c)^\star},
\end{equation*}
结合 \eqref{eq:11.10}, 给出 \eqref{eq:11.9}.
\end{proof}

\begin{proposition}
\label{prop:11.21}
空间 $\ell^1$, $\ell^\infty$, $c$ 与 $c_0$ 不是自反的.
\end{proposition}

\begin{proof}
由命题  \ref{prop:11.19} 与  \ref{prop:11.18} 我们知道, $(c_0)^\star$ 是
$\ell^1$, 且 $(\ell^1)^\star$ 是 $\ell^\infty$.因此由 $c_0$ 到 $\ell^\infty$
中的恒等映射对应于 \S \ref{sec:1.3} 中定义的典范单射 $J:c_0\to(c_0)^{\star\star}$.
由于它不是满射, 我们得出 $c_0$ 不是自反的.应用推论  \ref{cor:3.21}, 我们
推出 $\ell^1$ 与 $\ell^\infty$ 也不是自反的.此外, $c$ 不可能是自反的;
否则, 作为 $c$ 的闭子空间的 $c_0$ 由命题  \ref{prop:3.20} 必自反.
\end{proof}

下表总结了上面讨论的主要性质:

\begin{table}[htbp]
	\centering
	\begin{tabular}{lccc}
		\hline
		空间 & 自反性 & 可分性 & 对偶空间 \\
		\hline
		$\ell^p$ ($1<p<\infty$) & 是 & 是 & $\ell^{p'}$ \\
		$\ell^1$ & 否 & 是 & $\ell^\infty$ \\
		$c_0$ & 否 & 是 & $\ell^1$ \\
		$c$ & 否 & 是 & $\ell^1\times\R$ \\
		$\ell^\infty$ & 否 & 否 & 严格大于 $\ell^1$ \\
		\hline
	\end{tabular}
\end{table}
\section{\texorpdfstring{$\CC$ 上的 Banach 空间: 哪些相似, 哪些不同?}{C 上的 Banach 空间: 哪些相似, 哪些不同?}}
\label{sec:11.4}

在本节中, 我们始终假设 $E$ 是 $\CC$ 上的向量空间.当然, 只须考虑把乘标量
$\lambda x$ 中的 $\lambda$ 限制在 $\R$ 中 ($x\in E$), 我们就可以与 $E$
联系一个 $\R$ 上的向量空间; 为了与 $E$ 区别, 相应的 $\R$ 上的向量空间常记作
$E_\R$.

线性子空间 $M\subset E$ 是满足对一切 $\lambda\in\CC$ 及一切 $x,y\in M$ 有
$\lambda x\in M$ 与 $x+y\in M$ 的子集 $M$.当然, $E$ 的线性子空间 $M$ 也是
$E_\R$ 的线性子空间.但逆命题不成立.例如, $\R^2$ 中包含 $0$ 的一条直线
$L$ 是 $\R^2$ 的线性子空间.然而, 若我们把 $\CC$ 与 $\R^2$ 等同起来, 则
直线 $L$ 不再是 $\CC$ 的线性子空间, 因为 $iL$ (即 $L$ 旋转 $\pi/2$) 不包含
在 $L$ 中.

$E$ 上的范数按定义是取值于 $[0,+\infty)$ 的函数 $\norm{\cdot}$, 使得
$\norm{x}=0\iff x=0$, $\norm{\lambda x}=\abs{\lambda}\norm{x}\ \forall\lambda\in\CC,
\ \forall x\in E$, 以及 $\norm{x+y}\le\norm{x}+\norm{y}$.显然
$\norm{\cdot}$ 也是 $E_\R$ 上的范数, 但逆命题不成立.

$E$ 上的线性泛函是映射 $f:E\to\CC$, 满足对一切 $\lambda\in\CC$ 与一切
$x,y\in E$, $f(\lambda x)=\lambda f(x)$ 且 $f(x+y)=f(x)+f(y)$.对偶空间
$E^\star$ 是 $E$ 上的所有连续线性泛函组成的空间; $E^\star$ 是 $\CC$ 上的
向量空间, 配以范数
\begin{equation*}
\norm{f}_{E^\star}=\sup_{\substack{x\in E\\\norm{x}\le1}}\abs{f(x)}.
\end{equation*}
复数 $f(x)$ 也记作 $\dual{f}{x}$, 显然对一切 $\lambda,\mu\in\CC$ 与一切
$x\in E$, 有 $\dual{\lambda f}{\mu x}=\lambda\mu\dual{f}{x}$.复对偶
$E^\star$ 与实对偶 $E^\star_\R$ 之间的对应由下述简单而富有启发性的结果给出.

\begin{proposition}
\label{prop:11.22}
映射
\begin{equation*}
I:f\in E^\star\mapsto\operatorname{Re}f\in E^\star_\R
\end{equation*}
是由 $E^\star$ 到 $E^\star_\R$ 上的双射等距.
\end{proposition}

\begin{proof}
显然,
\begin{equation*}
\abs{\operatorname{Re}\dual{f}{x}}\le\abs{\dual{f}{x}}\le\norm{f}_{E^\star}\norm{x},
\end{equation*}
从而
\begin{equation}
\norm{I(f)}_{E^\star_\R}\le\norm{f}_{E^\star}.
\label{eq:11.11}
\end{equation}
$I$ 是单射也是显然的, 因为 $\operatorname{Re}\dual{f}{x}=0\ \forall x\in E$
蕴含 $\operatorname{Re}\dual{f}{ix}=0\ \forall x\in E$, 即
$\operatorname{Im}\dual{f}{x}=0\ \forall x\in E$, 从而 $f=0$.接下来我们
断言 $I$ 是满射.事实上, 给定 $\varphi\in E^\star_\R$, 令
\begin{equation}
f(x)=\varphi(x)-i\varphi(ix)\quad\forall x\in E
\label{eq:11.12}
\end{equation}
[警告: $\varphi(ix)$ 不等于 $i\varphi(x)$, 因为 $\varphi(x)$ 与
$\varphi(ix)$ 都属于 $\R$].容易验证 $f\in E^\star$, 即对一切
$\lambda\in\CC$ 与一切 $x\in E$, $f(\lambda x)=\lambda f(x)$ (请读者验证!),
且 $I(f)=\operatorname{Re}f=\varphi$.由 \eqref{eq:11.11} 我们有
$\norm{\varphi}_{E^\star_\R}\le\norm{f}_{E^\star}$.由 \eqref{eq:11.12}
还清楚地看出
\begin{equation*}
\abs{f(x)}\le\left(\abs{\varphi(x)}^2+\abs{\varphi(ix)}^2\right)^{1/2}
\le\sqrt{2}\norm{\varphi}_{E^\star_\R}\norm{x}
\end{equation*}
(因为 $\norm{ix}=\norm{x}$).但我们可以做得更好.设 $f(x)\ne0$, 并令
$\lambda=\dfrac{f(x)}{\abs{f(x)}}\in\CC$.则
\begin{equation*}
\abs{f(x)}=\frac1\lambda f(x)=f\left(\frac{x}{\lambda}\right)
=\varphi\left(\frac{x}{\lambda}\right)-i\varphi\left(\frac{ix}{\lambda}\right).
\end{equation*}
由于 $\abs{f(x)}\in\R$, $\varphi\left(\dfrac{x}{\lambda}\right)\in\R$, 且
$\varphi\left(\dfrac{ix}{\lambda}\right)\in\R$, 可见
$\varphi\left(\dfrac{ix}{\lambda}\right)=0$, 从而
$\abs{f(x)}=\varphi\left(\dfrac{x}{\lambda}\right)$.因此
\begin{equation*}
\abs{f(x)}\le\norm{\varphi}_{E^\star_\R}\norm{\frac{x}{\lambda}}
=\frac1{\abs{\lambda}}\norm{\varphi}_{E^\star_\R}\norm{x}
=\norm{\varphi}_{E^\star_\R}\norm{x}.
\end{equation*}
从而 $\norm{f}_{E^\star}\le\norm{\varphi}_{E^\star_\R}
=\norm{I(f)}_{E^\star_\R}$.结合 \eqref{eq:11.11}, 我们得出 $I$ 是等距.
\end{proof}

命题  \ref{prop:11.22} 意味着: 当我们处理 $\CC$ 上的向量空间时, 第 1--5 章
中几乎没有任何改变, 只不过我们需要对 Hahn--Banach 定理稍加小心 (见下文).
重大差别出现在第~
\ref{chap:6}~章中我们处理特征值与谱的时候.这在有限维中就已经能看
出来: 任何复系数 $n\times n$ 矩阵 $M$ 在 $\CC$ 中必有特征值; 但即使 $M$
的元素属于 $\R$, 它在 $\R$ 中也可能没有特征值.我们现在逐章叙述需要作出的
修改.

\runin{第~
\ref{chap:1}~章.} 我们挑选几个例子来说明: 某些断言保持不变, 而另一些断言
需要微小的修改.

\begin{proposition}
\label{prop:11.23}
设 $G\subset E$ 为线性子空间.若 $g:G\to\CC$ 为连续线性泛函, 则存在
延拓 $g$ 的 $f\in E^\star$, 使得
\begin{equation*}
\norm{f}_{E^\star}=\norm{g}_{G^\star}.
\end{equation*}
\end{proposition}

\begin{proof}
令 $\psi=\operatorname{Re}g$, 于是 $\psi$ 是 $G^\star_\R$ 的元素, 且
$\norm{\psi}_{G^\star_\R}=\norm{g}_{G^\star}$.由推论  \ref{cor:1.2}, 存在
延拓 $\psi$ 的 $\phi\in E^\star_\R$, 使得
$\norm{\phi}_{E^\star_\R}=\norm{\psi}_{G^\star_\R}$.应用命题
 \ref{prop:11.22}, 可见存在 $f\in E^\star$ 使得
$\phi=\operatorname{Re}f$ 且
$\norm{f}_{E^\star}=\norm{\phi}_{E^\star_\R}=\norm{\psi}_{G^\star_\R}
=\norm{g}_{G^\star}$.此外, 在 $G$ 上有
$\phi=\operatorname{Re}f=\psi=\operatorname{Re}g$, 即对一切 $x\in G$,
$\operatorname{Re}f(x)=\operatorname{Re}g(x)$; 用 $ix$ 替代 $x$ 给出对一切
$x\in G$, $\operatorname{Im}f(x)=\operatorname{Im}g(x)$, 从而在 $G$ 上
$f=g$.
\end{proof}

接下来, 我们叙述 Hahn--Banach 定理的几何形式之一.$E$ 中的\emph{闭实超平面}
$H$ 是形如
\begin{equation*}
H=\{x\in E;\ \operatorname{Re}\dual{f}{x}=\alpha\}
=[\operatorname{Re}f=\alpha]
\end{equation*}
的集合, 其中 $f\in E^\star$, $f\ne0$, 且 $\alpha\in\R$.我们再次提醒读者:
若 $\alpha=0$, 则 $H$ 是 $E_\R$ 的线性子空间, 但它不是 $E$ 在 $\CC$ 上的
线性子空间; 例如, 在 $E=\CC$ 中, $H$ 是一条直线 (而直线不是 $E$ 的复线性
子空间).若对一切 $x\in A$, $\operatorname{Re}\dual{f}{x}\le\alpha$, 且对
一切 $x\in B$, $\operatorname{Re}\dual{f}{x}\ge\alpha$, 则称 $H$ \emph{分离}
$A,B\subset E$.

\begin{proposition}
\label{prop:11.24}
设 $A,B\subset E$ 为 $E$ 的两个非空凸子集, 满足 $A\cap B=\emptyset$.假设
其中之一是开集.则存在分离 $A$ 与 $B$ 的闭实超平面.
\end{proposition}

\begin{proof}
把定理  \ref{thm:1.6} 应用于 $E_\R$, 得到在通常意义下分离 $A$ 与 $B$ 的
超平面 $H=[\phi=\alpha]$ ($\phi\in E^\star_\R$).然后利用命题
 \ref{prop:11.23} 断言对某个 $f\in E^\star$, 有 $\phi=\operatorname{Re}f$.
\end{proof}

$E$ 的线性子空间 $M$ 的正交空间 $M^\perp$ 的定义不变:
\begin{equation*}
M^\perp=\{f\in E^\star;\ \dual{f}{x}=0\ \forall x\in M\},
\end{equation*}
显然我们有 $M^\perp=\{f\in E^\star;\ \operatorname{Re}\dual{f}{x}=0\
\forall x\in M\}$ (因为我们可以用 $ix$ 代替 $x$).容易看出 $M^{\perp\perp}=M$.

给定函数 $\varphi:E\to(-\infty,+\infty]$, 我们在 $E^\star$ 上定义其共轭
函数 $\varphi^\star$:
\begin{equation*}
\varphi^\star(f)=\sup_{x\in E}\{\operatorname{Re}\dual{f}{x}-\varphi(x)\}.
\end{equation*}
采用显而易见的记号, 我们有
\begin{equation*}
\varphi^\star(f)=\varphi^\star_\R(If)\quad\forall f\in E^\star.
\end{equation*}

\begin{proposition}
\label{prop:11.25}
假设 $\varphi:E\to(-\infty,+\infty]$ 是凸的, 下半连续的, 且
$\varphi\not\equiv+\infty$.则 $\varphi^{\star\star}=\varphi$.
\end{proposition}

\begin{proof}
有两种方法.或者把定理  \ref{thm:1.11} 应用于看作 $E_\R$ 上的
$\tilde\varphi=\varphi$, 并结合命题  \ref{prop:11.22}.或者重复定理
 \ref{thm:1.11} 的证明; 当使用 Hahn--Banach 定理时, 如上所述用实超平面
分离凸集 $A$ 与 $B$.
\end{proof}

示性函数 $I_K$ 的定义不变.若 $M$ 是 $E$ 的线性子空间且 $\varphi=I_M$, 则
\begin{equation*}
\varphi^\star(f)=\sup_{x\in M}\operatorname{Re}\dual{f}{x}=I_{M^\perp}.
\end{equation*}
事实上, 若 $f\in M^\perp$, 则对一切 $x\in M$ 有 $\dual{f}{x}=0$, 从而
$\varphi^\star(f)=0$.反之, 若 $f\notin M^\perp$, 则存在某个 $x_0\in M$
使得 $\dual{f}{x_0}\ne0$.必要时用 $ix_0$ 代替 $x_0$, 可以假设
$\operatorname{Re}\dual{f}{x_0}\ne0$.必要时用 $-x_0$ 代替 $x_0$, 可以
假设 $\operatorname{Re}\dual{f}{x_0}>0$, 从而
$\sup_{\lambda>0}\dual{f}{\lambda x_0}=+\infty$.

\runin{第~
\ref{chap:2}~章.} 所有断言均保持不变 (在推论~
\ref{cor:2.4} 与 \ref{cor:2.5} 中把 $\R$ 换成 $\CC$).
某些证明依赖于 $\R$ 的结构 (例如定理~
\ref{thm:2.16} 的证明中的公式 \eqref{eq:2.21}).它们容易
改编到 $\CC$; 另外, 也可以通过把 $\R$ 版本应用于 $E_\R$ 来建立 $\CC$ 上的
断言.

\runin{第~
\ref{chap:3}~章.} 所有断言均保持不变 (在引理~
\ref{lem:3.2} 与 \ref{lem:3.3} 中把 $\R$ 换成 $\CC$).
某些证明需要显而易见的修改 (例如命题~
\ref{prop:3.11} 的证明).

\runin{第~
\ref{chap:4}~章.} 完全不变.

\runin{第~
\ref{chap:5}~章.} $\CC$ 上的 Hilbert 空间是配以标量积 $(u,v)\in\CC$ 的 $\CC$
上的向量空间.这是由 $H\times H$ 到 $\CC$ 的映射, 满足
\begin{align*}
(u,v)&=\overline{(v,u)} &&\forall u,v\in H,\\
\text{对每个}\ v\in H,\ u&\mapsto(u,v)\ \text{是线性的},\\
(u,u)&>0 &&\forall u\ne0.
\end{align*}
特别地, 我们有
\begin{equation*}
(\lambda u,\mu v)=\lambda\bar\mu(u,v)\quad\forall\lambda,\mu\in\CC,\
\forall u,v\in H.
\end{equation*}
量 $\abs{u}=(u,u)^{1/2}$ 是一个范数; 我们有
\begin{equation*}
\abs{u+v}^2=\abs{u}^2+2\operatorname{Re}(u,v)+\abs{v}^2\quad\forall u,v\in H,
\end{equation*}
且 Cauchy--Schwarz 不等式采取如下形式:
\begin{equation*}
\abs{(u,v)}\le\abs{u}\abs{v}\quad\forall u,v\in H.
\end{equation*}
一个典型的例子是配以标量积
\begin{equation*}
(u,v)=\int_\Omega u(x)\overline{v(x)}d\mu
\end{equation*}
的 $L^2(\Omega;\CC)$.

$\R$ 上的 Hilbert 空间与 $\CC$ 上的 Hilbert 空间之间的联系如下: 假设 $H$
是 $\CC$ 上的 Hilbert 空间.则配以标量积 $\operatorname{Re}(u,v)$ 的 $H_\R$
成为 $\R$ 上的 Hilbert 空间.因此第~
\ref{chap:5}~章的所有断言对 $H_\R$ 都适用.下面是
一些例子.

\begin{proposition}
\label{prop:11.26}
设 $K\subset H$ 为非空闭凸集.则对每个 $f\in H$, 存在唯一的元素 $u\in K$
使得
\begin{equation*}
\abs{f-u}=\min_{v\in K}\abs{f-v}=\dist(f,K).
\end{equation*}
此外, $u$ 由下述性质表征:
\begin{equation*}
u\in K\quad\text{且}\quad\operatorname{Re}(f-u,v-u)\le0\quad\forall v\in K.
\end{equation*}
\end{proposition}

\begin{proposition}
\label{prop:11.27}
给定任何 $\phi\in H^\star$, 存在唯一的 $f\in H$ 使得
\begin{equation*}
\phi(u)=(u,f)\quad\forall u\in H.
\end{equation*}
此外,
\begin{equation*}
\abs{f}=\norm{\phi}_{H^\star}.
\end{equation*}
\end{proposition}

\begin{proof}
把定理  \ref{thm:5.5} 应用于 $H_\R$ 中的 $\operatorname{Re}\phi$, 我们找到
某个 $f\in H$ 使得
\begin{equation*}
\operatorname{Re}\phi(u)=\operatorname{Re}(u,f)\quad\forall u\in H.
\end{equation*}
把它应用于 $iu$ 给出 $\operatorname{Im}\phi(u)=\operatorname{Im}(u,f)$,
从而 $\phi(u)=(u,f)\ \forall u\in H$.
\end{proof}

现在考虑满足下述条件的函数 $a(u,v):H\times H\to\CC$:
\begin{equation}
\forall v\in H,\ u\mapsto a(u,v)\ \text{是线性的, 且}\
\forall u\in H,\ v\mapsto\overline{a(u,v)}\ \text{是线性的},
\label{eq:11.13}
\end{equation}
\begin{equation}
a\ \text{连续, 即}\ \abs{a(u,v)}\le C\abs{u}\abs{v}\quad\forall u,v\in H,
\label{eq:11.14}
\end{equation}
\begin{equation}
a\ \text{强制, 即}\ \operatorname{Re}a(u,u)\ge\alpha\abs{u}^2
\quad\forall u\in H,\ \text{对某个}\ \alpha>0.
\label{eq:11.15}
\end{equation}

\begin{proposition}
\label{prop:11.28}
假设 $a$ 满足 \eqref{eq:11.13}, \eqref{eq:11.14} 与 \eqref{eq:11.15}.
设 $K$ 为 $H$ 中的非空闭凸集.则给定任何 $\phi\in H^\star$, 存在唯一的
$u\in K$ 使得
\begin{equation}
\operatorname{Re}a(u,v-u)\ge\operatorname{Re}\dual{\phi}{v-u}
\quad\forall v\in K.
\label{eq:11.16}
\end{equation}
此外, 若对一切 $v,w\in H$ 有 $a(v,w)=\overline{a(w,v)}$, 则 $u$ 由下述
性质表征:
\begin{equation*}
u\in K\quad\text{且}\quad
\frac12a(u,u)-\operatorname{Re}\dual{\phi}{u}
=\min_{v\in K}\left\{\frac12a(v,v)-\operatorname{Re}\dual{\phi}{v}\right\}.
\end{equation*}
\end{proposition}

当 $K=H$ 时, \eqref{eq:11.16} 变为 $a(u,v)=\dual{\phi}{v}\ \forall v\in H$.
特别地, 我们推出任何满足
\begin{equation}
\operatorname{Re}(Tu,u)\ge\alpha\abs{u}^2\quad\forall u\in H,\
\text{对某个}\ \alpha>0
\label{eq:11.17}
\end{equation}
的算子 $T\in\LL(H)$ 是由 $H$ 到自身的双射.有一个看上去稍微更一般的变体
(然而见下面的注记 1).

\begin{proposition}[Lax--Milgram]
\label{prop:11.29}
假设 $T\in\LL(H)$ 满足
\begin{equation}
\abs{(Tu,u)}\ge\alpha\abs{u}^2\quad\forall u\in H,\ \text{对某个}\ \alpha>0.
\label{eq:11.18}
\end{equation}
则 $T$ 是双射.
\end{proposition}

\begin{proof}
见第~
\ref{chap:5}~章注记 8.
\end{proof}

\begin{remark}
\label{rem:11.1}
显然 \eqref{eq:11.17} 蕴含 \eqref{eq:11.18}.反之, 假设 \eqref{eq:11.18}
成立.则存在某个满足 $\abs{\xi}=1$ 的 $\xi\in\CC$, 使得
\begin{equation}
\operatorname{Re}(\xi Tu,u)\ge\alpha\abs{u}^2\quad\forall u\in H.
\label{eq:11.19}
\end{equation}
事实上, 数值域
\begin{equation*}
W(T)=\{(Tu,u);\ u\in H,\ \abs{u}=1\}
\end{equation*}
是凸集 (由下面的命题  \ref{prop:11.33}).此外, 由 \eqref{eq:11.18} 我们
知道 $0\notin W(T)$, 且事实上 $\dist(0,W(T))\ge\alpha$.设 $p$ 表示 $0$
在 $W(T)$ 上的投影 (在 $\CC\simeq\R^2$ 中).经过平面内的一个把 $p$ 转到
$x$ 轴上的点 $(0,\abs{p})$ 的旋转 (即乘以 $\xi\in\CC$, $\abs{\xi}=1$)
之后, 我们得出 \eqref{eq:11.19} 成立.
\end{remark}

\runin{第~
\ref{chap:6}~章.} \S\ref{sec:6.1} 与 \S\ref{sec:6.2} 完全不变.主要区别出现在 \S\ref{sec:6.3} 中.

设 $E$ 为 $\CC$ 上的 Banach 空间, $T\in\LL(E)$.预解集定义为
\begin{equation*}
\rho(T)=\{\lambda\in\CC;\ (T-\lambda I)\ \text{是由}\ E\ \text{到}\ E\
\text{上的双射}\}.
\end{equation*}
谱是 $\rho(T)$ 的余集, 即 $\sigma(T)=\CC\setminus\rho(T)$.若相应的特征
子空间 $N(T-\lambda I)\ne\{0\}$, 则称数 $\lambda\in\CC$ 为\emph{特征值},
全体特征值的集合记作 $\operatorname{EV}(T)$.显然
$\operatorname{EV}(T)\subset\sigma(T)$.可能出现
$\operatorname{EV}(T)=\emptyset$ 的情形 (例如右移位算子
$Tu=(0,u_1,u_2,\dots)$).然而, $\sigma(T)$ 绝非空集.

\begin{proposition}
\label{prop:11.30}
谱 $\sigma(T)$ 是非空紧集, 且
\begin{equation*}
\sigma(T)\subset\{\lambda\in\CC;\ \abs{\lambda}\le\norm{T}\}.
\end{equation*}
\end{proposition}

\begin{proof}
主要的新颖之处在于 $\sigma(T)$ 非空.证明依赖于 $\CC$ 上的解析函数理论
(更精确地说是 Liouville 定理), 我们这里不予给出.有兴趣的读者可查阅
A.~Taylor--D.~Lay \bref{Taylor}{1}, W.~Rudin \bref{Rudin}{2}, 或
A.~Knapp \bref{Knapp}{2}.
\end{proof}

估计 $\abs{\lambda}\le\norm{T}\ \forall\lambda\in\sigma(T)$ 通常不是精确的.
例如, 在 $\CC^2$ 中算子 $T(u_1,u_2)=(u_2,0)$ 满足 $\sigma(T)=\{0\}$ 且
$\norm{T}=1$.最优界是用\emph{谱半径}给出的.我们已知 (见习题 6.23): 对每个
算子 $T\in\LL(E)$,
\begin{equation*}
r(T)=\lim_{n\to\infty}\norm{T^n}^{1/n}\ \text{存在},
\end{equation*}
且显然 $r(T)\le\norm{T}$; $r(T)$ \index{pubanjing@谱半径 (spectral radius)}
称为\emph{谱半径}.

\begin{proposition}
\label{prop:11.31}
对每个 $T\in\LL(E)$, 我们有
\begin{equation*}
r(T)=\max\{\abs{\lambda};\ \lambda\in\sigma(T)\}.
\end{equation*}
\end{proposition}

关于证明, 我们再次引证 A.~Taylor--D.~Lay \bref{Taylor}{1},
W.~Rudin \bref{Rudin}{2}, 或 A.~Knapp \bref{Knapp}{2}.论证通过 $\CC$ 上的
幂级数理论在很大程度上依赖于 $E$ 是 $\CC$ 上的 Banach 空间这一事实.当 $E$
是 $\R$ 上的 Banach 空间时, 我们只能说
$\max\{\abs{\lambda};\ \lambda\in\sigma(T)\}\le r(T)$, 且即使
$\sigma(T)$ 非空, 不等式也可以是严格的 (见习题 6.23).

实空间与复空间之间的另一个有趣区别与\index{puyingshedingli@谱映射定理 (spectral mapping theorem)}
所谓的\emph{谱映射定理}有关.先考虑
实情形: 设 $Q(t)=\sum_{k=0}^pa_kt^k$ 为实系数多项式, 设 $T\in\LL(E)$,
其中 $E$ 是 $\R$ 上的 Banach 空间.我们知道 (见习题 6.22)
\begin{equation}
Q(\operatorname{EV}(T))\subset\operatorname{EV}(Q(T))
\quad\text{且}\quad
Q(\sigma(T))\subset\sigma(Q(T)),
\label{eq:11.20}
\end{equation}
且这些包含关系可以是严格的 (但在 Hilbert 空间中当 $T^\star=T$ 时除外).
在复情形, 这些包含关系成为等式: 假设 $Q(t)$ 是复系数多项式, 设
$T\in\LL(E)$, 其中 $E$ 是 $\CC$ 上的 Banach 空间.

\begin{proposition}
\label{prop:11.32}
我们有
\begin{equation}
Q(\operatorname{EV}(T))=\operatorname{EV}(Q(T))
\label{eq:11.21}
\end{equation}
以及
\begin{equation}
Q(\sigma(T))=\sigma(Q(T)).
\label{eq:11.22}
\end{equation}
\end{proposition}

\begin{proof}
我们已知 \eqref{eq:11.20} 成立 (论证与习题 6.22 相同).反设包含关系是
严格的.则存在 $\mu\in\operatorname{EV}(Q(T))$ 使得
$\mu\notin Q(\operatorname{EV}(T))$.写出
\begin{equation*}
Q(t)-\mu=\alpha(t-t_1)(t-t_2)\dots(t-t_p),
\end{equation*}
其中 $\alpha\ne0$, 且 $t_i\notin\operatorname{EV}(T)\ \forall i$.此外,
存在某个 $x\ne0$ 使得 $Q(T)x=\mu x$.由于 $(T-t_1I)$ 是单射, 我们推出
$(T-t_2I)\dots(T-t_pI)x=0$, 重复同样的论证给出 $x=0$.矛盾.

类似地, 设 $\mu\in\sigma(Q(T))$ 满足 $\mu\notin Q(\sigma(T))$.如上写出
$Q(t)-\mu$, 其中 $t_i\notin\sigma(T)\ \forall i$.则 $Q(T)-\mu I$ 可以
写成双射算子的乘积.因此 $Q(T)-\mu I$ 是双射, 即
$\mu\in\rho(Q(T))$.矛盾.
\end{proof}

在 Hilbert 空间中, 研究谱的一个有用工具是数值域.设 $H$ 为 $\CC$ 上的
Hilbert 空间; \index{shuzhiyu@数值域 (numerical range)}
算子 $T\in\LL(H)$ 的\emph{数值域}定义为
\begin{equation*}
W(T)=\{(Tu,u);\ u\in H\ \text{且}\ \abs{u}=1\}.
\end{equation*}

\begin{proposition}
\label{prop:11.33}
我们有
\begin{equation*}
\sigma(T)\subset\overline{W(T)},
\end{equation*}
更精确地说, 若 $\lambda\notin\overline{W(T)}$, 则 $\lambda\in\rho(T)$,
满足
\begin{equation}
\norm{(T-\lambda I)^{-1}}\le\frac{1}{\dist(\lambda,W(T))}.
\label{eq:11.23}
\end{equation}
此外, $W(T)$ 是凸的.
\end{proposition}

\begin{proof}
假设 $\lambda\notin\overline{W(T)}$, 并令 $\alpha=\dist(\lambda,W(T))$.
对满足 $\abs{u}=1$ 的一切 $u\in H$, 有
\begin{equation*}
\abs{(Tu,u)-\lambda}\ge\alpha.
\end{equation*}
从而
\begin{equation*}
\abs{(Tu-\lambda u,u)}\ge\alpha\abs{u}^2\quad\forall u\in H.
\end{equation*}
应用 Lax--Milgram 定理 (命题  \ref{prop:11.29}), 我们得出 $(T-\lambda I)$
是双射, 且 $\abs{Tu-\lambda u}\ge\alpha\abs{u}\ \forall u\in H$, 即
$\norm{(T-\lambda I)^{-1}}\le1/\alpha$.

$W(T)$ 的凸性是由于 Toeplitz 与 Hausdorff 的一个违反直觉的事实.证明
参见 P.~R.~Halmos \bref{Halmos}{2}.
\end{proof}

一般说来, 数值域 $W(T)$ 可能比谱大得多.例如, 取 $H=\CC^2$ 与
$T(u_1,u_2)=(u_2,0)$, 我们有
$\operatorname{EV}(T)=\sigma(T)=\{0\}$, 而
$W(T)=\{\lambda\in\CC;\ \abs{\lambda}\le1/2\}$.然而, 若 $T$ 自伴, 或更
一般地为\index{zhengguisuanzi@正规算子 (normal operator)}
正规算子 (见下文), 则 $W(T)=\operatorname{conv}\sigma(T)$, 即
$\sigma(T)$ 的凸包 (见 P.~R.~Halmos \bref{Halmos}{2} 与下面的注记 2).

当 $H$ 是 $\CC$ 上的 Hilbert 空间且 $T\in\LL(H)$ 时, 对伴随算子 $T^\star$
的概念需要提醒读者注意.依照一般的步骤, 算子 $T\in\LL(H)$ 的伴随算子通过
关系式
\begin{equation*}
\dual{T^\star f}{u}_{H^\star,H}=\dual{f}{Tu}_{H^\star,H}
\quad\forall f\in H^\star,\ \forall u\in H
\end{equation*}
定义, 此时 $T^\star\in\LL(H^\star)$ (我们强调: 对一切 $\lambda\in\CC$ 及
一切 $f\in H^\star$, $T^\star(\lambda f)=\lambda T^\star f$).此外,
对一切 $\lambda\in\CC$ 有 $(\lambda T)^\star=\lambda T^\star$ (因为
$f:H\to\CC$ 是线性的).

另一方面, 我们也可以把 $H^\star$ 与 $H$ 等同起来 (通过命题
 \ref{prop:11.27} 中的同构), 并把 $T^\star$ 看作由 $H$ 到自身的算子, 它
通过关系式
\begin{equation*}
(Tu,v)=(u,T^\star v)\quad\forall u,v\in H
\end{equation*}
定义, 此时 $T^\star\in\LL(H)$ (我们强调: 对一切 $\lambda\in\CC$ 及一切
$v\in H$, $T^\star(\lambda v)=\lambda T^\star v$).然而, 我们现在有
\begin{equation}
(\lambda T)^\star=\bar\lambda T^\star\quad\forall\lambda\in\CC
\label{eq:11.24}
\end{equation}
(这容易验证).通常采用这一约定, 使得 $T^\star$ 与 $T$ 生活在同一个世界中:
可以比较 $T^\star$ 与 $T$, 作复合 $T^\star\circ T$ 等.

若 $T^\star=T$, 即
\begin{equation*}
(Tu,v)=(u,Tv)\quad\forall u,v\in H,
\end{equation*}
则称算子 $T\in\LL(H)$ 是\emph{自伴的} (或 \emph{Hermite 的}).若 $T$
自伴, 则对一切 $u\in H$ 有 $(Tu,u)=(u,Tu)=\overline{(Tu,u)}$, 从而对一切
$u\in H$ 有 $(Tu,u)\in\R$.特别地, 数值域 $W(T)$ 是 $\R$ 的子集, 从而
$\sigma(T)\subset\R$.

紧自伴算子的谱分解与第~
\ref{chap:6}~章中完全相同.

\begin{proposition}
\label{prop:11.34}
设 $H$ 为 $\CC$ 上的可分 Hilbert 空间, 设 $T$ 为紧自伴算子.则存在由 $T$
的特征向量组成的 Hilbert 基 (且相应的特征值为实数).
\end{proposition}

若算子 $T\in\LL(H)$ 满足 $T^\star\circ T=T\circ T^\star$, 则称它为
\emph{正规的}.当底层空间 $H$ 是 $\R$ 上的 Hilbert 空间时, 正规算子的各种
性质在问题 43 中作了讨论; 当 $H$ 是 $\CC$ 上的 Hilbert 空间时它们仍然成立.
但现在我们有更多结果:

\begin{proposition}
\label{prop:11.35}
设 $H$ 为 $\CC$ 上的 Hilbert 空间, 设 $T$ 为正规算子.则
\begin{equation}
\max\{\abs{\lambda};\ \lambda\in\sigma(T)\}=\norm{T}.
\label{eq:11.25}
\end{equation}
\end{proposition}

\begin{proof}
由于 $T$ 正规, 对每个整数 $p\ge1$ 有
\begin{equation*}
\norm{T^p}=\norm{T}^p.
\end{equation*}
这在 $H$ 是 $\R$ 上的 Hilbert 空间时已在问题 43 中证明, 同样的论证当 $H$
是 $\CC$ 上的 Hilbert 空间时仍然有效 (或者也可以把实结果应用于 $H_\R$ 上的
$T$).因此 $r(T)=\lim_{n\to\infty}\norm{T^n}^{1/n}=\norm{T}$.结合命题
 \ref{prop:11.31}, 给出 \eqref{eq:11.25}.
\end{proof}

\begin{proposition}
\label{prop:11.36}
设 $H$ 为 $\CC$ 上的可分 Hilbert 空间, 设 $T$ 为紧正规算子.则存在由 $T$
的特征向量组成的 Hilbert 基 (但相应的特征值不必为实数).
\end{proposition}

\begin{proof}
若 $T$ 正规, 则对任何 $\lambda\in\CC$, $(T-\lambda I)$ 也是正规的.因此
(如同问题 43 中那样) 我们有
$N(T-\lambda I)=N((T-\lambda I)^\star)=N(T^\star-\bar\lambda I)$.由此
推出当 $\lambda\ne\mu$ 时, $N(T-\lambda I)$ 与 $N(T-\mu I)$ 正交.我们
就可以完全像定理  \ref{thm:6.11} 的证明那样进行.我们在 $F^\perp$ 上得到
紧正规算子 $T_0$, 满足 $\sigma(T_0)=\{0\}$.我们不再援引推论  \ref{cor:6.10}
来得出 $T_0=0$, 而是应用命题  \ref{prop:11.35} 推出 $T_0=0$.正是在这里,
我们利用了 $H$ 是 $\CC$ 上的空间这一事实 (同样的结论在实空间中不成立).
\end{proof}

\begin{remark}
\label{rem:11.2}
当 $T$ 是紧正规算子时, 容易由命题  \ref{prop:11.36} 推出
$W(T)=\operatorname{conv}\sigma(T)$.事实上, 如命题  \ref{prop:11.36} 中
那样选取基 $(e_i)$.给定满足 $\abs{u}=1$ 的 $u\in H$, 把 $u$ 写成
$u=\sum u_ie_i$, 其中 $\sum\abs{u_i}^2=1$.则 $Tu=\sum\lambda_iu_ie_i$ 且
$(Tu,u)=\sum\lambda_i\abs{u_i}^2$.对任何正规算子 $T$ (不必紧),
$W(T)=\operatorname{conv}\sigma(T)$ 仍然成立; 参见 P.~R.~Halmos \bref{Halmos}{2}.
\end{remark}

设 $H$ 为 $\CC$ 上的 Hilbert 空间.若 $\abs{Tu}=\abs{u}\ \forall u\in H$,
则称算子 $T\in\LL(H)$ 为\emph{等距}; 若 $T$ 是等距且又是满射, 则称 $T$
为\emph{\index{yousuanzi@酉算子 (unitary operator)}
酉算子}.当底层空间 $H$ 是 $\R$ 上的 Hilbert 空间时, 等距与酉算子
的各种性质在问题 44 中作了讨论; 当 $H$ 是 $\CC$ 上的 Hilbert 空间时它们中
的大多数仍然成立, 但关于谱的断言除外, 它需要作如下修改:

\begin{proposition}
\label{prop:11.37}
设 $T$ 为等距.则
\begin{equation*}
\operatorname{EV}(T)\subset S^1=\{\lambda\in\CC;\ \abs{\lambda}=1\}.
\end{equation*}
若 $T$ 是酉算子, 则
\begin{equation*}
\sigma(T)\subset S^1,
\end{equation*}
若 $T$ 不是酉算子, 则
\begin{equation*}
\sigma(T)=\{\lambda\in\CC;\ \abs{\lambda}\le1\}.
\end{equation*}
\end{proposition}

证明是问题 44 第 6 问的解答中给出的证明的一个容易的改编.

若 $T^\star=-T$, 则\index{xiezibansuanzi@斜自伴算子 (skew-adjoint operator)}
称算子 $T\in\LL(H)$ 为\emph{斜自伴的} (或
\emph{反对称的}).显然, $T$ 斜自伴当且仅当 $iT$ 自伴 (这由 \eqref{eq:11.24}
推出).因此, 对任何斜自伴算子, 我们有
$\operatorname{EV}(T)\subset\sigma(T)\subset\overline{W(T)}\subset i\R$.


\backmatter
% chapters/solutions.tex
% 部分习题解答 (Solutions of Some Exercises)
% 源: 原书 PDF 第 375--437 页 (印刷页 371--434)

\chapter*{部分习题解答}
\addcontentsline{toc}{chapter}{部分习题解答}
\markboth{部分习题解答}{部分习题解答}

在本节中, 为了避免与前面各节的公式混淆, 公式用 (S1), (S2) 等编号.

\solutionhead{1.1}

1. 等式 $\dual{f}{x}=\norm{x}^2$ 蕴含 $\norm{x}\le\norm{f}$.推论  \ref{cor:1.3} 蕴含
$F(x)$ 非空.由 $F(x)$ 的第二种形式显然可见 $F(x)$ 是闭凸集.

2. 在严格凸赋范空间中, 包含在球面内的任何非空凸集都退化为单点.

3. 注意
\begin{equation*}
\dual{f}{y}\le\norm{f}\norm{y}\le\frac12\norm{f}^2+\frac12\norm{y}^2.
\end{equation*}
反之, 设 $f$ 满足
\begin{equation}
\frac12\norm{y}^2-\frac12\norm{x}^2\ge\dual{f}{y-x}\quad\forall y\in E.
\tag{S1}
\label{eq:S1.1}
\end{equation}
先在 \eqref{eq:S1.1} 中取 $y=\lambda x$ ($\lambda\in\R$); 变动 $\lambda$
可见 $\dual{f}{x}=\norm{x}^2$.接着在 \eqref{eq:S1.1} 中选取满足
$\norm{y}=\delta>0$ 的 $y$; 由此推出
\begin{equation*}
\dual{f}{y}\le\frac12\delta^2+\frac12\norm{x}^2.
\end{equation*}
因此我们得到
\begin{equation*}
\delta\norm{f}=\sup_{\substack{y\in E\\\norm{y}=\delta}}\dual{f}{y}
\le\frac12\delta^2+\frac12\norm{x}^2.
\end{equation*}
选取 $\delta=\norm{x}$ 即得出结论.

4. 若 $f\in F(x)$, 则有
\begin{equation*}
\frac12\norm{y}^2-\frac12\norm{x}^2\ge\dual{f}{y-x},
\end{equation*}
若 $g\in F(y)$, 则有
\begin{equation*}
\frac12\norm{x}^2-\frac12\norm{y}^2\ge\dual{g}{x-y}.
\end{equation*}
把这两个不等式相加得出 $\dual{f-g}{x-y}\ge0$.另一方面, 注意
\begin{equation*}
\dual{f-g}{x-y}=\norm{x}^2+\norm{y}^2-\dual{f}{y}-\dual{g}{x}
\ge\norm{x}^2+\norm{y}^2-2\norm{x}\norm{y}.
\end{equation*}

5. 由第 4 问我们已知 $\norm{x}=\norm{y}$.另一方面, 我们有
\begin{equation*}
\dual{F(x)-F(y)}{x-y}=[\norm{x}^2-\dual{F(x)}{y}]+[\norm{y}^2-\dual{F(y)}{x}],
\end{equation*}
括号内的两项都 $\ge0$.由此推出
$\norm{x}^2=\norm{y}^2=\dual{F(x)}{y}=\dual{F(y)}{x}$, 这蕴含
$F(x)\in F(y)$, 从而由第 2 问 $F(x)=F(y)$.

\solutionhead{1.2}

1(a). $\norm{f}_{E^\star}=\max_{1\le i\le n}\abs{f_i}$.

1(b). $f\in F(x)$ 当且仅当对每个 $1\le i\le n$, 有
\begin{equation*}
f_i=\begin{cases}
(\operatorname{sign}x_i)\norm{x}_1 & \text{若}\ x_i\ne0,\\
\text{区间}\ [-\norm{x}_1,+\norm{x}_1]\ \text{中的任意数} & \text{若}\ x_i=0.
\end{cases}
\end{equation*}

2(a). $\norm{f}_{E^\star}=\sum_{i=1}^n\abs{f_i}$.

2(b). 给定 $x\in E$, 考虑集合
\begin{equation*}
I=\{1\le i\le n;\ \abs{x_i}=\norm{x}_\infty\}.
\end{equation*}
则 $f\in F(x)$ 当且仅当成立
\begin{enumerate}[label=(\roman*), itemsep=0.1em]
	\item 对一切 $i\notin I$, $f_i=0$,
	\item 对一切 $i\in I$, $f_ix_i\ge0$, 且 $\sum_{i\in I}\abs{f_i}=\norm{x}_\infty$.
\end{enumerate}

3. $\norm{f}_{E^\star}=\left(\sum_{i=1}^n\abs{f_i}^2\right)^{1/2}$, 且
$f\in F(x)$ 当且仅当对一切 $i=1,2,\dots,n$, $f_i=x_i$.更一般地,
\begin{equation*}
\norm{f}_{E^\star}=\left(\sum_{i=1}^n\abs{f_i}^{p'}\right)^{1/p'},
\end{equation*}
其中 $1/p+1/p'=1$, 且 $f\in F(x)$ 当且仅当对一切 $i=1,2,\dots,n$,
$f_i=\abs{x_i}^{p-2}x_i/\norm{x}_p^{p-2}$.

\solutionhead{1.3}

1. $\norm{f}_{E^\star}=1$ (注意对一切 $\alpha>0$, $f(t^\alpha)=1/(1+\alpha)$).

2. 若存在这样的 $u$, 则会有 $\int_0^1(1-u)dt=0$, 从而 $u\equiv1$; 矛盾.

\solutionhead{1.5}

1. 用 $\mathcal{P}$ 记 $E$ 的所有线性无关子集的族.容易看出 $\mathcal{P}$
(按通常的包含关系偏序) 是归纳的.Zorn 引理蕴含 $\mathcal{P}$ 有极大元,
记作 $(e_i)_{i\in I}$, 它显然是一个代数基.由于对一切 $i\in I$, $e_i\ne0$,
经单位化后可以假设对一切 $i\in I$, $\norm{e_i}=1$.

2. 由于 $E$ 无限维, 可以假设 $\N\subset I$.存在 $E$ 上的 (唯一) 线性泛函,
使得当 $i\in\N$ 时 $f(e_i)=i$, 而当 $i\in I\setminus\N$ 时 $f(e_i)=0$.

3. 假设 $I$ 是可数的, 即 $I=\N$.考虑由 $(e_i)_{1\le i\le n}$ 张成的向量
空间 $F_n$.$F_n$ 是闭的 (见 \S \ref{sec:11.1}), 并且
$\bigcup_{n=1}^\infty F_n=E$.由 Baire 纲定理推出存在某个 $n_0$ 使得
$\operatorname{Int}(F_{n_0})\ne\emptyset$.因此 $E=F_{n_0}$; 矛盾.

\solutionhead{1.7}

1. 设 $x,y\in\overline C$, 从而 $x=\lim x_n$, $y=\lim y_n$, 其中
$x_n,y_n\in C$.于是 $tx+(1-t)y=\lim[tx_n+(1-t)y_n]$, 因此对一切
$t\in[0,1]$, $tx+(1-t)y\in\overline C$.

设 $x,y\in\operatorname{Int}C$, 从而存在某个 $r>0$ 使得 $B(x,r)\subset C$
且 $B(y,r)\subset C$.由此推出对一切 $t\in[0,1]$,
\begin{equation*}
tB(x,r)+(1-t)B(y,r)\subset C.
\end{equation*}
但 $tB(x,r)+(1-t)B(y,r)=B(tx+(1-t)y,r)$ (为什么?).

2. 设 $r>0$ 使得 $B(y,r)\subset C$.对一切 $t\in[0,1]$, 有
\begin{equation*}
tx+(1-t)B(y,r)\subset C,
\end{equation*}
从而 $B(tx+(1-t)y,(1-t)r)\subset C$.由此推出对一切 $t\in[0,1)$,
$tx+(1-t)y\in\operatorname{Int}C$.

3. 固定任何 $y_0\in\operatorname{Int}C$.给定 $x\in\overline C$, 有
$x=\lim_{n\to\infty}\left[\left(1-\frac1n\right)x+\frac1ny_0\right]$.但
$\left(1-\frac1n\right)x+\frac1ny_0\in\operatorname{Int}C$, 因此
$x\in\overline{\operatorname{Int}C}$.这证明了
$\overline C\subset\overline{\operatorname{Int}C}$, 从而
$\overline C=\overline{\operatorname{Int}C}$.

\solutionhead{1.8}

1. 我们已知对一切 $\lambda>0$ 与一切 $x\in E$, $p(\lambda x)=\lambda p(x)$,
且对一切 $x,y\in E$, $p(x+y)\le p(x)+p(y)$.还须验证:

(i) 对一切 $x\in E$, $p(-x)=p(x)$, 这由 $C$ 的对称性推出.

(ii) $p(x)=0\Rightarrow x=0$, 这由 $C$ 有界的事实推出.更精确地说, 设
$L>0$ 使得对一切 $x\in C$, $\norm{x}\le L$.容易看出对一切 $x\in E$,
$p(x)\ge\frac1L\norm{x}$.

2. $C$ 不是有界的.例如考虑序列 $u_n(t)=\sqrt{n}/(1+nt)$, 验证 $u_n\in C$,
而 $\norm{u_n}=\sqrt{n}$.这里 $p(u)=\left(\int_0^1\abs{u(t)}^2dt\right)^{1/2}$
是一个与 $\norm{u}$ 不等价的范数.

\solutionhead{1.9}

1. 令
\begin{equation*}
P=\left\{\lambda=(\lambda_1,\lambda_2,\dots,\lambda_n)\in\R^n;\
\lambda_i\ge0\ \forall i\ \text{且}\ \sum_{i=1}^n\lambda_i=1\right\},
\end{equation*}
于是 $P$ 是 $\R^n$ 的紧子集, 且 $C_n$ 是 $P$ 在连续映射
$\lambda\mapsto\sum_{i=1}^n\lambda_ix_i$ 下的像.

2. 对 $C_n$ 与 $\{0\}$ 应用 Hahn--Banach 第二几何形式.把与分离 $C_n$ 和
$\{0\}$ 的超平面相关联的线性泛函单位化.

4. 对 $C=A-B$ 应用上述构造.

\solutionhead{1.10}

(A) $\Rightarrow$ (B) 是显然的.

(B) $\Rightarrow$ (A). 设 $G$ 为由 $x_i$ ($i\in I$) 张成的向量空间.给定
$x\in G$, 把 $x$ 写成 $x=\sum_{i\in J}\beta_ix_i$, 并令
$g(x)=\sum_{i\in J}\beta_i\alpha_i$.假设 (B) 蕴含这个定义是有意义的, 且
对一切 $x\in G$, $\abs{g(x)}\le M\norm{x}$.接着用推论  \ref{cor:1.2} 把
$g$ 延拓到整个 $E$.

\solutionhead{1.11}

(A) $\Rightarrow$ (B) 再次是显然的.

(B) $\Rightarrow$ (A). 先设 $f_i$ 线性无关 ($1\le i\le n$).令
$\alpha=(\alpha_1,\alpha_2,\dots,\alpha_n)\in\R^n$.考虑由
\begin{equation*}
\phi(x)=(\dual{f_1}{x},\dots,\dual{f_n}{x})
\end{equation*}
定义的映射 $\phi:E\to\R^n$.令 $C=\{x\in E;\ \norm{x}\le M+\varepsilon\}$.
须证明 $\alpha\in\phi(C)$.反设 $\alpha\notin\phi(C)$, 并分离 $\phi(C)$ 与
$\{\alpha\}$ (见习题 1.9).因此存在某个 $\beta=(\beta_1,\dots,\beta_n)\in\R^n$,
$\beta\ne0$, 使得对一切 $x\in C$, $\beta\cdot\phi(x)\le\beta\cdot\alpha$, 即
\begin{equation*}
\dual{\sum\beta_if_i}{x}\le\sum\beta_i\alpha_i\quad\forall x\in C.
\end{equation*}
由此推出 $(M+\varepsilon)\norm{\sum\beta_if_i}\le\sum\beta_i\alpha_i$.利用
假设 (B), 发现 $\sum\beta_if_i=0$.由于 $f_i$ 线性无关, 得出 $\beta=0$;
矛盾.

在一般情形下, 把上述结果应用于 $(f_i)_{1\le i\le n}$ 的极大线性无关子集.

\solutionhead{1.15}

1. 显然 $C\subset C^{\star\star}$ 且 $C^{\star\star}$ 是闭的.反之, 设
$x_0\in C^{\star\star}$ 且 $x_0\notin C$.可以严格分离 $\{x_0\}$ 与 $C$,
从而存在某个 $f_0\in E^\star$ 与某个 $\alpha_0\in\R$ 使得
\begin{equation*}
\dual{f_0}{x}<\alpha_0<\dual{f_0}{x_0}\quad\forall x\in C.
\end{equation*}
由于 $0\in C$, 推出 $\alpha_0>0$; 令 $f=(1/\alpha_0)f_0$, 有
\begin{equation*}
\dual{f}{x}<1<\dual{f}{x_0}\quad\forall x\in C.
\end{equation*}
因此 $f\in C^\star$, 我们导出了矛盾, 因为 $x_0\in C^{\star\star}$.

2. 若 $C$ 是线性子空间, 则
\begin{equation*}
C^\star=\{f\in E^\star;\ \dual{f}{x}=0\ \forall x\in C\}=C^\perp.
\end{equation*}

\solutionhead{1.18}

(a) $\varphi^\star(f)=\begin{cases}-b & \text{若}\ f=a,\\ +\infty & \text{若}\ f\ne a.\end{cases}$

(b) $\varphi^\star(f)=\begin{cases}f\log f-f & \text{若}\ f>0,\\ 0 & \text{若}\ f=0,\\ +\infty & \text{若}\ f<0.\end{cases}$

(c) $\varphi^\star(f)=\abs{f}$.

(d) $\varphi^\star(f)=0$.

(e) $\varphi^\star(f)=\begin{cases}+\infty & \text{若}\ f\ge0,\\ -1-\log\abs{f} & \text{若}\ f<0.\end{cases}$

(f) $\varphi^\star(f)=(1+f^2)^{1/2}$.

(g) $\varphi^\star(f)=\begin{cases}\frac12f^2 & \text{若}\ \abs{f}\le1,\\ +\infty & \text{若}\ \abs{f}>1.\end{cases}$

(h) $\varphi^\star(f)=\frac{1}{p'}\abs{f}^{p'}$, 其中 $\dfrac1p+\dfrac{1}{p'}=1$.

(i) $\varphi^\star(f)=\begin{cases}0 & \text{若}\ 0\le f\le1,\\ +\infty & \text{其他}.\end{cases}$

(j) $\varphi^\star(f)=\begin{cases}\frac{1}{p'}f^{p'} & \text{若}\ f\ge0,\\ 0 & \text{若}\ f<0.\end{cases}$

(k) $\varphi^\star(f)=\begin{cases}+\infty & \text{若}\ f\ge0,\\ -\frac{1}{p'}\abs{f}^{p'} & \text{若}\ f<0.\end{cases}$

(l) $\varphi^\star(f)=\abs{f}+\frac{1}{p'}\abs{f}^{p'}$.

\solutionhead{1.20}

共轭函数定义在 $\ell^{p'}$ 上 ($\dfrac1p+\dfrac{1}{p'}=1$), 由下式给出:

(a) $\varphi^\star(f)=\begin{cases}\dfrac14\sum_{k=1}^\infty\dfrac1k\abs{f_k}^2 & \text{若}\ \sum_{k=1}^\infty\dfrac1k\abs{f_k}^2<+\infty,\\ +\infty & \text{其他}.\end{cases}$

(b) $\varphi^\star(f)=\begin{cases}\sum_{k=2}^{+\infty}a_k\abs{f_k}^{k/(k-1)} & \text{若}\ \sum_{k=2}^{+\infty}a_k\abs{f_k}^{k/(k-1)}<+\infty,\\ +\infty & \text{其他},\end{cases}$
其中 $a_k=(k-1)k^{-k/(k-1)}$.

(c) $\varphi^\star(f)=\begin{cases}0 & \text{若}\ \norm{f}_{\ell^\infty}\le1,\\ +\infty & \text{其他}.\end{cases}$

\solutionhead{1.21}

2. $\varphi^\star=I_A$, 其中 $A=\{[f_1,f_2];\ f_1\le0,\ f_2\le0,\ \text{且}\ 4f_1f_2\ge1\}$.

3. 有
\begin{equation*}
\inf_{x\in E}\{\varphi(x)+\psi(x)\}=0
\end{equation*}
且 $\varphi^\star=I_{D^\perp}$, 其中 $D^\perp=\{[f_1,f_2];\ f_2=0\}$.
由此推出对一切 $f\in E^\star$, $\varphi^\star(-f)+\psi^\star(f)=+\infty$, 从而
\begin{equation*}
\sup_{f\in E^\star}\{-\varphi^\star(-f)-\psi^\star(f)\}=-\infty.
\end{equation*}

4. 定理  \ref{thm:1.12} 的假设不满足: 不存在元素 $x_0\in E$ 使得
$\varphi(x_0)<+\infty$, $\psi(x_0)<+\infty$, 且 $\varphi$ 在 $x_0$ 处连续.

\solutionhead{1.22}

1. 写出
\begin{equation*}
\norm{x-a}\le\norm{x-y}+\norm{y-a}.
\end{equation*}
取 $\inf_{a\in A}$ 给出 $\varphi(x)\le\norm{x-y}+\varphi(y)$.然后交换 $x$
与 $y$.

2. 固定 $x,y\in E$ 与 $t\in[0,1]$.给定 $\varepsilon>0$, 存在 $a\in A$ 与
$b\in A$ 使得
\begin{equation*}
\norm{x-a}\le\varphi(x)+\varepsilon\quad\text{且}\quad
\norm{y-b}\le\varphi(y)+\varepsilon.
\end{equation*}
因此
\begin{equation*}
\norm{tx+(1-t)y-[ta+(1-t)b]}\le t\varphi(x)+(1-t)\varphi(y)+\varepsilon.
\end{equation*}
但 $ta+(1-t)b\in A$, 从而对一切 $\varepsilon>0$,
\begin{equation*}
\varphi(tx+(1-t)y)\le t\varphi(x)+(1-t)\varphi(y)+\varepsilon.
\end{equation*}

3. 由于 $A$ 闭, 有 $A=\{x\in E;\ \varphi(x)\le0\}$, 因此若 $\varphi$ 是
凸的, 则 $A$ 是凸的.

4. 有
\begin{align*}
\varphi^\star(f)&=\sup_{x\in E}\{\dual{f}{x}-\inf_{a\in A}\norm{x-a}\}
=\sup_{x\in E}\sup_{a\in A}\{\dual{f}{x}-\norm{x-a}\}\\
&=\sup_{a\in A}\sup_{x\in E}\{\dual{f}{x}-\norm{x-a}\}
=(I_A)^\star(f)+I_{B_{E^\star}}(f).
\end{align*}

\solutionhead{1.23}

1. 设 $f\in D(\varphi^\star)\cap D(\psi^\star)$.对每个 $x,y\in E$, 有
\begin{align*}
\dual{f}{x-y}-\varphi(x-y)&\le\varphi^\star(f),\\
\dual{f}{y}-\psi(y)&\le\psi^\star(f).
\end{align*}
把这两个不等式相加得出
\begin{equation*}
(\varphi\mathbin{\triangledown}\psi)(x)\ge\dual{f}{x}-\varphi^\star(f)-\psi^\star(f).
\end{equation*}
特别地, $(\varphi\mathbin{\triangledown}\psi)(x)>-\infty$.此外, 我们有
\begin{align*}
(\varphi\mathbin{\triangledown}\psi)^\star(f)
&=\sup_{x\in E}\{\dual{f}{x}-\inf_{y\in E}[\varphi(x-y)+\psi(y)]\}\\
&=\sup_{x\in E}\sup_{y\in E}\{\dual{f}{x}-\varphi(x-y)-\psi(y)\}\\
&=\sup_{y\in E}\sup_{x\in E}\{\dual{f}{x}-\varphi(x-y)-\psi(y)\}\\
&=\varphi^\star(f)+\psi^\star(f).
\end{align*}

2. 须验证对一切 $f,g\in E^\star$ 及一切 $x\in E$,
\begin{equation*}
\dual{f}{x}-\varphi(x)-\psi(x)\le\varphi^\star(f-g)+\psi^\star(g).
\end{equation*}
写成 $\dual{f}{x}=\dual{f-g}{x}+\dual{g}{x}$ 即显然.

3. 给定 $f\in E^\star$, 须证明
\begin{equation}
\sup_{x\in E}\{\dual{f}{x}-\varphi(x)-\psi(x)\}
=\inf_{g\in E^\star}\{\varphi^\star(f-g)+\psi^\star(g)\}.
\tag{S1}
\label{eq:S1.23}
\end{equation}
注意
\begin{equation*}
\sup_{x\in E}\{\dual{f}{x}-\varphi(x)-\psi(x)\}
=-\inf_{x\in E}\{\tilde\varphi(x)+\psi(x)\},
\end{equation*}
其中 $\tilde\varphi(x)=\varphi(x)-\dual{f}{x}$.把定理  \ref{thm:1.12}
应用于函数 $\tilde\varphi$ 与 $\psi$ 给出
\begin{equation*}
\inf_{x\in E}\{\tilde\varphi(x)+\psi(x)\}
=\sup_{g\in E^\star}\{-\tilde\varphi^\star(-g)-\psi^\star(g)\},
\end{equation*}
它恰好对应于 \eqref{eq:S1.23}.

4. 显然有
\begin{align*}
(\varphi^\star\mathbin{\triangledown}\psi^\star)^\star(x)
&=\sup_{f\in E^\star}\{\dual{f}{x}-\inf_{g\in E^\star}[\varphi^\star(f-g)+\psi^\star(g)]\}\\
&=\sup_{f\in E^\star}\sup_{g\in E^\star}\{\dual{f}{x}-\varphi^\star(f-g)-\psi^\star(g)\}\\
&=\sup_{g\in E^\star}\sup_{f\in E^\star}\{\dual{f}{x}-\varphi^\star(f-g)-\psi^\star(g)\}\\
&=\varphi^{\star\star}(x)+\psi^{\star\star}(x).
\end{align*}

\solutionhead{1.24}

1. 我们知道 (命题  \ref{prop:1.10}) 存在某个 $f\in E^\star$ 与常数 $C$ 使得
对一切 $y\in E$, $\varphi(y)\ge\dual{f}{y}-C$.选取 $n\ge\norm{f}$, 则有
$\varphi_n(x)\ge-n\norm{x}-C>-\infty$.

2. 函数 $\varphi_n$ 是两个凸函数的下卷积; 因此 $\varphi_n$ 是凸的 (见习题
1.23 第 7 问).为证明 $\abs{\varphi_n(x_1)-\varphi_n(x_2)}\le n\norm{x_1-x_2}$,
使用与习题 1.22 第 1 问相同的论证.

3. $(\varphi_n)^\star=I_{nB_{E^\star}}+\varphi^\star$ (由习题 1.23 第 1 问).

5. 由第 1 问我们对一切 $y\in E$ 有 $\varphi(y)\ge-\norm{f}\norm{y}-C$, 这给出
\begin{equation*}
n\norm{x-y_n}\le\norm{f}\norm{y_n}+C+\varphi(x)+1/n.
\end{equation*}
由此推出当 $n\to\infty$ 时 $\norm{y_n}$ 保持有界, 因此
$\lim_{n\to\infty}\norm{x-y_n}=0$.另一方面, 我们有
$\varphi_n(x)\ge\varphi(y_n)-1/n$, 由于 $\varphi$ 下半连续, 我们得出
$\liminf_{n\to\infty}\varphi_n(x)\ge\varphi(x)$.

6. 反设存在常数 $C$ 使得沿一个仍记作 $\varphi_n(x)$ 的子序列有
$\varphi_n(x)\le C$.如第 5 问那样选取 $y_n$, 可见 $y_n\to x$.此外,
$\varphi(y_n)\le C+1/n$, 从而
$\varphi(x)\le\liminf_{n\to\infty}\varphi(y_n)\le C$; 矛盾.

\solutionhead{1.25}

4. 对每个固定的 $t>0$, 函数
\begin{equation*}
y\mapsto\frac{1}{2t}\left(\norm{x+ty}^2-\norm{x}^2\right)
\end{equation*}
是凸的.因此作为凸函数的极限, 函数 $y\mapsto[x,y]$ 是凸的.另一方面,
\begin{equation*}
G(x,y)=\sup_{t>0}\left\{-\frac{1}{2t}\left(\norm{x+ty}^2-\norm{x}^2\right)\right\}
\end{equation*}
作为连续函数的上确界是下半连续的.

5. 我们已知 (见习题 1.1 第 3 问)
\begin{equation*}
\frac12\norm{x+ty}^2-\frac12\norm{x}^2\ge\dual{f}{ty},
\end{equation*}
因此对一切 $x,y\in E$ 及一切 $f\in F(x)$,
\begin{equation*}
[x,y]\ge\dual{f}{y}.
\end{equation*}
另一方面, 有
\begin{equation*}
\varphi^\star(f)=\frac12\norm{f}^2-\dual{f}{x}+\frac12\norm{x}^2,
\end{equation*}
以及
\begin{equation*}
\psi^\star(f)=\begin{cases}
0 & \text{若}\ \dual{f}{y}+\alpha\le0,\\
+\infty & \text{若}\ \dual{f}{y}+\alpha>0.
\end{cases}
\end{equation*}
容易验证 $\inf_{z\in E}\{\varphi(z)+\psi(z)\}=0$.由定理  \ref{thm:1.12}
推出存在某个 $f_0\in E^\star$ 使得 $\varphi^\star(f_0)+\psi^\star(-f_0)=0$,
即 $\dual{f_0}{y}\ge\alpha$ 且
$\frac12\norm{f_0}^2-\dual{f_0}{x}+\frac12\norm{x}^2=0$.从而我们有
$\norm{f_0}=\norm{x}$ 且 $\dual{f_0}{x}=\norm{x}^2$, 即 $f_0\in F(x)$.

6. (a) $1<p<\infty$, $[x,y]=\dfrac{\sum\abs{x_i}^{p-2}x_iy_i}{\norm{x}_p^{p-2}}$.

(b) $p=1$, $[x,y]=\norm{x}_1\left(\sum_{x_i\ne0}(\operatorname{sign}x_i)y_i
+\sum_{x_i=0}\abs{y_i}\right)$.

(c) $p=\infty$, $[x,y]=\max_{i\in I}\{x_iy_i\}$, 其中
$I=\{1\le i\le n;\ \abs{x_i}=\norm{x}_\infty\}$.

\solutionhead{1.27}

设 $\widetilde T:E\to F$ 为 $T$ 的连续线性延拓.容易验证
$E=N(\widetilde T)+G$ 且 $N(\widetilde T)\cap G=\{0\}$, 从而
$N(\widetilde T)$ 是 $G$ 的补空间; 矛盾.
\solutionhead{2.1}

不失一般性可以假设 $x_0=0$.

1. 设 $X=\{x\in E;\ \norm{x}\le\rho\}$, 其中 $\rho>0$ 足够小使得 $X\subset D(\varphi)$.
集合 $F_n$ 是闭的且 $\bigcup_{n=1}^\infty F_n=X$.由 Baire 纲定理存在某个 $n_0$
使得 $\operatorname{Int}(F_{n_0})\ne\emptyset$.设 $x_1\in E$ 与 $\rho_1>0$
使得 $B(x_1,\rho_1)\subset F_{n_0}$.给定任何满足 $\norm{x}<\rho_1/2$ 的
$x\in E$, 写出 $x=\frac12(x_1+2x)+\frac12(-x_1)$, 得出
$\varphi(x)\le\frac12n_0+\frac12\varphi(-x_1)$.

2. 存在某个 $\xi\in E$ 与某个常数 $t\in[0,1]$ 使得 $\norm{\xi}=R$ 且
$x_2=tx_1+(1-t)\xi$.由此推出
\begin{equation*}
\varphi(x_2)\le t\varphi(x_1)+(1-t)M,
\end{equation*}
从而 $\varphi(x_2)-\varphi(x_1)\le(1-t)[M-\varphi(x_1)]$.但
$x_2-x_1=(1-t)(\xi-x_1)$, 从而 $\norm{x_2-x_1}\ge(1-t)(R-r)$.因此我们有
\begin{equation*}
\varphi(x_2)-\varphi(x_1)\le\frac{\norm{x_2-x_1}}{R-r}[M-\varphi(x_1)].
\end{equation*}
另一方面, 若 $x_2=0$, 得到 $t\norm{x_1}=(1-t)R$, 因此
$(1-t)=\frac{\norm{x_1}}{\norm{x_1}+R}\le\frac12$.由此推出
$\varphi(0)-\varphi(x_1)\le\frac12[M-\varphi(x_1)]$, 从而
$M-\varphi(x_1)\le2[M-\varphi(0)]$.

\solutionhead{2.2}

由 (i) 我们有 $p(0)\le p(x_n)+p(-x_n)\to0$, 从而 $p(0)\le0$.另一方面由 (i)
$p(0)\le2p(0)$.因此 $p(0)=0$.

接着证明 $p(\alpha_nx_n)\to0$.反设沿某个子序列对某个 $\varepsilon>0$
有 $\abs{p(\alpha_nx_n)}>2\varepsilon$.抽取进一步的子序列, 可以假设对某个
$\alpha\in\R$ 有 $\alpha_n\to\alpha$.为简便起见, 仍用 $(x_n)$ 与 $(\alpha_n)$
记相应的序列.

集合 $F_n$ 是闭的且 $\bigcup_{n\ge1}F_n=\R$.应用 Baire 纲定理, 我们找到某个
$n_0$ 使得 $\operatorname{Int}F_{n_0}\ne\emptyset$.因此存在某个 $\lambda_0\in\R$
与某个 $\delta>0$ 使得对一切 $k\ge n_0$ 及一切满足 $\abs{t}<\delta$ 的 $t$,
有 $\abs{p((\lambda_0+t)x_k)}\le\varepsilon$.另一方面, 注意
\begin{align*}
p(\alpha_kx_k)&\le p((\lambda_0+\alpha_k-\alpha)x_k)+p((\alpha-\lambda_0)x_k),\\
-p(\alpha_kx_k)&\le-p((\lambda_0+\alpha_k-\alpha)x_k)+p((\lambda_0-\alpha)x_k).
\end{align*}
因此对足够大的 $k$, 我们得到 $\abs{p(\alpha_kx_k)}\le2\varepsilon$.矛盾.

最后, 写出
\begin{equation*}
p(\alpha_nx_n)-p(\alpha x)\le p(\alpha_n(x_n-x))+p(\alpha_nx)-p(\alpha x)\to0
\end{equation*}
以及
\begin{equation*}
p(\alpha_nx)\le p(\alpha_n(x-x_n))+p(\alpha_nx_n),
\end{equation*}
从而
\begin{equation*}
p(\alpha_nx_n)-p(\alpha x)\ge-p(\alpha_n(x_n-x))+p(\alpha_nx)-p(\alpha x)\to0.
\end{equation*}

\solutionhead{2.4}

由 (i), 存在线性算子 $T:E\to F^\star$ 使得对一切 $x,y$,
$a(x,y)=\dual{Tx}{y}_{F^\star,F}$.目标是证明 $T$ 是有界算子, 即 $T(B_E)$
在 $F^\star$ 中有界.由推论  \ref{cor:2.5}, 只须固定 $y\in F$ 并验证
$\dual{T(B_E)}{y}$ 有界.这由 (ii) 推出.

\solutionhead{2.6}

1. 有 $\dual{Ax_n-A(x_0+x)}{x_n-x_0-x}\ge0$, 从而
$\dual{Ax_n}{x}\le\varepsilon_n\norm{Ax_n}+C(x)$, 其中
$\varepsilon_n=\norm{x_n-x_0}$, $C(x)=\norm{A(x_0+x)}(1+\norm{x})$ (假设对
一切 $n$, $\varepsilon_n\le1$).由习题 2.5 推出 $(Ax_n)$ 有界; 矛盾.

2. 假设存在 $D(A)$ 中的序列 $(x_n)$ 使得 $x_n\to x_0$ 且
$\norm{Ax_n}\to\infty$.选取 $r>0$ 使得 $B(x_0,r)\subset\operatorname{conv}D(A)$.
对每个满足 $\norm{x}<r$ 的 $x\in E$, 写出
\begin{equation*}
x_0+x=\sum_{i=1}^m t_iy_i,\quad t_i\ge0\ \forall i,\quad
\sum_{i=1}^m t_i=1,\quad y_i\in D(A)\ \forall i
\end{equation*}
(当然 $t_i,y_i$ 与 $m$ 依赖于 $x$).我们有 $\dual{Ax_n-Ay_i}{x_n-y_i}\ge0$,
从而 $t_i\dual{Ax_n}{x_n-y_i}\ge t_i\dual{Ay_i}{x_n-y_i}$.由此推出
\begin{equation*}
\dual{Ax_n}{x_n-x_0-x}\ge\sum_{i=1}^m t_i\dual{Ay_i}{x_n-y_i},
\end{equation*}
这导出 $\dual{Ax_n}{x}\le\varepsilon_n\norm{Ax_n}+C(x)$, 其中
$\varepsilon_n=\norm{x_n-x_0}$,
$C(x)=\sum_{i=1}^m t_i\norm{Ay_i}(1+\norm{x_0-y_i})$.

3. 设 $x_0\in\operatorname{Int}D(A)$.沿用第 1 问中的相同论证, 可证明存在
两个常数 $R>0$ 与 $C$, 使得对一切满足 $\norm{x-x_0}<R$ 的 $x\in D(A)$ 及
一切 $f\in Ax$, 有 $\norm{f}\le C$.

\solutionhead{2.7}

对每个 $x\in\ell^p$ 令 $T_nx=\sum_{i=1}^n\alpha_ix_i$, 于是对每个
$x\in\ell^p$, $T_nx$ 收敛于一个极限.由推论  \ref{cor:2.3} 推出存在常数 $C$
使得对一切 $x\in\ell^p$ 及一切 $n$,
\begin{equation*}
\abs{T_nx}\le C\norm{x}_{\ell^p}.
\end{equation*}
适当地选取 $x$, 可见 $\alpha\in\ell^{p'}$ 且 $\norm{\alpha}_{p'}\le C$.

\solutionhead{2.8}

方法 (ii). 让我们验证 $T$ 的图是闭的.设 $(x_n)$ 为 $E$ 中的序列使得
$x_n\to x$ 且 $Tx_n\to f$.在不等式 $\dual{Tx_n-Ty}{x_n-y}\ge0$ 中取极限
得出
\begin{equation*}
\dual{f-Ty}{x-y}\ge0\quad\forall y\in E.
\end{equation*}
选取 $y=x+tz$ ($t\in\R,z\in E$), 可见 $f=Tx$.

\solutionhead{2.10}

1. 若 $T(M)$ 闭, 则 $M+N(T)=T^{-1}(T(M))$ 也闭.反之, 设 $M+N(T)$ 闭.
由于 $T$ 满, 有 $T((M+N(T))^c)=(T(M))^c$.开映射定理蕴含 $T((M+N(T))^c)$ 开,
从而 $T(M)$ 闭.

2. 若 $M$ 是任意闭子空间且 $N$ 是任意有限维空间, 则 $M+N$ 闭 (见 \S \ref{sec:11.1}).

\solutionhead{2.11}

由开映射定理, 存在常数 $c>0$ 使得 $T(B_E)\supset cB_F$.用 $(e_n)$ 记
$\ell^1$ 的典范基, 即 $e_n=(0,\dots,0,\underset{(n)}{1},0,\dots)$.存在某个
$u_n\in E$ 使得 $\norm{u_n}\le1/c$ 且 $T(u_n)=e_n$.给定
$y=(y_1,y_2,\dots,y_n,\dots)\in\ell^1$, 令 $Sy=\sum_{i=1}^\infty y_iu_i$.
显然级数收敛, 且 $S$ 具有所有要求的性质.

\solutionhead{2.12}

不失一般性可以假设 $T$ 满 (否则用 $R(T)$ 替代 $E$).反设存在 $E$ 中的序列
$(x_n)$ 使得 $\norm{x_n}_E=1$ 且 $\norm{Tx_n}_F+\abs{x_n}<1/n$.由开映射
定理存在常数 $c>0$ 使得 $T(B_E)\supset cB_F$.由于 $\norm{Tx_n}_F<1/n$,
存在某个 $y_n\in E$ 使得 $Tx_n=Ty_n$ 且 $\norm{y_n}_E<1/nc$.把 $x_n$ 写成
$x_n=y_n+z_n$, 其中 $z_n\in N(T)$, $\norm{y_n}_E\to0$ 且 $\norm{z_n}_E\to1$.
另一方面, $\abs{x_n}<1/n$; 因此
$\abs{z_n}<(1/n)+\abs{y_n}\le(1/n)+M\norm{y_n}_E$, 从而 $\abs{z_n}\to0$.
这是不可能的, 因为范数 $\norm{\cdot}_E$ 与 $\abs{\cdot}$ 在有限维空间 $N(T)$
上等价.

\solutionhead{2.13}

首先, 设 $T\in\mathcal{O}$, 从而 $T^{-1}\in\LL(F,E)$ (由推论  \ref{cor:2.7}).
则对满足 $\norm{U}$ 足够小的每个 $U\in\LL(E,F)$, $T+U\in\mathcal{O}$.
事实上, 方程 $Tx+Ux=f$ 可以写成 $x=T^{-1}(f-Ux)$; 只要
$\norm{T^{-1}}\norm{U}<1$, 它 (对每个 $f\in F$) 就有唯一解 (由 Banach
不动点定理; 见定理  \ref{thm:5.7}).

接着, 设 $T\in\mathcal{L}$.由定理  \ref{thm:2.13}, $R(T)$ 闭且在 $F$ 中有
补空间.设 $P:F\to R(T)$ 为连续投影.算子 $PT$ 是由 $E$ 到 $R(T)$ 上的
双射, 从而上面的分析适用.设 $U\in\LL(E,F)$ 满足 $\norm{U}<\delta$; 若
$\delta$ 足够小, 算子 $(PT+PU):E\to R(T)$ 是双射, 从而 $(PT+PU)^{-1}$
作为 $\LL(R(T),E)$ 的元素是良定的.令 $S=(PT+PU)^{-1}P$.显然
$S\in\LL(F,E)$ 且 $S(T+U)=I_E$.

\solutionhead{2.14}

1. 考虑商空间 $\widetilde E=E/N(T)$ 与典范满射 $\pi:E\to\widetilde E$, 使得
对一切 $x\in E$, $\norm{\pi x}_{\widetilde E}=\dist(x,N(T))$.$T$ 在
$\widetilde E$ 上诱导一个单射算子 $\widetilde T$.更精确地说, 把 $T$ 写成
$T=\widetilde T\circ\pi$, 其中 $\widetilde T\in\LL(\widetilde E,F)$, 从而
$R(T)=R(\widetilde T)$.另一方面, 推论  \ref{cor:2.7} 表明 $R(\widetilde T)$
闭当且仅当存在常数 $C$ 使得对一切 $y\in\widetilde E$,
$\norm{y}_{\widetilde E}\le C\norm{\widetilde Ty}$, 或等价地对一切 $x\in E$,
$\norm{\pi x}_{\widetilde E}\le C\norm{\widetilde T\pi x}$.最后一个不等式
即
\begin{equation*}
\dist(x,N(T))\le C\norm{Tx}\quad\forall x\in E.
\end{equation*}

\solutionhead{2.15}

算子 $T:E_1\times E_2\to F$ 是线性的, 有界的且是满射.此外,
$N(T)=N(T_1)\times N(T_2)$ (因为 $R(T_1)\cap R(T_2)=\{0\}$).把习题 2.10
应用于 $M=E_1\times\{0\}$, 可见只要 $M+N(T)$ 闭, $T(M)=R(T_1)$ 就是闭的.
但 $M+N(T)=E_1\times N(T_2)$ 确实是闭的.

\solutionhead{2.16}

设 $\pi$ 为由 $E$ 到 $E/L$ 上的典范满射 (见 \S \ref{sec:11.2}).考虑由
$Tx=\pi x$ ($x\in G$) 定义的算子 $T:G\to E/L$.我们有
\begin{equation*}
\dist(x,N(T))=\dist(x,G\cap L)\le C\dist(x,L)=C\norm{Tx}\quad\forall x\in G.
\end{equation*}
由此推出 (见习题 2.14) $R(T)=\pi(G)$ 是闭的.因此 $\pi^{-1}[\pi(G)]=G+L$
是闭的.

\solutionhead{2.19}

回忆 $N(A^\star)=R(A)^\perp$.

1. 设 $u\in N(A)$ 且 $v\in D(A)$; 对一切 $t\in\R$, 有
\begin{equation*}
\dual{A(u+tv)}{u+tv}\ge-C\norm{A(u+tv)}^2,
\end{equation*}
这蕴含 $\dual{Av}{u}=0$.因此 $N(A)\subset R(A)^\perp$.

2. 配以图范数的 $D(A)$ 是 Banach 空间.配以 $E^\star$ 范数的 $R(A)$ 是
Banach 空间.算子 $A:D(A)\to R(A)$ 满足开映射定理的假设.因此存在常数 $C$,
使得对一切 $f\in R(A)$, 存在 $v\in D(A)$ 满足 $Av=f$ 且
$\norm{v}_{D(A)}\le C\norm{f}$.特别地, $\norm{v}\le C\norm{f}$.给定
$u\in D(A)$, 把上述结果应用于 $f=Au$, 可见存在某个 $v\in D(A)$ 使得
$Au=Av$ 且 $\norm{v}\le C\norm{Au}$.由于 $u-v\in N(A)\subset R(A)^\perp$,
我们有
\begin{equation*}
\dual{Au}{u}=\dual{Av}{u}=\dual{Av}{v}\ge-\norm{Av}\norm{v}\ge-C\norm{Au}^2.
\end{equation*}

\solutionhead{2.21}

1. 区分两种情形: 情形 (i): $f(a)=1$.则 $N(A)=\R a$ 且 $R(A)=N(f)$.
情形 (ii): $f(a)\ne1$.则 $N(A)=\{0\}$ 且 $R(A)=E$.

2. $A$ 不是闭的.否则闭图像定理将蕴含 $A$ 有界, 从而 $f$ 连续.

3. $D(A^\star)=\{u\in E^\star;\ \dual{u}{a}=0\}$, 且对一切
$u\in D(A^\star)$, $A^\star u=u$.

4. $N(A^\star)=\{0\}$ 且 $R(A^\star)=\{u\in E^\star;\ \dual{u}{a}=0\}$.

5. $R(A)^\perp=\{0\}$ 且 $R(A^\star)^\perp=\R a$ (注意 $N(f)$ 在 $E$ 中稠密;
见习题 1.6).由此推出 $N(A^\star)=R(A)^\perp$ 且
$N(A)\subset R(A^\star)^\perp$.注意到在情形 (ii) 中,
$N(A)\ne R(A^\star)^\perp$.

6. 若 $A$ 不是闭的, 可能出现 $N(A)\ne R(A^\star)^\perp$ 的情形.

\solutionhead{2.22}

1. 显然 $D(A)$ 在 $E$ 中稠密.为了验证 $A$ 是闭的, 设 $(u^j)$ 为 $D(A)$
中的序列使得在 $E$ 中 $u^j\to u$ 且 $Au^j\to f$.由此推出对每个 $n$, 当
$j\to\infty$ 时 $u_n^j\to u_n$, 且 $nu_n^j\to f_n$.因此对一切 $n$,
$nu_n=f_n$, 从而 $u\in D(A)$ 且 $Au=f$.

2. $D(A^\star)=\{v=(v_n)\in\ell^\infty;\ (nv_n)\in\ell^\infty\}$,
$A^\star v=(nv_n)$, 且 $\overline{D(A^\star)}=c_0$.

\solutionhead{2.24}

1. 我们有 $D(B^\star)=\{v\in G^\star;\ T^\star v\in D(A^\star)\}$ 且
$B^\star=A^\star T^\star$.

2. 若 $D(A)\ne E$ 且 $T=0$, 则 $B$ 不是闭的.事实上, 设 $(u_n)$ 为 $D(A)$
中的序列使得 $u_n\to u$, 其中 $u\notin D(A)$.则 $Bu_n\to0$, 但 $u\notin D(B)$.

\solutionhead{2.25}

2. 由推论  \ref{cor:2.7}, $T^{-1}\in\LL(F,E)$.由于 $T^{-1}T=I_E$ 且
$TT^{-1}=I_F$, 推出 $T^\star(T^{-1})^\star=I_{E^\star}$ 以及
$(T^{-1})^\star T^\star=I_{F^\star}$.

\solutionhead{2.26}

我们有
\begin{equation*}
\varphi^\star(T^\star f)=\sup_{x\in E}\{\dual{T^\star f}{x}-\varphi(x)\}
=\sup_{y\in R(T)}\{\dual{f}{y}-\psi(y)\}
=-\inf_{y\in F}\{\psi(y)+\zeta(y)\},
\end{equation*}
其中 $\zeta(y)=-\dual{f}{y}+I_{R(T)}(y)$.应用定理  \ref{thm:1.12}, 我们得到
\begin{equation*}
\varphi^\star(T^\star f)=\min_{g\in F^\star}\{\zeta^\star(g)+\psi^\star(-g)\}.
\end{equation*}
但
\begin{equation*}
\zeta^\star(g)=\begin{cases}
0 & \text{若}\ f+g\in R(T)^\perp,\\
+\infty & \text{若}\ f+g\notin R(T)^\perp,
\end{cases}
\end{equation*}
从而
\begin{equation*}
\varphi^\star(T^\star f)=\min_{f+g\in N(T^\star)}\psi^\star(-g)
=\min_{h\in N(T^\star)}\psi^\star(f-h).
\end{equation*}

\solutionhead{2.27}

令 $G=E\times X$ 并考虑算子 $S(x,y)=Tx+y:G\to F$.应用开映射定理, 我们知道
$S$ 是开映射, 从而 $S(E\times(X\setminus\{0\}))=R(T)+(X\setminus\{0\})$ 在
$F$ 中开.因此它的补集 $R(T)$ 是闭的.
\solutionhead{3.1}

应用推论  \ref{cor:2.4}.

\solutionhead{3.2}

注意对一切 $f\in E^\star$, $\dual{f}{\sigma_n}=\dfrac1n\sum_{i=1}^n\dual{f}{x_i}$.
由于 $\dual{f}{x_n}\to\dual{f}{x}$, 推出 $\dual{f}{\sigma_n}\to\dual{f}{x}$.

\solutionhead{3.4}

1. 令 $G_n=\operatorname{conv}\left(\bigcup_{i=n}^\infty\{x_i\}\right)$.由于对拓扑
$\sigma(E,E^\star)$ 有 $x_n\rightharpoonup x$, 由此推出对一切 $n$,
$x\in\overline{G_n}^{\sigma(E,E^\star)}$.另一方面, $G_n$ 是凸集, 其弱拓扑
$\sigma(E,E^\star)$ 闭包与强拓扑闭包重合 (见习题 3.3).因此对一切 $n$,
$x\in\overline{G_n}$ ($G_n$ 的强闭包), 从而存在序列 $(y_n)$ 使得对一切 $n$,
$y_n\in G_n$ 且 $y_n\to x$ 强收敛.

2. 存在 $E$ 中的序列 $(u_k)$ 使得 $u_k\to x$ 且对一切 $k$,
$u_k\in\operatorname{conv}\left(\bigcup_{i=1}^\infty\{x_i\}\right)$.因此存在
严格递增的整数序列 $(n_k)$ 使得对一切 $k$,
$u_k\in\operatorname{conv}\left(\bigcup_{i=1}^{n_k}\{x_i\}\right)$.由
$n_k\le n<n_{k+1}$ 时 $z_n=u_k$ (以及 $1\le n<n_1$ 时 $z_n=x_1$) 定义的序列
$(z_n)$ 具有所有要求的性质.

\solutionhead{3.7}

1. 设 $x\notin A+B$.我们将构造 $0$ 的关于 $\sigma(E,E^\star)$ 的邻域 $W$
使得 $(x+W)\cap(A+B)=\emptyset$.

对每个 $y\in B$, 存在 $0$ 的凸邻域 $V(y)$ 使得 $(x+V(y))\cap(A+y)=\emptyset$
(因为 $A+y$ 闭且 $x\notin A+y$).显然
\begin{equation*}
B\subset\bigcup_{y\in B}\left(y-\frac12V(y)\right),
\end{equation*}
由于 $B$ 是紧集, 存在某个有限集 $I$ 使得
$B\subset\bigcup_{i\in I}\left(y_i-\frac12V(y_i)\right)$, 其中 $y_i\in B$.
令
\begin{equation*}
W=\frac12\bigcap_{i\in I}V(y_i).
\end{equation*}

我们断言 $(x+W)\cap(A+B)=\emptyset$.事实上, 反设存在某个 $w\in W$ 使得
$x+w\in(A+B)$.从而存在某个 $i\in I$ 使得
\begin{equation*}
x+w\in A+y_i-\frac12V(y_i).
\end{equation*}
由于 $V(y_i)$ 是凸的, 推出存在某个 $w'\in V(y_i)$ 使得 $x+w'\in A+y_i$.
因此 $(x+V(y_i))\cap(A+y_i)\ne\emptyset$; 矛盾.

\runin{注记.} 若 $E^\star$ 可分且 $A$ 有界, 则可以使用序列来证明 $A+B$
是闭的, 因为弱拓扑在有界集上是可度量化的 (见定理  \ref{thm:3.29}).这使论证
稍微更容易一些.事实上, 设 $x_n=a_n+b_n$ 为使得 $x_n\rightharpoonup x$
弱收敛的序列, 其中 $a_n\in A,b_n\in B$.由于 $B$ 弱紧 (且可度量化), 存在
子序列使得 $b_{n_k}\rightharpoonup b$, 其中 $b\in B$.从而
$a_{n_k}\rightharpoonup x-b$.但 $A$ 是弱闭的, 因此 $x-b\in A$, 即
$x\in A+B$.

2. 由第 1 问, $(A-B)$ 是弱闭的, 因此是强闭的.于是可以严格分离 $\{0\}$ 与
$(A-B)$.

\solutionhead{3.8}

1. 由于 $V_k$ 是 $0$ 关于 $\sigma(E,E^\star)$ 的邻域, 可以假设 (见命题
 \ref{prop:3.4}) $V_k$ 具有形式
\begin{equation*}
V_k=\{x\in E;\ \abs{\dual{f}{x}}<\varepsilon_k\ \forall f\in F_k\},
\end{equation*}
其中 $\varepsilon_k>0$, $F_k$ 是 $E^\star$ 的有限子集.因此集合
$F=\bigcup_{k=1}^\infty F_k$ 是可数的.我们断言任何 $g\in E^\star$ 可以写成
$F$ 中元素的有限线性组合.事实上, 令
\begin{equation*}
V=\{x\in E;\ \abs{\dual{g}{x}}<1\}.
\end{equation*}
由于 $V$ 是 $0$ 的关于 $\sigma(E,E^\star)$ 的邻域, 存在某个整数 $m$ 使得
$\{x\in E;\ d(x,0)<1/m\}\subset V$, 从而 $V_m\subset V$.设 $x\in E$ 使得
对一切 $f\in F_m$, $\dual{f}{x}=0$.则对一切 $t\in\R$, $tx\in V_m$, 从而
对一切 $t\in\R$, $tx\in V$, 即 $\dual{g}{x}=0$.应用引理  \ref{lem:3.2},
可见 $g$ 是 $F_m$ 中元素的线性组合.

2. 使用与习题 1.5 第 3 问相同的方法.

3. 若 $\dim E^\star<\infty$, 则 $\dim E^{\star\star}<\infty$; 从而
$\dim E<\infty$ (因为存在由 $E$ 到 $E^{\star\star}$ 中的典范单射).

4. 应用下述引理 (它是引理  \ref{lem:3.2} 的容易推论): 假设
$x_1,x_2,\dots,x_k,y\in E$ 满足
\begin{equation*}
[f\in E^\star;\ \dual{f}{x_i}=0\ \forall i]\implies[\dual{f}{y}=0].
\end{equation*}
则存在常数 $\lambda_1,\lambda_2,\dots,\lambda_k$ 使得
$y=\sum_{i=1}^k\lambda_ix_i$.

\solutionhead{3.9}

1. 应用定理  \ref{thm:1.12}, 取
\begin{equation*}
\varphi(x)=\dual{f_0}{x}+I_{B_E}(x)\quad\text{以及}\quad\psi(x)=I_M(x).
\end{equation*}

2. 注意 $B_{E^\star}$ 关于 $\sigma(E^\star,E)$ 是紧的, 而 $M^\perp$ 关于
$\sigma(E^\star,E)$ 是闭的 (为什么?).

\solutionhead{3.11}

只须对序列作论证 (为什么?).设在 $E$ 中 $x_n\to x$ 强收敛, 且
$Ax_n\not\rightharpoonup Ax$ 关于 $\sigma(E^\star,E)$, 即存在某个 $y\in E$
使得 $\dual{Ax_n}{y}\not\to\dual{Ax}{y}$.由习题 2.6 我们已知 $(Ax_n)$
有界.因此存在子序列使得 $\dual{Ax_{n_k}}{y}\to\ell\ne\dual{Ax}{y}$.应用
$A$ 的单调性, 我们有
\begin{equation*}
\dual{Ax_{n_k}-A(x+ty)}{x_{n_k}-x-ty}\ge0.
\end{equation*}
取极限, 得到
\begin{equation*}
-t\ell+t\dual{A(x+ty)}{y}\ge0,
\end{equation*}
这蕴含 $\ell=\dual{Ax}{y}$; 矛盾.

\solutionhead{3.12}

1. 假设 (A) 蕴含对一切 $f\in E^\star$,
$\varphi^\star(f)\ge R\norm{f}+\dual{f}{x_0}-M$.反之, 设 (B) 成立, 并令
$\psi(f)=\varphi^\star(f)-\dual{f}{x_0}$.我们断言存在常数 $k>0$ 与 $C$ 使得
\begin{equation}
\psi(f)\ge k\norm{f}-C\quad\forall f\in E^\star.
\tag{S1}
\label{eq:S3.12}
\end{equation}
平移后总可以假设 $\psi(0)<\infty$ (见命题  \ref{prop:1.10}).固定
$\alpha>\psi(0)$.利用假设 (B), 我们可以找到某个 $r>0$ 使得对一切满足
$\norm{g}\ge r$ 的 $g\in E^\star$, $\psi(g)\ge\alpha$.给定满足
$\norm{f}\ge r$ 的 $f\in E^\star$, 取 $t=r/\norm{f}$, 写出
$\psi(tf)\le t\psi(f)+(1-t)\psi(0)$.由于 $\norm{tf}=r$, 这导出
\begin{equation*}
\alpha-\psi(0)\le\frac{r}{\norm{f}}(\psi(f)-\psi(0)),
\end{equation*}
这建立了断言 \eqref{eq:S3.12}.对 \eqref{eq:S3.12} 取共轭即得到 (A).

2. 函数 $\psi$ 是凸的且关于弱$^\star$拓扑下半连续 (为什么?).假设 (B) 说
对每个 $\lambda\in\R$, 集合 $\{f\in E^\star;\ \psi(f)\le\lambda\}$ 是有界的.
因此它是弱$^\star$紧的 (由定理  \ref{thm:3.16}), 从而 $\inf_{E^\star}\psi$
可以达到.另一方面,
\begin{equation*}
\inf_{E^\star}\psi=-\sup_{f\in E^\star}\{\dual{f}{x_0}-\varphi^\star(f)\}
=-\varphi^{\star\star}(x_0)=-\varphi(x_0).
\end{equation*}
另外, 也可以把定理  \ref{thm:1.12} 应用于函数 $\varphi$ 与 $I_{\{x_0\}}$
(注意 $\varphi$ 在 $x_0$ 处连续; 见习题 2.1).

\solutionhead{3.13}

1. 对每个固定的 $p$, 对一切 $n$ 有 $x_{p+n}\in K_p$.取极限 ($n\to\infty$),
可见 $x\in K_p$, 因为 $K_p$ 是弱闭的 (见定理  \ref{thm:3.7}).另一方面, 设
$V$ 为 $x$ 的关于拓扑 $\sigma(E,E^\star)$ 的凸邻域.存在整数 $N$ 使得对
一切 $n\ge N$, $x_n\in V$.从而对一切 $n\ge N$, $K_n\subset V$, 因此
$\bigcap_{n=1}^\infty K_n\subset V$.由于这对 $x$ 的任何凸邻域 $V$ 都成立,
推出 $\bigcap_{n=1}^\infty K_n\subset\{x\}$ (为什么?).

2. 设 $V$ 为 $x$ 关于拓扑 $\sigma(E,E^\star)$ 的开邻域.令
$K'_n=K_n\cap(V^c)$.由于 $K_n$ 关于 $\sigma(E,E^\star)$ 是紧的 (为什么?),
推出 $K'_n$ 也关于 $\sigma(E,E^\star)$ 是紧的.另一方面,
$\bigcap_{n=1}^\infty K'_n=\emptyset$, 从而存在某个整数 $N$ 使得
$K'_N=\emptyset$, 即 $K_N\subset V$.

3. 可以假设 $x=0$.考虑渐近锥 $C_n=\bigcap_{\lambda>0}\lambda K_n$.由于
$C_n\subset K_n$, 我们推出 $\bigcap_{n=1}^\infty C_n=\{0\}$.令
$S_E=\{x\in E;\ \norm{x}=1\}$.序列 $(C_n\cap S)$ 递减且
$\bigcap_{n=1}^\infty(C_n\cap S)=\emptyset$.从而由紧性, 对某个 $n_0$ 有
$C_{n_0}\cap S=\emptyset$.因此 $C_{n_0}=\{0\}$, 从而 $K_{n_0}$ 有界
(为什么?).因此 $(x_n)$ 有界, 我们归结为第 2 问.

4. 考虑当 $n$ 为奇数时 $x_n=(0,\dots,0,\underset{(n)}{n},0,\dots)$, 当
$n$ 为偶数时 $x_n=0$ 的序列.

\solutionhead{3.18}

2. 反设在 $\ell^1$ 中对拓扑 $\sigma(\ell^1,\ell^\infty)$ 有
$e_{n_k}\rightharpoonup a$ ($k\to\infty$).于是对一切 $\xi\in\ell^\infty$,
有 $\dual{\xi}{e_{n_k}}\to\dual{\xi}{a}$.考虑由
\begin{equation*}
\xi=(0,\dots,0,\underset{(n_1)}{-1},0,\dots,\underset{(n_2)}{1},0,\dots,\underset{(n_3)}{-1},0,\dots)
\end{equation*}
定义的元素 $\xi\in\ell^\infty$.注意 $\dual{\xi}{e_{n_k}}=(-1)^k$ 当
$k\to\infty$ 时不收敛; 矛盾.

3. 设 $E=\ell^\infty$, 从而 $\ell^1\subset E^\star$.令 $f_n=e_n$, 看作
$E^\star$ 中的序列.我们断言 $(f_n)$ 没有关于 $\sigma(E^\star,E)$ 收敛的
子序列.反设对 $\sigma(E^\star,E)$ 有
$f_{n_k}\overset{*}{\rightharpoonup}f$, 即对一切 $\eta\in E$,
$\dual{f_{n_k}}{\eta}\to\dual{f}{\eta}$.如第 2 问中那样选取 $\eta=\xi$,
可见 $\dual{f_{n_k}}{\xi}=(-1)^k$ 不收敛; 矛盾.这里, 配以拓扑
$\sigma(E^\star,E)$ 的集合 $B_{E^\star}$ 是紧的 (由定理  \ref{thm:3.16}),
但它不是可度量化的.应用定理  \ref{thm:3.28}, 我们也可以说 $E=\ell^\infty$
不可分 (另一种证明见第~
\ref{chap:4}~章注记 8 与命题  \ref{prop:11.17}).

\solutionhead{3.19}

1. 注意若在 $\ell^p$ 中 $x_n\to x$ 强收敛, 则对一切 $\varepsilon>0$, 存在
$I$ 使得对一切 $n$, $\sum_{i=I}^\infty\abs{x_i^n}^p\le\varepsilon^p$.

2. 应用习题 3.17.

3. 空间 $B_E$ 对拓扑 $\sigma(E,E^\star)$ 是可度量化的 (由定理  \ref{thm:3.29}).
因此只须在序列上验证 $A$ 的连续性.

\solutionhead{3.20}

1. 考虑由 $(Tx)(t)=\dual{t}{x}$ ($x\in E,t\in B_{E^\star}=K$) 定义的映射
$T:E\to C(K)$.显然
\begin{equation*}
\norm{Tx}=\sup_{t\in K}\abs{(Tx)(t)}=\norm{x}.
\end{equation*}

2. 由于 $K=B_{E^\star}$ 关于 $\sigma(E^\star,E)$ 是可度量化且紧的, $K$ 中
存在稠密的可数子集 $(t_n)$.考虑由
\begin{equation*}
Sx=(\dual{t_1}{x},\dual{t_2}{x},\dots,\dual{t_n}{x},\dots)
\end{equation*}
定义的映射 $S:E\to\ell^\infty$.验证 $\norm{Sx}_{\ell^\infty}=\norm{x}$.

\solutionhead{3.21}

设 $(a_i)$ 为 $E$ 的稠密可数子集.选取第一个子序列使得当 $k\to\infty$ 时
$\dual{f_{n_k}}{a_1}$ 收敛于一个极限.然后从 $(n_k)$ 中抽取子序列使得
$\dual{f_{n'_k}}{a_2}$ 收敛, 等等.

由标准的对角线法, 可以从序列 $(f_n)$ 中抽取子序列 $(g_k)$ 使得对一切 $i$,
$\dual{g_k}{a_i}\to\ell_i$.由于集合 $(a_i)$ 在 $E$ 中稠密, 容易得到对
一切 $a\in E$, $\dual{g_k}{a}\to\ell_a$.由此推出 $g_k$ 对
$\sigma(E^\star,E)$ 收敛于某个 $g$ (见习题 3.16).

\solutionhead{3.22}

(a) $B_E$ 关于 $\sigma(E,E^\star)$ 是可度量化的 (由定理  \ref{thm:3.29}), 且
$0$ 属于 $S=\{x\in E;\ \norm{x}=1\}$ 关于 $\sigma(E,E^\star)$ 的闭包 (见第
3 章注记 2).

(b) 由于 $\dim E=\infty$, $E$ 中存在可分且满足 $\dim E_0=\infty$ 的闭子空间
$E_0$ (为什么?).注意 $E_0$ 自反, 应用情形 (a) (并结合推论  \ref{cor:3.27}).

\solutionhead{3.25}

反设 $C(K)$ 自反.则 $E=\{u\in C(K);\ u(a)=0\}$ 也自反, 且
$\sup_{u\in B_E}f(u)$ 可以达到.

另一方面, 我们断言 $\sup_{u\in B_E}f(u)=1$.事实上, 对一切 $N$, 存在
$u\in E$ 使得在 $K$ 上 $0\le u\le1$ 且对一切 $i=1,\dots,N$, $u(a_i)=1$
(例如应用 Tietze--Urysohn 定理; 参见例如 J.~Munkres \bref{Munkres}{1}).
因此存在某个 $u\in B_E$ 使得 $f(u)=1$.这导出对一切 $n$, $u(a_n)=1$ 且
$u(a)=0$; 矛盾.

\solutionhead{3.26}

1. 给定 $y\in B_F$, 存在某个整数 $n_1$ 使得 $\norm{y-a_{n_1}}<1/2$.由于集合
$\frac12(a_i)_{i>n_1}$ 在 $\frac12B_F$ 中稠密, 存在某个 $n_2>n_1$ 使得
$\norm{y-a_{n_1}-\frac12a_{n_2}}<1/4$.由归纳法构造递增的整数序列
$n_k\uparrow\infty$ 使得
\begin{equation*}
y=a_{n_1}+\frac12a_{n_2}+\frac14a_{n_3}+\dots+\frac{1}{2^{k-1}}a_{n_k}+\dots.
\end{equation*}

2. 反设 $S\in\LL(F,\ell^1)$ 使得 $TS=I_F$.设 $(y_n)$ 为 $F$ 中的任何序列
满足对一切 $n$, $\norm{y_n}=1$ 且 $y_n\rightharpoonup0$ 弱收敛于 $0$.从而
关于 $\sigma(\ell^1,\ell^\infty)$ 有 $Sy_n\rightharpoonup0$, 进而 $Sy_n\to0$
在 $\ell^1$ 中强收敛 (见问题 8).由此推出 $y_n=TSy_n\to0$; 矛盾.

3. 使用定理  \ref{thm:2.12}.

4. $T^\star:F^\star\to\ell^\infty$ 由
$T^\star v=(\dual{v}{a_1},\dual{v}{a_2},\dots,\dual{v}{a_n},\dots)$ 定义.

\solutionhead{3.27}

$B_{E^\star}$ 关于 $\sigma(E^\star,E)$ 是紧的且可度量化的.因此 $B_{E^\star}$
中存在关于 $\sigma(E^\star,E)$ 稠密的可数子集.

1. 显然对一切 $f\in E^\star$, $\norm{f}\le\norm{f}_1\le\sqrt{2}\norm{f}$.

2. 令 $\abs{f}^2=\sum_{n=1}^\infty\frac{1}{2^n}\abs{\dual{f}{a_n}}^2$.注意
范数 $\abs{\cdot}$ 与标量积相关联 (为什么?), 因此它是严格凸的, 即函数
$f\mapsto\abs{f}^2$ 严格凸.更精确地说, 对一切 $t\in[0,1]$ 与一切
$f,g\in E^\star$, 有
\begin{equation}
\abs{tf+(1-t)g}^2+t(1-t)\abs{f-g}^2=t\abs{f}^2+(1-t)\abs{g}^2.
\tag{S1}
\label{eq:S3.27}
\end{equation}
因此, 函数 $f\mapsto\norm{f}^2+\abs{f}^2$ 也是严格凸的.

3. 与第 2 问相同的方法.注意若对一切 $n$, $\dual{b_n}{x}=0$, 则 $x=0$
(为什么?).

4. 给定 $x\in E$, 令
$[x]=\left\{\sum_{n=1}^\infty\frac{1}{2^n}\abs{\dual{b_n}{x}}^2\right\}^{1/2}$,
并用 $[f]$ 记 $[\cdot]$ 在 $E^\star$ 上的对偶范数.注意 $[f]$ 也满足恒等式
\eqref{eq:S3.27}.事实上, 我们有
\begin{align*}
\frac12[tf+(1-t)g]^2&=\sup_{x\in E}\left\{\dual{tf+(1-t)g}{x}-\frac12[x]^2\right\},\\
\frac12[f-g]^2&=\sup_{y\in E}\left\{\dual{f-g}{y}-\frac12[y]^2\right\},
\end{align*}
从而
\begin{align*}
&\frac12[tf+(1-t)g]^2+\frac12t(1-t)[f-g]^2\\
&=\sup_{x,y}\left\{\dual{tf+(1-t)g}{x}+t(1-t)\dual{f-g}{y}-\frac12[x]^2-\frac12t(1-t)[y]^2\right\}.
\end{align*}
作变量代换 $x=t\xi+(1-t)\eta$ 与 $y=\xi-\eta$, 我们得出 \eqref{eq:S3.27}
成立.应用习题 1.23 第 3 问, 可见
\begin{equation*}
\norm{f}_2^2=\inf_{h\in E^\star}\{\norm{f-h}_1^2+[h]^2\}=\min_{h\in E^\star}\{\norm{f-h}_1^2+[h]^2\}.
\end{equation*}
我们断言函数 $f\mapsto\norm{f}_2^2$ 严格凸.事实上, 给定 $f,g\in E^\star$,
固定 $h_1,h_2\in E^\star$ 使得 $\norm{f}_2^2=\norm{f-h_1}_1^2+[h_1]^2$,
$\norm{g}_2^2=\norm{g-h_2}_1^2+[h_2]^2$.对每个 $t\in(0,1)$, 有
\begin{align*}
\norm{tf+(1-t)g}_2^2&\le\norm{tf+(1-t)g-(th_1+(1-t)h_2)}_1^2+[th_1+(1-t)h_2]^2\\
&<t\norm{f}_2^2+(1-t)\norm{g}_2^2,
\end{align*}
除非 $f-h_1=g-h_2$ 且 $h_1=h_2$, 即 $f=g$.

\solutionhead{3.28}

由于 $E$ 自反, $\sup_{x\in B_E}\dual{f}{x}$ 由某个唯一点 $x_0\in B_E$ 达到.
则 $x=x_0\norm{f}$ 满足 $f\in F(x)$.

另外, 也可以考虑由 $E^\star$ 到 $E^{\star\star}$ 中的对偶映射 $F^\star$.
集合 $F^\star(f)$ 非空 (由推论  \ref{cor:1.3}).固定任何 $\xi\in F^\star(f)$.
由于 $E$ 自反, 存在某个 $x\in E$ 使得 $Jx=\xi$ ($J$ 是由 $E$ 到
$E^{\star\star}$ 中的典范单射).我们有 $\norm{\xi}=\norm{f}=\norm{x}$ 以及
$\dual{\xi}{f}=\norm{f}^2=\dual{f}{x}$.因此 $f\in F(x)$.

\runin{唯一性.} 设 $x_1,x_2$ 使得 $f\in F(x_1)$ 且 $f\in F(x_2)$.则
$\norm{x_1}=\norm{x_2}=\norm{f}$, 因此若 $x_1\ne x_2$, 有
$\norm{\dfrac{x_1+x_2}{2}}<\norm{f}$.另一方面,
$\dual{f}{x_1}=\dual{f}{x_2}=\norm{f}^2$, 从而若 $x_1\ne x_2$, 有
$\norm{f}^2=\dual{f}{\dfrac{x_1+x_2}{2}}<\norm{f}^2$, 矛盾.

\solutionhead{3.29}

1. 反设存在 $M_0>0$, $\varepsilon_0>0$ 与两个序列 $(x_n),(y_n)$ 使得
$\norm{x_n}\le M$, $\norm{y_n}\le M$, $\norm{x_n-y_n}>\varepsilon_0$, 以及
\begin{equation}
\norm{\frac{x_n+y_n}{2}}^2>\frac12\norm{x_n}^2+\frac12\norm{y_n}^2-\frac1n.
\tag{S1}
\label{eq:S3.29}
\end{equation}
考虑仍记作 $(x_n)$ 与 $(y_n)$ 的子序列使得 $\norm{x_n}\to a$ 且
$\norm{y_n}\to b$.我们发现 $a+b\ge\varepsilon_0$ 且
$\frac12a^2+\frac12b^2\le\left(\frac{a+b}{2}\right)^2$.因此 $a=b\ne0$.

令 $x'_n=\dfrac{x_n}{\norm{x_n}}$ 与 $y'_n=\dfrac{y_n}{\norm{y_n}}$.对足够大的
$n$, 有 $\norm{x'_n-y'_n}\ge(\varepsilon_0/a)+o(1)$.由一致凸性存在
$\delta_0>0$ 使得 $\norm{\dfrac{x'_n+y'_n}{2}}\le1-\delta_0$.从而
\begin{equation*}
\norm{\frac{x_n+y_n}{2}}\le a(1-\delta_0)+o(1).
\end{equation*}
由 \eqref{eq:S3.29} 我们有 $\norm{\dfrac{x_n+y_n}{2}}^2\ge a^2+o(1)$.因此
$a^2\le a^2(1-\delta_0)^2+o(1)$; 矛盾.

\solutionhead{3.32}

1. 下确界可以达到, 因为 $E$ 自反, 我们可以应用推论  \ref{cor:3.23}.唯一性
来自空间 $E$ 严格凸的事实, 从而函数 $y\mapsto\norm{y-x}^2$ 严格凸.

2. 设 $(y_n)$ 为极小化序列; 令 $d_n=\norm{x-y_n}$ 与 $d=\inf_{y\in C}\norm{x-y}$,
从而 $d_n\to d$.设 $(y_{n_k})$ 为使得 $y_{n_k}\rightharpoonup z$ 弱收敛的序列.
从而 $z\in C$ 且 $\norm{x-z}\le d$ (为什么?).由此推出
$x-y_{n_k}\rightharpoonup x-z$ 弱收敛且 $\norm{x-y_{n_k}}\to d=\norm{x-z}$,
因此 (见命题  \ref{prop:3.32}) $y_{n_k}\to z$ 强收敛.极限的唯一性蕴含整个序列
$(y_n)$ 强收敛于 $P_Cx$.这个论证是标准的, 我们将多次使用它.为方便读者起见,
我们回忆它: 反设 $(y_n)$ 不收敛于 $y=P_Cx$.则存在 $\varepsilon>0$ 与子序列
$(y_{m_j})$ 使得对一切 $j$, $\norm{y_{m_j}-y}\ge\varepsilon$.从 $(y_{m_j})$
中我们抽取 (按上述论证) 进一步的子序列 $(y_{n_k})$ 使得 $y_{n_k}\to P_Cx$.
由于 $(y_{n_k})$ 是 $(y_{m_j})$ 的子序列, 对一切 $k$ 有
$\norm{y_{n_k}-y}\ge\varepsilon$, 从而 $\norm{P_Cx-y}\ge\varepsilon$.矛盾.

3 与 4. 反设存在某个 $\varepsilon_0>0$ 与序列 $(x_n),(y_n)$ 使得
$\norm{x_n}\le M$, $\norm{y_n}\le M$, $\norm{x_n-y_n}\to0$, 且对一切 $n$,
$\norm{P_Cx_n-P_Cy_n}\ge\varepsilon_0$.我们有
\begin{equation*}
\norm{x_n-P_Cx_n}\le\norm{\frac{x_n-P_Cx_n+P_Cy_n}{2}}\le\norm{\frac{x_n+y_n}{2}-\frac{P_Cx_n+P_Cy_n}{2}}+o(1),
\end{equation*}
同理
\begin{equation*}
\norm{y_n-P_Cy_n}\le\norm{\frac{x_n+y_n}{2}-\frac{P_Cx_n+P_Cy_n}{2}}+o(1).
\end{equation*}
由此推出
\begin{equation}
\frac12\norm{x_n-P_Cx_n}^2+\frac12\norm{y_n-P_Cy_n}^2\le\norm{\frac{x_n+y_n}{2}-\frac{P_Cx_n+P_Cy_n}{2}}^2+o(1).
\tag{S1}
\label{eq:S3.32a}
\end{equation}
另一方面, 若令 $a_n=x_n-P_Cx_n$ 与 $b_n=y_n-P_Cy_n$, 则
$\norm{a_n-b_n}\ge\varepsilon_0+o(1)$ 且 $\norm{a_n}\le M', \norm{b_n}\le M'$.
利用习题 3.29, 我们知道存在某个 $\delta_0>0$ 使得
\begin{equation*}
\norm{\frac{a_n+b_n}{2}}^2\le\frac12\norm{a_n}^2+\frac12\norm{b_n}^2-\delta_0,
\end{equation*}
即
\begin{equation}
\norm{\frac{x_n+y_n}{2}-\frac{P_Cx_n+P_Cy_n}{2}}^2\le\frac12\norm{x_n-P_Cx_n}^2+\frac12\norm{y_n-P_Cy_n}^2-\delta_0.
\tag{S2}
\label{eq:S3.32b}
\end{equation}
结合 \eqref{eq:S3.32a} 与 \eqref{eq:S3.32b} 导出矛盾.

5. 与第 1 问相同的论证.

6. 我们有
\begin{equation}
n\norm{y_n-x}^2+\varphi(y_n)\le n\norm{y-x}^2+\varphi(y)\quad\forall y\in D(\varphi).
\tag{S3}
\label{eq:S3.32c}
\end{equation}
由于 $\varphi$ 被一个连续仿射函数由下方控制 (见命题  \ref{prop:1.10}), 我们看到
当 $n\to\infty$ 时 $(y_n)$ 保持有界.设 $(y_{n_k})$ 为使得
$y_{n_k}\rightharpoonup z$ 弱收敛的子序列.注意 $z\in\overline{D(\varphi)}$
(为什么?).由 \eqref{eq:S3.32c} 我们得到对一切 $y\in D(\varphi)$,
$\norm{z-x}\le\norm{y-x}$, 从而对一切 $y\in\overline{D(\varphi)}$ 成立.因此
$z=P_Cx$, 其中 $C=\overline{D(\varphi)}$.再次使用 \eqref{eq:S3.32c} 给出
\begin{equation*}
\limsup_{n\to\infty}\norm{y_n-x}\le\norm{y-x}\quad\forall y\in D(\varphi),
\end{equation*}
特别地对 $y=z$ 成立.我们得出 $y_{n_k}\to z$ 强收敛, 最后极限的唯一性表明
(如上) 整个序列 $(y_n)$ 强收敛于 $P_Cx$.
\solutionhead{4.3}

2. 注意 $h_n=\frac12(\abs{f_n-g_n}+f_n+g_n)$.

3. 注意 $f_ng_n-fg=(f_n-f)g_n+f(g_n-g)$, 且由控制收敛定理, 在 $L^p(\Omega)$
中 $f(g_n-g)\to0$.

\solutionhead{4.5}

1. 回忆 $L^1(\Omega)\cap L^\infty(\Omega)\subset L^p(\Omega)$, 更精确地
$\norm{f}_p^p\le\norm{f}_\infty^{p-1}\norm{f}_1$.由于 $\Omega$ 是 $\sigma$-有限
的, 可以写成 $\Omega=\bigcup_n\Omega_n$, 其中对一切 $n$, $\abs{\Omega_n}<\infty$.
给定 $f\in L^p(\Omega)$, 验证 $f_n=\chi_{\Omega_n}T_nf\in L^1(\Omega)\cap L^\infty(\Omega)$,
且当 $n\to\infty$ 时在 $L^p(\Omega)$ 中 $f_n\to f$.

2. 设 $(f_n)$ 为 $L^p(\Omega)\cap L^q(\Omega)$ 中的序列, 使得在 $L^p(\Omega)$
中 $f_n\to f$ 且 $\norm{f_n}_q\le1$.我们假设 (通过抽取子序列) $f_n\to f$
几乎处处 (见定理  \ref{thm:4.9}).由 Fatou 引理推出 $f\in L^q(\Omega)$ 且
$\norm{f}_q\le1$.

3. 由第 2 问我们已知 $f\in L^q(\Omega)$, 从而对介于 $p$ 与 $q$ 之间的每个
$r$, $f\in L^r(\Omega)$.另一方面, 可以写出 $\dfrac1r=\dfrac\alpha p+\dfrac{1-\alpha}{q}$,
其中 $0<\alpha\le1$, 我们得到
\begin{equation*}
\norm{f_n-f}_r\le\norm{f_n-f}_p^\alpha\norm{f_n-f}_q^{1-\alpha}
\le\norm{f_n-f}_p^\alpha(2C)^{1-\alpha}.
\end{equation*}

\solutionhead{4.6}

1. 我们有 $\norm{f}_p\le\norm{f}_\infty\abs{\Omega}^{1/p}$, 从而
$\limsup_{p\to\infty}\norm{f}_p\le\norm{f}_\infty$.另一方面, 固定
$0<k<\norm{f}_\infty$, 并令 $A=\{x\in\Omega;\ \abs{f(x)}>k\}$.显然
$\abs{A}\ne0$ 且 $\norm{f}_p\ge k\abs{A}^{1/p}$.由此推出
$\liminf_{p\to\infty}\norm{f}_p\ge k$, 因此
$\liminf_{p\to\infty}\norm{f}_p\ge\norm{f}_\infty$.

2. 固定 $k>C$, 设 $A$ 如上定义.则 $k^p\abs{A}\le\norm{f}_p^p\le C^p$, 从而对
一切 $p\ge1$, $\abs{A}\le(C/k)^p$.令 $p\to\infty$, 可见 $\abs{A}=0$.

3. $f(x)=\log\abs{x}$.

\solutionhead{4.7}

考虑由 $Tu=au$ 定义的算子 $T:L^p(\Omega)\to L^q(\Omega)$.我们断言 $T$ 的
图是闭的.事实上, 设 $(u_n)$ 为 $L^p(\Omega)$ 中的序列使得在 $L^p(\Omega)$
中 $u_n\to u$ 且在 $L^q(\Omega)$ 中 $au_n\to f$.抽取子序列可以假设
$u_n\to u$ 几乎处处且 $au_n\to f$ 几乎处处.从而 $f=au$ 几乎处处, 于是
$f=Tu$.由闭图像定理 (定理  \ref{thm:2.9}) 推出 $T$ 有界, 因此存在常数 $C$
使得
\begin{equation}
\norm{au}_q\le C\norm{u}_p\quad\forall u\in L^p(\Omega).
\tag{S1}
\label{eq:S4.7}
\end{equation}

情形 1: $p<\infty$.由 \eqref{eq:S4.7} 推出
\begin{equation*}
\int_\Omega\abs{a}^q\abs{v}\le C^q\norm{v}_{p/q}\quad\forall v\in L^{p/q}(\Omega).
\end{equation*}
因此映射 $v\mapsto\int_\Omega\abs{a}^qv$ 是 $L^{p/q}(\Omega)$ 上的连续线性
泛函, 从而 $\abs{a}^q\in L^{(p/q)'}(\Omega)$.

情形 2: $p=\infty$.在 \eqref{eq:S4.7} 中选取 $u\equiv1$.

\solutionhead{4.8}

1. 配以范数 $\norm{\cdot}_1$ 的 $X$ 是 Banach 空间.对每个 $n$, $X_n$ 是
$X$ 的闭子集 (见习题 4.5).另一方面, $X=\bigcup_nX_n$.事实上, 对每个
$f\in X$, 存在某个 $q>1$ 使得 $f\in L^q(\Omega)$.因此只要
$1+(1/n)\le q$, 就有 $f\in L^{1+1/n}(\Omega)$, 此外
\begin{equation*}
\norm{f}_{1+1/n}\le\norm{f}_1^{\alpha_n}\norm{f}_q^{1-\alpha_n},
\end{equation*}
其中 $\dfrac{1}{1+(1/n)}=\dfrac{\alpha_n}{1}+\dfrac{1-\alpha_n}{q}$.由
Baire 纲定理推出存在某个整数 $n_0$ 使得 $\operatorname{Int}X_{n_0}\ne\emptyset$.
因此 $X\subset L^{1+1/n_0}(\Omega)$.

2. 恒等映射 $I:X\to L^p(\Omega)$ 是图为闭的线性算子.因此它是有界算子.

\solutionhead{4.9}

对每个 $t\in\R$, 在 $\Omega$ 上几乎处处有 $f(x)t\le j(f(x))+j^\star(t)$,
积分得到
\begin{equation*}
\left(\frac{1}{\abs{\Omega}}\int_\Omega f\right)t
\le\frac{1}{\abs{\Omega}}\int_\Omega j(f)+j^\star(t).
\end{equation*}
因此
\begin{equation*}
j\left(\frac{1}{\abs{\Omega}}\int_\Omega f\right)
=\sup_{t\in\R}\left\{\left(\frac{1}{\abs{\Omega}}\int_\Omega f\right)t-j^\star(t)\right\}
\le\frac{1}{\abs{\Omega}}\int_\Omega j(f).
\end{equation*}

\solutionhead{4.10}

1. 设 $u_1,u_2\in D(J)$ 且 $t\in[0,1]$.函数 $x\mapsto j(tu_1(x)+(1-t)u_2(x))$
是可测的 (因为 $j$ 连续).另一方面, $j(tu_1+(1-t)u_2)\le tj(u_1)+(1-t)j(u_2)$.
回忆存在常数 $a$ 与 $b$ 使得对一切 $s\in\R$, $j(s)\ge as+b$ (见命题  \ref{prop:1.10}).
由此推出 $j(tu_1+(1-t)u_2)\in L^1(\Omega)$, 且
$J(tu_1+(1-t)u_2)\le tJ(u_1)+(1-t)J(u_2)$.

2. 先设 $j\ge0$.我们断言对每个 $\lambda\in\R$, 集合
$\{u\in L^p(\Omega);\ J(u)\le\lambda\}$ 是闭的.事实上, 设 $(u_n)$ 为
$L^p(\Omega)$ 中的序列使得在 $L^p(\Omega)$ 中 $u_n\to u$ 且 $\int_\Omega j(u_n)\le\lambda$.
抽取子序列可以假设 $u_n\to u$ 几乎处处.由 Fatou 引理推出 $j(u)\in L^1(\Omega)$
且 $\int_\Omega j(u)\le\lambda$.因此 $J$ 下半连续.在一般情形下, 设
$\tilde j(s)=j(s)-(as+b)\ge0$.我们已知 $\widetilde J$ 下半连续, 从而
$J(u)=\widetilde J(u)+a\int_\Omega u+b\abs{\Omega}$ 也下半连续.

3. 我们先断言: 对使 $j^\star(f)\in L^1(\Omega)$ 的每个 $f\in L^{p'}(\Omega)$, 有
\begin{equation*}
J^\star(f)\le\int_\Omega j^\star(f).
\end{equation*}
事实上, 在 $\Omega$ 上几乎处处有 $fu-j(u)\le j^\star(f)\ \forall u\in L^p(\Omega)$,
从而
\begin{equation*}
\sup_{u\in D(J)}\left\{\int_\Omega fu-J(u)\right\}\le\int_\Omega j^\star(f).
\end{equation*}

反向不等式的证明更精细, 需要某种``正则化''步骤.先设 $1<p<\infty$, 并令
\begin{equation*}
j_n(t)=j(t)+\frac1n\abs{t}^p,\quad t\in\R.
\end{equation*}
我们断言对一切 $f\in L^{p'}(\Omega)$,
\begin{equation}
J_n^\star(f)=\int_\Omega j_n^\star(f).
\tag{S1}
\label{eq:S4.10}
\end{equation}
事实上, 设 $f\in L^{p'}(\Omega)$.对几乎处处固定的 $x\in\Omega$,
$\sup_{u\in\R}\left\{f(x)u-j(u)-\frac1n\abs{u}^p\right\}$ 由某个唯一的元素
$u=u(x)$ 达到.显然我们有
$j(u(x))+\frac1n\abs{u(x)}^p-f(x)u(x)\le j(0)$.由此推出 $u\in L^p(\Omega)$
且 $j(u)\in L^1(\Omega)$ (为什么?).我们得出
\begin{equation*}
J_n^\star(f)=\sup_{v\in D(J)}\left\{\int_\Omega fv-J_n(v)\right\}
\ge\int_\Omega\left(fu-j(u)-\frac1n\abs{u}^p\right)=\int_\Omega j_n^\star(f).
\end{equation*}
由于我们已经建立了反向不等式, 可见 \eqref{eq:S4.10} 成立.

接着令 $n\uparrow\infty$.显然 $J\le J_n$, 从而 $J_n^\star\le J^\star$, 即
$\int_\Omega j_n^\star(f)\le J^\star(f)$.我们断言对每个 $s\in\R$, 当
$n\uparrow\infty$ 时 $j_n^\star(s)\uparrow j^\star(s)$.事实上, 我们知道
$j_n^\star=j^\star\mathbin{\triangledown}\left(\frac1n\abs{\cdot}^p\right)^\star$
(见习题 1.23), 然后可以像习题 1.24 那样论证.由单调收敛定理得出: 若
$f\in D(J^\star)$, 则 $j^\star(f)\in L^1(\Omega)$ 且
$\int_\Omega j^\star(f)\le J^\star(f)$.

最后, 若 $p=1$, 可以例如使用 $j_n(t)=j(t)+\frac1nt^2$ 来修改上述方法.

4. 先假设在 $\Omega$ 上几乎处处有 $f(x)\in\partial j(u(x))$, 我们有
\begin{equation*}
j(v)-j(u(x))\ge f(x)(v-u(x))\quad\forall v\in\R,\ \text{在}\ \Omega\ \text{上几乎处处}.
\end{equation*}
选取 $v=0$, 可见 $j(u)\in L^1(\Omega)$, 从而对一切 $v\in D(J)$,
$J(v)-J(u)\ge\int_\Omega f(v-u)$.

反之, 假设 $f\in\partial J(u)$.则我们有 $J(u)+J^\star(f)=\int_\Omega fu$.
因此 $j(u)\in L^1(\Omega)$, $j^\star(f)\in L^1(\Omega)$, 且
$\int_\Omega\{j(u)+j^\star(f)-fu\}=0$.由于几乎处处有 $j(u)+j^\star(f)-fu\ge0$,
我们发现几乎处处有 $j(u)+j^\star(f)-fu=0$, 即几乎处处 $f(x)\in\partial j(u(x))$.

\solutionhead{4.11}

令 $f=u^\alpha$, $g=v^\alpha$, 以及 $p=1/\alpha$.须证明
\begin{equation}
\left(\int_\Omega f\right)^p+\left(\int_\Omega g\right)^p
\le\left[\int_\Omega(f^p+g^p)^{1/p}\right]^p.
\tag{S1}
\label{eq:S4.11}
\end{equation}
令 $a=\int_\Omega f$ 与 $b=\int_\Omega g$, 于是我们有
\begin{align*}
a^p+b^p&=\int_\Omega(a^{p-1}f+b^{p-1}g)\le\int_\Omega(a^p+b^p)^{1/p'}(f^p+g^p)^{1/p}\\
&=(a^p+b^p)^{1/p'}\int_\Omega(f^p+g^p)^{1/p}.
\end{align*}
由此推出 $(a^p+b^p)^{1/p}\le\int_\Omega(f^p+g^p)^{1/p}$, 即 \eqref{eq:S4.11}.

\solutionhead{4.12}

1. 只须证明
\begin{equation*}
\inf_{t\in[-1,+1]}\left\{\frac{(\abs{t}^p+1)^{1-s}(\abs{t}^p+1-2\abs{\frac{t+1}{2}}^p)^s}{\abs{t-1}^p}\right\}>0,
\end{equation*}
或等价地证明
\begin{equation*}
\inf_{t\in[-1,+1]}\left\{\frac{\abs{t}^p+1-2\abs{\frac{t+1}{2}}^p}{\abs{t-1}^2}\right\}>0.
\end{equation*}
但函数 $\phi(t)=\abs{t}^p+1-2\abs{\frac{t+1}{2}}^p$ 满足对一切 $t\in[-1,+1)$,
$\phi(t)>0$, 且 $\phi(1)=\phi'(1)=0$, $\phi''(1)>0$.

\solutionhead{4.14}

1 与 2. 令 $g_n=\chi_{S_n(\alpha)}$, 从而 $g_n\to0$ 几乎处处且 $\abs{g_n}\le1$.
由控制收敛定理 (因为 $\abs{\Omega}<\infty$) 推出 $\int_\Omega g_n\to0$, 即
$\abs{S_n(\alpha)}\to0$.

3. 给定任何整数 $m\ge1$, 我们可以把第 2 问应用于 $\alpha=1/m$, 找到整数
$N_m$ 使得 $\abs{S_{N_m}(1/m)}<\delta/2^m$.令 $\Omega_m=S_{N_m}(1/m)$,
我们得到对一切 $k\ge N_m$ 及一切 $x\in\Omega\setminus\Omega_m$,
$\abs{f_k(x)-f(x)}\le1/m$.最后, 令 $A=\bigcup_{m=1}^\infty\Omega_m$, 从而
$\abs{A}<\delta$.我们断言 $f_n\to f$ 在 $\Omega\setminus A$ 上一致收敛.
事实上, 给定 $\varepsilon>0$, 固定满足 $m_0>1/\varepsilon$ 的整数 $m_0$.显然
对一切 $k\ge N_{m_0}$ 及一切 $x\in\Omega\setminus\Omega_{m_0}$,
$\abs{f_k(x)-f(x)}<\varepsilon$, 从而对一切 $k\ge N_{m_0}$ 及一切
$x\in\Omega\setminus A$, $\abs{f_k(x)-f(x)}<\varepsilon$.

4. 给定 $\varepsilon>0$, 先利用 (i) 固定某个 $\delta>0$, 然后利用第 3 问
固定某个 $A$.我们得到对一切 $n$, $\int_A\abs{f_n}^p\le\varepsilon$, 且
$f_n\to f$ 在 $\Omega\setminus A$ 上一致收敛.由 Fatou 引理推出
$\int_A\abs{f}^p\le\varepsilon$, 从而
\begin{equation*}
\int_\Omega\abs{f_n-f}^p=\int_A\abs{f_n-f}^p+\int_{\Omega\setminus A}\abs{f_n-f}^p
\le2^p\varepsilon+\abs{\Omega}\norm{f_n-f}_{L^\infty(\Omega\setminus A)}^p.
\end{equation*}

\solutionhead{4.15}

1(iv). 注意对一切 $\phi\in C_c(\Omega)$, $\int_\Omega f_n\phi\to0$.反设
$f_{n_k}\rightharpoonup f$ 弱收敛于 $\sigma(L^1,L^\infty)$.由此推出对一切
$\phi\in C_c(\Omega)$, $\int_\Omega f\phi=0$, 从而 (由推论  \ref{cor:4.24})
$f=0$ 几乎处处.此外 $\int_\Omega f_{n_k}\to\int_\Omega f=0$; 但
$\int_\Omega f_{n_k}=\int_0^{n_k}e^{-t}dt\to1$; 矛盾.

2(iv). 注意对一切 $\phi\in C_c(\Omega)$, $\int_\Omega g_n\phi\to0$, 并利用
$C_c(\Omega)$ 在 $L^{p'}(\Omega)$ 中稠密的事实 (因为 $p'<\infty$).

\solutionhead{4.16}

1. 我们先验证: 若序列 $(f_n)$ 满足
\begin{equation}
f_n\rightharpoonup\tilde f\quad\text{弱收敛于}\ \sigma(L^p,L^{p'})
\tag{S1}
\label{eq:S4.16a}
\end{equation}
且
\begin{equation}
f_n\to f\quad\text{几乎处处},
\tag{S2}
\label{eq:S4.16b}
\end{equation}
则 $f=\tilde f$ 几乎处处.

事实上, 由习题 3.4 我们知道存在 $L^p(\Omega)$ 中的序列 $(g_n)$ 使得
$g_n\in\operatorname{conv}\{f_n,f_{n+1},\dots\}$ 且在 $L^p(\Omega)$ 中
$g_n\to\tilde f$ 强收敛.由 \eqref{eq:S4.16b} 推出 $g_n\to f$ 几乎处处.
另一方面 (由定理  \ref{thm:4.9}), 存在子序列 $(g_{n_k})$ 使得
$g_{n_k}\to\tilde f$ 几乎处处.因此 $f=\tilde f$ 几乎处处.

现在在假设 (i) 与 (ii) 下验证 $f_n\rightharpoonup f$ 弱收敛于 $\sigma(L^p,L^{p'})$.
存在弱收敛于某个极限 $\tilde f$ 的子序列 $(f_{n_k})$.由前面的讨论我们知道
$f=\tilde f$ 几乎处处.``极限的唯一性''蕴含整个序列 $(f_n)$ 弱收敛于 $f$.

3. 第一种方法. 写出
\begin{equation}
\norm{f_n-f}_q\le\norm{f_n-T_kf_n}_q+\norm{T_kf_n-T_kf}_q+\norm{T_kf-f}_q.
\tag{S3}
\label{eq:S4.16c}
\end{equation}
注意对每个 $k>0$, 有
\begin{equation*}
\int_\Omega\abs{f_n-T_kf_n}^q\le\int_{[\abs{f_n}\ge k]}\abs{f_n}^q.
\end{equation*}
另一方面, $\int_\Omega\abs{f_n}^p\le C^p$, 从而
\begin{equation*}
k^{p-q}\int_{[\abs{f_n}\ge k]}\abs{f_n}^q\le C^p.
\end{equation*}
由此推出对一切 $n$,
\begin{equation}
\norm{f_n-T_kf_n}_q\le\left(\frac{C^p}{k^{p-q}}\right)^{1/q}.
\tag{S4}
\label{eq:S4.16d}
\end{equation}
取极限 ($n\to\infty$), 借助 Fatou 引理我们得到
$\norm{f-T_kf}_q\le\left(\dfrac{C^p}{k^{p-q}}\right)^{1/q}$.给定
$\varepsilon>0$, 固定足够大的 $k$ 使得 $(C^p/k^{p-q})^{1/q}<\varepsilon$.
显然 (由控制收敛定理) 当 $n\to\infty$ 时 $\norm{T_kf_n-T_kf}_q\to0$, 从而
存在某个整数 $N$ 使得对一切 $n\ge N$, $\norm{T_kf_n-T_kf}_q<\varepsilon$.
结合 \eqref{eq:S4.16c} 与 \eqref{eq:S4.16d}, 可见对一切 $n\ge N$,
$\norm{f_n-f}_q<3\varepsilon$.

第二种方法. 由 Egorov 定理, 给定 $\delta>0$, 存在某个 $A\subset\Omega$ 使得
$\abs{A}<\delta$ 且 $f_n\to f$ 在 $\Omega\setminus A$ 上一致收敛.写出
\begin{align*}
\int_\Omega\abs{f_n-f}^q&=\int_{\Omega\setminus A}+\int_A
\le\norm{f_n-f}_{L^\infty(\Omega\setminus A)}^q\abs{\Omega}+\norm{f_n-f}_p^q\abs{A}^{1-(q/p)}\\
&\le\norm{f_n-f}_{L^\infty(\Omega\setminus A)}^q\abs{\Omega}+(2C)^q\delta^{1-(q/p)},
\end{align*}
这导出对一切 $\delta>0$, $\limsup_{n\to\infty}\int_\Omega\abs{f_n-f}^q\le(2C)^q\delta^{1-(q/p)}$.

\solutionhead{4.17}

1. 由齐次性, 只须验证
\begin{equation*}
\sup_{t\in[-1,+1]}\left\{\frac{\abs{\abs{t+1}^p-\abs{t}^p-1}}{\abs{t}^{p-1}+\abs{t}}\right\}<\infty.
\end{equation*}

\solutionhead{4.18}

1. 首先容易验证对一切 $a,b\in(0,1)$, $\int_a^bu_n(t)dt\to(b-a)\overline f$.
只要 $1<p\le\infty$, 这就蕴含 $u_n\rightharpoonup\overline f$ 弱收敛于
$\sigma(L^p,L^{p'})$ (因为 $p'<\infty$, 阶梯函数在 $L^{p'}$ 中稠密).
当 $p=1$ 即 $f\in L^1_{\text{loc}}(\R)$ 时, 存在一个 $T$-周期函数
$g\in L^\infty(\R)$ 使得 $\frac1T\int_0^T\abs{f-g}<\varepsilon$ (其中
$\varepsilon>0$ 任意固定).令 $v_n(x)=g(nx)$, $x\in(0,1)$, 并设 $\phi\in L^\infty(0,1)$.
我们有
\begin{equation*}
\abs{\int u_n\phi-\overline f\int\phi}\le3\varepsilon\norm{\phi}_\infty+\abs{\int v_n\phi-\overline g\int\phi},
\end{equation*}
从而对一切 $\varepsilon>0$, $\limsup_{n\to\infty}\abs{\int u_n\phi-\overline f\int\phi}\le3\varepsilon\norm{\phi}_\infty$.
由此推出 $u_n\rightharpoonup\overline f$ 弱收敛于 $\sigma(L^1,L^\infty)$.

2. $\lim_{n\to\infty}\norm{u_n-\overline f}_p=\left(\frac1T\int_0^T\abs{f-\overline f}^p\right)^{1/p}$.

3. (i) $u_n\overset{*}{\rightharpoonup}0$ 关于 $\sigma(L^\infty,L^1)$.

(ii) $u_n\overset{*}{\rightharpoonup}\frac12(\alpha+\beta)$ 关于 $\sigma(L^\infty,L^1)$.

\solutionhead{4.20}

1. 设 $(u_n)$ 为 $L^p(\Omega)$ 中的序列使得在 $L^p(\Omega)$ 中 $u_n\to u$
强收敛.存在子序列使得 $u_{n_k}(x)\to u(x)$ 几乎处处, 且对一切 $k$,
$\abs{u_{n_k}}\le v$, 其中 $v\in L^p(\Omega)$ (见定理  \ref{thm:4.9}).由
控制收敛定理推出在 $L^q(\Omega)$ 中 $Au_{n_k}\to Au$ 强收敛.``极限的唯一性''
蕴含整个序列 $(Au_n)$ 在 $L^q(\Omega)$ 中强收敛于 $Au$.

2. 考虑习题 4.18 第 3(ii) 问中定义的序列 $(u_n)$.注意 $u_n\rightharpoonup\frac12(\alpha+\beta)$,
而 $Au_n\rightharpoonup\frac12(a(\alpha)+a(\beta))$.由此推出对一切 $\alpha,\beta\in\R$,
$a\left(\frac{\alpha+\beta}{2}\right)=\frac12(a(\alpha)+a(\beta))$, 因此 $a$
必为仿射函数.

\solutionhead{4.21}

1. 验证对每个有界区间 $I$, $\int_I u_n(t)dt\to0$.然后利用 (具有紧支集的)
阶梯函数在 $L^{p'}(\R)$ 中的稠密性.

2. 我们再次断言对每个有界区间 $I$, $\int_I u_n(t)dt\to0$.事实上, 给定
$\varepsilon>0$, 固定 $\delta>0$ 使得 $\delta(\norm{u_0}_\infty+\abs{I})<\varepsilon$.
令 $E=[\abs{u_0}>\delta]$, 写出
\begin{equation*}
\int_I u_n(t)dt=\int_{I+n}u_0(t)dt=\int_{(I+n)\cap E}u_0+\int_{(I+n)\cap E^c}u_0.
\end{equation*}
选取充分大的 $N$ 使得对一切 $n>N$, $\abs{(I+n)\cap E}<\delta$.我们得到
对一切 $n\ge N$, $\abs{\int_I u_n(t)dt}\le\delta\norm{u_0}_\infty+\delta\abs{I}<\varepsilon$.
然后利用 (具有紧支集的) 阶梯函数在 $L^1(\R)$ 中的稠密性.

3. 反设 $u_{n_k}\rightharpoonup u$ 弱收敛于 $\sigma(L^1,L^\infty)$.考虑由
$f=\sum_i(-1)^i\chi_{(-n_i,-n_{i+1})}$ 定义的函数 $f\in L^\infty(\R)$.注意
$\int u_{n_k}f=(-1)^k$ 不收敛.

\solutionhead{4.22}

1. 为了证明 (B) $\Rightarrow$ (A), 利用由形如 $\chi_E$ (其中 $E$ 可测且
$\abs{E}<\infty$) 的函数张成的向量空间在 $L^{p'}(\Omega)$ 中稠密的事实
(只要 $p'<\infty$).

2. 利用由形如 $\chi_E$ (其中 $E\subset\Omega$ 且 $E$ 可测) 的函数张成的向量
空间在 $L^\infty(\Omega)$ 中稠密的事实.

4. 给定 $\varepsilon>0$, 固定某个满足 $\abs{\omega}<\infty$ 的可测子集
$\omega\subset\Omega$ 使得
\begin{equation}
\int_{\omega^c}f<\varepsilon.
\tag{S1}
\label{eq:S4.22}
\end{equation}
我们有 $\int_{\omega^c}f_n=\int_{\omega^c}f+\left(\int_\omega f-\int_\omega f_n\right)+\left(\int_\Omega f_n-\int_\Omega f\right)$,
从而由 (b) 与 (c), $\int_{\omega^c}f_n=\int_{\omega^c}f+o(1)$.另一方面,
$\int_F f_n=\int_{F\cap\omega}f_n+\int_{F\cap\omega^c}f_n=\int_{F\cap\omega}f+\int_{F\cap\omega^c}f_n+o(1)$,
从而 $\int_F f_n-\int_F f=\int_{F\cap\omega^c}(f_n-f)+o(1)$.结合 \eqref{eq:S4.22},
我们得到 $\abs{\int_F f_n-\int_F f}\le2\varepsilon+o(1)$.由此推出
$\int_F f_n\to\int_F f$.最后, 我们利用由函数 $\chi_F$ ($F\subset\Omega$ 可测且
$\abs{F}\le\infty$) 张成的向量空间在 $L^\infty(\Omega)$ 中稠密的事实.

\solutionhead{4.23}

1. 设 $(u_n)$ 为 $C$ 中的序列使得在 $L^p(\Omega)$ 中 $u_n\to u$ 强收敛.存在
子序列 $(u_{n_k})$ 使得 $u_{n_k}\to u$ 几乎处处.从而 $u\ge f$ 几乎处处.

2. 假设 $u\in L^\infty(\Omega)$ 满足对满足 $f\phi\in L^1(\Omega)$ 且 $\phi\ge0$
的一切 $\phi\in L^1(\Omega)$, $\int_\Omega u\phi\ge\int_\Omega f\phi$.我们
断言 $u\ge f$ 几乎处处.事实上, 写出 $\Omega=\bigcup_n\Omega_n$, 其中
$\abs{\Omega_n}<\infty$, 并令 $\Omega'_n=\Omega_n\cap[\abs{f}<n]$, 从而
$\bigcup_n\Omega'_n=\Omega$.令 $A=[u<f]$.选取 $\phi=\chi_{A\cap\Omega'_n}$,
我们发现 $\int_{A\cap\Omega'_n}\abs{f-u}\le0$, 从而对一切 $n$,
$\abs{A\cap\Omega'_n}=0$.由此推出 $\abs{A}=0$.

3. 注意若固定满足 $f\phi\in L^1(\Omega)$ 的 $\phi\in L^1(\Omega)$, 则集合
$\{u\in L^\infty(\Omega);\ \int_\Omega u\phi\ge\int_\Omega f\phi\}$ 关于拓扑
$\sigma(L^\infty,L^1)$ 是闭的.

\solutionhead{4.24}

1. 对每个 $\phi\in L^1(\R^N)$, 我们有
\begin{equation*}
\int v_n\phi=\int u\zeta_n(\check\rho_n\star\phi)
=\int u\zeta_n(\check\rho_n\star\phi-\phi)+\int u\zeta_n\phi,
\end{equation*}
从而
\begin{equation*}
\abs{\int v_n\phi-\int v\phi}\le\norm{u}_\infty\norm{\check\rho_n\star\phi-\phi}_1
+\norm{u}_\infty\norm{(\zeta_n-\zeta)\phi}_1.
\end{equation*}
右边的第一项由定理  \ref{thm:4.22} 趋于零, 而右边的第二项由控制收敛定理趋于零.

2. 设 $B=B(x_0,R)$, 设 $\chi$ 表示 $B(x_0,R+1)$ 的特征函数.令
$\tilde v_n=\rho_n\star(\zeta_n\chi u)$.注意在 $B(x_0,R)$ 上 $\tilde v_n=v_n$,
因为 $\operatorname{supp}(\tilde v_n-v_n)\subset B(0,1/n)+B(x_0,R+1)^c$.另一方面,
我们有
\begin{align*}
\int_B\abs{v_n-v}&=\int_B\abs{\tilde v_n-\chi v}\le\int_{\R^N}\abs{\tilde v_n-\chi v}\\
&\le\int_{\R^N}\abs{\rho_n\star(\zeta_n-\zeta)\chi u}+\int_{\R^N}\abs{(\rho_n\star\chi v)-\chi v}\\
&\le\int_{\R^N}\abs{(\zeta_n-\zeta)\chi u}+\int_{\R^N}\abs{(\rho_n\star\chi v)-\chi v}\to0.
\end{align*}

\solutionhead{4.25}

1. 用 $\bar u$ 记 $u$ 在 $\Omega$ 外延拓为 $0$ 所得的函数.令
$\Omega_n=\{x\in\Omega;\ \dist(x,\partial\Omega)>2/n\ \text{且}\ \abs{x}<n\}$.
设 $\zeta_n$ (相应地 $\zeta$) 表示 $\Omega_n$ (相应地 $\Omega$) 的特征函数,
从而在 $\R^N$ 上 $\zeta_n\to\zeta$.令 $v_n=\rho_n\star(\zeta_n\bar u)$.我们
知道 $v_n\in C_c^\infty(\Omega)$, 且对每个球 $B$, $\int_B\abs{v_n-u}\to0$ (由习题
4.24).因此对每个球 $B$, 存在 (依赖于 $B$ 的) 子序列在 $B$ 上几乎处处收敛于
$u$.由对角线法我们可以构造在 $\R^N$ 上几乎处处收敛于 $u$ 的子序列 $(v_{n_k})$.

\solutionhead{4.26}

1. 假设 $A<\infty$.证明 $f\in L^1(\Omega)$ 且 $\norm{f}_1\le A$.对一切
$\phi\in C_c(\Omega)$, 我们有 $\abs{\int_\Omega f\phi}\le A\norm{\phi}_\infty$.
设 $K\subset\Omega$ 为任何紧子集, 设 $\psi\in C_c(\Omega)$ 为满足在 $\Omega$
上 $0\le\psi\le1$ 且在 $K$ 上 $\psi=1$ 的函数.设 $u$ 为 $L^\infty(\Omega)$
中的任何函数.利用习题 4.25 我们可以构造 $C_c(\Omega)$ 中的序列 $(u_n)$ 使得
$\norm{u_n}_\infty\le\norm{u}_\infty$ 且在 $\Omega$ 上 $u_n\to u$ 几乎处处.
我们有 $\abs{\int_\Omega f\psi u_n}\le A\norm{u}_\infty$.由控制收敛定理令
$n\to\infty$, 得到对一切 $u\in L^\infty(\Omega)$, $\abs{\int_\Omega f\psi u}\le A\norm{u}_\infty$.
选取 $u=\operatorname{sign}(f)$, 发现对 $\Omega$ 的每个紧子集 $K$,
$\int_K\abs{f}\le A$.由此推出 $f\in L^1(\Omega)$ 且 $\norm{f}_1\le A$.

2. 假设 $B<\infty$.对满足 $\phi\ge0$ 的一切 $\phi\in C_c(\Omega)$, 有
$\int_\Omega f\phi\le B\norm{\phi}_\infty$.使用与第 1 问相同的方法, 得到
对满足 $u\ge0$ 的一切 $u\in L^\infty(\Omega)$, $\int_\Omega f\psi u\le B\norm{u}_\infty$.
选取 $u=\chi_{[f>0]}$ 发现 $\int_K f^+\le B$.

\solutionhead{4.27}

先考察一个抽象框架.设 $E$ 为向量空间, 设 $f,g$ 为 $E$ 上的两个线性泛函,
满足 $f\not\equiv0$.假设
\begin{equation*}
[\phi\in E\ \text{且}\ f(\phi)>0]\implies[g(\phi)\ge0].
\end{equation*}
我们断言存在常数 $\lambda\ge0$ 使得 $g=\lambda f$.事实上, 固定任何满足
$f(\phi_0)=1$ 的 $\phi_0\in E$.对每个 $\phi\in E$ 与每个 $\varepsilon>0$,
有 $f(\phi-f(\phi)\phi_0+\varepsilon\phi_0)=\varepsilon>0$, 从而
$g(\phi-f(\phi)\phi_0+\varepsilon\phi_0)\ge0$.由此推出对一切 $\phi\in E$,
$g(\phi)\ge\lambda f(\phi)$, 从而 $g=\lambda f$, 其中 $\lambda=g(\phi_0)\ge0$.

\runin{应用.} $E=C_c^\infty(\Omega)$, $f(\phi)=\int_\Omega u\phi$, 以及
$g(\phi)=\int_\Omega v\phi$.

\solutionhead{4.30}

1 与 2. 注意 $\dfrac{1}{p'}+\dfrac{1}{q'}+\dfrac1r=1$ 以及 $(1-\alpha)r=p$,
$(1-\beta)r=q$.对几乎所有的 $x\in\R^N$, 写出
\begin{equation*}
\abs{f(x-y)g(y)}=\phi_1(y)\phi_2(y)\phi_3(y),
\end{equation*}
其中 $\phi_1(y)=\abs{f(x-y)}^\alpha$, $\phi_2(y)=\abs{g(y)}^\beta$, 以及
$\phi_3(y)=\abs{f(x-y)}^{1-\alpha}\abs{g(y)}^{1-\beta}$.显然
$\phi_1\in L^{q'}(\R^N)$ 且 $\phi_2\in L^{p'}(\R^N)$.另一方面,
$\abs{\phi_3(y)}^r=\abs{f(x-y)}^p\abs{g(y)}^q$.由定理  \ref{thm:4.15} 推出:
对几乎所有的 $x\in\R^N$, 函数 $y\mapsto\abs{\phi_3(y)}^r$ 可积.由 Hölder
不等式 (见习题 4.4) 推出: 对几乎所有的 $x\in\R^N$, 函数 $y\mapsto\abs{f(x-y)g(y)}$
可积, 且
\begin{equation*}
\int_{\R^N}\abs{f(x-y)}\abs{g(y)}dy\le\norm{f}_p^\alpha\norm{g}_q^\beta
\left(\int_{\R^N}\abs{f(x-y)}^p\abs{g(y)}^qdy\right)^{1/r}.
\end{equation*}
因此
\begin{equation*}
\abs{(f\star g)(x)}^r\le\norm{f}_p^{\alpha r}\norm{g}_q^{\beta r}
\int_{\R^N}\abs{f(x-y)}^p\abs{g(y)}^qdy,
\end{equation*}
从而
\begin{equation*}
\int_{\R^N}\abs{(f\star g)(x)}^rdx\le\norm{f}_p^{\alpha r}\norm{g}_q^{\beta r}
\norm{f}_p^p\norm{g}_q^q=\norm{f}_p^r\norm{g}_q^r.
\end{equation*}

3. 若 $1<p<\infty$ 且 $1<q<\infty$, 存在 $C_c(\R^N)$ 中的序列 $(f_n)$ 与
$(g_n)$ 使得在 $L^p(\R^N)$ 中 $f_n\to f$, 在 $L^q(\R^N)$ 中 $g_n\to g$.
则 $f_n\star g_n\in C_c(\R^N)$, 此外 $\norm{(f_n\star g_n)-(f\star g)}_\infty\to0$.
由此推出当 $\abs{x}\to\infty$ 时 $(f\star g)(x)\to0$.

\solutionhead{4.34}

给定任何 $\varepsilon>0$, 存在 $\mathcal{F}$ 在 $L^p(\R^N)$ 中被半径为
$\varepsilon$ 的球的有限覆盖, 设 $\mathcal{F}\subset\bigcup_{i=1}^kB(f_i,\varepsilon)$.

2. 对每个 $i$, 存在某个 $\delta_i>0$ 使得对一切满足 $\abs{h}<\delta_i$ 的
$h\in\R^N$, $\norm{\tau_hf_i-f_i}_{L^p(\R^N)}<\varepsilon$ (见引理  \ref{lem:4.3}).
令 $\delta=\min_{1\le i\le k}\delta_i$.容易验证对一切 $f\in\mathcal{F}$ 及一切
满足 $\abs{h}<\delta$ 的 $h\in\R^N$, $\norm{\tau_hf-f}_p<3\varepsilon$.

3. 对每个 $i$, 存在某个有界开集 $\Omega_i\subset\R^N$ 使得
$\norm{f_i}_{L^p(\R^N\setminus\Omega_i)}<\varepsilon$.令
$\Omega=\bigcup_{i=1}^k\Omega_i$, 验证对一切 $f\in\mathcal{F}$,
$\norm{f}_{L^p(\R^N\setminus\Omega)}<2\varepsilon$.

\solutionhead{4.37}

1. 写出
\begin{align*}
\int_I u_n(x)\phi(x)dx&=\int_{-n}^{+n}f(t)\left[\phi\left(\frac{t}{n}\right)-\phi(0)\right]dt
+\phi(0)\int_{-n}^{+n}f(t)dt\\
&=A_n+B_n;
\end{align*}
由 Lebesgue 定理 $A_n\to0$, 且由于 $\int_{-\infty}^{+\infty}f(t)dt=0$,
$B_n\to0$.

2. 注意对所有 $\delta>0$,
\begin{equation*}
\int_0^\delta\abs{u_n(x)}dx=\int_0^{n\delta}\abs{f(t)}dt\to\int_0^\infty\abs{f(t)}dt>0.
\end{equation*}

3. 反设成立.我们将对一切 $\phi\in C([-1,+1])$ 有 $\int_I u\phi=0$, 从而
(由推论  \ref{cor:4.24}) $u\equiv0$.另一方面, 若选取 $\phi=\chi_{(0,1)}$,
我们得到 $\int_I u_n\phi=\int_0^nf(t)dt\to\int_0^{+\infty}f(t)dt>0$.矛盾.

\solutionhead{4.38}

2. 验证对一切 $\phi\in C^1([0,1])$, 当 $n\to\infty$ 时
$\int_I u_n\phi=\int_I\phi+O(1/n)$.然后利用 $\norm{u_n}_1$ 有界以及
$C^1([0,1])$ 在 $C([0,1])$ 中稠密的事实.

3. 序列 $(u_n)$ 不能是等度可积的, 因为 $\abs{\operatorname{supp}u_n}\to0$ 且
$1=\int_I u_n=\int_{\operatorname{supp}u_n}\abs{u_n}$.

4. 若 $u_{n_k}\rightharpoonup u$ 弱收敛于 $\sigma(L^1,L^\infty)$, 则由第
2 问与推论  \ref{cor:4.24}, 我们会有 $u\equiv1$.选取进一步的子序列 $(u_{n'_k})$
使得 $\sum_k\abs{\operatorname{supp}u_{n'_k}}<1$.设 $\phi=\chi_A$, 其中
$A=I\setminus\left(\bigcup_k\operatorname{supp}u_{n'_k}\right)$, 从而
$\abs{A}>0$.我们对一切 $k$ 有 $\int_I u_{n'_k}\phi=0$, 从而
$0=\int_I\phi=\abs{A}$.矛盾.

5. 考虑满足 $\sum_k\abs{\operatorname{supp}u_{n_k}}<\infty$ 的子序列 $(u_{n_k})$.
令 $B_k=\bigcup_{j\ge k}(\operatorname{supp}u_{n_j})$ 且 $B=\bigcap_kB_k$.
显然当 $k\to\infty$ 时 $\abs{B_k}\to0$, 从而 $\abs{B}=0$.若 $x\notin B$,
则存在某个 $k_0$ 使得对一切 $k\ge k_0$, $u_{n_k}(x)=0$.
\solutionhead{5.1}

1. 对 $a=u+v$ 与 $b=v$ 应用平行四边形法则给出 $(u,2v)=2(u,v)$.

2. 计算 (i) $-$ (ii) $+$ (iii).

3. 注意由 $(\ ,\ )$ 的定义, 映射 $\lambda\in\R\mapsto(\lambda u,v)$ 是连续的.

\solutionhead{5.2}

设 $A$ 为满足 $0<\abs{A}<\abs{\Omega}$ 的可测集, 选取满足 $A\cap B=\emptyset$ 且
$0<\abs{B}<\abs{\Omega}$ 的可测集 $B$.令 $u=\chi_A$ 与 $v=\chi_B$.

先设 $1\le p<\infty$.我们有
\begin{equation*}
\norm{u+v}_p^p=\norm{u-v}_p^p=\abs{A}+\abs{B},
\end{equation*}
从而 $\norm{u+v}_p^2+\norm{u-v}_p^2=2(\abs{A}+\abs{B})^{2/p}$.另一方面,
我们有 $2(\norm{u}_p^2+\norm{v}_p^2)=2(\abs{A}^{2/p}+\abs{B}^{2/p})$.最后,
注意对一切 $\alpha,\beta>0$, 当 $p<2$ 时
$(\alpha+\beta)^{2/p}>\alpha^{2/p}+\beta^{2/p}$, 当 $p>2$ 时
$(\alpha+\beta)^{2/p}<\alpha^{2/p}+\beta^{2/p}$.用相同的函数 $u$ 与 $v$
考察 $p=\infty$ 的情形.

\solutionhead{5.3}

验证
\begin{equation}
2(t_nu_n-t_mu_m,u_n-u_m)=(t_n+t_m)\abs{u_n-u_m}^2+(t_n-t_m)(\abs{u_n}^2-\abs{u_m}^2),
\tag{S1}
\label{eq:S5.3}
\end{equation}
这蕴含对一切 $m,n$, $(t_n-t_m)(\abs{u_n}^2-\abs{u_m}^2)\le0$.

1. 设 $n>m$, 从而 $t_n\ge t_m$, 进而 $\abs{u_n}\le\abs{u_m}$ (注意若
$t_n=t_m$, 由 \eqref{eq:S5.3} 有 $u_n=u_m$).另一方面, 对 $n>m$, 我们有
\begin{equation*}
(t_n+t_m)\abs{u_n-u_m}^2\le(t_n-t_m)(\abs{u_m}^2-\abs{u_n}^2)\le t_n(\abs{u_m}^2-\abs{u_n}^2),
\end{equation*}
从而 $\abs{u_n-u_m}^2\le\abs{u_m}^2-\abs{u_n}^2$.由此推出当 $n\uparrow\infty$
时 $\abs{u_n}\downarrow\ell$, 且 $(u_n)$ 是 Cauchy 序列.

2. 设 $n>m$, 从而 $t_m\ge t_n$ 且 $\abs{u_m}\le\abs{u_n}$.对 $n>m$, 有
\begin{equation*}
(t_n+t_m)\abs{u_n-u_m}^2\le(t_m-t_n)(\abs{u_n}^2-\abs{u_m}^2)\le t_m(\abs{u_n}^2-\abs{u_m}^2),
\end{equation*}
从而 $\abs{u_n-u_m}^2\le\abs{u_n}^2-\abs{u_m}^2$.我们有下述抉择:

(i) 或者当 $n\uparrow\infty$ 时 $\abs{u_n}\uparrow\infty$,

(ii) 或者当 $n\uparrow\infty$ 时 $\abs{u_n}\uparrow\ell<\infty$, 此时 $(u_n)$ 是 Cauchy 序列.

另一方面, 令 $v_n=t_nu_n$ 与 $s_n=1/t_n$, 得到
\begin{equation*}
(s_nv_n-s_mv_m,v_n-v_m)\le0,
\end{equation*}
从而由第 1 问 $(v_n)$ 收敛于一个极限.由此推出若 $t_n\to t>0$, 则 $(u_n)$
也收敛于一个极限.最后若 $t_n\to0$, 两种情形 (i) 与 (ii) 均可出现.例如取
$H=\R$, 情形 (i) 取 $u_n=C/t_n$, 情形 (ii) 取 $u_n=C$.

\solutionhead{5.4}

注意
\begin{equation*}
\abs{v-u}^2=\abs{v-f}^2-\abs{u-f}^2+2(f-u,v-u).
\end{equation*}

\solutionhead{5.5}

1. 设 $K=\overline{\bigcup_nK_n}$.我们断言 $u_n\to u=P_Kf$.首先注意序列
$d_n=\abs{f-u_n}=\dist(f,K_n)$ 单调递增且有上界.从而当 $n\uparrow\infty$
时 $d_n\uparrow\ell<\infty$.接着应用平行四边形法则 (取 $a=f-u_n,b=f-u_m$),
得到
\begin{equation*}
\abs{f-\frac{u_n+u_m}{2}}^2+\abs{\frac{u_n-u_m}{2}}^2
=\frac12(\abs{f-u_n}^2+\abs{f-u_m}^2).
\end{equation*}
由此推出若 $m\ge n$, 则 $\abs{u_n-u_m}^2\le2(d_m^2-d_n^2)$.因此 $(u_n)$
收敛于极限 $u$, 且显然 $u\in K$.另一方面, 对一切 $v\in K_n$ 有
$\abs{f-u_n}\le\abs{f-v}$, 特别地对一切 $v\in K$ 有 $\abs{f-u_n}\le\abs{f-v}$.
取极限得到对一切 $v\in K$, $\abs{f-u}\le\abs{f-v}$.

2. 显然 $K=\overline{\bigcup_nK_n}$ 是凸的 (为什么?).我们断言
$u_n\to u=P_Kf$.首先注意序列 $d_n=\abs{f-u_n}=\dist(f,K_n)$ 单调递减,
从而 $d_n\to\ell$.接着 (用上面的相同方法), 若 $m\ge n$, 有
$\abs{u_n-u_m}^2\le2(d_n^2-d_m^2)$.因此 $(u_n)$ 收敛于极限 $u$, 且显然
$u\in K$.最后注意只要 $m\ge n$, 对一切 $v\in K_n$ 就有 $\abs{f-u_m}\le\abs{f-v}$.
取极限 ($m\to\infty$) 导出对一切 $v\in\bigcup_nK_n$ 有 $\abs{f-u}\le\abs{f-v}$,
由稠密性对一切 $v\in K$ 成立.

序列 $(\alpha_n)$ 单调递减, 从而收敛于极限 $\alpha$.我们断言
$\alpha=\inf_K\varphi$.首先, 显然 $\inf_K\varphi\le\alpha_n$, 从而
$\inf_K\varphi\le\alpha$.另一方面, 设 $u$ 为 $K$ 中的任意元素, 设
$u_n=P_{K_n}u$.在不等式 $\alpha_n\le\varphi(u_n)$ 中取极限, 得到
$\alpha\le\varphi(u)$ (因为 $u_n\to u$).由此推出 $\alpha\le\inf_K\varphi$.

\solutionhead{5.6}

1. 例如考虑 $\norm{u}\ge1$ 且 $\norm{v}\le1$ 的情形.我们有
\begin{equation*}
\norm{Tu-Tv}=\norm{\frac{u}{\norm{u}}-v}=\frac{\norm{(u-v)+(v-v\norm{u})}}{\norm{u}}
\le\norm{u-v}+\norm{u}-1\le2\norm{u-v},
\end{equation*}
因为 $\norm{u}\le\norm{u-v}+\norm{v}\le\norm{u-v}+1$.

2. 设 $u=(1,0)$ 与 $v=(1,\alpha)$.则有
$\norm{Tu-Tv}=2\abs{\alpha}/(1+\abs{\alpha})$, 而 $\norm{u-v}=\abs{\alpha}$.
选取 $\alpha\ne0$ 且任意小即得出结论.

3. $T$ 与 $P_{B_E}$ 重合.只须验证若 $\norm{u}\ge1$, 则对一切 $v\in B_E$,
\begin{equation*}
\left(u-\frac{u}{\norm{u}},v-\frac{u}{\norm{u}}\right)\le0.
\end{equation*}

\solutionhead{5.10}

(i) $\Rightarrow$ (ii). 写出对一切 $t\in(0,1)$ 与一切 $v\in K$,
$F(u)\le F((1-t)u+tv)$, 这蕴含
\begin{equation*}
\frac1t[F(u+t(v-u))-F(u)]\ge0.
\end{equation*}
令 $t\to0$ 即得 (ii).

(ii) $\Rightarrow$ (i). 我们断言对一切 $u,v\in H$,
$F(v)-F(u)\ge(F'(u),v-u)$.事实上, 函数 $t\in\R\mapsto\phi(t)=F(u+t(v-u))$
属于 $C^1$ 类且是凸的.因此 $\phi(1)-\phi(0)\ge\phi'(0)$.

\solutionhead{5.12}

$T$ 满当且仅当 $E$ 完备.

1. 通过令对一切 $u,v\in E$,
\begin{equation*}
((T(u),T(v)))=(u,v),
\end{equation*}
把 $E$ 的标量积转移到 $R(T)$ 上.注意对一切 $f,g\in R(T)$,
$\abs{((f,g))}\le\norm{f}_{E^\star}\norm{g}_{E^\star}$.标量积 $(( , ))$
可以通过连续性与稠密性延拓到 $\overline{R(T)}$ 上, 它现在配有 Hilbert
空间的结构.

2. 固定任何 $f\in E^\star$.映射 $g\in R(T)\mapsto\dual{f}{T^{-1}(g)}$ 是
$R(T)$ 上的连续线性泛函.它可以 (通过连续性) 延拓到 $\overline{R(T)}$.在
$\overline{R(T)}$ 中应用 Riesz--Fréchet 表示定理, 得到某个元素
$h\in\overline{R(T)}$ 使得对一切 $g\in R(T)$, $((h,g))=\dual{f}{T^{-1}(g)}$.
从而对一切 $v\in E$, $((h,T(v)))=\dual{f}{v}$.另一方面, 对一切
$h\in\overline{R(T)}$ 及一切 $v\in E$, 有 $((h,Tv))=\dual{h}{v}$ (当
$h\in R(T)$ 时这显然).由此推出 $f=h$, 从而 $f\in R(T)$, 即 $R(T)=E^\star$.

3. 我们构造了等距映射 $T:E\to E^\star$, 且 $R(T)$ 在 $E^\star$ 中稠密.由于
$E^\star$ 完备, 我们得出 (在同构意义下) $E^\star$ 是 $E$ 的完备化.

\solutionhead{5.13}

1. 我们断言平行四边形法则成立.事实上, 设 $f\in F(u)$ 与 $g\in F(v)$.则
$f\pm g\in F(u\pm v)$, 从而我们有
\begin{equation*}
\dual{f+g}{u+v}=\norm{u+v}^2\quad\text{与}\quad\dual{f-g}{u-v}=\norm{u-v}^2.
\end{equation*}
把这些关系式相加得出 $2(\norm{u}^2+\norm{v}^2)=\norm{u+v}^2+\norm{u-v}^2$.

2. 设 $T:E\to E^\star$ 为习题 5.12 中引入的映射.我们断言 $F(u)=\{T(u)\}$.
显然 $T(u)\in F(u)$.另一方面, 我们知道 $E^\star$ 关于对偶范数
$\norm{\cdot}_{E^\star}$ 是 Hilbert 空间.特别地, $E^\star$ 严格凸, 从而
(见习题 1.1) $F(u)$ 退化为单点.

\solutionhead{5.14}

凸性不等式 $a(tu+(1-t)v,tu+(1-t)v)\le ta(u,u)+(1-t)a(v,v)$ 等价于
$t(1-t)a(u-v,u-v)\ge0$.

考虑由 $a(u,v)=(Au,v)\ \forall u,v\in H$ 定义的算子 $A\in\LL(H)$.则
$F'(u)=Au+A^\star u$, 因为我们有
\begin{equation*}
F(u+h)-F(u)=(Au+A^\star u,h)+a(h,h).
\end{equation*}

\solutionhead{5.15}

首先通过连续性把 $S$ 延拓为算子 $\overline S:G\to F$.接着令
$T=\overline S\circ P_G$, 其中 $P_G$ 表示由 $H$ 到 $G$ 上的投影算子.

\solutionhead{5.18}

(ii) $\Rightarrow$ (i). 假设 (ii) 蕴含 $T$ 是单射且 $R(T)$ 闭.从而 $R(T)$
有补空间 (因为 $H$ 是 Hilbert 空间).由定理  \ref{thm:2.13} 我们推出 $T$ 有左逆.

(i) $\Rightarrow$ (ii). 假设 (i) 蕴含 $T$ 是单射且 $R(T)$ 闭.然后使用定理~
\ref{thm:2.21}.

\solutionhead{5.19}

注意
\begin{equation*}
\limsup_{n\to\infty}\abs{u_n-u}^2=\limsup_{n\to\infty}(\abs{u_n}^2-2(u_n,u)+\abs{u}^2)\le0.
\end{equation*}

\solutionhead{5.20}

1. 若 $u\in N(S)$, 则对一切 $v\in H$, $(Sv,v-u)\ge0$; 用 $tv$ 替换 $v$,
可见对一切 $v\in H$, $(Sv,u)=0$.反之, 若 $u\in R(S)^\perp$, 则对一切
$v\in H$, $(Sv-Su,v)\ge0$; 用 $tv$ 替换 $v$, 可见对一切 $v\in H$, $(Su,v)=0$
(另见问题 16).

2. 应用推论  \ref{cor:5.8} (Lax--Milgram).

3. 方法 (a). 令 $u_t=(I+tS)^{-1}f$.

若 $f\in N(S)$, 则对一切 $t>0$, $u_t=f$.

若 $f\in R(S)$, 写出 $f=Sv$, 从而 $u_t+S(tu_t-v)=0$.由此推出
$(u_t,tu_t-v)\le0$, 从而 $\abs{u_t}\le(1/t)\abs{v}$.因此当 $t\to\infty$
时 $u_t\to0$.由稠密性, 仍可证明对每个 $f\in\overline{R(S)}$, 当 $t\to\infty$
时 $u_t\to0$.

在一般情形 $f\in H$ 下, 把 $f$ 写成 $f=f_1+f_2$, 其中 $f_1=P_{N(S)}f$ 且
$f_2=P_{\overline{R(S)}}f$.

方法 (b). 我们有 $u_t+tSu_t=f$, 从而 $\abs{u_t}\le\abs{f}$.抽取子序列
$t_n\to\infty$, 可以假设 $u_{t_n}\rightharpoonup u$ 弱收敛且 $Su=0$ (为什么?),
即 $u\in N(S)$.由第 1 问我们知道对一切 $v\in N(S)$, $(Su_t,v)=0$, 从而对一切
$v\in N(S)$, $(f-u_t,v)=0$.取极限, 发现对一切 $v\in N(S)$, $(f-u,v)=0$.
因此 $u=P_{N(S)}f$, ``极限的唯一性''蕴含当 $t\to\infty$ 时 $u_t\rightharpoonup u$
弱收敛.另一方面, 我们有 $(Su_t,u_t)\ge0$, 即 $(f-u_t,u_t)\ge0$, 从而
$\limsup_{t\to\infty}\abs{u_t}^2\le(f,u)=\abs{u}^2$.由此推出当 $t\to\infty$
时 $u_t\to u$ 强收敛.

\solutionhead{5.21}

1. 令 $S=I-T$ 并应用习题 5.20 第 1 问.

2. 写出 $f=u-Tu$, 注意 $\sigma_n(f)=\dfrac1n(u-T^nu)$.

3. 首先验证对一切 $f\in R(I-T)$, $\lim_{n\to\infty}\sigma_n(f)=0$.接着把一般的
$f\in H$ 分解为 $f=f_1+f_2$, 其中 $f_1\in N(I-T)$ 且
$f_2\in N(I-T)^\perp=\overline{R(I-T)}$.则我们有
$\sigma_n(f)=\sigma_n(f_1)+\sigma_n(f_2)=f_1+\sigma_n(f_2)$.

4. 依次对 $u,Su,S^2u,\dots,S^iu,\dots$ 应用不等式 (1), 并把所得的不等式相加.
注意对一切 $i=0,1,\dots,n$,
\begin{equation*}
\abs{S^nu-S^{n+1}u}\le\abs{S^iu-S^{i+1}u}.
\end{equation*}

5. 写出 $f=u-Tu=2(u-Su)$, 我们得到 $\abs{\mu_n(f)}\le2\abs{u}/\sqrt{n+1}$.

6. 使用与第 3 问相同的方法.

\solutionhead{5.25}

2. 设 $m>n$.对 $f=u_m$ 与 $v=P_Ku_n$ 应用习题 5.4, 得到
\begin{align*}
\abs{P_Ku_n-P_Ku_m}^2&\le\abs{P_Ku_n-u_m}^2-\abs{P_Ku_m-u_m}^2\\
&\le\abs{P_Ku_n-u_n}^2-\abs{P_Ku_m-u_m}^2.
\end{align*}
因此 $(P_Ku_n)$ 是 Cauchy 序列.

3. 可以假设 $u_{n_k}\rightharpoonup u$ 弱收敛.回忆对一切 $v\in K$,
$(u_n-P_Ku_n,v-P_Ku_n)\le0$.沿子序列 $n_k$ 取极限导出对一切 $v\in K$,
$(u-\ell,v-\ell)\le0$.由于 $u\in K$, 可以取 $v=u$ 并得出 $u=\ell$.再次,
``极限的唯一性''蕴含 $u_n\rightharpoonup\ell$ 弱收敛.

4. 对每个 $v\in K$, $\lim_{n\to\infty}\abs{u_n-v}^2$ 存在, 从而对每个
$v,w\in K$, $\lim_{n\to\infty}(u_n,v-w)$ 也存在.由此推出对每个 $z\in H$,
$\phi(z)=\lim_{n\to\infty}(u_n,z)$ 存在.由 Riesz--Fréchet 表示定理我们可以把
$\phi(z)$ 写成 $\phi(z)=(u,z)$, 其中 $u\in H$.最后注意对一切 $v\in K$,
$(u-\ell,v-\ell)\le0$, 从而 $\ell=P_Ku$.

5. 经平移与伸缩后总可以假设 $K=B_H$.从而 $\abs{u_n}\downarrow\alpha$.
若 $\alpha<1$, 则对充分大的 $n$, $u_n=P_Ku_n$ (我们已知 $P_Ku_n$ 强收敛).
若 $\alpha\ge1$, 则 $P_Ku_n=u_n/\abs{u_n}$ 强收敛, 从而 $u_n$ 也强收敛.

6. 回忆对一切 $v\in K$, $(u_n-P_Ku_n,v-P_Ku_n)\le0$, 从而对一切 $v\in K$,
$(u_n-\ell,v-\ell)\le\varepsilon_n$, 其中 $\varepsilon_n\to0$ ($\varepsilon_n$
依赖于 $v$).把这些不等式相加得出对一切 $v\in K$,
$(\sigma_n-\ell,v-\ell)\le\varepsilon'_n$, 其中 $\varepsilon'_n\to0$.假设
$\sigma_{n_k}\rightharpoonup\sigma$ 弱收敛, 则 $\sigma\in K$ 满足对一切
$v\in K$, $(\sigma-\ell,v-\ell)\le0$.因此 $\sigma=\ell$, ``极限的唯一性''
蕴含 $\sigma_n\rightharpoonup\ell$ 弱收敛.

\solutionhead{5.26}

3. 注意 $\sqrt{n}u_n$ 有界, 且对每个固定的 $j$, 当 $n\to\infty$ 时
$(\sqrt{n}u_n,e_j)\to0$.

\solutionhead{5.27}

设 $F$ 为由 $E_n$ 张成的向量空间的闭包.我们知道 (见定理  \ref{thm:5.9} 的证明)
对一切 $u\in H$, $\sum_{n=1}^\infty\abs{P_{E_n}u}^2=\abs{P_Fu}^2$, 从而对一切
$u\in D$, $\abs{P_Fu}=\abs{u}$.由此推出对一切 $u\in D$, $\abs{P_Fu}^2=\abs{u}^2$,
因此对一切 $u\in D$, $P_{F^\perp}u=0$.从而对一切 $u\in H$, $P_{F^\perp}u=0$,
即 $F^\perp=\{0\}$, 于是 $F=H$.

\solutionhead{5.28}

1. 由命题  \ref{prop:3.25}, $V$ 可分.考虑 $V$ 的稠密可数子集 $(v_n)$, 并像
定理  \ref{thm:5.11} 的证明那样得出结论.

\solutionhead{5.29}

2. 若 $2<p<\infty$, 使用不等式 $\norm{u}_p\le\norm{u}_\infty^{1-2/p}\norm{u}_2^{2/p}$.
注意每个无限维 Hilbert 空间 (可分或不可分) 都包含一个无限标准正交序列.

6. 在 $\Omega$ 上积分, 我们发现 $k\le M^2\abs{\Omega}$, 这给出了 $E$ 的维数的
一个上界.

\solutionhead{5.30}

1. 对每个固定的 $t\in[0,1]$, 考虑函数 $u_t(s)=p(s)\chi_{[0,t]}(s)$, 并写出
$\sum_{n=1}^\infty\abs{(u_t,e_n)}^2\le\norm{u_t}_2^2$.

2. (2) 中的等号蕴含对几乎所有的 $t\in[0,1]$, (1) 中等号成立.因此对几乎所有的
$t\in[0,1]$, $u_t=\sum_{n=1}^\infty(u_t,e_n)e_n$, 从而
$u_t\in E=\overline{\operatorname{span}\{e_n\}}$.还须验证由函数 $(u_t)$ 张成的
空间在 $L^2$ 中稠密.设 $f\in L^2$ 使得对几乎所有的 $t$, $\int_0^1fu_t=0$.由此
推出对一切 $t\in[0,1]$, $\int_0^tfp=0$, 从而 $fp=0$ 几乎处处.

\solutionhead{5.31}

容易验证当 $i\ne j$ 时 $(\phi_i,\phi_j)=0$.令 $n=2^{p+1}-1$.设 $E$ 为由
$\{\phi_0,\phi_1,\dots,\phi_n\}$ 张成的空间, $F$ 为由区间
$(\frac{i}{2^{p+1}},\frac{i+1}{2^{p+1}})$ ($0\le i\le2^{p+1}-1$) 的特征函数张成的
空间.显然 $E\subset F$, $\dim E=n+1=2^{p+1}$, 且 $\dim F=2^{p+1}$.因此 $E=F$.

\solutionhead{5.32}

2. 函数 $u=r_1r_2$ 与所有函数 $(r_i)_{i\ge0}$ 正交, 且 $u\ne0$.因此
$(r_i)_{i\ge0}$ 不是基.

3. 容易验证 $(w_n)_{n\ge0}$ 是标准正交系, 且 $w_0=r_0$, 对一切 $\ell\ge0$,
$w_{2^\ell}=r_{\ell+1}$.为了证明 $(w_n)_{n\ge0}$ 是基, 可以使用与习题 5.31
相同的论证.
\solutionhead{6.2}

3. 考虑定义在 $[0,1]$ 上的函数序列
\begin{equation*}
u_n(t)=\begin{cases}
0 & \text{若}\ 0\le t\le\frac12,\\
n\left(t-\frac12\right) & \text{若}\ \frac12<t\le\frac12+\frac1n,\\
1 & \text{若}\ \frac12+\frac1n<t\le1.
\end{cases}
\end{equation*}
注意 $T(u_n)\to f$, 但 $f\notin T(B_E)$, 因为 $f\notin C^1([0,1])$.

\solutionhead{6.3}

反证法.若结论不成立, 则存在某个 $\delta>0$ 使得对一切 $u\in E$,
$\norm{Tu}_F\ge\delta\norm{u}_E$.因此 $R(T)$ 闭.考虑由 $T_0=T$ 定义的
算子 $T_0:E\to R(T)$.显然 $T_0$ 是双射.由推论 2.6,
% 原书此处印作 "By Corollary 2.6" (印刷页 413); 第~
% \ref{chap:2}~章并无编号 2.6 的推论
% (2.6 是定理~
\ref{thm:2.6} 开映射定理), 有界逆定理是推论~
\ref{cor:2.7}. 依技能规范保留原样
% 并记录, 不擅自更正, 故此处不用 \ref 以免产生未定义引用.
$T_0^{-1}\in\LL(R(T),E)$.另一方面, $T_0\in\KK(E,R(T))$.因此 $B_E$ 紧
且 $\dim E<\infty$.

\solutionhead{6.5}

设 $T:V\to\ell^2$ 为由
\begin{equation*}
Tu=(\sqrt{\lambda_1}u_1,\sqrt{\lambda_2}u_2,\dots,\sqrt{\lambda_n}u_n,\dots)
\end{equation*}
定义的算子.显然对一切 $u\in V$, $\abs{Tu}_{\ell^2}=\norm{u}_V$, 且 $T$ 是由
$V$ 到 $\ell^2$ 上的满射.由于 $\ell^2$ 完备, 推出 $V$ 也完备.

考虑由 $J_nu=(u_1,u_2,\dots,u_n,0,0,\dots)$ 定义的算子 $J_n:V\to\ell^2$.容易
验证 $\norm{J_n-I}_{\LL(V,\ell^2)}\to0$, 从而由 $V$ 到 $\ell^2$ 中的典范
单射是紧的.

\solutionhead{6.7}

1. 假设 $T$ 是由 $E$ 弱到 $F$ 强的连续映射.则对每个 $\varepsilon>0$, 存在
$0$ 在 $E$ 弱中的邻域 $V$ 使得 $x\in V\Rightarrow\norm{Tx}<\varepsilon$.
可以假设 $V$ 具有形式
\begin{equation*}
V=\{x\in E;\ \abs{\dual{f_i}{x}}<\delta\ \forall i=1,2,\dots,n\},
\end{equation*}
其中 $f_1,\dots,f_n\in E^\star$, $\delta>0$.

设 $M=\{x\in E;\ \dual{f_i}{x}=0\ \forall i=1,\dots,n\}$, 从而对一切 $x\in M$,
$Tx=0$.另一方面, $M$ 具有有限余维 (见 \S \ref{sec:2.4} 例 2).因此
$E=M+N$, 其中 $\dim N<\infty$.由此推出 $R(T)=T(N)$ 是有限维的.

2. 注意若在 $E$ 中 $u_n\rightharpoonup u$ 弱收敛, 则在 $F$ 中
$Tu_n\rightharpoonup Tu$ 弱收敛.另一方面, $(Tu_n)$ 在 $F$ 中有紧闭包
(关于强拓扑).因此 $Tu_n\to Tu$ (参见例如习题 3.5).

6. 注意 $T^\star\in\LL(E^\star,(c_0)^\star)$.但 $(c_0)^\star=\ell^1$ (见
\S \ref{sec:11.3}).由于 $E^\star$ 自反, 由第 5 问推出 $T^\star$ 是紧的.
因此 (由定理  \ref{thm:6.4}) $T$ 紧.

\solutionhead{6.8}

1. 存在常数 $c$ 使得 $B_{R(T)}\subset cT(B_E)$, 因此 $R(T)$ 的单位球是紧的.

2. 设 $E_0$ 为 $N(T)$ 的补空间.则 $T_0=T|_{E_0}$ 是由 $E_0$ 到 $R(T)$ 上的
双射.因此 $\dim E_0=\dim R(T)<\infty$.

\solutionhead{6.9}

1. (A) $\Rightarrow$ (B):
设 $E_0$ 为 $N(T)$ 的补空间, 设 $P:E\to N(T)$ 为相应的投影算子.则
$T_0=T|_{E_0}$ 是由 $E_0$ 到 $R(T)$ 上的双射.由开映射定理存在常数 $C$
使得对一切 $u\in E_0$, $\norm{u}_E\le C\norm{Tu}_F$.由此推出对一切 $u\in E$,
\begin{equation*}
\norm{u}_E\le\norm{u-Pu}_E+\norm{Pu}_E\le C\norm{Tu}_F+\norm{Pu}_E.
\end{equation*}

(C) $\Rightarrow$ (A):
(i) 为了验证 $N(T)$ 中的单位球是紧的, 设 $(u_n)$ 为 $N(T)$ 中的序列使得
$\norm{u_n}_E\le1$.由于 $(Q(u_n))$ 在 $G$ 中有紧闭包, 可以抽取在 $G$ 中收敛的
子序列 $(Q(u_{n_k}))$.应用 (C), 可见 $(u_{n_k})$ 是 Cauchy 序列.

(ii) 引入 $N(T)$ 的补空间, 还可以额外假设 $T$ 是单射.设 $(u_n)$ 为 $E$ 中的
序列使得 $Tu_n\to f$.先验证 $(u_n)$ 有界.若不然, 令 $v_n=u_n/\norm{u_n}$.
应用 (C), 可见子序列 $(v_{n_k})$ 是 Cauchy 序列.设 $v_{n_k}\to v$, 其中
$v\in N(T)$ 且 $\norm{v}=1$; 不可能.因此 $(u_n)$ 有界, 我们可以抽取在 $G$ 中
收敛的子序列 $(Q(u_{n_k}))$.再次应用 (C), 发现 $(u_{n_k})$ 是 Cauchy 序列.

为了恢复习题 2.12 中的结果, 写出
\begin{equation*}
\norm{u}_E\le C(\norm{Tu}_F+\norm{Pu}_E)\le C(\norm{Tu}_F+\abs{Pu}),
\end{equation*}
因为 $N(T)$ 上的所有范数等价.此外,
\begin{equation*}
\abs{Pu}\le\abs{u-Pu}+\abs{u}\le C\norm{u-Pu}_E+\abs{u}\le C\norm{Tu}_F+\abs{u}.
\end{equation*}

2. 注意
\begin{equation*}
\norm{u}_E\le C(\norm{Tu}_F+\norm{Pu}_E)\le C(\norm{(T+S)u}_F+\norm{Pu}_E+\norm{Su}_F),
\end{equation*}
并考虑由 $Qu=[Pu,Su]$ 定义的紧算子 $Q:E\to E\times F$.

\solutionhead{6.10}

1. 注意对一切 $u\in E$,
\begin{align*}
\abs{Q(1)}\norm{u}&\le\norm{Q(1)u-Q(T)u}+\norm{Q(T)u}\\
&=\norm{\widetilde Q(T)(u-Tu)}+\norm{Q(T)u}\le C(\norm{u-Tu}+\norm{Q(T)u}).
\end{align*}

2. 蕴含关系 $N(I-T)=\{0\}\Rightarrow R(I-T)=E$ 的证明: 反设
$R(I-T)=E_1\ne E$.令 $E_n=(I-T)^nE$.则 $(E_n)$ 是闭子空间的递减序列.选取
$u_n\in E_n$ 使得 $\norm{u_n}=1$ 且 $\dist(u_n,E_{n+1})\ge1/2$.写出
\begin{equation*}
Q(T)u_n-Q(T)u_m=Q(T)u_n-Q(1)u_n+Q(1)u_n-Q(1)u_m+Q(1)u_m-Q(T)u_m.
\end{equation*}
因此对 $m>n$, 有 $\norm{Q(T)u_n-Q(T)u_m}\ge\abs{Q(1)}/2$, 这是不可能的.

关于逆命题, 沿用定理  \ref{thm:6.6} 的证明中叙述的论证.

3. 沿用定理  \ref{thm:6.6} 的证明中的相同记号, 写出 $S=T+\widetilde\Lambda\circ P$.
这里 $S\notin\KK(E)$, 但 $\widetilde\Lambda\circ P\in\KK(E)$.因此
$Q(S)\in\KK(E)$ (为什么?).然后像定理  \ref{thm:6.6} 的证明那样继续.

\solutionhead{6.11}

1. 存在整数 $n_0\ge1$ 使得 $\operatorname{Int}F_{n_0}\ne\emptyset$, 从而
$B(u_0,\rho)\subset F_{n_0}$.对每个 $u\in F$ 及满足 $\abs{\lambda}<\rho/\norm{u}$
的 $\lambda$, 有 $u_0+\lambda u\in F_{n_0}$.因此
\begin{equation*}
\abs{\lambda}\abs{u(x)-u(y)}\le\abs{u_0(x)-u_0(y)}+n_0d(x,y)^{1/n_0}
\le2n_0d(x,y)^{1/n_0}.
\end{equation*}
由此推出对一切 $x,y\in K$,
\begin{equation*}
\abs{u(x)-u(y)}\le\frac{2n_0}{\rho}\norm{u}d(x,y)^{1/n_0}.
\end{equation*}

2. Ascoli--Arzelà 定理蕴含 $B_F$ 是紧的.

\solutionhead{6.13}

反设存在某个 $\varepsilon_0>0$ 与序列 $(u_n)$ 使得 $\norm{u_n}_E=1$ 且
$\norm{Tu_n}_F\ge\varepsilon_0+n\abs{u_n}$.则 $\abs{u_n}\to0$, 可以假设
$u_{n_k}\rightharpoonup u$ 弱收敛.但函数 $u\mapsto\abs{u}$ 是凸且连续的.
因此它对弱拓扑下半连续, 从而 $u=0$.由此推出 $Tu_n\to0$.不可能.

\solutionhead{6.15}

1. 若 $u=f+\lambda(T-\lambda I)^{-1}f$, 则 $\lambda u=T(u-f)$, 从而
$\abs{\lambda}\norm{u}\le\norm{T}(\norm{u}+\norm{f})$.

2. 由命题  \ref{prop:6.7} 的证明我们知道: 若 $\mu\in\R$ 使得
$\abs{\mu-\lambda}\norm{(T-\lambda I)^{-1}}<1$, 则 $\mu\in\rho(T)$.因此
$\dist(\lambda,\sigma(T))\ge1/\norm{(T-\lambda I)^{-1}}$.

4. $(U-I)^{-1}=\frac12(T-I)$.

6. 注意关系式 $Uu-\lambda u=f$ 等价于
\begin{equation*}
Tu-\frac{\lambda+1}{\lambda-1}u=\frac{1}{\lambda-1}(f-Tf).
\end{equation*}

\solutionhead{6.16}

2. $(T-\lambda I)^{-1}=\dfrac{1}{1-\lambda^n}\sum_{i=0}^{n-1}\lambda^{n-i-1}T^i$.

3. $(T-\lambda I)^{-1}=-\sum_{i=0}^{n-1}\lambda^{-i-1}T^i$.

4. $(I-T)^{-1}=(I-T^n)^{-1}\sum_{i=0}^{n-1}T^i$.

\solutionhead{6.18}

1. $\norm{S_r}=\norm{S_\ell}=1$.注意 $S_\ell\circ S_r=I$, 从而
$S_r\notin\KK(E),S_\ell\notin\KK(E)$.

3. 对每个 $\lambda\in[-1,+1]$, 算子 $(S_r-\lambda I)$ 不是满射: 例如若
$f=(-1,0,0,\dots)$, 则方程 $S_rx-\lambda x=f$ 没有解 $x\in\ell^2$.

4. $N(S_\ell-\lambda I)=\R(1,\lambda,\lambda^2,\dots)$.

6. $S_r^\star=S_\ell$ 且 $S_\ell^\star=S_r$.

7. 写出 $S_rx-\lambda x=f$, 我们有
\begin{equation*}
\abs{x}=\abs{S_rx}=\abs{\lambda x+f}\le\abs{\lambda}\abs{x}+\abs{f}.
\end{equation*}
从而 $\abs{S_rx-\lambda x}\ge(1-\abs{\lambda})\abs{x}$, 因此 $R(S_r-\lambda I)$ 闭.
应用定理  \ref{thm:2.19} 给出
\begin{equation*}
R(S_r-\lambda I)=N(S_\ell-\lambda I)^\perp
=\left\{x\in\ell^2;\ \sum_{i=1}^\infty\lambda^{i-1}x_i=0\right\}
\end{equation*}
以及 $R(S_\ell-\lambda I)=N(S_r-\lambda I)^\perp=E$.

8. 我们有 $R(S_r\pm I)=N(S_\ell\pm I)^\perp=E$ 以及
$R(S_\ell\pm I)=N(S_r\pm I)^\perp=E$.由第 3 问已知 $R(S_r\pm I)\ne E$.
另一方面, $R(S_\ell\pm I)\ne E$; 否则, 由于 $S_\ell\pm I$ 是单射, 会有
$\pm1\in\rho(S_\ell)$, 矛盾.

9. 若对一切 $n$, $\alpha_n\ne0$, 则 $\operatorname{EV}(S_r\circ M)=\emptyset$;
若对某个 $n$, $\alpha_n=0$, 则 $\operatorname{EV}(S_r\circ M)=\{0\}$.

10. 总可以假设 $\alpha\ne0$; 否则 $S_r\circ M$ 紧, 结论显然.

我们证明对满足 $\abs{\lambda}>\abs{\alpha}$ 的每个 $\lambda$, $(T-\lambda I)$
是双射.注意 $M=\alpha I+K$, 其中 $K$ 是紧算子.令 $T=S_r\circ M$, 得到
$T=\alpha S_r+K_1$ 且 $(T-\lambda I)=(\alpha S_r-\lambda I)+K_1=J\circ(I+K_2)$,
其中 $J=(\alpha S_r-\lambda I)$ 是双射, $K_1,K_2$ 紧.应用定理  \ref{thm:6.6}(c),
只须验证 $N(T-\lambda I)=\{0\}$.这已在第 9 问中建立.

我们证明对 $\abs{\lambda}\le\abs{\alpha}$, $(T-\lambda I)$ 不是双射.反设
$(T-\lambda I)$ 是双射.写出
$(S_r-\frac\lambda\alpha I)=\frac1\alpha(T-\lambda I)-\frac1\alpha K_1=J'\circ(I+K_3)$,
其中 $J'$ 是双射且 $K_3$ 紧.再次应用定理  \ref{thm:6.6}(c), 可见
$(S_r-\frac\lambda\alpha I)$ 单射当且仅当它满射.但由第 2 与第 3 问已知
对 $\abs{\lambda}\le\abs{\alpha}$, $(S_r-\frac\lambda\alpha I)$ 是单射且不是满射.
矛盾.

11. $\sigma(S_r\circ M)=[-\sqrt{\abs{ab}},+\sqrt{\abs{ab}}]$.事实上, 若
$\abs{\lambda}\le\sqrt{\abs{ab}}$, 算子 $(S_r\circ M-\lambda I)$ 不是满射,
因为例如 $f=(-1,0,0,\dots)\notin R(S_r\circ M-\lambda I)$.另一方面, 若
$\abs{\lambda}>\sqrt{\abs{ab}}$, 算子 $(S_r\circ M-\lambda I)$ 是双射, 因为
$(S_r\circ M)^2=abS_r^2$.从而 $\norm{(S_r\circ M)^2}\le\abs{ab}$, 可以应用
习题 6.16 第 4 问.

\solutionhead{6.20}

1. 注意 $\abs{Tu(x)-Tu(y)}\le\abs{x-y}^{1/p'}\norm{u}_p$.若 $1<p<\infty$,
可以应用 Ascoli--Arzelà 定理得出 $T(B_E)$ 在 $C([0,1])$ 中有紧闭包, 从而在
$L^p(0,1)$ 中有紧闭包.若 $p=1$, 应用定理  \ref{thm:4.26}.

2. $\operatorname{EV}(T)=\emptyset$.首先注意 $0\notin\operatorname{EV}(T)$.
事实上, 方程 $Tu=0$ 蕴含对一切 $a,b\in[0,1]$, $\int_0^1u\chi_{[a,b]}=0$.
若 $1<p<\infty$, 利用阶梯函数在 $L^{p'}$ 中的稠密性得出 $u\equiv0$.当 $p=1$
时, 通过用阶梯函数一致逼近 $\phi$, 证明对一切 $\phi\in C([0,1])$ 有
$\int_0^1u\phi=0$.由推论  \ref{cor:4.24} 得出 $u\equiv0$.

3. 对 $\lambda\ne0$ 与 $f\in C([0,1])$, 令 $u=(T-\lambda I)^{-1}f$.则
$v(x)=\int_0^xu(t)dt$ 满足 $v\in C^1([0,1])$ 且 $v-\lambda v'=f$, $v(0)=0$.
因此
\begin{equation*}
u(x)=-\frac1\lambda f(x)-\frac{1}{\lambda^2}\int_0^xe^{(x-t)/\lambda}f(t)dt.
\end{equation*}
同样的公式对 $f\in L^p$ 仍然有效 (用稠密性论证).

4. $(T^\star v)(x)=\int_x^1v(t)dt$.

\solutionhead{6.22}

2. 反设存在某个 $\mu\in Q(\sigma(T))$ 使得 $\mu\notin\sigma(Q(T))$.则
$\mu=Q(\lambda)$, 其中 $\lambda\in\sigma(T)$, 且 $Q(T)-Q(\lambda)I=S$ 是双射.
我们可以写出对一切 $t\in\R$, $Q(t)-Q(\lambda)=(t-\lambda)\widetilde Q(t)$,
从而 $(T-\lambda I)\widetilde Q(T)=\widetilde Q(T)(T-\lambda I)=S$.因此
$T-\lambda I$ 是双射, $\lambda\in\rho(T)$; 矛盾.

3. 取 $E=\R^2$, $T=\begin{pmatrix}0&1\\-1&0\end{pmatrix}$ 且 $Q(t)=t^2$.
则 $\operatorname{EV}(T)=\sigma(T)=\emptyset$ 且
$\operatorname{EV}(T^2)=\sigma(T^2)=\{-1\}$.

4. $T^2+I$ 由 Lax--Milgram 定理是双射.每个没有实根的 2 次多项式可以写成
$Q(t)=t^2+at+b=(t+a/2)^2+b-a^2/4$, 其中 $b-a^2/4>0$, 我们可以再次应用
Lax--Milgram 定理.若多项式 $Q(t)$ 没有实根, 则其根为共轭复数.我们可以写出
$Q(t)=Q_1(t)Q_2(t)\dots Q_\ell(t)$, 其中每个 $Q_i(t)$ 是没有实根的 2 次
多项式.由于 $Q_i(T)$ 是双射, 对 $Q(T)$ 也成立同样结论.

5. (i) 反设 $\mu\in\operatorname{EV}(Q(T))$ 且 $\mu\notin Q(\operatorname{EV}(T))$.
则存在 $u\ne0$ 使得 $Q(T)u=\mu u$.写出
\begin{equation*}
Q(t)-\mu=(t-t_1)(t-t_2)\dots(t-t_q)\widetilde Q(t),
\end{equation*}
其中 $t_i$ 是多项式 $Q(t)-\mu$ 的实根, $\widetilde Q$ 没有实根.由于
$\mu\notin Q(\operatorname{EV}(T))$, 对一切 $i$ 有 $t_i\notin\operatorname{EV}(T)$.
我们有 $(T-t_1I)\dots(T-t_kI)\widetilde Q(T)u=0$.由于这个乘积中的每个因子
都是单射, 我们得出 $u=0$.矛盾.

(ii) 像 (i) 那样论证.

\solutionhead{6.23}

3. 在 $E=\R^2$ 中取 $T(u_1,u_2)=(u_2,0)$.则 $T^2=0$, 从而 $r(T)=0$, 而 $\norm{T}=1$.

5. 在 $E=\R^3$ 中取 $T(u_1,u_2,u_3)=(u_2,-u_1,0)$.则 $\sigma(T)=\{0\}$.
利用 $T^3=-T$ 的事实, 容易看出 $r(T)=1$.

\runin{评注.} 若在 $\CC$ 上的 Banach 空间中工作, 情况完全不同; 见 \S \ref{sec:11.4}.
在那里总有 $r(T)=\max\{\abs{\lambda};\ \lambda\in\sigma(T)\}$.在本例中取
$E=\CC^3$, 我们有 $\sigma(T)=\{0,+i,-i\}$, 从而 $r(T)=1$.

6. 假设该公式对 $T^n$ 成立, 我们有
\begin{align*}
(T^{n+1}u)(t)&=\frac{1}{(n-1)!}\int_0^tds\int_0^s(s-\tau)^{n-1}u(\tau)d\tau\\
&=\frac{1}{(n-1)!}\int_0^tu(\tau)\left[\int_\tau^t(s-\tau)^{n-1}ds\right]d\tau
=\frac{1}{n!}\int_0^t(t-\tau)^nu(\tau)d\tau.
\end{align*}

7. 考虑定义在 $\R$ 上的函数 $f$ 与 $g$:
\begin{equation*}
f(t)=\begin{cases}\frac{1}{(n-1)!}t^{n-1} & \text{若}\ 0\le t\le1,\\ 0 & \text{其他},\end{cases}
\qquad
g(t)=\begin{cases}u(t) & \text{若}\ 0\le t\le1,\\ 0 & \text{其他},\end{cases}
\end{equation*}
于是对 $0\le t\le1$, 有
$(f\star g)(t)=\int_0^1(t-\tau)u(\tau)d\tau=(T^nu)(t)$.我们推出
\begin{equation*}
\norm{f\star g}_{L^p(0,1)}\le\norm{f\star g}_{L^p(\R)}\le\norm{f}_{L^1(\R)}\norm{g}_{L^p(\R)}
=\frac{1}{n!}\norm{u}_{L^p(0,1)}.
\end{equation*}

8. 应用 Stirling 公式.

\solutionhead{6.24}

2. (v) $\Rightarrow$ (vi). 对每个 $\varepsilon>0$, $T_\varepsilon=T+\varepsilon I$
是双射且 $\sigma(T_\varepsilon)\subset[\varepsilon,1+\varepsilon]$.从而
$\sigma(T_\varepsilon^{-1})\subset\left[\frac{1}{1+\varepsilon},\frac1\varepsilon\right]$.
对 $T_\varepsilon^{-1}$ 应用命题  \ref{prop:6.9} 给出对一切 $v\in H$,
$(T_\varepsilon^{-1}v,v)\ge\frac{1}{1+\varepsilon}\abs{v}^2$, 即对一切 $u\in H$,
$(T_\varepsilon u,u)\ge\frac{1}{1+\varepsilon}\abs{T_\varepsilon u}^2$.

3. 令 $U=2T-I$.显然 (vii) 等价于对一切 $u\in H$, $\abs{u}\le\abs{Uu}$ (记作 (vii$'$)).
应用定理  \ref{thm:2.20}, 可见 (vii) $\Rightarrow$ $(-1,+1)\subset\rho(U)=2\rho(T)-1$.
从而 (vii) $\Rightarrow$ (viii).反之, (viii) $\Rightarrow$ $(-1,+1)\subset\rho(U)$.
因此 $\sigma(U)\subset(-\infty,-1]\cup[1,+\infty)$ 且 $\sigma(U^{-1})\subset[-1,+1]$.
由命题  \ref{prop:6.9} 我们知道 $\norm{U^{-1}}\le1$, 即 (vii$'$) 成立.

\solutionhead{6.25}

由构造, 在 $X$ 上有 $M\circ(I+K)=I$, 在 $R(I+K)$ 上有 $(I+K)\circ M=I$.
给定任何 $x\in E$, 写出 $x=x_1+x_2$, 其中 $x_1\in X,x_2\in N(I+K)$.则
$M\circ(I+K)(x)=x_1=x-Px$, 其中 $P$ 是到 $N(I+K)$ 上的投影算子.
对任何 $x\in E$, 有
\begin{equation*}
(I+K)\circ\widetilde M(x)=(I+K)\circ M\circ Q(x)=Qx=x-\widetilde Px,
\end{equation*}
其中 $\widetilde P$ 是到 $R(I+K)$ 在 $E$ 中的补空间上的有限秩投影算子.
\solutionhead{8.8}

4. 我们有
\begin{equation*}
u'_n=\zeta_nu'+\zeta'_nu.
\end{equation*}
显然由控制收敛定理, 在 $L^p$ 中 $\zeta_nu'\to u'$.还须证明在 $L^p$ 中
$\zeta'_nu\to0$.注意
\begin{equation*}
\norm{\zeta'_nu}_p^p\le C\int_{1/n}^{2/n}n^p\abs{u(x)}^pdx,
\end{equation*}
其中 $C=\norm{\zeta'}_{L^\infty}^p$.

当 $p=1$ 时, 由于 $u\in C([0,1])$ 且 $u(0)=0$, 有
\begin{equation*}
n\int_{1/n}^{2/n}\abs{u(x)}dx\le\max_{x\in[1/n,2/n]}\abs{u(x)}\to0\quad(n\to\infty).
\end{equation*}
当 $p>1$ 时, 我们有
\begin{equation*}
n^p\int_{1/n}^{2/n}\abs{u(x)}^pdx=n^p\int_{1/n}^{2/n}x^p\frac{\abs{u(x)}^p}{x^p}dx
\le2^p\int_{1/n}^{2/n}\frac{\abs{u(x)}^p}{x^p}dx\to0
\end{equation*}
(由第 1 问).

\solutionhead{8.9}

1. 由习题 8.8 第 1 问已知 $\dfrac{u'(x)}{x}\in L^p$.另一方面,
\begin{equation*}
u(x)=\int_0^xu'(t)dt=xu'(x)-\int_0^xu''(t)tdt,
\end{equation*}
从而
\begin{equation*}
\frac{u(x)}{x^2}=\frac{u'(x)}{x}-\frac{1}{x^2}\int_0^xu''(t)tdt.
\end{equation*}
但
\begin{equation*}
\frac{1}{x^2}\abs{\int_0^xu''(t)tdt}\le\frac1x\int_0^x\abs{u''(t)}dt\in L^p
\end{equation*}
(同上).

2. 我们有 $v\in C^1((0,1))$ 且
\begin{equation*}
v'(x)=-\frac{u(x)}{x^2}+\frac{u'(x)}{x}\in L^p
\end{equation*}
(由第 1 问).此外, 当 $x\to0$ 时
\begin{equation*}
v(x)=\frac{u(x)}{x}=\frac1x\int_0^xu'(t)dt\to0
\end{equation*}
(因为 $u\in C^1([0,1])$ 且 $u'(0)=0$).

3. 只须证明当 $n\to0$ 时 $\norm{\zeta'_nu'}_p+\norm{\zeta''_nu}_p\to0$.但
\begin{equation*}
\norm{\zeta'_nu'}_p^p\le Cn^p\int_{1/n}^{2/n}\abs{u'(x)}^pdx
\le2^pC\int_{1/n}^{2/n}\frac{\abs{u'(x)}^p}{x^p}dx,
\end{equation*}
以及
\begin{equation*}
\norm{\zeta''_nu}_p^p\le Cn^{2p}\int_{1/n}^{2/n}\abs{u(x)}^pdx
\le4^pC\int_{1/n}^{2/n}\frac{\abs{u(x)}^p}{x^{2p}}dx,
\end{equation*}
由 $\dfrac{u'(x)}{x}\in L^p$ 与 $\dfrac{u(x)}{x^2}\in L^p$ 即得出结论.

4. 设 $u\in X_m$.由归纳假设有 $u'\in X_{m-1}$ 且 $\dfrac{u'(x)}{x^{m-1}}\in L^p(I)$.
接着注意到
\begin{equation*}
\frac{u(x)}{x^m}=\frac{1}{x^m}\int_0^x\frac{u'(t)}{t^{m-1}}t^{m-1}dt.
\end{equation*}
再次应用 Hardy 不等式 (见问题 34 第 C 部分) 得到
\begin{equation*}
\frac{\abs{u(x)}}{x^m}\le\frac1x\int_0^x\frac{\abs{u'(t)}}{t^{m-1}}dt\in L^p(I).
\end{equation*}
为了证明 $\dfrac{u(x)}{x^{m-1}}\in X_1$, 注意
\begin{equation*}
D\left(\frac{u(x)}{x^{m-1}}\right)=\frac{Du(x)}{x^{m-1}}-(m-1)\frac{u(x)}{x^m}\in L^p(I),
\end{equation*}
且当 $x\to0$ 时
\begin{equation*}
\frac{\abs{u(x)}}{x^{m-1}}\le\frac{1}{x^{m-1}}\int_0^x\frac{\abs{u'(t)}}{t^{m-1}}t^{m-1}dt
\le\int_0^x\frac{\abs{u'(t)}}{t^{m-1}}dt\to0
\end{equation*}
(因为 $\dfrac{u'(t)}{t^{m-1}}\in L^p(I)$).

5. 只须对满足 $0\le\ell\le k-1$ 的每个整数 $\ell$ 验证 $D^\ell v\in X_1$.
但 $D^\ell v$ 是形如 $\dfrac{D^{j+\alpha}u(x)}{x^{m-j-k+\ell-\alpha}}$ (其中
$\alpha$ 是满足 $0\le\alpha\le\ell$ 的整数) 的函数的线性组合.然后使用第 4 问.

6. 只须证明当 $\alpha+\beta=m$ 且 $1\le\alpha\le m$ 时, 在 $L^p(I)$ 中
$(D^\alpha\zeta_n)(D^\beta u)\to0$.但 $\abs{D^\alpha\zeta_n(x)}\le Cn^\alpha$,
从而
\begin{align*}
\int_0^1\abs{D^\alpha\zeta_n(x)}^p\abs{D^\beta u(x)}^pdx
&\le Cn^{\alpha p}\int_{1/n}^{2/n}\abs{\frac{D^\beta u(x)}{x^\alpha}}^px^{\alpha p}dx\\
&\le C\int_{1/n}^{2/n}\abs{\frac{D^\beta u(x)}{x^\alpha}}^pdx\to0,
\end{align*}
因为由第 4 问 $\dfrac{D^\beta u(x)}{x^\alpha}\in L^p(I)$.

8. 为证明 $v\in C([0,1])$, 注意 $v(x)=\dfrac1x\int_0^xu'(t)dt$ 且 $u'\in C([0,1])$
满足 $u'(0)=0$.接着证明 $v\in W^{1,1}(I)$.分部积分可见
\begin{equation*}
v'(x)=\frac{1}{x^2}\int_0^xu''(t)tdt,
\end{equation*}
直接计算给出
\begin{equation*}
\norm{v'}_1\le\int_0^1\abs{u''(t)}(1-t)dt\le\norm{u''}_1.
\end{equation*}

9. 令 $u(x)=\int_0^x(1+\abs{\log t})^{-1}dt$.显然 $u\in W^{2,1}(I)$ 满足
$u(0)=u'(0)=0$, 此外 $\dfrac{u'(x)}{x}\notin L^1(I)$.关系式
$\dfrac{u(x)}{x^2}=\dfrac{u'(x)}{x}-v'(x)$ 结合第 8 问表明
$\dfrac{u(x)}{x^2}\notin L^1(I)$.

\solutionhead{8.10}

4. 显然当 $n\to\infty$ 时, $v'_n(x)=G'(nu(x))u'(x)\to f(x)$ 几乎处处, 其中
\begin{equation*}
f(x)=\begin{cases}
0 & \text{若}\ u(x)\ne0,\\
u'(x) & \text{若}\ u(x)=0.
\end{cases}
\end{equation*}

6. 对一切 $\phi\in C_c^1(I)$, 有 $\int_0^1v_n\phi'=-\int_0^1v'_n\phi$.令
$n\to\infty$ 取极限得出对一切 $\phi\in C_c^1(I)$, $\int_0^1f\phi=0$, 因此在
$I$ 上 $f=0$ 几乎处处, 即在 $[u=0]$ 上 $u'(x)=0$ 几乎处处.

\solutionhead{8.12}

1. 使用习题 8.2 以及 $\liminf_{n\to\infty}\norm{u'_n}_{L^p}\ge\norm{u'}_{L^p}$
的事实.

2. 考虑习题 8.2 中的序列 $(u_n)$.我们有 $\norm{u_n}_{L^1}\le1/2$ 且
$\norm{u'_n}_{L^1}=1$.因此 $\frac23u_n\in B_1$.另一方面, 在 $L^1$ 中
$\frac23u_n\to\frac23u$, 其中
\begin{equation*}
u(x)=\begin{cases}
0 & \text{若}\ x\in(0,1/2),\\
1 & \text{若}\ x\in(1/2,1).
\end{cases}
\end{equation*}
但 $u\notin W^{1,1}$.因此 $B_1$ 在 $L^1$ 中不是闭的.

\solutionhead{8.16}

2. $R(A)=L^p(0,1)$ 且 $N(A)=\{0\}$.

3. $v\in D(A^\star)$ 当且仅当 $v\in L^{p'}$ 且存在常数 $C$ 使得对一切
$u\in D(A)$, $\abs{\int_0^1vu'}\le C\norm{u}_p$.特别地,
$v\in D(A^\star)\Rightarrow v\in W^{1,p'}$, 此时
\begin{equation}
\abs{u(1)v(1)-\int_0^1uv'}\le C\norm{u}_p\quad\forall u\in D(A).
\tag{S1}
\label{eq:S8.16}
\end{equation}
由 \eqref{eq:S8.16} 推出对一切 $u\in D(A)$,
\begin{equation*}
\abs{u(1)}\abs{v(1)}\le(C+\norm{v'}_{p'})\norm{u}_p.
\end{equation*}
由于存在 $D(A)$ 中的序列 $(u_n)$ 使得 $u_n(1)=1$ 且 $\norm{u_n}_p\to0$,
由此推出 $v(1)=0$.于是我们证明了 $v\in D(A^\star)\Rightarrow v\in W^{1,p'}$
且 $v(1)=0$.容易推出
\begin{equation*}
D(A^\star)=\{v\in W^{1,p'};\ v(1)=0\},
\end{equation*}
且 $A^\star v=-v'$.

4. 我们有 $N(\widetilde A)=\{0\}$,
$R(\widetilde A)=\left\{f\in L^p;\ \int_0^1f(t)dt=0\right\}$, 以及
$(\widetilde A)^\star v=-v'$, 定义域为 $D((\widetilde A)^\star)=W^{1,p'}$.

\solutionhead{8.17}

在确定 $D(A^\star)$ 时, 记住下述事实是有用的: 设 $I=(0,1)$, $1<p\le\infty$.
假设 $u\in L^p(I)$ 满足
\begin{equation}
\abs{\int_Iu\phi'}\le C\norm{\phi}_{p'}\quad\forall\phi\in C_c^1(I)\ \text{使得}\ \int_I\phi=0,
\tag{S1}
\label{eq:S8.17}
\end{equation}
则 $u\in W^{1,p}(I)$.

事实上, 固定满足 $\int_I\psi_0=1$ 的函数 $\psi_0\in C_c^1(I)$.设 $\zeta$
为 $C_c^1(I)$ 中的任何函数.把 $\phi=\zeta-(\int_I\zeta)\psi_0$ 代入
\eqref{eq:S8.17}, 我们得到
\begin{equation*}
\abs{\int_Iu\zeta'}\le C\norm{\zeta}_{p'}+C'\abs{\int_I\zeta},
\end{equation*}
其中 $C'$ 只依赖于 $u$ 与 $\psi_0$.因此由命题  \ref{prop:8.3}, $u\in W^{1,p}(I)$.

当 $Au=u''-xu'$ 时, 我们有 $A^\star v=v''+xv'+v$.注意下述恒等式
\begin{equation*}
A^\star(e^{-x^2/2}u)=e^{-x^2/2}Au\quad\forall u\in H^2(I),
\end{equation*}
这允许在各种边界条件下计算 $N(A^\star)$.

\solutionhead{8.19}

给定 $f\in L^2(0,1)$, 令 $F(x)=\int_0^xf(t)dt$.则
\begin{equation*}
\varphi^\star(f)=\begin{cases}
\dfrac12\int_0^1F^2(x)dx & \text{若}\ \int_0^1f(t)dt=0,\\
+\infty & \text{其他}.
\end{cases}
\end{equation*}
事实上, 若 $\int_0^1f(t)dt=0$, 则对一切 $v\in H^1(0,1)$,
$\int_0^1fv=\int_0^1F'v=-\int_0^1Fv'$, 且
\begin{align*}
\varphi^\star(f)&=\sup_{v\in H^1}\left\{\int_0^1fv-\frac12\int_0^1v'^2\right\}
=\sup_{v\in H^1}\left\{-\int_0^1Fv'-\frac12\int_0^1v'^2\right\}\\
&=\sup_{w\in L^2}\left\{-\int_0^1Fw-\frac12\int_0^1w^2\right\}
=\frac12\int_0^1F^2.
\end{align*}

\solutionhead{8.21}

2. 设 $U$ 为满足 $-(pU')'+qU=f$ 在 $(0,1)$ 上且 $U(1)=0$ 的任何函数.则
\begin{equation*}
\int_0^1fv_0=p(0)(U'(0)-k_0U(0)).
\end{equation*}
因此若 $\int_0^1fv_0=0$, 任何这样的函数 $U$ 满足 $U'(0)=k_0U(0)$.由于
$U(0)$ 可以任意选取, 我们看到解集是一维的.

\solutionhead{8.22}

1. 函数 $\rho(x)=x$ 属于 $H^1(0,1)$, 但 $\sqrt{\rho(x)}=\sqrt{x}\notin H^1(0,1)$.

2. 对在 $(0,1)$ 上满足 $\rho\ge0$ 的每个 $\rho\in H^1(0,1)$, 令
$\gamma_\varepsilon=\sqrt{\rho+\varepsilon}$.由于函数 $t\mapsto\sqrt{t+\varepsilon}$
在 $[0,+\infty)$ 上属于 $C^1$ 类, 我们推出 $\gamma_\varepsilon\in H^1(0,1)$,
此外
\begin{equation*}
\gamma'_\varepsilon=\frac{\rho'}{2\sqrt{\rho+\varepsilon}},
\end{equation*}
从而在集合 $[\rho>0]$ 上 $\abs{\gamma'_\varepsilon}\le\mu$.另一方面, 我们
知道在集合 $[\rho=0]$ 上 $\rho'=0$ 几乎处处 (见习题 8.10), 从而在 $[\rho=0]$
上 $\abs{\gamma'_\varepsilon}\le\mu$ 几乎处处.因此在 $(0,1)$ 上几乎处处
$\abs{\gamma'_\varepsilon}\le\mu$.

由此推出若 $\mu\in L^2$, 则当 $\varepsilon\to0$ 时
$\norm{\gamma'_\varepsilon}_{L^2}\le C$.由于当 $\varepsilon\to0$ 时在
$C([0,1])$ 中 $\gamma_\varepsilon\to\sqrt{\rho}$, 且在 $L^2(0,1)$ 中
$\gamma'_\varepsilon\to\mu$, 我们得出 (见习题 8.2) $\sqrt{\rho}\in H^1(0,1)$
且 $(\sqrt{\rho})'=\mu$.

反之, 若 $\sqrt{\rho}\in H^1(0,1)$, 令 $\gamma=\sqrt{\rho}$, 从而 $\rho=\gamma^2$
且 $\rho'=2\gamma\gamma'$.因此在 $[\rho>0]$ 上 $\mu=\gamma'$ 几乎处处,
此外在 $[\rho=0]$ 上 $\mu=\gamma'$ 几乎处处 (因为在 $[\gamma=0]=[\rho=0]$ 上
$\gamma'=0$ 几乎处处).

\solutionhead{8.24}

1. 可以选取 $C_\varepsilon=1+1/\varepsilon$.

2. 弱表述为: $u\in H^1(I)$,
\begin{equation*}
a(u,v)=\int_0^1(u'v'+kuv)-u(1)v(1)=\int_0^1fv\quad\forall v\in H^1(I).
\end{equation*}
显然 $a(u,v)$ 是 $H^1(0,1)$ 上的连续双线性形式.由第 1 问, 例如当 $k>2$ 时
它是强制的.相应的极小化问题为
\begin{equation*}
\min_{v\in H^1}\left\{\frac12\int_0^1(v'^2+kv^2)-\frac12v(1)^2-\int_0^1fv\right\}.
\end{equation*}

3. 设 $g\in L^2(I)$, 设 $v\in H^2(I)$ 为 (1) 的相应解 (其中 $f$ 换成 $g$).
我们有
\begin{align*}
(Tf,g)_{L^2}&=\int_0^1ug=\int_0^1u(-v''+kv)\\
&=-u(1)v'(1)+u(0)v'(0)+u'(1)v(1)-u'(0)v(0)+\int_0^1(-u''+ku)v\\
&=-u(1)v(1)+u(1)v(1)+\int_0^1fv=(f,Tg)_{L^2}.
\end{align*}
因此 $T$ 自伴.它紧是因为它是由 $L^2(I)$ 到 $H^1(I)$ 中的有界算子, 且
$H^1(I)\subset L^2(I)$ 紧嵌入 (见定理  \ref{thm:8.8}).

4. 由 \S \ref{sec:6.4} 的结果我们知道存在 $L^2(I)$ 中的序列 $(u_n)$ 满足
$Tu_n=\mu_nu_n$, 其中 $\norm{u_n}_{L^2}=1$, 对一切 $n$ 有 $\mu_n>0$, 且
$\mu_n\to0$.因此我们有 $-u''_n+ku_n=\dfrac{1}{\mu_n}u_n$, 从而在 $I$ 上
$-u''_n=\left(\dfrac{1}{\mu_n}-k\right)u_n$.

5. 排除了 $\lambda=0$ 的值 (为什么?).若 $\lambda>0$, 有
$u(x)=A\cos\sqrt{\lambda}x+B\sin\sqrt{\lambda}x$, 其中常数 $A$ 与 $B$ 经
调整以满足边界条件, 即 $B=0$ 且 $A(\cos\sqrt{\lambda}+\sqrt{\lambda}\sin\sqrt{\lambda})=0$,
从而 $A\ne0$ 当且仅当 $\sqrt{\lambda}$ 是方程 $\tan t=-1/t$ 的解 (通过考察
图形可见它有无限正解序列 $t_n\to\infty$).若 $\lambda<0$, 有
$u(x)=Ae^{\sqrt{\abs{\lambda}}x}+Be^{-\sqrt{\abs{\lambda}}x}$.把它与边界
条件结合给出 $A=B$ 且
$A\sqrt{\abs{\lambda}}e^{\sqrt{\abs{\lambda}}}-B\sqrt{\abs{\lambda}}e^{-\sqrt{\abs{\lambda}}}=Ae^{\sqrt{\abs{\lambda}}}+Be^{-\sqrt{\abs{\lambda}}}$.
为了有某个 $u\not\equiv0$, $\lambda$ 必须满足
$\sqrt{\abs{\lambda}}(e^{\sqrt{\abs{\lambda}}}-e^{-\sqrt{\abs{\lambda}}})=e^{\sqrt{\abs{\lambda}}}+e^{-\sqrt{\abs{\lambda}}}$,
即 $t=\sqrt{\abs{\lambda}}$ 是方程 $e^{2t}=\dfrac{t+1}{t-1}$ 的解.考察图形表明
存在唯一解 $t_0>1$, 此时 $\lambda=-t_0^2$.

\solutionhead{8.25}

2. 反设存在 $H^1(I)$ 中的序列 $(u_n)$ 使得 $a(u_n,u_n)\to0$ 且
$\norm{u_n}_{H^1(I)}=1$.抽取子序列 $(u_{n_k})$ 可以假设在 $L^2$ 中
$u'_{n_k}\rightharpoonup u'$ 弱收敛且 $u_n\to u$ 强收敛.由下半连续性 (见命题
 \ref{prop:3.5}) 有 $\liminf\int_I(u'_{n_k})^2\ge\int_I(u')^2$, 从而
$a(u,u)=0$, 因此 $u=0$.但 $\int_I(u'_{n_k})^2=1-\int_I u_{n_k}^2$, 从而
\begin{equation*}
a(u_{n_k},u_{n_k})=\int_I(u'_{n_k})^2+\left(\int_0^1u_{n_k}\right)^2
=1-\int_I u_{n_k}^2+\left(\int_0^1u_{n_k}\right)^2\to1.
\end{equation*}
矛盾.

4. 对一切 $v\in H^1(I)$, 有 $\int_Iu'v'=\int_Igv$, 其中
$g=f-(\int_0^1u)\chi_{(0,1)}$.因此 $u\in H^2(I)$ 且满足
\begin{equation*}
\begin{cases}
-u''+(\int_0^1u)\chi_{(0,1)}=f & \text{在}\ I\ \text{上},\\
u'(0)=u'(2)=0.
\end{cases}
\end{equation*}

5. $u\in C^2(I)$ 当且仅当 $\int_0^1u=0$.这发生当且仅当 $\int_If=0$.

8. $T$ 的特征值为正, 且若 $1/\lambda$ 是特征值, 必有非平凡函数 $u\not\equiv0$
满足
\begin{equation*}
\begin{cases}
-u''+\int_0^1u=\lambda u & \text{在}\ (0,1)\ \text{上},\\
-u''=\lambda u & \text{在}\ (1,2)\ \text{上},\\
u'(0)=u'(2)=0,\\
u(1-)=u(1+),\quad u'(1-)=u'(1+).
\end{cases}
\end{equation*}
因此在 $(0,1)$ 上 $u(x)=\dfrac k\lambda+A\cos(\sqrt{\lambda}x)$, 在 $(1,2)$
上 $u(x)=A'\cos(\sqrt{\lambda}(x-2))$, 其中常数 $k,A,A'$ 由关系式
$u(1-)=u(1+)$, $u'(1-)=u'(1+)$ 以及 $k=\int_0^1u$ 确定.我们得出或者
$\sin\sqrt{\lambda}=0$, 即 $\lambda=n^2\pi^2$ ($n=1,2,\dots$); 或者 $\lambda$
是方程 $\tan\sqrt{\lambda}=2\sqrt{\lambda}(1-\lambda)$ 的解.

\solutionhead{8.26}

3. 令 $a(v,v)=\int_I pv'^2+qv^2$.我们有 $(S_Nf-S_Df,f)=\int_If(u_N-u_D)$.
由极小化性质已知 $\frac12a(u_N,u_N)-\int_Ifu_N\le\frac12a(u_D,u_D)-\int_Ifu_D$.
另一方面, $a(u_N,u_N)=\int_Ifu_N$ 且 $a(u_D,u_D)=\int_Ifu_D$.因此
$\int_If(u_N-u_D)\ge0$.

6. 令 $a_i(v,v)=a(v,v)+k_iv^2(0)$ ($i=1,2$), 以及 $u_{k_1}=u_1,u_{k_2}=u_2$.
由于 $u_i$ 是 $\frac12a_i(v,v)-\int_Ifv$ 在 $V=\{v\in H^1(I);\ v(1)=0\}$
上的极小元, 我们有
\begin{equation*}
\frac12a(u_2,u_2)+\frac12k_2u_2^2(0)-\int_Ifu_2
\le\frac12a(u_1,u_1)+\frac12k_2u_1^2(0)-\int_Ifu_1.
\end{equation*}
另一方面, $a(u_1,u_1)+k_1u_1^2(0)=\int_Ifu_1$, 以及
$a(u_2,u_2)+k_2u_2^2(0)=\int_Ifu_2$.因此
\begin{equation*}
-\frac12\int_Ifu_2\le\frac12\int_Ifu_1+\frac12(k_2-k_1)u_1^2(0)-\int_Ifu_1,
\end{equation*}
从而
\begin{equation*}
(S_{k_2}f-S_{k_1}f,f)=\int_If(u_2-u_1)\ge(k_1-k_2)u_1^2(0)\ge0.
\end{equation*}

\solutionhead{8.27}

4. 问题 $-\phi''+\phi=1$ 在 $I$ 上, $\phi(-1)=\phi(1)=0$ 的解由
$\phi(x)=1+A(e^x+e^{-x})$ 给出, 其中 $A=-e/(e^2+1)$.由 $\phi$ 的唯一性必有
$u=\lambda u(0)\phi$.因此
\begin{equation*}
\lambda_0=\frac{1}{\phi(0)}=\frac{e^2+1}{(e-1)^2}.
\end{equation*}

5. 方程 (1) 变为 $u=S(f+\lambda u(0))=Sf+\lambda u(0)S1=Sf+\lambda u(0)\phi$.
从而 $u(0)(1-\lambda\phi(0))=(Sf)(0)$, 即
\begin{equation*}
u(0)=\frac{\lambda_0(Sf)(0)}{\lambda_0-\lambda},
\end{equation*}
且 $u=Sf+\dfrac{\lambda\lambda_0(Sf)(0)\phi}{\lambda_0-\lambda}$ 是所求的解.

6. 当 $\lambda=\lambda_0$ 时, 解 $u$ 的存在性蕴含 $(Sf)(0)=0$.反之, 设
$(Sf)(0)=0$.则 (1) 的解必须具有形式 $u=Sf+A\phi$ (对某个常数 $A$).直接计算
表明任何这样的 $u$ 满足 $-u''+u=f+A$.但 $u(0)=(Sf)(0)+\frac{A}{\lambda_0}=\frac{A}{\lambda_0}$.
因此我们有 $-u''+u=f+\lambda_0u(0)$, 即对任何 $A$, (1) 都成立.因此当
$\lambda=\lambda_0$ 时 (1) 的所有解的集合为 $Sf+\R\phi$.

\solutionhead{8.29}

2. 解 $u\in H^1(0,1)$ 的存在性与唯一性来自 Lax--Milgram 定理.特别地, $u$
满足对一切 $v\in H_0^1(0,1)$, $\int_0^1u'v'=\int_0^1(f-u)v$, 因此
$u'\in H^1(0,1)$, 即 $u\in H^2(0,1)$; 此外在 $(0,1)$ 上 $-u''+u=f$.利用
$u\in H^2(0,1)$ 的信息, 我们现在可以写出
\begin{align*}
a(u,v)&=\int_0^1(-u''+u)v+u'(1)v(1)-u'(0)v(0)+(u(1)-u(0))(v(1)-v(0))\\
&=\int_0^1fv\quad\forall v\in H^1(0,1).
\end{align*}
从而对一切 $v\in H^1(0,1)$,
\begin{equation*}
(u'(1)+u(1)-u(0))v(1)-(u'(0)+u(1)-u(0))v(0)=0.
\end{equation*}
由于 $v(0)$ 与 $v(1)$ 是任意的, 我们得出
\begin{equation*}
u'(1)+u(1)-u(0)=0\quad\text{且}\quad u'(0)+u(1)-u(0)=0.
\end{equation*}

5. 使用定理  \ref{thm:8.19} 的证明中的相同函数 $G$, 取 $v=G(-u)$, 由于
$f\ge0$ 与 $G\ge0$, 有 $a(u,G(-u))=\int_0^1fG(-u)\ge0$.另一方面,
\begin{align*}
a(u,G(-u))&=-\int_0^1u'^2G'(-u)-\int_0^1(-u)G(-u)+(u(1)-u(0))(G(-u(1))-G(-u(0)))\\
&\le-\int_0^1(-u)G(-u)
\end{align*}
(因为 $G$ 单调递增).由此推出 $\int_0^1(-u)G(-u)\le0$, 从而 $-u\le0$.

7. 设 $1/\lambda$ 为特征值, $u$ 为相应的特征函数.则 $-u''+u=\lambda u$ 在
$(0,1)$ 上, $u'(0)=u(0)-u(1)$, $u'(1)=u(0)-u(1)$.由于
$a(u,u)=\lambda\int_0^1u^2\ge\int_0^1u^2$, 可见 $\lambda\ge1$.此外 $\lambda=1$
是对应于 $u=\text{常数}$ 的特征值.现设 $\lambda>1$ 并令 $\alpha=\sqrt{\lambda-1}$.
必有 $u(x)=A\cos\alpha x+B\sin\alpha x$.为了满足边界条件, 需要满足方程组
\begin{align*}
B\alpha&=A-A\cos\alpha-B\sin\alpha,\\
-A\alpha\sin\alpha+B\alpha\cos\alpha&=A-A\cos\alpha-B\sin\alpha.
\end{align*}
这个方程组有非平凡解当且仅当 $2(1-\cos\alpha)+\alpha\sin\alpha=0$, 即
$\sin(\alpha/2)=0$ 或 $(\alpha/2)+\tan(\alpha/2)=0$.

\solutionhead{8.34}

1. 设 $u$ 为古典解.则对一切 $v\in H^1(0,1)$,
\begin{equation*}
-u'(1)v(1)+u'(0)v(0)+\int_0^1(u'v'+uv)=\int_0^1fv.
\end{equation*}
令 $V=\{v\in H^1(0,1);\ v(0)=v(1)\}$.若 $v\in V$, 我们得到
\begin{equation*}
a(u,v)=\int_0^1(u'v'+uv)=\int_0^1fv+kv(0).
\end{equation*}
弱表述为: $u\in V$ 且对一切 $v\in V$, $a(u,v)=\int_0^1fv+kv(0)$.

2. 由 Lax--Milgram 定理存在唯一的弱解 $u\in V$, 相应的极小化问题为
\begin{equation*}
\min_{v\in V}\left\{\frac12\int_0^1(v'^2+v^2)-\int_0^1fv-kv(0)\right\}.
\end{equation*}

3. 显然任何弱解 $u$ 属于 $H^2(0,1)$ 且满足在 $(0,1)$ 上几乎处处 $-u''+u=f$,
对一切 $v\in V$ 有 $u'(1)v(1)-u'(0)v(0)=kv(0)$, 即 $u'(1)-u'(0)=k$.

5. $T$ 的特征值由 $\lambda_k=1/\mu_k$ 给出, 其中 $\mu_k$ 对应于在 $(0,1)$
上几乎处处
\begin{equation*}
-u''+u=\mu_ku,\qquad u(1)=u(0),\ u'(1)=u'(0)
\end{equation*}
的非平凡解.因此 $\mu_k\ge1$ 且 $u$ 由
\begin{equation*}
u(x)=A\sin(\sqrt{\mu_k-1}x)+B\cos(\sqrt{\mu_k-1}x)
\end{equation*}
给出, 其中 $\sqrt{\mu_k-1}=2\pi k$, $k=0,1,\dots$.

\solutionhead{8.38}

2. 设 $Tu=\lambda u$, 其中 $u\in H^2(\R)$ 且 $u\not\equiv0$.显然 $\lambda\ne0$
且 $u$ 在 $\R$ 上满足 $-u''+u=\dfrac1\lambda u$.若 $\lambda=1$, 有
$u(x)=Ax+B$.由于 $u\in L^2(\R)$, 推出 $A=B=0$.因此 $1\notin\operatorname{EV}(T)$.
若 $(1/\lambda-1)>0$, 有 $u(x)=A\sin\alpha x+B\cos\alpha x$, 其中
$\alpha=\sqrt{1/\lambda-1}$.条件 $u\in L^2(\R)$ 再次给出 $A=B=0$.类似地若
$(1/\lambda-1)<0$, 除 $u\equiv0$ 外没有解.因此 $\operatorname{EV}(T)=\emptyset$.
$T$ 不能是紧算子.否则由定理  \ref{thm:6.8} 我们会有 $\sigma(T)=\{0\}$, 进而由
推论  \ref{cor:6.10} $T\equiv0$.但显然 $T\not\equiv0$.

3. 若 $\lambda<0$, 例如由 Lax--Milgram 定理以及对一切 $f\in H$ 有 $(Tf,f)\ge0$
的事实, $(T-\lambda I)$ 是由 $H=L^2(\R)$ 到自身的双射.因此 $\lambda\in\rho(T)$.

4. 若 $\lambda>1\ge\norm{T}$, 由命题  \ref{prop:6.7} 我们有 $\lambda\in\rho(T)$.

6. $T$ 不是满射, 因为 $R(T)\subset H^2(\R)$.

7. $T-I$ 不是满射.事实上, 若试图对给定的 $\phi\in L^2(\R)$ 求解 $Tf-f=\phi$,
我们被引向 $-u''+u=f$ (令 $u=Tf$) 与 $u=f+\phi$.因此 $u''=\phi$ 允许解
$u\in H^2$.例如设 $\operatorname{supp}\phi\subset[0,1]$.直接计算给出对一切
$x\le0$ 有 $u(x)=0$, 对一切 $x\ge1$ 有 $u(x)=0$.从而 $u'(0)=u'(1)=0$.由此
推出 $0=u'(1)-u'(0)=\int_0^1\phi$.因此当 $\int_0^1\phi\ne0$ 时方程 $Tf-f=\phi$
没有解 $f\in L^2(\R)$.从而 $T-I$ 不能是满射.

8. $T-\lambda I$ 不是满射.事实上, 若试图求解 $Tf-\lambda f=\phi$, 我们被
引向 $-u''+u=f$ (令 $u=Tf$) 与 $u=\lambda f+\phi$.因此 $-u''+u=\dfrac1\lambda(u-\phi)$.
再次假设 $\operatorname{supp}\phi\subset[0,1]$.在 $[0,1]$ 外会有 $u''=-\mu^2u$,
其中 $\mu=\sqrt{1/\lambda-1}$.因此在 $[0,1]$ 外 $u\equiv0$, 从而
$u(0)=u'(0)=u(1)=u'(1)=0$.方程 $-u''+(1-1/\lambda)u=-\dfrac1\lambda\phi$ 蕴含
对在 $(0,1)$ 上的 $-v''=\mu^2v$ 的任何解 $v$, $\int_0^1\phi v=0$; 例如
$\int_0^1\phi(x)\sin\mu x=0$.因此当 $\int_0^1\phi(x)\sin\mu x\ne0$ 时方程
$Tf-\lambda f=\phi$ 没有解 $f\in L^2(\R)$.从而 $(T-\lambda I)$ 不是满射.

\solutionhead{8.39}

2. 对一切 $v\in\R$, 有 $v^2\le\frac12v^4+\frac12$, 从而
\begin{equation*}
\varphi(v)\ge\frac12\norm{v}_{H^1}^2-\frac14-\norm{f}_{L^2}\norm{v}_{H^1}.
\end{equation*}
因此当 $\norm{v}_{H^1}\to\infty$ 时 $\varphi(v)\to\infty$.

3. 唯一性来自 $\varphi$ 在 $H^1(0,1)$ 上严格凸的事实; 这是函数 $t\mapsto t^4$
在 $\R$ 上严格凸的推论.

4. 我们有
\begin{align*}
\varphi(u+\varepsilon v)&=\frac12\int_0^1(u'^2+2\varepsilon u'v'+\varepsilon^2v'^2)\\
&\quad+\frac14\int_0^1(u^4+4\varepsilon u^3v+6\varepsilon^2u^2v^2+4\varepsilon^3uv^3+\varepsilon^4v^4)-\int_0^1f(u+\varepsilon v).
\end{align*}
写出 $\varphi(u)\le\varphi(u+\varepsilon v)$ 给出
$\int_0^1(u'v'+u^3v-fv)+A_\varepsilon\ge0$, 其中当 $\varepsilon\to0$ 时
$A_\varepsilon\to0$.令 $\varepsilon\to0$ 取极限并选取 $\pm v$ 给出对一切
$v\in H^1(0,1)$,
\begin{equation*}
\int_0^1(u'v'+u^3v-fv)=0.
\end{equation*}

6. 由函数 $t\mapsto t^4$ 的凸性, 对一切 $u,v\in\R$ 有
$\frac14v^4-\frac14u^4\ge u^3(v-u)$.另一方面, 显然对一切 $u,v\in H^1(0,1)$,
在 $(0,1)$ 上几乎处处有 $\frac12v'^2-\frac12u'^2\ge u'(v'-u')$.从而对一切
$u,v\in H^1(0,1)$,
\begin{equation*}
\varphi(v)-\varphi(u)\ge\int_0^1u'(v'-u')+\int_0^1u^3(v-u)-\int_0^1f(v-u).
\end{equation*}
若 $u$ 是 (3) 的解, 则对一切 $v\in H^1(0,1)$ 有
$\int_0^1u'(v'-u')+\int_0^1u^3(v-u)=\int_0^1f(v-u)$, 从而对一切
$v\in H^1(0,1)$, $\varphi(u)\le\varphi(v)$.

9. 我们断言当 $\norm{v}_{H^1}\to\infty$ 时 $\psi(v)\to+\infty$.事实上, 这
归结为证明对每个常数 $C$, 集合 $\{v\in H^1(0,1);\ \psi(v)\le C\}$ 在
$H^1(0,1)$ 中有界.若 $\psi(v)\le C$, 写出
\begin{equation*}
\int_0^1fv=\int_0^1f(v-v(0))+v(0)\le\norm{f}_{L^2}(\norm{v'}_{L^2}+\abs{v(0)}),
\end{equation*}
从而 $\norm{v'}_{L^2}$ 与 $\abs{v(0)}$ 有界 (为什么?).因此
$\norm{v}_{L^2}\le\norm{v'}_{L^2}+\abs{v(0)}$ 也有界, 从而 $\norm{v}_{H^1}$
有界.极小元的唯一性: 验证
$\psi\left(\dfrac{u_1+u_2}{2}\right)\le\frac12(\psi(u_1)+\psi(u_2))$, 且
等号成立当且仅当 $u'_1=u'_2$ 且 $u_1(0)=u_2(0)$, 即 $u_1=u_2$.

我们有
\begin{align*}
\psi(u+\varepsilon v)&=\frac12\int_0^1(u'^2+2\varepsilon u'v'+\varepsilon^2v'^2)\\
&\quad+\frac14(u^4(0)+4\varepsilon u^3(0)v(0)+\dots+\varepsilon^4v^4(0))-\int_0^1f(u+\varepsilon v).
\end{align*}
若 $u$ 是 $\psi$ 的极小元, 写出 $\psi(u)\le\psi(u+\varepsilon v)$, 得到
$\int_0^1(u'v'-fv)+u^3(0)v(0)+B_\varepsilon\ge0$, 其中当 $\varepsilon\to0$ 时
$B_\varepsilon\to0$.令 $\varepsilon\to0$ 取极限并选取 $\pm v$ 给出
\begin{equation}
\int_0^1(u'v'-fv)+u^3(0)v(0)=0\quad\forall v\in H^1(0,1).
\tag{S1}
\label{eq:S8.39a}
\end{equation}
由此推出 $u\in H^2(0,1)$ 满足在 $(0,1)$ 上几乎处处
\begin{equation}
-u''=f.
\tag{S2}
\label{eq:S8.39b}
\end{equation}
回到 \eqref{eq:S8.39a} 并利用 \eqref{eq:S8.39b} 给出对一切 $v\in H^1(0,1)$,
$u'(1)v(1)-u'(0)v(0)+u^3(0)v(0)=0$, 从而
\begin{equation}
u'(1)=0,\qquad u'(0)=u^3(0).
\tag{S3}
\label{eq:S8.39c}
\end{equation}
反之, 满足 \eqref{eq:S8.39b} 与 \eqref{eq:S8.39c} 的任何函数 $u$ 是 $\psi$
的极小元: 论证与第 6 问相同.在这种情形下我们有显式解.\eqref{eq:S8.39b} 的
通解由 $u(x)=-\int_0^x(x-t)f(t)dt+Ax+B$ 给出, 此时 \eqref{eq:S8.39c} 等价于
$A=\int_0^1f(t)dt$ 且 $A=B^3$.

\solutionhead{8.42}

2. 对方程
\begin{equation}
v(x)=p^{1/4}(t)u(t)
\tag{S1}
\label{eq:S8.42a}
\end{equation}
关于 $t$ 求导给出
\begin{equation*}
v'(x)p^{-1/2}(t)=\frac14p^{-3/4}(t)p'(t)u(t)+p^{1/4}(t)u'(t).
\end{equation*}
从而
\begin{equation}
p(t)u'(t)=v'(x)p^{1/4}(t)-\frac14p'(t)u(t)=v'(x)p^{1/4}(t)-\frac14p'(t)p^{-1/4}(t)v(x).
\tag{S2}
\label{eq:S8.42b}
\end{equation}
对 \eqref{eq:S8.42b} 关于 $t$ 求导给出
\begin{equation}
(pu')'=v''(x)p^{-1/4}(t)-\frac14p''(t)p^{-1/4}(t)v(x)+\frac{1}{16}p'(t)^2p^{-5/4}(t)v(x).
\tag{S3}
\label{eq:S8.42c}
\end{equation}
把 \eqref{eq:S8.42c} 与 $(0,1)$ 上的方程 $-(pu')'+qu=\mu u$ 结合得出
\begin{equation}
v''(x)p^{-1/4}(t)-\frac14p''(t)p^{-1/4}(t)v(x)+\frac{1}{16}p'(t)^2p^{-5/4}(t)v(x)=(q(t)-\mu)p^{-1/4}(t)v(x).
\tag{S4}
\label{eq:S8.42d}
\end{equation}
因此 $v$ 满足在 $(0,L)$ 上 $-v''+a(x)v=\mu v$, 其中
\begin{equation*}
a(x)=q(t)+\frac14p''(t)-\frac{1}{16}p'(t)^2p^{-1}(t).
\end{equation*}

% chapters/problems.tex
% 问题 (Problems)
% 源: 原书 PDF 第 439--524 页 (印刷页 435--520)

\chapter*{问题}
\addcontentsline{toc}{chapter}{问题}
\markboth{问题}{问题}

括号中的数字表示求解该问题所需掌握的原书各章知识.

\begin{problem}[第~
\ref{chap:1}~章, 仅第 9 问需第~
\ref{chap:4}~章]
\index{KreinMilmandingli@Krein--Milman 定理 (Krein--Milman theorem)}
\index{jizhidian@极值点 (extreme point)}
\label{prob:1}
\runin{极值点; Krein--Milman 定理}
设 $E$ 为赋范向量空间, $K\subset E$ 为凸子集.若对满足 $x\ne y$ 的一切 $x,y\in K$ 及
一切 $t\in(0,1)$,
\begin{equation*}
tx+(1-t)y\ne a,
\end{equation*}
则称点 $a\in K$ 为\emph{极值点}.

1. 验证: $a\in K$ 是极值点当且仅当集合 $K\setminus\{a\}$ 是凸集.

2. 设 $a$ 为 $K$ 的极值点.设 $(x_i)_{1\le i\le n}$ 为 $K$ 中的有限列,
$(\alpha_i)_{1\le i\le n}$ 为满足对一切 $i$ 有 $\alpha_i>0$, $\sum\alpha_i=1$ 且
$\sum\alpha_ix_i=a$ 的实数有限列.证明对一切 $i$, $x_i=a$.

在下文中, 我们假设 $K\subset E$ 是 $E$ 的非空紧凸子集.若子集 $M\subset K$
非空, 闭, 且只要对某个 $t\in(0,1)$ 有 $tx+(1-t)y\in M$ ($x,y\in K$), 就有
$x\in M$ 且 $y\in M$, 则称 $M$ 为\emph{极端集}.

3. 设 $a\in K$.验证: $a$ 是极值点当且仅当 $\{a\}$ 是极端集.

我们的第一个目标是证明每个极端集包含至少一个极值点.

4. 设 $A\subset K$ 为极端集, $f\in E^\star$.令
\begin{equation*}
B=\left\{x\in A;\ \dual{f}{x}=\max_{y\in A}\dual{f}{y}\right\}.
\end{equation*}
证明 $B$ 是 $K$ 的极端子集.

5. 设 $M\subset K$ 为 $K$ 的极端集.考虑包含在 $M$ 内的 $K$ 的所有极端集的族
$\mathcal{F}$; $\mathcal{F}$ 配以下述偏序:
\begin{equation*}
A\le B\iff B\subset A.
\end{equation*}
证明 $\mathcal{F}$ 有极大元 $M_0$.

6. 证明 $M_0$ 退化为单点.\\
{[提示: 使用 Hahn--Banach 定理与第 4 问.]}

7. 得出结论.

8. 证明 $K$ 与其所有极值点的闭凸包重合.\\
{[提示: 用反证法并使用 Hahn--Banach 定理.]}

9. 确定在下述情形下 $B_E$ ($=E$ 的闭单位球) 的所有极值点的集合 $\mathcal{E}$:
\begin{enumerate}[label=(\alph*), itemsep=0.1em]
	\item $E=\ell^\infty$,
	\item $E=c$,
	\item $E=c_0$,
	\item $E=\ell^1$,
	\item $E=\ell^p$, 其中 $1<p<\infty$,
	\item $E=L^1(\R)$.
\end{enumerate}
[关于记号参见 \S \ref{sec:11.3}].
\end{problem}

\begin{problem}[第~
\ref{chap:1}~章, 仅第 B4 问需第~
\ref{chap:2}~章]
\index{ciweifen@次微分 (subdifferential)}
\label{prob:2}
\runin{凸函数的次微分}
设 $E$ 为赋范向量空间, $\varphi:E\to(-\infty,+\infty]$ 为满足
$\varphi\not\equiv+\infty$ 的凸函数.对每个 $x\in E$, $\varphi$ 的次微分定义为
\begin{equation*}
\begin{cases}
\partial\varphi(x)=\{f\in E^\star;\ \varphi(y)-\varphi(x)\ge\dual{f}{y-x}\ \forall y\in E\} & \text{若}\ x\in D(\varphi),\\
\partial\varphi(x)=\emptyset & \text{若}\ x\notin D(\varphi),
\end{cases}
\end{equation*}
我们令 $D(\partial\varphi)=\{x\in E;\ \partial\varphi(x)\ne\emptyset\}$, 从而
$D(\partial\varphi)\subset D(\varphi)$.构造一个使此包含关系为严格包含的例子.

\runin{-- A --}

1. 证明 $\partial\varphi(x)$ 是 $E^\star$ 的闭凸子集.

2. 设 $x_1,x_2\in D(\partial\varphi)$, $f_1\in\partial\varphi(x_1)$,
$f_2\in\partial\varphi(x_2)$.证明
\begin{equation*}
\dual{f_1-f_2}{x_1-x_2}\ge0.
\end{equation*}

3. 证明
\begin{equation*}
f\in\partial\varphi(x)\iff\varphi(x)+\varphi^\star(f)=\dual{f}{x}.
\end{equation*}

4. 在下述情形下确定 $\partial\varphi$:
\begin{enumerate}[label=(\alph*), itemsep=0.1em]
	\item $\varphi(x)=\frac12\norm{x}^2$,
	\item $\varphi(x)=\norm{x}$,
	\item $\varphi(x)=I_K(x)$ ($K$ 的示性函数), 其中 $K\subset E$ 为非空凸集 (相应地线性子空间),
	\item $\varphi(x)$ 为 $E$ 上的可微凸函数.
\end{enumerate}
[提示: 在情形 (a), (b) 中, $\partial\varphi$ 与第~
\ref{chap:1}~章注记 2 中定义的对偶映射 $F$
有关; 另见习题 1.1.]

5. 设 $\psi:E\to(-\infty,+\infty]$ 为另一个满足 $\psi\not\equiv+\infty$ 的凸函数.
假设 $D(\varphi)\cap D(\psi)\ne\emptyset$.证明对一切 $x\in E$,
\begin{equation*}
\partial\varphi(x)+\partial\psi(x)\subset\partial(\varphi+\psi)(x)
\end{equation*}
(约定若 $A=\emptyset$ 或 $B=\emptyset$, 则 $A+B=\emptyset$).构造一个使此包含
关系为严格包含的例子.

\runin{-- B --}

在整个第 B 部分中, 我们假设 $x_0\in E$ 满足条件
\begin{equation}
\exists M\in\R\ \text{与}\ \exists R>0\ \text{使得}\ \varphi(x)\le M
\quad\forall x\in E\ \text{满足}\ \norm{x-x_0}\le R.
\label{eq:prob2.1}
\end{equation}

1. 证明 $\partial\varphi(x_0)\ne\emptyset$.\\
{[提示: 在 $E\times\R$ 中使用 Hahn--Banach 定理.]}

2. 证明对一切 $f\in\partial\varphi(x_0)$,
$\norm{f}\le\frac1R(M-\varphi(x_0))$.

3. 由此推出: 对一切 $r<R$, 存在 $L\ge0$ 使得对一切满足 $\norm{x_i-x_0}\le r$
($i=1,2$) 的 $x_1,x_2\in E$,
\begin{equation*}
\abs{\varphi(x_1)-\varphi(x_2)}\le L\norm{x_1-x_2}.
\end{equation*}
[另一种证明另见习题 2.1.]

4. 这里假设 $E$ 为 Banach 空间且 $\varphi$ 下半连续.证明
\begin{equation*}
\operatorname{Int}D(\partial\varphi)=\operatorname{Int}D(\varphi).
\end{equation*}

5. 证明对每个 $y\in E$,
\begin{equation*}
\lim_{t\downarrow0}\frac{\varphi(x_0+ty)-\varphi(x_0)}{t}
=\max_{f\in\partial\varphi(x_0)}\dual{f}{y}.
\end{equation*}
[提示: 参看习题 1.25 第 5 问.]

6. 设 $\psi:E\to(-\infty,+\infty]$ 为满足 $x_0\in D(\psi)$ 的凸函数.证明对一切
$x\in E$,
\begin{equation*}
\partial\varphi(x)+\partial\psi(x)=\partial(\varphi+\psi)(x).
\end{equation*}
[提示: 给定 $f_0\in\partial(\varphi+\psi)(x)$, 对函数
$\tilde\varphi(y)=\varphi(y)-\varphi(x)-\dual{f_0}{y-x}$ 与
$\tilde\psi(y)=\psi(y)-\psi(x)$ 应用定理  \ref{thm:1.12}.]

\runin{-- C --}

1. 设 $\varphi:E\to\R$ 为满足对某个常数 $k\ge0$ 与 $C$, 对一切 $x\in E$
有 $\varphi(x)\le k\norm{x}+C$ 的凸函数.证明对一切 $x_1,x_2\in E$,
\begin{equation*}
\abs{\varphi(x_1)-\varphi(x_2)}\le k\norm{x_1-x_2}.
\end{equation*}
关于 $D(\varphi^\star)$ 能得出什么结论?

2. 设 $A\subset\R^n$ 为开凸集.设 $\varphi:A\to\R$ 为凸函数.证明 $\varphi$
在 $A$ 上连续.

\runin{-- D --}

设 $\varphi:E\to\R$ 为连续凸函数, 设 $C=\{x\in E;\ \varphi(x)\le0\}$.假设存在
某个 $x_0\in E$ 使得 $\varphi(x_0)<0$.给定 $x\in C$, 证明
$f\in\partial I_C(x)$ 当且仅当存在某个 $\lambda\in\R$ 使得
$f\in\lambda\partial\varphi(x)$, 其中若 $\varphi(x)<0$ 则 $\lambda=0$,
若 $\varphi(x)=0$ 则 $\lambda\ge0$.
\end{problem}

\begin{problem}[第~
\ref{chap:1}~章]
\index{Ekelandbianfenyuanli@Ekeland 变分原理 (Ekeland variational principle)}
\index{BishopPhelpsdingli@Bishop--Phelps 定理 (Bishop--Phelps theorem)}
\label{prob:3}
\runin{Ekeland, Brönsted--Rockafellar 与 Bishop--Phelps 定理; $\varepsilon$-次微分}

\runin{-- A --}

设 $M$ 为配以距离 $d(x,y)$ 的非空完备度量空间.设 $\psi:M\to(-\infty,+\infty]$
为有下界的, 下半连续的且满足 $\psi\not\equiv+\infty$ 的函数.我们的目标是
证明存在某个 $a\in M$ 使得
\begin{equation*}
\psi(x)-\psi(a)+d(x,a)\ge0\quad\forall x\in M.
\end{equation*}
给定 $x\in M$, 令
\begin{equation*}
S(x)=\{y\in M;\ \psi(y)-\psi(x)+d(x,y)\le0\}.
\end{equation*}

1. 验证 $x\in S(x)$, 且 $y\in S(x)\Rightarrow S(y)\subset S(x)$.

2. 固定实数的任意序列 $(\varepsilon_n)$, 满足对一切 $n$, $\varepsilon_n>0$,
且 $\varepsilon_n\to0$.给定 $x_0\in M$, 由归纳法构造序列 $(x_n)$ 如下: 一旦
已知 $x_n$, 选取满足下式的任意元素 $x_{n+1}$:
\begin{equation*}
\begin{cases}
x_{n+1}\in S(x_n),\\
\psi(x_{n+1})\le\inf_{x\in S(x_n)}\psi(x)+\varepsilon_{n+1}.
\end{cases}
\end{equation*}
验证对一切 $n$, $S(x_{n+1})\subset S(x_n)$, 且对一切 $n$ 与一切 $p$,
\begin{equation*}
\psi(x_{n+p})-\psi(x_n)+d(x_n,x_{n+p})\le0.
\end{equation*}
由此推出 $(x_n)$ 是 Cauchy 序列, 从而它收敛于一个极限, 记作 $a$.

3. 证明 $a$ 满足所要求的性质.\\
{[提示: 给定 $x\in M$, 考虑两种情形: 或者对一切 $n$, $x\in S(x_n)$; 或者存在
$N$ 使得 $x\notin S(x_N)$.]}

4. 给出几何解释.

\runin{-- B --}

设 $E$ 为 Banach 空间, $\varphi:E\to(-\infty,+\infty]$ 为满足
$\varphi\not\equiv+\infty$ 的凸下半连续函数.给定 $\varepsilon>0$ 与
$x\in D(\varphi)$, 令
\begin{equation*}
\partial_\varepsilon\varphi(x)=\{f\in E^\star;\ \varphi(x)+\varphi^\star(f)-\dual{f}{x}\le\varepsilon\}.
\end{equation*}
验证 $\partial_\varepsilon\varphi(x)\ne\emptyset$.

我们的目的是证明: 给定任何 $x_0\in D(\varphi)$ 与任何
$f_0\in\partial_\varepsilon\varphi(x_0)$, 下述性质成立:
\begin{equation*}
\forall\lambda>0,\ \exists x_1\in D(\varphi)\ \text{与}\ \exists f_1\in E^\star\
\text{满足}\ f_1\in\partial\varphi(x_1),\ \text{使得}\
\norm{x_1-x_0}\le\varepsilon/\lambda\ \text{且}\ \norm{f_1-f_0}\le\lambda.
\end{equation*}
(次微分 $\partial\varphi$ 在问题 2 中定义; 建议在做本题之前先做问题 2.)

1. 考虑由 $\psi(x)=\varphi(x)+\varphi^\star(f_0)-\dual{f_0}{x}$ 定义的函数 $\psi$.
证明存在某个 $x_1\in E$ 使得 $\norm{x_1-x_0}\le\varepsilon/\lambda$ 且对一切
$x\in E$,
\begin{equation*}
\psi(x)-\psi(x_1)+\lambda\norm{x-x_1}\ge0.
\end{equation*}
[提示: 在集合 $M=\{x\in E;\ \psi(x)\le\psi(x_0)-\lambda\norm{x-x_0}\}$ 上应用
第 A 部分的结果.]

2. 得出结论.\\
{[提示: 使用问题 2 第 B6 问的结果.]}

3. 由此推出 $\overline{D(\partial\varphi)}=\overline{D(\varphi)}$ 且
$\overline{R(\partial\varphi)}=\overline{D(\varphi^\star)}$, 其中
$R(\partial\varphi)=\{f\in E^\star;\ \exists x\in D(\partial\varphi)\ \text{使得}\ f\in\partial\varphi(x)\}$.

\runin{-- C --}

设 $E$ 为 Banach 空间, $C\subset E$ 为非空闭凸集.

1. 假设 $C$ 也是有界的, 证明集合
\begin{equation*}
\left\{f\in E^\star;\ \sup_{x\in C}\dual{f}{x}\ \text{可达到}\right\}
\end{equation*}
在 $E^\star$ 中稠密.\\
{[提示: 对函数 $\varphi=I_C$ 应用第 B 部分的结果.]}

2. 若 $E$ 的闭超平面 $H$ 分离 $C$ 与 $\{x\}$, 则称 $H$ 为 $C$ 在点 $x\in C$
处的\emph{支撑超平面}.证明: $C$ 中允许支撑超平面的点集在 $C$ 的边界
$\partial C$ ($=C\setminus\operatorname{Int}C$) 中稠密.
\end{problem}

\begin{problem}[第~
\ref{chap:1}~章]
\label{prob:4}
\runin{Asplund 定理与严格凸范数}
设 $E$ 为赋范向量空间, $\varphi_0,\psi_0:E\to[0,\infty)$ 为两个凸函数, 满足
$\varphi_0(0)=\psi_0(0)=0$ 以及对一切 $x\in E$, $0\le\psi_0(x)\le\varphi_0(x)$.
从 $\varphi_0$ 与 $\psi_0$ 出发, 由归纳法定义两个函数序列 $(\varphi_n)$ 与
$(\psi_n)$ 如下:
\begin{align*}
\varphi_{n+1}(x)&=\frac12(\varphi_n(x)+\psi_n(x)),\\
\psi_{n+1}(x)&=\frac12\inf_{y\in E}\{\varphi_n(x+y)+\psi_n(x-y)\}
=\frac12(\varphi_n\mathbin{\triangledown}\psi_n)(2x).
\end{align*}
[在开始本题之前, 先做涉及下卷积 $\triangledown$ 的习题 1.23.]

\runin{-- A --}

1. 验证对一切 $x\in E$ 与一切 $n$, $0\le\psi_n(x)\le\varphi_n(x)$, 且
$\varphi_n(0)=\psi_n(0)=0$.

2. 验证 $\varphi_n$ 与 $\psi_n$ 是凸的.

3. 证明序列 $(\varphi_n)$ 单调递减, 序列 $(\psi_n)$ 单调递增.由此推出
$(\varphi_n)$ 与 $(\psi_n)$ 有公共的极限, 记作 $\theta$, 满足
$\psi_0\le\theta\le\varphi_0$, 且 $\theta$ 是凸的.

4. 证明 $\varphi_n^\star\uparrow\theta^\star$.

5. 证明 $\psi_{n+1}^\star=\frac12(\varphi_n^\star+\psi_n^\star)$, 并由此推出
当 $D(\psi_0^\star)=E^\star$ 时 $\psi_n^\star\downarrow\theta^\star$.

6. 假设存在某个 $x_0\in D(\varphi_0)$ 使得 $\varphi_0$ 在 $x_0$ 处连续.
证明 $\varphi_n$ 与 $\psi_n$ 在 $x_0$ 处也连续.\\
{[提示: 应用习题 2.1 第 2 问.]}\\
由此推出 $\varphi_{n+1}^\star(f)=\frac12\inf_{g\in E^\star}\{\varphi_n^\star(f+g)+\psi_n^\star(f-g)\}$.

\runin{-- B --}

设 $\varphi:E\to[0,+\infty)$ 为二阶齐次的凸函数, 即对一切 $\lambda\in\R$ 与
一切 $x\in E$, $\varphi(\lambda x)=\lambda^2\varphi(x)$.证明对一切
$x,y\in E$ 与一切 $t\in(0,1)$,
\begin{equation*}
\varphi(x+y)\le\frac1t\varphi(x)+\frac{1}{1-t}\varphi(y).
\end{equation*}
由此推出函数 $x\mapsto\sqrt{\varphi(x)}$ 是半范数, 反之亦然.并建立
\begin{equation}
4\varphi(x)\le\frac1t\varphi(x+y)+\frac{1}{1-t}\varphi(x-y)
\quad\forall x,y\in E,\ \forall t\in(0,1).
\label{eq:prob4.1}
\end{equation}

在下文中, 我们额外假设 $\varphi_0$ 与 $\psi_0$ 是二阶齐次的, 且存在常数
$C>0$ 使得对一切 $x\in E$, $\varphi_0(x)\le(1+C)\psi_0(x)$.

1. 验证 $\varphi_n,\psi_n$ 与 $\theta$ 是二阶齐次的.

2. 证明对每个 $n$, 对一切 $x\in E$ 有
\begin{equation*}
\varphi_n(x)\le\left(1+\frac{C}{4^n}\right)\psi_n(x).
\end{equation*}
[提示: 用数学归纳法并使用 \eqref{eq:prob4.1}.]

3. 假设 $\varphi_0$ 或 $\psi_0$ 严格凸, 证明 $\theta$ 严格凸 (关于严格凸函数的
定义见习题 1.26).\\
{[提示: 使用第 B2 问中建立的不等式.把 $\varphi_n$ 分解为
$\varphi_n=\theta_n+\frac{1}{2^n}\varphi_0$ 是方便的, 其中 $\theta_n$ 是某个
凸函数, 不应试图把它显式写出.注意
$\theta_n+\left(\frac{1}{2^n}-\frac{C}{4^n}\right)\varphi_0\le\theta\le\theta_n+\frac{1}{2^n}\varphi_0$.]}

\runin{-- C --}

假设在 $E$ 上存在两个等价范数, 记作 $\norm{\cdot}_1$ 与 $\norm{\cdot}_2$.
用 $\norm{\cdot}_1^\star$ 与 $\norm{\cdot}_2^\star$ 记 $E^\star$ 上的相应
对偶范数.假设范数 $\norm{\cdot}_1$ 与 $\norm{\cdot}_2^\star$ 严格凸.
利用上述结果, 证明存在与 $\norm{\cdot}_1$ (从而与 $\norm{\cdot}_2$) 等价的
第三个范数 $\norm{\cdot}$, 它与其对偶范数 $\norm{\cdot}^\star$ 都严格凸.
\end{problem}

\begin{problem}[第 1, 2 章]
\label{prob:5}
\runin{正线性泛函}
设 $E$ 为赋范向量空间, $P$ 为以 $0$ 为顶点的凸锥, 即对一切 $x,y\in P$ 及一切
$\lambda,\mu>0$, $\lambda x+\mu y\in P$.令 $F=P-P$, 从而 $F$ 是线性子空间.
考虑下述两条性质:
\begin{enumerate}[label=(\roman*), itemsep=0.1em]
	\item $E$ 上满足对一切 $x\in P$ 有 $f(x)\ge0$ 的每个线性泛函 $f$ 在 $E$ 上连续.
	\item $F$ 是有限余维的闭子空间.
\end{enumerate}
本问题的目标是证明 (i) $\Rightarrow$ (ii); 反之, 当 $E$ 为 Banach 空间且 $P$
闭时, (ii) $\Rightarrow$ (i).

\runin{-- A --}

在整个第 A 部分中, 我们假设 (i) 成立.

1. 证明 $F$ 是闭的.\\
{[提示: 给定任何 $x_0\notin F$, 在 $E$ 上构造一个线性泛函 $f$, 满足 $f(x_0)=1$
且在 $F$ 上 $f=0$.]}

2. 设 $M$ 为 $E$ 的满足 $M\cap F=\{0\}$ 的任意线性子空间.证明 $\dim M<+\infty$.\\
{[提示: 使用习题 1.5.]}

3. 由此推出 (i) $\Rightarrow$ (ii).

\runin{-- B --}

在整个第 B 部分中, 我们假设 $E$ 为 Banach 空间且 $P$ 是闭的.

1. 这里额外假设 $P-P=E$.证明存在常数 $C>0$ 使得每个 $x\in E$ 有分解
$x=y-z$, 其中 $y,z\in P$, $\norm{y}\le C\norm{x}$ 且 $\norm{z}\le C\norm{x}$.\\
{[提示: 考虑集合 $K=\{x=y-z;\ y,z\in P,\ \norm{y}\le1,\ \norm{z}\le1\}$, 遵循
开映射定理 (定理  \ref{thm:2.6}) 证明的思想.]}

2. 由此推出 $P-P=E\Rightarrow\text{(i)}$.\\
{[提示: 用反证法, 考虑 $E$ 中的序列 $(x_n)$ 满足 $\norm{x_n}\le1/2^n$ 且
$f(x_n)\ge1$.然后使用第 B1 问的结果.]}

3. 证明 (ii) $\Rightarrow$ (i).

\runin{-- C --}

在下述例子中确定 $F=P-P$ 并检验 (i) 或 (ii) 是否成立:
\begin{enumerate}[label=(\alph*), itemsep=0.1em]
	\item 配以通常范数的 $E=C([0,1])$ 与 $P=\{u\in E;\ u(t)\ge0\ \forall t\in[0,1]\}$.
	\item 配以通常范数的 $E=C([0,1])$ 与 $P=\{u\in E;\ u(t)\ge0\ \forall t\in[0,1]\ \text{且}\ u(0)=u(1)=0\}$.
	\item 配以通常范数的 $E=\{u\in C^1([0,1]);\ u(0)=u(1)=0\}$ 与 $P=\{u\in E;\ u(t)\ge0\ \forall t\in[0,1]\}$.
\end{enumerate}
\end{problem}

\begin{problem}[第 1, 2 章]
\label{prob:6}
设 $E$ 为 Banach 空间, 设 $A:D(A)\subset E\to E^\star$ 为满足对一切
$x\in D(A)$, $\dual{Ax}{x}\ge0$ 的闭无界算子.

\runin{-- A --}

我们的目的是证明下述性质等价:
\begin{enumerate}[label=(\roman*), itemsep=0.1em]
	\item 对一切 $x\in D(A)$, 存在 $C(x)\in\R$ 使得对一切 $y\in D(A)$, $\dual{Ay}{y-x}\ge C(x)$,
	\item 存在 $k\ge0$ 使得对一切 $x,y\in D(A)$,
	\begin{equation*}
	\abs{\dual{Ay}{x}}\le k(\norm{x}+\norm{Ax})\sqrt{\dual{Ay}{y}}.
	\end{equation*}
\end{enumerate}

1. 证明 (ii) $\Rightarrow$ (i).

反之, 假设 (i) 成立.

2. 证明存在两个常数 $R>0$ 与 $M\ge0$ 使得对一切 $y\in D(A)$ 及一切满足
$\norm{x}+\norm{Ax}\le R$ 的 $x\in D(A)$, $\dual{Ay}{x-y}\le M$.\\
{[提示: 考虑函数 $\varphi(x)=\sup_{y\in D(A)}\dual{Ay}{x-y}$ 并应用习题 2.1.]}

3. 由此推出对一切 $y\in D(A)$ 及一切满足 $\norm{x}+\norm{Ax}\le R$ 的
$x\in D(A)$, $\abs{\dual{Ay}{x}}^2\le4M\dual{Ay}{y}$.

4. 得出结论.

\runin{-- B --}

在下文中假设 $D(A)=E$.设 $\alpha>0$.

1. 证明下述性质等价:
\begin{equation}
\norm{Ay}\le\alpha\sqrt{\dual{Ay}{y}}\quad\forall y\in E,
\tag{iii}
\label{eq:prob6.iii}
\end{equation}
\begin{equation}
\dual{Ay}{y-x}\ge-\frac1{4\alpha^2}\norm{x}^2\quad\forall x,y\in E.
\tag{iv}
\label{eq:prob6.iv}
\end{equation}
[提示: 使用与第 A 部分相同的方法.]

2. 设 $A^\star\in\LL(E^{\star\star},E^\star)$ 为 $A$ 的伴随算子.证明 \eqref{eq:prob6.iv} 等价于
\begin{equation}
\dual{A^\star y}{y-x}\ge-\frac1{4\alpha^2}\norm{x}^2\quad\forall x,y\in E.
\tag{iv$^\star$}
\label{eq:prob6.ivstar}
\end{equation}

3. 由此推出 \eqref{eq:prob6.iii} 等价于
\begin{equation}
\norm{A^\star y}\le\alpha\sqrt{\dual{A^\star y}{y}}\quad\forall y\in E.
\tag{iii$^\star$}
\label{eq:prob6.iiistar}
\end{equation}
\end{problem}

\begin{problem}[第 1, 2 章]
\label{prob:7}
\runin{两个无界线性算子之和的伴随算子}
设 $E$ 为 Banach 空间.给定 $E$ 中的两个闭线性子空间 $M$ 与 $N$, 令
\begin{equation*}
\rho(M,N)=\sup_{\substack{x\in M\\\norm{x}\le1}}\dist(x,N).
\end{equation*}

\runin{-- A --}

1. 验证 $\rho(M,N)\le1$; 若此外 $N\subset M$ 且 $N\ne M$, 证明 $\rho(M,N)=1$.\\
{[提示: 使用引理  \ref{lem:6.1}.]}

2. 设 $L,M$ 与 $N$ 为三个闭线性子空间.令 $a=\rho(M,N)$ 与 $b=\rho(N,L)$.
证明 $\rho(M,L)\le a+b+ab$.由此推出若 $L\subset M$, $a\le1/3$ 且 $b\le1/3$,
则 $L=M$.

3. 证明 $\rho(M,N)=\rho(N^\perp,M^\perp)$.\\
{[提示: 借助定理  \ref{thm:1.12} 验证对一切 $x\in E$ 与一切 $f\in E^\star$,
$\dist(x,N)=\sup_{g\in N^\perp,\norm{g}\le1}\dual{g}{x}$, 以及
$\dist(f,M^\perp)=\sup_{y\in M,\norm{y}\le1}\dual{f}{y}$.]}

\runin{-- B --}

设 $E$ 与 $F$ 为两个 Banach 空间; $E\times F$ 配以范数
$\norm{[u,v]}_{E\times F}=\norm{u}_E+\norm{v}_F$.给定稠密定义且闭的两个
无界算子 $A:D(A)\subset E\to F$ 与 $B:D(B)\subset E\to F$, 令
\begin{equation*}
\rho(A,B)=\rho(G(A),G(B)).
\end{equation*}

1. 证明 $\rho(A,B)=\rho(B^\star,A^\star)$.

2. 证明若 $D(A)\cap D(B)$ 在 $E$ 中稠密, 则
\begin{equation*}
A^\star+B^\star\subset(A+B)^\star.
\end{equation*}
[回忆 $D(A+B)=D(A)\cap D(B)$ 以及 $D(A^\star+B^\star)=D(A^\star)\cap D(B^\star)$.]
可能出现包含关系 $A^\star+B^\star\subset(A+B)^\star$ 是严格包含的情形——构造
这样的一个例子.我们的目的是证明在某些附加假设下等号成立.

3. 假设
\begin{equation}
D(A)\subset D(B)\ \text{且存在常数}\ k\in[0,1)\ \text{与}\ C\ge0\ \text{使得}\
\norm{Bu}\le k\norm{Au}+C\norm{u}\quad\forall u\in D(A).
\tag{H}
\label{eq:prob7.H}
\end{equation}
证明 $A+B$ 是闭的, 且 $\rho(A,A+B)\le k+C$.

4. 除 \eqref{eq:prob7.H} 外, 假设还成立
\begin{equation}
D(A^\star)\subset D(B^\star)\ \text{且存在常数}\ k^\star\in[0,1)\ \text{与}\
C^\star\ge0\ \text{使得}\ \norm{B^\star v}\le k^\star\norm{A^\star v}+C^\star\norm{v}\quad\forall v\in D(A^\star).
\tag{H$^\star$}
\label{eq:prob7.Hstar}
\end{equation}
设 $\varepsilon>0$ 使得 $\varepsilon(k+C)\le1/3$ 且 $\varepsilon(k^\star+C^\star)\le1/3$.
证明 $(A+\varepsilon B)^\star=A^\star+\varepsilon B^\star$.

5. 在假设 \eqref{eq:prob7.H} 与 \eqref{eq:prob7.Hstar} 下证明
$(A+B)^\star=A^\star+B^\star$.\\
{[提示: 使用相继步骤.验证下述不等式对一切 $t\in[0,1]$ 成立:
$\norm{Bu}\le\dfrac{k}{1-k}\norm{Au+tBu}+\dfrac{C}{1-k}\norm{u}\ \forall u\in D(A)$.]}
\end{problem}

\begin{problem}[第 2, 3 章, 仅第 6 问需第~
\ref{chap:4}~章]
\label{prob:8}
\runin{$\ell^1$ 中的弱收敛. Schur 定理}
设 $E=\ell^1$, 从而 $E^\star=\ell^\infty$ (见 \S \ref{sec:11.3}).给定 $x\in E$,
记 $x=(x_1,x_2,\dots,x_i,\dots)$, $\norm{x}_1=\sum_{i=1}^\infty\abs{x_i}$;
给定 $f\in E^\star$, 记 $f=(f_1,f_2,\dots,f_i,\dots)$,
$\norm{f}_\infty=\sup_i\abs{f_i}$.

设 $(x_n)$ 为 $E$ 中的序列, 使得 $x_n\rightharpoonup0$ 弱收敛于 $\sigma(E,E^\star)$.
我们的目标是证明 $\norm{x_n}_1\to0$.

1. 给定 $f,g\in B_{E^\star}$ (即 $\norm{f}_\infty\le1$ 且 $\norm{g}_\infty\le1$), 令
\begin{equation*}
d(f,g)=\sum_{i=1}^\infty\frac{1}{2^i}\abs{f_i-g_i}.
\end{equation*}
验证 $d$ 是 $B_{E^\star}$ 上的度量, 且 $B_{E^\star}$ 关于相应的拓扑是紧的.

2. 给定 $\varepsilon>0$, 令
\begin{equation*}
F_k=\{f\in B_{E^\star};\ \abs{\dual{f}{x_n}}\le\varepsilon\ \forall n\ge k\}.
\end{equation*}
证明存在某个 $f^0\in B_{E^\star}$, 常数 $\rho>0$, 以及整数 $k_0$, 使得
\begin{equation*}
[f\in B_{E^\star}\ \text{且}\ d(f,f^0)<\rho]\implies[f\in F_{k_0}].
\end{equation*}
[提示: 使用 Baire 纲定理.]

3. 固定满足 $(1/2^{N-1})<\rho$ 的整数 $N$.证明对一切 $n\ge k_0$,
\begin{equation*}
\norm{x_n}_1\le\varepsilon+2\sum_{i=1}^N\abs{x_i^n}.
\end{equation*}

4. 得出结论.

5. 用类似的方法证明: 若 $(x_n)$ 是 $\ell^1$ 中的序列使得对每个 $f\in\ell^\infty$,
序列 $(\dual{f}{x_n})$ 收敛于某个极限, 则 $(x_n)$ 在 $\ell^1$ 中强收敛于一个极限.

6. 考虑 $E=L^1(0,1)$, 从而 $E^\star=L^\infty(0,1)$.在 $E$ 中构造一个序列
$(u_n)$ 使得 $u_n\rightharpoonup0$ 弱收敛于 $\sigma(E,E^\star)$, 且对一切 $n$,
$\norm{u_n}_1=1$.
\end{problem}

\begin{problem}[第 1, 2, 3 章]
\label{prob:9}
\runin{弱$^\star$拓扑的 Hahn--Banach 定理及其应用}
设 $E$ 为 Banach 空间.

\runin{-- A --}

1. 设 $A\subset E^\star$ 与 $B\subset E^\star$ 为满足 $A\cap B=\emptyset$ 的两个
非空凸集.假设 $A$ 在拓扑 $\sigma(E^\star,E)$ 中开.证明存在某个 $x\in E$,
$x\ne0$, 以及常数 $\alpha$, 使得超平面 $\{f\in E^\star;\ \dual{f}{x}=\alpha\}$
分离 $A$ 与 $B$.

2. 假设 $A\subset E^\star$ 在 $\sigma(E^\star,E)$ 中闭且 $B\subset E^\star$ 在
$\sigma(E^\star,E)$ 中紧.证明 $A+B$ 在 $\sigma(E^\star,E)$ 中闭.

3. 设 $A\subset E^\star$ 与 $B\subset E^\star$ 为满足 $A\cap B=\emptyset$ 的两个
非空凸集.假设 $A$ 在 $\sigma(E^\star,E)$ 中闭且 $B$ 在 $\sigma(E^\star,E)$ 中紧.
证明存在某个 $x\in E$, $x\ne0$, 以及常数 $\alpha$, 使得超平面
$\{f\in E^\star;\ \dual{f}{x}=\alpha\}$ 严格分离 $A$ 与 $B$.

4. 设 $A\subset E^\star$ 为凸集.证明 $A$ 在 $\sigma(E^\star,E)$ 中的闭包
$\overline A^{\sigma(E^\star,E)}$ 是凸的.

\runin{-- B --}

下面是上述结果的若干应用:

1. 设 $N\subset E^\star$ 为线性子空间.回忆
\begin{equation*}
N^\perp=\{x\in E;\ \dual{f}{x}=0\ \forall f\in N\}\quad\text{且}\quad
N^{\perp\perp}=\{f\in E^\star;\ \dual{f}{x}=0\ \forall x\in N^\perp\}.
\end{equation*}
证明 $N^{\perp\perp}=\overline N^{\sigma(E^\star,E)}$.若 $E$ 自反, 能得出什么结论?
由此推出 $c_0$ 在拓扑 $\sigma(\ell^\infty,\ell^1)$ 下在 $\ell^\infty$ 中稠密.

2. 设 $\varphi:E\to(-\infty,+\infty]$ 为凸下半连续函数, $\varphi\not\equiv+\infty$.
证明 $\psi=\varphi^\star$ 在拓扑 $\sigma(E^\star,E)$ 下下半连续.

反之, 给定关于拓扑 $\sigma(E^\star,E)$ 下半连续的凸函数
$\psi:E^\star\to(-\infty,+\infty]$ ($\psi\not\equiv+\infty$), 证明存在凸下半
连续函数 $\varphi:E\to(-\infty,+\infty]$ ($\varphi\not\equiv+\infty$) 使得
$\psi=\varphi^\star$.

3. 设 $F$ 为另一个 Banach 空间, 设 $A:D(A)\subset E\to F$ 为稠密定义且闭的
无界线性算子.证明:
\begin{enumerate}[label=(\roman*), itemsep=0.1em]
	\item $\overline{R(A^\star)}^{\sigma(E^\star,E)}=N(A)^\perp$,
	\item $\overline{D(A^\star)}^{\sigma(F^\star,F)}=F^\star$.
\end{enumerate}
若 $E$ (相应地 $F$) 自反, 能得出什么结论?

4. 不借助引理~
\ref{lem:3.3} 证明: $J(B_E)$ 在拓扑 $\sigma(E^{\star\star},E^\star)$ 下在
$B_{E^{\star\star}}$ 中稠密 (见引理  \ref{lem:3.4}).

5. 设 $A:B_E\to E^\star$ 为单调映射, 即对一切 $x,y\in B_E$,
$\dual{Ax-Ay}{x-y}\ge0$.令 $S_E=\{x\in E;\ \norm{x}=1\}$.证明
$A(B_E)\subset\overline{\operatorname{conv}A(S_E)}^{\sigma(E^\star,E)}$.
\end{problem}

\begin{problem}[第~
\ref{chap:3}~章]
\index{Eberleinmuliandingli@Eberlein--Šmulian 定理 (Eberlein--Smulian theorem)}
\label{prob:10}
\runin{Eberlein--Šmulian 定理}
设 $E$ 为 Banach 空间, $A\subset E$.令 $B=\overline A^{\sigma(E,E^\star)}$.
本问题的目标是证明下述性质等价:
\begin{enumerate}[label=(\Alph*), itemsep=0.1em]
	\item[(P)] $B$ 在拓扑 $\sigma(E,E^\star)$ 中是紧的.
	\item[(Q)] $A$ 中的每个序列 $(x_n)$ 都有弱收敛的子序列.
\end{enumerate}
此外, (P) (或 (Q)) 蕴含下述性质:
\begin{enumerate}[label=(\Alph*), itemsep=0.1em]
	\item[(R)] 对每个 $y\in B$, 存在序列 $(y_n)\subset A$ 使得 $y_n\rightharpoonup y$ 弱收敛于 $\sigma(E,E^\star)$.
\end{enumerate}

\runin{-- A --}

断言 (P) $\Rightarrow$ (Q) 的证明.

1. 在 $E$ 可分的额外假设下证明 (P) $\Rightarrow$ (Q).\\
{[提示: 考虑 $B_{E^\star}$ 中的一个可数且对拓扑 $\sigma(E^\star,E)$ 在
$B_{E^\star}$ 中稠密的集合 $(b_k)$ (为什么这样的集合存在?).验证量
$d(x,y)=\sum_{k=1}^\infty\frac{1}{2^k}\abs{\dual{b_k}{x-y}}$ 是度量, 并由此推出
$B$ 对 $\sigma(E,E^\star)$ 是可度量化的.]}

2. 在一般情形下证明 (P) $\Rightarrow$ (Q).\\
{[提示: 使用第 A1 问.]}

\runin{-- B --}

后文将需要下述引理:

\runin{引理.} 设 $F$ 为赋范向量空间, $M\subset F^\star$ 为有限维向量空间.则
在 $B_F$ 中存在有限子集 $(a_i)_{1\le i\le k}$ 使得对一切 $g\in M$,
\begin{equation*}
\max_{1\le i\le k}\dual{g}{a_i}\ge\frac12\norm{g}.
\end{equation*}
[提示: 先在 $S_M$ 中选取点 $(g_i)_{1\le i\le k}$ 使得
$S_M\subset\bigcup_{i=1}^kB(g_i,1/4)$, 其中 $S_M=\{g\in M;\ \norm{g}=1\}$.]

\runin{-- C --}

设 $\xi\in E^{\star\star}$ 使得 $\xi\in\overline A^{\sigma(E^{\star\star},E^\star)}$.
利用假设 (Q), 我们将证明 $\xi\in B$ 且存在序列 $(y_k)\subset A$ 使得在
$\sigma(E,E^\star)$ 中 $y_k\rightharpoonup\xi$.

1. 令 $n_1=1$, 固定任何 $f_1\in B_{E^\star}$.证明存在某个 $x_1\in A$ 使得
$\abs{\dual{\xi}{f_1}-\dual{f_1}{x_1}}<1$.

2. 设 $M_1=[\xi,x_1]$ 为由 $\xi$ 与 $x_1$ 张成的线性空间.证明在 $B_{E^\star}$
中存在 $(f_i)_{1<i\le n_2}$ 使得对一切 $\eta\in M_1$,
$\max_{1<i\le n_2}\dual{\eta}{f_i}\ge\frac12\norm{\eta}$.证明存在某个
$x_2\in A$ 使得对一切 $1\le i\le n_2$, $\abs{\dual{\xi}{f_i}-\dual{f_i}{x_2}}<1/2$.

3. 重复上述构造, 我们得到两个序列 $(x_k)\subset A$ 与 $(f_i)\subset B_{E^\star}$,
以及递增的整数序列 $(n_k)$, 使得
\begin{align*}
\max_{n_k<i\le n_{k+1}}\dual{\eta}{f_i}&\ge\frac12\norm{\eta}\quad\forall\eta\in M_k=[\xi,x_1,x_2,\dots,x_k], \tag{a}\\
\abs{\dual{\xi}{f_i}-\dual{f_i}{x_{k+1}}}&<\frac{1}{k+1}\quad\forall i,\ 1\le i\le n_{k+1}. \tag{b}
\end{align*}

4. 由 (a) 推出对一切 $\eta\in\bigcup_{k=1}^\infty M_k=\mathcal{M}$,
$\sup_{i\ge1}\dual{\eta}{f_i}\ge\frac12\norm{\eta}$, 进而对一切 $\eta\in\overline{\mathcal{M}}$,
$\sup_{i\ge1}\dual{\eta}{f_i}\ge\frac12\norm{\eta}$, 其中 $\overline{\mathcal{M}}$
表示 $\mathcal{M}$ 在 $E^{\star\star}$ 中关于强拓扑的闭包.

5. 利用 (b) 与假设 (Q), 证明存在某个 $x\in B\cap\overline{\mathcal{M}}$ 使得
对一切 $i\ge1$, $\dual{\xi}{f_i}=\dual{f_i}{x}$.由此推出 $\xi=x$ 并得出结论.

\runin{-- D --}

1. 证明 (Q) $\Rightarrow$ (P).

2. 证明 (Q) $\Rightarrow$ (R).
\end{problem}
\begin{problem}[第~
\ref{chap:3}~章]
\label{prob:11}
\runin{Banach--Dieudonné--Krein--Šmulian 定理}
设 $E$ 为 Banach 空间, $C\subset E^\star$ 为凸集.假设对每个整数 $n$, 集合
$C\cap(nB_{E^\star})$ 在拓扑 $\sigma(E^\star,E)$ 中是闭的.本问题的目标是证明 $C$
在拓扑 $\sigma(E^\star,E)$ 中是闭的.

\runin{-- A --}

额外假设 $0\notin C$.我们将证明存在 $E$ 中的序列 $(x_n)$ 使得
\begin{equation}
\norm{x_n}\to0\quad\text{且}\quad\sup_n\dual{f}{x_n}>1\quad\forall f\in C.
\label{eq:prob11.1}
\end{equation}
令 $d=\dist(0,C)$, 考虑满足 $d_1>d$ 的递增序列 $d_n\uparrow+\infty$.令
$C_k=\{f\in C;\ \norm{f}\le d_k\}$.

1. 验证集合 $C_k$ 在拓扑 $\sigma(E^\star,E)$ 中是紧的.证明存在某个 $f_0\in C$
使得 $d=\norm{f_0}>0$.

2. 证明存在某个 $x_1\in E$ 使得对一切 $f\in C_1$, $\dual{f}{x_1}>1$.\\
{[提示: 使用弱$^\star$拓扑的 Hahn--Banach 定理; 见问题 9 第 A3 问.]}

3. 令 $A_1=\{x_1\}$.证明存在有限子集 $A_2\subset E$ 使得
$A_2\subset\frac{1}{d_1}B_E$ 且对一切 $f\in C_2$, $\sup_{x\in A_1\cup A_2}\dual{f}{x}>1$.\\
{[提示: 对满足 $A\subset\frac{1}{d_1}B_E$ 的每个有限子集 $A\subset E$, 考虑集合
$Y_A=\left\{f\in C_2;\ \sup_{x\in A_1\cup A}\dual{f}{x}\le1\right\}$, 先证明
$\bigcap_A Y_A=\emptyset$.]}

4. 由归纳法构造有限子集 $A_k\subset E$, 使得
$A_k\subset\frac{1}{d_{k-1}}B_E$ 且对一切 $f\in C_k$, $\sup_{x\in\bigcup_{i=1}^kA_i}\dual{f}{x}>1$.

5. 构造满足 \eqref{eq:prob11.1} 的序列 $(x_n)$.

\runin{-- B --}

1. 再次假设 $0\notin C$.证明存在某个 $x\in E$ 使得对一切 $f\in C$, $\dual{f}{x}\ge1$.\\
{[提示: 设 $(x_n)$ 为满足 \eqref{eq:prob11.1} 的序列.考虑由 $T(f)=(\dual{f}{x_n})_n$
定义的算子 $T:E^\star\to c_0$, 并在 $c_0$ 中分离 $T(C)$ 与 $c_0$ 的开单位球.]}

2. 得出结论.
\end{problem}

\begin{problem}[第 1, 2, 3 章]
\label{prob:12}
在开始本题之前, 必须先做习题 1.23.

设 $E$ 为自反 Banach 空间, $\varphi,\psi:E\to(-\infty,+\infty]$ 为满足
$D(\varphi)\cap D(\psi)\ne\emptyset$ 的凸下半连续函数.令
$\theta=\varphi^\star\mathbin{\triangledown}\psi^\star$.

\runin{-- A --}

我们断言 $D((\varphi+\psi)^\star)=D(\varphi^\star)+D(\psi^\star)$.

1. 证明 $D(\varphi^\star)+D(\psi^\star)\subset D((\varphi+\psi)^\star)$.

2. 证明 $\theta$ 将 $E^\star$ 映入 $(-\infty,+\infty]$, $\theta$ 是凸的,
$D(\theta)=D(\varphi^\star)+D(\psi^\star)$, 且 $\theta^\star=\varphi+\psi$.

3. 由此推出 $D((\varphi+\psi)^\star)=D(\theta)$ 并得出结论.

\runin{-- B --}

额外假设 $\varphi$ 与 $\psi$ 满足
\begin{equation}
\bigcup_{\lambda\ge0}\lambda(D(\varphi)-D(\psi))=E.
\tag{H}
\label{eq:prob12.H}
\end{equation}
我们断言:
\begin{enumerate}[label=(\roman*), itemsep=0.1em]
	\item $(\varphi+\psi)^\star=\varphi^\star\mathbin{\triangledown}\psi^\star$,
	\item $\inf_{x\in E}\{\varphi(x)+\psi(x)\}=\max_{g\in E^\star}\{-\varphi^\star(-g)-\psi^\star(g)\}$,
	\item $D((\varphi+\psi)^\star)=D(\varphi^\star)+D(\psi^\star)$.
\end{enumerate}

1. 证明对每个固定的 $f\in E^\star$ 与 $\alpha\in\R$, 集合
$M=\{g\in E^\star;\ \varphi^\star(f-g)+\psi^\star(g)\le\alpha\}$ 是有界的.\\
{[提示: 使用假设 \eqref{eq:prob12.H} 与推论  \ref{cor:2.5}.]}

2. 设 $\alpha\in\R$ 固定.设 $(f_n)$ 与 $(g_n)$ 为 $E^\star$ 中的两个序列使得
$(f_n)$ 有界且 $\varphi^\star(f_n-g_n)+\psi^\star(g_n)\le\alpha$.证明 $(g_n)$ 有界.

3. 由此推出 $\theta$ 是下半连续的.

4. 证明 (i), (ii) 与 (iii).将这些结果与习题 1.23 第 3 问以及定理  \ref{thm:1.12} 作比较.
\end{problem}

\begin{problem}[第 1, 3 章]
\label{prob:13}
\runin{对偶映射的性质. 一致凸性. 范数的可微性}
设 $E$ 为 Banach 空间.回忆对偶映射的定义 (见第~
\ref{chap:1}~章注记 2): 对每个 $x\in E$,
\begin{equation*}
F(x)=\{f\in E^\star;\ \norm{f}=\norm{x}\ \text{且}\ \dual{f}{x}=\norm{x}^2\}.
\end{equation*}
在开始本题之前, 先做习题 1.1 与 1.25 是有益的.

\runin{-- A --}

假设 $E^\star$ 严格凸, 从而 $F(x)$ 只由单个元素组成.

1. 验证对一切 $x,y\in E$,
\begin{equation*}
\lim_{\lambda\to0}\frac{1}{2\lambda}(\norm{x+\lambda y}^2-\norm{x}^2)=\dual{Fx}{y}.
\end{equation*}
[提示: 应用习题 1.25 的结果; 区分 $\lambda>0$ 与 $\lambda<0$ 的情形.]

2. 证明对一切 $x,y\in E$, 映射 $t\in\R\mapsto\dual{F(x+ty)}{y}$ 在 $t=0$ 处连续.\\
{[提示: 对 $u=x+ty$ 与 $v=x+\lambda y$ 使用不等式 $\frac12(\norm{v}^2-\norm{u}^2)\ge\dual{Fu}{v-u}$.]}

3. 由此推出 $F$ 从 $E$ 强到 $E^\star$ 弱$^\star$连续.\\
{[提示: 使用习题 3.11 的结果.]}\\
当 $E$ 自反或可分时, 用简捷的直接方法证明同样的结果.

4. 验证对一切 $x,y\in E$,
\begin{equation*}
\dual{Fx+Fy}{x+y}+\dual{Fx-Fy}{x-y}=2(\norm{x}^2+\norm{y}^2).
\end{equation*}
由此推出对一切满足 $\norm{x}=\norm{y}=1$ 的 $x,y\in E$,
$\norm{Fx+Fy}+\norm{x-y}\ge2$.

5. 额外假设 $E$ 自反且严格凸.证明 $F$ 是由 $E$ 到 $E^\star$ 上的双射.
验证 $F^{-1}$ 与 $E^\star$ 的对偶映射重合.

\runin{-- B --}

在本部分中假设 $E^\star$ 一致凸.

1. 证明 $F$ 从 $E$ 强到 $E^\star$ 强连续.

2. 更精确地说, 证明 $F$ 在 $E$ 的有界集上一致连续.\\
{[提示: 用反证法并应用第 A4 问.]}

3. 由此推出函数 $\varphi(x)=\frac12\norm{x}^2$ 可微且其微分是 $F$, 即对每个 $x_0\in E$,
\begin{equation*}
\lim_{\substack{x\to x_0\\x\ne x_0}}\frac{\varphi(x)-\varphi(x_0)-\dual{Fx_0}{x-x_0}}{\norm{x-x_0}}=0.
\end{equation*}

\runin{-- C --}

反之, 假设对每个 $x\in E$, 集合 $Fx$ 只包含单个元素, 且 $F$ 在 $E$ 的有界集上一致连续.
证明 $E^\star$ 一致凸.\\
	{[提示: 先证明对一切 $y\in E$ 与一切 $f,g\in E^\star$,
	\begin{equation*}
	\norm{f+g}\le\frac12\norm{f}^2+\frac12\norm{g}^2-\dual{f-g}{y}+\sup_{\substack{x\in E\\\norm{x}\le1}}\{\varphi(x+y)+\varphi(x-y)\}.
	\end{equation*}]}
\end{problem}

\begin{problem}[第 1, 3 章]
\label{prob:14}
\runin{通过下卷积对凸函数作正则化}
设 $E$ 为使得 $E^\star$ 一致凸的 Banach 空间.假设 $\varphi:E\to(-\infty,+\infty]$
为凸下半连续函数且 $\varphi\not\equiv+\infty$.本问题的目标是证明存在可微凸函数序列
$(\varphi_n)$ 使得当 $n\uparrow+\infty$ 时 $\varphi_n\uparrow\varphi$.

\runin{-- A --}

对每个固定的 $x\in E$, 考虑由
\begin{equation*}
\Phi_x(f)=\frac12\norm{f}^2+\varphi^\star(f)-\dual{f}{x},\quad f\in E^\star
\end{equation*}
定义的函数 $\Phi_x:E^\star\to(-\infty,+\infty]$.

1. 验证存在唯一的元素 $f_x\in E^\star$ 使得
$\Phi_x(f_x)=\inf_{f\in E^\star}\Phi_x(f)$.令 $Sx=f_x$.

2. 证明映射 $x\mapsto Sx$ 从 $E$ 强到 $E^\star$ 强连续.\\
{[提示: 先证明 $S$ 从 $E$ 强到 $E^\star$ 弱$^\star$连续.]}

\runin{-- B --}

考虑由 $\psi(x)=\inf_{y\in E}\left\{\frac12\norm{x-y}^2+\varphi(y)\right\}$ 定义的
函数 $\psi:E\to\R$.我们断言 $\psi$ 是凸的, 可微的, 且其微分与 $S$ 重合.

1. 验证 $\psi$ 是凸的, 且对一切 $x\in E$,
\begin{equation}
\psi(x)=-\min_{f\in E^\star}\left\{\frac12\norm{f}^2+\varphi^\star(f)-\dual{f}{x}\right\},
\tag{i}
\end{equation}
对一切 $f\in E^\star$,
\begin{equation}
\psi^\star(f)=\frac12\norm{f}^2+\varphi^\star(f).
\tag{ii}
\end{equation}
[提示: 应用定理  \ref{thm:1.12}.]

2. 由此推出对一切 $x\in E$, $\psi(x)+\psi^\star(Sx)=\dual{Sx}{x}$, 且对一切 $x,y\in E$,
\begin{equation*}
\abs{\psi(y)-\psi(x)-\dual{Sx}{y-x}}\le\norm{Sy-Sx}\norm{y-x}.
\end{equation*}

3. 得出结论.

\runin{-- C --}

对每个整数 $n\ge1$ 与每个 $x\in E$, 令
\begin{equation*}
\varphi_n(x)=\inf_{y\in E}\left\{\frac n2\norm{x-y}^2+\varphi(y)\right\}.
\end{equation*}
证明 $\varphi_n$ 是凸的, 可微的, 且对每个 $x\in E$, 当 $n\uparrow+\infty$ 时
$\varphi_n(x)\uparrow\varphi(x)$.\\
{[提示: 使用与习题 1.24 相同的方法.]}
\end{problem}

\begin{problem}[第~
\ref{chap:1}~章, 仅第 B6 问需第~
\ref{chap:5}~章]
\label{prob:15}
\runin{Chebyshev 意义下的集合中心. 正规结构. Edelstein 意义下的序列渐近中心. Kirk, Browder, Göhde 与 Edelstein 的紧致映射不动点定理}
在开始本题之前, 先做习题 3.29 是有益的.

设 $E$ 为一致凸 Banach 空间, $C\subset E$ 为非空闭凸集.

\runin{-- A --}

设 $A\subset C$ 为非空有界集.对每个 $x\in E$, 定义
\begin{equation*}
\phi(x)=\sup_{y\in A}\norm{x-y}.
\end{equation*}

1. 验证 $\phi$ 是凸函数, 且对一切 $x_1,x_2\in E$,
$\abs{\phi(x_1)-\phi(x_2)}\le\norm{x_1-x_2}$.

2. 证明存在唯一的元素 $c\in C$ 使得 $\phi(c)=\inf_{x\in C}\phi(x)$.点 $c$
称为 $A$ 的\emph{中心}, 记作 $c=\sigma(A)$.

3. 证明若 $A$ 不退化为单点, 则
\begin{equation*}
\phi(\sigma(A))<\operatorname{diam}A=\sup_{x,y\in A}\norm{x-y}.
\end{equation*}

\runin{-- B --}

设 $(a_n)$ 为 $C$ 中的有界序列; 令 $A_n=\bigcup_{i=n}^\infty\{a_i\}$,
对 $x\in E$ 令 $\phi_n(x)=\sup_{y\in A_n}\norm{x-y}$.

1. 对每个 $x\in E$, 考虑 $\phi(x)=\lim_{n\to+\infty}\phi_n(x)$.证明这个极限存在,
且 $\phi$ 在 $E$ 上是凸的, 连续的.

2. 证明存在唯一的元素 $\sigma\in C$ 使得 $\phi(\sigma)=\inf_{x\in C}\phi(x)$.
点 $\sigma$ 称为序列 $(a_n)$ 的\emph{渐近中心}.

3. 设 $\sigma_n=\sigma(A_n)$ 为在第 A2 问意义下集合 $A_n$ 的中心.证明
\begin{equation*}
\lim_{n\to\infty}\phi_n(\sigma_n)=\lim_{n\to\infty}\phi(\sigma_n)=\phi(\sigma),
\end{equation*}
且 $\sigma_n\rightharpoonup\sigma$ 弱收敛于 $\sigma(E,E^\star)$.

4. 由此推出 $\sigma_n\to\sigma$ 强收敛.\\
{[提示: 用反证法并应用习题 3.29 的结果.]}

5. 假设 $a_n\to a$ 强收敛.确定序列 $(a_n)$ 的渐近中心.

6. 这里假设 $E$ 为 Hilbert 空间且 $a_n\rightharpoonup a$ 弱收敛于 $\sigma(E,E^\star)$.
计算 $\phi(x)$ 并确定序列 $(a_n)$ 的渐近中心.\\
{[提示: 展开范数的平方.]}

\runin{-- C --}

假设 $T:C\to C$ 是压缩映射, 即对一切 $x,y\in C$, $\norm{Tx-Ty}\le\norm{x-y}$.
设 $a\in C$ 给定, 设 $a_n=T^na$ 为其迭代序列.假设序列 $(a_n)$ 有界.设 $\sigma$
为序列 $(a_n)$ 的渐近中心.

1. 证明 $\sigma$ 是 $T$ 的不动点, 即 $T\sigma=\sigma$.

2. 验证 $T$ 的不动点集是闭凸集.
\end{problem}
\begin{problem}[第 2, 3 章]
\label{prob:16}
\runin{线性极大单调算子的表征}
设 $E$ 为 Banach 空间, 设 $A:D(A)\subset E\to E^\star$ 为满足单调性条件
\begin{equation}
\dual{Au}{u}\ge0\quad\forall u\in D(A)
\tag{M}
\label{eq:prob16.M}
\end{equation}
的无界线性算子.我们用 (P) 记下述性质:
\begin{equation}
\text{若}\ x\in E\ \text{且}\ f\in E^\star\ \text{满足}\
\dual{Au-f}{u-x}\ge0\quad\forall u\in D(A),
\ \text{则}\ x\in D(A)\ \text{且}\ Ax=f.
\tag{P}
\label{eq:prob16.P}
\end{equation}

\runin{-- A --}

1. 证明若 (P) 成立, 则 $D(A)$ 在 $E$ 中稠密.\\
{[提示: 证明若 $f\in E^\star$ 且对一切 $u\in D(A)$ 有 $\dual{f}{u}=0$, 则 $f=0$.]}

2. 证明若 (P) 成立, 则 $A$ 是闭的.

3. 证明函数 $u\in D(A)\mapsto\dual{Au}{u}$ 是凸的.

4. 证明 $N(A)\subset R(A)^\perp$.由此推出若 $D(A)$ 在 $E$ 中稠密, 则 $N(A)\subset N(A^\star)$.

5. 证明若 $D(A)=E$, 则 (P) 成立.

在整个本题的其余部分中, 我们额外假设:
\begin{enumerate}[label=(\roman*), itemsep=0.1em]
	\item $E$ 自反, $E$ 与 $E^\star$ 严格凸,
	\item $D(A)$ 在 $E$ 中稠密且 $A$ 是闭的,
\end{enumerate}
从而 $A^\star:D(A^\star)\subset E\to E^\star$ 且 $D(A^\star)$ 在 $E$ 中稠密 (为什么?).
本问题的目标是建立等价性 (P) $\iff$ (M$^\star$), 其中 (M$^\star$) 记下述性质:
\begin{equation}
\dual{A^\star v}{v}\ge0\quad\forall v\in D(A^\star).
\tag{M$^\star$}
\label{eq:prob16.Mstar}
\end{equation}

\runin{-- B --}

在本节中假设 (P) 成立.

1. 证明对一切 $v\in D(A)\cap D(A^\star)$, $\dual{A^\star v}{v}\ge0$.

2. 设 $v\in D(A^\star)$ 满足 $v\notin D(A)$.证明对一切 $f\in E^\star$, 存在 $u\in D(A)$
使得 $\dual{Au-f}{u-v}<0$.选取 $f=-A^\star v$, 证明 $\dual{A^\star v}{v}>0$.
由此推出 (M$^\star$) 成立.

3. 证明 $N(A)=N(A^\star)$ 且 $R(A)=R(A^\star)$.

\runin{-- C --}

在本部分中假设 (M$^\star$) 成立.

1. 验证配以图范数 $\norm{u}_{D(A)}=\norm{u}_E+\norm{Au}_{E^\star}$ 的空间 $D(A)$
是自反的.

2. 给定 $x\in E$ 与 $f\in E^\star$, 考虑在 $D(A)$ 上由
\begin{equation*}
\phi(u)=\frac12\norm{Au-f}^2+\frac12\norm{u-x}^2+\dual{Au-f}{u-x}
\end{equation*}
定义的函数 $\phi$.证明 $\phi$ 在 $D(A)$ 上是凸的, 连续的.证明当 $\norm{u}_{D(A)}\to\infty$
时 $\phi(u)\to+\infty$.

3. 由此推出存在某个 $u_0\in D(A)$ 使得对一切 $u\in D(A)$, $\phi(u_0)\le\phi(u)$.
通过选取 $u=u_0+tv$ ($v\in D(A),t>0$) 并令 $t\to0$, 能得到关于 $A$ 与 $A^\star$ 的
什么方程? \\
{[提示: 应用问题 13 第 A 部分的结果.]}

4. 证明 (M$^\star$) $\Rightarrow$ (P).

5. 由此推出 $A^\star$ 也满足性质 (P).
\end{problem}

\begin{problem}[第 1, 3, 4 章]
\label{prob:17}
\runin{-- A --}

设 $E$ 为自反 Banach 空间, $M$ 为 $E$ 的闭线性子空间.设 $C$ 为 $E^\star$ 的凸子集.
对每个 $u\in E$ 令
\begin{equation*}
\phi(u)=\sup_{g\in C}\dual{g}{u}.
\end{equation*}

1. 证明对每个 $f\in M^\perp+C$, 对一切 $u\in M$ 有 $\phi(u)\ge\dual{f}{u}$.\\
{[提示: 先从 $f\in M^\perp+C$ 的情形开始.]}

2. 反之, 设 $f\in E^\star$ 满足对一切 $u\in M$, $\phi(u)\ge\dual{f}{u}$.证明
$f\in M^\perp+C$.\\
{[提示: 使用 Hahn--Banach 定理.]}

3. 假设 $C$ 是闭且有界的, 证明 $M^\perp+C$ 是闭的.

\runin{-- B --}

在本节中假设 $E=L^p(\Omega)$ ($1<p<\infty$),
\begin{equation*}
M=\left\{u\in L^p(\Omega);\ \int_\Omega ju=0\right\},
\end{equation*}
以及
\begin{equation*}
C=\{g\in L^{p'}(\Omega);\ \abs{g(x)}\le k(x)\ \text{在}\ \Omega\ \text{上几乎处处}\},
\end{equation*}
其中 $j$ 与 $k\ge0$ 是 $L^{p'}(\Omega)$ 中的给定函数.

1. 验证 $M$ 是闭线性子空间, 且 $C$ 是凸集, 闭集和有界集.

2. 确定 $M^\perp$.

3. 对每个 $u\in L^p(\Omega)$ 确定 $\phi(u)$.

4. 由此推出若 $f\in L^{p'}(\Omega)$ 满足对一切 $u\in M$, $\int_\Omega k\abs{u}\ge\int_\Omega fu$,
则存在常数 $\lambda\in\R$ 与函数 $g\in C$ 使得 $f=\lambda j+g$.

5. 证明逆命题也成立.

\runin{-- C --}

设 $M\subset L^1(\Omega)$ 为线性子空间.设 $f,g\in L^\infty(\Omega)$ 满足在 $\Omega$
上几乎处处 $f\le g$.证明下述性质等价:
\begin{enumerate}[label=(\roman*), itemsep=0.1em]
	\item 存在 $\phi\in M^\perp$ 使得在 $\Omega$ 上几乎处处 $f\le\phi\le g$,
	\item 对一切 $u\in M$, $\int_\Omega(fu^+-gu^-)\le0$,
\end{enumerate}
其中 $u^+=\max\{u,0\}$, $u^-=\max\{-u,0\}$.\\
{[提示: 假设 (ii), 验证对一切 $u\in M$, $\int_\Omega(g+f)u\le\int_\Omega(g-f)\abs{u}$,
并应用定理  \ref{thm:1.12} 找到满足 $\abs{\psi}\le g-f$ 的某个 $\psi\in L^\infty(\Omega)$
使得 $\psi-(g+f)\in M^\perp$.取 $\phi=\frac12(g+f-\psi)$.]}
\end{problem}

\begin{problem}[第 3, 4 章]
\label{prob:18}
设 $\Omega$ 为有限测度的测度空间.设 $1<p<\infty$.设 $g:\R\to\R$ 为连续单调不减函数,
满足对某个常数 $C$, 对一切 $t\in\R$, $\abs{g(t)}\le C(\abs{t}^{p-1}+1)$.令
$G(t)=\int_0^tg(s)ds$.

1. 验证对每个 $u\in L^p(\Omega)$, 有 $g(u)\in L^{p'}(\Omega)$ 且 $G(u)\in L^1(\Omega)$.

设 $(u_n)$ 为 $L^p(\Omega)$ 中的序列, $u\in L^p(\Omega)$ 满足:
(i) $u_n\rightharpoonup u$ 弱收敛于 $\sigma(L^p,L^{p'})$, 且
(ii) $\limsup_{n\to\infty}\int_\Omega G(u_n)\le\int_\Omega G(u)$.
本问题的目的是建立下述性质:
\begin{equation}
\text{对每个}\ q\in[1,p'),\ g(u_n)\to g(u)\ \text{在}\ L^q\ \text{中强收敛},
\tag{1}
\label{eq:prob18.1}
\end{equation}
\begin{equation}
\text{若此外假设}\ g\ \text{严格递增, 则对每个}\ q\in[1,p),\ u_n\to u\ \text{在}\ L^q\ \text{中强收敛}.
\tag{2}
\label{eq:prob18.2}
\end{equation}

2. 验证对一切 $a,b\in\R$, $G(a)-G(b)-g(b)(a-b)\ge0$.若 $G(a)-G(b)-g(b)(a-b)=0$,
能得出什么结论?

3. 设 $(a_n)$ 为 $\R$ 中的序列, $b\in\R$ 使得
$\lim_{n\to\infty}[G(a_n)-G(b)-g(b)(a_n-b)]=0$.证明 $g(a_n)\to g(b)$.

4. 证明 $\int_\Omega\abs{G(u_n)-G(u)-g(u)(u_n-u)}\to0$.由此推出存在子序列 $(u_{n_k})$
使得在 $\Omega$ 上几乎处处 $G(u_{n_k})-G(u)-g(u)(u_{n_k}-u)\to0$, 从而在 $\Omega$
上几乎处处 $g(u_{n_k})\to g(u)$.

5. 证明 \eqref{eq:prob18.1} 成立 (对整个序列验证 \eqref{eq:prob18.1}, 而不仅对子序列).

6. 证明 \eqref{eq:prob18.2} 成立.

在下文中, 我们额外假设存在常数 $\alpha>0$ 与 $C$ 使得
\begin{equation}
\abs{g(t)}\ge\alpha\abs{t}^{p-1}-C\quad\forall t\in\R.
\tag{3}
\label{eq:prob18.3}
\end{equation}

7. 证明在 $L^{p'}$ 中 $g(u_n)\to g(u)$ 强收敛.

8. 没有假设 \eqref{eq:prob18.3} 能得出同样的结论吗?

9. 若此外 $g$ 严格递增, 证明在 $L^p$ 中 $u_n\to u$ 强收敛.
\end{problem}

\begin{problem}[第 3, 4 章]
\label{prob:19}
设 $E$ 为配以范数 $\norm{u}_E=\norm{u}_1+\norm{u}_2$ 的空间 $L^1(\R)\cap L^2(\R)$.

1. 验证 $E$ 是 Banach 空间.设 $f(x)=f_1(x)+f_2(x)$, 其中 $f_1\in L^\infty(\R)$
且 $f_2\in L^2(\R)$.验证映射 $u\mapsto\int_\R f(x)u(x)dx$ 是 $E$ 上的连续线性泛函.

2. 设 $0<\alpha<1/2$; 验证映射 $u\mapsto\int_\R\dfrac{1}{\abs{x}^\alpha}u(x)dx$
是 $E$ 上的连续线性泛函.\\
{[提示: 把积分拆成两部分: $[\abs{x}>M]$ 与 $[\abs{x}\le M]$.]}

3. 令
\begin{equation*}
K=\left\{u\in E;\ u\ge0\ \text{几乎处处在}\ \R\ \text{上且}\ \int_\R u(x)dx\le1\right\}.
\end{equation*}
验证 $K$ 是 $E$ 的闭凸子集.

4. 设 $(u_n)$ 为 $K$ 中的序列, 设 $u\in K$ 使得在 $L^2(\R)$ 中 $u_n\rightharpoonup u$
弱收敛.验证 $u\in K$ 并证明
\begin{equation*}
\int_\R\frac{1}{\abs{x}^\alpha}u_n(x)dx\to\int_\R\frac{1}{\abs{x}^\alpha}u(x)dx.
\end{equation*}

考虑对每个 $u\in E$ 由 $J(u)=\int_\R u^2(x)dx-\int_\R\dfrac{1}{\abs{x}^\alpha}u(x)dx$
定义的函数 $J$.

5. 验证存在常数 $C$ 使得对一切 $u\in K$, $J(u)\ge C$.我们断言 $m=\inf_{u\in K}J(u)$
可以达到.

6. 设 $(u_n)$ 为 $K$ 中的序列使得 $J(u_n)\to m$.证明 $\norm{u_n}_E$ 有界.

7. 设 $(u_{n_k})$ 为使得在 $L^2(\R)$ 中 $u_{n_k}\rightharpoonup u$ 弱收敛的子序列.
证明 $J(u)=m$.

8. $E$ 是自反空间吗?
\end{problem}

\begin{problem}[第~
\ref{chap:4}~章]
\label{prob:20}
\runin{Clarkson 不等式. $L^p$ 的一致凸性}

\runin{-- A --}

在本部分中假设 $2\le p<\infty$, 我们将建立下述不等式:
\begin{equation}
\abs{x+y}^p+\abs{x-y}^p\le2(\abs{x}^{p'}+\abs{y}^{p'})^{p/p'}\quad\forall x,y\in\R,
\label{eq:prob20.1}
\end{equation}
\begin{equation}
2(\abs{x}^{p'}+\abs{y}^{p'})^{p/p'}\le2^{p-1}(\abs{x}^p+\abs{y}^p)\quad\forall x,y\in\R.
\label{eq:prob20.2}
\end{equation}

1. 证明 \eqref{eq:prob20.2}.\\
{[提示: 利用函数 $g(t)=\abs{t}^{p/p'}$ 的凸性.]}

2. 令 $f(x)=(1+x^{1/p})^p+(1-x^{1/p})^p$, $x\in(0,1)$.证明对一切 $x\in(0,1)$,
$f''(x)\le0$.由此推出
\begin{equation}
f(x)\le f(y)+(x-y)f'(y)\quad\forall x,y\in(0,1).
\label{eq:prob20.3}
\end{equation}

3. 证明对一切 $x\in(0,1)$, $f(x)\le2(1+x^{p'/p})^{p/p'}$.\\
{[提示: 在 \eqref{eq:prob20.3} 中取 $y=x^{p'}$.]}

4. 推出 \eqref{eq:prob20.1}.

在下文中, $\Omega$ 表示 $\sigma$-有限测度空间.

\runin{-- B --}

在本部分中再次假设 $2\le p<\infty$.

1. 证明下述不等式:
\begin{equation}
\norm{f+g}_p^p+\norm{f-g}_p^p\le2(\norm{f}_p^{p'}+\norm{g}_p^{p'})^{p/p'}\quad\forall f,g\in L^p(\Omega),
\label{eq:prob20.4}
\end{equation}
\begin{equation}
2(\norm{f}_p^{p'}+\norm{g}_p^{p'})^{p/p'}\le2^{p-1}(\norm{f}_p^p+\norm{g}_p^p)\quad\forall f,g\in L^p(\Omega).
\label{eq:prob20.5}
\end{equation}

2. 推出 Clarkson 第一不等式 (见定理  \ref{thm:4.10}).

\runin{-- C --}

在本部分中假设 $1<p\le2$.

1. 建立下述不等式:
\begin{equation}
\norm{f+g}_p^{p'}+\norm{f-g}_p^{p'}\le2(\norm{f}_p^p+\norm{g}_p^p)^{p'/p}\quad\forall f,g\in L^p(\Omega).
\label{eq:prob20.6}
\end{equation}
不等式 \eqref{eq:prob20.6} 称为 Clarkson 第二不等式.\\
{[提示: 有两种不同的方法:
(i) 由 \eqref{eq:prob20.4} 通过对偶性, 注意到
$\sup_{\phi,\psi\in L^{p'}}\dfrac{\int(u\phi+v\psi)}{[\norm{\phi}_{p'}^p+\norm{\psi}_{p'}^p]^{1/p}}=(\norm{u}_p^{p'}+\norm{v}_p^{p'})^{1/p'}$.
(ii) 直接由 \eqref{eq:prob20.1} 结合习题 4.11 的结果.]}

2. 由此推出对 $1<p\le2$, $L^p(\Omega)$ 是一致凸的.
\end{problem}
\begin{problem}[第~
\ref{chap:4}~章]
\index{Marcinkiewiczkongjian@Marcinkiewicz 空间 (Marcinkiewicz space)}
\index{fenbuhanshu@分布函数 (distribution function)}
\label{prob:21}
\runin{分布函数. Marcinkiewicz 空间}
在整个本题中, $\Omega$ 表示具有有限测度 $\mu$ 的测度空间.给定可测函数
$f:\Omega\to\R$, 我们定义其\emph{分布函数} $\alpha$ 为
\begin{equation*}
\alpha(t)=\abs{[\abs{f}>t]}=\operatorname{meas}\{x\in\Omega;\ \abs{f(x)}>t\}\quad\forall t\ge0.
\end{equation*}

\runin{-- A --}

1. 验证 $\alpha$ 单调不增.证明对一切 $t\ge0$, $\alpha(t+0)=\alpha(t)$.构造一个简单的例子,
使得对某个 $t>0$, $\alpha(t-0)\ne\alpha(t)$.

2. 设 $(f_n)$ 为可测函数序列, 使得在 $\Omega$ 上几乎处处 $f_n\to f$.用 $(\alpha_n)$ 与
$\alpha$ 记相应的分布函数.证明对一切 $t>0$,
\begin{equation*}
\alpha(t)\le\liminf_{n\to\infty}\alpha_n(t)\le\limsup_{n\to\infty}\alpha_n(t)\le\alpha(t-0).
\end{equation*}
由此推出 $\alpha_n(t)\to\alpha(t)$ 几乎处处成立.

\runin{-- B --}

1. 设 $g\in L^1_{\text{loc}}(\R)$ 为满足几乎处处 $g\ge0$ 的函数.令 $G(t)=\int_0^tg(s)ds$.
证明对每个可测函数 $f$,
\begin{equation*}
\int_\Omega G(\abs{f(x)})d\mu<\infty\iff\int_0^\infty\alpha(t)g(t)dt<\infty,
\end{equation*}
且
\begin{equation*}
\int_\Omega G(\abs{f(x)})d\mu=\int_0^\infty\alpha(t)g(t)dt.
\end{equation*}
[提示: 使用 Fubini 与 Tonelli 定理.]

2. 更一般地, 证明对一切 $\lambda\ge0$,
\begin{equation*}
\int_{[\abs{f}>\lambda]}G(\abs{f(x)})d\mu=\alpha(\lambda)G(\lambda)+\int_\lambda^\infty\alpha(t)g(t)dt.
\end{equation*}

3. 由此推出对 $1\le p<\infty$,
\begin{equation*}
f\in L^p(\Omega)\iff\int_0^\infty\alpha(t)t^{p-1}dt<\infty,
\end{equation*}
且对一切 $\lambda\ge0$,
\begin{equation*}
\int_{[\abs{f}>\lambda]}\abs{f(x)}^pd\mu=\alpha(\lambda)\lambda^p+p\int_\lambda^\infty\alpha(t)t^{p-1}dt.
\end{equation*}
验证若 $f\in L^p(\Omega)$, 则 $\lim_{t\to+\infty}\alpha(t)t^p=0$.

\runin{-- C --}

设 $1<p<\infty$.对每个 $f\in L^1(\Omega)$, 定义
\begin{equation*}
[f]_p=\sup\left\{\abs{A}^{-1/p'}\int_A\abs{f};\ A\subset\Omega\ \text{可测},\ \abs{A}>0\right\}\le\infty,
\end{equation*}
并考虑集合 $M^p(\Omega)=\{f\in L^1(\Omega);\ [f]_p<\infty\}$, 称为 $p$ 阶
\emph{Marcinkiewicz 空间}.空间 $M^p$ 也称为\emph{弱 $L^p$ 空间}, 但这个术语
容易引起混淆, 因为``弱''字已经用于弱拓扑.

1. 验证 $M^p(\Omega)$ 是线性空间, 且 $[\cdot]_p$ 是范数.证明 $L^p(\Omega)\subset M^p(\Omega)$,
且对每个 $f\in L^p(\Omega)$, $[f]_p\le\norm{f}_p$.

2. 证明配以范数 $[\cdot]_p$ 的 $M^p(\Omega)$ 是 Banach 空间.验证对 $1<q\le p$,
$M^p(\Omega)\subset M^q(\Omega)$ 连续嵌入.

我们断言: $f\in M^p(\Omega)$ 当且仅当 $f$ 可测且 $\sup_{t>0}t^p\alpha(t)<\infty$.

3. 证明若 $f\in M^p(\Omega)$, 则对一切 $t>0$, $t^p\alpha(t)\le[f]_p^p$.

4. 反之, 设 $f$ 为满足 $\sup_{t>0}t^p\alpha(t)<\infty$ 的可测函数.证明存在常数 $C_p$
(只依赖于 $p$) 使得
\begin{equation*}
[f]_p^p\le C_p\sup_{t>0}t^p\alpha(t).
\end{equation*}
[提示: 使用第 B3 问并写出 $\int_A\abs{f}=\int_{A\cap[\abs{f}>\lambda]}\abs{f}+\int_{A\cap[\abs{f}\le\lambda]}\abs{f}$; 然后变动 $\lambda$.]

5. 证明对 $1\le q<p$, $M^p(\Omega)\subset L^q(\Omega)$ 连续嵌入.

6. 设 $1<q<r<\infty$ 与 $\theta\in(0,1)$; 令 $\dfrac1p=\dfrac\theta q+\dfrac{1-\theta}{r}$.
证明存在只依赖于 $q,r$ 与 $\theta$ 的常数 $C$ 使得对一切 $f\in M^r(\Omega)$,
\begin{equation*}
\norm{f}_p\le C[f]_q^\theta[f]_r^{1-\theta}.
\end{equation*}

7. 设 $\Omega=\{x\in\R^N;\ \abs{x}<1\}$ 配以 Lebesgue 测度, 设 $f(x)=\abs{x}^{-N/p}$
($1<p<\infty$).验证 $f\in M^p(\Omega)$, 而 $f\notin L^p(\Omega)$.
\end{problem}

\begin{problem}[第~
\ref{chap:4}~章]
\label{prob:22}
\runin{一个插值定理 (Schur, Riesz, Thorin, Marcinkiewicz)}
设 $\Omega$ 为有限测度的测度空间.设 $T:L^1(\Omega)\to L^1(\Omega)$ 为有界线性算子,
其范数记作 $N_1=\norm{T}_{\LL(L^1,L^1)}$.我们假设 $T(L^\infty(\Omega))\subset L^\infty(\Omega)$.

1. 证明 $T$ 是由 $L^\infty(\Omega)$ 到自身中的有界算子.令 $N_\infty=\norm{T}_{\LL(L^\infty,L^\infty)}$.

本问题的目标是证明对每个 $1<p<\infty$,
\begin{equation*}
T(L^p(\Omega))\subset L^p(\Omega),
\end{equation*}
且 $T:L^p(\Omega)\to L^p(\Omega)$ 是有界算子, 其范数 $N_p=\norm{T}_{\LL(L^p,L^p)}$
满足不等式 $N_p\le2N_1^{1/p}N_\infty^{1/p'}$.

为简单起见, 先设 $N_\infty=1$.给定函数 $u\in L^1(\Omega)$, 对每个 $\lambda>0$, 我们令
$u=v_\lambda+w_\lambda$, 其中 $v_\lambda=u\chi_{[\abs{u}>\lambda]}$ 且
$w_\lambda=u\chi_{[\abs{u}\le\lambda]}$, $f=Tu$, $g_\lambda=Tv_\lambda$, 以及
$h_\lambda=Tw_\lambda$, 从而 $f=g_\lambda+h_\lambda$.

2. 验证对一切 $\lambda>0$, $\norm{g_\lambda}_1\le N_1\int_{[\abs{u}>\lambda]}\abs{u(x)}d\mu$,
且 $\norm{h_\lambda}_\infty\le\lambda$.

3. 考虑分布函数 $\alpha(t)=\abs{[\abs{u}>t]}$, $\beta(t)=\abs{[\abs{f}>t]}$, 以及
$\gamma_\lambda(t)=\abs{[\abs{g_\lambda}>t]}$.证明对一切 $\lambda>0$,
\begin{equation*}
\int_0^\infty\gamma_\lambda(t)dt\le N_1\left[\alpha(\lambda)\lambda+\int_\lambda^\infty\alpha(t)dt\right],
\end{equation*}
且对一切 $\lambda>0$ 及一切 $t>\lambda$, $\beta(t)\le\gamma_\lambda(t-\lambda)$.\\
{[提示: 应用问题 21 第 B 部分的结果.]}

4. 假设 $u\in L^p(\Omega)$, 证明 $f\in L^p(\Omega)$ 且 $\norm{f}_p\le2N_1^{1/p}\norm{u}_p$.

5. 在 $N_\infty\ne1$ 的一般情形下得出结论.

\runin{注记.} 事实上通过不同的论证可以证明 $N_p\le N_1^{1/p}N_\infty^{1/p'}$;
参见例如 Bergh--Löfström \bref{BerghLofstrom}{1} 及其书第~
\ref{chap:1}~章注记中的文献.
\end{problem}

\begin{problem}[第 3, 4 章]
\label{prob:23}
\runin{$L^1$ 的弱紧子集与等度可积族. Hahn--Vitali--Saks, Dunford--Pettis 与 de la Vallée-Poussin 定理}
设 $\Omega$ 为 $\sigma$-有限测度空间.回忆 (见习题 4.36), 若子集 $\mathcal{F}\subset L^1(\Omega)$
满足下述性质, 则称它是\emph{等度可积}的:
\begin{enumerate}[label=(\alph*), itemsep=0.1em]
	\item $\mathcal{F}$ 在 $L^1(\Omega)$ 中有界,
	\item 对一切 $\varepsilon>0$, 存在 $\delta>0$ 使得对满足 $\abs{A}<\delta$ 的一切可测集 $A\subset\Omega$ 及一切 $f\in\mathcal{F}$, $\int_A\abs{f}<\varepsilon$,
	\item 对一切 $\varepsilon>0$, 存在满足 $\abs{\omega}<\infty$ 的可测集 $\omega\subset\Omega$ 使得对一切 $f\in\mathcal{F}$, $\int_{\Omega\setminus\omega}\abs{f}<\varepsilon$.
\end{enumerate}
本问题的第一个目标是建立给定集合 $\mathcal{F}\subset L^1(\Omega)$ 的下述性质等价:
\begin{enumerate}[label=(\roman*), itemsep=0.1em]
	\item $\mathcal{F}$ 包含在 $L^1(\Omega)$ 的弱 (关于 $\sigma(L^1,L^\infty)$) 紧子集中,
	\item $\mathcal{F}$ 是等度可积的.
\end{enumerate}

\runin{-- A --}

蕴含关系 (i) $\Rightarrow$ (ii).

1. 设 $(f_n)$ 为 $L^1(\Omega)$ 中的序列, 使得对满足 $\abs{A}<\infty$ 的一切可测集
$A\subset\Omega$, $\int_A f_n\to0$.证明 $(f_n)$ 满足性质 (b).\\
{[提示: 考虑由 $X=\{\chi_A;\ A\subset\Omega\ \text{可测且}\ \abs{A}<\infty\}$ 定义的
子集 $X\subset L^1(\Omega)$.验证 $X$ 在 $L^1(\Omega)$ 中闭, 并对序列
$X_n=\{\chi_A\in X;\ \abs{\int_A f_k}\le\varepsilon\ \forall k\ge n\}$ 应用 Baire 纲定理,
其中 $\varepsilon>0$ 固定.]}

2. 设 $(f_n)$ 为 $L^1(\Omega)$ 中的序列, 使得对满足 $\abs{A}\le\infty$ 的一切可测集
$A\subset\Omega$, $\int_A f_n\to0$.证明 $(f_n)$ 满足性质 (c).\\
{[提示: 设 $(\Omega_i)$ 为满足 $\Omega=\bigcup_i\Omega_i$ 的测度有限的可测集的单调递增序列.
在 $L^\infty(\Omega)$ 上考虑由 $d(f,g)=\sum_i\frac{1}{2^i\abs{\Omega_i}}\int_{\Omega_i}\abs{f-g}$
定义的度量 $d$.令 $Y=\{\chi_A;\ A\subset\Omega\ \text{可测}\}$.验证 $Y$ 关于度量 $d$
完备, 并对序列 $Y_n=\{\chi_A\in Y;\ \abs{\int_A f_k}\le\varepsilon\ \forall k\ge n\}$
应用 Baire 纲定理.]}

3. 由此推出: 若 $(f_n)$ 是 $L^1(\Omega)$ 中的序列且 $f_n\rightharpoonup f$ 弱收敛于
$\sigma(L^1,L^\infty)$, 则 $(f_n)$ 是等度可积的.

4. 证明 (i) $\Rightarrow$ (ii).\\
{[提示: 用反证法并应用 Eberlein--Šmulian 定理; 见问题 10.]}

5. 仅假设对满足 $\abs{A}<\infty$ (相应地 $\abs{A}\le\infty$) 的每个可测集 $A\subset\Omega$,
$\int_A f_n$ 收敛于有限极限 $\ell(A)$, 重新讨论第 1 问 (相应地第 2 问).

\runin{-- B --}

蕴含关系 (ii) $\Rightarrow$ (i).

1. 设 $E$ 为 Banach 空间, $F\subset E$.假设对一切 $\varepsilon>0$, 存在弱
(关于 $\sigma(E,E^\star)$) 紧集 $F_\varepsilon\subset E$ 使得
$F\subset F_\varepsilon+\varepsilon B_E$.证明 $F$ 包含在 $E$ 的弱紧子集中.\\
{[提示: 考虑 $G=\overline F^{\sigma(E^{\star\star},E^\star)}$.]}

2. 由此推出 (ii) $\Rightarrow$ (i).\\
{[提示: 考虑族 $(\chi_\omega T_nf)_{f\in\mathcal{F}}$, 其中 $\abs{\omega}<\infty$,
$T_n$ 是定理  \ref{thm:4.12} 证明中的截断算子.]}

\runin{-- C --}

若干应用.

1. 设 $(f_n)$ 为 $L^1(\Omega)$ 中的序列使得 $f_n\rightharpoonup f$ 弱收敛于
$\sigma(L^1,L^\infty)$ 且 $f_n\to f$ 几乎处处.证明 $\norm{f_n-f}_1\to0$.\\
{[提示: 应用习题 4.14.]}

2. 设 $u_1,u_2\in L^1(\Omega)$ 满足几乎处处 $u_1\le u_2$.证明集合
$K=\{f\in L^1(\Omega);\ u_1\le f\le u_2\ \text{几乎处处}\}$ 在弱拓扑
$\sigma(L^1,L^\infty)$ 中是紧的.

3. 设 $(f_n)$ 为 $L^1(\Omega)$ 中的等度可积序列.证明存在弱收敛于 $\sigma(L^1,L^\infty)$
的子序列 $(f_{n_k})$.

4. 设 $(f_n)$ 为 $L^1(\Omega)$ 中的有界序列使得对每个可测集 $A\subset\Omega$,
$\int_A f_n$ 收敛于有限极限 $\ell(A)$.证明存在某个 $f\in L^1(\Omega)$ 使得
$f_n\rightharpoonup f$ 弱收敛于 $\sigma(L^1,L^\infty)$.

5. 设 $g:\R\to\R$ 为满足对一切 $t\in\R$ 有 $\abs{g(t)}\le C$ 的连续单调递增函数.
令 $G(t)=\int_0^tg(s)ds$.设 $(f_n)$ 为 $L^1(\Omega)$ 中的序列使得
$f_n\rightharpoonup f$ 弱收敛于 $\sigma(L^1,L^\infty)$ 且
$\limsup\int_\Omega G(f_n)\le\int_\Omega G(f)$.证明 $\norm{f_n-f}_1\to0$.\\
{[提示: 参看问题 18.]}

\runin{-- D --}

在本部分中假设 $\abs{\Omega}<\infty$.设 $\mathcal{F}\subset L^1(\Omega)$.

1. 设 $G:[0,+\infty)\to[0,+\infty)$ 为满足 $\lim_{t\to+\infty}G(t)/t=+\infty$
的连续函数.假设存在常数 $C$ 使得对一切 $f\in\mathcal{F}$, $\int_\Omega G(\abs{f})\le C$.
证明 $\mathcal{F}$ 是等度可积的.

2. 反之, 假设 $\mathcal{F}$ 是等度可积的.证明存在凸单调递增函数
$G:[0,+\infty)\to[0,+\infty)$ 满足 $\lim_{t\to+\infty}G(t)/t=+\infty$, 使得对某个
常数 $C$, 对一切 $f\in\mathcal{F}$ 有 $\int_\Omega G(\abs{f})\le C$.\\
{[提示: 使用分布函数; 见问题 21.]}
\end{problem}

\begin{problem}[第 1, 3, 4 章]
\index{Radonceduo@Radon 测度 (Radon measure)}
\label{prob:24}
\runin{Radon 测度}
设 $K$ 为带有距离 $d$ 的紧度量空间, 设 $E=C(K)$ 配以通常的范数
$\norm{f}=\max_{x\in K}\abs{f(x)}$.对偶空间 $E^\star$ 记作 $\mathcal{M}(K)$,
称为 $K$ 上的 \emph{Radon 测度空间}.空间 $\mathcal{M}(K)$ 配以对偶范数, 记作
$\norm{\cdot}_{\mathcal{M}}$ 或简记为 $\norm{\cdot}$.本问题的目的是介绍
$\mathcal{M}(K)$ 的若干性质.

\runin{-- A --}

这里我们证明 $C(K)$ 可分.给定 $\delta>0$, 设 $\bigcup_{j\in J}B(a_j,\delta/2)$ 为
$K$ 的有限覆盖.令
\begin{align*}
q_j(x)&=\max\{0,\delta-d(x,a_j)\},\quad j\in J,\ x\in K,\\
q(x)&=\sum_{j\in J}q_j(x).
\end{align*}

1. 验证函数 $(q_j)_{j\in J}$ 与 $q$ 在 $K$ 上连续.证明对一切 $x\in K$,
$q(x)\ge\delta/2$.

2. 令 $\theta_j(x)=\dfrac{q_j(x)}{q(x)}$ ($j\in J,x\in K$).证明函数
$(\theta_j)_{j\in J}$ 在 $K$ 上连续, $[\theta_j(x)\ne0]\iff[d(x,a_j)<\delta]$,
且对一切 $x\in K$, $\sum_{j\in J}\theta_j(x)=1$.函数族 $(\theta_j)_{j\in J}$
称为 (从属于开覆盖 $\bigcup_{j\in J}B(a_j,2\delta)$ 的) \emph{单位分解}.

3. 给定 $f\in C(K)$, 令 $\tilde f(x)=\sum_{j\in J}f(a_j)\theta_j(x)$.证明
\begin{equation*}
\norm{f-\tilde f}\le\sup_{\substack{x,y\in K\\d(x,y)<\delta}}\abs{f(x)-f(y)}.
\end{equation*}

4. 选取 $\delta=1/n$, 上述构造给出一个有限集 $J_n$, 点的有限族 $(a_j)_{j\in J_n}$,
以及函数的有限族 $(\theta_j)_{j\in J_n}$.证明由函数 $(\theta_j)$ ($j\in J_n,n=1,2,3,\dots$)
张成的向量空间在 $C(K)$ 中稠密.

5. 由此推出 $C(K)$ 可分.

\runin{-- B --}

在本部分中假设 $K=\overline\Omega$, 其中 $\Omega$ 是 $\R^N$ 中的有界开集.通过
由对一切 $u\in L^1(\Omega)$ 与一切 $f\in C(\overline\Omega)$ 由
$\dual{Tu}{f}=\int_\Omega uf$ 定义的嵌入 $T:L^1(\Omega)\to\mathcal{M}(\overline\Omega)$,
把 $L^1(\Omega)$ 视为 $\mathcal{M}(\overline\Omega)$ 的子空间是方便的.

1. 验证对一切 $u\in L^1(\Omega)$, $\norm{Tu}_{\mathcal{M}}=\norm{u}_{L^1}$.\\
{[提示: 使用习题 4.26.]}

2. 设 $(v_n)$ 为 $L^1(\Omega)$ 中的有界序列.证明存在子序列 $(v_{n_k})$ 与某个
$\mu\in\mathcal{M}(\overline\Omega)$ 使得在 $\mathcal{M}(\overline\Omega)$ 中
$v_{n_k}\overset{*}{\rightharpoonup}\mu$ 且
$\norm{\mu}_{\mathcal{M}}\le\liminf_{k\to\infty}\norm{v_{n_k}}_{L^1}$.\\
{[提示: 使用推论  \ref{cor:3.30}.]}

现在的目标是证明: 给定任何 $\mu_0\in\mathcal{M}(\overline\Omega)$, 存在 $C_c^\infty(\Omega)$
中的序列 $(u_n)$ 使得
\begin{equation}
\int_\Omega u_nf\to\dual{\mu_0}{f}\quad\forall f\in C(\overline\Omega)
\label{eq:prob24.1}
\end{equation}
且
\begin{equation}
\norm{u_n}_{L^1}=\norm{\mu_0}_{\mathcal{M}}\quad\forall n.
\label{eq:prob24.2}
\end{equation}
不失一般性可以假设 $\norm{\mu_0}_{\mathcal{M}}=1$.令
$A=\{u\in C_c^\infty(\Omega);\ \norm{u}_{L^1}\le1\}$.

3. 证明 $\mu_0\in\overline A^{\sigma(E^\star,E)}$.\\
{[提示: 在 $E^\star$ 中对弱$^\star$拓扑 $\sigma(E^\star,E)$ 应用 Hahn--Banach 定理;
见问题 9.然后使用推论  \ref{cor:4.23}.]}

4. 由此推出存在 $A$ 中的序列 $(v_n)$ 使得在 $\sigma(E^\star,E)$ 中
$v_n\overset{*}{\rightharpoonup}\mu_0$.验证 $\lim_{n\to\infty}\norm{v_n}_{L^1}=1$.

5. 得出序列 $u_n=v_n/\norm{v_n}_{L^1}$ 满足 \eqref{eq:prob24.1} 与 \eqref{eq:prob24.2}.

若对满足 $f\ge0$ 的一切 $f\in C(\overline\Omega)$ 有 $\dual{\mu}{f}\ge0$, 则称
$\mu\ge0$.

6. 验证若 $\mu\ge0$, 则 $\dual{\mu}{1}=\norm{\mu}$, 其中 $1$ 表示恒等于 $1$ 的函数.

7. 假设 $\mu_0\in\mathcal{M}(\overline\Omega)$ 满足 $\mu_0\ge0$ 且 $\norm{\mu_0}=1$
(这样的测度称为\emph{概率测度}).在 $C_c^\infty(\Omega)$ 中构造满足 \eqref{eq:prob24.1}, \eqref{eq:prob24.2} 且此外对一切 $n$ 与一切 $x\in\Omega$ 满足 $u_n(x)\ge0$ 的序列 $(u_n)$.

8. 计算 $\norm{u+\delta_a}_{\mathcal{M}}$, 其中 $u\in L^1$, $\delta_a$ ($a\in\Omega$)
在下文定义.

\runin{-- C --}

现在回到一般框架, 用 $\delta_a$ 记点 $a\in K$ 处的 Dirac 质量, 即由
$\dual{\delta_a}{f}=f(a)$ ($\forall f\in C(K)$) 定义的测度.令
\begin{equation*}
\mathcal{D}=\left\{\mu=\sum_{j\in J}\alpha_j\delta_{a_j};\ J\ \text{有限},\ \alpha_j\in\R,\ \text{且点}\ a_j\ \text{互不相同}\right\}.
\end{equation*}

1. 证明若 $\mu\in\mathcal{D}$, 则 $\norm{\mu}=\sum_{j\in J}\abs{\alpha_j}$,
且 $[\mu\ge0]\iff[\alpha_j\ge0\ \forall j]$.令 $\mathcal{D}_1=\{\mu\in\mathcal{D};\ \norm{\mu}\le1\}$.

2. 证明满足 $\norm{\mu_0}\le1$ 的任何测度 $\mu_0\in\mathcal{M}(K)$ 属于
$\overline{\mathcal{D}_1}^{\sigma(E^\star,E)}$.\\
{[提示: 使用与第 B 部分相同的技巧.]}

3. 由此推出对给定的任何测度 $\mu_0\in\mathcal{M}(\overline\Omega)$, 存在 $\mathcal{D}$
中的序列 $(\nu_n)$ 使得 $\nu_n\overset{*}{\rightharpoonup}\mu_0$ 且对一切 $n$,
$\norm{\nu_n}=\norm{\mu_0}$.

4. 设 $\mu_0$ 为概率测度.证明存在 $\mathcal{D}$ 中的概率测度序列 $(\nu_n)$
使得 $\nu_n\overset{*}{\rightharpoonup}\mu_0$.

\runin{-- D --}

本部分的目的是证明每个 $\mu\in\mathcal{M}(K)$ 允许唯一的分解 $\mu=\mu_1-\mu_2$,
其中 $\mu_1,\mu_2\in\mathcal{M}(K)$, $\mu_1,\mu_2\ge0$, 且 $\norm{\mu_1}+\norm{\mu_2}=\norm{\mu}$
(测度 $\mu_1$ 与 $\mu_2$ 常记作 $\mu^+$ 与 $\mu^-$).

给定满足 $f\ge0$ 的 $f\in C(K)$, 令
\begin{equation*}
L(f)=\sup\{\dual{\mu}{g};\ g\in C(K)\ \text{且在}\ K\ \text{上}\ 0\le g\le f\}.
\end{equation*}

1. 验证 $0\le L(f)\le\norm{\mu}\norm{f}$, 对一切 $\lambda\ge0$ 有 $L(\lambda f)=\lambda L(f)$,
且对满足 $f_1\ge0$ 与 $f_2\ge0$ 的一切 $f_1,f_2\in C(K)$, $L(f_1+f_2)=L(f_1)+L(f_2)$.

给定任何 $f\in C(K)$, 令 $\mu_1(f)=L(f^+)-L(f^-)$, 其中 $f^+=\max\{f,0\}$ 与
$f^-=\max\{-f,0\}$.

2. 证明映射 $f\mapsto\mu_1(f)$ 在 $C(K)$ 上是线性的且对一切 $f\in C(K)$ 满足
$\abs{\mu_1(f)}\le\norm{\mu}\norm{f}$, 从而 $\mu_1\in\mathcal{M}(K)$.验证 $\mu_1\ge0$.

3. 令 $\mu_2=\mu_1-\mu$ 并验证 $\mu_2\ge0$.证明 $\norm{\mu}=\norm{\mu_1}+\norm{\mu_2}$.

4. 设 $\nu\in\mathcal{M}(K)$ 满足 $\nu\ge0$ 且 $\nu\ge\mu$.证明 $\nu\ge\mu_1$.
类似地若 $\nu\in\mathcal{M}(K)$ 且 $\nu\ge-\mu$, 证明 $\nu\ge\mu_2$.由此推出分解的唯一性.

\runin{-- E --}

证明当把空间 $E=C(\overline\Omega)$ 换成子空间
$E_0=\{f\in C(\overline\Omega);\ f=0\ \text{在}\ \partial\Omega\ \text{上}\}$ 时,
上述所有结果 (除第 B6 问外) 仍然成立.$E_0$ 的对偶常记作 $\mathcal{M}_0(\Omega)$
或 $\mathcal{M}(\Omega)$.

\runin{-- F --}

重访 Dunford--Pettis 定理.设 $(f_n)$ 为 $L^1(\Omega)$ 中的序列.回忆若 $(f_n)$
满足对一切 $\varepsilon>0$, 存在 $\delta>0$ 使得对满足 $\abs{A}<\delta$ 的一切可测集
$A\subset\Omega$ 及一切 $n$, $\int_A\abs{f_n}<\varepsilon$, 则称 $(f_n)$ 是等度可积的.
目标是证明每个等度可积序列 $(f_n)$ 允许子序列 $(f_{n_j})$ 弱收敛于某个 $f\in L^1(\Omega)$.

1. 证明 $(f_n)$ 在 $L^1(\Omega)$ 中有界.

2. 验证对一切 $n,k$, $\int_\Omega\abs{f_n-T_kf_n}\le\int_{[\abs{f_n}>k]}\abs{f_n}$.

3. 由此推出对一切 $\varepsilon>0$, 存在 $k>0$ 使得对一切 $n$, $\int_\Omega\abs{f_n-T_kf_n}\le\varepsilon$.

抽取子序列 (仍记作 $(f_n)$), 可以假设在 $\mathcal{M}(\overline\Omega)$ 中
$f_n\overset{*}{\rightharpoonup}\mu$ 弱$^\star$收敛于某个测度 $\mu$.

4. 证明对一切 $\varepsilon>0$, 存在 $g=g_\varepsilon\in L^\infty(\Omega)$ 使得
$\norm{\mu-g_\varepsilon}_{\mathcal{M}}\le\varepsilon$.

5. 由此推出 $\mu\in L^1(\Omega)$.

6. 证明 $f_n\rightharpoonup\mu$ 弱收敛于 $\sigma(L^1,L^\infty)$.
\end{problem}

\begin{problem}[第 1, 5 章]
\label{prob:25}
设 $H$ 为 Hilbert 空间, $C\subset H$ 为以 $0$ 为顶点的凸锥, 即 $0\in C$, 且对一切
$\lambda,\mu>0$ 与一切 $u,v\in C$, $\lambda u+\mu v\in C$.我们假设 $C$ 非空, 开,
且 $C\ne H$.

验证 $0\notin C$ 且 $0\in\overline C$.考虑集合
\begin{equation*}
\Gamma=\{u\in H;\ (u,v)\le0\ \forall v\in C\}.
\end{equation*}

1. 验证 $\Gamma$ 是以 $0$ 为顶点的凸锥, $\Gamma$ 闭, 且 $0\in\Gamma$.证明
$C=\{v\in H;\ (u,v)<0\ \forall u\in\Gamma\setminus\{0\}\}$, 并由此推出 $\Gamma$
不退化为 $\{0\}$.\\
{[提示: 使用 Hahn--Banach 定理.]}

2. 设 $\omega\in C$ 固定, 考虑集合 $K=\{u\in\Gamma;\ (u,\omega)=-1\}$.证明 $K$
是非空, 有界, 闭的凸集, 使得 $0\notin K$ 且
$\Gamma\setminus\{0\}=\bigcup_{\lambda>0}\lambda K$.画出示意图.\\
{[提示: 考虑以 $\omega$ 为中心, 包含在 $C$ 内的半径为 $\rho>0$ 的球.]}

3. 设 $a=P_K0$.证明 $a\in(-C)\cap\Gamma$.

4. 用简捷的方法直接证明 $(-C)\cap\Gamma\ne\emptyset$.

5. 设 $D\subset H$ 为非空开凸集, $x_0\notin D$.证明存在某个 $w_0\in D$ 使得
对一切 $w\in D$, $(w_0-x_0,w-x_0)>0$.给出几何解释.\\
{[提示: 考虑集合 $C=\bigcup_{\mu>0}\mu(D-x_0)$.]}
\end{problem}
\begin{problem}[第 1, 5 章]
\index{MoreauProxyingshe@Moreau Prox 映射 (Moreau prox map)}
\label{prob:26}
\runin{Moreau 意义下的 Prox 映射}
设 $H$ 为 Hilbert 空间, $\varphi:H\to(-\infty,+\infty]$ 为满足 $\varphi\not\equiv+\infty$
的凸下半连续函数.

1. 证明对每个 $f\in H$, 存在某个 $u\in D(\varphi)$ 使得
\begin{equation}
\frac12\abs{f-u}^2+\varphi(u)=\inf_{v\in H}\left\{\frac12\abs{f-v}^2+\varphi(v)\right\}\equiv I.
\tag{P}
\label{eq:prob26.P}
\end{equation}
[提示: 先验证 $I>-\infty$.然后使用 Cauchy 序列论证或 $H$ 自反的事实.]

2. 验证 $u$ 满足 \eqref{eq:prob26.P} 当且仅当
\begin{equation}
u\in D(\varphi)\quad\text{且}\quad(u,v-u)+\varphi(v)-\varphi(u)\ge(f,v-u)\quad\forall v\in D(\varphi).
\tag{Q}
\label{eq:prob26.Q}
\end{equation}

3. 证明若 $u$ 与 $\bar u$ 分别是对应于 $f$ 与 $\bar f$ 的 \eqref{eq:prob26.P} 的解,
则 $\abs{u-\bar u}\le\abs{f-\bar f}$.由此推出 \eqref{eq:prob26.P} 的解的唯一性.

4. 考察 $\varphi=I_K$ 为闭凸集 $K$ 的示性函数的特殊情形.

5. 设 $\varphi^\star$ 为 $\varphi$ 的共轭函数, 考虑问题
\begin{equation}
\frac12\abs{f-u^\star}^2+\varphi^\star(u^\star)=\inf_{v\in H}\left\{\frac12\abs{f-v}^2+\varphi^\star(v)\right\}=I^\star.
\tag{P$^\star$}
\label{eq:prob26.Pstar}
\end{equation}
证明 \eqref{eq:prob26.P} 的解 $u$ 与 \eqref{eq:prob26.Pstar} 的解 $u^\star$ 满足
\begin{equation*}
u+u^\star=f\quad\text{且}\quad I+I^\star=\frac12\abs{f}^2.
\end{equation*}

6. 给定 $f\in H$ 与 $\lambda>0$, 设 $u_\lambda$ 表示问题
\begin{equation}
\frac12\abs{f-u_\lambda}^2+\lambda\varphi(u_\lambda)=\inf_{v\in H}\left\{\frac12\abs{f-v}^2+\lambda\varphi(v)\right\}
\tag{P$_\lambda$}
\label{eq:prob26.Plambda}
\end{equation}
的解.证明 $\lim_{\lambda\to0}u_\lambda=P_{\overline{D(\varphi)}}f$ ($=f$ 在
$\overline{D(\varphi)}$ 上的投影).\\
{[提示: 从弱收敛开始或使用习题 5.3.]}

7. 设 $K=\{v\in D(\varphi);\ \varphi(v)=\inf_H\varphi\}$ 并假设 $K\ne\emptyset$.
验证 $K$ 是闭凸集并证明 $\lim_{\lambda\to+\infty}u_\lambda=P_Kf$.当 $K=\emptyset$
时, 当 $\lambda\to+\infty$ 时 $(u_\lambda)$ 发生什么?

8. 证明 $\lim_{\lambda\to+\infty}\dfrac1\lambda u_\lambda=-P_{\overline{D(\varphi^\star)}}0$.\\
{[提示: 从 $f=0$ 的情形开始并应用第 5 与第 6 问.]}
\end{problem}

\begin{problem}[第~
\ref{chap:5}~章]
\label{prob:27}
\runin{交替投影}
设 $H$ 为 Hilbert 空间, $K\subset H$ 为非空闭凸集.验证对一切 $u,v\in H$,
\begin{equation*}
\abs{P_Ku-P_Kv}^2\le(P_Ku-P_Kv,u-v)\le\abs{u-v}^2.
\end{equation*}

设 $K_1\subset H$ 与 $K_2\subset H$ 为两个非空闭凸集.令 $P_1=P_{K_1}$ 与
$P_2=P_{K_2}$.给定 $u\in H$, 由归纳法定义序列 $(u_n)$ 如下:
\begin{equation*}
u_0=u,\quad u_1=P_1u_0,\quad u_2=P_2u_1,\quad\dots,\quad u_{2n-1}=P_1u_{2n-2},\quad u_{2n}=P_2u_{2n-1},\quad\dots
\end{equation*}

\runin{-- A --}

本部分的目的是证明序列 $(u_{2n}-u_{2n-1})$ 收敛于 $P_K0$, 其中 $K=K_2-K_1$
(注意 $K$ 是凸集, 为什么?).

1. 给定 $v\in H$, 考虑以 $v_0=v$ 开始由上述相同迭代定义的序列 $(v_n)$.验证
\begin{align*}
\abs{u_{2n}-v_{2n}}^2&\le(u_{2n}-v_{2n},u_{2n-1}-v_{2n-1})\le\abs{u_{2n-1}-v_{2n-1}}^2,\\
\abs{u_{2n+1}-v_{2n+1}}^2&\le(u_{2n+1}-v_{2n+1},u_{2n}-v_{2n})\le\abs{u_{2n}-v_{2n}}^2.
\end{align*}

2. 由此推出序列 $(\abs{u_n-v_n})$ 单调不增, 从而收敛于一个极限, 记作 $\ell$.
证明
\begin{equation*}
\lim_{n\to\infty}\abs{(u_{2n}-v_{2n})-(u_{2n-1}-v_{2n-1})}^2=0.
\end{equation*}

3. 验证序列 $(\abs{u_{2n}-u_{2n-1}})$ 单调不增.令
$d=\dist(K_1,K_2)=\inf\{\abs{a_1-a_2};\ a_1\in K_1,a_2\in K_2\}$.我们断言
$\lim_{n\to\infty}\abs{u_{2n}-u_{2n-1}}=d$.

4. 给定 $\varepsilon>0$, 选取 $v\in K_2$ 使得 $\dist(v,K_1)\le d+\varepsilon$.
证明对一切 $n$, $\abs{v_{2n}-v_{2n-1}}\le d+\varepsilon$.

5. 由此推出 $\lim_{n\to\infty}\abs{u_{2n}-u_{2n-1}}=d$.令 $z=P_K0$.

6. 验证 $\abs{z}=d$ 且对一切 $w\in K_2-K_1$, $\abs{z}^2\le(z,w)$.

7. 证明序列 $(u_{2n}-u_{2n-1})$ 收敛于 $z$.\\
{[提示: 使用上述结果估计 $\abs{z-(u_{2n}-u_{2n-1})}^2$.]}

8. 给出几何解释.

\runin{-- B --}

在整个本题的其余部分中, 我们假设 $z=P_K0\in K_2-K_1$ (例如若集合 $K_1$ 或 $K_2$
之一有界, 这个假设就成立, 为什么?).

我们断言存在 $a_1\in K_1$ 与 $a_2\in K_2$ 满足 $a_2-a_1=z$, 使得
$u_{2n}\rightharpoonup a_2$ 且 $u_{2n-1}\rightharpoonup a_1$ 弱收敛.注意 $a_1$
与 $a_2$ 可能依赖于初值 $u_0=u$ 的选取.画出示意图.

1. 考虑配以其自然标量积的 Hilbert 空间 $\mathcal{H}=H\times H$.令
$\mathcal{K}=\{[b_1,b_2]\in\mathcal{H};\ b_1\in K_1,b_2\in K_2\ \text{且}\ b_2-b_1=z\}$.
验证 $\mathcal{K}$ 是非空闭凸集.

2. 设 $b=[b_1,b_2]\in\mathcal{K}$.确定对应于 $v_0=b_1$ 的序列 $(v_n)$.
由此推出序列 $(\abs{u_{2n-1}-b_1})$ 与 $(\abs{u_{2n}-b_2})$ 单调不增.

3. 令 $x_n=[u_{2n-1},u_{2n}]$, 证明序列 $(x_n)$ 满足下述性质: 对弱收敛于某个
元素 $\bar x\in\mathcal{H}$ 的每个子序列 $(x_{n_k})$, 有 $\bar x\in\mathcal{K}$.

4. 应用 Opial 引理 (见习题 5.25 第 3 问) 并得出结论.
\end{problem}

\begin{problem}[第~
\ref{chap:5}~章]
\label{prob:28}
\runin{投影与正交投影}
设 $H$ 为 Hilbert 空间.满足 $P^2=P$ 的算子 $P\in\LL(H)$ 称为\emph{投影算子}.
验证投影算子满足下述性质:
(a) $I-P$ 是投影算子,
(b) $N(I-P)=R(P)$ 且 $N(P)=R(I-P)$,
(c) $N(P)\cap N(I-P)=\{0\}$,
(d) $H=N(P)+N(I-P)$.

\runin{-- A --}

若存在闭线性子空间 $M$ 使得 $P=P_M$ (其中 $P_M$ 在推论  \ref{cor:5.4} 中定义),
则称算子 $P\in\LL(H)$ 为\emph{正交投影算子}.验证每个正交投影算子是投影算子.

1. 给定投影算子 $P$, 证明下述性质等价:
(a) $P$ 是正交投影算子,
(b) $P^\star=P$,
(c) $\norm{P}\le1$,
(d) $N(P)\perp N(I-P)$,
其中记号 $X\perp Y$ 意味着对一切 $x\in X$ 及一切 $y\in Y$, $(x,y)=0$.

2. 设 $T\in\LL(H)$ 为满足 $T^\star=T$ 且 $T^2=I$ 的算子.证明 $P=\frac12(I-T)$
是正交投影算子.证明逆命题.额外假设对一切 $u\in H$ 有 $(Tu,u)\ge0$, 证明 $T=I$.

\runin{-- B --}

在整个本部分中, $M$ 与 $N$ 表示 $H$ 的两个闭线性子空间.令 $P=P_M$ 与 $Q=P_N$.

1. 证明下述性质等价:
(a) $PQ=QP$,
(b) $PQ$ 是投影算子,
(c) $QP$ 是投影算子.
在这时, 验证: (i) $PQ$ 是到 $M\cap N$ 上的正交投影算子, (ii) $(P+Q-PQ)$ 是到
$M+N$ 上的正交投影算子.

2. 证明下述性质等价:
(a) $M\perp N$,
(b) $PQ=0$,
(c) $QP=0$,
(d) 对一切 $u\in H$, $\abs{Pu}^2+\abs{Qu}^2\le\abs{u}^2$,
(e) 对一切 $u\in H$, $\abs{Pu}\le\abs{u-Qu}$,
(f) 对一切 $u\in H$, $\abs{Qu}\le\abs{u-Pu}$,
(g) $P+Q$ 是投影算子.
在这时, 验证 $(P+Q)$ 是到 $M+N$ 上的正交投影算子 (注意 $M+N$ 是闭的; 为什么?).

3. 证明下述性质等价:
(a) $M\subset N$,
(b) $PQ=P$,
(c) $QP=P$,
(d) 对一切 $u\in H$, $\abs{Pu}\le\abs{Qu}$,
(e) $Q-P$ 是投影算子.
在这时, 验证 $Q-P$ 是到 $M^\perp\cap N$ 上的正交投影算子.
\end{problem}

\begin{problem}[第~
\ref{chap:5}~章]
\label{prob:29}
\runin{非线性压缩映射的迭代. Opial 与 Baillon 遍历定理}
设 $H$ 为 Hilbert 空间, 设 $T:H\to H$ 为非线性压缩映射, 即对一切 $u,v\in H$,
$\abs{Tu-Tv}\le\abs{u-v}$.我们假设不动点集 $K=\{u\in H;\ Tu=u\}$ 非空.验证
$K$ 是闭凸集.给定 $f\in H$, 令
\begin{align*}
\sigma_n&=\frac1n(f+Tf+T^2f+\dots+T^{n-1}f),\\
\mu_n&=\left(\frac{I+T}{2}\right)^nf.
\end{align*}
本问题的目标是证明下述结论:
\begin{enumerate}[label=(\Alph*), itemsep=0.1em]
	\item 序列 $(\sigma_n)$ 与 $(\mu_n)$ 各自弱收敛于 $T$ 的某个不动点.
	\item 若此外 $T$ 是奇的, 即对一切 $v\in H$, $T(-v)=-Tv$, 则 $(\sigma_n)$ 与 $(\mu_n)$ 强收敛.
\end{enumerate}
在开始本题之前, 建议先做习题 5.22 与 5.25.在 $T$ 为线性的特殊情形下, 另见习题 5.21.

\runin{-- A --}

令 $u_n=T^nf$.

1. 验证对每个 $v\in K$, 序列 $(\abs{u_n-v})$ 单调不增.由此推出序列 $(\sigma_n)$
与 $(T\sigma_n)$ 有界.

2. 证明对一切 $n\ge1$, $\abs{\sigma_n-T\sigma_n}\le\frac{1}{\sqrt{n}}\abs{f-T\sigma_n}$.\\
{[提示: 注意 $\abs{T\sigma_n-Tu_i}^2\le\abs{\sigma_n-u_i}^2$, 并对 $0\le i\le n-1$ 将这些不等式相加.]}

3. 由此推出序列 $(\sigma_n)$ 满足习题 5.25 的性质 (P).得出结论: $\sigma_n\rightharpoonup\sigma$
弱收敛, 其中 $\sigma\in K$.

令 $S=\frac12(I+T)$.

4. 证明对一切 $u,v\in H$,
\begin{equation*}
\abs{(u-Su)-(v-Sv)}^2+\abs{Su-Sv}^2\le\abs{u-v}^2.
\end{equation*}

5. 由此推出对每个 $v\in K$, $\sum_{n=0}^\infty\abs{\mu_n-\mu_{n+1}}^2\le\abs{f-v}^2$,
从而对一切 $n$, $\abs{\mu_n-S\mu_n}\le\frac{1}{\sqrt{n+1}}\abs{f-v}$.

6. 得出结论: $\mu_n\rightharpoonup\mu$ 弱收敛, 其中 $\mu\in K$.

\runin{-- B --}

在整个本题的其余部分中, 我们假设 $T$ 是奇的, 即对一切 $v\in H$, $T(-v)=-Tv$.

1. 证明对每个整数 $p$, 对一切 $u,v\in H$,
\begin{equation*}
2\abs{(u,v)-(T^pu,T^pv)}\le\abs{u}^2+\abs{v}^2-\abs{T^pu}^2-\abs{T^pv}^2.
\end{equation*}
[提示: 从对一切 $u,v\in H$ 有 $\abs{T^pu-T^pv}^2\le\abs{u-v}^2$ 的不等式开始.]

2. 由此推出对每个固定的整数 $i\ge0$, $\ell(i)=\lim_{n\to\infty}(u_n,u_{n+i})$ 存在.
证明这个收敛关于 $i$ 一致成立, 即对一切 $i$ 与一切 $n$,
\begin{equation}
\abs{(u_n,u_{n+i})-\ell(i)}\le\varepsilon_n,
\label{eq:prob29.1}
\end{equation}
其中 $\lim_{n\to\infty}\varepsilon_n=0$.

3. 同理证明对每个固定的整数 $i\ge0$, $m(i)=\lim_{n\to\infty}(\mu_n,\mu_{n+i})$ 存在.
证明 $m(0)=m(1)=m(2)=\dots$.\\
{[提示: 使用第 A5 问的结果.]}

4. 由此推出 $\mu_n\to\mu$ 强收敛.

我们现在断言 $\sigma_n\to\sigma$ 强收敛.

5. 令 $X_p=\frac1p\sum_{i=0}^{p-1}\ell(i)$.证明对一切 $n$ 与一切 $p$,
\begin{equation*}
\abs{(u_n,\sigma_{n+p})-X_p}\le\varepsilon_n+\frac{2n}{p}\abs{f}^2.
\end{equation*}
[提示: 使用 \eqref{eq:prob29.1}.]

6. 由此推出: (i) $X=\lim_{p\to\infty}X_p$ 存在, (ii) 对一切 $n$, $\abs{(u_n,\sigma)-X}\le\varepsilon_n$,
(iii) $\abs{\sigma}^2=X$.

7. 证明 $\abs{\sigma_n}^2\le\frac{2}{n^2}\sum_{i=0}^{n-1}(n-i)\ell(i)+\frac2n\sum_{i=0}^{n-1}\varepsilon_i$.

8. 由此推出 $\limsup_{n\to\infty}\abs{\sigma_n}^2\le X$ 并得出结论.
\end{problem}

\begin{problem}[第 3, 5 章]
\label{prob:30}
\runin{Stampacchia 定理的变体. von Neumann 极小化极大定理}
设 $H$ 为 Hilbert 空间.

\runin{-- A --}

设 $a(u,v):H\times H\to\R$ 为满足对一切 $v\in H$ 有 $a(v,v)\ge0$ 的连续双线性形式.
设 $K\subset H$ 为非空闭凸集.设 $f\in H$.假设存在某个 $v_0\in K$ 使得集合
$\{u\in K;\ a(u,v_0-u)\ge(f,v_0-u)\}$ 是有界的.

1. 证明存在某个 $u\in K$ 使得对一切 $v\in K$, $a(u,v-u)\ge(f,v-u)$.\\
{[提示: 令 $f_\varepsilon=f+\varepsilon v_0$, 并考虑双线性形式
$a_\varepsilon(u,v)=a(u,v)+\varepsilon(u,v)$ ($\varepsilon>0$).然后利用习题 5.14
令 $\varepsilon\to0$ 取极限.]}

2. 恢复 Stampacchia 定理.

3. 在 $K$ 有界且对一切 $u,v\in H$ 有 $a(u,v)=0$ 的情形下给出几何解释.

\runin{-- B --}

设 $b(u,v):H\times H\to\R$ 为连续且强制的双线性形式.设 $\varphi:H\to(-\infty,+\infty]$
为满足 $\varphi\not\equiv+\infty$ 的凸下半连续函数.

1. 证明存在唯一的 $u\in D(\varphi)$ 使得对一切 $v\in D(\varphi)$,
\begin{equation*}
b(u,v-u)+\varphi(v)-\varphi(u)\ge0.
\end{equation*}
[提示: 在空间 $H\times\R$ 中应用第 A1 问的结果, 取 $K=\operatorname{epi}\varphi$,
$f=[0,-1]$, $a(U,V)=b(u,v)$, 其中 $U=[u,\lambda],V=[v,\mu]$.注意 $a$ 不是强制的.]

2. 恢复 Stampacchia 定理.

\runin{-- C --}

设 $H_1$ 与 $H_2$ 为两个 Hilbert 空间, 设 $A\subset H_1,B\subset H_2$ 为两个非空, 有界, 闭的凸集.

1. 设 $F(\lambda,\mu):H_1\times H_2\to\R$ 为连续双线性形式.证明存在 $\bar\lambda\in A$
与 $\bar\mu\in B$ 使得对一切 $\lambda\in A$ 与一切 $\mu\in B$,
\begin{equation}
F(\bar\lambda,\mu)\le F(\bar\lambda,\bar\mu)\le F(\lambda,\bar\mu).
\label{eq:prob30.1}
\end{equation}
[提示: 在 $H=H_1\times H_2$, $\mathcal{K}=A\times B$ 上应用第 A1 问, 取
$a(u,v)=F(\bar\lambda,\mu)-F(\lambda,\bar\mu)$, 其中 $u=[\bar\lambda,\bar\mu],v=[\lambda,\mu]$.]

2. 由此推出
\begin{equation}
\min_{\lambda\in A}\max_{\mu\in B}F(\lambda,\mu)=\max_{\mu\in B}\min_{\lambda\in A}F(\lambda,\mu).
\label{eq:prob30.2}
\end{equation}
注意所有的极小与极大都可以达到 (为什么?).\\
{[提示: 验证没有任何附加假设下就有 $\max\min\le\min\max$; 使用 \eqref{eq:prob30.1}
证明反向不等式.]}

3. 证明 \eqref{eq:prob30.2} 蕴含满足 \eqref{eq:prob30.1} 的某个 $\bar\lambda\in A$
与 $\bar\mu\in B$ 的存在性.

\runin{-- D --}

设 $E$ 与 $F$ 为两个自反 Banach 空间; 设 $A\subset E$ 与 $B\subset F$ 为两个非空, 有界, 闭的凸集.设 $K:E\times F\to\R$ 为满足下述假设的函数:
\begin{enumerate}[label=(\alph*), itemsep=0.1em]
	\item 对每个固定的 $v\in B$, 函数 $u\mapsto K(u,v)$ 是凸的且下半连续.
	\item 对每个固定的 $u\in A$, 函数 $v\mapsto K(u,v)$ 是凹的且上半连续, 即函数 $v\mapsto-K(u,v)$ 是凸的且下半连续.
\end{enumerate}
我们的目标是证明
\begin{equation*}
\min_{u\in A}\max_{v\in B}K(u,v)=\max_{v\in B}\min_{u\in A}K(u,v).
\end{equation*}
我们将用反证法并假设存在常数 $\gamma$ 使得
\begin{equation*}
\max_{v\in B}\min_{u\in A}K(u,v)<\gamma<\min_{u\in A}\max_{v\in B}K(u,v).
\end{equation*}

1. 对每个 $u\in A$, 令 $B_u=\{v\in B;\ K(u,v)\ge\gamma\}$; 对每个 $v\in B$, 令
$A_v=\{u\in A;\ K(u,v)\le\gamma\}$.验证 $\bigcap_{u\in A}B_u=\emptyset$ 且
$\bigcap_{v\in B}A_v=\emptyset$.

2. 选取 $u_1,\dots,u_n\in A$ 与 $v_1,\dots,v_m\in B$ 使得 $\bigcap_{i=1}^nB_{u_i}=\emptyset$
且 $\bigcap_{j=1}^mA_{v_j}=\emptyset$ (说明理由).在 $H_1=\R^n,H_2=\R^m$,
$A'=\left\{\lambda=(\lambda_1,\dots,\lambda_n);\ \lambda_i\ge0\ \forall i,\ \sum_{i=1}^n\lambda_i=1\right\}$,
$B'=\left\{\mu=(\mu_1,\dots,\mu_m);\ \mu_j\ge0\ \forall j,\ \sum_{j=1}^m\mu_j=1\right\}$,
以及 $F(\lambda,\mu)=\sum_{i,j}\lambda_i\mu_jK(u_i,v_j)$ 上应用第 C1 问的结果.
令 $\bar u=\sum_i\bar\lambda_iu_i$ 与 $\bar v=\sum_j\bar\mu_jv_j$.证明对一切
$k=1,\dots,n$ 与一切 $\ell=1,\dots,m$,
\begin{equation*}
K(\bar u,v_\ell)\le K(u_k,\bar v).
\end{equation*}

3. 验证另一方面有 $\min_k K(u_k,\bar v)<\gamma$ 以及 $\max_\ell K(\bar u,v_\ell)>\gamma$.\\
{[提示: 用反证法.]}

4. 得出结论.
\end{problem}
\begin{problem}[第 3, 5 章]
\label{prob:31}
\runin{单调算子. Minty--Browder 定理}
设 $E$ 为自反 Banach 空间.(非线性) 映射 $A:D(A)\subset E\to E^\star$ 若满足
\begin{equation*}
\dual{Au-Av}{u-v}\ge0\quad\forall u,v\in D(A),
\end{equation*}
则称它是\emph{单调的} (这里 $D(A)$ 表示 $E$ 的任意子集).

\runin{-- A --}

设 $A:D(A)\subset E\to E^\star$ 为单调映射, 设 $K\subset E$ 为非空, 有界, 闭的凸集.
我们的目标是证明存在某个 $u\in K$ 使得
\begin{equation*}
\dual{Av}{u-v}\ge0\quad\forall v\in D(A)\cap K.
\end{equation*}
为此, 对每个 $v\in D(A)\cap K$, 令 $K_v=\{u\in K;\ \dual{Av}{v-u}\ge0\}$.须证明
$\bigcap_{v\in D(A)\cap K}K_v\ne\emptyset$; 我们用反证法并假设
$\bigcap_{v\in D(A)\cap K}K_v=\emptyset$.

1. 验证 $K_v$ 是闭凸集.

2. 由此推出存在 $v_1,v_2,\dots,v_n\in D(A)\cap K$ 使得 $\bigcap_{i=1}^nK_{v_i}=\emptyset$.
令
\begin{equation*}
B=\left\{\lambda=(\lambda_1,\dots,\lambda_n);\ \lambda_i\ge0\ \forall i,\ \sum_{i=1}^n\lambda_i=1\right\},
\end{equation*}
并考虑由 $F(\lambda,\mu)=\sum_{i,j=1}^n\lambda_i\mu_j\dual{Av_j}{v_i-v_j}$ 定义的双线性形式
$F:\R^n\times\R^n\to\R$.

3. 验证对一切 $\lambda\in\R^n$, $F(\lambda,\lambda)\le0$.

4. 证明存在某个 $\bar\lambda\in B$ 使得对一切 $\mu\in B$, $F(\bar\lambda,\mu)\le0$.\\
{[提示: 应用问题 30 第 C1 问.]}

5. 令 $\bar u=\sum_{i=1}^n\bar\lambda_iv_i$, 证明对一切 $j=1,\dots,n$,
$\dual{Av_j}{\bar u-v_j}\le0$.

6. 得出结论.

\runin{-- B --}

在整个本题的其余部分中, 我们假设 $D(A)=E$, $A:E\to E^\star$ 单调, 且 $A$ 连续.

1. 设 $K\subset E$ 为非空, 有界, 闭的凸集.证明存在某个 $u\in K$ 使得对一切
$w\in K$, $\dual{Au}{w-u}\ge0$.\\
{[提示: 考虑 $v_t=(1-t)u+tw$, 其中 $t\in(0,1),w\in K$.]}

2. 设 $K$ 为包含 $0$ 的闭凸集 ($K$ 不必有界).假设集合 $\{u\in K;\ \dual{Au}{u}\le0\}$
是有界的.证明存在某个 $u\in K$ 使得对一切 $v\in K$, $\dual{Au}{v-u}\ge0$.\\
{[提示: 对充分大的 $R$, 把 B1 应用于集合 $K_R=\{v\in K;\ \norm{v}\le R\}$.]}

3. 假设 $\lim_{\norm{v}\to\infty}\dfrac{\dual{Av}{v}}{\norm{v}}=+\infty$.证明 $A$ 是满射.

4. 这里假设 $E$ 是与 $E^\star$ 等同的 Hilbert 空间.证明 $I+A$ 是由 $E$ 到自身上的双射.
\end{problem}

\begin{problem}[第~
\ref{chap:5}~章]
\label{prob:32}
\runin{压缩映射的延拓. Kirszbraun--Valentine 定理的 Schoenberg 方法}
设 $H$ 为 Hilbert 空间, $I$ 为有限指标集.

\runin{-- A --}

设 $(y_i)_{i\in I}$ 为 $H$ 中的元素, $(c_i)_{i\in I}$ 为 $\R$ 中的元素.令
\begin{equation*}
\phi(u)=\max_{i\in I}\{\abs{u-y_i}^2-c_i\},\quad u\in H,
\end{equation*}
以及 $J(u)=\{i\in I;\ \abs{u-y_i}^2-c_i=\phi(u)\}$.

1. 验证 $\inf_{u\in H}\phi(u)$ 由某个唯一的元素 $u_0\in H$ 达到.

2. 证明对一切 $v\in H$, $\max_{i\in J(u_0)}(v,u_0-y_i)\ge0$.

3. 由此推出
\begin{equation}
u_0\in\operatorname{conv}\left(\bigcup_{i\in J(u_0)}\{y_i\}\right).
\label{eq:prob32.1}
\end{equation}

4. 反之, 若 $u_0\in H$ 满足 \eqref{eq:prob32.1}, 证明 $\phi(u_0)=\inf_{u\in H}\phi(u)$.

5. 把这一结果推广到 $\phi(u)=\max_{i\in I}\{f_i(u)\}$ 的情形, 其中每个 $f_i:H\to\R$
是凸的 $C^1$ 函数.

\runin{-- B --}

设 $(x_i)_{i\in I}$ 与 $(y_i)_{i\in I}$ 为 $H$ 中的元素, 满足对一切 $i,j\in I$,
$\abs{y_i-y_j}\le\abs{x_i-x_j}$.我们断言: 给定任何 $p\in H$, 存在某个
$q\in\operatorname{conv}\left(\bigcup_{i\in I}\{y_i\}\right)$ 使得对一切 $i\in I$,
$\abs{q-y_i}\le\abs{p-x_i}$.

1. 令 $P=\{\lambda=(\lambda_i)_{i\in I};\ \lambda_i\ge0\ \forall i\ \text{且}\ \sum_{i\in I}\lambda_i=1\}$.
证明对每个 $p\in H$ 与每个 $\lambda\in P$,
\begin{equation*}
\sum_{j\in I}\lambda_j\abs{\sum_{i\in I}\lambda_iy_i-y_j}^2\le\sum_{j\in I}\lambda_j\abs{p-x_j}^2.
\end{equation*}
[提示: 验证 $\sum_{j\in I}\lambda_j\abs{\sum_{i\in I}\lambda_iy_i-y_j}^2=\frac12\sum_{i,j\in I}\lambda_i\lambda_j\abs{y_i-y_j}^2$.]

2. 考虑函数 $\phi(u)=\max_{i\in I}\{\abs{u-y_i}^2-\abs{p-x_i}^2\}$.设 $u_0\in H$
使得 $\phi(u_0)=\inf_{u\in H}\phi(u)$.证明 $\phi(u_0)\le0$.\\
{[提示: 应用第 A3 与 B1 问.]}

3. 得出结论.

\runin{-- C --}

1. 把第 B 部分的结果推广到 $I$ 为无限指标集的情形.

2. 设 $D\subset H$ 为 $H$ 的任意子集, 设 $S:D\to H$ 为压缩映射, 即对一切 $u,v\in D$,
$\abs{Su-Sv}\le\abs{u-v}$.证明存在定义在整个 $H$ 上的压缩映射 $T$, 延拓 $S$ 且满足
$T(H)\subset\overline{\operatorname{conv}S(D)}$.\\
{[提示: 使用 Zorn 引理与第 C1 问.]}
\end{problem}

\begin{problem}[第 4, 6 章]
\label{prob:33}
\runin{$L^p$ 中的乘法算子}
设 $\Omega$ 为测度空间 (有限或无限测度).令 $E=L^p(\Omega)$, 其中 $1<p<\infty$.
设 $a:\Omega\to\R$ 为可测函数.考虑由
\begin{equation*}
D(A)=\{u\in L^p(\Omega);\ au\in L^p(\Omega)\}\quad\text{以及}\quad Au=au
\end{equation*}
定义的无界线性算子 $A:D(A)\subset E\to E$.

1. 证明 $D(A)$ 在 $E$ 中稠密.\\
{[提示: 给定 $u\in E$, 考虑序列 $u_n(x)=(1+n^{-1}\abs{a(x)})^{-1}u(x)$.]}

2. 证明 $A$ 是闭的.

3. 证明 $D(A)=E$ 当且仅当 $a\in L^\infty(\Omega)$.\\
{[提示: 应用闭图像定理.]}

4. 确定 $N(A)$ 与 $N(A)^\perp$.

5. 确定 $D(A^\star)$, $A^\star$, $N(A^\star)$ 与 $N(A^\star)^\perp$.

6. 证明 $A$ 是满射当且仅当存在 $\alpha>0$ 使得在 $\Omega$ 上几乎处处 $\abs{a(x)}\ge\alpha$.\\
{[提示: 使用第 3 问.]}

在下文中假设 $a\in L^\infty(\Omega)$.

7. 确定 $A$ 的特征值与谱.验证 $\sigma(A)\subset[\operatorname{ess\,inf}_\Omega a,\operatorname{ess\,sup}_\Omega a]$,
且 $\operatorname{ess\,inf}_\Omega a\in\sigma(A)$, $\operatorname{ess\,sup}_\Omega a\in\sigma(A)$.

8. 若 $\Omega$ 是 $\R^N$ 中的开集 (配以 Lebesgue 测度) 且 $a\in C(\Omega)\cap L^\infty(\Omega)$,
证明 $\sigma(A)=\overline{a(\Omega)}$.

9. 证明 $\sigma(A)=\{0\}$ 当且仅当在 $\Omega$ 上几乎处处 $a=0$.

10. 假设 $\Omega$ 没有原子.证明 $A$ 紧当且仅当在 $\Omega$ 上几乎处处 $a=0$.
\end{problem}

\begin{problem}[第 4, 6 章]
\label{prob:34}
\runin{Hardy 算子 $Tu(x)=\frac1x\int_0^xu(t)dt$ 的谱分析}

\runin{-- A --}

设 $E=C([0,1])$ 配以范数 $\norm{u}=\sup_{t\in[0,1]}\abs{u(t)}$.给定 $u\in E$, 在 $[0,1]$
上定义函数 $Tu$:
\begin{equation*}
Tu(x)=\begin{cases}
\dfrac1x\int_0^xu(t)dt & \text{若}\ x\in(0,1],\\
u(0) & \text{若}\ x=0.
\end{cases}
\end{equation*}
验证 $Tu\in E$ 且对一切 $u\in E$, $\norm{Tu}\le\norm{u}$, 从而 $T\in\LL(E)$.

1. 证明 $\operatorname{EV}(T)=(0,1]$ 并确定相应的特征函数.

2. 验证 $\norm{T}_{\LL(E)}=1$.$T$ 是由 $E$ 到自身中的紧算子吗?

3. 证明 $\sigma(T)=[0,1]$.当 $\lambda\in\rho(T)$ 时给出 $(T-\lambda I)^{-1}$ 的显式公式.
证明对每个不属于 $\{0,1\}$ 的 $\lambda\in\R$, $(T-\lambda I)$ 是由 $E$ 到 $E$ 上的满射.
验证 $T$ 与 $(T-I)$ 不是满射.

4. 在本问中把 $T$ 看作由 $E=C([0,1])$ 到 $F=L^q(0,1)$ ($1\le q<\infty$) 中的有界算子.
证明 $T\in\KK(E,F)$.\\
{[提示: 考虑由 $(T_\varepsilon u)(x)=\dfrac{1}{x+\varepsilon}\int_0^xu(t)dt$ ($\varepsilon>0$)
定义的算子, 并估计当 $\varepsilon\to0$ 时 $\norm{T_\varepsilon-T}_{\LL(E,F)}$.]}

\runin{-- B --}

在本部分中设 $E=C^1([0,1])$ 配以范数 $\norm{u}=\sup_{t\in[0,1]}\abs{u(t)}+\sup_{t\in[0,1]}\abs{u'(t)}$.
给定 $u\in C^1([0,1])$, 如第 A 部分中那样定义 $Tu$.

1. 验证若 $u\in C^1([0,1])$, 则 $Tu\in C^1([0,1])$ 且对一切 $u\in E$, $\norm{Tu}\le\norm{u}$.

2. 证明 $\operatorname{EV}(T)=(0,1/2]\cup\{1\}$.

3. 证明 $\sigma(T)=[0,1/2]\cup\{1\}$.

\runin{-- C --}

在本部分中设 $E=L^p(0,1)$ ($1<p<\infty$).给定 $u\in L^p(0,1)$, 对 $x\in(0,1]$ 由
$Tu(x)=\frac1x\int_0^xu(t)dt$ 定义 $Tu$.验证 $Tu\in C((0,1])$ 且对每个 $q<p$,
$Tu\in L^q(0,1)$.我们的目标是证明 $Tu\in L^p(0,1)$ 且
\begin{equation}
\norm{Tu}_{L^p(0,1)}\le\frac{p}{p-1}\norm{u}_{L^p(0,1)}\quad\forall u\in E.
\label{eq:prob34.1}
\end{equation}

1. 证明当 $u\in C_c((0,1))$ 时 \eqref{eq:prob34.1} 成立.\\
{[提示: 令 $\phi(x)=\int_0^xu(t)dt$; 验证 $\abs{\phi}^p\in C^1([0,1])$ 并计算其导数.
使用公式 $\int_0^1\abs{Tu(x)}^pdx=\frac{1}{p-1}\int_0^1\abs{\phi(x)}^pd\left(-\frac{1}{x^{p-1}}\right)$
并分部积分估计 $\norm{Tu}_{L^p}$.]}

2. 证明对每个 $u\in E$, \eqref{eq:prob34.1} 成立.

在下文中把 $T$ 看作由 $E$ 到自身中的有界算子.

3. 证明 $\operatorname{EV}(T)=\left(0,\frac{p}{p-1}\right)$.

4. 由此推出 $\norm{T}_{\LL(E)}=\frac{p}{p-1}$.$T$ 是由 $E$ 到自身中的紧算子吗?

5. 证明 $\sigma(T)=\left[0,\frac{p}{p-1}\right]$.

6. 确定 $T^\star$.

7. 在本问中把 $T$ 看作由 $E=L^p(0,1)$ 到 $F=L^q(0,1)$ ($1\le q<p<\infty$) 中的有界算子.
证明 $T\in\KK(E,F)$.
\end{problem}

\begin{problem}[第~
\ref{chap:6}~章]
\label{prob:35}
\runin{Cotlar 引理}
设 $H$ 为与其对偶空间等同的 Hilbert 空间.

\runin{-- A --}

假设 $T\in\LL(H)$, 从而 $T^\star\in\LL(H)$.

1. 证明 $\norm{T^\star T}=\norm{T}^2$.

2. 在本问中假设 $T$ 自伴.证明对每个整数 $N$, $\norm{T^N}=\norm{T}^N$.

3. 由此推出 (对一般的 $T\in\LL(H)$), 对每个整数 $N$, $\norm{(T^\star T)^N}=\norm{T}^{2N}$.

\runin{-- B --}

设 $(T_j)_{1\le j\le m}$ 为 $\LL(H)$ 中的算子的有限族.假设对一切 $j,k\in\{1,2,\dots,m\}$,
\begin{align}
\norm{T_j^\star T_k}^{1/2}&\le\omega(j-k), \label{eq:prob35.1}\\
\norm{T_kT_j^\star}^{1/2}&\le\omega(k-j), \label{eq:prob35.2}
\end{align}
其中 $\omega:\Z\to[0,\infty)$.令 $\sigma=\sum_{i=-(m-1)}^{m-1}\omega(i)$.
本问题的目标是证明
\begin{equation}
\norm{\sum_{j=1}^m T_j}\le\sigma.
\label{eq:prob35.3}
\end{equation}
令 $U=\sum_{j=1}^mT_j$, 并固定整数 $N$.

1. 证明对整数 $j_1,k_1,\dots,j_N,k_N\in\{1,2,\dots,m\}$ 的任何选取,
\begin{equation*}
\norm{T_{j_1}^\star T_{k_1}T_{j_2}^\star T_{k_2}\dots T_{j_N}^\star T_{k_N}}
\le\sigma\omega(j_1-k_1)\omega(k_1-j_2)\dots\omega(k_{N-1}-j_N)\omega(j_N-k_N).
\end{equation*}

2. 由此推出 $\sum_{j_1}\sum_{k_1}\dots\sum_{j_N}\sum_{k_N}\norm{T_{j_1}^\star T_{k_1}\dots T_{j_N}^\star T_{k_N}}\le m\sigma^{2N}$.

3. 证明 $\norm{(U^\star U)^N}\le m\sigma^{2N}$, 并由此推出 \eqref{eq:prob35.3} 成立.
\end{problem}
\begin{problem}[第~
\ref{chap:6}~章]
\label{prob:36}
\runin{关于 Riesz--Fredholm 理论的补充}
设 $E$ 为 Banach 空间, $T\in\KK(E)$.对每个整数 $k\ge1$, 令
\begin{equation*}
N_k=N((I-T)^k)\quad\text{以及}\quad R_k=R((I-T)^k).
\end{equation*}

1. 验证对一切 $k\ge1$, $R_{k+1}\subset R_k$, $R_k$ 闭, $T(R_k)\subset R_k$, 以及
$(I-T)R_k\subset R_{k+1}$.

2. 证明存在整数 $p\ge1$ 使得对一切 $k<p$, $R_{k+1}\ne R_k$ (若 $p=1$ 则无条件),
且对一切 $k\ge p$, $R_{k+1}=R_k$.

3. 验证对一切 $k\ge1$, $N_k\subset N_{k+1}$, $\dim N_k<\infty$, $T(N_k)\subset N_k$,
以及 $(I-T)N_{k+1}\subset N_k$.

4. 证明对一切 $k\ge1$, $\operatorname{codim}R_k=\dim N_k$, 并由此推出对一切 $k<p$,
$N_{k+1}\ne N_k$ (若 $p=1$ 则无条件), 且对一切 $k\ge p$, $N_{k+1}=N_k$.

5. 证明 $R_p\cap N_p=\{0\}$ 且 $R_p+N_p=E$.

6. 证明 $(I-T)$ 限制在 $R_p$ 上是由 $R_p$ 到自身上的双射.

7. 这里额外假设 $E$ 是 Hilbert 空间且 $T$ 自伴.证明 $p=1$.
\end{problem}

\begin{problem}[第~
\ref{chap:6}~章]
\index{Rayleighshang@Rayleigh 商 (Rayleigh quotient)}
\index{CourantFischerjixiaohuajidayuanli@Courant--Fischer 极小化极大原理 (Courant--Fischer min-max principle)}
\label{prob:37}
\runin{Courant--Fischer 极小化极大原理. Rayleigh--Ritz 方法}
设 $H$ 为无限维可分 Hilbert 空间.设 $T$ 为由 $H$ 到自身的自伴紧算子, 满足对一切
$x\in H$, $(Tx,x)\ge0$.用 $(\mu_k)$ ($k\ge1$) 记其按不增顺序排列的, 计入重数的
特征值:
\begin{equation*}
\mu_1\ge\mu_2\ge\dots\ge0.
\end{equation*}
设 $(e_j)$ 为由特征向量组成的相应标准正交基.设 $E_k$ 为由 $\{e_1,e_2,\dots,e_k\}$
张成的空间.对 $x\ne0$, 我们定义 \emph{Rayleigh 商}
\begin{equation*}
R(x)=\frac{(Tx,x)}{\abs{x}^2}.
\end{equation*}

1. 证明对一切 $k\ge1$, $\min_{\substack{x\in E_k\\x\ne0}}R(x)=\mu_k$.

2. 证明对一切 $k\ge2$, $\max_{\substack{x\in E_{k-1}^\perp\\x\ne0}}R(x)=\mu_k$,
且 $\max_{\substack{x\in H\\x\ne0}}R(x)=\mu_1$.

3. 设 $\Sigma$ 为 $H$ 的任何 $k$ 维子空间 ($k\ge1$).证明
$\min_{\substack{x\in\Sigma\\x\ne0}}R(x)\le\mu_k$.\\
{[提示: 若 $k\ge2$, 证明 $\Sigma\cap E_{k-1}^\perp\ne\{0\}$ 并应用第 2 问.]}

4. 由此推出对一切 $k\ge1$,
\begin{equation*}
\max_{\substack{\Sigma\subset H\\\dim\Sigma=k}}\min_{\substack{x\in\Sigma\\x\ne0}}R(x)=\mu_k.
\end{equation*}

5. 设 $\Sigma$ 为 $H$ 的任何 $(k-1)$ 维子空间 ($k\ge2$).证明
$\max_{\substack{x\in\Sigma^\perp\\x\ne0}}R(x)\ge\mu_k$.\\
{[提示: 证明 $\Sigma^\perp\cap E_k\ne\{0\}$.]}

6. 由此推出对一切 $k\ge1$,
\begin{equation*}
\min_{\substack{\Sigma\subset H\\\dim\Sigma=k-1}}\max_{\substack{x\in\Sigma^\perp\\x\ne0}}R(x)=\mu_k.
\end{equation*}

7. 这里假设 $N(T)=\{0\}$, 从而对一切 $x\ne0$, $R(x)\ne0$, 或等价地对一切 $k$,
$\mu_k>0$.证明对一切 $k\ge1$,
\begin{equation*}
\min_{\substack{\Sigma\subset H\\\dim\Sigma=k}}\max_{\substack{x\in\Sigma\\x\ne0}}\frac{1}{R(x)}=\frac{1}{\mu_k},
\end{equation*}
以及
\begin{equation*}
\max_{\substack{M\subset H\\\operatorname{codim}M=k-1}}\min_{\substack{x\in M\\x\ne0}}\frac{1}{R(x)}=\frac{1}{\mu_k},
\end{equation*}
其中 $\Sigma$ 与 $M$ 是 $H$ 的闭子空间.特别地, 对 $k=1$,
$\min_{\substack{x\in H\\x\ne0}}\frac{1}{R(x)}=\frac{1}{\mu_1}$; 此外对一切 $k\ge2$,
$\min_{\substack{x\in E_{k-1}^\perp\\x\ne0}}\frac{1}{R(x)}=\frac{1}{\mu_k}$.

8. 设 $V$ 为 $H$ 的闭子空间 (有限维或无限维).设 $P_V$ 为由 $H$ 到 $V$ 上的正交
投影算子, 考虑由 $S=P_V\circ T|_V$ 定义的算子 $S:V\to V$.验证 $S$ 是由 $V$ 到
自身的自伴紧算子, 且对一切 $x\in V$, $(Sx,x)\ge0$.

9. 用 $(\nu_k)$ ($k\ge1$) 记 $S$ 的按不增顺序排列的, 计入重数的特征值.证明对
满足 $1\le k\le\dim V$ 的一切 $k$,
\begin{equation*}
\max_{\substack{\Sigma\subset V\\\dim\Sigma=k}}\min_{\substack{x\in\Sigma\\x\ne0}}R(x)=\nu_k.
\end{equation*}
由此推出对满足 $1\le k\le\dim V$ 的一切 $k$, $\nu_k\le\mu_k$.

10. 现在考虑 $H$ 的闭子空间的递增序列 $V^{(n)}$, 满足 $\overline{\bigcup_nV^{(n)}}=H$.
令 $S^{(n)}=P_{V^{(n)}}\circ T|_{V^{(n)}}$, 并设 $(\nu_k^{(n)})$ 记如第 9 问中那样排列的
$S^{(n)}$ 的特征值.证明对每个固定的 $k$, 序列 $n\mapsto\nu_k^{(n)}$ 单调不减, 且当
$n\to\infty$ 时收敛于 $\mu_k$.
\end{problem}

\begin{problem}[第 2, 6, 11 章]
\index{FredholmNoethersuanzi@Fredholm--Noether 算子 (Fredholm--Noether operator)}
\label{prob:38}
\runin{Fredholm--Noether 算子}
设 $E$ 与 $F$ 为 Banach 空间, $T\in\LL(E,F)$.

\runin{-- A --}

第 A 部分的目标是证明下述条件等价:
\begin{equation}
\text{(a)}\ R(T)\ \text{闭且在}\ F\ \text{中有有限余维,}\quad
\text{(b)}\ N(T)\ \text{在}\ E\ \text{中有补空间}.
\tag{1}
\label{eq:prob38.1}
\end{equation}
\begin{equation}
\text{存在}\ S\in\LL(F,E)\ \text{与}\ K\in\KK(F,F)\ \text{使得}\ T\circ S=I_F+K.
\tag{2}
\label{eq:prob38.2}
\end{equation}
\begin{equation}
\text{存在}\ U\in\LL(F,E)\ \text{与}\ F\ \text{中的有限秩投影}\ P\ \text{使得}\ T\circ U=I_F-P.
\tag{3}
\label{eq:prob38.3}
\end{equation}
此外, 可以选取 $U$ 与 $P$ 使得 $\dim R(P)=\operatorname{codim}R(T)$.

1. 证明 \eqref{eq:prob38.1} $\Rightarrow$ \eqref{eq:prob38.3}.\\
{[提示: 设 $X$ 为 $N(T)$ 在 $E$ 中的补空间.则 $T|_X$ 是由 $X$ 到 $R(T)$ 上的双射.
用 $U_0$ 记其逆算子.设 $Q$ 为由 $F$ 到 $R(T)$ 上的投影算子, 令 $U=U_0\circ Q$.]}

2. 证明 \eqref{eq:prob38.2} $\Rightarrow$ \eqref{eq:prob38.3}.\\
{[提示: 使用习题 6.25.]}

3. 证明 \eqref{eq:prob38.3} $\Rightarrow$ \eqref{eq:prob38.1}.\\
{[提示: 为建立 \eqref{eq:prob38.1} 的 (a) 部分, 注意 $R(T)\supset R(I_F-P)$ 并应用
命题  \ref{prop:11.5}.同理证明 $R(U^\star)$ 闭, 从而 $R(U)$ 也闭.最后证明在 $E$ 中
存在有限维空间 $\Sigma_1$ 与 $\Sigma_2$ 使得 $N(T)+R(U)+\Sigma_1=E$ 且
$N(T)\cap R(U)\subset\Sigma_2$.然后应用命题  \ref{prop:11.7}.]}

4. 得出结论.

\runin{-- B --}

证明下述条件等价:
\begin{equation}
\text{(a)}\ R(T)\ \text{闭且有补空间,}\quad\text{(b)}\ \dim N(T)<\infty.
\tag{4}
\label{eq:prob38.4}
\end{equation}
\begin{equation}
\text{存在}\ \widetilde S\in\LL(F,E)\ \text{与}\ \widetilde K\in\KK(E,E)\ \text{使得}\ \widetilde S\circ T=I_E+\widetilde K.
\tag{5}
\label{eq:prob38.5}
\end{equation}
\begin{equation}
\text{存在}\ \widetilde U\in\LL(F,E)\ \text{与}\ E\ \text{中的有限秩投影}\ \widetilde P\ \text{使得}\ \widetilde U\circ T=I_E-\widetilde P.
\tag{6}
\label{eq:prob38.6}
\end{equation}

\runin{-- C --}

若算子 $T\in\LL(E,F)$ 满足
\begin{equation}
\text{(a)}\ R(T)\ \text{闭且有有限余维,}\quad\text{(b)}\ \dim N(T)<\infty,
\tag{FN}
\label{eq:prob38.FN}
\end{equation}
则称它为 \emph{Fredholm} 算子 (或 \emph{Noether} 算子) ($R(T)$ 闭的性质可由其他
假设推出; 见习题 2.27).满足 \eqref{eq:prob38.FN} 的算子类记作 $\Phi(E,F)$.$T$ 的
\emph{指标}按定义为
\begin{equation*}
\operatorname{ind}T=\dim N(T)-\operatorname{codim}R(T).
\end{equation*}

1. 设 $T\in\Phi(E,F)$.证明存在 $U\in\LL(F,E)$ 以及 $F$ (相应地 $E$) 中的有限秩投影
$P$ (相应地 $\widetilde P$), 使得
\begin{equation}
T\circ U=I_F-P,\qquad U\circ T=I_E-\widetilde P,
\label{eq:prob38.7}
\end{equation}
满足 $\dim R(P)=\operatorname{codim}R(T)$, $\dim R(\widetilde P)=\dim N(T)$.\\
{[提示: 使用第 A1 问中构造的算子 $U$.]}

满足 $T\circ V=I_F+K$ 且 $V\circ T=I_E+\widetilde K$ ($K\in\KK(F)$, $\widetilde K\in\KK(E)$)
的算子 $V\in\LL(F,E)$ 称为 $T$ 的\emph{拟逆} (或模紧算子的逆).

2. 证明任何拟逆 $V$ 属于 $\Phi(F,E)$.

3. 证明算子 $T\in\LL(E,F)$ 属于 $\Phi(E,F)$ 当且仅当 $R(T)$ 闭, $\dim N(T)<\infty$,
且 $\dim N(T^\star)<\infty$.此外,
\begin{equation*}
\operatorname{ind}T=\dim N(T)-\dim N(T^\star).
\end{equation*}
[提示: 应用命题  \ref{prop:11.14} 与命题 2.18.]
% 原书此处印作 "Propositions 11.14 and 2.18" (印刷页 494); 第~
% \ref{chap:2}~章并无编号 2.18
% 的命题 (2.18 是推论~
% \ref{cor:2.18}). 依技能规范保留原样并记录, 不擅自更正, 故此处不用
% \ref 以免产生未定义引用.

4. 设 $T\in\Phi(E,F)$.证明 $T^\star\in\Phi(F^\star,E^\star)$ 且
$\operatorname{ind}T^\star=-\operatorname{ind}T$.\\
{[提示: 应用命题  \ref{prop:11.13} 与定理  \ref{thm:2.19}.]}

5. 反之, 设 $T\in\LL(E,F)$ 满足 $T^\star\in\Phi(F^\star,E^\star)$.证明 $T\in\Phi(E,F)$.

6. 假设 $J\in\LL(E,F)$ 是双射且 $K\in\KK(E,F)$.证明 $T=J+K$ 属于 $\Phi(E,F)$ 且
$\operatorname{ind}T=0$.反之, 若 $T\in\Phi(E,F)$ 且 $\operatorname{ind}T=0$, 证明
$T$ 可以写成 $T=J+K$, 其中 $J$ 与 $K$ 如上 (甚至可以选取 $K$ 为有限秩算子).\\
{[提示: 应用定理  \ref{thm:6.6}, 证明 $I_E+J^{-1}\circ K$ 属于 $\Phi(E,E)$ 且指标为零.
关于逆命题, 考虑由 $N(T)$ 到 $R(T)$ 的补空间 $Y$ 上的同构.]}

7. 设 $T\in\Phi(E,F)$ 且 $K\in\KK(E,F)$.证明 $T+K\in\Phi(E,F)$.

8. 在前一问的假设下, 证明 $\operatorname{ind}(T+K)=\operatorname{ind}T$.\\
{[提示: 令 $\widetilde E=E\times Y$, $\widetilde F=F\times N(T)$, 以及由
$\widetilde T(x,y)=(Tx+Kx,0)$ 定义的 $\widetilde T:\widetilde E\to\widetilde F$.
证明 $\widetilde T=\widetilde J+\widetilde K$, 其中 $\widetilde J$ 是由 $\widetilde E$
到 $\widetilde F$ 上的双射, 且 $\widetilde K\in\KK(\widetilde E,\widetilde F)$.
然后应用第 6 问.]}

9. 设 $T\in\Phi(E,F)$.证明存在 $\varepsilon>0$ (依赖于 $T$), 使得对满足
$\norm{M}<\varepsilon$ 的每个 $M\in\LL(E,F)$, 有 $T+M\in\Phi(E,F)$.证明
$\operatorname{ind}(T+M)=\operatorname{ind}T$.\\
{[提示: 设 $V$ 为 $T$ 的拟逆.则若 $\norm{M}<\norm{V}^{-1}$, $W=I_E+(V\circ M)$
是双射.验证 $T+M=(T\circ W)+\text{紧算子}$; 然后应用前一问.]}

10. 设 $(H_t)_{t\in[0,1]}$ 为 $\LL(E,F)$ 中的算子族.假设 $t\mapsto H_t$ 从 $[0,1]$
到 $\LL(E,F)$ 连续, 且对一切 $t\in[0,1]$, $H_t\in\Phi(E,F)$.证明 $\operatorname{ind}H_t$
在 $[0,1]$ 上为常数.

11. 设 $E_1,E_2$ 与 $E_3$ 为 Banach 空间, 设 $T_1\in\Phi(E_1,E_2)$, $T_2\in\Phi(E_2,E_3)$.
证明 $T_2\circ T_1\in\Phi(E_1,E_3)$.

12. 沿用上面的相同记号, 证明
\begin{equation*}
\operatorname{ind}(T_2\circ T_1)=\operatorname{ind}T_1+\operatorname{ind}T_2.
\end{equation*}
[提示: 考虑用矩阵记号对 $t\in[0,1]$ 由
\begin{equation*}
H_t=\begin{pmatrix}I&0\\0&T_2\end{pmatrix}\begin{pmatrix}(1-t)I&tI\\-tI&(1-t)I\end{pmatrix}\begin{pmatrix}T_1&0\\0&I\end{pmatrix}
\end{equation*}
定义的算子族 $H_t:E_1\times E_2\to E_2\times E_3$, 其中 $I$ 是 $E_2$ 中的恒等算子.
验证 $t\mapsto H_t$ 从 $[0,1]$ 到 $\LL(E_1\times E_2,E_2\times E_3)$ 连续.利用前一问
证明对每个 $t$, $H_t\in\Phi(E_1\times E_2,E_2\times E_3)$.计算 $\operatorname{ind}H_0$
与 $\operatorname{ind}H_1$.]

13. 设 $T\in\Phi(E,F)$.计算 $T$ 的任何拟逆 $V$ 的指标.

\runin{-- D --}

在本部分中我们研究两个简单的例子.

1. 假设 $\dim E<\infty$ 且 $\dim F<\infty$.证明从 $E$ 到 $F$ 的任何线性算子 $T$
属于 $\Phi(E,F)$ 并计算其指标.

2. 设 $E=F=\ell^2$.考虑习题 6.18 中定义的移位算子 $S_r$ 与 $S_\ell$.证明对满足
$\lambda\ne\pm1$ 的每个 $\lambda\in\R$, 有 $S_r-\lambda I\in\Phi(\ell^2,\ell^2)$
且 $S_\ell-\lambda I\in\Phi(\ell^2,\ell^2)$.计算它们的指标.证明 $S_r\pm I$,
$S_\ell\pm I$ 不属于 $\Phi(\ell^2,\ell^2)$.\\
{[提示: 使用习题 6.18 的结果.]}
\end{problem}

\begin{problem}[第 5, 6 章]
\label{prob:39}
\runin{自伴非负算子的平方根}
设 $H$ 为 Hilbert 空间.设 $S\in\LL(H)$; 若对一切 $x\in H$, $(Sx,x)\ge0$, 则称 $S$ 是
\emph{非负的}, 记作 $S\ge0$.当 $S_1,S_2\in\LL(H)$ 时, 若 $S_1-S_2\ge0$, 则记
$S_1\ge S_2$ (或 $S_2\le S_1$).

\runin{-- A --}

1. 设 $S\in\LL(H)$ 满足 $S^\star=S$ 且 $0\le S\le I$.证明 $\norm{S^2}=\norm{S}^2\le1$,
且 $0\le S^2\le S\le I$.\\
{[提示: 使用习题 6.24.]}

2. 设 $S\in\LL(H)$ 满足 $S^\star=S$ 且 $S\ge0$.设 $P(t)=\sum a_kt^k$ 为对一切 $k$
满足 $a_k\ge0$ 的多项式.证明 $[P(S)]^\star=P(S)$ 且 $P(S)\ge0$.

3. 设 $(S_n)$ 为 $\LL(H)$ 中的序列使得对一切 $n$, $S_n^\star=S_n$, 且对一切 $n$,
$S_{n+1}\le S_n$.假设存在常数 $M$ 使得对一切 $n$, $\norm{S_n}\le M$.证明对每个
$x\in H$, 当 $n\to\infty$ 时 $S_nx$ 收敛于一个极限, 记作 $Sx$, 且 $S\in\LL(H)$
满足 $S^\star=S$.\\
{[提示: 设 $n\ge m$.使用习题 6.24 证明 $\abs{S_nx-S_mx}^2\le2M(S_mx-S_nx,x)$.]}

\runin{-- B --}

假设 $T\in\LL(H)$ 满足 $T^\star=T$, $T\ge0$, 且 $\norm{T}\le1$.考虑从 $S_0=I$ 开始
由下式定义的序列 $(S_n)$:
\begin{equation*}
S_{n+1}=S_n+\frac12(T-S_n^2),\quad n\ge0.
\end{equation*}

1. 证明对一切 $n\ge0$, $S_n^\star=S_n$.

2. 证明 $I-S_{n+1}=\frac12(I-S_n)^2+\frac12(I-T)$, 并由此推出对一切 $n$, $I-S_n\ge0$.

3. 证明对一切 $n$, $S_n\ge0$.\\
{[提示: 利用第 A1 与 B2 问由数学归纳法证明 $I-S_n\le I$.]}

4. 由此推出对一切 $n$, $\norm{S_n}\le1$.

5. 证明对一切 $n$, $S_n-S_{n+1}=\frac12[(I-S_n)+(I-S_{n-1})]\circ(S_{n-1}-S_n)$,
并由此推出对一切 $n$, $S_{n-1}-S_n\ge0$.\\
{[提示: 由数学归纳法证明 $(I-S_n)=P_n(I-T)$ 与 $(S_{n-1}-S_n)=Q_n(I-T)$, 其中 $P_n$
与 $Q_n$ 是具有非负系数的多项式.]}

6. 证明 $\lim_{n\to\infty}S_nx=Sx$ 存在.证明 $S\in\LL(H)$ 满足 $S^\star=S$, $S\ge0$,
$\norm{S}\le1$, 且 $S^2=T$.

\runin{-- C --}

1. 设 $U\in\LL(H)$ 满足 $U^\star=U$ 且 $U\ge0$.证明存在 $V\in\LL(H)$ 使得 $V^\star=V$,
$V\ge0$, 且 $V^2=U$.\\
{[提示: 对 $T=U/\norm{U}$ 应用第 B 部分的构造.]}

接着我们证明 $V$ 的唯一性.更精确地说, 若 $W$ 是满足 $W^\star=W$, $W\ge0$ 且 $W^2=U$
的任何算子 $W\in\LL(H)$, 则 $W=V$.算子 $V$ 称为 $U$ 的\emph{平方根}, 记作 $U^{1/2}$.

2. 证明上述构造的算子 $V$ 与每个与 $U$ 对易的算子 $X$ 对易 (即 $X\circ U=U\circ X$
蕴含 $X\circ V=V\circ X$).

3. 证明 $W$ 与 $U$ 对易, 并由此推出 $V$ 与 $W$ 对易.

4. 验证 $(V-W)\circ(V+W)=0$, 并由此推出在 $R(V+W)$ 上 $V=W$.证明
$N(V)=N(W)=N(U)=N(V+W)$.得出在 $H$ 上 $V=W$ 的结论.\\
{[提示: 注意在 $\overline{R(V+W)}=N(U)^\perp$ 上 $V=W$, 且在 $N(U)$ 上 $V=W=0$.]}

5. 证明 $\norm{U^{1/2}}=\norm{U}^{1/2}$.

6. 设 $U_1,U_2\in\LL(H)$ 满足 $U_1^\star=U_1,U_2^\star=U_2,U_1\ge0,U_2\ge0$, 且
$U_1\circ U_2=U_2\circ U_1$.证明 $U_1\circ U_2\ge0$.\\
{[提示: 引入 $U_1^{1/2}$ 与 $U_2^{1/2}$.]}

\runin{-- D --}

设 $U\in\KK(H)$ 满足 $U^\star=U$ 且 $U\ge0$.证明其平方根 $V$ 属于 $\KK(H)$.假设
$H$ 可分, 在由 $U$ 的特征向量组成的 Hilbert 基上计算 $V$.求 $V$ 的特征值.
\end{problem}

\begin{problem}[第 4, 5, 6 章]
\index{HilbertSchmidtsuanzi@Hilbert--Schmidt 算子 (Hilbert--Schmidt operator)}
\label{prob:40}
\runin{Hilbert--Schmidt 算子}

\runin{-- A --}

设 $E$ 与 $F$ 为可分 Hilbert 空间, 都与其对偶空间等同.$E$ 与 $F$ 上的范数用相同符号
$\abs{\cdot}$ 表示.设 $T\in\LL(E,F)$, 从而 $T^\star\in\LL(F,E)$.

1. 设 $(e_k)$ (相应地 $(f_k)$) 为 $E$ (相应地 $F$) 的任意标准正交基.证明
$\sum_{k=1}^\infty\abs{T(e_k)}^2<\infty$ 当且仅当 $\sum_{k=1}^\infty\abs{T^\star(f_k)}^2<\infty$,
且 $\sum_{k=1}^\infty\abs{T(e_k)}^2=\sum_{k=1}^\infty\abs{T^\star(f_k)}^2$.

2. 设 $(e_k)$ 与 $(\tilde e_k)$ 为 $E$ 的两个标准正交基.证明
$\sum_{k=1}^\infty\abs{T(e_k)}^2<\infty$ 当且仅当 $\sum_{k=1}^\infty\abs{T(\tilde e_k)}^2<\infty$,
且 $\sum_{k=1}^\infty\abs{T(e_k)}^2=\sum_{k=1}^\infty\abs{T(\tilde e_k)}^2$.

若存在 $E$ 的某个标准正交基 $(e_k)$ 使得 $\sum_{k=1}^\infty\abs{T(e_k)}^2<\infty$,
则称 $T\in\LL(E,F)$ 为 \emph{Hilbert--Schmidt} 算子, 记作 $T\in\mathcal{HS}(E,F)$.

3. 证明 $\mathcal{HS}(E,F)$ 是 $\LL(E,F)$ 的线性子空间, 且
\begin{equation*}
\norm{T}_{\mathcal{HS}}=\left(\sum_{k=1}^\infty\abs{T(e_k)}^2\right)^{1/2}
\end{equation*}
定义了 $\mathcal{HS}(E,F)$ 上的范数.设 $\norm{\cdot}$ 表示 $\LL(E,F)$ 上的通常范数.
证明对一切 $T\in\mathcal{HS}(E,F)$, $\norm{T}\le\norm{T}_{\mathcal{HS}}$.

4. 证明配以范数 $\norm{\cdot}_{\mathcal{HS}}$ 的 $\mathcal{HS}(E,F)$ 是 Banach 空间.
证明事实上它是 Hilbert 空间.

5. 证明 $\mathcal{HS}(E,F)\subset\KK(E,F)$.\\
{[提示: 给定 $x\in E$, 写出 $x=\sum_{k=1}^\infty x_ke_k$, 并令
$T_n(x)=\sum_{k=1}^nx_kT(e_k)$.证明当 $n\to\infty$ 时 $\norm{T_n-T}\to0$.]}

6. 证明由 $E$ 到 $F$ 中的任何有限秩算子属于 $\mathcal{HS}(E,F)$.

7. 设 $T\in\LL(E,F)$.证明 $T\in\mathcal{HS}(E,F)\iff T^\star\in\mathcal{HS}(F,E)$,
且 $\norm{T^\star}_{\mathcal{HS}(F,E)}=\norm{T}_{\mathcal{HS}(E,F)}$.

8. 假设 $T\in\KK(E,E)$ 满足 $T^\star=T$, 设 $(\lambda_k)$ 记 $T$ 的特征值序列.
证明 $T\in\mathcal{HS}(E,E)$ 当且仅当 $\sum_{k=1}^\infty\lambda_k^2<\infty$, 且
$\norm{T}_{\mathcal{HS}}^2=\sum_{k=1}^\infty\lambda_k^2$.构造一个在 $E=\ell^2$ 上的
算子 $T\in\KK(E,E)$ 使得 $T\notin\mathcal{HS}(E,E)$.

9. 设 $G$ 为另一个可分 Hilbert 空间.设 $T_1\in\LL(E,F)$ 且 $T_2\in\LL(F,G)$.
证明若 $T_1$ 或 $T_2$ 之一属于 $\mathcal{HS}$, 则 $T_2\circ T_1\in\mathcal{HS}(E,G)$.

10. 设 $T\in\mathcal{HS}(E,E)$ 并假设 $N(I+T)=\{0\}$.证明 $(I+T)$ 是双射, 且
$(I+T)^{-1}=I+S$, 其中 $S\in\mathcal{HS}(E,E)$.

11. 设 $(e_k)$ (相应地 $(f_\ell)$) 为 $E$ (相应地 $F$) 的标准正交基.考虑由
$T_{k,\ell}(x)=(x,e_k)f_\ell$ 定义的算子 $T_{k,\ell}:E\to F$.证明 $(T_{k,\ell})$
是 $\mathcal{HS}(E,F)$ 的标准正交基.

\runin{-- B --}

假设 $\Omega$ 是 $\R^N$ 中的开子集.在下文中取 $E=F=L^2(\Omega)$.设
$K\in L^2(\Omega\times\Omega)$, 考虑算子
\begin{equation}
(Tu)(x)=\int_\Omega K(x,y)u(y)dy.
\label{eq:prob40.1}
\end{equation}

1. 证明 $T\in\LL(E,E)$ 且 $\norm{T}_{\LL(E,E)}\le\norm{K}_{L^2(\Omega\times\Omega)}$.

2. 证明 $T\in\mathcal{HS}(E,E)$ 且 $\norm{T}_{\mathcal{HS}(E,E)}\le\norm{K}_{L^2(\Omega\times\Omega)}$.\\
{[提示: 设 $(e_j)$ 为 $L^2(\Omega)$ 的标准正交基.验证族 $e_{j,k}=e_j\otimes e_k$
(其中 $(e_j\otimes e_k)(x,y)=e_j(x)e_k(y)$) 是 $L^2(\Omega\times\Omega)$ 的标准正交基.
然后写出 $\norm{T(e_k)}_{L^2(\Omega)}^2=\sum_{j=1}^\infty\abs{(T(e_k),e_j)}^2=\sum_{j=1}^\infty\abs{(K,e_j\otimes e_k)}$.]}

3. 反之, 设 $T\in\mathcal{HS}(E,E)$.证明存在唯一的函数 $K\in L^2(\Omega\times\Omega)$
使得 \eqref{eq:prob40.1} 成立.$K$ 称为 $T$ 的\emph{积分核}.\\
{[提示: 令 $t_{j,k}=(Te_k,e_j)$ 并验证 $\sum_{j,k=1}^\infty\abs{t_{j,k}}^2<\infty$.
定义 $K=\sum_{j,k=1}^\infty t_{j,k}e_j\otimes e_k$ 并证明 \eqref{eq:prob40.1} 成立.]}

4. 假设 $\Omega=(0,1)$, $E=L^2(\Omega)$, 考虑算子 $(Tu)(x)=\int_0^xu(t)dt$.证明
$T\in\mathcal{HS}(E,E)$ 并计算 $\norm{T}_{\mathcal{HS}}$.
\end{problem}
\begin{problem}[第 1, 6 章]
\label{prob:41}
\runin{Krein--Rutman 定理}
设 $E$ 为 Banach 空间, $P\subset E$ 为包含 $0$ 的闭凸集.假设 $P$ 是以 $0$ 为顶点的
凸锥, 即对一切 $\lambda>0,\mu>0,x\in P,y\in P$, $\lambda x+\mu y\in P$.

假设
\begin{equation}
\operatorname{Int}P\ne\emptyset\quad\text{且}\quad P\ne E.
\label{eq:prob41.1}
\end{equation}
设 $T\in\KK(E)$ 满足
\begin{equation}
T(P\setminus\{0\})\subset\operatorname{Int}P.
\label{eq:prob41.2}
\end{equation}

\runin{-- A --}

1. 证明 $(\operatorname{Int}P)\cap(-P)=\emptyset$.\\
{[提示: 使用习题 1.7.]}

在下文中我们固定某个 $u\in\operatorname{Int}P$.

2. 证明存在 $\alpha>0$ 使得对一切 $x\in P$, $\norm{x+u}\ge\alpha$.\\
{[提示: 用反证法并得出 $-u\in P$.]}

3. 验证存在 $r>0$ 使得 $Tu-ru\in P$.

4. 假设某个 $x\in P$ 满足对某个 $\lambda\in\R$, $T(x+u)=\lambda x$.证明 $\lambda\ge r$.\\
{[提示: 在 $E$ 上引入由 $y\ge z\iff y-z\in P$ 定义的序关系是方便的.由数学归纳法证明
$(\lambda/r)^nx\ge u$, $n=1,2,\dots$.]}

5. 考虑非线性映射 $F(x)=T\left(\dfrac{x+u}{\norm{x+u}}\right)$, $x\in P$.证明
$F:P\to P$ 连续且对 $E$ 的某个紧子集 $K$ 有 $F(P)\subset K$.由此推出存在某个
$x_1\in P$ 使得 $T(x_1+u)=\lambda_1x_1$, 其中 $\lambda_1=\norm{x_1+u}\ge r$.\\
{[提示: 应用 Schauder 不动点定理; 见习题 6.26.]}

6. 由此推出对每个 $\varepsilon>0$, 存在 $x_\varepsilon\in P$ 使得
$T(x_\varepsilon+\varepsilon u)=\lambda_\varepsilon x_\varepsilon$, 其中
$\lambda_\varepsilon=\norm{x_\varepsilon+\varepsilon u}\ge r$.

7. 证明存在 $x_0\in\operatorname{Int}P$ 与 $\mu_0>0$ 使得 $Tx_0=\mu_0x_0$.\\
{[提示: 证明 $(x_\varepsilon)$ 有界.由此推出存在满足所要求性质的序列
$\varepsilon_n\to0$ 使得 $x_{\varepsilon_n}\to x_0$ 且 $\lambda_{\varepsilon_n}\to\mu_0$.]}

\runin{-- B --}

1. 给定两点 $a\in\operatorname{Int}P$ 与 $b\in E\setminus P$, 证明存在唯一的
$\sigma\in(0,1)$ 使得对一切 $t\in[0,\sigma)$, $(1-t)a+tb\in\operatorname{Int}P$;
$(1-\sigma)a+\sigma b\in P$; 且对一切 $t\in(\sigma,1]$, $(1-t)a+tb\notin P$.
然后令 $\tau(a,b)=\sigma/(1-\sigma)$, 满足 $0<\tau(a,b)<\infty$.

2. 设 $x\in P\setminus\{0\}$ 满足对某个 $\mu\in\R$, $Tx=\mu x$.证明 $\mu=\mu_0$
且对某个 $m>0$ 有 $x=mx_0$, 其中 $\mu_0$ 与 $x_0$ 是在第 A7 问中构造的.\\
{[提示: 反设对一切 $m>0$ 都有 $x\ne mx_0$.证明 $\mu>0,x\in\operatorname{Int}P$,
且 $-x\notin P$.令 $y=x_0-\tau_0x$, 其中 $\tau_0=\tau(x_0,-x)$.计算 $Ty$ 并得出
$\mu<\mu_0$.然后互换 $x_0$ 与 $x$ 的角色.]}

3. 设 $x\in E\setminus\{0\}$ 满足对某个 $\mu\in\R$, $Tx=\mu x$.证明或者 $\mu=\mu_0$
且 $x=mx_0$ ($m\in\R,m\ne0$), 或者 $\abs{\mu}<\mu_0$.\\
{[提示: 由第 2 问可以假设 $x\notin P$ 且 $-x\notin P$.若 $\mu>0$ 考虑 $\tau(x_0,x)$;
若 $\mu<0$ 则同时考虑 $\tau(x_0,x)$ 与 $\tau(x_0,-x)$.]}

4. 由此推出 $N(T-\mu_0I)=\R x_0$.换言之, 特征值 $\mu_0$ 的几何重数为 1.

5. 证明对所有 $k\ge2$, $N((T-\mu_0I)^k)=\R x_0$.换言之, 特征值 $\mu_0$ 的代数重数也是 1.\\
{[提示: 鉴于问题 36, 只须证明 $N((T-\mu_0I)^2)=\R x_0$.]}
\end{problem}

\begin{problem}[第~
\ref{chap:6}~章]
\label{prob:42}
\runin{关于不变子空间的 Lomonosov 定理}
设 $E$ 为无限维 Banach 空间, 设 $T\in\KK(E)$, $T\ne0$.第 A 部分的目标是证明
存在 $T$ 的非平凡闭不变子空间 $Z$, 即 $T(Z)\subset Z$, 其中 $Z\ne\{0\}$ 且 $Z\ne E$.

\runin{-- A --}

令
\begin{align*}
\mathcal{A}&=\operatorname{span}\{I,T,T^2,\dots\}\\
&=\left\{\sum_{i\in I}\lambda_iT^i;\ \lambda_i\in\R,\ I\ \text{是}\ \{0,1,2,\dots\}\ \text{的有限子集}\right\}.
\end{align*}
对每个 $y\in E$, 令 $\mathcal{A}y=\{Sy;\ S\in\mathcal{A}\}$.显然 $y\in\mathcal{A}y$,
因此对每个 $y\ne0$, $\mathcal{A}y\ne\{0\}$.此外 $\mathcal{A}y$ 是 $E$ 的子空间
且 $T(\mathcal{A}y)\subset\mathcal{A}y$, 从而 $T(\overline{\mathcal{A}y})\subset\overline{\mathcal{A}y}$.
若对某个 $y\ne0$ 有 $\overline{\mathcal{A}y}\ne E$, 则 $\overline{\mathcal{A}y}$
是 $T$ 的非平凡闭不变子空间.因此我们可以假设对一切 $y\in E$ ($y\ne0$),
\begin{equation}
\overline{\mathcal{A}y}=E.
\label{eq:prob42.1}
\end{equation}

由于 $T\ne0$, 我们可以固定某个 $x_0\in E$ 使得 $Tx_0\ne0$, 并固定 $r$ 使得
$0<r\le\dfrac{\norm{Tx_0}}{2\norm{T}}\le\dfrac{\norm{x_0}}{2}$.令
$C=\{x\in E;\ \norm{x-x_0}\le r\}$.

1. 验证 $0\notin C$, 且对一切 $x\in C$, $\norm{Tx-Tx_0}\le\frac12\norm{Tx_0}$,
从而对一切 $x\in C$, $\norm{Tx}\ge\frac12\norm{Tx_0}$.由此推出 $0\notin\overline{T(C)}$.

2. 证明对每个 $y\in E$ ($y\ne0$), 存在某个 $S\in\mathcal{A}$ (记作 $S_y$), 使得
$\norm{S_yy-x_0}\le r/2$.\\
{[提示: 使用假设 \eqref{eq:prob42.1}.]}

3. 由此推出对每个 $y\in E$ ($y\ne0$), 存在某个 $\varepsilon_y>0$ (依赖于 $y$),
使得对一切 $z\in B(y,\varepsilon_y)$, $\norm{S_yz-x_0}\le r$.

4. 考虑 $T(C)$ 被球 $B(y_j,\frac12\varepsilon_{y_j})$ ($j\in J,J$ 有限) 的有限覆盖.
对 $j\in J$ 与 $x\in E$, 令 $q_j(x)=\max\{0,\varepsilon_{y_j}-\norm{Tx-y_j}\}$ 以及
$q(x)=\sum_{j\in J}q_j(x)$.验证函数 $q_j$ 与 $q$ 在 $E$ 上连续.证明对一切 $x\in C$,
$q(x)\ge\min_{j\in J}\frac12\varepsilon_{y_j}>0$.对 $x\in C$, 令
\begin{equation*}
F(x)=\frac{1}{q(x)}\sum_{j\in J}q_j(x)S_{y_j}(Tx).
\end{equation*}

5. 证明 $F$ 从 $C$ 到 $E$ 连续, 且对一切 $x\in C$, $\norm{F(x)-x_0}\le r$.\\
{[提示: 使用第 3 问.]}

6. 证明 $F(C)\subset K$, 其中 $K$ 是 $C$ 的紧子集.由此推出存在 $\xi\in C$ 使得
$F(\xi)=\xi$.\\
{[提示: 应用 Schauder 不动点定理; 见习题 6.26.]}

7. 令 $U=\dfrac{1}{q(\xi)}\sum_{j\in J}q_j(\xi)(S_{y_j}\circ T)$, 其中 $\xi$ 如第 6 问中那样.
证明 $U\in\KK(E)$.由此推出 $Z=N(I-U)$ 是有限维的; 验证 $\xi\in Z$.

8. 证明 $T(Z)\subset Z$ 并得出结论.\\
{[提示: 证明 $U\in\mathcal{A}$, 由此推出 $T\circ U=U\circ T$.]}

9. 构造没有非平凡不变子空间的线性算子 $T:\R^2\to\R^2$.

\runin{-- B --}

现在建立上述结果的一个更强版本.假设 $T\in\KK(E)$ 且 $T\ne0$.设 $R\in\LL(E)$ 满足
$R\circ T=T\circ R$.证明 $R$ 允许非平凡闭不变子空间.\\
{[提示: 令 $\mathcal{B}=\operatorname{span}\{I,R,R^2,\dots\}$ 且
$\mathcal{B}y=\{Sy;\ S\in\mathcal{B}\}$.验证第 A 部分中的所有步骤在把 $\mathcal{A}$
换成 $\mathcal{B}$, 把 $\mathcal{A}y$ 换成 $\mathcal{B}y$ 之后仍然成立.]}
\end{problem}

\begin{problem}[第 2, 4, 5, 6 章]
\label{prob:43}
\runin{正规算子}
设 $H$ 为与其对偶空间等同的 Hilbert 空间.若算子 $T\in\LL(H)$ 满足
$T\circ T^\star=T^\star\circ T$, 则称它为\emph{正规算子}.

1. 证明 $T$ 正规当且仅当它满足对一切 $u\in H$, $\abs{Tu}=\abs{T^\star u}$.\\
{[提示: 计算 $\abs{T(u+v)}^2$.]}

在整个本题的其余部分中, 我们假设 $T$ 是正规的.

2. 假设 $u\in N(T-\lambda I)$ 且 $v\in N(T-\mu I)$, 其中 $\lambda\ne\mu$.证明
$(u,v)=0$.\\
{[提示: 利用第 1 问证明 $N(T^\star-\mu I)=N(T-\mu I)$, 并计算 $(Tu,v)$.]}

3. 证明 $\overline{R(T)}=\overline{R(T^\star)}=N(T)^\perp=N(T^\star)^\perp$.

4. 设 $f\in R(T)$.验证存在 $u\in R(T^\star)$ 满足 $f=Tu$.\\
{[提示: 注意 $H=\overline{R(T)}\oplus N(T)$.]}

5. 考虑 $R(T^\star)$ 中的序列 $(u_n)$ 使得当 $n\to\infty$ 时 $u_n\to u$.写出
$u_n=T^\star y_n$ ($y_n\in H$).证明当 $n\to\infty$ 时 $Ty_n$ 收敛于极限 $z\in H$,
满足 $T^\star z=f$.\\
{[提示: 使用第 1 问与 Cauchy 序列论证.]}

6. 由此推出 $R(T)=R(T^\star)$.\\
{[提示: 利用 $N(T)=N(T^\star)$ 的事实.]}

7. 证明 $\norm{T^2}=\norm{T}^2$.\\
{[提示: 写出 $\abs{Tu}^2\le\abs{T^\star Tu}\abs{u}=\abs{T^2u}\abs{u}$.]}

8. 由此推出对每个整数 $p\ge1$, $\norm{T^p}=\norm{T}^p$.\\
{[提示: 先考虑 $p=2^k$ 的情形.对一般的整数 $p$, 选取满足 $2^k\ge p$ 的任何 $k$,
并写出 $\norm{T}^{2^k}=\norm{T^{2^k}}=\norm{T^{2^k-p}T^p}$.]}

9. 证明 $N(T^2)=N(T)$, 并由此推出对每个整数 $p\ge1$, $N(T^p)=N(T)$.\\
{[提示: 注意若 $T^2u=0$, 则 $Tu\in N(T)\cap R(T)$.]}
\end{problem}

\begin{problem}[第 5, 6 章]
\label{prob:44}
\runin{等距与酉算子. 斜自伴算子. 极分解与 Cayley 变换}
设 $H$ 为与其对偶空间等同的 Hilbert 空间, 设 $T\in\LL(H)$.称
\begin{enumerate}[label=(\roman*), itemsep=0.1em]
	\item $T$ 为等距, 若对一切 $u\in H$, $\abs{Tu}=\abs{u}$,
	\item $T$ 为酉算子, 若 $T$ 是等距且又是满射,
	\item $T$ 为斜自伴的 (或反对称的), 若 $T^\star=-T$.
\end{enumerate}

\runin{-- A --}

1. 假设 $T$ 是等距.验证 $\norm{T}=1$.

2. 证明 $T\in\LL(H)$ 是等距当且仅当 $T^\star\circ T=I$.

3. 假设 $T\in\LL(H)$ 是等距.证明下述条件等价:
(a) $T$ 是酉算子,
(b) $T^\star$ 是单射,
(c) $T\circ T^\star=I$,
(d) $T^\star$ 是等距,
(e) $T^\star$ 是酉算子.

4. 举出一个不是酉算子的等距的例子.\\
{[提示: 使用习题 6.18.]}

5. 假设 $T$ 是等距.证明 $R(T)$ 闭且 $T\circ T^\star=P_{R(T)}$ ($=R(T)$ 上的正交投影算子).

6. 假设 $T$ 是等距.证明或者 $T$ 是酉算子此时 $\sigma(T)\subset\{-1,+1\}$, 或者 $T$ 不是酉算子此时 $\sigma(T)=[-1,+1]$.

7. 假设 $T\in\KK(H)$ 是等距.证明 $\dim H<\infty$.

8. 证明 $T\in\LL(H)$ 斜自伴当且仅当对一切 $u\in H$, $(Tu,u)=0$.

9. 假设 $T\in\LL(H)$ 斜自伴.证明 $\sigma(T)\subset\{0\}$.\\
{[提示: 使用 Lax--Milgram 定理.]}

10. 假设 $T\in\LL(H)$ 斜自伴.令 $U=(T+I)\circ(T-I)^{-1}$.验证 $U$ 良定,
$U=(T-I)^{-1}\circ(T+I)$, 且 $U\circ T=T\circ U$.证明 $U$ 是酉算子 ($U$ 称为 $T$
的 \emph{Cayley 变换}).

11. 反之, 设 $T\in\LL(H)$ 满足 $1\notin\sigma(T)$.假设 $U=(T+I)\circ(T-I)^{-1}$
是等距.证明 $T$ 斜自伴.

\runin{-- B --}

若算子 $T\in\LL(H)$ 满足
\begin{equation}
\text{存在由}\ N(T)\ \text{到}\ N(T^\star)\ \text{中的等距}\ J,
\tag{1}
\label{eq:prob44.1}
\end{equation}
则称它满足性质 \eqref{eq:prob44.1}.

第 B 部分的目标是证明满足性质 \eqref{eq:prob44.1} 的每个算子 $T\in\LL(H)$ 可以
分解为 $T=U\circ P$, 其中 $U\in\LL(H)$ 是等距且 $P\in\LL(H)$ 是自伴非负算子
(回忆非负意味着对一切 $u\in H$, $(Pu,u)\ge0$).这样的分解称为 $T$ 的\emph{极分解}.
此外, $P$ 在 $H$ 上唯一确定, $U$ 在 $N(T)^\perp$ 上唯一确定 (但在 $H$ 上不唯一).

1. 验证在下述情形下假设 \eqref{eq:prob44.1} 成立:
(i) $T$ 是单射,
(ii) $\dim H<\infty$,
(iii) $T$ 正规 (见问题 43),
(iv) $T=I-K$, 其中 $K\in\KK(H)$.

2. 举出不满足 \eqref{eq:prob44.1} 的例子.\\
{[提示: 使用习题 6.18.]}

3. 假设有极分解 $T=U\circ P$.证明 $P^2=T^\star\circ T$.

4. 由此推出 $P$ 在 $H$ 上唯一确定.\\
{[提示: 使用问题 39.]}

5. 设 $T=U\circ P$ 为 $T$ 的极分解.证明 $U$ 在 $N(T)^\perp$ 上唯一确定.

6. 假设 $T$ 允许极分解.证明 \eqref{eq:prob44.1} 成立.\\
{[提示: 令 $J=U|_{N(T)}$.]}

7. 证明满足 \eqref{eq:prob44.1} 的每个算子 $T\in\LL(H)$ 允许极分解.

8. 假设 $T$ 满足更强的假设
\begin{equation}
\text{存在由}\ N(T)\ \text{到}\ N(T^\star)\ \text{上的等距}\ J.
\tag{2}
\label{eq:prob44.2}
\end{equation}
证明 $T$ 允许极分解 $T=U\circ P$, 其中 $U$ 是酉算子.

9. 由此推出每个正规算子 $T\in\LL(H)$ 允许极分解 $T=U\circ P$, 其中 $U$ 是酉算子且
$U\circ P=P\circ U$.

10. 证明满足 \eqref{eq:prob44.2} 的每个算子 $T\in\LL(H)$ 可以分解为 $T=P\circ U$,
其中 $U\in\LL(H)$ 是酉算子, $P\in\LL(H)$ 是自伴非负算子.\\
{[提示: 对 $T^\star$ 应用第 8 问.]}

11. 证明满足 \eqref{eq:prob44.1} 的每个算子 $T\in\KK(H)$ 允许极分解 $T=U\circ P$,
其中 $P\in\KK(H)$.

12. 假设 $H$ 可分且 $T\in\KK(H)$ ($T$ 不必满足 \eqref{eq:prob44.1}).证明存在 $H$ 的
两个标准正交基 $(e_n)$ 与 $(f_n)$ 使得对一切 $u\in H$,
\begin{equation*}
Tu=\sum_{n=1}^\infty\alpha_n(u,e_n)f_n,
\end{equation*}
其中 $(\alpha_n)$ 是满足对一切 $n$ 有 $\alpha_n\ge0$ 且当 $n\to\infty$ 时 $\alpha_n\to0$
的序列.计算 $T^\star$.反之, 证明这种形式的任何算子必为紧算子.
\end{problem}

\begin{problem}[第~
\ref{chap:8}~章]
\label{prob:45}
\runin{强极大值原理}
考虑双线性形式
\begin{equation*}
a(u,v)=\int_0^1(pu'v'+quv),
\end{equation*}
其中 $p\in C^1([0,1])$, 在 $(0,1)$ 上 $p\ge\alpha>0$, 且 $q\in C([0,1])$.我们假设
$a$ 在 $H_0^1(0,1)$ 上强制 (但不对 $q$ 作任何符号假设).

给定 $f\in L^2(0,1)$, 设 $u\in H^2(0,1)$ 为
\begin{equation}
\begin{cases}
-(pu')'+qu=f & \text{在}\ (0,1)\ \text{上},\\
u(0)=u(1)=0
\end{cases}
\label{eq:prob45.1}
\end{equation}
的解.假设在 $(0,1)$ 上几乎处处 $f\ge0$ 且 $f\not\equiv0$.我们的目标是证明
\begin{equation}
u'(0)>0,\quad u'(1)<0\quad\text{且在}\ (0,1)\ \text{上}\ u(x)>0.
\label{eq:prob45.2}
\end{equation}

1. 假设 $\psi\in H^1(0,1)$ 满足对一切在 $(0,1)$ 上满足 $v\ge0$ 的 $v\in H_0^1(0,1)$,
$a(\psi,v)\le0$, 且 $\psi(0)\le0,\psi(1)\le0$.证明在 $(0,1)$ 上 $\psi\le0$.\\
{[提示: 取 $v=\psi^+$ 并使用习题 8.11.]}

考虑问题
\begin{equation}
\begin{cases}
-(p\zeta')'+q\zeta=0 & \text{在}\ (0,1)\ \text{上},\\
\zeta(0)=0,\quad\zeta(1)=1.
\end{cases}
\label{eq:prob45.5}
\end{equation}

2. 证明 \eqref{eq:prob45.5} 有唯一解 $\zeta$, 且在 $(0,1)$ 上 $\zeta\ge0$.

3. 验证在 $(0,1)$ 上 $u\ge0$, 并由此推出 $u'(0)\ge0$ 且 $u'(1)\le0$.

4. 证明 $p(1)\abs{u'(1)}=\int_0^1f\zeta$.\\
{[提示: 把 \eqref{eq:prob45.1} 乘以 $\zeta$, 把 \eqref{eq:prob45.5} 乘以 $u$.]}

令 $\phi(x)=(e^{Bx}-1)$, $B>0$.

5. 验证若 $B$ 充分大 (只依赖于 $p$ 与 $q$), 则在 $(0,1)$ 上
$-(p\phi')'+q\phi\le0$.在下文中固定使此式成立的 $B$.

6. 设 $A=(e^B-1)^{-1}$.证明在 $(0,1)$ 上 $\zeta\ge A\phi$.\\
{[提示: 对 $\psi=A\phi-\zeta$ 应用第 1 问.]}

7. 由此推出 $u'(1)<0$.\\
{[提示: 应用第 4 问.]}

8. 验证 $u'(0)>0$.\\
{[提示: 将 $t$ 换为 $1-t$.]}

9. 固定充分小的 $\delta\in(0,1/2)$ 使得对一切 $x\in(0,\delta)$, $\dfrac{u(x)}{x}\ge\frac12u'(0)$,
且对一切 $x\in(1-\delta,1)$, $\dfrac{u(x)}{1-x}\ge\frac12\abs{u'(1)}$.为什么这样的 $\delta$
存在? 设 $v$ 为在 $(\delta,1-\delta)$ 上 $-(pv')'+qv=0$, $v(\delta)=v(1-\delta)=\gamma$
的解, 其中 $\gamma=\dfrac\delta2\min\{u'(0),\abs{u'(1)}\}$.证明在 $(\delta,1-\delta)$
上 $u\ge v\ge0$.

10. 证明在 $(\delta,1-\delta)$ 上 $v>0$.\\
{[提示: 反设对某个 $x_0\in(\delta,1-\delta)$ 有 $v(x_0)=0$, 像习题 8.33 那样应用定理
 \ref{thm:7.3} (Cauchy--Lipschitz--Picard).]}

11. 由此推出对一切 $x\in(0,1)$, $u(x)>0$.

最后, 我们给出强极大值原理的一个更强的形式.

12. 证明存在常数 $a>0$ (只依赖于 $p$ 与 $q$) 使得
\begin{equation*}
u(x)\ge ax(1-x)\int_0^1f(t)t(1-t)dt.
\end{equation*}
[提示: 先从 $p\equiv1$ 且 $q\equiv k^2$ (正实数) 的情形开始; 使用 \eqref{eq:prob45.1}
的显式解.接着考虑 $p\equiv1$ 且不对 $q$ 作进一步假设的情形.最后通过变量代换把一般
情形归结为前述情形.]
\end{problem}

\begin{problem}[第~
\ref{chap:8}~章]
\label{prob:46}
\runin{上下解方法}
设 $h(t):[0,+\infty)\to[0,+\infty)$ 为连续单调不减函数.假设存在两个函数
$v,w\in C^2([0,1])$ 满足
\begin{equation}
\begin{cases}
0\le v\le w & \text{在}\ I=(0,1)\ \text{上},\\
-v''+v\le h(v) & \text{在}\ I\ \text{上},\\
v(0)=v(1)=0,\\
-w''+w\ge h(w) & \text{在}\ I\ \text{上},\\
w(0)\ge0,\quad w(1)\ge0
\end{cases}
\label{eq:prob46.1}
\end{equation}
($v$ 称为下解, $w$ 称为上解).目标是证明问题
\begin{equation}
\begin{cases}
-u''+u=h(u) & \text{在}\ I\ \text{上},\\
u(0)=u(1)=0,\\
v\le u\le w & \text{在}\ I\ \text{上}
\end{cases}
\label{eq:prob46.2}
\end{equation}
存在解 $u\in C^2([0,1])$.

考虑从 $u_0=w$ 开始由下式归纳定义的序列 $(u_n)_{n\ge1}$:
\begin{equation}
\begin{cases}
-u''_n+u_n=h(u_{n-1}) & \text{在}\ I\ \text{上},\ n\ge1,\\
u_n(0)=u_n(1)=0.
\end{cases}
\label{eq:prob46.3}
\end{equation}

1. 证明在 $I$ 上 $v\le u_1\le w$.\\
{[提示: 对 $(u_1-w)$ 与 $(u_1-v)$ 应用极大值原理.]}

2. 由归纳法证明对每个 $n\ge1$, 在 $I$ 上 $v\le u_n$ 且 $u_{n+1}\le u_n$.

3. 由此推出序列 $(u_n)$ 在 $L^2(I)$ 中收敛于极限 $u$, 且在 $L^2(I)$ 中 $h(u_n)\to h(u)$.

4. 证明 $u\in H_0^1(I)$, 且对一切 $\phi\in H_0^1(I)$,
\begin{equation*}
\int_0^1u'\phi'+\int_0^1u\phi=\int_0^1h(u)\phi.
\end{equation*}

5. 得出 $u\in C^2([0,1])$ 是 \eqref{eq:prob46.2} 的古典解.

在下文中我们选取 $h(t)=t^\alpha$, 其中 $0<\alpha<1$.目标是证明存在唯一的函数
$u\in C^2([0,1])$ 满足
\begin{equation}
\begin{cases}
-u''+u=u^\alpha & \text{在}\ I\ \text{上},\\
u(0)=u(1)=0,\\
u(x)>0 & \forall x\in I.
\end{cases}
\label{eq:prob46.4}
\end{equation}

6. 设 $v(x)=\varepsilon\sin(\pi x)$ 且 $w(x)\equiv1$.证明若 $\varepsilon$ 充分小,
假设 \eqref{eq:prob46.1} 满足.由此推出 \eqref{eq:prob46.4} 存在解.

现在转向唯一性问题.设 $u$ 为由上述方法从 $u_0\equiv1$ 出发得到的 \eqref{eq:prob46.4}
的解.设 $\tilde u\in C^2([0,1])$ 为 \eqref{eq:prob46.4} 的另一个解.

7. 证明在 $I$ 上 $\tilde u\le1$.\\
{[提示: 考虑 $\tilde u$ 达到其最大值的点 $x_0\in[0,1]$.]}

8. 证明从 $u_0\equiv1$ 开始由 \eqref{eq:prob46.3} 定义的序列 $(u_n)_{n\ge1}$ 满足在
$I$ 上 $\tilde u\le u_n$, 并由此推出在 $I$ 上 $\tilde u\le u$.

9. 证明 $\int_0^1(\tilde u^\alpha u-u^\alpha\tilde u)=0$.

10. 得出在 $I$ 上 $\tilde u=u$.\\
{[提示: 写出 $\tilde u^\alpha u-u^\alpha\tilde u=u\tilde u(\tilde u^{\alpha-1}-u^{\alpha-1})$,
并注意 $u^{\alpha-1}\le\tilde u^{\alpha-1}$.]}

我们现在给出存在性的另一种证明.对每个 $u\in H_0^1(I)$, 令
\begin{equation*}
F(u)=\frac12\int_0^1(u'^2+u^2)-\int_0^1g(u),
\end{equation*}
其中 $g(t)=\frac{1}{\alpha+1}(t^+)^\alpha+1$, $0<\alpha<1$, $t^+=\max(t,0)$.

11. 证明存在常数 $C$ 使得对一切 $u\in H_0^1(I)$,
$F(u)\ge\frac12\norm{u}_{H^1}^2-C\norm{u}_{H^1}^{\alpha+1}$.

12. 由此推出 $m=\inf_{v\in H_0^1(I)}F(v)>-\infty$, 且下确界可以达到.\\
{[提示: 设 $(u_n)$ 为极小化序列.验证子序列 $(u_{n_k})$ 在 $H_0^1(I)$ 中弱收敛于极限 $u$,
且 $\int_0^1g(u_{n_k})\to\int_0^1g(u)$.提醒读者注意泛函 $F$ 不是凸的; 为什么?]}

13. 证明 $m<0$.\\
{[提示: 证明对满足 $v^+\not\equiv0$ 的一切 $v\in H_0^1(I)$ 及充分小的 $\varepsilon$,
$F(\varepsilon v)<0$.]}

14. 验证对一切 $a,b\in\R$, $g(b)-g(a)\ge(a^+)^\alpha(b-a)$.

15. 设 $u\in H_0^1(I)$ 为 $F$ 在 $H_0^1(I)$ 上的极小元.证明对一切 $v\in H_0^1(I)$,
\begin{equation*}
\int_0^1(u'v'+uv)=\int_0^1(u^+)^\alpha v.
\end{equation*}
[提示: 写出 $F(u)\le F(u+tv)$, 应用第 14 问, 并令 $t\to0$.]

16. 由此推出 $u\in C^2([0,1])$ 是 $-u''+u=(u^+)^\alpha$ 在 $I$ 上, $u(0)=u(1)=0$
的解.证明在 $I$ 上 $u\ge0$ 且 $u\not\equiv0$.

17. 利用强极大值原理得出在 $I$ 上 $u>0$ (见问题 45).
\end{problem}
\begin{problem}[第~
\ref{chap:8}~章]
\label{prob:47}
\runin{Poincaré--Wirtinger 不等式}
设 $I=(0,1)$.

\runin{-- A --}

1. 证明对一切 $u\in W^{1,1}(I)$,
\begin{equation}
\norm{u-\bar u}_{L^\infty(I)}\le\norm{u'}_{L^1(I)},
\label{eq:prob47.1}
\end{equation}
其中 $\bar u=\int_I u$.\\
{[提示: 注意对某个 $x_0\in[0,1]$, $\bar u=u(x_0)$.]}

2. 证明 \eqref{eq:prob47.1} 中的常数 $1$ 是最优的, 即
\begin{equation}
\sup\{\norm{u-\bar u}_{L^\infty(I)};\ u\in W^{1,1}(I)\ \text{且}\ \norm{u'}_{L^1(I)}=1\}=1.
\label{eq:prob47.2}
\end{equation}
[提示: 考虑 $[0,1]$ 上的光滑函数序列 $(u_n)$, 满足在 $(0,1)$ 上对一切 $n$ 有
$u'_n\ge0$, $u_n(1)=1$, 以及在 $[0,1-1/n]$ 上 $u_n(x)=0$.]

3. 证明 \eqref{eq:prob47.2} 中的上确界达不到, 即不存在函数 $u\in W^{1,1}(I)$
使得 $\norm{u-\bar u}_{L^\infty(I)}=1$ 且 $\norm{u'}_{L^1(I)}=1$.

4. 证明对一切 $u\in W_0^{1,1}(I)$,
\begin{equation}
\norm{u}_{L^\infty(I)}\le\frac12\norm{u'}_{L^1(I)}.
\label{eq:prob47.3}
\end{equation}
[提示: 写出 $\abs{u(x)-u(0)}\le\int_0^x\abs{u'(t)}dt$ 以及
$\abs{u(x)-u(1)}\le\int_x^1\abs{u'(t)}dt$.]

5. 证明 $1/2$ 是 \eqref{eq:prob47.3} 中的最佳常数.它能达到吗? \\
{[提示: 固定 $a\in(0,1)$, 考虑在 $(0,a)$ 上递增, 在 $(a,1)$ 上递减且 $u(a)=1$ 的函数 $u\in W_0^{1,1}(I)$.]}

6. 由此推出下述不等式成立:
\begin{equation}
\norm{u-\bar u}_{L^q(I)}\le C\norm{u'}_{L^p(I)}\quad\forall u\in W^{1,p}(I),
\label{eq:prob47.4}
\end{equation}
以及
\begin{equation}
\norm{u}_{L^q(I)}\le C\norm{u'}_{L^p(I)}\quad\forall u\in W_0^{1,p}(I),
\label{eq:prob47.5}
\end{equation}
其中 $1\le q\le\infty$ 且 $1\le p\le\infty$.证明当 $1\le q\le\infty$ 且 $1<p\le\infty$ 时,
\eqref{eq:prob47.4} 与 \eqref{eq:prob47.5} 中的最佳常数可以达到.\\
{[提示: 在满足 $\norm{u-\bar u}_{L^q(I)}=1$ 的 $u\in W^{1,p}(I)$ 类 (相应地
$u\in W_0^{1,p}(I)$ 且 $\norm{u}_{L^q(I)}=1$) 中极小化 $\norm{u'}_{L^p(I)}$.]}

\runin{-- B --}

接下来的目标是求 $p=q=2$ 时 \eqref{eq:prob47.4} 中的最佳常数, 即
\begin{equation}
\norm{u-\bar u}_{L^2(I)}\le C\norm{u'}_{L^2(I)}\quad\forall u\in H^1(I).
\label{eq:prob47.6}
\end{equation}
令 $H=\{f\in L^2(I);\ \int_I f=0\}$ 与 $V=\{v\in H^1(I);\ \int_I v=0\}$.

1. 验证对每个 $f\in H$, 存在唯一的 $u\in V$ 使得对一切 $v\in V$, $\int_I u'v'=\int_I fv$.

2. 证明 $u\in H^2(I)$ 且满足在 $I$ 上几乎处处 $-u''=f$, $u'(0)=u'(1)=0$.

3. 证明由 $Tf=u$ 定义的算子 $T:H\to H$ 是自伴紧算子, 且对一切 $f\in H$, $\int_I fTf\ge0$.

4. 设 $\lambda_1$ 为 $T$ 的最大特征值.证明 \eqref{eq:prob47.6} 在 $C=\sqrt{\lambda_1}$
下成立, 且 $\sqrt{\lambda_1}$ 是 \eqref{eq:prob47.6} 中的最佳常数.\\
{[提示: 使用习题 6.24.]}

5. 显式计算 \eqref{eq:prob47.6} 中的最佳常数.

\runin{-- C --}

1. 证明对一切 $u\in W^{1,1}(I)$,
\begin{equation}
\norm{u-\bar u}_{L^1(I)}\le2\int_I\abs{u'(t)}t(1-t)dt.
\label{eq:prob47.7}
\end{equation}

2. 由此推出对一切 $u\in W^{1,1}(I)$,
\begin{equation}
\norm{u-\bar u}_{L^1(I)}\le\frac12\norm{u'}_{L^1(I)}.
\label{eq:prob47.8}
\end{equation}

3. 证明 \eqref{eq:prob47.8} 中的常数 $1/2$ 是最优的, 即
\begin{equation}
\sup\{\norm{u-\bar u}_{L^1(I)};\ u\in W^{1,1}(I)\ \text{且}\ \norm{u'}_{L^1(I)}=1\}=\frac12.
\label{eq:prob47.9}
\end{equation}

4. \eqref{eq:prob47.9} 中的上确界能达到吗?
\end{problem}

\begin{problem}[第~
\ref{chap:8}~章]
\label{prob:48}
\runin{一个非线性问题}
设 $j:[-1,+1]\to[0,+\infty)$ 为连续凸函数, 满足 $j\in C^2((-1,+1))$, $j(0)=0$,
$j'(0)=0$, 且
\begin{equation*}
\lim_{t\uparrow+1}j'(t)=+\infty,\qquad\lim_{t\downarrow-1}j'(t)=-\infty.
\end{equation*}
(一个应牢记的好例子是 $j(t)=1-\sqrt{1-t^2}$, $t\in[-1,+1]$.) 给定 $f\in L^2(0,1)$,
定义函数 $\phi:H_0^1(0,1)\to(-\infty,+\infty]$:
\begin{equation*}
\phi(v)=\begin{cases}
\dfrac12\int_0^1v'^2+\int_0^1j(v)-\int_0^1fv & \text{若}\ v\in H_0^1(0,1)\ \text{且}\ \norm{v}_{L^\infty}\le1,\\
+\infty & \text{其他}.
\end{cases}
\end{equation*}

1. 验证 $\phi$ 在 $H_0^1(0,1)$ 上是凸的, 下半连续的, 且 $\lim_{\norm{v}_{H_0^1}\to+\infty}\phi(v)=+\infty$.

2. 由此推出存在唯一的 $u\in H_0^1(0,1)$ 使得 $\phi(u)=\min_{v\in H_0^1(0,1)}\phi(v)$.

目标是证明若 $f\in L^\infty(0,1)$, 则 $\norm{u}_{L^\infty(0,1)}<1$, $u\in H^2(0,1)$,
且 $u$ 满足
\begin{equation}
\begin{cases}
-u''+j'(u)=f & \text{在}\ (0,1)\ \text{上},\\
u(0)=u(1)=0.
\end{cases}
\label{eq:prob48.1}
\end{equation}

3. 验证对一切 $t\in[-1,+1]$ 与一切 $a\in(-1,+1)$, $j(t)-j(a)\ge j'(a)(t-a)$.\\
{[提示: 使用 $j$ 的凸性.]}

固定 $a\in[0,1)$.

4. 令 $v=\min(u,a)$.证明 $v\in H_0^1(0,1)$ 且在 $[u\le a]$ 上几乎处处 $v'=u'$,
在 $[u>a]$ 上几乎处处 $v'=0$.\\
{[提示: 写出 $v=a-(a-u)^+$ 并使用习题 8.11.]}

5. 证明 $\frac12\int_{[u>a]}u'^2\le\int_{[u>a]}(f-j'(a))(u-a)$.\\
{[提示: 写出 $\phi(u)\le\phi(v)$, 其中 $v$ 在第 4 问中定义.然后使用第 3 问.]}

6. 选取 $a\in[0,1)$ 使得对一切 $x\in[0,1]$, $f(x)\le j'(a)$, 并证明对一切 $x\in[0,1]$,
$u(x)\le a$.\\
{[提示: 证明 $\int_0^1w'^2=0$, 其中 $w=(u-a)^+$ 属于 $H_0^1(0,1)$; 为什么?]}

7. 得出结论: $\norm{u}_{L^\infty(0,1)}<1$.\\
{[提示: 应用前面的论证, 将 $u$ 换为 $-u$, $j(t)$ 换为 $j(-t)$, $f$ 换为 $-f$.]}

8. 由此推出 $u$ 属于 $H^2(0,1)$ 且满足 \eqref{eq:prob48.1}.\\
{[提示: 写出 $\phi(u)\le\phi(u+\varepsilon v)$, 其中 $v\in H_0^1(0,1)$, $\varepsilon$ 很小.]}

9. 验证若 $f\in C([0,1])$, 则 $u\in C^2([0,1])$.

10. 反之, 证明满足 $\norm{u}_{L^\infty(0,1)}<1$ 且满足 \eqref{eq:prob48.1} 的任何函数
$u\in C^2([0,1])$ 是 $\phi$ 在 $H_0^1(0,1)$ 上的极小元.\\
{[提示: 在第 3 问中取 $t=v(x)$ 与 $a=u(x)$.]}

现在假设 $f\in L^2(0,1)$.令 $f_n=T_nf$, 其中 $T_n$ 是截断算子 (定义在第~
\ref{chap:4}~章定理~
\ref{thm:4.12}
之后).设 $u_n$ 为对应于 $f_n$ 的 \eqref{eq:prob48.1} 的解.

11. 证明当 $n\to\infty$ 时 $\norm{j'(u_n)}_{L^2(0,1)}\le C$.\\
{[提示: 把 \eqref{eq:prob48.1} 乘以 $j'(u_n)$.]}

12. 由此推出 $\norm{u_n}_{H^2(0,1)}\le C$.

13. 证明某个子序列 $(u_{n_k})$ 在 $H^2(0,1)$ 中弱收敛于极限 $u\in H^2(0,1)$, 满足在
$C^1([0,1])$ 中 $u_{n_k}\to u$.证明在 $(0,1)$ 上几乎处处 $\abs{u(x)}<1$, 且
$j'(u)\in L^2(0,1)$.\\
{[提示: 对序列 $j'(u_{n_k})^2$ 应用 Fatou 引理.]}

14. 证明 $j'(u_{n_k})$ 在 $L^2(0,1)$ 中弱收敛于 $j'(u)$, 并由此推出 \eqref{eq:prob48.1} 成立.\\
{[提示: 应用习题 4.16.]}

15. 若额外假设
\begin{equation*}
\liminf_{t\uparrow1}j'(t)(1-t)^{1/3}>0\quad\text{且}\quad\limsup_{t\downarrow-1}j'(t)(1+t)^{1/3}<0,
\end{equation*}
由此推出 $\norm{u}_{L^\infty(0,1)}<1$.\\
{[提示: 反设对某个 $x_0\in(0,1)$ 有 $u(x_0)=1$.验证对一切 $x\in(0,1)$,
$\abs{u'(x)}\le\abs{x-x_0}^{1/2}\norm{u''}_{L^2}$ 以及
$\abs{u(x)-1}\le\frac23\abs{x-x_0}^{3/2}\norm{u''}_{L^2}$.由此推出 $j'(u)\notin L^2(0,1)$.]}
\end{problem}

\begin{problem}[第~
\ref{chap:8}~章]
\label{prob:49}
\runin{Sturm--Liouville 算子特征值的极小化极大原理}
考虑 $(0,1)$ 上的 Sturm--Liouville 算子 $Au=-(pu')'+qu$, 带 Dirichlet 边界条件
$u(0)=u(1)=0$.假设 $p\in C^1([0,1])$, 对一切 $x\in[0,1]$ 有 $p(x)\ge\alpha>0$,
且 $q\in C([0,1])$.令
\begin{equation*}
a(u,v)=\int_0^1(pu'v'+quv)\quad\forall u,v\in H_0^1(0,1).
\end{equation*}
注意我们不对 $q$ 作进一步的假设, 从而双线性形式 $a$ 不必是强制的.固定充分大的 $M$
使得 $\tilde a(u,v)=a(u,v)+M\int_0^1uv$ 是强制的 (例如 $M>-\min_{x\in[0,1]}q(x)$).
设 $(\lambda_k)$ 为 $A$ 的特征值序列.空间 $H=H_0^1(0,1)$ 配以标量积 $\tilde a(u,v)$
(记作 $(u,v)_H$) 以及相应的范数 $\abs{u}_H=\tilde a(u,u)^{1/2}$.给定任何 $f\in L^2(0,1)$,
设 $u\in H_0^1(0,1)$ 为问题 $\tilde a(u,v)=\int_0^1fv$ ($\forall v\in H_0^1(0,1)$) 的
唯一解.令 $u=Tf$, 并把 $T$ 看作由 $H$ 到自身中的算子.

1. 证明 $T$ 自伴且紧.\\
{[提示: 回忆由 $H$ 到 $L^2(0,1)$ 中的恒等映射是紧的.]}

2. 设 $(\lambda_k)$ 为 $A$ 的特征值序列 (在定理  \ref{thm:8.22} 的意义下), 相应的特征
函数为 $(e_k)$, 设 $(\mu_k)$ 为 $T$ 的特征值序列.验证对一切 $k$, $\mu_k>0$, 并证明
对一切 $k$, $\lambda_k=\dfrac{1}{\mu_k}-M$ 且 $T(e_k)=\mu_ke_k$.

3. 证明对一切 $w\in H$, $(Tw,w)_H=\int_0^1w^2$, 并由此推出对一切 $w\in H$ ($w\ne0$),
\begin{equation*}
\frac{1}{R(w)}=\frac{a(w,w)}{\int_0^1w^2}+M,
\end{equation*}
其中 $R$ 是与 $T$ 相关联的 Rayleigh 商, 即 $R(w)=\dfrac{(Tw,w)_H}{\abs{w}_H^2}$
(见问题 37).

4. 证明
\begin{equation}
\lambda_1=\min_{\substack{w\in H_0^1\\w\ne0}}\left\{\frac{a(w,w)}{\int_0^1w^2}\right\},
\label{eq:prob49.1}
\end{equation}
且对一切 $k\ge2$,
\begin{equation*}
\lambda_k=\min\left\{\frac{a(w,w)}{\int_0^1w^2};\ w\in H_0^1(0,1),\ w\ne0,\ \int_0^1we_j=0\ \forall j=1,\dots,k-1\right\}.
\end{equation*}
[提示: 应用问题 37 第 2 问并证明 $(w,e_j)_H=0\iff\int_0^1we_j=0$.]

5. 证明对一切 $k\ge1$,
\begin{equation*}
\lambda_k=\min_{\substack{\Sigma\subset H_0^1(0,1)\\\dim\Sigma=k}}\max_{\substack{u\in\Sigma\\u\ne0}}\left\{\frac{a(u,u)}{\int_0^1u^2}\right\},
\end{equation*}
以及
\begin{equation*}
\lambda_k=\max_{\substack{M\subset H_0^1(0,1)\\\operatorname{codim}M=k-1}}\min_{\substack{u\in M\\u\ne0}}\left\{\frac{a(u,u)}{\int_0^1u^2}\right\},
\end{equation*}
其中 $\Sigma$ 与 $M$ 是 $H_0^1(0,1)$ 的闭子空间.\\
{[提示: 应用问题 37 第 7 问.]}

6. 对带 Neumann 边界条件的 Sturm--Liouville 算子证明类似的结果.

现在我们回到公式 \eqref{eq:prob49.1} 并讨论对应于第一个特征值 $\lambda_1$ 的特征函数
的进一步性质.特别地, 我们将看到存在一个生成与 $\lambda_1$ 相关联的特征子空间的正特征函数.

7. 设 $w_0\in H_0^1(0,1)$ 为 \eqref{eq:prob49.1} 的满足 $\int_0^1w_0^2=1$ 的极小元.
证明在 $(0,1)$ 上 $Aw_0=\lambda_1w_0$.

8. 令 $w_1=\abs{w_0}$.验证 $w_1$ 也是 \eqref{eq:prob49.1} 的极小元, 并由此推出在
$(0,1)$ 上 $Aw_1=\lambda_1w_1$.\\
{[提示: 使用习题 8.11.]}

9. 证明在 $(0,1)$ 上 $w_1>0$, $w'_1(0)>0$, 且 $w'_1(1)<0$.\\
{[提示: 对算子 $A+M$ 应用强极大值原理; 见问题 45.]}

10. 假设 $w\in H_0^1(0,1)$ 在 $(0,1)$ 上满足 $Aw=\lambda_1w$.证明 $w$ 是 $w_1$ 的倍数.\\
{[提示: 回忆特征值是单的; 见习题 8.33.找出不依赖特征值单性的另一种证明; 在 (2) 中使用
$w^2/w_1$ 作为检验函数.]}

11. 证明在 $(0,1)$ 上满足 $A\psi=\mu\psi$, $\psi\ge0$ 且 $\int_0^1\psi^2=1$ ($\mu\in\R$)
的任何函数 $\psi\in H_0^1(0,1)$ 必与 $w_1$ 重合.\\
{[提示: 若 $\mu\ne\lambda_1$, 验证 $\int_0^1\psi w_1=0$.由此推出 $\mu=\lambda_1$.]}
\end{problem}

\begin{problem}[第~
\ref{chap:8}~章]
\label{prob:50}
\runin{另一个非线性问题}
设 $q\in C([0,1])$ 并考虑双线性形式 $a(u,v)=\int_0^1(u'v'+quv)$ ($u,v\in H_0^1(0,1)$).
假设存在 $v_1\in H_0^1(0,1)$ 使得
\begin{equation}
a(v_1,v_1)<0.
\label{eq:prob50.1}
\end{equation}

1. 验证假设 \eqref{eq:prob50.1} 等价于 $\lambda_1(A)<0$, 其中 $\lambda_1(A)$ 是带零
Dirichlet 条件的算子 $Au=-u''+qu$ 的第一个特征值.

2. 验证
\begin{equation}
-\infty<m=\inf_{u\in H_0^1(0,1)}\left\{\frac12a(u,u)+\frac14\int_0^1\abs{u}^4\right\}<0.
\label{eq:prob50.3}
\end{equation}
[提示: 使用充分小的 $\varepsilon>0$ 的 $u=\varepsilon v_1$.]

3. 证明 \eqref{eq:prob50.3} 中的下确界由某个 $u_0$ 达到.\\
{[警告: \eqref{eq:prob50.3} 中的泛函不是凸的; 为什么?]}

我们的目标是证明 \eqref{eq:prob50.3} 恰好有两个极小元.

4. 证明 $u_0$ 属于 $C^2([0,1])$ 且满足
\begin{equation}
\begin{cases}
-u''+qu+u^3=0 & \text{在}\ (0,1)\ \text{上},\\
u(0)=u(1)=0.
\end{cases}
\label{eq:prob50.4}
\end{equation}

5. 令 $u_1=\abs{u_0}$.证明 $u_1$ 也是 \eqref{eq:prob50.3} 的极小元.由此推出 $u_1$
满足 \eqref{eq:prob50.4}.\\
{[提示: 应用习题 8.11.]}

6. 证明对一切 $x\in(0,1)$, $u_1(x)>0$, $u'_1(0)>0$, 且 $u'_1(1)<0$.\\
{[提示: 选取充分大的常数 $a$ 使得 $-u''_1+a^2u_1=f\ge0,f\not\equiv0$.然后使用强极大值原理.]}

7. 设 $u_0\in H_0^1(0,1)$ 再次为 \eqref{eq:prob50.3} 中的任意极小元.证明或者对一切
$x\in(0,1)$ 有 $u_0(x)>0$, 或者对一切 $x\in(0,1)$ 有 $u_0(x)<0$.\\
{[提示: 验证对一切 $x\in(0,1)$, $\abs{u_0(x)}>0$.]}

8. 设 $U_1$ 为满足在 $[0,1]$ 上 $U_1\ge0$ 且 $U_1\not\equiv0$ 的 \eqref{eq:prob50.4}
的任何解.令 $\rho_1=U_1^2$.考虑在集合
\begin{equation*}
K=\{\rho\in H_0^1(0,1);\ \rho\ge0\ \text{在}\ (0,1)\ \text{上且}\ \sqrt{\rho}\in H_0^1(0,1)\}
\end{equation*}
上定义的泛函 $\Phi(\rho)=\int_0^1\left[\abs{(\sqrt{\rho})'}^2+q\rho+\frac12\rho^2\right]$.
证明对一切 $\rho\in K$,
\begin{equation}
\Phi(\rho)-\Phi(\rho_1)\ge\frac12\int_0^1(\rho-\rho_1)^2.
\label{eq:prob50.5}
\end{equation}
[提示: 设 $u\in C_c^1((0,1))$.注意在 $(0,1)$ 上
$\dfrac{2U'_1uu'}{U_1}\le u'^2+\dfrac{U_1'^2u^2}{U_1^2}$, 由此推出 (利用分部积分)
对一切 $u\in H_0^1(0,1)$, $\int_0^1(u'^2-U_1'^2)\ge-\int_0^1\dfrac{U''_1}{U_1}(u^2-U_1^2)$.
然后应用方程 \eqref{eq:prob50.4} 建立 \eqref{eq:prob50.5}.]

9. 由此推出存在 \eqref{eq:prob50.4} 的恰好一个在 $[0,1]$ 上满足 $u\ge0$ 的非平凡解.
记作 $U_0$.\\
{[评注: 一般说来 \eqref{eq:prob50.4} 存在许多变号解.]}

10. 证明 \eqref{eq:prob50.3} 恰好存在两个极小元: $U_0$ 与 $-U_0$.
\end{problem}

\begin{problem}[第~
\ref{chap:8}~章]
\label{prob:51}
\runin{调和振子. Hermite 多项式}
设 $p\in C(\R)$ 满足在 $\R$ 上 $p\ge0$.考虑空间
\begin{equation*}
V=\left\{v\in H^1(\R);\ \int_{-\infty}^{+\infty}pv^2<\infty\right\},
\end{equation*}
它配以标量积 $(u,v)_V=\int_{-\infty}^{+\infty}(u'v'+uv+puv)$ 以及相应的范数
$\abs{u}_V=(u,u)_V^{1/2}$.

1. 验证 $V$ 是可分 Hilbert 空间.

2. 证明 $C_c^\infty(\R)$ 在 $V$ 中稠密.\\
{[提示: 设 $\zeta_n$ 为定理  \ref{thm:8.7} 证明中的截断函数序列.给定 $u\in V$, 考虑
$\zeta_nu$, 然后使用卷积.]}

考虑双线性形式 $a(u,v)=\int_{-\infty}^{+\infty}(u'v'+puv)$, $u,v\in V$.
在下文中我们假设存在常数 $\delta>0$ 与 $A>0$ 使得
\begin{equation}
p(x)\ge\delta\quad\forall x\in\R\ \text{满足}\ \abs{x}\ge A.
\label{eq:prob51.1}
\end{equation}

3. 证明 $a$ 在 $V$ 上是强制的.由此推出对每个 $f\in L^2(\R)$, 存在问题
\begin{equation}
a(u,v)=\int_{-\infty}^{+\infty}fv\quad\forall v\in V
\label{eq:prob51.2}
\end{equation}
的唯一解 $u\in V$.

4. 假设 $f\in L^2(\R)\cap C(\R)$, 证明 $u$ 满足
\begin{equation}
\begin{cases}
u\in C^2(\R),\\
-u''+pu=f & \text{在}\ \R\ \text{上},\\
\lim_{\abs{x}\to\infty}u(x)=0.
\end{cases}
\label{eq:prob51.3}
\end{equation}

5. 反之, 证明 \eqref{eq:prob51.3} 的任何解 $u$ 属于 $V$ 且满足 \eqref{eq:prob51.2}.\\
{[提示: 把方程 $-u''+pu=f$ 乘以 $\zeta_n^2u$ 并使用 $a$ 强制的事实.]}

在下文中假设
\begin{equation}
\lim_{\abs{x}\to\infty}p(x)=+\infty.
\label{eq:prob51.4}
\end{equation}

6. 给定 $f\in L^2(\R)$, 令 $u=Tf$, 其中 $u$ 是 \eqref{eq:prob51.2} 的解.证明
$T:L^2(\R)\to L^2(\R)$ 自伴且紧.\\
{[提示: 使用推论  \ref{cor:4.27} 验证 $V\subset L^2(\R)$ 紧嵌入.]}

7. 由此推出存在满足当 $n\to\infty$ 时 $\lambda_n\to\infty$ 的正数序列 $(\lambda_n)$,
以及满足下式的 $L^2(\R)$ 的 Hilbert 基 $(e_n)$:
\begin{equation}
\begin{cases}
e_n\in V\cap C^2(\R),\\
-e''_n+pe_n=\lambda_ne_n & \text{在}\ \R\ \text{上}.
\end{cases}
\label{eq:prob51.5}
\end{equation}

在下文中取 $p(x)=x^2$.

8. 验证 \eqref{eq:prob51.5} 允许形如 $e_n(x)=e^{-x^2/2}P_n(x)$ 的解, 其中 $\lambda_n=2n+1$,
$P_n(x)$ 是 $n$ 次多项式.
\end{problem}

% chapters/partial_solutions.tex
% 部分问题解答 (Partial Solutions of the Problems)
% 源: 原书 PDF 第 525--586 页 (印刷页 521--582)

\chapter*{部分问题解答}
\addcontentsline{toc}{chapter}{部分问题解答}
\markboth{部分问题解答}{部分问题解答}

\problemsolutionhead{1}

5. 鉴于 Zorn 引理 (引理  \ref{lem:1.1}), 只须验证 $\mathcal{F}$ 是归纳的.设 $(A_i)_{i\in I}$
为 $\mathcal{F}$ 的全序子集.令 $A=\bigcap_{i\in I}A_i$, 验证 $A$ 非空, $A$ 是 $K$ 的
极端集, $A\in\mathcal{F}$, 且 $A$ 是 $(A_i)_{i\in I}$ 的上界.

6. 反设存在 $M_0$ 中的两个相异点 $a,b$.由 Hahn--Banach 定理 (定理  \ref{thm:1.7}), 存在
某个 $f\in E^\star$ 使得 $\dual{f}{a}\ne\dual{f}{b}$.令
\begin{equation*}
M_1=\left\{x\in M_0;\ \dual{f}{x}=\max_{y\in M_0}\dual{f}{y}\right\}.
\end{equation*}
显然 $M_1\in\mathcal{F}$ 且 $M_0\le M_1$.由于 $M_0$ 极大, 推出 $M_1=M_0$.这是荒谬的,
因为点 $a$ 与 $b$ 不能同时属于 $M_1$.

8. 设 $K_1$ 为 $K$ 的所有极值点的闭凸包.反设存在某个点 $a\in K$ 使得 $a\notin K_1$.
则存在严格分离 $\{a\}$ 与 $K_1$ 的超平面.设 $f\in E^\star$ 使得对一切 $x\in K_1$,
$\dual{f}{x}<\dual{f}{a}$.注意
\begin{equation*}
B=\left\{x\in K;\ \dual{f}{x}=\max_{y\in K}\dual{f}{y}\right\}
\end{equation*}
是 $K$ 的满足 $B\cap K_1=\emptyset$ 的极端集.但 $B$ 包含 $K$ 的至少一个极值点; 矛盾.

9. (a) $\mathcal{E}=\{x=(x_i);\ \abs{x_i}=1\ \forall i\}$,

(b) $\mathcal{E}=\{x=(x_i);\ \abs{x_i}=1\ \forall i,\ \text{且对大的}\ i,\ x_i\ \text{平稳}\}$,

(c) $\mathcal{E}=\emptyset$,

(d) $\mathcal{E}=\{x=(x_i);\ \exists j\ \text{使}\ \abs{x_j}=1,\ \text{且对一切}\ i\ne j,\ x_i=0\}$,

(e) $\mathcal{E}=\{x=(x_i);\ \sum\abs{x_i}^p=1\}$,

(f) $\mathcal{E}=\emptyset$.为看出情形 (f) 中 $\mathcal{E}=\emptyset$, 设 $f\in L^1(\R)$
为满足 $\int_\R\abs{f}=1$ 的任意函数.经平移后总可以假设 $\int_{-\infty}^0\abs{f}=\int_0^\infty\abs{f}=1/2$.
然后写出 $f=(g+h)/2$, 其中在 $(-\infty,0)$ 上 $g=2f$, 在 $(0,+\infty)$ 上 $g=0$;
在 $(-\infty,0)$ 上 $h=0$, 在 $(0,+\infty)$ 上 $h=2f$.

\problemsolutionhead{2}

确定对由 $x\ge0$ 时 $\varphi(x)=-\sqrt{x}$, $x<0$ 时 $\varphi(x)=+\infty$ 定义的函数
$\varphi$ 的 $\partial\varphi(x)$.

\runin{-- A --}

4. (a) $\partial\varphi(x)=F(x)$,

(b) 当 $x\ne0$ 时 $\partial\varphi(x)=\dfrac{1}{\norm{x}}F(x)$, 且 $\partial\varphi(0)=B_{E^\star}$,

(c) 若 $x\in\operatorname{Int}K$, $\partial\varphi(x)=\{0\}$; 若 $x\in\partial K$,
$\partial\varphi(x)$ 为 $K$ 在 $x$ 处的外法锥; 若 $K$ 是线性子空间, $\partial\varphi(x)=K^\perp$,

(d) $\partial\varphi(x)=D\varphi(x)=\varphi\ \text{在}\ x\ \text{处的微分}$.

5. 研究下述例子: 在 $E=\R^2$ (配以 Euclidean 范数) 中, $\varphi=I_C$, 其中
$C=\{[x_1,x_2];\ (x_1-1)^2+x_2^2\le1\}$, 以及 $\psi=I_D$, 其中 $D=\{[x_1,x_2];\ x_1=0\}$.

\runin{-- B --}

1. 设 $C=\operatorname{epi}\varphi$.对 $A=\operatorname{Int}C$ 与 $B=[x_0,\varphi(x_0)]$
应用 Hahn--Banach 定理 (第一几何形式).注意 $A\ne\emptyset$ (为什么?).因此存在某个
$f\in E^\star$ 与常数 $k$ 与 $a$, 满足 $\norm{f}+\abs{k}\ne0$, 使得对一切 $x\in D(\varphi)$
及一切 $\lambda\ge\varphi(x)$,
\begin{equation*}
\dual{f}{x}+k\lambda\ge a\ge\dual{f}{x_0}+k\varphi(x_0).
\end{equation*}
验证 $k>0$, 并由此推出 $-\frac1kf\in\partial\varphi(x_0)$.

6. 注意 $\inf_E(\tilde\varphi+\tilde\psi)=0$, 从而存在某个 $g\in E^\star$ 使得
$\tilde\varphi^\star(-g)+\tilde\psi^\star(g)=0$.验证 $f_0-g\in\partial\varphi(x)$ 且
$g\in\partial\psi(x)$; 因此 $f_0\in\partial\varphi(x)+\partial\psi(x)$.

\runin{-- C --}

1. 对每个 $R>0$ 与每个 $x_0\in E$, 对一切满足 $\norm{x-x_0}\le R$ 的 $x\in E$, 有
\begin{equation*}
\varphi(x)\le k(\norm{x_0}+R)+C\equiv M(R).
\end{equation*}
从而对一切 $f\in\partial\varphi(x_0)$, $\norm{f}\le\frac1R(k\norm{x_0}+kR+C-\varphi(x_0))$.
令 $R\to\infty$, 可见对一切 $f\in\partial\varphi(x_0)$, $\norm{f}\le k$, 由此推出对一切
$x,x_0\in E$, $\varphi(x)-\varphi(x_0)\ge-k\norm{x-x_0}$.我们有 $D(\varphi^\star)\subset kB_{E^\star}$.
事实上, 若 $f\in D(\varphi^\star)$, 写出 $\dual{f}{x}\le\varphi(x)+\varphi^\star(f)\le k\norm{x}+C+\varphi^\star(f)$.
选取 $\norm{x}=R$ 给出对一切 $R>0$, $R\norm{f}\le kR+C+\varphi^\star(f)$, 令 $R\to\infty$
即得出结论.

2. 借助 $\R^n$ 的基验证每个点 $x_0\in A$ 满足假设 (1).

\runin{-- D --}

主要难点在于证明: 若 $f\in\partial I_C(x)$, 其中 $\varphi(x)=0$ 且 $f\ne0$, 则存在某个
$\lambda>0$ 使得 $f\in\lambda\partial\varphi(x)$.在 $E\times\R$ 中对凸集
$A=\operatorname{Int}(\operatorname{epi}\varphi)$ 与 $B=\{[y,0]\in E\times\R;\ \dual{f}{y-x}\ge0\}$
应用 Hahn--Banach 定理 (第一几何形式) (验证 $A\cap B=\emptyset$).因此存在某个 $g\in E^\star$
与常数 $k$, 满足 $\norm{g}+\abs{k}\ne0$, 使得对一切 $[y,\mu]\in\operatorname{epi}\varphi$
及一切 $[z,0]\in B$, $\dual{g}{y}+k\mu\ge\dual{g}{z}$.特别地, 由此推出 $k\ge0$ 且对一切
$y\in E$, $\dual{g}{y}+k\varphi(y)\ge\dual{g}{x}$.事实上 $k\ne0$ (因为 $k=0$ 会蕴含 $g=0$).
因此 $-\dfrac gk\in\partial\varphi(x)$ (因为 $\varphi(x)=0$).此外 $g\ne0$ (为什么?).
最后, 对一切 $[z,0]\in B$ 我们有 $\dual{g}{x}\ge\dual{g}{z}$, 从而对一切满足 $\dual{f}{u}\ge0$
的 $u\in E$ 有 $\dual{g}{u}\le0$.由此推出在集合 $f^{-1}(\{0\})$ 上 $g=0$.得出存在常数
$\theta<0$ 使得 $g=\theta f$ (见引理  \ref{lem:3.2}).

\problemsolutionhead{3}

\runin{-- A --}

3. 或者对一切 $n$ 有 $x\in S(x_n)$, 此时对一切 $n$, $\psi(x_{n+1})\le\psi(x)+\varepsilon_{n+1}$.
取极限得出 $\psi(a)\le\psi(x)$, 更有 $\psi(x)-\psi(a)+d(x,a)\ge0$.

或者存在 $N$ 使得 $x\notin S(x_N)$, 此时对一切 $n\ge N$, $x\notin S(x_n)$.由此推出对
一切 $n\ge N$, $\psi(x)-\psi(x_n)+d(x,x_n)>0$.取极限同样得出 $\psi(x)-\psi(a)+d(x,a)\ge0$.

\runin{-- B --}

1. 配以距离 $d(x,y)=\lambda\norm{x-y}$ 的集合 $M$ 是完备的 (因为 $\psi$ 下半连续) 且非空
($x_0\in M$).注意 $\psi\ge0$.由第 A 部分的结果, 存在某个 $x_1\in M$ 使得对一切 $x\in M$,
$\psi(x)-\psi(x_1)+\lambda\norm{x-x_1}\ge0$.若 $x\notin M$, 由 $M$ 的定义有
$\psi(x)>\psi(x_0)-\lambda\norm{x_0-x}$, 而由于 $x_1\in M$,
$\psi(x_0)-\lambda\norm{x_0-x}\ge\psi(x_1)-\lambda\norm{x-x_1}$.结合这两种情形, 可见对一切
$x\in E$, $\psi(x)-\psi(x_1)+\lambda\norm{x-x_1}\ge0$.另一方面由于 $x_1\in M$, 我们有
$\psi(x_1)\le\psi(x_0)-\lambda\norm{x_0-x_1}$.但 $\psi(x_0)\le\varepsilon$ 且 $\psi(x_1)\ge0$.
因此 $\norm{x_0-x_1}\le\varepsilon/\lambda$.

2. 考虑函数 $\omega(x)=\psi(x)-\psi(x_1)$ 与 $\theta(x)=\lambda\norm{x-x_1}$, 从而
$0\in\partial(\omega+\theta)(x_1)$.我们已知 $\partial(\omega+\theta)=\partial\omega+\partial\theta$,
且 $\partial\theta(x_1)=\lambda B_{E^\star}$.由此推出 $0\in\partial\varphi(x_1)-f_0+\lambda B_{E^\star}$.

3. 验证 $D(\varphi)\subset\overline{D(\partial\varphi)}$.给定任何 $x_0\in D(\varphi)$,
由前面的各问已知对一切 $\varepsilon>0$ 与一切 $\lambda>0$, 存在 $x_1\in D(\partial\varphi)$
使得 $\norm{x_1-x_0}<\varepsilon/\lambda$.显然 $R(\partial\varphi)\subset D(\varphi^\star)$.
反之验证 $D(\varphi^\star)\subset\overline{R(\partial\varphi)}$.给定任何 $f_0\in D(\varphi^\star)$,
我们已知对一切 $\varepsilon>0$, 存在 $x_0\in D(\varphi)$ 使得 $f_0\in\partial_\varepsilon\varphi(x_0)$,
从而对一切 $\lambda>0$, 存在 $f_1\in R(\partial\varphi)$ 使得 $\norm{f_1-f_0}<\lambda$.

\runin{-- C --}

1. 设 $f_0\in E^\star$.由于 $(I_C)^\star(f_0)<\infty$, 我们知道对一切 $\varepsilon>0$, 存在
$x_0\in C$ 使得 $f_0\in\partial_\varepsilon I_C(x_0)$.由此推出对一切 $\lambda>0$, 存在
$x_1\in C$ 与 $f_1\in\partial I_C(x_1)$, 满足 $\norm{f_1-f_0}\le\lambda$.显然我们有
$\sup_{x\in C}\dual{f_1}{x}=\dual{f_1}{x_1}$.

2. 设 $x_0$ 为 $C$ 的边界点.则对一切 $\varepsilon>0$, 存在 $a\in E\setminus C$ 使得
$\norm{a-x_0}<\varepsilon$.用闭超平面分离 $C$ 与 $\{a\}$, 得到满足 $f_0\ne0$ 且对一切
$x\in C$ 有 $\dual{f_0}{x-a}\le0$ 的某个 $f_0\in E^\star$.当然可以假设 $\norm{f_0}=1$.
从而对一切 $x\in C$, $\dual{f_0}{x-x_0}\le\varepsilon$, 进而 $f_0\in\partial_\varepsilon I_C(x_0)$.
把第 B 部分的结果应用于 $\lambda=\sqrt{\varepsilon}$, 找到某个 $x_1\in C$ 与某个
$f_1\in\partial I_C(x_1)$ 使得 $\norm{x_1-x_0}\le\sqrt{\varepsilon}$ 且
$\norm{f_1-f_0}\le\sqrt{\varepsilon}$.只要 $\varepsilon<1$, 就有 $f_1\ne0$, 可见存在
$C$ 在 $x_1$ 处的支撑超平面.

\problemsolutionhead{4}

2. 用数学归纳法并应用习题 1.23 第 7 问.

3. 注意 $x=\frac12[(x+y)+(x-y)]$, 从而由凸性, 对一切 $x,y$,
\begin{equation*}
\psi_n(x)\le\frac12[\psi_n(x+y)+\psi_n(x-y)]\le\frac12[\varphi_n(x+y)+\psi_n(x-y)].
\end{equation*}
因此 $\psi_n(x)\le\psi_{n+1}(x)$.我们有 $\varphi_n\downarrow\theta$,
$\psi_n\uparrow\tilde\theta$, 且 $\varphi_{n+1}=\frac12(\varphi_n+\psi_n)$.因此 $\theta=\tilde\theta$.

4. 序列 $(\varphi_n^\star)$ 单调不减且收敛于极限 $\omega$.由于 $\theta\le\varphi_n$,
推出 $\varphi_n^\star\le\theta^\star$ 且 $\omega\le\theta^\star$.另一方面, 对一切 $x\in E$
及一切 $f\in E^\star$, $\dual{f}{x}-\varphi_n(x)\le\varphi_n^\star(f)$.从而对一切 $x\in E$
及一切 $f\in E^\star$, $\dual{f}{x}-\theta(x)\le\omega(f)$, 即 $\theta^\star\le\omega$.
我们得出 $\omega=\theta^\star$.

5. 应用习题 1.23 第 1 问, 可见 $\psi_{n+1}^\star=\frac12(\varphi_n^\star+\psi_n^\star)$.
序列 $(\psi_n^\star)$ 单调不增, 从而收敛于满足 $\zeta=\frac12(\theta^\star+\zeta)$ 的
极限 $\zeta$.由此推出 $\zeta=\theta^\star$ (因为 $\zeta<\infty$).

\runin{-- B --}

由 $\varphi$ 的凸性与齐次性得到
\begin{equation*}
\varphi(x+y)=\varphi\left(t\frac xt+(1-t)\frac{y}{1-t}\right)
\le t\varphi\left(\frac xt\right)+(1-t)\varphi\left(\frac{y}{1-t}\right)
=\frac1t\varphi(x)+\frac{1}{1-t}\varphi(y).
\end{equation*}
为了建立 (1), 选取 $x=\frac12(X+Y)$ 与 $y=\frac12(X-Y)$.

2. 利用 (1) 发现对一切 $x,y\in E$ 及一切 $t\in(0,1)$,
\begin{align*}
\varphi_{n+1}(x)&=\frac12\{\varphi_n(x)+\psi_n(x)\}\\
&\le\frac12\left\{\frac{1}{4t}\varphi_n(x+y)+\frac{1}{4(1-t)}\varphi_n(x-y)
+\frac{1}{4t}\psi_n(x+y)+\frac{1}{4(1-t)}\psi_n(x-y)\right\}.
\end{align*}
应用 A1 与归纳假设, 对一切 $x,y\in E$ 与一切 $t\in(0,1)$, 有
\begin{equation*}
\varphi_{n+1}(x)\le\frac12\left\{\frac{2}{4t}\varphi_{n+1}(x+y)
+\frac{1}{4(1-t)}\left(2+\frac{C}{4^n}\right)\psi_n(x-y)\right\}.
\end{equation*}
选取 $t$ 使得 $\frac{2}{4t}=\frac{1}{4(1-t)}\left(2+\frac{C}{4^n}\right)$, 即
$t=1/\left(2\left(1+\frac{C}{4^{n+1}}\right)\right)$, 我们得出对一切 $x,y\in E$,
\begin{equation*}
\varphi_{n+1}(x)\le\frac12\left(1+\frac{C}{4^{n+1}}\right)\{\varphi_n(x+y)+\psi_n(x-y)\}.
\end{equation*}
由此推出对一切 $x\in E$, $\varphi_{n+1}(x)\le\frac12\left(1+\frac{C}{4^{n+1}}\right)\psi_{n+1}(x)$.

3. 设 $x\ne y$ 且 $t\in(0,1)$, 写出
\begin{align*}
\theta(tx+(1-t)y)&\le\theta_n(tx+(1-t)y)+\frac{1}{2^n}\varphi_0(tx+(1-t)y)\\
&\le t\theta_n(x)+(1-t)\theta_n(y)+\frac{1}{2^n}\varphi_0(tx+(1-t)y)\\
&\le t\theta(x)+(1-t)\theta(y)\\
&\quad+\frac{1}{2^n}\left\{\varphi_0(tx+(1-t)y)-t\varphi_0(x)-(1-t)\varphi_0(y)+\frac{C}{2^n}(t\varphi_0(x)+(1-t)\varphi_0(y))\right\}\\
&<t\theta(x)+(1-t)\theta(y),
\end{align*}
对充分大的 $n$ 成立, 因为 $\varphi_0$ 严格凸.

\runin{-- C --}

取 $\varphi_0(x)=\frac12\norm{x}_1^2$ 与 $\psi_0(x)=\frac12\alpha^2\norm{x}_2^2$, 其中
$\alpha>0$ 充分小.范数 $\norm{\cdot}$ 通过关系式 $\theta(x)=\frac12\norm{x}^2$ 定义.

\problemsolutionhead{5}

\runin{-- B --}

1. 只须证明存在常数 $c>0$ 使得 $B(0,c)\subset K$.由 (iii) 有 $\bigcup_{n=1}^\infty(nK)=E$,
从而 $\bigcup_{n=1}^\infty\overline{nK}=E$.应用 Baire 定理, 可见
$\operatorname{Int}(\overline K)\ne\emptyset$, 从而存在某个 $y_0\in E$ 与常数 $c>0$
使得 $B(y_0,4c)\subset\overline K$.由于 $K$ 凸且对称, 推出 $B(0,2c)\subset\overline K$.

我们断言 $B(0,c)\subset K$.固定满足 $\norm{x}<c$ 的 $x\in E$.存在 $y_1,z_1\in P$
使得 $\norm{y_1}\le1/2$, $\norm{z_1}\le1/2$, 且 $\norm{x-(y_1-z_1)}<c/2$.接着存在
$y_2,z_2\in P$ 使得 $\norm{y_2}\le1/4$, $\norm{z_2}\le1/4$, 且
$\norm{x-(y_1-z_1)-(y_2-z_2)}<c/4$.重复这一构造, 得到 $P$ 中的序列 $(y_n)$ 与 $(z_n)$
使得 $\norm{y_n}\le1/2^n$, $\norm{z_n}\le1/2^n$, 且
$\norm{x-\sum_{i=1}^n(y_i-z_i)}<c/2^n$.然后把 $x$ 写成 $x=y-z$, 其中
$y=\sum_{i=1}^\infty y_i,z=\sum_{i=1}^\infty z_i$, 从而 $x\in K$.

2. 把 $x_n$ 写成 $x_n=y_n-z_n$, 其中 $y_n,z_n\in P$, $\norm{y_n}\le C/2^n$,
$\norm{z_n}\le C/2^n$.则 $1\le f(x_n)\le f(y_n)$.令 $u_n=\sum_{i=1}^ny_i$ 与
$u=\sum_{i=1}^\infty y_i$.一方面 $f(u_n)\ge n$, 另一方面 $f(u-u_n)\ge0$.由此推出
对一切 $n$, $f(u)\ge n$; 矛盾.

3. 考虑 $F$ 的补空间 (见 \S \ref{sec:2.4}).

\runin{-- C --}

(a) 有 $F=P-P=E$; 也可以直接验证 (i): 若在 $P$ 上 $f\ge0$, 则对一切 $u\in E$,
$\abs{f(u)}\le\norm{u}_\infty f(1)$.

(b) 这里 $F=\{u\in E;\ u(0)=u(1)=0\}$ 是有限余维的闭子空间.

(c) 有 $F=E$.事实上若 $u\in E$, 存在常数 $c>0$ 使得对一切 $t\in[0,1]$,
$\abs{u(t)}\le ct(1-t)$, 从而可以把 $u$ 写成 $u=v-w$, 其中 $w=ct(1-t)$, $v=u+ct(1-t)$.
\problemsolutionhead{7}

\runin{-- A --}

2. 固定满足 $\norm{x}\le1$ 的 $x\in M$.设 $\varepsilon>0$.由于 $\dist(x,N)\le a$,
存在某个 $y\in N$ 使得 $\norm{x-y}\le a+\varepsilon$, 从而 $\norm{y}\le 1+a+\varepsilon$.
另一方面, $\dist\left(\dfrac{y}{\norm{y}+\varepsilon},L\right)\le b$, 从而
$\dist(y,L)\le b(\norm{y}+\varepsilon)\le b(1+a+2\varepsilon)$.由此推出对一切 $\varepsilon>0$,
$\dist(x,L)\le a+\varepsilon+b(1+a+2\varepsilon)$.

\runin{-- B --}

为了构造使 $A^\star+B^\star\ne(A+B)^\star$ 的例子, 只须考虑满足 $D(A)\ne E$ 的
任何稠密定义且闭的无界算子 $A:D(A)\subset E\to F$.然后取 $B=-A$.我们有
$(A+B)^\star=0$, 其定义域为 $D((A+B)^\star)=F^\star$, 而 $A^\star+B^\star=0$,
其定义域为 $D(A^\star+B^\star)=D(A^\star)$.[注意 $D(A^\star)\ne F^\star$; 为什么?].

3. $A+B$ 是闭的; 事实上, 设 $(u_n)$ 为 $E$ 中的序列使得在 $E$ 中 $u_n\to u$ 且在
$F$ 中 $(A+B)u_n\to f$.注意对一切 $u\in D(A)$,
\begin{equation*}
\norm{Bu}\le k\norm{Au+Bu}+k\norm{Bu}+C\norm{u},
\end{equation*}
从而对一切 $u\in D(A)$,
\begin{equation*}
\norm{Bu}\le\frac{k}{1-k}\norm{Au+Bu}+\frac{C}{1-k}\norm{u}.
\end{equation*}
由此推出 $(Bu_n)$ 是 Cauchy 序列.设 $Bu_n\to g$, 于是 $u\in D(B)$ 且 $Bu=g$.
另一方面, $Au_n\to f-Bu$, 从而 $u\in D(A)$ 且 $Au+Bu=f$.显然有
\begin{align*}
\rho(A,A+B)&=\sup_{\substack{u\in D(A)\\\norm{u}+\norm{Au}\le1}}\inf_{v\in D(A)}\{\norm{u-v}+\norm{Au-(Av+Bv)}\}\\
&\le\sup_{\substack{u\in D(A)\\\norm{u}+\norm{Au}\le1}}\norm{Bu}\le k+C.
\end{align*}

4. 同样的论证表明在假设 \eqref{eq:prob7.Hstar} 下, 有 $\rho(A^\star,A^\star+B^\star)\le k^\star+C^\star$.

5. 设 $t\in[0,1]$.对每个 $u\in D(A)$, 有
\begin{equation*}
\norm{Bu}\le k\norm{Au}+C\norm{u}\le k(\norm{Au+tBu}+t\norm{Bu})+C\norm{u},
\end{equation*}
从而 $\norm{Bu}\le\dfrac{k}{1-k}\norm{Au+tBu}+\dfrac{C}{1-k}\norm{u}$.固定满足
$1/\varepsilon=n$ 为整数, $\dfrac{\varepsilon(k+C)}{1-k}\le\frac13$ 且
$\dfrac{\varepsilon(k^\star+C^\star)}{1-k^\star}\le\frac13$ 的任何 $\varepsilon>0$.
令 $A_1=A+\varepsilon B$, 从而 $A_1^\star=A^\star+\varepsilon B^\star$, 此外对一切
$u\in D(A)$, $\norm{Bu}\le\dfrac{k}{1-k}\norm{A_1u}+\dfrac{C}{1-k}\norm{u}$, 且对一切
$v\in D(A^\star)$, $\norm{B^\star v}\le\dfrac{k^\star}{1-k^\star}\norm{A_1^\star v}+\dfrac{C^\star}{1-k^\star}\norm{v}$.
由此推出 $(A_1+\varepsilon B)^\star=A_1^\star+\varepsilon B^\star$, 即
$(A+2\varepsilon B)^\star=A^\star+2\varepsilon B^\star$, 依此类推逐步推进到
$A_j=A+j\varepsilon B$ ($j\le n-1$).

\problemsolutionhead{8}

1. 设 $\mathcal{T}$ 为对应于度量 $d$ 的拓扑.由于配以拓扑 $\sigma(E,E^\star)$ 的
$B_{E^\star}$ 是紧的, 只须验证典范单射 $(B_{E^\star},\sigma(E^\star,E))\to(B_{E^\star},\mathcal{T})$
是连续的.这归结为证明对每个 $f_0\in B_{E^\star}$ 与每个 $\varepsilon>0$, 存在
$f_0$ 关于 $\sigma(E^\star,E)$ 的邻域 $V(f_0)$ 使得
\begin{equation*}
V(f_0)\cap B_{E^\star}\subset\{f\in B_{E^\star};\ d(f,f_0)<\varepsilon\}.
\end{equation*}
设 $(e_i)$ 为 $\ell^1$ 的典范基.选取满足 $\delta+(1/2^{n-1})<\varepsilon$ 的
$V(f_0)=\{f\in E^\star;\ \abs{\dual{f-f_0}{e_i}}<\delta\ \forall i=1,2,\dots,n\}$.

2. 注意 $(B_{E^\star},d)$ 是完备度量空间 (因为它紧).集合 $F_k$ 关于拓扑 $\mathcal{T}$
是闭的, 此外 $\bigcup_{k=1}^\infty F_k=B_{E^\star}$ (因为 $\dual{f}{x_n}\to0$).应用
Baire 纲定理即得出结论.

3. 取满足 $(1/2^{N-1})<\rho$ 的整数 $N$.对任何 $f\in B_{E^\star}$, 构造满足
$g_i=f_i$ ($1\le i\le N$) 且在 $i>N$ 时使得 $d(g,f_0)<\rho$ 的 $g\in B_{E^\star}$.
由于 $g\in F_{k_0}$, 对一切 $n\ge k_0$, $\abs{\dual{g}{x_n}}\le\varepsilon$.由此推出
\begin{equation*}
\abs{\sum_{i=1}^Nf_ix_i^n}\le\varepsilon+\sum_{i=N+1}^\infty\abs{x_i^n}\le\varepsilon+\norm{x_n}_1.
\end{equation*}
在 $B_{E^\star}$ 上取上确界得出 $\norm{x_n}_1\le\varepsilon+2\sum_{i=1}^N\abs{x_i^n}$.

4. 由于对每个固定的 $i$, 当 $n\to\infty$ 时 $x_i^n\to0$, 由第 3 问推出
$\limsup_{n\to\infty}\norm{x_n}_1\le\varepsilon$.令 $\varepsilon\to0$ 给出 $\norm{x_n}_1\to0$.

5. 序列 $(x_n)$ 是弱 Cauchy 列.由第 1--4 问的方法, 证明它是强 Cauchy 列, 从而强收敛.

6. 在 $E=L^1(0,1)$ 中取 Rademacher 函数列 $u_n(t)=\operatorname{sign}(\sin(2^n\pi t))$.
则 $\norm{u_n}_1=1$ 且由 Riemann--Lebesgue 引理, $u_n\rightharpoonup0$ 弱收敛于 $\sigma(L^1,L^\infty)$.

\problemsolutionhead{9}

\runin{-- A --}

1. 拓扑 $\sigma(E^\star,E)$ 的开凸集可由 Hahn--Banach 定理分离.因为 $\sigma(E^\star,E)$
上的连续线性泛函全都是由 $E$ 的元素给出的评估泛函 (见定理  \ref{thm:3.10}), 分离超平面必由
某个 $x\in E$ 给出.

2. 设 $f_0\notin A+B$.对每个 $y\in B$, $f_0-y\notin A$.由于 $A$ 闭, 存在 $0$ 在
$\sigma(E^\star,E)$ 中的邻域 $V(y)$ 使得 $(f_0-y+V(y))\cap A=\emptyset$.利用 $B$ 紧,
提取有限覆盖即证 $A+B$ 闭.

3. 由第 2 问, $A-B$ 在 $\sigma(E^\star,E)$ 中闭且不含 $0$.利用第 1 问严格分离 $\{0\}$
与 $A-B$.

4. 凸集的闭包仍为凸集 (由弱$^\star$拓扑下加法与数乘的连续性).

\runin{-- B --}

1. 显然 $N\subset N^{\perp\perp}$, 且 $N^{\perp\perp}$ 在 $\sigma(E^\star,E)$ 中闭,
从而 $\overline N^{\sigma(E^\star,E)}\subset N^{\perp\perp}$.反之若 $f_0\notin\overline N^{\sigma(E^\star,E)}$,
由第 A3 问存在 $x_0\in E$ 严格分离 $\{f_0\}$ 与 $\overline N^{\sigma(E^\star,E)}$.由于
$N$ 是子空间, 必有在 $N$ 上 $\dual{f}{x_0}=0$, 从而 $x_0\in N^\perp$, 但 $\dual{f_0}{x_0}\ne0$,
从而 $f_0\notin N^{\perp\perp}$.这证明了 $N^{\perp\perp}=\overline N^{\sigma(E^\star,E)}$.
若 $E$ 自反, 则所有闭子空间关于弱$^\star$拓扑都闭, 此时 $N^{\perp\perp}=N$.在
$E=\ell^1$ 中取 $N=c_0$, 则 $N^\perp=\{0\}$, 从而 $N^{\perp\perp}=\ell^\infty$, 这证明了
$c_0$ 在拓扑 $\sigma(\ell^\infty,\ell^1)$ 下在 $\ell^\infty$ 中稠密.

2. 函数 $\varphi^\star$ 是弱$^\star$连续的仿射泛函 $f\mapsto\dual{f}{x}-\varphi(x)$
的上确界, 因而关于 $\sigma(E^\star,E)$ 下半连续.反之, 设 $\psi$ 是凸且在 $\sigma(E^\star,E)$
下下半连续的, 其上图在 $E^\star\times\R$ 中关于弱$^\star$拓扑闭, 从而由其仿射下支撑泛函
的上确界给出, 这些泛函来自 $E$, 故 $\psi=\varphi^\star$, 其中 $\varphi=\psi^\star|_E$.

3. (i) 应用第 B1 问于 $N=R(A^\star)$, 并利用 $N(A)=R(A^\star)^\perp$.
(ii) 应用第 B1 问于 $N=D(A^\star)$, 并利用 $D(A^\star)^\perp=\{0\}$.

4. 反设存在 $\xi_0\in B_{E^{\star\star}}\setminus\overline{J(B_E)}^{\sigma(E^{\star\star},E^\star)}$.
由弱$^\star$ Hahn--Banach 定理, 存在 $f_0\in E^\star$ 严格分离 $\{\xi_0\}$ 与 $J(B_E)$,
得到 $\dual{\xi_0}{f_0}>\sup_{x\in B_E}\dual{f_0}{x}=\norm{f_0}$, 矛盾于 $\norm{\xi_0}\le1$.

5. 设 $x_0\in B_E$ 且 $f_0=Ax_0$.若 $f_0\notin\overline{\operatorname{conv}A(S_E)}^{\sigma(E^\star,E)}$,
用弱$^\star$超平面分离得到某个 $y\in E$ 使得对一切 $z\in S_E$, $\dual{Az}{y}\le\alpha<\dual{f_0}{y}$.
构造沿射线趋于边界的点引出矛盾.

\problemsolutionhead{10}

\runin{-- A --}

1. 若 $E$ 可分, $B_{E^\star}$ 在拓扑 $\sigma(E^\star,E)$ 下包含稠密可数集 $(b_k)$.
量 $d(x,y)=\sum_{k=1}^\infty\frac{1}{2^k}\abs{\dual{b_k}{x-y}}$ 是 $B$ 上的相容度量,
从而 $B$ 是紧度量空间.紧度量空间是列紧的, 故 (P) $\Rightarrow$ (Q).

2. 在一般情形下, 设 $(x_n)$ 为 $A$ 中的序列.令 $E_0=\overline{\operatorname{span}\{x_n\}}$,
则 $E_0$ 可分.对 $E_0$ 应用第 A1 问.

\runin{-- B --}

由于 $S_M$ 是有限维球面的紧集, 存在有限 $1/4$-网 $(g_i)_{1\le i\le k}$.对每个 $g_i$,
由对偶范数的定义存在 $a_i\in B_F$ 使得 $\dual{g_i}{a_i}\ge3/4$.对任何 $g\in S_M$,
选取最近的 $g_i$ 使得 $\norm{g-g_i}<1/4$, 则 $\dual{g}{a_i}\ge\dual{g_i}{a_i}-\norm{g-g_i}>1/2$.

\runin{-- C --}

1. 由于 $\xi\in\overline A^{\sigma(E^{\star\star},E^\star)}$, 在 $A$ 中存在满足
$\abs{\dual{\xi}{f_1}-\dual{f_1}{x_1}}<1$ 的 $x_1$.

2. 对 $M_1=[\xi,x_1]$ 应用第 B 部分的引理, 找到有限集 $(f_i)_{1<i\le n_2}$.然后利用
$\xi$ 在弱$^\star$闭包中选取 $x_2\in A$.

3. 归纳构造出序列 $(x_k)$ 与 $(f_i)$ 以及指标 $(n_k)$.

4. 由构造知在 $\mathcal{M}$ 及其闭包上 $\sup_{i\ge1}\dual{\eta}{f_i}\ge\frac12\norm{\eta}$.

5. 由假设 (Q), $(x_k)$ 有弱收敛子序列 $x_{k_j}\rightharpoonup x\in B$.取极限得到
$\dual{\xi}{f_i}=\dual{f_i}{x}$ 对一切 $i\ge1$.若 $\xi\ne x$, 取 $\eta=\xi-x\in\overline{\mathcal{M}}$,
由第 4 问知 $\sup_i\dual{\eta}{f_i}>0$, 矛盾.故 $\xi=x\in B$.

\runin{-- D --}

1. 由第 C 部分, $\overline A^{\sigma(E^{\star\star},E^\star)}\subset E$, 从而
$B=\overline A^{\sigma(E,E^\star)}$ 在 $\sigma(E^{\star\star},E^\star)$ 中闭,
进而由于有界而弱紧.

2. 第 C 部分的构造直接给出了弱收敛于 $y$ 的序列 $(y_k)\subset A$.

\problemsolutionhead{11}

\runin{-- A --}

1. 集合 $C_k=C\cap(d_kB_{E^\star})$ 按假设在 $\sigma(E^\star,E)$ 中闭, 且包含在弱$^\star$
紧球中, 故弱$^\star$紧.由于 $0\notin C$, 存在 $f_0\in C$ 使 $d=\norm{f_0}>0$.

2. 由于 $C_1$ 弱$^\star$紧且 $0\notin C_1$, 由问题 9 第 A3 问严格分离 $\{0\}$ 与 $C_1$,
存在 $x_1\in E$ 使对一切 $f\in C_1$, $\dual{f}{x_1}>1$.

3. 若对一切有限集 $A\subset\frac{1}{d_1}B_E$ 都有 $Y_A\ne\emptyset$, 由有限交性质
在紧集 $C_2$ 中存在交集点 $f\in\bigcap_AY_A$.对该 $f$ 有 $\norm{f}\le d_2$ 且对一切
$x\in\frac{1}{d_1}B_E$, $\dual{f}{x}\le1$, 导出 $\norm{f}\le d_1$, 从而 $f\in C_1$,
但这与 $\dual{f}{x_1}\le1$ 矛盾.

4. 归纳地构造满足条件的有限集 $A_k$.

5. 把所有有限集 $A_k$ 中的元素排成序列 $(x_n)$, 则由 $A_k\subset\frac{1}{d_{k-1}}B_E$
知 $\norm{x_n}\to0$, 且满足 \eqref{eq:prob11.1}.

\runin{-- B --}

1. 设 $T:E^\star\to c_0$ 由 $T(f)=(\dual{f}{x_n})_n$ 定义.由于对每个 $f\in C$,
$\sup_n\dual{f}{x_n}>1$, $T(C)$ 与 $c_0$ 的开单位球不相交.用 Hahn--Banach 定理在
$c_0$ 中分离它们.由 $(c_0)^\star=\ell^1$ 得出存在 $x\in E$ 满足对一切 $f\in C$,
$\dual{f}{x}\ge1$.

2. 若 $C$ 不是弱$^\star$闭的, 设 $f_0\in\overline C^{\sigma(E^\star,E)}\setminus C$.
经平移后设 $0\notin C-f_0$, 应用第 B1 问得到严格分离超平面, 导出矛盾.

\problemsolutionhead{12}

\runin{-- A --}

1. 设 $f=f_1+f_2$, 其中 $f_1\in D(\varphi^\star),f_2\in D(\psi^\star)$.对一切 $x\in E$,
$\dual{f}{x}-\varphi(x)-\psi(x)=(\dual{f_1}{x}-\varphi(x))+(\dual{f_2}{x}-\psi(x))\le\varphi^\star(f_1)+\psi^\star(f_2)<\infty$,
故 $f\in D((\varphi+\psi)^\star)$.

2. 由习题 1.23, $\theta=\varphi^\star\mathbin{\triangledown}\psi^\star$ 是凸的且
$D(\theta)=D(\varphi^\star)+D(\psi^\star)$.此外 $\theta^\star=\varphi^{\star\star}+\psi^{\star\star}=\varphi+\psi$.

3. 由此推出 $D((\varphi+\psi)^\star)=D(\theta)$.

\runin{-- B --}

1. 设 $g\in M$.对一切 $x\in D(\varphi)$ 与 $y\in D(\psi)$, 由 Fenchel--Young 不等式可得
$\dual{g}{x-y}\le\text{常数}$.利用假设 \eqref{eq:prob12.H} 与一致有界原理 (推论  \ref{cor:2.5}),
推出 $M$ 有界.

2. 类似地利用假设 \eqref{eq:prob12.H} 可证 $(g_n)$ 有界.

3. 证明水平集有界且弱闭, 从而下半连续.

4. 极小元达到, 从而 $(\varphi+\psi)^\star=\varphi^\star\mathbin{\triangledown}\psi^\star$,
且上确界可达到.

\problemsolutionhead{13}

\runin{-- A --}

1. 结合习题 1.25 的单侧差商性质, 当 $\lambda\to0$ 时两侧差商极限重合且等于 $\dual{Fx}{y}$.

2. 函数 $t\mapsto\frac12\norm{x+ty}^2$ 是凸的且处处可微, 其导数连续.

3. 由弱$^\star$拓扑的定义与差商连续性, 得到强到弱$^\star$连续.

4. 展开恒等式并使用范数三角不等式.

5. 当 $E$ 自反且严格凸时, 对偶映射是一对一且满射的, 其逆与 $E^\star$ 的对偶映射重合.

\runin{-- B --}

1. 若 $x_n\to x$, 则 $\norm{Fx_n}=\norm{x_n}\to\norm{x}=\norm{Fx}$ 且 $Fx_n\overset{*}{\rightharpoonup}Fx$.
由 $E^\star$ 一致凸 (从而满足 Radon--Riesz 性质), 强收敛 $Fx_n\to Fx$ 成立.

2. 用反证法, 结合第 A4 问的不等式证明有界集上的一致连续性.

3. 差商关于方向一致收敛, 从而 Fréchet 可微且微分为 $F$.

\runin{-- C --}

由范数次微分的一致连续性反推出对偶球的凸模估计, 证明 $E^\star$ 一致凸.

\problemsolutionhead{14}

\runin{-- A --}

1. 函数 $\Phi_x$ 严格凸且在自反空间中强制, 故存在唯一的极小元 $f_x$.

2. 极小元的严格极小性与一致凸性保证了 $x\mapsto Sx$ 的强连续性.

\runin{-- B --}

1. 应用定理  \ref{thm:1.12} 计算下卷积的共轭函数.

2. 验证次微分包含关系并建立 Lipschitz 估计.

3. 得出 $\psi$ 具有连续微分 $S$.

\runin{-- C --}

由习题 1.24 的论证, $\varphi_n$ 逐点递增收敛于 $\varphi$.

\problemsolutionhead{15}

\runin{-- A --}

1. 作为一族范数的上确界, $\phi$ 凸且满足 Lipschitz 条件.

2. 由于 $E$ 一致凸且 $C$ 闭凸, 极小化序列强收敛于唯一点 $c=\sigma(A)$.

3. 若 $A$ 不退化为单点, 由一致凸性在弦的中点产生严格缩减.

\runin{-- B --}

1. 极限序列单调递减且有下界, 故极限存在, 且保持凸性与连续性.

2. 类似第 A2 问, 渐近中心 $\sigma$ 存在且唯一.

3. 由各集合的中心逼近性质建立弱收敛.

4. 应用习题 3.29 得出强收敛.

5. 若 $a_n\to a$, 则渐近中心为 $a$.

6. 在 Hilbert 空间中展开平方和, 渐近中心为弱极限 $a$.

\runin{-- C --}

1. 利用压缩映射性质证明 $T\sigma$ 也是渐近中心, 由唯一性 $T\sigma=\sigma$.

2. 不动点集的闭性来自连续性, 凸性来自压缩性质.
\problemsolutionhead{16}

\runin{-- A --}

1. 我们对一切 $u\in D(A)$ 有 $\dual{Au-f}{u}\ge0$, 利用 (P) 可知 $f=A0=0$.

2. 设 $(u_n)$ 为 $D(A)$ 中的序列使得在 $E$ 中 $u_n\to x$ 且在 $E^\star$ 中 $Au_n\to f$.
对一切 $u\in D(A)$ 有 $\dual{Au-Au_n}{u-u_n}\ge0$.取极限得到对一切 $u\in D(A)$,
$\dual{Au-f}{u-x}\ge0$.由 (P) 推出 $x\in D(A)$ 且 $Ax=f$.

3. 容易验证对 $t\in(0,1)$, 凸性不等式
$\dual{A(tu+(1-t)v)}{tu+(1-t)v}\le t\dual{Au}{u}+(1-t)\dual{Av}{v}$ 等价于
$\dual{Au-Av}{u-v}\ge0$.

4. 设 $u\in N(A)$; 对一切 $v\in D(A)$, 有 $\dual{Av}{v-u}\ge0$.用 $\lambda v$ 替换 $v$,
可见对一切 $v\in D(A)$, $\dual{Av}{u}=0$; 即 $u\in R(A)^\perp$.

\runin{-- B --}

1. 注意对一切 $v\in D(A)\cap D(A^\star)$, $\dual{A^\star v}{v}=\dual{Av}{v}\ge0$.

2. 第一项断言是 (P) 以及假设 $v\notin D(A)$ 的直接推论.选取 $f=-A^\star v$, 存在
某个 $u\in D(A)$ 使得 $\dual{Au+A^\star v}{u-v}<0$, 从而
$\dual{A^\star v}{v}>\dual{Au}{u}\ge0$.

3. 把第 A4 问应用于 $A^\star$ (这是允许的, 因为 $A^\star$ 单调), 可见
$N(A^\star)\subset N(A^{\star\star})=N(A)$; 因此 $N(A)=N(A^\star)$.我们总有
$\overline{R(A)}=N(A^\star)^\perp$, 由于 $E$ 自反, 还有 $\overline{R(A^\star)}=N(A)^\perp$.

\runin{-- C --}

1. 映射 $u\in D(A)\mapsto[u,Au]$ 是配以图范数的 $D(A)$ 到 $G(A)$ 上的等距, 而 $G(A)$
是 $E\times E^\star$ 的闭子空间.

2. 注意 $\dual{Au-f}{u-x}\ge-\norm{Au}\norm{x}-\norm{f}\norm{u}+\dual{f}{x}$.

3. 利用对偶映射的微分性质 (见问题 13), 发现对一切 $v\in D(A)$,
\begin{equation*}
\dual{Av}{F^{-1}(Au_0-f)}+\dual{F(u_0-x)}{v}+\dual{Au_0-f}{v}+\dual{Av}{u_0-x}=0.
\end{equation*}
由此推出 $F^{-1}(Au_0-f)+u_0-x\in D(A^\star)$ 且
\begin{equation}
A^\star[F^{-1}(Au_0-f)+u_0-x]+(Au_0-f)+F(u_0-x)=0.
\tag{S1}
\label{eq:prob16.S1}
\end{equation}

4. 设 $x\in E$ 与 $f\in E^\star$ 满足对一切 $u\in D(A)$, $\dual{Au-f}{u-x}\ge0$.须证明
$x\in D(A)$ 且 $Ax=f$.我们知道存在满足 \eqref{eq:prob16.S1} 的某个 $u_0\in D(A)$.
应用 (M$^\star$) 给出
\begin{equation*}
\dual{Au_0-f+F(u_0-x)}{F^{-1}(Au_0-f)+u_0-x}\le0,
\end{equation*}
展开得出 $\norm{Au_0-f}^2+\norm{u_0-x}^2\le\norm{u_0-x}\norm{Au_0-f}$; 因此
$u_0=x$ 且 $Au_0=f$.

5. 对算子 $A^\star$ 应用蕴含关系 (M$^\star$) $\Rightarrow$ (P).

\problemsolutionhead{18}

2. $[G(a)-G(b)-g(b)(a-b)=0]\iff[g(a)=g(b)]$.

3. 抽取子序列可以假设 $a_{n_k}\to a$ (可能为 $\pm\infty$).有
$\int_b^a(g(t)-g(b))dt=0$, 因此 $g(a)=g(b)$.由此推出 $g(a_{n_k})\to g(b)$.

4. 注意
\begin{equation*}
0\le\int_\Omega\abs{G(u_n)-G(u)-g(u)(u_n-u)}=\int_\Omega G(u_n)-\int_\Omega G(u)-\int_\Omega g(u)(u_n-u),
\end{equation*}
使用假设 (ii).然后应用定理  \ref{thm:4.9}.

5. 由于 $g(u_n)$ 在 $L^{p'}$ 中有界, 由习题 4.16 推出对每个 $q\in[1,p')$, 在 $L^q$ 中
$g(u_{n_k})\to g(u)$ 强收敛.极限的唯一性蕴含整个序列 $g(u_n)\to g(u)$.

6. 若 $g$ 严格递增, 则 $u_{n_k}\to u$ 几乎处处, 再次使用习题 4.16 可见对每个 $q\in[1,p)$,
在 $L^q$ 中 $u_{n_k}\to u$ 强收敛.

7 与 9. 应用第 4 问与定理  \ref{thm:4.9}, 存在 $L^1$ 中的函数 $f$ 使得对一切 $k$,
\begin{equation*}
\abs{G(u_{n_k})-G(u)-g(u)(u_{n_k}-u)}\le f.
\end{equation*}
结合假设 (3) 控制住被积函数并应用控制收敛定理.

\problemsolutionhead{19}

1. 验证 $E$ 完备.注意 $\abs{\int_\R fu}\le\norm{f_1}_\infty\norm{u}_1+\norm{f_2}_2\norm{u}_2\le\max\{\norm{f_1}_\infty,\norm{f_2}_2\}\norm{u}_E$.

2. 把积分分解为 $[\abs{x}>1]$ 与 $[\abs{x}\le1]$, 在第一部分用 Hölder 不等式与 $L^1$ 控制,
在第二部分用 Hölder 不等式与 $L^2$ 控制 (因为当 $\alpha<1/2$ 时 $\abs{x}^{-\alpha}\in L^2([-1,1])$).

3. $K$ 是两个闭凸集的交, 因而闭且凸.

4. 弱收敛蕴含范数有界且点态收敛在紧集上一致, 结合奇异因子的积分估计可得收敛.

5. 建立下界估计 $J(u)\ge\norm{u}_2^2-C\norm{u}_2\ge-C^2/4$.

6. 极小化序列在 $L^2$ 中有界, 并在 $K$ 的约束下保持 $L^1$ 有界.

7. 弱下半连续性保证弱极限达到极小值.

8. $E$ 不是自反空间 (因为它包含与 $L^1$ 类似的子结构).

\problemsolutionhead{20}

\runin{-- A --}

1. 利用函数 $g(t)=\abs{t}^{p/p'}$ 在 $\R$ 上的凸性:
$\abs{\frac{a+b}{2}}^{p/p'}\le\frac12\abs{a}^{p/p'}+\frac12\abs{b}^{p/p'}$.

2. 计算 $f''(x)$, 验证在 $(0,1)$ 上 $f''(x)\le0$.由凹性得出切线不等式 \eqref{eq:prob20.3}.

3. 在 \eqref{eq:prob20.3} 中取 $y=x^{p'}$, 计算得出结果.

4. 设 $x,y\in\R$, 不失一般性设 $\abs{y}\le\abs{x}$, 考虑比值并应用第 3 问的不等式.

\runin{-- B --}

1. 对点态不等式 \eqref{eq:prob20.1} 与 \eqref{eq:prob20.2} 在 $\Omega$ 上积分并应用 Minkowski 不等式.

2. 由 \eqref{eq:prob20.4} 直接推出 Clarkson 第一不等式.

\runin{-- C --}

1. 通过对偶性 (方法 (i)) 或直接由 \eqref{eq:prob20.1} 结合习题 4.11 的反向 Minkowski 不等式 (方法 (ii)) 建立.

2. Clarkson 第二不等式直接给出当 $\norm{f-g}_p>\varepsilon$ 时
$\norm{\frac{f+g}{2}}_p\le(1-\delta)$, 从而 $L^p(\Omega)$ 一致凸.

\problemsolutionhead{21}

\runin{-- A --}

1. 用单调收敛定理证明 $\alpha(t+0)=\alpha(t)$.注意若 $f=\chi_\omega$, 则
$\alpha(1-0)=\abs{\omega}$, 而 $\alpha(1)=0$.

2. 令 $\omega_n=[\abs{f_n}>t]$, $\omega=[\abs{f}>t]$, $\chi_n=\chi_{\omega_n}$, $\chi=\chi_\omega$.
几乎处处有 $\chi(x)\le\liminf\chi_n(x)$.由 Fatou 引理可得 $\alpha(t)\le\liminf\alpha_n(t)$.
另一方面对 $\delta\in(0,t)$, 利用控制收敛定理估计得 $\limsup\alpha_n(t)\le\alpha(t-\delta)$.

\runin{-- B --}

1. 考虑函数 $H(x,t)=g(t)\chi_{[\abs{f(x)}>t]}$.应用 Fubini 与 Tonelli 定理交换积分次序.

2. 对截断函数 $\tilde f=f\chi_{[\abs{f}>\lambda]}$ 应用第 B1 问.

3. 取 $g(t)=pt^{p-1}$, 结合第 B1 与 B2 问得出结论.

\runin{-- C --}

3. 在不等式 $\int_A\abs{f}\le\abs{A}^{1/p'}[f]_p$ 中取 $A=[\abs{f}>t]$, 注意 $\int_A\abs{f}\ge t\alpha(t)$.

4. 设 $C=\sup_{t>0}t^p\alpha(t)$.写出 $\int_A\abs{f}\le\alpha(\lambda)\lambda+\int_\lambda^\infty\alpha(t)dt+\lambda\abs{A}$,
选取 $\lambda=\abs{A}^{-1/p}$ 得出估计.

5. 由分布函数积分公式与 $M^p$ 控制建立 $L^q$ 范数估计.

6. 将积分在 $(0,\lambda)$ 与 $(\lambda,\infty)$ 上分段并适当选取 $\lambda$.

7. 直接计算积分知在原点附近 $\abs{x}^{-N/p}$ 属于 $M^p$ 但不属于 $L^p$.

\problemsolutionhead{22}

1. 应用闭图像定理.

3. 由问题 21 第 B3 问, $\norm{g_\lambda}_1=\int_0^\infty\gamma_\lambda(t)dt\le N_1[\alpha(\lambda)\lambda+\int_\lambda^\infty\alpha(t)dt]$.
由 $\norm{f-g_\lambda}_\infty\le\lambda$ 得到集合包含关系.

4. 对第 3 问的不等式乘以 $\lambda^{p-2}$ 并在 $(0,\infty)$ 上积分, 得到
$\norm{f}_p^p\le pN_1\norm{u}_p^p$, 并利用 $p^{1/p}\le2$.

5. 将算子按 $T/N_\infty$ 归一化后应用上述结果.

\problemsolutionhead{23}

\runin{-- A --}

1. 集合 $X_n$ 在 $L^1$ 中闭且覆盖 $X$.由 Baire 纲定理存在内点, 从而在测度充分小的集上建立一致估计.

2. 在度量空间 $(Y,d)$ 上应用 Baire 纲定理, 得到在无穷远处尾部的控制.

3. 弱收敛序列在有界集与尾部均满足等度可积性条件.

4. 用反证法结合 Eberlein--Šmulian 定理.

\runin{-- B --}

1. $F_\varepsilon+\varepsilon B_E$ 在双对偶中弱$^\star$紧, 证明其闭包落在 $E$ 中.

2. 截断函数族在 $L^\infty$ 中有界, 从而弱紧, 结合等度可积性逼近 $F$.

\runin{-- C --}

1. 应用习题 4.14 (Vitali 收敛定理).

2. 阶梯函数控制下的函数集是等度可积的, 因而弱紧.

3. 等度可积列包含在弱紧集中, 故有弱收敛子序列.

4. 由第 A5 问与 B2 问, 序列等度可积且弱列紧.

5. 结合问题 18 的论证得出强收敛.

\runin{-- D --}

1. 凸函数的超线性增长给出等度可积性.

2. 利用分布函数构造满足超线性增长的凸函数 $G$.

\problemsolutionhead{24}

\runin{-- A --}

1--3. 标准的连续单位分解构造与一致逼近.

4--5. 可数个步长的连续分解生成稠密子空间, 故 $C(K)$ 可分.

\runin{-- B --}

1. 由习题 4.26, $L^1$ 范数与诱导的测度范数一致.

2. 由推论  \ref{cor:3.30} 提取弱$^\star$收敛子序列.

3--5. 用弱$^\star$ Hahn--Banach 定理证明 $C_c^\infty(\Omega)$ 的球在测度球中弱$^\star$稠密.

6--7. 对非负概率测度建立保号逼近列.

8. 直接检验得 $\norm{u+\delta_a}_{\mathcal{M}}=\norm{u}_{L^1}+1$.

\runin{-- C --}

1--4. Dirac 质量的有限凸组合在概率测度集中弱$^\star$稠密.

\runin{-- D --}

1--4. 构造 Jordan 分解的正部与负部并证明唯一性.

\runin{-- E --}

用逼近恒等函数 $\theta_n\uparrow1$ 替代 $1$.

\runin{-- F --}

用测度空间中的弱$^\star$收敛与 Egorov 定理重访 Dunford--Pettis 定理.

\problemsolutionhead{25}

1. 对分离超平面应用 Hahn--Banach 定理.

2. 利用开球性质证明 $K$ 有界闭凸.

3. 投影的极值性质给出 $a\in(-C)\cap\Gamma$.

4. 用反证法分离凸锥 $(-C)$ 与 $\Gamma$.

5. 考虑锥 $C=\bigcup_{\mu>0}\mu(D-x_0)$ 并应用前述结果.

\problemsolutionhead{26}

1. 目标泛函强制且严格凸, 极小元存在且唯一.

2. 次微分刻画极小性条件.

3. 投影收缩性给出解对数据的 Lipschitz 连续性.

5. 共轭函数的 Moreau 恒等式 $u+u^\star=f$.

6--7. 变动参数 $\lambda$ 的渐近分析: $\lambda\to0$ 趋于投影, $\lambda\to\infty$ 趋于极小值集上的投影.

8. 缩放后利用对偶关系求极限.

\problemsolutionhead{27}

\runin{-- A --}

1--2. 投影算子的非扩张性导出收敛性.

3--7. 极小化序列趋于距离向量 $z=P_K0$.

\runin{-- B --}

为看出 $a_1$ 与 $a_2$ 可能依赖于 $u_0$, 取如图  \ref{fig:9} 所示的凸集 $K_1$ 与 $K_2$.

\begin{figure}[tbp]
	\centering
	\begin{tikzpicture}[scale=1.0]
		\draw (0,-1.0) arc (-90:-270:1.8cm and 1.0cm) -- (0,1.0) -- cycle;
		\node at (-0.8,0) {$K_1$};
		\draw (0.6,-1.5) -- (0.6,1.5) arc (90:-90:3.0cm and 1.5cm) -- cycle;
		\node at (1.8,0) {$K_2$};
	\end{tikzpicture}
	\caption{}
	\label{fig:9}
\end{figure}

1--4. 在乘积空间中应用 Opial 引理证明弱收敛.

\problemsolutionhead{28}

\runin{-- A --}

1. 正交投影的等价刻画: 自伴性, 范数不超过 1, 核与像正交.

2. 对合算子与正交投影的对应关系.

\runin{-- B --}

1--3. 两个正交投影的可易性与子空间交, 和的正交投影.

\problemsolutionhead{29}

\runin{-- A --}

1--3. 遍历均值序列 $(\sigma_n)$ 弱收敛于不动点.

4--6. 算术平均迭代 $(\mu_n)$ 弱收敛于不动点.

\runin{-- B --}

1--8. 奇算子情形下的 Baillon 强收敛遍历定理.

\problemsolutionhead{30}

\runin{-- A --}

1--3. 摄动双线性形式取极限证明变分不等式解的存在性.

\runin{-- B --}

对上图应用变分不等式得出下半连续凸泛函上的 Stampacchia 定理.

\runin{-- C --}

双线性鞍点问题的 von Neumann 极小化极大定理.

\runin{-- D --}

凸-凹鞍点问题的一般极小化极大定理.
\problemsolutionhead{31}

\runin{-- A --}

3. 注意 $\sum_{i,j}\lambda_i\lambda_j\dual{Av_j}{v_i-v_j}=\frac12\sum_{i,j}\lambda_i\lambda_j\dual{Av_j-Av_i}{v_i-v_j}\le0$.

\runin{-- B --}

2. 对每个 $R>0$, 存在某个 $u_R\in K_R$ 使得对一切 $v\in K_R$, $\dual{Au_R}{v-u_R}\ge0$.
选取 $v=0$, 可见存在与 $R$ 无关的常数 $M$ 使得对一切 $R$, $\norm{u_R}\le M$.固定任何
$R>M$.给定 $w\in K$, 对充分小的 $t>0$ 取 $v=(1-t)u_R+tw$.

3. 取 $K=E$.先证明存在某个 $u\in E$ 使得 $Au=0$.然后用映射 $v\mapsto Av-f$ ($f\in E^\star$ 固定)
代替 $A$.

\problemsolutionhead{32}

2. 对充分小的 $\varepsilon>0$,
\begin{align*}
\phi(u_0)&\le\phi(u_0+\varepsilon v)=\max_{i\in I}\{\abs{u_0-y_i}^2-c_i+2\varepsilon(v,u_0-y_i)\}+O(\varepsilon^2)\\
&=\max_{i\in J(u_0)}\{\abs{u_0-y_i}^2-c_i+2\varepsilon(v,u_0-y_i)\}+O(\varepsilon^2)\\
&\le\phi(u_0)+2\varepsilon\max_{i\in J(u_0)}\{(v,u_0-y_i)\}+O(\varepsilon^2).
\end{align*}

3. 用反证法并应用 Hahn--Banach 定理.

4. 注意对一切 $u,v\in H$,
$\phi(v)-\phi(u)\ge\max_{i\in J(u)}\{\abs{v-y_i}^2-\abs{u-y_i}^2\}\ge2\max_{i\in J(u)}\{(u-y_i,v-u)\}$.

5. 条件 (1) 替换为 $0\in\operatorname{conv}\left(\bigcup_{i\in J(u_0)}\{f'_i(u_0)\}\right)$.

\runin{-- B --}

1. 令 $\sigma_x=\sum_{i\in I}\lambda_ix_i$ 与 $\sigma_y=\sum_{i\in I}\lambda_iy_i$, 我们得到
\begin{align*}
\sum_{j\in I}\lambda_j\abs{\sigma_y-y_j}^2&=-\abs{\sigma_y}^2+\sum_{j\in I}\lambda_j\abs{y_j}^2
=\frac12\sum_{i,j\in I}\lambda_i\lambda_j\abs{y_i-y_j}^2\\
&\le\frac12\sum_{i,j\in I}\lambda_i\lambda_j\abs{x_i-x_j}^2=-\abs{\sigma_x}^2+\sum_{j\in I}\lambda_j\abs{x_j}^2\\
&=-\abs{\sigma_x-p}^2+\sum_{j\in I}\lambda_j\abs{x_j-p}^2.
\end{align*}

2. 把 $u_0$ 写成 $u_0=\sum_{i\in J(u_0)}\lambda_jy_j$.由 B1 的结果,
$\sum_{j\in J(u_0)}\lambda_j\abs{u_0-y_j}^2\le\sum_{j\in J(u_0)}\lambda_j\abs{p-x_j}^2$.由此推出
$\sum_{j\in J(u_0)}\lambda_j\phi(u_0)\le0$, 从而 $\phi(u_0)\le0$.

\runin{注记.} 也可以通过对函数
\begin{equation*}
K(\lambda,\mu)=\sum_{j\in I}\mu_j\abs{\sum_{i\in I}\lambda_iy_i-y_j}^2-\sum_{j\in I}\mu_j\abs{p-x_j}^2
\end{equation*}
应用 von Neumann 极小化极大定理 (见问题 30 第 D 部分) 来建立 $q$ 的存在性.

\runin{-- C --}

1. 令 $K_i=\{z\in H;\ \abs{z-y_i}\le\abs{p-x_i}\}$ 与 $K=\overline{\operatorname{conv}\left(\bigcup_{i\in I}\{y_i\}\right)}$.
须证明 $\left(\bigcap_{i\in I}K_i\right)\cap K\ne\emptyset$.用反证法并归结为有限集 $I$.

2. 考虑由延拓 $S$ 且满足 $T(D(T))\subset\overline{\operatorname{conv}S(D)}$ 的所有压缩映射
$T:D(T)\subset H\to H$ 组成的偏序集.由 Zorn 引理它有极大元 $T_0$, 且 $D(T_0)=H$.

\problemsolutionhead{33}

1. 注意 $\abs{au_n}\le n\abs{u}$, 从而 $u_n\in D(A)$.此外 $\abs{u_n}\le\abs{u}$ 且 $u_n\to u$ 几乎处处.

2. 设在 $L^p$ 中 $u_n\to u$ 且 $au_n\to f$.抽取子序列可以假设 $u_n\to u$ 几乎处处且 $au_n\to f$ 几乎处处.从而 $au=f$.

3. 若 $D(A)=E$, 闭图像定理 (定理  \ref{thm:2.9}) 蕴含存在常数 $C$ 使得对一切 $u\in L^p$,
$\int_\Omega\abs{au}^p\le C\int_\Omega\abs{u}^p$.因此映射 $v\mapsto\int_\Omega\abs{a}^pv$ 是 $L^1$ 上的
连续线性泛函.由定理  \ref{thm:4.14} 存在 $f\in L^\infty$ 使得对一切 $v\in L^1$,
$\int_\Omega\abs{a}^pv=\int_\Omega fv$.因此 $a\in L^\infty$.

4. $N(A)=\{u\in L^p;\ u=0\ \text{几乎处处在}\ [a\ne0]\ \text{上}\}$ 以及
$N(A)^\perp=\{f\in L^{p'};\ f=0\ \text{几乎处处在}\ [a=0]\ \text{上}\}$.

5. $D(A^\star)=\{v\in L^{p'};\ av\in L^{p'}\}$ 且 $A^\star v=av$.

6. 结合闭图像定理与 $a$ 在几乎处处的下界控制.

7. $\sigma(A)$ 落在实质值域闭包中, 两个端点必定属于谱.

8. 若 $a$ 连续, 实质值域与通常值域闭包一致.

9--10. 谱为单点 $\{0\}$ 或算子紧均强迫 $a=0$ 几乎处处.

\problemsolutionhead{34}

\runin{-- A --}

1. 显然 $0\notin\operatorname{EV}(T)$.设 $\lambda\in\operatorname{EV}(T)$ 且 $\lambda\ne0$.
设 $u$ 为相应的特征函数, 从而 $\frac1x\int_0^xu(t)dt=\lambda u(x)$.由此推出 $u\in C^1((0,1])$
且满足 $u=\lambda u+\lambda xu'$.积分这个常微分方程, 可见存在常数 $C$ 使得 $u(x)=Cx^{-1+1/\lambda}$.
由于 $u\in C([0,1])$, 必有 $0<\lambda\le1$.反之任何 $\lambda\in(0,1]$ 是特征值, 相应的特征子空间为
$\CC x^{-1+1/\lambda}$.

3. 我们已知 $[0,1]\subset\sigma(T)\subset[-1,+1]$.对任何满足 $\lambda\notin\{0,1\}$ 的
$\lambda\in\R$, 方程 $Tu-\lambda u=f$ 的显式解由积分给出.

4. 直接计算给出当 $q>1$ 时 $\norm{T_\varepsilon u-Tu}_{L^q(0,1)}\le\dfrac{\varepsilon^{1/q}}{(q-1)^{1/q}}\norm{u}_{L^\infty(0,1)}$,
当 $q=1$ 时 $\norm{T_\varepsilon u-Tu}_{L^1(0,1)}\le\varepsilon\log(1+1/\varepsilon)\norm{u}_{L^\infty(0,1)}$.
从而 $\norm{T_\varepsilon-T}\to0$.显然 $T_\varepsilon\in\KK(E,F)$, 从而 $T\in\KK(E,F)$.

\runin{-- B --}

1. 把 $Tu(x)$ 写成 $Tu(x)=\int_0^1u(xs)ds$, 从而 $(Tu)'(x)=\int_0^1u'(xs)sds$.

2. $\operatorname{EV}(T)=(0,1/2]\cup\{1\}$.

3. $\sigma(T)=[0,1/2]\cup\{1\}$.

\runin{-- C --}

1. 分部积分并应用 Hölder 不等式建立 Hardy 不等式.

2. 由稠密性推广到一般的 $u\in E$.

3. $\operatorname{EV}(T)=(0,p/(p-1))$.

4. $\norm{T}_{\LL(E)}=\dfrac{p}{p-1}$.$T$ 不是紧算子.

5. $\sigma(T)=[0,p/(p-1)]$.

6. $(T^\star v)(x)=\int_x^1\dfrac{v(t)}{t}dt$.

7. 由 Ascoli 定理与范数逼近证明 $T\in\KK(E,F)$.

\problemsolutionhead{35}

\runin{-- A --}

1. 显然 $\norm{T^\star T}\le\norm{T}^2$.另一方面, $\abs{Tx}^2=(Tx,Tx)=(T^\star Tx,x)\le\norm{T^\star T}\abs{x}^2$.
故 $\norm{T}^2\le\norm{T^\star T}$.

2. 由归纳法 $\norm{T^{2^k}}=\norm{T}^{2^k}$.对任何整数 $N$, 选取满足 $N\le2^k$ 的 $k$,
写出 $\norm{T}^{2^k}=\norm{T^{2^k}}\le\norm{T^N}\norm{T}^{2^k-N}$, 从而 $\norm{T}^N\le\norm{T^N}$.

\runin{-- B --}

1. 利用假设 (1) 与 (2) 交替控制相邻算子对的范数.

2--3. 求和并取 $2N$ 次方根, 令 $N\to\infty$ 建立 Cotlar 引理的不等式.

\problemsolutionhead{36}

1. $R_k$ 闭是因为 $(I-T)^k=I-S$, 其中 $S\in\KK(E)$, 应用定理  \ref{thm:6.6}.

2. 降链平稳性.

4. 对偶关系的维数等式 $\operatorname{codim}R_k=\dim N_k$.

5. $R_p\cap N_p=\{0\}$ 且 $R_p+N_p=E$.

6. $(I-T)$ 限制在 $R_p$ 上是满同构.

7. 自伴算子下平方等于零蕴含本身为零, 故 $p=1$.

\problemsolutionhead{37}

10. 由第 9 问, $\nu_k^{(n)}$ 关于 $n$ 单调不减且 $\nu_k^{(n)}\le\mu_k$.利用投影算子强收敛到恒等
算子的性质, 逼近特征向量张成的有限维子空间, 建立 $\liminf\nu_k^{(n)}\ge\mu_k$.

\problemsolutionhead{38}

\runin{-- A --}

2. 对算子 $(I_F+K)$ 应用第 1 问.

3. 由有限余维子空间性质与对偶闭域定理证明.

\runin{-- B --}

对合伴随算子对偶转化.

\runin{-- C --}

1--13. Fredholm 算子的指标性质: 伴随算子的反号指标, 有限秩与紧摄动不变性, 乘积指标可加性
$\operatorname{ind}(T_2\circ T_1)=\operatorname{ind}T_1+\operatorname{ind}T_2$.

\runin{-- D --}

1. 有限维情形 $\operatorname{ind}T=\dim E-\dim F$.

2. 移位算子的 Fredholm 性与指标计算.
\problemsolutionhead{41}

\runin{-- A --}

1. 反设 $a\in(\operatorname{Int}P)\cap(-P)$.由习题 1.7, $0=\frac12a+\frac12(-a)\in\operatorname{Int}P$,
这蕴含 $P=E$, 与假设 (2) 矛盾.

2. 反设存在 $P$ 中的序列 $(x_n)$ 使得 $x_n+u\to0$.由于 $(x_n+u)-u=x_n\in P$, 取极限得到
$-u\in P$, 矛盾于第 1 问.

3. 显然 $u\ne0$ (因为由 (2), $0\notin\operatorname{Int}P$).由 (3) 有 $Tu\in\operatorname{Int}P$,
从而对某个 $\rho>0$ 有 $B(Tu,\rho)\subset P$.选取 $0<r<\rho/\norm{u}$.

4. 由于 $\lambda x=T(x+u)\ge Tu\ge ru$, 有 $\frac\lambda r x\ge u$.若归纳假设
$(\lambda/r)^nx\ge u$, 则有 $(\lambda/r)^nTx\ge Tu$, 从而 $(\lambda/r)^n(\lambda x-Tu)\ge Tu\ge ru$.
因此 $(\lambda/r)^n\lambda x\ge ru$, 即 $(\lambda/r)^{n+1}x\ge u$.另一方面,
$\lambda x=T(x+u)\in\operatorname{Int}P$, 这蕴含 $\lambda>0$ (由第 1 问).若 $0<\lambda<r$,
令 $n\to\infty$ 取极限将得到 $-u\in P$, 这不可能.

5. 映射 $x\mapsto(x+u)/\norm{x+u}$ 在 $P$ 上显然连续 (由第 2 问).由于 $T\in\KK(E)$,
$F(P)\subset T(B_E)\subset K$.应用 Schauder 不动点定理 (见习题 6.26).

6. 把 $u$ 换成 $\varepsilon u$ 时, 第 2 问中的常数 $\alpha$ 可能改变, 但第 3 问中的常数 $r$ 保持不变.

7. 有 $\lambda_\varepsilon\norm{x_\varepsilon}=\norm{T(x_\varepsilon+\varepsilon u)}\le\norm{T}\norm{x_\varepsilon+\varepsilon u}$,
因此 $\norm{x_\varepsilon}\le\norm{T}$.抽取子序列 $\varepsilon_n\to0$, 设 $\lambda_{\varepsilon_n}\to\mu_0$,
$Tx_{\varepsilon_n}\to\ell$, 得出 $x_{\varepsilon_n}\to x_0$, 其中 $x_0\in P$,
$\mu_0=\norm{x_0}\ge r$ 且 $Tx_0=\mu_0x_0$, 由 (3) 知 $x_0\in\operatorname{Int}P$.

\runin{-- B --}

1. 集合 $\Sigma=\{s\in[0,1];\ (1-s)a+sb\in P\}$ 是闭区间.由习题 1.7, $\sigma=\max\{s;\ s\in\Sigma\}$
具有所要求的性质.

2. 不能有 $\mu=0$ (否则 $0\in\operatorname{Int}P$), 也不能有 $\mu<0$ (否则 $-x\in(\operatorname{Int}P)\cap(-P)$).
故 $\mu>0$, 从而 $x\in\operatorname{Int}P$ 且 $-x\notin P$.注意 $x_0$ 与 $x$ 扮演对称角色.
令 $y=x_0-\tau_0x$, 其中 $\tau_0=\tau(x_0,-x)$.则 $y\in P$ 且 $y\ne0$.因此 $Ty\in\operatorname{Int}P$.
但 $Ty=\mu_0x_0-\tau_0\mu x$, 从而 $x_0+\dfrac{\tau_0\mu}{\mu_0}(-x)\in\operatorname{Int}P$.由 $\tau_0$
的定义得出 $\dfrac{\tau_0\mu}{\mu_0}<\tau_0$, 从而 $\mu<\mu_0$.互换 $x_0$ 与 $x$ 的角色导出 $\mu_0<\mu$.
矛盾.因此存在 $m>0$ 使得 $x=mx_0$, 进而 $\mu=\mu_0$.

3. 若 $x\in P$ 或 $-x\in P$, 由第 2 问可得结论的第一种情形.故可设 $x\notin P$ 且 $-x\notin P$.
通过分别对 $\mu>0$ 与 $\mu<0$ 考虑 $\tau_0$ 与 $\tilde\tau_0$, 证明 $\abs{\mu}<\mu_0$.

4. 几何重数为 1.

5. 鉴于问题 36, 只须证明 $N((T-\mu_0I)^2)=N(T-\mu_0I)$.反设存在满足 $Tx-\mu_0x=\alpha x_0$
($\alpha\ne0$) 的 $x$, 迭代计算得出当 $n\to\infty$ 时产生矛盾.

\problemsolutionhead{42}

\runin{-- A --}

1. 验证 $0\notin C$, 且对一切 $x\in C$, $\norm{Tx-Tx_0}\le\frac12\norm{Tx_0}$, 从而
$\norm{Tx}\ge\frac12\norm{Tx_0}$.由此推出 $0\notin\overline{T(C)}$.

2--3. 利用假设 \eqref{eq:prob42.1} 在每个非零点附近找到满意的算子 $S_y$ 并由连续性拓展到邻域.

4--6. 提取有限覆盖, 构造连续非线性映射 $F$, 并由 Schauder 不动点定理得出不动点 $\xi$.

7--8. 由不动点算子构造紧算子 $U$, 证明其核有限维且非平凡, 并在 $T$ 下不变.

9. 在 $\R^2$ 中考虑旋转算子.

\runin{-- B --}

把单算子代数 $\mathcal{A}$ 换成与 $T$ 对易的算子 $R$ 生成的代数 $\mathcal{B}$, 所有论证平行成立.

\problemsolutionhead{43}

1. 计算 $\abs{T(u+v)}^2$ 展开并作内积实部消去.

2. 利用正规性得出 $N(T^\star-\mu I)=N(T-\mu I)$, 进而正交.

3. 闭域定理与正规性给出正交补恒等式.

4--6. 证明 $R(T)=R(T^\star)$.

7--8. 归纳证明高阶幂范数等式 $\norm{T^p}=\norm{T}^p$.

9. 核的降链在一步即平稳: $N(T^2)=N(T)$.

\problemsolutionhead{44}

\runin{-- A --}

2. 由等距定义展开 $\abs{T(u+v)}^2$ 证明 $T^\star\circ T=I$.

3. 酉算子, 伴随算子等距性与可逆性之间的五重等价.

4. 经典反例: $\ell^2$ 上的右移位算子.

5. $T\circ T^\star=P_{R(T)}$.

6. 酉算子谱落在 $\{-1,+1\}$ 中, 非酉等距谱填满整个闭区间 $[-1,+1]$.

7. 紧等距强迫空间有限维.

8--11. 斜自伴算子性质与其 Cayley 变换的酉性.

\runin{-- B --}

1--12. 极分解定理 $T=U\circ P$ 的存在性, 唯一性与紧算子的谱分解展开.

\problemsolutionhead{45}

12. 对常系数热方程显式求解 Green 函数, 利用初等函数估计下界, 再通过变量代换推广到一般 Sturm--Liouville 算子.

\problemsolutionhead{46}

12. 极小化序列在 $H_0^1(I)$ 中有界, 弱极限与紧嵌入结合证明下确界可达.

\problemsolutionhead{47}

\runin{-- A --}

2--3. 平移尖峰逼近列给出上确界为 1, 但无法由 $W^{1,1}$ 函数达到.

4--5. 两端积分控制给出 $W_0^{1,1}$ 的最佳常数 $1/2$.

6. $p>1$ 时通过弱紧性证明最佳常数可达.

\runin{-- B --}

1--5. 通过正交特征函数展开与 Sturm--Liouville 特征值求解出最佳常数为 $1/\pi$.

\runin{-- C --}

1--4. Fubini 定理给出加权导数积分控制, 证明常数 $1/2$ 是最优的但不可达.

\problemsolutionhead{49}

7--11. 极小化比值确定第一特征函数 $w_0$, 取绝对值后由强极大值原理证明严格正性, 并由检验函数法证明特征子空间的一维性.

\problemsolutionhead{50}

2--6. 变分极小化构造非平凡非负解, 强极大值原理给出正性.

8--10. 严格凸泛函变换证明正解唯一性, 从而原问题恰有两个对称极小元.

\problemsolutionhead{51}

1--3. 调和振子二次增长势能 $p(x)=x^2$ 的完备化空间 $V$, 证明形式强制.

4--6. 证明 Green 算子自伴紧, 从而存在离散谱.

7--8. Hermite 多项式基与其特征值 $\lambda_n=2n+1$.

% chapters/notation.tex
% 符号表
% 源: 原书 PDF 第 587--588 页 (印刷页 583--584)

\chapter*{符号表}
\addcontentsline{toc}{chapter}{符号表}
\markboth{符号表}{符号表}

\section*{一般符号}

\begin{longtable}{@{}p{3.8cm}p{9.5cm}@{}}
$A^c$ & 集合 $A$ 的补集 \\
$E^\star$ & 对偶空间 \\
$\dual{\cdot}{\cdot}$ & 对偶 $E^\star,E$ 中的标量积 (对偶配对) \\
$[f=\alpha]=\{x;\ f(x)=\alpha\}$ & 泛函 $f$ 的等水平集 \\
$B(x_0,r)$ & 以 $x_0$ 为中心, $r$ 为半径的开球 \\
$B_E=\{x\in E;\ \norm{x}\le1\}$ & 空间 $E$ 中的闭单位球 \\
$\operatorname{epi}\varphi=\{[x,\lambda];\ \varphi(x)\le\lambda\}$ & 函数 $\varphi$ 的上图 \\
$\varphi^\star$ & 共轭函数 \\
$\LL(E,F)$ & 从 $E$ 到 $F$ 的有界线性算子空间 \\
$M^\perp$ & $M$ 的正交空间 \\
$D(A)$ & 算子 $A$ 的定义域 \\
$G(A)$ & 算子 $A$ 的图 \\
$N(A)$ & 算子 $A$ 的核 (零空间) \\
$R(A)$ & 算子 $A$ 的值域 \\
$\sigma(E,E^\star)$ & $E$ 上的弱拓扑 \\
$\sigma(E^\star,E)$ & $E^\star$ 上的弱$^\star$拓扑 \\
$\weakto$ & 弱收敛 \\
$J$ & 从 $E$ 到 $E^{\star\star}$ 的典范单射 \\
$p'$ & $p$ 的共轭指数, 即 $\dfrac{1}{p}+\dfrac{1}{p'}=1$ \\
几乎处处 (a.e.) & 在零测度集外 \\
$\abs{A}$ & 集合 $A$ 的测度 \\
$\operatorname{supp}f$ & 函数 $f$ 的支集 \\
$f\star g$ & $f$ 与 $g$ 的卷积 \\
$\rho_n$ & 磨光核序列 \\
$(\tau_hf)(x)=f(x+h)$ & 函数 $f$ 的平移 \\
$\omega\subset\subset\Omega$ & $\omega$ 紧含于 $\Omega$, 即 $\overline\omega$ 为紧集且 $\overline\omega\subset\Omega$ \\
$P_K$ & 到闭凸集 $K$ 上的投影算子 \\
$\abs{\cdot}$ & Hilbert 范数 \\
$\rho(T)$ & 算子 $T$ 的预解集 \\
$\sigma(T)$ & 算子 $T$ 的谱 \\
$\operatorname{EV}(T)$ & 算子 $T$ 的全体特征值的集合 \\
$J_\lambda=(I+\lambda A)^{-1}$ & 算子 $A$ 的预解式 \\
$A_\lambda=AJ_\lambda$ & 算子 $A$ 的 Yosida 逼近 \\
$\nabla u=\left(\dfrac{\partial u}{\partial x_1},\dfrac{\partial u}{\partial x_2},\dots,\dfrac{\partial u}{\partial x_N}\right)$ & 函数 $u$ 的梯度 \\
$D^\alpha u=\dfrac{\partial^{\abs{\alpha}}u}{\partial x_1^{\alpha_1}\partial x_2^{\alpha_2}\dots\partial x_N^{\alpha_N}}$ & 多重指标 $\alpha=(\alpha_1,\dots,\alpha_N)$, $\abs{\alpha}=\sum_{i=1}^N\alpha_i$ \\
$\Delta u=\sum_{i=1}^N\dfrac{\partial^2u}{\partial x_i^2}$ & $u$ 的 Laplace 算子 \\
$\R^N_+=\{x=(x',x_N)\in\R^{N-1}\times\R;\ x_N>0\}$ & 上半空间 \\
$Q=\{x=(x',x_N)\in\R^N\times\R;\ \abs{x'}<1\ \text{且}\ \abs{x_N}<1\}$ & 单位立方体 \\
$Q_+=Q\cap\R^N_+$ & 上半立方体 \\
$Q_0=\{x\in Q;\ x_N=0\}$ & 立方体的底面 \\
$(D_hu)(x)=\dfrac{1}{\abs{h}}(u(x+h)-u(x))$ & 差商算子 \\
$\dfrac{\partial u}{\partial n}$ & 外法向导数 \\
\end{longtable}

\section*{函数空间}

\begin{longtable}{@{}p{3.8cm}p{9.5cm}@{}}
$\Omega\subset\R^N$ & $\R^N$ 中的开集 \\
$\partial\Omega=\Gamma$ & $\Omega$ 的边界 \\
$L^p(\Omega)$ & $\{u:\Omega\to\R;\ u\ \text{可测且}\ \int_\Omega\abs{u}^p<\infty\}$, $1\le p<\infty$ \\
$L^\infty(\Omega)$ & $\{u:\Omega\to\R;\ u\ \text{可测且存在常数}\ C\ \text{使}\ \abs{u(x)}\le C\ \text{几乎处处}\}$ \\
$C_c(\Omega)$ & $\Omega$ 中具有紧支集的连续函数空间 \\
$C^k(\Omega)$ & $\Omega$ 上的 $k$ 次连续可微函数空间, $k\ge0$ \\
$C^\infty(\Omega)=\bigcap_{k\ge0}C^k(\Omega)$ & 光滑函数空间 \\
$C^k(\overline\Omega)$ & $C^k(\Omega)$ 中使得对满足 $\abs{\alpha}\le k$ 的每个多重指标 $\alpha$, 函数 $x\mapsto D^\alpha u(x)$ 可连续延拓到 $\overline\Omega$ 上的函数组成的集合 \\
$C^\infty(\overline\Omega)=\bigcap_{k\ge0}C^k(\overline\Omega)$ & 直到边界都光滑的函数空间 \\
$C^{0,\alpha}(\overline\Omega)$ & $\left\{u\in C(\overline\Omega);\ \sup_{\substack{x,y\in\Omega\\x\ne y}}\dfrac{\abs{u(x)-u(y)}}{\abs{x-y}^\alpha}<\infty\right\}$, $0<\alpha<1$ (Hölder 空间) \\
$C^{k,\alpha}(\overline\Omega)$ & $\{u\in C^k(\overline\Omega);\ D^ju\in C^{0,\alpha}(\overline\Omega)\ \forall j,\ \abs{j}\le k\}$ \\
$W^{1,p}(\Omega)$, $W_0^{1,p}(\Omega)$, $W^{m,p}(\Omega)$ & Sobolev 空间 \\
$H^1(\Omega)$, $H_0^1(\Omega)$, $H^m(\Omega)$ & 相应的 $L^2$ 型 Sobolev 空间 \\
\end{longtable}

% chapters/references.tex
% 参考文献 (References)
% 源: 原书 PDF 第 589--598 页 (印刷页 585--594)

\chapter*{参考文献}
\addcontentsline{toc}{chapter}{参考文献}
\markboth{参考文献}{参考文献}

\noindent\small

\noindent\hangindent=1.5em\hangafter=1 Adams, R. A., \hypertarget{ref:Adams-1}{[1]} Sobolev spaces, Academic Press, 1975.\par\addvspace{0.35\baselineskip}
\noindent\hangindent=1.5em\hangafter=1 Agmon, S., \hypertarget{ref:Agmon-1}{[1]} Lectures on Elliptic BoundaryValue Problems,Van Nostrand, 1965, \hypertarget{ref:Agmon-2}{[2]} On positive solutions of elliptic equations with periodic coefficients in $\R^n$, spectral results and extensions operators on Riemannian manifolds in Differential Equations (Knowles, I. W. and Lewis, R. T., eds.), North-Holland, 1984, pp. 7--17.\par\addvspace{0.35\baselineskip}
\noindent\hangindent=1.5em\hangafter=1 Agmon, S., Douglis,A. and Nirenberg, L., \hypertarget{ref:AgmonDouglisNirenberg-1}{[1]} Estimates near the boundary for solutions of elliptic partial differential equations satisfying general boundary value conditions I, Comm. Pure Appl. Math. 12 (1959), pp. 623--727.\par\addvspace{0.35\baselineskip}
\noindent\hangindent=1.5em\hangafter=1 Akhiezer, N. and Glazman, I., \hypertarget{ref:AkhiezerGlazman-1}{[1]} Theory of Linear Operators in Hilbert Space, Pitman, 1980.\par\addvspace{0.35\baselineskip}
\noindent\hangindent=1.5em\hangafter=1 Albiac, F. and Kalton, N., \hypertarget{ref:Albiac-1}{[1]} Topics in Banach Space Theory, Springer, 2006.\par\addvspace{0.35\baselineskip}
\noindent\hangindent=1.5em\hangafter=1 Alexits, G., \hypertarget{ref:Alexits-1}{[1]} Convergence Problems of Orthogonal Series, Pergamon Press, 1961.\par\addvspace{0.35\baselineskip}
\noindent\hangindent=1.5em\hangafter=1 Ambrosetti, A. and Prodi, G., \hypertarget{ref:AmbrosettiProdi-1}{[1]} A Primer of Nonlinear Analysis, Cambridge University Press, 1993.\par\addvspace{0.35\baselineskip}
\noindent\hangindent=1.5em\hangafter=1 Ambrosio, L., Fusco, N., and Pallara, D., \hypertarget{ref:Ambrosio-1}{[1]} Functions of Bounded Variation and Free Disconti- nuity Problems, Oxford University Press, 2000.\par\addvspace{0.35\baselineskip}
\noindent\hangindent=1.5em\hangafter=1 Apostol, T. M., \hypertarget{ref:Apostol-1}{[1]} Mathematical Analysis (2nd ed.), Addison--Wesley, 1974.\par\addvspace{0.35\baselineskip}
\noindent\hangindent=1.5em\hangafter=1 Ash, J. M. (ed.), \hypertarget{ref:Ash-1}{[1]} Studies in Harmonic Analysis, Mathematical Association of America, Wash- ington D.C., 1976.\par\addvspace{0.35\baselineskip}
\noindent\hangindent=1.5em\hangafter=1 Aubin, J. P., \hypertarget{ref:Aubin-1}{[1]} Mathematical Methods of Game and Economic Theory, North-Holland, 1979, \hypertarget{ref:Aubin-2}{[2]} Applied Functional Analysis, Wiley, 1999, \hypertarget{ref:Aubin-3}{[3]} Optima and Equilibria, Springer, 1993.\par\addvspace{0.35\baselineskip}
\noindent\hangindent=1.5em\hangafter=1 Aubin, J. P. and Ekeland, I., \hypertarget{ref:AubinEkeland-1}{[1]} Applied Nonlinear Analysis, Wiley, 1984.\par\addvspace{0.35\baselineskip}
\noindent\hangindent=1.5em\hangafter=1 Aubin, Th., \hypertarget{ref:Aubin-1}{[1]} Problèmes isopérimétriques et espaces de Sobolev, C. R. Acad. Sci. Paris 280 (1975), pp. 279--281, and J. Diff. Geom. 11 (1976), pp. 573--598, \hypertarget{ref:Aubin-2}{[2]} Nonlinear Analysis on Man- ifolds, Monge--Ampère equations, Springer, 1982, \hypertarget{ref:Aubin-3}{[3]} Some Nonlinear Problems in Riemannian Geometry, Springer, 1998.\par\addvspace{0.35\baselineskip}
\noindent\hangindent=1.5em\hangafter=1 Auchmuty, G., \hypertarget{ref:Auchmuty-1}{[1]} Duality for non-convex variational principles, J. Diff. Eq 50 (1983), pp. 80--145.\par\addvspace{0.35\baselineskip}
\noindent\hangindent=1.5em\hangafter=1 Azencott, R., Guivarc'h,Y., and Gundy, R. F. (Hennequin, P. L., ed.), \hypertarget{ref:Azencott-1}{[1]} Ecole d'été de probabilités de Saint-Flour VIII-1978, Springer, 1980.\par\addvspace{0.35\baselineskip}
\noindent\hangindent=1.5em\hangafter=1 Bachman, G., Narici, L., and Beckenstein, E., \hypertarget{ref:BachmanNariciBeckenstein-1}{[1]} Fourier and Wavelet Analysis, Springer, 2000.\par\addvspace{0.35\baselineskip}
\noindent\hangindent=1.5em\hangafter=1 Baiocchi, C. and Capelo, A., \hypertarget{ref:BaiocchiCapelo-1}{[1]} Variational and Quasivariational Inequalities. Applications to Free Boundary Problems, Wiley, 1984.\par\addvspace{0.35\baselineskip}
\noindent\hangindent=1.5em\hangafter=1 Balakrishnan, A., \hypertarget{ref:Balakrishnan-1}{[1]} Applied Functional Analysis, Springer, 1976.\par\addvspace{0.35\baselineskip}
\noindent\hangindent=1.5em\hangafter=1 Baouendi,M.S.andGoulaouic,C.,\hypertarget{ref:BaouendiGoulaouic-1}{[1]}Régularitéetthéoriespectralepouruneclassed'opérateurs elliptiques dégénérés, Archive Rat. Mech. Anal. 34 (1969), pp. 361--379.\par\addvspace{0.35\baselineskip}
\noindent\hangindent=1.5em\hangafter=1 Barbu, V., \hypertarget{ref:Barbu-1}{[1]} Nonlinear Semigroups and Differential Equations in Banach Spaces, Noordhoff, 1976, \hypertarget{ref:Barbu-2}{[2]} Optimal Control of Variational Inequalities, Pitman, 1984.\par\addvspace{0.35\baselineskip}
\noindent\hangindent=1.5em\hangafter=1 Barbu, V. and Precupanu, I., \hypertarget{ref:BarbuPrecupanu-1}{[1]} Convexity and Optimization in Banach Spaces, Noordhoff, 1978.\par\addvspace{0.35\baselineskip}
\noindent\hangindent=1.5em\hangafter=1 Beauzamy, B., \hypertarget{ref:Beauzamy-1}{[1]} Introduction to Banach Spaces and Their Geometry, North-Holland, 1983.\par\addvspace{0.35\baselineskip}
\noindent\hangindent=1.5em\hangafter=1 Benedetto, J. J. and Frazier, M.W. (eds.), \hypertarget{ref:BenedettoFrazier-1}{[1]}Wavelets: Mathematics and Applications, CRC Press, 1994.\par\addvspace{0.35\baselineskip}
\noindent\hangindent=1.5em\hangafter=1 Bensoussan,A. and Lions, J. L., \hypertarget{ref:BensoussanLions-1}{[1]} Applications of Variational Inequalities in Stochastic Control, North-Holland, Elsevier, 1982.\par\addvspace{0.35\baselineskip}
\noindent\hangindent=1.5em\hangafter=1 Benyamini,Y. and Lindenstrauss, J., \hypertarget{ref:Benyamini-1}{[1]} Geometric Functional Analysis, American Mathematical Society, 2000.\par\addvspace{0.35\baselineskip}
\noindent\hangindent=1.5em\hangafter=1 Bérard, P., \hypertarget{ref:Berard-1}{[1]} Spectral Geometry: Direct and Inverse Problems, Lecture Notes in Math., 1207, Springer, 1986.\par\addvspace{0.35\baselineskip}
\noindent\hangindent=1.5em\hangafter=1 Berger, Marcel, \hypertarget{ref:Berger-1}{[1]} Geometry of the spectrum, in Differential Geometry (Chern, S. S. and Osser- man, R., eds.), Proc. Sympos. Pure Math., Vol. 27, Part 2;, American Mathematical Society, 1975, pp. 129--152.\par\addvspace{0.35\baselineskip}
\noindent\hangindent=1.5em\hangafter=1 Berger, Melvyn, \hypertarget{ref:Berger-1}{[1]} Nonlinearity and Functional Analysis, Academic Press, 1977.\par\addvspace{0.35\baselineskip}
\noindent\hangindent=1.5em\hangafter=1 Bergh, J. and Löfstróm J., \hypertarget{ref:BerghLofstrom-1}{[1]} Interpolation Spaces: An Introduction, Springer, 1976.\par\addvspace{0.35\baselineskip}
\noindent\hangindent=1.5em\hangafter=1 Bers, L., John, F. and Schechter, M., \hypertarget{ref:BersJohnSchechter-1}{[1]} Partial Differential Equations, American Mathematical Society, 1979.\par\addvspace{0.35\baselineskip}
\noindent\hangindent=1.5em\hangafter=1 Bombieri, E., \hypertarget{ref:Bombieri-1}{[1]} Variational problems and elliptic equations, in Mathematical Developments Arising from Hilbert Problems (Browder, F., ed.), Proc. Sympos. Pure Math., Vol. 28, Part 2;, American Mathematical Society, 1977, pp. 525--536.\par\addvspace{0.35\baselineskip}
\noindent\hangindent=1.5em\hangafter=1 Bonsall, F.F., \hypertarget{ref:Bonsall-1}{[1]} Linear operators in complete positive cones, Proc. London Math. Soc. 8 (1958), pp. 53--75.\par\addvspace{0.35\baselineskip}
\noindent\hangindent=1.5em\hangafter=1 Bourbaki, N., \hypertarget{ref:Bourbaki-1}{[1]} Topological Vector Spaces, Springer, 1987.\par\addvspace{0.35\baselineskip}
\noindent\hangindent=1.5em\hangafter=1 Bourgain, J. and Brezis, H., \hypertarget{ref:BourgainBrezis-1}{[1]} On the equation divY = f and application to control of phases, J. Amer. Math. Soc. 16 (2003), pp. 393--426, \hypertarget{ref:BourgainBrezis-2}{[2]} New estimates for elliptic equations and Hodge type systems, J. European Math. Soc. 9 (2007), pp. 277--315.\par\addvspace{0.35\baselineskip}
\noindent\hangindent=1.5em\hangafter=1 Brandt, A., \hypertarget{ref:Brandt-1}{[1]} Interior estimates for second order elliptic differential (or finite-difference) equa- tions via the maximum principle, Israel J. Math. 76 (1969), pp. 95--121, \hypertarget{ref:Brandt-2}{[2]} Interior Schauder estimates for parabolic differential (or difference) equations, Israel J. Math. 7, (1969), pp. 254-- 262.\par\addvspace{0.35\baselineskip}
\noindent\hangindent=1.5em\hangafter=1 Bressan, A., \hypertarget{ref:Bressan-1}{[1]} BV solutions to systems of conservation laws, in Hyperbolic Systems of Balance Laws, Lectures Given at the C.I.M.E. Summer School Held in Cetraro, Italy, July 14-21, 2003 (Marcati, P., ed.), Lecture Notes in Math., 1911, Springer, 2007, pp. 1--78.\par\addvspace{0.35\baselineskip}
\noindent\hangindent=1.5em\hangafter=1 Brezis, H., \hypertarget{ref:Brezis-1}{[1]} Opérateurs maximaux monotones et semi-groupes de, contractions dans les espaces de Hilbert North-Holland, 1973, \hypertarget{ref:Brezis-2}{[2]} Periodic solutions of nonlinear vibrating strings and duality principles, Bull. Amer. Math. Soc. 8 (1983), 409--426.\par\addvspace{0.35\baselineskip}
\noindent\hangindent=1.5em\hangafter=1 Brezis, H. and Browder, F., \hypertarget{ref:Browder-1}{[1]} Partial differential equations in the 20th century, Advances in Math. 135 (1998), pp. 76--144.\par\addvspace{0.35\baselineskip}
\noindent\hangindent=1.5em\hangafter=1 Brezis, H., Coron, J. M. and Nirenberg, L., \hypertarget{ref:BrezisCoronNirenberg-1}{[1]} Free vibrations for a nonlinear wave equation and a theorem of P. Rabinowitz, Comm. Pure Appl. Math. 33 (1980), pp. 667--689.\par\addvspace{0.35\baselineskip}
\noindent\hangindent=1.5em\hangafter=1 Brezis, H. and Lieb E. H., \hypertarget{ref:Brezis-1}{[1]} A relation between pointwise convergence of functions and conver- gence of functionals, Proc. Amer. Math. Soc. 88 (1983), pp. 486--490.\par\addvspace{0.35\baselineskip}
\noindent\hangindent=1.5em\hangafter=1 Browder, F., \hypertarget{ref:Browder-1}{[1]} Nonlinear operators and nonlinear equations of evolution, Proc. Sympos. Pure Math., Vol. 18, Part 2, American Mathematical Society, 1976.\par\addvspace{0.35\baselineskip}
\noindent\hangindent=1.5em\hangafter=1 Bui, H. D., \hypertarget{ref:Bui-1}{[1]} Duality and symmetry lost in solid mechanics, C. R. Mécanique 336 (2008), pp. 12--23.\par\addvspace{0.35\baselineskip}
\noindent\hangindent=1.5em\hangafter=1 Buttazzo, G., Giaquinta, M. and Hildebrandt S., \hypertarget{ref:ButtazzoGiaquintaHildebrandt-1}{[1]} One-Dimensional Variational Problems. An Introduction, Oxford University Press, 1998.\par\addvspace{0.35\baselineskip}
\noindent\hangindent=1.5em\hangafter=1 CabréX.andCaffarelliL.,\hypertarget{ref:CabreCaffarelli-1}{[1]}FullyNonlinearEllipticEquations,AmericanMathematicalSociety, 1995.\par\addvspace{0.35\baselineskip}
\noindent\hangindent=1.5em\hangafter=1 Caffarelli, L., Nirenberg, L. and Spruck, J., \hypertarget{ref:CaffarelliNirenbergSpruck-1}{[1]} The Dirichlet problem for nonlinear second-order elliptic equations I: Monge-Ampère equations, Comm. Pure Appl. Math. 37 (1984), pp. 339--402.\par\addvspace{0.35\baselineskip}
\noindent\hangindent=1.5em\hangafter=1 Caffarelli L. and Salsa S., \hypertarget{ref:CaffarelliSalsa-1}{[1]} A Geometric Approach to Free Boundary Problems, American Mathematical Society, 2005.\par\addvspace{0.35\baselineskip}
\noindent\hangindent=1.5em\hangafter=1 Carleson, L., \hypertarget{ref:Carleson-1}{[1]} On convergence and growth of partial sums of Fourier series, Acta Math. 116 (1966), pp. 135--157.\par\addvspace{0.35\baselineskip}
\noindent\hangindent=1.5em\hangafter=1 Cazenave, T. and Haraux, A., \hypertarget{ref:CazenaveHaraux-1}{[1]} An Introduction to Semilinear Evolution Equations, Oxford University Press, 1998.\par\addvspace{0.35\baselineskip}
\noindent\hangindent=1.5em\hangafter=1 Chae, S. B., \hypertarget{ref:Chae-1}{[1]} Lebesgue Integration, Dekker, 1980.\par\addvspace{0.35\baselineskip}
\noindent\hangindent=1.5em\hangafter=1 Chan, T. F. and Shen, J., \hypertarget{ref:ChanShen-1}{[1]} Image Processing and Analysis - Variational, PDE, Wavelet, and Stochastic Methods, SIAM Publisher, 2005.\par\addvspace{0.35\baselineskip}
\noindent\hangindent=1.5em\hangafter=1 Chavel, I., \hypertarget{ref:Chavel-1}{[1]} Eigenvalues in Riemannian Geometry, Academic Press, 1994.\par\addvspace{0.35\baselineskip}
\noindent\hangindent=1.5em\hangafter=1 Chen, Ya-Zhe and Wu, Lan-Cheng, \hypertarget{ref:ChenWu-1}{[1]} Second Order Elliptic Equations and Elliptic Systems, American Mathematical Society, 1998.\par\addvspace{0.35\baselineskip}
\noindent\hangindent=1.5em\hangafter=1 Choquet, G, \hypertarget{ref:Choquet-1}{[1]} Topology,Academic Press, 1966, \hypertarget{ref:Choquet-2}{[2]} Lectures on Analysis (3 volumes), Benjamin, 1969.\par\addvspace{0.35\baselineskip}
\noindent\hangindent=1.5em\hangafter=1 Choquet-Bruhat, Y., Dewitt-Morette, C., and Dillard-Bleick, M., \hypertarget{ref:ChoquetDewittDillard-1}{[1]} Analysis, Manifolds and Physics, North-Holland, 1977.\par\addvspace{0.35\baselineskip}
\noindent\hangindent=1.5em\hangafter=1 Chui, C. K., \hypertarget{ref:Chui-1}{[1]} An Introduction to Wavelets, Academic Press, 1992.\par\addvspace{0.35\baselineskip}
\noindent\hangindent=1.5em\hangafter=1 Ciarlet, Ph., \hypertarget{ref:Ciarlet-1}{[1]} The Finite Element Method for Elliptic Problems (2nd ed.), North-Holland, 1979.\par\addvspace{0.35\baselineskip}
\noindent\hangindent=1.5em\hangafter=1 Clarke, F., \hypertarget{ref:Clarke-1}{[1]} Optimization and Nonsmooth Analysis, Wiley, 1983.\par\addvspace{0.35\baselineskip}
\noindent\hangindent=1.5em\hangafter=1 Clarke, F. and Ekeland, I., \hypertarget{ref:ClarkeEkeland-1}{[1]} Hamiltonian trajectories having prescribed minimal period, Comm. Pure Appl. Math. 33 (1980), pp. 103--116.\par\addvspace{0.35\baselineskip}
\noindent\hangindent=1.5em\hangafter=1 Coddington, E. and Levinson, N., \hypertarget{ref:CoddingtonLevinson-1}{[1]} Theory of Ordinary Differential Equations, McGraw Hill, 1955.\par\addvspace{0.35\baselineskip}
\noindent\hangindent=1.5em\hangafter=1 Cohn, D. L., \hypertarget{ref:Cohn-1}{[1]} Measure Theory, Birkhäuser, 1980.\par\addvspace{0.35\baselineskip}
\noindent\hangindent=1.5em\hangafter=1 Coifman, R. and Meyer, Y., \hypertarget{ref:CoifmanMeyer-1}{[1]} Wavelets: Calderón--Zygmund and Multilinear Operators, Cam- bridge University Press, 1997.\par\addvspace{0.35\baselineskip}
\noindent\hangindent=1.5em\hangafter=1 Conway, J. B., \hypertarget{ref:Conway-1}{[1]} A Course in Functional Analysis, Springer, 1990.\par\addvspace{0.35\baselineskip}
\noindent\hangindent=1.5em\hangafter=1 Courant, R. and Hilbert, D., \hypertarget{ref:CourantHilbert-1}{[1]} Methods of Mathematical Physics (2 volumes), Interscience, 1962.\par\addvspace{0.35\baselineskip}
\noindent\hangindent=1.5em\hangafter=1 Crank J., \hypertarget{ref:Crank-1}{[1]} Free and Moving Boundary Problems, Oxford University Press, 1987.\par\addvspace{0.35\baselineskip}
\noindent\hangindent=1.5em\hangafter=1 Croke C.B., Lasiecka I., Uhlmann G., and Vogelius M., eds, \hypertarget{ref:CrokeEtAl-1}{[1]} Geometric Methods in Inverse Problems and PDE Control, Springer, 2004.\par\addvspace{0.35\baselineskip}
\noindent\hangindent=1.5em\hangafter=1 Damlamian, A., \hypertarget{ref:Damlamian-1}{[1]} Application de la dualité non convexe à un problème non linéaire à frontière libre (équilibre d'un plasma confiné), C. R. Acad. Sci. Paris 286 (1978), pp. 153--155.\par\addvspace{0.35\baselineskip}
\noindent\hangindent=1.5em\hangafter=1 Daubechies, I., \hypertarget{ref:Daubechies-1}{[1]} Ten Lectures on Wavelets, CBMS-NSF, Reg. Conf. Ser. Appl. Math., Vol. 61, SIAM, 1992.\par\addvspace{0.35\baselineskip}
\noindent\hangindent=1.5em\hangafter=1 Dautray, R. and Lions, J.-L., \hypertarget{ref:DautrayLions-1}{[1]} Mathematical Analysis and Numerical Methods for Science and Technology (6 volumes), Springer, 1988.\par\addvspace{0.35\baselineskip}
\noindent\hangindent=1.5em\hangafter=1 David, G., \hypertarget{ref:David-1}{[1]} Wavelets and Singular Integrals on Curves and Surfaces, Lecture Notes in Math., 1465, Springer, 1992.\par\addvspace{0.35\baselineskip}
\noindent\hangindent=1.5em\hangafter=1 Davies, E., \hypertarget{ref:Davies-1}{[1]} One parameter semigroups, Academic Press, 1980, \hypertarget{ref:Davies-2}{[2]} Linear operators and their spectra, Cambridge University Press, 2007, \hypertarget{ref:Davies-3}{[3]} Heat kernels and Spectral Theory, Cambridge University Press, 1989.\par\addvspace{0.35\baselineskip}
\noindent\hangindent=1.5em\hangafter=1 de Figueiredo, D. G. and Karlovitz, L., \hypertarget{ref:De-1}{[1]} On the radial projection in normed spaces, Bull. Amer. Math. Soc. 73 (1967), pp. 364--368.\par\addvspace{0.35\baselineskip}
\noindent\hangindent=1.5em\hangafter=1 Deimling, K., \hypertarget{ref:Deimling-1}{[1]} Nonlinear Functional Analysis, Springer, 1985.\par\addvspace{0.35\baselineskip}
\noindent\hangindent=1.5em\hangafter=1 Dellacherie, C. and Meyer, P.-A., \hypertarget{ref:DellacherieMeyer-1}{[1]} Probabilités et potentiel, Hermann, Paris, 1983.\par\addvspace{0.35\baselineskip}
\noindent\hangindent=1.5em\hangafter=1 Deville, R., Godefroy, G., and Zizler, V., \hypertarget{ref:Deville-1}{[1]} Smoothness and Renormings in Banach Spaces, Longman, 1993.\par\addvspace{0.35\baselineskip}
\noindent\hangindent=1.5em\hangafter=1 DeVito, C., \hypertarget{ref:DeVito-1}{[1]} Functional Analysis, Academic Press, 1978.\par\addvspace{0.35\baselineskip}
\noindent\hangindent=1.5em\hangafter=1 DiBenedetto, E., \hypertarget{ref:DiBenedetto-1}{[1]} Partial Differential Equations (2nd ed.), Birkhäuser, 2009.\par\addvspace{0.35\baselineskip}
\noindent\hangindent=1.5em\hangafter=1 Diestel, J., \hypertarget{ref:Diestel-1}{[1]} Geometry of Banach spaces: Selected Topics, Springer, 1975, \hypertarget{ref:Diestel-2}{[2]} Sequences and Series in Banach Spaces, Springer, 1984.\par\addvspace{0.35\baselineskip}
\noindent\hangindent=1.5em\hangafter=1 Dieudonné, J., \hypertarget{ref:Dieudonne-1}{[1]} Treatise on Analysis (8 volumes), Academic Press, 1969, \hypertarget{ref:Dieudonne-2}{[2]} History of Func- tional Analysis, North-Holland, 1981.\par\addvspace{0.35\baselineskip}
\noindent\hangindent=1.5em\hangafter=1 Dixmier, J., \hypertarget{ref:Dixmier-1}{[1]} General Topology, Springer, 1984.\par\addvspace{0.35\baselineskip}
\noindent\hangindent=1.5em\hangafter=1 Dugundji, J., \hypertarget{ref:Dugundji-1}{[1]} Topology, Allyn and Bacon, 1965.\par\addvspace{0.35\baselineskip}
\noindent\hangindent=1.5em\hangafter=1 Dunford, N. and Schwartz, J. T, \hypertarget{ref:DunfordSchwartz-1}{[1]} Linear Operators (3 volumes), Interscience (1972).\par\addvspace{0.35\baselineskip}
\noindent\hangindent=1.5em\hangafter=1 Duvaut, G. and Lions, J. L., \hypertarget{ref:DuvautLions-1}{[1]} Inequalities in Mechanics and Physics, Springer, 1976.\par\addvspace{0.35\baselineskip}
\noindent\hangindent=1.5em\hangafter=1 Dym, H. and McKean, H. P., \hypertarget{ref:DymMcKean-1}{[1]} Fourier Series and Integrals, Academic Press, 1972,.\par\addvspace{0.35\baselineskip}
\noindent\hangindent=1.5em\hangafter=1 Edwards, R., \hypertarget{ref:Edwards-1}{[1]} Functional Analysis, Holt--Rinehart--Winston, 1965.\par\addvspace{0.35\baselineskip}
\noindent\hangindent=1.5em\hangafter=1 Ekeland, I., \hypertarget{ref:Ekeland-1}{[1]} Convexity Methods in Hamiltonian mechanics, Springer, 1990.\par\addvspace{0.35\baselineskip}
\noindent\hangindent=1.5em\hangafter=1 Ekeland, I. and Temam, R., \hypertarget{ref:EkelandTemam-1}{[1]} Convex Analysis and Variational Problems, North-Holland, 1976.\par\addvspace{0.35\baselineskip}
\noindent\hangindent=1.5em\hangafter=1 Ekeland, I. and Turnbull, T., \hypertarget{ref:EkelandTurnbull-1}{[1]} Infinite-Dimensional Optimization and Convexity, The University of Chicago Press, 1983.\par\addvspace{0.35\baselineskip}
\noindent\hangindent=1.5em\hangafter=1 Enflo, P., \hypertarget{ref:Enflo-1}{[1]} A counterexample to the approximation property in Banach spaces, Acta Math. 130 (1973), pp. 309--317.\par\addvspace{0.35\baselineskip}
\noindent\hangindent=1.5em\hangafter=1 Evans, L. C., \hypertarget{ref:Evans-1}{[1]} Partial Differential Equations, American Mathematical Society, 1998.\par\addvspace{0.35\baselineskip}
\noindent\hangindent=1.5em\hangafter=1 Evans, L. C. and Gariepy, R. F., \hypertarget{ref:EvansGariepy-1}{[1]} Measure Theory and Fine Properties of Functions, CRC Press, 1992.\par\addvspace{0.35\baselineskip}
\noindent\hangindent=1.5em\hangafter=1 Fife, P., \hypertarget{ref:Fife-1}{[1]} Mathematical aspects of reacting and diffusing systems, Lecture Notes in Biomathe- matics, 28, Springer, 1979.\par\addvspace{0.35\baselineskip}
\noindent\hangindent=1.5em\hangafter=1 Folland, G.B., \hypertarget{ref:Folland-1}{[1]} Introduction to Partial Differential Equations, Princeton University Press, 1976, \hypertarget{ref:Folland-2}{[2]} Real Analysis, Wiley, 1984.\par\addvspace{0.35\baselineskip}
\noindent\hangindent=1.5em\hangafter=1 Fonseca, I., Leoni, G., \hypertarget{ref:FonsecaLeoni-1}{[1]} Modern Methods in the Calculus of Variations: $L^p$ Spaces, Springer, 2007.\par\addvspace{0.35\baselineskip}
\noindent\hangindent=1.5em\hangafter=1 Franklin, J., \hypertarget{ref:Franklin-1}{[1]} Methods of Mathematical Economics: linear and Nonlinear Programming, Fixed Point Theorems, SIAM, 2002.\par\addvspace{0.35\baselineskip}
\noindent\hangindent=1.5em\hangafter=1 Friedman,A., \hypertarget{ref:Friedman-1}{[1]} Partial Differential Equations of Parabolic Type, Prentice Hall, 1964, \hypertarget{ref:Friedman-2}{[2]} Partial, differential Equations, Holt--Rinehart--Winston, 1969, \hypertarget{ref:Friedman-3}{[3]} Foundations of Modern Analysis, Holt-- Rinehart--Winston, 1970, \hypertarget{ref:Friedman-4}{[4]} Variational Principles and Free Boundary Problems, Wiley, 1982.\par\addvspace{0.35\baselineskip}
\noindent\hangindent=1.5em\hangafter=1 Garabedian, P., \hypertarget{ref:Garabedian-1}{[1]} Partial Differential equations, Wiley, 1964.\par\addvspace{0.35\baselineskip}
\noindent\hangindent=1.5em\hangafter=1 Germain, P., \hypertarget{ref:Germain-1}{[1]} Duality and convection in continuum mechanics, in Trends in Applications of Pure Mathematics to Mechanics (Fichera, G., ed.), Monographs and Studies in Math., Vol. 2, Pitman, 1976, pp. 107--128, \hypertarget{ref:Germain-2}{[2]} Functional concepts in continuum mechanics, Meccanica 33 (1998), pp. 433--444.\par\addvspace{0.35\baselineskip}
\noindent\hangindent=1.5em\hangafter=1 Giaquinta, M., \hypertarget{ref:Giaquinta-1}{[1]} Multiple Integrals in the Calculus of Variations and Nonlinear Elliptic Systems, Princeton University Press, 1983.\par\addvspace{0.35\baselineskip}
\noindent\hangindent=1.5em\hangafter=1 Gilbarg, D. andTrudinger, N., \hypertarget{ref:GilbargTrudinger-1}{[1]} Elliptic Partial Differential Equations of Second Order, Springer, 1977.\par\addvspace{0.35\baselineskip}
\noindent\hangindent=1.5em\hangafter=1 Gilkey, P., \hypertarget{ref:Gilkey-1}{[1]} The Index theorem and the Heat Equation, Publish or Perish, 1974.\par\addvspace{0.35\baselineskip}
\noindent\hangindent=1.5em\hangafter=1 Giusti, E., \hypertarget{ref:Giusti-1}{[1]} Minimal Surfaces and Functions of Bounded Variation, Birkhäuser, 1984, \hypertarget{ref:Giusti-2}{[2]} Direct Methods in the Calculus of Variations, World Scientific, 2003.\par\addvspace{0.35\baselineskip}
\noindent\hangindent=1.5em\hangafter=1 Goldstein, J., \hypertarget{ref:Goldstein-1}{[1]} Semigroups of Operators and Applications, Oxford University Press, 1985.\par\addvspace{0.35\baselineskip}
\noindent\hangindent=1.5em\hangafter=1 Gordon, C., Webb, L., and Wolpert, S., \hypertarget{ref:GordonWebbWolpert-1}{[1]} One cannot hear the shape of a drum, Bull. Amer. Math. Soc. 27 (1992), pp. 134--138.\par\addvspace{0.35\baselineskip}
\noindent\hangindent=1.5em\hangafter=1 Grafakos, L., \hypertarget{ref:Grafakos-1}{[1]} Classical and Modern Fourier Analysis (2nd ed.), Springer, 2008.\par\addvspace{0.35\baselineskip}
\noindent\hangindent=1.5em\hangafter=1 Granas, A. and Dugundji, J., \hypertarget{ref:GranasDugundji-1}{[1]} Fixed Point Theory, Springer, 2003.\par\addvspace{0.35\baselineskip}
\noindent\hangindent=1.5em\hangafter=1 Grisvard, P., \hypertarget{ref:Grisvard-1}{[1]} Équations différentielles abstraites, Ann. Sci. ENS 2 (1969), pp. 311--395.\par\addvspace{0.35\baselineskip}
\noindent\hangindent=1.5em\hangafter=1 Gurtin, M., \hypertarget{ref:Gurtin-1}{[1]} An Introduction to Continuum Mechanics, Academic Press, 1981.\par\addvspace{0.35\baselineskip}
\noindent\hangindent=1.5em\hangafter=1 Halmos, P., \hypertarget{ref:Halmos-1}{[1]} Measure Theory, Van Nostrand, 1950, \hypertarget{ref:Halmos-2}{[2]} A Hilbert Space Problem Book (2nd ed.), Springer, 1982.\par\addvspace{0.35\baselineskip}
\noindent\hangindent=1.5em\hangafter=1 Han, Q. and Lin, F. H., \hypertarget{ref:HanLin-1}{[1]} Elliptic Partial Differential Equations,American Mathematical Society, 1997.\par\addvspace{0.35\baselineskip}
\noindent\hangindent=1.5em\hangafter=1 Harmuth, H. F., \hypertarget{ref:Harmuth-1}{[1]} Transmission of Information by Orthogonal Functions, Springer, 1972.\par\addvspace{0.35\baselineskip}
\noindent\hangindent=1.5em\hangafter=1 Hartman, Ph., \hypertarget{ref:Hartman-1}{[1]} Ordinary Differential Equations, Wiley, 1964.\par\addvspace{0.35\baselineskip}
\noindent\hangindent=1.5em\hangafter=1 Henry, D., \hypertarget{ref:Henry-1}{[1]} Geometric Theory of Semilinear Parabolic Equations, Springer, 1981.\par\addvspace{0.35\baselineskip}
\noindent\hangindent=1.5em\hangafter=1 Hernandez, E. and Weiss, G., \hypertarget{ref:HernandezWeiss-1}{[1]} A First Course on Wavelets, CRC Press, 1996.\par\addvspace{0.35\baselineskip}
\noindent\hangindent=1.5em\hangafter=1 Hewitt, E. and Stromberg, K., \hypertarget{ref:HewittStromberg-1}{[1]} Real and Abstract Analysis, Springer, 1965.\par\addvspace{0.35\baselineskip}
\noindent\hangindent=1.5em\hangafter=1 Hille, E., \hypertarget{ref:Hille-1}{[1]} Methods in Classical and Functional Analysis, Addison--Wesley, 1972.\par\addvspace{0.35\baselineskip}
\noindent\hangindent=1.5em\hangafter=1 Hiriart-Urruty, J. B. and Lemaréchal, C., \hypertarget{ref:HiriartUrrutyLemarechal-1}{[1]} Convex Analysis and Minimization Algorithms (2 volumes), Springer, 1993.\par\addvspace{0.35\baselineskip}
\noindent\hangindent=1.5em\hangafter=1 Holmes, R., \hypertarget{ref:Holmes-1}{[1]} Geometric Functional Analysis and Its Applications, Springer, 1975.\par\addvspace{0.35\baselineskip}
\noindent\hangindent=1.5em\hangafter=1 Hörmander, L., \hypertarget{ref:Hormander-1}{[1]} Linear Partial Differential Operators, Springer, 1963, \hypertarget{ref:Hormander-2}{[2]} The Analysis of Lin- ear Partial Differenetial Operators, (4 volumes), Springer, 1983--1985, \hypertarget{ref:Hormander-3}{[3]} Notions of Convexity, Birkhäuser, 1994.\par\addvspace{0.35\baselineskip}
\noindent\hangindent=1.5em\hangafter=1 Horvath, J., \hypertarget{ref:Horvath-1}{[1]} Topological Vector Spaces and Distributions, Addison--Wesley, 1966.\par\addvspace{0.35\baselineskip}
\noindent\hangindent=1.5em\hangafter=1 Ince, E. L., \hypertarget{ref:Ince-1}{[1]} Ordinary Differential Equations, Dover Publications, 1944.\par\addvspace{0.35\baselineskip}
\noindent\hangindent=1.5em\hangafter=1 James, R. C., \hypertarget{ref:James-1}{[1]} A non-reflexive Banach space isometric with its second conjugate space, Proc. Nat. Acad. Sci. USA 37 (1951), pp. 174--177.\par\addvspace{0.35\baselineskip}
\noindent\hangindent=1.5em\hangafter=1 Jost, J., \hypertarget{ref:Jost-1}{[1]} Partial Differential Equations, Springer, 2002.\par\addvspace{0.35\baselineskip}
\noindent\hangindent=1.5em\hangafter=1 Kac, M., \hypertarget{ref:Kac-1}{[1]} Can one hear the shape of a drum? Amer. Math. Monthly 73 (1966), pp. 1--23.\par\addvspace{0.35\baselineskip}
\noindent\hangindent=1.5em\hangafter=1 Kahane, J. P. and Lemarié-Rieusset, P. G., \hypertarget{ref:KahaneLemarie-1}{[1]} Fourier Series and Wavelets, Gordon Breach, 1996.\par\addvspace{0.35\baselineskip}
\noindent\hangindent=1.5em\hangafter=1 Kaiser, G., \hypertarget{ref:Kaiser-1}{[1]} A Friendly Guide to Wavelets, Birkhäuser, 1994.\par\addvspace{0.35\baselineskip}
\noindent\hangindent=1.5em\hangafter=1 Kakutani, S., \hypertarget{ref:Kakutani-1}{[1]} Some characterizations of Euclidean spaces, Jap. J. Math. 16 (1939), pp. 93--97.\par\addvspace{0.35\baselineskip}
\noindent\hangindent=1.5em\hangafter=1 Karlin, S., \hypertarget{ref:Karlin-1}{[1]} Mathematical Methods and Theory in Games, Programming, and Economics (2 vol- umes), Addison--Wesley, 1959.\par\addvspace{0.35\baselineskip}
\noindent\hangindent=1.5em\hangafter=1 Kato, T., \hypertarget{ref:Kato-1}{[1]} Perturbation Theory for Linear Operators, Springer, 1976.\par\addvspace{0.35\baselineskip}
\noindent\hangindent=1.5em\hangafter=1 Katznelson, Y., \hypertarget{ref:Katznelson-1}{[1]} An Introduction to Harmonic Analysis, Dover Publications, 1976.\par\addvspace{0.35\baselineskip}
\noindent\hangindent=1.5em\hangafter=1 Kelley, J. L., \hypertarget{ref:Kelley-1}{[1]} Banach spaces with the extension property, Trans. Amer. Math. Soc. 72 (1952), pp. 323--326.\par\addvspace{0.35\baselineskip}
\noindent\hangindent=1.5em\hangafter=1 Kelley, J. and Namioka, I., \hypertarget{ref:KelleyNamioka-1}{[1]} Linear Topological Spaces (2nd ed.), Springer, 1976.\par\addvspace{0.35\baselineskip}
\noindent\hangindent=1.5em\hangafter=1 Kinderlehrer, D. and Stampacchia, G., \hypertarget{ref:KinderlehrerStampacchia-1}{[1]} An Introduction to Variational Inequalities and Their Applications, Academic Press, 1980.\par\addvspace{0.35\baselineskip}
\noindent\hangindent=1.5em\hangafter=1 Klainerman, S., \hypertarget{ref:Klainerman-1}{[1]} Introduction to Analysis, Lecture Notes, Princeton University, 2008.\par\addvspace{0.35\baselineskip}
\noindent\hangindent=1.5em\hangafter=1 Knapp, A., \hypertarget{ref:Knapp-1}{[1]} Basic Real Analysis, Birkhäuser, 2005, \hypertarget{ref:Knapp-2}{[2]} Advanced Real Analysis, Birkhäuser, 2005.\par\addvspace{0.35\baselineskip}
\noindent\hangindent=1.5em\hangafter=1 Knerr, B., \hypertarget{ref:Knerr-1}{[1]} Parabolic interior Schauder estimates by the maximum principle, Archive Rat. Mech. Anal. 75 (1980), pp. 51--58.\par\addvspace{0.35\baselineskip}
\noindent\hangindent=1.5em\hangafter=1 Kohn, J.J. and Nirenberg, L., \hypertarget{ref:KohnNirenberg-1}{[1]} Degenerate elliptic parabolic equations of second order, Comm. Pure Appl. Math. 20 (1967), pp. 797--872.\par\addvspace{0.35\baselineskip}
\noindent\hangindent=1.5em\hangafter=1 Kolmogorov, A. and Fomin, S., \hypertarget{ref:KolmogorovFomin-1}{[1]} Introductory Real Analysis, Prentice-Hall, 1970.\par\addvspace{0.35\baselineskip}
\noindent\hangindent=1.5em\hangafter=1 Köthe, G., \hypertarget{ref:Kothe-1}{[1]} Topological Vector Spaces (2 volumes), Springer, 1969--1979.\par\addvspace{0.35\baselineskip}
\noindent\hangindent=1.5em\hangafter=1 Krasnoselskii, M., \hypertarget{ref:Krasnoselskii-1}{[1]} Topological Methods in the Theory of Nonlinear Integral Equations, Mcmil- lan, 1964.\par\addvspace{0.35\baselineskip}
\noindent\hangindent=1.5em\hangafter=1 Kreyszig, E., \hypertarget{ref:Kreyszig-1}{[1]} Introductory Functional Analysis with Applications, Wiley, 1978.\par\addvspace{0.35\baselineskip}
\noindent\hangindent=1.5em\hangafter=1 Krylov, N. V., \hypertarget{ref:Krylov-1}{[1]} Lectures on Elliptic and Parabolic Equations in Hölder Spaces, American Mathematical Society, 1996, \hypertarget{ref:Krylov-2}{[2]} Lectures on Elliptic and Parabolic Equations in Sobolev Spaces, American Mathematical Society, 2008.\par\addvspace{0.35\baselineskip}
\noindent\hangindent=1.5em\hangafter=1 Ladyzhenskaya, O. and Uraltseva, N., \hypertarget{ref:Ladyzhenskaya-1}{[1]} Linear and quasilinear elliptic equations, Academic Press, 1968.\par\addvspace{0.35\baselineskip}
\noindent\hangindent=1.5em\hangafter=1 Ladyzhenskaya, O., Solonnikov, V., and Uraltseva, N., \hypertarget{ref:LadyzhenskayaSolonnikovUraltseva-1}{[1]} Linear and Quasilinear Equations of Parabolic Type, American Mathematical Society, 1968.\par\addvspace{0.35\baselineskip}
\noindent\hangindent=1.5em\hangafter=1 Lang, S., \hypertarget{ref:Lang-1}{[1]} Real and Functional Analysis, Springer, 1993.\par\addvspace{0.35\baselineskip}
\noindent\hangindent=1.5em\hangafter=1 Larsen, R., \hypertarget{ref:Larsen-1}{[1]} Functional Analysis: An Introduction, Dekker, 1973.\par\addvspace{0.35\baselineskip}
\noindent\hangindent=1.5em\hangafter=1 Lax, P., \hypertarget{ref:Lax-1}{[1]} Functional Analysis, Wiley, 2002.\par\addvspace{0.35\baselineskip}
\noindent\hangindent=1.5em\hangafter=1 Levitan, B. M., \hypertarget{ref:Levitan-1}{[1]} Inverse Sturm-Liouville problems, VNU Science Press, 1987.\par\addvspace{0.35\baselineskip}
\noindent\hangindent=1.5em\hangafter=1 Lieb, E., \hypertarget{ref:Lieb-1}{[1]} Sharp constants in the Hardy-Littlewood-Sobolev and related inequalities,Ann. Math. 118 (1983), pp. 349--374.\par\addvspace{0.35\baselineskip}
\noindent\hangindent=1.5em\hangafter=1 Lieb, E. H. and Loss, M., \hypertarget{ref:Lieb-1}{[1]} Analysis, American Mathematical Society, 1997.\par\addvspace{0.35\baselineskip}
\noindent\hangindent=1.5em\hangafter=1 Lin, F. H. and Yang, X. P., \hypertarget{ref:LinYang-1}{[1]} Geometric Measure Theory -- An Introduction, International Press, 2002.\par\addvspace{0.35\baselineskip}
\noindent\hangindent=1.5em\hangafter=1 Lindenstrauss, J. and Tzafriri, L., \hypertarget{ref:LindenstraussTzafriri-1}{[1]} On the complemented subspaces problem, Israel J. Math. 9 (1971), pp. 263--269, \hypertarget{ref:LindenstraussTzafriri-2}{[2]} Classical Banach Spaces, (2 volumes), Springer, 1973--1979.\par\addvspace{0.35\baselineskip}
\noindent\hangindent=1.5em\hangafter=1 Lions, J. L., \hypertarget{ref:Lions-1}{[1]} Problèmes aux limites dans les équations aux dérivées partielles, Presses de l'Université de Montreal, 1965, \hypertarget{ref:Lions-2}{[2]} Optimal Control of Systems Governed by Partial Differential Equations, Springer, 1971, \hypertarget{ref:Lions-3}{[3]} Quelques méthodes de résolution des problèmes aux limites non linéaires, Dunod-Gauthier Villars, 1969.\par\addvspace{0.35\baselineskip}
\noindent\hangindent=1.5em\hangafter=1 Lions, J. L. and Magenes, E., \hypertarget{ref:LionsMagenes-1}{[1]} Non-homogeneous Boundary Value Problems and Applications (3 volumes), Springer, 1972.\par\addvspace{0.35\baselineskip}
\noindent\hangindent=1.5em\hangafter=1 Magenes, E., \hypertarget{ref:Magenes-1}{[1]} Topics in parabolic equations: some typical free boundary problems, in Boundary Value Problems for Linear Evolution Partial Differential Equations (Garnir, H. G., ed.), Reidel, 1977.\par\addvspace{0.35\baselineskip}
\noindent\hangindent=1.5em\hangafter=1 Mallat, S., \hypertarget{ref:Mallat-1}{[1]} A Wavelet Tour of Signal Processing, Academic Press, 1999.\par\addvspace{0.35\baselineskip}
\noindent\hangindent=1.5em\hangafter=1 Malliavin, P., \hypertarget{ref:Malliavin-1}{[1]} Integration and Probabilities, Springer, 1995.\par\addvspace{0.35\baselineskip}
\noindent\hangindent=1.5em\hangafter=1 Martin, R. H., \hypertarget{ref:Martin-1}{[1]} Nonlinear Operators and Differential Equations in Banach Spaces, Wiley, 1976.\par\addvspace{0.35\baselineskip}
\noindent\hangindent=1.5em\hangafter=1 Mawhin, J., R. Ortega, R., and Robles-Pérez,A. M., \hypertarget{ref:MawhinOrtegaRoblesPerez-1}{[1]} Maximum principles for bounded solutions of the telegraph equation in space dimensions two and three and applications, J. Differential Equations 208 (2005), pp. 42--63.\par\addvspace{0.35\baselineskip}
\noindent\hangindent=1.5em\hangafter=1 Mawhin, J. and Willem, M., \hypertarget{ref:MawhinWillem-1}{[1]} Critical Point Theory and Hamiltonian Systems, Springer, 1989.\par\addvspace{0.35\baselineskip}
\noindent\hangindent=1.5em\hangafter=1 Maz'ja, V. G., \hypertarget{ref:Mazja-1}{[1]} Sobolev Spaces, Springer, 1985.\par\addvspace{0.35\baselineskip}
\noindent\hangindent=1.5em\hangafter=1 Meyer,Y., \hypertarget{ref:Meyer-1}{[1]}Wavelets and Operators, Cambridge University Press, 1992. \hypertarget{ref:Meyer-2}{[2]}Wavelets,Vibrations and Scaling, American Mathematical Society, 1998, \hypertarget{ref:Meyer-3}{[3]} Oscillating Patterns in Image Processing and Nonlinear Evolution Equations, American Mathematical Society, 2001.\par\addvspace{0.35\baselineskip}
\noindent\hangindent=1.5em\hangafter=1 Mikhlin, S., \hypertarget{ref:Mikhlin-1}{[1]} An Advanced Course of Mathematical Physics, North-Holland, 1970.\par\addvspace{0.35\baselineskip}
\noindent\hangindent=1.5em\hangafter=1 Miranda, C., \hypertarget{ref:Miranda-1}{[1]} Partial Differential Equations of Elliptic Type, Springer, 1970.\par\addvspace{0.35\baselineskip}
\noindent\hangindent=1.5em\hangafter=1 Mizohata, S., \hypertarget{ref:Mizohata-1}{[1]} The Theory of Partial Differential Equations, Cambridge University Press, 1973.\par\addvspace{0.35\baselineskip}
\noindent\hangindent=1.5em\hangafter=1 Morawetz, C., \hypertarget{ref:Morawetz-1}{[1]} $L^p$ inequalities, Bull. Amer. Math. Soc. 75 (1969), pp. 1299--1302.\par\addvspace{0.35\baselineskip}
\noindent\hangindent=1.5em\hangafter=1 Moreau, J. J, \hypertarget{ref:Moreau-1}{[1]} Fonctionnelles convexes, Séminaire Leray, Collège de France, 1966. \hypertarget{ref:Moreau-2}{[2]} Appli- cations of convex analysis to the treatment of elastoplastic systems, in Applications of methods of FunctionalAnalysis to Problems in Mechanics, Sympos. IUTAM/IMU (Germain, P. and Nayroles, B., eds.), Springer, 1976.\par\addvspace{0.35\baselineskip}
\noindent\hangindent=1.5em\hangafter=1 Morrey, C., \hypertarget{ref:Morrey-1}{[1]} Multiple Integrals in the Calculus of Variations, Springer, 1966.\par\addvspace{0.35\baselineskip}
\noindent\hangindent=1.5em\hangafter=1 Morton K.W. and Mayers D.F., \hypertarget{ref:MortonMayers-1}{[1]} Numerical Solutions of Partial Differential Equations (2nd ed.), Cambridge University Press, 2005.\par\addvspace{0.35\baselineskip}
\noindent\hangindent=1.5em\hangafter=1 Munkres, J. R., \hypertarget{ref:Munkres-1}{[1]} Topology, A First Course, Prentice Hall, 1975.\par\addvspace{0.35\baselineskip}
\noindent\hangindent=1.5em\hangafter=1 Nachbin, L., \hypertarget{ref:Nachbin-1}{[1]} A theorem of the Hahn-Banach type for linear transformations, Trans. Amer. Math. Soc. 68 (1950), pp. 28--46.\par\addvspace{0.35\baselineskip}
\noindent\hangindent=1.5em\hangafter=1 Neˇcas, J. and Hlavaˇcek, L., \hypertarget{ref:NecasHlavacek-1}{[1]} Mathematical Theory of Elastic and Elastoplastic Bodies. An Introduction, Elsevier, 1981.\par\addvspace{0.35\baselineskip}
\noindent\hangindent=1.5em\hangafter=1 Neveu, J., \hypertarget{ref:Neveu-1}{[1]} , Mathematical Foundations of the Calculus of Probability, Holden-Day, 1965.\par\addvspace{0.35\baselineskip}
\noindent\hangindent=1.5em\hangafter=1 Nirenberg, L., \hypertarget{ref:Nirenberg-1}{[1]} On elliptic partial differential equations, Ann. Sc. Norm. Sup. Pisa 13 (1959), pp. 116--162, \hypertarget{ref:Nirenberg-2}{[2]} Topics in Nonlinear Functional Analysis, New York University Lecture Notes, 1974, reprinted by the American Mathematical Society, 2001, \hypertarget{ref:Nirenberg-3}{[3]} Variational and topological methods in nonlinear problems, Bull. Amer. Math. Soc. 4, (1981), pp. 267--302.\par\addvspace{0.35\baselineskip}
\noindent\hangindent=1.5em\hangafter=1 Nussbaum, R., \hypertarget{ref:Nussbaum-1}{[1]} Eigenvectors of nonlinear positive operators and the linear Krein-Rutman theorem, in Fixed Point Theory (Fadell, E. and Fournier, G., eds.), Lecture Notes in Math., 886, Springer, 1981, pp. 309--330.\par\addvspace{0.35\baselineskip}
\noindent\hangindent=1.5em\hangafter=1 Oleinik, O. and Radkevitch, E., \hypertarget{ref:OleinikRadkevitch-1}{[1]} Second Order Equations with Non-negative Characteristic Form, Plenum, 1973.\par\addvspace{0.35\baselineskip}
\noindent\hangindent=1.5em\hangafter=1 Osserman, R., \hypertarget{ref:Osserman-1}{[1]} Isoperimetric inequalities and eigenvalues of the Laplacian, in Proc. Int. Congress of Math. Helsinski, 1978, and Bull. Amer. Math. Soc. 84 (1978), pp. 1182--1238.\par\addvspace{0.35\baselineskip}
\noindent\hangindent=1.5em\hangafter=1 Pazy, A., \hypertarget{ref:Pazy-1}{[1]} Semigroups of Linear Operators and Applications to Partial Differential Equations, Springer, 1983.\par\addvspace{0.35\baselineskip}
\noindent\hangindent=1.5em\hangafter=1 Pearcy, C. M., \hypertarget{ref:Pearcy-1}{[1]} Some Recent Developments in Operator Theory, CBMS Reg. Conf. Series, American Mathematical Society, 1978.\par\addvspace{0.35\baselineskip}
\noindent\hangindent=1.5em\hangafter=1 Phelps, R., \hypertarget{ref:Phelps-1}{[1]} Lectures on Choquet's Theorem,Van Nostrand, 1966, \hypertarget{ref:Phelps-2}{[2]} Convex Functions, Mono- tone Operators and Differentiability, Lecture Notes in Math., 1364, Springer, 1989.\par\addvspace{0.35\baselineskip}
\noindent\hangindent=1.5em\hangafter=1 Pietsch, A., \hypertarget{ref:Pietsch-1}{[1]} History of Banach Spaces and Linear Operators, Birkhäuser, 2007.\par\addvspace{0.35\baselineskip}
\noindent\hangindent=1.5em\hangafter=1 Pinchover,Y. and Rubinstein, J., \hypertarget{ref:PinchoverRubinstein-1}{[1]} An Introduction to Partial Differential Equations, Cambridge University Press, 2005.\par\addvspace{0.35\baselineskip}
\noindent\hangindent=1.5em\hangafter=1 Protter, M. and Weinberger, H., \hypertarget{ref:ProtterWeinberger-1}{[1]} Maximum Principles in Differential Equations, Prentice-Hall, 1967.\par\addvspace{0.35\baselineskip}
\noindent\hangindent=1.5em\hangafter=1 Pucci, P. and Serrin, J., \hypertarget{ref:PucciSerrin-1}{[1]} The Maximum Principle, Birkhäuser, 2007.\par\addvspace{0.35\baselineskip}
\noindent\hangindent=1.5em\hangafter=1 Rabinowitz, P., \hypertarget{ref:Rabinowitz-1}{[1]} Variational methods for nonlinear eigenvalue problems, in Eigenvalues of Nonlinear Problems CIME, Cremonese, 1974, \hypertarget{ref:Rabinowitz-2}{[2]} Théorie du degré topologique et applications à des problèmes aux limites non linéaires, Lecture Notes written by H. Berestycki, Laboratoire d'Analyse Numérique, Université Paris VI, 1975.\par\addvspace{0.35\baselineskip}
\noindent\hangindent=1.5em\hangafter=1 Reed, M. and Simon, B., \hypertarget{ref:ReedSimon-1}{[1]} Methods of Modern Mathematical Physics (4 volumes), Academic Press, 1972--1979.\par\addvspace{0.35\baselineskip}
\noindent\hangindent=1.5em\hangafter=1 Rees, C., Shah, S., and Stanojevic, C., \hypertarget{ref:ReesShahStanojevic-1}{[1]} Theory and Applications of Fourier Series, Dekker, 1981.\par\addvspace{0.35\baselineskip}
\noindent\hangindent=1.5em\hangafter=1 Rockafellar, R. T., \hypertarget{ref:Rockafellar-1}{[1]} Convex Analysis, Princeton University Press, 1970, \hypertarget{ref:Rockafellar-2}{[2]} The Theory of Subgradients and ItsApplications to Problems of Optimization: Convex and Nonconvex Functions, Helderman Verlag, Berlin, 1981.\par\addvspace{0.35\baselineskip}
\noindent\hangindent=1.5em\hangafter=1 Royden, H. L., \hypertarget{ref:Royden-1}{[1]} Real Analysis, Macmillan, 1963.\par\addvspace{0.35\baselineskip}
\noindent\hangindent=1.5em\hangafter=1 Rudin, W., \hypertarget{ref:Rudin-1}{[1]} Functional Analysis, McGraw Hill, 1973, \hypertarget{ref:Rudin-2}{[2]} Real and Complex Analysis, McGraw Hill, 1987.\par\addvspace{0.35\baselineskip}
\noindent\hangindent=1.5em\hangafter=1 Ruskai, M. B. et al. (eds.), \hypertarget{ref:RuskaiEtAl-1}{[1]} Wavelets and Their Applications, Jones and Bartlett, 1992.\par\addvspace{0.35\baselineskip}
\noindent\hangindent=1.5em\hangafter=1 Sadosky, C., \hypertarget{ref:Sadosky-1}{[1]} Interpolation of Operators and Singular Integrals, Dekker, 1979.\par\addvspace{0.35\baselineskip}
\noindent\hangindent=1.5em\hangafter=1 Schaefer, H., \hypertarget{ref:Schaefer-1}{[1]} Topological Vector Spaces, Springer, 1971.\par\addvspace{0.35\baselineskip}
\noindent\hangindent=1.5em\hangafter=1 Schechter, M., \hypertarget{ref:Schechter-1}{[1]} Principles of Functional Analysis,Academic Press, 1971, \hypertarget{ref:Schechter-2}{[2]} Operator Methods in Quantum Mechanics, North-Holland, 1981.\par\addvspace{0.35\baselineskip}
\noindent\hangindent=1.5em\hangafter=1 Schwartz, J. T., \hypertarget{ref:Schwartz-1}{[1]} Nonlinear Functional Analysis, Gordon Breach, 1969.\par\addvspace{0.35\baselineskip}
\noindent\hangindent=1.5em\hangafter=1 Schwartz, L., \hypertarget{ref:Schwartz-1}{[1]} Théorie des distributions, Hermann, 1973, \hypertarget{ref:Schwartz-2}{[2]} Geometry and Probability in Banach spaces, Bull. Amer. Math. Soc., 4 (1981), 135--141, and Lecture Notes in Math., 852, Springer, 1981, \hypertarget{ref:Schwartz-3}{[3]} Fonctions mesurables et $\star$-scalairement mesurables, propriété de Radon- Nikodym, Exposés 4, 5 et 6, Séminaire Maurey-Schwartz, École Polytechnique (1974--1975).\par\addvspace{0.35\baselineskip}
\noindent\hangindent=1.5em\hangafter=1 Serrin, J., \hypertarget{ref:Serrin-1}{[1]} The solvability of boundary value problems, in Mathematical Developments Arising from Hilbert Problems (Browder, F., ed.), Proc. Sympos. Pure Math., Vol. 28, Part 2, American Mathematical Society, 1977, pp. 507--525.\par\addvspace{0.35\baselineskip}
\noindent\hangindent=1.5em\hangafter=1 Simon, L., \hypertarget{ref:Simon-1}{[1]} Lectures on Geometric Measure Theory, Australian National University, Center for Mathematical Analysis, Canberra, 1983, \hypertarget{ref:Simon-2}{[2]} Schauder estimates by scaling, Calc. Var. Partial Differential Equations 5 (1997), pp. 391--407.\par\addvspace{0.35\baselineskip}
\noindent\hangindent=1.5em\hangafter=1 Singer, I., \hypertarget{ref:Singer-1}{[1]} Bases in Banach Spaces (2 volumes) Springer, 1970.\par\addvspace{0.35\baselineskip}
\noindent\hangindent=1.5em\hangafter=1 Singer, I. M., \hypertarget{ref:Singer-1}{[1]} Eigenvalues of the Laplacian and invariants, of manifolds, in Proc. Int. Congress of Math. Vancouver, 1974.\par\addvspace{0.35\baselineskip}
\noindent\hangindent=1.5em\hangafter=1 Sperb, R., \hypertarget{ref:Sperb-1}{[1]} Maximum Principles and Their Applications, Academic Press, 1981.\par\addvspace{0.35\baselineskip}
\noindent\hangindent=1.5em\hangafter=1 Stampacchia, G., \hypertarget{ref:Stampacchia-1}{[1]} Equations elliptiques du second ordre à coefficients discontinus, Presses de l'Université de Montreal, 1966.\par\addvspace{0.35\baselineskip}
\noindent\hangindent=1.5em\hangafter=1 Stein, E., \hypertarget{ref:Stein-1}{[1]} Singular Integrals and Differentiability Properties of Functions, Princeton University Press, 1970.\par\addvspace{0.35\baselineskip}
\noindent\hangindent=1.5em\hangafter=1 Stein, E. and Weiss, G., \hypertarget{ref:SteinWeiss-1}{[1]} Introduction to Fourier Analysis on Euclidean Spaces, Princeton University Press, 1971.\par\addvspace{0.35\baselineskip}
\noindent\hangindent=1.5em\hangafter=1 Stoer, J. and Witzgall, C., \hypertarget{ref:StoerWitzgall-1}{[1]} Convexity and Optimization in Finite Dimensions, Springer, 1970.\par\addvspace{0.35\baselineskip}
\noindent\hangindent=1.5em\hangafter=1 Strauss, W., \hypertarget{ref:Strauss-1}{[1]} Partial Differential Equations: An Introduction, Wiley, 1992.\par\addvspace{0.35\baselineskip}
\noindent\hangindent=1.5em\hangafter=1 Stroock, D., \hypertarget{ref:Stroock-1}{[1]} An Introduction to the Theory of Large Deviations, Springer, 1984.\par\addvspace{0.35\baselineskip}
\noindent\hangindent=1.5em\hangafter=1 Stroock, D. and Varadhan, S., \hypertarget{ref:StroockVaradhan-1}{[1]} Multidimensional Diffusion Processes, Springer, 1979.\par\addvspace{0.35\baselineskip}
\noindent\hangindent=1.5em\hangafter=1 Struwe, M., \hypertarget{ref:Struwe-1}{[1]} Variational Methods: Applications to Nonlinear Partial Differential Equations and Hamiltonian Systems, Springer, 1990.\par\addvspace{0.35\baselineskip}
\noindent\hangindent=1.5em\hangafter=1 Szankowski, A., \hypertarget{ref:Szankowski-1}{[1]} B(H) does not have the approximation property, Acta Math. 147 (1981), pp. 89--108.\par\addvspace{0.35\baselineskip}
\noindent\hangindent=1.5em\hangafter=1 Talenti, G., \hypertarget{ref:Talenti-1}{[1]} Best constants in Sobolev inequality,Ann. Mat. PuraAppl. 110 (1976), pp. 353--372.\par\addvspace{0.35\baselineskip}
\noindent\hangindent=1.5em\hangafter=1 Tanabe, H., \hypertarget{ref:Tanabe-1}{[1]} Equations of Evolution, Pitman, 1979.\par\addvspace{0.35\baselineskip}
\noindent\hangindent=1.5em\hangafter=1 Taylor, A. and Lay, D., \hypertarget{ref:TaylorLay-1}{[1]} Introduction to Functional Analysis, Wiley, 1980.\par\addvspace{0.35\baselineskip}
\noindent\hangindent=1.5em\hangafter=1 Taylor, M., \hypertarget{ref:Taylor-1}{[1]} Partial Differential Equations, vols. I--III, Springer, 1996.\par\addvspace{0.35\baselineskip}
\noindent\hangindent=1.5em\hangafter=1 Temam, T, \hypertarget{ref:Temam-1}{[1]} Navier--Stokes Equations, North-Holland, 1979.\par\addvspace{0.35\baselineskip}
\noindent\hangindent=1.5em\hangafter=1 Temam, R. and Strang, G., \hypertarget{ref:TemamStrang-1}{[1]} Duality and relaxation in the variational problems of plasticity, J. de Mécanique 19 (1980), pp. 493--528, \hypertarget{ref:TemamStrang-2}{[2]} Functions of bounded deformation,Archive Rat. Mech. Anal. 75 (1980), pp. 7--21.\par\addvspace{0.35\baselineskip}
\noindent\hangindent=1.5em\hangafter=1 Toland, J. F., \hypertarget{ref:Toland-1}{[1]} Duality in nonconvex optimization, J. Math. Anal. Appl. 66 (1978), pp. 399--415, \hypertarget{ref:Toland-2}{[2]} A duality principle for nonconvex optimization and the calculus of variations,Arch. Rat. Mech. Anal. 71, (1979), 41--61, \hypertarget{ref:Toland-3}{[3]} Stability of heavy rotating chains, J. Diff. Eq. 32 (1979), pp. 15--31, \hypertarget{ref:Toland-4}{[4]} Self-adjoint operators and cones, J. London Math. Soc. 53 (1996), pp. 167--183.\par\addvspace{0.35\baselineskip}
\noindent\hangindent=1.5em\hangafter=1 Treves, F., \hypertarget{ref:Treves-1}{[1]} Topological Vector Spaces, Distributions and Kernels, Academic Press, 1967, \hypertarget{ref:Treves-2}{[2]} Linear Partial Differential Equations with Constant Coefficients, Gordon Breach, 1967, \hypertarget{ref:Treves-3}{[3]}Locally Convex Spaces and Linear Partial Differential Equations, Springer, 1967, \hypertarget{ref:Treves-4}{[4]} Basic Linear Partial Differential Equations, Academic Press, 1975.\par\addvspace{0.35\baselineskip}
\noindent\hangindent=1.5em\hangafter=1 Triebel, H., \hypertarget{ref:Triebel-1}{[1]} Theory of function spaces, (3 volumes), Birkhäuser, 1983--2006.\par\addvspace{0.35\baselineskip}
\noindent\hangindent=1.5em\hangafter=1 Uhlmann, G. (ed.), \hypertarget{ref:Uhlmann-1}{[1]} Inside Out: Inverse Problems and Applications, MSRI Publications, Vol- ume 47, Cambridge University Press, 2003.\par\addvspace{0.35\baselineskip}
\noindent\hangindent=1.5em\hangafter=1 Volpert,A. I., \hypertarget{ref:Volpert-1}{[1]}The spaces BV and quasilinear equations, Mat. USSR-Sbornik 2 (1967), pp. 225-- 267.\par\addvspace{0.35\baselineskip}
\noindent\hangindent=1.5em\hangafter=1 Weinberger,H.,\hypertarget{ref:Weinberger-1}{[1]}AFirstCourseinPartialDifferentialEquations,Blaisdell,1965,\hypertarget{ref:Weinberger-2}{[2]}Variational Methods for Eigenvalue Approximation, Reg. Conf. Appl. Math., SIAM, 1974.\par\addvspace{0.35\baselineskip}
\noindent\hangindent=1.5em\hangafter=1 Weidmann, J., \hypertarget{ref:Weidmann-1}{[1]} Linear Operators and Hilbert Spaces, Springer, 1980.\par\addvspace{0.35\baselineskip}
\noindent\hangindent=1.5em\hangafter=1 Weir, A. J., \hypertarget{ref:Weir-1}{[1]} General Integration and Measure, Cambridge University Press, 1974.\par\addvspace{0.35\baselineskip}
\noindent\hangindent=1.5em\hangafter=1 Wheeden, R. and Zygmund, A., \hypertarget{ref:WheedenZygmund-1}{[1]} Measure and Integral, Dekker, 1977.\par\addvspace{0.35\baselineskip}
\noindent\hangindent=1.5em\hangafter=1 Willem, M., \hypertarget{ref:Willem-1}{[1]} Minimax Theorems, Birkhäuser, 1996.\par\addvspace{0.35\baselineskip}
\noindent\hangindent=1.5em\hangafter=1 Wojtaszczyk P., \hypertarget{ref:Wojtaszczyk-1}{[1]} A Mathematical Introduction to Wavelets, Cambridge University Press, 1997.\par\addvspace{0.35\baselineskip}
\noindent\hangindent=1.5em\hangafter=1 Yau, S. T., \hypertarget{ref:Yau-1}{[1]} The role of partial differential equations in differential geometry, in Proc. Int. Congress of Math. Helsinki, 1978.\par\addvspace{0.35\baselineskip}
\noindent\hangindent=1.5em\hangafter=1 Yosida, K, \hypertarget{ref:Yosida-1}{[1]} Functional Analysis, Springer, 1965.\par\addvspace{0.35\baselineskip}
\noindent\hangindent=1.5em\hangafter=1 Zeidler, E., \hypertarget{ref:Zeidler-1}{[1]} Nonlinear Functional Analysis and Its Applications, Springer, 1988.\par\addvspace{0.35\baselineskip}
\noindent\hangindent=1.5em\hangafter=1 Zettl, A., \hypertarget{ref:Zettl-1}{[1]} Sturm--Liouville Theory, American Mathematical Society, 2005.\par\addvspace{0.35\baselineskip}
\noindent\hangindent=1.5em\hangafter=1 Ziemer, W., \hypertarget{ref:Ziemer-1}{[1]} Weakly Differentiable Functions, Springer, 1989.\par\addvspace{0.35\baselineskip}


\cleardoublepage
\phantomsection
\addcontentsline{toc}{chapter}{名词索引}
\printindex

\end{document}
